\input{preamble}

% OK, start here.
%
\begin{document}

\title{Algèbre commutative}


\maketitle

\phantomsection
\label{section-phantom}

\tableofcontents





\section{Introduction}
\label{section-introduction}

\noindent
Les notions de base d'algèbre commutative seront exposées dans ce document.
Une référence est \cite{MatCA}.






\section{Conventions}
\label{section-conventions}

\noindent
Un anneau est commutatif et possède un élément unité $1$. L'anneau nul est un anneau. En fait, c'est
le seul anneau qui ne possède pas d'idéal premier. Nous utiliserons le symbole
de Kronecker $\delta_{ij}$. Si $R \to S$ est un morphisme d'anneaux et si
$\mathfrak q$ est un idéal premier de $S$, nous employons alors la notation
``$\mathfrak p = R \cap \mathfrak q$''
pour désigner l'idéal premier qui est l'image réciproque de $\mathfrak q$ par
$R \to S$, même si $R$ n'est pas un sous-anneau de $S$ et même si $R \to S$
n'est pas injectif.






\section{Notions de base}
\label{section-rings-basic}

\noindent
Voici une liste de notions de base d'algèbre commutative. Certaines de ces
notions sont étudiées plus en détail dans la suite du texte et certaines sont
définies dans la liste, mais d'autres sont considérées comme élémentaires et ne seront pas définies.
Si la plupart des notions en italique ne vous sont pas familières, nous vous conseillons
de consulter un texte d'introduction à l'algèbre avant de poursuivre.

\begin{enumerate}
\item $R$ est un {\it anneau},
\label{item-ring}
\item $x\in R$ est {\it nilpotent},
\label{item-ring-element-nilpotent}
\item $x\in R$ est un {\it diviseur de zéro},
\label{item-ring-element-zerodivisor}
\item $x\in R$ est une {\it unité},
\label{item-ring-element-unit}
\item $e \in R$ est un {\it idempotent},
\label{item-ring-element-idempotent}
\item un idempotent $e \in R$ est dit {\it trivial} si $e = 1$ ou $e = 0$,
\label{item-idempotent-trivial}
\item $\varphi : R_1 \to R_2$ est un {\it homomorphisme d'anneaux},
\label{item-ring-homomorphism}
\item
\label{item-ring-homomorphism-finite-presentation}
$\varphi : R_1 \to R_2$ est {\it de présentation finie}, ou
{\it $R_2$ est une $R_1$-algèbre de présentation finie},
voir la Définition \ref{definition-finite-type},
\item
\label{item-ring-homomorphism-finite-type}
$\varphi : R_1 \to R_2$ est {\it de type fini}, ou
{\it $R_2$ est une $R_1$-algèbre de type fini},
voir la Définition \ref{definition-finite-type},
\item
\label{item-ring-homomorphism-finite}
$\varphi : R_1 \to R_2$ est {\it fini}, ou
{\it $R_2$ est une $R_1$-algèbre finie},
\item $R$ est un {\it anneau intègre},
\label{item-ring-domain}
\item $R$ est {\it réduit},
\label{item-ring-reduced}
\item $R$ est {\it noethérien},
\label{item-ring-Noetherian}
\item $R$ est un {\it anneau principal} ou un {\it PID},
\label{item-ring-PID}
\item $R$ est un {\it anneau euclidien},
\label{item-ring-Euclidean}
\item $R$ est un {\it anneau factoriel} ou un {\it UFD},
\label{item-ring-UFD}
\item $R$ est un {\it anneau de valuation discrète} ou un {\it dvr},
\label{item-ring-dvr}
\item $K$ est un {\it corps},
\label{item-field}
\item $L/K$ est une {\it extension de corps},
\label{item-field-extension}
\item $L/K$ est une {\it extension algébrique de corps},
\label{item-field-extension-algebraic}
\item $\{t_i\}_{i\in I}$ est une {\it base de transcendance} de $L$ sur $K$,
\label{item-transcendence-basis}
\item le {\it degré de transcendance} $\text{trdeg}(L/K)$ de $L$ sur $K$,
\label{item-transcendence-degree}
\item le corps $k$ est {\it algébriquement clos},
\label{item-algebraically-closed}
\item
\label{item-extend-into-algebraically-closed}
si $L/K$ est algébrique et si $\Omega/K$ est une extension avec $\Omega$
algébriquement clos,
alors il existe un morphisme d'anneaux $L \to \Omega$ qui prolonge le morphisme sur $K$,
\item $I \subset R$ est un {\it idéal},
\label{item-ideal}
\item $I \subset R$ est {\it radical},
\label{item-ideal-radical}
\item si $I$ est un idéal, nous disposons de son {\it radical} $\sqrt{I}$,
\label{item-radical-ideal}
\item
\label{item-ideal-nilpotent}
$I \subset R$ est {\it nilpotent} signifie que $I^n = 0$ pour
un certain $n \in \mathbf{N}$,
\item
\label{item-ideal-locally-nilpotent}
$I \subset R$ est {\it localement nilpotent} signifie que tout
élément de $I$ est nilpotent,
\item $\mathfrak p \subset R$ est un {\it idéal premier},
\label{item-prime-ideal}
\item
\label{item-prime-product-ideals}
si $\mathfrak p \subset R$ est premier, si $I, J \subset R$
sont des idéaux et si $IJ\subset \mathfrak p$, alors
$I \subset \mathfrak p$ ou $J \subset \mathfrak p$.
\item $\mathfrak m \subset R$ est un {\it idéal maximal},
\label{item-maximal-ideal}
\item tout anneau non nul possède un idéal maximal,
\label{item-exists-maximal-ideal}
\item
\label{item-jacobson-radical}
le {\it radical de Jacobson} de $R$ est $\text{rad}(R) =
\bigcap_{\mathfrak m \subset R} \mathfrak m$, l'intersection
de tous les idéaux maximaux de $R$,
\item l'idéal $(T)$ {\it engendré} par un sous-ensemble $T \subset R$,
\label{item-ideal-generated-by}
\item l'{\it anneau quotient} $R/I$,
\label{item-quotient-ring}
\item un idéal $I$ de l'anneau $R$ est premier si et seulement si $R/I$ est intègre,
\label{item-characterize-prime-ideal}
\item
\label{item-characterize-maximal-ideal}
un idéal $I$ de l'anneau $R$ est maximal si et seulement si
l'anneau $R/I$ est un corps,
\item
\label{item-inverse-image-ideal}
si $\varphi : R_1 \to R_2$ est un homomorphisme d'anneaux et si
$I \subset R_2$ est un idéal, alors $\varphi^{-1}(I)$ est un
idéal de $R_1$,
\item
\label{item-image-ideal}
si $\varphi : R_1 \to R_2$ est un homomorphisme d'anneaux et si
$I \subset R_1$ est un idéal, alors $\varphi(I) \cdot R_2$ (parfois
noté $I \cdot R_2$ ou $IR_2$) est l'idéal de $R_2$ engendré
par $\varphi(I)$,
\item
\label{item-inverse-image-prime}
si $\varphi : R_1 \to R_2$ est un homomorphisme d'anneaux et si
$\mathfrak p \subset R_2$ est un idéal premier, alors
$\varphi^{-1}(\mathfrak p)$ est un idéal premier de $R_1$,
\item $M$ est un {\it $R$-module},
\label{item-module}
\item
\label{item-annihilator}
pour $m \in M$, l'{\it annulateur}
$I = \{f \in R \mid fm = 0\}$ de $m$ dans $R$,
\item $N \subset M$ est un {\it $R$-sous-module},
\label{item-submodule}
\item $M$ est un {\it $R$-module noethérien},
\label{item-Noetherian-module}
\item $M$ est un {\it $R$-module de type fini},
\label{item-finite-module}
\item $M$ est un {\it $R$-module engendré par un nombre fini d'éléments},
\label{item-finitely-generated-module}
\item $M$ est un {\it $R$-module de présentation finie},
\label{item-finitely-presented-module}
\item $M$ est un {\it $R$-module libre},
\label{item-free-module}
\item
\label{item-extension-free}
si $0 \to K \to L \to M \to 0$ est une suite exacte courte
de $R$-modules et si $K$, $M$ sont libres, alors $L$ est libre,
\item si $N \subset M \subset L$ sont des $R$-modules, alors $L/M = (L/N)/(M/N)$,
\label{item-isomorphism-theorem}
\item $S$ est une {\it partie multiplicative de $R$},
\label{item-multiplicative-subset}
\item la {\it localisation} $R \to S^{-1}R$ de $R$,
\label{item-localization-ring}
\item
\label{item-localization-zero}
si $R$ est un anneau et si $S$ est une partie multiplicative
de $R$, alors $S^{-1}R$ est l'anneau nul si et seulement si $S$ contient $0$,
\item
\label{item-localize-nonzerodivisors}
si $R$ est un anneau et si la partie multiplicative $S$
est entièrement constituée d'éléments non diviseurs de zéro, alors $R \to S^{-1}R$
est injectif,
\item si $\varphi : R_1 \to R_2$ est un homomorphisme d'anneaux et si
$S$ est une partie multiplicative de $R_1$, alors $\varphi(S)$ est
une partie multiplicative de $R_2$,
\item
\label{item-products-multiplicative-subsets}
si $S$, $S'$ sont des parties multiplicatives de $R$,
et si $SS'$ désigne l'ensemble des produits $SS' =
\{r \in R \mid \exists s\in S, \exists s' \in S', r = ss'\}$,
alors $SS'$ est une partie multiplicative de $R$,
\item
\label{item-localization-localization}
si $S$, $S'$ sont des parties multiplicatives de $R$,
et si $\overline{S}$ désigne l'image de $S$ dans $(S')^{-1}R$,
alors $(SS')^{-1}R = \overline{S}^{-1}((S')^{-1}R)$,
\item la {\it localisation} $S^{-1}M$ du $R$-module $M$,
\label{item-localization-module}
\item
\label{item-localization-exact}
le foncteur $M \mapsto S^{-1}M$ préserve les applications injectives,
les applications surjectives et l'exactitude,
\item
\label{item-localization-localization-module}
si $S$, $S'$ sont des parties multiplicatives de $R$,
et si $M$ est un $R$-module, alors
$(SS')^{-1}M = S^{-1}((S')^{-1}M)$,
\item
\label{item-localize-ideal}
si $R$ est un anneau, si $I$ est un idéal de $R$ et si $S$ est une partie multiplicative
de $R$, alors $S^{-1}I$ est un idéal de $S^{-1}R$, et nous avons
$S^{-1}R/S^{-1}I = \overline{S}^{-1}(R/I)$, où $\overline{S}$
est l'image de $S$ dans $R/I$,
\item
\label{item-ideal-in-localization}
si $R$ est un anneau et si $S$ est une partie multiplicative
de $R$, alors tout idéal $I'$ de $S^{-1}R$ est
de la forme $S^{-1}I$, où l'on peut prendre pour $I$
l'image réciproque de $I'$ dans $R$,
\item
\label{item-submodule-in-localization}
si $R$ est un anneau, si $M$ est un $R$-module et si $S$ est une partie multiplicative
de $R$, alors tout sous-module $N'$ de $S^{-1}M$ est de la forme
$S^{-1}N$ pour un certain sous-module $N \subset M$, où
l'on peut prendre pour $N$ l'image réciproque de $N'$ dans $M$,
\item si $S = \{1, f, f^2, \ldots\}$, alors $R_f = S^{-1}R$ et $M_f = S^{-1}M$,
\label{item-localize-f}
\item
\label{item-localize-p}
si $S = R \setminus \mathfrak p = \{x\in R \mid x\not\in \mathfrak p\}$
pour un certain idéal premier $\mathfrak p$,
alors on note usuellement $R_{\mathfrak p} = S^{-1}R$
et $M_{\mathfrak p} = S^{-1}M$,
\item un {\it anneau local} est un anneau qui possède exactement un idéal maximal,
\label{item-local-ring}
\item un {\it anneau semi-local} est un anneau qui possède un nombre fini d'idéaux maximaux,
\label{item-semi-local-ring}
\item
\label{item-localize-p-local-ring}
si $\mathfrak p$ est un idéal premier de $R$, alors $R_{\mathfrak p}$ est
un anneau local d'idéal maximal $\mathfrak p R_{\mathfrak p}$,
\item
\label{item-residue-field}
le {\it corps résiduel}, noté $\kappa(\mathfrak p)$,
de l'idéal premier $\mathfrak p$ de l'anneau $R$ est le
corps des fractions de l'anneau intègre $R/\mathfrak p$ ;
il est égal à $R_\mathfrak p/\mathfrak pR_\mathfrak p
= (R \setminus \mathfrak p)^{-1}R/\mathfrak p$,
\item étant donnés $R$ et $M_1$, $M_2$, le {\it produit tensoriel} $M_1 \otimes_R M_2$,
\label{item-tensor-product}
\item
\label{item-cauchy-binet}
étant données des matrices $A$ et $B$ dans un anneau $R$, de tailles $m \times n$ et
$n \times m$, nous avons $\det(AB) = \sum \det(A_S)\det({}_SB)$ dans $R$, où
la somme porte sur les sous-ensembles $S \subset \{1, \ldots, n\}$ de cardinal $m$,
et $A_S$ est la sous-matrice $m \times m$ de $A$ formée des colonnes
correspondant à $S$, tandis que ${}_SB$ est la sous-matrice $m \times m$ de $B$
formée des lignes correspondant à $S$,
\item etc.
\end{enumerate}




\section{Lemme du serpent}
\label{section-snake}

\noindent
Le lemme du serpent et ses variantes sont étudiés dans le cadre des
catégories abéliennes dans
Homologie, section \ref{homology-section-abelian-categories}.

\begin{lemma}
\label{lemma-snake}
\begin{reference}
\cite[III, Lemma 3.3]{Cartan-Eilenberg}
\end{reference}
Étant donné un diagramme commutatif
$$
\xymatrix{
& X \ar[r] \ar[d]^\alpha &
Y \ar[r] \ar[d]^\beta &
Z \ar[r] \ar[d]^\gamma &
0 \\
0 \ar[r] & U \ar[r] & V \ar[r] & W
}
$$
de groupes abéliens dont les lignes sont exactes, il existe une suite exacte canonique
$$
\Ker(\alpha) \to \Ker(\beta) \to \Ker(\gamma)
\to
\Coker(\alpha) \to \Coker(\beta) \to \Coker(\gamma)
$$
De plus, si $X \to Y$ est injectif, la première application est
injective ; si $V \to W$ est surjectif, la dernière
application est surjective.
\end{lemma}

\begin{proof}
L'application $\partial : \Ker(\gamma) \to \Coker(\alpha)$ est définie
comme suit. Soit $z \in \Ker(\gamma)$. Choisissons $y \in Y$ d'image $z$.
Alors $\beta(y) \in V$ s'envoie sur zéro dans $W$. Ainsi, $\beta(y)$ est l'image d'un
élément $u \in U$. Posons $\partial z = \overline{u}$, la classe de $u$ dans le
conoyau de $\alpha$. La démonstration de l'exactitude est omise.
\end{proof}

\section{Modules de type fini et modules de présentation finie}
\label{section-module-finite-type}

\noindent
Voici simplement quelques notations et lemmes élémentaires.

\begin{definition}
\label{definition-module-finite-type}
Soit $R$ un anneau. Soit $M$ un $R$-module.
\begin{enumerate}
\item Nous disons que $M$ est un {\it $R$-module de type fini}, ou un {\it
$R$-module engendré par un nombre fini d'éléments}, s'il existe $n \in \mathbf{N}$ et $x_1, \ldots, x_n \in M$
tels que tout élément de $M$ soit une combinaison $R$-linéaire des $x_i$.
De manière équivalente, cela signifie qu'il existe une surjection
$R^{\oplus n} \to M$ pour un certain $n \in \mathbf{N}$.
\item Nous disons que $M$ est un {\it $R$-module de présentation finie} ou un
{\it $R$-module de présentation finie} s'il existe des entiers
$n, m \in \mathbf{N}$ et une suite exacte
$$
R^{\oplus m} \longrightarrow R^{\oplus n} \longrightarrow M \longrightarrow 0
$$
\end{enumerate}
\end{definition}

\noindent
De manière informelle, $M$ est un $R$-module de présentation finie si et seulement s'il
est de type fini et si le module des relations entre ces
générateurs est lui aussi de type fini.
Le choix d'une suite exacte comme dans la définition s'appelle une
{\it présentation} de $M$.

\begin{lemma}
\label{lemma-lift-map}
Soit $R$ un anneau. Soient $\alpha : R^{\oplus n} \to M$ et $\beta : N \to M$ des
homomorphismes de modules. Si $\Im(\alpha) \subset \Im(\beta)$, il existe alors
un homomorphisme de $R$-modules $\gamma : R^{\oplus n} \to N$ tel que
$\alpha = \beta \circ \gamma$.
\end{lemma}

\begin{proof}
Soit $e_i = (0, \ldots, 0, 1, 0, \ldots, 0)$ le $i$-ème vecteur de base
de $R^{\oplus n}$. Soit $x_i \in N$ un élément tel que
$\alpha(e_i) = \beta(x_i)$ ; il existe par hypothèse. Posons
$\gamma(a_1, \ldots, a_n) = \sum a_i x_i$. Par construction,
$\alpha = \beta \circ \gamma$.
\end{proof}

\begin{lemma}
\label{lemma-extension}
Soit $R$ un anneau.
Soit
$$
0 \to M_1 \to M_2 \to M_3 \to 0
$$
une suite exacte courte de $R$-modules.
\begin{enumerate}
\item Si $M_1$ et $M_3$ sont des $R$-modules de type fini, alors $M_2$ est un
$R$-module de type fini.
\item Si $M_1$ et $M_3$ sont des $R$-modules de présentation finie, alors $M_2$
est un $R$-module de présentation finie.
\item Si $M_2$ est un $R$-module de type fini, alors $M_3$ est un $R$-module de type fini.
\item Si $M_2$ est un $R$-module de présentation finie et si $M_1$ est un
$R$-module de type fini, alors $M_3$ est un $R$-module de présentation finie.
\item Si $M_3$ est un $R$-module de présentation finie et si $M_2$ est un
$R$-module de type fini, alors $M_1$ est un $R$-module de type fini.
\end{enumerate}
\end{lemma}

\begin{proof}
Démonstration de (1). Si $x_1, \ldots, x_n$ sont des générateurs de $M_1$ et si
$y_1, \ldots, y_m \in M_2$ sont des éléments dont les images dans $M_3$ sont des
générateurs de $M_3$, alors $x_1, \ldots, x_n, y_1, \ldots, y_m$
engendrent $M_2$.

\medskip\noindent
L'assertion (3) découle immédiatement de la définition.

\medskip\noindent
Démonstration de (5). Supposons $M_3$ de présentation finie et $M_2$ de type fini.
Choisissons une présentation
$$
R^{\oplus m} \to R^{\oplus n} \to M_3 \to 0
$$
D'après le Lemme \ref{lemma-lift-map}, il existe une application
$R^{\oplus n} \to M_2$ telle que
le diagramme en traits pleins
$$
\xymatrix{
& R^{\oplus m} \ar[r] \ar@{..>}[d] & R^{\oplus n} \ar[r] \ar[d] &
M_3 \ar[r] \ar[d]^{\text{id}} & 0 \\
0 \ar[r] & M_1 \ar[r] & M_2 \ar[r] & M_3 \ar[r] & 0
}
$$
soit commutatif. On obtient ainsi la flèche en pointillé. Par le lemme du serpent
(Lemme \ref{lemma-snake}), on voit que l'on obtient un isomorphisme
$$
\Coker(R^{\oplus m} \to M_1)
\cong
\Coker(R^{\oplus n} \to M_2)
$$
En particulier, nous en déduisons que $\Coker(R^{\oplus m} \to M_1)$
est un $R$-module de type fini. Puisque $\Im(R^{\oplus m} \to M_1)$
est de type fini d'après (3), on voit que $M_1$ est de type fini d'après l'assertion (1).

\medskip\noindent
Démonstration de (4). Supposons $M_2$ de présentation finie et $M_1$ de type fini.
Choisissons une présentation $R^{\oplus m} \to R^{\oplus n} \to M_2 \to 0$.
Choisissons une surjection $R^{\oplus k} \to M_1$. D'après le Lemme \ref{lemma-lift-map},
il existe une factorisation $R^{\oplus k} \to R^{\oplus n} \to M_2$
de la composée $R^{\oplus k} \to M_1 \to M_2$. Alors
$R^{\oplus k + m} \to R^{\oplus n} \to M_3 \to 0$
est une présentation.

\medskip\noindent
Démonstration de (2). Supposons $M_1$ et $M_3$ de présentation finie.
L'argument de la démonstration de l'assertion (1) fournit un diagramme commutatif
$$
\xymatrix{
0 \ar[r] & R^{\oplus n} \ar[d] \ar[r] & R^{\oplus n + m} \ar[d] \ar[r] &
R^{\oplus m} \ar[d] \ar[r] & 0 \\
0 \ar[r] & M_1 \ar[r] & M_2 \ar[r] & M_3 \ar[r] & 0
}
$$
dont les flèches verticales sont surjectives. Par le lemme du serpent, nous obtenons une suite
exacte courte
$$
0 \to \Ker(R^{\oplus n} \to M_1) \to
\Ker(R^{\oplus n + m} \to M_2) \to
\Ker(R^{\oplus m} \to M_3) \to 0
$$
D'après l'assertion (5), les deux modules extrêmes sont de type fini. Par conséquent, celui du
milieu l'est aussi. D'après (4), on voit que $M_2$ est de présentation finie.
\end{proof}

\begin{lemma}
\label{lemma-trivial-filter-finite-module}
\begin{slogan}
Les modules de type fini admettent des filtrations dont les quotients successifs sont des
modules cycliques.
\end{slogan}
Soit $R$ un anneau et soit $M$ un $R$-module de type fini.
Il existe une filtration par des $R$-sous-modules de type fini
$$
0 = M_0 \subset M_1 \subset \ldots \subset M_n = M
$$
telle que chaque quotient $M_i/M_{i - 1}$ soit isomorphe
à $R/I_i$ pour un certain idéal $I_i$ de $R$.
\end{lemma}

\begin{proof}
Nous raisonnons par récurrence sur le nombre de générateurs de $M$. Soient
$x_1, \ldots, x_r \in M$ des générateurs.
Posons $M' = Rx_1 \subset M$. Alors $M/M'$ possède $r - 1$ générateurs
et l'hypothèse de récurrence s'applique. De plus, il est clair que $M' \cong R/I_1$
avec $I_1 = \{f \in R \mid fx_1 = 0\}$.
\end{proof}

\begin{lemma}
\label{lemma-finite-over-subring}
Soit $R \to S$ un morphisme d'anneaux.
Soit $M$ un $S$-module.
Si $M$ est de type fini comme $R$-module, alors $M$ est de type fini comme $S$-module.
\end{lemma}

\begin{proof}
En effet, tout système générateur sur $R$ de $M$ est aussi un système générateur sur $S$ de
$M$, puisque la structure de $R$-module est induite par l'image de $R$ dans $S$.
\end{proof}




\section{Morphismes d'anneaux de type fini et de présentation finie}
\label{section-finite-type}

\begin{definition}
\label{definition-finite-type}
Soit $R \to S$ un morphisme d'anneaux.
\begin{enumerate}
\item Nous disons que $R \to S$ est de {\it type fini}, ou que {\it $S$ est une
$R$-algèbre de type fini}, s'il existe un $n \in \mathbf{N}$ et une surjection
de $R$-algèbres $R[x_1, \ldots, x_n] \to S$.
\item Nous disons que $R \to S$ est de {\it présentation finie} s'il
existe des entiers $n, m \in \mathbf{N}$, des polynômes
$f_1, \ldots, f_m \in R[x_1, \ldots, x_n]$
et un isomorphisme de $R$-algèbres
$R[x_1, \ldots, x_n]/(f_1, \ldots, f_m) \cong S$.
\end{enumerate}
\end{definition}

\noindent
De manière informelle, $R \to S$ est de présentation finie si et seulement si
$S$ est de type fini comme $R$-algèbre
et si l'idéal des relations entre les générateurs est de type fini.
Le choix d'une surjection $R[x_1, \ldots, x_n] \to S$ comme dans la définition
est parfois appelé une {\it présentation} de $S$.

\begin{lemma}
\label{lemma-compose-finite-type}
Les notions de type fini et de présentation finie possèdent les propriétés
de permanence suivantes.
\begin{enumerate}
\item Une composée de morphismes d'anneaux de type fini est de type fini.
\item Une composée de morphismes d'anneaux de présentation finie est de présentation
finie.
\item Étant donné $R \to S' \to S$ avec $R \to S$ de type fini,
le morphisme $S' \to S$ est alors de type fini.
\item Étant donné $R \to S' \to S$, avec $R \to S$ de présentation finie
et $R \to S'$ de type fini, le morphisme $S' \to S$ est alors de présentation finie.
\end{enumerate}
\end{lemma}

\begin{proof}
Nous ne démontrons que la dernière assertion.
Écrivons $S = R[x_1, \ldots, x_n]/(f_1, \ldots, f_m)$
et $S' = R[y_1, \ldots, y_a]/I$. Supposons que la classe
$\bar y_i$ de $y_i$ s'envoie
sur $h_i \bmod (f_1, \ldots, f_m)$ dans $S$.
Il est alors clair que
$S = S'[x_1, \ldots, x_n]/(f_1, \ldots, f_m,
h_1 - \bar y_1, \ldots, h_a - \bar y_a)$.
\end{proof}

\begin{lemma}
\label{lemma-finite-presentation-independent}
Soit $R \to S$ un morphisme d'anneaux de présentation finie.
Pour toute surjection $\alpha : R[x_1, \ldots, x_n] \to S$, le
noyau de $\alpha$ est un idéal de type fini de $R[x_1, \ldots, x_n]$.
\end{lemma}

\begin{proof}
Écrivons $S = R[y_1, \ldots, y_m]/(f_1, \ldots, f_k)$.
Choisissons des $g_i \in R[y_1, \ldots, y_m]$ qui relèvent
les $\alpha(x_i)$. On voit alors que $S = R[x_i, y_j]/(f_l, x_i - g_i)$.
Choisissons $h_j \in R[x_1, \ldots, x_n]$ tel que $\alpha(h_j)$
corresponde à $y_j \bmod (f_1, \ldots, f_k)$. Considérons
l'application $\psi : R[x_i, y_j] \to R[x_i]$, $x_i \mapsto x_i$,
$y_j \mapsto h_j$. Le noyau de $\alpha$
est alors l'image de $(f_l, x_i - g_i)$ par $\psi$, ce qui conclut.
\end{proof}

\begin{lemma}
\label{lemma-finitely-presented-over-subring}
Soit $R \to S$ un morphisme d'anneaux.
Soit $M$ un $S$-module.
Supposons $R \to S$ de type fini et
$M$ de présentation finie comme $R$-module.
Alors $M$ est de présentation finie comme $S$-module.
\end{lemma}

\begin{proof}
La démonstration est semblable à celle de l'assertion (4) du
Lemme \ref{lemma-compose-finite-type}.
Nous pouvons supposer $S = R[x_1, \ldots, x_n]/J$.
Choisissons $y_1, \ldots, y_m \in M$ qui engendrent $M$ comme $R$-module
et des relations $\sum a_{ij} y_j = 0$, $i = 1, \ldots, t$, qui
engendrent le noyau de $R^{\oplus m} \to M$. Pour tous
$i = 1, \ldots, n$ et $j = 1, \ldots, m$, écrivons
$$
x_i y_j = \sum a_{ijk} y_k
$$
pour certains $a_{ijk} \in R$. Considérons le $S$-module $N$ engendré par
$y_1, \ldots, y_m$ et soumis aux relations
$\sum a_{ij} y_j = 0$, $i = 1, \ldots, t$, ainsi que
$x_i y_j = \sum a_{ijk} y_k$, $i = 1, \ldots, n$ et $j = 1, \ldots, m$.
Alors $N$ possède une présentation
$$
S^{\oplus nm + t} \longrightarrow S^{\oplus m} \longrightarrow N
\longrightarrow 0
$$
Par construction, il existe une application surjective $\varphi : N \to M$.
Pour achever la démonstration, montrons que $\varphi$ est injective.
Supposons $z = \sum b_j y_j \in N$ pour certains $b_j \in S$.
Nous pouvons considérer les $b_j$ comme des polynômes en $x_1, \ldots, x_n$
à coefficients dans $R$.
En appliquant les relations de la forme $x_i y_j = \sum a_{ijk} y_k$,
nous pouvons abaisser par récurrence le degré des polynômes.
Nous voyons donc que $z = \sum c_j y_j$ pour certains $c_j \in R$.
Par conséquent, si $\varphi(z) = 0$, le vecteur $(c_1, \ldots, c_m)$
est une combinaison $R$-linéaire des vecteurs $(a_{i1}, \ldots, a_{im})$,
et nous concluons que $z = 0$, comme souhaité.
\end{proof}




\section{Morphismes d'anneaux finis}
\label{section-finite}

\noindent
Voici la définition.

\begin{definition}
\label{definition-finite-ring-map}
Soit $\varphi : R \to S$ un morphisme d'anneaux. Nous disons que $\varphi : R \to S$ est
{\it fini} si $S$ est de type fini comme $R$-module.
\end{definition}

\begin{lemma}
\label{lemma-finite-module-over-finite-extension}
Soit $R \to S$ un morphisme d'anneaux fini.
Soit $M$ un $S$-module.
Alors $M$ est de type fini comme $R$-module si et seulement si $M$ est de type fini
comme $S$-module.
\end{lemma}

\begin{proof}
L'une des implications résulte du
Lemme \ref{lemma-finite-over-subring}.
Pour voir l'autre, supposons que $M$ soit de type fini comme $S$-module.
Choisissons $x_1, \ldots, x_n \in S$ qui engendrent $S$ comme $R$-module.
Choisissons $y_1, \ldots, y_m \in M$ qui engendrent $M$ comme $S$-module.
Alors les $x_i y_j$ engendrent $M$ comme $R$-module.
\end{proof}

\begin{lemma}
\label{lemma-finite-transitive}
Supposons que $R \to S$ et $S \to T$ soient des morphismes d'anneaux finis.
Alors $R \to T$ est fini.
\end{lemma}

\begin{proof}
Si les $t_i$ engendrent $T$ comme $S$-module et si les $s_j$ engendrent $S$ comme
$R$-module, alors les $t_i s_j$ engendrent $T$ comme $R$-module.
(Cela résulte également du
Lemme \ref{lemma-finite-module-over-finite-extension}.)
\end{proof}

\begin{lemma}
\label{lemma-finite-finite-type}
Soit $\varphi : R \to S$ un morphisme d'anneaux.
\begin{enumerate}
\item Si $\varphi$ est fini, alors $\varphi$ est de type fini.
\item Si $S$ est de présentation finie comme $R$-module, alors
$\varphi$ est de présentation finie.
\end{enumerate}
\end{lemma}

\begin{proof}
Pour (1), si $x_1, \ldots, x_n \in S$ engendrent $S$ comme $R$-module,
alors $x_1, \ldots, x_n$ engendrent $S$ comme $R$-algèbre. Pour (2),
supposons que $\sum r_j^ix_i = 0$, $j = 1, \ldots, m$, soit un système
de générateurs des relations entre les $x_i$ considérés comme des
générateurs du $R$-module $S$. Écrivons en outre
$1 = \sum r_ix_i$ pour certains $r_i \in R$, et
$x_ix_j = \sum r_{ij}^k x_k$ pour certains $r_{ij}^k \in R$.
Alors
$$
S = R[t_1, \ldots, t_n]/
(\sum r_j^it_i,\ 1 - \sum r_it_i,\ t_it_j - \sum r_{ij}^k t_k)
$$
comme $R$-algèbre, ce qui démontre (2).
\end{proof}

\noindent
Pour plus de renseignements sur les morphismes d'anneaux finis, voir
la section \ref{section-finite-ring-extensions}.


\section{Colimites}
\label{section-colimits}

\noindent
Une partie du contenu de cette section recoupe l'étude générale
des colimites dans
Catégories, sections \ref{categories-section-limits} --
\ref{categories-section-posets-limits}.
La notion d'ensemble préordonné est définie dans
Catégories, Définition \ref{categories-definition-directed-set}.
Elle est légèrement plus faible que celle d'ensemble partiellement ordonné.

\begin{definition}
\label{definition-directed-system}
Soit $(I, \leq)$ un ensemble préordonné.
Un {\it système $(M_i, \mu_{ij})$ de $R$-modules sur $I$}
est constitué d'une famille de $R$-modules $\{M_i\}_{i\in I}$ indexée
par $I$ et d'une famille d'homomorphismes de $R$-modules $\{\mu_{ij} : M_i \to M_j\}_{i \leq j}$
telles que, pour tous $i \leq j \leq k$,
$$
\mu_{ii} = \text{id}_{M_i}\quad
\mu_{ik} = \mu_{jk}\circ \mu_{ij}
$$
Nous disons que $(M_i, \mu_{ij})$ est un {\it système filtrant} si $I$ est un ensemble filtrant.
\end{definition}

\noindent
C'est la même notion que celle définie dans Catégories,
Définition \ref{categories-definition-system-over-poset}
et section \ref{categories-section-posets-limits}.
Nous renvoyons à Catégories, Définition \ref{categories-definition-colimit},
pour la définition d'une colimite d'un diagramme ou d'un système dans une
catégorie quelconque.

\begin{lemma}
\label{lemma-colimit}
Soit $(M_i, \mu_{ij})$ un système de $R$-modules sur l'ensemble préordonné $I$.
La colimite du système $(M_i, \mu_{ij})$ est le $R$-module quotient
$(\bigoplus_{i\in I} M_i) /Q$, où $Q$ est le
$R$-sous-module engendré par tous les éléments
$$
\iota_i(x_i) - \iota_j(\mu_{ij}(x_i))
$$
où $\iota_i : M_i \to \bigoplus_{i\in I} M_i$
est l'inclusion canonique. Nous notons la colimite
$M = \colim_i M_i$. Nous notons
$\pi : \bigoplus_{i\in I} M_i \to M$ la
projection et
$\phi_i = \pi \circ \iota_i : M_i \to M$.
\end{lemma}

\begin{proof}
Ce lemme est un cas particulier de
Catégories, Lemme \ref{categories-lemma-colimits-coproducts-coequalizers},
mais nous allons aussi le démontrer directement dans ce cas.
Remarquons en effet que $\phi_i = \phi_j\circ \mu_{ij}$ dans la
construction ci-dessus. Pour montrer que la paire $(M, \phi_i)$ est la colimite, nous devons
montrer qu'elle satisfait la propriété universelle : pour toute autre paire de ce type
$(Y, \psi_i)$ avec $\psi_i : M_i \to
Y$, $\psi_i = \psi_j\circ \mu_{ij}$, il existe un unique homomorphisme de
$R$-modules $g : M \to Y$ tel que le
diagramme suivant soit commutatif :
$$
\xymatrix{
M_i \ar[rr]^{\mu_{ij}} \ar[dr]^{\phi_i} \ar[ddr]_{\psi_i} & &
M_j\ar[dl]_{\phi_j} \ar[ddl]^{\psi_j} \\
& M \ar[d]^{g}\\
& Y
}
$$
Cela est clair, car nous pouvons définir $g$ en prenant
l'application $\psi_i$ sur le facteur $M_i$ de la somme directe
$\bigoplus M_i$.
\end{proof}

\begin{lemma}
\label{lemma-directed-colimit}
Soit $(M_i, \mu_{ij})$ un système de $R$-modules sur
l'ensemble préordonné $I$. Supposons que $I$ soit filtrant.
La colimite du système $(M_i, \mu_{ij})$ est canoniquement
isomorphe au module $M$ défini comme suit :
\begin{enumerate}
\item comme ensemble, posons
$$
M = \left(\coprod\nolimits_{i \in I} M_i\right)/\sim
$$
où, pour $m \in M_i$ et $m' \in M_{i'}$, nous avons
$$
m \sim m' \Leftrightarrow
\mu_{ij}(m) = \mu_{i'j}(m')\text{ pour un certain }j \geq i, i'
$$
\item comme groupe abélien, pour $m \in M_i$ et $m' \in M_{i'}$,
nous définissons la somme des classes de $m$ et $m'$ dans $M$
comme étant la classe de $\mu_{ij}(m) + \mu_{i'j}(m')$, où
$j \in I$ est un indice quelconque tel que $i \leq j$ et $i' \leq j$, et
\item comme $R$-module, pour $m \in M_i$ et $x \in R$, définissons
le produit de $x$ par la classe de $m$ dans $M$ comme étant la
classe de $xm$ dans $M$.
\end{enumerate}
Les applications canoniques $\phi_i : M_i \to M$ sont induites par les applications
canoniques $M_i \to \coprod_{i \in I} M_i$.
\end{lemma}

\begin{proof}
Démonstration omise. Comparer à
Catégories, section \ref{categories-section-directed-colimits}.
\end{proof}

\begin{lemma}
\label{lemma-zero-directed-limit}
Soit $(M_i, \mu_{ij})$ un système filtrant.
Soit $M = \colim M_i$ avec $\mu_i : M_i \to M$.
Alors $\mu_i(x_i) = 0$ pour $x_i \in M_i$ si et seulement s'il
existe $j \geq i$ tel que $\mu_{ij}(x_i) = 0$.
\end{lemma}

\begin{proof}
Cela résulte clairement de la description de la colimite filtrante
du Lemme \ref{lemma-directed-colimit}.
\end{proof}

\begin{example}
\label{example-zero-colimit-different}
Considérons l'ensemble partiellement ordonné $I = \{a, b, c\}$ avec
$a < b$ et $a < c$, et aucune autre inégalité stricte.
Un système $(M_a, M_b, M_c, \mu_{ab}, \mu_{ac})$
sur $I$ est constitué de trois $R$-modules $M_a, M_b, M_c$
et de deux homomorphismes de $R$-modules $\mu_{ab} : M_a \to M_b$ et
$\mu_{ac} : M_a \to M_c$.
La colimite du système est simplement
$$
M := \colim_{i \in I} M_i = \Coker(M_a \to M_b \oplus M_c)
$$
où l'application est $\mu_{ab} \oplus -\mu_{ac}$. Le noyau de
l'application canonique $M_a \to M$ est donc $\Ker(\mu_{ab}) + \Ker(\mu_{ac})$.
Et le noyau de l'application canonique $M_b \to M$ est l'image de
$\Ker(\mu_{ac})$ par l'application $\mu_{ab}$. Par conséquent, il est clair que
le résultat du Lemme \ref{lemma-zero-directed-limit} est faux pour les
systèmes généraux.
\end{example}

\begin{definition}
\label{definition-homomorphism-directed-systems}
Soient $(M_i, \mu_{ij})$, $(N_i, \nu_{ij})$ des
systèmes de $R$-modules sur le même ensemble préordonné $I$.
Un {\it homomorphisme de systèmes} $\Phi$ de $(M_i, \mu_{ij})$ vers
$(N_i, \nu_{ij})$ est, par définition, une famille d'homomorphismes de $R$-modules
$\phi_i : M_i \to N_i$
tels que $\phi_j \circ \mu_{ij} = \nu_{ij} \circ \phi_i$
pour tous $i \leq j$.
\end{definition}

\noindent
C'est la même notion que celle de transformation naturelle de foncteurs
entre les diagrammes associés $M : I \to \text{Mod}_R$
et $N : I \to \text{Mod}_R$, dans le langage des
catégories.
Le lemme suivant est un cas particulier de
Catégories, Lemme \ref{categories-lemma-functorial-colimit}.

\begin{lemma}
\label{lemma-homomorphism-limit}
Soient $(M_i, \mu_{ij})$, $(N_i, \nu_{ij})$ des
systèmes de $R$-modules sur le même ensemble préordonné.
Un morphisme de systèmes $\Phi = (\phi_i)$ de $(M_i, \mu_{ij})$ vers
$(N_i, \nu_{ij})$ induit un unique homomorphisme
$$
\colim \phi_i : \colim M_i \longrightarrow \colim N_i
$$
tel que
$$
\xymatrix{
M_i \ar[r] \ar[d]_{\phi_i} & \colim M_i \ar[d]^{\colim \phi_i} \\
N_i \ar[r] & \colim N_i
}
$$
soit commutatif pour tout $i \in I$.
\end{lemma}

\begin{proof}
Écrivons $M = \colim M_i$, $N = \colim N_i$ et $\phi = \colim \phi_i$
(ce dernier restant à construire). Nous utiliserons sans autre mention la description explicite de $M$
et de $N$ donnée au Lemme \ref{lemma-colimit}.
La condition du lemme équivaut à ce que
$$
\xymatrix{
\bigoplus_{i\in I} M_i \ar[r] \ar[d]_{\bigoplus\phi_i} & M \ar[d]^\phi \\
\bigoplus_{i\in I} N_i \ar[r] & N
}
$$
soit commutatif. Il est donc clair que, si $\phi$ existe, il est unique.
Pour voir que $\phi$ existe, il suffit de montrer que le noyau de la
flèche horizontale supérieure est envoyé par $\bigoplus \phi_i$ dans le noyau
de la flèche horizontale inférieure. Pour cela, soient $j \leq k$ et
$x_j \in M_j$. Alors
$$
(\bigoplus \phi_i)(x_j - \mu_{jk}(x_j))
=
\phi_j(x_j) - \phi_k(\mu_{jk}(x_j))
=
\phi_j(x_j) - \nu_{jk}(\phi_j(x_j))
$$
appartient au noyau de la flèche horizontale inférieure, comme voulu.
\end{proof}

\begin{lemma}
\label{lemma-directed-colimit-exact}
\begin{slogan}
Les colimites filtrantes sont exactes. Les colimites sur un ensemble filtrant sont exactes.
\end{slogan}
Soit $I$ un ensemble filtrant.
Soient $(L_i, \lambda_{ij})$, $(M_i, \mu_{ij})$ et
$(N_i, \nu_{ij})$ des systèmes de $R$-modules sur $I$.
Soient $\varphi_i : L_i \to M_i$ et $\psi_i : M_i \to N_i$ des
morphismes de systèmes sur $I$. Supposons que, pour tout $i \in I$, la
suite de $R$-modules
$$
\xymatrix{
L_i \ar[r]^{\varphi_i} &
M_i \ar[r]^{\psi_i} &
N_i
}
$$
soit un complexe d'homologie $H_i$.
Alors les $R$-modules $H_i$ forment un système sur $I$,
la suite de $R$-modules
$$
\xymatrix{
\colim_i L_i \ar[r]^\varphi &
\colim_i M_i \ar[r]^\psi &
\colim_i N_i
}
$$
est elle aussi un complexe et, en notant $H$ son homologie, nous avons
$$
H = \colim_i H_i.
$$
\end{lemma}

\begin{proof}
Il est clair que
$
\xymatrix{
\colim_i L_i \ar[r]^\varphi &
\colim_i M_i \ar[r]^\psi &
\colim_i N_i
}
$
est un complexe. Pour chaque $i \in I$, il existe un morphisme canonique de
$R$-modules $H_i \to H$ (qui envoie chaque
$[m] \in H_i = \Ker(\psi_i) / \Im(\varphi_i)$ sur la
classe dans $H = \Ker(\psi) / \Im(\varphi)$
de l'image de $m$ dans $\colim_i M_i$). Ces morphismes donnent lieu
à un morphisme $\colim_i H_i \to H$. Il reste à
montrer que ce morphisme est surjectif et injectif.

\medskip\noindent
Nous allons utiliser à plusieurs reprises, sans autre mention, la description des colimites sur $I$
donnée au Lemme \ref{lemma-directed-colimit}.
Soit $h \in H$.
Puisque $H = \Ker(\psi)/\Im(\varphi)$, on voit que
$h$ est la classe modulo $\Im(\varphi)$ d'un élément $[m]$
de $\Ker(\psi) \subset \colim_i M_i$. Choisissons un
$i$ tel que $[m]$ provienne d'un élément $m \in M_i$. Choisissons
$j \geq i$ tel que $\nu_{ij}(\psi_i(m)) = 0$, ce qui est possible
puisque $[m] \in \Ker(\psi)$. Après avoir remplacé $i$ par $j$ et
$m$ par $\mu_{ij}(m)$, nous voyons que nous pouvons supposer $m \in \Ker(\psi_i)$.
Cela montre que l'application $\colim_i H_i \to H$ est surjective.

\medskip\noindent
Supposons qu'un $h_i \in H_i$ ait une image nulle dans $H$. Puisque
$H_i = \Ker(\psi_i)/\Im(\varphi_i)$, nous pouvons représenter
$h_i$ par un élément $m \in \Ker(\psi_i) \subset M_i$.
L'hypothèse d'annulation de $h_i$ dans $H$ signifie que
la classe de $m$ dans $\colim_i M_i$ appartient à l'image
de $\varphi$. Il existe donc un $j \geq i$ et un $l \in L_j$
tels que $\varphi_j(l) = \mu_{ij}(m)$. Cela montre clairement que
l'image de $h_i$ dans $H_j$ est nulle. Ceci démontre
l'injectivité de $\colim_i H_i \to H$.
\end{proof}

\begin{example}
\label{example-colimit-not-exact}
La prise des colimites n'est pas exacte en général.
Considérons l'ensemble partiellement ordonné $I = \{a, b, c\}$ avec
$a < b$ et $a < c$, et aucune autre inégalité stricte,
comme dans l'Exemple \ref{example-zero-colimit-different}.
Considérons l'application de systèmes
$(0, \mathbf{Z}, \mathbf{Z}, 0, 0) \to
(\mathbf{Z}, \mathbf{Z}, \mathbf{Z}, 1, 1)$.
D'après la description de la colimite dans
l'Exemple \ref{example-zero-colimit-different},
nous voyons que l'application associée entre les colimites n'est pas injective,
bien que l'application de systèmes soit injective sur chaque objet.
Par conséquent, le résultat du Lemme \ref{lemma-directed-colimit-exact}
est faux pour les systèmes généraux.
\end{example}

\begin{lemma}
\label{lemma-almost-directed-colimit-exact}
Soit $\mathcal{I}$ une catégorie d'indices satisfaisant aux hypothèses de
Catégories, Lemme \ref{categories-lemma-split-into-directed}.
Alors la prise des colimites de diagrammes de groupes abéliens sur $\mathcal{I}$
est exacte (c'est-à-dire que l'analogue du
Lemme \ref{lemma-directed-colimit-exact}
vaut dans cette situation).
\end{lemma}

\begin{proof}
D'après
Catégories, Lemme \ref{categories-lemma-split-into-directed},
nous pouvons écrire $\mathcal{I} = \coprod_{j \in J} \mathcal{I}_j$, où chaque
$\mathcal{I}_j$ est une catégorie filtrante et où $J$ peut être vide. D'après
Catégories, Lemme \ref{categories-lemma-directed-category-system},
prendre les colimites sur les catégories d'indices $\mathcal{I}_j$ revient à
prendre la colimite sur un certain ensemble filtrant. Le
Lemme \ref{lemma-directed-colimit-exact}
s'applique donc à ces colimites. Le problème se ramène ainsi à montrer que
les coproduits dans la catégorie des $R$-modules indexés par l'ensemble $J$ sont exacts.
Autrement dit, étant données des suites exactes
$L_j \to M_j \to N_j$ de $R$-modules, nous devons montrer que
$$
\bigoplus\nolimits_{j \in J} L_j
\longrightarrow
\bigoplus\nolimits_{j \in J} M_j
\longrightarrow
\bigoplus\nolimits_{j \in J} N_j
$$
est exacte. Cela se vérifie directement et reste vrai même si $J$ est vide.
\end{proof}






\section{Localisation}
\label{section-localization}

\begin{definition}
\label{definition-multiplicative-subset}
Soit $R$ un anneau et soit $S$ un sous-ensemble de $R$.
Nous disons que $S$ est une {\it partie multiplicative de $R$} si
$1\in S$ et si $S$ est stable
par multiplication, c'est-à-dire si $s, s' \in S \Rightarrow ss' \in S$.
\end{definition}

\noindent
Étant donnés un anneau $A$ et une partie multiplicative $S$, nous
définissons une relation sur $A \times S$ comme suit :
$$
(x, s) \sim (y, t)
\Leftrightarrow
\exists u \in S \text{ tel que } (xt-ys)u = 0
$$
On vérifie aisément qu'il s'agit d'une relation d'équivalence.
Notons $x/s$ (ou $\frac{x}{s}$) la classe d'équivalence de $(x, s)$ et
$S^{-1}A$ l'ensemble de toutes les classes d'équivalence. Définissons l'addition
et la multiplication dans $S^{-1}A$ comme suit :
$$
x/s + y/t = (xt + ys)/st, \quad
x/s \cdot y/t = xy/st
$$
On vérifie que ces opérations font de $S^{-1}A$ un anneau.

\begin{definition}
\label{definition-localization}
Cet anneau est appelé la {\it localisation de $A$ relativement à $S$}.
\end{definition}

\noindent
Il existe un morphisme naturel d'anneaux de $A$ vers sa localisation $S^{-1}A$,
$$
A \longrightarrow S^{-1}A, \quad x \longmapsto x/1
$$
parfois appelé le {\it morphisme de localisation}. En général, le
morphisme de localisation n'est pas injectif, à moins que $S$ ne contienne aucun diviseur de zéro.
En effet, si $x/1 = 0$, il existe un $u\in S$ tel que $xu = 0$ dans $A$, et
donc $x = 0$ puisqu'il n'y a pas de diviseur de zéro dans $S$.
La localisation d'un anneau possède la propriété universelle suivante.

\begin{proposition}
\label{proposition-universal-property-localization}
Soit $f : A \to B$ un morphisme d'anneaux qui envoie chaque élément de $S$ sur une unité
de $B$. Il existe alors un unique homomorphisme $g : S^{-1}A \to B$ tel
que le diagramme suivant soit commutatif.
$$
\xymatrix{
A \ar[rr]^{f} \ar[dr] & & B \\
& S^{-1}A \ar[ur]_g
}
$$
\end{proposition}

\begin{proof}
Existence. Définissons une application $g$ comme suit. Pour $x/s\in S^{-1}A$, posons
$g(x/s) = f(x)f(s)^{-1}\in B$. La définition permet de vérifier aisément
qu'il s'agit d'un morphisme d'anneaux bien défini. Il est également clair que
ce morphisme rend le diagramme commutatif.

\medskip\noindent
Unicité. Montrons maintenant que, si $g' : S^{-1}A \to B$
satisfait $g'(x/1) = f(x)$, alors $g = g'$. La commutativité du diagramme donne donc $f(s) = g'(s/1)$ pour
$s \in S$. Mais alors $g'(1/s)f(s) = 1$
dans $B$, ce qui implique $g'(1/s) = f(s)^{-1}$, puis
$g'(x/s) = g'(x/1)g'(1/s) = f(x)f(s)^{-1} = g(x/s)$.
\end{proof}

\begin{lemma}
\label{lemma-localization-zero}
La localisation $S^{-1}A$ est l'anneau nul si et seulement si $0\in S$.
\end{lemma}

\begin{proof}
Si $0\in S$, toute paire $(a, s)\sim (0, 1)$ par définition.
Si $0\not \in S$, alors il est clair que $1/1 \neq 0/1$ dans $S^{-1}A$.
\end{proof}

\begin{lemma}
\label{lemma-localization-and-modules}
Soit $R$ un anneau. Soit $S \subset R$ une partie multiplicative.
La catégorie des $S^{-1}R$-modules est équivalente à la catégorie
des $R$-modules $N$ ayant la propriété que tout $s \in S$ agit sur
$N$ par un automorphisme.
\end{lemma}

\begin{proof}
Le foncteur qui définit l'équivalence associe à un $S^{-1}R$-module
$M$ le même module, considéré maintenant comme un $R$-module au moyen du morphisme de localisation
$R \to S^{-1}R$. Réciproquement, si $N$ est un $R$-module tel que tout
$s \in S$ agisse par un automorphisme $s_N$, nous pouvons considérer $N$ comme un
$S^{-1}R$-module en faisant agir $x/s$ par $x_N  \circ s_N^{-1}$.
Nous omettons de vérifier que ces deux foncteurs sont quasi-inverses l'un de
l'autre.
\end{proof}

\noindent
La notion de localisation d'un anneau se généralise à la
localisation d'un module. Soient $A$ un anneau, $S$ une partie multiplicative
de $A$ et $M$ un $A$-module. Nous définissons une relation sur
$M \times S$ comme suit :
$$
(m, s) \sim (n, t)
\Leftrightarrow
\exists u\in S \text{ tel que } (mt-ns)u = 0
$$
Il s'agit clairement d'une relation d'équivalence. Notons $m/s$ (ou
$\frac{m}{s}$) la classe d'équivalence de $(m, s)$ et $S^{-1}M$
l'ensemble de toutes les classes d'équivalence. Définissons l'addition et la multiplication
scalaire comme suit :
$$
m/s + n/t = (mt + ns)/st,\quad
m/s\cdot n/t = mn/st
$$
Il est clair que cela fait de $S^{-1}M$ un $S^{-1}A$-module.

\begin{definition}
\label{definition-localization-module}
Le $S^{-1}A$-module $S^{-1}M$ est appelé la {\it localisation} de $M$
relativement à $S$.
\end{definition}

\noindent
Remarquons qu'il existe un homomorphisme de $A$-modules $M \to S^{-1}M$, $m \mapsto m/1$,
parfois
appelé le {\it morphisme de localisation}. Il satisfait la propriété
universelle suivante.

\begin{lemma}
\label{lemma-universal-property-localization-module}
Soit $R$ un anneau. Soit $S \subset R$ une partie multiplicative. Soient $M$, $N$
des $R$-modules. Supposons que tous les éléments de $S$ agissent sur $N$ par des automorphismes.
Alors l'application canonique
$$
\Hom_R(S^{-1}M, N) \longrightarrow \Hom_R(M, N)
$$
induite par le morphisme de localisation est un isomorphisme.
\end{lemma}

\begin{proof}
Il est clair que l'application est bien définie et $R$-linéaire.
Injectivité : soit $\alpha \in \Hom_R(S^{-1}M, N)$ et prenons un
élément arbitraire $m/s \in S^{-1}M$. Puisque $s \cdot \alpha(m/s) = \alpha(m/1)$,
nous avons $ \alpha(m/s) =s^{-1}(\alpha (m/1))$ ; ainsi, $\alpha$ est entièrement
déterminé par ses valeurs sur l'image de $M$ dans $S^{-1}M$.
Surjectivité : soit $\beta : M \rightarrow N$ une application R-linéaire donnée.
Nous devons montrer qu'elle peut être ``prolongée'' à $S^{-1}M$. Définissons une application d'ensembles
$$
M \times S \rightarrow N,\quad
(m,s) \mapsto s^{-1}\beta(m)
$$
Cette application respecte clairement la relation d'équivalence ci-dessus ; elle
passe donc au quotient en une application bien définie $\alpha : S^{-1}M \rightarrow N$.
Il reste à montrer que cette application est $R$-linéaire ; prenons donc
$r, r' \in R$, ainsi que $s, s' \in S$ et
$m, m' \in M$. Alors
\begin{align*}
\alpha(r \cdot m/s + r' \cdot m' /s')
& =
\alpha((r \cdot s' \cdot m + r' \cdot s \cdot m') /(ss')) \\
& =
(ss')^{-1}\beta(r \cdot s' \cdot m + r' \cdot s \cdot m') \\
& =
(ss')^{-1} (r \cdot s' \beta (m) + r' \cdot s \beta (m')) \\
& =
r \alpha (m/s) + r' \alpha (m' /s')
\end{align*}
ce qui conclut.
\end{proof}

\begin{example}
\label{example-localize-at-prime}
Soit $A$ un anneau et soit $M$ un $A$-module.
Voici quelques exemples importants de localisations.
\begin{enumerate}
\item Étant donné un idéal premier $\mathfrak p$ de $A$, considérons
$S = A\setminus\mathfrak p$. On vérifie
immédiatement que $S$ est une partie multiplicative. Dans ce cas,
nous notons $A_\mathfrak p$ et $M_\mathfrak p$ les localisations de
$A$ et de $M$ relativement à $S$, respectivement. Elles sont
appelées la {\it localisation de $A$, resp.\ de $M$, en $\mathfrak p$}.
\item Soit $f\in A$. Considérons $S = \{1, f, f^2, \ldots\}$.
C'est clairement une partie multiplicative de $A$.
Dans ce cas, nous notons $A_f$
(resp. $M_f$) la localisation $S^{-1}A$ (resp. $S^{-1}M$).
Elle est appelée la {\it localisation de $A$, resp.\ de $M$, relativement
à $f$}.
Remarquons que $A_f = 0$ si et seulement si $f$ est nilpotent dans $A$.
\item Soit $S = \{f \in A \mid f \text{ n'est pas un diviseur de zéro dans }A\}$.
C'est une partie multiplicative de $A$. Dans ce cas, l'anneau
$Q(A) = S^{-1}A$ est appelé soit l'{\it anneau total des quotients},
soit l'{\it anneau total des fractions}
de $A$.
\item Si $A$ est un anneau intègre, alors l'anneau total des quotients $Q(A)$ est
le corps des fractions de $A$. Voir
Corps, Exemple \ref{fields-example-quotient-field}.
\end{enumerate}
\end{example}

\begin{lemma}
\label{lemma-localization-colimit}
Soit $R$ un anneau.
Soit $S \subset R$ une partie multiplicative.
Soit $M$ un $R$-module.
Alors
$$
S^{-1}M = \colim_{f \in S} M_f
$$
où le préordre sur $S$ est donné par
$f \geq f' \Leftrightarrow f = f'f''$ pour un certain $f'' \in R$,
auquel cas l'application $M_{f'} \to M_f$ est donnée
par $m/(f')^e \mapsto m(f'')^e/f^e$.
\end{lemma}

\begin{proof}
Démonstration omise. Indication : utiliser la propriété universelle du
Lemme \ref{lemma-universal-property-localization-module}.
\end{proof}

\noindent
Dans le paragraphe suivant,
$A$ désigne un anneau,
et $M, N$ désignent des modules sur $A$.

\medskip\noindent
Si $S$ et $S'$ sont des parties multiplicatives de $A$, il est alors
clair que
$$
SS' = \{ss' : s\in S, \ s'\in S'\}
$$
est également une partie multiplicative de $A$. On a alors le résultat suivant.

\begin{proposition}
\label{proposition-localize-twice}
Soit $\overline{S}$ l'image de $S$ dans $S'^{-1}A$ ; alors
$(SS')^{-1}A$ est isomorphe à $\overline{S}^{-1}(S'^{-1}A)$.
\end{proposition}

\begin{proof}
L'application qui envoie $x\in A$ sur $x/1\in (SS')^{-1}A$ induit, par propriété
universelle, une application qui envoie $x/s\in S'^{-1}A$ sur $x/s \in (SS')^{-1}A$.
Les images des éléments de $\overline{S}$ sont inversibles
dans $(SS')^{-1}A$. La propriété universelle fournit une application
$f : \overline{S}^{-1}(S'^{-1}A) \to (SS')^{-1}A$ qui envoie
$(x/s')/(s/1)$ sur $x/ss'$.

\medskip\noindent
D'autre part, l'application de $A$ vers $\overline{S}^{-1}(S'^{-1}A)$
qui envoie $x\in A$ sur $(x/1)/(1/1)$ induit elle aussi, par la propriété universelle,
une application $g : (SS')^{-1}A \to \overline{S}^{-1}(S'^{-1}A)$ qui envoie $x/ss'$
sur $(x/s')/(s/1)$. On vérifie
immédiatement que $f$ et $g$ sont inverses l'un de l'autre ;
ce sont donc tous deux des isomorphismes.
\end{proof}

\noindent
Pour le module $M$, nous avons

\begin{proposition}
\label{proposition-localize-twice-module}
Considérons $S'^{-1}M$ comme un $A$-module ; alors $S^{-1}(S'^{-1}M)$ est
isomorphe à $(SS')^{-1}M$.
\end{proposition}

\begin{proof}
Remarquons qu'étant donné un $A$-module M, nous n'avons démontré aucune
propriété universelle pour $S^{-1}M$. Nous ne pouvons donc pas raisonner
comme dans la démonstration précédente ; nous devons construire explicitement l'isomorphisme.

\medskip\noindent
Définissons les applications comme suit :
\begin{align*}
& f : S^{-1}(S'^{-1}M) \longrightarrow (SS')^{-1}M, \quad \frac{x/s'}{s}\mapsto
x/ss'\\
& g : (SS')^{-1}M \longrightarrow S^{-1}(S'^{-1}M), \quad x/t\mapsto
\frac{x/s'}{s}\ \text{pour certains }s\in S, s'\in S', \text{ et }
t = ss'
\end{align*}
Nous devons vérifier que ces homomorphismes sont bien définis, c'est-à-dire
indépendants du choix de la fraction. Cela se vérifie aisément et il est également
immédiat qu'ils sont inverses l'un de l'autre.
\end{proof}

\noindent
Si $u : M \to N$ est un homomorphisme de $A$-modules, la localisation
induit bien un homomorphisme de $S^{-1}A$-modules $S^{-1}u : S^{-1}M \to
S^{-1}N$ qui envoie $x/s$ sur $u(x)/s$. On vérifie immédiatement que
cette construction est fonctorielle, si bien que $S^{-1}$
est effectivement un foncteur de la catégorie des $A$-modules vers la
catégorie des $S^{-1}A$-modules. De plus, ce foncteur est exact,
comme le montre la proposition suivante.

\begin{proposition}
\label{proposition-localization-exact}
\begin{slogan}
La localisation est exacte.
\end{slogan}
Soit $L\xrightarrow{u} M\xrightarrow{v} N$ une suite exacte
de $R$-modules. Alors
$S^{-1}L \to S^{-1}M \to S^{-1}N$ est également exacte.
\end{proposition}

\begin{proof}
Tout d'abord, il est clair que $S^{-1}L \to S^{-1}M \to S^{-1}N$ est un complexe,
puisque la localisation est un foncteur. Supposons ensuite que $x/s$ s'envoie sur zéro
dans $S^{-1}N$ pour un certain $x/s \in S^{-1}M$. Il existe alors, par définition, un
$t\in S$ tel que $v(xt) = v(x)t = 0$ dans $N$, ce qui signifie que
$xt \in \Ker(v)$. Par l'exactitude de $L \to M \to N$, nous avons
$xt = u(y)$ pour un certain $y$ dans $L$. Ainsi, $x/s$ est l'image de $y/st$.
Cela démontre l'exactitude.
\end{proof}

\begin{lemma}
\label{lemma-localize-quotient-modules}
La localisation respecte les quotients : si $N$ est un sous-module de
$M$, alors $S^{-1}(M/N)\simeq (S^{-1}M)/(S^{-1}N)$.
\end{lemma}

\begin{proof}
La suite exacte
$$
0 \longrightarrow N \longrightarrow M \longrightarrow M/N \longrightarrow 0
$$
donne
$$
0 \longrightarrow S^{-1}N \longrightarrow S^{-1}M
\longrightarrow S^{-1}(M/N) \longrightarrow 0
$$
L'énoncé en résulte.
\end{proof}

\noindent
Si, dans le lemme précédent, nous prenons $N = I$ et $M = A$ pour un idéal $I$ de
$A$, nous voyons que $S^{-1}A/S^{-1}I \simeq S^{-1}(A/I)$ comme $A$-modules. La
proposition suivante montre qu'ils sont isomorphes comme anneaux.

\begin{proposition}
\label{proposition-localize-quotient}
Soient $I$ un idéal de $A$ et $S$ une partie multiplicative de $A$. Alors
$S^{-1}I$ est un idéal de $S^{-1}A$ et $\overline{S}^{-1}(A/I)$ est
isomorphe à $S^{-1}A/S^{-1}I$, où $\overline{S}$ est
l'image de $S$ dans $A/I$.
\end{proposition}

\begin{proof}
Le fait que $S^{-1}I$ soit un idéal est clair puisque $I$ est lui-même un
idéal. Définissons
$$
f : S^{-1}A\longrightarrow \overline{S}^{-1}(A/I), \quad x/s\mapsto
\overline{x}/\overline{s}
$$
où $\overline{x}$ et $\overline{s}$ sont les images de $x$ et de
$s$ dans $A/I$. Nous conserverons des notations semblables dans cette démonstration.
Cette application est bien définie par la propriété universelle de
$S^{-1}A$, et $S^{-1}I$ est contenu dans son noyau ;
elle induit donc une application
$$
\overline{f} : S^{-1}A/S^{-1}I \longrightarrow \overline{S}^{-1}(A/I), \quad
\overline{x/s}\mapsto \overline{x}/\overline{s}
$$

\medskip\noindent
D'autre part, l'application $A \to S^{-1}A/S^{-1}I$ qui envoie $x$ sur
$\overline{x/1}$ induit une application $A/I \to S^{-1}A/S^{-1}I$ qui envoie
$\overline{x}$ sur $\overline{x/1}$. L'image de $\overline{S}$ est
inversible dans $S^{-1}A/S^{-1}I$ ; elle induit donc une application
$$
g : \overline{S}^{-1}(A/I) \longrightarrow S^{-1}A/S^{-1}I, \quad
\frac{\overline{x}}{\overline{s}}\mapsto \overline{x/s}
$$
par la propriété universelle. Il est alors clair que $\overline{f}$ et $g$
sont inverses l'un de l'autre ; ce sont donc tous deux des isomorphismes.
\end{proof}

\noindent
Étudions maintenant le comportement des sous-modules par localisation.

\begin{lemma}
\label{lemma-submodule-localization}
Tout sous-module $N'$ de $S^{-1}M$ est de la forme $S^{-1}N$ pour un certain
$N\subset M$. On peut en effet prendre pour $N$ l'image réciproque de
$N'$ dans $M$.
\end{lemma}

\begin{proof}
Soit $N$ l'image réciproque de $N'$ dans $M$. On voit alors que
$S^{-1}N\supset N'$. Pour montrer qu'ils sont égaux, prenons $x/s$ dans
$S^{-1}N$, où $s\in S$ et $x\in N$. Il s'ensuit que $x/1\in
N'$. Puisque $N'$ est un $S^{-1}R$-sous-module, nous avons
$x/s = x/1\cdot 1/s\in N'$. Cela achève la démonstration.
\end{proof}

\noindent
En prenant $M = A$ et $N = I$, où I est un idéal de $A$, nous obtenons le
corollaire suivant, que l'on peut considérer comme une réciproque de la première partie de la
Proposition \ref{proposition-localize-quotient}.

\begin{lemma}
\label{lemma-ideal-in-localization}
\begin{slogan}
Les idéaux de la localisation d'un anneau sont les localisés d'idéaux.
\end{slogan}
Tout idéal $I'$ de $S^{-1}A$ est de la forme $S^{-1}I$, où l'on peut
prendre pour $I$ l'image réciproque de $I'$ dans $A$.
\end{lemma}

\begin{proof}
Cela résulte immédiatement du Lemme \ref{lemma-submodule-localization}.
\end{proof}









\section{Hom interne}
\label{section-hom}

\noindent
Si $R$ est un anneau et si $M$, $N$ sont des $R$-modules, alors
$$
\Hom_R(M, N) = \{ \varphi : M \to N\}
$$
est l'ensemble des applications $R$-linéaires de $M$ vers $N$. Cet ensemble est muni
d'une structure de groupe abélien en posant, comme d'habitude,
$(\varphi + \psi)(m) = \varphi(m) + \psi(m)$.
En fait, $\Hom_R(M, N)$ est aussi un $R$-module par la règle
$(x \varphi)(m) = x \varphi(m) = \varphi(xm)$.

\medskip\noindent
Étant donnés des homomorphismes $a : M \to M'$ et $b : N \to N'$ de $R$-modules, nous pouvons
précomposer et postcomposer les homomorphismes par $a$ et $b$. On obtient
le diagramme commutatif suivant :
$$
\xymatrix{
\Hom_R(M', N) \ar[d]_{- \circ a} \ar[r]_{b \circ -} &
\Hom_R(M', N') \ar[d]^{- \circ a} \\
\Hom_R(M, N) \ar[r]^{b \circ -} &
\Hom_R(M, N')
}
$$
En fait, les applications de ce diagramme sont des homomorphismes de $R$-modules.
Ainsi, $\Hom_R$ définit un foncteur additif
$$
\text{Mod}_R^{opp} \times \text{Mod}_R \longrightarrow \text{Mod}_R, \quad
(M, N) \longmapsto \Hom_R(M, N)
$$

\begin{lemma}
\label{lemma-hom-exact}
Exactitude et $\Hom_R$. Soit $R$ un anneau. Soient $M_1$, $M_2$, $M_3$
des $R$-modules. Soient $M_1 \to M_2$ et $M_2 \to M_3$ des homomorphismes de $R$-modules.
\begin{enumerate}
\item $M_1 \to M_2 \to M_3 \to 0$ est exacte si et seulement si
$0 \to \Hom_R(M_3, N) \to \Hom_R(M_2, N) \to \Hom_R(M_1, N)$
est exacte pour tout $R$-module $N$.
\item $0 \to M_1 \to M_2 \to M_3$ est exacte si et seulement si
$0 \to \Hom_R(N, M_1) \to \Hom_R(N, M_2) \to \Hom_R(N, M_3)$
est exacte pour tout $R$-module $N$.
\end{enumerate}
\end{lemma}

\begin{proof}
Démonstration omise.
\end{proof}

\begin{lemma}
\label{lemma-hom-from-finitely-presented}
Soit $R$ un anneau. Soit $M$ un $R$-module de présentation finie.
Soit $N$ un $R$-module.
\begin{enumerate}
\item Pour $f \in R$, nous avons
$\Hom_R(M, N)_f = \Hom_{R_f}(M_f, N_f) = \Hom_R(M_f, N_f)$,
\item pour une partie multiplicative $S$ de $R$, nous avons
$$
S^{-1}\Hom_R(M, N) = \Hom_{S^{-1}R}(S^{-1}M, S^{-1}N) =
\Hom_R(S^{-1}M, S^{-1}N).
$$
\end{enumerate}
\end{lemma}

\begin{proof}
L'assertion (1) est un cas particulier de l'assertion (2).
La seconde égalité dans (2) résulte du
Lemme \ref{lemma-localization-and-modules}.
Choisissons une présentation
$$
\bigoplus\nolimits_{j = 1, \ldots, m} R
\longrightarrow
\bigoplus\nolimits_{i = 1, \ldots, n} R
\to M \to 0.
$$
D'après le
Lemme \ref{lemma-hom-exact},
nous obtenons une suite exacte
$$
0 \to
\Hom_R(M, N) \to
\bigoplus\nolimits_{i = 1, \ldots, n} N
\longrightarrow
\bigoplus\nolimits_{j = 1, \ldots, m} N.
$$
En inversant $S$ et en utilisant la Proposition \ref{proposition-localization-exact},
nous obtenons une suite exacte
$$
0 \to
S^{-1}\Hom_R(M, N) \to
\bigoplus\nolimits_{i = 1, \ldots, n} S^{-1}N
\longrightarrow
\bigoplus\nolimits_{j = 1, \ldots, m} S^{-1}N
$$
et le résultat en découle puisque $S^{-1}M$ s'insère dans
une suite exacte
$$
\bigoplus\nolimits_{j = 1, \ldots, m} S^{-1}R
\longrightarrow
\bigoplus\nolimits_{i = 1, \ldots, n} S^{-1}R \to S^{-1}M \to 0
$$
qui induit (d'après le Lemme \ref{lemma-hom-exact})
la suite exacte
$$
0 \to
\Hom_{S^{-1}R}(S^{-1}M, S^{-1}N) \to
\bigoplus\nolimits_{i = 1, \ldots, n} S^{-1}N
\longrightarrow
\bigoplus\nolimits_{j = 1, \ldots, m} S^{-1}N
$$
qui est la même que ci-dessus.
\end{proof}




\section{Caractérisation des modules de type fini et de présentation finie}
\label{section-colim-and-hom}

\noindent
Étant donné un module $N$ sur un anneau $R$, on peut caractériser si
$N$ est ou non un module de type fini ou un module de présentation finie
à l'aide du foncteur $\Hom_R(N, -)$.

\begin{lemma}
\label{lemma-characterize-finite-module-hom}
Soit $R$ un anneau. Soit $N$ un $R$-module. Les assertions suivantes sont équivalentes :
\begin{enumerate}
\item $N$ est un $R$-module de type fini,
\item pour toute colimite filtrante $M = \colim M_i$ de $R$-modules, l'application
$\colim \Hom_R(N, M_i) \to \Hom_R(N, M)$ est injective.
\end{enumerate}
\end{lemma}

\begin{proof}
Supposons (1) et choisissons des générateurs $x_1, \ldots, x_m$ de $N$.
Si $N \to M_i$ est un homomorphisme de modules et si la composée
$N \to M_i \to M$ est nulle, alors, puisque $M = \colim_{i' \geq i} M_{i'}$,
pour tout $j \in \{1, \ldots, m\}$, nous pouvons trouver un $i' \geq i$ tel que
$x_j$ s'envoie sur zéro dans $M_{i'}$. Comme les $x_j$ sont en nombre fini,
nous pouvons trouver un seul $i'$ qui convienne à tous.
La composée $N \to M_i \to M_{i'}$ est alors nulle, et nous en concluons que
l'application est injective, c'est-à-dire que l'assertion (2) vaut.

\medskip\noindent
Supposons (2). Pour un sous-ensemble fini $E \subset N$, notons $N_E \subset N$
le $R$-sous-module engendré par les éléments de $E$. Alors
$0 = \colim N/N_E$ est une colimite filtrante. Nous voyons donc que
$\text{id} : N \to N$ se factorise par $N_E$ pour un certain $E$, c'est-à-dire que $N$
est de type fini.
\end{proof}

\noindent
Pour référence, nous définissons ce que signifie l'existence d'une relation
entre des éléments d'un module.

\begin{definition}
\label{definition-relation}
Soit $R$ un anneau. Soit $M$ un $R$-module.
Soient $n \geq 0$ et $x_i \in M$ pour $i = 1, \ldots, n$.
Une {\it relation} entre $x_1, \ldots, x_n$ dans $M$ est une
suite d'éléments $f_1, \ldots, f_n \in R$ telle que
$\sum_{i = 1, \ldots, n} f_i x_i = 0$.
\end{definition}

\begin{lemma}
\label{lemma-module-colimit-fp}
Soit $R$ un anneau et soit $M$ un $R$-module.
Alors $M$ est la colimite d'un système filtrant
$(M_i, \mu_{ij})$ de $R$-modules
dont tous les $M_i$ sont des $R$-modules de présentation finie.
\end{lemma}

\begin{proof}
Considérons un sous-ensemble fini quelconque $S \subset M$ et une famille finie
de relations $E$ entre les éléments
de $S$. Ainsi, chaque $s \in S$ correspond à un $x_s \in M$ et
chaque $e \in E$ est constitué d'un vecteur
d'éléments $f_{e, s} \in R$ tel que $\sum f_{e, s} x_s = 0$.
Soit $M_{S, E}$ le conoyau de l'application
$$
R^{\# E} \longrightarrow R^{\# S}, \quad
(g_e)_{e\in E} \longmapsto (\sum g_e f_{e, s})_{s\in S}.
$$
Il existe des applications canoniques $M_{S, E} \to M$.
Si $S \subset S'$ et si les éléments de
$E$ correspondent, par cette application, à des relations
de $E'$, il existe alors une application évidente
$M_{S, E} \to M_{S', E'}$ qui commute aux
applications vers $M$. Soit $I$ l'ensemble des paires
$(S, E)$, ordonné par l'inclusion décrite ci-dessus.
Il est clair que la colimite de ce système filtrant est $M$.
\end{proof}

\begin{lemma}
\label{lemma-characterize-finitely-presented-module-hom}
Soit $R$ un anneau. Soit $N$ un $R$-module. Les assertions suivantes sont équivalentes :
\begin{enumerate}
\item $N$ est un $R$-module de présentation finie,
\item pour toute colimite filtrante $M = \colim M_i$ de $R$-modules, l'application
$\colim \Hom_R(N, M_i) \to \Hom_R(N, M)$ est bijective.
\end{enumerate}
\end{lemma}

\begin{proof}
Supposons (1) et choisissons une suite exacte $F_{-1} \to F_0 \to N \to 0$
où les $F_i$ sont libres de type fini. Nous avons alors une suite exacte
$$
0 \to \Hom_R(N, M) \to \Hom_R(F_0, M) \to \Hom_R(F_{-1}, M)
$$
fonctorielle en le $R$-module $M$. Les foncteurs $\Hom_R(F_i, M)$ commutent
aux colimites filtrantes puisque $\Hom_R(R^{\oplus n}, M) = M^{\oplus n}$.
Comme les colimites filtrantes sont exactes
(Lemme \ref{lemma-directed-colimit-exact}),
nous voyons que (2) vaut.

\medskip\noindent
Supposons (2). D'après le Lemme \ref{lemma-module-colimit-fp},
nous pouvons écrire $N = \colim N_i$ comme une colimite
filtrante telle que $N_i$ soit de présentation finie pour tout $i$.
Ainsi, $\text{id}_N$ se factorise par $N_i$ pour un certain $i$.
Cela signifie que $N$ est facteur direct d'un module de
présentation finie sur $R$ (à savoir $N_i$), et est donc de présentation finie.
\end{proof}












\section{Produits tensoriels}
\label{section-tensor-product}

\begin{definition}
\label{definition-bilinear}
Soient $R$ un anneau et $M, N, P$ trois $R$-modules.
Une application $f : M \times N \to P$ (où $M \times N$
est considéré uniquement comme le produit cartésien de deux $R$-modules) est dite
{\it $R$-bilinéaire} si, pour tout $x \in M$,
l'application $y\mapsto f(x, y)$ de $N$ dans $P$ est $R$-linéaire et si, pour tout
$y\in N$, l'application $x\mapsto f(x, y)$ est également $R$-linéaire.
\end{definition}

\begin{lemma}
\label{lemma-tensor-product}
Soient $M, N$ des $R$-modules. Il existe alors une paire $(T, g)$,
où $T$ est un $R$-module et
$g : M \times N \to T$ une application $R$-bilinéaire,
qui possède la propriété universelle suivante :
pour tout $R$-module $P$ et toute application $R$-bilinéaire
$f : M \times N \to P$, il
existe une unique application $R$-linéaire
$\tilde{f} : T \to P$ telle que $f = \tilde{f} \circ g$.
Autrement dit, le diagramme suivant est commutatif :
$$
\xymatrix{
M \times N \ar[rr]^f \ar[dr]_g & & P\\
& T \ar[ur]_{\tilde f}
}
$$
De plus, si $(T, g)$ et $(T', g')$
sont deux paires possédant cette propriété, il
existe un unique isomorphisme
$j : T \to T'$ tel que $j\circ g = g'$.
\end{lemma}

\noindent
Le $R$-module $T$ qui satisfait à la propriété universelle ci-dessus est appelé
le {\it produit tensoriel} des $R$-modules $M$ et $N$, et est noté
$M \otimes_R N$.

\begin{proof}
Démontrons d'abord l'existence d'un tel $R$-module $T$.
Soient $M, N$ des $R$-modules.
Soit $T$ le module quotient
$P/Q$, où $P$ est le $R$-module libre $R^{(M \times N)}$ et où $Q$ est le
$R$-module engendré par tous les éléments
des types suivants : ($x\in M, y\in N$)
\begin{align*}
(x + x', y) - (x, y) - (x', y), \\
(x, y + y') - (x, y) - (x, y'), \\
(ax, y) - a(x, y), \\
(x, ay) - a(x, y)
\end{align*}
Notons $\pi : M \times N \to T$ l'application canonique.
Cette application est $R$-bilinéaire, comme
le montrent les relations ci-dessus
lorsque l'on vérifie les conditions de bilinéarité. Notons l'image
$\pi(x, y) = x \otimes
y$ ; ces éléments engendrent alors
$T$. Soit maintenant $f : M \times N \to P$ une application $R$-bilinéaire ;
nous pouvons définir
$f' : T \to P$ en prolongeant l'application
$f'(x \otimes y) = f(x, y)$. Il est clair que $f = f'\circ \pi$. De plus, $f'$ est
uniquement déterminée par ses valeurs sur le
système générateur $\{x \otimes y : x\in M, y\in N\}$.
Supposons qu'il existe une autre paire $(T', g')$ satisfaisant aux mêmes propriétés.
Il existe alors un unique $j : T \to T'$, ainsi
qu'un $j' : T' \to T$, tels que $g' = j\circ g$, $g = j'\circ g'$.
Mais les deux applications $(j\circ j') \circ g$ et $g$
satisfont alors aux propriétés universelles ; par unicité, elles sont donc égales,
et par conséquent $j'\circ j$ est l'identité sur $T$.
De même, $(j'\circ j) \circ g' = g'$ et $j\circ j'$ est l'identité sur $T'$.
Ainsi, $j$ est un isomorphisme.
\end{proof}

\begin{lemma}
\label{lemma-flip-tensor-product}
Soient $M, N, P$ des $R$-modules ; les applications bilinéaires
\begin{align*}
(x, y) & \mapsto y \otimes x\\
(x + y, z) & \mapsto x \otimes z + y \otimes z\\
(r, x) & \mapsto rx
\end{align*}
induisent des isomorphismes uniques
\begin{align*}
M \otimes_R N & \to N \otimes_R M, \\
(M\oplus N)\otimes_R P & \to (M \otimes_R P)\oplus(N \otimes_R P),  \\
R \otimes_R M & \to M
\end{align*}
\end{lemma}

\begin{proof}
Démonstration omise.
\end{proof}

\noindent
Nous pouvons généraliser le produit tensoriel de deux $R$-modules à un nombre fini de
$R$-modules et établir une
correspondance entre le produit multitensoriel et les applications multilinéaires.
Une construction presque identique
permet de démontrer le résultat suivant :

\begin{lemma}
\label{lemma-multilinear}
Soient $M_1, \ldots, M_r$ des $R$-modules. Il existe alors une paire $(T, g)$
constituée d'un $R$-module T et d'une application $R$-multilinéaire
$g : M_1\times \ldots \times M_r \to T$ possédant la propriété
universelle suivante : pour toute application $R$-multilinéaire
$f : M_1\times \ldots \times M_r \to P$, il existe un unique homomorphisme de $R$-modules
$f' : T \to P$ tel que $f'\circ g = f$.
Un tel module $T$ est unique à isomorphisme unique près. Nous le notons
$M_1\otimes_R \ldots \otimes_R M_r$ et nous notons l'application
multilinéaire universelle $(m_1, \ldots, m_r) \mapsto m_1 \otimes \ldots \otimes m_r$.
\end{lemma}

\begin{proof}
Démonstration omise.
\end{proof}

\begin{lemma}
\label{lemma-transitive}
Les homomorphismes
$$
(M \otimes_R N)\otimes_R P \to
M \otimes_R N \otimes_R P \to
M \otimes_R (N \otimes_R P)
$$
tels que
$f((x \otimes y)\otimes z) = x \otimes y \otimes z$
et $g(x \otimes y \otimes z) = x \otimes (y \otimes z)$,
$x\in M, y\in N, z\in P$, sont bien définis et sont des isomorphismes.
\end{lemma}

\begin{proof}
Nous allons montrer que $f$ est bien défini et est un isomorphisme ; la démonstration
s'applique de manière analogue à $g$. Fixons un
$z\in P$ quelconque. L'application $(x, y)\mapsto x \otimes y \otimes z$,
$x\in M, y\in N$, est alors $R$-bilinéaire en $x$ et $y$,
et induit donc un homomorphisme $f_z : M \otimes N \to M \otimes N \otimes P$
qui vérifie
$f_z(x \otimes y) = x \otimes y \otimes z$.
Considérons ensuite $(M \otimes N)\times P \to M \otimes N \otimes P$ donnée par
$(w, z)\mapsto f_z(w)$. Cette application est
$R$-bilinéaire et induit donc
$f : (M \otimes_R N)\otimes_R P \to M \otimes_R N \otimes_R P$,
avec $f((x \otimes y)\otimes z) = x \otimes y \otimes z$.
Pour construire l'inverse, remarquons que l'application
$\pi : M \times N \times P \to (M \otimes N)\otimes P$ est
$R$-trilinéaire.
Elle induit donc une application $R$-linéaire
$h : M \otimes N \otimes P \to (M \otimes N)\otimes P$ qui
s'accorde avec la propriété universelle. Nous voyons ici que
$h(x \otimes y \otimes z) = (x \otimes y)\otimes z$.
D'après les expressions explicites de $f$ et de $h$, $f\circ h$ et $h\circ f$ sont
les applications identiques de $M \otimes N \otimes
P$ et de $(M \otimes N)\otimes P$, respectivement ; ainsi, $f$ est bien
l'isomorphisme recherché.
\end{proof}

\noindent
En raisonnant par récurrence, nous voyons que cela s'étend aux produits multitensoriels. En combinant
ce résultat avec le Lemme \ref{lemma-flip-tensor-product}, nous voyons que
l'opération de produit tensoriel sur la catégorie des $R$-modules est associative,
commutative et distributive.

\begin{definition}
\label{definition-bimodule}
Un groupe abélien $N$ est appelé un {\it $(A, B)$-bimodule} s'il est à la fois un
$A$-module et un $B$-module et si, pour tous $a \in A$ et $b \in B$, les
multiplications par $a$ et par $b$ commutent, de sorte que $b(an) = a(bn)$ pour tout $n \in N$.
Dans cette situation, nous écrivons habituellement l'action de $B$ à droite : ainsi,
pour $b \in B$ et $n \in N$, le résultat de la multiplication de $n$ par $b$
est noté $nb$. Avec cette convention, la compatibilité ci-dessus s'écrit
$(ax)b = a(xb)$ pour tous $a\in A, b\in B, x\in N$.
La notation abrégée $_AN_B$ désigne un $(A, B)$-bimodule $N$.
\end{definition}

\begin{lemma}
\label{lemma-tensor-with-bimodule}
Pour un $A$-module $M$, un $B$-module $P$ et un $(A, B)$-bimodule $N$, les modules
$(M \otimes_A N)\otimes_B P$ et $M \otimes_A(N \otimes_B P)$ peuvent tous deux être
munis d'une structure de $(A, B)$-bimodule,
et de plus
$$
(M \otimes_A N)\otimes_B P \cong M \otimes_A(N \otimes_B P).
$$
\end{lemma}

\begin{proof}
A priori, $M \otimes_A N$ est un $A$-module, mais nous pouvons le munir d'une
structure de $B$-module en posant
$$
(x \otimes y)b = x \otimes yb, \quad x\in M, y\in N, b\in B
$$
Ainsi, $M \otimes_A N$ devient un $(A, B)$-bimodule. Il en va de même pour
$N \otimes_B P$, et donc pour
$(M \otimes_A N)\otimes_B P$ et $M \otimes_A(N \otimes_B P)$. D'après le
Lemme \ref{lemma-transitive}, ces deux
modules sont isomorphes à la fois comme $A$-modules et comme $B$-modules par la même
application.
\end{proof}

\begin{lemma}
\label{lemma-hom-from-tensor-product}
Pour trois $R$-modules quelconques $M, N, P$,
$$
\Hom_R(M \otimes_R N, P) \cong \Hom_R(M, \Hom_R(N, P))
$$
\end{lemma}

\begin{proof}
Une application $R$-linéaire $\hat{f}\in \Hom_R(M \otimes_R N, P)$ correspond à une
application $R$-bilinéaire $f : M \times N \to P$. Pour
tout $x\in M$, l'application $y\mapsto f(x, y)$ est $R$-linéaire par la propriété
universelle. Ainsi, $f$ correspond à une
application $\phi_f : M \to \Hom_R(N, P)$. Cette application est $R$-linéaire puisque
$$
\phi_f(ax + y)(z) =
f(ax + y, z) = af(x, z)+f(y, z) =
(a\phi_f(x)+\phi_f(y))(z),
$$
pour tous $a \in R$, $x \in M$, $y \in M$ et
$z \in N$. Réciproquement, tout
$f \in \Hom_R(M, \Hom_R(N, P))$ définit une application $R$-bilinéaire
$M \times N \to P$, à savoir $(x, y)\mapsto f(x)(y)$.
Nous obtenons donc une bijection naturelle entre les
deux modules
$\Hom_R(M \otimes_R N, P)$ et $\Hom_R(M, \Hom_R(N, P))$.
\end{proof}

\begin{lemma}[Les produits tensoriels commutent aux colimites]
\label{lemma-tensor-products-commute-with-limits}
Soit $(M_i, \mu_{ij})$ un système sur l'ensemble préordonné $I$.
Soit $N$ un $R$-module. Alors
$$
\colim (M_i \otimes N) \cong (\colim M_i)\otimes N.
$$
De plus, l'isomorphisme est induit par les homomorphismes
$\mu_i \otimes 1: M_i \otimes N \to M \otimes N$,
où $M = \colim_i M_i$ et où $\mu_i : M_i \to M$ sont les applications canoniques.
\end{lemma}

\begin{proof}
Première démonstration. Le foncteur $M' \mapsto M' \otimes_R N$ est adjoint à gauche
du foncteur $N' \mapsto \Hom_R(N, N')$ d'après le
Lemme \ref{lemma-hom-from-tensor-product}. Ainsi, $M' \mapsto M' \otimes_R N$
commute à toutes les colimites ; voir
Catégories, Lemme \ref{categories-lemma-adjoint-exact}.

\medskip\noindent
Seconde démonstration directe. Soit $P = \colim (M_i \otimes N)$, muni des coprojections
$\lambda_i : M_i \otimes N \to P$. Soit $M = \colim M_i$, muni des coprojections
$\mu_i : M_i \to M$. Alors, pour tous $i\leq j$, le diagramme suivant est commutatif :
$$
\xymatrix{
M_i \otimes N \ar[r]_{\mu_i \otimes 1} \ar[d]_{\mu_{ij} \otimes 1} &
M \otimes N \ar[d]^{\text{id}} \\
M_j \otimes N \ar[r]^{\mu_j \otimes 1} &
M \otimes N
}
$$
D'après le Lemme \ref{lemma-homomorphism-limit}, ces applications induisent un unique homomorphisme
$\psi : P \to M \otimes N$ tel que
$\mu_i \otimes 1 = \psi \circ \lambda_i$.

\medskip\noindent
Pour construire l'application inverse, il existe, pour tout $i\in I$, l'application canonique
$R$-bilinéaire $g_i : M_i \times N \to
M_i \otimes N$. Celle-ci induit une unique application
$\widehat{\phi} : M \times N \to P$
telle que $\widehat{\phi} \circ (\mu_i \times 1) = \lambda_i \circ g_i$.
Elle est $R$-bilinéaire et induit donc une
application $R$-linéaire $\phi : M \otimes N \to P$.
Le diagramme commutatif ci-dessous :
$$
\xymatrix{
M_i \times N \ar[r]^{g_i} \ar[d]^{\mu_i \times \text{id}} &
M_i \otimes N\ar[r]_{\text{id}} \ar[d]_{\lambda_i} &
M_i \otimes N \ar[d]_{\mu_i \otimes \text{id}} \ar[rd]^{\lambda_i} \\
M \times N \ar[r]^{\widehat{\phi}} &
P \ar[r]^{\psi} & M \otimes N \ar[r]^{\phi} & P
}
$$
montre que $\psi\circ\widehat{\phi} = g$, où g est l'application $R$-bilinéaire canonique
$g : M \times N \to M \otimes N$. Ainsi,
$\psi\circ\phi$ est l'identité de $M \otimes N$. Le carré et le
triangle de droite montrent que $\phi\circ\psi$ est également
l'identité de $P$.
\end{proof}

\begin{lemma}
\label{lemma-tensor-product-exact}
Soit
\begin{align*}
M_1\xrightarrow{f} M_2\xrightarrow{g} M_3 \to 0
\end{align*}
une suite exacte de $R$-modules et d'homomorphismes, et soit $N$ un
$R$-module quelconque. Alors la suite
\begin{equation}
\label{equation-2ndex}
M_1\otimes N\xrightarrow{f \otimes 1} M_2\otimes N \xrightarrow{g \otimes 1}
M_3\otimes N \to 0
\end{equation}
est exacte. Autrement dit, le foncteur $- \otimes_R N$ est
{\it exact à droite}, au sens où tensoriser
chaque terme de la suite exacte à droite initiale préserve l'exactitude.
\end{lemma}

\begin{proof}
Pour tout $R$-module $P$,
appliquons le foncteur $\Hom(-, \Hom(N, P))$ à la première suite
exacte. Nous obtenons
$$
0 \to
\Hom(M_3, \Hom(N, P)) \to
\Hom(M_2, \Hom(N, P)) \to
\Hom(M_1, \Hom(N, P))
$$
qui est exacte d'après le Lemme \ref{lemma-hom-exact} (1).
D'après le Lemme \ref{lemma-hom-from-tensor-product}, cette suite devient
$$
0 \to \Hom(M_3 \otimes N, P) \to
\Hom(M_2 \otimes N, P) \to \Hom(M_1 \otimes N, P)
$$
et est donc elle aussi exacte. En appliquant de nouveau le
Lemme \ref{lemma-hom-exact} (1), nous obtenons la suite exacte recherchée.
\end{proof}

\begin{remark}
\label{remark-tensor-product-not-exact}
Cependant, le produit tensoriel NE préserve PAS les suites exactes en général.
Autrement dit, si $M_1 \to M_2 \to M_3$ est
exacte, il n'est pas nécessaire que
$M_1 \otimes N \to M_2 \otimes N \to M_3 \otimes N$
soit exacte pour un $R$-module $N$ arbitraire.
\end{remark}

\begin{example}
\label{example-tensor-product-not-exact}
Considérons l'application injective $2 : \mathbf{Z}\to \mathbf{Z}$
comme une application de $\mathbf{Z}$-modules.
Soit $N = \mathbf{Z}/2$. L'application induite
$\mathbf{Z} \otimes \mathbf{Z}/2 \to \mathbf{Z} \otimes \mathbf{Z}/2$
n'est PAS injective. En effet, pour
$x \otimes y\in \mathbf{Z} \otimes \mathbf{Z}/2$,
$$
(2 \otimes 1)(x \otimes y) = 2x \otimes y = x \otimes 2y = x \otimes 0 = 0
$$
L'application induite est donc l'application nulle, tandis que $\mathbf{Z} \otimes N\neq 0$.
\end{example}

\begin{remark}
\label{remark-flat-module}
Pour un $R$-module $N$, si le
foncteur $-\otimes_R N$ est exact, c'est-à-dire si la tensorisation
par $N$ préserve toutes les suites
exactes, alors $N$ est appelé un $R$-module {\it plat}.
Nous étudierons cette notion plus loin dans la section \ref{section-flat}.
\end{remark}

\begin{lemma}
\label{lemma-tensor-finiteness}
Soit $R$ un anneau. Soient $M$ et $N$ des $R$-modules.
\begin{enumerate}
\item Si $N$ et $M$ sont de type fini, alors $M \otimes_R N$ l'est aussi.
\item Si $N$ et $M$ sont de présentation finie, alors $M \otimes_R N$ l'est aussi.
\end{enumerate}
\end{lemma}

\begin{proof}
Supposons $M$ de type fini. Choisissons alors une présentation
$0 \to K \to R^{\oplus n} \to M \to 0$. Elle donne une suite exacte
$K \otimes_R N \to N^{\oplus n} \to M \otimes_R N \to 0$ d'après le
Lemme \ref{lemma-tensor-product-exact}.
Nous en concluons que, si $N$ est également de type fini, alors $M \otimes_R N$
est un quotient d'un module de type fini, donc est de type fini ; voir le
Lemme \ref{lemma-extension}.
De même, si $N$ et $M$ sont tous deux de présentation finie,
nous voyons que $K$ est de type fini et que $M \otimes_R N$
est le quotient du module de présentation finie $N^{\oplus n}$ par
un module de type fini, à savoir $K \otimes_R N$, et est donc de présentation finie ; voir le
Lemme \ref{lemma-extension}.
\end{proof}

\begin{lemma}
\label{lemma-tensor-localization}
Soit $M$ un $R$-module. Alors les $S^{-1}R$-modules $S^{-1}M$
et $S^{-1}R \otimes_R M$ sont canoniquement isomorphes, et
l'isomorphisme canonique $f : S^{-1}R \otimes_R M \to S^{-1}M$
est donné par
$$
f((a/s) \otimes m) = am/s, \forall a \in R, m \in M, s \in S
$$
\end{lemma}

\begin{proof}
Il est évident que l'application
$f' : S^{-1}R \times M \to S^{-1}M$ donnée par $f'(a/s, m) = am/s$ est
bilinéaire ; par la
propriété universelle, elle induit donc un unique homomorphisme de $S^{-1}R$-modules
$f : S^{-1}R \otimes_R M \to S^{-1}M$ comme dans l'énoncé du lemme.
En fait, tout élément de $S^{-1}M$ est de la forme $m/s$, $m\in M, s\in S$, et
tout élément de
$S^{-1}R \otimes_R M$ est de la forme $1/s \otimes m$. Pour voir ce dernier fait,
écrivons un élément de
$S^{-1}R \otimes_R M$ sous la forme
$$
\sum_k \frac{a_k}{s_k} \otimes m_k =
\sum_k \frac{a_k t_k}{s} \otimes m_k =
\frac{1}{s} \otimes \sum_k {a_k t_k}m_k = \frac{1}{s} \otimes m
$$
où $m = \sum_k {a_k t_k}m_k$. Il est alors évident que $f$ est surjectif,
et si $f(\frac{1}{s} \otimes m) = m/s = 0$, il existe $t \in S$ tel que
$tm = 0$ dans $M$. Nous avons alors
$$
\frac{1}{s} \otimes m = \frac{1}{st} \otimes tm = \frac{1}{st} \otimes 0 = 0
$$
Par conséquent, $f$ est injectif.
\end{proof}

\begin{lemma}
\label{lemma-tensor-product-localization}
Soient $M, N$ des $R$-modules. Il existe alors un
isomorphisme canonique de $S^{-1}R$-modules
$f : S^{-1}M \otimes_{S^{-1}R}S^{-1}N \to S^{-1}(M \otimes_R N)$,
donné par
$$
f((m/s)\otimes(n/t)) = (m \otimes n)/st
$$
\end{lemma}

\begin{proof}
Nous pouvons utiliser à plusieurs reprises le Lemme \ref{lemma-tensor-with-bimodule}
et le Lemme \ref{lemma-tensor-localization} pour
voir que ces deux
$S^{-1}R$-modules sont isomorphes, en remarquant que $S^{-1}R$ est un
$(R, S^{-1}R)$-bimodule :
\begin{align*}
S^{-1}(M \otimes_R N) & \cong S^{-1}R \otimes_R (M \otimes_R N)\\
 & \cong S^{-1}M \otimes_R N\\
 & \cong (S^{-1}M \otimes_{S^{-1}R}S^{-1}R)\otimes_R N\\
 & \cong S^{-1}M \otimes_{S^{-1}R}(S^{-1}R \otimes_R N)\\
 & \cong S^{-1}M \otimes_{S^{-1}R}S^{-1}N
\end{align*}
On vérifie aisément que cet isomorphisme est bien celui énoncé dans le lemme.
\end{proof}

\section{Algèbre tensorielle}
\label{section-tensor-algebra}

\noindent
Soit $R$ un anneau. Soit $M$ un $R$-module.
Nous définissons l'{\it algèbre tensorielle de $M$ sur $R$}
comme la $R$-algèbre non commutative
$$
\text{T}(M) = \text{T}_R(M) =
\bigoplus\nolimits_{n \geq 0} \text{T}^n(M)
$$
où
$\text{T}^0(M) = R$,
$\text{T}^1(M) = M$,
$\text{T}^2(M) = M \otimes_R M$,
$\text{T}^3(M) = M \otimes_R M \otimes_R M$, et ainsi de suite.
La multiplication est définie par la règle suivante sur les tenseurs purs :
$$
(x_1 \otimes x_2 \otimes \ldots \otimes x_n)
\cdot
(y_1 \otimes y_2 \otimes \ldots \otimes y_m)
=
x_1 \otimes x_2 \otimes \ldots \otimes x_n \otimes
y_1 \otimes y_2 \otimes \ldots \otimes y_m
$$
et nous l'étendons par linéarité.

\medskip\noindent
Nous définissons l'{\it algèbre extérieure $\wedge(M)$ de $M$ sur $R$}
comme le quotient de $\text{T}(M)$ par l'idéal bilatère
engendré par les éléments $x \otimes x \in \text{T}^2(M)$.
L'image d'un tenseur pur $x_1 \otimes \ldots \otimes x_n$
dans $\wedge^n(M)$ est notée $x_1 \wedge \ldots \wedge x_n$.
Ces éléments engendrent $\wedge^n(M)$, ils dépendent $R$-linéairement
de chaque $x_i$ et ils sont nuls lorsque deux des $x_i$ sont égaux
(autrement dit, ils sont alternés en
$x_1, x_2, \ldots, x_n$). La multiplication sur $\wedge(M)$ est
commutative au sens gradué, c'est-à-dire que tous $x \in M$ et $y \in M$
vérifient $x \wedge y = - y \wedge x$.

\medskip\noindent
Considérons le cas où $M = Rx_1 \oplus \ldots \oplus Rx_n$
est un module libre de type fini. Alors $\wedge(M)$ est libre sur
$R$, de base formée des éléments
$$
x_{i_1} \wedge \ldots \wedge x_{i_r}
$$
où $0 \leq r \leq n$ et $1 \leq i_1 < i_2 < \ldots < i_r \leq n$.

\medskip\noindent
Nous définissons l'{\it algèbre symétrique $\text{Sym}(M)$ de $M$ sur $R$}
comme le quotient de $\text{T}(M)$ par l'idéal bilatère
engendré par les éléments $x \otimes y - y \otimes x \in \text{T}^2(M)$.
L'image d'un tenseur pur $x_1 \otimes \ldots \otimes x_n$
dans $\text{Sym}^n(M)$ est simplement notée $x_1 \ldots x_n$.
Ces éléments engendrent $\text{Sym}^n(M)$, ils dépendent $R$-linéairement
de chaque $x_i$, et $x_1 \ldots x_n = x_1' \ldots x_n'$ si la
suite d'éléments $x_1, \ldots, x_n$ est une permutation de la
suite $x_1', \ldots, x_n'$. Ainsi, $\text{Sym}(M)$
est commutative.

\medskip\noindent
Considérons le cas où $M = Rx_1 \oplus \ldots \oplus Rx_n$
est un module libre de type fini. Alors
$\text{Sym}(M) = R[x_1, \ldots, x_n]$ est une algèbre polynomiale.

\begin{lemma}
\label{lemma-free-tensor-algebra}
Soit $R$ un anneau. Soit $M$ un $R$-module.
Si $M$ est un $R$-module libre, chacune de ses puissances symétriques
et extérieures l'est aussi.
\end{lemma}

\begin{proof}
La démonstration est omise ; voir toutefois ci-dessus le cas libre de type fini.
\end{proof}

\begin{lemma}
\label{lemma-presentation-sym-exterior}
Soit $R$ un anneau.
Soit $M_2 \to M_1 \to M \to 0$ une suite exacte de $R$-modules.
Il existe des suites exactes
$$
M_2 \otimes_R \text{Sym}^{n - 1}(M_1)
\to
\text{Sym}^n(M_1)
\to
\text{Sym}^n(M)
\to
0
$$
et, de même,
$$
M_2 \otimes_R \wedge^{n - 1}(M_1)
\to
\wedge^n(M_1)
\to
\wedge^n(M)
\to
0
$$
\end{lemma}

\begin{proof}
La démonstration est omise.
\end{proof}

\begin{lemma}
\label{lemma-present-sym-wedge}
Soit $R$ un anneau.
Soit $M$ un $R$-module.
Soit $x_i$, $i \in I$, une famille donnée de générateurs de
$M$ comme $R$-module. Soit $n \geq 2$.
Il existe une suite exacte canonique
$$
\bigoplus_{1 \leq j_1 < j_2 \leq n}
\bigoplus_{i_1, i_2 \in I}
\text{T}^{n - 2}(M)
\oplus
\bigoplus_{1 \leq j_1 < j_2 \leq n}
\bigoplus_{i \in I}
\text{T}^{n - 2}(M)
\to
\text{T}^n(M)
\to
\wedge^n(M)
\to
0
$$
où le tenseur pur $m_1 \otimes \ldots \otimes m_{n - 2}$ du premier
facteur direct est envoyé sur
\begin{align*}
\underbrace{
m_1 \otimes \ldots \otimes x_{i_1} \otimes \ldots
\otimes x_{i_2} \otimes \ldots \otimes m_{n - 2}
}_{\text{avec } x_{i_1} \text{ et } x_{i_2}
\text{ occupant les places } j_1 \text{ et } j_2
\text{ dans le tenseur}} \\
+
\underbrace{
m_1 \otimes \ldots \otimes x_{i_2} \otimes \ldots
\otimes x_{i_1} \otimes \ldots \otimes m_{n - 2}
}_{\text{avec } x_{i_2} \text{ et } x_{i_1}
\text{ occupant les places } j_1 \text{ et } j_2
\text{ dans le tenseur}}
\end{align*}
et $m_1 \otimes \ldots \otimes m_{n - 2}$ du second
facteur direct est envoyé sur
$$
\underbrace{
m_1 \otimes \ldots \otimes x_i \otimes \ldots
\otimes x_i \otimes \ldots \otimes m_{n - 2}
}_{\text{avec } x_{i} \text{ et } x_{i}
\text{ occupant les places } j_1 \text{ et } j_2
\text{ dans le tenseur}}
$$
Il existe aussi une suite exacte canonique
$$
\bigoplus_{1 \leq j_1 < j_2 \leq n}
\bigoplus_{i_1, i_2 \in I}
\text{T}^{n - 2}(M)
\to
\text{T}^n(M)
\to
\text{Sym}^n(M)
\to
0
$$
où le tenseur pur $m_1 \otimes \ldots \otimes m_{n - 2}$ est envoyé sur
\begin{align*}
\underbrace{
m_1 \otimes \ldots \otimes x_{i_1} \otimes \ldots
\otimes x_{i_2} \otimes \ldots \otimes m_{n - 2}
}_{\text{avec } x_{i_1} \text{ et } x_{i_2}
\text{ occupant les places } j_1 \text{ et } j_2
\text{ dans le tenseur}} \\
-
\underbrace{
m_1 \otimes \ldots \otimes x_{i_2} \otimes \ldots
\otimes x_{i_1} \otimes \ldots \otimes m_{n - 2}
}_{\text{avec } x_{i_2} \text{ et } x_{i_1}
\text{ occupant les places } j_1 \text{ et } j_2
\text{ dans le tenseur}}
\end{align*}
\end{lemma}

\begin{proof}
La démonstration est omise.
\end{proof}

\begin{lemma}
\label{lemma-present-wedge}
Soit $A \to B$ un morphisme d'anneaux. Soit $M$ un $B$-module.
Soit $n > 1$. Le noyau de l'application $A$-linéaire
$M \otimes_A \ldots \otimes_A M \to \wedge^n_B(M)$
est engendré comme $A$-module par les éléments
$m_1 \otimes \ldots \otimes m_n$ tels que $m_i = m_j$
pour $i \not = j$, avec $m_1, \ldots, m_n \in M$, et par les éléments
$m_1 \otimes \ldots \otimes bm_i \otimes \ldots \otimes m_n -
m_1 \otimes \ldots \otimes bm_j \otimes \ldots \otimes m_n$
pour $i \not = j$, $m_1, \ldots, m_n \in M$ et $b \in B$.
\end{lemma}

\begin{proof}
La démonstration est omise.
\end{proof}

\begin{lemma}
\label{lemma-colimit-tensor-algebra}
\begin{slogan}
La formation des algèbres tensorielles commute aux colimites filtrantes.
\end{slogan}
Soit $R$ un anneau. Supposons que les $M_i$ forment un système filtrant de
$R$-modules. Alors
$\colim_i \text{T}(M_i) = \text{T}(\colim_i M_i)$,
et il en va de même pour les algèbres symétriques et extérieures.
\end{lemma}

\begin{proof}
La démonstration est omise. Indication : appliquer le lemme \ref{lemma-tensor-products-commute-with-limits}.
\end{proof}

\begin{lemma}
\label{lemma-tensor-algebra-localization}
Soit $R$ un anneau et soit $S \subset R$ une partie multiplicative.
Alors $S^{-1}T_R(M) = T_{S^{-1}R}(S^{-1}M)$ pour tout $R$-module $M$.
Il en va de même pour les algèbres symétriques et extérieures.
\end{lemma}

\begin{proof}
La démonstration est omise. Indication : appliquer le lemme \ref{lemma-tensor-product-localization}.
\end{proof}

\section{Changement de base}
\label{section-base-change}

\noindent
Nous introduisons formellement le changement de base en algèbre comme suit.

\begin{definition}
\label{definition-base-change}
Soit $\varphi : R \to S$ un morphisme d'anneaux. Soit $M$ un $S$-module.
Soit $R \to R'$ un morphisme d'anneaux quelconque. Le {\it changement de base}
de $\varphi$ par $R \to R'$ est le morphisme d'anneaux $R' \to S \otimes_R R'$.
Dans cette situation, nous écrivons souvent $S' = S \otimes_R R'$.
Le {\it changement de base} du $S$-module $M$ est le $S'$-module
$M \otimes_R R'$.
\end{definition}

\noindent
Si $S = R[x_i]/(f_j)$ pour une famille de variables $x_i$, $i \in I$,
et une famille de polynômes $f_j \in R[x_i]$, $j \in J$, alors
$S \otimes_R R' = R'[x_i]/(f'_j)$, où $f'_j \in R'[x_i]$ est l'image
de $f_j$ par le morphisme $R[x_i] \to R'[x_i]$ induit par $R \to R'$.
Cette remarque simple est la clé de la compréhension du changement de base.

\begin{lemma}
\label{lemma-base-change-finiteness}
Soit $R \to S$ un morphisme d'anneaux. Soit $M$ un $S$-module.
Soit $R \to R'$ un morphisme d'anneaux, et soient
$S' = S \otimes_R R'$ et $M' = M \otimes_R R'$ les changements de base.
\begin{enumerate}
\item Si $M$ est un $S$-module de type fini, alors son changement de base
$M'$ est un $S'$-module de type fini.
\item Si $M$ est un $S$-module de présentation finie, alors son
changement de base $M'$ est un $S'$-module de présentation finie.
\item Si $R \to S$ est de type fini, alors son changement de base
$R' \to S'$ est de type fini.
\item Si $R \to S$ est de présentation finie, alors son
changement de base $R' \to S'$ est de présentation finie.
\end{enumerate}
\end{lemma}

\begin{proof}
Démonstration de (1). Choisissons une application $S$-linéaire surjective
$S^{\oplus n} \to M \to 0$.
Par les lemmes \ref{lemma-flip-tensor-product} et \ref{lemma-tensor-product-exact},
sa tensorisation par $R^\prime$ donne une surjection
${S^\prime}^{\oplus n} \to M^\prime \rightarrow 0$,
de sorte que $M^\prime$ est un $S^\prime$-module de type fini.
Démonstration de (2). Choisissons une présentation
$S^{\oplus m} \to S^{\oplus n} \to M \to 0$.
Par les lemmes \ref{lemma-flip-tensor-product} et \ref{lemma-tensor-product-exact},
sa tensorisation par $R^\prime$ donne une présentation finie
${S^\prime}^{\oplus m} \to {S^\prime}^{\oplus n} \to M^\prime \to 0$
du $S^\prime$-module $M^\prime$. Démonstration de (3). Cela résulte
de la remarque précédant le lemme, puisque nous pouvons supposer $I$ fini.
Démonstration de (4). Cela résulte de la remarque précédant le lemme,
puisque nous pouvons supposer $I$ et $J$ finis.
\end{proof}

\noindent
Soit $\varphi : R \to S$ un morphisme d'anneaux. Étant donné un $S$-module
$N$, nous obtenons un $R$-module $N_R$ par la règle
$r \cdot n = \varphi(r)n$.
On l'appelle parfois la {\it restriction des scalaires} de $N$ à $R$.

\begin{lemma}
\label{lemma-adjoint-tensor-restrict}
Soit $R \to S$ un morphisme d'anneaux. Les foncteurs
$\text{Mod}_S \to \text{Mod}_R$, $N \mapsto N_R$ (restriction des scalaires),
et $\text{Mod}_R \to \text{Mod}_S$, $M \mapsto M \otimes_R S$
(changement de base) sont adjoints. Plus précisément,
$$
\Hom_R(M, N_R) = \Hom_S(M \otimes_R S, N)
$$
\end{lemma}

\begin{proof}
Si $\alpha : M \to N_R$ est un morphisme de $R$-modules, nous définissons
$\alpha' : M \otimes_R S \to N$ par la règle
$\alpha'(m \otimes s) = s\alpha(m)$. Si $\beta : M \otimes_R S \to N$
est un morphisme de $S$-modules, nous définissons
$\beta' : M \to N_R$ par la règle
$\beta'(m) = \beta(m \otimes 1)$.
Nous omettons la vérification que ces constructions sont inverses l'une de l'autre.
\end{proof}

\noindent
Le lemme précédent nous dit que la restriction des scalaires admet un adjoint
à gauche, à savoir le changement de base. Elle admet aussi un adjoint à droite.

\begin{lemma}
\label{lemma-adjoint-hom-restrict}
Soit $R \to S$ un morphisme d'anneaux. Les foncteurs
$\text{Mod}_S \to \text{Mod}_R$, $N \mapsto N_R$ (restriction des scalaires),
et $\text{Mod}_R \to \text{Mod}_S$, $M \mapsto \Hom_R(S, M)$
sont adjoints. Plus précisément,
$$
\Hom_R(N_R, M) = \Hom_S(N, \Hom_R(S, M))
$$
\end{lemma}

\begin{proof}
Si $\alpha : N_R \to M$ est un morphisme de $R$-modules, nous définissons
$\alpha' : N \to \Hom_R(S, M)$ par la règle
$\alpha'(n) = (s \mapsto \alpha(sn))$. Si $\beta : N \to \Hom_R(S, M)$
est un morphisme de $S$-modules, nous définissons
$\beta' : N_R \to M$ par la règle
$\beta'(n) = \beta(n)(1)$.
Nous omettons la vérification que ces constructions sont inverses l'une de l'autre.
\end{proof}

\begin{lemma}
\label{lemma-hom-from-tensor-product-variant}
Soit $R \to S$ un morphisme d'anneaux. Étant donnés des $S$-modules $M, N$
et un $R$-module $P$, nous avons
$$
\Hom_R(M \otimes_S N, P) = \Hom_S(M, \Hom_R(N, P))
$$
\end{lemma}

\begin{proof}
On peut le démontrer directement, mais c'est aussi une conséquence
des lemmes \ref{lemma-adjoint-hom-restrict} et \ref{lemma-hom-from-tensor-product}.
En effet, nous avons
\begin{align*}
\Hom_R(M \otimes_S N, P)
& =
\Hom_S(M \otimes_S N, \Hom_R(S, P)) \\
& =
\Hom_S(M, \Hom_S(N, \Hom_R(S, P))) \\
& =
\Hom_S(M, \Hom_R(N, P))
\end{align*}
comme voulu.
\end{proof}

\section{Divers}
\label{section-miscellany}

\noindent
Les démonstrations de cette section ne doivent faire appel à aucun résultat
autre que ceux de la section sur les notions de base, la section \ref{section-rings-basic}.

\begin{lemma}
\label{lemma-product-ideals-in-prime}
Soient $R$ un anneau, $I$ et $J$ deux idéaux, et $\mathfrak p$ un idéal premier
contenant le produit $IJ$. Alors $\mathfrak{p}$ contient $I$ ou $J$.
\end{lemma}

\begin{proof}
Supposons le contraire et choisissons $x \in I \setminus \mathfrak p$ et
$y \in J \setminus \mathfrak p$. Leur produit appartient à
$IJ \subset \mathfrak p$, ce qui contredit le fait que
$\mathfrak p$ est premier.
\end{proof}

\begin{lemma}[Évitement des idéaux premiers]
\label{lemma-silly}
\begin{slogan}
1. Dans un schéma affine, tout ensemble fini de points contenu dans un
ouvert est contenu dans un ouvert principal plus petit.
2. Les ouverts affines sont cofinaux parmi les voisinages d'une partie finie
donnée d'un schéma affine.
\end{slogan}
Soit $R$ un anneau. Soient $I_i \subset R$, $i = 1, \ldots, r$,
et $J \subset R$ des idéaux. Supposons que
\begin{enumerate}
\item $J \not\subset I_i$ pour $i = 1, \ldots, r$, et
\item tous les $I_i$ sauf deux sont des idéaux premiers.
\end{enumerate}
Alors il existe $x \in J$ tel que $x\not\in I_i$ pour tout $i$.
\end{lemma}

\begin{proof}
Le résultat est vrai pour $r = 1$. Si $r = 2$, choisissons $x, y \in J$ tels que
$x \not \in I_1$ et $y \not \in I_2$. Nous avons conclu, sauf si $x \in I_2$
et $y \in I_1$. Mais alors l'élément $x + y$ ne peut appartenir à $I_1$
(car cela entraînerait $x + y - y \in I_1$), ni à $I_2$.

\medskip\noindent
Pour $r \geq 3$, supposons le résultat vrai pour $r - 1$. Après renumérotation,
nous pouvons supposer que $I_r$ est premier. Nous pouvons aussi supposer
qu'il n'existe aucune inclusion entre les $I_i$. Choisissons $x \in J$ tel que
$x \not \in I_i$ pour tout $i = 1, \ldots, r - 1$. Si $x \not\in I_r$,
nous avons conclu. Supposons donc $x \in I_r$. Si
$J I_1 \ldots I_{r - 1} \subset I_r$, alors
$J \subset I_r$ (par le lemme \ref{lemma-product-ideals-in-prime}), contradiction.
Choisissons $y \in J I_1 \ldots I_{r - 1}$ tel que $y \not \in I_r$.
Alors $x + y$ convient.
\end{proof}

\begin{lemma}
\label{lemma-silly-silly}
Soit $R$ un anneau. Soient $x \in R$, $I \subset R$ un idéal, et
$\mathfrak p_i$, $i = 1, \ldots, r$, des idéaux premiers.
Supposons que $x + I \not \subset \mathfrak p_i$ pour
$i = 1, \ldots, r$. Alors il existe $y \in I$
tel que $x + y \not \in \mathfrak p_i$ pour tout $i$.
\end{lemma}

\begin{proof}
Nous pouvons supposer qu'il n'existe aucune inclusion entre les $\mathfrak p_i$.
Après réordonnement, nous pouvons supposer $x \not \in \mathfrak p_i$ pour $i < s$
et $x \in \mathfrak p_i$ pour $i \geq s$. Si $s = r + 1$, nous avons conclu.
Sinon, nous pouvons trouver $y \in I$ tel que $y \not \in \mathfrak p_s$.
Choisissons $f \in \bigcap_{i < s} \mathfrak p_i$ tel que
$f \not \in \mathfrak p_s$.
Alors $x + fy$ n'appartient à aucun de $\mathfrak p_1, \ldots, \mathfrak p_s$.
Nous concluons donc par récurrence sur $s$.
\end{proof}

\begin{lemma}[Théorème des restes chinois]
\label{lemma-chinese-remainder}
Soit $R$ un anneau.
\begin{enumerate}
\item Si $I_1, \ldots, I_r$ sont des idéaux tels que $I_a + I_b = R$
lorsque $a \not = b$, alors $I_1 \cap \ldots \cap I_r =
I_1I_2\ldots I_r$ et $R/(I_1I_2\ldots I_r)
\cong R/I_1 \times \ldots \times R/I_r$.
\item Si $\mathfrak m_1, \ldots, \mathfrak m_r$ sont des idéaux maximaux
deux à deux distincts, alors $\mathfrak m_a + \mathfrak m_b = R$ pour
$a \not = b$, et ce qui précède s'applique.
\end{enumerate}
\end{lemma}

\begin{proof}
Démontrons d'abord $I_1 \cap \ldots \cap I_r = I_1 \ldots I_r$,
car cela entraînera aussi l'injectivité du morphisme d'anneaux induit
$R/(I_1 \ldots I_r) \rightarrow R/I_1 \times \ldots \times R/I_r$.
L'inclusion $I_1 \cap \ldots \cap I_r \supset I_1 \ldots I_r$ est toujours
vraie, puisque les idéaux sont stables par multiplication par des éléments
quelconques de l'anneau. Pour démontrer l'autre inclusion, affirmons que les idéaux
$$
I_1 \ldots \hat I_i \ldots I_r,\quad i = 1, \ldots, r
$$
engendrent l'anneau $R$. Nous le démontrons par récurrence sur $r$.
C'est vrai pour $r = 2$. Si $r > 2$, nous voyons que $R$ est la somme des idéaux
$I_1 \ldots \hat I_i \ldots I_{r - 1}$, $i = 1, \ldots, r - 1$.
Par conséquent, $I_r$ est la somme des idéaux
$I_1 \ldots \hat I_i \ldots I_r$, $i = 1, \ldots, r - 1$.
En appliquant le même argument aux idéaux pris dans l'ordre inverse,
nous voyons que $I_1$ est la somme des idéaux
$I_1 \ldots \hat I_i \ldots I_r$, $i = 2, \ldots, r$.
Puisque $R = I_1 + I_r$ par hypothèse, nous voyons que $R$ est la somme
des idéaux affichés ci-dessus. Nous pouvons donc trouver des éléments
$a_i \in I_1 \ldots \hat I_i \ldots I_r$
dont la somme vaut un. En multipliant cette égalité par un élément
de $I_1 \cap \ldots \cap I_r$, nous obtenons l'autre inclusion.
Il reste à montrer que l'application canonique
$R/(I_1 \ldots I_r) \rightarrow R/I_1 \times \ldots \times R/I_r$
est surjective. Pour cela, considérons son action sur l'égalité
$1 = \sum_{i=1}^r a_i$ obtenue ci-dessus. D'une part, un
morphisme d'anneaux envoie 1 sur 1 ; d'autre part, l'image de tout
$a_i$ est nulle dans $R/I_j$ pour $j \neq i$. L'image de $a_i$
dans $R/I_i$ est donc l'élément unité. Ainsi, pour tout élément
$(\bar{b_1}, \ldots, \bar{b_r}) \in R/I_1 \times \ldots \times R/I_r$,
l'élément $\sum_{i=1}^r a_i \cdot b_i$ en est un antécédent dans $R$.

\medskip\noindent
Pour voir (2), la définition même d'idéaux maximaux distincts donne
$\mathfrak{m}_a + \mathfrak{m}_b = R$ pour $a \neq b$,
et ce qui précède s'applique.
\end{proof}

\begin{lemma}
\label{lemma-matrix-left-inverse}
Soit $R$ un anneau. Soit $n \geq m$. Soit $A$ une matrice
$n \times m$ à coefficients dans $R$. Soit $J \subset R$
l'idéal engendré par les mineurs $m \times m$ de $A$.
\begin{enumerate}
\item Pour tout $f \in J$, il existe une matrice $m \times n$, $B$,
telle que $BA = f 1_{m \times m}$.
\item Si $f \in R$ et $BA = f 1_{m \times m}$ pour une matrice
$m \times n$, $B$, alors $f^m \in J$.
\end{enumerate}
\end{lemma}

\begin{proof}
Pour $I \subset \{1, \ldots, n\}$ avec $|I| = m$, notons
$E_I$ la matrice $m \times n$ de la projection
$$
R^{\oplus n} = \bigoplus\nolimits_{i \in \{1, \ldots, n\}} R
\longrightarrow \bigoplus\nolimits_{i \in I} R
$$
et posons $A_I = E_I A$ ; ainsi, $A_I$ est la matrice $m \times m$
dont les lignes sont celles de $A$ d'indices appartenant à $I$.
Soit $B_I$ la matrice adjointe (transposée de la matrice
des cofacteurs) de $A_I$, c'est-à-dire telle que
$A_I B_I = B_I A_I = \det(A_I) 1_{m \times m}$.
Les mineurs $m \times m$ de $A$ sont les déterminants $\det A_I$
pour tous les $I \subset \{1, \ldots, n\}$ tels que $|I| = m$.
Si $f \in J$, nous pouvons écrire $f = \sum c_I \det(A_I)$ avec
$c_I \in R$. Posons $B = \sum c_I B_I E_I$ ; cela démontre (1).

\medskip\noindent
Si $f 1_{m \times m} = BA$, la formule de Cauchy-Binet
(\ref{item-cauchy-binet}) donne
$f^m = \sum b_I \det(A_I)$, où $b_I$ est le déterminant
de la matrice $m \times m$ dont les colonnes sont celles de $B$
d'indices appartenant à $I$.
\end{proof}

\begin{lemma}
\label{lemma-matrix-right-inverse}
Soit $R$ un anneau. Soit $n \geq m$. Soit $A = (a_{ij})$ une matrice
$n \times m$ à coefficients dans $R$, écrite par blocs sous la forme
$$
A =
\left(
\begin{matrix}
A_1 \\
A_2
\end{matrix}
\right)
$$
où $A_1$ est de taille $m \times m$. Soit $B$ la matrice adjointe
(transposée de la matrice des cofacteurs) de $A_1$. Alors
$$
AB = 
\left(
\begin{matrix}
f 1_{m \times m} \\
C
\end{matrix}
\right)
$$
où $f = \det(A_1)$ et où $c_{ij}$ est, au signe près, le déterminant du
mineur $m \times m$ de $A$ correspondant aux lignes
$1, \ldots, \hat j, \ldots, m, i$.
\end{lemma}

\begin{proof}
Puisque la matrice adjointe vérifie $A_1B = B A_1 = f$, le premier bloc
de l'expression de $AB$ est correct. Remarquons que
$$
c_{ij} = \sum\nolimits_k a_{ik}b_{kj} = \sum (-1)^{j + k}a_{ik} \det(A_1^{jk})
$$
où $A_1^{ij}$ désigne la matrice $A_1$ privée de sa $j$-ième ligne
et de sa $k$-ième colonne. Cette dernière expression est le développement
suivant une ligne du déterminant de la matrice de l'énoncé du lemme.
\end{proof}

\begin{lemma}
\label{lemma-map-cannot-be-injective}
\begin{slogan}
Un morphisme de modules libres de type fini ne peut être injectif si la source
est de rang strictement supérieur à celui du but.
\end{slogan}
Soit $R$ un anneau non nul. Soit $n \geq 1$. Soit $M$ un $R$-module engendré
par un nombre d'éléments $< n$. Alors tout morphisme de $R$-modules
$f : R^{\oplus n} \to M$ a un noyau non nul.
\end{lemma}

\begin{proof}
Choisissons une surjection $R^{\oplus n - 1} \to M$.
Nous pouvons relever le morphisme $f$ en un morphisme
$f' : R^{\oplus n} \to R^{\oplus n - 1}$
(lemme \ref{lemma-lift-map}).
Il suffit de démontrer que $f'$ a un noyau non nul.
Le morphisme $f' : R^{\oplus n} \to R^{\oplus n - 1}$ est donné par une
matrice $A = (a_{ij})$. Si l'un des $a_{ij}$ n'est pas nilpotent, disons
que $a = a_{ij}$ ne l'est pas, nous pouvons remplacer $R$ par la localisation
$R_a$ et supposer que $a_{ij}$ est une unité. En effet, si nous trouvons un
noyau non nul après localisation, il existait déjà un noyau non nul,
car la localisation est exacte ; voir la proposition
\ref{proposition-localization-exact}.
Dans ce cas, nous pouvons effectuer un changement de base à la fois dans
$R^{\oplus n}$ et dans $R^{\oplus n - 1}$ et nous ramener au cas où
$$
A =
\left(
\begin{matrix}
1 & 0 & 0 & \ldots \\
0 & a_{22} & a_{23} & \ldots \\
0 & a_{32} & \ldots \\
\ldots & \ldots
\end{matrix}
\right)
$$
Nous concluons alors par récurrence sur $n$. Sinon, chaque
$a_{ij}$ est nilpotent. Posons $I = (a_{ij}) \subset R$. Remarquons que
$I^{m + 1} = 0$ pour un certain $m \geq 0$. Soit $m$ le plus grand entier
tel que $I^m \not = 0$. Alors $(I^m)^{\oplus n}$ est
contenu dans le noyau du morphisme, ce qui conclut.
\end{proof}

\begin{lemma}
\label{lemma-rank}
\begin{slogan}
Le rang d'un module libre de type fini est bien défini.
\end{slogan}
Soit $R$ un anneau non nul. Soient $n, m \geq 0$ des entiers.
Si $R^{\oplus n}$ est isomorphe à $R^{\oplus m}$ comme
$R$-modules, alors $n = m$.
\end{lemma}

\begin{proof}
C'est immédiat d'après le lemme \ref{lemma-map-cannot-be-injective}.
\end{proof}

\section{Cayley-Hamilton}
\label{section-cayley-hamilton}

\begin{lemma}
\label{lemma-charpoly}
Soit $R$ un anneau. Soit $A = (a_{ij})$ une matrice $n \times n$
à coefficients dans $R$. Soit $P(x) \in R[x]$
le polynôme caractéristique de $A$ (défini
comme $\det(x\text{id}_{n \times n} - A)$).
Alors $P(A) = 0$ dans $\text{Mat}(n \times n, R)$.
\end{lemma}

\begin{proof}
Nous ramenons la question au théorème de Cayley-Hamilton bien connu
d'algèbre linéaire, en plusieurs étapes :
\begin{enumerate}
\item Si $\phi :S \rightarrow R$ est un morphisme d'anneaux et si les $b_{ij}$
sont des antécédents des $a_{ij}$ par ce morphisme, il suffit de démontrer
l'énoncé pour $S$ et $(b_{ij})$, puisque $\phi$ est un morphisme d'anneaux.
\item Si $\psi :R \hookrightarrow S$ est un morphisme injectif d'anneaux,
il suffit clairement de démontrer le résultat pour $S$ et pour les $a_{ij}$
considérés comme éléments de $S$.
\item Nous pouvons donc d'abord nous ramener au cas $R = \mathbf{Z}[X_{ij}]$,
$a_{ij} = X_{ij}$, d'un anneau de polynômes, puis au
cas $R = \mathbf{Q}(X_{ij})$, où nous pouvons enfin appliquer Cayley-Hamilton.
\end{enumerate}
\end{proof}

\begin{lemma}
\label{lemma-charpoly-module}
Soit $R$ un anneau.
Soit $M$ un $R$-module de type fini.
Soit $\varphi : M \to M$ un endomorphisme.
Alors il existe un polynôme unitaire $P \in R[T]$ tel que
$P(\varphi) = 0$ comme endomorphisme de $M$.
\end{lemma}

\begin{proof}
Choisissons un morphisme surjectif de $R$-modules $R^{\oplus n} \to M$, donné par
$(a_1, \ldots, a_n) \mapsto \sum a_ix_i$, pour des générateurs $x_i \in M$.
Choisissons $(a_{i1}, \ldots, a_{in}) \in R^{\oplus n}$ tels que
$\varphi(x_i) = \sum a_{ij} x_j$. Autrement dit, le diagramme
$$
\xymatrix{
R^{\oplus n} \ar[d]_A \ar[r] & M \ar[d]^\varphi \\
R^{\oplus n} \ar[r] & M
}
$$
est commutatif, où $A = (a_{ij})$. D'après
le lemme \ref{lemma-charpoly},
il existe un polynôme unitaire $P$ tel que $P(A) = 0$.
Il s'ensuit que $P(\varphi) = 0$.
\end{proof}

\begin{lemma}
\label{lemma-charpoly-module-ideal}
Soit $R$ un anneau. Soit $I \subset R$ un idéal.
Soit $M$ un $R$-module de type fini.
Soit $\varphi : M \to M$ un endomorphisme tel
que $\varphi(M) \subset IM$.
Alors il existe un polynôme unitaire
$P = t^n + a_1 t^{n - 1} + \ldots + a_n \in R[T]$
tel que $a_j \in I^j$ et que $P(\varphi) = 0$ comme endomorphisme de $M$.
\end{lemma}

\begin{proof}
Choisissons un morphisme surjectif de $R$-modules $R^{\oplus n} \to M$, donné par
$(a_1, \ldots, a_n) \mapsto \sum a_ix_i$, pour des générateurs $x_i \in M$.
Choisissons $(a_{i1}, \ldots, a_{in}) \in I^{\oplus n}$ tels que
$\varphi(x_i) = \sum a_{ij} x_j$. Autrement dit, le diagramme
$$
\xymatrix{
R^{\oplus n} \ar[d]_A \ar[r] & M \ar[d]^\varphi \\
I^{\oplus n} \ar[r] & M
}
$$
est commutatif, où $A = (a_{ij})$. D'après
le lemme \ref{lemma-charpoly},
le polynôme
$P(t) = \det(t\text{id}_{n \times n} - A)$
possède toutes les propriétés voulues.
\end{proof}

\noindent
À titre d'application plaisante, démontrons le lemme surprenant suivant.

\begin{lemma}
\label{lemma-fun}
Soit $R$ un anneau.
Soit $M$ un $R$-module de type fini.
Soit $\varphi : M \to M$ un morphisme surjectif de $R$-modules.
Alors $\varphi$ est un isomorphisme.
\end{lemma}

\begin{proof}[Première démonstration]
Écrivons $R' = R[x]$ et considérons $M$ comme un $R'$-module de type fini,
$x$ agissant par $\varphi$. Posons $I = (x) \subset R'$. L'hypothèse de
surjectivité de $\varphi$ donne $IM = M$. Nous pouvons donc appliquer
le lemme \ref{lemma-charpoly-module-ideal}
à $M$ considéré comme un $R'$-module, à l'idéal $I$ et à l'endomorphisme
$\text{id}_M$.
Nous en déduisons que
$(1 + a_1 + \ldots + a_n)\text{id}_M = 0$ avec $a_j \in I$.
Écrivons $a_j = b_j(x)x$ pour un certain $b_j(x) \in R[x]$.
En revenant à $\varphi$, nous voyons que
$\text{id}_M = -(\sum_{j = 1, \ldots, n} b_j(\varphi)) \varphi$,
et donc que $\varphi$ est inversible.
\end{proof}

\begin{proof}[Seconde démonstration]
Nous raisonnons par récurrence sur le nombre de générateurs de $M$ sur $R$. Si
$M$ est engendré par un élément, alors $M \cong R/I$ pour un certain idéal
$I \subset R$. Nous pouvons alors remplacer $R$ par $R/I$, de sorte que $M = R$.
Dans ce cas, $\varphi : R \to R$ est donné par la multiplication sur $M$ par un
élément $r \in R$. La surjectivité de $\varphi$ force $r$ à être inversible,
puisque $\varphi$ doit atteindre $1$ ; il s'ensuit que $\varphi$ est
inversible.

\medskip\noindent
Supposons maintenant le lemme démontré pour les modules
engendrés par $n - 1$ éléments, et considérons un module $M$ engendré
par $n$ éléments. Notons $A$ l'anneau $R[t]$ et regardons le module
$M$ comme un $A$-module en faisant agir $t$ par $\varphi$ ; comme $M$ est
de type fini sur $R$, il l'est aussi sur $R[t]$. Puisque nous cherchons
à démontrer que $\varphi$ est injectif, ce qui est une propriété ensembliste,
nous pouvons tout aussi bien démontrer que l'endomorphisme
$t : M \to M$ sur $A$ est injectif. Nous avons ramené le problème au cas
où notre endomorphisme est la multiplication par un élément de l'anneau de base.
Notons $M' \subset M$ le sous-$A$-module engendré par les $n - 1$ premiers
générateurs de $M$, et considérons le diagramme
$$
\xymatrix{
0 \ar[r] & M' \ar[r]\ar[d]^{\varphi\mid_{M'}} & M\ar[d]^\varphi \ar[r] &
M/M' \ar[d]^{\varphi \bmod M'} \ar[r] & 0 \\
0 \ar[r] & M' \ar[r]                                  & M \ar[r]
           & M/M' \ar[r]                                  & 0,
}
$$
où la restriction de $\varphi$ à $M'$ et l'application induite par $\varphi$
sur le quotient $M/M'$ sont bien définies, puisque $\varphi$ est la
multiplication par un élément de la base, et que $M'$ et $M/M'$ sont
eux-mêmes des $A$-modules. Par le cas $n = 1$, l'application
$M/M' \to M/M'$ est un isomorphisme. Une chasse au diagramme montre que
$\varphi|_{M'}$ est surjective ; par récurrence, $\varphi|_{M'}$ est donc
un isomorphisme. Le lemme du serpent impose alors à la colonne du milieu
d'être un isomorphisme.
\end{proof}

\section{Le spectre d'un anneau}
\label{section-spectrum-ring}

\noindent
Nous décidons arbitrairement que le spectre d'un anneau, considéré comme espace
topologique, relève du chapitre d'algèbre, tandis qu'un schéma affine relève du
chapitre sur les schémas.

\begin{definition}
\label{definition-spectrum-ring}
Soit $R$ un anneau.
\begin{enumerate}
\item Le {\it spectre} de $R$ est l'ensemble des idéaux premiers de $R$.
On le note habituellement $\Spec(R)$.
\item Étant donnée une partie $T \subset R$, nous notons
$V(T) \subset \Spec(R)$ l'ensemble des idéaux premiers contenant $T$,
c'est-à-dire $V(T) = \{ \mathfrak p \in
\Spec(R) \mid \forall f\in T, f\in \mathfrak p\}$.
\item Étant donné un élément $f \in R$, nous notons $D(f) \subset \Spec(R)$
l'ensemble des idéaux premiers ne contenant pas $f$.
\end{enumerate}
\end{definition}

\begin{lemma}
\label{lemma-Zariski-topology}
Soit $R$ un anneau.
\begin{enumerate}
\item Le spectre d'un anneau $R$ est vide si et seulement si $R$
est l'anneau nul.
\item Tout anneau non nul possède un idéal maximal.
\item Tout anneau non nul possède un idéal premier minimal.
\item Étant donnés un idéal $I \subset R$ et un idéal premier vérifiant
$I \subset \mathfrak p$, il existe un idéal premier
$I \subset \mathfrak q \subset \mathfrak p$ tel que
$\mathfrak q$ soit minimal au-dessus de $I$.
\item Si $T \subset R$ et si $(T)$ est l'idéal engendré par
$T$ dans $R$, alors $V((T)) = V(T)$.
\item Si $I$ est un idéal et si $\sqrt{I}$ est son radical,
voir la notion de base (\ref{item-radical-ideal}), alors $V(I) = V(\sqrt{I})$.
\item Pour tout idéal $I$ de $R$, nous avons $\sqrt{I} =
\bigcap_{I \subset \mathfrak p} \mathfrak p$.
\item Si $I$ est un idéal, alors $V(I) = \emptyset$ si et seulement
si $I$ est l'idéal unité.
\item Si $I$ et $J$ sont des idéaux de $R$, alors $V(I) \cup V(J) =
V(I \cap J)$.
\item Si $(I_a)_{a\in A}$ est une famille d'idéaux de $R$, alors
$\bigcap_{a\in A} V(I_a) = V(\bigcup_{a\in A} I_a)$.
\item Si $f \in R$, alors $D(f) \amalg V(f) = \Spec(R)$.
\item Si $f \in R$, alors $D(f) = \emptyset$ si et seulement si $f$
est nilpotent.
\item Si $f = u f'$ pour une unité $u \in R$, alors $D(f) = D(f')$.
\item Si $I \subset R$ est un idéal et si $\mathfrak p$ est un idéal premier de
$R$ tel que $\mathfrak p \not\in V(I)$, alors il existe $f \in R$
tel que $\mathfrak p \in D(f)$ et $D(f) \cap V(I) = \emptyset$.
\item Si $f, g \in R$, alors $D(fg) = D(f) \cap D(g)$.
\item Si $f_i \in R$ pour $i \in I$, alors
$\bigcup_{i\in I} D(f_i)$ est le complémentaire de
$V(\{f_i \}_{i\in I})$ dans $\Spec(R)$.
\item Si $f \in R$ et $D(f) = \Spec(R)$, alors $f$ est une unité.
\end{enumerate}
\end{lemma}

\begin{proof}
Nous traitons chaque assertion dans l'item correspondant ci-dessous.
\begin{enumerate}
\item C'est une conséquence directe de (2) ou de (3).
\item Soit $\mathfrak{A}$ l'ensemble de tous les idéaux propres de $R$.
Cet ensemble est ordonné par inclusion et n'est pas vide, puisque
$(0) \in \mathfrak{A}$ est un idéal propre.
Soit $A$ une partie totalement ordonnée de $\mathfrak A$.
Alors $\bigcup_{I \in A} I$ est en fait un idéal. Comme 1 $\notin I$
pour tout $I \in A$, la réunion ne contient pas 1 et est donc propre.
Ainsi, $\bigcup_{I \in A} I$ appartient à $\mathfrak{A}$ et majore
la partie $A$. Le lemme de Zorn assure donc à $\mathfrak{A}$
l'existence d'un élément maximal, qui est l'idéal maximal recherché.
\item Puisque $R$ est non nul, il contient un idéal maximal, donc premier.
Ainsi, l'ensemble $\mathfrak{A}$ de tous les idéaux premiers de $R$
n'est pas vide. $\mathfrak{A}$ est ordonné par l'inclusion opposée.
Soit $A$ une partie totalement ordonnée de $\mathfrak{A}$. Il est clair que
$J = \bigcap_{I \in A} I$ est un idéal. Il est moins clair qu'il soit premier.
Soit $xy \in J$. Alors $xy \in I$ pour tout $I \in A$.
Posons $B = \{I \in A | y \in I\}$ et $K
= \bigcap_{I \in B} I$. Puisque $A$ est totalement ordonné, ou bien
$K = J$ (et nous avons conclu, car alors $y \in J$), ou bien $K \supset J$
et, pour tout $I \in A$ tel que $I$ soit strictement contenu dans $K$,
nous avons $y \notin I$. Mais cela signifie que, pour tous ces idéaux
$I, x \in I$, puisqu'ils sont premiers. Ainsi $x \in J$. Dans les deux cas,
$J$ est premier, comme voulu. Le lemme de Zorn fournit donc un élément
maximal, qui est ici un idéal premier minimal.
\item C'est exactement le même argument que dans (3), en ne considérant que
les idéaux premiers contenus dans $\mathfrak{p}$ et contenant $I$.
\item $(T)$ est le plus petit idéal contenant $T$. Par conséquent, si
$T \subset I$ pour un certain idéal, alors $(T) \subset I$ également.
Ainsi, si $I \in V(T)$, alors $I \in V((T))$ également.
L'autre inclusion est évidente.
\item Puisque $I \subset \sqrt{I}, V(\sqrt{I}) \subset V(I)$, prenons
$\mathfrak{p} \in V(I)$. Soit $x \in \sqrt{I}$. Alors $x^n \in I$
pour un certain $n$. Donc $x^n \in \mathfrak{p}$. Comme $\mathfrak{p}$
est premier, une récurrence immédiate donne $x \in \mathfrak{p}$.
Ainsi $\sqrt{I} \subset \mathfrak{p}$ et
$\mathfrak{p} \in V(\sqrt{I})$.
\item Soit $f \in R \setminus \sqrt{I}$. Alors $f^n \notin I$
pour tout $n$. Par conséquent,
$S = \{1, f, f^2, \ldots\}$ est une partie multiplicative ne contenant pas $0$.
Choisissons un idéal premier
$\bar{\mathfrak{p}} \subset S^{-1}R$ contenant $S^{-1}I$.
Alors $\mathfrak{p}$, image réciproque dans $R$ de $\bar{\mathfrak{p}}$,
est un idéal premier contenant $I$ et ne rencontrant pas $S$.
Cela montre que $\bigcap_{I \subset
\mathfrak p} \mathfrak p \subset \sqrt{I}$. Maintenant, si
$a \in \sqrt{I}$, alors $a^n \in I$ pour un certain $n$.
Ainsi, si $I \subset \mathfrak{p}$, alors $a^n \in
\mathfrak{p}$. Comme $\mathfrak{p}$ est premier, nous avons
$a \in \mathfrak{p}$. L'égalité est donc démontrée.
\item $I$ n'est pas l'idéal unité si et seulement si $I$
est contenu dans un idéal maximal (pour le voir, appliquer (2) à l'anneau
$R/I$), lequel est donc premier.
\item Si $\mathfrak{p} \in V(I) \cup V(J)$, alors $I \subset \mathfrak{p}$
ou $J \subset \mathfrak{p}$, ce qui signifie que
$I \cap J \subset \mathfrak{p}$. Réciproquement, si
$I \cap J \subset \mathfrak{p}$, alors $IJ \subset \mathfrak{p}$ ;
ainsi, soit $I$, soit $J$, est contenu dans $\mathfrak{p}$,
puisque $\mathfrak{p}$ est premier.
\item $\mathfrak{p} \in \bigcap_{a \in A} V(I_a) \Leftrightarrow
I_a \subset \mathfrak{p}, \forall a \in A \Leftrightarrow
\mathfrak{p} \in V(\bigcup_{a\in A} I_a)$
\item Si $\mathfrak{p}$ est un idéal premier et $f \in R$, alors ou bien
$f \in \mathfrak{p}$, ou bien $f \notin \mathfrak{p}$, de façon exclusive ;
c'est précisément ce qu'exprime la réunion disjointe.
\item Si $a \in R$ est nilpotent, alors $a^n = 0$ pour un certain $n$.
Ainsi $a^n \in \mathfrak{p}$ pour tout idéal premier. Une récurrence donne
donc $a \in \mathfrak{p}$ et $D(a) = \emptyset$. Réciproquement, comme
nous l'avons montré en (7), si $a \in R$ n'est pas nilpotent, il existe
un idéal premier qui ne le contient pas.
\item $f \in \mathfrak{p} \Leftrightarrow uf \in \mathfrak{p}$,
puisque $u$ est inversible.
\item Si $\mathfrak{p} \notin V(I)$, alors
$\exists f \in I \setminus \mathfrak{p}$. Ainsi
$f \notin \mathfrak{p}$, donc $\mathfrak{p} \in D(f)$. En outre, si
$\mathfrak{q} \in D(f)$, alors $f \notin \mathfrak{q}$, si bien que $I$
n'est pas contenu dans $\mathfrak{q}$. Ainsi
$D(f) \cap V(I) = \emptyset$.
\item Si $fg \in \mathfrak{p}$, alors $f \in \mathfrak{p}$ ou
$g \in \mathfrak{p}$. Par conséquent, si $f \notin \mathfrak{p}$ et
$g \notin \mathfrak{p}$, alors $fg \notin \mathfrak{p}$.
Puisque $\mathfrak{p}$ est un idéal, si $fg \notin
\mathfrak{p}$, alors $f \notin \mathfrak{p}$ et $g \notin \mathfrak{p}$.
\item $\mathfrak{p} \in \bigcup_{i \in I} D(f_i) \Leftrightarrow \exists i \in
I, f_i \notin \mathfrak{p} \Leftrightarrow \mathfrak{p} \in \Spec(R)
\setminus V(\{f_i\}_{i \in I})$
\item Si $D(f) = \Spec(R)$, alors $V(f) = \emptyset$ et donc
$fR = R$, de sorte que $f$ est une unité.
\end{enumerate}
\end{proof}

\noindent
Le lemme entraîne que les parties $V(T)$ de la
définition \ref{definition-spectrum-ring} sont les fermés
d'une topologie sur $\Spec(R)$. Il montre aussi que
les ensembles $D(f)$ sont ouverts et forment une base de cette
topologie.

\begin{definition}
\label{definition-Zariski-topology}
Soit $R$ un anneau.
La topologie sur $\Spec(R)$ dont les fermés sont les
ensembles $V(T)$ est appelée topologie {\it de Zariski}. Les ouverts
$D(f)$ sont appelés les {\it ouverts principaux} de $\Spec(R)$.
\end{definition}

\noindent
Le contexte indiquera clairement si nous considérons $\Spec(R)$
seulement comme un ensemble ou comme un espace topologique.

\begin{lemma}
\label{lemma-spec-functorial}
\begin{slogan}
Fonctorialité du spectre
\end{slogan}
Supposons que $\varphi : R \to R'$ soit un morphisme d'anneaux.
L'application induite
$$
\Spec(\varphi) : \Spec(R') \longrightarrow \Spec(R),
\quad
\mathfrak p' \longmapsto \varphi^{-1}(\mathfrak p')
$$
est continue pour les topologies de Zariski. En fait, pour tout
élément $f \in R$, nous avons
$\Spec(\varphi)^{-1}(D(f)) = D(\varphi(f))$.
\end{lemma}

\begin{proof}
La notion de base (\ref{item-inverse-image-prime}) assure que
$\mathfrak p := \varphi^{-1}(\mathfrak p')$
est bien un idéal premier de $R$. La dernière assertion
du lemme résulte directement des définitions
et entraîne la première.
\end{proof}

\noindent
Si $\varphi' : R' \to R''$ est un second morphisme d'anneaux,
alors la composée
$$
\Spec(R'')
\longrightarrow
\Spec(R')
\longrightarrow
\Spec(R)
$$
est égale à $\Spec(\varphi' \circ \varphi)$. Autrement dit,
$\Spec$ est un foncteur contravariant de la catégorie des anneaux
vers la catégorie des espaces topologiques.

\begin{lemma}
\label{lemma-spec-localization}
Soit $R$ un anneau. Soit $S \subset R$ une partie multiplicative.
Le morphisme $R \to S^{-1}R$ induit, par fonctorialité de $\Spec$,
un homéomorphisme
$$
\Spec(S^{-1}R)
\longrightarrow
\{\mathfrak p \in \Spec(R) \mid S \cap \mathfrak p = \emptyset \}
$$
où le membre de droite est muni de la topologie induite par la
topologie de Zariski sur $\Spec(R)$. L'application inverse est donnée
par $\mathfrak p \mapsto S^{-1}\mathfrak p = \mathfrak p(S^{-1}R)$.
\end{lemma}

\begin{proof}
Notons $D$ le membre de droite de la flèche du lemme.
Choisissons un idéal premier $\mathfrak p' \subset S^{-1}R$ et soit
$\mathfrak p$ l'image réciproque de $\mathfrak p'$ dans $R$.
Puisque $\mathfrak p'$ ne contient pas $1$, nous voyons que $\mathfrak p$
ne contient aucun élément de $S$. Ainsi $\mathfrak p \in D$, et nous voyons que
l'image est contenue dans $D$. Soit $\mathfrak p \in D$.
Par hypothèse, l'image $\overline{S}$ ne contient pas $0$.
D'après la notion de base (\ref{item-localization-zero}),
$\overline{S}^{-1}(R/\mathfrak p)$ n'est pas l'anneau nul.
D'après la notion de base (\ref{item-localize-ideal}), nous voyons que
$S^{-1}R / S^{-1}\mathfrak p = \overline{S}^{-1}(R/\mathfrak p)$
est un anneau intègre, et donc que $S^{-1}\mathfrak p$ est premier.
L'égalité d'anneaux montre aussi que l'image réciproque de
$S^{-1}\mathfrak p$ dans $R$ est égale à $\mathfrak p$,
car $R/\mathfrak p \to \overline{S}^{-1}(R/\mathfrak p)$
est injective d'après la notion de base
(\ref{item-localize-nonzerodivisors}).
Cela démontre que l'application $\Spec(S^{-1}R) \to \Spec(R)$
est une bijection sur $D$, dont l'inverse est celui indiqué.
Elle est continue d'après le lemme \ref{lemma-spec-functorial}.
Enfin, soit $D(g) \subset \Spec(S^{-1}R)$ un ouvert principal.
Écrivons $g = h/s$ pour certains $h\in R$ et $s\in S$.
Puisque $g$ s'obtient à partir de $h/1$ par multiplication par une unité, nous avons
$D(g) = D(h/1)$ dans $\Spec(S^{-1}R)$.
Le lemme \ref{lemma-spec-functorial} et la bijectivité ci-dessus montrent donc
que l'image de $D(g) = D(h/1)$ est $D \cap D(h)$.
Cela prouve que l'application est également ouverte.
\end{proof}

\begin{lemma}
\label{lemma-standard-open}
Soit $R$ un anneau. Soit $f \in R$.
Le morphisme $R \to R_f$ induit, par fonctorialité de
$\Spec$, un homéomorphisme
$$
\Spec(R_f) \longrightarrow D(f) \subset \Spec(R).
$$
L'application inverse est donnée par
$\mathfrak p \mapsto \mathfrak p \cdot R_f$.
\end{lemma}

\begin{proof}
C'est un cas particulier du lemme \ref{lemma-spec-localization}.
\end{proof}

\noindent
Tout « ouvert affine » d'un spectre n'est pas nécessairement
un ouvert principal. Voir l'exemple
\ref{example-affine-open-not-standard}.

\begin{lemma}
\label{lemma-spec-closed}
Soit $R$ un anneau. Soit $I \subset R$ un idéal.
Le morphisme $R \to R/I$ induit, par fonctorialité de
$\Spec$, un homéomorphisme
$$
\Spec(R/I) \longrightarrow V(I) \subset \Spec(R).
$$
L'application inverse est donnée par
$\mathfrak p \mapsto \mathfrak p / I$.
\end{lemma}

\begin{proof}
Il est immédiat que l'image est contenue dans $V(I)$.
D'autre part, si $\mathfrak p \in V(I)$, alors
$\mathfrak p \supset I$, et nous pouvons considérer
l'idéal $\mathfrak p /I \subset R/I$. La notion de base
(\ref{item-isomorphism-theorem}) montre que
$(R/I)/(\mathfrak p/I) = R/\mathfrak p$ est un anneau intègre,
et donc que $\mathfrak p/I$ est un idéal premier. Il est alors
immédiat que l'image de $D(f + I)$ est
$D(f) \cap V(I)$, et donc que l'application est un homéomorphisme.
\end{proof}



\begin{lemma}
\label{lemma-quasi-compact}
\begin{slogan}
Le spectre d'un anneau est quasi-compact
\end{slogan}
Soit $R$ un anneau. L'espace $\Spec(R)$ est quasi-compact.
\end{lemma}

\begin{proof}
Il suffit de démontrer que de tout recouvrement de $\Spec(R)$
par des ouverts principaux on peut extraire un recouvrement fini.
Supposons donc que $\Spec(R) = \cup D(f_i)$
pour une famille d'éléments $\{f_i\}_{i\in I}$ de $R$. Cela signifie que
$\cap V(f_i) = \emptyset$. D'après le lemme
\ref{lemma-Zariski-topology}, cela signifie que
$V(\{f_i \}) = \emptyset$. D'après le
même lemme, cela signifie que l'idéal engendré
par les $f_i$ est l'idéal unité de $R$. Nous pouvons donc
écrire $1$ comme une somme {\it finie} :
$1 = \sum_{i \in J} r_i f_i$ avec $J \subset I$ fini.
Il s'ensuit alors que $\Spec(R)
= \cup_{i \in J} D(f_i)$.
\end{proof}

\begin{lemma}
\label{lemma-topology-spec}
Soit $R$ un anneau.
La topologie sur $X = \Spec(R)$ possède les propriétés suivantes :
\begin{enumerate}
\item $X$ est quasi-compact,
\item $X$ possède une base de la topologie formée d'ouverts quasi-compacts, et
\item l'intersection de deux ouverts quasi-compacts quelconques est quasi-compacte.
\end{enumerate}
\end{lemma}

\begin{proof}
Le spectre d'un anneau est quasi-compact ; voir le
lemme \ref{lemma-quasi-compact}.
Il possède une base de la topologie formée des ouverts principaux
$D(f) = \Spec(R_f)$
(lemme \ref{lemma-standard-open}),
qui sont quasi-compacts d'après la première assertion.
L'intersection de deux ouverts principaux est quasi-compacte,
car $D(f) \cap D(g) = D(fg)$. Étant donnés deux ouverts quasi-compacts
$U, V \subset X$, nous pouvons écrire $U = D(f_1) \cup \ldots \cup D(f_n)$
et $V = D(g_1) \cup \ldots \cup D(g_m)$. Alors
$U \cap V = \bigcup D(f_ig_j)$, qui est quasi-compact.
\end{proof}

\section{Anneaux locaux}
\label{section-local-rings}

\noindent
Les anneaux locaux sont au cœur de la géométrie algébrique.

\begin{definition}
\label{definition-local-ring}
Un {\it anneau local} est un anneau qui possède exactement un idéal maximal.
Si $R$ est un anneau local, on note souvent son idéal maximal
$\mathfrak m_R$, et le corps $R/\mathfrak m_R$ est appelé le
{\it corps résiduel} de l'anneau local $R$.
Nous disons souvent « soit $(R, \mathfrak m)$ un anneau local »
ou « soit $(R, \mathfrak m, \kappa)$ un anneau local »
pour indiquer que $R$ est local, que $\mathfrak m$ est son unique
idéal maximal et que $\kappa = R/\mathfrak m$ est son corps résiduel.
Un {\it morphisme local d'anneaux locaux} est un morphisme d'anneaux
$\varphi : R \to S$ tel que $R$ et $S$ soient des anneaux locaux et
que $\varphi(\mathfrak m_R) \subset \mathfrak m_S$.
S'il est entendu que $R$ et $S$ sont des anneaux locaux, l'expression
« {\it morphisme local d'anneaux $\varphi : R \to S$} » signifie que
$\varphi$ est un morphisme local d'anneaux locaux.
\end{definition}

\noindent
Un corps est un anneau local. Tout morphisme d'anneaux entre corps est
un morphisme local d'anneaux locaux.

\medskip\noindent
La localisation $R_\mathfrak p$ d'un anneau $R$ en un idéal premier
$\mathfrak p$ est un anneau local d'idéal maximal
$\mathfrak p R_\mathfrak p$. En effet, d'après le
lemme \ref{lemma-spec-localization}, tout idéal premier de $R_\mathfrak p$
est contenu dans l'idéal premier $\mathfrak p R_\mathfrak p$
(qui est donc un idéal maximal, et l'unique idéal maximal de $R_\mathfrak p$).
Le corps résiduel de $R_\mathfrak p$ est noté $\kappa(\mathfrak p)$ ;
nous l'appelons le {\it corps résiduel de $\mathfrak p$}. D'après la
proposition \ref{proposition-localize-quotient}, nous pouvons identifier
$\kappa(\mathfrak p)$ au corps des fractions de l'anneau intègre
$R/\mathfrak p$. Par la composée
$$
\Spec(\kappa(\mathfrak p)) \to \Spec(R_\mathfrak p) \to \Spec(R)
$$
l'unique point de la source est envoyé sur le point $\mathfrak p$ du but.

\medskip\noindent
Soit $\varphi : R \to S$ un morphisme d'anneaux. Soit
$\mathfrak q \subset S$ un idéal premier, et considérons l'idéal premier
$\mathfrak p = \varphi^{-1}(\mathfrak q)$ de $R$.
Puisque $\varphi(\mathfrak p) \subset \mathfrak q$, le morphisme d'anneaux induit
$$
R_\mathfrak p \to S_\mathfrak q,\quad r/g \mapsto \varphi(r)/\varphi(g)
$$
est un morphisme local d'anneaux locaux, et nous obtenons un morphisme induit
entre les corps résiduels $\kappa(\mathfrak p) \to \kappa(\mathfrak q)$.


\begin{example}
\label{example-not-local}
Si $R$ est un anneau local et si $\mathfrak p \subset R$ est un idéal premier
non maximal, alors $R \to R_\mathfrak p$ n'est pas un morphisme local.
\end{example}

\begin{lemma}
\label{lemma-characterize-local-ring}
Soit $R$ un anneau. Les conditions suivantes sont équivalentes :
\begin{enumerate}
\item $R$ est un anneau local,
\item $\Spec(R)$ possède exactement un point fermé,
\item $R$ possède un idéal maximal $\mathfrak m$
et tout élément de $R \setminus \mathfrak m$
est une unité, et
\item $R$ n'est pas l'anneau nul et, pour tout $x \in R$, ou bien $x$,
ou bien $1 - x$, ou bien les deux sont inversibles.
\end{enumerate}
\end{lemma}

\begin{proof}
Soient $R$ un anneau et $\mathfrak m$ un idéal maximal.
Si $x \in R \setminus \mathfrak m$ et si $x$ n'est pas une unité,
alors il existe un idéal maximal $\mathfrak m'$ contenant $x$.
Ainsi $R$ possède au moins deux idéaux maximaux. Réciproquement,
si $\mathfrak m'$ est un autre idéal maximal, choisissons
$x \in \mathfrak m'$ tel que $x \not \in \mathfrak m$. Clairement,
$x$ n'est pas une unité. Cela démontre l'équivalence de (1) et de (3).
L'équivalence de (1) et de (2) est tautologique.
Si $R$ est local, alors (4) est vraie, puisque $x$ appartient ou non à
$\mathfrak m$. Si (4) est vraie et si $\mathfrak m$, $\mathfrak m'$
sont des idéaux maximaux distincts, nous pouvons choisir $x \in R$ tel que
$x \bmod \mathfrak m' = 0$ et $x \bmod \mathfrak m = 1$
par le théorème des restes chinois
(lemme \ref{lemma-chinese-remainder}).
Cet élément $x$ n'est pas inversible, et $1 - x$ ne l'est pas davantage,
ce qui est une contradiction. Ainsi (4) et (1) sont équivalentes.
\end{proof}

\begin{lemma}
\label{lemma-characterize-local-ring-map}
Soit $\varphi : R \to S$ un morphisme d'anneaux. Supposons que $R$ et $S$
soient des anneaux locaux. Les conditions suivantes sont équivalentes :
\begin{enumerate}
\item $\varphi$ est un morphisme local d'anneaux,
\item $\varphi(\mathfrak m_R) \subset \mathfrak m_S$,
\item $\varphi^{-1}(\mathfrak m_S) = \mathfrak m_R$, et
\item pour tout $x \in R$, si $\varphi(x)$ est inversible dans $S$, alors $x$
est inversible dans $R$.
\end{enumerate}
\end{lemma}

\begin{proof}
Les conditions (1) et (2) sont équivalentes par définition.
Si (3) est vraie, alors (2) l'est. Réciproquement, si (2) est vraie,
$\varphi^{-1}(\mathfrak m_S)$ est un idéal premier contenant
l'idéal maximal $\mathfrak m_R$, donc
$\varphi^{-1}(\mathfrak m_S) = \mathfrak m_R$. Enfin, (4) est la
contraposée de (2) d'après le lemme \ref{lemma-characterize-local-ring}.
\end{proof}

\begin{remark}
\label{remark-fundamental-diagram}
Un diagramme commutatif fondamental associé à un morphisme d'anneaux
$\varphi : R \to S$ et à un idéal premier $\mathfrak p \subset R$ est le suivant :
$$
\xymatrix{
\kappa(\mathfrak p) \otimes_R S =
S_{\mathfrak p}/{\mathfrak p}S_{\mathfrak p}
&
S_{\mathfrak p} \ar[l]
&
S \ar[r] \ar[l]
&
S/\mathfrak pS \ar[r]
&
(R \setminus \mathfrak p)^{-1}S/\mathfrak pS
\\
\kappa(\mathfrak p) =
R_{\mathfrak p}/{\mathfrak p}R_{\mathfrak p} \ar[u]
&
R_{\mathfrak p} \ar[u] \ar[l]
&
R \ar[u] \ar[r] \ar[l]
&
R/\mathfrak p \ar[u] \ar[r]
&
\kappa(\mathfrak p) \ar[u]
}
$$
Dans ce diagramme, les colonnes extrêmes gauche et droite sont identiques.
Sur les spectres, les morphismes horizontaux induisent des homéomorphismes
sur leurs images, et les carrés induisent des carrés cartésiens d'espaces
topologiques (voir les lemmes \ref{lemma-spec-localization}
et \ref{lemma-spec-closed}).
Cela montre que $\mathfrak p$ appartient à l'image de l'application
entre spectres si et seulement si
$S \otimes_R \kappa(\mathfrak p)$ n'est pas l'anneau nul.
S'il existe un idéal premier $\mathfrak q \subset S$ au-dessus de
$\mathfrak p$, c'est-à-dire tel que
$\mathfrak p = \varphi^{-1}(\mathfrak q)$,
alors nous pouvons compléter le diagramme en le diagramme suivant :
$$
\xymatrix{
\kappa(\mathfrak q) = S_{\mathfrak q}/{\mathfrak q}S_{\mathfrak q}
&
S_{\mathfrak q} \ar[l]
&
S \ar[r] \ar[l]
&
S/\mathfrak q \ar[r]
&
\kappa(\mathfrak q)
\\
\kappa(\mathfrak p) \otimes_R S =
S_{\mathfrak p}/{\mathfrak p}S_{\mathfrak p} \ar[u]
&
S_{\mathfrak p} \ar[u] \ar[l]
&
S \ar[u] \ar[r] \ar[l]
&
S/\mathfrak pS \ar[u] \ar[r]
&
(R \setminus \mathfrak p)^{-1}S/\mathfrak pS \ar[u]
\\
\kappa(\mathfrak p) =
R_{\mathfrak p}/{\mathfrak p}R_{\mathfrak p} \ar[u]
&
R_{\mathfrak p} \ar[u] \ar[l]
&
R \ar[u] \ar[r] \ar[l]
&
R/\mathfrak p \ar[u] \ar[r]
&
\kappa(\mathfrak p) \ar[u]
}
$$
Dans ce diagramme encore, les colonnes extrêmes gauche et droite sont
identiques et, sur les spectres, les morphismes horizontaux induisent
des homéomorphismes sur leurs images.
\end{remark}

\begin{lemma}
\label{lemma-in-image}
Soit $\varphi : R \to S$ un morphisme d'anneaux. Soit $\mathfrak p$
un idéal premier de $R$. Les conditions suivantes sont équivalentes :
\begin{enumerate}
\item $\mathfrak p$ appartient à l'image de
$\Spec(S) \to \Spec(R)$,
\item $S \otimes_R \kappa(\mathfrak p) \not = 0$,
\item $S_{\mathfrak p}/\mathfrak p S_{\mathfrak p} \not = 0$,
\item $(S/\mathfrak pS)_{\mathfrak p} \not = 0$, et
\item $\mathfrak p = \varphi^{-1}(\mathfrak pS)$.
\end{enumerate}
\end{lemma}

\begin{proof}
Nous avons déjà vu l'équivalence des deux premières conditions
dans la remarque \ref{remark-fundamental-diagram}. Les autres
n'en sont que des reformulations.
\end{proof}

\begin{remark}
\label{remark-local-ring-fibre}
Soit $R \to S$ un morphisme d'anneaux. Soit $\mathfrak q$ un idéal premier
de $S$ au-dessus de l'idéal premier $\mathfrak p$ de $R$.
D'après la remarque \ref{remark-fundamental-diagram},
l'idéal premier $\mathfrak q$ correspond à un unique idéal premier
$\overline{\mathfrak q}$ de l'anneau de la fibre
$F = S \otimes_R \kappa(\mathfrak p)$.
Alors nous avons
$$
F_{\overline{\mathfrak q}}
\cong S_\mathfrak q \otimes_{R_\mathfrak p} \kappa(\mathfrak p)
\cong S_\mathfrak q/\mathfrak p S_\mathfrak q
$$
En effet, il existe un morphisme d'anneaux évident
$F \to S_\mathfrak q \otimes_{R_\mathfrak p} \kappa(\mathfrak p)$
qui s'identifie aisément à $F \to F_{\overline{\mathfrak q}}$.
Le second isomorphisme résulte du fait que $\kappa(\mathfrak p)$
est le quotient de $R_\mathfrak p$ par $\mathfrak pR_\mathfrak p$.
\end{remark}

\section{Le radical de Jacobson d'un anneau}
\label{section-radical}

\noindent
Rappelons que le {\it radical de Jacobson} $\text{rad}(R)$ d'un anneau $R$
est l'intersection de tous les idéaux maximaux de $R$. Si $R$ est local,
alors $\text{rad}(R)$ est l'idéal maximal de $R$.

\begin{lemma}
\label{lemma-contained-in-radical}
Soit $R$ un anneau de radical de Jacobson $\text{rad}(R)$.
Soit $I \subset R$ un idéal. Les conditions suivantes sont
équivalentes :
\begin{enumerate}
\item $I \subset \text{rad}(R)$, et
\item tout élément de $1 + I$ est une unité de $R$.
\end{enumerate}
Dans ce cas, tout élément de $R$ dont l'image est une unité de $R/I$
est une unité.
\end{lemma}

\begin{proof}
Si $f \in \text{rad}(R)$, alors $f \in \mathfrak m$ pour tout
idéal maximal $\mathfrak m$ de $R$. Ainsi $1 + f \not \in \mathfrak m$
pour tout idéal maximal $\mathfrak m$ de $R$. Le fermé
$V(1 + f)$ de $\Spec(R)$ est donc vide. Cela entraîne
que $1 + f$ est une unité ; voir le lemme \ref{lemma-Zariski-topology}.

\medskip\noindent
Réciproquement, supposons que $1 + f$ soit une unité pour tout $f \in I$.
Si $\mathfrak m$ est un idéal maximal et si $I \not \subset \mathfrak m$,
alors $I + \mathfrak m = R$. Par conséquent,
$1 = f + g$ pour un certain $g \in \mathfrak m$ et un certain
$f \in I$. Mais $g = 1 + (-f)$ n'est pas une unité, contradiction.

\medskip\noindent
Pour la dernière assertion, soit $f \in R$ dont l'image est une unité de $R/I$.
Nous pouvons trouver $g \in R$ dont l'image est l'inverse multiplicatif
de $f \bmod I$. Alors $fg = 1 \bmod I$. Ainsi $fg$ est une unité de $R$
d'après (2), ce qui entraîne que $f$ est une unité.
\end{proof}

\begin{lemma}
\label{lemma-surjective-on-spec-units}
Soit $\varphi : R \to S$ un morphisme d'anneaux tel que l'application induite
$\Spec(S) \to \Spec(R)$ soit surjective. Alors un élément $x \in R$
est une unité si et seulement si $\varphi(x) \in S$ est une unité.
\end{lemma}

\begin{proof}
Si $x$ est une unité, alors $\varphi(x)$ l'est aussi. Réciproquement, si
$\varphi(x)$ est une unité, alors
$\varphi(x) \not \in \mathfrak q$ pour tout
$\mathfrak q \in \Spec(S)$. Par conséquent,
$x \not \in \varphi^{-1}(\mathfrak q) = \Spec(\varphi)(\mathfrak q)$
pour tout $\mathfrak q \in \Spec(S)$. Puisque $\Spec(\varphi)$ est surjective,
nous concluons que $x$ est une unité d'après
la partie (17) du lemme \ref{lemma-Zariski-topology}.
\end{proof}












\section{Le lemme de Nakayama}
\label{section-nakayama}

\noindent
Nous citons \cite{MatCA} : « Ce lemme simple mais
important est dû à T.~Nakayama, G.~Azumaya et W.~Krull. La question de la priorité
est obscure et, bien qu'il soit habituellement appelé lemme de Nakayama, feu
le professeur Nakayama n'aimait pas ce nom. »

\begin{lemma}[Lemme de Nakayama]
\label{lemma-NAK}
\begin{reference}
\cite[1.M Lemme (NAK), page 11]{MatCA}
\end{reference}
\begin{history}
Nous citons \cite{MatCA} : « Ce lemme simple mais
important est dû à T.~Nakayama, G.~Azumaya et W.~Krull. La question de la priorité
est obscure et, bien qu'il soit habituellement appelé lemme de Nakayama, feu
le professeur Nakayama n'aimait pas ce nom. »
\end{history}
Soit $R$ un anneau de radical de Jacobson $\text{rad}(R)$.
Soit $M$ un $R$-module. Soit $I \subset R$
un idéal.
\begin{enumerate}
\item
\label{item-nakayama}
Si $IM = M$ et si $M$ est de type fini, alors il existe $f \in 1 + I$ tel que
$fM = 0$.
\item Si $IM = M$, si $M$ est de type fini et si $I \subset \text{rad}(R)$, alors $M = 0$.
\item Si $N, N' \subset M$, si $M = N + IN'$ et si $N'$ est de type fini,
alors il existe $f \in 1 + I$ tel que $fM \subset N$ et $M_f = N_f$.
\item Si $N, N' \subset M$, si $M = N + IN'$, si $N'$ est de type fini et si
$I \subset \text{rad}(R)$, alors $M = N$.
\item Si $N \to M$ est un morphisme de modules, si $N/IN \to M/IM$ est
surjectif et si $M$ est de type fini, alors il existe $f \in 1 + I$
tel que $N_f \to M_f$ soit surjectif.
\item Si $N \to M$ est un morphisme de modules, si $N/IN \to M/IM$ est
surjectif, si $M$ est de type fini et si $I \subset \text{rad}(R)$,
alors $N \to M$ est surjectif.
\item Si $x_1, \ldots, x_n \in M$ engendrent $M/IM$ et si $M$ est de type fini,
alors il existe $f \in 1 + I$ tel que $x_1, \ldots, x_n$
engendrent $M_f$ comme $R_f$-module.
\item Si $x_1, \ldots, x_n \in M$ engendrent $M/IM$, si $M$ est de type fini et si
$I \subset \text{rad}(R)$, alors $M$ est engendré par $x_1, \ldots, x_n$.
\item Si $IM = M$ et si $I$ est nilpotent, alors $M = 0$.
\item Si $N, N' \subset M$, si $M = N + IN'$ et si $I$ est nilpotent, alors $M = N$.
\item Si $N \to M$ est un morphisme de modules, si $I$ est nilpotent et si $N/IN \to M/IM$
est surjectif, alors $N \to M$ est surjectif.
\item Si $\{x_\alpha\}_{\alpha \in A}$ est un ensemble d'éléments de $M$
qui engendrent $M/IM$ et si $I$ est nilpotent, alors $M$ est engendré
par les $x_\alpha$.
\end{enumerate}
\end{lemma}

\begin{proof}
Démonstration de (\ref{item-nakayama}). Choisissons des générateurs $y_1, \ldots, y_m$ de $M$
comme $R$-module. Pour chaque $i$, nous pouvons écrire $y_i = \sum z_{ij} y_j$ avec
$z_{ij} \in I$ (puisque $M = IM$).
Autrement dit, $\sum_j (\delta_{ij} - z_{ij})y_j = 0$.
Soit $f$ le déterminant de la matrice $m \times m$
$A = (\delta_{ij} - z_{ij})$. Notons que $f \in 1 + I$
(puisque la matrice $A$ est congrue terme à terme à la
matrice identité $m \times m$ modulo $I$).
D'après le lemme \ref{lemma-matrix-left-inverse} (1),
il existe une matrice $m \times m$
$B$ telle que $BA = f 1_{m \times m}$. En développant, nous voyons que
$\sum_{i} b_{hi} a_{ij} = f \delta_{hj}$ pour tous
$h$ et $j$ ; ainsi, $\sum_{i, j} b_{hi} a_{ij} y_j
= \sum_{j} f \delta_{hj} y_j = f y_h$ pour tout $h$.
Autrement dit, $0 = f y_h$ pour tout $h$ (puisque chaque
$i$ vérifie $\sum_j a_{ij} y_j = 0$).
Cela entraîne que $f$ annule $M$.

\medskip\noindent
D'après le lemme \ref{lemma-contained-in-radical}, tout élément de $1 + \text{rad}(R)$ est
un élément inversible de $R$. Nous voyons donc que (\ref{item-nakayama}) entraîne
(2). Nous obtenons (3) en appliquant (1) à $M/N$, qui est de type fini puisque $N'$ l'est.
Nous obtenons (4) en appliquant (2) à $M/N$, qui est de type fini puisque $N'$ l'est.
Nous obtenons (5) en appliquant (3) à $M$ et aux sous-modules $\Im(N \to M)$
et $M$. Nous obtenons (6) en appliquant (4) à $M$ et aux sous-modules
$\Im(N \to M)$ et $M$.
Nous obtenons (7) en appliquant (5) au morphisme $R^{\oplus n} \to M$,
$(a_1, \ldots, a_n) \mapsto a_1x_1 + \ldots + a_nx_n$.
Nous obtenons (8) en appliquant (6) au morphisme $R^{\oplus n} \to M$,
$(a_1, \ldots, a_n) \mapsto a_1x_1 + \ldots + a_nx_n$.

\medskip\noindent
L'assertion (9) est vraie car, si $M = IM$, alors $M = I^nM$ pour tout $n \geq 0$,
et la nilpotence de $I$ signifie que $I^n = 0$ pour un certain $n \gg 0$. Les assertions
(10), (11) et (12) résultent de (9) par les arguments utilisés ci-dessus.
\end{proof}

\begin{lemma}
\label{lemma-NAK-localization}
Soit $R$ un anneau, soit $S \subset R$ une partie multiplicative,
soit $I \subset R$ un idéal et soit $M$ un $R$-module de type fini.
Si $x_1, \ldots, x_r \in M$ engendrent $S^{-1}(M/IM)$
comme $S^{-1}(R/I)$-module, alors il existe $f \in S + I$
tel que $x_1, \ldots, x_r$ engendrent $M_f$ comme
$R_f$-module.\footnote{Cas particuliers : (I) $I = 0$. Le lemme affirme que,
si $x_1, \ldots, x_r$ engendrent $S^{-1}M$, alors $x_1, \ldots, x_r$
engendrent $M_f$ pour un certain $f \in S$. (II) $I = \mathfrak p$ est
un idéal premier et $S = R \setminus \mathfrak p$. Le lemme affirme que, si
$x_1, \ldots, x_r$ engendrent $M \otimes_R \kappa(\mathfrak p)$,
alors $x_1, \ldots, x_r$ engendrent $M_f$ pour un certain
$f \in R$, $f \not \in \mathfrak p$.}
\end{lemma}

\begin{proof}
Cas particulier $I = 0$. Soient $y_1, \ldots, y_s$ des générateurs de $M$ comme $R$-module.
Puisque $S^{-1}M$ est engendré par $x_1, \ldots, x_r$, pour chaque $i$,
nous pouvons écrire $y_i = \sum (a_{ij}/s_{ij})x_j$ dans $S^{-1}M$
pour certains $a_{ij} \in R$ et $s_{ij} \in S$. En multipliant par le produit
$s \in S$ des $s_{ij}$, nous voyons que $sy_i = \sum a'_{ij}x_j$
dans $S^{-1}M$ pour certains $a'_{ij} \in R$.
Cela signifie à son tour qu'il existe des $t_i \in S$ tels que
$t_isy_i = \sum t_ia'_{ij}x_j$ dans $M$. Ainsi, si $t \in S$
est le produit des $t_i$, alors nous
voyons que $y_i$ appartient au sous-$R_{st}$-module engendré par $x_1, \ldots, x_r$
dans $M_{st}$. Par conséquent, $x_1, \ldots, x_r$ engendrent $M_{st}$.

\medskip\noindent
Cas général. D'après le cas particulier, nous pouvons trouver $s \in S$
tel que $x_1, \ldots, x_r$ engendrent $(M/IM)_s$ comme $(R/I)_s$-module.
D'après le lemme \ref{lemma-NAK}, nous pouvons trouver $g \in 1 + I_s \subset R_s$
tel que $x_1, \ldots, x_r$ engendrent $(M_s)_g$ comme $(R_s)_g$-module.
Écrivons $g = 1 + i/s'$. Alors $f = ss' + is$ convient ; nous omettons les détails.
\end{proof}

\begin{lemma}
\label{lemma-when-surjective-local}
Soit $A \to B$ un morphisme local d'anneaux locaux.
Supposons que
\begin{enumerate}
\item $B$ soit de type fini comme $A$-module,
\item $\mathfrak m_B$ soit un idéal de type fini,
\item $A \to B$ induise un isomorphisme sur les corps résiduels, et
\item $\mathfrak m_A/\mathfrak m_A^2 \to \mathfrak m_B/\mathfrak m_B^2$
soit surjectif.
\end{enumerate}
Alors $A \to B$ est surjectif.
\end{lemma}

\begin{proof}
Pour montrer que $A \to B$ est surjectif, considérons-le comme un morphisme de $A$-modules
et appliquons le lemme \ref{lemma-NAK} (6). Nous en concluons qu'il suffit
de montrer que $A/\mathfrak m_A \to B/\mathfrak m_AB$ est surjectif.
Comme $A/\mathfrak m_A = B/\mathfrak m_B$, il suffit de montrer que
$\mathfrak m_AB \to \mathfrak m_B$ est surjectif. Considérons
$\mathfrak m_AB \to \mathfrak m_B$ comme un morphisme de $B$-modules et appliquons
le lemme \ref{lemma-NAK} (6). Nous en concluons qu'il suffit de voir que
$\mathfrak m_AB/\mathfrak m_A\mathfrak m_B \to \mathfrak m_B/\mathfrak m_B^2$
est surjectif. Cela résulte de l'hypothèse (4).
\end{proof}







\section{Sous-ensembles ouverts et fermés des spectres}
\label{section-open-and-closed}

\noindent
Il se trouve que les sous-ensembles à la fois ouverts et fermés d'un spectre
correspondent aux idempotents de l'anneau.

\begin{lemma}
\label{lemma-idempotent-spec}
Soit $R$ un anneau. Soit $e \in R$ un idempotent.
Dans ce cas,
$$
\Spec(R) = D(e) \amalg D(1-e).
$$
\end{lemma}

\begin{proof}
Notons qu'un idempotent $e$ d'un anneau intègre est égal à $1$ ou à $0$.
Nous voyons donc que
\begin{eqnarray*}
D(e)
& = &
\{ \mathfrak p \in \Spec(R)
\mid
e \not\in \mathfrak p \} \\
& = &
\{ \mathfrak p \in \Spec(R)
\mid
e \not = 0\text{ dans }\kappa(\mathfrak p) \} \\
& = &
\{ \mathfrak p \in \Spec(R)
\mid
e = 1\text{ dans }\kappa(\mathfrak p) \}
\end{eqnarray*}
De même, nous avons
\begin{eqnarray*}
D(1-e)
& = &
\{ \mathfrak p \in \Spec(R)
\mid
1 - e \not\in \mathfrak p \} \\
& = &
\{ \mathfrak p \in \Spec(R)
\mid
e \not = 1\text{ dans }\kappa(\mathfrak p) \} \\
& = &
\{ \mathfrak p \in \Spec(R)
\mid
e = 0\text{ dans }\kappa(\mathfrak p) \}
\end{eqnarray*}
Puisque l'image de $e$ dans tout corps résiduel est égale à $1$ ou à $0$,
nous en déduisons que $D(e)$ et $D(1-e)$ recouvrent tout $\Spec(R)$.
\end{proof}

\begin{lemma}
\label{lemma-spec-product}
Soient $R_1$ et $R_2$ des anneaux.
Soit $R = R_1 \times R_2$.
Les morphismes $R \to R_1$, $(x, y) \mapsto x$ et $R \to R_2$,
$(x, y) \mapsto y$
induisent des applications continues $\Spec(R_1) \to \Spec(R)$ et
$\Spec(R_2) \to \Spec(R)$.
L'application ainsi obtenue
$$
\Spec(R_1) \amalg \Spec(R_2)
\longrightarrow
\Spec(R)
$$
est un homéomorphisme. Autrement dit,
le spectre de $R = R_1\times R_2$ est la
réunion disjointe du spectre de $R_1$ et du
spectre de $R_2$.
\end{lemma}

\begin{proof}
Écrivons $1 = e_1 + e_2$ avec $e_1 = (1, 0)$ et $e_2 = (0, 1)$.
Notons que $e_1$ et $e_2 = 1 - e_1$ sont des idempotents.
Nous laissons au lecteur le soin de montrer que
$R_1 = R_{e_1}$ est la localisation de $R$ en $e_1$.
Il en va de même pour $e_2$.
Ainsi, l'assertion du lemme résulte du lemme
\ref{lemma-idempotent-spec} combiné au lemme
\ref{lemma-standard-open}.
\end{proof}

\noindent
Nous redémontrerons le lemme suivant plus loin, après avoir introduit
un lemme de recollement pour les fonctions. Voir la section
\ref{section-tilde-module-sheaf}.

\begin{lemma}
\label{lemma-disjoint-decomposition}
Soit $R$ un anneau. Pour tout $U \subset \Spec(R)$
qui est à la fois ouvert et fermé,
il existe un unique idempotent $e \in R$ tel que
$U = D(e)$. Cela induit une correspondance bijective entre les
sous-ensembles à la fois ouverts et fermés $U \subset \Spec(R)$ et les
idempotents $e \in R$.
\end{lemma}

\begin{proof}
Soit $U \subset \Spec(R)$ un sous-ensemble ouvert et fermé.
Puisque $U$ est fermé, il est quasi-compact d'après le
lemme \ref{lemma-quasi-compact}, et il en va de même pour
son complémentaire.
Écrivons $U = \bigcup_{i = 1}^n D(f_i)$ comme réunion finie d'ouverts principaux.
De même, écrivons $\Spec(R) \setminus U = \bigcup_{j = 1}^m D(g_j)$
comme réunion finie d'ouverts principaux. Puisque $\emptyset =
D(f_i) \cap D(g_j) = D(f_i g_j)$, nous voyons que $f_i g_j$ est
nilpotent d'après le lemme \ref{lemma-Zariski-topology}.
Soit $I = (f_1, \ldots, f_n) \subset R$ et soit
$J = (g_1, \ldots, g_m) \subset R$.
Notons que $V(J)$ est égal à $U$, que $V(I)$
est égal au complémentaire de $U$, de sorte que $\Spec(R) = V(I) \amalg V(J)$.
D'après la remarque ci-dessus concernant la nilpotence,
nous voyons que $(IJ)^N = (0)$ pour un entier $N$ suffisamment grand.
Puisque $\bigcup D(f_i) \cup \bigcup D(g_j) = \Spec(R)$,
nous voyons que $I + J = R$ ; voir le lemme \ref{lemma-Zariski-topology}.
En élevant cette égalité à la puissance $2N$, nous en concluons que
$I^N + J^N = R$. Écrivons $1 = x + y$ avec $x \in I^N$ et $y \in J^N$.
Alors $0 = xy = x(1 - x)$ puisque $I^N J^N = (0)$. Ainsi $x = x^2$
est idempotent et appartient
à $I^N \subset I$. L'idempotent $y = 1 - x$ appartient à $J^N \subset J$.
Cela montre que l'idempotent $x$ a pour image $1$ dans tout corps résiduel
$\kappa(\mathfrak p)$ pour $\mathfrak p \in V(J)$ et que $x$ a pour image $0$
dans $\kappa(\mathfrak p)$ pour tout $\mathfrak p \in V(I)$.

\medskip\noindent
Pour voir l'unicité, supposons que $e_1, e_2$ soient des
idempotents distincts de $R$. Nous devons montrer qu'il
existe un idéal premier $\mathfrak p$ tel que $e_1 \in \mathfrak p$
et $e_2 \not \in \mathfrak p$, ou l'inverse.
Écrivons $e_i' = 1 - e_i$. Si $e_1 \not = e_2$, alors
$0 \not = e_1 - e_2  = e_1(e_2 + e_2') - (e_1 + e_1')e_2
= e_1 e_2' - e_1' e_2$. Par conséquent, soit l'idempotent
$e_1 e_2' \not = 0$, soit $e_1' e_2 \not = 0$. Un idempotent non nul
n'est pas nilpotent ; nous trouvons donc un idéal premier
$\mathfrak p$ tel que soit $e_1e_2' \not \in \mathfrak p$,
soit $e_1'e_2 \not \in \mathfrak p$, d'après le lemme \ref{lemma-Zariski-topology}.
Il est facile de voir que cela fournit l'idéal premier recherché.
\end{proof}

\begin{lemma}
\label{lemma-characterize-spec-connected}
Soit $R$ un anneau non nul. Alors $\Spec(R)$ est
connexe si et seulement si $R$ ne possède aucun idempotent
non trivial.
\end{lemma}

\begin{proof}
Cela résulte immédiatement du lemme \ref{lemma-disjoint-decomposition}
et de la définition d'un espace topologique connexe.
\end{proof}

\begin{lemma}
\label{lemma-ideal-is-squared-union-connected}
Soit $I \subset R$ un idéal de type fini d'un anneau $R$
tel que $I = I^2$. Alors
\begin{enumerate}
\item il existe un idempotent $e \in R$ tel que $I = (e)$,
\item $R/I \cong R_{e'}$ pour l'idempotent $e' = 1 - e \in R$, et
\item $V(I)$ est ouvert et fermé dans $\Spec(R)$.
\end{enumerate}
\end{lemma}

\begin{proof}
D'après le lemme de Nakayama \ref{lemma-NAK}, il existe un élément
$f = 1 + i$, $i \in I$ tel que $fI = 0$. Alors $f^2 = f + fi = f$
est un idempotent. Considérons l'idempotent $e = 1 - f = -i \in I$.
Pour $j \in I$, nous avons $ej = j - fj = j$, donc $I = (e)$.
Cela démontre (1).

\medskip\noindent
Les parties (2) et (3) résultent de (1). En effet, nous avons
$V(I) = V(e) = \Spec(R) \setminus D(e)$, qui est ouvert et
fermé d'après le lemme \ref{lemma-idempotent-spec} ou le
lemme \ref{lemma-disjoint-decomposition}. Cela démontre (3).
Pour (2), observons que le morphisme $R \to R_{e'}$ est surjectif
puisque $x/(e')^n = x/e' = xe'/(e')^2 = xe'/e' = x/1$ dans $R_{e'}$.
Le noyau du morphisme $R \to R_{e'}$ est l'ensemble des éléments de
$R$ annulés par une puissance strictement positive de $e'$. Puisque $e'$ est
idempotent, c'est l'idéal des éléments annulés par $e'$,
qui est l'idéal $I = (e)$ puisque $e + e' = 1$ est une paire
d'idempotents orthogonaux. Cela démontre (2).
\end{proof}






\section{Composantes connexes des spectres}
\label{section-connected-components}

\noindent
Les composantes connexes des spectres ne sont pas aussi faciles à comprendre
qu'on pourrait le penser de prime abord. Cela vient de ce que nous sommes habitués à la topologie des
espaces localement connexes, alors que le spectre d'un anneau n'est en général
pas localement connexe.

\begin{lemma}
\label{lemma-closed-union-connected-components}
Soit $R$ un anneau. Soit $T \subset \Spec(R)$ un sous-ensemble du spectre.
Les conditions suivantes sont équivalentes :
\begin{enumerate}
\item $T$ est fermé et est une réunion de composantes connexes de
$\Spec(R)$,
\item $T$ est une intersection de sous-ensembles ouverts et fermés de
$\Spec(R)$, et
\item $T = V(I)$, où $I \subset R$ est un idéal engendré par des idempotents.
\end{enumerate}
De plus, l'idéal de (3), s'il existe, est unique.
\end{lemma}

\begin{proof}
D'après
le lemme \ref{lemma-topology-spec}
et
Topologie, lemme \ref{topology-lemma-closed-union-connected-components},
nous voyons que (1) et (2) sont équivalentes.
Supposons (2) et écrivons $T = \bigcap U_\alpha$ avec
$U_\alpha \subset \Spec(R)$ ouvert et fermé.
Alors $U_\alpha = D(e_\alpha)$ pour un certain idempotent $e_\alpha \in R$ d'après le
lemme \ref{lemma-disjoint-decomposition}.
En posant alors $I = (1 - e_\alpha)$, nous voyons que $T = V(I)$, c'est-à-dire que (3) est satisfaite.
Enfin, supposons (3). Écrivons $T = V(I)$ et $I = (e_\alpha)$ pour une certaine
famille d'idempotents $e_\alpha$. Il est alors clair que
$T = \bigcap V(e_\alpha) = \bigcap D(1 - e_\alpha)$.

\medskip\noindent
Supposons que $I$ soit un idéal engendré par des idempotents.
Soit $e \in R$ un idempotent tel que $V(I) \subset V(e)$. Alors, d'après le
lemme \ref{lemma-Zariski-topology},
nous voyons que $e^n \in I$ pour un certain $n \geq 1$. Comme $e$ est un idempotent,
cela signifie que $e \in I$. Nous voyons donc que $I$ est engendré par
exactement les idempotents $e$ tels que $T \subset V(e)$.
Autrement dit, l'idéal $I$ est entièrement déterminé par le
sous-ensemble fermé $T$, ce qui démontre l'unicité.
\end{proof}

\begin{lemma}
\label{lemma-connected-component}
Soit $R$ un anneau.
Une composante connexe de
$\Spec(R)$ est de la forme $V(I)$,
où $I$ est un idéal engendré par des idempotents
tel que tout idempotent de $R$ ait pour image
$0$ ou $1$ dans $R/I$.
\end{lemma}

\begin{proof}
Soit $\mathfrak p$ un idéal premier de $R$. D'après le
lemme \ref{lemma-topology-spec},
nous voyons que les hypothèses de
Topologie, lemme \ref{topology-lemma-connected-component-intersection},
sont satisfaites pour l'espace topologique $\Spec(R)$.
Par conséquent, la composante connexe de $\mathfrak p$ dans $\Spec(R)$
est l'intersection des sous-ensembles ouverts et fermés de $\Spec(R)$
qui contiennent $\mathfrak p$. Elle est donc égale à $V(I)$, où
$I$ est engendré par les idempotents $e \in R$, tels que $e$ ait pour image $0$
dans $\kappa(\mathfrak p)$ ; voir le
lemme \ref{lemma-disjoint-decomposition}.
Tout idempotent $e$ qui n'appartient pas à cette famille a manifestement pour image $1$
dans $R/I$.
\end{proof}









\section{Propriétés de recollement}
\label{section-more-glueing}

\noindent
Dans cette section, nous rassemblons plusieurs résultats classiques de la
forme suivante : si une propriété est vraie sur tous les membres d'un
recouvrement par des ouverts principaux, alors elle est vraie. En fait, il suffit souvent de vérifier
les propriétés au niveau des anneaux locaux, comme dans le lemme suivant.

\begin{lemma}
\label{lemma-characterize-zero-local}
Soit $R$ un anneau.
\begin{enumerate}
\item Pour un élément $x$ d'un $R$-module $M$, les conditions suivantes sont équivalentes :
\begin{enumerate}
\item $x = 0$,
\item l'image de $x$ dans $M_\mathfrak p$ est nulle pour tout $\mathfrak p \in \Spec(R)$,
\item l'image de $x$ dans $M_{\mathfrak m}$ est nulle pour tout idéal maximal
$\mathfrak m$ de $R$.
\end{enumerate}
Autrement dit, le morphisme $M \to \prod_{\mathfrak m} M_{\mathfrak m}$
est injectif.
\item Pour un $R$-module $M$, les conditions suivantes sont équivalentes :
\begin{enumerate}
\item $M$ est nul,
\item $M_{\mathfrak p}$ est nul pour tout $\mathfrak p \in \Spec(R)$,
\item $M_{\mathfrak m}$ est nul pour tout idéal maximal $\mathfrak m$ de $R$.
\end{enumerate}
\item Pour un complexe $M_1 \to M_2 \to M_3$
de $R$-modules, les conditions suivantes sont équivalentes :
\begin{enumerate}
\item $M_1 \to M_2 \to M_3$ est exact,
\item pour tout idéal premier $\mathfrak p$ de $R$, le complexe localisé
$M_{1, \mathfrak p} \to M_{2, \mathfrak p} \to M_{3, \mathfrak p}$
est exact,
\item pour tout idéal maximal $\mathfrak m$ de $R$, le complexe localisé
$M_{1, \mathfrak m} \to M_{2, \mathfrak m} \to M_{3, \mathfrak m}$
est exact.
\end{enumerate}
\item Pour un morphisme $f : M \to M'$ de $R$-modules, les conditions suivantes sont équivalentes :
\begin{enumerate}
\item $f$ est injectif,
\item $f_{\mathfrak p} : M_\mathfrak p \to M'_\mathfrak p$ est injectif
pour tout idéal premier $\mathfrak p$ de $R$,
\item $f_{\mathfrak m} : M_\mathfrak m \to M'_\mathfrak m$ est injectif
pour tout idéal maximal $\mathfrak m$ de $R$.
\end{enumerate}
\item Pour un morphisme $f : M \to M'$ de $R$-modules, les conditions suivantes sont équivalentes :
\begin{enumerate}
\item $f$ est surjectif,
\item $f_{\mathfrak p} : M_\mathfrak p \to M'_\mathfrak p$ est surjectif
pour tout idéal premier $\mathfrak p$ de $R$,
\item $f_{\mathfrak m} : M_\mathfrak m \to M'_\mathfrak m$ est surjectif
pour tout idéal maximal $\mathfrak m$ de $R$.
\end{enumerate}
\item Pour un morphisme $f : M \to M'$ de $R$-modules, les conditions suivantes sont équivalentes :
\begin{enumerate}
\item $f$ est bijectif,
\item $f_{\mathfrak p} : M_\mathfrak p \to M'_\mathfrak p$ est bijectif
pour tout idéal premier $\mathfrak p$ de $R$,
\item $f_{\mathfrak m} : M_\mathfrak m \to M'_\mathfrak m$ est bijectif
pour tout idéal maximal $\mathfrak m$ de $R$.
\end{enumerate}
\end{enumerate}
\end{lemma}

\begin{proof}
Soit $x \in M$ comme dans (1). Soit $I = \{f \in R \mid fx = 0\}$.
Il est facile de voir que $I$ est un idéal (c'est
l'annulateur de $x$). La condition (1)(c) signifie que, pour
tout idéal maximal $\mathfrak m$, il existe
$f \in R \setminus \mathfrak m$ tel que $fx =0$.
Autrement dit, $V(I)$ ne contient aucun point fermé.
D'après le lemme \ref{lemma-Zariski-topology}, nous voyons que $I$ est l'idéal unité.
Ainsi $x$ est nul, c'est-à-dire que (1)(a) est satisfaite. Cela démontre (1).

\medskip\noindent
La partie (2) résulte de l'application simultanée de (1) à tous les éléments de $M$.

\medskip\noindent
Démonstration de (3). Soit $H$ l'homologie de la suite, c'est-à-dire
$H = \Ker(M_2 \to M_3)/\Im(M_1 \to M_2)$. D'après la
proposition \ref{proposition-localization-exact},
$H_\mathfrak p$ est l'homologie de la suite
$M_{1, \mathfrak p} \to M_{2, \mathfrak p} \to M_{3, \mathfrak p}$.
Ainsi (3) résulte de (2).

\medskip\noindent
Les parties (4) et (5) sont des cas particuliers de (3). La partie (6) résulte
formellement de la combinaison de (4) et (5).
\end{proof}

\begin{lemma}
\label{lemma-cover}
\begin{slogan}
Propriétés locales pour Zariski des modules et des algèbres
\end{slogan}
Soit $R$ un anneau. Soit $M$ un $R$-module. Soit $S$ une $R$-algèbre.
Supposons que $f_1, \ldots, f_n$ soit une liste finie
d'éléments de $R$ telle que $\bigcup D(f_i) = \Spec(R)$,
autrement dit $(f_1, \ldots, f_n) = R$.
\begin{enumerate}
\item Si chaque $M_{f_i} = 0$, alors $M = 0$.
\item Si chaque $M_{f_i}$ est un $R_{f_i}$-module de type fini,
alors $M$ est un $R$-module de type fini.
\item Si chaque $M_{f_i}$ est un $R_{f_i}$-module de présentation finie,
alors $M$ est un $R$-module de présentation finie.
\item Soit $M \to N$ un morphisme de $R$-modules. Si $M_{f_i} \to N_{f_i}$
est un isomorphisme pour tout $i$, alors $M \to N$ est un isomorphisme.
\item Soit $0 \to M'' \to M \to M' \to 0$ un complexe de $R$-modules.
Si $0 \to M''_{f_i} \to M_{f_i} \to M'_{f_i} \to 0$ est exact pour tout $i$,
alors $0 \to M'' \to M \to M' \to 0$ est exact.
\item Si chaque $R_{f_i}$ est noethérien, alors $R$ est noethérien.
\item Si chaque $S_{f_i}$ est une $R_{f_i}$-algèbre de type fini, alors
$S$ est une $R$-algèbre de type fini.
\item Si chaque $S_{f_i}$ est de présentation finie sur $R_{f_i}$, alors
$S$ est une $R$-algèbre de présentation finie.
\end{enumerate}
\end{lemma}

\begin{proof}
Démontrons successivement chacune des parties.
\begin{enumerate}
\item D'après la proposition \ref{proposition-localize-twice},
cela entraîne que $M_\mathfrak p = 0$ pour tout $\mathfrak p \in \Spec(R)$ ;
nous concluons donc par le lemme \ref{lemma-characterize-zero-local}.
\item Pour chaque $i$, choisissons un ensemble fini de générateurs $X_i$ de $M_{f_i}$.
Nous pouvons supposer sans perte de généralité que les éléments de $X_i$
appartiennent à l'image du morphisme de localisation $M \rightarrow M_{f_i}$ ; prenons donc
un ensemble fini $Y_i$ d'antécédents des éléments de $X_i$ dans $M$. Soit $Y$
la réunion de ces ensembles. C'est encore un ensemble fini.
Considérons le morphisme $R$-linéaire évident $R^Y \rightarrow M$ qui envoie l'élément
de base $e_y$ sur $y$. D'après l'hypothèse, ce morphisme est surjectif après localisation
en tout idéal premier $\mathfrak p$ de $R$ ; il est donc surjectif d'après le
lemme \ref{lemma-characterize-zero-local},
et $M$ est de type fini.
\item D'après (2), nous avons une suite exacte courte
$$
0 \rightarrow K \rightarrow R^n \rightarrow M \rightarrow 0
$$
Puisque la localisation est un foncteur exact et que $M_{f_i}$ est de présentation
finie, nous voyons que $K_{f_i}$ est de type fini pour tout
$1 \leq i \leq n$ d'après le lemme \ref{lemma-extension}.
D'après (2), cela entraîne que $K$ est un $R$-module de type fini et donc que
$M$ est de présentation finie.
\item D'après la proposition \ref{proposition-localize-twice},
l'hypothèse entraîne que le morphisme induit
sur les localisés en tous les idéaux premiers est un isomorphisme ; nous concluons donc
par le lemme \ref{lemma-characterize-zero-local}.
\item D'après la proposition \ref{proposition-localize-twice}, l'hypothèse
entraîne que la
suite des localisés en tous les idéaux premiers est exacte courte ; nous
concluons donc par le lemme \ref{lemma-characterize-zero-local}.
\item Nous allons montrer que tout idéal de $R$ possède un ensemble fini de générateurs.
Pour cela, soit $I \subset R$ un idéal arbitraire. D'après la
proposition \ref{proposition-localization-exact},
chaque $I_{f_i} \subset R_{f_i}$ est un idéal. Ils sont tous
de type fini par hypothèse, et nous concluons donc par (2).
\item Pour chaque $i$, choisissons un ensemble fini de générateurs $X_i$ de $S_{f_i}$.
Nous pouvons supposer sans perte de généralité que les éléments de $X_i$
appartiennent à l'image du morphisme de localisation $S \rightarrow S_{f_i}$ ; prenons donc
un ensemble fini $Y_i$ d'antécédents des éléments de $X_i$ dans
$S$. Soit $Y$ la réunion de ces ensembles. C'est encore un ensemble fini.
Considérons le morphisme d'algèbres $R[X_y]_{y \in Y} \rightarrow S$
induit par $Y$. Puisqu'il s'agit d'un morphisme d'algèbres, l'image $T$
est un $R$-sous-module du $R$-module $S$ ; nous pouvons donc considérer le
module quotient $S/T$. D'après l'hypothèse, celui-ci est nul après localisation
en chacun des $f_i$ ; il est donc nul d'après (1), et $S$ est par conséquent une
$R$-algèbre de type fini.
\item D'après le point précédent, il existe un morphisme surjectif de $R$-algèbres
$R[X_1, \ldots, X_n] \rightarrow S$. Soit $K$ le noyau
de ce morphisme. C'est un idéal de $R[X_1, \ldots, X_n]$, de type fini
après chacune des localisations en $f_i$. Puisque les $f_i$ engendrent l'idéal unité
dans $R$, ils engendrent également l'idéal unité dans $R[X_1, \ldots, X_n]$ ; une
application de (2) achève donc la démonstration.
\end{enumerate}
\end{proof}

\begin{lemma}
\label{lemma-cover-upstairs}
Soit $R \to S$ un morphisme d'anneaux.
Supposons que $g_1, \ldots, g_n$ soit une liste finie
d'éléments de $S$ telle que $\bigcup D(g_i) = \Spec(S)$,
autrement dit $(g_1, \ldots, g_n) = S$.
\begin{enumerate}
\item Si chaque $S_{g_i}$ est de type fini sur $R$, alors $S$ est
de type fini sur $R$.
\item Si chaque $S_{g_i}$ est de présentation finie sur $R$,
alors $S$ est de présentation finie sur $R$.
\end{enumerate}
\end{lemma}

\begin{proof}
Choisissons $h_1, \ldots, h_n \in S$ tels que $\sum h_i g_i = 1$.

\medskip\noindent
Démonstration de (1). Pour chaque $i$, choisissons une liste finie d'éléments
$x_{i, j} \in S_{g_i}$, $j = 1, \ldots, m_i$,
qui engendrent $S_{g_i}$ comme $R$-algèbre.
Écrivons $x_{i, j} = y_{i, j}/g_i^{n_{i, j}}$ pour un certain $y_{i, j} \in S$ et
un certain $n_{i, j} \ge 0$. Considérons la $R$-sous-algèbre $S' \subset S$
engendrée par $g_1, \ldots, g_n$, $h_1, \ldots, h_n$ et
les $y_{i, j}$, $i = 1, \ldots, n$, $j = 1, \ldots, m_i$.
Puisque la localisation est exacte (proposition \ref{proposition-localization-exact}),
nous voyons que $S'_{g_i} \to S_{g_i}$ est injectif.
D'autre part, il est surjectif par notre choix des $y_{i, j}$.
Les éléments $g_1, \ldots, g_n$ engendrent l'idéal unité dans $S'$
puisque $h_1, \ldots, h_n \in S'$.
Ainsi, $S' \to S$, considéré comme un morphisme de $S'$-modules, est un isomorphisme
d'après le lemme \ref{lemma-cover}.

\medskip\noindent
Démonstration de (2). Nous savons déjà que $S$ est de type fini.
Écrivons $S = R[x_1, \ldots, x_m]/J$ pour un certain idéal $J$.
Pour chaque $i$, choisissons un relèvement $g'_i \in R[x_1, \ldots, x_m]$ de $g_i$
et choisissons un relèvement $h'_i \in R[x_1, \ldots, x_m]$ de $h_i$.
Nous voyons alors que
$$
S_{g_i} = R[x_1, \ldots, x_m, y_i]/(J_i + (1 - y_ig'_i))
$$
où $J_i$ est l'idéal de $R[x_1, \ldots, x_m, y_i]$
engendré par $J$. Nous omettons un petit détail. D'après le
lemme \ref{lemma-finite-presentation-independent},
nous pouvons choisir une liste finie d'éléments
$f_{i, j} \in J$, $j = 1, \ldots, m_i$,
telle que les images des $f_{i, j}$ dans $J_i$ ainsi que $1 - y_ig'_i$
engendrent l'idéal $J_i + (1 - y_ig'_i)$.
Posons
$$
S' = R[x_1, \ldots, x_m]/\left(\sum h'_ig'_i - 1, f_{i, j}; 
i = 1, \ldots, n, j = 1, \ldots, m_i\right)
$$
Il existe un morphisme surjectif de $R$-algèbres $S' \to S$.
Les classes des éléments $g'_1, \ldots, g'_n$ dans $S'$
engendrent l'idéal unité et, par construction, les morphismes
$S'_{g'_i} \to S_{g_i}$ sont injectifs.
Nous concluons donc comme dans la partie (1).
\end{proof}





\section{Recollement de fonctions}
\label{section-tilde-module-sheaf}

\noindent
Dans cette section, nous montrons qu'étant donné un recouvrement ouvert
$$
\Spec(R) = \bigcup\nolimits_{i = 1}^n D(f_i)
$$
par des ouverts principaux, et étant donné un élément $h_i \in R_{f_i}$
pour chaque $i$, tel que $h_i = h_j$ comme éléments de $R_{f_i f_j}$,
il existe un unique $h \in R$ tel que l'image de
$h$ dans $R_{f_i}$ soit $h_i$. Ce résultat peut s'interpréter
de deux manières :
\begin{enumerate}
\item La règle $D(f) \mapsto R_f$ est un faisceau d'anneaux
sur les ouverts principaux ; voir Faisceaux, section \ref{sheaves-section-bases}.
\item Si nous considérons les éléments de $R_f$ comme les fonctions « algébriques »
ou « régulières » sur $D(f)$, alors celles-ci se recollent
comme le feraient les fonctions continues, resp.\ différentiables,
sur une variété topologique, resp.\ différentiable.
\end{enumerate}

\begin{lemma}
\label{lemma-cover-module}
Soit $R$ un anneau. Soient $f_1, \ldots, f_n$ des éléments de $R$
qui engendrent l'idéal unité. Soit $M$ un $R$-module.
La suite
$$
0 \to
M \xrightarrow{\alpha}
\bigoplus\nolimits_{i = 1}^n M_{f_i} \xrightarrow{\beta}
\bigoplus\nolimits_{i, j = 1}^n M_{f_i f_j}
$$
est exacte, où $\alpha(m) = (m/1, \ldots, m/1)$
et $\beta(m_1/f_1^{e_1}, \ldots, m_n/f_n^{e_n})
= (m_i/f_i^{e_i} - m_j/f_j^{e_j})_{(i, j)}$.
\end{lemma}

\begin{proof}
Il suffit de montrer que la localisation de la suite en
tout idéal maximal $\mathfrak m$ est exacte ; voir le
lemme \ref{lemma-characterize-zero-local}.
Puisque $f_1, \ldots, f_n$ engendrent l'idéal unité,
il existe $i$ tel que $f_i \not \in \mathfrak m$.
Après renumérotation, nous pouvons supposer que $i = 1$.
Notons que $(M_{f_i})_\mathfrak m = (M_\mathfrak m)_{f_i}$
et que $(M_{f_if_j})_\mathfrak m = (M_\mathfrak m)_{f_if_j}$ ; voir la
proposition \ref{proposition-localize-twice-module}.
En particulier, $(M_{f_1})_\mathfrak m = M_\mathfrak m$ et
$(M_{f_1 f_i})_\mathfrak m = (M_\mathfrak m)_{f_i}$, car
$f_1$ est une unité.
Notons que les morphismes de la suite sont les morphismes canoniques
provenant du
lemme \ref{lemma-universal-property-localization-module}
et du morphisme identité de $M$.
Cela étant dit, après avoir remplacé $R$ par $R_\mathfrak m$,
$M$ par $M_\mathfrak m$, les $f_i$ par leurs images dans $R_\mathfrak m$
et $f_1$ par $1 \in R_\mathfrak m$,
nous nous ramenons au cas où $f_1 = 1$.

\medskip\noindent
Supposons $f_1 = 1$. L'injectivité de $\alpha$ est alors immédiate. Soit
$m = (m_i) \in \bigoplus_{i = 1}^n M_{f_i}$ un élément du noyau de $\beta$.
Alors $m_1 \in M_{f_1} = M$. De plus, $\beta(m) = 0$
entraîne que $m_1$ et $m_i$ ont la même image dans
$M_{f_1f_i} = M_{f_i}$. Ainsi $\alpha(m_1) = m$, ce qui
achève la démonstration.
\end{proof}

\begin{lemma}
\label{lemma-standard-covering}
Soit $R$ un anneau, et soient $f_1, f_2, \ldots f_n\in R$ des éléments qui engendrent
l'idéal unité de $R$.
Alors la suite suivante est exacte :
$$
0 \longrightarrow
R \longrightarrow
\bigoplus\nolimits_i R_{f_i} \longrightarrow
\bigoplus\nolimits_{i, j}R_{f_if_j}
$$
où les morphismes $\alpha : R \longrightarrow \bigoplus_i R_{f_i}$
et $\beta : \bigoplus_i R_{f_i} \longrightarrow \bigoplus_{i, j} R_{f_if_j}$
sont définis par
$$
\alpha(x) = \left(\frac{x}{1}, \ldots, \frac{x}{1}\right)
\text{ et }
\beta\left(\frac{x_1}{f_1^{r_1}}, \ldots, \frac{x_n}{f_n^{r_n}}\right)
=
\left(\frac{x_i}{f_i^{r_i}}-\frac{x_j}{f_j^{r_j}}~\text{dans}~R_{f_if_j}\right).
$$
\end{lemma}

\begin{proof}
C'est un cas particulier du lemme \ref{lemma-cover-module}.
\end{proof}

\noindent
Nous avons déjà vu le résultat suivant ci-dessus, mais nous l'énonçons explicitement ici
pour plus de commodité.

\begin{lemma}
\label{lemma-disjoint-implies-product}
Soit $R$ un anneau.
Si $\Spec(R) = U \amalg V$, où $U$ et $V$ sont tous deux ouverts,
alors $R \cong R_1 \times R_2$, avec $U \cong \Spec(R_1)$
et $V \cong \Spec(R_2)$ via les morphismes du lemme \ref{lemma-spec-product}.
De plus, $R_1$ et $R_2$ sont tous deux des localisations ainsi que des quotients
de l'anneau $R$.
\end{lemma}

\begin{proof}
D'après le lemme \ref{lemma-disjoint-decomposition}, nous avons
$U = D(e)$ et $V = D(1-e)$ pour un certain idempotent $e$.
D'après le lemme \ref{lemma-standard-covering}, nous voyons que
$R \cong R_e \times R_{1 - e}$ (puisque, manifestement, $R_{e(1-e)} = 0$,
de sorte que la condition de recollement est triviale ; bien entendu, il est
également immédiat de démontrer directement la décomposition en produit dans ce
cas). Le lemme en résulte.
\end{proof}

\begin{lemma}
\label{lemma-when-injective-covering}
Soit $R$ un anneau.
Soient $f_1, \ldots, f_n \in R$.
Soit $M$ un $R$-module.
Alors $M \to \bigoplus M_{f_i}$ est injectif si et seulement si
$$
M \longrightarrow \bigoplus\nolimits_{i = 1, \ldots, n} M, \quad
m \longmapsto (f_1m, \ldots, f_nm)
$$
est injectif.
\end{lemma}

\begin{proof}
Le morphisme $M \to \bigoplus M_{f_i}$ est injectif si et seulement si,
pour tous $m \in M$ et $e_1, \ldots, e_n \geq 1$ tels que
$f_i^{e_i}m = 0$, $i = 1, \ldots, n$, nous avons $m = 0$.
Cela entraîne manifestement que le morphisme affiché est injectif.
Réciproquement, supposons le morphisme affiché injectif et que
$m \in M$ et $e_1, \ldots, e_n \geq 1$ soient tels que
$f_i^{e_i}m = 0$, $i = 1, \ldots, n$. Si $e_i = 1$ pour tout $i$,
alors l'injectivité du morphisme affiché donne immédiatement $m = 0$.
Démontrons ensuite que cela vaut pour toutes données de ce type
par récurrence sur $e = \sum e_i$. Le cas initial est $e = n$, et nous venons
de le traiter. Si un certain $e_i > 1$, posons $m' = f_im$.
Par récurrence, nous voyons que $m' = 0$. Ainsi $f_i m = 0$,
c'est-à-dire que nous pouvons prendre $e_i = 1$, ce qui diminue $e$ et achève la récurrence.
\end{proof}

\noindent
Le lemme suivant s'énonce et se démontre mieux dans le cadre plus général
de la descente plate. Il est néanmoins naturel de l'énoncer ici,
car il s'accorde bien avec ce qui précède.

\begin{lemma}
\label{lemma-glue-modules}
Soit $R$ un anneau. Soient $f_1, \ldots, f_n \in R$. Supposons données
les données suivantes :
\begin{enumerate}
\item Pour chaque $i$, un $R_{f_i}$-module $M_i$.
\item Pour chaque paire $i, j$, un isomorphisme de $R_{f_if_j}$-modules
$\psi_{ij} : (M_i)_{f_j} \to (M_j)_{f_i}$.
\end{enumerate}
qui satisfont la « condition de cocycle » selon laquelle tous les diagrammes
$$
\xymatrix{
(M_i)_{f_jf_k}
\ar[rd]_{\psi_{ij}}
\ar[rr]^{\psi_{ik}}
& &
(M_k)_{f_if_j} \\
&
(M_j)_{f_if_k} \ar[ru]_{\psi_{jk}}
}
$$
sont commutatifs (pour tous les triplets $i, j, k$). À partir de ces données, définissons
$$
M = \Ker\left(
\bigoplus\nolimits_{1 \leq i \leq n} M_i
\longrightarrow
\bigoplus\nolimits_{1 \leq i, j \leq n} (M_i)_{f_j}
\right)
$$
où $(m_1, \ldots, m_n)$ est envoyé sur l'élément dont
la composante d'indice $(i, j)$ est $m_i/1 - \psi_{ji}(m_j/1)$.
Alors le morphisme naturel $M \to M_i$ induit un isomorphisme
$M_{f_i} \to M_i$. De plus, $\psi_{ij}(m/1) = m/1$
pour tout $m \in M$ (avec les notations évidentes).
\end{lemma}

\begin{proof}
Pour montrer que $M_{f_1} \to M_1$ est un isomorphisme, il suffit
de montrer que sa localisation en tout idéal premier $\mathfrak p'$
de $R_{f_1}$ est un isomorphisme ; voir le
lemme \ref{lemma-characterize-zero-local}.
Écrivons $\mathfrak p' = \mathfrak p R_{f_1}$
pour un certain idéal premier $\mathfrak p \subset R$, avec $f_1 \not \in \mathfrak p$ ; voir le
lemme \ref{lemma-standard-open}.
Puisque la localisation est exacte
(proposition \ref{proposition-localization-exact}),
nous voyons que
\begin{align*}
(M_{f_1})_{\mathfrak p'} & =
M_\mathfrak p \\
& =
\Ker\left(
\bigoplus\nolimits_{1 \leq i \leq n} M_{i, \mathfrak p}
\longrightarrow
\bigoplus\nolimits_{1 \leq i, j \leq n} ((M_i)_{f_j})_\mathfrak p
\right) \\
& =
\Ker\left(
\bigoplus\nolimits_{1 \leq i \leq n} M_{i, \mathfrak p}
\longrightarrow
\bigoplus\nolimits_{1 \leq i, j \leq n} (M_{i, \mathfrak p})_{f_j}
\right)
\end{align*}
Nous avons également utilisé ici la proposition \ref{proposition-localize-twice-module}.
Puisque $f_1$ est une unité de $R_\mathfrak p$, nous nous ramenons au cas
où $f_1 = 1$ en remplaçant $R$ par $R_\mathfrak p$, les $f_i$ par les
images des $f_i$ dans $R_\mathfrak p$, $M$ par $M_\mathfrak p$ et
$f_1$ par $1$.

\medskip\noindent
Supposons $f_1 = 1$. Alors $\psi_{1j} : (M_1)_{f_j} \to M_j$
est un isomorphisme pour $j = 2, \ldots, n$. Si nous utilisons ces
isomorphismes pour identifier $M_j = (M_1)_{f_j}$, nous voyons alors
que $\psi_{ij} : (M_1)_{f_if_j} \to (M_1)_{f_if_j}$ est
l'identification canonique. Ainsi, le complexe
$$
0 \to M_1 \to
\bigoplus\nolimits_{1 \leq i \leq n} (M_1)_{f_i}
\longrightarrow
\bigoplus\nolimits_{1 \leq i, j \leq n}
(M_1)_{f_if_j}
$$
est exact d'après le lemme \ref{lemma-cover-module}.
Ainsi, le premier morphisme identifie $M_1$ à $M$ dans ce cas,
et tout est clair.
\end{proof}












\section{Diviseurs de zéro et anneaux totaux des fractions}
\label{section-total-quotient-ring}

\noindent
L'anneau local en un idéal premier minimal possède les propriétés suivantes.

\begin{lemma}
\label{lemma-minimal-prime-reduced-ring}
Soit $\mathfrak p$ un idéal premier minimal d'un anneau $R$.
Tout élément de l'idéal maximal de $R_{\mathfrak p}$
est nilpotent. Si $R$ est réduit, alors $R_{\mathfrak p}$
est un corps.
\end{lemma}

\begin{proof}
Si un élément $x$ de ${\mathfrak p}R_{\mathfrak p}$
n'est pas nilpotent, alors $D(x) \not = \emptyset$ ; voir le
lemme \ref{lemma-Zariski-topology}. Cela contredit
la minimalité de $\mathfrak p$. Si $R$ est réduit,
alors ${\mathfrak p}R_{\mathfrak p} = 0$ et
l'anneau local est donc un corps.
\end{proof}

\begin{lemma}
\label{lemma-reduced-ring-sub-product-fields}
Soit $R$ un anneau réduit. Alors
\begin{enumerate}
\item $R$ est un sous-anneau d'un produit de corps,
\item $R \to \prod_{\mathfrak p\text{ minimal}} R_{\mathfrak p}$
est un plongement dans un produit de corps,
\item $\bigcup_{\mathfrak p\text{ minimal}} \mathfrak p$ est l'ensemble
des diviseurs de zéro de $R$.
\end{enumerate}
\end{lemma}

\begin{proof}
D'après le lemme \ref{lemma-minimal-prime-reduced-ring}, chacun des anneaux
$R_\mathfrak p$ est un corps. En particulier, le noyau du morphisme d'anneaux
$R \to R_\mathfrak p$ est $\mathfrak p$.
D'après le lemme \ref{lemma-Zariski-topology},
nous avons $\bigcap_{\mathfrak p} \mathfrak p = (0)$.
Ainsi (2) et (1) sont vraies. Si $x y = 0$ et $y \not = 0$, alors
$y \not \in \mathfrak p$ pour un certain idéal premier minimal $\mathfrak p$.
Par conséquent, $x \in \mathfrak p$. Ainsi, tout diviseur de zéro de $R$ appartient
à $\bigcup_{\mathfrak p\text{ minimal}} \mathfrak p$.
Réciproquement, supposons que $x \in \mathfrak p$ pour un certain idéal
premier minimal $\mathfrak p$. Alors l'image de $x$ dans $R_\mathfrak p$ est nulle ;
il existe donc $y \in R$, $y \not \in \mathfrak p$ tel que
$xy = 0$. Autrement dit, $x$ est un diviseur de zéro. Cela achève la
démonstration de (3) et du lemme.
\end{proof}

\noindent
L'anneau total des fractions $Q(R)$ d'un anneau $R$ a été introduit dans
l'exemple \ref{example-localize-at-prime}.

\begin{lemma}
\label{lemma-total-ring-fractions}
Soit $R$ un anneau.
Soit $S \subset R$ une partie multiplicative formée d'éléments non diviseurs de zéro.
Alors $Q(R) \cong Q(S^{-1}R)$.
En particulier, $Q(R) \cong Q(Q(R))$.
\end{lemma}

\begin{proof}
Si $x \in S^{-1}R$ est un élément non diviseur de zéro et si
$x = r/f$ pour certains $r \in R$, $f \in S$, alors
$r$ est un élément non diviseur de zéro de $R$. Le lemme en résulte.
\end{proof}

\noindent
Nous pouvons appliquer les résultats de recollement pour démontrer une propriété des
anneaux totaux des fractions $Q(R)$, que nous avons introduits dans
l'exemple \ref{example-localize-at-prime}.

\begin{lemma}
\label{lemma-total-ring-fractions-no-embedded-points}
Soit $R$ un anneau.
Supposons que $R$ possède un nombre fini d'idéaux premiers minimaux
$\mathfrak q_1, \ldots, \mathfrak q_t$, et que
$\mathfrak q_1 \cup \ldots \cup \mathfrak q_t$ soit l'ensemble
des diviseurs de zéro de $R$.
Alors l'anneau total des fractions $Q(R)$ est égal à
$R_{\mathfrak q_1} \times \ldots \times R_{\mathfrak q_t}$.
\end{lemma}

\begin{proof}
Il existe des morphismes naturels $Q(R) \to R_{\mathfrak q_i}$ puisque
tout élément non diviseur de zéro appartient à $R \setminus \mathfrak q_i$.
Il existe donc un morphisme naturel
$Q(R) \to R_{\mathfrak q_1} \times \ldots \times R_{\mathfrak q_t}$.
Pour tout idéal premier non minimal $\mathfrak p \subset R$, nous voyons que
$\mathfrak p \not \subset \mathfrak q_1 \cup \ldots \cup \mathfrak q_t$
d'après le lemme \ref{lemma-silly}. Par conséquent,
$\Spec(Q(R)) = \{\mathfrak q_1, \ldots, \mathfrak q_t\}$
(comme sous-ensembles de $\Spec(R)$ ; voir le lemme \ref{lemma-spec-localization}).
Ainsi $\Spec(Q(R))$ est un ensemble discret fini et
il s'ensuit que $Q(R) = A_1 \times \ldots \times A_t$,
avec $\Spec(A_i) = \{\mathfrak{q}_i\}$ ; voir le
lemme \ref{lemma-disjoint-implies-product}.
De plus, $A_i$ est un anneau local qui est une localisation
de $R$. Ainsi $A_i \cong R_{\mathfrak q_i}$.
\end{proof}

















\section{Composantes irréductibles des spectres}
\label{section-irreducible}

\noindent
Nous montrons que les composantes irréductibles du
spectre d'un anneau correspondent aux
idéaux premiers minimaux de l'anneau.

\begin{lemma}
\label{lemma-irreducible}
Soit $R$ un anneau.
\begin{enumerate}
\item Pour un idéal premier $\mathfrak p \subset R$, l'adhérence
de $\{\mathfrak p\}$ dans la topologie de Zariski est $V(\mathfrak p)$.
Sous forme de formule, $\overline{\{\mathfrak p\}} = V(\mathfrak p)$.
\item Les sous-ensembles fermés irréductibles de $\Spec(R)$ sont
exactement les sous-ensembles $V(\mathfrak p)$, où $\mathfrak p \subset R$
est un idéal premier.
\item Les composantes irréductibles (voir Topologie,
définition \ref{topology-definition-irreducible-components})
de $\Spec(R)$ sont exactement les sous-ensembles $V(\mathfrak p)$,
où $\mathfrak p \subset R$ est un idéal premier minimal.
\end{enumerate}
\end{lemma}

\begin{proof}
Notons que, si $ \mathfrak p \in V(I)$, alors
$I \subset \mathfrak p$. Par conséquent,
il est clair que $\overline{\{\mathfrak p\}} = V(\mathfrak p)$.
En particulier, $V(\mathfrak p)$ est l'adhérence d'un
singleton et est donc irréductible.
La deuxième assertion entraîne la troisième.
Pour démontrer la deuxième, soit
$V(I) \subset \Spec(R)$, où $I$ est un idéal radical.
Si $I$ n'est pas premier, choisissons $a, b\in R$, $a, b\not \in I$
avec $ab\in I$. Dans ce cas, $V(I, a) \cup V(I, b) = V(I)$,
mais ni $V(I, b) = V(I)$ ni $V(I, a) = V(I)$, d'après le
lemme \ref{lemma-Zariski-topology}. Ainsi $V(I)$ n'est pas
irréductible.
\end{proof}

\noindent
Autrement dit, ce lemme montre que tout sous-ensemble fermé irréductible
de $\Spec(R)$ est de la forme $V(\mathfrak p)$ pour
un certain idéal premier $\mathfrak p$. Puisque $V(\mathfrak p) = \overline{\{\mathfrak p\}}$,
nous voyons que tout sous-ensemble fermé irréductible possède un unique point générique ;
voir Topologie, définition \ref{topology-definition-generic-point}.
En particulier, $\Spec(R)$ est un espace topologique sobre.
Consignons ce fait dans le lemme suivant.

\begin{lemma}
\label{lemma-spec-spectral}
Le spectre d'un anneau est un espace spectral ; voir Topologie, définition
\ref{topology-definition-spectral-space}.
\end{lemma}

\begin{proof}
Formellement, cela résulte du lemme \ref{lemma-irreducible} et du
lemme \ref{lemma-topology-spec}. Voir aussi la discussion ci-dessus.
\end{proof}

\begin{lemma}
\label{lemma-irreducible-components-containing-x}
Soit $R$ un anneau. Soit $\mathfrak p \subset R$ un idéal premier.
\begin{enumerate}
\item L'ensemble des sous-ensembles fermés irréductibles de $\Spec(R)$
passant par $\mathfrak p$ est en correspondance bijective avec les
idéaux premiers $\mathfrak q \subset R_{\mathfrak p}$.
\item L'ensemble des composantes irréductibles de $\Spec(R)$ passant par
$\mathfrak p$ est en correspondance bijective avec les idéaux
premiers minimaux $\mathfrak q \subset R_{\mathfrak p}$.
\end{enumerate}
\end{lemma}

\begin{proof}
Cela résulte du lemme \ref{lemma-irreducible}
et de la description de $\Spec(R_\mathfrak p)$ donnée dans le
lemme \ref{lemma-spec-localization}, qui montre que
$\Spec(R_\mathfrak p)$ correspond aux idéaux premiers $\mathfrak q$ de $R$
tels que $\mathfrak q \subset \mathfrak p$.
\end{proof}

\begin{lemma}
\label{lemma-standard-open-containing-maximal-point}
Soit $R$ un anneau.
Soit $\mathfrak p$ un idéal premier minimal de $R$.
Soit $W \subset \Spec(R)$ un ouvert quasi-compact
qui ne contient pas le point $\mathfrak p$. Alors il
existe $f \in R$, $f \not \in \mathfrak p$, tel
que $D(f) \cap W = \emptyset$.
\end{lemma}

\begin{proof}
Puisque $W$ est quasi-compact, nous pouvons l'écrire comme réunion finie
d'ouverts affines principaux $D(g_i)$, $i = 1, \ldots, n$.
Puisque $\mathfrak p \not \in W$, nous avons $g_i \in \mathfrak p$ pour
tout $i$. D'après le lemme \ref{lemma-minimal-prime-reduced-ring},
chaque $g_i$ est nilpotent dans $R_{\mathfrak p}$. Nous pouvons donc trouver
$f \in R$, $f \not \in \mathfrak p$, tel que, pour tout $i$, nous ayons
$f g_i^{n_i} = 0$ pour un certain $n_i > 0$. Alors $D(f)$ convient.
\end{proof}

\begin{lemma}
\label{lemma-ring-with-only-minimal-primes}
Soit $R$ un anneau. Soit $X = \Spec(R)$ comme espace topologique.
Les conditions suivantes sont équivalentes :
\begin{enumerate}
\item $X$ est profini,
\item $X$ est séparé,
\item $X$ est totalement discontinu.
\item tout ouvert quasi-compact de $X$ est fermé,
\item il n'existe aucune inclusion non triviale entre ses idéaux premiers,
\item tout idéal premier est un idéal maximal,
\item tout idéal premier est minimal,
\item tout ouvert principal $D(f) \subset X$ est fermé, et
\item ajouter d'autres conditions ici.
\end{enumerate}
\end{lemma}

\begin{proof}
Première démonstration. Il est clair que (5), (6) et (7) sont équivalentes.
Il est clair que (4) et (8) sont équivalentes, puisque tout ouvert quasi-compact
est une réunion finie d'ouverts principaux.
L'implication (7) $\Rightarrow$ (4) résulte du
lemme \ref{lemma-standard-open-containing-maximal-point}.
Supposons (4). Soient $\mathfrak p, \mathfrak p'$ deux idéaux premiers
distincts de $R$. Choisissons $f \in \mathfrak p'$, $f \not \in \mathfrak p$
(au besoin, échangeons $\mathfrak p$ et $\mathfrak p'$).
Alors $\mathfrak p' \not \in D(f)$ et $\mathfrak p \in D(f)$.
D'après (4), l'ouvert $D(f)$ est également fermé.
Ainsi $\mathfrak p$ et $\mathfrak p'$ appartiennent à des voisinages ouverts
disjoints dont la réunion est $X$. Par conséquent, $X$ est séparé et totalement
discontinu. Ainsi (4) $\Rightarrow$ (2) et (3).
Si (3) est satisfaite, il ne peut exister aucune spécialisation
entre les points de $\Spec(R)$, et nous voyons que (5) est satisfaite.
Si $X$ est séparé, tout point est fermé, donc (2) entraîne (6).
Ainsi (2), (3), (4), (5), (6), (7) et (8) sont équivalentes.
Tout espace profini est séparé, donc (1) entraîne (2).
Si $X$ satisfait (2) et (3), alors $X$ (qui est quasi-compact d'après le
lemme \ref{lemma-quasi-compact}) est profini d'après
Topologie, lemme \ref{topology-lemma-profinite}.

\medskip\noindent
Deuxième démonstration. Outre l'équivalence de (4) et (8), cela résulte
du lemme \ref{lemma-spec-spectral} et de faits purement topologiques ; voir
Topologie, lemme \ref{topology-lemma-characterize-profinite-spectral}.
\end{proof}













\section{Exemples de spectres d'anneaux}
\label{section-examples-spectra}

\noindent
Dans cette section, nous donnons quelques exemples de spectres.

\begin{example}
\label{example-spec-Zxmodx2minus4}
Dans cet exemple, nous décrivons $X = \Spec(\mathbf{Z}[x]/(x^2 - 4))$.
Soit $\mathfrak{p}$ un idéal premier arbitraire de $X$.
Soit $\phi : \mathbf{Z} \to \mathbf{Z}[x]/(x^2 - 4)$ le morphisme d'anneaux naturel.
Alors $ \phi^{-1}(\mathfrak p)$ est un idéal premier de $\mathbf{Z}$.
Si $ \phi^{-1}(\mathfrak p) = (2)$, alors, puisque $\mathfrak p$ contient $2$,
il correspond à un idéal premier de
$\mathbf{Z}[x]/(x^2 - 4, 2) \cong (\mathbf{Z}/2\mathbf{Z})[x]/(x^2)$
via le morphisme $ \mathbf{Z}[x]/(x^2 - 4) \to  \mathbf{Z}[x]/(x^2 - 4, 2)$.
Tout idéal premier de $(\mathbf{Z}/2\mathbf{Z})[x]/(x^2)$ correspond à un idéal premier
de $(\mathbf{Z}/2\mathbf{Z})[x]$ contenant $(x^2)$. Un tel idéal premier
contient alors $x$. Puisque
$(\mathbf{Z}/2\mathbf{Z}) \cong (\mathbf{Z}/2\mathbf{Z})[x]/(x)$ est un corps,
$(x)$ est un idéal maximal. Puisque tout idéal premier contient $(x)$ et que $(x)$ est
maximal, l'anneau ne possède qu'un seul idéal premier, $(x)$. Ainsi, dans ce cas,
$\mathfrak p = (2, x)$. Maintenant, si $ \phi^{-1}(\mathfrak p) = (q)$ pour
$q > 2$, alors, puisque $\mathfrak p$ contient $q$, il correspond à un
idéal premier de
$\mathbf{Z}[x]/(x^2 - 4, q) \cong (\mathbf{Z}/q\mathbf{Z})[x]/(x^2 - 4)$
via le morphisme $ \mathbf{Z}[x]/(x^2 - 4) \to  \mathbf{Z}[x]/(x^2 - 4, q)$.
Tout idéal premier de $(\mathbf{Z}/q\mathbf{Z})[x]/(x^2 - 4)$ correspond à un
idéal premier de $(\mathbf{Z}/q\mathbf{Z})[x]$ contenant $(x^2 - 4) = (x -2)(x + 2)$.
Par conséquent, ces idéaux premiers doivent contenir soit $x -2$, soit $x + 2$. Puisque
$(\mathbf{Z}/q\mathbf{Z})[x]$ est un anneau principal, tous les idéaux premiers
non nuls sont maximaux, et il existe donc
exactement deux idéaux premiers de $(\mathbf{Z}/q\mathbf{Z})[x]$ contenant
$(x-2)(x + 2)$, à savoir $(x-2)$ et $(x + 2)$. En conclusion, il existe deux
idéaux premiers $(q, x-2)$ et $(q, x + 2)$ puisque $2 \neq -2 \in \mathbf{Z}/(q)$.
Enfin, traitons le cas où $\phi^{-1}(\mathfrak p) = (0)$. Remarquons
que $\mathfrak p$ correspond à un idéal premier de $\mathbf{Z}[x]$ qui
contient $(x^2 - 4) = (x -2)(x + 2)$. Ainsi, $\mathfrak p$ contient soit
$(x-2)$, soit $(x + 2)$. Par conséquent, $\mathfrak p$ correspond à un idéal premier de
$\mathbf{Z}[x]/(x - 2)$ ou de $\mathbf{Z}[x]/(x + 2)$ dont l'intersection
avec $\mathbf{Z}$ est réduite à $0$, par hypothèse. Puisque
$\mathbf{Z}[x]/(x - 2) \cong \mathbf{Z}$ et
$\mathbf{Z}[x]/(x + 2) \cong \mathbf{Z}$, cela signifie que $\mathfrak p$
doit correspondre à $0$ dans l'un de ces anneaux. Ainsi,
$\mathfrak p = (x - 2)$ ou $\mathfrak p = (x + 2)$ dans l'anneau initial.
\end{example}

\begin{example}
\label{example-spec-Zx}
Dans cet exemple, nous décrivons $X = \Spec(\mathbf{Z}[x])$.
Fixons $\mathfrak p \in X$.
Soit $\phi : \mathbf{Z} \to \mathbf{Z}[x]$ et remarquons
que $\phi^{-1}(\mathfrak p) \in \Spec(\mathbf{Z})$.
Si $\phi^{-1}(\mathfrak p) = (q)$ pour $q$ un nombre premier tel que $q > 0$,
alors $\mathfrak p$ correspond à un idéal premier de $(\mathbf{Z}/(q))[x]$,
qui doit être engendré par un polynôme irréductible dans
$(\mathbf{Z}/(q))[x]$. Si nous choisissons un représentant de ce polynôme
de degré minimal, il sera également irréductible dans $\mathbf{Z}[x]$.
Ainsi, dans ce cas, $\mathfrak p = (q, f_q)$, où $f_q$ est un polynôme
irréductible de $\mathbf{Z}[x]$ qui reste irréductible lorsqu'on le considère
dans $(\mathbf{Z}/(q) [x])$. Supposons maintenant que $\phi^{-1}(\mathfrak p) = (0)$.
Dans ce cas, si $\mathfrak p$ est non nul, il doit être engendré par des polynômes non constants
qui, puisque $\mathfrak p$ est premier, peuvent être supposés irréductibles dans
$\mathbf{Z}[x]$. D'après le lemme de Gauss, ces polynômes sont aussi irréductibles
dans $\mathbf{Q}[x]$. Puisque $\mathbf{Q}[x]$ est un anneau euclidien, s'il
existe au moins deux polynômes irréductibles distincts $f, g$ qui engendrent $\mathfrak p$,
alors $1 = af + bg$ pour certains $a, b \in \mathbf{Q}[x]$. En multipliant par
un dénominateur commun, nous voyons que $m = \bar{a}f + \bar{b} g$ pour
certains $\bar{a}, \bar{b} \in \mathbf{Z}[x]$ et un certain $m \in \mathbf{Z}$ non nul.
C'est une contradiction. Ainsi, lorsqu'il est non nul, $\mathfrak p$ est engendré par un
polynôme irréductible de $\mathbf{Z}[x]$.
\end{example}

\begin{example}
\label{example-spec-kxy}
Dans cet exemple, nous décrivons $X = \Spec(k[x, y])$
lorsque $k$ est un corps arbitraire.
Il est clair que $(0)$ est premier, et tout idéal principal engendré par un
polynôme irréductible est également premier puisque $k[x, y]$ est un
anneau factoriel. Supposons maintenant que $\mathfrak p$ soit un
élément de $X$ qui n'est pas principal. Puisque $k[x, y]$ est un
anneau factoriel noethérien, l'idéal premier $\mathfrak p$ peut être engendré
par un nombre fini de polynômes irréductibles $(f_1, \ldots, f_n)$.
Affirmons maintenant que, si $f, g$ sont des polynômes irréductibles de $k[x, y]$
qui ne sont pas associés, alors $(f, g) \cap k[x] \neq 0$. Pour le montrer,
il suffit d'établir que $f$ et $g$ sont premiers entre eux lorsqu'on les
considère dans $k(x)[y]$. Dans ce cas, $k(x)[y]$ est un anneau euclidien ;
en appliquant l'algorithme d'Euclide et en chassant les dénominateurs, nous
obtenons $p = af + bg$ pour $p, a, b \in k[x]$. Supposons donc que ce ne soit pas
le cas, c'est-à-dire qu'un certain élément non inversible $h \in k(x)[y]$ divise à la fois
$f$ et $g$. Alors, d'après le lemme de Gauss, pour certains $a, b \in k(x)$, nous
avons $ah | f$ et $bh | g$ avec $ah, bh \in k[x]$. Par
irréductibilité, $ah = f$ et
$bh = g$ (puisque $h \notin k(x)$). Ainsi, de retour dans $k(x)[y]$, $f, g $
sont associés, car $\frac{a}{b} g = f$. Puisque
$k(x)$ est le corps des fractions de $k[x]$, nous pouvons écrire $g = \frac{r}{s} f $
pour des éléments $r , s \in k[x]$ sans facteur commun. Cela
entraîne que $sg = rf$ dans $k[x, y]$, et donc que $s$ doit diviser $f$,
puisque $k[x, y]$ est un anneau factoriel. Ainsi, $s = 1$ ou $s = f$. Si $s = f$,
alors $r = g$, ce qui entraîne $f, g \in k[x]$ ; ils doivent donc être des unités de
$k(x)$ et être premiers entre eux dans $k(x)[y]$, en contradiction avec notre
hypothèse. Si $s = 1$, alors $g = rf$, autre contradiction.
Ainsi, $f, g$ doivent être premiers entre eux dans $k(x)[y]$, qui est un
anneau euclidien. Nous nous sommes donc ramenés au cas où $\mathfrak p$
contient un polynôme irréductible $p \in k[x] \subset k[x, y]$.
D'après ce qui précède, $\mathfrak p$ correspond à un idéal premier de l'anneau
$k[x, y]/(p) = k(\alpha)[y]$, où $\alpha$ est un élément
algébrique sur $k$ de polynôme minimal $p$. Cet anneau est
principal, et tout idéal premier correspond donc à $(0)$ ou à un
polynôme irréductible de $k(\alpha)[y]$. Ainsi, $\mathfrak p$
est de la forme $(p)$ ou $(p, f)$, où $f$ est un
polynôme de $k[x, y]$ qui est irréductible dans le quotient
$k[x, y]/(p)$.
\end{example}

\begin{example}
\label{example-affine-open-not-standard}
Considérons l'anneau
$$
R = \{ f \in \mathbf{Q}[z]\text{ vérifiant }f(0) = f(1) \}.
$$
Considérons le morphisme
$$
\varphi : \mathbf{Q}[A, B] \to R
$$
défini par $\varphi(A) = z^2-z$ et $\varphi(B) = z^3-z^2$. On vérifie
aisément que $(A^3 - B^2 + AB) \subset \Ker(\varphi)$ et que
$A^3 - B^2 + AB$ est irréductible. Supposons que $\varphi$ soit surjectif ;
alors, puisque $R$ est un anneau intègre (c'est un sous-anneau d'un anneau
intègre), $\Ker(\varphi)$ doit être un idéal premier de $\mathbf{Q}[A, B]$.
Les idéaux premiers qui contiennent $(A^3-B^2 + AB)$ sont $(A^3-B^2 + AB)$
lui-même et tout idéal maximal $(f, g)$, avec $f, g\in\mathbf{Q}[A, B]$,
tel que $f$ soit irréductible modulo $g$. Mais $R$ n'est pas un corps, donc le
noyau doit être $(A^3-B^2 + AB)$ ; ainsi, $\varphi$ donne un isomorphisme
$R \to \mathbf{Q}[A, B]/(A^3-B^2 + AB)$.

\medskip\noindent
Pour voir que $\varphi$ est surjectif, nous devons exprimer tout
$f\in R$ comme un polynôme à coefficients dans $\mathbf{Q}$ en $A(z) = z^2-z$
et $B(z) = z^3-z^2$. Notons la relation $zA(z) = B(z)$. Soit
$a = f(0) = f(1)$. Alors $z(z-1)$ doit diviser $f(z)-a$, et nous pouvons donc écrire
$f(z) = z(z-1)g(z)+a = A(z)g(z)+a$. Si $\deg(g) < 2$, alors
$h(z) = c_1z + c_0$ et $f(z) = A(z)(c_1z + c_0)+a = c_1B(z)+c_0A(z)+a$, et nous
avons terminé. Si $\deg(g)\geq 2$, alors, d'après l'algorithme de division
euclidienne des polynômes, nous pouvons écrire $g(z) = A(z)h(z)+b_1z + b_0$
($\deg(h)\leq\deg(g)-2$), de sorte que $f(z) = A(z)^2h(z)+b_1B(z)+b_0A(z)$.
En appliquant la division à $h(z)$ et en itérant, nous obtenons une expression
de $f(z)$ comme polynôme en $A(z)$ et $B(z)$ ; ainsi, $\varphi$ est
surjectif.

\medskip\noindent
Soit maintenant $a \in \mathbf{Q}$, $a \neq 0, \frac{1}{2}, 1$, et
considérons
$$
R_a = \{ f \in \mathbf{Q}[z, \frac{1}{z-a}]\text{ vérifiant }f(0) = f(1)
\}.
$$
C'est encore une $\mathbf{Q}$-algèbre de type fini : on vérifie
aisément que les fonctions $z^2-z$, $z^3-z$ et
$\frac{a^2-a}{z-a}+z$ engendrent $R_a$ comme $\mathbf{Q}$-algèbre. Nous
avons les inclusions suivantes :
$$
R\subset R_a\subset\mathbf{Q}[z, \frac{1}{z-a}], \quad
R\subset\mathbf{Q}[z]\subset\mathbf{Q}[z, \frac{1}{z-a}].
$$
Rappelons (lemme \ref{lemma-spec-localization}) que, pour un anneau T et une
partie multiplicative $S\subset T$, le morphisme d'anneaux $T \to S^{-1}T$
induit un morphisme sur les spectres $\Spec(S^{-1}T) \to \Spec(T)$
qui est un homéomorphisme sur le sous-ensemble
$$
\{\mathfrak p \in \Spec(T) \mid S \cap \mathfrak p = \emptyset\}
\subset \Spec(T).
$$
Lorsque $S = \{ 1, f, f^2, \ldots\}$ pour un certain $f\in T$, il s'agit
de l'ouvert $D(f)\subset T$. Vérifions maintenant une propriété analogue
pour le morphisme d'anneaux $R \to R_a$ : nous allons montrer que le morphisme
$\theta : \Spec(R_a) \to \Spec(R)$ induit par l'inclusion
$R\subset R_a$ est un homéomorphisme sur un sous-ensemble ouvert de
$\Spec(R)$ en vérifiant que $\theta$ est un homéomorphisme local
injectif. Nous le faisons relativement à un recouvrement ouvert de
$\Spec(R_a)$ par deux ouverts principaux, que nous décrivons maintenant.
Pour tout $r\in\mathbf{Q}$, soit $\text{ev}_r : R \to \mathbf{Q}$ le
morphisme donné par l'évaluation en $r$. Notons que, pour $r = 0$ et
$r = 1-a$, celui-ci se prolonge en un morphisme
$\text{ev}_r' : R_a \to \mathbf{Q}$ (dans le second cas, parce que $\frac{1}{z-a}$
est bien défini en $z = 1-a$, puisque $a\neq\frac{1}{2}$). Toutefois,
$\text{ev}_a$ ne se prolonge pas à $R_a$. Posons
$\mathfrak{m}_r = \Ker(\text{ev}_r)$. Nous avons
$$
\mathfrak{m}_0 = (z^2-z, z^3-z),
$$
$$
\mathfrak{m}_a = ((z-1 + a)(z-a), (z^2-1 + a)(z-a)), \text{ et}
$$
$$
\mathfrak{m}_{1-a} = ((z-1 + a)(z-a), (z-1 + a)(z^2-a)).
$$
Pour le vérifier, notons que les membres de droite sont clairement contenus dans
les membres de gauche. Vérifions ensuite que les membres de droite sont des
idéaux maximaux en écrivant les générateurs en fonction de $A$ et $B$,
et en considérant $R$ comme $\mathbf{Q}[A, B]/(A^3-B^2 + AB)$. Notons que
$\mathfrak{m}_a$ n'appartient pas à l'image de $\theta$ : nous avons
$$
(z^2 - z)^2(z - a)\left(\frac{a^2 - a}{z - a} + z\right) =
(z^2 - z)^2(a^2 - a) + (z^2 - z)^2(z - a)z
$$
Le membre de gauche appartient à $\mathfrak m_a R_a$ parce que
$(z^2 - z)(z - a)$ appartient à $\mathfrak m_a$ et que
$(z^2 - z)(\frac{a^2 - a}{z - a} + z)$ appartient à $R_a$. De même,
l'élément $(z^2 - z)^2(z - a)z$ appartient à $\mathfrak m_a R_a$
parce que $(z^2 - z)$ appartient à $R_a$ et que $(z^2 - z)(z - a)$ appartient à $\mathfrak m_a$.
Comme $a \not \in \{0, 1\}$, nous en concluons que
$(z^2 - z)^2 \in \mathfrak m_a R_a$. Par conséquent,
aucun idéal $I$ de $R_a$ ne peut vérifier $I \cap R = \mathfrak m_a$, car un tel
$I$ devrait contenir $(z^2 - z)^2$, qui appartient à $R$ mais pas à
$\mathfrak m_a$. L'ouvert principal
$D((z-1 + a)(z-a))\subset\Spec(R)$ est égal au complémentaire du
fermé $\{\mathfrak{m}_a, \mathfrak{m}_{1-a}\}$.
Vérifions alors que $R_{(z-1 + a)(z-a)} = (R_a)_{(z-1 + a)(z-a)}$ ; en
notant cet anneau localisé $R'$, il s'ensuit que le morphisme $R \to R'$
se factorise en $R \to R_a \to R'$. D'après le lemme
\ref{lemma-spec-localization}, ces morphismes présentent alors
$\Spec(R') \subset \Spec(R_a)$ et
$\Spec(R') \subset \Spec(R)$ comme des sous-ensembles ouverts ; ainsi,
$\theta : \Spec(R_a) \to \Spec(R)$, restreint à
$D((z-1 + a)(z-a))$, est un homéomorphisme sur un sous-ensemble ouvert.
De même, la restriction de $\theta$ à
$D((z^2 + z + 2a-2)(z-a)) \subset \Spec(R_a)$ est un homéomorphisme
sur le sous-ensemble ouvert $D((z^2 + z + 2a-2)(z-a)) \subset \Spec(R)$.
Selon que $z^2 + z + 2a-2$ est irréductible ou non sur
$\mathbf{Q}$, le complémentaire de ce dernier ouvert principal
est constitué d'un ou deux points fermés, ainsi que du point fermé
$\mathfrak{m}_a$. De plus, l'idéal de $R_a$ engendré par les
éléments $(z^2 + z + 2a-a)(z-a)$ et $(z-1 + a)(z-a)$ est égal à $R_a$ tout entier ; ainsi,
ces deux ouverts principaux recouvrent $\Spec(R_a)$. Par conséquent, pour
montrer que $\theta$ est un homéomorphisme sur
$\Spec(R)-\{\mathfrak{m}_a\}$, il suffit de montrer
que ces un ou deux points ne peuvent jamais être égaux à $\mathfrak{m}_{1-a}$.
Et c'est bien le cas, puisque $1-a$ est une racine de $z^2 + z + 2a-2$
si et seulement si $a = 0$ ou $a = 1$, deux cas qui ne se produisent pas.

\medskip\noindent
Malgré cet homéomorphisme, qui imite le comportement d'une
localisation en un élément de $R$, tandis que
$\mathbf{Q}[z, \frac{1}{z-a}]$ est la localisation de $\mathbf{Q}[z]$
relativement à l'élément $(z-a)$, l'anneau $R_a$ n'est {\it pas} une
localisation de $R$ : toute localisation $S^{-1}R$ produit davantage
d'unités que l'anneau initial $R$. Les unités de $R$ sont
$\mathbf{Q}^\times$, les unités de $\mathbf{Q}$. En fait, il est
facile de voir que les unités de $R_a$ sont $\mathbf{Q}^*$.
En effet, les unités de $\mathbf{Q}[z, \frac{1}{z - a}]$ sont
$c (z - a)^n$ pour $c \in \mathbf{Q}^*$ et $n \in \mathbf{Z}$,
et il est clair que celles-ci n'appartiennent à $R_a$ que si $n = 0$.
Ainsi, $R_a$ ne possède pas davantage d'unités que
$R$ et ne peut donc pas être une localisation de $R$.

\medskip\noindent
Nous avons utilisé le fait que $a\neq 0, 1$ pour garantir que
$\frac{1}{z-a}$ est défini en $z = 0, 1$. Nous avons utilisé le fait que
$a\neq 1/2$ en plusieurs endroits : (1) pour pouvoir parler du
noyau de $\text{ev}_{1-a}$ sur $R_a$, ce qui garantit que
$\mathfrak{m}_{1-a}$ est un point de $R_a$ (c'est-à-dire que $R_a$
ne possède qu'un point de moins que $R$). (2) À la fin, pour conclure
que $(z-a)^{k + \ell}$ ne peut appartenir à $R$ que pour $k = \ell = 0$ ; en effet, si
$a = 1/2$, cet élément appartient à $R$ dès que $k + \ell$ est pair. Il
y aurait alors effectivement davantage d'unités dans $R_a$ que dans $R$, et $R_a$
pourrait éventuellement être une localisation de $R$.
\end{example}





\section{Une méta-observation sur les idéaux premiers}
\label{section-oka-families}

\noindent
Cette section est tirée du projet CRing. Soit $R$ un anneau et
soit $S \subset R$ une partie multiplicative.
Une conséquence du
lemme \ref{lemma-spec-localization}
est qu'un idéal $I \subset R$ maximal parmi ceux qui ne
rencontrent pas $S$ est premier. La raison en est que $I = R \cap \mathfrak m$
pour un certain idéal maximal $\mathfrak m$ de l'anneau $S^{-1}R$.
Il s'avère que, pour de nombreuses propriétés d'idéaux, les éléments maximaux
sont premiers. Une méthode générale permettant de le voir a été développée dans \cite{Lam-Reyes}.
Dans cette section, nous faisons une digression pour expliquer ce phénomène.

\medskip\noindent
Soit $R$ un anneau. Si $I$ est un idéal de $R$ et $a \in R$, nous
définissons
$$
(I : a) = \left\{ x \in R \mid xa \in I\right\}.
$$
Plus généralement, si $J \subset R$ est un idéal, nous définissons
$$
(I : J) = \left\{ x \in R \mid xJ \subset I\right\}.
$$

\begin{lemma}
\label{lemma-colon}
Soit $R$ un anneau. Pour un idéal principal $J \subset R$ et pour tout idéal
$I \subset J$, nous avons $I = J (I : J)$.
\end{lemma}

\begin{proof}
Écrivons $J = (a)$. Alors $(I : J) = (I : a)$.
Puisque $I \subset J$, nous voyons que tout $y \in I$ est de la forme
$y = xa$ pour un certain $x \in (I : a)$. Ainsi, $I \subset J (I : J)$.
Réciproquement, si $x \in (I : a)$, alors $xJ = (xa) \subset I$, ce qui
démontre l'autre inclusion.
\end{proof}

\noindent
Soit $\mathcal{F}$ une collection d'idéaux de $R$. Nous cherchons des
conditions garantissant que les éléments maximaux du complémentaire
de $\mathcal{F}$ sont premiers.

\begin{definition}
\label{definition-oka-family}
Soit $R$ un anneau. Soit $\mathcal{F}$ un ensemble d'idéaux de $R$. Nous disons
que $\mathcal{F}$ est une {\it famille d'Oka} si $R \in \mathcal{F}$ et si,
chaque fois que $I \subset R$ est un idéal et que $(I : a), (I, a) \in \mathcal{F}$
pour un certain $a \in R$, alors $I \in \mathcal{F}$.
\end{definition}

\noindent
Donnons quelques exemples de familles d'Oka. Le premier est l'exemple
fondamental présenté dans l'introduction de cette section.

\begin{example}
\label{example-oka-family-not-meet-multiplicative-set}
Soit $R$ un anneau et soit $S$ une partie multiplicative de $R$.
Nous affirmons que $\mathcal{F} = \{I \subset R \mid I \cap S \not = \emptyset\}$
est une famille d'Oka. En effet, supposons que $(I : a), (I, a) \in \mathcal{F}$
pour un certain $a \in R$. Choisissons alors $s \in (I, a) \cap S$ et
$s' \in (I : a) \cap S$. Alors $ss' \in I \cap S$, et donc
$I \in \mathcal{F}$. Ainsi, $\mathcal{F}$ est une famille d'Oka.
\end{example}

\begin{example}
\label{example-oka-family-finitely-generated}
Soient $R$ un anneau, $I \subset R$ un idéal et $a \in R$.
Si $(I : a)$ est engendré par $a_1, \ldots, a_n$ et si $(I, a)$
est engendré par $a, b_1, \ldots, b_m$, avec
$b_1, \ldots, b_m \in I$, alors $I$ est engendré par
$aa_1, \ldots, aa_n, b_1, \ldots, b_m$.
Pour le voir, notons que, si $x \in I$, alors $x \in (I, a)$
est une combinaison linéaire de $a, b_1, \ldots, b_m$, mais le
coefficient de $a$ doit appartenir à $(I:a)$.
Nous en déduisons que
la famille des idéaux de type fini est une famille d'Oka.
\end{example}

\begin{example}
\label{example-oka-family-principal}
Montrons que la famille des idéaux principaux d'un anneau $R$ est une famille d'Oka.
En effet, supposons que $I \subset R$ soit un idéal, que $a \in R$ et que $(I, a)$ et
$(I : a)$ soient principaux. Notons que $(I : a) = (I : (I, a))$.
En posant $J = (I, a)$, nous constatons que $J$ est principal et que $(I : J)$ l'est aussi. D'après
le lemme \ref{lemma-colon},
nous avons $I = J (I : J)$.
Ainsi, dans notre situation, puisque $J = (I, a)$ et $(I : J)$
sont principaux, $I$ est principal.
\end{example}

\begin{example}
\label{example-oka-family-bound-cardinality}
Soit $R$ un anneau.
Soit $\kappa$ un cardinal infini.
La famille des idéaux qui peuvent être engendrés par au plus $\kappa$ éléments
est une famille d'Oka. L'argument est analogue à celui de
l'exemple \ref{example-oka-family-finitely-generated}
et est omis.
\end{example}

\begin{example}
\label{example-oka-family-property-modules}
Soient $A$ un anneau, $I \subset A$ un idéal et $a \in A$ un élément.
Il existe une suite exacte courte $0 \to A/(I : a) \to A/I \to A/(I, a) \to 0$
dont la première flèche est donnée par la multiplication par $a$. Ainsi, si $P$
est une propriété des $A$-modules qui est stable par extensions et qui est vérifiée
par $0$, alors la famille des idéaux $I$ tels que $A/I$ possède $P$
est une famille d'Oka.
\end{example}

\begin{proposition}
\label{proposition-oka}
Si $\mathcal{F}$ est une famille d'Oka d'idéaux, alors tout élément maximal du
complémentaire de $\mathcal{F}$ est premier.
\end{proposition}

\begin{proof}
Supposons que $I \not \in \mathcal{F}$ soit maximal parmi les idéaux qui n'appartiennent pas à
$\mathcal{F}$, mais que $I$ ne soit pas premier. Notons que $I \not = R$, car
$R \in \mathcal{F}$. Puisque $I$ n'est pas premier, nous pouvons trouver $a, b \in R - I$
tels que $ab \in I$. Il s'ensuit que $(I, a) \neq I$ et que $(I : a)$ contient
$b \not \in I$, donc que $(I : a) \neq I$ également. Ainsi, $(I : a), (I, a)$ contiennent tous deux
strictement $I$ et doivent donc appartenir à $\mathcal{F}$.
Par la condition d'Oka, nous avons $I \in \mathcal{F}$, ce qui est une contradiction.
\end{proof}

\noindent
Nous pouvons maintenant convertir la plupart des exemples ci-dessus en
un lemme sur les idéaux premiers d'un anneau.

\begin{lemma}
\label{lemma-simple}
Soit $R$ un anneau. Soit $S$ une partie multiplicative de $R$.
Un idéal $I \subset R$ maximal parmi ceux qui vérifient
$I \cap S = \emptyset$ est premier.
\end{lemma}

\begin{proof}
C'est l'exemple présenté dans l'introduction de cette section.
Pour une autre démonstration, combiner
l'exemple \ref{example-oka-family-not-meet-multiplicative-set}
avec la
proposition \ref{proposition-oka}.
\end{proof}

\begin{lemma}
\label{lemma-cohen}
Soit $R$ un anneau.
\begin{enumerate}
\item Un idéal $I \subset R$ maximal parmi ceux qui ne sont pas
de type fini est premier.
\item Si tout idéal premier de $R$ est
de type fini, alors
tout idéal de $R$ est de type fini\footnote{Nous dirons plus loin
que $R$ est noethérien.}.
\end{enumerate}
\end{lemma}

\begin{proof}
La première assertion est une conséquence immédiate de
l'exemple \ref{example-oka-family-finitely-generated} et de la
proposition \ref{proposition-oka}. Pour la seconde,
supposons qu'il existe un idéal $I \subset R$ qui ne soit pas de type
fini. La réunion d'une chaîne totalement ordonnée $\left\{I_\alpha\right\}$
d'idéaux qui ne sont pas de type fini n'est pas de type fini ;
en effet, si $I = \bigcup I_\alpha$ était engendré par
$a_1, \ldots, a_n$, alors tous les générateurs appartiendraient à un certain
$I_\alpha $ et l'engendreraient par conséquent.
D'après le lemme de Zorn, il existe un idéal maximal parmi ceux qui ne sont pas de type
fini. D'après la première partie, cet idéal est premier.
\end{proof}

\begin{lemma}
\label{lemma-primes-principal}
Soit $R$ un anneau.
\begin{enumerate}
\item Un idéal $I \subset R$ maximal parmi ceux qui ne sont pas
principaux est premier.
\item Si tout idéal premier de $R$ est principal, alors
tout idéal de $R$ est principal.
\end{enumerate}
\end{lemma}

\begin{proof}
La première partie résulte de
l'exemple \ref{example-oka-family-principal} et de la
proposition \ref{proposition-oka}.
Pour la seconde, supposons qu'il existe un idéal $I \subset R$
qui ne soit pas principal. La réunion d'une chaîne totalement ordonnée
$\left\{I_\alpha\right\}$ d'idéaux qui ne sont pas principaux n'est pas principale ;
en effet, si $I = \bigcup I_\alpha$ était engendré par
$a$, alors $a$ appartiendrait à un certain $I_\alpha $ et $a$ l'engendrerait.
D'après le lemme de Zorn, il existe un idéal maximal parmi ceux qui ne sont pas
principaux. Cet idéal est nécessairement premier d'après la première partie.
\end{proof}

\begin{lemma}
\label{lemma-characterize-domain}
Soit $R$ un anneau.
\begin{enumerate}
\item Un idéal maximal parmi les idéaux qui ne contiennent aucun
élément non diviseur de zéro est premier.
\item Si $R$ est non nul et si tout idéal premier non nul de $R$
contient un élément non diviseur de zéro, alors $R$ est un anneau intègre.
\end{enumerate}
\end{lemma}

\begin{proof}
Considérons l'ensemble $S$ des éléments non diviseurs de zéro. C'est une partie
multiplicative de $R$. Ainsi, tout idéal maximal parmi ceux qui ne rencontrent pas
$S$ est premier, voir le
lemme \ref{lemma-simple}.
Par conséquent, si tout idéal premier non nul contient un élément non diviseur de zéro, alors
$(0)$ est premier, c'est-à-dire que $R$ est un anneau intègre.
\end{proof}

\begin{remark}
\label{remark-cohen-bound-cardinality}
Soit $R$ un anneau. Soit $\kappa$ un cardinal infini.
En appliquant
l'exemple \ref{example-oka-family-bound-cardinality} et la
proposition \ref{proposition-oka},
nous voyons que tout idéal maximal parmi ceux qui ne peuvent pas être
engendrés par $\kappa$ éléments est premier. Ce résultat n'est pas très
utile, car il existe un anneau $R$ dont tout idéal premier
peut être engendré par $\aleph_0$ éléments, mais dont un certain
idéal ne le peut pas. En effet, soit $k$ un corps, soit $T$ un ensemble dont
la cardinalité est strictement supérieure à $\aleph_0$, et posons
$$
R = k[\{x_n\}_{n \geq 1}, \{z_{t, n}\}_{t \in T, n \geq 0}]/
(x_n^2, z_{t, n}^2, x_n z_{t, n} - z_{t, n - 1})
$$
C'est un anneau local dont l'unique idéal premier est
$\mathfrak m = (x_n)$. Mais l'idéal $(z_{t, n})$ ne peut pas
être engendré par une famille dénombrable d'éléments.
\end{remark}

\begin{example}
\label{example-noetherian-topology-on-spec}
\begin{reference}
Commentaire de Lukas Heger du 12 novembre 2020.
\end{reference}
Soient $R$ un anneau et $X = \Spec(R)$. Puisque les fermés de $X$ correspondent aux
idéaux radicaux de $R$ (lemme \ref{lemma-Zariski-topology}), nous voyons que $X$ est un
espace topologique noethérien si et seulement si les idéaux radicaux vérifient la condition de chaîne ascendante.
Cela équivaut à dire que tout idéal radical est le radical d'un idéal de type
fini (détails omis). Posons
$$
\mathcal{F} = \{I \subset R \mid
\sqrt{I} = \sqrt{(f_1, \ldots, f_n)}\text{ pour un certain }n
\text{ et }f_1, \ldots, f_n \in R\}.
$$
Le lecteur peut montrer que $\mathcal{F}$ est une famille d'Oka en utilisant l'identité
$$
\sqrt{I} = \sqrt{(I, a)(I : a)}
$$
qui vaut pour tout idéal $I \subset R$ et tout élément $a \in R$.
D'autre part, si nous avons une chaîne totalement ordonnée d'idéaux
$\{I_\alpha\}$ dont aucun n'appartient à $\mathcal{F}$, alors la réunion
$I = \bigcup I_\alpha$ ne peut pas non plus appartenir à $\mathcal{F}$. Sinon,
$\sqrt{I} = \sqrt{(f_1, \ldots, f_n)}$, puis $f_i^e \in I$ pour un certain $e$,
puis $f_i^e \in I_\alpha$ pour un certain $\alpha$ indépendant de $i$, et enfin
$\sqrt{I_\alpha} = \sqrt{(f_1, \ldots, f_n)}$, contradiction.
Ainsi, si l'ensemble des idéaux qui n'appartiennent pas à $\mathcal{F}$ est non vide, il
possède des éléments maximaux et,
exactement comme dans le lemme \ref{lemma-cohen}, nous concluons que $X$ est un
espace topologique noethérien si et seulement si tout idéal premier de $R$
est égal à $\sqrt{(f_1, \ldots, f_n)}$ pour certains $f_1, \ldots, f_n \in R$.
Si ce résultat nous est un jour nécessaire, nous l'énoncerons et le démontrerons
ici avec soin.
\end{example}








\section{Images des morphismes d'anneaux de présentation finie}
\label{section-images-finite-presentation}

\noindent
Dans cette section, nous démontrons quelques résultats sur la
topologie des applications $\Spec(S) \to \Spec(R)$
induites par des morphismes d'anneaux $R \to S$, principalement le théorème de Chevalley.
Pour ce faire, nous utiliserons les notions de parties constructibles,
de parties quasi-compactes, de parties rétrocompactes, etc.,
qui sont définies dans Topologie, section \ref{topology-section-constructible}.

\begin{lemma}
\label{lemma-qc-open}
Soit $U \subset \Spec(R)$ un ouvert. Les conditions suivantes
sont équivalentes :
\begin{enumerate}
\item $U$ est rétrocompact dans $\Spec(R)$,
\item $U$ est quasi-compact,
\item $U$ est une réunion finie d'ouverts principaux, et
\item il existe un idéal de type fini $I \subset R$ tel
que $X \setminus V(I) = U$.
\end{enumerate}
\end{lemma}

\begin{proof}
Nous avons (1) $\Rightarrow$ (2), car $\Spec(R)$ est quasi-compact, voir le
lemme \ref{lemma-quasi-compact}. Nous avons (2) $\Rightarrow$ (3), car les
ouverts principaux forment une base de la topologie. Démontrons (3) $\Rightarrow$ (1).
Soit $U = \bigcup_{i = 1\ldots n} D(f_i)$. Pour montrer que $U$ est rétrocompact
dans $\Spec(R)$, il suffit de montrer que $U \cap V$ est quasi-compact pour tout
ouvert quasi-compact $V$ de $\Spec(R)$. Écrivons
$V = \bigcup_{j = 1\ldots m} D(g_j)$, ce qui est possible par (2) $\Rightarrow$
(3). Tout ouvert principal est homéomorphe au spectre d'un anneau et est donc
quasi-compact, voir les lemmes \ref{lemma-standard-open} et
\ref{lemma-quasi-compact}. Ainsi,
$U \cap V =
(\bigcup_{i = 1\ldots n} D(f_i)) \cap (\bigcup_{j = 1\ldots m} D(g_j))
= \bigcup_{i, j} D(f_i g_j)$ est une réunion finie d'ouverts quasi-compacts,
donc est quasi-compact. Pour achever la démonstration, notons
que (4) équivaut à (3) d'après le
lemme \ref{lemma-Zariski-topology}.
\end{proof}

\begin{lemma}
\label{lemma-affine-map-quasi-compact}
Soit $\varphi : R \to S$ un morphisme d'anneaux.
L'application continue induite $f : \Spec(S) \to \Spec(R)$
est quasi-compacte. Pour toute partie constructible $E \subset \Spec(R)$,
l'image réciproque $f^{-1}(E)$ est constructible dans $\Spec(S)$.
\end{lemma}

\begin{proof}
Montrons d'abord que l'image réciproque de tout ouvert quasi-compact
$U \subset \Spec(R)$ est quasi-compacte. D'après le
lemme \ref{lemma-qc-open}, nous pouvons écrire $U$ comme une réunion
finie d'ouverts principaux. Le lemme \ref{lemma-spec-functorial}
montre donc que $f^{-1}(U)$ est une réunion finie d'ouverts principaux.
Ainsi, $f^{-1}(U)$ est quasi-compact d'après le lemme \ref{lemma-qc-open}.
La seconde assertion résulte alors de Topologie, lemme
\ref{topology-lemma-inverse-images-constructibles}.
\end{proof}

\begin{lemma}
\label{lemma-constructible}
Soit $R$ un anneau. Une partie de $\Spec(R)$ est constructible si et seulement
si elle peut s'écrire comme une réunion finie de parties de la forme
$D(f) \cap V(g_1, \ldots, g_m)$, où $f, g_1, \ldots, g_m \in R$.
\end{lemma}

\begin{proof}
D'après le lemme \ref{lemma-qc-open}, la partie $D(f)$ et le complémentaire de
$V(g_1, \ldots, g_m)$ sont des ouverts rétrocompacts. Ainsi,
$D(f) \cap V(g_1, \ldots, g_m)$ est une partie constructible, de même que
toute réunion finie de telles parties. Réciproquement, soit $T \subset \Spec(R)$ une
partie constructible. D'après Topologie, définition \ref{topology-definition-constructible},
nous pouvons supposer que $T = U \cap V^c$, où $U, V \subset \Spec(R)$
sont des ouverts rétrocompacts. D'après le lemme \ref{lemma-qc-open}, nous pouvons écrire
$U = \bigcup_{i = 1, \ldots, n} D(f_i)$ et
$V = \bigcup_{j = 1, \ldots, m} D(g_j)$. Alors
$T = \bigcup_{i = 1, \ldots, n} \big(D(f_i) \cap V(g_1, \ldots, g_m)\big)$.
\end{proof}

\begin{lemma}
\label{lemma-constructible-is-image}
Soit $R$ un anneau et soit $T \subset \Spec(R)$
une partie constructible. Il existe alors un morphisme d'anneaux $R \to S$ de
présentation finie tel que $T$ soit l'image de
$\Spec(S)$ dans $\Spec(R)$.
\end{lemma}

\begin{proof}
Le spectre d'un produit fini d'anneaux
est la réunion disjointe des spectres, voir le
lemme \ref{lemma-spec-product}. Par conséquent, si $T = T_1 \cup T_2$
et si le résultat vaut pour $T_1$ et $T_2$, alors il
vaut pour $T$.
D'après le lemme \ref{lemma-constructible}, nous pouvons supposer
que $T = D(f) \cap V(g_1, \ldots, g_m)$.
Dans ce cas, $T$ est l'image du morphisme
$\Spec((R/(g_1, \ldots, g_m))_f) \to \Spec(R)$, voir les lemmes
\ref{lemma-standard-open} et \ref{lemma-spec-closed}.
\end{proof}

\begin{lemma}
\label{lemma-open-fp}
Soit $R$ un anneau.
Soit $f$ un élément de $R$.
Soit $S = R_f$.
Alors l'image d'une partie constructible de $\Spec(S)$
est constructible dans $\Spec(R)$.
\end{lemma}

\begin{proof}
Nous utilisons à plusieurs reprises le lemme \ref{lemma-qc-open} sans le mentionner.
Soient $U, V$ des ouverts quasi-compacts de $\Spec(S)$.
Nous allons montrer que l'image de $U \cap V^c$ est constructible.
Sous l'identification
$\Spec(S) = D(f)$ du lemme \ref{lemma-standard-open},
les parties $U, V$ correspondent à des ouverts quasi-compacts
$U', V'$ de $\Spec(R)$.
Il suffit donc de montrer que $U' \cap (V')^c$
est constructible dans $\Spec(R)$, ce qui est clair.
\end{proof}

\begin{lemma}
\label{lemma-closed-fp}
Soit $R$ un anneau.
Soit $I$ un idéal de type fini de $R$.
Soit $S = R/I$.
Alors l'image d'une partie constructible de $\Spec(S)$
est constructible dans $\Spec(R)$.
\end{lemma}

\begin{proof}
Si $I = (f_1, \ldots, f_m)$, alors nous voyons que
$V(I)$ est le complémentaire de $\bigcup D(f_i)$,
voir le lemme \ref{lemma-Zariski-topology}.
Il est donc constructible, d'après le lemme \ref{lemma-qc-open}.
Notons le morphisme $R \to S$ par $f \mapsto \overline{f}$.
Nous devons montrer que, si $\overline{U}, \overline{V}$
sont des ouverts rétrocompacts de $\Spec(S)$, alors l'image de
$\overline{U} \cap \overline{V}^c$
dans $\Spec(R)$ est constructible.
D'après le lemme \ref{lemma-qc-open}, nous pouvons écrire
$\overline{U} = \bigcup D(\overline{g_i})$.
En posant $U = \bigcup D({g_i})$, nous voyons que $\overline{U}$
a pour image $U \cap V(I)$, qui est constructible dans
$\Spec(R)$. De même, l'image de $\overline{V}$ est égale à
$V \cap V(I)$ pour un certain ouvert rétrocompact $V$ de $\Spec(R)$.
Ainsi, l'image de $\overline{U} \cap \overline{V}^c$
est égale à $U \cap V(I) \cap V^c$, comme voulu.
\end{proof}

\begin{lemma}
\label{lemma-affineline-open}
Soit $R$ un anneau. Le morphisme $\Spec(R[x]) \to \Spec(R)$
est ouvert, et l'image de tout ouvert principal est un ouvert
quasi-compact.
\end{lemma}

\begin{proof}
Il suffit de montrer que l'image d'un ouvert principal
$D(f)$, $f\in R[x]$, est un ouvert quasi-compact.
L'image de $D(f)$ est celle du morphisme
$\Spec(R[x]_f) \to \Spec(R)$.
Soit $\mathfrak p \subset R$ un idéal premier.
Soit $\overline{f}$ l'image de $f$ dans
$\kappa(\mathfrak p)[x]$.
Rappelons, voir le lemme \ref{lemma-in-image},
que $\mathfrak p$ appartient à l'image
si et seulement si $R[x]_f \otimes_R \kappa(\mathfrak p) =
\kappa(\mathfrak p)[x]_{\overline{f}}$ n'est pas l'anneau
nul. C'est exactement la condition que $f$ ne s'envoie pas
sur zéro dans $\kappa(\mathfrak p)[x]$, autrement dit qu'un
coefficient de $f$ n'appartienne pas à $\mathfrak p$.
Ainsi, si $f = a_d x^d + \ldots + a_0$, alors
l'image de $D(f)$ est $D(a_d) \cup \ldots \cup D(a_0)$.
\end{proof}

\noindent
Nous démontrons une propriété des polynômes caractéristiques qui
sera utilisée ci-dessous.

\begin{lemma}
\label{lemma-characteristic-polynomial-prime}
Soit $R \to A$ un morphisme d'anneaux.
Supposons que $A \cong R^{\oplus n}$ comme $R$-module.
Soit $f \in A$. L'application de multiplication $m_f: A
\to A$ est $R$-linéaire et possède donc
un polynôme caractéristique
$P(T) = T^n + r_{n-1}T^{n-1} + \ldots + r_0 \in R[T]$.
Pour tout idéal premier
$\mathfrak{p} \in \Spec(R)$, $f$ agit de manière nilpotente sur $A
\otimes_R \kappa(\mathfrak{p})$ si et seulement si $\mathfrak p \in
V(r_0, \ldots, r_{n-1})$.
\end{lemma}

\begin{proof}
Cela résulte assez facilement du fait que le polynôme
caractéristique $\bar P(T) \in \kappa(\mathfrak p)[T]$ de l'application
de multiplication $m_{\bar f}: A \otimes_R \kappa(\mathfrak p) \to
A \otimes_R \kappa(\mathfrak p)$, qui multiplie les éléments de $A
\otimes_R \kappa(\mathfrak p)$ par $\bar f$, l'image de $f$ dans cette algèbre sur
$\kappa(\mathfrak p)$, est précisément l'image de $P(T)$ dans
$\kappa(\mathfrak p)[T]$. Soit $(a_{ij})$ la matrice de l'application
$m_f$, à coefficients dans $R$, relativement à une base $e_1, \ldots, e_n$
de $A$ comme $R$-module.
Alors $A \otimes_R \kappa(\mathfrak p) \cong (R \otimes_R
\kappa(\mathfrak p))^{\oplus n} = \kappa(\mathfrak p)^n$, qui est
un espace vectoriel de dimension $n$ sur $\kappa(\mathfrak p)$, de
base $e_1 \otimes 1, \ldots, e_n \otimes 1$. L'image est $\bar f = f
\otimes 1$, et l'application de multiplication $m_{\bar f}$ a donc pour matrice
$(a_{ij} \otimes 1)$. Ainsi, le polynôme caractéristique est
précisément l'image de $P(T)$.

\medskip\noindent
Nous savons par l'algèbre linéaire qu'un endomorphisme agit
de manière nilpotente sur un espace vectoriel de dimension $n$ si et seulement si son
polynôme caractéristique est $T^n$ (car le polynôme caractéristique
divise une puissance du polynôme minimal). Ainsi,
$f$ agit de manière nilpotente sur $A \otimes_R \kappa(\mathfrak p)$ si et
seulement si $\bar P(T) = T^n$. Cela se produit si et seulement si $r_i \in
\mathfrak p$ pour tout $0 \leq i \leq n - 1$, c'est-à-dire lorsque $\mathfrak p \in
V(r_0, \ldots, r_{n - 1}).$
\end{proof}

\begin{lemma}
\label{lemma-affineline-special}
Soit $R$ un anneau. Soient $f, g \in R[x]$ des polynômes.
Supposons que le coefficient dominant de $g$ soit une unité de $R$.
Il existe des éléments $r_i\in R$, $i = 1\ldots, n$, tels que
l'image de $D(f) \cap V(g)$ dans $\Spec(R)$ soit
$\bigcup_{i = 1, \ldots, n} D(r_i)$.
\end{lemma}

\begin{proof}
Écrivons $g = ux^d + a_{d-1}x^{d-1} + \ldots + a_0$, où
$d$ est le degré de $g$, et donc $u \in R^*$.
Considérons l'anneau $A = R[x]/(g)$.
Comme $R$-module, il est libre de type fini, de base les images
de $1, x, \ldots, x^{d-1}$. Considérons la multiplication
par (l'image de) $f$ sur $A$. C'est un morphisme de $R$-modules.
Nous pouvons donc noter $P(T) \in R[T]$ le polynôme caractéristique
de ce morphisme. Écrivons $P(T) = T^d + r_{d-1} T^{d-1} + \ldots + r_0$.
Nous affirmons que $r_0, \ldots, r_{d-1}$ ont la propriété voulue.
Nous utiliserons ci-dessous la propriété des polynômes caractéristiques
selon laquelle
$$
\mathfrak p \in V(r_0, \ldots, r_{d-1})
\Leftrightarrow
\text{la multiplication par }f\text{ est nilpotente sur }
A \otimes_R \kappa(\mathfrak p).
$$
Elle a été démontrée dans le lemme \ref{lemma-characteristic-polynomial-prime}.

\medskip\noindent
Supposons que $\mathfrak q\in D(f) \cap V(g)$, et soit
$\mathfrak p = \mathfrak q \cap R$. Il existe alors un morphisme non nul
$A \otimes_R \kappa(\mathfrak p) \to \kappa(\mathfrak q)$ qui
est compatible avec la multiplication par $f$.
Et $f$ agit comme une unité sur $\kappa(\mathfrak q)$.
Nous en concluons que $\mathfrak p \not \in  V(r_0, \ldots, r_{d-1})$.

\medskip\noindent
D'autre part, supposons que $r_i \not\in \mathfrak p$ pour un certain
idéal premier $\mathfrak p$ de $R$ et un certain $0 \leq i \leq d - 1$.
Alors la multiplication par $f$ n'est pas nilpotente sur l'algèbre
$A \otimes_R \kappa(\mathfrak p)$.
Il existe donc un idéal premier $\overline{\mathfrak q} \subset
A \otimes_R \kappa(\mathfrak p)$ qui ne contient pas l'image de $f$.
L'image réciproque de $\overline{\mathfrak q}$ dans $R[x]$
est un élément de $D(f) \cap V(g)$ qui s'envoie sur $\mathfrak p$.
\end{proof}

\begin{theorem}[Théorème de Chevalley]
\label{theorem-chevalley}
Supposons que $R \to S$ soit de présentation finie.
L'image d'une partie constructible de
$\Spec(S)$ dans $\Spec(R)$ est constructible.
\end{theorem}

\begin{proof}
Écrivons $S = R[x_1, \ldots, x_n]/(f_1, \ldots, f_m)$.
Nous pouvons factoriser $R \to S$ sous la forme $R \to R[x_1] \to R[x_1, x_2]
\to \ldots \to R[x_1, \ldots, x_{n-1}] \to S$. Nous pouvons donc
supposer que $S = R[x]/(f_1, \ldots, f_m)$.
Dans ce cas, nous factorisons le morphisme sous la forme $R \to R[x] \to S$
et, grâce au lemme \ref{lemma-closed-fp}, nous nous ramenons au
cas $S = R[x]$. D'après le lemme \ref{lemma-qc-open}, il suffit
de montrer que, si
$T = (\bigcup_{i = 1\ldots n} D(f_i)) \cap V(g_1, \ldots, g_m)$
pour $f_i , g_j \in R[x]$, alors l'image dans $\Spec(R)$ est
constructible. Puisque les réunions finies de parties constructibles
sont constructibles, il suffit de traiter le cas $n = 1$,
c'est-à-dire celui où $T = D(f) \cap V(g_1, \ldots, g_m)$.

\medskip\noindent
Notons que, si $c \in R$, alors
$$
\Spec(R) =
V(c) \amalg D(c) =
\Spec(R/(c)) \amalg \Spec(R_c),
$$
et, de façon correspondante, $\Spec(R[x]) =
V(c) \amalg D(c) = \Spec(R/(c)[x]) \amalg
\Spec(R_c[x])$. L'intersection de $T = D(f) \cap V(g_1, \ldots, g_m)$
avec chaque partie conserve la même forme, où $f$, $g_i$ sont remplacés
par leurs images dans $R/(c)[x]$ et $R_c[x]$, respectivement.
Notons que l'image de $T$
dans $\Spec(R)$ est la réunion des images de
$T \cap V(c)$ et de $T \cap D(c)$. D'après les lemmes \ref{lemma-open-fp}
et \ref{lemma-closed-fp}, il suffit de montrer que les images des deux
parties sont constructibles dans $\Spec(R/(c))$ et
$\Spec(R_c)$, respectivement.

\medskip\noindent
Supposons que $T = D(f) \cap V(g_1, \ldots, g_m)$
comme ci-dessus, avec $\deg(g_1) \leq \deg(g_2) \leq \ldots \leq \deg(g_m)$.
Nous allons raisonner par récurrence sur $m$ et sur les
degrés des $g_i$. Soit $d_1 = \deg(g_1)$, c'est-à-dire $g_1 = c x^{d_1} + l.o.t$
avec $c \in R$ non nul. En décomposant $R$ en les deux morceaux
$R/(c)$ et $R_c$, soit nous abaissons le degré de $g_1$ (et ce
cas est couvert par récurrence),
soit nous nous ramenons au cas où $c$ est inversible.
Si $c$ est inversible et si $m > 1$, écrivons
$g_2 = c' x^{d_2} + l.o.t$. Dans ce cas, considérons
$g_2' = g_2 - (c'/c) x^{d_2 - d_1} g_1$. Puisque les idéaux
$(g_1, g_2, \ldots, g_m)$ et $(g_1, g_2', g_3, \ldots, g_m)$
sont égaux, nous voyons que $T = D(f) \cap V(g_1, g_2', g_3\ldots, g_m)$.
Mais ici, le degré de $g_2'$ est strictement inférieur à celui
de $g_2$, et ce cas est donc couvert par récurrence.

\medskip\noindent
Les cas de base de la récurrence ci-dessus sont les cas
(a) $T = D(f) \cap V(g)$ où le coefficient dominant
de $g$ est inversible, et (b) $T = D(f)$. Ces deux cas
sont traités par les lemmes \ref{lemma-affineline-special}
et \ref{lemma-affineline-open}.
\end{proof}











\section{Compléments sur les images}
\label{section-more-images}

\noindent
Dans cette section, nous rassemblons quelques lemmes supplémentaires concernant l'image
au niveau de $\Spec$ pour les morphismes d'anneaux. Voir aussi, par exemple, la section
\ref{section-going-up}.

\begin{lemma}
\label{lemma-generic-finite-presentation}
Soit $R \subset S$ une inclusion d'anneaux intègres.
Supposons que $R \to S$ soit de type fini.
Il existe un élément non nul $f \in R$ et un élément non nul $g \in S$
tels que $R_f \to S_{fg}$ soit de présentation finie.
\end{lemma}

\begin{proof}
Nous raisonnons par récurrence sur le nombre de générateurs de $S$ sur $R$.
Au cours de la démonstration, nous pouvons remplacer $R$ par $R_f$ et $S$ par $S_f$
pour un certain élément non nul $f \in R$.

\medskip\noindent
Supposons que $S$ soit engendré par un seul élément
sur $R$. Alors $S = R[x]/\mathfrak q$ pour un certain
idéal premier $\mathfrak q \subset R[x]$. Si $\mathfrak q = (0)$,
il n'y a rien à démontrer. Si $\mathfrak q \not = (0)$,
soit $h \in \mathfrak q$ un élément non nul de degré minimal
en $x$. Écrivons $h = f x^d + a_{d - 1} x^{d - 1} + \ldots + a_0$,
avec $a_i \in R$ et $f \not = 0$. Après avoir inversé $f$
dans $R$ et $S$, nous pouvons supposer que $h$ est unitaire. Nous obtenons
un morphisme surjectif de $R$-algèbres $R[x]/(h) \to S$.
Nous avons $R[x]/(h) = R \oplus Rx \oplus \ldots \oplus Rx^{d - 1}$
comme $R$-module et, par minimalité de $d$, nous voyons que
$R[x]/(h)$ s'injecte dans $S$. Ainsi, $R[x]/(h) \cong S$
est de présentation finie sur $R$.

\medskip\noindent
Supposons que $S$ soit engendré par $n > 1$ éléments sur $R$.
Disons que $x_1, \ldots, x_n \in S$ engendrent $S$. Notons $S' \subset S$
la sous-$R$-algèbre engendrée par $x_1, \ldots, x_{n-1}$. Par l'hypothèse
de récurrence, nous voyons qu'il existe des éléments non nuls $f\in R$ et $g \in S'$
tels que $R_f \to S'_{fg}$ soit de présentation finie.
Nous appliquons ensuite l'hypothèse de récurrence à $S'_{fg} \to S_{fg}$
pour voir qu'il existe $f' \in S'_{fg}$ et
$g' \in S_{fg}$ tels que $S'_{fgf'} \to S_{fgf'g'}$
soit de présentation finie. Nous laissons au lecteur le soin de conclure.
\end{proof}

\begin{lemma}
\label{lemma-characterize-image-finite-type}
Soit $R \to S$ un morphisme d'anneaux de type fini.
Notons $X = \Spec(R)$ et $Y = \Spec(S)$.
Notons $f : Y \to X$ le morphisme induit
sur les spectres. Soit $E \subset Y = \Spec(S)$ une
partie constructible.
Si un point $\xi \in X$ appartient à $f(E)$, alors
$\overline{\{\xi\}} \cap f(E)$ contient un sous-ensemble ouvert
dense de $\overline{\{\xi\}}$.
\end{lemma}

\begin{proof}
Soit $\xi \in X$ un point de $f(E)$. Choisissons un point $\eta \in E$
qui s'envoie sur $\xi$. Soit $\mathfrak p \subset R$ l'idéal premier
correspondant à $\xi$, et soit $\mathfrak q \subset S$ l'idéal premier
correspondant à $\eta$. Considérons le diagramme
$$
\xymatrix{
\eta \ar[r] \ar@{|->}[d] & E \cap Y' \ar[r] \ar[d] &
Y' = \Spec(S/\mathfrak q) \ar[r] \ar[d] &
Y \ar[d] \\
\xi \ar[r] & f(E) \cap X' \ar[r] &
X' = \Spec(R/\mathfrak p) \ar[r] &
X
}
$$
D'après le lemme \ref{lemma-affine-map-quasi-compact}, la partie $E \cap Y'$
est constructible dans $Y'$.
Il s'ensuit que nous pouvons remplacer $X$ par $X'$ et
$Y$ par $Y'$. Nous pouvons donc supposer que $R \subset S$ est une
inclusion d'anneaux intègres, que $\xi$ est le point générique
de $X$ et que $\eta$ est le point générique de $Y$.
En combinant le lemme \ref{lemma-generic-finite-presentation}
avec le théorème de Chevalley
(théorème \ref{theorem-chevalley}),
nous voyons qu'il existe des ouverts denses $U \subset X$ et
$V \subset Y$ tels que $f(V) \subset U$ et
que $f : V \to U$ envoie les parties constructibles
sur des parties constructibles. Notons que $E \cap V$ est
constructible dans $V$, voir Topologie,
lemme \ref{topology-lemma-open-immersion-constructible-inverse-image}.
Ainsi, $f(E \cap V)$ est constructible dans $U$ et contient $\xi$.
D'après Topologie, lemme \ref{topology-lemma-generic-point-in-constructible},
nous voyons que $f(E \cap V)$ contient un ouvert dense $U' \subset U$.
\end{proof}

\noindent
À la fin de cette section, nous présentons quelques résultats supplémentaires sur les
images des applications entre spectres qui n'ont rien à voir avec les parties
constructibles.

\begin{lemma}
\label{lemma-surjective-spec-radical-ideal}
Soit $\varphi : R \to S$ un morphisme d'anneaux.
Les conditions suivantes sont équivalentes :
\begin{enumerate}
\item L'application $\Spec(S) \to \Spec(R)$ est surjective.
\item Pour tout idéal $I \subset R$,
l'image réciproque de $\sqrt{IS}$ dans $R$ est égale à $\sqrt{I}$.
\item Pour tout idéal radical $I \subset R$, l'image réciproque
de $IS$ dans $R$ est égale à $I$.
\item Pour tout idéal premier $\mathfrak p$ de $R$, l'image
réciproque de $\mathfrak p S$ dans $R$ est $\mathfrak p$.
\end{enumerate}
Dans ce cas, il en va de même après tout changement de base : étant donné un morphisme d'anneaux
$R \to R'$, le morphisme d'anneaux $R' \to R' \otimes_R S$ possède lui aussi les
propriétés équivalentes (1), (2), (3).
\end{lemma}

\begin{proof}
Si $J \subset S$ est un idéal, alors
$\sqrt{\varphi^{-1}(J)} = \varphi^{-1}(\sqrt{J})$. Cela montre que (2)
et (3) sont équivalentes.
L'implication (3) $\Rightarrow$ (4) est immédiate.
Si $I \subset R$ est un idéal radical, alors le
lemme \ref{lemma-Zariski-topology}
garantit que $I = \bigcap_{I \subset \mathfrak p} \mathfrak p$.
Ainsi, (4) $\Rightarrow$ (2). D'après le
lemme \ref{lemma-in-image},
nous avons $\mathfrak p = \varphi^{-1}(\mathfrak p S)$ si et seulement si
$\mathfrak p$ appartient à l'image. Ainsi, (1) $\Leftrightarrow$ (4).
Les conditions (1), (2), (3) et (4) sont donc équivalentes.

\medskip\noindent
Supposons que (1) soit vérifiée. Soit $R \to R'$ un morphisme d'anneaux. Soit
$\mathfrak p' \subset R'$ un idéal premier au-dessus de l'idéal premier
$\mathfrak p$ de $R$. Pour voir que $\mathfrak p'$ appartient à l'image
de $\Spec(R' \otimes_R S) \to \Spec(R')$, nous devons montrer
que $(R' \otimes_R S) \otimes_{R'} \kappa(\mathfrak p')$ n'est pas nul, voir le
lemme \ref{lemma-in-image}.
Mais nous avons
$$
(R' \otimes_R S) \otimes_{R'} \kappa(\mathfrak p') =
S \otimes_R \kappa(\mathfrak p)
\otimes_{\kappa(\mathfrak p)} \kappa(\mathfrak p')
$$
qui n'est pas nul, car $S \otimes_R \kappa(\mathfrak p)$ n'est pas nul
par hypothèse et $\kappa(\mathfrak p) \to \kappa(\mathfrak p')$ est
une extension de corps.
\end{proof}

\begin{lemma}
\label{lemma-domain-image-dense-set-points-generic-point}
Soit $R$ un anneau intègre. Soit $\varphi : R \to S$ un morphisme d'anneaux.
Les conditions suivantes sont équivalentes :
\begin{enumerate}
\item Le morphisme d'anneaux $R \to S$ est injectif.
\item L'image de $\Spec(S) \to \Spec(R)$
contient un ensemble dense de points.
\item Il existe un idéal premier $\mathfrak q \subset S$
dont l'image réciproque dans $R$ est $(0)$.
\end{enumerate}
\end{lemma}

\begin{proof}
Soit $K$ le corps des fractions de l'anneau intègre $R$.
Supposons que $R \to S$ soit injectif. Puisque la localisation
est exacte, nous voyons que $K \to S \otimes_R K$ est injectif.
Il existe donc un idéal premier s'envoyant sur $(0)$ d'après le
lemme \ref{lemma-in-image}.

\medskip\noindent
Notons que $(0)$ est dense dans $\Spec(R)$, de sorte que la
dernière condition implique la seconde.

\medskip\noindent
Supposons que la seconde condition soit vérifiée. Soit $f \in R$,
$f \not = 0$. Puisque $R$ est un anneau intègre, nous voyons que $V(f)$
est un fermé propre du spectre de $R$. Par hypothèse,
il existe un idéal premier $\mathfrak q$
de $S$ tel que $\varphi(f) \not \in \mathfrak q$.
Ainsi, $\varphi(f) \not = 0$.
Le morphisme $R \to S$ est donc injectif.
\end{proof}

\begin{lemma}
\label{lemma-injective-minimal-primes-in-image}
Soit $R \subset S$ un morphisme d'anneaux injectif.
Alors $\Spec(S) \to \Spec(R)$
atteint tous les idéaux premiers minimaux.
\end{lemma}

\begin{proof}
Soit $\mathfrak p \subset R$ un idéal premier minimal.
Dans ce cas, $R_{\mathfrak p}$ possède un unique idéal premier.
Il suffit donc de montrer que $S_{\mathfrak p}$ n'est pas nul.
Cela résulte du fait que la localisation est exacte,
voir la proposition \ref{proposition-localization-exact}.
\end{proof}

\begin{lemma}
\label{lemma-image-dense-generic-points}
Soit $R \to S$ un morphisme d'anneaux. Les conditions suivantes sont équivalentes :
\begin{enumerate}
\item Le noyau de $R \to S$ est constitué d'éléments nilpotents.
\item Les idéaux premiers minimaux de $R$ appartiennent à l'image de
$\Spec(S) \to \Spec(R)$.
\item L'image de $\Spec(S) \to \Spec(R)$ est dense
dans $\Spec(R)$.
\end{enumerate}
\end{lemma}

\begin{proof}
Soit $I = \Ker(R \to S)$. Notons que
$\sqrt{(0)} = \bigcap_{\mathfrak q \subset S} \mathfrak q$, voir le
lemme \ref{lemma-Zariski-topology}.
Ainsi, $\sqrt{I} = \bigcap_{\mathfrak q \subset S} R \cap \mathfrak q$.
Donc $V(I) = V(\sqrt{I})$ est l'adhérence de l'image de
$\Spec(S) \to \Spec(R)$.
Cela montre que (1) équivaut à (3). Il est clair que
(2) implique (3). Enfin, supposons (1). Nous pouvons remplacer
$R$ par $R/I$ et $S$ par $S/IS$ sans changer la topologie
des spectres ni l'application. Ainsi, l'implication (1) $\Rightarrow$ (2)
résulte du lemme \ref{lemma-injective-minimal-primes-in-image}.
\end{proof}

\begin{lemma}
\label{lemma-minimal-prime-image-minimal-prime}
Soit $R \to S$ un morphisme d'anneaux. Si un idéal premier minimal $\mathfrak p \subset R$
appartient à l'image de $\Spec(S) \to \Spec(R)$, alors il est l'image
d'un idéal premier minimal.
\end{lemma}

\begin{proof}
Écrivons $\mathfrak p = \mathfrak q \cap R$. Choisissons alors un
idéal premier minimal $\mathfrak r \subset S$ tel que $\mathfrak r \subset \mathfrak q$, voir le
lemme \ref{lemma-Zariski-topology}.
Par minimalité de $\mathfrak p$, nous voyons que
$\mathfrak p = \mathfrak r \cap R$.
\end{proof}

\begin{lemma}
\label{lemma-intersection-not-zero}
Soit $A \subset B$ une inclusion d'anneaux intègres induisant une extension algébrique
des corps des fractions. Si $J \subset B$ est un idéal non nul, alors $A \cap J$
est également non nul. Ainsi, l'image d'un fermé propre de $\Spec(B)$
n'est pas dense dans $\Spec(A)$.
\end{lemma}

\begin{proof}
Soit $x \in J$ un élément non nul. Puisque $x$ est algébrique sur le corps des
fractions de $A$, il existe un $d \geq 1$ et des $a_0, \ldots, a_d \in A$
tels que $a_0, a_d \not = 0$ et que
$a_d x^d + a_{d - 1} x^{d - 1} + \ldots + a_0 = 0$ dans $B$.
Alors $a_0 \in A \cap J$.
\end{proof}




\section{Anneaux noethériens}
\label{section-Noetherian}

\noindent
Un anneau $R$ est {\it noethérien} si tout idéal de $R$ est de type fini.
Cela équivaut clairement à la condition de chaîne ascendante pour les idéaux de $R$.
D'après le
lemme \ref{lemma-cohen},
il suffit de vérifier que tout idéal premier de $R$ est de type fini.

\begin{lemma}
\label{lemma-Noetherian-permanence}
\begin{slogan}
La propriété noethérienne est stable par extension de type fini
et par localisation.
\end{slogan}
Tout anneau de type fini sur un anneau noethérien
est noethérien. Toute localisation d'un anneau noethérien
est noethérienne.
\end{lemma}

\begin{proof}
L'assertion sur les localisations résulte du fait
que tout idéal $J \subset S^{-1}R$ est de la forme
$I \cdot S^{-1}R$. Tout quotient $R/I$ d'un anneau noethérien
$R$ est noethérien, car tout idéal $\overline{J} \subset R/I$
est de la forme $J/I$ pour un certain idéal $I \subset J \subset R$.
Il suffit donc de montrer que, si $R$ est noethérien, alors
$R[X]$ l'est aussi. Supposons que $J_1 \subset J_2 \subset \ldots$ soit une
chaîne ascendante d'idéaux de $R[X]$. Considérons les idéaux $I_{i, d}$
définis comme les idéaux constitués des éléments de $R$ qui apparaissent comme coefficients
dominants de polynômes de degré $d$ appartenant à $J_i$.
Il est clair que $I_{i, d} \subset I_{i', d'}$ dès que
$i \leq i'$ et $d \leq d'$. Par la condition de chaîne ascendante
dans $R$, il n'y a qu'un nombre fini d'idéaux distincts parmi tous les
$I_{i, d}$.
(Indication : tout ensemble infini d'éléments de
$\mathbf{N} \times \mathbf{N}$ contient une suite infinie
croissante.)
Prenons $i_0$ assez grand pour que $I_{i, d} = I_{i_0, d}$
pour tous $i \geq i_0$ et tout $d$. Supposons que $f \in J_i$ pour un certain $i \geq i_0$.
Par récurrence sur le degré $d = \deg(f)$, montrons que $f \in J_{i_0}$.
En effet, il existe un $g\in J_{i_0}$ de degré $d$ ayant
le même coefficient dominant que $f$. Par récurrence,
$f - g \in J_{i_0}$, et le résultat s'ensuit.
\end{proof}

\begin{lemma}
\label{lemma-Noetherian-power-series}
Si $R$ est un anneau noethérien, alors l'anneau de séries
formelles $R[[x_1, \ldots, x_n]]$ l'est également.
\end{lemma}

\begin{proof}
Puisque $R[[x_1, \ldots, x_{n + 1}]] \cong R[[x_1, \ldots, x_n]][[x_{n + 1}]]$,
il suffit de montrer que $R[[x]]$ est noethérien lorsque
$R$ l'est. Soit $I \subset R[[x]]$ un idéal.
Nous devons montrer que $I$ est un idéal de type fini.
Pour tout entier naturel
$d$, notons $I_d = \{a \in R \mid ax^d + \text{termes d'ordre supérieur} \in I\}$.
Alors nous voyons que $I_0 \subset I_1 \subset \ldots$ se stabilise, puisque $R$
est noethérien. Choisissons $d_0$ tel que $I_{d_0} = I_{d_0 + 1} = \ldots$.
Pour tout $d \leq d_0$, choisissons des éléments $f_{d, j} \in I \cap (x^d)$,
$j = 1, \ldots, n_d$, tels que, si nous écrivons
$f_{d, j} = a_{d, j}x^d + \text{termes d'ordre supérieur}$, alors $I_d = (a_{d, j})$.
Notons $I' = (\{f_{d, j}\}_{d = 0, \ldots, d_0, j = 1, \ldots, n_d})$.
Il est alors clair que $I' \subset I$. Choisissons $f \in I$.
Nous pouvons d'abord choisir des $c_{d, i} \in R$ tels que
$$
f - \sum c_{d, i} f_{d, i} \in (x^{d_0 + 1}) \cap I.
$$
Ensuite, nous pouvons choisir $c_{i, 1} \in R$, $i = 1, \ldots, n_{d_0}$, tels que
$$
f - \sum c_{d, i} f_{d, i} - \sum c_{i, 1}xf_{d_0, i} \in (x^{d_0 + 2}) \cap I.
$$
Ensuite, nous pouvons choisir $c_{i, 2} \in R$, $i = 1, \ldots, n_{d_0}$, tels que
$$
f - \sum c_{d, i} f_{d, i} - \sum c_{i, 1}xf_{d_0, i}
- \sum c_{i, 2}x^2f_{d_0, i}
\in (x^{d_0 + 3}) \cap I.
$$
Et ainsi de suite. Finalement, nous voyons que
$$
f = \sum c_{d, i} f_{d, i} +
\sum\nolimits_i (\sum\nolimits_e c_{i, e} x^e)f_{d_0, i}
$$
appartient à $I'$, comme voulu.
\end{proof}

\noindent
Le lemme suivant, quoique facile, est utile, car les
$\mathbf{Z}$-algèbres de type fini apparaissent assez souvent dans
une technique appelée « réduction noethérienne absolue ».

\begin{lemma}
\label{lemma-obvious-Noetherian}
Toute algèbre de type fini sur un corps est noethérienne.
Toute algèbre de type fini sur $\mathbf{Z}$ est noethérienne.
\end{lemma}

\begin{proof}
Cela résulte immédiatement du lemme \ref{lemma-Noetherian-permanence}
et du fait que les corps sont des anneaux noethériens et que
$\mathbf{Z}$ est un anneau noethérien (car c'est un
anneau principal).
\end{proof}

\begin{lemma}
\label{lemma-Noetherian-finite-type-is-finite-presentation}
Soit $R$ un anneau noethérien.
\begin{enumerate}
\item Tout $R$-module de type fini est de présentation finie.
\item Tout sous-module d'un $R$-module de type fini est de type fini.
\item Toute $R$-algèbre de type fini est de présentation finie sur $R$.
\end{enumerate}
\end{lemma}

\begin{proof}
Soit $M$ un $R$-module de type fini. D'après le
lemme \ref{lemma-trivial-filter-finite-module},
nous pouvons trouver une filtration finie de $M$ dont les quotients successifs sont
de la forme $R/I$. Puisque tout idéal est de type fini, chacun des
quotients $R/I$ est de présentation finie. Ainsi, $M$ est de présentation
finie d'après le
lemme \ref{lemma-extension}.
Cela démontre (1).

\medskip\noindent
Soit $N \subset M$ un sous-module. Puisque $M$ est de type fini, le quotient
$M/N$ est de type fini. Ainsi, $M/N$ est de présentation finie d'après la partie (1).
Nous voyons donc que $N$ est de type fini d'après le lemme \ref{lemma-extension}, partie (5).
Cela démontre la partie (2).

\medskip\noindent
Pour voir (3), notons que tout idéal de
$R[x_1, \ldots, x_n]$ est de type fini d'après le
lemme \ref{lemma-Noetherian-permanence}.
\end{proof}

\begin{lemma}
\label{lemma-Noetherian-topology}
Si $R$ est un anneau noethérien, alors $\Spec(R)$
est un espace topologique noethérien, voir Topologie,
définition \ref{topology-definition-noetherian}.
\end{lemma}

\begin{proof}
Cela vient du fait que tout fermé de $\Spec(R)$
s'écrit de manière unique sous la forme $V(I)$, où $I$ est un idéal radical,
voir le lemme \ref{lemma-Zariski-topology}.
Et cette correspondance renverse les inclusions.
Le résultat découle donc des définitions.
\end{proof}

\begin{lemma}
\label{lemma-Noetherian-irreducible-components}
\begin{slogan}
Un schéma affine noethérien possède un nombre fini de points génériques.
\end{slogan}
Si $R$ est un anneau noethérien, alors $\Spec(R)$
possède un nombre fini de composantes irréductibles. Autrement dit,
$R$ possède un nombre fini d'idéaux premiers minimaux.
\end{lemma}

\begin{proof}
D'après le lemme \ref{lemma-Noetherian-topology} et
Topologie, lemme \ref{topology-lemma-Noetherian},
nous voyons qu'il n'y a qu'un nombre fini de composantes irréductibles.
D'après le lemme \ref{lemma-irreducible}, celles-ci correspondent aux
idéaux premiers minimaux de $R$.
\end{proof}

\begin{lemma}
\label{lemma-Noetherian-base-change-finite-type}
Soit $R \to S$ un morphisme d'anneaux. Supposons que $R \to R'$ soit de type fini.
Si $S$ est noethérien, alors le changement de base $S' = R' \otimes_R S$
est noethérien.
\end{lemma}

\begin{proof}
D'après le lemme \ref{lemma-base-change-finiteness}, le caractère de type fini est stable par
changement de base. Ainsi, $S \to S'$ est de type fini. Puisque $S$ est noethérien, nous
pouvons appliquer le lemme \ref{lemma-Noetherian-permanence}.
\end{proof}

\begin{lemma}
\label{lemma-Noetherian-field-extension}
Soit $k$ un corps et soit $R$ une $k$-algèbre noethérienne.
Si $K/k$ est une extension de corps de type fini, alors
$K \otimes_k R$ est noethérien.
\end{lemma}

\begin{proof}
Puisque $K/k$ est une extension de corps de type fini, il existe
une $k$-algèbre de type fini $B \subset K$ telle que $K$ soit
le corps des fractions de $B$. Autrement dit, $K = S^{-1}B$,
où $S = B \setminus \{0\}$. Alors $K \otimes_k R = S^{-1}(B \otimes_k R)$.
Puis, $B \otimes_k R$ est noethérien d'après le
lemme \ref{lemma-Noetherian-base-change-finite-type}.
Enfin, $K \otimes_k R = S^{-1}(B \otimes_k R)$ est noethérien d'après le
lemme \ref{lemma-Noetherian-permanence}.
\end{proof}

\noindent
Voici quelques lemmes amusants qui sont parfois utiles.

\begin{lemma}
\label{lemma-subring-of-local-ring}
Soient $R$ un anneau et $\mathfrak p \subset R$ un idéal premier.
Il existe un $f \in R$, $f \not \in \mathfrak p$, tel
que $R_f \to R_\mathfrak p$ soit injectif dans chacun des
cas suivants :
\begin{enumerate}
\item $R$ est un anneau intègre,
\item $R$ est noethérien, ou
\item $R$ est réduit et possède un nombre fini d'idéaux premiers minimaux.
\end{enumerate}
\end{lemma}

\begin{proof}
Si $R$ est un anneau intègre, alors $R \subset R_\mathfrak p$, et donc $f = 1$ convient.
Si $R$ est noethérien, alors le noyau $I$ de $R \to R_\mathfrak p$
est un idéal de type fini, et nous pouvons trouver
$f \in R$, $f \not \in \mathfrak p$, tel que $IR_f = 0$.
Pour ce $f$, le morphisme $R_f \to R_\mathfrak p$ est injectif,
et $f$ convient. Si $R$ est réduit et possède un nombre
fini d'idéaux premiers minimaux $\mathfrak p_1, \ldots, \mathfrak p_n$,
alors nous pouvons choisir
$f \in \bigcap_{\mathfrak p_i \not \subset \mathfrak p} \mathfrak p_i$,
$f \not \in \mathfrak p$. En effet, si $\mathfrak{p}_i\not\subset
\mathfrak{p}$, alors il existe $f_i \in \mathfrak{p}_i$,
$f_i \not\in \mathfrak{p}$, et $f = \prod f_i$ convient.
Pour ce $f$, nous avons $R_f \subset R_\mathfrak p$, car les idéaux premiers
minimaux de $R_f$ correspondent aux idéaux premiers minimaux de $R_\mathfrak p$
et nous pouvons appliquer le lemme \ref{lemma-reduced-ring-sub-product-fields}
(certains détails sont omis).
\end{proof}

\begin{lemma}
\label{lemma-surjective-endo-noetherian-ring-is-iso}
Tout endomorphisme surjectif d'un anneau noethérien est un isomorphisme.
\end{lemma}

\begin{proof}
Si $f : R \to R$ était un tel endomorphisme sans être injectif, alors
$$
\Ker(f) \subset \Ker(f \circ f) \subset
\Ker(f \circ f \circ f) \subset \ldots
$$
serait une chaîne strictement croissante d'idéaux.
\end{proof}







\section{Idéaux localement nilpotents}
\label{section-locally-nilpotent}

\noindent
Voici la définition.

\begin{definition}
\label{definition-locally-nilpotent-ideal}
Soit $R$ un anneau. Soit $I \subset R$ un idéal.
Nous disons que $I$ est {\it localement nilpotent} si, pour tout
$x \in I$, il existe un $n \in \mathbf{N}$ tel
que $x^n = 0$. Nous disons que $I$ est {\it nilpotent} s'il
existe un $n \in \mathbf{N}$ tel que $I^n = 0$.
\end{definition}

\begin{example}
\label{example-locally-nilpotent-not-nilpotent}
Soit $R = k[x_n | n \in \mathbf{N}]$ l'anneau de polynômes en une infinité
de variables sur un corps $k$. Soit $I$ l'idéal engendré par
les éléments $x_n^n$ pour $n \in \mathbf{N}$, et soit $S = R/I$. Alors l'idéal
$J \subset S$ engendré par les images des $x_n$, $n \in \mathbf{N}$,
est localement nilpotent, mais n'est pas nilpotent. En effet, puisque les combinaisons
$S$-linéaires d'éléments nilpotents sont nilpotentes, pour prouver que $J$ est localement
nilpotent, il suffit d'observer que tous ses générateurs sont nilpotents
(ce qu'ils sont manifestement). D'autre part, pour tout $n \in \mathbf{N}$,
nous avons $x_{n + 1}^n \not \in I$, de sorte que $J^n \not = 0$.
Il s'ensuit que $J$ n'est pas nilpotent.
\end{example}

\begin{lemma}
\label{lemma-locally-nilpotent}
Soit $R \to R'$ un morphisme d'anneaux et soit $I \subset R$ un idéal localement
nilpotent. Alors $IR'$ est un idéal localement nilpotent de $R'$.
\end{lemma}

\begin{proof}
Cela résulte du fait que, si $x, y \in R'$ sont nilpotents, alors
$x + y$ l'est aussi. En effet, si $x^n = 0$ et $y^m = 0$, alors
$(x + y)^{n + m - 1} = 0$.
\end{proof}

\begin{lemma}
\label{lemma-locally-nilpotent-unit}
Soit $R$ un anneau et soit $I \subset R$ un idéal localement
nilpotent.
Un élément $x$ de $R$ est une unité si et seulement si l'image de $x$
dans $R/I$ est une unité.
\end{lemma}

\begin{proof}
Si $x$ est une unité de $R$, alors son image est clairement une unité de $R/I$.
Il reste à démontrer la réciproque.
Supposons que l'image de $y \in R$ dans $R/I$ soit l'inverse de l'image de $x$.
Alors $xy = 1 - z$ pour un certain $z \in I$.
Cela signifie que $1\equiv z$ modulo $xR$.
Puisque $z$ appartient à l'idéal localement nilpotent
$I$, nous avons $z^N = 0$ pour un certain $N$ assez grand.
Il s'ensuit que $1 = 1^N \equiv z^N = 0$ modulo $xR$.
Autrement dit, $x$ divise $1$ et est donc une unité.
\end{proof}

\begin{lemma}
\label{lemma-Noetherian-power}
\begin{slogan}
Un idéal d'un anneau noethérien est nilpotent si chacun de ses éléments
est nilpotent.
\end{slogan}
Soit $R$ un anneau noethérien. Soient $I, J$ des idéaux de $R$.
Supposons que $J \subset \sqrt{I}$. Alors $J^n \subset I$ pour un certain $n$.
En particulier, dans un anneau noethérien, les notions d'
« idéal localement nilpotent »
et d'« idéal nilpotent » coïncident.
\end{lemma}

\begin{proof}
Écrivons $J = (f_1, \ldots, f_s)$.
Par hypothèse, $f_i^{d_i} \in I$.
Prenons $n = d_1 + d_2 + \ldots + d_s + 1$.
\end{proof}

\begin{lemma}
\label{lemma-lift-idempotents}
Soit $R$ un anneau. Soit $I \subset R$ un idéal localement nilpotent.
Alors $R \to R/I$ induit une bijection sur les idempotents.
\end{lemma}

\begin{proof}[Première démonstration du lemme \ref{lemma-lift-idempotents}]
Puisque $I$ est localement nilpotent, il est contenu dans tout idéal premier.
Ainsi, $\Spec(R/I) = V(I) = \Spec(R)$. Le
lemme résulte donc du lemme \ref{lemma-disjoint-decomposition}.
\end{proof}

\begin{proof}[Seconde démonstration du lemme \ref{lemma-lift-idempotents}]
Supposons que $\overline{e} \in R/I$ soit un idempotent.
Nous devons relever $\overline{e}$ en un idempotent de $R$.

\medskip\noindent
Choisissons d'abord un relèvement quelconque $f \in R$ de $\overline{e}$ et posons
$x = f^2 - f$. Alors $x \in I$, donc $x$ est nilpotent (puisque $I$
est localement nilpotent). Soit maintenant $J$ l'idéal de $R$ engendré
par $x$. Alors $J$ est nilpotent (et pas seulement localement nilpotent),
puisqu'il est engendré par l'élément nilpotent $x$.

\medskip\noindent
Supposons maintenant que nous ayons trouvé un relèvement $e \in R$ de $\overline{e}$
tel que $e^2 - e \in J^k$ pour un certain $k \geq 1$.
Soit $e' = e - (2e - 1)(e^2 - e) = 3e^2 - 2e^3$, qui est un autre
relèvement de $\overline{e}$ (car l'idempotence de $\overline{e}$
donne $e^2 - e \in I$). Alors
$$
(e')^2 - e' = (4e^2 - 4e - 3)(e^2 - e)^2 \in J^{2k}
$$
par un calcul direct.

\medskip\noindent
Nous sommes donc partis d'un relèvement $e$ de $\overline{e}$ tel
que $e^2 - e \in J^k$, et nous avons obtenu un relèvement $e'$ de
$\overline{e}$ tel que $(e')^2 - e' \in J^{2k}$.
Nous pouvons ainsi améliorer successivement l'approximation
(en commençant par $e = f$, qui convient pour $k = 1$).
Nous finissons par atteindre un stade où $J^k = 0$, et à ce
stade nous disposons d'un relèvement $e$ de $\overline{e}$ tel que
$e^2 - e \in J^k = 0$, c'est-à-dire que cet $e$ est idempotent.

\medskip\noindent
Nous avons ainsi vu que, si $\overline{e} \in R/I$ est un
idempotent quelconque, il existe un relèvement de $\overline{e}$
qui est un idempotent de $R$.
Il reste à montrer que ce relèvement est unique. En effet, soient
$e_1$ et $e_2$ deux tels relèvements. Nous
devons montrer que $e_1 = e_2$.

\medskip\noindent
Par définition de $e_1$ et de $e_2$, nous avons $e_1 \equiv e_2
\mod I$, et $e_1$ et $e_2$ sont tous deux idempotents. De
$e_1 \equiv e_2 \mod I$, nous déduisons que $e_1 - e_2 \in I$,
de sorte que $e_1 - e_2$ est nilpotent (puisque $I$ est localement nilpotent).
Un calcul direct
(utilisant l'idempotence de $e_1$ et de $e_2$)
montre que $(e_1 - e_2)^3 = e_1 - e_2$. En utilisant cela et
une récurrence, nous obtenons $(e_1 - e_2)^k = e_1 - e_2$ pour tout
entier impair positif $k$. Comme tout $k$ assez grand vérifie
$(e_1 - e_2)^k = 0$ (puisque $e_1 - e_2$ est nilpotent),
cela montre que $e_1 - e_2 = 0$, et donc que $e_1 = e_2$, ce qui
achève la démonstration.
\end{proof}

\begin{lemma}
\label{lemma-lift-idempotents-noncommutative}
Soit $A$ une algèbre éventuellement non commutative.
Soit $e \in A$ un élément tel que $x = e^2 - e$ soit nilpotent.
Il existe alors un idempotent de la forme
$e' = e + x(\sum a_{i, j}e^ix^j) \in A$,
avec $a_{i, j} \in \mathbf{Z}$.
\end{lemma}

\begin{proof}
Considérons l'anneau $R_n = \mathbf{Z}[e]/((e^2 - e)^n)$. Il est clair que,
si nous pouvons démontrer le résultat pour chaque $R_n$, le lemme en résulte.
Dans $R_n$, considérons l'idéal $I = (e^2 - e)$ et appliquons le
lemme \ref{lemma-lift-idempotents}.
\end{proof}

\begin{lemma}
\label{lemma-lift-nth-roots}
Soit $R$ un anneau. Soit $I \subset R$ un idéal localement nilpotent.
Soit $n \geq 1$ un entier inversible dans $R/I$. Alors
\begin{enumerate}
\item l'application d'élévation à la puissance $n$-ième $1 + I \to 1 + I$, $1 + x \mapsto (1 + x)^n$
est une bijection,
\item une unité de $R$ est une puissance $n$-ième si et seulement si son image dans $R/I$
est une puissance $n$-ième.
\end{enumerate}
\end{lemma}

\begin{proof}
Soit $a \in R$ une unité dont l'image dans $R/I$ est égale à l'image
de $b^n$, avec $b \in R$. Alors $b$ est une unité
(lemme \ref{lemma-locally-nilpotent-unit}) et
$ab^{-n} = 1 + x$ pour un certain $x \in I$. Ainsi, $ab^{-n} = c^n$ d'après
la partie (1). La partie (2) résulte donc de la partie (1).

\medskip\noindent
Démontrons (1). Cela est vrai parce que l'application
$1 + x \mapsto (1 + x)^n$ possède une inverse. Nous pouvons en effet considérer l'application
qui envoie $1 + x$ sur
\begin{align*}
(1 + x)^{1/n}
& =
1 + {1/n \choose 1}x +
{1/n \choose 2}x^2 +
{1/n \choose 3}x^3 + \ldots \\
& =
1 + \frac{1}{n} x + \frac{1 - n}{2n^2}x^2 +
\frac{(1 - n)(1 - 2n)}{6n^3}x^3 + \ldots
\end{align*}
comme en calcul élémentaire. Cette expression a un sens, car la série est finie,
puisque $x^k = 0$ pour tout $k \gg 0$, et que chaque coefficient
$ {1/n \choose k} \in \mathbf{Z}[1/n]$ (détails omis ; observons que
$n$ est inversible dans $R$ d'après le lemme \ref{lemma-locally-nilpotent-unit}).
\end{proof}






\section{Curiosité}
\label{section-curiosity}

\noindent
Le lemme \ref{lemma-disjoint-implies-product}
explique ce qui se passe lorsque $V(I)$ est ouvert pour un certain idéal $I \subset R$.
Mais qu'en est-il si $\Spec(S^{-1}R)$ est fermé dans $\Spec(R)$ ?
Les deux lemmes suivants donnent une réponse partielle. Pour plus d'informations, voir la
section \ref{section-pure-ideals}.

\begin{lemma}
\label{lemma-invert-closed-quotient}
Soit $R$ un anneau. Soit $S \subset R$ une partie multiplicative.
Supposons que l'image du morphisme $\Spec(S^{-1}R) \to \Spec(R)$
soit fermée. Alors $S^{-1}R \cong R/I$ pour un certain idéal $I \subset R$.
\end{lemma}

\begin{proof}
Soit $I = \Ker(R \to S^{-1}R)$, de sorte que $V(I)$ contient l'image.
Disons que l'image est le fermé $V(I') \subset \Spec(R)$ pour
un certain idéal $I' \subset R$. Ainsi, $V(I') \subset V(I)$.
Pour $f \in I'$, nous voyons que $f/1 \in S^{-1}R$
appartient à tout idéal premier. Ainsi, $f^n$ s'envoie sur zéro dans $S^{-1}R$
pour un certain $n \geq 1$ (lemme \ref{lemma-Zariski-topology}).
Par conséquent, $V(I') = V(I)$.
Cela implique que tout $g \in S$ est inversible modulo $I$.
Nous obtenons donc des morphismes d'anneaux $R/I \to S^{-1}R$ et $S^{-1}R \to R/I$.
Le premier morphisme est injectif par le choix de $I$.
Le second est le morphisme $S^{-1}R \to S^{-1}(R/I) = R/I$, qui
a pour noyau $S^{-1}I$, car la localisation est exacte.
Puisque $S^{-1}I = 0$, nous voyons que le second morphisme est lui aussi injectif.
Ainsi, $S^{-1}R \cong R/I$.
\end{proof}

\begin{lemma}
\label{lemma-invert-closed-split}
Soit $R$ un anneau. Soit $S \subset R$ une partie multiplicative.
Supposons que l'image du morphisme $\Spec(S^{-1}R) \to \Spec(R)$
soit fermée. Si $R$ est noethérien, ou si $\Spec(R)$ est un
espace topologique noethérien, ou si $S$ est de type fini comme monoïde,
alors $R \cong S^{-1}R \times R'$ pour un certain anneau $R'$.
\end{lemma}

\begin{proof}
D'après le lemme \ref{lemma-invert-closed-quotient}, nous avons $S^{-1}R \cong R/I$
pour un certain idéal $I \subset R$. D'après le lemme \ref{lemma-disjoint-implies-product},
il suffit de montrer que $V(I)$ est ouvert.
Si $R$ est noethérien, alors $\Spec(R)$ est un espace topologique
noethérien, voir le lemme \ref{lemma-Noetherian-topology}.
Si $\Spec(R)$ est un espace topologique noethérien,
alors le complémentaire $\Spec(R) \setminus V(I)$ est quasi-compact, voir
Topologie, lemme \ref{topology-lemma-Noetherian-quasi-compact}.
Il existe donc un nombre fini d'éléments $f_1, \ldots, f_n \in I$ tels
que $V(I) = V(f_1, \ldots, f_n)$.
Puisque chaque $f_i$ s'envoie sur zéro dans $S^{-1}R$,
il existe un $g \in S$ tel que $gf_i = 0$ pour
$i = 1, \ldots, n$. Ainsi, $D(g) = V(I)$, comme voulu.
Dans le cas où $S$ est de type fini comme monoïde, disons que $S$ est engendré
par $g_1, \ldots, g_m$, nous avons $S^{-1}R \cong R_{g_1 \ldots g_m}$,
et nous concluons que $V(I) = D(g_1 \ldots g_m)$.
\end{proof}














\section{Nullstellensatz de Hilbert}
\label{section-nullstellensatz}

\begin{theorem}[Nullstellensatz de Hilbert]
\label{theorem-nullstellensatz}
Soit $k$ un corps.
\begin{enumerate}
\item
\label{item-finite-kappa}
Pour tout idéal maximal $\mathfrak m \subset k[x_1, \ldots, x_n]$,
l'extension de corps $\kappa(\mathfrak m)/k$ est finie.
\item
\label{item-polynomial-ring-Jacobson}
Tout idéal radical $I \subset k[x_1, \ldots, x_n]$
est l'intersection des idéaux maximaux qui le contiennent.
\end{enumerate}
Il en va de même dans toute $k$-algèbre de type fini.
\end{theorem}

\begin{proof}
Il suffit de démontrer la partie (\ref{item-finite-kappa}) du
théorème dans le cas d'une algèbre polynomiale
$k[x_1, \ldots, x_n]$, car toute $k$-algèbre de type fini
est un quotient d'une telle algèbre polynomiale.
Nous le démontrons par récurrence sur $n$. Le cas $n = 0$ est clair.
Supposons que $\mathfrak m$ soit un idéal maximal de $k[x_1, \ldots, x_n]$.
Soit $\mathfrak p \subset k[x_n]$ l'intersection
de $\mathfrak m$ avec $k[x_n]$.

\medskip\noindent
Si $\mathfrak p \not = (0)$,
alors $\mathfrak p$ est maximal et engendré par un polynôme
irréductible unitaire $P$ (grâce à l'algorithme d'Euclide
dans $k[x_n]$). Alors
$k' = k[x_n]/\mathfrak p$ est une extension finie de $k$
et est contenu dans $\kappa(\mathfrak m)$. Dans ce cas,
nous obtenons une surjection
$$
k'[x_1, \ldots, x_{n-1}]
\to
k'[x_1, \ldots, x_n] =
k' \otimes_k k[x_1, \ldots, x_n]
\longrightarrow
\kappa(\mathfrak m)
$$
et nous voyons donc que $\kappa(\mathfrak m)$ est une extension finie
de $k'$ par l'hypothèse de récurrence. Ainsi, $\kappa(\mathfrak m)$
est également finie sur $k$.

\medskip\noindent
Si $\mathfrak p = (0)$, nous considérons l'extension d'anneaux
$k[x_n] \subset k[x_1, \ldots, x_n]/\mathfrak m$.
Cette extension d'anneaux est de type fini, donc
de présentation finie d'après les
lemmes \ref{lemma-obvious-Noetherian} et
\ref{lemma-Noetherian-finite-type-is-finite-presentation}.
Ainsi, l'image de $\Spec(k[x_1, \ldots, x_n]/\mathfrak m)$
dans $\Spec(k[x_n])$ est constructible d'après le
théorème \ref{theorem-chevalley}. Puisque l'image
contient $(0)$, nous concluons qu'elle contient un ouvert
standard $D(f)$ pour un certain $f\in k[x_n]$ non nul. Comme
$D(f)$ est manifestement infini, nous obtenons une contradiction avec
l'hypothèse que $k[x_1, \ldots, x_n]/\mathfrak m$ est
un corps (et a donc un spectre constitué d'un seul point).

\medskip\noindent
Démonstration de (\ref{item-polynomial-ring-Jacobson}). Soit
$I \subset R$ un idéal radical, où $R$ est de type fini sur $k$.
Soit $f \in R$, $f \not \in I$. Nous devons trouver un idéal maximal
$\mathfrak m \subset R$ tel que $I \subset \mathfrak m$ et
$f \not \in \mathfrak m$. L'anneau $(R/I)_f$ est non nul, car
$1 = 0$ dans cet anneau signifierait que $f^n \in I$ et, puisque $I$ est
radical, cela signifierait que $f \in I$, contrairement à notre hypothèse sur $f$.
Nous pouvons donc choisir un idéal maximal $\mathfrak m'$
de $(R/I)_f$, voir le lemme \ref{lemma-Zariski-topology}.
Soit $\mathfrak m \subset R$
l'image réciproque de $\mathfrak m'$ dans $R$. Nous voyons que
$I \subset \mathfrak m$
et $f \not \in \mathfrak m$. Si nous montrons que $\mathfrak m$ est un idéal
maximal de $R$, alors nous avons terminé. Nous avons manifestement
$$
k \subset R/\mathfrak m \subset \kappa(\mathfrak m').
$$
D'après la partie (\ref{item-finite-kappa}), l'extension de corps
$\kappa(\mathfrak m')/k$ est finie. Ainsi,
$R/\mathfrak m$ est un corps d'après Corps, lemme
\ref{fields-lemma-subalgebra-algebraic-extension-field}.
Par conséquent, $\mathfrak m$ est maximal, ce qui achève la démonstration.
\end{proof}

\begin{lemma}
\label{lemma-field-finite-type-over-domain}
Soit $R$ un anneau. Soit $K$ un corps.
Si $R \subset K$ et si $K$ est de type fini sur $R$,
alors il existe un $f \in R$ tel que $R_f$ soit un corps
et que $K/R_f$ soit une extension finie de corps.
\end{lemma}

\begin{proof}
D'après le lemme \ref{lemma-characterize-image-finite-type}, il
existe un ouvert non vide $U \subset \Spec(R)$
contenu dans l'image $\{(0)\}$ de $\Spec(K) \to \Spec(R)$.
Choisissons $f \in R$, $f \not = 0$, tel que $D(f) \subset U$, c'est-à-dire
$D(f) = \{(0)\}$. Alors $R_f$ est un anneau intègre dont le spectre a exactement un
point, et $R_f$ est un corps. Ensuite, $K$ est une algèbre de type fini
sur le corps $R_f$, donc une extension finie de
$R_f$ d'après le Nullstellensatz de Hilbert (théorème \ref{theorem-nullstellensatz}).
\end{proof}





















\section{Anneaux de Jacobson}
\label{section-ring-jacobson}

\noindent
Soit $R$ un anneau. Les points fermés de $\Spec(R)$ sont les
idéaux maximaux de $R$. Souvent, les anneaux qui apparaissent naturellement en
géométrie algébrique ont beaucoup d'idéaux maximaux. C'est par exemple le cas des algèbres
de type fini sur un corps ou sur $\mathbf{Z}$. Nous montrerons que celles-ci
sont des exemples d'anneaux de Jacobson.

\begin{definition}
\label{definition-ring-jacobson}
Soit $R$ un anneau. Nous disons que $R$ est un
{\it anneau de Jacobson} si tout idéal
radical $I$ est l'intersection des
idéaux maximaux qui le contiennent.
\end{definition}

\begin{lemma}
\label{lemma-finite-type-field-Jacobson}
Toute algèbre de type fini sur un corps est de Jacobson.
\end{lemma}

\begin{proof}
Cela résulte du théorème \ref{theorem-nullstellensatz}
et de la définition \ref{definition-ring-jacobson}.
\end{proof}

\begin{lemma}
\label{lemma-jacobson-prime}
Soit $R$ un anneau. Si tout idéal premier de $R$ est
l'intersection des idéaux maximaux qui le contiennent,
alors $R$ est de Jacobson.
\end{lemma}

\begin{proof}
C'est immédiat du fait que
chaque idéal radical $I \subset R$ est
l'intersection des idéaux premiers qui le contiennent.
Voir le lemme \ref{lemma-Zariski-topology}.
\end{proof}

\begin{lemma}
\label{lemma-jacobson}
Un anneau $R$ est de Jacobson si et seulement si $\Spec(R)$
est un espace de Jacobson, voir Topologie,
définition \ref{topology-definition-space-jacobson}.
\end{lemma}

\begin{proof}
Supposons que $R$ soit de Jacobson. Soit $Z \subset \Spec(R)$
un fermé. Nous devons montrer que l'ensemble des points
fermés de $Z$ est dense dans $Z$. Soit $U \subset \Spec(R)$
un ouvert tel que $U \cap Z$ soit non vide.
Nous devons montrer que $Z \cap U$ contient un point fermé
de $\Spec(R)$. Nous pouvons
supposer que $U = D(f)$, car les ouverts standards forment une base de la
topologie de $\Spec(R)$. D'après le
lemme \ref{lemma-Zariski-topology}, nous pouvons supposer que
$Z = V(I)$, où $I$ est un idéal radical. Nous voyons également
que $f \not \in I$. Par hypothèse, il existe un
idéal maximal $\mathfrak m \subset R$ tel que
$I \subset \mathfrak m$ mais $f \not\in \mathfrak m$.
Ainsi, $\mathfrak m \in D(f) \cap V(I) = U \cap Z$, comme voulu.

\medskip\noindent
Réciproquement, supposons que $\Spec(R)$ soit un espace de Jacobson.
Soit $I \subset R$ un idéal radical. Soit
$J = \cap_{I \subset \mathfrak m} \mathfrak m$
l'intersection des idéaux maximaux qui contiennent $I$.
Manifestement, $J$ est un idéal radical, $V(J) \subset V(I)$, et
$V(J)$ est le plus petit fermé de $V(I)$ qui contient
tous les points fermés de $V(I)$. Par hypothèse, nous voyons que
$V(J) = V(I)$. Mais le lemme \ref{lemma-Zariski-topology}
montre qu'il existe une bijection entre les fermés de Zariski
et les idéaux radicaux ; ainsi, $I = J$, comme voulu.
\end{proof}

\begin{lemma}
\label{lemma-characterize-jacobson}
Soit $R$ un anneau. Si $R$ n'est pas de Jacobson, il existe
un idéal premier $\mathfrak p \subset R$ et un élément $f \in R$
tels que les propriétés suivantes soient satisfaites :
\begin{enumerate}
\item $\mathfrak p$ n'est pas un idéal maximal,
\item $f \not \in \mathfrak p$,
\item $V(\mathfrak p) \cap D(f) = \{\mathfrak p\}$, et
\item $(R/\mathfrak p)_f$ est un corps.
\end{enumerate}
En revanche, si $R$ est de Jacobson, alors, pour toute paire $(\mathfrak p, f)$
qui satisfait (1) et (2), l'ensemble $V(\mathfrak p) \cap D(f)$ est
infini.
\end{lemma}

\begin{proof}
Supposons que $R$ ne soit pas de Jacobson.
D'après le lemme \ref{lemma-jacobson}, cela signifie qu'il existe un
fermé $T \subset \Spec(R)$
dont l'ensemble $T_0 \subset T$ des points fermés n'est pas dense dans $T$.
Choisissons un $f \in R$ tel que $T_0 \subset V(f)$ mais
$T \not \subset V(f)$. Notons que $T \cap D(f)$
est homéomorphe à $\Spec((R/I)_f)$ si $T = V(I)$, voir les
lemmes \ref{lemma-spec-closed} et \ref{lemma-standard-open}.
Comme tout anneau possède un idéal maximal
(lemme \ref{lemma-Zariski-topology}), nous pouvons choisir un point fermé $t$ de
l'espace $T \cap D(f)$. Alors $t$ correspond à un idéal premier
$\mathfrak p \subset R$ qui n'est pas maximal (car $t \not \in T_0$).
Ainsi, (1) est satisfaite. Par construction, $f \not \in \mathfrak p$, d'où (2).
Comme $t$ est un point fermé de $T \cap D(f)$, nous voyons que
$V(\mathfrak p) \cap D(f) = \{\mathfrak p\}$, c'est-à-dire que (3) est satisfaite. Nous
concluons donc que $(R/\mathfrak p)_f$ est un anneau intègre dont le
spectre a un seul point, d'où (4)
(on peut par exemple combiner les lemmes \ref{lemma-characterize-local-ring} et
\ref{lemma-minimal-prime-reduced-ring}).

\medskip\noindent
Réciproquement, supposons que $R$ soit de Jacobson et que $(\mathfrak p, f)$
satisfasse (1) et (2). Si
$V(\mathfrak p) \cap D(f) =
\{\mathfrak p, \mathfrak q_1, \ldots, \mathfrak q_t\}$,
alors $\mathfrak p \not = \mathfrak q_i$
implique qu'il existe un élément $g \in R$ tel que $g \not \in \mathfrak p$
mais $g \in \mathfrak q_i$ pour tout $i$. Ainsi,
$V(\mathfrak p) \cap D(fg) = \{\mathfrak p\}$, ce qui
est impossible puisque tout sous-ensemble localement fermé de $\Spec(R)$
contient au moins un point fermé, car $\Spec(R)$ est
un espace topologique de Jacobson.
\end{proof}

\begin{lemma}
\label{lemma-pid-jacobson}
L'anneau $\mathbf{Z}$ est un anneau de Jacobson.
Plus généralement, soit $R$ un anneau tel que
\begin{enumerate}
\item $R$ soit un anneau intègre,
\item $R$ soit noethérien,
\item tout idéal premier non nul soit un idéal maximal, et
\item $R$ possède une infinité d'idéaux maximaux.
\end{enumerate}
Alors $R$ est un anneau de Jacobson.
\end{lemma}

\begin{proof}
Soit $R$ satisfaisant (1), (2), (3) et (4). L'assertion
signifie que $(0) = \bigcap_{\mathfrak m \subset R} \mathfrak m$.
Puisque $R$ possède une infinité d'idéaux maximaux, il suffit de
montrer que tout $x \in R$ non nul appartient à seulement
un nombre fini d'idéaux maximaux, autrement dit que $V(x)$ est fini.
D'après le lemme \ref{lemma-spec-closed},
nous voyons que $V(x)$ est homéomorphe à $\Spec(R/xR)$.
Par l'hypothèse (3), tout idéal premier de $R/xR$ est minimal et
correspond donc à une composante irréductible de $\Spec(R/xR)$
(lemme \ref{lemma-irreducible}).
Comme $R/xR$ est noethérien, l'espace topologique $\Spec(R/xR)$
est noethérien (lemme \ref{lemma-Noetherian-topology})
et possède un nombre fini de composantes irréductibles
(Topologie, lemme \ref{topology-lemma-Noetherian}).
Ainsi, $V(x)$ est fini, comme voulu.
\end{proof}

\begin{example}
\label{example-infinite-product-fields-jacobson}
Soit $A$ un ensemble infini.
Pour chaque $\alpha \in A$, soit $k_\alpha$ un corps.
Nous affirmons que $R = \prod_{\alpha\in A} k_\alpha$ est de Jacobson.
Remarquons d'abord que tout élément $f \in R$ est de la forme
$f = ue$, où $u \in R$ est une unité et $e\in R$ un idempotent
(laissé au lecteur).
Ainsi, $D(f) = D(e)$, et $R_f = R_e = R/(1-e)$ est un quotient de $R$.
En fait, tout anneau possédant cette propriété est de Jacobson.
En effet, soit $\mathfrak p \subset R$ un idéal premier
et $f \in R$, $f \not \in \mathfrak p$. Nous devons trouver
un idéal maximal $\mathfrak m$ de $R$ tel que
$\mathfrak p \subset \mathfrak m$ et $f \not\in \mathfrak m$.
Puisque $R_f$ est un quotient de $R$, nous voyons que tout idéal
maximal de $R_f$ correspond à un idéal maximal de $R$
ne contenant pas $f$. Le résultat s'obtient donc
en choisissant un idéal maximal de $R_f$ contenant $\mathfrak p R_f$.
\end{example}

\begin{example}
\label{example-not-jacobson}
Un anneau intègre $R$ possédant un nombre fini d'idéaux maximaux
$\mathfrak m_i$, $i = 1, \ldots, n$, n'est pas un
anneau de Jacobson, sauf s'il s'agit d'un corps.
En effet, dans ce cas, $(0)$ n'est pas l'intersection
des idéaux maximaux : $(0) \not =
\mathfrak m_1 \cap \mathfrak m_2 \cap \ldots \cap \mathfrak m_n
\supset \mathfrak m_1 \cdot \mathfrak m_2 \cdot \ldots
\cdot \mathfrak m_n \not = 0$.
En particulier, un anneau de valuation discrète, ou tout anneau local ayant
au moins deux idéaux premiers, n'est pas un anneau
de Jacobson.
\end{example}

\begin{lemma}
\label{lemma-finite-residue-extension-closed}
Soit $R \to S$ un morphisme d'anneaux.
Soit $\mathfrak m \subset R$ un idéal maximal.
Soit $\mathfrak q \subset S$ un idéal premier
au-dessus de $\mathfrak m$ tel que $\kappa(\mathfrak q)/\kappa(\mathfrak m)$
soit une extension algébrique de corps.
Alors $\mathfrak q$ est un idéal maximal de $S$.
\end{lemma}

\begin{proof}
Considérons le diagramme
$$
\xymatrix{
S \ar[r] & S/\mathfrak q \ar[r] & \kappa(\mathfrak q) \\
R \ar[r] \ar[u] & R/\mathfrak m \ar[u]
}
$$
Nous voyons que $\kappa(\mathfrak m) \subset S/\mathfrak q \subset
\kappa(\mathfrak q)$. Puisque l'extension de corps
$\kappa(\mathfrak m) \subset \kappa(\mathfrak q)$
est algébrique, tout anneau compris entre $\kappa(\mathfrak m)$
et $\kappa(\mathfrak q)$ est un corps
(Corps, lemme \ref{fields-lemma-subalgebra-algebraic-extension-field}).
Ainsi, $S/\mathfrak q$ est un corps, et est donc égal
à $\kappa(\mathfrak q)$.
\end{proof}

\begin{lemma}
\label{lemma-dimension}
Supposons que $k$ soit un corps et que $V$ soit un espace vectoriel non nul
sur $k$. Supposons que la dimension de $V$ (qui est un nombre cardinal)
soit strictement inférieure au cardinal de $k$. Alors, pour tout endomorphisme
$T : V \to V$, il existe un polynôme unitaire $P(t) \in k[t]$ tel que
$P(T)$ ne soit pas inversible.
\end{lemma}

\begin{proof}
Sinon, $V$ hérite d'une structure d'espace vectoriel sur
le corps $k(t)$. Or la dimension de $k(t)$ sur $k$ est
au moins égale au cardinal de $k$, par exemple parce que les éléments
$\frac{1}{t - \lambda}$ sont linéairement indépendants sur $k$.
\end{proof}

\noindent
Voici une autre version du Nullstellensatz de Hilbert.

\begin{theorem}
\label{theorem-uncountable-nullstellensatz}
Soit $k$ un corps. Soit $S$ une $k$-algèbre engendrée sur $k$
par les éléments $\{x_i\}_{i \in I}$. Supposons que le cardinal de $I$
soit strictement inférieur au cardinal de $k$. Alors
\begin{enumerate}
\item pour tout idéal maximal $\mathfrak m \subset S$, l'extension de corps
$\kappa(\mathfrak m)/k$
est algébrique, et
\item $S$ est un anneau de Jacobson.
\end{enumerate}
\end{theorem}

\begin{proof}
Si $I$ est fini, le résultat découle du Nullstellensatz de Hilbert,
théorème \ref{theorem-nullstellensatz}. Dans la suite de la démonstration, nous supposons
$I$ infini. Il suffit de démontrer le résultat lorsque
$\mathfrak m \subset k[\{x_i\}_{i \in I}]$ est maximal dans l'anneau
de polynômes en les variables $x_i$, puisque $S$ en est un quotient.
Puisque $I$ est infini, l'ensemble
des monômes $x_{i_1}^{e_1} \ldots x_{i_r}^{e_r}$, $i_1, \ldots, i_r \in I$
et $e_1, \ldots, e_r \geq 0$, est de cardinal inférieur ou égal au
cardinal de $I$. En effet, le cardinal de $I \times \ldots \times I$
est le cardinal de $I$, et le cardinal de
$\bigcup_{n \geq 0} I^n$ est lui aussi le même.
(Si $I$ est fini, cela n'est plus vrai et,
dans ce cas, cette démonstration ne fonctionne que si $k$ est non dénombrable.)

\medskip\noindent
Pour obtenir une contradiction, choisissons $T \in \kappa(\mathfrak m)$ transcendant
sur $k$. Notons que l'application $k$-linéaire
$T : \kappa(\mathfrak m) \to \kappa(\mathfrak m)$
donnée par la multiplication par $T$ possède la propriété que $P(T)$ est inversible
pour tout polynôme unitaire $P(t) \in k[t]$.
De plus, $\kappa(\mathfrak m)$ est de dimension au plus égale au cardinal de $I$
sur $k$, puisqu'il s'agit d'un quotient de l'espace vectoriel
$k[\{x_i\}_{i \in I}]$ sur $k$ (dont la dimension vaut $\# I$, comme vu ci-dessus).
C'est impossible d'après le lemme \ref{lemma-dimension}.

\medskip\noindent
Pour montrer que $S$ est de Jacobson, raisonnons comme suit. Sinon,
il existe un idéal premier $\mathfrak q \subset S$ et un élément $f \in S$,
$f \not \in \mathfrak q$, tels que $\mathfrak q$ ne soit pas maximal
et que $(S/\mathfrak q)_f$ soit un corps, voir le
lemme \ref{lemma-characterize-jacobson}.
Mais remarquons que $(S/\mathfrak q)_f$ est engendré par au plus $\# I + 1$ éléments.
Ainsi, l'extension de corps $(S/\mathfrak q)_f/k$ est algébrique
(d'après la première partie de la démonstration).
Cela implique que $\kappa(\mathfrak q)$ est une extension algébrique de $k$ ;
donc $\mathfrak q$ est maximal d'après le
lemme \ref{lemma-finite-residue-extension-closed}. Cette contradiction
achève la démonstration.
\end{proof}

\begin{lemma}
\label{lemma-base-change-Jacobson}
Soit $k$ un corps. Soit $S$ une $k$-algèbre.
Pour toute extension de corps $K/k$ dont le cardinal est strictement supérieur
au cardinal de $S$, nous avons :
\begin{enumerate}
\item pour tout idéal maximal $\mathfrak m$ de $S_K$, le corps
$\kappa(\mathfrak m)$ est algébrique sur $K$, et
\item $S_K$ est un anneau de Jacobson.
\end{enumerate}
\end{lemma}

\begin{proof}
Considérons $k \subset K$ tel que le cardinal de $K$ soit supérieur
au cardinal de $S$. Puisque les éléments de $S$ engendrent
la $K$-algèbre $S_K$, nous voyons que le
théorème \ref{theorem-uncountable-nullstellensatz}
s'applique.
\end{proof}

\begin{example}
\label{example-countable-trick-does-not-work}
L'astuce employée dans la démonstration du
théorème \ref{theorem-uncountable-nullstellensatz}
ne fonctionne réellement pas si $k$ est un corps dénombrable et $I$ est lui aussi dénombrable.
Soit $k$ un corps dénombrable. Soit $x$ une indéterminée,
et soit $k(x)$ le corps des fractions rationnelles en $x$.
Considérons l'algèbre polynomiale $R = k[x, \{x_f\}_{f \in k[x]-\{0\}}]$.
Soit $I = (\{fx_f - 1\}_{f\in k[x] - \{0\}})$. Remarquons que
$I$ est un idéal propre de $R$.
Choisissons un idéal maximal $I \subset \mathfrak m$.
Alors $k \subset R/\mathfrak m$ est isomorphe à
$k(x)$, et n'est pas algébrique sur $k$.
\end{example}

\begin{lemma}
\label{lemma-Jacobson-invert-element}
Soit $R$ un anneau de Jacobson. Soit $f \in R$. L'anneau $R_f$ est de Jacobson, et
les idéaux maximaux de $R_f$ correspondent aux idéaux maximaux de $R$ qui ne contiennent pas $f$.
\end{lemma}

\begin{proof}
D'après Topologie, lemme \ref{topology-lemma-jacobson-inherited},
nous voyons que $D(f) = \Spec(R_f)$ est de Jacobson et
que les points fermés de $D(f)$
correspondent aux points fermés de $\Spec(R)$
qui appartiennent à $D(f)$. Ainsi, $R_f$ est de Jacobson d'après le
lemme \ref{lemma-jacobson}.
\end{proof}

\begin{example}
\label{example-localize-not-preserve-closed-points}
Voici un exemple simple qui montre que le lemme \ref{lemma-Jacobson-invert-element}
est faux si $R$ n'est pas de Jacobson.
Considérons l'anneau $R = \mathbf{Z}_{(2)}$, c'est-à-dire la localisation
de $\mathbf{Z}$ en l'idéal premier $(2)$. La localisation de $R$ par
l'élément $2$ est isomorphe à $\mathbf{Q}$ ; sous forme d'une formule :
$R_2 \cong \mathbf{Q}$. Manifestement, le morphisme $R \to R_2$ envoie le
point fermé de $\Spec(\mathbf{Q})$ sur le point générique
de $\Spec(R)$.
\end{example}

\begin{example}
\label{example-infinite-localize-not-preserve-closed-points}
Voici un exemple simple qui montre que le
lemme \ref{lemma-Jacobson-invert-element}
est faux si $R$ est de Jacobson mais que nous localisons par une infinité
d'éléments.
En effet, soit $R = \mathbf{Z}$ et considérons la localisation
$(R \setminus \{0\})^{-1}R \cong \mathbf{Q}$
de $R$ par l'ensemble de tous les éléments non nuls.
Manifestement, le morphisme $\mathbf{Z} \to \mathbf{Q}$ envoie le
point fermé de $\Spec(\mathbf{Q})$ sur le point générique
de $\Spec(\mathbf{Z})$.
\end{example}

\begin{lemma}
\label{lemma-Jacobson-mod-ideal}
Soit $R$ un anneau de Jacobson. Soit $I \subset R$ un idéal.
L'anneau $R/I$ est de Jacobson, et les idéaux maximaux
de $R/I$ correspondent aux idéaux maximaux de $R$ qui contiennent $I$.
\end{lemma}

\begin{proof}
La démonstration est la même que celle du
lemme \ref{lemma-Jacobson-invert-element}.
\end{proof}

\begin{lemma}
\label{lemma-silly-jacobson}
Soit $R$ un anneau de Jacobson. Soit $K$ un corps. Supposons que $R \subset K$ et
que $K$ soit de type fini sur $R$. Alors $R$ est un corps et $K/R$
est une extension finie de corps.
\end{lemma}

\begin{proof}
Remarquons d'abord que $R$ est un anneau intègre.
D'après le lemme \ref{lemma-field-finite-type-over-domain},
nous voyons que $R_f$ est un corps et que $K/R_f$ est une extension finie de corps
pour un certain $f \in R$ non nul. Ainsi, $(0)$ est un idéal maximal de $R_f$
et, d'après le
lemme \ref{lemma-Jacobson-invert-element},
nous concluons que $(0)$ est un idéal maximal de $R$.
\end{proof}

\begin{proposition}
\label{proposition-Jacobson-permanence}
Soit $R$ un anneau de Jacobson. Soit $R \to S$ un
morphisme d'anneaux de type fini. Alors :
\begin{enumerate}
\item L'anneau $S$ est de Jacobson.
\item Le morphisme $\Spec(S) \to \Spec(R)$ envoie
les points fermés sur des points fermés.
\item Pour $\mathfrak m' \subset S$ maximal au-dessus de $\mathfrak m \subset R$,
l'extension de corps $\kappa(\mathfrak m')/\kappa(\mathfrak m)$
est finie.
\end{enumerate}
\end{proposition}

\begin{proof}
Soit $\mathfrak m' \subset S$ un idéal maximal, et
$R \cap \mathfrak m' = \mathfrak m$.
Alors $R/\mathfrak m \to S/\mathfrak m'$ satisfait
les conditions du lemme \ref{lemma-silly-jacobson}
d'après le lemme \ref{lemma-Jacobson-mod-ideal}.
Ainsi, $R/\mathfrak m$ est un corps,
$\mathfrak m$ est un idéal maximal et l'extension induite
des corps résiduels est finie. Cela démontre (2) et (3).  

\medskip\noindent
Si $S$ n'est pas de Jacobson, alors, d'après le lemme \ref{lemma-characterize-jacobson}, il
existe un idéal premier non maximal $\mathfrak q$ de $S$ et un
$g \in S$, $g \not\in \mathfrak q$, tels que $(S/\mathfrak q)_g$ soit un corps.
Pour obtenir une contradiction, nous montrons que $\mathfrak q$ est un idéal maximal.
Soit $\mathfrak p = \mathfrak q \cap R$. Alors
$R/\mathfrak p \to (S/\mathfrak q)_g$ satisfait les conditions du
lemme \ref{lemma-silly-jacobson} d'après le
lemme \ref{lemma-Jacobson-mod-ideal}.
Ainsi, $R/\mathfrak p$ est un corps et l'extension de corps
$\kappa(\mathfrak p) \to (S/\mathfrak q)_g = \kappa(\mathfrak q)$ est
finie, donc algébrique. Alors $\mathfrak q$ est un idéal maximal de $S$ d'après le
lemme \ref{lemma-finite-residue-extension-closed}. Contradiction.
\end{proof}

\begin{lemma}
\label{lemma-corollary-jacobson}
Toute algèbre de type fini sur $\mathbf{Z}$ est de Jacobson.
\end{lemma}

\begin{proof}
Combiner le lemme \ref{lemma-pid-jacobson} et la
proposition \ref{proposition-Jacobson-permanence}.
\end{proof}

\begin{lemma}
\label{lemma-image-finite-type-map-Jacobson-rings}
Soit $R \to S$ un morphisme de type fini entre anneaux de Jacobson.
Posons $X = \Spec(R)$ et $Y = \Spec(S)$.
Notons $f : Y \to X$ le morphisme induit
sur les spectres. Soit $E \subset Y = \Spec(S)$ un ensemble
constructible. Notons par un indice ${}_0$ l'ensemble
des points fermés d'un espace topologique.
\begin{enumerate}
\item Nous avons $f(E)_0 = f(E_0) = X_0 \cap f(E)$.
\item Un point $\xi \in X$ appartient à $f(E)$ si et seulement si
$\overline{\{\xi\}} \cap f(E_0)$ est dense dans $\overline{\{\xi\}}$.
\end{enumerate}
\end{lemma}

\begin{proof}
Nous avons un diagramme commutatif d'applications continues
$$
\xymatrix{
E \ar[r] \ar[d] & Y \ar[d] \\
f(E) \ar[r] & X
}
$$
Supposons que $x \in f(E)$ soit fermé dans $f(E)$. Alors $f^{-1}(\{x\})\cap E$
est non vide et fermé dans $E$. En appliquant le
lemme \ref{topology-lemma-jacobson-inherited} de Topologie
aux deux inclusions
$$
f^{-1}(\{x\}) \cap E \subset E \subset Y
$$
nous trouvons qu'il existe un point $y \in f^{-1}(\{x\}) \cap E$ qui est
fermé dans $Y$. Autrement dit, il existe $y \in Y_0$ et $y \in E_0$
qui s'envoie sur $x$. Ainsi, $x \in f(E_0)$.
Cela démontre que $f(E)_0 \subset f(E_0)$.
La proposition \ref{proposition-Jacobson-permanence} implique que
$f(E_0) \subset X_0 \cap f(E)$. L'inclusion
$X_0 \cap f(E) \subset f(E)_0$ est triviale. Cela démontre la
première assertion.

\medskip\noindent
Supposons que $\xi \in f(E)$. D'après le
lemme \ref{lemma-characterize-image-finite-type},
l'ensemble $f(E) \cap \overline{\{\xi\}}$ contient un ouvert
dense de $\overline{\{\xi\}}$. Puisque $X$ est de Jacobson,
nous concluons que $f(E) \cap \overline{\{\xi\}}$ contient un
ensemble dense de points fermés, voir Topologie,
lemme \ref{topology-lemma-jacobson-inherited}.
Nous concluons par la partie (1) du lemme.

\medskip\noindent
Réciproquement, supposons que $\overline{\{\xi\}} \cap f(E_0)$
soit dense dans $\overline{\{\xi\}}$. D'après le
lemme \ref{lemma-constructible-is-image},
il existe un morphisme d'anneaux $S \to S'$ de présentation finie
tel que $E$ soit l'image de $Y' := \Spec(S') \to Y$.
Alors $E_0$ est l'image de $Y'_0$ d'après la première partie du
lemme appliquée au morphisme d'anneaux $S \to S'$. Nous pouvons donc supposer que
$E = Y$ en remplaçant $S$ par $S'$. Supposons que $\xi$ corresponde
à $\mathfrak p \subset R$. Considérons le diagramme
$$
\xymatrix{
S \ar[r] & S/\mathfrak p S \\
R \ar[r] \ar[u] & R/\mathfrak p \ar[u]
}
$$
Ce diagramme et la densité de $f(Y_0) \cap V(\mathfrak p)$
dans $V(\mathfrak p)$
montrent que le morphisme $R/\mathfrak p \to S/\mathfrak p S$
satisfait la condition (2) du
lemme \ref{lemma-domain-image-dense-set-points-generic-point}.
Nous en concluons
qu'il existe un idéal premier $\overline{\mathfrak q} \subset S/\mathfrak pS$
qui s'envoie sur $(0)$. Autrement dit, l'image réciproque $\mathfrak q$
de $\overline{\mathfrak q}$ dans $S$ s'envoie sur $\mathfrak p$, comme voulu.
\end{proof}

\noindent
La conclusion du lemme ci-dessus est que nous pouvons déterminer
l'image de $f$ à partir de l'ensemble des points fermés de l'image.
C'est un peu plus commode lorsque le morphisme est de présentation finie,
car nous savons alors que l'image d'une partie constructible est constructible.
Avant de l'énoncer, introduisons quelques notations.
Notons $\text{Constr}(X)$ l'ensemble des parties constructibles.
Soit $R \to S$ un morphisme d'anneaux.
Posons $X = \Spec(R)$ et $Y = \Spec(S)$.
Notons $f : Y \to X$ le morphisme induit sur les spectres.
Notons par un indice ${}_0$ l'ensemble
des points fermés d'un espace topologique.


\begin{lemma}
\label{lemma-conclude-jacobson-Noetherian}
Avec les notations ci-dessus. Supposons que $R$ soit un anneau de Jacobson noethérien.
Supposons en outre que $R \to S$ soit de type fini.
Il existe un diagramme commutatif
$$
\xymatrix{
\text{Constr}(Y) \ar[r]^{E \mapsto E_0} \ar[d]^{E \mapsto f(E)} &
\text{Constr}(Y_0) \ar[d]^{E \mapsto f(E)} \\
\text{Constr}(X) \ar[r]^{E \mapsto E_0} &
\text{Constr}(X_0)
}
$$
où les flèches horizontales sont les bijections du
lemme \ref{topology-lemma-jacobson-equivalent-constructible} de Topologie.
\end{lemma}

\begin{proof}
Puisque $R \to S$ est de type fini, il est de présentation finie,
voir le lemme \ref{lemma-Noetherian-finite-type-is-finite-presentation}.
Ainsi, l'image d'une partie constructible dans $X$ est constructible
dans $Y$ d'après le théorème de Chevalley
(théorème \ref{theorem-chevalley}).
Combiné au
lemme \ref{lemma-image-finite-type-map-Jacobson-rings},
cela démontre le lemme.
\end{proof}

\noindent
Pour illustrer l'emploi des anneaux de Jacobson, donnons les deux exemples suivants.

\begin{example}
\label{example-product-matrices-zero}
Soit $k$ un corps. L'espace $\Spec(k[x, y]/(xy))$
possède deux composantes irréductibles : l'axe des $x$ et
l'axe des $y$. Comme généralisation, soit
$$
R = k[x_{11}, x_{12}, x_{21}, x_{22}, y_{11}, y_{12}, y_{21}, y_{22}]/
\mathfrak a,
$$
où $\mathfrak a$ est l'idéal de
$k[x_{11}, x_{12}, x_{21}, x_{22}, y_{11}, y_{12}, y_{21}, y_{22}]$
engendré par les coefficients du produit des matrices $2 \times 2$
$$
\left(
\begin{matrix}
x_{11} & x_{12}\\
x_{21} & x_{22}
\end{matrix}
\right)
 \left(
\begin{matrix}
y_{11} & y_{12}\\
y_{21} & y_{22}
\end{matrix}
\right).
$$
Dans cet exemple, nous allons décrire $\Spec(R)$.

\medskip\noindent
Pour démontrer l'assertion concernant $\Spec(k[x, y]/(xy))$, raisonnons comme suit.
Si $\mathfrak p \subset k[x, y]$ est un idéal premier contenant $xy$, alors
soit $x$, soit $y$ appartient à $\mathfrak p$. Les idéaux premiers minimaux
ayant cette propriété sont donc simplement $(x)$ et $(y)$. Lorsque $k$ est
algébriquement clos, le $\text{max-Spec}$ de ces composantes
se visualise alors comme les ensembles de points des axes des $y$ et des $x$.

\medskip\noindent
Pour la généralisation, remarquons que nous pouvons identifier les points
fermés du spectre de
$k[x_{11}, x_{12}, x_{21}, x_{22}, y_{11}, y_{12}, y_{21}, y_{22}])$
à l'espace des matrices
$$
\left\{ (X, Y) \in \text{Mat}(2, k)\times \text{Mat}(2, k) \mid
X = \left(
\begin{matrix}
x_{11} & x_{12}\\
x_{21} & x_{22}
\end{matrix}
\right),
Y= \left(
\begin{matrix}
y_{11} & y_{12}\\
y_{21} & y_{22}
\end{matrix}
\right)
\right\}
$$
du moins si $k$ est algébriquement clos.
Définissons maintenant une action du groupe
$\text{GL}(2, k)\times \text{GL}(2, k)\times \text{GL}(2, k)$
sur l'espace des matrices $\{(X, Y)\}$ par
$$
(g_1, g_2, g_3) \times (X, Y) \mapsto ((g_1Xg_2^{-1}, g_2Yg_3^{-1})).
$$
Observons également que l'ensemble algébrique
$$
\text{GL}(2, k)\times \text{GL}(2, k)\times \text{GL}(2, k) \subset
\text{Mat}(2, k)\times \text{Mat}(2, k) \times \text{Mat}(2, k)
$$
est irréductible, puisqu'il est le spectre maximal de l'anneau intègre
$$
k[x_{11}, x_{12}, \ldots, z_{21}, z_{22}, (x_{11}x_{22}-x_{12}x_{21})^{-1}
, (y_{11}y_{22}-y_{12}y_{21})^{-1}, (z_{11}z_{22}-z_{12}z_{21})^{-1}].
$$
Puisque l'image d'un ensemble algébrique irréductible est encore
irréductible, il suffit de classifier les orbites de l'ensemble
$\{(X, Y)\in \text{Mat}(2, k)\times \text{Mat}(2, k)|XY = 0\}$ et de prendre leurs
adhérences. L'algèbre linéaire usuelle nous ramène aux
trois cas suivants :
\begin{enumerate}
\item $\exists (g_1, g_2)$ tel que $g_1Xg_2^{-1} = I_{2\times 2}$.
Alors $Y$ vaut nécessairement $0$ ; cet ensemble algébrique est
invariant sous l'action du groupe. Il s'ensuit que cette orbite est
contenue dans l'ensemble algébrique irréductible défini par l'idéal premier
$(y_{11}, y_{12}, y_{21}, y_{22})$. En prenant l'adhérence, nous voyons
que $(y_{11}, y_{12}, y_{21}, y_{22})$ est effectivement une composante.
\item $\exists (g_1, g_2)$ tel que
$$
g_1Xg_2^{-1} = \left(
\begin{matrix}
1 & 0 \\
0 & 0
\end{matrix}
\right).
$$
Ce cas se produit si et seulement si $X$ est une matrice de rang 1,
et, de plus, $Y$ est annulée par une telle matrice $X$ si et seulement si
$$
x_{11}y_{11}+x_{12}y_{21} = 0; \quad x_{11}y_{12}+x_{12}y_{22} = 0;
$$
$$
x_{21}y_{11}+x_{22}y_{21} = 0; \quad x_{21}y_{12}+x_{22}y_{22} = 0.
$$
Fixons une matrice $X$ de rang 1 ; les matrices $Y$ non nulles qui satisfont les
équations ci-dessus forment un ensemble algébrique irréductible pour la
raison suivante ($Y = 0$ relève du cas précédent) :
$0 = g_1Xg_2^{-1}g_2Y$ implique que
$$
g_2Y = \left(
\begin{matrix}
0 & 0 \\
y_{21}' & y_{22}'
\end{matrix}
\right).
$$
À l'aide d'une autre action de $\text{GL}(2, k)$ à droite par $g_3$,
$g_2Y$ peut être amenée à la forme
$$
g_2Yg_3^{-1} = \left(
\begin{matrix}
0 & 0 \\
0 & 1
\end{matrix}
\right),
$$
et ces matrices $Y$ forment donc un ensemble algébrique irréductible
isomorphe à l'image de $\text{GL}(2, k)$ sous cette action. Enfin,
remarquons que la condition « de rang 1 » sur les matrices $X$ définit un ouvert dense
de l'ensemble algébrique irréductible
$\det X = x_{11}x_{22} - x_{12}x_{21} = 0$.
Il s'ensuit maintenant que les cinq équations définissent une composante irréductible
$(x_{11}y_{11}+x_{12}y_{21}, x_{11}y_{12}+x_{12}y_{22}, x_{21}y_{11}
+x_{22}y_{21}, x_{21}y_{12}+x_{22}y_{22}, x_{11}x_{22}-x_{12}x_{21})$
dans l'ouvert de l'espace des paires de matrices non nulles.
On peut montrer que les deux équations
$\det X = 0$, $\det Y = 0$ découpent $\Spec(R)$ suivant une composante irréductible
dont le lieu ci-dessus est un ouvert dense.
\item $\exists (g_1, g_2)$ tel que $g_1Xg_2^{-1} = 0$, ou,
de manière équivalente, $X = 0$. Alors $Y$ peut être arbitraire, et cette composante
est donc définie par $(x_{11}, x_{12}, x_{21}, x_{22})$.
\end{enumerate}
\end{example}


\begin{example}
\label{example-idempotent-matrices}
Pour un autre exemple, considérons
$R = k[\{t_{ij}\}_{i, j = 1}^{n}]/\mathfrak a$, où $\mathfrak a$ est
l'idéal engendré par les coefficients de la matrice produit $T^2-T$,
$T = (t_{ij})$. L'algèbre linéaire nous apprend que, sous
l'action de $GL(n, k)$ définie par $g, T \mapsto gTg^{-1}$, $T$ est
classifiée par son rang, et chaque $T$ est conjuguée à une matrice
$\text{diag}(1, \ldots, 1, 0, \ldots, 0)$, qui possède $r$ coefficients égaux à 1 et $n-r$ coefficients égaux à 0.
Ainsi, chaque orbite d'une telle matrice $\text{diag}(1, \ldots, 1, 0, \ldots, 0)$ sous
l'action du groupe forme une composante irréductible, et toute matrice
idempotente appartient à l'une de ces orbites. Montrons maintenant que deux
orbites distinctes sont nécessairement disjointes. Pour cela, il
suffit de construire des fonctions polynomiales qui prennent des
valeurs différentes sur des orbites différentes. En caractéristique 0, une telle
fonction peut être
$f(t_{ij}) = trace(T) = \sum_{i = 1}^nt_{ii}$. En caractéristique
positive, les choses sont un peu plus délicates, car nous
pouvons avoir $trace(T) = 0$ même si $T \neq 0$. Par exemple, $char = 3$
$$
trace\left(
\begin{matrix}
1 & & \\
& 1 & \\
& & 1
\end{matrix}
\right) = 3 = 0
$$
Quoi qu'il en soit, ces composantes peuvent être séparées à l'aide d'autres fonctions.
Par exemple, en caractéristique 3, $tr(\wedge^3T)$ prend
la valeur 1 sur les composantes correspondant à $diag(1, 1, 1)$ et 0 sur
les autres composantes.
\end{example}





































\section{Extensions d'anneaux finies et entières}
\label{section-finite-ring-extensions}

\noindent
Lemmes élémentaires concernant les morphismes d'anneaux finis et entiers.
Rappelons la définition.

\begin{definition}
\label{definition-integral-ring-map}
Soit $\varphi : R \to S$ un morphisme d'anneaux.
\begin{enumerate}
\item Un élément $s \in S$
est {\it entier sur $R$} s'il existe un polynôme
unitaire $P(x) \in R[x]$ tel que
$P^\varphi(s) = 0$, où $P^\varphi(x) \in S[x]$
est l'image de $P$ par $\varphi : R[x] \to S[x]$.
\item Le morphisme d'anneaux $\varphi$ est {\it entier}
si tout $s \in S$ est entier sur $R$.
\end{enumerate}
\end{definition}

\begin{lemma}
\label{lemma-characterize-integral-element}
Soit $\varphi : R \to S$ un morphisme d'anneaux. Soit $y \in S$. S'il existe un
sous-$R$-module de type fini $M$ de $S$ tel que $1 \in M$ et $yM \subset M$,
alors $y$ est entier sur $R$.
\end{lemma}

\begin{proof}
Considérons l'application $\varphi : M \to M$, $x \mapsto y \cdot x$.
D'après le lemme \ref{lemma-charpoly-module}, il existe un polynôme unitaire
$P \in R[T]$ tel que $P(\varphi) = 0$. Dans l'anneau $S$, nous obtenons
$P(y) = P(y) \cdot 1 = P(\varphi)(1) = 0$.
\end{proof}

\begin{lemma}
\label{lemma-finite-is-integral}
Un morphisme d'anneaux fini est entier.
\end{lemma}

\begin{proof}
Soit $R \to S$ fini. Soit $y \in S$. Appliquons le
lemme \ref{lemma-characterize-integral-element}
à $M = S$ pour voir que $y$ est entier sur $R$.
\end{proof}

\begin{lemma}
\label{lemma-characterize-integral}
Soit $\varphi : R \to S$ un morphisme d'anneaux. Soient $s_1, \ldots, s_n$
un ensemble fini d'éléments de $S$.
Alors les $s_i$ sont entiers sur $R$ pour tout $i = 1, \ldots, n$
si et seulement s'il
existe une $R$-sous-algèbre $S' \subset S$ finie sur $R$
qui contient tous les $s_i$.
\end{lemma}

\begin{proof}
Si chaque $s_i$ est entier, alors la sous-algèbre
engendrée par $\varphi(R)$ et les $s_i$ est finie
sur $R$. En effet, si $s_i$ satisfait une équation unitaire
de degré $d_i$ sur $R$, alors cette sous-algèbre est engendrée comme
$R$-module par les éléments $s_1^{e_1} \ldots s_n^{e_n}$
avec $0 \leq e_i \leq d_i - 1$.
Réciproquement, supposons donnée une $R$-sous-algèbre finie
$S'$ qui contient tous les $s_i$. Alors tous les
$s_i$ sont entiers d'après le lemme \ref{lemma-finite-is-integral}.
\end{proof}

\begin{lemma}
\label{lemma-characterize-finite-in-terms-of-integral}
Soit $R \to S$ un morphisme d'anneaux. Les propriétés suivantes sont équivalentes :
\begin{enumerate}
\item $R \to S$ est fini,
\item $R \to S$ est entier et de type fini, et
\item il existe $x_1, \ldots, x_n \in S$ qui engendrent $S$ comme
algèbre sur $R$ et tels que chaque $x_i$ soit entier sur $R$.
\end{enumerate}
\end{lemma}

\begin{proof}
C'est clair d'après le lemme \ref{lemma-characterize-integral}.
\end{proof}

\begin{lemma}
\label{lemma-integral-transitive}
\begin{slogan}
Une composée de morphismes d'anneaux entiers est entière
\end{slogan}
Supposons que $R \to S$ et $S \to T$ soient des
morphismes d'anneaux entiers. Alors $R \to T$ est entier.
\end{lemma}

\begin{proof}
Soit $t \in T$. Soit $P(x) \in S[x]$ un
polynôme unitaire tel que $P(t) = 0$.
Appliquons le lemme \ref{lemma-characterize-integral}
à l'ensemble fini des coefficients de $P$.
Ainsi, $t$ est entier sur une certaine sous-algèbre
$S' \subset S$ finie sur $R$. Appliquons à nouveau le lemme
\ref{lemma-characterize-integral} pour trouver
une sous-algèbre $T' \subset T$ finie sur $S'$ et
contenant $t$. Le lemme \ref{lemma-finite-transitive}
appliqué à $R \to S' \to T'$ montre que $T'$ est finie
sur $R$. L'intégralité de $t$ sur $R$
résulte alors du lemme \ref{lemma-finite-is-integral}.
\end{proof}

\begin{lemma}
\label{lemma-integral-closure-is-ring}
Soit $R \to S$ un morphisme d'anneaux.
L'ensemble
$$
S' = \{s \in S \mid s\text{ est entier sur }R\}
$$
est une $R$-sous-algèbre de $S$.
\end{lemma}

\begin{proof}
C'est clair d'après les lemmes \ref{lemma-characterize-integral}
et \ref{lemma-finite-is-integral}.
\end{proof}

\begin{lemma}
\label{lemma-finite-product-integral}
Soient $R_i\to S_i$ des morphismes d'anneaux, $i = 1, \ldots, n$.
Notons respectivement $R$ et $S$ les produits des $R_i$ et des $S_i$.
Alors un élément $s = (s_1, \ldots, s_n) \in S$ est entier sur $R$
si et seulement si chaque $s_i$ est entier sur $R_i$.
\end{lemma}

\begin{proof}
Omis.
\end{proof}

\begin{definition}
\label{definition-integral-closure}
Soit $R \to S$ un morphisme d'anneaux.
L'anneau $S' \subset S$ des éléments entiers sur
$R$, voir le lemme \ref{lemma-integral-closure-is-ring},
est appelé la {\it clôture intégrale} de $R$
dans $S$. Si $R \subset S$, nous disons que $R$ est
{\it intégralement clos} dans $S$ si $R = S'$.
\end{definition}

\noindent
En particulier, nous voyons que $R \to S$ est entier si et seulement
si la clôture intégrale de $R$ dans $S$ est $S$ tout entier.

\begin{lemma}
\label{lemma-finite-product-integral-closure}
Soient $R_i\to S_i$ des morphismes d'anneaux, $i = 1, \ldots, n$.
Notons la clôture intégrale de $R_i$ dans $S_i$ par $S'_i$.
Notons en outre respectivement $R$ et $S$ les produits des $R_i$ et des $S_i$.
Alors la clôture intégrale de $R$ dans $S$
est le produit des $S'_i$. En particulier, $R \to S$ est
intégralement clos si et seulement si chaque $R_i \to S_i$ est intégralement clos.
\end{lemma}

\begin{proof}
Cela résulte immédiatement du lemme \ref{lemma-finite-product-integral}.
\end{proof}

\begin{lemma}
\label{lemma-integral-closure-localize}
La clôture intégrale commute à la localisation : si $A \to B$ est un
morphisme d'anneaux et $S \subset A$ une partie multiplicative, alors la clôture
intégrale de $S^{-1}A$ dans $S^{-1}B$ est $S^{-1}B'$, où $B' \subset B$
est la clôture intégrale de $A$ dans $B$.
\end{lemma}

\begin{proof}
Puisque la localisation est exacte, nous voyons que $S^{-1}B' \subset S^{-1}B$.
Supposons que $x \in B'$ et $f \in S$. Alors
$x^d + \sum_{i = 1, \ldots, d} a_i x^{d - i} = 0$
dans $B$ pour certains $a_i \in A$. Ainsi,
$$
(x/f)^d + \sum\nolimits_{i = 1, \ldots, d} a_i/f^i (x/f)^{d - i} = 0
$$
dans $S^{-1}B$. Nous voyons ainsi que $S^{-1}B'$ est contenu dans
la clôture intégrale de $S^{-1}A$ dans $S^{-1}B$. Réciproquement, supposons
que $x/f \in S^{-1}B$ soit entier sur $S^{-1}A$. Alors nous avons
$$
(x/f)^d + \sum\nolimits_{i = 1, \ldots, d} (a_i/f_i) (x/f)^{d - i} = 0
$$
dans $S^{-1}B$ pour certains $a_i \in A$ et $f_i \in S$. Cela signifie que
$$
(f'f_1 \ldots f_d x)^d +
\sum\nolimits_{i = 1, \ldots, d}
f^i(f')^if_1^i \ldots f_i^{i - 1} \ldots f_d^i a_i
(f'f_1 \ldots f_dx)^{d - i} = 0
$$
pour un $f' \in S$ convenable. Ainsi, $f'f_1\ldots f_dx \in B'$ et donc
$x/f \in S^{-1}B'$, comme voulu.
\end{proof}

\begin{lemma}
\label{lemma-integral-closure-stalks}
\begin{slogan}
Un élément d'une algèbre sur un anneau est entier sur cet anneau
si et seulement s'il est localement entier en tout idéal premier de l'anneau.
\end{slogan}
Soit $\varphi : R \to S$ un morphisme d'anneaux.
Soit $x \in S$. Les propriétés suivantes sont équivalentes :
\begin{enumerate}
\item $x$ est entier sur $R$, et
\item pour tout idéal premier $\mathfrak p \subset R$, l'élément
$x \in S_{\mathfrak p}$ est entier sur $R_{\mathfrak p}$.
\end{enumerate}
\end{lemma}

\begin{proof}
Il est clair que (1) implique (2). Supposons (2). Considérons la $R$-algèbre
$S' \subset S$ engendrée par $\varphi(R)$ et $x$. Soit $\mathfrak p$ un
idéal premier de $R$. Nous savons alors que
$x^d + \sum_{i = 1, \ldots, d} \varphi(a_i) x^{d - i} = 0$
dans $S_{\mathfrak p}$ pour certains $a_i \in R_{\mathfrak p}$. En examinant
les dénominateurs qui interviennent, nous voyons donc que,
pour un certain $f \in R$, $f \not \in \mathfrak p$, nous avons
$a_i \in R_f$ et
$x^d + \sum_{i = 1, \ldots, d} \varphi(a_i) x^{d - i} = 0$
dans $S_f$. Cela implique que $S'_f$ est fini sur $R_f$.
Puisque $\mathfrak p$ était arbitraire et que $\Spec(R)$ est quasi-compact
(lemme \ref{lemma-quasi-compact}), nous pouvons trouver un nombre fini d'éléments
$f_1, \ldots, f_n \in R$
qui engendrent l'idéal unité de $R$ tels que $S'_{f_i}$ soit fini
sur $R_{f_i}$. Nous concluons donc du lemme \ref{lemma-cover} que
$S'$ est fini sur $R$. Ainsi, $x$ est entier sur $R$ d'après le
lemme \ref{lemma-characterize-integral}.
\end{proof}

\begin{lemma}
\label{lemma-base-change-integral}
\begin{slogan}
L'intégralité et la finitude sont préservées par changement de base.
\end{slogan}
Soient $R \to S$ et $R \to R'$ des morphismes d'anneaux.
Posons $S' = R' \otimes_R S$.
\begin{enumerate}
\item Si $R \to S$ est entier, alors $R' \to S'$ l'est aussi.
\item Si $R \to S$ est fini, alors $R' \to S'$ l'est aussi.
\end{enumerate}
\end{lemma}

\begin{proof}
Nous démontrons (1).
Soient $s_i \in S$ des générateurs de $S$ sur $R$.
Chacun d'eux satisfait une équation polynomiale unitaire $P_i$
sur $R$. Ainsi, les éléments $1 \otimes s_i \in S'$ engendrent
$S'$ sur $R'$ et satisfont le polynôme correspondant
$P_i'$ sur $R'$. Puisque ces éléments engendrent $S'$ sur $R'$,
nous voyons que $S'$ est entier sur $R'$.
Démonstration de (2) omise.
\end{proof}

\begin{lemma}
\label{lemma-integral-local}
Soit $R \to S$ un morphisme d'anneaux.
Soient $f_1, \ldots, f_n \in R$ engendrant l'idéal unité.
\begin{enumerate}
\item Si chaque $R_{f_i} \to S_{f_i}$ est entier, alors $R \to S$ l'est aussi.
\item Si chaque $R_{f_i} \to S_{f_i}$ est fini, alors $R \to S$ l'est aussi.
\end{enumerate}
\end{lemma}

\begin{proof}
Démonstration de (1).
Soit $s \in S$. Considérons l'idéal $I \subset R[x]$ des
polynômes $P$ tels que $P(s) = 0$. Notons $J \subset R$
l'idéal (!) des coefficients dominants des éléments de $I$.
Par hypothèse et en chassant les dénominateurs,
nous voyons que $f_i^{n_i} \in J$ pour tout $i$
et certains $n_i \geq 0$. Ainsi, $J$ contient $1$, et nous voyons que
$s$ est entier sur $R$. Démonstration de (2) omise.
\end{proof}

\begin{lemma}
\label{lemma-integral-permanence}
Soient $A \to B \to C$ des morphismes d'anneaux.
\begin{enumerate}
\item Si $A \to C$ est entier, alors $B \to C$ l'est aussi.
\item Si $A \to C$ est fini, alors $B \to C$ l'est aussi.
\end{enumerate}
\end{lemma}

\begin{proof}
Omis.
\end{proof}

\begin{lemma}
\label{lemma-integral-closure-transitive}
Soient $A \to B \to C$ des morphismes d'anneaux.
Soit $B'$ la clôture intégrale de $A$ dans $B$,
et soit $C'$ la clôture intégrale de $B'$ dans $C$. Alors
$C'$ est la clôture intégrale de $A$ dans $C$.
\end{lemma}

\begin{proof}
Omis.
\end{proof}

\begin{lemma}
\label{lemma-integral-overring-surjective}
Supposons que $R \to S$ soit une extension
d'anneaux entière avec $R \subset S$.
Alors $\varphi : \Spec(S) \to \Spec(R)$
est surjective.
\end{lemma}

\begin{proof}
Soit $\mathfrak p \subset R$ un idéal premier.
Nous devons montrer que $\mathfrak pS_{\mathfrak p} \not = S_{\mathfrak p}$, voir le
lemme \ref{lemma-in-image}.
La localisation $R_{\mathfrak p} \to S_{\mathfrak p}$ est injective
(car la localisation est exacte) et entière d'après le
lemme \ref{lemma-integral-closure-localize} ou
\ref{lemma-base-change-integral}.
Nous pouvons donc remplacer $R$, $S$ par $R_{\mathfrak p}$, $S_{\mathfrak p}$ et
supposer que $R$ est local d'idéal maximal $\mathfrak m$ ; il
suffit alors de montrer que $\mathfrak mS \not = S$.
Supposons que $1 = \sum f_i s_i$, avec $f_i \in \mathfrak m$
et $s_i \in S$, afin d'obtenir une contradiction.
Soit $R \subset S' \subset S$
tel que $R \to S'$ soit fini et que $s_i \in S'$, voir le
lemme \ref{lemma-characterize-integral}.
L'équation $1 = \sum f_i s_i$ implique que
le $R$-module de type fini $S'$ satisfait $S' = \mathfrak m S'$. Le
lemme de Nakayama \ref{lemma-NAK} donne donc
$S' = 0$. Contradiction.
\end{proof}

\begin{lemma}
\label{lemma-integral-under-field}
Soit $R$ un anneau. Soit $K$ un corps.
Si $R \subset K$ et si $K$ est entier sur $R$,
alors $R$ est un corps et $K$ est une extension algébrique.
Si $R \subset K$ et si $K$ est fini sur $R$,
alors $R$ est un corps et $K$ est une extension algébrique finie.
\end{lemma}

\begin{proof}
Supposons que $R \subset K$ soit entier.
D'après le lemme \ref{lemma-integral-overring-surjective}, nous voyons que
$\Spec(R)$ possède $1$ point. Comme $R$ est manifestement un anneau intègre, nous voyons
que $R = R_{(0)}$ est un corps (lemme \ref{lemma-minimal-prime-reduced-ring}).
Les autres assertions en résultent immédiatement.
\end{proof}

\begin{lemma}
\label{lemma-integral-over-field}
Soit $k$ un corps. Soit $S$ une $k$-algèbre sur $k$.
\begin{enumerate}
\item Si $S$ est un anneau intègre de dimension finie sur $k$,
alors $S$ est un corps.
\item Si $S$ est entier sur $k$ et est un anneau intègre,
alors $S$ est un corps.
\item Si $S$ est entier sur $k$, alors tout idéal premier de
$S$ est un idéal maximal (voir le
lemme \ref{lemma-ring-with-only-minimal-primes}
pour d'autres conséquences).
\end{enumerate}
\end{lemma}

\begin{proof}
L'assertion sur les idéaux premiers résulte de l'assertion
« entier $+$ intègre $\Rightarrow$ corps ».
Soit $S$ entier sur $k$ et supposons que $S$ soit un anneau intègre.
Prenons $s \in S$. D'après le lemme
\ref{lemma-characterize-integral}, nous pouvons trouver une
$k$-sous-algèbre de dimension finie $k \subset S' \subset S$
qui contient $s$. Ainsi, $S$ est un corps si nous pouvons démontrer la
première assertion. Supposons $S$ de dimension finie
sur $k$ et intègre. Choisissons $s\in S$ non nul.
Puisque $S$ est un anneau intègre, l'application de multiplication
$s : S \to S$ est surjective pour des raisons de
dimension. Il existe donc un élément $s_1 \in S$
tel que $ss_1 = 1$. Ainsi, $S$ est un corps.
\end{proof}

\begin{lemma}
\label{lemma-integral-no-inclusion}
Supposons que $R \to S$ soit entier.
Soient $\mathfrak q, \mathfrak q' \in \Spec(S)$
deux idéaux premiers distincts
ayant la même image dans $\Spec(R)$.
Alors ni $\mathfrak q \subset \mathfrak q'$
ni $\mathfrak q' \subset \mathfrak q$.
\end{lemma}

\begin{proof}
Soit $\mathfrak p \subset R$ l'image.
D'après la remarque \ref{remark-fundamental-diagram},
les idéaux premiers $\mathfrak q, \mathfrak q'$
correspondent à des idéaux de
$S \otimes_R \kappa(\mathfrak p)$.
Le lemme résulte donc du lemme \ref{lemma-integral-over-field}.
\end{proof}

\begin{lemma}
\label{lemma-finite-finite-fibres}
Supposons que $R \to S$ soit fini.
Alors les fibres de $\Spec(S) \to \Spec(R)$ sont finies.
\end{lemma}

\begin{proof}
D'après la discussion de la
remarque \ref{remark-fundamental-diagram},
les fibres sont les spectres des anneaux $S \otimes_R \kappa(\mathfrak p)$.
Comme $R \to S$ est fini, ces anneaux fibres sont finis sur
$\kappa(\mathfrak p)$, donc noethériens d'après le
lemme \ref{lemma-Noetherian-permanence}.
D'après le
lemme \ref{lemma-integral-no-inclusion},
tout idéal premier de $S \otimes_R \kappa(\mathfrak p)$ est
minimal. Ainsi, d'après le
lemme \ref{lemma-Noetherian-irreducible-components},
il n'y en a qu'un nombre fini.
\end{proof}

\begin{lemma}
\label{lemma-integral-going-up}
Soit $R \to S$ un morphisme d'anneaux tel que
$S$ soit entier sur $R$.
Soient $\mathfrak p \subset \mathfrak p' \subset R$
des idéaux premiers. Soit $\mathfrak q$ un idéal premier de $S$ qui s'envoie
sur $\mathfrak p$. Alors il existe un idéal premier $\mathfrak q'$
tel que $\mathfrak q \subset \mathfrak q'$
et qui s'envoie sur $\mathfrak p'$.
\end{lemma}

\begin{proof}
Nous pouvons remplacer $R$ par $R/\mathfrak p$ et $S$ par $S/\mathfrak q$.
Cela nous ramène à une extension entière d'anneaux
intègres $R \subset S$ et à un idéal premier $\mathfrak p' \subset R$.
Le lemme \ref{lemma-integral-overring-surjective} permet de conclure.
\end{proof}

\noindent
La propriété exprimée dans le lemme ci-dessus est appelée
la « propriété de montée » pour le morphisme d'anneaux $R \to S$,
voir la définition \ref{definition-going-up-down}.

\begin{lemma}
\label{lemma-finite-finitely-presented-extension}
Soit $R \to S$ un morphisme d'anneaux fini et de présentation finie.
Soit $M$ un $S$-module.
Alors $M$ est de présentation finie comme $R$-module si et seulement si
$M$ est de présentation finie comme $S$-module.
\end{lemma}

\begin{proof}
L'une des implications résulte du
lemme \ref{lemma-finitely-presented-over-subring}.
Pour voir l'autre, supposons que $M$ soit de présentation finie comme $S$-module.
Choisissons une présentation
$$
S^{\oplus m} \longrightarrow
S^{\oplus n} \longrightarrow
M \longrightarrow 0
$$
Comme $S$ est fini comme $R$-module, le noyau de
$S^{\oplus n} \to M$ est un $R$-module de type fini. Ainsi, le
lemme \ref{lemma-extension}
montre qu'il suffit de démontrer que $S$ est de présentation finie comme
$R$-module.

\medskip\noindent
Choisissons $y_1, \ldots, y_n \in S$ tels que $y_1, \ldots, y_n$ engendrent $S$
comme $R$-module. D'après le lemme \ref{lemma-characterize-integral-element},
chaque $y_i$ est entier sur $R$. Choisissons des polynômes unitaires
$P_i(x) \in R[x]$ tels que $P_i(y_i) = 0$. Considérons l'anneau
$$
S' = R[x_1, \ldots, x_n]/(P_1(x_1), \ldots, P_n(x_n))
$$
Nous voyons alors que $S$ est de présentation finie comme $S'$-algèbre
d'après le lemme \ref{lemma-compose-finite-type}. Puisque $S' \to S$ est surjectif,
le noyau $J = \Ker(S' \to S)$ est engendré par un nombre fini d'éléments comme idéal d'après le
lemme \ref{lemma-finite-presentation-independent}. Ainsi, $J$ est un
$S'$-module de type fini (cela découle immédiatement des définitions).
Par conséquent, $S = \Coker(J \to S')$ est de présentation finie comme $S'$-module
d'après le lemme \ref{lemma-extension}.
Ainsi, en raisonnant comme dans le premier paragraphe, il suffit de montrer que $S'$ est
de présentation finie comme $R$-module. En fait, $S'$ est libre comme
$R$-module, de base les monômes $x_1^{e_1} \ldots x_n^{e_n}$
pour $0 \leq e_i < \deg(P_i)$. En effet, écrivons $R \to S'$ comme la composée
$$
R \to R[x_1]/(P_1(x_1)) \to R[x_1, x_2]/(P_1(x_1), P_2(x_2)) \to
\ldots \to S'
$$
Cela montre que le $i$-ème anneau de cette suite est libre comme module sur le
$(i - 1)$-ième, de base $1, x_i, \ldots, x_i^{\deg(P_i) - 1}$. Le résultat
s'en déduit aisément par récurrence. Certains détails sont omis.
\end{proof}

\begin{lemma}
\label{lemma-silly-normal}
Soit $R$ un anneau. Soient $x, y \in R$ des non-diviseurs de zéro.
Soit $R[x/y] \subset R_{xy}$ la $R$-sous-algèbre engendrée
par $x/y$, et de même pour les sous-algèbres $R[y/x]$ et $R[x/y, y/x]$.
Si $R$ est intégralement clos dans $R_x$ ou $R_y$, alors la suite
$$
0 \to R \xrightarrow{(-1, 1)} R[x/y] \oplus R[y/x] \xrightarrow{(1, 1)}
R[x/y, y/x] \to 0
$$
est une suite exacte courte de $R$-modules.
\end{lemma}

\begin{proof}
Puisque $x/y \cdot y/x = 1$, il est clair que l'application
$R[x/y] \oplus R[y/x] \to R[x/y, y/x]$ est surjective.
Soit $\alpha \in R[x/y] \cap R[y/x]$. Pour montrer l'exactitude au milieu,
nous devons prouver que $\alpha \in R$. Par hypothèse, nous pouvons écrire
$$
\alpha = a_0 + a_1 x/y + \ldots + a_n (x/y)^n =
b_0 + b_1 y/x + \ldots + b_m(y/x)^m
$$
pour certains $n, m \geq 0$ et $a_i, b_j \in R$.
Choisissons un $N > \max(n, m)$.
Considérons le $R$-sous-module de type fini $M$ de $R_{xy}$ engendré par les éléments
$$
(x/y)^N, (x/y)^{N - 1}, \ldots, x/y, 1, y/x, \ldots, (y/x)^{N - 1}, (y/x)^N
$$
Nous affirmons que $\alpha M \subset M$. En effet, il est clair que
$(x/y)^i (b_0 + b_1 y/x + \ldots + b_m(y/x)^m) \in M$ pour
$0 \leq i \leq N$ et que
$(y/x)^i (a_0 + a_1 x/y + \ldots + a_n(x/y)^n) \in M$ pour
$0 \leq i \leq N$. Ainsi, $\alpha$ est entier sur $R$ d'après le
lemme \ref{lemma-characterize-integral-element}. Remarquons que
$\alpha \in R_x$ ; si $R$ est intégralement clos dans $R_x$,
alors $\alpha \in R$, comme voulu.
\end{proof}







\section{Anneaux normaux}
\label{section-normal-rings}

\noindent
Introduisons d'abord la notion d'anneau intègre normal, puis
la notion (très générale) d'anneau normal.

\begin{definition}
\label{definition-domain-normal}
Un anneau intègre $R$ est dit {\it normal} s'il est intégralement
clos dans son corps des fractions.
\end{definition}

\begin{lemma}
\label{lemma-integral-closure-in-normal}
Soit $R \to S$ un morphisme d'anneaux.
Si $S$ est un anneau intègre normal, alors la clôture intégrale de $R$
dans $S$ est un anneau intègre normal.
\end{lemma}

\begin{proof}
Omis.
\end{proof}

\noindent
La notion suivante est parfois utile dans l'étude
de la normalité.

\begin{definition}
\label{definition-almost-integral}
Soit $R$ un anneau intègre.
\begin{enumerate}
\item Un élément $g$ du corps des
fractions de $R$ est dit {\it presque entier sur $R$}
s'il existe un élément $r \in R$, $r\not = 0$,
tel que $rg^n \in R$ pour tout $n \geq 0$.
\item L'anneau intègre $R$ est dit {\it complètement normal} si tout
élément presque entier du corps des fractions de $R$
appartient à $R$.
\end{enumerate}
\end{definition}

\noindent
Le lemme suivant montre qu'un anneau intègre noethérien est normal
si et seulement s'il est complètement normal.

\begin{lemma}
\label{lemma-almost-integral}
Soit $R$ un anneau intègre de corps des fractions $K$.
Si $u, v \in K$ sont presque entiers sur $R$, alors
$u + v$ et $uv$ le sont aussi. Tout élément $g \in K$ qui est entier sur $R$
est presque entier sur $R$. Si $R$ est noethérien,
alors la réciproque est également vraie.
\end{lemma}

\begin{proof}
Si $ru^n \in R$ pour tout $n \geq 0$ et
$v^nr' \in R$ pour tout $n \geq 0$, alors
$(uv)^nrr'$ et $(u + v)^nrr'$ appartiennent à $R$ pour
tout $n \geq 0$. D'où la première assertion.
Supposons que $g \in K$ soit entier sur $R$.
Dans ce cas, il existe un $d > 0$ tel que
l'anneau $R[g]$ soit engendré par $1, g, \ldots, g^d$ comme $R$-module.
Soit $r \in R$ un dénominateur commun des éléments
$1, g, \ldots, g^d \in K$. Il s'ensuit que $rR[g] \subset R$,
et donc que $g$ est presque entier sur $R$.

\medskip\noindent
Supposons que $R$ soit noethérien et que $g \in K$ soit presque entier sur $R$.
Soit $r \in R$, $r\not = 0$, comme dans la définition.
Alors $R[g] \subset \frac{1}{r}R$ comme $R$-module.
Puisque $R$ est noethérien, cela implique que $R[g]$ est
fini sur $R$. Ainsi, $g$ est entier sur $R$, voir le
lemme \ref{lemma-finite-is-integral}.
\end{proof}

\begin{lemma}
\label{lemma-localize-normal-domain}
Toute localisation d'un anneau intègre normal est normale.
\end{lemma}

\begin{proof}
Soit $R$ un anneau intègre normal, et soit $S \subset R$ une
partie multiplicative. Supposons que $g$ soit un élément
du corps des fractions de $R$ qui soit entier sur $S^{-1}R$.
Soit $P = x^d + \sum_{j < d} a_j x^j$ un polynôme
avec $a_i \in S^{-1}R$ tel que $P(g) = 0$.
Choisissons $s \in S$ tel que $sa_i \in R$ pour tout $i$.
Alors $sg$ satisfait le polynôme unitaire
$x^d + \sum_{j < d} s^{d-j}a_j x^j$, dont les coefficients
$s^{d-j}a_j$ appartiennent à $R$. Ainsi, $sg \in R$ puisque $R$ est normal.
Donc $g \in S^{-1}R$.
\end{proof}

\begin{lemma}
\label{lemma-PID-normal}
Un anneau principal est normal.
\end{lemma}

\begin{proof}
Soit $R$ un anneau principal.
Soit $g = a/b$ un élément du corps des fractions
de $R$ qui soit entier sur $R$. Puisque $R$ est un anneau principal,
nous pouvons diviser $a$ et $b$ par un facteur commun
et supposer que $(a, b) = R$. Dans ce cas, toute équation
$(a/b)^n + r_{n-1} (a/b)^{n-1} + \ldots + r_0 = 0$
avec $r_i \in R$ impliquerait que $a^n \in (b)$. Cela
contredit $(a, b) = R$, sauf si $b$ est une unité de $R$.
\end{proof}

\begin{lemma}
\label{lemma-prepare-polynomial-ring-normal}
Soit $R$ un anneau intègre de corps des fractions $K$.
Supposons que $f = \sum \alpha_i x^i$ soit un
élément de $K[x]$.
\begin{enumerate}
\item Si $f$ est entier sur $R[x]$,
alors tous les $\alpha_i$ sont entiers sur $R$, et
\item Si $f$ est presque entier sur $R[x]$,
alors tous les $\alpha_i$ sont presque entiers sur $R$.
\end{enumerate}
\end{lemma}

\begin{proof}
Démontrons d'abord la seconde assertion.
Écrivons $f = \alpha_0 + \alpha_1 x + \ldots + \alpha_r x^r$
avec $\alpha_r \not = 0$.
Par hypothèse, il existe $h = b_0 + b_1 x + \ldots + b_s x^s \in R[x]$,
$b_s \not = 0$, tel que $f^n h \in R[x]$ pour tout
$n \geq 0$. Cela implique que $b_s \alpha_r^n \in R$
pour tout $n \geq 0$. Ainsi, $\alpha_r$ est presque
entier sur $R$. Puisque l'ensemble des éléments presque
entiers forme un sous-anneau (lemme \ref{lemma-almost-integral}), nous en déduisons que
$f - \alpha_r x^r = \alpha_0 + \alpha_1 x + \ldots + \alpha_{r - 1} x^{r - 1}$
est presque entier sur $R[x]$. Une récurrence sur $r$ permet de conclure.

\medskip\noindent
Pour démontrer la première assertion, nous utiliserons une réduction
noethérienne absolue. Écrivons $\alpha_i = a_i / b_i$ et
soit $P(t) = t^d + \sum_{j < d} f_j t^j$ un polynôme
à coefficients $f_j \in R[x]$ tel que $P(f) = 0$.
Soit $f_j = \sum f_{ji}x^i$. Considérons le sous-anneau
$R_0 \subset R$ engendré par la liste finie des éléments
$a_i, b_i, f_{ji}$ de $R$. C'est un anneau intègre ; soit
$K_0$ son corps des fractions. Puisque $R_0$ est une
$\mathbf{Z}$-algèbre de type fini, il est noethérien, voir le
lemme \ref{lemma-obvious-Noetherian}. Il est toujours
vrai que $f \in K_0[x]$ est entier sur $R_0[x]$,
car toutes les identités dans $R$
entre les éléments $a_i, b_i, f_{ji}$ sont également valables dans $R_0$.
D'après le lemme \ref{lemma-almost-integral}, l'élément
$f$ est presque entier sur $R_0[x]$. D'après la seconde assertion du
lemme, les éléments $\alpha_i$ sont presque entiers
sur $R_0$. Et puisque $R_0$ est noethérien, ils sont
entiers sur $R_0$, voir le lemme \ref{lemma-almost-integral}.
Ils sont alors bien sûr entiers sur $R$.
\end{proof}

\begin{lemma}
\label{lemma-polynomial-domain-normal}
Soit $R$ un anneau intègre normal.
Alors $R[x]$ est un anneau intègre normal.
\end{lemma}

\begin{proof}
Le résultat est vrai si $R$ est un corps $K$, car
$K[x]$ est un anneau euclidien, donc un anneau principal,
et donc normal d'après le lemme \ref{lemma-PID-normal}.
Soit $g$ un élément du corps des fractions de
$R[x]$ qui soit entier sur $R[x]$. Puisque $g$
est entier sur $K[x]$, où $K$ est le corps des
fractions de $R$, nous pouvons écrire $g = \alpha_d x^d + \alpha_{d-1}x^{d-1} +
\ldots + \alpha_0$ avec $\alpha_i \in K$.
D'après le lemme \ref{lemma-prepare-polynomial-ring-normal},
les éléments $\alpha_i$ sont entiers sur $R$ et
appartiennent donc à $R$.
\end{proof}

\begin{lemma}
\label{lemma-power-series-over-Noetherian-normal-domain}
Soit $R$ un anneau intègre noethérien normal. Alors $R[[x]]$ est
un anneau intègre noethérien normal.
\end{lemma}

\begin{proof}
L'anneau de séries formelles est noethérien d'après le
lemme \ref{lemma-Noetherian-power-series}.
Soient $f, g \in R[[x]]$ des éléments non nuls tels que
$w = f/g$ soit entier sur $R[[x]]$.
Soit $K$ le corps des fractions de $R$. Puisque l'anneau des séries de Laurent
$K((x)) = K[[x]][1/x]$ est un corps, nous pouvons écrire
$w = a_n x^n + a_{n + 1} x^{n + 1} + \ldots$
pour certains $n \in \mathbf{Z}$, $a_i \in K$, et $a_n \not = 0$.
D'après le lemme \ref{lemma-almost-integral}, nous voyons qu'il existe un
élément non nul $h = b_m x^m + b_{m + 1} x^{m + 1} + \ldots$
de $R[[x]]$, avec $b_m \not = 0$, tel que
$w^e h \in R[[x]]$ pour tout $e \geq 1$. Nous concluons que $n \geq 0$ et que
$b_m a_n^e \in R$ pour tout $e \geq 1$.
Puisque $R$ est noethérien, le même lemme implique que $a_n \in R$.
Maintenant, si $a_n, a_{n + 1}, \ldots, a_{N - 1} \in R$,
nous pouvons appliquer le même argument à
$w - a_n x^n - \ldots - a_{N - 1} x^{N - 1} = a_N x^N + \ldots$.
Nous voyons ainsi que tous les $a_i \in R$, ce qui démontre le lemme.
\end{proof}

\begin{lemma}
\label{lemma-normality-is-local}
Soit $R$ un anneau intègre. Les propriétés suivantes sont équivalentes :
\begin{enumerate}
\item L'anneau intègre $R$ est normal,
\item pour tout idéal premier $\mathfrak p \subset R$, l'anneau local
$R_{\mathfrak p}$ est un anneau intègre normal, et
\item pour tout idéal maximal $\mathfrak m$, l'anneau $R_{\mathfrak m}$
est un anneau intègre normal.
\end{enumerate}
\end{lemma}

\begin{proof}
Nous déduisons (1) $\Rightarrow$ (2) du lemme \ref{lemma-localize-normal-domain}.
L'implication (2) $\Rightarrow$ (3) est immédiate. L'implication
(3) $\Rightarrow$ (1) résulte du fait que, pour tout anneau intègre $R$, nous avons
$$
R = \bigcap\nolimits_{\mathfrak m} R_{\mathfrak m}
$$
dans le corps des fractions de $R$. En effet, si $g$ est un élément du
membre de droite, alors l'idéal $I = \{x \in R \mid xg \in R\}$
n'est contenu dans aucun idéal maximal $\mathfrak m$, d'où $I = R$.
\end{proof}

\noindent
Le lemme \ref{lemma-normality-is-local} montre que la définition suivante
est compatible avec la définition \ref{definition-domain-normal}. (Il s'agit de la
définition d'EGA -- voir \cite[IV, 5.13.5 et 0, 4.1.4]{EGA}.)

\begin{definition}
\label{definition-ring-normal}
Un anneau $R$ est dit {\it normal} si, pour tout idéal premier
$\mathfrak p \subset R$, la localisation $R_{\mathfrak p}$ est
un anneau intègre normal (voir la définition \ref{definition-domain-normal}).
\end{definition}

\noindent
Remarquons qu'un anneau normal est réduit, puisque $R$ est un sous-anneau du produit
de ses localisations en tous les idéaux premiers (voir par exemple le
lemme \ref{lemma-characterize-zero-local}).

\begin{lemma}
\label{lemma-normal-ring-integrally-closed}
Un anneau normal est intégralement clos dans son anneau total des fractions.
\end{lemma}

\begin{proof}
Soit $R$ un anneau normal. Soit $x \in Q(R)$ un élément de l'anneau total
des fractions de $R$ entier sur $R$. Posons $I = \{f \in R, fx \in R\}$. Soit
$\mathfrak p \subset R$ un idéal premier. Comme $R \to R_{\mathfrak p}$ est
plat, nous voyons que $R_{\mathfrak p} \subset Q(R) \otimes_R R_{\mathfrak p}$. Comme
$R_{\mathfrak p}$ est un anneau intègre normal, nous voyons que $x \otimes 1$ est un élément de
$R_{\mathfrak p}$. Nous pouvons donc trouver $a, f \in R$, $f \not \in \mathfrak p$,
tels que $x \otimes 1 = a \otimes 1/f$. Cela signifie que $fx - a$ s'envoie sur
zéro dans $Q(R) \otimes_R R_{\mathfrak p} = Q(R)_{\mathfrak p}$, ce qui
signifie à son tour qu'il existe un $f' \in R$, $f' \not \in \mathfrak p$,
tel que $f'fx = f'a$ dans $R$. Autrement dit, $ff' \in I$. Ainsi, $I$
est un idéal qui n'est contenu dans aucun idéal premier de $R$, c'est-à-dire
$I = R$ et $x \in R$.
\end{proof}

\begin{lemma}
\label{lemma-localization-normal-ring}
Une localisation d'un anneau normal est un anneau normal.
\end{lemma}

\begin{proof}
Omis.
\end{proof}

\begin{lemma}
\label{lemma-polynomial-ring-normal}
Soit $R$ un anneau normal. Alors $R[x]$ est un anneau normal.
\end{lemma}

\begin{proof}
Soit $\mathfrak q$ un idéal premier de $R[x]$. Posons $\mathfrak p = R \cap \mathfrak q$.
Nous voyons alors que $R_{\mathfrak p}[x]$ est un anneau intègre normal d'après le
lemme \ref{lemma-polynomial-domain-normal}.
Ainsi, $(R[x])_{\mathfrak q}$ est un anneau intègre normal d'après le
lemme \ref{lemma-localize-normal-domain}.
\end{proof}

\begin{lemma}
\label{lemma-finite-product-normal}
Un produit fini d'anneaux normaux est normal.
\end{lemma}

\begin{proof}
Il suffit de montrer que le produit de deux anneaux normaux, disons $R$ et $S$, est
normal. D'après le lemme \ref{lemma-disjoint-decomposition}, les idéaux premiers de
$R\times S$ sont de la forme $\mathfrak{p}\times S$ et $R\times
\mathfrak{q}$, où $\mathfrak{p}$ et $\mathfrak{q}$ sont des idéaux premiers respectivement de $R$
et de $S$. La localisation donne
$(R\times S)_{\mathfrak{p}\times S}=R_{\mathfrak{p}}$, qui est un anneau intègre normal
par hypothèse. De même pour $S$.
\end{proof}

\begin{lemma}
\label{lemma-characterize-reduced-ring-normal}
Soit $R$ un anneau. Supposons que $R$ soit réduit et possède un nombre fini
d'idéaux premiers minimaux. Alors les propriétés suivantes sont équivalentes :
\begin{enumerate}
\item $R$ est un anneau normal,
\item $R$ est intégralement clos dans son anneau total des fractions, et
\item $R$ est un produit fini d'anneaux intègres normaux.
\end{enumerate}
\end{lemma}

\begin{proof}
Les implications (1) $\Rightarrow$ (2) et
(3) $\Rightarrow$ (1) sont vraies en général,
voir les lemmes \ref{lemma-normal-ring-integrally-closed} et
\ref{lemma-finite-product-normal}.

\medskip\noindent
Soient $\mathfrak p_1, \ldots, \mathfrak p_n$ les idéaux premiers minimaux de $R$.
D'après les lemmes \ref{lemma-reduced-ring-sub-product-fields} et
\ref{lemma-total-ring-fractions-no-embedded-points}, nous avons
$Q(R) = R_{\mathfrak p_1} \times \ldots \times R_{\mathfrak p_n}$, et,
d'après le lemme \ref{lemma-minimal-prime-reduced-ring}, chaque facteur est un corps.
Notons $e_i = (0, \ldots, 0, 1, 0, \ldots, 0)$ le $i$-ème idempotent
de $Q(R)$.

\medskip\noindent
Si $R$ est intégralement clos dans $Q(R)$, alors il contient en particulier
les idempotents $e_i$, et nous voyons que $R$ est un produit de $n$
anneaux intègres (voir les sections \ref{section-connected-components} et
\ref{section-tilde-module-sheaf}). Chaque facteur est de la forme
$R/\mathfrak p_i$, de corps des fractions $R_{\mathfrak p_i}$.
D'après le lemme \ref{lemma-finite-product-integral-closure}, chaque morphisme
$R/\mathfrak p_i \to R_{\mathfrak p_i}$ est intégralement clos.
Ainsi, $R$ est un produit fini d'anneaux intègres normaux.
\end{proof}

\begin{lemma}
\label{lemma-colimit-normal-ring}
Soit $(R_i, \varphi_{ii'})$ un système inductif
(Catégories, définition \ref{definition-directed-system})
d'anneaux. Si chaque $R_i$ est un anneau normal, alors
$R = \colim_i R_i$ l'est aussi.
\end{lemma}

\begin{proof}
Soit $\mathfrak p \subset R$ un idéal premier.
Posons $\mathfrak p_i = R_i \cap \mathfrak p$ (abus de notation usuel).
Nous voyons alors que
$R_{\mathfrak p} = \colim_i (R_i)_{\mathfrak p_i}$.
Puisque chaque $(R_i)_{\mathfrak p_i}$ est un anneau intègre normal, nous
nous ramenons à démontrer l'énoncé du lemme pour des anneaux intègres
normaux. Si $a, b \in R$ et si $a/b$ satisfait un polynôme unitaire
$P(T) \in R[T]$, alors nous pouvons trouver un $i \in I$ (suffisamment grand)
tel que $a, b$ proviennent d'éléments $a_i, b_i$ sur $R_i$, que $P$ provienne d'un
polynôme unitaire $P_i\in R_i[T]$ et que $P_i(a_i/b_i)=0$. Puisque $R_i$ est normal, nous
voyons que $a_i/b_i \in R_i$, et donc aussi que $a/b \in R$.
\end{proof}







\section{Descente pour une extension entière au-dessus d'un anneau normal}
\label{section-going-down-integral-over-normal}

\noindent
Commençons par manipuler un peu la notion d'éléments
entiers sur un idéal, puis démontrons le théorème mentionné
dans le titre de la section.

\begin{definition}
\label{definition-integral-over-ideal}
Soit $\varphi : R \to S$ un morphisme d'anneaux.
Soit $I \subset R$ un idéal.
Nous disons qu'un élément $g \in S$ est
{\it entier sur $I$} s'il
existe un polynôme unitaire
$P = x^d + \sum_{j < d} a_j x^j$
à coefficients $a_j \in I^{d-j}$ tel
que $P^\varphi(g) = 0$ dans $S$.
\end{definition}

\noindent
Ceci est surtout utilisé lorsque $\varphi = \text{id}_R : R \to R$.
Dans ce cas, l'ensemble $I'$ des éléments entiers sur $I$ est appelé
la {\it clôture intégrale de $I$}. Nous verrons que $I'$ est
un idéal de $R$ (et, bien sûr, $I \subset I'$).

\begin{lemma}
\label{lemma-characterize-integral-ideal}
Soit $\varphi : R \to S$ un morphisme d'anneaux.
Soit $I \subset R$ un idéal.
Soit $A = \sum I^nt^n \subset R[t]$ le
sous-anneau de l'anneau de polynômes
engendré par $R \oplus It \subset R[t]$.
Un élément $s \in S$ est entier sur $I$ si
et seulement si l'élément $st \in S[t]$
est entier sur $A$.
\end{lemma}

\begin{proof}
Supposons que $st$ soit entier sur $A$.
Soit $P = x^d + \sum_{j < d} a_j x^j$
un polynôme unitaire à coefficients dans $A$
tel que $P^\varphi(st) = 0$. Soit $a_j' \in A$
la partie de degré $d-j$ de $a_j$, autrement
dit $a_j' = a_j'' t^{d-j}$ avec $a_j'' \in I^{d-j}$.
Pour des raisons de degré, nous avons encore
$(st)^d + \sum_{j < d} \varphi(a_j'') t^{d-j} (st)^j = 0$.
Ainsi, $s^d + \sum_{j < d} \varphi(a_j'') s^j = 0$,
et nous voyons que $s$ est entier sur $I$.

\medskip\noindent
Supposons que $s$ soit entier sur $I$.
Écrivons $P = x^d + \sum_{j < d} a_j x^j$
avec $a_j \in I^{d-j}$. Nous trouvons alors immédiatement un
polynôme $Q = x^d + \sum_{j < d} (a_j t^{d-j}) x^j$
à coefficients dans $A$ qui démontre que
$st$ est entier sur $A$.
\end{proof}

\begin{lemma}
\label{lemma-integral-over-ideal-is-submodule}
Soit $\varphi : R \to S$ un morphisme d'anneaux.
Soit $I \subset R$ un idéal.
L'ensemble des éléments de $S$ qui sont entiers
sur $I$ forme un $R$-sous-module de $S$.
En outre, si $s \in S$ est entier sur
$R$ et si $s'$ est entier sur $I$, alors
$ss'$ est entier sur $I$.
\end{lemma}

\begin{proof}
Nous utiliserons le lemme \ref{lemma-integral-closure-is-ring} sans
autre mention. La stabilité par addition résulte clairement de la
caractérisation du lemme \ref{lemma-characterize-integral-ideal},
dont nous adoptons les notations.
Tout élément $s \in S$ qui est entier sur
$R$ correspond à l'élément de degré $0$, $s$, de $S[t]$,
qui est entier sur $A$ (car $R \subset A$).
Nous voyons donc que la multiplication par $s$ sur $S[t]$
préserve la propriété d'être entier sur $A$.
\end{proof}

\begin{lemma}
\label{lemma-integral-integral-over-ideal}
Supposons que $\varphi : R \to S$ soit entier.
Supposons que $I \subset R$ soit un idéal.
Alors tout élément de $IS$ est entier sur $I$.
\end{lemma}

\begin{proof}
C'est immédiat d'après le lemme \ref{lemma-integral-over-ideal-is-submodule}.
\end{proof}

\begin{lemma}
\label{lemma-polynomials-divide}
Soit $K$ un corps. Soient $n, m \in \mathbf{N}$ et
$a_0, \ldots, a_{n - 1}, b_0, \ldots, b_{m - 1} \in K$.
Si le polynôme $x^n + a_{n - 1}x^{n - 1} + \ldots + a_0$
divise le polynôme $x^m + b_{m - 1} x^{m - 1} + \ldots + b_0$
dans $K[x]$, alors :
\begin{enumerate}
\item $a_0, \ldots, a_{n - 1}$ sont entiers sur tout sous-anneau
$R_0$ de $K$ contenant les éléments $b_0, \ldots, b_{m - 1}$, et
\item chaque $a_i$ appartient à $\sqrt{(b_0, \ldots, b_{m-1})R}$
pour tout sous-anneau $R \subset K$ contenant les éléments
$a_0, \ldots, a_{n - 1}, b_0, \ldots, b_{m - 1}$.
\end{enumerate}
\end{lemma}

\begin{proof}
Soit $L/K$ une extension de corps telle que nous puissions écrire
$x^m + b_{m - 1} x^{m - 1} + \ldots + b_0 =
\prod_{i = 1}^m (x - \beta_i)$ avec $\beta_i \in L$.
Voir Corps, section \ref{fields-section-splitting-fieds}.
Chaque $\beta_i$ est entier sur $R_0$.
Puisque chaque $a_i$ est un polynôme homogène en $\beta_1, \ldots, \beta_m$,
nous en déduisons la même propriété pour les $a_i$
(utiliser le lemme \ref{lemma-integral-closure-is-ring}). Cela démontre (1).

\medskip\noindent
Soit $R$ comme dans (2). Choisissons $c_0, \ldots, c_{m - n - 1} \in K$ tels que
$$
\begin{matrix}
x^m + b_{m - 1} x^{m - 1} + \ldots + b_0 =  \\
(x^n + a_{n - 1}x^{n - 1} + \ldots + a_0)
(x^{m - n} + c_{m - n - 1}x^{m - n - 1}+ \ldots + c_0).
\end{matrix}
$$
Cette équation implique
\begin{align*}
c_{m - n - 1} & = b_{m - 1} - a_{n - 1}, \\
c_{m - n - 2} & = b_{m - 2} - a_{n - 2} - a_{n - 1}c_{m - n - 1}, \\
\ldots
\end{align*}
Ainsi, $c_j \in R$ pour tout $j$. En quotientant par le radical
$\sqrt{(b_0, \ldots, b_{m - 1})}$, nous obtenons un anneau réduit $\overline{R}$.
Nous devons montrer que les images $\overline{a}_i \in \overline{R}$
sont nulles. Et, dans
$\overline{R}[x]$, nous avons la relation
$$
\begin{matrix}
x^m = x^m + \overline{b}_{m - 1} x^{m - 1} + \ldots + \overline{b}_0 = \\
(x^n + \overline{a}_{n - 1}x^{n - 1} + \ldots + \overline{a}_0)
(x^{m - n} + \overline{c}_{m - n - 1}x^{m - n - 1}+ \ldots + \overline{c}_0).
\end{matrix}
$$
Il est facile de voir que cela implique $\overline{a}_i = 0$ pour tout $i$. En effet,
d'après le lemme \ref{lemma-minimal-prime-reduced-ring}, la localisation de
$\overline{R}$ en un idéal premier minimal $\mathfrak{p}$ est un corps et
$\overline{R}_{\mathfrak p}[x]$ est un anneau factoriel. Ainsi,
$f = x^n + \sum \overline{a}_i x^i$
est associé à $x^n$ et, puisque $f$ est unitaire, $f = x^n$
dans $\overline{R}_{\mathfrak p}[x]$.
Il existe alors un $s \in \overline{R}$, $s \not\in \mathfrak p$,
tel que $s(f - x^n) = 0$. Par conséquent, tous les $\overline{a}_i$ appartiennent
à $\mathfrak p$, et nous concluons par le
lemme \ref{lemma-reduced-ring-sub-product-fields}.
\end{proof}

\begin{lemma}
\label{lemma-minimal-polynomial-normal-domain}
Soit $R \subset S$ une inclusion d'anneaux intègres.
Supposons que $R$ soit normal. Soit $g \in S$ entier
sur $R$. Alors le polynôme minimal de $g$
est à coefficients dans $R$.
\end{lemma}

\begin{proof}
Soit $P = x^m + b_{m-1} x^{m-1} + \ldots + b_0$
un polynôme à coefficients dans $R$
tel que $P(g) = 0$. Soit $Q = x^n + a_{n-1}x^{n-1} + \ldots + a_0$
le polynôme minimal de $g$ sur le corps des fractions
$K$ de $R$. Alors $Q$ divise $P$ dans $K[x]$. D'après le lemme
\ref{lemma-polynomials-divide}, nous voyons que les $a_i$ sont
entiers sur $R$. Puisque $R$ est normal, cela
signifie qu'ils appartiennent à $R$.
\end{proof}

\begin{proposition}
\label{proposition-going-down-normal-integral}
Soit $R \subset S$ une inclusion d'anneaux intègres.
Supposons que $R$ soit normal et que $S$ soit entier sur $R$.
Soient $\mathfrak p \subset \mathfrak p' \subset R$
des idéaux premiers. Soit $\mathfrak q'$ un idéal premier de $S$
tel que $\mathfrak p' = R \cap \mathfrak q'$.
Alors il existe un idéal premier $\mathfrak q$
tel que $\mathfrak q \subset \mathfrak q'$
et $\mathfrak p = R \cap \mathfrak q$. Autrement dit,
la propriété de descente est satisfaite pour $R \to S$, voir la
définition \ref{definition-going-up-down}.
\end{proposition}

\begin{proof}
Soient $\mathfrak p$, $\mathfrak p'$ et $\mathfrak q'$
comme dans l'énoncé. Nous devons montrer qu'il existe un idéal premier
$\mathfrak q$, avec $\mathfrak q \subset \mathfrak q'$ et
$R \cap \mathfrak q = \mathfrak p$. Cela revient
à trouver un idéal premier de
$S_{\mathfrak q'}$ qui s'envoie sur $\mathfrak p$.
D'après le lemme \ref{lemma-in-image}, nous devons montrer
que $\mathfrak p S_{\mathfrak q'} \cap R
= \mathfrak p$. Choisissons $z \in \mathfrak p S_{\mathfrak q'} \cap R$.
Nous pouvons écrire $z = y/g$ avec $y \in \mathfrak pS$ et
$g \in S$, $g \not\in \mathfrak q'$. Autrement dit,
nous avons $zg = y$.

\medskip\noindent
D'après le lemme \ref{lemma-integral-integral-over-ideal},
il existe un polynôme unitaire
$P = x^m + b_{m-1} x^{m-1} + \ldots + b_0$
avec $b_i \in \mathfrak p$ tel que $P(y) = 0$.

\medskip\noindent
D'après le lemme \ref{lemma-minimal-polynomial-normal-domain},
le polynôme minimal de $g$ sur $K$ est à coefficients
dans $R$. Écrivons-le $Q = x^n + a_{n-1} x^{n-1} + \ldots
+ a_0$. Remarquons que les $a_i$, $i = n-1, \ldots, 0$,
n'appartiennent pas tous à $\mathfrak p$, car cela impliquerait que
$g^n = \sum_{j < n} a_j g^j \in \mathfrak pS
\subset \mathfrak p'S
\subset \mathfrak q'$,
ce qui est une contradiction.

\medskip\noindent
Puisque $y = zg$, ce qui précède montre immédiatement
que $Q' = x^n + za_{n-1} x^{n-1} + \ldots + z^{n}a_0$
est le polynôme minimal de $y$. Ainsi,
$Q'$ divise $P$ et, d'après le lemme \ref{lemma-polynomials-divide},
nous voyons que $z^ja_{n - j} \in \sqrt{(b_0, \ldots, b_{m-1})}
\subset \mathfrak p$, $j =  1, \ldots, n$.
Puisque les $a_i$, $i = n-1, \ldots, 0$,
n'appartiennent pas tous à $\mathfrak p$, nous concluons que $z \in \mathfrak p$,
comme voulu.
\end{proof}














\section{Modules plats et morphismes d'anneaux plats}
\label{section-flat}

\noindent
Un résultat souvent utilisé est le suivant : si $M = \colim_{i\in \mathcal{I}} M_i$
est une colimite de $R$-modules et si $N$ est un $R$-module, alors
$$
M \otimes N
=
\colim_{i\in \mathcal{I}} M_i \otimes_R N,
$$
voir le lemme \ref{lemma-tensor-products-commute-with-limits}.
On exprime habituellement cette propriété en disant
que {\it $\otimes$ commute aux colimites}.
Un autre résultat souvent utilisé est le suivant : si $0 \to N_1 \to N_2 \to N_3 \to 0$
est une suite exacte et si $M$ est un $R$-module quelconque, alors
$$
M \otimes_R N_1
\to
M \otimes_R N_2
\to
M \otimes_R N_3
\to
0
$$
est encore exacte, voir le lemme \ref{lemma-tensor-product-exact}.
Ces deux propriétés nous disent que le foncteur
$N \mapsto M \otimes_R N$ {\it est exact à droite}.
Voir Catégories, section \ref{categories-section-exact-functor},
et Homologie, section \ref{homology-section-functors}.
Un $R$-module $M$ est plat si $N \mapsto N \otimes_R M$ est aussi exact à gauche,
c'est-à-dire s'il est exact. Voici la définition précise.

\begin{definition}
\label{definition-flat}
Soit $R$ un anneau.
\begin{enumerate}
\item Un $R$-module $M$ est dit {\it plat} si, chaque fois que
$N_1 \to N_2 \to N_3$ est une suite exacte de $R$-modules,
la suite $M \otimes_R N_1 \to M \otimes_R N_2 \to M \otimes_R N_3$
est également exacte.
\item Un $R$-module $M$ est dit {\it fidèlement plat} si le
complexe de $R$-modules
$N_1 \to N_2 \to N_3$ est exact si et seulement si
la suite $M \otimes_R N_1 \to M \otimes_R N_2 \to M \otimes_R N_3$
est exacte.
\item Un morphisme d'anneaux $R \to S$ est dit {\it plat} si
$S$ est plat comme $R$-module.
\item Un morphisme d'anneaux $R \to S$ est dit {\it fidèlement plat} si
$S$ est fidèlement plat comme $R$-module.
\end{enumerate}
\end{definition}

\noindent
Voici un exemple d'utilisation de la condition de platitude.

\begin{lemma}
\label{lemma-flat-intersect-ideals}
Soit $R$ un anneau. Soient $I, J \subset R$ des idéaux. Soit $M$ un
$R$-module plat. Alors $IM \cap JM = (I \cap J)M$.
\end{lemma}

\begin{proof}
Considérons la suite exacte $0 \to I \cap J \to R \to R/I \oplus R/J$.
En tensorisant par le module plat $M$, nous obtenons une suite exacte
$$
0 \to (I \cap J) \otimes_R M \to M \to M/IM \oplus M/JM
$$
Puisque le noyau de $M \to M/IM \oplus M/JM$ est égal à
$IM \cap JM$, nous concluons.
\end{proof}

\begin{lemma}
\label{lemma-colimit-flat}
Soit $R$ un anneau. Soit $\{M_i, \varphi_{ii'}\}$ un système filtrant de
$R$-modules plats. Alors $\colim_i M_i$ est un $R$-module plat.
\end{lemma}

\begin{proof}
Cela résulte du fait que $\otimes$ commute aux colimites et que
les colimites filtrantes sont exactes, voir le
lemme \ref{lemma-directed-colimit-exact}.
\end{proof}

\begin{lemma}
\label{lemma-composition-flat}
Une composée de morphismes d'anneaux (fidèlement) plats est
(fidèlement) plate.
Si $R \to R'$ est (fidèlement) plat et si $M'$ est un
$R'$-module (fidèlement) plat, alors $M'$ est un
$R$-module (fidèlement) plat.
\end{lemma}

\begin{proof}
La première assertion du lemme est un cas particulier de la
seconde, de sorte qu'il suffit clairement de démontrer cette dernière. Soit
$R \to R'$ un morphisme d'anneaux plat et soit $M'$ un $R'$-module plat.
Nous devons démontrer que $M'$ est un $R$-module plat. Soit
$N_1 \to N_2 \to N_3$ un complexe exact de $R$-modules.
Alors le complexe $R' \otimes_R N_1 \to
R' \otimes_R N_2 \to R' \otimes_R N_3$ est exact (puisque $R'$
est plat comme $R$-module), et le complexe
$M' \otimes_{R'} \left(R' \otimes_R N_1\right)
\to M' \otimes_{R'} \left(R' \otimes_R N_2\right)
\to M' \otimes_{R'} \left(R' \otimes_R N_3\right)$ est donc
exact (puisque $M'$ est un $R'$-module plat). Comme
$M' \otimes_{R'} \left(R' \otimes_R N\right)
\cong \left(M' \otimes_{R'} R'\right) \otimes_R N
\cong M' \otimes_R N$ pour tout $R$-module $N$, fonctoriellement
(d'après les lemmes \ref{lemma-tensor-with-bimodule} et
\ref{lemma-flip-tensor-product}), ce complexe est isomorphe
au complexe
$M' \otimes_R N_1 \to M' \otimes_R N_2 \to M' \otimes_R N_3$,
qui est donc lui aussi exact. Cela montre que $M'$ est un
$R$-module plat. En remontant cet argument, nous pouvons montrer
que, si $R \to R'$ est fidèlement plat et si $M'$ est
fidèlement plat comme $R'$-module, alors $M'$ est fidèlement
plat comme $R$-module.
\end{proof}

\begin{lemma}
\label{lemma-flat}
Soit $M$ un $R$-module. Les assertions suivantes sont équivalentes :
\begin{enumerate}
\item
\label{item-flat}
$M$ est plat sur $R$.
\item
\label{item-injective}
pour toute injection de $R$-modules $N \subset N'$,
l'application $N \otimes_R M \to N'\otimes_R M$ est injective.
\item
\label{item-f-ideal}
pour tout idéal $I \subset R$, l'application
$I \otimes_R M \to R \otimes_R M = M$ est injective.
\item
\label{item-ffg-ideal}
pour tout idéal de type fini $I \subset R$,
l'application $I \otimes_R M \to R \otimes_R M = M$ est injective.
\end{enumerate}
\end{lemma}

\begin{proof}
Les implications successives de (\ref{item-flat}) à (\ref{item-injective}), puis
à (\ref{item-f-ideal}) et à (\ref{item-ffg-ideal}), sont immédiates. Nous démontrons
donc que (\ref{item-ffg-ideal}) implique (\ref{item-flat}).
Supposons que $N_1 \to N_2 \to N_3$ soit exacte.
Posons $K = \Ker(N_2 \to N_3)$ et $Q = \Im(N_2 \to N_3)$.
Nous obtenons alors des applications
$$
N_1 \otimes_R M \to
K \otimes_R M \to
N_2 \otimes_R M \to
Q \otimes_R M \to
N_3 \otimes_R M
$$
Remarquons que les première et troisième flèches sont surjectives. Ainsi, si nous montrons
que les deuxième et quatrième flèches sont injectives, nous aurons
conclu\footnote{Voici l'argument plus en détail :
supposons que nous sachions que les deuxième et quatrième flèches sont
injectives. Le lemme \ref{lemma-tensor-product-exact} (appliqué
à la suite exacte $K \to N_2 \to Q \to 0$) donne que
la suite $K \otimes_R M \to N_2 \otimes_R M \to
Q \otimes_R M \to 0$ est exacte. Par conséquent,
$\Ker \left(N_2 \otimes_R M \to Q \otimes_R M\right)
= \Im \left(K \otimes_R M \to N_2 \otimes_R M\right)$.
Comme
$\Im \left(K \otimes_R M \to N_2 \otimes_R M\right)
= \Im \left(N_1 \otimes_R M \to N_2 \otimes_R M\right)$
(en raison de la surjectivité de $N_1 \otimes_R M \to
K \otimes_R M$) et
$\Ker \left(N_2 \otimes_R M \to Q \otimes_R M\right)
= \Ker \left(N_2 \otimes_R M \to N_3 \otimes_R M\right)$
(en raison de l'injectivité de $Q \otimes_R M \to
N_3 \otimes_R M$), cela devient
$\Ker \left(N_2 \otimes_R M \to N_3 \otimes_R M\right)
= \Im \left(N_1 \otimes_R M \to N_2 \otimes_R M\right)$,
ce qui montre que le foncteur $- \otimes_R M$ est exact,
et donc que $M$ est plat.}.
Il suffit donc de montrer que $- \otimes_R M$ transforme
les applications injectives de $R$-modules en applications injectives de $R$-modules.

\medskip\noindent
Supposons que $K \to N$ soit une application injective de $R$-modules et
soit $x \in \Ker(K \otimes_R M \to N \otimes_R M)$.
Nous devons montrer que $x$ est nul.
Le $R$-module $K$ est la réunion de ses
$R$-sous-modules de type fini ; par conséquent, $K \otimes_R M$ est
la colimite des $R$-modules de la forme
$K_i \otimes_R M$, où $K_i$ parcourt tous les
$R$-sous-modules de type fini de $K$
(car le produit tensoriel commute aux colimites).
Ainsi, pour un certain $i$, notre $x$ provient d'un élément
$x_i \in K_i \otimes_R M$. Nous pouvons donc supposer que $K$
est un $R$-module de type fini. Faisons cette hypothèse. Nous considérons
l'injection $K \to N$ comme une inclusion, de sorte que
$K \subset N$.

\medskip\noindent
Le $R$-module $N$ est la réunion de ses
$R$-sous-modules de type fini qui contiennent $K$. Par conséquent, $N \otimes_R M$
est la colimite des $R$-modules de la forme
$N_i \otimes_R M$, où $N_i$ parcourt tous les
$R$-sous-modules de type fini de $N$ qui contiennent $K$
(de nouveau parce que le produit tensoriel commute aux colimites).
Remarquons qu'il s'agit d'une colimite sur un système filtrant
(puisque la somme de deux sous-modules de type fini de $N$ est
encore de type fini).
Par conséquent, (d'après le lemme \ref{lemma-zero-directed-limit}),
l'élément $x \in K \otimes_R M$ s'envoie sur
zéro dans au moins un de ces $R$-modules
$N_i \otimes_R M$ (puisque $x$ s'envoie sur zéro
dans $N \otimes_R M$).
Nous pouvons donc supposer que $N$ est un $R$-module de type fini.

\medskip\noindent
Supposons $N$ de type fini sur $R$. Écrivons $N = R^{\oplus n}/L$ et $K = L'/L$
pour certains $L \subset L' \subset R^{\oplus n}$.
Pour tout $R$-sous-module $G \subset R^{\oplus n}$,
nous disposons d'une application canonique $G \otimes_R M \to M^{\oplus n}$
obtenue en composant
$G \otimes_R M \to R^n \otimes_R M = M^{\oplus n}$.
Il suffit de prouver que $L \otimes_R M \to M^{\oplus n}$
et $L' \otimes_R M \to M^{\oplus n}$ sont injectives.
En effet, si tel est le cas, nous voyons que
$K \otimes_R M = L' \otimes_R M/L \otimes_R M \to M^{\oplus n}/L \otimes_R M$
est elle aussi injective\footnote{Cela devient évident si nous
identifions $L' \otimes_R M$ et $L \otimes_R M$ à des
sous-modules de $M^{\oplus n}$ (ce qui est légitime puisque
les applications $L \otimes_R M \to M^{\oplus n}$
et $L' \otimes_R M \to M^{\oplus n}$ sont injectives et
commutent avec l'application évidente $L' \otimes_R M \to L \otimes_R M$).}.

\medskip\noindent
Il suffit donc de montrer que $L \otimes_R M \to M^{\oplus n}$
est injective lorsque $L \subset R^{\oplus n}$ est un $R$-sous-module.
Nous le faisons par récurrence sur $n$. Nous traitons ci-dessous le cas initial $n = 1$.
Pour l'étape de récurrence, supposons $n > 1$ et posons
$L' = L \cap R \oplus 0^{\oplus n - 1}$. Alors $L'' = L/L'$ est un sous-module
de $R^{\oplus n - 1}$. Nous obtenons un diagramme
$$
\xymatrix{
&
L' \otimes_R M \ar[r] \ar[d] &
L \otimes_R M \ar[r] \ar[d] &
L'' \otimes_R M \ar[r] \ar[d] &
0 \\
0 \ar[r] &
M \ar[r] &
M^{\oplus n} \ar[r] &
M^{\oplus n - 1} \ar[r] & 0
}
$$
D'après l'hypothèse de récurrence et le cas initial, les flèches verticales
de gauche et de droite sont injectives. Les lignes sont exactes. Il s'ensuit que la flèche verticale
du milieu est elle aussi injective.

\medskip\noindent
Le cas initial de la récurrence ci-dessus est celui où $L \subset R$ est un idéal.
Autrement dit, nous devons montrer que $I \otimes_R M \to M$ est injective
pour tout idéal $I$ de $R$. Nous savons que cela est vrai lorsque $I$ est de type
fini. Toutefois, $I = \bigcup I_\alpha$ est la réunion des
idéaux de type fini $I_\alpha$ qu'il contient. Autrement dit,
$I = \colim I_\alpha$. Puisque $\otimes$ commute aux colimites, nous voyons que
$I \otimes_R M = \colim I_\alpha \otimes_R M$ et, puisque toutes
les applications $I_\alpha \otimes_R M \to M$ sont injectives par
hypothèse, il en va de même de $I \otimes_R M \to M$.
\end{proof}

\begin{lemma}
\label{lemma-colimit-rings-flat}
Soit $\{R_i, \varphi_{ii'}\}$ un système d'anneaux sur l'ensemble filtrant $I$.
Soit $R = \colim_i R_i$.
\begin{enumerate}
\item Si $M$ est un $R$-module tel que $M$ soit plat comme $R_i$-module
pour tout $i$, alors $M$ est plat comme $R$-module.
\item Pour $i \in I$, soit $M_i$ un $R_i$-module plat et,
pour $i' \geq i$, soit $f_{ii'} : M_i \to M_{i'}$ une application
$\varphi_{ii'}$-linéaire telle que $f_{i' i''} \circ f_{i i'} = f_{i i''}$.
Alors $M = \colim_{i \in I} M_i$ est un $R$-module plat.
\end{enumerate}
\end{lemma}

\begin{proof}
La partie (1) est un cas particulier de la partie (2), avec $M_i = M$ pour tout $i$
et $f_{i i'} = \text{id}_M$. Démontrons (2).
Soit $\mathfrak a \subset R$ un idéal de type fini. D'après le
lemme \ref{lemma-flat},
il suffit de montrer que $\mathfrak a \otimes_R M \to M$ est
injective. Nous pouvons trouver $i \in I$ et un idéal de type fini
$\mathfrak a' \subset R_i$ tels que $\mathfrak a = \mathfrak a'R$.
Alors $\mathfrak a = \colim_{i' \geq i} \mathfrak a'R_{i'}$.
Puisque $\otimes$ commute aux colimites, l'application
$\mathfrak a \otimes_R M \to M$ est la colimite des applications
$$
\mathfrak a'R_{i'} \otimes_{R_{i'}} M_{i'} \longrightarrow M_{i'}
$$
Toutes ces applications sont injectives par hypothèse. Puisque les colimites sur $I$
sont exactes d'après le lemme \ref{lemma-directed-colimit-exact}, nous concluons.
\end{proof}

\begin{lemma}
\label{lemma-flat-base-change}
Supposons que $M$ soit (fidèlement) plat sur $R$ et que $R \to R'$
soit un morphisme d'anneaux. Alors $M \otimes_R R'$ est (fidèlement) plat sur $R'$.
\end{lemma}

\begin{proof}
Pour tout $R'$-module $N$, nous avons un
isomorphisme canonique $N \otimes_{R'} (R'\otimes_R M)
= N \otimes_R M$. Les propriétés d'exactitude voulues du foncteur
$-\otimes_{R'}(R'\otimes_R M)$ résultent donc
des propriétés d'exactitude correspondantes du foncteur $-\otimes_R M$.
\end{proof}

\begin{lemma}
\label{lemma-flatness-descends}
Soit $R \to R'$ un morphisme d'anneaux fidèlement plat.
Soit $M$ un module sur $R$, et posons $M' = R' \otimes_R M$.
Alors $M$ est plat sur $R$ si et seulement si $M'$ est plat sur $R'$.
\end{lemma}

\begin{proof}
D'après le lemme \ref{lemma-flat-base-change}, si $M$ est plat,
alors $M'$ est plat. Réciproquement, supposons $M'$ plat.
Soit $N_1 \to N_2 \to N_3$ une suite exacte de $R$-modules.
Nous voulons montrer que $N_1 \otimes_R M \to N_2 \otimes_R M \to N_3 \otimes_R M$
est exacte. Nous savons que
$N_1 \otimes_R R' \to N_2 \otimes_R R' \to N_3 \otimes_R R'$ est
exacte, car $R \to R'$ est plat. La platitude de $M'$ implique que
$N_1 \otimes_R R' \otimes_{R'} M'
\to N_2 \otimes_R R' \otimes_{R'} M'
\to N_3 \otimes_R R' \otimes_{R'} M'$ est exacte.
Nous pouvons l'écrire
$N_1 \otimes_R M \otimes_R R'
\to N_2 \otimes_R M \otimes_R R'
\to N_3 \otimes_R M \otimes_R R'$.
Enfin, la fidèle platitude implique que
$N_1 \otimes_R M \to N_2 \otimes_R M \to N_3 \otimes_R M$
est exacte.
\end{proof}

\begin{lemma}
\label{lemma-flatness-descends-more-general}
Soit $R$ un anneau. Soit $S \to S'$ un morphisme plat de $R$-algèbres.
Soit $M$ un module sur $S$, et posons $M' = S' \otimes_S M$.
\begin{enumerate}
\item Si $M$ est plat sur $R$, alors $M'$ est plat sur $R$.
\item Si $S \to S'$ est fidèlement plat, alors $M$ est plat
sur $R$ si et seulement si $M'$ est plat sur $R$.
\end{enumerate}
\end{lemma}

\begin{proof}
Soit $N \to N'$ une injection de $R$-modules. Par la platitude
de $S \to S'$, nous avons
$$
\Ker(N \otimes_R M \to N' \otimes_R M) \otimes_S S'
=
\Ker(N \otimes_R M' \to N' \otimes_R M')
$$
Si $M$ est plat sur $R$, alors le membre de gauche est nul et
nous trouvons que $M'$ est plat sur $R$ par la deuxième caractérisation
de la platitude dans le lemme \ref{lemma-flat}.
Si $M'$ est plat sur $R$, le membre de droite est nul et, si
en outre $S \to S'$ est fidèlement plat, cela implique que
$\Ker(N \otimes_R M \to N' \otimes_R M)$ est nul, ce qui
montre à son tour que $M$ est plat sur $R$.
\end{proof}

\begin{lemma}
\label{lemma-flat-permanence}
Soit $R \to S$ un morphisme d'anneaux. Soit $M$ un $S$-module.
Si $M$ est plat comme $R$-module et fidèlement plat comme $S$-module,
alors $R \to S$ est plat.
\end{lemma}

\begin{proof}
Soit $N_1 \to N_2 \to N_3$ une suite exacte de $R$-modules.
Par hypothèse, $N_1 \otimes_R M \to N_2 \otimes_R M \to N_3 \otimes_R M$
est exacte. Nous pouvons l'écrire
$$
N_1 \otimes_R S \otimes_S M
\to
N_2 \otimes_R S \otimes_S M
\to
N_3 \otimes_R S \otimes_S M.
$$
Par fidèle platitude de $M$ sur $S$, nous concluons que
$N_1 \otimes_R S \to N_2 \otimes_R S \to N_3 \otimes_R S$ est exacte.
Ainsi $R \to S$ est plat.
\end{proof}

\noindent
Soit $R$ un anneau.
Soit $M$ un $R$-module.
Soit $\sum f_i x_i = 0$ une relation dans $M$.
Nous disons que la relation $\sum f_i x_i$
est {\it triviale} s'il existe un entier $m \geq 0$,
des éléments $y_j \in M$, $j = 1, \ldots, m$, et des éléments $a_{ij} \in R$,
$i = 1, \ldots, n$, $j = 1, \ldots, m$, tels que
$$
x_i = \sum\nolimits_j a_{ij} y_j, \forall i,
\quad\text{et}\quad
0 = \sum\nolimits_i f_ia_{ij}, \forall j.
$$

\begin{lemma}[Critère équationnel de platitude]
\label{lemma-flat-eq}
Un module $M$ sur $R$ est plat si et seulement si
toute relation dans $M$ est triviale.
\end{lemma}

\begin{proof}
Supposons $M$ plat et soit $\sum f_i x_i = 0$ une relation dans $M$.
Posons $I = (f_1, \ldots, f_n)$ et
$K = \Ker(R^n \to I, (a_1, \ldots, a_n) \mapsto \sum_i a_i f_i)$.
Nous avons donc la suite exacte courte
$0 \to K \to R^n \to I \to 0$. Alors $\sum f_i \otimes x_i$
est un élément de $I \otimes_R M$ qui s'envoie
sur zéro dans $R \otimes_R M = M$. Par platitude,
$\sum f_i \otimes x_i$ est nul dans $I \otimes_R M$.
Il existe donc un élément de $K \otimes_R M$ dont l'image
est $\sum e_i \otimes x_i \in R^n \otimes_R M$, où $e_i$
est le $i$-ième vecteur de base de $R^n$.
Écrivons cet élément sous la forme $\sum k_j \otimes y_j$,
puis l'image de $k_j$ dans $R^n$ sous la forme
$\sum a_{ij} e_i$, ce qui donne le résultat.

\medskip\noindent
Supposons toute relation triviale, soit $I$
un idéal de type fini, et soit $x = \sum f_i \otimes x_i$
un élément de $I \otimes_R M$ qui s'envoie sur zéro dans $R \otimes_R M = M$.
Cela signifie exactement que $\sum f_i x_i$ est une relation dans
$M$. Le fait qu'elle soit triviale implique aisément que
$x$ est nul, car
$$
x
=
\sum f_i \otimes x_i
=
\sum f_i \otimes \left(\sum a_{ij}y_j\right)
=
\sum \left(\sum f_i a_{ij}\right) \otimes y_j
=
0
$$
\end{proof}

\begin{lemma}
\label{lemma-flat-tor-zero}
Supposons que $R$ soit un anneau, que $0 \to M'' \to M' \to M \to 0$
soit une suite exacte courte et que $N$ soit un $R$-module. Si $M$ est plat,
alors $N \otimes_R M'' \to N \otimes_R M'$ est injective, c'est-à-dire que la
suite
$$
0 \to N \otimes_R M'' \to N \otimes_R M' \to N \otimes_R M \to 0
$$
est une suite exacte courte.
\end{lemma}

\begin{proof}
Soit $R^{(I)} \to N$ une surjection d'un module libre
sur $N$, de noyau $K$. Le résultat découle
du lemme du serpent appliqué au diagramme suivant
$$
\begin{matrix}
 & & 0 & & 0 & & 0 & & \\
 & & \uparrow & & \uparrow & & \uparrow & & \\
 & & M''\otimes_R N & \to & M' \otimes_R N & \to & M \otimes_R N & \to & 0 \\
 & & \uparrow & & \uparrow & & \uparrow & & \\
0 & \to & (M'')^{(I)} & \to & (M')^{(I)} & \to & M^{(I)} & \to & 0 \\
 & & \uparrow & & \uparrow & & \uparrow & & \\
 & & M''\otimes_R K & \to & M' \otimes_R K & \to & M \otimes_R K & \to & 0 \\
 & & & & & & \uparrow & & \\
 & & & & & & 0 & &
\end{matrix}
$$
dont les lignes et les colonnes sont exactes. La ligne du milieu est exacte car tensoriser
par le module libre $R^{(I)}$ est exact.
\end{proof}


\begin{lemma}
\label{lemma-flat-ses}
Supposons que $0 \to M' \to M \to M'' \to 0$ soit
une suite exacte courte de $R$-modules.
Si $M'$ et $M''$ sont plats, alors $M$ l'est aussi.
Si $M$ et $M''$ sont plats, alors $M'$ l'est aussi.
\end{lemma}

\begin{proof}
Nous utiliserons le critère suivant : un module $N$ est plat si, pour
tout idéal $I \subset R$, l'application $N \otimes_R I \to N$ est injective ;
voir le lemme \ref{lemma-flat}.
Considérons un idéal $I \subset R$.
Considérons le diagramme
$$
\begin{matrix}
0 & \to & M' & \to & M & \to & M'' & \to & 0 \\
& & \uparrow & & \uparrow & & \uparrow & & \\
& & M'\otimes_R I & \to & M \otimes_R I & \to & M''\otimes_R I & \to & 0
\end{matrix}
$$
dont les lignes sont exactes. Cela démontre immédiatement la première assertion.
La seconde résulte du fait que, si $M''$ est plat, alors la flèche horizontale
inférieure gauche est injective d'après le lemme \ref{lemma-flat-tor-zero}.
\end{proof}

\begin{lemma}
\label{lemma-easy-ff}
Soit $R$ un anneau.
Soit $M$ un $R$-module.
Les assertions suivantes sont équivalentes :
\begin{enumerate}
\item $M$ est fidèlement plat, et
\item $M$ est plat et, pour tous les homomorphismes de $R$-modules $\alpha : N \to N'$,
nous avons $\alpha = 0$ si et seulement si $\alpha \otimes \text{id}_M = 0$.
\end{enumerate}
\end{lemma}

\begin{proof}
Si $M$ est fidèlement plat, alors
$0 \to \Ker(\alpha) \to N \to N'$ est exacte si et seulement si elle le reste
après tensorisation par $M$. Cela démontre que (1) implique (2).
Réciproquement, supposons (2). Soit $N_1 \to N_2 \to N_3$
un complexe et supposons que le complexe
$N_1 \otimes_R M \to N_2 \otimes_R M \to N_3\otimes_R M$
soit exact. Prenons $x \in \Ker(N_2 \to N_3)$
et considérons l'application $\alpha : R \to N_2/\Im(N_1)$,
$r \mapsto rx + \Im(N_1)$. Par l'exactitude
du complexe $-\otimes_R M$, nous voyons que $\alpha \otimes
\text{id}_M$ est nulle. Par hypothèse, nous obtenons que $\alpha$ est
nulle. Ainsi $x $ appartient à l'image de $N_1 \to N_2$.
\end{proof}

\begin{lemma}
\label{lemma-ff}
\begin{slogan}
Un module plat est fidèlement plat si et seulement si ses fibres sont non nulles.
\end{slogan}
Soit $M$ un $R$-module plat.
Les assertions suivantes sont équivalentes :
\begin{enumerate}
\item $M$ est fidèlement plat,
\item pour tout $R$-module non nul $N$, le produit tensoriel $M \otimes_R N$
est non nul,
\item pour tout $\mathfrak p \in \Spec(R)$,
le produit tensoriel $M \otimes_R \kappa(\mathfrak p)$ est non nul, et
\item pour tout idéal maximal $\mathfrak m$ de $R$,
le produit tensoriel $M \otimes_R \kappa(\mathfrak m) = M/{\mathfrak m}M$
est non nul.
\end{enumerate}
\end{lemma}

\begin{proof}
Supposons $M$ fidèlement plat et $N \not = 0$. D'après le lemme \ref{lemma-easy-ff},
l'application non nulle $1 : N \to N$ induit une application non nulle
$M \otimes_R N \to M \otimes_R N$, donc $M \otimes_R N \not = 0$.
Ainsi (1) implique (2). Les implications (2) $\Rightarrow$ (3) $\Rightarrow$ (4)
sont immédiates.

\medskip\noindent
Supposons (4). Supposons que $N_1 \to N_2 \to N_3$ soit un complexe et
que $N_1 \otimes_R M \to N_2\otimes_R M \to
N_3\otimes_R M$ soit exact. Soit $H$ l'homologie du complexe,
donc $H = \Ker(N_2 \to N_3)/\Im(N_1 \to N_2)$. Pour achever la preuve,
nous allons montrer que $H = 0$. Par platitude, nous voyons que $H \otimes_R M = 0$.
Prenons $x \in H$ et soit $I = \{f \in R \mid fx = 0 \}$
son annulateur. Puisque $R/I \subset H$, nous obtenons
$M/IM \subset H \otimes_R M = 0$ par platitude de $M$.
Si $I \not =  R$, nous pouvons choisir
un idéal maximal $I \subset \mathfrak m \subset R$.
Cela donne immédiatement une contradiction.
\end{proof}

\begin{lemma}
\label{lemma-ff-rings}
Soit $R \to S$ un morphisme d'anneaux plat.
Les assertions suivantes sont équivalentes :
\begin{enumerate}
\item $R \to S$ est fidèlement plat,
\item l'application induite sur $\Spec$ est surjective, et
\item tout point fermé $x \in \Spec(R)$ appartient
à l'image de l'application $\Spec(S) \to \Spec(R)$.
\end{enumerate}
\end{lemma}

\begin{proof}
Cela résulte rapidement du lemme \ref{lemma-ff}, car nous
avons vu dans la remarque \ref{remark-fundamental-diagram}
que $\mathfrak p$ appartient à l'image
si et seulement si l'anneau $S \otimes_R \kappa(\mathfrak p)$
est non nul.
\end{proof}

\begin{lemma}
\label{lemma-local-flat-ff}
Un homomorphisme local plat d'anneaux locaux est fidèlement plat.
\end{lemma}

\begin{proof}
Cela résulte immédiatement du lemme \ref{lemma-ff-rings}.
\end{proof}

\noindent
La platitude se comporte bien avec la localisation.

\begin{lemma}
\label{lemma-flat-localization}
Soit $R$ un anneau. Soit $S \subset R$ une partie multiplicative.
\begin{enumerate}
\item La localisation $S^{-1}R$ est une $R$-algèbre plate.
\item Si $M$ est un $S^{-1}R$-module, alors $M$ est un $R$-module plat
si et seulement si $M$ est un $S^{-1}R$-module plat.
\item Supposons que $M$ soit un $R$-module. Alors
$M$ est un $R$-module plat si et seulement si $M_{\mathfrak p}$ est un
$R_{\mathfrak p}$-module plat pour tout idéal premier $\mathfrak p$ de $R$.
\item Supposons que $M$ soit un $R$-module. Alors $M$ est un $R$-module plat si
et seulement si $M_{\mathfrak m}$ est un
$R_{\mathfrak m}$-module plat pour tout idéal maximal $\mathfrak m$ de $R$.
\item Supposons que $R \to A$ soit un morphisme d'anneaux, que $M$ soit un $A$-module,
et que $g_1, \ldots, g_m \in A$ soient des éléments qui engendrent l'idéal
unité de $A$. Alors $M$ est plat sur $R$ si et seulement si chaque localisation
$M_{g_i}$ est plate sur $R$.
\item Supposons que $R \to A$ soit un morphisme d'anneaux et que $M$ soit un $A$-module.
Alors $M$ est un $R$-module plat si et seulement si la localisation
$M_{\mathfrak q}$ est un $R_{\mathfrak p}$-module plat
(où $\mathfrak p$ est l'idéal premier de $R$ situé sous $\mathfrak q$)
pour tout idéal premier $\mathfrak q$ de $A$.
\item Supposons que $R \to A$ soit un morphisme d'anneaux et que $M$ soit un $A$-module.
Alors $M$ est un $R$-module plat si et seulement si la localisation
$M_{\mathfrak m}$ est un $R_{\mathfrak p}$-module plat
(où $\mathfrak p = R \cap \mathfrak m$)
pour tout idéal maximal $\mathfrak m$ de $A$.
\end{enumerate}
\end{lemma}

\begin{proof}
Démontrons la dernière assertion du lemme.
Dans la preuve, nous utiliserons à plusieurs reprises le fait que la localisation est exacte
et commute au produit tensoriel, voir les sections \ref{section-localization}
et \ref{section-tensor-product}.

\medskip\noindent
Supposons que $R \to A$ soit un morphisme d'anneaux et que $M$ soit un $A$-module.
Supposons que $M_{\mathfrak m}$ soit un $R_{\mathfrak p}$-module plat
pour tout idéal maximal $\mathfrak m$ de $A$ (avec
$\mathfrak p = R \cap \mathfrak m$). Soit $I \subset R$ un idéal.
Nous devons montrer que l'application $I \otimes_R M \to M$ est injective.
Nous pouvons la considérer comme une application de $A$-modules.
Par hypothèse, la localisation
$(I \otimes_R M)_{\mathfrak m} \to M_{\mathfrak m}$ est injective
car
$(I \otimes_R M)_{\mathfrak m} =
I_{\mathfrak p} \otimes_{R_{\mathfrak p}} M_{\mathfrak m}$.
Ainsi, le noyau de $I \otimes_R M \to M$ est nul d'après le
lemme \ref{lemma-characterize-zero-local}.
Par conséquent, $M$ est plat sur $R$.

\medskip\noindent
Réciproquement, supposons $M$ plat sur $R$. Prenons un idéal premier
$\mathfrak q$ de $A$ situé au-dessus de l'idéal premier $\mathfrak p$ de $R$.
Supposons que $I \subset R_{\mathfrak p}$ soit un idéal. Nous devons montrer que
$I \otimes_{R_{\mathfrak p}} M_{\mathfrak q} \to M_{\mathfrak q}$
est injective. Nous pouvons écrire $I = J_{\mathfrak p}$ pour un certain
idéal $J \subset R$. Alors l'application
$I \otimes_{R_{\mathfrak p}} M_{\mathfrak q} \to M_{\mathfrak q}$
est simplement la localisation (en $\mathfrak q$) de l'application
$J \otimes_R M \to M$, qui est injective. Puisque la localisation est exacte,
nous voyons que $M_{\mathfrak q}$ est un $R_{\mathfrak p}$-module plat.

\medskip\noindent
Cela démontre (7) et (6). Les autres assertions résultent de manière immédiate
de la dernière assertion (preuves omises).
\end{proof}

\begin{lemma}
\label{lemma-flat-going-down}
Soit $R \to S$ plat. Soient $\mathfrak p \subset \mathfrak p'$
des idéaux premiers de $R$. Soit $\mathfrak q' \subset S$ un idéal premier de $S$
qui s'envoie sur $\mathfrak p'$. Alors il existe un idéal premier
$\mathfrak q \subset \mathfrak q'$ qui s'envoie sur $\mathfrak p$.
\end{lemma}

\begin{proof}
D'après le lemme \ref{lemma-flat-localization}, le morphisme d'anneaux locaux
$R_{\mathfrak p'} \to S_{\mathfrak q'}$ est plat.
D'après le lemme \ref{lemma-local-flat-ff}, ce morphisme d'anneaux locaux est fidèlement
plat. D'après le lemme \ref{lemma-ff-rings}, il existe un idéal premier qui s'envoie sur
$\mathfrak p R_{\mathfrak p'}$. L'image réciproque de cet
idéal premier dans $S$ convient.
\end{proof}

\noindent
La propriété de $R \to S$ décrite dans le lemme est appelée
« propriété de descente ». Voir la définition \ref{definition-going-up-down}.

\begin{lemma}
\label{lemma-colimit-faithfully-flat}
Soit $R$ un anneau. Soit $\{S_i, \varphi_{ii'}\}$ un système filtrant de
$R$-algèbres fidèlement plates. Alors $S = \colim_i S_i$ est une
$R$-algèbre fidèlement plate.
\end{lemma}

\begin{proof}
D'après le lemme \ref{lemma-colimit-flat}, nous voyons que $S$ est plate.
Soit $\mathfrak m \subset R$ un idéal maximal. D'après le
lemme \ref{lemma-ff-rings},
aucun des anneaux $S_i/\mathfrak m S_i$ n'est nul.
Ainsi $S/\mathfrak mS = \colim S_i/\mathfrak mS_i$ est lui aussi non nul,
car $1$ n'est pas égal à zéro. L'image de
$\Spec(S) \to \Spec(R)$ contient donc $\mathfrak m$, et nous voyons que
$R \to S$ est fidèlement plat d'après le lemme \ref{lemma-ff-rings}.
\end{proof}





\section{Supports et annulateurs}
\label{section-supp-and-ann}

\noindent
Quelques définitions et lemmes très élémentaires.

\begin{definition}
\label{definition-support-module}
Soit $R$ un anneau et soit $M$ un $R$-module.
Le {\it support de $M$} est l'ensemble
$$
\text{Supp}(M)
=
\{
\mathfrak p \in \Spec(R)
\mid
M_{\mathfrak p} \not = 0
\}
$$
\end{definition}

\begin{lemma}
\label{lemma-support-zero}
\begin{slogan}
Un module sur un anneau est de support vide si et seulement s'il est nul.
\end{slogan}
Soit $R$ un anneau. Soit $M$ un $R$-module. Alors
$$
M = (0) \Leftrightarrow \text{Supp}(M) = \emptyset.
$$
\end{lemma}

\begin{proof}
En fait, le
lemme \ref{lemma-characterize-zero-local}
montre même que $\text{Supp}(M)$ contient toujours un idéal maximal
si $M$ n'est pas nul.
\end{proof}

\begin{definition}
\label{definition-annihilator}
Soit $R$ un anneau. Soit $M$ un $R$-module.
\begin{enumerate}
\item Étant donné un élément $m \in M$, l'{\it annulateur de $m$}
est l'idéal
$$
\text{Ann}_R(m) = \text{Ann}(m) = \{f \in R \mid fm = 0\}.
$$
\item L'{\it annulateur de $M$}
est l'idéal
$$
\text{Ann}_R(M) = \text{Ann}(M) = \{f \in R \mid fm = 0\ \forall m \in M\}.
$$
\end{enumerate}
\end{definition}

\begin{lemma}
\label{lemma-annihilator-flat-base-change}
Soit $R \to S$ un morphisme d'anneaux plat. Soit $M$ un $R$-module et
$m \in M$. Alors $\text{Ann}_R(m) S = \text{Ann}_S(m \otimes 1)$.
Si $M$ est un $R$-module de type fini, alors
$\text{Ann}_R(M) S = \text{Ann}_S(M \otimes_R S)$.
\end{lemma}

\begin{proof}
Posons $I = \text{Ann}_R(m)$. Par définition, il existe une suite exacte
$0 \to I \to R \to M$, où l'application $R \to M$ envoie $f$ sur $fm$. En utilisant
la platitude, nous obtenons une suite exacte
$0 \to I \otimes_R S \to S \to M \otimes_R S$, ce qui démontre la première
assertion. Si $m_1, \ldots, m_n$ est un système de générateurs de $M$,
alors $\text{Ann}_R(M) = \bigcap \text{Ann}_R(m_i)$. De même,
$\text{Ann}_S(M \otimes_R S) = \bigcap \text{Ann}_S(m_i \otimes 1)$.
Posons $I_i = \text{Ann}_R(m_i)$. Il suffit alors de montrer que
$\bigcap_{i = 1, \ldots, n} (I_i S) = (\bigcap_{i = 1, \ldots, n} I_i)S$.
C'est le lemme \ref{lemma-flat-intersect-ideals}.
\end{proof}

\begin{lemma}
\label{lemma-support-closed}
Soit $R$ un anneau et soit $M$ un $R$-module.
Si $M$ est de type fini, alors $\text{Supp}(M)$ est fermé.
Plus précisément, si $I = \text{Ann}(M)$ est l'annulateur de $M$, alors
$V(I) = \text{Supp}(M)$.
\end{lemma}

\begin{proof}
Nous allons montrer que $V(I) = \text{Supp}(M)$.

\medskip\noindent
Supposons $\mathfrak p \in \text{Supp}(M)$. Alors $M_{\mathfrak p} \not = 0$.
Choisissons un élément $m \in M$ dont l'image dans $M_\mathfrak p$ est non nulle.
L'annulateur de $m$ est alors contenu dans $\mathfrak p$ par construction
de la localisation $M_\mathfrak p$. Par conséquent,
à plus forte raison, $I = \text{Ann}(M)$ doit être contenu dans $\mathfrak p$.

\medskip\noindent
Réciproquement, supposons que $\mathfrak p \not \in \text{Supp}(M)$.
Alors $M_{\mathfrak p} = 0$.
Soient $x_1, \ldots, x_r \in M$ des générateurs.
D'après le lemme \ref{lemma-localization-colimit}, il existe
$f \in R$, $f\not\in \mathfrak p$, tel que
$x_i/1 = 0$ dans $M_f$. Ainsi $f^{n_i} x_i = 0$ pour un certain $n_i \geq 1$.
Par conséquent $f^nM = 0$ pour $n = \max\{n_i\}$, comme voulu.
\end{proof}

\begin{lemma}
\label{lemma-support-base-change}
Soit $R \to R'$ un morphisme d'anneaux et soit $M$ un $R$-module de type fini.
Alors $\text{Supp}(M \otimes_R R')$ est l'image réciproque de
$\text{Supp}(M)$.
\end{lemma}

\begin{proof}
Soit $\mathfrak p \in \text{Supp}(M)$. D'après le lemme de Nakayama
(lemme \ref{lemma-NAK}), nous voyons que
$$
M \otimes_R \kappa(\mathfrak p) = M_\mathfrak p/\mathfrak p M_\mathfrak p
$$
est un espace vectoriel non nul sur $\kappa(\mathfrak p)$.
Ainsi, pour tout idéal premier $\mathfrak p' \subset R'$ situé
au-dessus de $\mathfrak p$, nous voyons que
$$
(M \otimes_R R')_{\mathfrak p'}/\mathfrak p' (M \otimes_R R')_{\mathfrak p'} =
(M \otimes_R R') \otimes_{R'} \kappa(\mathfrak p') =
M \otimes_R \kappa(\mathfrak p) \otimes_{\kappa(\mathfrak p)}
\kappa(\mathfrak p')
$$
est non nul. Cela implique $\mathfrak p' \in \text{Supp}(M \otimes_R R')$.
Réciproquement, si $\mathfrak p' \subset R'$ est un idéal premier situé
au-dessus d'un idéal premier quelconque $\mathfrak p \subset R$, alors
$$
(M \otimes_R R')_{\mathfrak p'} =
M_\mathfrak p \otimes_{R_\mathfrak p} R'_{\mathfrak p'}.
$$
Par conséquent, si $\mathfrak p' \in \text{Supp}(M \otimes_R R')$
est situé au-dessus de l'idéal premier $\mathfrak p \subset R$, alors
$\mathfrak p \in \text{Supp}(M)$.
\end{proof}

\begin{lemma}
\label{lemma-support-element}
Soit $R$ un anneau, soit $M$ un $R$-module et soit $m \in M$.
Alors $\mathfrak p \in V(\text{Ann}(m))$ si et seulement si
$m$ ne s'envoie pas sur zéro dans $M_\mathfrak p$.
\end{lemma}

\begin{proof}
Nous pouvons remplacer $M$ par $Rm \subset M$. Alors (1) $\text{Ann}(m) = \text{Ann}(M)$
et (2) $m$ ne s'envoie pas sur zéro dans $M_\mathfrak p$ si et seulement si
$\mathfrak p \in \text{Supp}(M)$.
Le résultat découle maintenant du lemme \ref{lemma-support-closed}.
\end{proof}

\begin{lemma}
\label{lemma-support-finite-presentation-constructible}
Soit $R$ un anneau et soit $M$ un $R$-module.
Si $M$ est un $R$-module de présentation finie, alors $\text{Supp}(M)$ est une
partie fermée de $\Spec(R)$ dont le complémentaire est quasi-compact.
\end{lemma}

\begin{proof}
Choisissons une présentation
$$
R^{\oplus m} \longrightarrow R^{\oplus n} \longrightarrow M \to 0
$$
Soit $A \in \text{Mat}(n \times m, R)$ la matrice de la première
application. D'après le lemme de Nakayama \ref{lemma-NAK}, nous voyons que
$$
M_{\mathfrak p} \not = 0 \Leftrightarrow
M \otimes \kappa(\mathfrak p) \not = 0 \Leftrightarrow
\text{rang}(A \bmod \mathfrak p) < n.
$$
Ainsi, si $I$ est l'idéal de $R$ engendré par les mineurs $n \times n$
de $A$, alors $\text{Supp}(M) = V(I)$. Comme $I$
est de type fini, disons $I = (f_1, \ldots, f_t)$,
nous voyons que $\Spec(R) \setminus V(I)$ est
une réunion finie des ouverts standards $D(f_i)$, donc est quasi-compact.
\end{proof}

\begin{lemma}
\label{lemma-support-quotient}
Soit $R$ un anneau et soit $M$ un $R$-module.
\begin{enumerate}
\item Si $M$ est de type fini, alors le support
de $M/IM$ est $\text{Supp}(M) \cap V(I)$.
\item Si $N \subset M$, alors $\text{Supp}(N) \subset
\text{Supp}(M)$.
\item Si $Q$ est un module quotient de $M$, alors $\text{Supp}(Q) \subset
\text{Supp}(M)$.
\item Si $0 \to N \to M \to Q \to 0$ est une suite exacte courte,
alors $\text{Supp}(M) = \text{Supp}(Q) \cup \text{Supp}(N)$.
\end{enumerate}
\end{lemma}

\begin{proof}
Les foncteurs $M \mapsto M_{\mathfrak p}$ sont exacts. Cela implique immédiatement
toutes les assertions sauf la première. Pour la première assertion,
nous devons montrer que $M_\mathfrak p \not = 0$ et
$I \subset \mathfrak p$ impliquent $(M/IM)_{\mathfrak p}
= M_\mathfrak p/IM_\mathfrak p \not = 0$. Cela résulte
du lemme de Nakayama \ref{lemma-NAK}.
\end{proof}





\section{Montée et descente}
\label{section-going-up}

\noindent
Supposons que $\mathfrak p$, $\mathfrak p'$ soient des idéaux premiers
de l'anneau $R$. Soit $X = \Spec(R)$ muni de la topologie de
Zariski. Notons $x \in X$ le point correspondant
à $\mathfrak p$ et $x' \in X$ le point correspondant
à $\mathfrak p'$. Nous avons alors :
$$
x' \leadsto x \Leftrightarrow \mathfrak p' \subset \mathfrak p.
$$
Autrement dit, $x$ est une spécialisation de $x'$ si et
seulement si $\mathfrak p' \subset \mathfrak p$.
Voir Topologie, section \ref{topology-section-specialization},
pour la terminologie et les notations.

\begin{definition}
\label{definition-going-up-down}
Soit $\varphi : R \to S$ un morphisme d'anneaux.
\begin{enumerate}
\item Nous disons que $\varphi : R \to S$ vérifie la {\it propriété de montée} si,
étant donnés des idéaux premiers $\mathfrak p \subset \mathfrak p'$ de $R$
et un idéal premier $\mathfrak q$ de $S$ situé au-dessus de $\mathfrak p$,
il existe un idéal premier $\mathfrak q'$ de $S$ tel que
(a) $\mathfrak q \subset \mathfrak q'$ et (b)
$\mathfrak q'$ soit situé au-dessus de $\mathfrak p'$.
\item Nous disons que $\varphi : R \to S$ vérifie la {\it propriété de descente} si,
étant donnés des idéaux premiers $\mathfrak p \subset \mathfrak p'$ de $R$
et un idéal premier $\mathfrak q'$ de $S$ situé au-dessus de $\mathfrak p'$,
il existe un idéal premier $\mathfrak q$ de $S$ tel que
(a) $\mathfrak q \subset \mathfrak q'$ et (b)
$\mathfrak q$ soit situé au-dessus de $\mathfrak p$.
\end{enumerate}
\end{definition}

\noindent
Nous avons jusqu'ici rencontré les cas suivants :
\begin{enumerate}
\item Un morphisme d'anneaux entier vérifie la propriété de montée, voir le
lemme \ref{lemma-integral-going-up}.
\item En particulier, les morphismes d'anneaux finis vérifient la propriété de montée.
\item En particulier, les morphismes quotients $R \to R/I$ vérifient la propriété de montée.
\item Un morphisme d'anneaux plat vérifie la propriété de descente, voir le
lemme \ref{lemma-flat-going-down}.
\item En particulier, toute localisation vérifie la propriété de descente.
\item Une extension $R \subset S$ d'anneaux intègres, où $R$ est normal
et $S$ est entier sur $R$, vérifie la propriété de descente, voir la
proposition \ref{proposition-going-down-normal-integral}.
\end{enumerate}
Voici un autre cas où la propriété de descente est vérifiée.

\begin{lemma}
\label{lemma-open-going-down}
Soit $R \to S$ un morphisme d'anneaux. Si l'application induite
$\varphi : \Spec(S) \to \Spec(R)$ est ouverte, alors
$R \to S$ vérifie la propriété de descente.
\end{lemma}

\begin{proof}
Supposons que $\mathfrak p \subset \mathfrak p' \subset R$ et qu'un
$\mathfrak q' \subset S$ soit situé au-dessus de $\mathfrak p'$. Comme $\varphi$ est ouverte,
pour tout $g \in S$, $g \not \in \mathfrak q'$, nous voyons que $\mathfrak p$
appartient à l'image de $D(g) \subset \Spec(S)$. Autrement dit,
$S_g \otimes_R \kappa(\mathfrak p)$ n'est pas nul. Puisque $S_{\mathfrak q'}$
est la colimite filtrante de ces $S_g$, cela implique
que $S_{\mathfrak q'} \otimes_R \kappa(\mathfrak p)$ n'est pas
nul, voir les
lemmes \ref{lemma-localization-colimit} et
\ref{lemma-tensor-products-commute-with-limits}.
Par conséquent, $\mathfrak p$ appartient à l'image de
$\Spec(S_{\mathfrak q'}) \to \Spec(R)$, comme voulu.
\end{proof}

\begin{lemma}
\label{lemma-going-up-down-specialization}
Soit $R \to S$ un morphisme d'anneaux.
\begin{enumerate}
\item $R \to S$ vérifie la propriété de descente si et seulement si
les généralisations se relèvent le long de l'application $\Spec(S) \to \Spec(R)$,
voir Topologie, définition \ref{topology-definition-lift-specializations}.
\item $R \to S$ vérifie la propriété de montée si et seulement si
les spécialisations se relèvent le long de l'application $\Spec(S) \to \Spec(R)$,
voir Topologie, définition \ref{topology-definition-lift-specializations}.
\end{enumerate}
\end{lemma}

\begin{proof}
Omis.
\end{proof}

\begin{lemma}
\label{lemma-going-up-down-composition}
Supposons que $R \to S$ et $S \to T$ soient des morphismes d'anneaux vérifiant
la propriété de descente. Alors $R \to T$ la vérifie aussi. Il en va de même pour la propriété de montée.
\end{lemma}

\begin{proof}
D'après le lemme \ref{lemma-going-up-down-specialization},
cela résulte du
lemme de Topologie \ref{topology-lemma-lift-specialization-composition}.
\end{proof}

\begin{lemma}
\label{lemma-image-stable-specialization-closed}
Soit $R \to S$ un morphisme d'anneaux. Soit $T \subset \Spec(R)$
l'image de $\Spec(S)$. Si $T$ est stable par spécialisation,
alors $T$ est fermé.
\end{lemma}

\begin{proof}
Nous donnons deux preuves.

\medskip\noindent
Première preuve. Soit $\mathfrak p \subset R$ un idéal premier tel que
le point correspondant de $\Spec(R)$ appartienne à l'adhérence
de $T$. Cela signifie que, pour tout $f \in R$, $f \not \in \mathfrak p$,
nous avons $D(f) \cap T \not = \emptyset$. Remarquons que $D(f) \cap T$
est l'image de $\Spec(S_f)$ dans $\Spec(R)$. Nous
concluons donc que $S_f \not = 0$. Autrement dit, $1 \not = 0$ dans
l'anneau $S_f$. Puisque $S_{\mathfrak p}$ est la colimite filtrante
des anneaux $S_f$, nous concluons que $1 \not = 0$ dans
$S_{\mathfrak p}$. Autrement dit, $S_{\mathfrak p} \not = 0$ et,
en considérant l'image de $\Spec(S_{\mathfrak p})
\to \Spec(S) \to \Spec(R)$, nous voyons qu'il existe
$\mathfrak p' \in T$ avec $\mathfrak p' \subset \mathfrak p$.
Comme nous avons supposé $T$ stable par spécialisation, nous concluons que
$\mathfrak p$ est un point de $T$, comme voulu.

\medskip\noindent
Deuxième preuve. Posons $I = \Ker(R \to S)$. Nous pouvons remplacer $R$ par $R/I$.
Dans ce cas, le morphisme d'anneaux $R \to S$ est injectif.
D'après le lemme \ref{lemma-injective-minimal-primes-in-image},
tous les idéaux premiers minimaux de $R$ appartiennent à l'image $T$. Par conséquent,
si $T$ est stable par spécialisation, alors il contient tous les idéaux premiers.
\end{proof}

\begin{lemma}
\label{lemma-going-up-closed}
Soit $R \to S$ un morphisme d'anneaux. Les assertions suivantes sont équivalentes :
\begin{enumerate}
\item La propriété de montée est vérifiée pour $R \to S$, et
\item l'application $\Spec(S) \to \Spec(R)$ est fermée.
\end{enumerate}
\end{lemma}

\begin{proof}
Les spécialisations se relèvent le long de toute
application fermée d'espaces topologiques ; voir le
lemme de Topologie \ref{topology-lemma-closed-open-map-specialization}.
La seconde condition implique donc la première.

\medskip\noindent
Supposons que la propriété de montée soit vérifiée pour $R \to S$.
Soit $V(I) \subset \Spec(S)$ une partie fermée.
Nous voulons montrer que l'image de $V(I)$ dans $\Spec(R)$ est fermée.
Le morphisme d'anneaux $S \to S/I$ vérifie évidemment la propriété de montée.
Ainsi $R \to S \to S/I$ vérifie la propriété de montée,
d'après le lemme \ref{lemma-going-up-down-composition}.
En remplaçant $S$ par $S/I$, il suffit de montrer que l'image $T$
de $\Spec(S)$ dans $\Spec(R)$ est fermée.
D'après les lemmes de Topologie \ref{topology-lemma-open-closed-specialization}
et \ref{topology-lemma-lift-specializations-images}, cette
image est stable par spécialisation. Le résultat découle donc
du lemme \ref{lemma-image-stable-specialization-closed}.
\end{proof}

\begin{lemma}
\label{lemma-constructible-stable-specialization-closed}
Soit $R$ un anneau. Soit $E \subset \Spec(R)$ une partie constructible.
\begin{enumerate}
\item Si $E$ est stable par spécialisation, alors $E$ est fermée.
\item Si $E$ est stable par généralisation, alors $E$ est ouverte.
\end{enumerate}
\end{lemma}

\begin{proof}
Première preuve. La première assertion
résulte du lemme \ref{lemma-image-stable-specialization-closed}
combiné au lemme \ref{lemma-constructible-is-image}.
La seconde résulte du fait que le complémentaire d'une partie constructible
est constructible
(voir Topologie, lemme \ref{topology-lemma-constructible}),
de la première partie du lemme et du
lemme de Topologie \ref{topology-lemma-open-closed-specialization}.

\medskip\noindent
Deuxième preuve. Puisque $\Spec(R)$ est un espace spectral d'après le
lemme \ref{lemma-spec-spectral}, il s'agit d'un cas particulier du
lemme de Topologie
\ref{topology-lemma-constructible-stable-specialization-closed}.
\end{proof}

\begin{proposition}
\label{proposition-fppf-open}
Soit $R \to S$ plat et de présentation finie.
Alors $\Spec(S) \to \Spec(R)$ est ouverte.
Plus généralement, cela vaut pour tout morphisme d'anneaux $R \to S$ de
présentation finie qui vérifie la propriété de descente.
\end{proposition}

\begin{proof}
Si $R \to S$ est plat, alors $R \to S$ vérifie la propriété de descente d'après le
lemme \ref{lemma-flat-going-down}. Pour démontrer la proposition, nous pouvons donc
supposer que $R \to S$ est de présentation finie et vérifie la propriété de descente.

\medskip\noindent
Puisque les ouverts standards $D(g) \subset \Spec(S)$, $g \in S$,
forment une base de la topologie, il suffit de prouver que
l'image de $D(g)$ est ouverte. Rappelons que $\Spec(S_g) \to \Spec(S)$
est un homéomorphisme de $\Spec(S_g)$ sur $D(g)$
(lemme \ref{lemma-standard-open}).
Puisque $S \to S_g$ vérifie la propriété de descente (voir ci-dessus), nous voyons que
$R \to S_g$ vérifie la propriété de descente d'après le
lemme \ref{lemma-going-up-down-composition}. Ainsi, après avoir remplacé
$S$ par $S_g$, nous voyons qu'il suffit de prouver que l'image est ouverte.
D'après le théorème de Chevalley (théorème \ref{theorem-chevalley}),
l'image est une partie constructible $E$. Et $E$ est stable
par généralisation parce que $R \to S$ vérifie la propriété de descente ;
voir les lemmes de Topologie \ref{topology-lemma-open-closed-specialization}
et \ref{topology-lemma-lift-specializations-images}.
Par conséquent, $E$ est ouverte d'après le
lemme \ref{lemma-constructible-stable-specialization-closed}.
\end{proof}

\begin{lemma}
\label{lemma-same-image}
Soit $k$ un corps, et soient $R$, $S$ des $k$-algèbres.
Soit $S' \subset S$ une sous-$k$-algèbre, et soit $f \in S' \otimes_k R$.
Dans le diagramme commutatif
$$
\xymatrix{
\Spec((S \otimes_k R)_f) \ar[rd] \ar[rr] & &
\Spec((S' \otimes_k R)_f) \ar[ld] \\
& \Spec(R) &
}
$$
les images des flèches diagonales sont les mêmes.
\end{lemma}

\begin{proof}
Soit $\mathfrak p \subset R$ dans l'image de la flèche
sud-ouest. Cela signifie (lemme \ref{lemma-in-image}) que
$$
(S' \otimes_k R)_f \otimes_R \kappa(\mathfrak p)
=
(S' \otimes_k \kappa(\mathfrak p))_f
$$
n'est pas l'anneau nul, c'est-à-dire que $S' \otimes_k \kappa(\mathfrak p)$
n'est pas l'anneau nul et que l'image de $f$ dans cet anneau n'est pas nilpotente.
Le morphisme d'anneaux
$S' \otimes_k \kappa(\mathfrak p) \to S \otimes_k \kappa(\mathfrak p)$
est injectif. Ainsi, $S \otimes_k \kappa(\mathfrak p)$
n'est pas non plus l'anneau nul et l'image de $f$ n'y est pas nilpotente.
Par conséquent, $(S \otimes_k R)_f \otimes_R \kappa(\mathfrak p)$
n'est pas l'anneau nul. Ainsi (lemme \ref{lemma-in-image}),
nous voyons que $\mathfrak p$ appartient à l'image de la flèche sud-est,
comme voulu.
\end{proof}

\begin{lemma}
\label{lemma-map-into-tensor-algebra-open}
Soit $k$ un corps.
Soient $R$ et $S$ des $k$-algèbres.
L'application $\Spec(S \otimes_k R) \to \Spec(R)$
est ouverte.
\end{lemma}

\begin{proof}
Soit $f \in S \otimes_k R$.
Il suffit de prouver que l'image de l'ouvert standard $D(f)$ est ouverte.
Soit $S' \subset S$ une sous-$k$-algèbre de type fini telle que
$f \in S' \otimes_k R$. L'application $R \to S' \otimes_k R$ est plate
et de présentation finie ; l'image $U$ de
$\Spec((S' \otimes_k R)_f) \to \Spec(R)$ est donc ouverte
d'après la proposition \ref{proposition-fppf-open}.
D'après le lemme \ref{lemma-same-image}, c'est aussi l'image de $D(f)$, et nous concluons.
\end{proof}

\noindent
Voici un lemme délicat qui est parfois utile.

\begin{lemma}
\label{lemma-unique-prime-over-localize-below}
Soit $R \to S$ un morphisme d'anneaux.
Soit $\mathfrak p \subset R$ un idéal premier.
Supposons que
\begin{enumerate}
\item il existe un unique idéal premier $\mathfrak q \subset S$ situé au-dessus de
$\mathfrak p$, et
\item soit
\begin{enumerate}
\item la propriété de montée est vérifiée pour $R \to S$, soit
\item la propriété de descente est vérifiée pour $R \to S$ et il y a au plus un idéal premier
de $S$ au-dessus de chaque idéal premier de $R$.
\end{enumerate}
\end{enumerate}
Alors $S_{\mathfrak p} = S_{\mathfrak q}$.
\end{lemma}

\begin{proof}
Considérons un idéal premier quelconque $\mathfrak q' \subset S$ qui correspond à
un point de $\Spec(S_{\mathfrak p})$. Cela signifie que
$\mathfrak p' = R \cap \mathfrak q'$ est contenu dans $\mathfrak p$.
Voici un diagramme
$$
\xymatrix{
\mathfrak q' \ar@{-}[d] \ar@{-}[r] & ? \ar@{-}[r] \ar@{-}[d] & S \ar@{-}[d] \\
\mathfrak p' \ar@{-}[r] & \mathfrak p \ar@{-}[r] & R
}
$$
Supposons (1) et (2)(a).
Par montée, il existe un idéal premier $\mathfrak q'' \subset S$
tel que $\mathfrak q' \subset \mathfrak q''$ et que $\mathfrak q''$
soit situé au-dessus de $\mathfrak p$. Par l'unicité de $\mathfrak q$, nous
concluons que $\mathfrak q'' = \mathfrak q$. Autrement dit,
$\mathfrak q'$ définit un point de $\Spec(S_{\mathfrak q})$.

\medskip\noindent
Supposons (1) et (2)(b).
Par descente, il existe un idéal premier $\mathfrak q'' \subset \mathfrak q$
situé au-dessus de $\mathfrak p'$. Par l'unicité des idéaux premiers situés au-dessus de
$\mathfrak p'$, nous voyons que $\mathfrak q' = \mathfrak q''$. Autrement dit,
$\mathfrak q'$ définit un point de $\Spec(S_{\mathfrak q})$.

\medskip\noindent
Dans les deux cas, nous concluons que l'application
$\Spec(S_{\mathfrak q}) \to \Spec(S_{\mathfrak p})$
est bijective. Cela signifie clairement que tous les éléments de $S - \mathfrak q$
sont inversibles dans $S_{\mathfrak p}$, autrement dit
$S_{\mathfrak p} = S_{\mathfrak q}$.
\end{proof}

\noindent
Le lemme suivant est une généralisation de la propriété de descente pour les
morphismes d'anneaux plats.

\begin{lemma}
\label{lemma-going-down-flat-module}
Soit $R \to S$ un morphisme d'anneaux. Soit $N$ un $S$-module de type fini plat sur $R$.
Munissons $\text{Supp}(N) \subset \Spec(S)$ de la topologie induite.
Alors les généralisations se relèvent le long de $\text{Supp}(N) \to \Spec(R)$.
\end{lemma}

\begin{proof}
L'énoncé signifie ce qui suit. Soient
$\mathfrak p \subset \mathfrak p' \subset R$ des idéaux premiers. Soit
$\mathfrak q' \subset S$ un idéal premier tel que $\mathfrak q' \in \text{Supp}(N)$.
Il existe alors un idéal premier $\mathfrak q \subset \mathfrak q'$,
$\mathfrak q \in \text{Supp}(N)$, situé au-dessus de $\mathfrak p$.
Comme $N$ est plat sur $R$, nous voyons que $N_{\mathfrak q'}$ est plat
sur $R_{\mathfrak p'}$, voir le lemme \ref{lemma-flat-localization}.
Comme $N_{\mathfrak q'}$ est de type fini sur $S_{\mathfrak q'}$
et non nul puisque $\mathfrak q' \in \text{Supp}(N)$, nous voyons
que $N_{\mathfrak q'} \otimes_{S_{\mathfrak q'}} \kappa(\mathfrak q')$
est non nul d'après le lemme de Nakayama \ref{lemma-NAK}.
Ainsi $N_{\mathfrak q'} \otimes_{R_{\mathfrak p'}} \kappa(\mathfrak p')$
n'est pas nul non plus. Nous déduisons du lemme \ref{lemma-ff}
que $N_{\mathfrak q'} \otimes_{R_{\mathfrak p'}} \kappa(\mathfrak p)$
est non nul. Soit
$J \subset S_{\mathfrak q'} \otimes_{R_{\mathfrak p'}} \kappa(\mathfrak p)$
l'annulateur du module non nul de type fini
$N_{\mathfrak q'} \otimes_{R_{\mathfrak p'}} \kappa(\mathfrak p)$.
Puisque $J$ est un idéal propre, nous pouvons choisir un idéal premier
$\mathfrak q \subset S$ qui correspond à un idéal premier de
$S_{\mathfrak q'} \otimes_{R_{\mathfrak p'}} \kappa(\mathfrak p)/J$.
Cet idéal premier appartient au support de $N$, est situé au-dessus de
$\mathfrak p$ et est contenu dans $\mathfrak q'$, comme voulu.
\end{proof}

\section{Extensions séparables}
\label{section-separability}

\noindent
Dans cette section, nous étudions la séparabilité pour les extensions de
corps non algébriques. Elle est étroitement liée à la notion d'algèbres
géométriquement réduites, voir la
Définition \ref{definition-geometrically-reduced}.

\begin{definition}
\label{definition-separable-field-extension}
Soit $K/k$ une extension de corps.
\begin{enumerate}
\item On dit que $K$ est {\it séparablement engendré sur $k$} s'il existe
une base de transcendance $\{x_i; i \in I\}$ de $K/k$ telle que l'extension
$K/k(x_i; i \in I)$ soit une extension algébrique séparable.
\item On dit que $K$ est {\it séparable sur $k$} si, pour toute sous-extension
$k \subset K' \subset K$ avec $K'$ de type fini
sur $k$, l'extension $K'/k$ est séparablement engendrée.
\end{enumerate}
\end{definition}

\noindent
Avec cette définition quelque peu maladroite, il n'est pas clair qu'une
extension de corps séparablement engendrée soit elle-même séparable.
Nous verrons que c'est bien le cas, voir le
Lemme \ref{lemma-separably-generated-separable}.

\begin{lemma}
\label{lemma-subextensions-are-separable}
Soit $K/k$ une extension de corps séparable.
Pour toute sous-extension $K/K'/k$, l'extension de corps
$K'/k$ est séparable.
\end{lemma}

\begin{proof}
C'est immédiat d'après la définition.
\end{proof}

\begin{lemma}
\label{lemma-generating-finitely-generated-separable-field-extensions}
Soit $K/k$ une extension de corps séparablement engendrée et de type fini.
Posons $r = \text{trdeg}_k(K)$. Alors il existe des éléments
$x_1, \ldots, x_{r + 1}$ de $K$ tels que
\begin{enumerate}
\item $x_1, \ldots, x_r$ soit une base de transcendance de $K$ sur $k$,
\item $K = k(x_1, \ldots, x_{r + 1})$, et
\item $x_{r + 1}$ soit séparable sur $k(x_1, \ldots, x_r)$.
\end{enumerate}
\end{lemma}

\begin{proof}
Il suffit de combiner la définition avec le
lemme \ref{fields-lemma-primitive-element} de Corps.
\end{proof}

\begin{lemma}
\label{lemma-make-separably-generated}
Soit $K/k$ une extension de corps de type fini.
Il existe un diagramme
$$
\xymatrix{
K \ar[r] & K' \\
k \ar[u] \ar[r] & k' \ar[u]
}
$$
où $k'/k$ et $K'/K$ sont des extensions de corps finies purement
inséparables, telles que $K'/k'$ soit une extension de corps
séparablement engendrée.
\end{lemma}

\begin{proof}
Ce lemme n'est intéressant que lorsque la caractéristique de $k$ est
$p > 0$. Choisissons $x_1, \ldots, x_r$ comme base de transcendance
de $K$ sur $k$. Puisque $K$ est de type fini sur $k$, l'extension
$k(x_1, \ldots, x_r) \subset K$ est finie.
Soit $K/K_{sep}/k(x_1, \ldots, x_r)$ la sous-extension trouvée dans le
lemme \ref{fields-lemma-separable-first} de Corps.
Si $K = K_{sep}$, c'est terminé.
Nous raisonnerons par récurrence sur $d = [K : K_{sep}]$.

\medskip\noindent
Supposons que $d > 1$. Choisissons un $\beta \in K$ tel que
$\alpha = \beta^p \in K_{sep}$ et $\beta \not \in K_{sep}$.
Soit $P = T^n + a_1T^{n - 1} + \ldots + a_n$
le polynôme minimal de $\alpha$ sur $k(x_1, \ldots, x_r)$.
Soit $k'/k$ une extension finie purement inséparable
obtenue en adjoignant des racines $p$-ièmes, de sorte que chaque $a_i$
soit une puissance $p$-ième dans $k'(x_1^{1/p}, \ldots, x_r^{1/p})$.
Une telle extension existe ; les détails sont omis.
Soit $L$ un corps s'insérant dans le diagramme
$$
\xymatrix{
K \ar[r] & L \\
k(x_1, \ldots, x_r) \ar[u] \ar[r] & k'(x_1^{1/p}, \ldots, x_r^{1/p}) \ar[u]
}
$$
Nous pouvons et allons supposer que $L$ est le compositum de $K$ et
de $k'(x_1^{1/p}, \ldots, x_r^{1/p})$.
Soit $L/L_{sep}/k'(x_1^{1/p}, \ldots, x_r^{1/p})$
la sous-extension trouvée dans le
lemme \ref{fields-lemma-separable-first} de Corps.
Alors $L_{sep}$ est le compositum de
$K_{sep}$ et de $k'(x_1^{1/p}, \ldots, x_r^{1/p})$.
L'élément $\alpha \in L_{sep}$ est une racine du polynôme
$P$, dont tous les coefficients sont des puissances $p$-ièmes dans
$k'(x_1^{1/p}, \ldots, x_r^{1/p})$ et dont les racines sont
deux à deux distinctes. D'après le
lemme \ref{fields-lemma-pth-root} de Corps,
nous voyons que $\alpha = (\alpha')^p$ pour un certain $\alpha' \in L_{sep}$.
Cela signifie clairement que $\beta$ s'envoie sur $\alpha' \in L_{sep}$.
Autrement dit, nous obtenons la tour de corps
$$
\xymatrix{
K \ar[r] & L \\
K_{sep}(\beta) \ar[r] \ar[u] & L_{sep} \ar[u] \\
K_{sep} \ar[r] \ar[u] & L_{sep} \ar@{=}[u] \\
k(x_1, \ldots, x_r) \ar[u] \ar[r] & k'(x_1^{1/p}, \ldots, x_r^{1/p}) \ar[u] \\
k \ar[r] \ar[u] & k' \ar[u]
}
$$
Ainsi cette construction conduit à une nouvelle situation avec
$[L : L_{sep}] < [K : K_{sep}]$. Par récurrence, nous pouvons trouver
$k' \subset k''$ et $L \subset L'$ comme dans le lemme pour
l'extension $L/k'$. Alors les extensions $k''/k$ et $L'/K$
conviennent pour l'extension $K/k$.
Ceci démontre le lemme.
\end{proof}

\section{Algèbres géométriquement réduites}
\label{section-geometrically-reduced}

\noindent
Le résultat principal sur les algèbres géométriquement réduites est le
Lemme \ref{lemma-geometrically-reduced-finite-purely-inseparable-extension}.
Nous suggérons au lecteur de passer directement à ce lemme après avoir lu la définition.

\begin{definition}
\label{definition-geometrically-reduced}
Soit $k$ un corps. Soit $S$ une $k$-algèbre.
On dit que $S$ est {\it géométriquement réduite sur $k$}
si, pour toute extension de corps $K/k$, la
$K$-algèbre $K \otimes_k S$ est réduite.
\end{definition}

\noindent
Soient $k$ un corps et $S$ une $k$-algèbre réduite.
Pour vérifier que $S$ est géométriquement réduite, il suffit
de vérifier que $\overline{k} \otimes_k S$ est réduite (où
$\overline{k}$ désigne la clôture algébrique de $k$).
En fait, il suffit de le vérifier pour les extensions de corps finies
purement inséparables $k'/k$. Voir le
Lemme \ref{lemma-geometrically-reduced-finite-purely-inseparable-extension}.

\begin{lemma}
\label{lemma-subalgebra-separable}
Propriétés élémentaires des algèbres géométriquement réduites.
Soit $k$ un corps. Soit $S$ une $k$-algèbre.
\begin{enumerate}
\item Si $S$ est géométriquement réduite sur $k$, toute
$k$-sous-algèbre.
\item Si toutes les $k$-sous-algèbres de type fini de $S$ sont
géométriquement réduites, alors $S$ est géométriquement réduite.
\item Une colimite filtrante de $k$-algèbres géométriquement réduites
est géométriquement réduite.
\item Si $S$ est géométriquement réduite sur $k$, toute localisation
de $S$ est géométriquement réduite sur $k$.
\end{enumerate}
\end{lemma}

\begin{proof}
Omis. Les deuxième et troisième propriétés résultent du fait que
le produit tensoriel commute aux colimites.
\end{proof}

\begin{lemma}
\label{lemma-geometrically-reduced-permanence}
Soit $k$ un corps.
Si $R$ est géométriquement réduite sur $k$
et si $S \subset R$ est une partie multiplicative, alors la localisation
$S^{-1}R$ est géométriquement réduite sur $k$.
Si $R$ est géométriquement réduite sur $k$, alors $R[x]$ est
géométriquement réduite sur $k$.
\end{lemma}

\begin{proof}
Omis. Indications : une localisation d'un anneau réduit est réduite,
et la localisation commute aux produits tensoriels.
\end{proof}

\noindent
Dans les preuves des lemmes suivants, nous utiliserons à plusieurs reprises
l'observation suivante : supposons que $R' \subset R$ et
$S' \subset S$ soient des inclusions de $k$-algèbres.
Alors l'application $R' \otimes_k S' \to R \otimes_k S$
est injective.

\begin{lemma}
\label{lemma-limit-argument}
Soit $k$ un corps. Soient $R$ et $S$ des $k$-algèbres.
\begin{enumerate}
\item Si $R \otimes_k S$ n'est pas réduit, il existe des
sous-algèbres de type fini $R' \subset R$ et
$S' \subset S$ telles que $R' \otimes_k S'$ ne soit pas réduit.
\item Si $R \otimes_k S$ contient un diviseur de zéro non nul, il existe des
sous-algèbres de type fini $R' \subset R$ et
$S' \subset S$ telles que $R' \otimes_k S'$ contienne un diviseur de zéro non nul.
\item Si $R \otimes_k S$ contient un idempotent non trivial, il existe des
sous-algèbres de type fini $R' \subset R$ et
$S' \subset S$ telles que $R' \otimes_k S'$ contienne un idempotent non trivial.
\end{enumerate}
\end{lemma}

\begin{proof}
Supposons que $z \in R \otimes_k S$ soit nilpotent. On peut écrire
$z = \sum_{i = 1, \ldots, n} x_i \otimes y_i$.
On peut donc prendre pour $R'$ la $k$-sous-algèbre engendrée par
les $x_i$, et pour $S'$ la $k$-sous-algèbre engendrée par les $y_i$.
Les deuxième et troisième assertions se démontrent de la même manière.
\end{proof}

\begin{lemma}
\label{lemma-geometrically-reduced-any-reduced-base-change}
Soit $k$ un corps.
Soit $S$ une $k$-algèbre géométriquement réduite.
Soit $R$ une $k$-algèbre réduite quelconque.
Alors $R \otimes_k S$ est réduite.
\end{lemma}

\begin{proof}
D'après le Lemme \ref{lemma-limit-argument},
nous pouvons supposer que $R$ est de type fini sur $k$.
Alors, comme anneau réduit noethérien, $R$ se plonge dans un produit
fini de corps
(voir les Lemmes \ref{lemma-total-ring-fractions-no-embedded-points},
\ref{lemma-Noetherian-irreducible-components} et
\ref{lemma-minimal-prime-reduced-ring}).
Nous pouvons donc supposer que $R$ est un produit fini de corps.
Dans ce cas, la
Définition \ref{definition-geometrically-reduced}
montre que $R \otimes_k S$ est réduite.
\end{proof}

\begin{lemma}
\label{lemma-separable-extension-preserves-reducedness}
Soit $k$ un corps.
Soit $S$ une $k$-algèbre réduite.
Soit $K/k$ une extension de corps séparable
ou une extension de corps séparablement engendrée.
Alors $K \otimes_k S$ est réduite.
\end{lemma}

\begin{proof}
Supposons que $k \subset K$ soit séparable.
D'après le Lemme \ref{lemma-limit-argument},
nous pouvons supposer que $S$ est de type fini sur $k$
et que $K$ est de type fini sur $k$.
Alors $S$ se plonge dans un produit fini de corps,
à savoir son anneau total des fractions (voir les
Lemmes \ref{lemma-minimal-prime-reduced-ring} et
\ref{lemma-total-ring-fractions-no-embedded-points}).
Nous pouvons donc en fait supposer que $S$ est un domaine.
Choisissons $x_1, \ldots, x_{r + 1} \in K$ comme dans le
Lemme \ref{lemma-generating-finitely-generated-separable-field-extensions}.
Soit $P \in k(x_1, \ldots, x_r)[T]$
le polynôme minimal de $x_{r + 1}$. C'est un polynôme séparable.
Il est facile de voir que
$k[x_1, \ldots, x_r] \otimes_k S = S[x_1, \ldots, x_r]$ est un domaine.
Il s'ensuit que $k(x_1, \ldots, x_r) \otimes_k S$ est un domaine,
car c'est une localisation de $S[x_1, \ldots, x_r]$.
L'extension d'anneaux
$k(x_1, \ldots, x_r) \otimes_k S \subset K \otimes_k S$
est engendrée par un seul élément $x_{r + 1}$ qui satisfait une seule
équation, à savoir $P$. Par conséquent, $K \otimes_k S$ se plonge dans
$F[T]/(P)$, où $F$ est le corps des fractions de
$k(x_1, \ldots, x_r) \otimes_k S$.
Comme $P$ est séparable, ceci est un produit fini de corps, et la conclusion s'ensuit.

\medskip\noindent
À ce stade, nous ne savons pas encore qu'une extension de corps
séparablement engendrée est séparable ; il faut donc aussi démontrer
le lemme dans ce cas.
Pour cela, supposons que $\{x_i\}_{i \in I}$ soit une base de
transcendance séparante de $K$ sur $k$. Pour tout ensemble fini
d'éléments $\lambda_j \in K$, il existe un sous-ensemble fini
$T \subset I$ tel que
$k(\{x_i\}_{i\in T}) \subset k(\{x_i\}_{i \in T} \cup \{\lambda_j\})$
soit finie et séparable. Il s'ensuit que $K$ est une colimite filtrante
d'extensions de $k$ de type fini et séparablement engendrées.
L'argument du paragraphe précédent s'applique donc également à ce cas.
\end{proof}

\begin{lemma}
\label{lemma-generic-points-geometrically-reduced}
Soit $k$ un corps et soit $S$ une $k$-algèbre. Supposons que
$S$ soit réduite et que $S_{\mathfrak p}$ soit géométriquement
réduite pour tout idéal premier minimal $\mathfrak p$ de $S$.
Alors $S$ est géométriquement réduite.
\end{lemma}

\begin{proof}
Comme $S$ est réduite, l'application
$S \to \prod_{\mathfrak p\text{ minimal}} S_{\mathfrak p}$
est injective, voir le
Lemme \ref{lemma-reduced-ring-sub-product-fields}.
Si $K/k$ est une extension de corps, alors les applications
$$
S \otimes_k K \to (\prod S_\mathfrak p) \otimes_k K \to
\prod S_\mathfrak p \otimes_k K
$$
sont injectives : la première parce que $k \to K$ est plate, et la seconde
par inspection puisque $K$ est un $k$-module libre. Comme $S_\mathfrak p$ est
géométriquement réduite, l'anneau de droite est réduit. Ainsi $S \otimes_k K$
est réduit comme sous-anneau d'un anneau réduit.
\end{proof}

\begin{lemma}
\label{lemma-separable-algebraic-diagonal}
Soit $k'/k$ une extension algébrique séparable.
Alors il existe une partie multiplicative $S \subset k' \otimes_k k'$
telle que l'application de multiplication $k' \otimes_k k' \to k'$
s'identifie à $k' \otimes_k k' \to S^{-1}(k' \otimes_k k')$.
\end{lemma}

\begin{proof}
Supposons d'abord que $k'/k$ soit finie et séparable. Alors
$k' = k(\alpha)$, voir le lemme \ref{fields-lemma-primitive-element} de Corps.
Soit $P \in k[x]$ le polynôme minimal de $\alpha$ sur $k$.
Alors $P$ est un polynôme irréductible, séparable et unitaire, voir
la section \ref{fields-section-separable-extensions} de Corps.
Alors $k'[x]/(P) \to k' \otimes_k k'$,
$\sum \alpha_i x^i \mapsto \alpha_i \otimes \alpha^i$ est un isomorphisme.
On peut factoriser $P = (x - \alpha) Q$ dans $k'[x]$ et, puisque $P$
est séparable, on a $Q(\alpha) \not = 0$.
Il est alors clair que la partie multiplicative $S'$ engendrée par
$Q$ dans $k'[x]/(P)$ convient, c'est-à-dire
$k' = (S')^{-1}(k'[x]/(P))$.
Par transport de structure, l'image $S$ de $S'$ dans $k' \otimes_k k'$
convient.

\medskip\noindent
Dans le cas général, écrivons $k' = \bigcup k_i$ comme réunion
de ses sous-extensions finies sur $k$. Pour chaque $i$, il existe une
partie multiplicative $S_i \subset k_i \otimes_k k_i$
telle que $k_i = S_i^{-1}(k_i \otimes_k k_i)$. Soit $S$ la clôture
multiplicative de $\bigcup S_i \subset k' \otimes_k k'$.
Il est clair que l'application de multiplication envoie chaque élément de $S$
sur un élément inversible de $k'$. Comme $k' \otimes_k k'$ est la réunion
des anneaux $k_i \otimes_k k_i$, tout élément du noyau de l'application
de multiplication est envoyé sur zéro dans $S^{-1}(k' \otimes_k k')$.
Par exactitude de la localisation, le résultat s'ensuit.
\end{proof}

\begin{lemma}
\label{lemma-geometrically-reduced-over-separable-algebraic}
Soit $k'/k$ une extension de corps algébrique séparable.
Soit $A$ une algèbre sur $k'$. Alors $A$ est géométriquement
réduite sur $k$ si et seulement si elle est géométriquement réduite sur $k'$.
\end{lemma}

\begin{proof}
Supposons que $A$ soit géométriquement réduite sur $k'$.
Soit $K/k$ une extension de corps. Alors $K \otimes_k k'$ est
un anneau réduit d'après le
Lemme \ref{lemma-separable-extension-preserves-reducedness}.
Ainsi, d'après le
Lemme \ref{lemma-geometrically-reduced-any-reduced-base-change},
nous obtenons que $K \otimes_k A = (K \otimes_k k') \otimes_{k'} A$ est réduit.

\medskip\noindent
Supposons que $A$ soit géométriquement réduite sur $k$. Soit $K/k'$ une extension
de corps. Alors
$$
K \otimes_{k'} A = (K \otimes_k A) \otimes_{(k' \otimes_k k')} k'
$$
Comme $k' \otimes_k k' \to k'$ est une localisation d'après le
Lemme \ref{lemma-separable-algebraic-diagonal}, nous voyons que
$K \otimes_{k'} A$ est une localisation d'une algèbre réduite, donc réduite.
\end{proof}

\section{Extensions séparables, suite}
\label{section-separability-continued}

\noindent
Dans cette section, nous poursuivons la discussion commencée dans la
section \ref{section-separability}.

\begin{lemma}
\label{lemma-mini-separability}
Soit $k$ un corps de caractéristique $p > 1$. Soit $K/k$
une extension de corps engendrée par $x_1, \ldots, x_{n + 1} \in K$ telle que
\begin{enumerate}
\item $\{x_1, \ldots, x_n\}$ soit une base de transcendance de $K/k$,
\item pour tout sous-ensemble $k$-linéairement indépendant
$\{a_1, \ldots, a_m\}$ de $K$, l'ensemble
$\{a^p_1, \ldots, a_m^p\}$ soit linéairement indépendant sur $k$.
\end{enumerate}
Alors il existe $1 \leq j \leq n+1$ tel que
$\{ x_1, \ldots, \widehat{x}_j, \ldots, x_{n+1}\}$
soit une base de transcendance séparante de $K / k$.
\end{lemma}

\begin{proof}
Par hypothèse, $x_{n + 1}$ est algébrique sur $k(x_1, \ldots, x_n)$,
et il existe donc un polynôme non nul $F \in k[X_1, \ldots, X_{n + 1}]$
tel que $F(x_1, \ldots, x_{n+1}) = 0$. Choisissons $F$ de degré total minimal.
Alors $F$ est irréductible, car au moins un facteur irréductible
doit également avoir cette propriété.

\medskip\noindent
Nous affirmons que, pour un certain $i$, toutes les puissances de $X_i$
qui apparaissent dans $F$ ne sont pas multiples de $p$. Supposons par
contradiction que tous les exposants apparaissant dans $F$ soient
des multiples de $p$ ; alors l'ensemble
$$
\{x_1^{\alpha_1} \ldots x^{\alpha_{n+1}}_{n+1} \mid \lambda_\alpha \neq 0\}
\subset
K
$$
est linéairement dépendant sur $k$, où les $\lambda_\alpha$ sont les
coefficients de $F$. Par l'hypothèse (2), nous concluons que l'ensemble
$$
\{x_1^{\alpha_1 / p} \ldots x^{\alpha_{n+1} / p}_{n+1}
\mid \lambda_\alpha \neq 0 \}
$$
est également linéairement dépendant sur $k$, ce qui contredit la minimalité
de $\deg(F)$.

\medskip\noindent
Choisissons $i$ tel qu'une puissance non-$p$-ième de $X_i$ apparaisse
dans $F$. Alors $x_i$ est algébrique sur
$L = k(x_1, \ldots, x_{i - 1}, x_{i + 1}, \ldots, x_{n+1})$.
D'après le lemme \ref{fields-lemma-transcendence-degree} de Corps,
nous voyons que $x_1, \ldots, x_{i - 1}, x_{i + 1}, \ldots, x_{n+1}$
est une base de transcendance de $K/k$. Ainsi $L$ est le corps des fractions
de l'anneau de polynômes sur $k$ en
$x_1, \ldots, x_{i - 1}, x_{i + 1}, \ldots, x_{n + 1}$.
Par le Lemme de Gauss, nous concluons que
$$
P(T) =
F(x_1, \ldots, x_{i - 1}, T, x_{i + 1}, \ldots, x_{n + 1}) \in L[T]
$$
est irréductible. Par construction, $P(T)$ n'est pas contenu dans $L[T^p]$.
Ainsi $K/L$ est séparable, comme requis.
\end{proof}

\noindent
Soit $p$ un nombre premier et soit $k$ un corps de caractéristique $p$.
Dans ce cas, nous écrivons $k^{1/p}$ pour l'extension de $k$ obtenue
en adjoignant les racines $p$-ièmes de tous les éléments de $k$
dans $k$.
(Autrement dit, c'est le sous-corps d'une clôture algébrique de
$k$ engendré par les racines $p$-ièmes des éléments de $k$.)

\begin{lemma}
\label{lemma-characterize-separable-field-extensions}
Soit $k$ un corps de caractéristique $p > 0$.
Soit $K/k$ une extension de corps.
Les assertions suivantes sont équivalentes :
\begin{enumerate}
\item $K$ est séparable sur $k$,
\item pour tout sous-ensemble $k$-linéairement indépendant
$\{a_1, \ldots, a_m\}$ de $K$, l'ensemble $\{a^p_1, \ldots, a_m^p\}$ est linéairement
indépendant sur $k$,
\item l'anneau $K \otimes_k k^{1/p}$ est réduit, et
\item $K$ est géométriquement réduit sur $k$.
\end{enumerate}
\end{lemma}

\begin{proof}
L'implication (1) $\Rightarrow$ (4) résulte du
Lemme \ref{lemma-separable-extension-preserves-reducedness}.
L'implication (4) $\Rightarrow$ (3) est immédiate.

\medskip\noindent
Supposons (3). Considérons l'homomorphisme d'anneaux
$m : K \otimes_k k^{1/p} \rightarrow K$ défini par
$$
\lambda \otimes \mu \rightarrow \lambda^p \mu^p
$$
Remarquons que $x^p = m(x) \otimes 1$ pour tout $x \in K \otimes_k k^{1/p}$.
Comme $K \otimes_k k^{1/p}$ est réduit, nous voyons que $m$ est injectif.
Si $\{a_1, \ldots, a_m\} \subset K$ est linéairement indépendant sur $k$,
alors $\{a_1 \otimes 1, \ldots, a_m \otimes 1\}$ est linéairement
indépendant sur $k^{1/p}$. Par injectivité de $m$, nous déduisons
qu'aucune combinaison $k$-linéaire non triviale de
$a_1^p, \ldots, a_m^p$ n'est nulle. Ainsi (3) implique (2).

\medskip\noindent
Supposons (2). Pour démontrer (1), nous pouvons supposer que $K$ est
de type fini sur $k$, et nous devons montrer que $K$ est séparablement
engendrée sur $k$.
Soit $\{x_1, \ldots, x_d\}$ une base de transcendance de $K/k$.
D'après le lemme \ref{fields-lemma-algebraic-finitely-generated} de Corps,
nous avons $[K : K'] < \infty$, où $K' = k(x_1, \ldots, x_d)$.
Choisissons la base de transcendance de sorte que le degré d'inséparabilité
$[K : K']_i$ soit minimal. Si $K / K'$ est séparable, la conclusion est acquise.
Supposons le contraire afin d'obtenir une contradiction. Alors il existe
$x_{d + 1} \in K$ qui n'est pas séparable sur $K'$ et, en particulier,
$[K'(x_{d+1}) : K']_i > 1$. Alors, d'après le
Lemme \ref{lemma-mini-separability}, il existe $1 \leq j \leq n + 1$
tel que $K'' = k(x_1, \ldots, \widehat{x}_j, \ldots, x_{d+1})$
satisfasse $[K'(x_{d+1}) : K'']_i = 1$. Par multiplicativité,
$[K : K'']_i < [K : K']_i$, ce qui donne la contradiction.
\end{proof}

\begin{lemma}
\label{lemma-separably-generated-separable}
Une extension de corps séparablement engendrée est séparable.
\end{lemma}

\begin{proof}
Combiner le Lemme \ref{lemma-separable-extension-preserves-reducedness}
avec le Lemme \ref{lemma-characterize-separable-field-extensions}.
\end{proof}

\noindent
Dans le lemme suivant, nous utiliserons la notion de clôture parfaite,
qui est définie dans la
Définition \ref{definition-perfection}.

\begin{lemma}
\label{lemma-geometrically-reduced-finite-purely-inseparable-extension}
Soit $k$ un corps. Soit $S$ une $k$-algèbre.
Les assertions suivantes sont équivalentes :
\begin{enumerate}
\item $k' \otimes_k S$ est réduit pour toute extension finie
purement inséparable $k'$ de $k$,
\item $k^{1/p} \otimes_k S$ est réduit,
\item $k^{perf} \otimes_k S$ est réduit, où $k^{perf}$ est la
clôture parfaite de $k$,
\item $\overline{k} \otimes_k S$ est réduit, où $\overline{k}$ est la
clôture algébrique de $k$, et
\item $S$ est géométriquement réduite sur $k$.
\end{enumerate}
\end{lemma}

\begin{proof}
Remarquons que toute extension finie purement inséparable $k'/k$ se plonge
dans $k^{perf}$. De plus, $k^{1/p}$ se plonge dans $k^{perf}$, qui se
plonge lui-même dans $\overline{k}$. Il est donc clair que
(5) $\Rightarrow$ (4) $\Rightarrow$ (3) $\Rightarrow$ (2)
et que (3) $\Rightarrow$ (1).

\medskip\noindent
Prouvons que (1) $\Rightarrow$ (5).
Supposons que $k' \otimes_k S$ soit réduit pour toute extension finie
purement inséparable $k'$ de $k$. Soit $K/k$ une extension de corps.
Nous devons montrer que $K \otimes_k S$ est réduit. D'après le
Lemme \ref{lemma-limit-argument}, nous nous ramenons au cas où
$K/k$ est une extension de corps de type fini. Choisissons un diagramme
$$
\xymatrix{
K \ar[r] & K' \\
k \ar[u] \ar[r] & k' \ar[u]
}
$$
comme dans le Lemme \ref{lemma-make-separably-generated}.
Par hypothèse, $k' \otimes_k S$ est réduit.
D'après le Lemme \ref{lemma-separable-extension-preserves-reducedness},
il s'ensuit que $K' \otimes_k S$ est réduit.
Nous concluons donc que $K \otimes_k S$ est réduit, comme voulu.

\medskip\noindent
Enfin, démontrons que (2) $\Rightarrow$ (5).
Supposons que $k^{1/p} \otimes_k S$ soit réduit. Alors $S$ est réduite.
De plus, pour chaque localisation $S_{\mathfrak p}$ en un idéal premier
minimal $\mathfrak p$, l'anneau $k^{1/p}\otimes_k S_{\mathfrak p}$
est une localisation de $k^{1/p} \otimes_k S$, donc est réduit.
Mais $S_{\mathfrak p}$ est un corps d'après le
Lemme \ref{lemma-minimal-prime-reduced-ring},
donc $S_{\mathfrak p}$ est géométriquement réduit d'après le
Lemme \ref{lemma-characterize-separable-field-extensions}.
Il s'ensuit du Lemme \ref{lemma-generic-points-geometrically-reduced}
que $S$ est géométriquement réduite.
\end{proof}

\section{Corps parfaits}
\label{section-perfect-fields}

\noindent
Voici la définition.

\begin{definition}
\label{definition-perfect}
Soit $k$ un corps. On dit que $k$ est {\it parfait}
si toute extension de corps de $k$ est séparable sur $k$.
\end{definition}

\begin{lemma}
\label{lemma-perfect}
Un corps $k$ est parfait si et seulement s'il est de caractéristique $0$
ou de caractéristique $p > 0$ et si tout élément possède une
racine $p$-ième.
\end{lemma}

\begin{proof}
Le cas de caractéristique nulle est clair.
Supposons que la caractéristique de $k$ soit $p > 0$.
Si $k$ est parfait, toutes les extensions de corps obtenues en adjoignant
une racine $p$-ième d'un élément de $k$ doivent être triviales ; tout
élément de $k$ possède donc une racine $p$-ième. Réciproquement, si tout
élément possède une racine $p$-ième, alors $k = k^{1/p}$ et toute extension
de corps de $k$ est séparable d'après le
Lemme \ref{lemma-characterize-separable-field-extensions}.
\end{proof}

\begin{lemma}
\label{lemma-make-separable}
Soit $K/k$ une extension de corps de type fini.
Il existe un diagramme
$$
\xymatrix{
K \ar[r] & K' \\
k \ar[u] \ar[r] & k' \ar[u]
}
$$
où $k'/k$ et $K'/K$ sont des extensions de corps finies
purement inséparables telles que $K'/k'$ soit une extension de corps séparable.
Dans cette situation, on peut supposer que $K' = k'K$ est le compositum,
et aussi que $K' = (k' \otimes_k K)_{red}$.
\end{lemma}

\begin{proof}
D'après le Lemme \ref{lemma-make-separably-generated},
nous pouvons trouver un tel diagramme avec $K'/k'$ séparablement engendrée.
D'après le
Lemme \ref{lemma-separably-generated-separable},
il s'ensuit que $K'$ est séparable sur $k'$.
Le compositum $k'K$ est une sous-extension de $K'/k'$ et donc
$k' \subset k'K$ est séparable d'après le
Lemme \ref{lemma-subextensions-are-separable}.
L'anneau $(k' \otimes_k K)_{red}$ est un domaine, car pour un certain
$n \gg 0$ l'application $x \mapsto x^{p^n}$ l'envoie dans $K$.
Il est donc un corps d'après le
Lemme \ref{lemma-integral-over-field}.
Ainsi $(k' \otimes_k K)_{red} \to K'$ l'envoie isomorphiquement sur $k'K$.
\end{proof}

\begin{lemma}
\label{lemma-perfection}
\begin{slogan}
Tout corps possède une clôture parfaite unique.
\end{slogan}
Pour tout corps $k$, il existe une extension purement inséparable
$k'/k$ telle que $k'$ soit parfaite. L'extension de corps
$k'/k$ est unique à isomorphisme unique près.
\end{lemma}

\begin{proof}
Si la caractéristique de $k$ est nulle, alors $k' = k$ est l'unique choix.
Supposons que la caractéristique de $k$ soit $p > 0$.
Pour tout $n > 0$, il existe une unique extension algébrique
$k \subset k^{1/p^n}$ telle que (a) tout élément $\lambda \in k$
possède une racine $p^n$-ième dans $k^{1/p^n}$ et (b) pour tout élément
$\mu \in k^{1/p^n}$, on ait $\mu^{p^n} \in k$.
C'est-à-dire, considérons l'application d'anneaux
$k \to k^{1/p^n} = k$, $x \mapsto x^{p^n}$.
Elle est injective et satisfait (a) et (b). Il est clair que
$k^{1/p^n} \subset k^{1/p^{n + 1}}$ comme extensions de $k$ via
l'application $y \mapsto y^p$. Nous pouvons donc prendre
$k' = \bigcup k^{1/p^n}$.
Quelques détails sont omis.
\end{proof}

\begin{definition}
\label{definition-perfection}
Soit $k$ un corps. L'extension de corps $k'/k$ du
Lemme \ref{lemma-perfection} s'appelle la {\it clôture parfaite} de $k$.
Notation : $k^{perf}/k$.
\end{definition}

\noindent
Remarquons que si $k'/k$ est une extension algébrique purement inséparable,
alors $k'$ est une sous-extension de $k^{perf}$, c'est-à-dire
$k^{perf}/k'/k$. En effet, $(k')^{perf}$ est isomorphe à $k^{perf}$
par l'unicité du
Lemme \ref{lemma-perfection}.

\begin{lemma}
\label{lemma-perfect-reduced}
Soit $k$ un corps parfait.
Toute $k$-algèbre réduite est géométriquement réduite sur $k$.
Soient $R$ et $S$ des $k$-algèbres.
Supposons que $R$ et $S$ soient toutes deux réduites.
Alors la $k$-algèbre $R \otimes_k S$ est réduite.
\end{lemma}

\begin{proof}
La première assertion résulte du
Lemme \ref{lemma-geometrically-reduced-finite-purely-inseparable-extension}.
Pour la seconde, utiliser la première assertion et le
Lemme \ref{lemma-geometrically-reduced-any-reduced-base-change}.
\end{proof}

\section{Homéomorphismes universels}
\label{section-universal-homeomorphism}

\noindent
Soit $k'/k$ une extension de corps algébrique purement inséparable.
Alors, pour toute $k$-algèbre $R$, l'application d'anneaux
$R \to k' \otimes_k R$ induit un homéomorphisme des spectres.
La raison en est le
Lemme \ref{lemma-p-ring-map} un peu plus général ci-dessous.

\begin{lemma}
\label{lemma-surjective-locally-nilpotent-kernel}
Soit $\varphi : R \to S$ un homomorphisme surjectif à noyau localement nilpotent.
Alors $\varphi$ induit un homéomorphisme des spectres et des isomorphismes
sur les corps résiduels. Pour toute application d'anneaux $R \to R'$, l'application
$R' \to R' \otimes_R S$ est surjective à noyau localement nilpotent.
\end{lemma}

\begin{proof}
D'après le Lemme \ref{lemma-spec-closed}, l'application
$\Spec(S) \to \Spec(R)$ est un homéomorphisme sur le sous-ensemble fermé
$V(\Ker(\varphi))$. Bien sûr, $V(\Ker(\varphi)) = \Spec(R)$, car tout idéal
premier de $R$ contient tout élément nilpotent de $R$. Cela implique également
l'assertion sur les corps résiduels. Par l'exactitude à droite du produit
tensoriel, nous voyons que $\Ker(\varphi)R'$ est le noyau de l'application
surjective $R' \to R' \otimes_R S$. La dernière assertion résulte donc du
Lemme \ref{lemma-locally-nilpotent}.
\end{proof}

\begin{lemma}
\label{lemma-powers-field}
\begin{reference}
\cite[Lemma 3.1.6]{Alper-adequate}
\end{reference}
Soit $k'/k$ une extension de corps. Les assertions suivantes sont équivalentes :
\begin{enumerate}
\item pour tout $x \in k'$, il existe un $n > 0$ tel que $x^n \in k$, et
\item $k' = k$, ou bien $k$ et $k'$ sont de caractéristique $p > 0$ et
soit $k'/k$ est une extension purement inséparable, soit
$k$ et $k'$ sont des extensions algébriques de $\mathbf{F}_p$.
\end{enumerate}
\end{lemma}

\begin{proof}
Remarquons que chacune des possibilités énumérées en (2) satisfait (1).
Supposons donc que $k'/k$ satisfasse (1), et montrons que l'un des cas de (2)
se produit. En écartant le cas $k = k'$, nous pouvons supposer que
$k' \not = k$. Il est clair que $k'/k$ est algébrique.
Nous pouvons donc supposer que $k'/k$ est une extension finie non triviale.
Soit $k'/k'_{sep}/k$ la sous-extension séparable
trouvée dans le lemme \ref{fields-lemma-separable-first} de Corps.
Nous devons montrer que $k = k'_{sep}$ ou que $k$ est algébrique sur
$\mathbf{F}_p$. Nous pouvons donc supposer que $k'/k$
est une extension finie séparable non triviale et nous devons montrer
que $k$ est algébrique sur $\mathbf{F}_p$.

\medskip\noindent
Choisissons $x \in k'$, $x \not \in k$. Choisissons $n, m > 0$ tels que
$x^n \in k$ et $(x + 1)^m \in k$. Soit $\overline{k}$ une clôture
algébrique de $k$. Nous pouvons choisir des plongements
$\sigma, \tau : k' \to \overline{k}$ tels que $\sigma(x) \not = \tau(x)$.
Cela résulte de la discussion de la
section \ref{fields-section-separable-extensions} de Corps
(plus précisément, après avoir remplacé $k'$ par l'extension de $k$
engendrée par $x$, cela résulte du
lemme \ref{fields-lemma-count-embeddings} de Corps).
Nous voyons alors que $\sigma(x) = \zeta \tau(x)$ pour une
racine $n$-ième de l'unité $\zeta$ dans $\overline{k}$.
De même, $\sigma(x + 1) = \zeta' \tau(x + 1)$
pour une racine $m$-ième de l'unité $\zeta' \in \overline{k}$.
Comme $\sigma(x + 1) \not = \tau(x + 1)$, nous voyons que $\zeta' \not = 1$.
Alors
$$
\zeta' (\tau(x) + 1) =
\zeta' \tau(x + 1) =
\sigma(x + 1) =
\sigma(x) + 1 =
\zeta \tau(x) + 1
$$
implique que
$$
\tau(x) (\zeta' - \zeta) = 1 - \zeta'
$$
d'où $\zeta' \not = \zeta$ et
$$
\tau(x) = (1 - \zeta')/(\zeta' - \zeta)
$$
Ainsi, tout élément de $k'$ qui n'appartient pas à $k$ est algébrique
sur le sous-corps premier. Comme $k'$ est engendré sur le sous-corps premier
par les éléments de $k'$ qui n'appartiennent pas à $k$, nous concluons que
$k'$ (et donc $k$) est algébrique sur le sous-corps premier.

\medskip\noindent
Enfin, si la caractéristique de $k$ est $0$, ce qui précède conduit à une
contradiction comme suit (nous encourageons le lecteur à trouver sa propre preuve).
Pour tout nombre rationnel $y$, nous obtenons de même une racine de l'unité
$\zeta_y$ telle que $\sigma(x + y) = \zeta_y\tau(x + y)$.
Nous trouvons alors
$$
\zeta \tau(x) + y = \zeta_y(\tau(x) + y)
$$
et, grâce à la formule ci-dessus pour $\tau(x)$, nous concluons que
$\zeta_y \in \mathbf{Q}(\zeta, \zeta')$. Comme le corps de nombres
$\mathbf{Q}(\zeta, \zeta')$ ne contient qu'un nombre fini de racines de
l'unité, il existe deux nombres rationnels distincts $y, y'$ tels que
$\zeta_y = \zeta_{y'}$. Nous concluons alors que
$$
y - y' =
\sigma(x + y) - \sigma(x + y') =
\zeta_y(\tau(x + y)) - \zeta_{y'}\tau(x + y') = \zeta_y(y - y')
$$
ce qui implique $\zeta_y = 1$, contradiction.
\end{proof}

\begin{lemma}
\label{lemma-powers}
Soit $\varphi : R \to S$ une application d'anneaux. Supposons que
\begin{enumerate}
\item pour tout $x \in S$, il existe $n > 0$ tel que
$x^n$ appartienne à l'image de $\varphi$, et
\item $\Ker(\varphi)$ soit localement nilpotent,
\end{enumerate}
alors $\varphi$ induit un homéomorphisme des spectres et induit des
extensions de corps résiduels satisfaisant aux conditions équivalentes du
Lemme \ref{lemma-powers-field}.
\end{lemma}

\begin{proof}
Supposons (1) et (2). Soient $\mathfrak q, \mathfrak q'$ des idéaux premiers
de $S$ au-dessus du même idéal premier $\mathfrak p$ de $R$. Supposons que
$x \in S$ vérifie $x \in \mathfrak q$, $x \not \in \mathfrak q'$. Alors
$x^n \in \mathfrak q$ et $x^n \not \in \mathfrak q'$ pour tout $n > 0$. Si
$x^n = \varphi(y)$ avec $y \in R$ pour un certain $n > 0$, alors
$$
x^n \in \mathfrak q \Rightarrow y \in \mathfrak p \Rightarrow
x^n \in \mathfrak q'
$$
ce qui est une contradiction. Il n'existe donc pas de $x$ comme ci-dessus,
et nous concluons que $\mathfrak q = \mathfrak q'$, c'est-à-dire que
l'application sur les spectres est injective. D'après (2), le noyau
$I = \Ker(\varphi)$ est contenu dans tout idéal premier, donc
$\Spec(R) = \Spec(R/I)$ comme espaces topologiques. Comme l'application
induite $R/I \to S$ est entière par (1), le
Lemme \ref{lemma-integral-overring-surjective} montre que
$\Spec(S) \to \Spec(R/I)$ est surjective. En combinant ce qui précède, nous
voyons que $\Spec(S) \to \Spec(R)$ est bijective.
Si $x \in S$ est quelconque et si nous choisissons $y \in R$ tel que
$\varphi(y) = x^n$ pour un certain $n > 0$, alors l'ouvert
$D(x) \subset \Spec(S)$ correspond à l'ouvert
$D(y) \subset \Spec(R)$ par la bijection ci-dessus. Ainsi
l'application $\Spec(S) \to \Spec(R)$ est un homéomorphisme.

\medskip\noindent
Pour l'assertion sur les corps résiduels, soit
$\mathfrak q \subset S$ un idéal premier au-dessus d'un idéal premier
$\mathfrak p \subset R$. Soit $x \in \kappa(\mathfrak q)$. Si nous considérons
$\kappa(\mathfrak q)$ comme le corps résiduel de l'anneau local
$S_\mathfrak q$, nous voyons que $x$ est l'image d'un certain
$y/z \in S_\mathfrak q$ avec $y \in S$, $z \in S$, $z \not \in \mathfrak q$.
Choisissons $n, m > 0$ tels que $y^n, z^m$ appartiennent à l'image de $\varphi$.
Alors $x^{nm}$ est le résidu de $(y/z)^{nm} = (y^n)^m/(z^m)^n$,
qui appartient à l'image de $R_\mathfrak p \to S_\mathfrak q$.
Ainsi $x^{nm}$ appartient à l'image de
$\kappa(\mathfrak p) \to \kappa(\mathfrak q)$.
\end{proof}

\begin{lemma}
\label{lemma-2-3-ring-map}
Soit $\varphi : R \to S$ une application d'anneaux. Supposons que
\begin{enumerate}
\item[(a)] $S$ soit engendrée comme $R$-algèbre par des éléments $x$ tels que
$x^2, x^3 \in \varphi(R)$, et
\item[(b)] $\Ker(\varphi)$ soit localement nilpotent,
\end{enumerate}
Alors $\varphi$ induit des isomorphismes sur les corps résiduels et
un homéomorphisme des spectres. Pour toute application d'anneaux $R \to R'$,
l'application $R' \to R' \otimes_R S$ satisfait également (a) et (b).
\end{lemma}

\begin{proof}
Supposons (a) et (b). L'application sur les spectres est fermée, car $S$ est
entière sur $R$, voir les Lemmes \ref{lemma-going-up-closed} et
\ref{lemma-integral-going-up}. L'image est dense d'après le
Lemme \ref{lemma-image-dense-generic-points}. Ainsi
$\Spec(S) \to \Spec(R)$ est surjective. Si
$\mathfrak q \subset S$ est un idéal premier au-dessus de
$\mathfrak p \subset R$, alors l'extension de corps
$\kappa(\mathfrak q)/\kappa(\mathfrak p)$ est engendrée par des éléments
$\alpha \in \kappa(\mathfrak q)$ dont le carré et le cube appartiennent à
$\kappa(\mathfrak p)$. Il est alors clair que
$\alpha \in \kappa(\mathfrak p)$ et nous obtenons
$\kappa(\mathfrak q) = \kappa(\mathfrak p)$.
Si $\mathfrak q, \mathfrak q'$ étaient deux idéaux premiers distincts
au-dessus de $\mathfrak p$, au moins l'un des générateurs $x$ de $S$ comme
dans (a) aurait des images distinctes dans
$\kappa(\mathfrak q) = \kappa(\mathfrak p)$ et
$\kappa(\mathfrak q') = \kappa(\mathfrak p)$.
Cela contredirait le fait que $x^2$ et $x^3$ ont bien la même image.
Cela montre que $\Spec(S) \to \Spec(R)$ est injective, donc un homéomorphisme
(d'après ce qui a déjà été démontré).

\medskip\noindent
Puisque $\varphi$ induit un homéomorphisme des spectres, elle est en particulier
surjective sur les spectres, propriété préservée par tout changement de base,
voir le Lemme \ref{lemma-surjective-spec-radical-ideal}.
Par conséquent, pour toute application $R \to R'$, le noyau de l'application
d'anneaux $R' \to R' \otimes_R S$ est constitué d'éléments nilpotents,
voir le Lemme \ref{lemma-image-dense-generic-points}; autrement dit, (b) est
satisfaite pour $R' \to R' \otimes_R S$.
Il est clair que (a) est préservée par changement de base.
\end{proof}

\begin{lemma}
\label{lemma-help-with-powers}
Soit $p$ un nombre premier. Soient $n, m > 0$ deux entiers. Il existe
un entier $a$ tel que
$(x + y)^{p^a}, p^a(x + y) \in \mathbf{Z}[x^{p^n}, p^nx, y^{p^m}, p^my]$.
\end{lemma}

\begin{proof}
C'est clair pour $p^a(x + y)$ dès que $a \geq n, m$.
En fait, choisissons $a \gg n, m$. Écrivons
$$
(x + y)^{p^a}  = \sum\nolimits_{i, j \geq 0, i + j = p^a}
{p^a \choose i, j} x^iy^j
$$
Pour tous $i, j \geq 0$ tels que $i + j = p^a$, écrivons
$i = q p^n + r$ avec $r \in \{0, \ldots, p^n - 1\}$ et
$j = q' p^m + r'$ avec $r' \in \{0, \ldots, p^m - 1\}$.
La condition $(x + y)^{p^a} \in \mathbf{Z}[x^{p^n}, p^nx, y^{p^m}, p^my]$
est satisfaite si
$$
p^{nr + mr'} \text{ divise } {p^a \choose i, j}
$$
Si $r = r' = 0$, la divisibilité est satisfaite. Si $r \not = 0$, nous
écrivons
$$
{p^a \choose i, j} = \frac{p^a}{i} {p^a - 1 \choose i - 1, j}
$$
Comme $r \not = 0$, le nombre rationnel $p^a/i$ a une valuation
$p$-adique au moins égale à $a - (n - 1)$ (car $i$ n'est pas divisible par
$p^n$). Ainsi ${p^a \choose i, j}$ est divisible par $p^{a - n + 1}$
dans ce cas. De même, si $r' \not = 0$, alors ${p^a \choose i, j}$ est
divisible par $p^{a - m + 1}$. Le choix
$a = np^n + mp^m + n + m$ convient.
\end{proof}

\begin{lemma}
\label{lemma-p-ring-map-field}
Soit $k'/k$ une extension de corps. Soit $p$ un nombre premier.
Les assertions suivantes sont équivalentes
\begin{enumerate}
\item $k'$ est engendré comme extension de corps de $k$ par des éléments
$x$ tels qu'il existe un $n > 0$ avec $x^{p^n} \in k$ et
$p^nx \in k$, et
\item $k = k'$ ou bien la caractéristique de $k$ et de $k'$ est $p$
et $k'/k$ est purement inséparable.
\end{enumerate}
\end{lemma}

\begin{proof}
Soit $x \in k'$. S'il existe un $n > 0$ tel que $x^{p^n} \in k$ et
$p^nx \in k$ et si la caractéristique n'est pas $p$, alors $x \in k$.
Si la caractéristique est $p$, alors $x^{p^n} \in k$ et donc $x$ est
purement inséparable sur $k$.
\end{proof}

\begin{lemma}
\label{lemma-p-ring-map}
Soit $\varphi : R \to S$ une application d'anneaux. Soit $p$ un nombre
premier. Supposons
\begin{enumerate}
\item[(a)] que $S$ soit engendrée comme $R$-algèbre par des éléments $x$
tels qu'il existe un $n > 0$ avec $x^{p^n} \in \varphi(R)$ et
$p^nx \in \varphi(R)$, et
\item[(b)] que $\Ker(\varphi)$ soit localement nilpotent,
\end{enumerate}
Alors $\varphi$ induit un homéomorphisme des spectres et des extensions
de corps résiduels satisfaisant aux conditions équivalentes du
Lemme \ref{lemma-p-ring-map-field}. Pour toute application d'anneaux
$R \to R'$, l'application d'anneaux $R' \to R' \otimes_R S$ satisfait
également (a) et (b).
\end{lemma}

\begin{proof}
Supposons (a) et (b). Remarquons que (b) équivaut à la condition (2)
du Lemme \ref{lemma-powers}. Soit $T \subset S$ l'ensemble des
éléments $x \in S$ tels qu'il existe un entier
$n > 0$ tel que $x^{p^n} , p^n x \in \varphi(R)$.
Nous affirmons que $T = S$. Cela montrera que la condition (1) du
Lemme \ref{lemma-powers} est satisfaite et donc que $\varphi$ induit
un homéomorphisme sur les spectres.
Par (a), il suffit de montrer que $T \subset S$ est une sous-algèbre sur
$R$ de l'algèbre ambiante. Si $x \in T$ et $y \in R$, il est clair que $yx \in T$.
Supposons $x, y \in T$ et $n, m > 0$ tels que
$x^{p^n}, y^{p^m}, p^n x, p^m y \in \varphi(R)$.
Alors $(xy)^{p^{n + m}}, p^{n + m}xy \in \varphi(R)$,
d'où $xy \in T$. Nous avons $x + y \in T$ par le
Lemme \ref{lemma-help-with-powers}, ce qui prouve l'affirmation.

\medskip\noindent
Puisque $\varphi$ induit un homéomorphisme sur les spectres, elle est en
particulier surjective sur les spectres, propriété préservée par tout
changement de base, voir le Lemme \ref{lemma-surjective-spec-radical-ideal}.
Par conséquent, pour toute application $R \to R'$, le noyau de
l'application d'anneaux $R' \to R' \otimes_R S$ est constitué d'éléments
nilpotents, voir le Lemme \ref{lemma-image-dense-generic-points};
autrement dit, (b) est satisfaite pour $R' \to R' \otimes_R S$.
Il est clair que (a) est préservée par changement de base.
Enfin, la condition sur les corps résiduels découle de (a), car les
générateurs de $S$ comme $R$-algèbre donnent des générateurs des
extensions de corps résiduels.
\end{proof}

\begin{lemma}
\label{lemma-radicial}
Soit $\varphi : R \to S$ une application d'anneaux. Supposons que
\begin{enumerate}
\item $\varphi$ induise une application injective sur les spectres,
\item $\varphi$ induise des extensions de corps résiduels purement
inséparables.
\end{enumerate}
Alors, pour toute application d'anneaux $R \to R'$, les propriétés (1) et
(2) sont vraies pour $R' \to R' \otimes_R S$.
\end{lemma}

\begin{proof}
Posons $S' = R' \otimes_R S$, de sorte que nous ayons un diagramme
commutatif d'applications continues de spectres d'anneaux
$$
\xymatrix{
\Spec(S') \ar[r] \ar[d] & \Spec(S) \ar[d] \\
\Spec(R') \ar[r] & \Spec(R)
}
$$
Soit $\mathfrak p' \subset R'$ un idéal premier au-dessus de
$\mathfrak p \subset R$. S'il n'existe aucun idéal premier de $S$
au-dessus de $\mathfrak p$, alors il n'existe aucun idéal premier de
$S'$ au-dessus de $\mathfrak p'$. Sinon, d'après la
Remarque \ref{remark-fundamental-diagram}, il existe un unique idéal
premier $\mathfrak r$ de $F = S \otimes_R \kappa(\mathfrak p)$ dont le
corps résiduel est une extension purement inséparable de
$\kappa(\mathfrak p)$. Considérons les applications d'anneaux
$$
\kappa(\mathfrak p) \to F \to \kappa(\mathfrak r)
$$
Par le Lemme \ref{lemma-minimal-prime-reduced-ring}, l'idéal
$\mathfrak r \subset F$ est localement nilpotent; nous pouvons donc
appliquer le Lemme \ref{lemma-surjective-locally-nilpotent-kernel}
à l'application d'anneaux $F \to \kappa(\mathfrak r)$.
Nous pouvons appliquer le Lemme \ref{lemma-p-ring-map}
à l'application $\kappa(\mathfrak p) \to \kappa(\mathfrak r)$.
Ainsi, la composée et la seconde flèche des applications
$$
\kappa(\mathfrak p') \to
\kappa(\mathfrak p') \otimes_{\kappa(\mathfrak p)} F \to
\kappa(\mathfrak p') \otimes_{\kappa(\mathfrak p)} \kappa(\mathfrak r)
$$
induisent des bijections sur les spectres et des extensions de corps
résiduels purement inséparables. Il en est donc de même pour la première
application. Comme
$$
\kappa(\mathfrak p') \otimes_{\kappa(\mathfrak p)} F =
\kappa(\mathfrak p') \otimes_{\kappa(\mathfrak p)}
\kappa(\mathfrak p) \otimes_R S =
\kappa(\mathfrak p') \otimes_R S =
\kappa(\mathfrak p') \otimes_{R'} R' \otimes_R S
$$
nous concluons par la discussion de la
Remarque \ref{remark-fundamental-diagram}.
\end{proof}

\begin{lemma}
\label{lemma-radicial-integral}
Soit $\varphi : R \to S$ une application d'anneaux. Supposons que
\begin{enumerate}
\item $\varphi$ soit entière,
\item $\varphi$ induise une application injective sur les spectres,
\item $\varphi$ induise des extensions de corps résiduels purement
inséparables.
\end{enumerate}
Alors $\varphi$ induit un homéomorphisme de $\Spec(S)$ sur un sous-ensemble
fermé de $\Spec(R)$ et, pour toute application d'anneaux
$R \to R'$, les propriétés (1), (2), (3) sont vraies pour
$R' \to R' \otimes_R S$.
\end{lemma}

\begin{proof}
L'application sur les spectres est fermée par les
Lemmes \ref{lemma-going-up-closed} et \ref{lemma-integral-going-up}.
Les propriétés sont préservées par changement de base, par les
Lemmes \ref{lemma-radicial} et \ref{lemma-base-change-integral}.
\end{proof}

\begin{lemma}
\label{lemma-radicial-integral-bijective}
Soit $\varphi : R \to S$ une application d'anneaux. Supposons que
\begin{enumerate}
\item $\varphi$ soit entière,
\item $\varphi$ induise une application bijective sur les spectres,
\item $\varphi$ induise des extensions de corps résiduels purement
inséparables.
\end{enumerate}
Alors $\varphi$ induit un homéomorphisme sur les spectres et, pour toute
application d'anneaux $R \to R'$, les propriétés (1), (2), (3) sont vraies
pour $R' \to R' \otimes_R S$.
\end{lemma}

\begin{proof}
Cela résulte des Lemmes \ref{lemma-radicial-integral} et
\ref{lemma-surjective-spec-radical-ideal}.
\end{proof}

\begin{lemma}
\label{lemma-universally-bijective}
Soit $\varphi : R \to S$ une application d'anneaux telle que
\begin{enumerate}
\item le noyau de $\varphi$ soit localement nilpotent, et
\item $S$ soit engendrée comme $R$-algèbre par des éléments $x$
tels qu'il existe un $n > 0$ et un polynôme $P(T) \in R[T]$
dont l'image dans $S[T]$ soit $(T - x)^n$.
\end{enumerate}
Alors $\Spec(S) \to \Spec(R)$ est un homéomorphisme et $R \to S$
induit des extensions de corps résiduels purement inséparables.
De plus, les conditions (1) et (2) restent vraies après tout changement
de base.
\end{lemma}

\begin{proof}
Nous pouvons remplacer $R$ par $R/\Ker(\varphi)$, voir le
Lemme \ref{lemma-surjective-locally-nilpotent-kernel}.
L'hypothèse (2) implique que $S$ est engendrée sur $R$ par des éléments
entiers sur $R$. Ainsi $R \subset S$ est une extension entière
(Lemme \ref{lemma-integral-closure-is-ring}).
En particulier, $\Spec(S) \to \Spec(R)$ est surjective et fermée
(Lemmes \ref{lemma-integral-overring-surjective},
\ref{lemma-going-up-closed} et \ref{lemma-integral-going-up}).

\medskip\noindent
Soit $x \in S$ l'un des générateurs de (2), c'est-à-dire qu'il existe un
$n > 0$ tel que $(T - x)^n \in R[T]$.
Soit $\mathfrak p \subset R$ un idéal premier.
L'anneau $\kappa(\mathfrak p) \otimes_R S$ est non nul d'après ce qui
précède et le Lemme \ref{lemma-in-image}.
Si la caractéristique de $\kappa(\mathfrak p)$ est nulle,
nous voyons que $nx \in R$ implique que $1 \otimes x$ appartient à l'image
de $\kappa(\mathfrak p) \to \kappa(\mathfrak p) \otimes_R S$.
Ainsi $\kappa(\mathfrak p) \to \kappa(\mathfrak p) \otimes_R S$
est un isomorphisme.
Si la caractéristique de $\kappa(\mathfrak p)$ est $p > 0$,
écrivons $n = p^k m$ avec $m$ premier à $p$.
Dans $\kappa(\mathfrak p) \otimes_R S[T]$, nous avons
$$
(T - 1 \otimes x)^n = ((T - 1 \otimes x)^{p^k})^m =
(T^{p^k} - 1 \otimes x^{p^k})^m
$$
et nous voyons que $mx^{p^k} \in R$. Cela implique que
$1 \otimes x^{p^k}$ appartient à l'image de
$\kappa(\mathfrak p) \to \kappa(\mathfrak p) \otimes_R S$.
Le Lemme \ref{lemma-p-ring-map} s'applique donc à
$\kappa(\mathfrak p) \to \kappa(\mathfrak p) \otimes_R S$.
Dans les deux cas, nous concluons que
$\kappa(\mathfrak p) \otimes_R S$ possède un unique idéal premier dont
le corps résiduel est une extension purement inséparable de
$\kappa(\mathfrak p)$. Par la Remarque \ref{remark-fundamental-diagram},
nous concluons que $\varphi$ est bijective sur les spectres.

\medskip\noindent
L'assertion sur le changement de base est immédiate.
\end{proof}

\section{Algèbres géométriquement irréductibles}
\label{section-algebras-over-fields}

\noindent
Une algèbre $S$ sur un corps $k$ est dite géométriquement irréductible si
l'algèbre $S \otimes_k k'$ possède un unique idéal premier minimal pour
toute extension de corps $k'/k$. Dans cette section, nous développons
quelques éléments de théorie pertinents pour cette notion.

\begin{lemma}
\label{lemma-flat-fibres-irreducible}
Soit $R \to S$ une application d'anneaux. Supposons que
\begin{enumerate}
\item[(a)] $\Spec(R)$ soit irréductible,
\item[(b)] $R \to S$ soit plate,
\item[(c)] $R \to S$ soit de présentation finie,
\item[(d)] les anneaux fibres $S \otimes_R \kappa(\mathfrak p)$ aient des
spectres irréductibles pour une famille dense d'idéaux premiers
$\mathfrak p$ de $R$.
\end{enumerate}
Alors $\Spec(S)$ est irréductible.
Cela reste vrai plus généralement si (b) $+$ (c) est remplacé par
« l'application $\Spec(S) \to \Spec(R)$ est ouverte ».
\end{lemma}

\begin{proof}
Les hypothèses (b) et (c) impliquent que l'application sur les spectres
est ouverte, voir la Proposition \ref{proposition-fppf-open}.
Le lemme résulte donc du
lemme \ref{topology-lemma-irreducible-on-top} de Topologie.
\end{proof}

\begin{lemma}
\label{lemma-separably-closed-irreducible}
Soit $k$ un corps séparablement clos.
Soient $R$ et $S$ des $k$-algèbres. Si $R$ et $S$ possèdent un unique
idéal premier minimal, il en est de même pour $R \otimes_k S$.
\end{lemma}

\begin{proof}
Soit $k \subset \overline{k}$ une clôture parfaite, voir la
Définition \ref{definition-perfection}.
D'après l'hypothèse, $\overline{k}$ est algébriquement clos.
Les applications d'anneaux $R \to R \otimes_k \overline{k}$ et
$S \to S \otimes_k \overline{k}$ ainsi que
$R \otimes_k S \to (R \otimes_k S) \otimes_k \overline{k}
= (R \otimes_k \overline{k}) \otimes_{\overline{k}} (S \otimes_k \overline{k})$
satisfont aux hypothèses du Lemme \ref{lemma-p-ring-map}.
Nous pouvons donc supposer que $k$ est algébriquement clos.

\medskip\noindent
Nous pouvons remplacer $R$ et $S$ par leurs réductions.
Nous pouvons donc supposer que $R$ et $S$ sont des domaines.
Par le Lemme \ref{lemma-perfect-reduced}, $R \otimes_k S$ est réduit.
Son spectre est donc réductible si et seulement s'il contient un diviseur
de zéro non nul. Par le Lemme \ref{lemma-limit-argument}, nous nous
ramenons au cas où $R$ et $S$ sont des domaines de type fini sur le corps
$k$ algébriquement clos.

\medskip\noindent
Remarquons que l'application d'anneaux $R \to R \otimes_k S$ est de
présentation finie et plate. De plus, pour tout idéal maximal
$\mathfrak m$ de $R$, nous avons
$(R \otimes_k S) \otimes_R R/\mathfrak m \cong S$ car
$k \cong R/\mathfrak m$ par le théorème des zéros de Hilbert,
Théorème \ref{theorem-nullstellensatz}. En outre, l'ensemble des idéaux
maximaux est dense dans le spectre de $R$ puisque $\Spec(R)$ est un
espace de Jacobson, voir le Lemme \ref{lemma-finite-type-field-Jacobson}.
Le Lemme \ref{lemma-flat-fibres-irreducible} s'applique donc à
l'application $R \to R \otimes_k S$, et nous concluons que le spectre de
$R \otimes_k S$ est irréductible, comme voulu.
\end{proof}

\begin{lemma}
\label{lemma-geometrically-irreducible}
Soit $k$ un corps.
Soit $R$ une $k$-algèbre.
Les assertions suivantes sont équivalentes
\begin{enumerate}
\item pour toute extension de corps $k'/k$, le spectre de
$R \otimes_k k'$ est irréductible,
\item pour toute extension de corps finie et séparable $k'/k$, le spectre de
$R \otimes_k k'$ est irréductible,
\item le spectre de $R \otimes_k \overline{k}$ est irréductible,
où $\overline{k}$ est la clôture algébrique séparable de $k$, et
\item le spectre de $R \otimes_k \overline{k}$ est irréductible,
où $\overline{k}$ est la clôture algébrique de $k$.
\end{enumerate}
\end{lemma}

\begin{proof}
Il est clair que (1) implique (2).

\medskip\noindent
Supposons (2) et soit $\overline{k}$ la clôture algébrique séparable de $k$.
Supposons que $\mathfrak q_i \subset R \otimes_k \overline{k}$, $i = 1, 2$,
soient deux idéaux premiers minimaux. Pour toute sous-extension finie
$\overline{k}/k'/k$, l'extension $k'/k$ est séparable et l'application
$R \otimes_k k' \to R \otimes_k \overline{k}$ est plate.
Ainsi $\mathfrak p_i = (R \otimes_k k') \cap \mathfrak q_i$ sont des
idéaux premiers minimaux (car nous avons la descente des idéaux premiers
pour les applications d'anneaux plates, par le
Lemme \ref{lemma-flat-going-down}). Nous avons donc
$\mathfrak p_1 = \mathfrak p_2$ par (2). Comme
$\overline{k} = \bigcup k'$, nous concluons que
$\mathfrak q_1 = \mathfrak q_2$. Ainsi $\Spec(R \otimes_k \overline{k})$
est irréductible.

\medskip\noindent
Supposons (3) et soit $\overline{k}$ la clôture algébrique de $k$.
Soit $\overline{k}/\overline{k}'/k$ la clôture algébrique séparable
correspondante de $k$. Alors $\overline{k}/\overline{k}'$ est purement
inséparable (en caractéristique positive) ou triviale.
Ainsi $R \otimes_k \overline{k}' \to R \otimes_k \overline{k}$
induit un homéomorphisme sur les spectres, par exemple par le
Lemme \ref{lemma-p-ring-map}. Nous obtenons donc (4).

\medskip\noindent
Supposons (4). Soit $k'/k$ une extension de corps quelconque et soit
$\overline{k}$ la clôture algébrique de $k$. Choisissons un corps $F$
tel que $k'$ et $\overline{k}$ soient tous deux isomorphes à des
sous-corps de $F$. Alors
$$
R \otimes_k F = (R \otimes_k \overline{k}) \otimes_{\overline{k}} F
$$
et, par le Lemme \ref{lemma-separably-closed-irreducible},
$R \otimes_k F$ possède un unique idéal premier minimal.
Enfin, l'application d'anneaux $R \otimes_k k' \to R \otimes_k F$ est
plate et injective, et tout idéal premier minimal de $R \otimes_k k'$ est
donc l'image d'un idéal premier minimal de $R \otimes_k F$ (par le
Lemme \ref{lemma-injective-minimal-primes-in-image} et la descente).
Nous concluons qu'il n'existe qu'un seul idéal premier minimal, ce qui
achève la démonstration.
\end{proof}

\begin{definition}
\label{definition-geometrically-irreducible}
Soit $k$ un corps.
Soit $S$ une $k$-algèbre.
Nous disons que $S$ est {\it géométriquement irréductible sur $k$}
si, pour toute extension de corps $k'/k$, le spectre de
$S \otimes_k k'$ est irréductible\footnote{Un espace irréductible est non vide.}.
\end{definition}

\noindent
Par le Lemme \ref{lemma-geometrically-irreducible}, il suffit de vérifier
cette propriété pour les extensions de corps finies et séparables $k'/k$
ou pour $k'$ égal à la clôture algébrique séparable de $k$.

\begin{lemma}
\label{lemma-separably-closed-irreducible-implies-geometric}
Soit $k$ un corps.
Soit $R$ une $k$-algèbre.
Si $k$ est séparablement algébriquement clos, alors $R$ est
géométriquement irréductible sur $k$ si et seulement si le spectre de
$R$ est irréductible.
\end{lemma}

\begin{proof}
Cela résulte immédiatement de la remarque qui suit la
Définition \ref{definition-geometrically-irreducible}.
\end{proof}

\begin{lemma}
\label{lemma-subalgebra-geometrically-irreducible}
Soit $k$ un corps. Soit $S$ une $k$-algèbre.
\begin{enumerate}
\item Si $S$ est géométriquement irréductible sur $k$, toute
$k$-sous-algèbre l'est également.
\item Si toutes les $k$-sous-algèbres de type fini de $S$ sont
géométriquement irréductibles, alors $S$ est géométriquement irréductible.
\item Une colimite filtrante de $k$-algèbres géométriquement
irréductibles est géométriquement irréductible.
\end{enumerate}
\end{lemma}

\begin{proof}
Soit $S' \subset S$ une sous-algèbre. Alors, pour toute extension $k'/k$,
l'application d'anneaux $S' \otimes_k k' \to S \otimes_k k'$ est également
injective. Ainsi (1) résulte du
Lemme \ref{lemma-injective-minimal-primes-in-image}
(et du fait que l'image d'un espace irréductible par une application
continue est irréductible). Les deuxième et troisième assertions résultent
du fait que le produit tensoriel commute aux colimites.
\end{proof}

\begin{lemma}
\label{lemma-geometrically-irreducible-any-base-change}
Soit $k$ un corps.
Soit $S$ une $k$-algèbre géométriquement irréductible.
Soit $R$ une $k$-algèbre quelconque.
L'application
$$
\Spec(R \otimes_k S) \longrightarrow \Spec(R)
$$
induit une bijection sur les composantes irréductibles.
\end{lemma}

\begin{proof}
Rappelons que les composantes irréductibles correspondent aux idéaux
premiers minimaux (Lemme \ref{lemma-irreducible}).
Comme $R \to R \otimes_k S$ est plate, la descente montre
(Lemme \ref{lemma-flat-going-down}) que tout idéal premier minimal de
$R \otimes_k S$ est au-dessus d'un idéal premier minimal de $R$.
Réciproquement, si $\mathfrak p \subset R$ est un idéal premier
(minimal), alors
$$
R \otimes_k S/\mathfrak p(R \otimes_k S)
=
(R/\mathfrak p) \otimes_k S
\subset
\kappa(\mathfrak p) \otimes_k S
$$
par platitude de $R \to R \otimes_k S$. L'anneau
$\kappa(\mathfrak p) \otimes_k S$ a un spectre irréductible
par hypothèse. Il s'ensuit que
$R \otimes_k S/\mathfrak p(R \otimes_k S)$ possède un unique idéal
premier minimal (Lemme \ref{lemma-injective-minimal-primes-in-image}).
Autrement dit, l'image inverse de l'ensemble irréductible $V(\mathfrak p)$
est irréductible. Le lemme en résulte.
\end{proof}

\noindent
Faisons maintenant quelques remarques sur la notion d'extensions de corps
géométriquement irréductibles.

\begin{lemma}
\label{lemma-field-extension-geometrically-irreducible}
Soit $K/k$ une extension de corps. Si $k$ est algébriquement clos dans $K$,
alors $K$ est géométriquement irréductible sur $k$.
\end{lemma}

\begin{proof}
Supposons que $k$ soit algébriquement clos dans $K$. Par la
Définition \ref{definition-geometrically-irreducible} et le
Lemme \ref{lemma-geometrically-irreducible}, il suffit de montrer que le
spectre de $K \otimes_k k'$ est irréductible pour toute extension séparable
finie $k'/k$. Disons que $k'$ est engendré par $\alpha \in k'$ sur $k$, voir
le lemme \ref{fields-lemma-primitive-element} de Corps. Soit
$P = T^d + a_1 T^{d - 1} + \ldots + a_d \in k[T]$ le polynôme minimal
de $\alpha$. Alors $K \otimes_k k' \cong K[T]/(P)$.
La seule façon pour que le spectre de $K[T]/(P)$ soit réductible est que
$P$ soit réductible dans $K[T]$. Supposons donc que $P = P_1 P_2$ soit une
factorisation non triviale dans $K[T]$, afin d'obtenir une contradiction.
Par le Lemme \ref{lemma-polynomials-divide}, les coefficients de $P_1$ et
$P_2$ sont algébriques sur $k$. Notre hypothèse implique que les
coefficients de $P_1$ et $P_2$ appartiennent à $k$, ce qui contredit le fait
que $P$ soit irréductible sur $k$.
\end{proof}

\begin{lemma}
\label{lemma-geometrically-irreducible-transitive}
Soit $K/k$ une extension de corps géométriquement irréductible.
Soit $S$ une $K$-algèbre géométriquement irréductible.
Alors $S$ est géométriquement irréductible sur $k$.
\end{lemma}

\begin{proof}
Par la Définition \ref{definition-geometrically-irreducible} et le
Lemme \ref{lemma-geometrically-irreducible}, il suffit de montrer que le
spectre de $S \otimes_k k'$ est irréductible pour toute extension séparable
finie $k'/k$. Comme $K$ est géométriquement irréductible sur $k$, nous
voyons que $K' = K \otimes_k k'$ est une extension de corps finie et
séparable de $K$. Par conséquent, le spectre de
$S \otimes_k k' = S \otimes_K K'$ est irréductible, puisque $S$ est,
par hypothèse, géométriquement irréductible sur $K$.
\end{proof}

\begin{lemma}
\label{lemma-geometrically-irreducible-base-change-transcendental}
Soit $K/k$ une extension de corps. Les assertions suivantes sont équivalentes
\begin{enumerate}
\item $K$ est géométriquement irréductible sur $k$, et
\item l'extension induite $K(t)/k(t)$ d'extensions purement transcendantes
est géométriquement irréductible.
\end{enumerate}
\end{lemma}

\begin{proof}
Supposons (1). Notons $\Omega$ une clôture algébrique de $k(t)$.
Par la Définition \ref{definition-geometrically-irreducible}, nous obtenons
que le spectre de
$$
K \otimes_k \Omega = K \otimes_k k(t) \otimes_{k(t)} \Omega
$$
est irréductible. Comme $K(t)$ est une localisation de
$K \otimes_k k(T)$, nous concluons que le spectre de
$K(t) \otimes_{k(t)} \Omega$ est irréductible. Ainsi, par le
Lemme \ref{lemma-geometrically-irreducible}, $K(t)/k(t)$ est
géométriquement irréductible.

\medskip\noindent
Supposons (2). Soit $k'/k$ une extension de corps.
Nous devons montrer que $K \otimes_k k'$ possède un unique idéal premier
minimal. Nous savons que le spectre de
$$
K(t) \otimes_{k(t)} k'(t)
$$
est irréductible, c'est-à-dire qu'il possède un unique idéal premier minimal.
Comme il existe une application injective
$K \otimes_k k' \to K(t) \otimes_{k(t)} k'(t)$ (détails omis), nous
concluons par les Lemmes \ref{lemma-injective-minimal-primes-in-image} et
\ref{lemma-minimal-prime-image-minimal-prime}.
\end{proof}

\begin{lemma}
\label{lemma-geometrically-irreducible-add-transcendental}
Soit $K/L/M$ une tour de corps telle que $L/M$ soit géométriquement
irréductible. Soit $x \in K$ transcendant sur $L$. Alors $L(x)/M(x)$ est
géométriquement irréductible.
\end{lemma}

\begin{proof}
Cela résulte du
Lemme \ref{lemma-geometrically-irreducible-base-change-transcendental},
car les corps $L(x)$ et $M(x)$ sont des extensions purement transcendantes
de $L$ et $M$.
\end{proof}

\begin{lemma}
\label{lemma-geometrically-irreducible-separable-elements}
Soit $K/k$ une extension de corps. Les assertions suivantes sont équivalentes
\begin{enumerate}
\item $K/k$ est géométriquement irréductible, et
\item tout élément $\alpha \in K$ qui est algébrique séparable sur $k$
appartient à $k$.
\end{enumerate}
\end{lemma}

\begin{proof}
Supposons (1) et soit $\alpha \in K$ algébrique séparable sur $k$.
Alors $k' = k(\alpha)$ est une extension finie et séparable de $k$ contenue
dans $K$. Par le Lemme \ref{lemma-subalgebra-geometrically-irreducible},
l'extension $k'/k$ est géométriquement irréductible.
En particulier, le spectre de $k' \otimes_k \overline{k}$ est
irréductible (et donc, s'il s'agit d'un produit de corps, il n'y a
exactement qu'un seul facteur). Par le
lemme \ref{fields-lemma-finite-separable-tensor-alg-closed} de Corps,
$\Hom_k(k', \overline{k})$ possède un seul élément, ce qui implique à son
tour que $k' = k$ par le
lemme \ref{fields-lemma-separable-equality} de Corps. L'assertion (2)
en résulte.

\medskip\noindent
Supposons (2). Soit $k' \subset K$ le sous-corps constitué des éléments
algébriques sur $k$. Par le
Lemme \ref{lemma-field-extension-geometrically-irreducible}, l'extension
$K/k'$ est géométriquement irréductible. Par hypothèse, $k'/k$ est une
extension purement inséparable. Par le Lemme \ref{lemma-p-ring-map},
l'extension $k'/k$ est géométriquement irréductible. Ainsi, par le
Lemme \ref{lemma-geometrically-irreducible-transitive}, $K/k$ est
géométriquement irréductible.
\end{proof}

\begin{lemma}
\label{lemma-make-geometrically-irreducible}
Soit $K/k$ une extension de corps. Considérons la sous-extension
$K/k'/k$ constituée des éléments algébriques séparables sur $k$.
Alors $K$ est géométriquement irréductible sur $k'$. Si $K/k$ est une
extension de corps de type fini, alors $[k' : k] < \infty$.
\end{lemma}

\begin{proof}
La première assertion résulte immédiatement du
Lemme \ref{lemma-geometrically-irreducible-separable-elements} et du fait
que les éléments algébriques séparables sur $k'$ appartiennent à $k'$ par
transitivité des extensions algébriques séparables, voir le
lemme \ref{fields-lemma-separable-permanence} de Corps.
Si $K/k$ est de type fini, alors $k'$ est une extension finie de $k$ par le
lemme \ref{fields-lemma-algebraic-closure-in-finitely-generated} de Corps.
\end{proof}

\begin{lemma}
\label{lemma-Galois-orbit}
Soit $K/k$ une extension de corps.
Soit $\overline{k}/k$ une clôture algébrique séparable.
Alors $\text{Gal}(\overline{k}/k)$ agit transitivement sur les idéaux
premiers de $\overline{k} \otimes_k K$.
\end{lemma}

\begin{proof}
Soit $K/k'/k$ la sous-extension obtenue dans le
Lemme \ref{lemma-make-geometrically-irreducible}.
Remarquons que, comme $k \subset \overline{k}$ est entière, tous les
idéaux premiers de $\overline{k} \otimes_k K$ et
$\overline{k} \otimes_k k'$ sont maximaux, voir le
Lemme \ref{lemma-integral-no-inclusion}.
Par le Lemme \ref{lemma-geometrically-irreducible-any-base-change},
l'application
$$
\Spec(\overline{k} \otimes_k K) \to \Spec(\overline{k} \otimes_k k')
$$
est bijective car (1) tous les idéaux premiers sont minimaux, (2)
$$
\overline{k} \otimes_k K = (\overline{k} \otimes_k k') \otimes_{k'} K
$$
et (3) $K$ est géométriquement irréductible sur $k'$.
Il suffit donc de démontrer le lemme pour l'action de
$\text{Gal}(\overline{k}/k)$ sur les idéaux premiers de
$\overline{k} \otimes_k k'$.

\medskip\noindent
Comme tout idéal premier de $\overline{k} \otimes_k k'$ est maximal, les
corps résiduels sont isomorphes à $\overline{k}$. Les idéaux premiers de
$\overline{k} \otimes_k k'$ correspondent donc bijectivement aux éléments
de $\Hom_k(k', \overline{k})$, où $\sigma \in \Hom_k(k', \overline{k})$
correspond au noyau $\mathfrak p_\sigma$ de
$1 \otimes \sigma : \overline{k} \otimes_k k' \to \overline{k}$.
En particulier, $\text{Gal}(\overline{k}/k)$ agit transitivement sur cet
ensemble, comme voulu.
\end{proof}

\section{Algèbres géométriquement connexes}
\label{section-geometrically-connected}

\begin{lemma}
\label{lemma-separably-closed-connected}
Soit $k$ un corps séparablement algébriquement clos.
Soient $R$ et $S$ des $k$-algèbres. Si $\Spec(R)$ et
$\Spec(S)$ sont connexes, il en est de même pour
$\Spec(R \otimes_k S)$.
\end{lemma}

\begin{proof}
Rappelons que $\Spec(R)$ est connexe si et seulement si
$R$ ne possède pas d'idempotent non trivial, voir le
Lemme \ref{lemma-characterize-spec-connected}.
Par conséquent, par le Lemme \ref{lemma-limit-argument}, nous pouvons
supposer que $R$ et $S$ sont de type fini sur $k$.
Dans ce cas, $R$ et $S$ sont noethériens et possèdent un nombre fini
d'idéaux premiers minimaux, voir le
Lemme \ref{lemma-Noetherian-irreducible-components}.
Nous pouvons donc raisonner par récurrence sur $n + m$, où $n$, resp.\ $m$,
est le nombre de composantes irréductibles de $\Spec(R)$, resp.\ $\Spec(S)$.
Le cas où $n$ ou $m$ est nul est bien entendu trivial. Si $n = m = 1$,
c'est-à-dire si $\Spec(R)$ et $\Spec(S)$ ont chacun une seule composante
irréductible, le résultat est donné par le
Lemme \ref{lemma-separably-closed-irreducible}.
Supposons $n > 1$. Soit $\mathfrak p \subset R$ un idéal premier minimal
correspondant au fermé irréductible $T \subset \Spec(R)$.
Soit $T' \subset \Spec(R)$ l'union des $n - 1$ autres composantes
irréductibles. Choisissons un idéal $I \subset R$ tel que
$T' = V(I) = \Spec(R/I)$ (Lemme \ref{lemma-spec-closed}).
En choisissant convenablement notre idéal premier minimal, nous pouvons
de plus faire en sorte que $T'$ soit connexe, voir le
lemme \ref{topology-lemma-remove-irreducible-connected} de Topologie.
Alors $T \cup T' = \Spec(R)$ et
$T \cap T' = V(\mathfrak p + I) = \Spec(R/(\mathfrak p + I))$
n'est pas vide puisque $\Spec(R)$ est supposé connexe.
L'image inverse de $T$ dans $\Spec(R \otimes_k S)$ est
$\Spec(R/\mathfrak p \otimes_k S)$, et l'image inverse de $T'$ dans
$\Spec(R \otimes_k S)$ est $\Spec(R/I \otimes_k S)$.
Par récurrence, ces deux espaces sont connexes. L'image inverse de
$T \cap T'$ est $\Spec(R/(\mathfrak p + I) \otimes_k S)$, qui n'est pas vide.
Ainsi $\Spec(R \otimes_k S)$ est connexe.
\end{proof}

\begin{lemma}
\label{lemma-geometrically-connected}
Soit $k$ un corps.
Soit $R$ une $k$-algèbre.
Les assertions suivantes sont équivalentes
\begin{enumerate}
\item pour toute extension de corps $k'/k$, le spectre de
$R \otimes_k k'$ est connexe, et
\item pour toute extension de corps finie et séparable $k'/k$, le spectre
de $R \otimes_k k'$ est connexe.
\end{enumerate}
\end{lemma}

\begin{proof}
Pour toute extension de corps $k'/k$, la connexité du spectre de
$R \otimes_k k'$ équivaut à l'absence d'idempotents non triviaux dans
$R \otimes_k k'$, voir le Lemme \ref{lemma-characterize-spec-connected}.
Supposons (2). Soit $k \subset \overline{k}$ une clôture algébrique
séparable de $k$. Par le Lemme \ref{lemma-limit-argument}, (2) équivaut
à l'absence d'idempotents non triviaux dans $R \otimes_k \overline{k}$.
Pour toute extension de corps $k'/k$, il existe une extension
$\overline{k}'/\overline{k}$ telle que $k' \subset \overline{k}'$.
Par le Lemme \ref{lemma-separably-closed-connected},
$R \otimes_k \overline{k}'$ ne possède pas d'idempotent non trivial.
Si $R \otimes_k k'$ possédait un idempotent non trivial, il en serait de
même dans $R \otimes_k \overline{k}'$, contradiction.
\end{proof}

\begin{definition}
\label{definition-geometrically-connected}
Soit $k$ un corps.
Soit $S$ une $k$-algèbre.
Nous disons que $S$ est {\it géométriquement connexe sur $k$}
si, pour toute extension de corps $k'/k$, le spectre
de $S \otimes_k k'$ est connexe.
\end{definition}

\noindent
Par le Lemme \ref{lemma-geometrically-connected}, il suffit de vérifier
cette propriété pour les extensions de corps finies et séparables $k'/k$.

\begin{lemma}
\label{lemma-separably-closed-connected-implies-geometric}
Soit $k$ un corps.
Soit $R$ une $k$-algèbre.
Si $k$ est séparablement algébriquement clos, alors $R$ est
géométriquement connexe sur $k$ si et seulement si le spectre de
$R$ est connexe.
\end{lemma}

\begin{proof}
Cela résulte immédiatement de la remarque qui suit la
Définition \ref{definition-geometrically-connected}.
\end{proof}

\begin{lemma}
\label{lemma-subalgebra-geometrically-connected}
Soit $k$ un corps. Soit $S$ une $k$-algèbre.
\begin{enumerate}
\item Si $S$ est géométriquement connexe sur $k$, toute
$k$-sous-algèbre l'est également.
\item Si toutes les $k$-sous-algèbres de type fini de $S$ sont
géométriquement connexes, alors $S$ est géométriquement connexe.
\item Une colimite filtrante de $k$-algèbres géométriquement connexes
est géométriquement connexe.
\end{enumerate}
\end{lemma}

\begin{proof}
Cela résulte de la caractérisation de la connexité par l'absence
d'idempotents non triviaux. Les deuxième et troisième assertions résultent
du fait que le produit tensoriel commute aux colimites.
\end{proof}

\noindent
Le lemme suivant sera remplacé par le résultat plus général du
lemme \ref{varieties-lemma-bijection-connected-components} de Variétés.

\begin{lemma}
\label{lemma-geometrically-connected-any-base-change}
Soit $k$ un corps.
Soit $S$ une $k$-algèbre géométriquement connexe.
Soit $R$ une $k$-algèbre quelconque.
L'application
$$
R \longrightarrow R \otimes_k S
$$
induit une bijection sur les idempotents, et l'application
$$
\Spec(R \otimes_k S) \longrightarrow \Spec(R)
$$
induit une bijection sur les composantes connexes.
\end{lemma}

\begin{proof}
La seconde assertion résulte de la première et du
Lemme \ref{lemma-connected-component}.
Par les Lemmes \ref{lemma-subalgebra-geometrically-connected} et
\ref{lemma-limit-argument}, nous pouvons supposer que $R$ et $S$ sont de
type fini sur $k$. Alors $R \otimes_k S$ est également de type fini sur $k$.
Remarquons que, dans ce cas, tous les anneaux sont noethériens et que leurs
spectres possèdent un nombre fini de composantes connexes (puisqu'ils ont
un nombre fini de composantes irréductibles, voir le
Lemme \ref{lemma-Noetherian-irreducible-components}).
En particulier, toutes les composantes connexes en question sont ouvertes !
Par le Lemme \ref{lemma-disjoint-implies-product}, la première assertion
du lemme est alors équivalente à la seconde. Démontrons cette dernière.
Comme l'algèbre $S$ est géométriquement connexe et non nulle, toutes les
fibres de $X = \Spec(R \otimes_k S) \to \Spec(R) = Y$ sont connexes
et non vides. De plus, comme $R \to R \otimes_k S$ est plate et de
présentation finie, l'application $X \to Y$ est ouverte
(Proposition \ref{proposition-fppf-open}).
Le Lemme \ref{topology-lemma-connected-fibres-connected-components} de
Topologie montre que $X \to Y$ induit une bijection sur les composantes
connexes.
\end{proof}

\section{Algèbres géométriquement intègres}
\label{section-geometrically-integral}

\noindent
Voici la définition.

\begin{definition}
\label{definition-geometrically-integral}
Soit $k$ un corps.
Soit $S$ une $k$-algèbre.
Nous disons que $S$ est {\it géométriquement intègre sur $k$}
si, pour toute extension de corps $k'/k$, l'anneau
$S \otimes_k k'$ est un domaine intègre.
\end{definition}

\noindent
Toute question concernant les algèbres géométriquement intègres peut être
traduite en une question concernant les algèbres géométriquement réduites
et irréductibles.

\begin{lemma}
\label{lemma-geometrically-integral}
Soit $k$ un corps.
Soit $S$ une $k$-algèbre.
Alors $S$ est géométriquement intègre sur $k$ si et seulement si
$S$ est à la fois géométriquement irréductible et géométriquement réduite
sur $k$.
\end{lemma}

\begin{proof}
Démonstration omise.
\end{proof}

\begin{lemma}
\label{lemma-characterize-geometrically-integral}
Soit $k$ un corps. Soit $S$ une $k$-algèbre.
Les assertions suivantes sont équivalentes
\begin{enumerate}
\item $S$ est géométriquement intègre sur $k$,
\item pour toute extension finie de corps $k'/k$, l'anneau
$S \otimes_k k'$ est un domaine intègre,
\item $S \otimes_k \overline{k}$ est un domaine intègre,
où $\overline{k}$ est la clôture algébrique de $k$.
\end{enumerate}
\end{lemma}

\begin{proof}
Cela résulte des Lemmes \ref{lemma-geometrically-integral},
\ref{lemma-geometrically-reduced-finite-purely-inseparable-extension} et
\ref{lemma-geometrically-irreducible}.
\end{proof}

\begin{lemma}
\label{lemma-geometrically-integral-any-integral-base-change}
Soit $k$ un corps. Soit $S$ une $k$-algèbre géométriquement intègre.
Soit $R$ une $k$-algèbre qui est un domaine intègre. Alors
$R \otimes_k S$ est un domaine intègre.
\end{lemma}

\begin{proof}
Par le Lemme \ref{lemma-geometrically-reduced-any-reduced-base-change},
l'anneau $R \otimes_k S$ est réduit et, par le
Lemme \ref{lemma-geometrically-irreducible-any-base-change}, l'anneau
$R \otimes_k S$ est irréductible (son spectre possède une seule composante
irréductible), donc $R \otimes_k S$ est un domaine intègre.
\end{proof}


\section{Anneaux de valuation}
\label{section-valuation-rings}

\noindent
Voici quelques définitions.

\begin{definition}
\label{definition-valuation-ring}
Anneaux de valuation.
\begin{enumerate}
\item Soit $K$ un corps. Soient $A$, $B$ des anneaux locaux contenus
dans $K$. Nous disons que $B$ {\it domine} $A$ si $A \subset B$
et si $\mathfrak m_A = A \cap \mathfrak m_B$.
\item Soit $A$ un anneau. Nous disons que $A$ est un {\it anneau de valuation}
si $A$ est un domaine local et si $A$ est maximal
pour la relation de domination parmi les anneaux locaux contenus dans
le corps des fractions de $A$.
\item Soit $A$ un anneau de valuation de corps des fractions $K$.
Si $R \subset K$ est un sous-anneau de $K$, nous disons que $A$
est {\it centré} sur $R$ si $R \subset A$.
\end{enumerate}
\end{definition}

\noindent
Avec cette définition, un corps est un anneau de valuation.

\begin{lemma}
\label{lemma-dominate}
Soit $K$ un corps. Soit $A \subset K$ un sous-anneau local.
Alors il existe un anneau de valuation de corps des fractions $K$
qui domine $A$.
\end{lemma}

\begin{proof}
Nous considérons l'ensemble des sous-anneaux locaux
de $K$ comme un ensemble partiellement ordonné par la relation de domination.
Supposons que $\{A_i\}_{i \in I}$ soit une famille totalement ordonnée
de sous-anneaux locaux de $K$. Alors $B = \bigcup A_i$
est un sous-anneau local qui domine tous les $A_i$.
Par conséquent, par le lemme de Zorn, il suffit de montrer que si $A \subset K$
est un anneau local dont le corps des fractions n'est pas $K$, alors il
existe un anneau local $B \subset K$, $B \not = A$, qui domine $A$.

\medskip\noindent
Choisissons $t \in K$ qui n'appartient pas au corps des fractions de $A$.
Si $t$ est transcendant sur $A$, alors $A[t] \subset K$
et donc $A[t]_{(t, \mathfrak m)} \subset K$ est un anneau local
distinct de $A$ qui domine $A$. Supposons que $t$ soit algébrique sur $A$.
Alors, pour un certain $a \in A$ non nul, l'élément $at$ est entier sur $A$.
Dans ce cas, le sous-anneau $A' \subset K$ engendré par $A$ et
$ta$ est fini sur $A$.
Par le Lemme \ref{lemma-integral-overring-surjective}, il existe
un idéal premier $\mathfrak m' \subset A'$ au-dessus de
$\mathfrak m$. Alors $A'_{\mathfrak m'}$ domine
$A$. Si $A = A'_{\mathfrak m'}$, alors $t$
appartient au corps des fractions de $A$, ce que nous avons supposé faux.
Ainsi $A \not = A'_{\mathfrak m'}$, comme voulu.
\end{proof}

\begin{lemma}
\label{lemma-valuation-ring-normal}
Soit $A$ un anneau de valuation.
Alors $A$ est un domaine normal.
\end{lemma}

\begin{proof}
Supposons que $x$ appartienne au corps des fractions de $A$ et soit entier sur $A$.
Soit $A'$ le sous-anneau de $K$ engendré par $A$ et $x$.
Comme $A\subset A'$ est une extension entière, le
Lemme \ref{lemma-integral-overring-surjective} montre qu'il existe
un idéal premier $\mathfrak m' \subset A'$ au-dessus de
$\mathfrak m$. Alors $A'_{\mathfrak m'}$ domine
$A$. Comme $A$ est un anneau de valuation, nous concluons que $A=A'_{\mathfrak m'}$
et donc que $x\in A$.
\end{proof}


\begin{lemma}
\label{lemma-valuation-ring-x-or-x-inverse}
Soit $A$ un anneau de valuation d'idéal maximal $\mathfrak m$ et de
corps des fractions $K$.
Soit $x \in K$. Alors $x \in A$ ou bien $x^{-1} \in A$, ou les deux.
\end{lemma}

\begin{proof}
Supposons que $x$ n'appartienne pas à $A$.
Soit $A'$ le sous-anneau de $K$ engendré par $A$ et $x$.
Comme $A$ est un anneau de valuation, aucun idéal premier de $A'$
n'est au-dessus de $\mathfrak m$. Comme $\mathfrak m$ est maximal,
nous avons $V(\mathfrak m A') = \emptyset$. Alors $\mathfrak m A' = A'$
par le Lemme \ref{lemma-Zariski-topology}.
Nous pouvons donc écrire
$1 = \sum_{i = 0}^d t_i x^i$ avec $t_i \in \mathfrak m$.
Cela implique que $(1 - t_0) (x^{-1})^d - \sum t_i (x^{-1})^{d - i} = 0$.
En particulier, $x^{-1}$ est entier sur $A$, et donc
$x^{-1} \in A$ par le
Lemme \ref{lemma-valuation-ring-normal}.
\end{proof}

\begin{lemma}
\label{lemma-x-or-x-inverse-valuation-ring}
Soit $A \subset K$ un sous-anneau d'un corps $K$ tel que, pour tout
$x \in K$, on ait $x \in A$ ou bien $x^{-1} \in A$, ou les deux.
Alors $A$ est un anneau de valuation de corps des fractions $K$.
\end{lemma}

\begin{proof}
Si $A$ n'est pas $K$, alors $A$ n'est pas un corps et possède un idéal
maximal non nul $\mathfrak m$.
Si $\mathfrak m'$ est un second idéal maximal, choisissons $x, y \in A$
tels que $x \in \mathfrak m$, $y \not \in \mathfrak m$,
$x \not \in \mathfrak m'$, et $y \in \mathfrak m'$.
Alors ni $x/y \in A$ ni $y/x \in A$,
ce qui contredit l'hypothèse du lemme. Ainsi $A$ est
un anneau local. Supposons que $A'$ soit un anneau local contenu dans $K$
qui domine $A$. Soit $x \in A'$. Nous devons montrer que $x \in A$. Sinon,
$x^{-1} \in A$, et bien sûr $x^{-1} \in \mathfrak m_A$. Mais alors
$x^{-1} \in \mathfrak m_{A'}$, ce qui contredit $x \in A'$.
\end{proof}


\begin{lemma}
\label{lemma-colimit-valuation-rings}
\begin{slogan}
Les anneaux de valuation sont stables par limites directes filtrantes
\end{slogan}
Soit $I$ un ensemble dirigé. Soit $(A_i, \varphi_{ij})$
un système d'anneaux de valuation indexé par $I$.
Alors $A = \colim A_i$ est un anneau de valuation.
\end{lemma}

\begin{proof}
Il est clair que $A$ est un domaine. Soient $a, b \in A$.
Le Lemme \ref{lemma-x-or-x-inverse-valuation-ring} nous dit qu'il faut
montrer que $a | b$ ou $b | a$ dans $A$. Choisissons $i$
assez grand pour qu'il existe $a_i, b_i \in A_i$ dont les images soient $a, b$.
Alors le Lemme \ref{lemma-valuation-ring-x-or-x-inverse}
appliqué à $a_i, b_i$ dans $A_i$ donne le résultat pour $a, b$ dans $A$.
\end{proof}

\begin{lemma}
\label{lemma-valuation-ring-cap-field}
Soit $L/K$ une extension de corps. Si $B \subset L$
est un anneau de valuation, alors $A = K \cap B$ est un anneau de valuation.
\end{lemma}

\begin{proof}
Nous pouvons remplacer $L$ par le corps des fractions $F$ de $B$ et $K$ par
$K \cap F$. Le lemme résulte alors d'une combinaison des
Lemmes \ref{lemma-valuation-ring-x-or-x-inverse} et
\ref{lemma-x-or-x-inverse-valuation-ring}.
\end{proof}

\begin{lemma}
\label{lemma-valuation-ring-cap-field-finite}
Soit $L/K$ une extension algébrique de corps. Si $B \subset L$
est un anneau de valuation de corps des fractions $L$ et n'est pas un corps, alors
$A = K \cap B$ est un anneau de valuation et n'est pas un corps.
\end{lemma}

\begin{proof}
Par le Lemme \ref{lemma-valuation-ring-cap-field}, l'anneau $A$ est un anneau de
valuation. Si $A$ est un corps, alors $A = K$. Alors $A = K \subset B$ est une extension
entière, donc il n'y a pas d'inclusions propres entre les idéaux premiers de $B$
(Lemme \ref{lemma-integral-no-inclusion}).
Cela contredit le fait que $B$ est un domaine local et n'est pas un corps.
\end{proof}

\begin{lemma}
\label{lemma-make-valuation-rings}
Soit $A$ un anneau de valuation. Pour tout idéal premier $\mathfrak p \subset A$, le
quotient $A/\mathfrak p$ est un anneau de valuation. Il en va de même pour la
localisation $A_\mathfrak p$ et, en fait, pour toute localisation de $A$.
\end{lemma}

\begin{proof}
Utiliser la caractérisation des anneaux de valuation donnée
dans le Lemme \ref{lemma-x-or-x-inverse-valuation-ring}.
\end{proof}

\begin{lemma}
\label{lemma-stack-valuation-rings}
Soit $A'$ un anneau de valuation de corps résiduel $K$.
Soit $A$ un anneau de valuation de corps des fractions $K$.
Alors
$C = \{\lambda \in A' \mid \lambda \bmod \mathfrak m_{A'} \in A\}$
est un anneau de valuation.
\end{lemma}

\begin{proof}
Remarquons que $\mathfrak m_{A'} \subset C$ et que $C/\mathfrak m_{A'} = A$.
En particulier, le corps des fractions de $C$ est égal au corps des fractions
de $A'$. Nous utiliserons le critère du
Lemme \ref{lemma-x-or-x-inverse-valuation-ring} pour démontrer le lemme.
Soit $x$ un élément du corps des fractions de $C$.
D'après le lemme, nous pouvons supposer que $x \in A'$. Si $x \in \mathfrak m_{A'}$,
alors $x \in C$. Sinon, $x$ est une unité de $A'$ et nous avons
également $x^{-1} \in A'$. Ainsi $x$ ou $x^{-1}$ s'envoie sur
un élément de $A$, à nouveau par le lemme.
\end{proof}

\begin{lemma}
\label{lemma-find-valuation-rings}
Soit $A$ un domaine normal de corps des fractions $K$.
\begin{enumerate}
\item Pour tout $x \in K$, $x \not \in A$, il existe un anneau de valuation
$A \subset V \subset K$ de corps des fractions $K$ tel que $x \not \in V$.
\item Si $A$ est local, nous pouvons de plus choisir $V$ qui domine $A$.
\end{enumerate}
Autrement dit, $A$ est l'intersection de tous les anneaux de valuation dans $K$
contenant $A$ et, si $A$ est local, $A$ est l'intersection de tous les anneaux de
valuation dans $K$ qui dominent $A$.
\end{lemma}

\begin{proof}
Supposons que $x \in K$, $x \not \in A$. Considérons $B = A[x^{-1}]$.
Alors $x \not \in B$. En effet, si $x = a_0 + a_1x^{-1} + \ldots + a_d x^{-d}$
alors $x^{d + 1} - a_0x^d - \ldots - a_d = 0$ et $x$ est entier
sur $A$, contradiction avec le fait que $A$ est normal.
Ainsi $x^{-1}$ n'est pas une unité de $B$. Donc $V(x^{-1}) \subset \Spec(B)$
n'est pas vide (Lemme \ref{lemma-Zariski-topology}), et nous pouvons choisir un idéal premier
$\mathfrak p \subset B$ tel que $x^{-1} \in \mathfrak p$.
Choisissons un anneau de valuation $V \subset K$ dominant $B_\mathfrak p$
(Lemme \ref{lemma-dominate}).
Alors $x \not \in V$ puisque $x^{-1} \in \mathfrak m_V$.

\medskip\noindent
Si $A$ est local, nous affirmons que $x^{-1} B + \mathfrak m_A B \not = B$.
En effet, si $1 = (a_0 + a_1x^{-1} + \ldots + a_d x^{-d})x^{-1} +
a'_0 + \ldots + a'_d x^{-d}$ avec $a_i \in A$ et $a'_i \in \mathfrak m_A$,
alors nous obtenons
$$
(1 - a'_0) x^{d + 1} - (a_0 + a'_1) x^d - \ldots - a_d = 0
$$
Comme $a'_0 \in \mathfrak m_A$, $1 - a'_0$ est une unité de $A$
et nous concluons que $x$ serait entier sur $A$, contradiction comme
précédemment. Choisissons alors l'idéal premier $\mathfrak p \supset x^{-1} B + \mathfrak m_A B$ ;
nous trouvons $V$ dominant $A$.
\end{proof}

\noindent
Un {\it groupe abélien totalement ordonné} est un couple $(\Gamma, \geq)$
constitué d'un groupe abélien $\Gamma$ muni d'un
ordre total $\geq$ tel que $\gamma \geq \gamma' \Rightarrow
\gamma + \gamma'' \geq \gamma' + \gamma''$ pour tous
$\gamma, \gamma', \gamma'' \in \Gamma$.

\begin{lemma}
\label{lemma-valuation-group}
Soit $A$ un anneau de valuation de corps des fractions $K$.
Posons $\Gamma = K^*/A^*$ (la loi de groupe étant écrite additivement).
Pour $\gamma, \gamma' \in \Gamma$,
définissons $\gamma \geq \gamma'$ si et seulement si
$\gamma - \gamma'$ appartient à l'image de $A - \{0\} \to \Gamma$.
Alors $(\Gamma, \geq)$ est un groupe abélien totalement ordonné.
\end{lemma}

\begin{proof}
Démonstration omise, mais le résultat découle facilement du
Lemme \ref{lemma-valuation-ring-x-or-x-inverse}.
Remarquons que dans le cas $A = K$, on obtient le groupe nul
$\Gamma = \{0\}$ muni de son unique ordre total.
\end{proof}

\begin{definition}
\label{definition-value-group}
Soit $A$ un anneau de valuation.
\begin{enumerate}
\item Le groupe abélien totalement ordonné $(\Gamma, \geq)$ du
Lemme \ref{lemma-valuation-group} est appelé le
{\it groupe des valeurs} de l'anneau de valuation $A$.
\item L'application $v : A - \{0\} \to \Gamma$ et aussi $v : K^* \to \Gamma$ est
appelée la {\it valuation} associée à $A$.
\item L'anneau de valuation $A$ est appelé un {\it anneau de valuation discrète}
si $\Gamma \cong \mathbf{Z}$.
\end{enumerate}
\end{definition}

\noindent
Remarquons que si $\Gamma \cong \mathbf{Z}$, il existe un unique isomorphisme de ce type
tel que $1 \geq 0$. Si l'isomorphisme est choisi de cette manière,
l'ordre devient l'ordre usuel des entiers.

\begin{lemma}
\label{lemma-properties-valuation}
Soit $A$ un anneau de valuation. La valuation $v : A -\{0\} \to \Gamma_{\geq 0}$
possède les propriétés suivantes :
\begin{enumerate}
\item $v(a) = 0 \Leftrightarrow a \in A^*$,
\item $v(ab) = v(a) + v(b)$,
\item $v(a + b) \geq \min(v(a), v(b))$ pourvu que $a + b \not = 0$.
\end{enumerate}
\end{lemma}

\begin{proof}
Démonstration omise.
\end{proof}

\begin{lemma}
\label{lemma-characterize-valuation-ring}
Soit $A$ un anneau. Les assertions suivantes sont équivalentes
\begin{enumerate}
\item $A$ est un anneau de valuation,
\item $A$ est un domaine local et tout idéal
de $A$ engendré par un nombre fini d'éléments est principal.
\end{enumerate}
\end{lemma}

\begin{proof}
Supposons que $A$ soit un anneau de valuation et soient $f_1, \ldots, f_n \in A$.
Choisissons $i$ tel que $v(f_i)$ soit minimal parmi les $v(f_j)$.
Alors $(f_i) = (f_1, \ldots, f_n)$. Réciproquement, supposons que $A$ soit
un domaine local et que tout idéal de $A$ engendré par un nombre fini d'éléments soit principal.
Prenons $f, g \in A$ et écrivons $(f, g) = (h)$. Alors $f = ah$ et $g = bh$
et $h = cf + dg$ pour certains $a, b, c, d \in A$. Ainsi $ac + bd = 1$
et nous voyons que $a$ ou $b$ est une unité, c'est-à-dire que $g/f$
ou $f/g$ appartient à $A$. Cela montre que $A$ est un anneau de valuation
par le Lemme \ref{lemma-x-or-x-inverse-valuation-ring}.
\end{proof}

\begin{lemma}
\label{lemma-valuation-valuation-ring}
Soit $(\Gamma, \geq)$ un groupe abélien totalement ordonné.
Soit $K$ un corps. Soit $v : K^* \to \Gamma$ un homomorphisme
de groupes abéliens tel que $v(a + b) \geq \min(v(a), v(b))$ pour
$a, b \in K$ avec $a, b, a + b$ non nuls. Alors
$$
A =
\{
x \in K \mid x = 0 \text{ ou } v(x) \geq 0
\}
$$
est un anneau de valuation de groupe des valeurs $\Im(v) \subset \Gamma$,
d'idéal maximal
$$
\mathfrak m =
\{
x \in K \mid x = 0 \text{ ou } v(x) > 0
\}
$$
et de groupe des unités
$$
A^* =
\{
x \in K^* \mid v(x) = 0
\}.
$$
\end{lemma}

\begin{proof}
Démonstration omise.
\end{proof}

\noindent
Soit $(\Gamma, \geq)$ un groupe abélien totalement ordonné.
Un {\it idéal de $\Gamma$} est un sous-ensemble $I \subset \Gamma$ tel
que tous les éléments de $I$ soient $\geq 0$ et que, si $\gamma \in I$,
$\gamma' \geq \gamma$ implique $\gamma' \in I$. Nous disons qu'un tel
idéal est {\it premier} si $0 \not \in I$ et si
$\gamma + \gamma' \in I, \gamma, \gamma' \geq 0
\Rightarrow \gamma \in I \text{ ou } \gamma' \in I$.

\begin{lemma}
\label{lemma-ideals-valuation-ring}
Soit $A$ un anneau de valuation.
Les idéaux de $A$ correspondent $1 - 1$ aux idéaux de $\Gamma$.
Cette bijection respecte l'inclusion et envoie les idéaux premiers
sur les idéaux premiers.
\end{lemma}

\begin{proof}
Démonstration omise.
\end{proof}

\begin{lemma}
\label{lemma-valuation-ring-Noetherian-discrete}
Un anneau de valuation est noethérien si et seulement si c'est
un anneau de valuation discrète ou un corps.
\end{lemma}

\begin{proof}
Supposons que $A$ soit un anneau de valuation discrète
de valuation $v : A \setminus \{0\} \to \mathbf{Z}$
normalisée de sorte que $\Im(v) = \mathbf{Z}_{\geq 0}$.
Par le Lemme \ref{lemma-ideals-valuation-ring}, les idéaux de $A$ sont les sous-ensembles
$I_n = \{0\} \cup v^{-1}(\mathbf{Z}_{\geq n})$. Il est clair
que tout élément $x \in A$ tel que $v(x) = n$ engendre $I_n$.
Ainsi $A$ est un PID, donc certainement noethérien.

\medskip\noindent
Supposons que $A$ soit un anneau de valuation noethérien de groupe des valeurs $\Gamma$.
Par le Lemme \ref{lemma-ideals-valuation-ring}, la condition de chaîne ascendante
est satisfaite pour les idéaux de $\Gamma$.
Nous pouvons supposer que $A$ n'est pas un corps, c'est-à-dire qu'il existe un $\gamma \in \Gamma$
tel que $\gamma > 0$. En appliquant la condition de chaîne ascendante aux sous-ensembles
$\gamma + \Gamma_{\geq 0}$ avec $\gamma > 0$, nous voyons
qu'il existe un plus petit élément $\gamma_0$ strictement supérieur à $0$.
Soit $\gamma \in \Gamma$ un élément tel que $\gamma > 0$. Considérons la suite
des éléments $\gamma$, $\gamma - \gamma_0$, $\gamma - 2\gamma_0$,
etc. Par la condition de chaîne ascendante, ils ne peuvent pas tous être $> 0$.
Soit $\gamma - n \gamma_0$ le dernier qui soit $\geq 0$. Par minimalité
de $\gamma_0$, nous voyons que $0 = \gamma - n \gamma_0$. Ainsi $\Gamma$
est un groupe cyclique, comme voulu.
\end{proof}


\section{Autres anneaux noethériens}
\label{section-Noetherian-again}


\begin{lemma}
\label{lemma-Noetherian-basic}
Soit $R$ un anneau noethérien.
Tout $R$-module fini est de présentation finie.
Tout sous-module d'un $R$-module fini est fini.
La condition de chaîne ascendante est satisfaite pour les sous-modules
$R$-d'un $R$-module fini.
\end{lemma}

\begin{proof}
Montrons d'abord que tout sous-module $N$ d'un $R$-module fini
$M$ est fini. Nous procédons par récurrence sur le nombre de
générateurs de $M$. Si ce nombre vaut $1$, alors $N = J/I \subset
M = R/I$ pour certains idéaux $I \subset J \subset R$. Ainsi la définition
de noethérien donne le résultat. Si le nombre de générateurs de
$M$ est supérieur à $1$, nous pouvons trouver une suite exacte courte
$0 \to M' \to M \to M'' \to 0$ où $M'$ et $M''$ ont moins de
générateurs. Remarquons qu'en posant $N' = M' \cap N$ et
$N'' = \Im(N \to M'')$, nous obtenons une suite exacte courte analogue pour $N$.
Le résultat découle donc de l'hypothèse de récurrence
puisque le nombre de générateurs de $N$ est au plus le nombre de
générateurs de $N'$ plus le nombre de générateurs de $N''$.

\medskip\noindent
Pour montrer que $M$ est de présentation finie, il suffit d'appliquer le résultat précédent
au noyau d'une présentation $R^n \to M$.

\medskip\noindent
Il est bien connu et facile de démontrer que la condition de chaîne ascendante pour les
sous-modules $R$ de $M$ est équivalente à la condition que tout sous-module
de $M$ soit un $R$-module fini. Nous omettons la démonstration.
\end{proof}

\begin{lemma}[Artin-Rees]
\label{lemma-Artin-Rees}
Supposons que $R$ soit noethérien et que $I \subset R$ soit un idéal.
Soient $N \subset M$ des $R$-modules finis.
Il existe une constante $c > 0$ telle que
$I^n M \cap N  =  I^{n-c}(I^cM \cap N)$ pour tout $n \geq c$.
\end{lemma}

\begin{proof}
Considérons l'anneau $S = R \oplus I \oplus I^2 \oplus \ldots
= \bigoplus_{n \geq 0} I^n$. Convention : $I^0 = R$.
La multiplication envoie $I^n \times I^m$
dans $I^{n + m}$ par la multiplication dans $R$.
Remarquons que si $I = (f_1, \ldots, f_t)$,
alors $S$ est un quotient de l'anneau noethérien $R[X_1, \ldots, X_t]$.
L'application envoie simplement le monôme $X_1^{e_1}\ldots X_t^{e_t}$
sur $f_1^{e_1}\ldots f_t^{e_t}$. Ainsi $S$ est noethérien.
De même, considérons le module $M \oplus IM \oplus I^2M \oplus \ldots
= \bigoplus_{n \geq 0} I^nM$. C'est un $S$-module de type fini.
En effet, si $x_1, \ldots, x_r$ engendrent $M$ sur $R$, ils engendrent aussi
$\bigoplus_{n \geq 0} I^nM$ sur $S$. Ensuite, considérons le
sous-module $\bigoplus_{n \geq 0} I^nM \cap N$.
Il s'agit d'un $S$-sous-module, comme on le vérifie facilement. Par le
Lemme \ref{lemma-Noetherian-basic}, il est de type fini comme
$S$-module,
disons engendré par $\xi_j \in \bigoplus_{n \geq 0} I^nM \cap N$, $j = 1, \ldots, s$.
Nous pouvons supposer, en décomposant chaque $\xi_j$ en ses composantes homogènes,
que $\xi_j \in I^{d_j}M \cap N$ pour un certain $d_j$.
Posons $c = \max\{d_j\}$. Alors, pour tout $n \geq c$, tout élément
de $I^nM \cap N$ est de la forme $\sum h_j \xi_j$ avec
$h_j \in I^{n - d_j}$. Le lemme résulte alors de cette observation et du fait trivial que
$I^{n-d_j}(I^{d_j}M \cap N) \subset I^{n-c}(I^cM \cap N)$.
\end{proof}

\begin{lemma}
\label{lemma-map-AR}
Supposons que $0 \to K \to M \xrightarrow{f} N$ soit une
suite exacte de modules de type fini
sur un anneau noethérien $R$. Soit $I \subset R$ un idéal.
Alors il existe un $c$ tel que
$$
f^{-1}(I^nN) = K + I^{n-c}f^{-1}(I^cN)
\quad\text{et}\quad
f(M) \cap I^nN \subset f(I^{n - c}M)
$$
pour tout $n \geq c$.
\end{lemma}

\begin{proof}
Appliquer le Lemme \ref{lemma-Artin-Rees} à
$\Im(f) \subset N$ et remarquer que
$f : I^{n-c}M \to I^{n-c}f(M)$ est surjectif.
\end{proof}

\begin{lemma}[théorème d'intersection de Krull]
\label{lemma-intersect-powers-ideal-module-zero}
Soit $R$ un anneau local noethérien. Soit $I \subset R$ un
idéal propre. Soit $M$ un $R$-module fini.
Alors $\bigcap_{n \geq 0} I^nM = 0$.
\end{lemma}

\begin{proof}
Soit $N = \bigcap_{n \geq 0} I^nM$.
Alors $N = I^nM \cap N$ pour tout $n \geq 0$.
Par le Lemme d'Artin-Rees \ref{lemma-Artin-Rees},
nous voyons que $N = I^nM \cap N \subset IN$
pour un certain $n$ suffisamment grand. Par le Lemme de Nakayama
\ref{lemma-NAK}, nous avons $N = 0$.
\end{proof}

\begin{lemma}
\label{lemma-intersection-powers-ideal-module}
Soit $R$ un anneau noethérien. Soit $I \subset R$ un idéal.
Soit $M$ un $R$-module fini. Soit $N = \bigcap_n I^n M$.
\begin{enumerate}
\item Pour tout idéal premier $\mathfrak p$, $I \subset \mathfrak p$, il
existe un $f \in R$, $f \not \in \mathfrak p$, tel que $N_f = 0$.
\item Si $I$ est contenu dans le radical de Jacobson
de $R$, alors $N = 0$.
\end{enumerate}
\end{lemma}

\begin{proof}
Démontrons (1). Soient $x_1, \ldots, x_n$ des générateurs du module $N$,
voir le Lemme \ref{lemma-Noetherian-basic}. Pour tout idéal premier
$\mathfrak p$, $I \subset \mathfrak p$, l'image de $N$ dans la localisation $M_{\mathfrak p}$
est nulle, par le Lemme \ref{lemma-intersect-powers-ideal-module-zero}.
Nous pouvons donc trouver $g_i \in R$, $g_i \not \in \mathfrak p$
tel que $x_i$ s'envoie sur zéro dans $N_{g_i}$. Ainsi
$N_{g_1g_2\ldots g_n} = 0$.

\medskip\noindent
La partie (2) découle de (1) et du Lemme \ref{lemma-characterize-zero-local}.
\end{proof}

\begin{remark}
\label{remark-intersection-powers-ideal}
Le Lemme \ref{lemma-intersect-powers-ideal-module-zero} implique en particulier
que $\bigcap_n I^n = (0)$ lorsque $I \subset R$ est un idéal non unité dans un
anneau local noethérien $R$. Plus généralement, soit $R$ un anneau noethérien et
$I \subset R$ un idéal. Supposons que $f \in \bigcap_{n \in \mathbf{N}} I^n$.
Alors le Lemme \ref{lemma-intersection-powers-ideal-module}
dit que, pour tout idéal premier $I \subset \mathfrak p$,
il existe un $g \in R$, $g \not \in \mathfrak p$,
tel que $f$ s'envoie sur zéro dans $R_g$. En géométrie algébrique, nous
exprimons cela en disant que « $f$ est nul dans un voisinage ouvert
de l'ensemble fermé $V(I)$ de $\Spec(R)$ ».
\end{remark}

\begin{lemma}[Artin-Tate]
\label{lemma-Artin-Tate}
Soit $R$ un anneau noethérien. Soit $S$ une $R$-algèbre
de type fini. Si $T \subset S$ est une $R$-sous-algèbre telle que
$S$ soit de type fini comme $T$-module, alors $T$ est de
type fini sur $R$.
\end{lemma}

\begin{proof}
Choisissons des éléments $x_1, \ldots, x_n \in S$ qui engendrent $S$ comme $R$-algèbre.
Choisissons $y_1, \ldots, y_m$ dans $S$ qui engendrent $S$ comme $T$-module.
Il existe donc des $a_{ij} \in T$ tels que
$x_i = \sum a_{ij} y_j$. Il existe également des $b_{ijk} \in T$ tels
que $y_i y_j = \sum b_{ijk} y_k$. Soit $T' \subset T$ la
$R$-sous-algèbre engendrée par les $a_{ij}$ et $b_{ijk}$. Il s'agit d'une $R$-algèbre
de type fini, donc noethérienne. Considérons l'algèbre
$$
S' = T'[Y_1, \ldots, Y_m]/(Y_i Y_j - \sum b_{ijk} Y_k).
$$
Remarquons que $S'$ est fini sur $T'$, car en tant que $T'$-module il est
engendré par les classes de $1, Y_1, \ldots, Y_m$.
Considérons l'homomorphisme de $T'$-algèbres $S' \to S$ qui envoie
$Y_i$ sur $y_i$. Comme $a_{ij} \in T'$, nous voyons que $x_j$ appartient à l'image de cette application. Ainsi $S' \to S$ est surjectif.
Par conséquent, $S$ est également fini sur $T'$. Comme $T'$ est noethérien,
nous concluons que $T \subset S$ est fini sur $T'$ et nous avons terminé.
\end{proof}


\section{Longueur}
\label{section-length}

\begin{definition}
\label{definition-length}
Soit $R$ un anneau. Pour tout $R$-module $M$,
nous définissons la {\it longueur} de $M$ sur $R$ par la
formule
$$
\text{length}_R(M)
=
\sup
\{
n
\mid
\exists\ 0 = M_0 \subset M_1 \subset \ldots \subset M_n = M,
\text{ }M_i \not = M_{i + 1}
\}.
$$
\end{definition}

\noindent
Autrement dit, c'est le supremum des longueurs des chaînes
de sous-modules. Il existe une notion évidente de chaîne
de sous-modules qui raffine une autre. Cela donne un
ordre partiel sur l'ensemble de toutes les chaînes de sous-modules,
la chaîne minimale ayant la forme $0 = M_0 \subset M_1 = M$
si $M$ n'est pas nul.
Remarquons le fait évident que si la longueur de
$M$ est finie, toute chaîne peut être raffinée en une
chaîne maximale. Mais il est moins évident que toutes les chaînes maximales
ont la même longueur (comme nous le verrons plus loin).

\begin{lemma}
\label{lemma-finite-length-finite}
\begin{slogan}
Les modules de longueur finie sont finis.
\end{slogan}
Soit $R$ un anneau.
Soit $M$ un $R$-module.
Si $\text{length}_R(M) < \infty$, alors $M$ est un $R$-module fini.
\end{lemma}

\begin{proof}
Démonstration omise.
\end{proof}

\begin{lemma}
\label{lemma-length-additive}
\begin{slogan}
La longueur est additive dans les suites exactes courtes.
\end{slogan}
Si $0 \to M' \to M \to M'' \to 0$
est une suite exacte courte de modules sur $R$, alors
la longueur de $M$ est la somme des
longueurs de $M'$ et $M''$.
\end{lemma}

\begin{proof}
Étant données des filtrations de $M'$ et $M''$ de longueurs $n', n''$,
il est facile de construire une filtration correspondante de $M$
de longueur $n' + n''$. Ainsi
$\text{length}_R M \geq \text{length}_R M' + \text{length}_R M''$.
Réciproquement, étant donnée une filtration
$M_0 \subset M_1 \subset \ldots \subset M_n$ de
$M$, considérons les filtrations induites
$M_i' = M_i \cap M'$ et $M_i'' = \Im(M_i \to M'')$.
Soit $n'$ (resp.\ $n''$) le nombre d'étapes de la filtration
$\{M'_i\}$ (resp.\ $\{M''_i\}$).
Si $M_i' = M_{i + 1}'$ et $M_i'' = M_{i + 1}''$, alors
$M_i = M_{i + 1}$. Nous concluons donc que $n' + n'' \geq n$.
Avec le résultat précédent, c'est démontré.
\end{proof}

\begin{lemma}
\label{lemma-length-infinite}
Soit $R$ un anneau local d'idéal maximal $\mathfrak m$.
Si $M$ est un $R$-module et $\mathfrak m^n M \not = 0$ pour tout $n \geq 0$,
alors $\text{length}_R(M) = \infty$. Autrement dit, si $M$
est de longueur finie, alors $\mathfrak m^nM = 0$ pour un certain $n$.
\end{lemma}

\begin{proof}
Supposons que $\mathfrak m^n M \not = 0$ pour tout $n\geq 0$.
Choisissons $x \in M$ et $f_1, \ldots, f_n \in \mathfrak m$
tels que $f_1f_2 \ldots f_n x \not = 0$.
Les $n$ premières étapes de la filtration
$$
0 \subset R f_1 \ldots f_n x \subset R f_1 \ldots f_{n - 1} x
\subset \ldots \subset R x \subset M
$$
sont distinctes. Par exemple, si
$R f_1 x = R f_1 f_2 x$, alors $f_1 x = g f_1 f_2 x$ pour un certain $g$,
d'où $(1 - gf_2) f_1 x = 0$, donc $f_1 x = 0$ puisque $1 - gf_2$ est une unité,
ce qui contredit le choix de $x$ et de $f_1, \ldots, f_n$.
Ainsi la longueur est infinie.
\end{proof}

\begin{lemma}
\label{lemma-length-independent}
Soit $R \to S$ un morphisme d'anneaux. Soit $M$ un $S$-module.
Nous avons toujours $\text{length}_R(M) \geq \text{length}_S(M)$.
Si $R \to S$ est surjectif, alors il y a égalité.
\end{lemma}

\begin{proof}
Une filtration de $M$ par des $S$-sous-modules donne une filtration
de $M$ par des $R$-sous-modules. Cela démontre l'inégalité.
Et si $R \to S$ est surjectif, tout $R$-sous-module
de $M$ est automatiquement un $S$-sous-module. Il y a donc égalité
dans ce cas.
\end{proof}

\begin{lemma}
\label{lemma-dimension-is-length}
Soit $R$ un anneau d'idéal maximal $\mathfrak m$.
Supposons que $M$ soit un $R$-module tel que
$\mathfrak m M  =  0$. Alors la longueur de $M$ comme
$R$-module coïncide avec la dimension de $M$ comme
espace vectoriel sur $R/\mathfrak m$.
La longueur est finie si et seulement si $M$ est un $R$-module fini.
\end{lemma}

\begin{proof}
La première assertion est un cas particulier du
Lemme \ref{lemma-length-independent}.
Ainsi la longueur est finie si et seulement si $M$ possède une base finie
comme espace vectoriel sur $R/\mathfrak m$, si et seulement si $M$ possède un ensemble fini
de générateurs comme $R$-module.
\end{proof}

\begin{lemma}
\label{lemma-length-localize}
Soit $R$ un anneau. Soit $M$ un $R$-module. Soit $S \subset R$ un
ensemble multiplicatif. Alors
$\text{length}_R(M) \geq \text{length}_{S^{-1}R}(S^{-1}M)$.
\end{lemma}

\begin{proof}
Tout sous-module $N' \subset S^{-1}M$ est de la forme
$S^{-1}N$ pour un certain sous-module $R$-module $N \subset M$, par le Lemme
\ref{lemma-submodule-localization}. Le lemme en découle.
\end{proof}

\begin{lemma}
\label{lemma-length-finite}
Soit $R$ un anneau dont l'idéal maximal $\mathfrak m$ est engendré par un nombre fini d'éléments.
(Par exemple, $R$ est noethérien.)
Supposons que $M$ soit un $R$-module fini tel que
$\mathfrak m^n M  =  0$ pour un certain $n$.
Alors $\text{length}_R(M) < \infty$.
\end{lemma}

\begin{proof}
Considérons la filtration
$0 = \mathfrak m^n M \subset
\mathfrak m^{n-1} M \subset
\ldots \subset \mathfrak m M \subset M$.
Tous les quotients successifs sont des $R$-modules de type fini
auxquels s'applique le Lemme \ref{lemma-dimension-is-length}. Nous concluons
par additivité, voir le Lemme \ref{lemma-length-additive}.
\end{proof}


\begin{definition}
\label{definition-simple-module}
Soit $R$ un anneau. Soit $M$ un $R$-module.
Nous disons que $M$ est {\it simple} si $M \not = 0$ et
tout sous-module de $M$ est soit égal à $M$, soit
à $0$.
\end{definition}

\begin{lemma}
\label{lemma-characterize-length-1}
Soit $R$ un anneau. Soit $M$ un $R$-module.
Les assertions suivantes sont équivalentes :
\begin{enumerate}
\item $M$ est simple,
\item $\text{length}_R(M) = 1$, et
\item $M \cong R/\mathfrak m$ pour un idéal maximal
$\mathfrak m \subset R$.
\end{enumerate}
\end{lemma}

\begin{proof}
Soit $\mathfrak m$ un idéal maximal de $R$.
Par le Lemme \ref{lemma-dimension-is-length}, le module
$R/\mathfrak m$ est de longueur $1$. L'équivalence des deux
premières assertions est tautologique.
Supposons que $M$ soit simple. Choisissons $x \in M$, $x \not = 0$.
Comme $M$ est simple, nous avons $M = R \cdot x$.
Soit $I \subset R$ l'annulateur de $x$, c'est-à-dire
$I = \{f \in R \mid fx = 0\}$. L'application $R/I \to M$,
$f \bmod I \mapsto fx$ est un isomorphisme, donc
$R/I$ est un $R$-module simple. Comme $R/I \not = 0$, nous voyons que $I \not = R$.
Soit $I \subset \mathfrak m$ un idéal maximal contenant $I$.
Si $I \not = \mathfrak m$, alors $\mathfrak m /I \subset R/I$
est un sous-module non trivial, ce qui contredit la simplicité
de $R/I$. Nous obtenons donc $I = \mathfrak m$, comme voulu.
\end{proof}

\begin{lemma}
\label{lemma-simple-pieces}
Soit $R$ un anneau. Soit $M$ un $R$-module de longueur finie.
Choisissons une chaîne maximale quelconque de sous-modules
$$
0 = M_0 \subset M_1 \subset M_2 \subset \ldots \subset M_n = M
$$
telle que $M_i \not = M_{i-1}$, $i = 1, \ldots, n$. Alors
\begin{enumerate}
\item $n = \text{length}_R(M)$,
\item chaque $M_i/M_{i-1}$ est simple,
\item chaque $M_i/M_{i-1}$ est de la forme
$R/\mathfrak m_i$ pour un certain idéal maximal $\mathfrak m_i$,
\item étant donné un idéal maximal $\mathfrak m \subset R$,
nous avons
$$
\# \{i \mid \mathfrak m_i = \mathfrak m\}
=
\text{length}_{R_{\mathfrak m}} (M_{\mathfrak m}).
$$
\end{enumerate}
\end{lemma}

\begin{proof}
Si $M_i/M_{i-1}$ n'est pas simple, nous pouvons raffiner la filtration
et la filtration n'est pas maximale. Ainsi $M_i/M_{i-1}$
est simple. Par le Lemme \ref{lemma-characterize-length-1}, les modules
$M_i/M_{i-1}$ sont de longueur $1$ et sont de la forme $R/\mathfrak m_i$
pour certains idéaux maximaux $\mathfrak m_i$. Par additivité de la longueur,
Lemme \ref{lemma-length-additive}, nous voyons que $n = \text{length}_R(M)$.
Comme la localisation est exacte, nous voyons que
$$
0 = (M_0)_{\mathfrak m}
\subset (M_1)_{\mathfrak m}
\subset (M_2)_{\mathfrak m}
\subset \ldots
\subset (M_n)_{\mathfrak m} = M_{\mathfrak m}
$$
est une filtration de $M_{\mathfrak m}$ dont les quotients successifs sont
$(M_i/M_{i-1})_{\mathfrak m}$. La dernière assertion découle alors
directement du fait que pour des idéaux maximaux $\mathfrak m$,
$\mathfrak m'$ de $R$, nous avons
$$
(R/\mathfrak m')_{\mathfrak m}
\cong
\left\{
\begin{matrix}
0 &
\text{si } \mathfrak m \not = \mathfrak m', \\
R_{\mathfrak m}/\mathfrak m R_{\mathfrak m} &
\text{si } \mathfrak m  = \mathfrak m'
\end{matrix}
\right.
$$
Nous laissons cela au lecteur.
\end{proof}

\begin{lemma}
\label{lemma-pushdown-module}
Soit $A$ un anneau local d'idéal maximal $\mathfrak m$.
Soit $B$ un anneau semi-local d'idéaux maximaux $\mathfrak m_i$,
$i = 1, \ldots, n$.
Supposons que $A \to B$ soit un homomorphisme tel que chaque $\mathfrak m_i$
soit au-dessus de $\mathfrak m$ et tel que
$$
[\kappa(\mathfrak m_i) : \kappa(\mathfrak m)] < \infty.
$$
Soit $M$ un $B$-module de longueur finie.
Alors
$$
\text{length}_A(M) = \sum\nolimits_{i = 1, \ldots, n}
[\kappa(\mathfrak m_i) : \kappa(\mathfrak m)]
\text{length}_{B_{\mathfrak m_i}}(M_{\mathfrak m_i}),
$$
en particulier $\text{length}_A(M) < \infty$.
\end{lemma}

\begin{proof}
Choisissons une chaîne maximale
$$
0 = M_0
\subset M_1
\subset M_2
\subset \ldots
\subset M_m = M
$$
par des $B$-sous-modules comme dans le Lemme \ref{lemma-simple-pieces}.
Alors chaque quotient $M_j/M_{j - 1}$ est isomorphe à
$\kappa(\mathfrak m_{i(j)})$ pour un certain $i(j) \in \{1, \ldots, n\}$.
De plus
$\text{length}_A(\kappa(\mathfrak m_i)) =
[\kappa(\mathfrak m_i) : \kappa(\mathfrak m)]$ par le
Lemme \ref{lemma-dimension-is-length}. Le lemme découle
de l'additivité des longueurs (Lemme \ref{lemma-length-additive}).
\end{proof}

\begin{lemma}
\label{lemma-pullback-module}
Soit $A \to B$ un homomorphisme local plat d'anneaux locaux.
Alors, pour tout $A$-module $M$, nous avons
$$
\text{length}_A(M) \text{length}_B(B/\mathfrak m_AB)
=
\text{length}_B(M \otimes_A B).
$$
En particulier, si $\text{length}_B(B/\mathfrak m_AB) < \infty$,
alors $M$ est de longueur finie si et seulement si $M \otimes_A B$ est de longueur finie.
\end{lemma}

\begin{proof}
Le morphisme d'anneaux $A \to B$ est fidèlement plat par le
Lemme \ref{lemma-local-flat-ff}.
Ainsi, si $0 = M_0 \subset M_1 \subset \ldots \subset M_n = M$
est une chaîne de longueur $n$ dans $M$, la chaîne correspondante
$0 = M_0 \otimes_A B \subset M_1 \otimes_A B \subset
\ldots \subset M_n \otimes_A B = M \otimes_A B$ est aussi de longueur $n$.
Cela démontre
$\text{length}_A(M) = \infty \Rightarrow
\text{length}_B(M \otimes_A B) = \infty$.
Supposons maintenant $\text{length}_A(M) < \infty$. Dans ce cas, $M$ possède une filtration de longueur $\ell = \text{length}_A(M)$
dont les quotients sont $A/\mathfrak m_A$. En raisonnant comme ci-dessus,
nous voyons que $M \otimes_A B$ possède une filtration de longueur $\ell$
dont les quotients sont isomorphes à
$B \otimes_A A/\mathfrak m_A =  B/\mathfrak m_AB$.
Le lemme en découle.
\end{proof}

\begin{lemma}
\label{lemma-pullback-transitive}
Soient $A \to B \to C$ des homomorphismes locaux plats d'anneaux locaux. Alors
$$
\text{length}_B(B/\mathfrak m_A B)
\text{length}_C(C/\mathfrak m_B C)
=
\text{length}_C(C/\mathfrak m_A C)
$$
\end{lemma}

\begin{proof}
Cela résulte du Lemme \ref{lemma-pullback-module} appliqué au morphisme d'anneaux
$B \to C$ et au $B$-module $M = B/\mathfrak m_A B$.
\end{proof}


\section{Anneaux artiniens}
\label{section-artinian}

\noindent
Les anneaux artiniens, et en particulier les anneaux artiniens locaux,
jouent un rôle important en géométrie algébrique, par exemple
dans la théorie des déformations.

\begin{definition}
\label{definition-artinian}
Un anneau $R$ est {\it artinien} s'il satisfait la
condition de chaîne descendante pour les idéaux.
\end{definition}

\begin{lemma}
\label{lemma-finite-dimensional-algebra}
Supposons que $R$ soit une algèbre de dimension finie sur un corps.
Alors $R$ est artinien.
\end{lemma}

\begin{proof}
La condition de chaîne descendante pour les idéaux est évidente.
\end{proof}

\begin{lemma}
\label{lemma-artinian-finite-nr-max}
Si $R$ est artinien, alors $R$ ne possède qu'un nombre fini d'idéaux maximaux.
\end{lemma}

\begin{proof}
Supposons que $\mathfrak m_i$, $i = 1, 2, 3, \ldots$, soient
des idéaux maximaux deux à deux distincts.
Alors $\mathfrak m_1 \supset \mathfrak m_1\cap \mathfrak m_2
\supset \mathfrak m_1 \cap \mathfrak m_2 \cap \mathfrak m_3 \supset \ldots$
est une suite descendante infinie (car, par le
théorème des restes chinois, toutes les applications
$R \to \oplus_{i = 1}^n R/\mathfrak m_i$ sont surjectives).
\end{proof}

\begin{lemma}
\label{lemma-artinian-radical-nilpotent}
Soit $R$ artinien. Le radical de Jacobson de $R$ est un idéal nilpotent.
\end{lemma}

\begin{proof}
Soit $I \subset R$ le radical de Jacobson.
Remarquons que $I \supset I^2 \supset I^3 \supset \ldots$ est une suite
descendante. Ainsi $I^n = I^{n + 1}$ pour un certain $n$.
Posons $J = \{ x\in R \mid xI^n = 0\}$. Nous devons montrer que $J = R$.
Sinon, choisissons un idéal $J' \not = J$, $J \subset J'$ minimal
(possible par la propriété artinienne). Alors $J'/J$ est un
$R$-module simple, donc isomorphe à $R/\mathfrak m$ pour un certain idéal maximal $\mathfrak m$,
voir le Lemme \ref{lemma-characterize-length-1}.
Alors $\mathfrak m I^n$ annule $J'$. Comme $I \subset \mathfrak m$,
nous concluons que $I^{n + 1} = I^n$ annule $J'$. Donc $J' = J$,
ce qui est une contradiction.
\end{proof}

\begin{lemma}
\label{lemma-product-local}
Tout anneau ayant un nombre fini d'idéaux maximaux et
un radical de Jacobson localement nilpotent est le produit de ses localisations
en ses idéaux maximaux. De plus, tous les idéaux premiers sont maximaux.
\end{lemma}

\begin{proof}
Soit $R$ un anneau ayant un nombre fini d'idéaux maximaux
$\mathfrak m_1, \ldots, \mathfrak m_n$.
Soit $I = \bigcap_{i = 1}^n \mathfrak m_i$
le radical de Jacobson de $R$. Supposons que $I$ soit localement nilpotent.
Soit $\mathfrak p$ un idéal premier de $R$.
Comme tout idéal premier contient tout élément
nilpotent de $R$, nous voyons
$ \mathfrak p \supset \mathfrak m_1 \cap \ldots \cap \mathfrak m_n$.
Comme $\mathfrak m_1 \cap \ldots \cap \mathfrak m_n \supset
\mathfrak m_1 \ldots \mathfrak m_n$,
nous concluons $\mathfrak p \supset \mathfrak m_1 \ldots \mathfrak m_n$.
Donc $\mathfrak p \supset \mathfrak m_i$ pour un certain $i$, et ainsi
$\mathfrak p = \mathfrak m_i$. Le spectre de $R$
est donc l'espace topologique discret $\{\mathfrak m_1, \ldots, \mathfrak m_n\}$.
Par le Lemme \ref{lemma-disjoint-implies-product}, appliqué $n - 1$ fois,
nous trouvons que $R = R_1 \times \ldots \times R_n$, où le
spectre de chaque $R_i$ est un singleton pour chaque $i$.
Ainsi $R_i$ est un anneau local et, comme il est une localisation de $R$
(par le lemme), il est l'un des anneaux locaux de $R$, comme voulu.
\end{proof}

\begin{lemma}
\label{lemma-artinian-finite-length}
Un anneau $R$ est artinien si et seulement s'il est de longueur finie
comme module sur lui-même. Un tel anneau $R$ est à la fois artinien et
noethérien, tout idéal premier de $R$ est maximal, et $R$ est égal
au produit (fini) de ses localisations en ses idéaux maximaux.
\end{lemma}

\begin{proof}
Si $R$ est de longueur finie sur lui-même, il satisfait à la fois
la condition de chaîne ascendante et la condition de chaîne descendante
pour les idéaux. Il est donc à la fois noethérien et artinien.
Tout anneau artinien est égal au produit de ses localisations
en ses idéaux maximaux par les Lemmes \ref{lemma-artinian-finite-nr-max},
\ref{lemma-artinian-radical-nilpotent} et \ref{lemma-product-local}.

\medskip\noindent
Supposons que $R$ soit artinien. Nous allons montrer que $R$ est de longueur finie
sur lui-même. Il suffit d'exhiber une chaîne de
sous-modules dont les quotients successifs sont de longueur finie.
D'après ce qui précède, nous pouvons supposer que $R$ est local, d'idéal maximal $\mathfrak m$.
Par le Lemme \ref{lemma-artinian-radical-nilpotent}, nous avons
$\mathfrak m^n =0$ pour un certain $n$.
Considérons la suite
$0 = \mathfrak m^n \subset \mathfrak m^{n-1} \subset
\ldots \subset \mathfrak m \subset R$. Par le Lemme
\ref{lemma-dimension-is-length}, la longueur de chaque quotient successif
$\mathfrak m^j/\mathfrak m^{j + 1}$ est sa dimension
comme espace vectoriel sur $\kappa(\mathfrak m)$. Celle-ci doit être
finie, car sinon nous aurions une chaîne descendante infinie
de sous-espaces vectoriels, qui correspondrait à une chaîne descendante infinie d'idéaux dans $R$.
\end{proof}

\section{Homomorphismes essentiellement de type fini}
\label{section-essentially-of-finite-type}

\noindent
Quelques remarques simples sur les localisations de morphismes d'anneaux de type fini.

\begin{definition}
\label{definition-essentially-finite-p-t}
Soit $R \to S$ un morphisme d'anneaux.
\begin{enumerate}
\item Nous disons que $R \to S$ est {\it essentiellement de type fini} si
$S$ est la localisation d'une $R$-algèbre de type fini.
\item Nous disons que $R \to S$ est {\it essentiellement de présentation finie} si
$S$ est la localisation d'une $R$-algèbre de présentation finie.
\end{enumerate}
\end{definition}

\begin{lemma}
\label{lemma-composition-essentially-of-finite-type}
La classe des morphismes d'anneaux qui sont essentiellement de type fini
est stable par composition. De même pour les morphismes essentiellement de
présentation finie.
\end{lemma}

\begin{proof}
Omis.
\end{proof}

\begin{lemma}
\label{lemma-base-change-essentially-of-finite-type}
La classe des morphismes d'anneaux qui sont essentiellement de type fini
est stable par changement de base. De même pour les morphismes essentiellement de
présentation finie.
\end{lemma}

\begin{proof}
Omis.
\end{proof}

\begin{lemma}
\label{lemma-essentially-of-finite-type-into-artinian-local}
Soit $R \to S$ un morphisme d'anneaux. Supposons que $S$ soit un anneau local
artinien d'idéal maximal $\mathfrak m$. Alors
\begin{enumerate}
\item $R \to S$ est fini si et seulement si $R \to S/\mathfrak m$ est fini,
\item $R \to S$ est de type fini si et seulement si $R \to S/\mathfrak m$
est de type fini.
\item $R \to S$ est essentiellement de type fini si et seulement si la
composition $R \to S/\mathfrak m$ est essentiellement de type fini.
\end{enumerate}
\end{lemma}

\begin{proof}
Si $R \to S$ est fini, alors $R \to S/\mathfrak m$
est fini par le Lemme \ref{lemma-finite-transitive}.
Réciproquement, supposons que $R \to S/\mathfrak m$ soit fini.
Comme $S$ est de longueur finie sur lui-même
(Lemme \ref{lemma-artinian-finite-length}), nous pouvons choisir une filtration
$$
0 \subset I_1 \subset \ldots \subset I_n = S
$$
par des idéaux telle que $I_i/I_{i - 1} \cong S/\mathfrak m$ comme
$S$-modules. Ainsi, $S$ possède une filtration par des $R$-sous-modules
$I_i$ dont chaque quotient successif est un $R$-module fini. Donc $S$ est
un $R$-module fini par le Lemme \ref{lemma-extension}.

\medskip\noindent
Si $R \to S$ est de type fini, alors $R \to S/\mathfrak m$
est de type fini par le Lemme \ref{lemma-compose-finite-type}.
Réciproquement, supposons que $R \to S/\mathfrak m$ soit de type fini.
Choisissons $f_1, \ldots, f_n \in S$ qui se projettent sur des générateurs
de $S/\mathfrak m$. Alors $A = R[x_1, \ldots, x_n] \to S$, $x_i \mapsto f_i$
est un morphisme d'anneaux tel que $A \to S/\mathfrak m$ soit surjectif
(en particulier fini). Ainsi $A \to S$ est fini par (1), et nous voyons que
$R \to S$ est de type fini par le Lemme \ref{lemma-compose-finite-type}.

\medskip\noindent
Si $R \to S$ est essentiellement de type fini, alors $R \to S/\mathfrak m$
est essentiellement de type fini par le
Lemme \ref{lemma-composition-essentially-of-finite-type}.
Réciproquement, supposons que $R \to S/\mathfrak m$ soit essentiellement
de type fini. Supposons que $S/\mathfrak m$ soit la localisation
de $R[x_1, \ldots, x_n]/I$. Choisissons $f_1, \ldots, f_n \in S$
dont les classes modulo $\mathfrak m$ correspondent aux classes de
$x_1, \ldots, x_n$ modulo $I$.
Considérons le morphisme $R[x_1, \ldots, x_n] \to S$, $x_i \mapsto f_i$
de noyau $J$. Posons $A = R[x_1, \ldots, x_n]/J \subset S$
et $\mathfrak p = A \cap \mathfrak m$. Remarquons que
$A/\mathfrak p \subset S/\mathfrak m$ est égal à l'image
de $R[x_1, \ldots, x_n]/I$ dans $S/\mathfrak m$. Ainsi
$\kappa(\mathfrak p) = S/\mathfrak m$. Donc $A_\mathfrak p \to S$
est fini par (1). Nous concluons que $S$ est essentiellement de type
fini par le Lemme \ref{lemma-composition-essentially-of-finite-type}.
\end{proof}

\noindent
Le lemme suivant peut être démontré en utilisant la propreté de l'espace
projectif au lieu de l'argument algébrique que nous donnons ici.

\begin{lemma}
\label{lemma-localization-at-closed-point-special-fibre}
Soit $\varphi : R \to S$ essentiellement de type fini, avec $R$ et $S$
locaux (mais $\varphi$ n'est pas nécessairement local). Alors il existe
un $n$ et un idéal maximal $\mathfrak m \subset R[x_1, \ldots, x_n]$
au-dessus de $\mathfrak m_R$ tels que $S$ soit une localisation d'un
quotient de $R[x_1, \ldots, x_n]_\mathfrak m$.
\end{lemma}

\begin{proof}
Nous pouvons écrire $S$ comme une localisation d'un quotient de
$R[x_1, \ldots, x_n]$. Il suffit donc de démontrer le lemme dans le cas où
$S = R[x_1, \ldots, x_n]_\mathfrak q$ pour un idéal premier
$\mathfrak q \subset R[x_1, \ldots, x_n]$.
Si $\mathfrak q + \mathfrak m_R R[x_1, \ldots, x_n] \not = R[x_1, \ldots, x_n]$
alors nous pouvons trouver un idéal maximal $\mathfrak m$ comme dans l'énoncé
du lemme avec $\mathfrak q \subset \mathfrak m$, et le résultat est clair.

\medskip\noindent
Choisissons un anneau de valuation $A \subset \kappa(\mathfrak q)$
qui domine l'image de $R \to \kappa(\mathfrak q)$
(Lemme \ref{lemma-dominate}). Si l'image $\lambda_i \in \kappa(\mathfrak q)$
de $x_i$ est contenue dans $A$, alors $\mathfrak q$ est contenue dans
l'image réciproque de $\mathfrak m_A$ par $R[x_1, \ldots, x_n] \to A$,
ce qui nous ramène au cas précédent.
Il existe donc un $i$ tel que $\lambda_i^{-1} \in A$ et tel que
$\lambda_j/\lambda_i \in A$ pour tout $j = 1, \ldots, n$
(car le groupe des valeurs de $A$ est totalement ordonné, voir le
Lemme \ref{lemma-valuation-group}). Nous considérons alors le morphisme
$$
R[y_0, y_1, \ldots, \hat{y_i}, \ldots, y_n]
\to R[x_1, \ldots, x_n]_\mathfrak q,\quad
y_0 \mapsto 1/x_i,\quad y_j \mapsto x_j/x_i
$$
Soit $\mathfrak q' \subset R[y_0, \ldots, \hat{y_i}, \ldots, y_n]$
l'image réciproque de $\mathfrak q$. Comme $y_0 \not \in \mathfrak q'$, il est
facile de voir que la flèche affichée définit un isomorphisme entre les
localisations. D'autre part, le résultat du premier paragraphe s'applique à
$R[y_0, \ldots, \hat{y_i}, \ldots, y_n]$
car $y_j$ s'envoie sur un élément de $A$.
Ceci achève la démonstration.
\end{proof}

\section{Groupes K}
\label{section-K-groups}

\noindent
Soit $R$ un anneau. Nous allons introduire deux groupes abéliens associés
à $R$. Le premier est noté $K'_0(R)$ et possède les propriétés suivantes
\footnote{La définition a un sens pour tout anneau, mais elle est rarement
utilisée lorsque $R$ n'est pas noethérien.} :
\begin{enumerate}
\item Pour tout $R$-module fini $M$, un élément $[M]$ de
$K'_0(R)$ est donné,
\item pour toute suite exacte courte $0 \to M' \to M \to M'' \to 0$
de $R$-modules finis, nous avons la relation $[M] = [M'] + [M'']$,
\item le groupe $K'_0(R)$ est engendré par les éléments $[M]$, et
\item toutes les relations dans $K'_0(R)$ entre les générateurs $[M]$
sont des combinaisons $\mathbf{Z}$-linéaires
des relations provenant des suites exactes comme ci-dessus.
\end{enumerate}
La construction effective est un peu plus délicate, car il faut prendre garde
à ce que la collection de tous les $R$-modules de type fini soit une classe
propre. Cependant, on peut surmonter ce problème en prenant comme ensemble
de générateurs du groupe $K'_0(R)$ les éléments $[R^n/K]$, où $n$ parcourt
les entiers et où $K$ parcourt tous les sous-modules $K \subset R^n$.
Les générateurs du sous-groupe des relations imposées sur ces éléments
seront les relations provenant des suites exactes courtes dont les termes
sont de la forme $R^n/K$. L'élément $[M]$ est défini en
choisissant $n$ et $K$ tels que $M \cong R^n/K$ et en posant
$[M] = [R^n/K]$. Les détails sont laissés au lecteur.

\begin{lemma}
\label{lemma-length-K0}
Si $R$ est un anneau local artinien, la fonction longueur
définit un homomorphisme naturel de groupes abéliens
$\text{length}_R : K'_0(R) \to \mathbf{Z}$.
\end{lemma}

\begin{proof}
La longueur de tout $R$-module fini est finie,
car il est quotient de $R^n$, qui est de longueur finie par
le Lemme \ref{lemma-artinian-finite-length}. Et la fonction longueur
est additive, voir le Lemme \ref{lemma-length-additive}.
\end{proof}

\noindent
Le second est noté $K_0(R)$ et possède les propriétés suivantes :
\begin{enumerate}
\item Pour tout $R$-module projectif fini $M$, un élément $[M]$ de
$K_0(R)$ est donné,
\item pour toute suite exacte courte $0 \to M' \to M \to M'' \to 0$
de $R$-modules projectifs finis, nous avons la relation $[M] = [M'] + [M'']$,
\item le groupe $K_0(R)$ est engendré par les éléments $[M]$, et
\item toutes les relations dans $K_0(R)$ sont des combinaisons
$\mathbf{Z}$-linéaires des relations provenant des suites exactes comme ci-dessus.
\end{enumerate}
La construction de ce groupe se fait comme ci-dessus.

\medskip\noindent
Nous notons qu'il existe un morphisme évident $K_0(R) \to K'_0(R)$
qui n'est pas un isomorphisme en général.

\begin{example}
\label{example-K0-field}
Notons que si $R = k$ est un corps, alors nous avons clairement
$K_0(k) = K'_0(k) \cong \mathbf{Z}$, l'isomorphisme
étant donné par la fonction dimension (qui est aussi la fonction longueur).
\end{example}

\begin{example}
\label{example-K0-PID}
Soit $R$ un anneau principal. Nous affirmons que $K_0(R) = K'_0(R) = \mathbf{Z}$.
En effet, tout $R$-module projectif fini est libre fini.
Un module libre fini possède un rang bien défini par le Lemme \ref{lemma-rank}.
Étant donnée une suite exacte courte de modules libres finis
$$
0 \to M' \to M \to M'' \to 0
$$
nous avons $\text{rank}(M) = \text{rank}(M') + \text{rank}(M'')$
car dans ce cas $M \cong M' \oplus M'$ (par exemple,
la suite se scinde par le Lemme \ref{lemma-lift-map}).
Nous en concluons que $K_0(R) = \mathbf{Z}$.

\medskip\noindent
Le théorème de structure des modules sur un anneau principal dit que
tout $R$-module de type fini est de la forme
$M = R^{\oplus r} \oplus R/(d_1) \oplus \ldots \oplus R/(d_k)$.
Considérons la suite exacte courte
$$
0 \to (d_i) \to R \to R/(d_i) \to 0
$$
Comme l'idéal $(d_i)$ est isomorphe à $R$ comme module
(il est libre de générateur $d_i$), dans $K'_0(R)$ nous avons
$[(d_i)] = [R]$. Alors $[R/(d_i)] = [(d_i)]-[R] = 0$. Il
s'ensuit qu'un module de torsion a une classe nulle dans $K'_0(R)$.
En utilisant le rang de la partie libre, on obtient une identification
$K'_0(R) = \mathbf{Z}$ et l'homomorphisme canonique
$K_0(R) \to K'_0(R)$ est un isomorphisme.
\end{example}

\begin{example}
\label{example-K0-polynomial-ring}
Soit $k$ un corps. Alors $K_0(k[x]) = K'_0(k[x]) = \mathbf{Z}$.
Ceci découle de l'Exemple \ref{example-K0-PID}, puisque
$R = k[x]$ est un anneau principal.
\end{example}

\begin{example}
\label{example-K0-node}
Soit $k$ un corps. Soit $R = \{f \in k[x] \mid f(0) = f(1)\}$,
voir l'Exemple \ref{example-affine-open-not-standard}.
Dans ce cas, $K_0(R) \cong k^* \oplus \mathbf{Z}$, mais
$K'_0(R) = \mathbf{Z}$.
\end{example}

\begin{lemma}
\label{lemma-K0-product}
Soit $R = R_1 \times R_2$. Alors $K_0(R) = K_0(R_1) \times K_0(R_2)$
et $K'_0(R) = K'_0(R_1) \times K'_0(R_2)$.
\end{lemma}

\begin{proof}
Omis.
\end{proof}

\begin{lemma}
\label{lemma-K0prime-Artinian}
Soit $R$ un anneau local artinien.
L'application $\text{length}_R : K'_0(R) \to \mathbf{Z}$
du Lemme \ref{lemma-length-K0} est un isomorphisme.
\end{lemma}

\begin{proof}
Omis.
\end{proof}

\begin{lemma}
\label{lemma-K0-local}
Soit $(R, \mathfrak m)$ un anneau local. Tout $R$-module projectif fini
est libre fini. L'application $\text{rank}_R : K_0(R) \to \mathbf{Z}$
donnée par $[M] \to \text{rank}_R(M)$ est bien définie
et est un isomorphisme.
\end{lemma}

\begin{proof}
Soit $P$ un $R$-module projectif fini. Choisissons des éléments
$x_1, \ldots, x_n \in P$ qui se projettent sur une base de $P/\mathfrak m P$.
Par le Lemme de Nakayama \ref{lemma-NAK}, ces éléments engendrent $P$.
La surjection correspondante $u : R^{\oplus n} \to P$ admet une section
car $P$ est projectif. Donc $R^{\oplus n} = P \oplus Q$ avec $Q = \Ker(u)$.
Il s'ensuit que $Q/\mathfrak m Q = 0$, donc $Q$ est nul par
le lemme de Nakayama. Nous voyons ainsi que tout $R$-module projectif
fini est libre fini.
Un module libre fini possède un rang bien défini par le Lemme \ref{lemma-rank}.
Étant donnée une suite exacte courte de $R$-modules libres finis
$$
0 \to M' \to M \to M'' \to 0
$$
nous avons $\text{rank}(M) = \text{rank}(M') + \text{rank}(M'')$
car dans ce cas $M \cong M' \oplus M'$ (par exemple,
la suite se scinde par le Lemme \ref{lemma-lift-map}).
Nous en concluons que $K_0(R) = \mathbf{Z}$.
\end{proof}

\begin{lemma}
\label{lemma-K0-and-K0prime-Artinian-local}
Soit $R$ un anneau artinien local. Il existe un diagramme
commutatif
$$
\xymatrix{
K_0(R) \ar[rr] \ar[d]_{\text{rank}_R} & &
K'_0(R) \ar[d]^{\text{length}_R} \\
\mathbf{Z} \ar[rr]^{\text{length}_R(R)} & & 
\mathbf{Z}
}
$$
où les applications verticales sont des isomorphismes par
les Lemmes \ref{lemma-K0prime-Artinian} et \ref{lemma-K0-local}.
\end{lemma}

\begin{proof}
Soit $P$ un $R$-module projectif fini. Nous devons montrer que
$\text{length}_R(P) = \text{rank}_R(P) \text{length}_R(R)$.
Par le Lemme \ref{lemma-K0-local}, le module $P$ est libre fini. Donc
$P \cong R^{\oplus n}$ pour un certain $n \geq 0$. Alors
$\text{rank}_R(P) = n$ et
$\text{length}_R(R^{\oplus n}) = n \text{length}_R(R)$
par additivité des longueurs (Lemme \ref{lemma-length-additive}).
Le résultat s'ensuit.
\end{proof}


\section{Anneaux gradués}
\label{section-graded}

\noindent
Un {\it anneau gradué} sera pour nous un anneau $S$ muni
d'une décomposition en somme directe $S = \bigoplus_{d \geq 0} S_d$
du groupe abélien sous-jacent, telle que $S_d \cdot S_e \subset S_{d + e}$.
Remarquons que nous n'autorisons pas d'éléments non nuls en degrés négatifs.
L'{\it idéal non pertinent} est l'idéal $S_{+} = \bigoplus_{d > 0} S_d$.
Un {\it module gradué}
sera un $S$-module $M$ muni d'une décomposition en somme directe
$M = \bigoplus_{n\in \mathbf{Z}} M_n$ du groupe abélien sous-jacent
telle que $S_d \cdot M_e \subset M_{d + e}$. Remarquons que pour les modules
nous autorisons des éléments non nuls en degrés négatifs.
Nous considérons $S$ comme un $S$-module gradué en posant $S_{-k} = (0)$
pour $k > 0$. Un élément $x$ (resp.\ $f$) de $M$ (resp.\ $S$) est appelé
{\it homogène}
si $x \in M_d$ (resp.\ $f \in S_d$) pour un certain $d$.
Un {\it morphisme de $S$-modules gradués} est un morphisme de $S$-modules
$\varphi : M \to M'$ tel que $\varphi(M_d) \subset M'_d$.
Nous n'autorisons pas les morphismes à décaler les degrés. Notons
$\text{GrHom}_0(M, N)$ le $S_0$-module des homomorphismes
de modules gradués de $M$ vers $N$.

\medskip\noindent
À ce stade interviennent les notions d'idéal gradué,
d'anneau quotient gradué, de sous-module gradué, de module quotient gradué,
de produit tensoriel gradué, etc. Nous laissons au lecteur le soin de trouver
les définitions et les lemmes pertinents. Par exemple : une suite exacte courte
de modules gradués est exacte courte à chaque degré.

\medskip\noindent
Étant donné un anneau gradué $S$, un $S$-module gradué $M$ et $n \in \mathbf{Z}$,
nous notons $M(n)$ le $S$-module gradué tel que $M(n)_d = M_{n + d}$.
C'est ce que l'on appelle le {\it tordu de $M$ par $n$}. En particulier, nous obtenons
les modules $S(n)$, $n \in \mathbf{Z}$, qui joueront un rôle important
dans l'étude des schémas projectifs. Il existe des isomorphismes
fonctoriels évidents tels que
$(M \oplus N)(n) = M(n) \oplus N(n)$,
$(M \otimes_S N)(n) = M \otimes_S N(n) = M(n) \otimes_S N$.
De plus, nous pouvons définir une structure de $S$-module gradué sur
le $S_0$-module
$$
\text{GrHom}(M, N) =
\bigoplus\nolimits_{n \in \mathbf{Z}} \text{GrHom}_n(M, N),
\quad
\text{GrHom}_n(M, N) = \text{GrHom}_0(M, N(n)).
$$
Nous omettons la définition de la multiplication.

\begin{lemma}
\label{lemma-graded-NAK}
Soit $S$ un anneau gradué. Soit $M$ un $S$-module gradué.
\begin{enumerate}
\item Si $S_+M = M$ et si $M$ est fini, alors $M = 0$.
\item Si $N, N' \subset M$ sont des sous-modules gradués,
$M = N + S_+N'$, et si $N'$ est fini, alors $M = N$.
\item Si $N \to M$ est un morphisme de modules gradués, si $N/S_+N \to M/S_+M$
est surjectif, et si $M$ est fini, alors $N \to M$ est surjectif.
\item Si $x_1, \ldots, x_n \in M$ sont homogènes, engendrent $M/S_+M$
et si $M$ est fini, alors $x_1, \ldots, x_n$ engendrent $M$.
\end{enumerate}
\end{lemma}

\begin{proof}
Preuve de (1). Choisissons des générateurs $y_1, \ldots, y_r$ de $M$ sur $S$.
Nous pouvons supposer que $y_i$ est homogène de degré $d_i$. Après
réindexation, nous pouvons supposer $d_r = \min(d_i)$. Alors la condition
$S_+M = M$ implique $y_r = 0$. Donc $M = 0$ par récurrence sur $r$.
La partie (2) s'obtient en appliquant (1) à $M/N$. La partie (3) s'obtient en
appliquant (2) aux sous-modules $\Im(N \to M)$ et $M$.
La partie (4) s'obtient en appliquant (3) au morphisme de modules
$\bigoplus S(-d_i) \to M$, $(s_1, \ldots, s_n) \mapsto \sum s_i x_i$.
\end{proof}

\noindent
Soit $S$ un anneau gradué. Soit $d \geq 1$ un entier.
Posons $S^{(d)} = \bigoplus_{n \geq 0} S_{nd}$. Nous considérons
$S^{(d)}$ comme un anneau gradué dont la composante de degré $n$ est
$(S^{(d)})_n = S_{nd}$. Étant donné un $S$-module gradué $M$, nous
pouvons de même considérer $M^{(d)} = \bigoplus_{n \in \mathbf{Z}} M_{nd}$,
qui est un $S^{(d)}$-module gradué.

\begin{lemma}
\label{lemma-uple-generated-degree-1}
Soit $S$ un anneau gradué, qui est de type fini sur $S_0$.
Alors, pour tout $d$ suffisamment divisible, l'algèbre
$S^{(d)}$ est engendrée en degré $1$ sur $S_0$.
\end{lemma}

\begin{proof}
Disons que $S$ est engendré par $f_1, \ldots, f_r \in S$ sur $S_0$.
Après remplacement des $f_i$ par leurs composantes homogènes, nous pouvons
supposer que $f_i$ est homogène de degré $d_i > 0$. Alors tout élément de
$S_n$ est une combinaison linéaire à coefficients dans $S_0$ de monômes
$f_1^{e_1} \ldots f_r^{e_r}$ tels que $\sum e_i d_i = n$.
Soit $m$ un multiple de $\text{lcm}(d_i)$. Pour tout $N \geq r$, si
$$
\sum e_i d_i = N m
$$
alors il existe un $i$ tel que $e_i \geq m/d_i$, par un argument élémentaire.
Ainsi tout monôme de degré $N m$ est le produit d'un monôme
de degré $m$, à savoir $f_i^{m/d_i}$, et d'un monôme de degré $(N - 1)m$.
Il s'ensuit que tout monôme de degré $nrm$, pour $n \geq 2$,
est un produit de monômes de degré $rm$. Ainsi $S^{(rm)}$ est engendrée
en degré $1$ sur $S_0$.
\end{proof}

\begin{lemma}
\label{lemma-integral-closure-graded}
\begin{reference}
\cite[Theorem 2.3.2]{Huneke-Swanson}
\end{reference}
Soit $R \to S$ un homomorphisme d'anneaux gradués.
Soit $S' \subset S$ la clôture intégrale de $R$ dans $S$.
Alors
$$
S' = \bigoplus\nolimits_{d \geq 0} S' \cap S_d,
$$
c'est-à-dire que $S'$ est une $R$-sous-algèbre graduée de $S$.
\end{lemma}

\begin{proof}
Nous devons montrer ce qui suit : si
$s = s_n + s_{n + 1} + \ldots + s_m \in S'$, alors chaque
composante homogène $s_j \in S'$. Nous le démontrerons par récurrence
sur $m - n$ pour tous les homomorphismes $R \to S$ d'anneaux gradués.
Tout d'abord, il est immédiat que $s_0$ est entier sur $R_0$ (et donc sur $R$),
car il existe un morphisme d'anneaux $S \to S_0$ compatible avec le morphisme
d'anneaux $R \to R_0$. Ainsi, après avoir remplacé $s$ par $s - s_0$, nous
pouvons supposer $n > 0$. Considérons l'extension d'anneaux gradués
$R[t, t^{-1}] \to S[t, t^{-1}]$ où $t$ est de degré $0$. Il existe un diagramme
commutatif
$$
\xymatrix{
S[t, t^{-1}] \ar[rr]_{s \mapsto t^{\deg(s)}s} & & S[t, t^{-1}] \\
R[t, t^{-1}] \ar[u] \ar[rr]^{r \mapsto t^{\deg(r)}r} & &  R[t, t^{-1}] \ar[u]
}
$$
où les morphismes horizontaux sont des automorphismes d'anneaux. Ainsi la clôture
intégrale $C$ de $S[t, t^{-1}]$ sur $R[t, t^{-1}]$ se transforme en elle-même.
Nous voyons donc que
$$
t^m(s_n + s_{n + 1} + \ldots + s_m) -
(t^ns_n + t^{n + 1}s_{n + 1} + \ldots + t^ms_m) \in C
$$
ce qui implique, par l'hypothèse de récurrence, que chaque $(t^m - t^i)s_i \in C$
pour $i = n, \ldots, m - 1$. Remarquons que pour tout anneau $A$ et
$m > i \geq n > 0$, nous avons
$A[t, t^{-1}]/(t^m - t^i - 1) \cong A[t]/(t^m - t^i - 1) \supset A$
car $t(t^{m - 1} - t^{i - 1}) = 1$ dans $A[t]/(t^m - t^i - 1)$.
Comme $t^m - t^i$ s'envoie sur $1$, nous voyons que l'image de $s_i$ dans
l'anneau $S[t]/(t^m - t^i - 1)$ est entière sur
$R[t]/(t^m - t^i - 1)$ pour $i = n, \ldots, m - 1$. Comme
$R \to R[t]/(t^m - t^i - 1)$ est fini, nous voyons que $s_i$ est entier
sur $R$ par transitivité, voir le Lemme \ref{lemma-integral-transitive}.
Enfin, nous concluons également que $s_m = s - \sum_{i = n, \ldots, m - 1} s_i$
est entier sur $R$.
\end{proof}

\section{Proj d'un anneau gradué}
\label{section-proj}

\noindent
Soit $S$ un anneau gradué.
Un {\it idéal homogène} est simplement un idéal
$I \subset S$ qui est aussi un sous-module gradué de $S$.
De manière équivalente, c'est un idéal engendré par des éléments homogènes.
De manière équivalente, si $f \in I$ et
$$
f = f_0 + f_1 + \ldots + f_n
$$
est la décomposition de $f$ en composantes homogènes dans $S$, alors $f_i \in I$
pour tout $i$. Pour vérifier qu'un idéal homogène $\mathfrak p$
est premier, il suffit de vérifier que si $ab \in \mathfrak p$
avec $a, b$ homogènes, alors soit $a \in \mathfrak p$, soit
$b \in \mathfrak p$.

\begin{definition}
\label{definition-proj}
Soit $S$ un anneau gradué.
Nous définissons $\text{Proj}(S)$ comme l'ensemble des idéaux
premiers homogènes $\mathfrak p$ de $S$ tels que
$S_{+} \not \subset \mathfrak p$.
L'ensemble $\text{Proj}(S)$ est un sous-ensemble de $\Spec(S)$
et nous le munissons de la topologie induite.
L'espace topologique $\text{Proj}(S)$ est appelé le
{\it spectre homogène} de l'anneau gradué $S$.
\end{definition}

\noindent
Remarquons que, par construction, il existe un morphisme continu
$$
\text{Proj}(S) \longrightarrow \Spec(S_0).
$$

\medskip\noindent
Soit $S = \oplus_{d \geq 0} S_d$ un anneau gradué.
Soit $f\in S_d$ et supposons que $d \geq 1$.
Nous définissons $S_{(f)}$ comme le sous-anneau de $S_f$
constitué des éléments de la forme $r/f^n$, avec $r$ homogène et
$\deg(r) = nd$. Si $M$ est un $S$-module gradué,
nous définissons le $S_{(f)}$-module $M_{(f)}$ comme le
sous-module de $M_f$ constitué des éléments de
la forme $x/f^n$, où $x$ est homogène de degré $nd$.

\begin{lemma}
\label{lemma-Z-graded}
Soit $S$ un anneau gradué par $\mathbf{Z}$ contenant un
élément inversible homogène de degré strictement positif. Alors l'ensemble
$G \subset \Spec(S)$ des idéaux premiers gradués par $\mathbf{Z}$ de $S$
(muni de la topologie induite) s'envoie homéomorphiquement sur $\Spec(S_0)$.
\end{lemma}

\begin{proof}
Montrons d'abord que le morphisme est une bijection en construisant un inverse.
Soit $f \in S_d$, $d > 0$, inversible dans $S$.
Si $\mathfrak p_0$ est un idéal premier de $S_0$, alors
$\mathfrak p_0S$ est un idéal gradué par $\mathbf{Z}$ de $S$ tel que
$\mathfrak p_0S \cap S_0 = \mathfrak p_0$. Et si $ab \in \mathfrak p_0S$
avec $a$, $b$ homogènes, alors
$a^db^d/f^{\deg(a) + \deg(b)} \in \mathfrak p_0$.
Ainsi, soit $a^d/f^{\deg(a)} \in \mathfrak p_0$, soit
$b^d/f^{\deg(b)} \in \mathfrak p_0$, autrement dit soit
$a^d \in \mathfrak p_0S$, soit $b^d \in \mathfrak p_0S$.
Il s'ensuit que $\sqrt{\mathfrak p_0S}$ est un idéal premier
gradué par $\mathbf{Z}$ de $S$ dont l'intersection avec $S_0$ est $\mathfrak p_0$.

\medskip\noindent
Pour montrer que le morphisme est un homéomorphisme, montrons que
l'image de $G \cap D(g)$ est ouverte. Si $g = \sum g_i$
avec $g_i \in S_i$, alors, d'après ce qui précède, $G \cap D(g)$
s'envoie sur l'ensemble $\bigcup D(g_i^d/f^i)$, qui est ouvert.
\end{proof}

\noindent
Pour $f \in S$ homogène de degré $> 0$, nous définissons
$$
D_{+}(f) = \{ \mathfrak p \in \text{Proj}(S) \mid f \not\in \mathfrak p \}.
$$
Enfin, pour un idéal homogène $I \subset S$, nous définissons
$$
V_{+}(I) = \{ \mathfrak p \in \text{Proj}(S) \mid I \subset \mathfrak p \}.
$$
Nous utiliserons plus généralement la notation $V_{+}(E)$ pour tout
ensemble $E$ d'éléments homogènes $E \subset S$.

\begin{lemma}[Topologie sur Proj]
\label{lemma-topology-proj}
Soit $S = \oplus_{d \geq 0} S_d$ un anneau gradué.
\begin{enumerate}
\item Les ensembles $D_{+}(f)$ sont ouverts dans $\text{Proj}(S)$.
\item Nous avons $D_{+}(ff') = D_{+}(f) \cap D_{+}(f')$.
\item Soit $g = g_0 + \ldots + g_m$ un élément
de $S$ avec $g_i \in S_i$. Alors
$$
D(g) \cap \text{Proj}(S) =
(D(g_0) \cap \text{Proj}(S))
\cup
\bigcup\nolimits_{i \geq 1} D_{+}(g_i).
$$
\item
Soit $g_0\in S_0$ un élément homogène de degré $0$. Alors
$$
D(g_0) \cap \text{Proj}(S)
=
\bigcup\nolimits_{f \in S_d, \ d\geq 1} D_{+}(g_0 f).
$$
\item Les ouverts $D_{+}(f)$ forment une
base de la topologie de $\text{Proj}(S)$.
\item Soit $f \in S$ homogène de degré positif.
L'anneau $S_f$ possède une graduation naturelle par $\mathbf{Z}$.
Les morphismes d'anneaux $S \to S_f \leftarrow S_{(f)}$ induisent
des homéomorphismes
$$
D_{+}(f)
\leftarrow
\{\mathbf{Z}\text{-idéaux premiers gradués de }S_f\}
\to
\Spec(S_{(f)}).
$$
\item Il existe un $S$ tel que $\text{Proj}(S)$ ne soit pas
quasi-compact.
\item Les ensembles $V_{+}(I)$ sont fermés.
\item Tout fermé $T \subset \text{Proj}(S)$ est de
la forme $V_{+}(I)$ pour un certain idéal homogène $I \subset S$.
\item Pour tout idéal gradué $I \subset S$, nous avons
$V_{+}(I) = \emptyset$ si et seulement si $S_{+} \subset \sqrt{I}$.
\end{enumerate}
\end{lemma}

\begin{proof}
Puisque $D_{+}(f) = \text{Proj}(S) \cap D(f)$, ces ensembles sont ouverts.
Ceci démontre (1). De même, (2) découle de $D(ff') = D(f) \cap D(f')$.
De manière analogue, les ensembles $V_{+}(I) = \text{Proj}(S) \cap V(I)$
sont fermés. Ceci démontre (8).

\medskip\noindent
Supposons que $T \subset \text{Proj}(S)$ soit fermé.
Alors nous pouvons écrire $T = \text{Proj}(S) \cap V(J)$ pour un
idéal $J \subset S$. Par définition d'un idéal homogène,
si $g \in J$, $g = g_0 + \ldots + g_m$
avec $g_d \in S_d$, alors $g_d \in \mathfrak p$ pour tout
$\mathfrak p \in T$. Ainsi, si $I \subset S$
est l'idéal engendré par les composantes homogènes des éléments
de $J$, nous avons $T = V_{+}(I)$. Ceci démontre (9).

\medskip\noindent
La formule pour $\text{Proj}(S) \cap D(g)$, avec $g \in S$, découle directement
des définitions. Ceci démontre (3).
Considérons la formule pour $\text{Proj}(S) \cap D(g_0)$.
L'inclusion du membre de droite dans le membre de gauche est
évidente. Pour l'autre inclusion, supposons $g_0 \not \in \mathfrak p$
avec $\mathfrak p \in \text{Proj}(S)$. Si $g_0f \in \mathfrak p$
pour tout $f$ homogène de degré positif, alors nous voyons que
$S_{+} \subset \mathfrak p$, ce qui est une contradiction. Cela donne
l'autre inclusion. Ceci démontre (4).

\medskip\noindent
La collection des ouverts $D(g) \cap \text{Proj}(S)$
forme une base de la topologie, puisque les ouverts standards
$D(g) \subset \Spec(S)$ forment une base de la topologie de
$\Spec(S)$. D'après les formules ci-dessus, nous pouvons exprimer
$D(g) \cap \text{Proj}(S)$ comme une réunion d'ouverts $D_{+}(f)$.
Ainsi la collection des ouverts $D_{+}(f)$ forme aussi une base de la topologie.
Ceci démontre (5).

\medskip\noindent
Preuve de (6). Notons d'abord que $D_{+}(f)$ peut être identifié
à un sous-ensemble (muni de la topologie induite) de $D(f) = \Spec(S_f)$
via le Lemme \ref{lemma-standard-open}. Remarquons que l'anneau
$S_f$ est gradué par $\mathbf{Z}$. Les éléments homogènes sont
de la forme $r/f^n$ avec $r \in S$ homogène et ont
pour degré $\deg(r/f^n) = \deg(r) - n\deg(f)$. Le sous-ensemble
$D_{+}(f)$ correspond exactement aux idéaux premiers
$\mathfrak p \subset S_f$ qui sont des idéaux gradués par $\mathbf{Z}$
(c'est-à-dire engendrés par des éléments homogènes). Il suffit donc de montrer que
l'ensemble des idéaux premiers gradués par $\mathbf{Z}$ de $S_f$ s'envoie
homéomorphiquement sur $\Spec(S_{(f)})$. Cela découle du Lemme \ref{lemma-Z-graded}.

\medskip\noindent
Soit $S = \mathbf{Z}[X_1, X_2, X_3, \ldots]$ avec une graduation telle que
chaque $X_i$ soit de degré $1$. Alors il est facile de voir que
$$
\text{Proj}(S) = \bigcup\nolimits_{i = 1}^\infty D_{+}(X_i)
$$
n'admet pas de sous-recouvrement fini. Ceci démontre (7).

\medskip\noindent
Soit $I \subset S$ un idéal gradué.
Si $\sqrt{I} \supset S_{+}$, alors $V_{+}(I) = \emptyset$, puisque
tout idéal premier $\mathfrak p \in \text{Proj}(S)$ ne contient pas
$S_{+}$ par définition. Réciproquement, supposons que
$S_{+} \not \subset \sqrt{I}$. Nous pouvons trouver un élément
$f \in S_{+}$ tel que $f$ ne soit pas nilpotent modulo $I$.
Cela signifie clairement que l'une des composantes homogènes de $f$
n'est pas nilpotente modulo $I$; autrement dit, nous pouvons (et allons)
supposer que $f$ est homogène. Cela implique que
$I S_f \not = S_f$, autrement dit que $(S/I)_f$ n'est pas
nul. Donc $(S/I)_{(f)} \not = 0$, puisque c'est un anneau qui
se plonge dans $(S/I)_f$. Choisissons un idéal premier
$\mathfrak q \subset (S/I)_{(f)}$. Celui-ci correspond à
un idéal premier gradué de $S/I$, qui ne contient pas l'idéal non pertinent
$(S/I)_{+}$. Cela correspond à son tour à un idéal premier
gradué $\mathfrak p$ de $S$, contenant $I$ mais ne contenant pas $S_{+}$,
comme voulu. Ceci démontre (10) et achève la démonstration.
\end{proof}

\begin{example}
\label{example-proj-polynomial-ring-1-variable}
Soit $R$ un anneau. Si $S = R[X]$ avec $\deg(X) = 1$, alors le morphisme naturel
$\text{Proj}(S) \to \Spec(R)$ est une bijection et même un homéomorphisme.
En effet, soit $\mathfrak p \in \text{Proj}(S)$. Puisque
$S_{+} \not \subset \mathfrak p$, nous voyons que $X \not \in \mathfrak p$.
Ainsi, si $aX^n \in \mathfrak p$ avec $a \in R$ et $n > 0$, alors
$a \in \mathfrak p$. Il s'ensuit que $\mathfrak p = \mathfrak p_0S$
avec $\mathfrak p_0 = \mathfrak p \cap R$.
\end{example}

\noindent
Si $\mathfrak p \in \text{Proj}(S)$, nous
définissons $S_{(\mathfrak p)}$ comme l'anneau dont les
éléments sont les fractions $r/f$, où $r, f \in S$ sont des
éléments homogènes de même degré tels que $f \not\in \mathfrak p$.
Comme d'habitude, nous disons que $r/f = r'/f'$ si et seulement s'il existe
un certain $f'' \in S$ homogène, $f'' \not \in \mathfrak p$, tel
que $f''(rf' - r'f) = 0$.
Étant donné un $S$-module gradué $M$, nous posons
$M_{(\mathfrak p)}$ égal au $S_{(\mathfrak p)}$-module
dont les éléments sont les fractions $x/f$, où $x \in M$
et $f \in S$ sont homogènes de même degré, avec
$f \not \in \mathfrak p$. Nous disons que $x/f = x'/f'$
si et seulement s'il existe un certain $f'' \in S$ homogène,
$f'' \not \in \mathfrak p$, tel que $f''(xf' - x'f) = 0$.

\begin{lemma}
\label{lemma-proj-prime}
Soit $S$ un anneau gradué. Soit $M$ un $S$-module gradué.
Soit $\mathfrak p$ un élément de $\text{Proj}(S)$.
Soit $f \in S$ un élément homogène de degré positif
tel que $f \not \in \mathfrak p$, c'est-à-dire $\mathfrak p \in D_{+}(f)$.
Soit $\mathfrak p' \subset S_{(f)}$ l'élément de
$\Spec(S_{(f)})$ correspondant à $\mathfrak p$ comme dans
le Lemme \ref{lemma-topology-proj}. Alors
$S_{(\mathfrak p)} = (S_{(f)})_{\mathfrak p'}$
et, de manière compatible,
$M_{(\mathfrak p)} = (M_{(f)})_{\mathfrak p'}$.
\end{lemma}

\begin{proof}
Nous définissons un morphisme $\psi : M_{(\mathfrak p)} \to (M_{(f)})_{\mathfrak p'}$.
Soit $x/g \in M_{(\mathfrak p)}$. Posons
$$
\psi(x/g) = (x g^{\deg(f) - 1}/f^{\deg(x)})/(g^{\deg(f)}/f^{\deg(g)}).
$$
Ceci a un sens puisque $\deg(x) = \deg(g)$ et puisque
$g^{\deg(f)}/f^{\deg(g)} \not \in \mathfrak p'$.
Nous omettons la vérification que $\psi$ est bien défini, est un morphisme
de modules et est un isomorphisme. Indication : l'inverse envoie
$(x/f^n)/(g/f^m)$ sur $(xf^m)/(g f^n)$.
\end{proof}

\noindent
Voici une variante graduée du Lemme \ref{lemma-silly}.

\begin{lemma}
\label{lemma-graded-silly}
Supposons que $S$ soit un anneau gradué, que les $\mathfrak p_i$, $i = 1, \ldots, r$,
soient des idéaux premiers homogènes et que $I \subset S_{+}$ soit un idéal gradué.
Supposons que $I \not\subset \mathfrak p_i$ pour tout $i$. Alors il
existe un élément homogène $x\in I$ de degré positif tel que
$x\not\in \mathfrak p_i$ pour tout $i$.
\end{lemma}

\begin{proof}
Nous pouvons supposer qu'il n'y a aucune inclusion entre les $\mathfrak p_i$.
Le résultat est vrai pour $r = 1$. Supposons le résultat vrai pour $r - 1$.
Choisissons $x \in I$ homogène de degré positif tel que
$x \not \in \mathfrak p_i$ pour tous les $i = 1, \ldots, r - 1$.
Si $x \not\in \mathfrak p_r$, c'est terminé. Supposons donc $x \in \mathfrak p_r$.
Si $I \mathfrak p_1 \ldots \mathfrak p_{r-1} \subset \mathfrak p_r$,
alors $I \subset \mathfrak p_r$, contradiction.
Choisissons $y \in I\mathfrak p_1 \ldots \mathfrak p_{r-1}$ homogène
et $y \not \in \mathfrak p_r$. Alors $x^{\deg(y)} + y^{\deg(x)}$ convient.
\end{proof}

\begin{lemma}
\label{lemma-smear-out}
Soit $S$ un anneau gradué.
Soit $\mathfrak p \subset S$ un idéal premier.
Soit $\mathfrak q$ l'idéal homogène de $S$ engendré par les
éléments homogènes de $\mathfrak p$. Alors $\mathfrak q$ est un
idéal premier de $S$.
\end{lemma}

\begin{proof}
Pour prouver que $\mathfrak q$ est premier, il suffit de vérifier que
si $f, g \in S$ sont homogènes et si $fg \in \mathfrak q$, alors
$f$ ou $g$ appartient à $\mathfrak q$. Alors $fg \in \mathfrak p$
car $\mathfrak p \subset \mathfrak q$. Comme $\mathfrak p$
est premier, nous voyons que soit $f \in \mathfrak p$, soit $g \in \mathfrak p$.
Comme $f$ et $g$ sont homogènes, il est alors clair que soit
$f \in \mathfrak q$, soit $g \in \mathfrak q$.
\end{proof}

\begin{lemma}
\label{lemma-graded-ring-minimal-prime}
Soit $S$ un anneau gradué.
\begin{enumerate}
\item Tout idéal premier minimal de $S$ est un idéal homogène de $S$.
\item Étant donné un idéal homogène $I \subset S$, tout idéal premier
minimal au-dessus de $I$ est homogène.
\end{enumerate}
\end{lemma}

\begin{proof}
La première assertion est vraie car l'idéal premier $\mathfrak q$ construit dans
le Lemme \ref{lemma-smear-out} vérifie $\mathfrak q \subset \mathfrak p$.
La seconde découle de ce que nous pouvons considérer $S/I$ et appliquer la première partie.
\end{proof}

\begin{lemma}
\label{lemma-dehomogenize-finite-type}
Soit $R$ un anneau. Soit $S$ une $R$-algèbre graduée. Soit $f \in S_{+}$
homogène. Supposons que $S$ soit de type fini sur $R$. Alors
\begin{enumerate}
\item l'anneau $S_{(f)}$ est de type fini sur $R$, et
\item pour tout $S$-module gradué fini $M$, le module $M_{(f)}$
est un $S_{(f)}$-module fini.
\end{enumerate}
\end{lemma}

\begin{proof}
Choisissons $f_1, \ldots, f_n \in S$ qui engendrent $S$ comme $R$-algèbre.
Nous pouvons supposer que chaque $f_i$ est homogène (en décomposant chaque $f_i$
en ses composantes homogènes). Un élément de $S_{(f)}$ est une somme
de la forme
$$
\sum\nolimits_{e\deg(f) =
\sum e_i\deg(f_i)} \lambda_{e_1 \ldots e_n} f_1^{e_1} \ldots f_n^{e_n}/f^e
$$
avec $\lambda_{e_1 \ldots e_n} \in R$. Ainsi $S_{(f)}$ est engendré
comme $R$-algèbre par les $f_1^{e_1} \ldots f_n^{e_n} /f^e$ vérifiant
$e\deg(f) = \sum e_i\deg(f_i)$. Si $e_i \geq \deg(f)$,
nous pouvons écrire ceci sous la forme
$$
f_1^{e_1} \ldots f_n^{e_n}/f^e =
f_i^{\deg(f)}/f^{\deg(f_i)} \cdot
f_1^{e_1} \ldots f_i^{e_i - \deg(f)} \ldots f_n^{e_n}/f^{e - \deg(f_i)}
$$
Ainsi, il ne nous faut que les éléments $f_i^{\deg(f)}/f^{\deg(f_i)}$, ainsi
que les éléments $f_1^{e_1} \ldots f_n^{e_n} /f^e$ tels que
$e \deg(f) = \sum e_i \deg(f_i)$ et $e_i < \deg(f)$.
Il s'agit d'une liste finie, et nous voyons que (1) est vraie.

\medskip\noindent
Pour voir (2), supposons que $M$ soit engendré par des éléments homogènes
$x_1, \ldots, x_m$. En raisonnant
comme ci-dessus, nous trouvons que $M_{(f)}$ est engendré comme $S_{(f)}$-module
par la liste finie des éléments de la forme
$f_1^{e_1} \ldots f_n^{e_n} x_j /f^e$
avec $e \deg(f) = \sum e_i \deg(f_i) + \deg(x_j)$ et
$e_i < \deg(f)$.
\end{proof}

\begin{lemma}
\label{lemma-homogenize}
Soit $R$ un anneau.
Soit $R'$ une $R$-algèbre de type fini, et soit $M$ un $R'$-module fini.
Il existe une $R$-algèbre graduée $S$, un $S$-module gradué $N$ et
un élément $f \in S$ homogène de degré $1$ tels que
\begin{enumerate}
\item $R' \cong S_{(f)}$ et $M \cong N_{(f)}$ (comme modules),
\item $S_0 = R$ et $S$ est engendré par un nombre fini d'éléments
de degré $1$ sur $R$, et
\item $N$ est un $S$-module fini.
\end{enumerate}
\end{lemma}

\begin{proof}
Nous pouvons écrire $R' = R[x_1, \ldots, x_n]/I$ pour un certain idéal $I$.
Pour un élément $g \in R[x_1, \ldots, x_n]$, notons
$\tilde g \in R[X_0, \ldots, X_n]$ l'élément homogène de degré minimal
tel que $g = \tilde g(1, x_1, \ldots, x_n)$.
Soit $\tilde I \subset R[X_0, \ldots, X_n]$ l'idéal engendré par tous les
éléments $\tilde g$, $g \in I$.
Posons $S = R[X_0, \ldots, X_n]/\tilde I$ et notons $f$ l'image
de $X_0$ dans $S$. Par construction, nous avons un isomorphisme
$$
S_{(f)} \longrightarrow R', \quad
X_i/X_0 \longmapsto x_i.
$$
Pour faire de même avec le module $M$, choisissons une présentation
$$
M = (R')^{\oplus r}/\sum\nolimits_{j \in J} R'k_j
$$
avec $k_j = (k_{1j}, \ldots, k_{rj})$. Posons $d_{ij} = \deg(\tilde k_{ij})$.
Posons $d_j = \max\{d_{ij}\}$. Posons $K_{ij} = X_0^{d_j - d_{ij}}\tilde k_{ij}$,
qui est homogène de degré $d_j$. Avec cette notation, posons
$$
N = \Coker\Big(
\bigoplus\nolimits_{j \in J} S(-d_j) \xrightarrow{(K_{ij})} S^{\oplus r}
\Big)
$$
ce qui convient. Certains détails sont omis.
\end{proof}

\section{Anneaux gradués noethériens}
\label{section-noetherian-graded}

\noindent
Quelques éléments de théorie des anneaux gradués noethériens,
notamment sur les polynômes de Hilbert.

\begin{lemma}
\label{lemma-S-plus-generated}
Soit $S$ un anneau gradué. Un ensemble d'éléments homogènes
$f_i \in S_{+}$ engendre $S$ comme algèbre sur $S_0$ si et seulement
si ces éléments engendrent $S_{+}$ comme idéal de $S$.
\end{lemma}

\begin{proof}
Si les $f_i$ engendrent $S$ comme algèbre sur $S_0$, tout élément
de $S_{+}$ est un polynôme sans terme constant en les $f_i$, et donc
$S_{+}$ est engendré par les $f_i$ comme idéal. Réciproquement, supposons que
$S_{+} = \sum Sf_i$. Nous montrerons que tout élément $f$ de $S$ s'écrit
comme un polynôme en les $f_i$ à coefficients dans $S_0$. Il suffit
de le faire pour les éléments homogènes. Supposons que $f$ soit de degré $d$.
Nous pouvons alors raisonner par récurrence sur $d$. Le cas $d = 0$ est immédiat.
Si $d > 0$, alors $f \in S_{+}$, donc nous pouvons écrire $f = \sum g_i f_i$
pour certains $g_i \in S$. Comme $S$ est gradué, nous pouvons remplacer
$g_i$ par sa composante homogène de degré $d - \deg(f_i)$. Par récurrence,
nous voyons que chaque $g_i$ est un polynôme en les $f_i$, et le résultat suit.
\end{proof}

\begin{lemma}
\label{lemma-graded-Noetherian}
Un anneau gradué $S$ est noethérien si et seulement si $S_0$ est
noethérien et si $S_{+}$ est de type fini comme idéal de $S$.
\end{lemma}

\begin{proof}
Il est clair que si $S$ est noethérien, alors $S_0 = S/S_{+}$ est noethérien
et $S_{+}$ est de type fini. Réciproquement, supposons $S_0$ noethérien
et $S_{+}$ de type fini comme idéal de $S$. Choisissons des générateurs
$S_{+} = (f_1, \ldots, f_n)$. En décomposant les $f_i$ en composantes homogènes,
nous pouvons supposer que chaque $f_i$ est homogène. Par le
Lemme \ref{lemma-S-plus-generated}, nous voyons que
$S_0[X_1, \ldots X_n] \to S$ qui envoie $X_i$ sur $f_i$
est surjectif. Ainsi $S$ est noethérien par le
Lemme \ref{lemma-Noetherian-permanence}.
\end{proof}

\begin{definition}
\label{definition-numerical-polynomial}
Soit $A$ un groupe abélien.
Nous disons qu'une fonction $f : n \mapsto f(n) \in A$
définie pour tous les entiers $n$ suffisamment grands est un
{\it polynôme numérique} s'il existe $r \geq 0$,
des éléments $a_0, \ldots, a_r\in A$ tels que
$$
f(n) = \sum\nolimits_{i = 0}^r \binom{n}{i} a_i
$$
pour tout $n \gg 0$.
\end{definition}

\noindent
La raison d'utiliser les coefficients binomiaux est le
fait élémentaire que tout polynôme $P \in \mathbf{Q}[T]$
dont toutes les valeurs aux points entiers sont entières est
égal à une somme $P(T) = \sum a_i \binom{T}{i}$ avec
$a_i \in \mathbf{Z}$. Remarquons qu'en particulier les
expressions $\binom{T + 1}{i + 1}$ sont de cette forme.

\begin{lemma}
\label{lemma-numerical-polynomial-functorial}
Si $A \to A'$ est un homomorphisme de groupes abéliens et si
$f : n \mapsto f(n) \in A$ est un polynôme numérique,
alors la composée l'est également.
\end{lemma}

\begin{proof}
Ceci est immédiat d'après les définitions.
\end{proof}

\begin{lemma}
\label{lemma-numerical-polynomial}
Supposons que $f: n \mapsto f(n) \in A$
soit définie pour tout $n$ suffisamment grand,
et supposons que $n \mapsto f(n) - f(n-1)$
soit un polynôme numérique. Alors $f$ est un
polynôme numérique.
\end{lemma}

\begin{proof}
Soit $f(n) - f(n-1) = \sum\nolimits_{i = 0}^r \binom{n}{i} a_i$
pour tout $n \gg 0$. Posons
$g(n) = f(n) - \sum\nolimits_{i = 0}^r \binom{n + 1}{i + 1} a_i$.
Alors $g(n) - g(n-1) = 0$ pour tout $n \gg 0$. Ainsi $g$ est
finalement constante, disons égale à $a_{-1}$. Nous laissons
au lecteur le soin de montrer que
$a_{-1} + \sum\nolimits_{i = 0}^r \binom{n + 1}{i + 1} a_i$
a la forme requise (voir la remarque au-dessus du lemme).
\end{proof}

\begin{lemma}
\label{lemma-graded-module-fg}
Si $M$ est un $S$-module gradué de type fini,
et si $S$ est de type fini sur $S_0$, alors
chaque $M_n$ est un $S_0$-module fini.
\end{lemma}

\begin{proof}
Supposons que les générateurs de $M$ soient les $m_i$ et les générateurs
de $S$ les $f_i$. En prenant les composantes homogènes, nous pouvons
supposer que les $m_i$ et les $f_i$ sont homogènes
et nous pouvons supposer $f_i \in S_{+}$. Dans ce cas, il est
clair que chaque $M_n$ est engendré sur $S_0$
par les « monômes » $\prod f_i^{e_i} m_j$ dont le
degré est $n$.
\end{proof}

\begin{proposition}
\label{proposition-graded-hilbert-polynomial}
Supposons que $S$ soit un anneau gradué noethérien
et $M$ un $S$-module gradué fini. Considérons la
fonction
$$
\mathbf{Z} \longrightarrow K'_0(S_0), \quad
n \longmapsto [M_n]
$$
voir le Lemme \ref{lemma-graded-module-fg}.
Si $S_{+}$ est engendré par des éléments de degré $1$,
alors cette fonction est un polynôme numérique.
\end{proposition}

\begin{proof}
Nous démontrons ceci par récurrence sur le nombre minimal de
générateurs de $S_1$. Si ce nombre est $0$, alors
$M_n = 0$ pour tout $n \gg 0$ et le résultat est vrai.
Pour l'étape de récurrence, soit $x\in S_1$
l'un des générateurs d'un ensemble minimal, tel que
l'hypothèse de récurrence s'applique à l'anneau gradué $S/(x)$.

\medskip\noindent
Montrons d'abord que le résultat est vrai si $x$ est nilpotent sur $M$.
Nous raisonnons par récurrence sur le plus petit entier $r$ tel que
$x^r M  = 0$. Si $r = 1$, alors $M$ est un module sur $S/xS$
et le résultat est vrai (par l'autre hypothèse de récurrence).
Si $r > 1$, nous pouvons trouver une suite exacte courte
$0 \to M' \to M \to M'' \to 0$ telle que les entiers
$r', r''$ soient strictement plus petits que $r$. Nous connaissons donc
le résultat pour $M''$ et $M'$. Ainsi
nous obtenons le résultat pour $M$ en raison de la relation
$
[M_d]  = [M'_d] + [M''_d]
$
dans $K'_0(S_0)$.

\medskip\noindent
Si $x$ n'est pas nilpotent sur $M$, soit $M' \subset M$ le
plus grand sous-module sur lequel $x$ est nilpotent.
Considérons la suite exacte $0 \to M' \to M \to M/M' \to 0$ :
nous voyons à nouveau qu'il suffit de prouver le résultat pour $M/M'$.
Autrement dit, nous pouvons supposer que la multiplication par $x$ est injective.

\medskip\noindent
Posons $\overline{M} = M/xM$. Remarquons que l'application $x : M \to M$
n'est {\it pas} une application de $S$-modules gradués, car elle
n'envoie pas $M_d$ dans $M_d$. En effet, pour chaque $d$, nous avons la
suite exacte courte suivante
$$
0 \to M_d \xrightarrow{x} M_{d + 1} \to \overline{M}_{d + 1} \to 0
$$
Ceci montre que $[M_{d + 1}] - [M_d] = [\overline{M}_{d + 1}]$.
Le résultat découle alors du Lemme \ref{lemma-numerical-polynomial}.
\end{proof}

\begin{remark}
\label{remark-period-polynomial}
Si $S$ est toujours noethérien mais $S$ n'est pas engendré en degré $1$,
alors la fonction associée à un $S$-module gradué est un
polynôme périodique (c'est-à-dire un polynôme numérique sur les
classes de congruence des entiers modulo $n$ pour un certain $n$).
\end{remark}

\begin{example}
\label{example-hilbert-function}
Supposons que $S = k[X_1, \ldots, X_d]$.
Par l'Exemple \ref{example-K0-field}, nous pouvons identifier
$K_0(k) = K'_0(k) = \mathbf{Z}$. Ainsi tout module gradué
de type fini sur $k[X_1, \ldots, X_d]$
donne lieu à un polynôme numérique
$n \mapsto \dim_k(M_n)$.
\end{example}

\begin{lemma}
\label{lemma-quotient-smaller-d}
Soit $k$ un corps. Supposons que $I \subset k[X_1, \ldots, X_d]$
soit un idéal gradué non nul. Soit $M = k[X_1, \ldots, X_d]/I$.
Alors le polynôme numérique $n \mapsto \dim_k(M_n)$ (voir
l'Exemple \ref{example-hilbert-function})
est de degré $ < d - 1$ (ou est nul si $d = 1$).
\end{lemma}

\begin{proof}
Le polynôme numérique associé au module gradué
$k[X_1, \ldots, X_d]$ est $n \mapsto \binom{n - 1 + d}{d - 1}$.
Pour tout élément homogène non nul $f \in I$ de degré $e$
et tout degré $n >> e$, nous avons $I_n \supset f \cdot k[X_1, \ldots, X_d]_{n-e}$
et donc $\dim_k(I_n) \geq \binom{n - e - 1 + d}{d - 1}$. Ainsi
$\dim_k(M_n) \leq \binom{n - 1 + d}{d - 1} - \binom{n - e - 1 + d}{d - 1}$.
Le résultat suit car la dernière expression
est de degré $ < d - 1$ (ou est nulle si $d = 1$).
\end{proof}

\section{Anneaux locaux noethériens}
\label{section-Noetherian-local}

\noindent
Dans toute cette section, $(R, \mathfrak m, \kappa)$ est un anneau local
noethérien. Nous y développons quelques éléments de théorie des fonctions
de Hilbert des modules. Soit $M$ un $R$-module de type fini. Nous définissons
la {\it fonction de Hilbert} de $M$ comme la fonction
$$
\varphi_M : n
\longmapsto
\text{longueur}_R(\mathfrak m^nM/{\mathfrak m}^{n + 1}M)
$$
définie pour tout entier $n \geq 0$. Un autre invariant important est la
fonction
$$
\chi_M : n
\longmapsto
\text{longueur}_R(M/{\mathfrak m}^{n + 1}M)
$$
définie pour tout entier $n \geq 0$.
D'après le Lemme \ref{lemma-length-additive}, nous avons
$$
\chi_M(n) = \sum\nolimits_{i = 0}^n \varphi_M(i).
$$
Il existe une variante de cette construction qui utilise un idéal de définition.

\begin{definition}
\label{definition-ideal-definition}
Soit $(R, \mathfrak m)$ un anneau local noethérien.
Un idéal $I \subset R$ tel que $\sqrt{I} = \mathfrak m$ est appelé
{\it idéal de définition de $R$}.
\end{definition}

\noindent
Soit $I \subset R$ un idéal de définition.
Comme $R$ est noethérien, cela implique que
$\mathfrak m^r \subset I$ pour un certain $r$; voir le Lemme
\ref{lemma-Noetherian-power}. Tout $R$-module de type fini
annulé par une puissance de $I$ est donc de longueur finie; voir le Lemme
\ref{lemma-length-finite}.
Il est ainsi légitime de définir
$$
\varphi_{I, M}(n) = \text{longueur}_R(I^nM/I^{n + 1}M)
\quad\text{et}\quad
\chi_{I, M}(n) = \text{longueur}_R(M/I^{n + 1}M)
$$
pour tout $n \geq 0$. Nous avons encore
$$
\chi_{I, M}(n) = \sum\nolimits_{i = 0}^n \varphi_{I, M}(i).
$$

\begin{lemma}
\label{lemma-differ-finite}
Supposons que $M' \subset M$ soient des $R$-modules de type fini
dont le quotient est de longueur finie. Il existe alors des
constantes $c_1, c_2$ telles que, pour tout $n \geq c_2$, nous ayons
$$
c_1 + \chi_{I, M'}(n - c_2) \leq \chi_{I, M}(n) \leq
c_1 + \chi_{I, M'}(n)
$$
\end{lemma}

\begin{proof}
Comme $M/M'$ est de longueur finie, il existe $c_2 \geq 0$ tel que
$I^{c_2}M \subset M'$. Posons $c_1 = \text{longueur}_R(M/M')$.
Pour $n \geq c_2$, nous avons
\begin{eqnarray*}
\chi_{I, M}(n)
& = &
\text{longueur}_R(M/I^{n + 1}M) \\
& = &
c_1 + \text{longueur}_R(M'/I^{n + 1}M) \\
& \leq &
c_1 + \text{longueur}_R(M'/I^{n + 1}M') \\
& = &
c_1 + \chi_{I, M'}(n)
\end{eqnarray*}
D'autre part, comme $I^{c_2}M \subset M'$,
nous avons $I^nM \subset I^{n - c_2}M'$ pour $n \geq c_2$.
Ainsi, pour $n \geq c_2$, nous obtenons
\begin{eqnarray*}
\chi_{I, M}(n)
& = &
\text{longueur}_R(M/I^{n + 1}M) \\
& = &
c_1 + \text{longueur}_R(M'/I^{n + 1}M) \\
& \geq &
c_1 + \text{longueur}_R(M'/I^{n + 1 - c_2}M') \\
& = &
c_1 + \chi_{I, M'}(n - c_2)
\end{eqnarray*}
ce qui achève la démonstration.
\end{proof}

\begin{lemma}
\label{lemma-hilbert-ses}
Supposons que $0 \to M' \to M \to M'' \to 0$
soit une suite exacte courte de $R$-modules de type fini.
Il existe alors un sous-module $N \subset M'$ dont le quotient est de
longueur finie $l$, et un entier $c \geq 0$, tels que
$$
\chi_{I, M}(n) = \chi_{I, M''}(n) + \chi_{I, N}(n - c) + l
$$
et
$$
\varphi_{I, M}(n) = \varphi_{I, M''}(n) + \varphi_{I, N}(n - c)
$$
pour tout $n \geq c$.
\end{lemma}

\begin{proof}
Remarquons que $M/I^nM \to M''/I^nM''$ est surjectif
et a pour noyau $M' / M' \cap I^nM$. D'après le lemme d'Artin--Rees
(Lemme \ref{lemma-Artin-Rees}), il existe une constante
$c$ telle que $M' \cap I^nM =
I^{n - c}(M' \cap I^cM)$. Posons $N = M' \cap I^cM$.
Remarquons que $I^c M' \subset N \subset M'$.
Ainsi $\text{longueur}_R(M' / M' \cap I^nM)
= \text{longueur}_R(M'/N) + \text{longueur}_R(N/I^{n - c}N)$ pour $n \geq c$.
La suite exacte courte
$$
0 \to M' / M' \cap I^nM \to M/I^nM \to M''/I^nM'' \to 0
$$
et l'additivité des longueurs (Lemme \ref{lemma-length-additive})
donnent l'égalité
$$
\chi_{I, M}(n - 1)
=
\chi_{I, M''}(n - 1)
+
\chi_{I, N}(n - c - 1)
+
\text{longueur}_R(M'/N)
$$
pour $n \geq c$. Nous avons
$\varphi_{I, M}(n) = \chi_{I, M}(n) - \chi_{I, M}(n - 1)$,
et de même pour les modules $M''$ et $N$. Nous obtenons donc
$\varphi_{I, M}(n) = \varphi_{I, M''}(n) + \varphi_{I, N}(n-c)$ pour
$n \geq c$.
\end{proof}

\begin{lemma}
\label{lemma-hilbert-change-I}
Supposons que $I$, $I'$ soient deux idéaux de définition
de l'anneau local noethérien $R$. Soit $M$ un
$R$-module de type fini. Il existe une constante $a$ telle que
$\chi_{I, M}(n) \leq \chi_{I', M}(an)$ pour $n \geq 1$.
\end{lemma}

\begin{proof}
Il existe un entier $c \geq 1$ tel que $(I')^c \subset I$.
Nous obtenons donc une surjection $M/(I')^{c(n + 1)}M \to M/I^{n + 1}M$.
Le résultat s'ensuit avec $a = 2c - 1$.
\end{proof}

\begin{proposition}
\label{proposition-hilbert-function-polynomial}
Soit $R$ un anneau local noethérien. Soit $M$ un $R$-module de type fini.
Soit $I \subset R$ un idéal de définition.
La fonction de Hilbert $\varphi_{I, M}$ et la fonction
$\chi_{I, M}$ sont des polynômes numériques.
\end{proposition}

\begin{proof}
Considérons l'anneau gradué $S = R/I \oplus I/I^2 \oplus I^2/I^3 \oplus
\ldots = \bigoplus_{d \geq 0} I^d/I^{d + 1}$. Considérons le
$S$-module gradué $N = M/IM \oplus IM/I^2M \oplus \ldots =
\bigoplus_{d \geq 0} I^dM/I^{d + 1}M$. La paire $(S, N)$ satisfait
aux hypothèses de la Proposition \ref{proposition-graded-hilbert-polynomial}.
Le résultat pour $\varphi_{I, M}$ découle donc de cette proposition et du
Lemme \ref{lemma-length-K0}. Le résultat pour $\chi_{I, M}$ découle
de ce qui précède et du Lemme \ref{lemma-numerical-polynomial}.
\end{proof}

\begin{definition}
\label{definition-hilbert-polynomial}
Soit $R$ un anneau local noethérien. Soit $M$ un $R$-module de type fini.
Le {\it polynôme de Hilbert} de $M$ sur $R$ est l'élément
$P(t) \in \mathbf{Q}[t]$ tel que $P(n) = \varphi_M(n)$ pour $n \gg 0$.
\end{definition}

\noindent
La Proposition \ref{proposition-hilbert-function-polynomial}
montre que le polynôme de Hilbert existe.

\begin{lemma}
\label{lemma-d-independent}
Soit $R$ un anneau local noethérien. Soit $M$ un $R$-module de type fini.
\begin{enumerate}
\item Le degré du polynôme numérique $\varphi_{I, M}$ est indépendant
de l'idéal de définition $I$.
\item Le degré du polynôme numérique $\chi_{I, M}$ est indépendant
de l'idéal de définition $I$.
\end{enumerate}
\end{lemma}

\begin{proof}
La partie (2) découle immédiatement du Lemme \ref{lemma-hilbert-change-I}.
La partie (1) découle de (2), car
$\varphi_{I, M}(n) = \chi_{I, M}(n) - \chi_{I, M}(n - 1)$
pour $n \geq 1$.
\end{proof}

\begin{definition}
\label{definition-d}
Soit $R$ un anneau local noethérien et $M$ un $R$-module de type fini.
On note {\it $d(M)$} l'élément de $\{-\infty, 0, 1, 2, \ldots \}$
défini comme suit :
\begin{enumerate}
\item si $M = 0$, nous posons $d(M) = -\infty$;
\item si $M \not = 0$, alors $d(M)$ est le degré du polynôme
numérique $\chi_M$.
\end{enumerate}
\end{definition}

\noindent
Si $\mathfrak m^nM \not = 0$ pour tout $n$, alors nous voyons que
$d(M)$ est le degré, augmenté d'une unité ($+1$), du polynôme de Hilbert
de $M$.

\begin{lemma}
\label{lemma-differ-finite-chi}
Soit $R$ un anneau local noethérien. Soit $I \subset R$ un idéal
de définition. Soit $M$ un $R$-module de type fini
qui n'est pas de longueur finie. Si $M' \subset M$ est un sous-module
dont le quotient est de longueur finie, alors $\chi_{I, M} - \chi_{I, M'}$
est un polynôme de degré strictement inférieur ($<$) à celui de
chacun des deux polynômes.
\end{lemma}

\begin{proof}
Cela découle du Lemme \ref{lemma-differ-finite} par un calcul élémentaire.
\end{proof}

\begin{lemma}
\label{lemma-hilbert-ses-chi}
Soit $R$ un anneau local noethérien. Soit $I \subset R$ un idéal de
définition. Soit $0 \to M' \to M \to M'' \to 0$ une suite exacte courte
de $R$-modules de type fini. Alors :
\begin{enumerate}
\item si $M'$ n'est pas de longueur finie, alors
$\chi_{I, M} - \chi_{I, M''} - \chi_{I, M'}$
est un polynôme numérique de degré strictement inférieur ($<$) à celui de
$\chi_{I, M'}$;
\item $\max\{ \deg(\chi_{I, M'}), \deg(\chi_{I, M''}) \} = \deg(\chi_{I, M})$;
et
\item $\max\{d(M'), d(M'')\} = d(M)$.
\end{enumerate}
\end{lemma}

\begin{proof}
Démontrons d'abord (1). Soit $N \subset M'$ comme dans le
Lemme \ref{lemma-hilbert-ses}.
D'après le Lemme \ref{lemma-differ-finite-chi}, le polynôme numérique
$\chi_{I, M'} - \chi_{I, N}$ est de degré strictement inférieur ($<$)
au degré commun de
$\chi_{I, M'}$ et $\chi_{I, N}$. D'après le Lemme \ref{lemma-hilbert-ses},
la différence
$$
\chi_{I, M}(n) - \chi_{I, M''}(n) - \chi_{I, N}(n - c)
$$
est constante pour $n \gg 0$. Par un calcul élémentaire, la différence
$\chi_{I, N}(n) - \chi_{I, N}(n - c)$ est de degré strictement
inférieur ($<$) au degré de
$\chi_{I, N}$, qui est strictement positif (voir ci-dessus). En réunissant
ces observations, nous obtenons (1).

\medskip\noindent
Remarquons que les coefficients dominants de $\chi_{I, M'}$ et
$\chi_{I, M''}$ sont non négatifs. Le degré de
$\chi_{I, M'} + \chi_{I, M''}$ est donc égal au maximum des degrés.
Ainsi, si $M'$ n'est pas de longueur finie, alors (2) découle de (1).
Si $M'$ est de longueur finie, alors
$I^nM \to I^nM''$ est un isomorphisme pour tout $n \gg 0$ d'après
Artin--Rees (Lemme \ref{lemma-Artin-Rees}). Ainsi
$M/I^nM \to M''/I^nM''$ est surjectif et a pour noyau $M'$ pour
$n \gg 0$, et nous voyons que
$\chi_{I, M}(n) - \chi_{I, M''}(n) = \text{longueur}(M')$
pour tout $n \gg 0$. La partie (2) est donc également vraie dans ce cas.

\medskip\noindent
Démonstration de (3). Cela découle de (2), sauf si l'un des modules
$M$, $M'$ ou $M''$ est nul. Nous omettons la démonstration dans ces cas particuliers.
\end{proof}

\section{Dimension}
\label{section-dimension}

\noindent
Comparer avec
Topologie, section \ref{topology-section-krull-dimension}.

\begin{definition}
\label{definition-chain-primes}
Soit $R$ un anneau. Une {\it chaîne d'idéaux premiers} est une suite
$\mathfrak p_0 \subset \mathfrak p_1 \subset \ldots \subset \mathfrak p_n$
d'idéaux premiers de $R$ telle que $\mathfrak p_i \not = \mathfrak p_{i + 1}$
pour $i = 0, \ldots, n - 1$. La {\it longueur} de cette chaîne d'idéaux
premiers est $n$.
\end{definition}

\noindent
Rappelons qu'il existe une bijection renversant les inclusions entre les
idéaux premiers d'un anneau $R$ et les sous-ensembles fermés irréductibles de
$\Spec(R)$; voir le Lemme \ref{lemma-irreducible}.

\begin{definition}
\label{definition-Krull}
La {\it dimension de Krull} de l'anneau $R$ est la dimension de Krull
de l'espace topologique $\Spec(R)$; voir
Topologie, Définition \ref{topology-definition-Krull}.
Autrement dit, c'est la borne supérieure des entiers $n\geq 0$
pour lesquels $R$ possède une chaîne d'idéaux premiers
$$
\mathfrak p_0
\subset
\mathfrak p_1
\subset
\ldots
\subset
\mathfrak p_n, \quad
\mathfrak p_i \not = \mathfrak p_{i + 1}.
$$
Cette chaîne est de longueur $n$.
\end{definition}

\begin{definition}
\label{definition-height}
La {\it hauteur} d'un idéal premier $\mathfrak p$ d'un anneau $R$
est la dimension de l'anneau local $R_{\mathfrak p}$.
\end{definition}

\begin{lemma}
\label{lemma-dimension-height}
La dimension de Krull de $R$ est la borne supérieure des hauteurs de ses
idéaux premiers (maximaux).
\end{lemma}

\begin{proof}
Cela résulte du fait que l'on peut toujours ajouter un idéal maximal à la fin
d'une chaîne d'idéaux premiers.
\end{proof}

\begin{lemma}
\label{lemma-Noetherian-dimension-0}
Un anneau noethérien de dimension $0$ est artinien.
Réciproquement, tout anneau artinien est noethérien de dimension zéro.
\end{lemma}

\begin{proof}
Supposons que $R$ soit un anneau noethérien de dimension $0$.
D'après le Lemme \ref{lemma-Noetherian-topology}, l'espace $\Spec(R)$
est noethérien. D'après Topologie, Lemme \ref{topology-lemma-Noetherian},
$\Spec(R)$ possède un nombre fini de composantes irréductibles, disons
$\Spec(R) = Z_1 \cup \ldots \cup Z_r$.
D'après le Lemme \ref{lemma-irreducible}, chaque $Z_i = V(\mathfrak p_i)$,
où $\mathfrak p_i$ est un idéal premier minimal. Comme la dimension est $0$,
ces $\mathfrak p_i$ sont également maximaux. Ainsi $\Spec(R)$ est l'espace
topologique discret dont les éléments sont les $\mathfrak p_i$.
Tout élément $f$ du radical de Jacobson $\bigcap \mathfrak p_i$ est nilpotent,
car sinon $R_f$ ne serait pas l'anneau nul et nous aurions un autre idéal premier.
D'après le Lemme \ref{lemma-product-local}, $R$ est égal à
$\prod R_{\mathfrak p_i}$. Comme $R_{\mathfrak p_i}$ est lui aussi noethérien
et de dimension $0$, les arguments précédents montrent que son radical
$\mathfrak p_iR_{\mathfrak p_i}$ est localement nilpotent.
Le Lemme \ref{lemma-Noetherian-power} donne
$\mathfrak p_i^nR_{\mathfrak p_i} = 0$ pour un certain $n \geq 1$.
D'après le Lemme \ref{lemma-length-finite}, nous en concluons que
$R_{\mathfrak p_i}$ est de longueur finie sur $R$. Le Lemme
\ref{lemma-artinian-finite-length} montre alors que $R$ est artinien.

\medskip\noindent
Si $R$ est un anneau artinien, alors le Lemme
\ref{lemma-artinian-finite-length} montre qu'il est noethérien. Tous ses
idéaux premiers sont maximaux par combinaison des Lemmes
\ref{lemma-artinian-finite-nr-max},
\ref{lemma-artinian-radical-nilpotent} et \ref{lemma-product-local}.
\end{proof}

\noindent
Dans la suite, nous utiliserons l'invariant $d(-)$ défini dans la
Définition \ref{definition-d}. Commençons par un lemme préparatoire.

\begin{lemma}
\label{lemma-dimension-0-d-0}
Soit $R$ un anneau local noethérien.
Alors $\dim(R) = 0 \Leftrightarrow d(R) = 0$.
\end{lemma}

\begin{proof}
En effet, $d(R) = 0$ si et seulement si $R$ est de longueur finie
comme $R$-module. Voir le Lemme \ref{lemma-artinian-finite-length}.
\end{proof}

\begin{proposition}
\label{proposition-dimension-zero-ring}
Soit $R$ un anneau. Les conditions suivantes sont équivalentes :
\begin{enumerate}
\item $R$ est artinien;
\item $R$ est noethérien et $\dim(R) = 0$;
\item $R$ est de longueur finie comme module sur lui-même;
\item $R$ est un produit fini d'anneaux locaux artiniens;
\item $R$ est noethérien et $\Spec(R)$ est un espace topologique discret fini;
\item $R$ est un produit fini d'anneaux locaux noethériens de dimension $0$;
\item $R$ est un produit fini d'anneaux locaux noethériens $R_i$
tels que $d(R_i) = 0$;
\item $R$ est un produit fini d'anneaux locaux noethériens $R_i$
dont les idéaux maximaux sont nilpotents;
\item $R$ est noethérien, possède un nombre fini d'idéaux maximaux
et son radical de Jacobson est nilpotent; et
\item $R$ est noethérien et il n'existe aucune inclusion stricte entre
ses idéaux premiers.
\end{enumerate}
\end{proposition}

\begin{proof}
C'est une combinaison des Lemmes
\ref{lemma-product-local},
\ref{lemma-artinian-finite-length},
\ref{lemma-Noetherian-dimension-0} et
\ref{lemma-dimension-0-d-0}.
\end{proof}

\begin{lemma}
\label{lemma-height-1}
Soit $R$ un anneau local noethérien.
Les conditions suivantes sont équivalentes :
\begin{enumerate}
\item
\label{item-dim-1}
$\dim(R) = 1$;
\item
\label{item-d-1}
$d(R) = 1$;
\item
\label{item-Vx}
il existe $x \in \mathfrak m$, avec $x$ non nilpotent, tel que
$V(x) = \{\mathfrak m\}$;
\item
\label{item-x}
il existe $x \in \mathfrak m$, avec $x$ non nilpotent, tel que
$\mathfrak m = \sqrt{(x)}$; et
\item
\label{item-ideal-1}
il existe un idéal de définition engendré par $1$ élément,
et aucun idéal de définition n'est engendré par $0$ élément.
\end{enumerate}
\end{lemma}

\begin{proof}
Supposons d'abord que $\dim(R) = 1$.
Soient $\mathfrak p_i$ les idéaux premiers minimaux de $R$.
Comme la dimension est $1$, le seul autre idéal premier de $R$ est
$\mathfrak m$.
D'après le Lemme \ref{lemma-Noetherian-irreducible-components},
ils sont en nombre fini. Nous pouvons donc trouver $x \in \mathfrak m$,
$x \not \in \mathfrak p_i$; voir le Lemme \ref{lemma-silly}.
Le seul idéal premier contenant $x$ est ainsi $\mathfrak m$, d'où
(\ref{item-Vx}).

\medskip\noindent
Si (\ref{item-Vx}) est satisfaite, alors
$\mathfrak m = \sqrt{(x)}$ d'après le Lemme
\ref{lemma-Zariski-topology}, d'où (\ref{item-x}).
La réciproque est également claire.
L'équivalence entre (\ref{item-x}) et (\ref{item-ideal-1}) découle
directement des définitions.

\medskip\noindent
Supposons (\ref{item-ideal-1}).
Soit $I = (x)$ un idéal de définition.
Remarquons que $I^n/I^{n + 1}$ est un quotient de $R/I$ par la multiplication
par $x^n$, et donc que $\text{longueur}_R(I^n/I^{n + 1})$ est bornée.
Ainsi $d(R) = 0$ ou $d(R) = 1$, mais $d(R) = 0$ est exclu par
l'hypothèse selon laquelle $0$ n'est pas un idéal de définition.

\medskip\noindent
Supposons (\ref{item-d-1}). Pour obtenir une contradiction, supposons qu'il
existe des idéaux premiers $\mathfrak p \subset \mathfrak q \subset \mathfrak m$,
les deux inclusions étant strictes. Choisissons un idéal de définition
$I \subset R$. Nous utiliserons plusieurs fois le
Lemme \ref{lemma-hilbert-ses-chi}. Tout d'abord, appliqué à la suite exacte
$0 \to \mathfrak p \to R \to R/\mathfrak p \to 0$,
il implique que $d(R/\mathfrak p) \leq 1$. Mais cette valeur ne peut
évidemment pas être nulle. Choisissons $x\in \mathfrak q$,
$x\not \in \mathfrak p$. Considérons la suite exacte courte
$$
0 \to R/\mathfrak p \to R/\mathfrak p \to R/(xR + \mathfrak p) \to 0.
$$
Cela implique que $\chi_{I, R/\mathfrak p} - \chi_{I, R/\mathfrak p}
- \chi_{I, R/(xR + \mathfrak p)} = - \chi_{I, R/(xR + \mathfrak p)}$
est de degré strictement inférieur ($ < 1$). Autrement dit,
$d(R/(xR + \mathfrak p)) = 0$, et donc
$\dim(R/(xR + \mathfrak p)) = 0$, d'après le
Lemme \ref{lemma-dimension-0-d-0}. Mais $R/(xR + \mathfrak p)$ possède les
idéaux premiers distincts $\mathfrak q/(xR + \mathfrak p)$ et
$\mathfrak m/(xR + \mathfrak p)$, ce qui donne la contradiction cherchée.
\end{proof}

\begin{proposition}
\label{proposition-dimension}
Soit $R$ un anneau local noethérien. Soit $d \geq 0$ un entier.
Les conditions suivantes sont équivalentes :
\begin{enumerate}
\item
\label{item-dim-d}
$\dim(R) = d$;
\item
\label{item-d-d}
$d(R) = d$;
\item
\label{item-ideal-d}
il existe un idéal de définition engendré par $d$ éléments,
et aucun idéal de définition n'est engendré par moins de $d$ éléments.
\end{enumerate}
\end{proposition}

\begin{proof}
Cette démonstration est essentiellement la même que celle du Lemme
\ref{lemma-height-1}. Nous procédons par récurrence sur $d$.
D'après les Lemmes \ref{lemma-dimension-0-d-0} et \ref{lemma-height-1},
nous pouvons supposer que $d > 1$. Notons le nombre minimal de
générateurs d'un idéal de définition de $R$ par $d'(R)$.
Nous allons démontrer les inégalités
$\dim(R) \geq d'(R) \geq d(R) \geq \dim(R)$,
ce qui montrera que ces nombres sont tous égaux.

\medskip\noindent
Supposons d'abord que $\dim(R) = d$.
Soient $\mathfrak p_i$ les idéaux premiers minimaux de $R$.
D'après le Lemme \ref{lemma-Noetherian-irreducible-components},
ils sont en nombre fini. Nous pouvons donc trouver $x \in \mathfrak m$,
$x \not \in \mathfrak p_i$; voir le Lemme \ref{lemma-silly}.
Toute chaîne maximale d'idéaux premiers commence par un certain
$\mathfrak p_i$; la dimension de $R/xR$ est donc au plus $d-1$.
Par récurrence, il existe $x_2, \ldots, x_d$ qui engendrent un idéal de
définition dans $R/xR$. Ainsi $R$ possède un idéal de définition engendré
par (au plus) $d$ éléments.

\medskip\noindent
Supposons $d'(R) = d$. Soit $I = (x_1, \ldots, x_d)$ un idéal de définition.
Remarquons que $I^n/I^{n + 1}$ est un quotient d'une somme directe de
$\binom{d + n - 1}{d - 1}$ copies de $R/I$, par multiplication par tous les
monômes de degré $n$ en $x_1, \ldots, x_d$.
Ainsi $\text{longueur}_R(I^n/I^{n + 1})$ est bornée par un polynôme
de degré $d-1$. Donc $d(R) \leq d$.

\medskip\noindent
Supposons $d(R) = d$. Considérons une chaîne d'idéaux premiers
$\mathfrak p \subset \mathfrak q \subset
\mathfrak q_2 \subset \ldots \subset \mathfrak q_e = \mathfrak m$,
où toutes les inclusions sont strictes et $e \geq 2$.
Choisissons un idéal de définition $I \subset R$.
Nous utiliserons plusieurs fois le Lemme \ref{lemma-hilbert-ses-chi}.
Tout d'abord, appliqué à la suite exacte
$0 \to \mathfrak p \to R \to R/\mathfrak p \to 0$,
il implique que $d(R/\mathfrak p) \leq d$. Mais cette valeur ne peut
évidemment pas être nulle. Choisissons $x\in \mathfrak q$,
$x\not \in \mathfrak p$. Considérons la suite exacte courte
$$
0 \to R/\mathfrak p \to R/\mathfrak p \to R/(xR + \mathfrak p) \to 0.
$$
Cela implique que $\chi_{I, R/\mathfrak p} - \chi_{I, R/\mathfrak p}
- \chi_{I, R/(xR + \mathfrak p)} = - \chi_{I, R/(xR + \mathfrak p)}$
est de degré strictement inférieur ($ < d$). Autrement dit,
$d(R/(xR + \mathfrak p)) \leq d - 1$, et donc
$\dim(R/(xR + \mathfrak p)) \leq d - 1$, par récurrence.
Or $R/(xR + \mathfrak p)$ possède la chaîne d'idéaux premiers
$\mathfrak q/(xR + \mathfrak p) \subset \mathfrak q_2/(xR + \mathfrak p)
\subset \ldots \subset \mathfrak q_e/(xR + \mathfrak p)$, ce qui donne
$e - 1 \leq d - 1$. Comme nous sommes partis d'une chaîne arbitraire
d'idéaux premiers, cela prouve que $\dim(R) \leq d(R)$.

\medskip\noindent
En relisant la démonstration, le lecteur constatera que nous avons obtenu
les inégalités cycliques voulues.
\end{proof}

\noindent
Soit $(R, \mathfrak m)$ un anneau local noethérien.
D'après ce qui précède, il est clair que $\mathfrak m$ ne peut être
engendré par moins de $\dim(R)$ éléments.
D'après le lemme de Nakayama \ref{lemma-NAK}, le nombre minimal de
générateurs de $\mathfrak m$ est égal à $\dim_{\kappa(\mathfrak m)}
\mathfrak m/\mathfrak m^2$. Nous avons donc l'inégalité fondamentale
$$
\dim(R) \leq \dim_{\kappa(\mathfrak m)} \mathfrak m/\mathfrak m^2.
$$
Il se trouve que les anneaux pour lesquels l'égalité est vérifiée
possèdent de nombreuses bonnes propriétés. On les appelle anneaux locaux
réguliers.

\begin{definition}
\label{definition-regular-local}
Soit $(R, \mathfrak m)$ un anneau local noethérien de dimension $d$.
\begin{enumerate}
\item Un {\it système de paramètres de $R$} est une suite d'éléments
$x_1, \ldots, x_d \in \mathfrak m$ qui engendre un idéal de définition
de $R$;
\item s'il existe $x_1, \ldots, x_d \in \mathfrak m$ tels que
$\mathfrak m = (x_1, \ldots, x_d)$, alors nous appelons $R$ un
{\it anneau local régulier}, et $x_1, \ldots, x_d$ un {\it système
régulier de paramètres}.
\end{enumerate}
\end{definition}

\noindent
Les lemmes suivants ressortent clairement des démonstrations des lemmes
et de la proposition ci-dessus, mais nous les énonçons explicitement afin
de disposer de références commodes.

\begin{lemma}
\label{lemma-minimal-over-1}
Soit $R$ un anneau noethérien. Soit $x \in R$.
\begin{enumerate}
\item Si $\mathfrak p$ est minimal au-dessus de $(x)$,
alors la hauteur de $\mathfrak p$ est $0$ ou $1$.
\item Si $\mathfrak p, \mathfrak q \in \Spec(R)$ et si $\mathfrak q$
est minimal au-dessus de $(\mathfrak p, x)$, alors il n'existe aucun idéal
premier strictement compris entre $\mathfrak p$ et $\mathfrak q$.
\end{enumerate}
\end{lemma}

\begin{proof}
Démonstration de (1). Si $\mathfrak p$ est minimal parmi les idéaux premiers
contenant $x$, alors le seul idéal premier de $R_\mathfrak p$ contenant $x$
est l'idéal maximal $\mathfrak p R_\mathfrak p$. En effet, les idéaux premiers
de $R_\mathfrak p$ correspondent de façon $1$ à $1$ aux idéaux premiers de
$R$ contenus dans $\mathfrak p$; voir le Lemme
\ref{lemma-spec-localization}.
Le Lemme \ref{lemma-height-1} montre donc que
$\dim(R_\mathfrak p) = 1$ si $x$ n'est pas nilpotent dans
$R_\mathfrak p$. Bien entendu, si $x$ est nilpotent dans
$R_\mathfrak p$, l'argument montre que $\mathfrak pR_\mathfrak p$ est le
seul idéal premier et que la hauteur est $0$.

\medskip\noindent
Démonstration de (2). D'après la partie (1),
$\mathfrak q/\mathfrak p$ est un idéal premier de hauteur $1$ ou $0$
dans $R/\mathfrak p$. Il en résulte immédiatement qu'il ne peut exister
d'idéal premier strictement compris entre $\mathfrak p$ et $\mathfrak q$.
\end{proof}

\begin{lemma}
\label{lemma-minimal-over-r}
Soit $R$ un anneau noethérien. Soient $f_1, \ldots, f_r \in R$.
\begin{enumerate}
\item Si $\mathfrak p$ est minimal au-dessus de $(f_1, \ldots, f_r)$,
alors la hauteur de $\mathfrak p$ est $\leq r$.
\item Si $\mathfrak p, \mathfrak q \in \Spec(R)$ et si
$\mathfrak q$ est minimal au-dessus de $(\mathfrak p, f_1, \ldots, f_r)$,
alors toute chaîne d'idéaux premiers entre $\mathfrak p$ et
$\mathfrak q$ est de longueur au plus $r$.
\end{enumerate}
\end{lemma}

\begin{proof}
Démonstration de (1). Si $\mathfrak p$ est minimal parmi les idéaux premiers
contenant $f_1, \ldots, f_r$, alors le seul idéal premier de
$R_\mathfrak p$ contenant $f_1, \ldots, f_r$ est l'idéal maximal
$\mathfrak p R_\mathfrak p$. En effet, les idéaux premiers de
$R_\mathfrak p$ correspondent de façon $1$ à $1$ aux idéaux premiers de
$R$ contenus dans $\mathfrak p$; voir le Lemme
\ref{lemma-spec-localization}. La Proposition \ref{proposition-dimension}
montre donc que $\dim(R_\mathfrak p) \leq r$.

\medskip\noindent
Démonstration de (2). D'après la partie (1),
$\mathfrak q/\mathfrak p$ est un idéal premier de hauteur $\leq r$.
Cela implique immédiatement l'assertion sur les chaînes d'idéaux premiers
entre $\mathfrak p$ et $\mathfrak q$.
\end{proof}

\begin{lemma}
\label{lemma-one-equation}
Supposons que $R$ soit un anneau local noethérien et que
$x\in \mathfrak m$ soit un élément de son idéal maximal. Alors
$\dim R \leq \dim R/xR + 1$.
Si $x$ n'appartient à aucun des idéaux premiers minimaux de $R$,
alors il y a égalité. (Par exemple, si $x$ est un non-diviseur de zéro.)
\end{lemma}

\begin{proof}
Si $x_1, \ldots, x_{\dim R/xR} \in R$ s'envoient sur des éléments de
$R/xR$ qui engendrent un idéal de définition de $R/xR$, alors
$x, x_1, \ldots, x_{\dim R/xR}$ engendrent un idéal de définition de $R$.
L'inégalité résulte donc de la Proposition \ref{proposition-dimension}.
D'autre part, si $x$ n'appartient à aucun idéal premier minimal de $R$,
alors toute chaîne d'idéaux premiers de $R/xR$ donne naissance à une
chaîne dans $R$ à laquelle il manque encore au moins une étape pour être
maximale.
\end{proof}

\begin{lemma}
\label{lemma-elements-generate-ideal-definition}
Soit $(R, \mathfrak m)$ un anneau local noethérien.
Supposons que $x_1, \ldots, x_d \in \mathfrak m$ engendrent un idéal
de définition et que $d = \dim(R)$. Alors
$\dim(R/(x_1, \ldots, x_i)) = d - i$ pour tout $i = 1, \ldots, d$.
\end{lemma}

\begin{proof}
Cela découle soit de la démonstration de la Proposition
\ref{proposition-dimension}, soit d'une récurrence sur $d$ utilisant le
Lemme \ref{lemma-one-equation}.
\end{proof}

\section{Applications de la théorie de la dimension}
\label{section-applications-dimension-theory}

\noindent
Nous pouvons utiliser les résultats sur la dimension pour montrer que certains
anneaux ont un spectre infini et pour construire de nouveaux anneaux de
Jacobson.

\begin{lemma}
\label{lemma-Noetherian-local-domain-dim-2-infinite-opens}
Soit $R$ un domaine local noethérien de dimension $\geq 2$.
Toute partie ouverte non vide $U \subset \Spec(R)$ est infinie.
\end{lemma}

\begin{proof}
Pour obtenir une contradiction, supposons que $U \subset \Spec(R)$ soit finie.
Dans ce cas, $(0) \in U$ et $\{(0)\}$ est une partie ouverte de $U$ (car le
complémentaire de $\{(0)\}$ est la réunion des adhérences des autres points).
Nous pouvons donc supposer que $U = \{(0)\}$.
Soit $\mathfrak m \subset R$ l'idéal maximal.
Nous pouvons trouver $x \in \mathfrak m$, $x \not = 0$, tel que
$V(x) \cup U = \Spec(R)$. Autrement dit,
$D(x) = \{(0)\}$. En particulier,
$\dim(R/xR) = \dim(R) - 1 \geq 1$, voir le Lemme \ref{lemma-one-equation}.
Soient $\overline{y}_2, \ldots, \overline{y}_{\dim(R)} \in R/xR$ des éléments
engendrant un idéal de définition de $R/xR$, voir la Proposition
\ref{proposition-dimension}. Choisissons des relèvements
$y_2, \ldots, y_{\dim(R)} \in R$, de sorte que
$x, y_2, \ldots, y_{\dim(R)}$ engendrent un idéal de définition de $R$.
Il s'ensuit que $\dim(R/(y_2)) = \dim(R) - 1$ et que
$\dim(R/(y_2, x)) = \dim(R) - 2$, voir le Lemme
\ref{lemma-elements-generate-ideal-definition}.
Il existe donc un idéal premier $\mathfrak p$ contenant $y_2$ mais pas $x$.
Cela contredit l'égalité $D(x) = \{(0)\}$.
\end{proof}

\noindent
Les anneaux $k[[t]]$, où $k$ est un corps, ainsi que l'anneau des nombres
$p$-adiques sont des anneaux noethériens de dimension $1$ qui n'ont qu'un
nombre fini d'idéaux premiers. C'est la plus grande dimension pour laquelle
cela puisse se produire.

\begin{lemma}
\label{lemma-Noetherian-finite-nr-primes}
Un anneau noethérien qui n'a qu'un nombre fini d'idéaux premiers est de
dimension $\leq 1$.
\end{lemma}

\begin{proof}
Soit $R$ un anneau noethérien qui n'a qu'un nombre fini d'idéaux premiers.
Si $R$ est un domaine local, le lemme découle du Lemme
\ref{lemma-Noetherian-local-domain-dim-2-infinite-opens}.
Si $R$ est un domaine, alors $R_\mathfrak m$ est de dimension $\leq 1$ pour
tout idéal maximal $\mathfrak m$, d'après le cas local.
Ainsi, $\dim(R) \leq 1$ d'après le Lemme \ref{lemma-dimension-height}.
Si $R$ est quelconque, alors $\dim(R/\mathfrak q) \leq 1$ pour tout idéal premier
minimal $\mathfrak q$ de $R$. Puisque tout idéal premier contient un idéal
premier minimal (Lemme \ref{lemma-Zariski-topology}), il s'ensuit que
$\dim(R) \leq 1$.
\end{proof}

\begin{lemma}
\label{lemma-finite-type-algebra-finite-nr-primes}
Soit $S$ une algèbre non nulle de type fini sur un corps $k$.
Les conditions suivantes sont équivalentes :
\begin{enumerate}
\item $\dim(S) = 0$,
\item $S$ n'a qu'un nombre fini d'idéaux premiers,
\item $S$ n'a qu'un nombre fini d'idéaux maximaux,
\item $\Spec(S)$ satisfait l'une des conditions équivalentes du
Lemme \ref{lemma-ring-with-only-minimal-primes},
\item $\dim_k(S) < \infty$,
\item $S$ est artinien,
\item $\Spec(S)$ est un espace topologique discret,
\item ajouter ici d'autres conditions.
\end{enumerate}
\end{lemma}

\begin{proof}
Il résulte immédiatement des définitions que (1) équivaut à (4), en
considérant l'assertion (5) du Lemme
\ref{lemma-ring-with-only-minimal-primes}.
Rappelons que $\Spec(S)$ est sobre, noethérien et de Jacobson, voir les Lemmes
\ref{lemma-spec-spectral}, \ref{lemma-Noetherian-topology},
\ref{lemma-finite-type-field-Jacobson} et \ref{lemma-jacobson}.
Si $S$ est de dimension $0$, chaque point définit une composante irréductible,
et il n'y a qu'un nombre fini de composantes irréductibles (Topologie,
Lemme \ref{topology-lemma-Noetherian}). Ainsi, (1) implique (2).
Il est clair que (2) implique (3). Si (3) est satisfaite, alors $\Spec(S)$ est
discret d'après Topologie, Lemme \ref{topology-lemma-finite-jacobson}, et la
dimension de $S$ est donc $0$.

\medskip\noindent
Nous savons à présent que (1) -- (4) sont équivalentes.
L'implication (5) $\Rightarrow$ (6) est le Lemme
\ref{lemma-finite-dimensional-algebra}.
L'implication (6) $\Rightarrow$ (7) résulte de la Proposition
\ref{proposition-dimension-zero-ring}.
L'implication (7) $\Rightarrow$ (4) est immédiate.
Réciproquement, si $S$ satisfait (1) -- (4), alors $S$ possède un nombre fini
d'idéaux premiers $\mathfrak m_1, \ldots, \mathfrak m_r$, tous maximaux.
Remarquons que $\kappa(\mathfrak m_i)$ est une extension finie de $k$ d'après
le Nullstellensatz de Hilbert (Théorème \ref{theorem-nullstellensatz}).
La Proposition \ref{proposition-dimension-zero-ring} montre également que
$S$ est artinien. Ensuite, le Lemme \ref{lemma-artinian-finite-length} donne
$\text{longueur}_S(S) < \infty$. Ainsi, $\dim_k(S) < \infty$ d'après le
Lemme \ref{lemma-pushdown-module}. Nous concluons que (1) -- (7) sont
équivalentes. (Remarque : une autre manière, plus classique, de démontrer que
$\dim(S) = 0 \Rightarrow \dim_k(S) < \infty$ consiste à utiliser la
normalisation de Noether, mais nous ne disposons pas encore de ce résultat.)
\end{proof}

\begin{lemma}
\label{lemma-noetherian-dim-1-Jacobson}
Anneaux de Jacobson noethériens.
\begin{enumerate}
\item Tout domaine noethérien $R$ de dimension $1$ qui possède une infinité
d'idéaux premiers est un anneau de Jacobson.
\item Tout anneau noethérien tel que chaque idéal premier $\mathfrak p$ soit
maximal ou soit contenu dans une infinité d'idéaux premiers est un anneau de
Jacobson.
\end{enumerate}
\end{lemma}

\begin{proof}
L'assertion (1) est une reformulation du Lemme \ref{lemma-pid-jacobson}.

\medskip\noindent
Soit $R$ un anneau noethérien tel que tout idéal premier non maximal
$\mathfrak p$ soit contenu dans une infinité d'idéaux premiers.
Supposons, pour obtenir une contradiction, que $\Spec(R)$ ne soit pas un
espace de Jacobson. Les Lemmes \ref{lemma-irreducible} et
\ref{lemma-Noetherian-topology} montrent que $\Spec(R)$ est un espace
topologique sobre et noethérien.
D'après Topologie, Lemme
\ref{topology-lemma-non-jacobson-Noetherian-characterize}, il existe un idéal
non maximal $\mathfrak p \subset R$ tel que $\{\mathfrak p\}$ soit une partie
localement fermée de $\Spec(R)$. Autrement dit, $\mathfrak p$ n'est pas
maximal et $\{\mathfrak p\}$ est une partie ouverte de $V(\mathfrak p)$.
Considérons un idéal premier $\mathfrak q \subset R$ tel que
$\mathfrak p \subset \mathfrak q$. Rappelons que la topologie du spectre de
$(R/\mathfrak p)_{\mathfrak q} = R_{\mathfrak q}/\mathfrak pR_{\mathfrak q}$
est induite par celle de $\Spec(R)$, voir les Lemmes
\ref{lemma-spec-localization} et \ref{lemma-spec-closed}.
Ainsi, $\{(0)\}$ est une partie localement fermée de
$\Spec((R/\mathfrak p)_{\mathfrak q})$. Le Lemme
\ref{lemma-Noetherian-local-domain-dim-2-infinite-opens} donne alors
$\dim((R/\mathfrak p)_{\mathfrak q}) = 1$.
Comme cela vaut pour tout $\mathfrak q \supset \mathfrak p$, nous concluons
que $\dim(R/\mathfrak p) = 1$. Utilisons maintenant l'hypothèse selon laquelle
$\mathfrak p$ est contenu dans une infinité d'idéaux premiers : elle montre
que $\Spec(R/\mathfrak p)$ est infini. L'assertion (1) du lemme implique donc
que $V(\mathfrak p) \cong \Spec(R/\mathfrak p)$ est l'adhérence de l'ensemble
de ses points fermés. C'est la contradiction recherchée, puisque cela signifie
que $\{\mathfrak p\} \subset V(\mathfrak p)$ ne peut pas être ouvert.
\end{proof}

\section{Support et dimension des modules}
\label{section-support}

\noindent
Quelques résultats élémentaires sur le support et la dimension des modules.

\begin{lemma}
\label{lemma-filter-Noetherian-module}
Soit $R$ un anneau noethérien et soit $M$ un $R$-module de type fini.
Il existe une filtration par des sous-$R$-modules
$$
0 = M_0 \subset M_1 \subset \ldots \subset M_n = M
$$
telle que chaque quotient $M_i/M_{i-1}$ soit isomorphe à
$R/\mathfrak p_i$ pour un certain idéal premier $\mathfrak p_i$ de $R$.
\end{lemma}

\begin{proof}[Première démonstration]
D'après le Lemme \ref{lemma-trivial-filter-finite-module}, il suffit de traiter
le cas $M = R/I$ pour un certain idéal $I$.
Considérons l'ensemble $S$ des idéaux $J$ tels que le lemme ne soit pas vrai
pour le module $R/J$, et ordonnons-le par inclusion. Pour obtenir une
contradiction, supposons que $S$ ne soit pas vide. Comme $R$ est noethérien,
$S$ possède un élément maximal $J$. Par définition de $S$, l'idéal $J$ ne
peut pas être premier. Choisissons $a, b\in R$ tels que $ab \in J$, mais que
ni $a \in J$ ni $b\in J$. Considérons la filtration
$0 \subset aR/(J \cap aR) \subset R/J$.
Remarquons que le sous-module $aR/(J \cap aR)$ et le module quotient
$(R/J)/(aR/(J \cap aR))$ sont tous deux cycliques ; écrivons-les respectivement
$R/J'$ et $R/J''$, de sorte que nous ayons une suite exacte courte
$0 \to R/J' \to R/J \to R/J'' \to 0$.
L'inclusion $J \subset J'$ est stricte puisque $b \in J'$, et l'inclusion
$J \subset J''$ est stricte puisque $a \in J''$.
Par maximalité de $J$, les modules $R/J'$ et $R/J''$ possèdent donc tous deux
une filtration comme ci-dessus, et il en va de même de $R/J$. Contradiction.
\end{proof}

\begin{proof}[Seconde démonstration]
Pour un $R$-module $M$, disons que $P(M)$ est satisfaite s'il existe une
filtration comme dans l'énoncé du lemme. Remarquons que $P$ est stable par
extensions et qu'elle est satisfaite pour $0$. D'après le Lemme
\ref{lemma-trivial-filter-finite-module}, il suffit de démontrer que $P(R/I)$
est satisfaite pour tout idéal $I$. Dans le cas contraire, comme $R$ est
noethérien, il existe un contre-exemple maximal $J$. D'après l'Exemple
\ref{example-oka-family-property-modules} et la Proposition
\ref{proposition-oka}, l'idéal $J$ est premier, ce qui est une contradiction.
\end{proof}

\begin{lemma}
\label{lemma-filter-primes-in-support}
Soient $R$, $M$, $M_i$ et $\mathfrak p_i$ comme dans le Lemme
\ref{lemma-filter-Noetherian-module}.
Alors $\text{Supp}(M) = \bigcup V(\mathfrak p_i)$ et, en particulier,
$\mathfrak p_i \in \text{Supp}(M)$.
\end{lemma}

\begin{proof}
Cela découle des Lemmes \ref{lemma-support-closed} et
\ref{lemma-support-quotient}.
\end{proof}

\begin{lemma}
\label{lemma-support-point}
Supposons que $R$ soit un anneau local noethérien d'idéal maximal
$\mathfrak m$. Soit $M$ un $R$-module non nul de type fini.
Alors $\text{Supp}(M) = \{ \mathfrak m\}$ si et seulement si $M$ est de
longueur finie sur $R$.
\end{lemma}

\begin{proof}
Supposons que $\text{Supp}(M) = \{ \mathfrak m\}$.
Il suffit de montrer que tous les idéaux premiers $\mathfrak p_i$ de la
filtration du Lemme \ref{lemma-filter-Noetherian-module} sont égaux à l'idéal
maximal. Cela résulte immédiatement du Lemme
\ref{lemma-filter-primes-in-support}.

\medskip\noindent
Supposons que $M$ soit de longueur finie sur $R$.
Alors $\mathfrak m^n M = 0$ d'après le Lemme \ref{lemma-length-infinite}.
Comme un élément de $\mathfrak m$ devient une unité dans $R_{\mathfrak p}$
pour tout idéal premier $\mathfrak p \not = \mathfrak m$ de $R$, nous voyons
que $M_{\mathfrak p} = 0$.
\end{proof}

\begin{lemma}
\label{lemma-Noetherian-power-ideal-kills-module}
Soit $R$ un anneau noethérien.
Soit $I \subset R$ un idéal.
Soit $M$ un $R$-module de type fini.
Alors $I^nM = 0$ pour un certain $n \geq 0$ si et seulement si
$\text{Supp}(M) \subset V(I)$.
\end{lemma}

\begin{proof}
En effet, $I^nM = 0$ équivaut à $I^n \subset \text{Ann}(M)$.
Puisque $R$ est noethérien, cela équivaut à
$I \subset \sqrt{\text{Ann}(M)}$, voir le Lemme
\ref{lemma-Noetherian-power}.
À son tour, cette condition équivaut à
$V(I) \supset V(\text{Ann}(M))$, voir le Lemme
\ref{lemma-Zariski-topology}.
D'après le Lemme \ref{lemma-support-closed}, elle équivaut à
$V(I) \supset \text{Supp}(M)$.
\end{proof}

\begin{lemma}
\label{lemma-filter-minimal-primes-in-support}
Soient $R$, $M$, $M_i$ et $\mathfrak p_i$ comme dans le Lemme
\ref{lemma-filter-Noetherian-module}.
Les éléments minimaux de l'ensemble $\{\mathfrak p_i\}$ sont les éléments
minimaux de $\text{Supp}(M)$. Le nombre de fois qu'un idéal premier minimal
$\mathfrak p$ apparaît est
$$
\#\{i \mid \mathfrak p_i = \mathfrak p\}
=
\text{longueur}_{R_\mathfrak p} M_{\mathfrak p}.
$$
\end{lemma}

\begin{proof}
La première assertion résulte de l'égalité
$\text{Supp}(M) = \bigcup V(\mathfrak p_i)$, voir le Lemme
\ref{lemma-filter-primes-in-support}.
Soit $\mathfrak p \in \text{Supp}(M)$ un élément minimal.
Le support de $M_{\mathfrak p}$ est l'ensemble réduit à l'idéal maximal
$\mathfrak p R_{\mathfrak p}$. D'après le Lemme \ref{lemma-support-point},
la longueur de $M_{\mathfrak p}$ est donc finie et $> 0$.
Remarquons ensuite que $M_{\mathfrak p}$ possède une filtration dont les
sous-quotients sont
$
(R/\mathfrak p_i)_{\mathfrak p}
=
R_{\mathfrak p}/{\mathfrak p_i}R_{\mathfrak p}
$.
Ils sont nuls si $\mathfrak p_i \not \subset \mathfrak p$, et ils sont égaux à
$\kappa(\mathfrak p)$ si $\mathfrak p_i \subset \mathfrak p$, car la
minimalité de $\mathfrak p$ donne alors $\mathfrak p_i = \mathfrak p$.
Le résultat en découle puisque $\kappa(\mathfrak p)$ est de longueur $1$.
\end{proof}

\begin{lemma}
\label{lemma-support-dimension-d}
Soit $R$ un anneau local noethérien.
Soit $M$ un $R$-module de type fini.
Alors $d(M) = \dim(\text{Supp}(M))$, où $d(M)$ est défini dans la
Définition \ref{definition-d}.
\end{lemma}

\begin{proof}
Soient $M_i, \mathfrak p_i$ comme dans le Lemme
\ref{lemma-filter-Noetherian-module}.
Le Lemme \ref{lemma-hilbert-ses-chi} donne l'égalité
$d(M) = \max \{ d(R/\mathfrak p_i) \}$. D'après la Proposition
\ref{proposition-dimension}, nous avons
$d(R/\mathfrak p_i) = \dim(R/\mathfrak p_i)$.
Il est clair que $\dim(R/\mathfrak p_i) = \dim V(\mathfrak p_i)$.
Comme tous les idéaux premiers minimaux de $\text{Supp}(M)$ figurent parmi les
$\mathfrak p_i$ (Lemme \ref{lemma-filter-minimal-primes-in-support}), le
résultat en découle.
\end{proof}

\begin{lemma}
\label{lemma-ses-dimension}
Soit $R$ un anneau noethérien. Soit
$0 \to M' \to M \to M'' \to 0$ une suite exacte courte de $R$-modules de
type fini. Alors
$\max\{\dim(\text{Supp}(M')), \dim(\text{Supp}(M''))\} =
\dim(\text{Supp}(M))$.
\end{lemma}

\begin{proof}
Si $R$ est local, cela découle immédiatement des Lemmes
\ref{lemma-support-dimension-d} et \ref{lemma-hilbert-ses-chi}.
Un argument plus élémentaire, qui vaut aussi lorsque $R$ n'est pas local,
consiste à utiliser le fait que $\text{Supp}(M')$, $\text{Supp}(M'')$ et
$\text{Supp}(M)$ sont fermés (Lemme \ref{lemma-support-closed}) et que
$\text{Supp}(M) = \text{Supp}(M') \cup \text{Supp}(M'')$
(Lemme \ref{lemma-support-quotient}).
\end{proof}

\section{Idéaux premiers associés}
\label{section-ass}

\noindent
Voici la définition usuelle. Pour les anneaux non noethériens et les modules
qui ne sont pas de type fini, il peut être plus approprié d'utiliser la
définition de la section \ref{section-weakly-ass}.

\begin{definition}
\label{definition-associated}
Soit $R$ un anneau. Soit $M$ un $R$-module.
Un idéal premier $\mathfrak p$ de $R$ est {\it associé} à $M$ s'il existe un
élément $m \in M$ dont l'annulateur est $\mathfrak p$.
L'ensemble de tous ces idéaux premiers se note $\text{Ass}_R(M)$ ou
$\text{Ass}(M)$.
\end{definition}

\begin{lemma}
\label{lemma-ass-support}
Soit $R$ un anneau. Soit $M$ un $R$-module.
Alors $\text{Ass}(M) \subset \text{Supp}(M)$.
\end{lemma}

\begin{proof}
Si $m \in M$ a pour annulateur $\mathfrak p$, alors aucun élément de
$R \setminus \mathfrak p$ n'annule $m$. Ainsi, $m$ est un élément non nul de
$M_{\mathfrak p}$, c'est-à-dire que $\mathfrak p \in \text{Supp}(M)$.
\end{proof}

\begin{lemma}
\label{lemma-ass}
Soit $R$ un anneau. Soit $0 \to M' \to M \to M'' \to 0$ une suite exacte
courte de $R$-modules. Alors $\text{Ass}(M') \subset \text{Ass}(M)$ et
$\text{Ass}(M) \subset \text{Ass}(M') \cup \text{Ass}(M'')$.
De plus, $\text{Ass}(M' \oplus M'') = \text{Ass}(M') \cup \text{Ass}(M'')$.
\end{lemma}

\begin{proof}
Si $m' \in M'$, l'annulateur de $m'$ considéré comme élément de $M'$ est le
même que l'annulateur de $m'$ considéré comme élément de $M$. On obtient donc
l'inclusion $\text{Ass}(M') \subset \text{Ass}(M)$. Soit $m \in M$ un élément
dont l'annulateur est un idéal premier $\mathfrak p$. S'il existe
$g \in R$, $g \not \in \mathfrak p$, tel que $m' = gm \in M'$, alors
l'annulateur de $m'$ est $\mathfrak p$. S'il n'existe aucun
$g \in R$, $g \not \in \mathfrak p$, tel que $gm \in M'$, alors l'annulateur
de l'image $m'' \in M''$ de $m$ est $\mathfrak p$. Cela démontre l'inclusion
$\text{Ass}(M) \subset \text{Ass}(M') \cup \text{Ass}(M'')$.
Nous omettons la démonstration de la dernière assertion.
\end{proof}

\begin{lemma}
\label{lemma-ass-filter}
Soit $R$ un anneau et soit $M$ un $R$-module.
Supposons qu'il existe une filtration par des sous-$R$-modules
$$
0 = M_0 \subset M_1 \subset \ldots \subset M_n = M
$$
telle que chaque quotient $M_i/M_{i-1}$ soit isomorphe à $R/\mathfrak p_i$
pour un certain idéal premier $\mathfrak p_i$ de $R$.
Alors $\text{Ass}(M) \subset \{\mathfrak p_1, \ldots, \mathfrak p_n\}$.
\end{lemma}

\begin{proof}
Raisonnons par récurrence sur la longueur $n$ de la filtration $\{ M_i \}$.
Choisissons $m \in M$ dont l'annulateur est un idéal premier $\mathfrak p$.
Si $m \in M_{n-1}$, le résultat découle de l'hypothèse de récurrence. Sinon,
$m$ s'envoie sur un élément non nul de
$M/M_{n-1} \cong R/\mathfrak p_n$. Nous avons donc
$\mathfrak p \subset \mathfrak p_n$.
S'il n'y a pas égalité, nous pouvons trouver $f \in \mathfrak p_n$,
$f \not\in \mathfrak p$. Dans ce cas, l'annulateur de $fm$ est encore
$\mathfrak p$ et $fm \in M_{n-1}$. L'hypothèse de récurrence conclut.
\end{proof}

\begin{lemma}
\label{lemma-finite-ass}
Soit $R$ un anneau noethérien.
Soit $M$ un $R$-module de type fini.
Alors $\text{Ass}(M)$ est fini.
\end{lemma}

\begin{proof}
Cela résulte immédiatement du Lemme \ref{lemma-ass-filter} et du Lemme
\ref{lemma-filter-Noetherian-module}.
\end{proof}

\begin{proposition}
\label{proposition-minimal-primes-associated-primes}
Soit $R$ un anneau noethérien.
Soit $M$ un $R$-module de type fini.
Les ensembles d'idéaux premiers suivants sont égaux :
\begin{enumerate}
\item les idéaux premiers minimaux du support de $M$ ;
\item les idéaux premiers minimaux de $\text{Ass}(M)$ ;
\item pour toute filtration $0 = M_0 \subset M_1 \subset \ldots
\subset M_{n-1} \subset M_n = M$ telle que
$M_i/M_{i-1} \cong R/\mathfrak p_i$, les idéaux premiers minimaux de
l'ensemble $\{\mathfrak p_i\}$.
\end{enumerate}
\end{proposition}

\begin{proof}
Choisissons une filtration comme en (3).
Le Lemme \ref{lemma-filter-minimal-primes-in-support} montre que les ensembles
de (1) et (3) sont égaux.

\medskip\noindent
Soit $\mathfrak p$ un élément minimal de l'ensemble $\{\mathfrak p_i\}$.
Choisissons $i$ minimal tel que $\mathfrak p = \mathfrak p_i$.
Prenons $m \in M_i$, $m \not \in M_{i-1}$. L'annulateur de $m$ est contenu
dans $\mathfrak p_i = \mathfrak p$ et contient
$\mathfrak p_1 \mathfrak p_2 \ldots \mathfrak p_i$. Par notre choix de
$i$ et de $\mathfrak p$, nous avons $\mathfrak p_j \not \subset \mathfrak p$
pour $j < i$, et donc
$\mathfrak p_1 \mathfrak p_2 \ldots \mathfrak p_{i - 1}
\not \subset \mathfrak p_i$. Choisissons
$f \in \mathfrak p_1 \mathfrak p_2 \ldots \mathfrak p_{i - 1}$,
$f \not \in \mathfrak p$. Alors $fm$ a pour annulateur $\mathfrak p$.
Ainsi, $\mathfrak p$ est un idéal premier associé à $M$.
Le Lemme \ref{lemma-ass-support} donne
$\text{Ass}(M) \subset \text{Supp}(M)$ ; par conséquent, $\mathfrak p$ est
minimal dans $\text{Ass}(M)$. L'ensemble des idéaux premiers de (1) est donc
contenu dans l'ensemble des idéaux premiers de (2).

\medskip\noindent
Soit $\mathfrak p$ un élément minimal de $\text{Ass}(M)$.
Puisque $\text{Ass}(M) \subset \text{Supp}(M)$, il existe un élément minimal
$\mathfrak q$ de $\text{Supp}(M)$ tel que
$\mathfrak q \subset \mathfrak p$. Nous venons de montrer que
$\mathfrak q \in \text{Ass}(M)$. Nous avons donc
$\mathfrak q = \mathfrak p$ par minimalité de $\mathfrak p$. Ainsi, l'ensemble des idéaux premiers de (2) est
contenu dans celui de (1).
\end{proof}

\begin{lemma}
\label{lemma-ass-zero}
\begin{slogan}
Sur un anneau noethérien, tout module non nul possède un idéal premier associé.
\end{slogan}
Soit $R$ un anneau noethérien. Soit $M$ un $R$-module.
Alors
$$
M = (0) \Leftrightarrow \text{Ass}(M) = \emptyset.
$$
\end{lemma}

\begin{proof}
Si $M = (0)$, alors $\text{Ass}(M) = \emptyset$ par définition.
Si $M \not = 0$, choisissons un sous-module non nul de type fini
$M' \subset M$, par exemple le sous-module engendré par un seul élément non
nul. Le Lemme \ref{lemma-support-zero} montre que $\text{Supp}(M')$ n'est pas
vide. D'après la Proposition
\ref{proposition-minimal-primes-associated-primes}, il en résulte que
$\text{Ass}(M')$ n'est pas vide. Le Lemme \ref{lemma-ass} implique alors que
$\text{Ass}(M) \not = \emptyset$.
\end{proof}

\begin{lemma}
\label{lemma-ass-minimal-prime-support}
Soit $R$ un anneau noethérien.
Soit $M$ un $R$-module.
Tout $\mathfrak p \in \text{Supp}(M)$ qui est minimal parmi les éléments de
$\text{Supp}(M)$ appartient à $\text{Ass}(M)$.
\end{lemma}

\begin{proof}
Si $M$ est un $R$-module de type fini, cela résulte de la Proposition
\ref{proposition-minimal-primes-associated-primes}.
Dans le cas général, écrivons $M = \bigcup M_\lambda$ comme réunion de ses
sous-modules de type fini et utilisons les égalités
$\text{Supp}(M) = \bigcup \text{Supp}(M_\lambda)$ et
$\text{Ass}(M) = \bigcup \text{Ass}(M_\lambda)$.
\end{proof}

\begin{lemma}
\label{lemma-ass-zero-divisors}
Soit $R$ un anneau noethérien.
Soit $M$ un $R$-module.
La réunion $\bigcup_{\mathfrak q \in \text{Ass}(M)} \mathfrak q$ est
l'ensemble des éléments de $R$ qui sont des diviseurs de zéro dans $M$.
\end{lemma}

\begin{proof}
Tout élément appartenant à un idéal premier associé est manifestement un
diviseur de zéro dans $M$. Réciproquement, supposons que $x \in R$ soit un
diviseur de zéro dans $M$. Considérons le sous-module
$N = \{m \in M \mid xm = 0\}$. Comme $N$ n'est pas nul, il possède un idéal
premier associé $\mathfrak q$ d'après le Lemme \ref{lemma-ass-zero}.
Alors $x \in \mathfrak q$, et $\mathfrak q$ est un idéal premier associé à
$M$ d'après le Lemme \ref{lemma-ass}.
\end{proof}

\begin{lemma}
\label{lemma-one-equation-module}
Soient $R$ un anneau local noethérien, $M$ un $R$-module de type fini et
$f \in \mathfrak m$ un élément de l'idéal maximal de $R$. Alors
$$
\dim(\text{Supp}(M/fM)) \leq
\dim(\text{Supp}(M)) \leq
\dim(\text{Supp}(M/fM)) + 1
$$
Si $f$ n'appartient à aucun des idéaux premiers minimaux du support de $M$
(par exemple, si $f$ est un non-diviseur de zéro dans $M$), alors l'inégalité
de droite est une égalité.
\end{lemma}

\begin{proof}
(L'assertion entre parenthèses résulte du Lemme
\ref{lemma-ass-zero-divisors}.) La première inégalité découle de
$\text{Supp}(M/fM) \subset \text{Supp}(M)$, voir le Lemme
\ref{lemma-support-quotient}. Pour la seconde, remarquons que
$\text{Supp}(M/fM) = \text{Supp}(M) \cap V(f)$, voir le Lemme
\ref{lemma-support-quotient}. Il en résulte, par exemple d'après le Lemme
\ref{lemma-filter-primes-in-support} et les propriétés élémentaires de la
dimension, qu'il suffit de montrer que
$\dim V(\mathfrak p) \leq \dim (V(\mathfrak p) \cap V(f)) + 1$
pour les idéaux premiers $\mathfrak p$ de $R$. C'est une conséquence du
Lemme \ref{lemma-one-equation}.
Enfin, si $f$ n'est contenu dans aucun idéal premier minimal du support de $M$,
alors toute chaîne d'idéaux premiers de $\text{Supp}(M/fM)$ donne naissance à
une chaîne dans $\text{Supp}(M)$ à laquelle il manque encore au moins une étape
pour être maximale.
\end{proof}

\begin{lemma}
\label{lemma-ass-functorial}
Soit $\varphi : R \to S$ un morphisme d'anneaux.
Soit $M$ un $S$-module.
Alors $\Spec(\varphi)(\text{Ass}_S(M)) \subset \text{Ass}_R(M)$.
\end{lemma}

\begin{proof}
Si $\mathfrak q \in \text{Ass}_S(M)$, il existe un élément $m$ de $M$ tel que
l'annulateur de $m$ dans $S$ soit $\mathfrak q$. L'annulateur de $m$ dans $R$
est alors $\mathfrak q \cap R$.
\end{proof}

\begin{remark}
\label{remark-ass-reverse-functorial}
Soit $\varphi : R \to S$ un morphisme d'anneaux.
Soit $M$ un $S$-module.
Il n'est pas toujours vrai que
$\Spec(\varphi)(\text{Ass}_S(M)) \supset \text{Ass}_R(M)$.
Par exemple, considérons le morphisme d'anneaux
$R = k \to S = k[x_1, x_2, x_3, \ldots]/(x_i^2)$ et $M = S$.
Alors $\text{Ass}_R(M)$ n'est pas vide, tandis que $\text{Ass}_S(S)$ est vide.
\end{remark}

\begin{lemma}
\label{lemma-ass-functorial-Noetherian}
Soit $\varphi : R \to S$ un morphisme d'anneaux.
Soit $M$ un $S$-module. Si $S$ est noethérien, alors
$\Spec(\varphi)(\text{Ass}_S(M)) = \text{Ass}_R(M)$.
\end{lemma}

\begin{proof}
Nous avons déjà vu dans le Lemme \ref{lemma-ass-functorial} que
$\Spec(\varphi)(\text{Ass}_S(M)) \subset \text{Ass}_R(M)$.
Pour l'inclusion réciproque, choisissons un idéal premier
$\mathfrak p \in \text{Ass}_R(M)$. Soit $m \in M$ un élément tel que
l'annulateur de $m$ dans $R$ est $\mathfrak p$. Soit
$I = \{g \in S \mid gm = 0\}$ l'annulateur de $m$ dans $S$.
Alors $R/\mathfrak p \subset S/I$ est injectif. En combinant les Lemmes
\ref{lemma-injective-minimal-primes-in-image} et
\ref{lemma-minimal-prime-image-minimal-prime}, nous voyons qu'il existe un
idéal premier $\mathfrak q \subset S$ minimal au-dessus de $I$ qui s'envoie
sur $\mathfrak p$. D'après la Proposition
\ref{proposition-minimal-primes-associated-primes}, $\mathfrak q$ est un idéal
premier associé à $S/I$ ; par le Lemme \ref{lemma-ass}, $\mathfrak q$ est donc
un idéal premier associé à $M$, ce qui conclut.
\end{proof}

\begin{lemma}
\label{lemma-ass-quotient-ring}
Soit $R$ un anneau.
Soit $I$ un idéal.
Soit $M$ un $R/I$-module.
Via l'injection canonique $\Spec(R/I) \to \Spec(R)$, nous avons
$\text{Ass}_{R/I}(M) = \text{Ass}_R(M)$.
\end{lemma}

\begin{proof}
Omis.
\end{proof}

\begin{lemma}
\label{lemma-associated-primes-localize}
Soit $R$ un anneau.
Soit $M$ un $R$-module.
Soit $\mathfrak p \subset R$ un idéal premier.
\begin{enumerate}
\item Si $\mathfrak p \in \text{Ass}(M)$, alors
$\mathfrak pR_{\mathfrak p} \in \text{Ass}(M_{\mathfrak p})$.
\item Si $\mathfrak p$ est de type fini, la réciproque est également vraie.
\end{enumerate}
\end{lemma}

\begin{proof}
Si $\mathfrak p \in \text{Ass}(M)$, il existe un élément $m \in M$ dont
l'annulateur est $\mathfrak p$. Comme la localisation est exacte
(Proposition \ref{proposition-localization-exact}), l'annulateur de $m/1$ dans
$M_{\mathfrak p}$ est $\mathfrak pR_{\mathfrak p}$ ; cela démontre (1).
Supposons $\mathfrak pR_{\mathfrak p} \in \text{Ass}(M_{\mathfrak p})$ et
$\mathfrak p = (f_1, \ldots, f_n)$. Soit $m/g$ un élément de
$M_{\mathfrak p}$ dont l'annulateur est $\mathfrak pR_{\mathfrak p}$.
Cela implique que l'annulateur de $m$ est contenu dans $\mathfrak p$.
Comme $f_i m/g = 0$ dans $M_{\mathfrak p}$, il existe
$g_i \in R$, $g_i \not \in \mathfrak p$, tel que $g_i f_i m = 0$ dans $M$.
En combinant ces faits, nous voyons que l'annulateur de $g_1\ldots g_nm$ est
$\mathfrak p$. Ainsi, $\mathfrak p \in \text{Ass}(M)$.
\end{proof}

\begin{lemma}
\label{lemma-localize-ass}
Soit $R$ un anneau. Soit $M$ un $R$-module.
Soit $S \subset R$ une partie multiplicative.
Via l'injection canonique $\Spec(S^{-1}R) \to \Spec(R)$, nous avons
\begin{enumerate}
\item $\text{Ass}_R(S^{-1}M) = \text{Ass}_{S^{-1}R}(S^{-1}M)$,
\item
$\text{Ass}_R(M) \cap \Spec(S^{-1}R) \subset \text{Ass}_R(S^{-1}M)$, et
\item si $R$ est noethérien, cette inclusion est une égalité.
\end{enumerate}
\end{lemma}

\begin{proof}
Pour $m \in S^{-1}M$, soient $I \subset R$ et $J \subset S^{-1}R$ les
annulateurs de $m$. Alors $I$ est l'image réciproque de $J$ par le morphisme
$R \to S^{-1}R$, et $J = S^{-1}I$. L'égalité de (1) résulte de la description
du morphisme $\Spec(S^{-1}R) \to \Spec(R)$ donnée dans le Lemme
\ref{lemma-spec-localization}.
Pour $m \in M$, soit $I \subset R$ l'annulateur de $m$ dans $R$, et soit
$J \subset S^{-1}R$ l'annulateur de $m/1 \in S^{-1}M$.
Nous avons $J = S^{-1}I$, ce qui implique (2). Dans le cas noethérien,
l'égalité résulte du Lemme \ref{lemma-associated-primes-localize}, car pour
$\mathfrak p \in R$, $S \cap \mathfrak p = \emptyset$, nous avons
$M_{\mathfrak p} = (S^{-1}M)_{S^{-1}\mathfrak p}$.
\end{proof}

\begin{lemma}
\label{lemma-localize-ass-nonzero-divisors}
Soit $R$ un anneau. Soit $M$ un $R$-module.
Soit $S \subset R$ une partie multiplicative.
Supposons que tout $s \in S$ soit un non-diviseur de zéro dans $M$.
Alors
$$
\text{Ass}_R(M) = \text{Ass}_R(S^{-1}M).
$$
\end{lemma}

\begin{proof}
Puisque $M \subset S^{-1}M$ par hypothèse, le Lemme \ref{lemma-ass} donne
l'inclusion $\text{Ass}(M) \subset \text{Ass}(S^{-1}M)$.
Réciproquement, supposons que $n/s \in S^{-1}M$ soit un élément dont
l'annulateur est un idéal premier $\mathfrak p$. Alors l'annulateur de
$n \in M$ est lui aussi $\mathfrak p$.
\end{proof}

\begin{lemma}
\label{lemma-ideal-nonzerodivisor}
Soit $R$ un anneau local noethérien d'idéal maximal $\mathfrak m$.
Soit $I \subset \mathfrak m$ un idéal. Soit $M$ un $R$-module de type fini.
Les conditions suivantes sont équivalentes :
\begin{enumerate}
\item Il existe $x \in I$ qui n'est pas un diviseur de zéro dans $M$.
\item Nous avons $I \not \subset \mathfrak q$ pour tout
$\mathfrak q \in \text{Ass}(M)$.
\end{enumerate}
\end{lemma}

\begin{proof}
S'il existe un non-diviseur de zéro $x$ dans $I$, alors $x$ ne peut
manifestement appartenir à aucun idéal premier associé à $M$.
Réciproquement, supposons $I \not \subset \mathfrak q$ pour tout
$\mathfrak q \in \text{Ass}(M)$. Nous pouvons alors choisir
$x \in I$, $x \not \in \mathfrak q$ pour tout
$\mathfrak q \in \text{Ass}(M)$, d'après les Lemmes
\ref{lemma-finite-ass} et \ref{lemma-silly}.
Le Lemme \ref{lemma-ass-zero-divisors} montre que l'élément $x$ n'est pas un
diviseur de zéro dans $M$.
\end{proof}

\begin{lemma}
\label{lemma-zero-at-ass-zero}
Soit $R$ un anneau. Soit $M$ un $R$-module. Si $R$ est noethérien,
l'application
$$
M
\longrightarrow
\prod\nolimits_{\mathfrak p \in \text{Ass}(M)} M_{\mathfrak p}
$$
est injective.
\end{lemma}

\begin{proof}
Soit $x \in M$ un élément du noyau de l'application.
Alors, si $\mathfrak p$ est un idéal premier associé à $Rx \subset M$, nous
voyons d'une part que $\mathfrak p \in \text{Ass}(M)$
(Lemme \ref{lemma-ass}) et d'autre part que
$(Rx)_{\mathfrak p} \subset M_{\mathfrak p}$ n'est pas nul.
Cette contradiction montre que $\text{Ass}(Rx) = \emptyset$.
Ainsi, $Rx = 0$ d'après le Lemme \ref{lemma-ass-zero}.
\end{proof}

\noindent
Ce lemme devrait probablement être placé ailleurs.

\begin{lemma}
\label{lemma-dim-not-zero-exists-nonzerodivisor-nonunit}
Soit $k$ un corps. Soit $S$ une $k$-algèbre de type fini.
Si $\dim(S) > 0$, alors il existe un élément $f \in S$ qui est un
non-diviseur de zéro et n'est pas inversible.
\end{lemma}

\begin{proof}
D'après le Lemme \ref{lemma-finite-ass}, l'anneau $S$ ne possède qu'un nombre
fini d'idéaux premiers associés. D'après le Lemme
\ref{lemma-finite-type-algebra-finite-nr-primes}, l'anneau $S$ possède une
infinité d'idéaux maximaux. Nous pouvons donc choisir un idéal maximal
$\mathfrak m \subset S$ qui ne soit pas un idéal premier associé à $S$.
Par évitement des idéaux premiers (Lemme \ref{lemma-silly}), nous pouvons
choisir un élément non nul $f \in \mathfrak m$ qui n'appartient à aucun des
idéaux premiers associés à $S$. D'après le Lemme
\ref{lemma-ass-zero-divisors}, l'élément $f$ est un non-diviseur de zéro ; et,
comme $f \in \mathfrak m$, nous voyons que $f$ n'est pas une unité.
\end{proof}

\section{Puissances symboliques}
\label{section-symbolic-power}

\noindent
Voici la définition.

\begin{definition}
\label{definition-symbolic-power}
Soit $R$ un anneau. Soit $\mathfrak p$ un idéal premier. Pour $n \geq 0$, la
$n$-ième {\it puissance symbolique} de $\mathfrak p$ est l'idéal
$\mathfrak p^{(n)} = \Ker(R \to R_\mathfrak p/\mathfrak p^nR_\mathfrak p)$.
\end{definition}

\noindent
Remarquons que $\mathfrak p^n \subset \mathfrak p^{(n)}$, mais qu'il n'y a
pas toujours égalité.

\begin{lemma}
\label{lemma-symbolic-power-associated}
Soit $R$ un anneau noethérien.
Soit $\mathfrak p$ un idéal premier.
Soit $n > 0$. Alors $\text{Ass}(R/\mathfrak p^{(n)}) = \{\mathfrak p\}$.
\end{lemma}

\begin{proof}
Si $\mathfrak q$ est un idéal premier associé à $R/\mathfrak p^{(n)}$, alors
$\mathfrak p \subset \mathfrak q$ manifestement.
D'autre part, tout élément $x \in R$, $x \not \in \mathfrak p$, est un
non-diviseur de zéro dans $R/\mathfrak p^{(n)}$.
En effet, si $y \in R$ et
$xy \in \mathfrak p^{(n)} = R \cap \mathfrak p^nR_{\mathfrak p}$, alors
$y \in \mathfrak p^nR_{\mathfrak p}$, d'où
$y \in \mathfrak p^{(n)}$. Le lemme en découle.
\end{proof}

\begin{lemma}
\label{lemma-symbolic-power-flat-extension}
Soit $R \to S$ un morphisme d'anneaux plat. Soit $\mathfrak p \subset R$ un
idéal premier tel que $\mathfrak q = \mathfrak p S$ soit un idéal premier de
$S$. Alors $\mathfrak p^{(n)} S = \mathfrak q^{(n)}$.
\end{lemma}

\begin{proof}
Puisque
$\mathfrak p^{(n)} = \Ker(R \to R_\mathfrak p/\mathfrak p^nR_\mathfrak p)$,
la platitude montre que $\mathfrak p^{(n)} S$ est le noyau du morphisme
$S \to S_\mathfrak p/\mathfrak p^nS_\mathfrak p$. D'autre part,
$\mathfrak q^{(n)}$ est le noyau du morphisme
$S \to S_\mathfrak q/\mathfrak q^nS_\mathfrak q =
S_\mathfrak q/\mathfrak p^nS_\mathfrak q$. Il suffit donc de montrer que
$$
S_\mathfrak p/\mathfrak p^nS_\mathfrak p
\longrightarrow
S_\mathfrak q/\mathfrak p^nS_\mathfrak q
$$
est injectif. Remarquons que le module de droite est le localisé du module de
gauche par les éléments $f \in S$, $f \not \in \mathfrak q$.
Il suffit donc de montrer que ces éléments sont des non-diviseurs de zéro dans
$S_\mathfrak p/\mathfrak p^nS_\mathfrak p$. Par platitude, le module
$S_\mathfrak p/\mathfrak p^nS_\mathfrak p$ possède une filtration finie dont
les sous-quotients sont
$$
\mathfrak p^iS_\mathfrak p/\mathfrak p^{i + 1}S_\mathfrak p
\cong \mathfrak p^iR_\mathfrak p/\mathfrak p^{i + 1}R_\mathfrak p
\otimes_{R_\mathfrak p} S_\mathfrak p \cong
V \otimes_{\kappa(\mathfrak p)} (S/\mathfrak q)_\mathfrak p
$$
où $V$ est un espace vectoriel sur $\kappa(\mathfrak p)$. Ainsi, $f$ agit de
manière inversible, comme voulu.
\end{proof}



\section{Assassin relatif}
\label{section-relative-assassin}

\noindent
Discussion des assassins relatifs. Soit $R \to S$ un morphisme d'anneaux.
Soit $N$ un $S$-module. Dans cette situation, nous pouvons introduire les
ensembles suivants d'idéaux premiers $\mathfrak q$ de $S$ :
\begin{enumerate}
\item $A$ : en posant $\mathfrak p = R \cap \mathfrak q$, on a
$\mathfrak q \in \text{Ass}_S(N \otimes_R \kappa(\mathfrak p))$,
\item $A'$ : en posant $\mathfrak p = R \cap \mathfrak q$, l'idéal
$\mathfrak q$ appartient à l'image de
$\text{Ass}_{S \otimes \kappa(\mathfrak p)}(N \otimes_R \kappa(\mathfrak p))$
par le morphisme canonique
$\Spec(S \otimes_R \kappa(\mathfrak p)) \to \Spec(S)$,
\item $A_{fin}$ : en posant $\mathfrak p = R \cap \mathfrak q$, on a
$\mathfrak q \in \text{Ass}_S(N/\mathfrak pN)$,
\item $A'_{fin}$ : il existe un idéal premier $\mathfrak p' \subset R$ tel que
$\mathfrak q \in \text{Ass}_S(N/\mathfrak p'N)$,
\item $B$ : il existe un $R$-module $M$ tel que
$\mathfrak q \in \text{Ass}_S(N \otimes_R M)$, et
\item $B_{fin}$ : il existe un $R$-module de type fini $M$ tel que
$\mathfrak q \in \text{Ass}_S(N \otimes_R M)$.
\end{enumerate}
Déterminons quelques-unes des relations entre ces ensembles.

\begin{lemma}
\label{lemma-compare-relative-assassins}
Soit $R \to S$ un morphisme d'anneaux. Soit $N$ un $S$-module.
Soient $A$, $A'$, $A_{fin}$, $B$ et $B_{fin}$ les sous-ensembles de
$\Spec(S)$ introduits ci-dessus.
\begin{enumerate}
\item Nous avons toujours $A = A'$.
\item Nous avons toujours $A_{fin} \subset A$,
$B_{fin} \subset B$, $A_{fin} \subset A'_{fin} \subset B_{fin}$
et $A \subset B$.
\item Si $S$ est noethérien, alors $A = A_{fin}$ et $B = B_{fin}$.
\item Si $N$ est plat sur $R$, alors $A = A_{fin} = A'_{fin}$ et $B = B_{fin}$.
\item Si $R$ est noethérien et si $N$ est plat sur $R$, alors tous ces
ensembles sont égaux, c'est-à-dire
$A = A' = A_{fin} = A'_{fin} = B = B_{fin}$.
\end{enumerate}
\end{lemma}

\begin{proof}
Certains arguments de la démonstration seront répétés dans les démonstrations
de lemmes ultérieurs, qui sont plus précis que celui-ci (car ils portent sur
un module $M$ donné ou sur un idéal premier $\mathfrak p$ donné, et non sur
l'ensemble de tous ces objets).

\medskip\noindent
Démonstration de (1). Soit $\mathfrak p$ un idéal premier de $R$. Nous avons
alors
$$
\text{Ass}_S(N \otimes_R \kappa(\mathfrak p)) =
\text{Ass}_{S/\mathfrak pS}(N \otimes_R \kappa(\mathfrak p)) =
\text{Ass}_{S \otimes_R \kappa(\mathfrak p)}(N \otimes_R \kappa(\mathfrak p))
$$
la première égalité d'après le
Lemme \ref{lemma-ass-quotient-ring}
et la seconde d'après le
Lemme \ref{lemma-localize-ass}, partie (1). Cela montre que $A = A'$.
L'inclusion $A_{fin} \subset A'_{fin}$ est claire.

\medskip\noindent
Démonstration de (2). Chacune des inclusions découle immédiatement des
définitions, à l'exception peut-être de $A_{fin} \subset A$, qui résulte du
Lemme \ref{lemma-localize-ass}
et du fait que nous imposons $\mathfrak p = R \cap \mathfrak q$ dans la
définition de $A_{fin}$.

\medskip\noindent
Démonstration de (3). L'égalité $A = A_{fin}$ découle du
Lemme \ref{lemma-localize-ass}, partie (3), si $S$ est noethérien.
Soit $\mathfrak q = (g_1, \ldots, g_m)$ un idéal premier de type fini de $S$.
Soit $z \in N \otimes_R M$ un élément dont l'annulateur est $\mathfrak q$.
Nous pouvons choisir un sous-module de type fini $M' \subset M$ tel que $z$
soit l'image d'un élément $z' \in N \otimes_R M'$. Alors
$\text{Ann}_S(z') \subset \mathfrak q = \text{Ann}_S(z)$.
Puisque $N \otimes_R -$ commute aux colimites et que $M$ est la colimite
filtrante de ses sous-$R$-modules de type fini, nous pouvons trouver
$M' \subset M'' \subset M$ tel que l'image
$z'' \in N \otimes_R M''$ soit annulée par $g_1, \ldots, g_m$. Ainsi,
$\text{Ann}_S(z'') = \mathfrak q$. Cela montre que $B = B_{fin}$ lorsque $S$
est noethérien.

\medskip\noindent
Démonstration de (4). Si $N$ est plat, le foncteur $N \otimes_R -$ est exact.
En particulier, si $M' \subset M$, alors
$N \otimes_R M' \subset N \otimes_R M$. Ainsi, si
$z \in N \otimes_R M$ est un élément dont l'annulateur
$\mathfrak q = \text{Ann}_S(z)$ est premier, nous pouvons choisir un
sous-$R$-module de type fini quelconque $M' \subset M$ tel que
$z \in N \otimes_R M'$ ; l'annulateur de $z$ comme élément de
$N \otimes_R M'$ est alors égal à $\mathfrak q$. Ainsi, $B = B_{fin}$.
Soit $\mathfrak p'$ un idéal premier de $R$, et soit $\mathfrak q$ un idéal
premier de $S$ associé à $N/\mathfrak p'N$. Cela implique que
$\mathfrak p'S \subset \mathfrak q$. Comme $N$ est plat sur $R$, le module
$N/\mathfrak p'N$ est plat sur l'anneau intègre $R/\mathfrak p'$.
Par conséquent, tout élément non nul de $R/\mathfrak p'$ est un
non-diviseur de zéro dans $N/\mathfrak p'$. L'image d'aucun de ces éléments
ne peut donc appartenir à $\mathfrak q$, et nous concluons que
$\mathfrak p' = R \cap \mathfrak q$. Ainsi, $A_{fin} = A'_{fin}$. Enfin, le
Lemme \ref{lemma-localize-ass-nonzero-divisors}
donne
$\text{Ass}_S(N/\mathfrak p'N) =
\text{Ass}_S(N \otimes_R \kappa(\mathfrak p'))$, c'est-à-dire
$A'_{fin} = A$.

\medskip\noindent
Démonstration de (5). Il suffit de prouver $A'_{fin} = B_{fin}$, puisque les
autres égalités ont été établies en (4). Pour cela, soit $M$ un $R$-module de
type fini. D'après le
Lemme \ref{lemma-filter-Noetherian-module},
il existe une filtration par des sous-$R$-modules
$$
0 = M_0 \subset M_1 \subset \ldots \subset M_n = M
$$
telle que chaque quotient $M_i/M_{i-1}$ soit isomorphe à
$R/\mathfrak p_i$ pour un certain idéal premier $\mathfrak p_i$ de $R$.
Puisque $N$ est plat, nous obtenons une filtration par des sous-$S$-modules
$$
0 = N \otimes_R M_0 \subset N \otimes_R M_1 \subset \ldots \subset
N \otimes_R M_n = N \otimes_R M
$$
dont chaque sous-quotient est isomorphe à $N/\mathfrak p_iN$. D'après le
Lemme \ref{lemma-ass},
nous en déduisons que
$\text{Ass}_S(N \otimes_R M) \subset \bigcup \text{Ass}_S(N/\mathfrak p_iN)$.
Ainsi, $B_{fin} \subset A'_{fin}$. L'inclusion réciproque faisant partie de
(2), l'égalité est démontrée.
\end{proof}

\noindent
Nous définissons l'assassin de $N$ relativement à $S/R$ comme
l'ensemble $A = A'$ ci-dessus. Pour motiver cette définition, remarquons
qu'il ne dépend que des modules fibres $N \otimes_R \kappa(\mathfrak p)$ sur
les anneaux fibres. Comme dans le cas de l'assassin d'un module, nous
avertissons le lecteur que cette notion est surtout pertinente lorsque les
anneaux fibres $S \otimes_R \kappa(\mathfrak p)$ sont noethériens, par exemple
lorsque $R \to S$ est de type fini.

\begin{definition}
\label{definition-relative-assassin}
Soit $R \to S$ un morphisme d'anneaux. Soit $N$ un $S$-module.
L'{\it assassin de $N$ relativement à $S/R$} est l'ensemble
$$
\text{Ass}_{S/R}(N)
=
\{ \mathfrak q \subset S \mid
\mathfrak q \in \text{Ass}_S(N \otimes_R \kappa(\mathfrak p))
\text{ avec }\mathfrak p = R \cap \mathfrak q\}.
$$
C'est l'ensemble noté $A$ dans le
Lemme \ref{lemma-compare-relative-assassins}.
\end{definition}

\noindent
L'esprit des quelques résultats suivants est qu'ils portent sur l'assassin
relatif, même si cela n'apparaît pas immédiatement.

\begin{lemma}
\label{lemma-bourbaki}
Soit $R \to S$ un morphisme d'anneaux.
Soit $M$ un $R$-module et soit $N$ un $S$-module.
Si $N$ est plat comme $R$-module, alors
$$
\text{Ass}_S(M \otimes_R N)
\supset
\bigcup\nolimits_{\mathfrak p \in \text{Ass}_R(M)} \text{Ass}_S(N/\mathfrak pN)
$$
et, si $R$ est noethérien, il y a égalité.
\end{lemma}

\begin{proof}
Si $\mathfrak p \in \text{Ass}_R(M)$, il existe une injection
$R/\mathfrak p \to M$. Comme $N$ est plat sur $R$, nous obtenons une
injection $R/\mathfrak p \otimes_R N \to M \otimes_R N$. Puisque
$R/\mathfrak p \otimes_R N = N/\mathfrak pN$, nous concluons que
$\text{Ass}_S(N/\mathfrak pN) \subset \text{Ass}_S(M \otimes_R N)$, voir le
Lemme \ref{lemma-ass}. Le membre de droite est donc contenu dans le membre de
gauche.

\medskip\noindent
Écrivons $M = \bigcup M_\lambda$ comme la réunion de ses sous-$R$-modules de
type fini. Alors $N \otimes_R M = \bigcup N \otimes_R M_\lambda$ également
(car $N$ est plat sur $R$). La définition des idéaux premiers associés donne
$\text{Ass}_S(N \otimes_R M) = \bigcup \text{Ass}_S(N \otimes_R M_\lambda)$
et $\text{Ass}_R(M) = \bigcup \text{Ass}(M_\lambda)$. Nous pouvons donc
supposer que $M$ est de type fini.

\medskip\noindent
Soit $\mathfrak q \in \text{Ass}_S(M \otimes_R N)$, et supposons que $R$ soit
noethérien et que $M$ soit un $R$-module de type fini. Pour achever la
démonstration, nous devons montrer que $\mathfrak q$ appartient au membre de
droite. Remarquons d'abord que
$\mathfrak qS_{\mathfrak q} \in
\text{Ass}_{S_{\mathfrak q}}((M \otimes_R N)_{\mathfrak q})$,
voir le Lemme \ref{lemma-associated-primes-localize}.
Soit $\mathfrak p$ l'idéal premier correspondant de $R$.
Remarquons que
$$
(M \otimes_R N)_{\mathfrak q} = M \otimes_R N_{\mathfrak q}
= M_{\mathfrak p} \otimes_{R_{\mathfrak p}} N_{\mathfrak q}
$$
Si
$\mathfrak pR_{\mathfrak p} \not \in
\text{Ass}_{R_{\mathfrak p}}(M_{\mathfrak p})$,
il existe un élément $x \in \mathfrak pR_{\mathfrak p}$ qui est un
non-diviseur de zéro dans $M_{\mathfrak p}$ (voir le
Lemme \ref{lemma-ideal-nonzerodivisor}). Comme
$N_{\mathfrak q}$ est plat sur $R_{\mathfrak p}$, l'image de $x$ dans
$\mathfrak qS_{\mathfrak q}$ est un non-diviseur de zéro dans
$(M \otimes_R N)_{\mathfrak q}$. Cela contredit l'hypothèse
$\mathfrak qS_{\mathfrak q} \in \text{Ass}_S((M \otimes_R N)_{\mathfrak q})$.
Nous concluons donc que $\mathfrak p$ est l'un des idéaux premiers associés à
$M$.

\medskip\noindent
Poursuivant l'argument, choisissons une filtration
$$
0 = M_0 \subset M_1 \subset \ldots \subset M_n = M
$$
telle que chaque quotient $M_i/M_{i-1}$ soit isomorphe à
$R/\mathfrak p_i$ pour un certain idéal premier $\mathfrak p_i$ de $R$, voir
le Lemme \ref{lemma-filter-Noetherian-module}.
(D'après le Lemme \ref{lemma-ass-filter}, nous avons
$\mathfrak p_i = \mathfrak p$ pour au moins un $i$.) Nous obtenons une
filtration
$$
0 = M_0 \otimes_R N \subset M_1 \otimes_R N \subset \ldots
\subset M_n \otimes_R N = M \otimes_R N
$$
dont les sous-quotients sont isomorphes à $N/\mathfrak p_iN$. Si
$\mathfrak p_i \not = \mathfrak p$, alors $\mathfrak q$ ne peut être associé
au module $N/\mathfrak p_iN$, d'après le résultat du paragraphe précédent
(car $\text{Ass}_R(R/\mathfrak p_i) = \{\mathfrak p_i\}$). Nous concluons
donc que $\mathfrak q$ est associé à $N/\mathfrak pN$, comme voulu.
\end{proof}

\begin{lemma}
\label{lemma-post-bourbaki}
Soit $R \to S$ un morphisme d'anneaux.
Soit $N$ un $S$-module.
Supposons que $N$ soit plat comme $R$-module et que $R$ soit un anneau
intègre de corps des fractions $K$. Alors
$$
\text{Ass}_S(N) =
\text{Ass}_S(N \otimes_R K) =
\text{Ass}_{S \otimes_R K}(N \otimes_R K)
$$
par l'inclusion canonique
$\Spec(S \otimes_R K) \subset \Spec(S)$.
\end{lemma}

\begin{proof}
Remarquons que $S \otimes_R K = (R \setminus \{0\})^{-1}S$ et
$N \otimes_R K = (R \setminus \{0\})^{-1}N$.
Pour tout élément non nul $x \in R$, la multiplication par $x$ dans $N$ est
injective, puisque $N$ est plat sur $R$. Le lemme résulte donc du
Lemme \ref{lemma-localize-ass-nonzero-divisors}
combiné avec le
Lemme \ref{lemma-localize-ass}, partie (1).
\end{proof}

\begin{lemma}
\label{lemma-bourbaki-fibres}
Soit $R \to S$ un morphisme d'anneaux.
Soit $M$ un $R$-module et soit $N$ un $S$-module.
Supposons que $N$ soit plat comme $R$-module. Alors
$$
\text{Ass}_S(M \otimes_R N)
\supset
\bigcup\nolimits_{\mathfrak p \in \text{Ass}_R(M)}
\text{Ass}_{S \otimes_R \kappa(\mathfrak p)}(N \otimes_R \kappa(\mathfrak p))
$$
où nous utilisons la
Remarque \ref{remark-fundamental-diagram}
pour considérer les spectres des anneaux fibres comme des sous-ensembles de
$\Spec(S)$. Si $R$ est noethérien, cette inclusion est une égalité.
\end{lemma}

\begin{proof}
Ceci équivaut au
Lemme \ref{lemma-bourbaki}
d'après les
Lemmes \ref{lemma-ass-quotient-ring},
\ref{lemma-flat-base-change} et
\ref{lemma-post-bourbaki}.
\end{proof}

\begin{remark}
\label{remark-bourbaki}
Soit $R \to S$ un morphisme d'anneaux. Soit $N$ un $S$-module.
Soit $\mathfrak p$ un idéal premier de $R$. Alors
$$
\text{Ass}_S(N \otimes_R \kappa(\mathfrak p)) =
\text{Ass}_{S/\mathfrak pS}(N \otimes_R \kappa(\mathfrak p)) =
\text{Ass}_{S \otimes_R \kappa(\mathfrak p)}(N \otimes_R \kappa(\mathfrak p)).
$$
La première égalité résulte du
Lemme \ref{lemma-ass-quotient-ring}
et la seconde du
Lemme \ref{lemma-localize-ass}, partie (1).
\end{remark}
















\section{Idéaux premiers faiblement associés}
\label{section-weakly-ass}

\noindent
Il s'agit d'une variante de la notion d'idéal premier associé qui est utile
pour les anneaux non noethériens et les modules qui ne sont pas de type fini.

\begin{definition}
\label{definition-weakly-associated}
Soit $R$ un anneau. Soit $M$ un $R$-module.
Un idéal premier $\mathfrak p$ de $R$ est {\it faiblement associé} à $M$
s'il existe un élément $m \in M$ tel que $\mathfrak p$ soit minimal
parmi les idéaux premiers contenant l'annulateur
$\text{Ann}(m) = \{f \in R \mid fm = 0\}$.
L'ensemble de tous ces idéaux premiers est noté $\text{WeakAss}_R(M)$
ou $\text{WeakAss}(M)$.
\end{definition}

\noindent
Ainsi, tout idéal premier associé est faiblement associé.
Voici une caractérisation à l'aide de la localisation en l'idéal premier.

\begin{lemma}
\label{lemma-weakly-ass-local}
Soit $R$ un anneau. Soit $M$ un $R$-module.
Soit $\mathfrak p$ un idéal premier de $R$.
Les assertions suivantes sont équivalentes :
\begin{enumerate}
\item $\mathfrak p$ est faiblement associé à $M$,
\item $\mathfrak pR_{\mathfrak p}$ est faiblement associé à $M_{\mathfrak p}$,
et
\item $M_{\mathfrak p}$ contient un élément dont
l'annulateur a pour radical $\mathfrak pR_{\mathfrak p}$.
\end{enumerate}
\end{lemma}

\begin{proof}
Supposons (1). Il existe alors un élément $m \in M$ tel que
$\mathfrak p$ soit minimal parmi les idéaux premiers contenant l'annulateur
$I = \{x \in R \mid xm = 0\}$ de $m$. Comme la localisation est exacte,
l'annulateur de $m$ dans $M_{\mathfrak p}$ est $I_{\mathfrak p}$.
Ainsi, $\mathfrak pR_{\mathfrak p}$ est un idéal premier minimal de
$R_{\mathfrak p}$ contenant l'annulateur $I_{\mathfrak p}$
de $m$ dans $M_{\mathfrak p}$. Ceci montre (2), ainsi que (3),
car $\sqrt{I_{\mathfrak p}} = \mathfrak pR_{\mathfrak p}$.

\medskip\noindent
En appliquant l'implication (1) $\Rightarrow$ (3) à $M_{\mathfrak p}$
sur $R_{\mathfrak p}$, on voit que (2) $\Rightarrow$ (3).

\medskip\noindent
Enfin, supposons (3). Il existe donc un élément
$m/f \in M_{\mathfrak p}$ dont l'annulateur a pour radical
$\mathfrak pR_{\mathfrak p}$. L'annulateur
$I = \{x \in R \mid xm = 0\}$ de $m$ dans $M$ vérifie alors
$\sqrt{I_{\mathfrak p}} = \mathfrak pR_{\mathfrak p}$. Cela signifie
clairement que $\mathfrak p$ contient $I$ et est minimal parmi les
idéaux premiers contenant $I$, c'est-à-dire que (1) est vérifiée.
\end{proof}

\begin{lemma}
\label{lemma-reduced-weakly-ass-minimal}
Pour un anneau réduit, les idéaux premiers faiblement associés à l'anneau sont
les idéaux premiers minimaux.
\end{lemma}

\begin{proof}
Soit $(R, \mathfrak m)$ un anneau local réduit.
Supposons que $x \in R$ soit un élément dont l'annulateur
a pour radical $\mathfrak m$. Si $\mathfrak m \not = 0$, alors $x$
ne peut pas être inversible, donc $x \in \mathfrak m$. En particulier, $x^{1 + n} = 0$
pour un certain $n \geq 0$. Ainsi $x = 0$, ce qui contredit l'hypothèse
selon laquelle l'annulateur de $x$ est contenu dans $\mathfrak m$.
On voit donc que $\mathfrak m = 0$, c'est-à-dire que $R$ est un corps.
Le Lemme \ref{lemma-weakly-ass-local} donne alors
l'assertion du lemme.
\end{proof}

\begin{lemma}
\label{lemma-weakly-ass}
Soit $R$ un anneau.
Soit $0 \to M' \to M \to M'' \to 0$ une suite exacte courte
de $R$-modules.
Alors $\text{WeakAss}(M') \subset \text{WeakAss}(M)$ et
$\text{WeakAss}(M) \subset \text{WeakAss}(M') \cup \text{WeakAss}(M'')$.
\end{lemma}

\begin{proof}
Nous utiliserons la caractérisation des idéaux premiers faiblement associés du
Lemme \ref{lemma-weakly-ass-local}.
Soit $\mathfrak p$ un idéal premier de $R$. L'exactitude de la localisation donne
la suite exacte courte
$0 \to M'_{\mathfrak p} \to M_{\mathfrak p} \to M''_{\mathfrak p} \to 0$.
Supposons que $m \in M_{\mathfrak p}$ soit un élément dont l'annulateur
a pour radical $\mathfrak pR_{\mathfrak p}$. Alors, ou bien l'image $\overline{m}$
de $m$ dans $M''_{\mathfrak p}$ est nulle et $m \in M'_{\mathfrak p}$, ou bien le
radical de l'annulateur de $\overline{m}$ est $\mathfrak pR_{\mathfrak p}$.
Ceci montre que
$\text{WeakAss}(M) \subset \text{WeakAss}(M') \cup \text{WeakAss}(M'')$.
L'inclusion $\text{WeakAss}(M') \subset \text{WeakAss}(M)$ résulte immédiatement
des définitions.
\end{proof}

\begin{lemma}
\label{lemma-weakly-ass-zero}
\begin{slogan}
Tout module non nul possède un idéal premier faiblement associé.
\end{slogan}
Soit $R$ un anneau. Soit $M$ un $R$-module. Alors
$$
M = (0) \Leftrightarrow \text{WeakAss}(M) = \emptyset
$$
\end{lemma}

\begin{proof}
Si $M = (0)$, alors $\text{WeakAss}(M) = \emptyset$ par définition.
Réciproquement, supposons $M \not = 0$. Choisissons un élément non nul $m \in M$.
Notons $I = \{x \in R \mid xm = 0\}$ l'annulateur de $m$.
Alors $R/I \subset M$. Ainsi $\text{WeakAss}(R/I) \subset \text{WeakAss}(M)$ par le
Lemme \ref{lemma-weakly-ass}.
Or, puisque $I \not = R$, l'ensemble $V(I) = \Spec(R/I)$ contient un idéal premier
minimal ; voir les
Lemmes \ref{lemma-Zariski-topology} et
\ref{lemma-spec-closed},
ce qui conclut.
\end{proof}

\begin{lemma}
\label{lemma-weakly-ass-support}
Soit $R$ un anneau. Soit $M$ un $R$-module. Alors
$$
\text{Ass}(M) \subset \text{WeakAss}(M) \subset \text{Supp}(M).
$$
\end{lemma}

\begin{proof}
La première inclusion résulte immédiatement des définitions.
Si $\mathfrak p \in \text{WeakAss}(M)$, alors, d'après le
Lemme \ref{lemma-weakly-ass-local},
on a $M_{\mathfrak p} \not = 0$, donc $\mathfrak p \in \text{Supp}(M)$.
\end{proof}

\begin{lemma}
\label{lemma-weakly-ass-zero-divisors}
Soit $R$ un anneau.
Soit $M$ un $R$-module.
La réunion $\bigcup_{\mathfrak q \in \text{WeakAss}(M)} \mathfrak q$
est l'ensemble des éléments de $R$ qui sont des diviseurs de zéro sur $M$.
\end{lemma}

\begin{proof}
Supposons $f \in \mathfrak q \in \text{WeakAss}(M)$.
Il existe un élément $m \in M$ tel que
$\mathfrak q$ soit minimal au-dessus de $I = \{x \in R \mid xm = 0\}$.
Il existe donc $g \in R$, avec $g \not \in \mathfrak q$, et $n > 0$
tels que $f^ngm = 0$. Remarquons que $gm \not = 0$ puisque $g \not \in I$.
Si l'on choisit $n$ minimal avec cette propriété, alors $f (f^{n - 1}gm) = 0$
et $f^{n - 1}gm \not = 0$ ; ainsi $f$ est un diviseur de zéro sur $M$.
Réciproquement, supposons que $f \in R$ soit un diviseur de zéro sur $M$.
Considérons le sous-module $N = \{m \in M \mid fm = 0\}$.
Comme $N$ n'est pas nul, il possède un idéal premier faiblement associé $\mathfrak q$ par le
Lemme \ref{lemma-weakly-ass-zero}.
On a clairement $f \in \mathfrak q$ et, d'après le
Lemme \ref{lemma-weakly-ass},
$\mathfrak q$ est un idéal premier faiblement associé à $M$.
\end{proof}

\begin{lemma}
\label{lemma-weakly-ass-minimal-prime-support}
Soit $R$ un anneau.
Soit $M$ un $R$-module.
Tout $\mathfrak p \in \text{Supp}(M)$ qui est minimal parmi les éléments
de $\text{Supp}(M)$ appartient à $\text{WeakAss}(M)$.
\end{lemma}

\begin{proof}
Remarquons que $\text{Supp}(M_{\mathfrak p}) = \{\mathfrak pR_{\mathfrak p}\}$
dans $\Spec(R_{\mathfrak p})$. En particulier, $M_{\mathfrak p}$
est non nul, et donc $\text{WeakAss}(M_{\mathfrak p}) \not = \emptyset$ d'après le
Lemme \ref{lemma-weakly-ass-zero}.
Puisque $\text{WeakAss}(M_{\mathfrak p}) \subset \text{Supp}(M_{\mathfrak p})$
d'après le
Lemme \ref{lemma-weakly-ass-support},
on en déduit que
$\text{WeakAss}(M_{\mathfrak p}) = \{\mathfrak pR_{\mathfrak p}\}$,
d'où $\mathfrak p \in \text{WeakAss}(M)$ par le
Lemme \ref{lemma-weakly-ass-local}.
\end{proof}

\begin{lemma}
\label{lemma-ass-weakly-ass}
Soit $R$ un anneau. Soit $M$ un $R$-module.
Soit $\mathfrak p$ un idéal premier de $R$ qui est de type fini.
Alors
$$
\mathfrak p \in \text{Ass}(M) \Leftrightarrow
\mathfrak p \in \text{WeakAss}(M).
$$
En particulier, si $R$ est noethérien, alors $\text{Ass}(M) = \text{WeakAss}(M)$.
\end{lemma}

\begin{proof}
Écrivons $\mathfrak p = (g_1, \ldots, g_n)$ avec $g_i \in R$.
Il suffit de démontrer l'implication « $\Leftarrow$ », puisque l'autre
implication est toujours vraie ; voir le
Lemme \ref{lemma-weakly-ass-support}.
Supposons $\mathfrak p \in \text{WeakAss}(M)$.
D'après le
Lemme \ref{lemma-weakly-ass-local},
il existe un élément $m \in M_{\mathfrak p}$ tel que
$I = \{x \in R_{\mathfrak p} \mid xm = 0\}$ ait pour radical
$\mathfrak pR_{\mathfrak p}$. Pour tout $i$, il existe donc
un plus petit $e_i > 0$ tel que $g_i^{e_i}m = 0$ dans $M_{\mathfrak p}$.
Si $e_i > 1$ pour un certain $i$, on peut remplacer $m$ par
$g_i^{e_i - 1} m \not = 0$ et diminuer $\sum e_i$.
On peut donc supposer que l'annulateur de $m \in M_{\mathfrak p}$ est
$(g_1, \ldots, g_n)R_{\mathfrak p} = \mathfrak p R_{\mathfrak p}$. Le
Lemme \ref{lemma-associated-primes-localize}
montre alors que $\mathfrak p \in \text{Ass}(M)$.
\end{proof}

\begin{remark}
\label{remark-weakly-ass-not-functorial}
Soit $\varphi : R \to S$ un homomorphisme d'anneaux. Soit $M$ un $S$-module.
On n'a pas toujours
$\Spec(\varphi)(\text{WeakAss}_S(M)) \subset \text{WeakAss}_R(M)$
contrairement au cas des idéaux premiers associés (voir le
Lemme \ref{lemma-ass-functorial}).
Un exemple est donné par l'homomorphisme d'anneaux
$$
R = k[x_1, x_2, x_3, \ldots] \to
S = k[x_1, x_2, x_3, \ldots, y_1, y_2, y_3, \ldots]/
(x_1y_1, x_2y_2, x_3y_3, \ldots)
$$
et par $M = S$. Dans ce cas, $\mathfrak q = \sum x_iS$ est un idéal premier minimal de
$S$, donc un idéal premier faiblement associé à $M = S$ (voir le
Lemme \ref{lemma-weakly-ass-minimal-prime-support}).
Mais, d'autre part, pour tout élément non nul de $S$, son annulateur dans $R$
est de type fini et n'a donc pas
pour radical $R \cap \mathfrak q = (x_1, x_2, x_3, \ldots)$
(détails omis).
\end{remark}

\begin{lemma}
\label{lemma-weakly-ass-reverse-functorial}
Soit $\varphi : R \to S$ un homomorphisme d'anneaux. Soit $M$ un $S$-module.
Alors on a
$\Spec(\varphi)(\text{WeakAss}_S(M)) \supset \text{WeakAss}_R(M)$.
\end{lemma}

\begin{proof}
Soit $\mathfrak p$ un élément de $\text{WeakAss}_R(M)$.
Il existe alors $m \in M_{\mathfrak p}$ dont l'annulateur
$I = \{x \in R_{\mathfrak p} \mid xm = 0\}$ a pour radical
$\mathfrak pR_{\mathfrak p}$. Considérons l'annulateur
$J = \{x \in S_{\mathfrak p} \mid xm = 0 \}$ de $m$ dans $S_{\mathfrak p}$.
Comme $IS_{\mathfrak p} \subset J$, tout idéal premier minimal
$\mathfrak q \subset S_{\mathfrak p}$ au-dessus de $J$ est au-dessus de $\mathfrak p$.
De plus, un tel $\mathfrak q$ correspond à un idéal premier faiblement associé
à $M$, par exemple d'après le
Lemme \ref{lemma-weakly-ass-local}.
\end{proof}

\begin{remark}
\label{remark-ass-functorial}
Soit $\varphi : R \to S$ un homomorphisme d'anneaux. Soit $M$ un $S$-module.
Notons $f : \Spec(S) \to \Spec(R)$ l'application induite sur les spectres.
Alors on a
$$
f(\text{Ass}_S(M)) \subset
\text{Ass}_R(M) \subset
\text{WeakAss}_R(M) \subset
f(\text{WeakAss}_S(M))
$$
voir les
Lemmes \ref{lemma-ass-functorial},
\ref{lemma-weakly-ass-reverse-functorial} et
\ref{lemma-weakly-ass-support}.
En général, toutes les inclusions peuvent être strictes ; voir les
Remarques \ref{remark-ass-reverse-functorial} et
\ref{remark-weakly-ass-not-functorial}.
Si $S$ est noethérien, alors toutes les inclusions sont des égalités puisque,
d'après le
Lemme \ref{lemma-ass-weakly-ass}, les deux extrêmes sont égaux.
\end{remark}

\begin{lemma}
\label{lemma-weakly-ass-finite-ring-map}
Soit $\varphi : R \to S$ un homomorphisme d'anneaux. Soit $M$ un $S$-module.
Notons $f : \Spec(S) \to \Spec(R)$ l'application induite sur les spectres.
Si $\varphi$ est un homomorphisme d'anneaux fini, alors
$$
\text{WeakAss}_R(M) = f(\text{WeakAss}_S(M)).
$$
\end{lemma}

\begin{proof}
L'une des inclusions a déjà été démontrée ; voir la
Remarque \ref{remark-ass-functorial}.
Pour démontrer l'autre, supposons $\mathfrak q \in \text{WeakAss}_S(M)$
et soit $\mathfrak p$ l'idéal premier correspondant de $R$. Soit $m \in M$
un élément tel que $\mathfrak q$ soit un idéal premier minimal au-dessus de
$J = \{g \in S \mid gm = 0\}$. Le radical de
$JS_{\mathfrak q}$ est donc $\mathfrak qS_{\mathfrak q}$.
Comme $R \to S$ est fini, il n'existe qu'un nombre
fini d'idéaux premiers
$\mathfrak q = \mathfrak q_1, \mathfrak q_2, \ldots, \mathfrak q_l$
au-dessus de $\mathfrak p$ ; voir le
Lemme \ref{lemma-finite-finite-fibres}.
Choisissons $x \in \mathfrak q$ tel que $x \not \in \mathfrak q_i$ pour $i > 1$ ; voir le
Lemme \ref{lemma-silly}.
D'après ce qui précède, il existe un élément $y \in S$, avec $y \not \in \mathfrak q$,
et un entier $t > 0$ tels que $y x^t m = 0$. Ainsi, l'élément
$ym \in M$ est annulé par $x^t$ ; ainsi $ym$ a une image nulle dans
$M_{\mathfrak q_i}$ pour $i = 2, \ldots, l$. En revanche, l'image de $ym$ n'est
pas nulle dans $S_{\mathfrak q}$.

\medskip\noindent
L'anneau $S_{\mathfrak p}$ est semi-local, d'idéaux maximaux
$\mathfrak q_i S_{\mathfrak p}$, par le théorème de montée pour les homomorphismes
d'anneaux finis ; voir le Lemme \ref{lemma-integral-going-up}.
Si $f \in \mathfrak pR_{\mathfrak p}$, une puissance de $f$ appartient
à $JS_{\mathfrak q}$ ; ainsi, pour un certain $t > 0$, l'image de
$f^t ym$ est nulle dans $M_{\mathfrak q}$. Comme $ym$ s'annule aux
autres idéaux maximaux de $S_{\mathfrak p}$, on en déduit que $f^t ym$ est nul
dans $M_{\mathfrak p}$ ; voir le
Lemme \ref{lemma-characterize-zero-local}.
On voit ainsi que $\mathfrak p$ est un idéal premier minimal au-dessus de
l'annulateur de $ym$ dans $R$, ce qui conclut.
\end{proof}

\begin{lemma}
\label{lemma-weakly-ass-quotient-ring}
Soit $R$ un anneau.
Soit $I$ un idéal.
Soit $M$ un $R/I$-module.
Par l'injection canonique
$\Spec(R/I) \to \Spec(R)$
on a $\text{WeakAss}_{R/I}(M) = \text{WeakAss}_R(M)$.
\end{lemma}

\begin{proof}
C'est un cas particulier du Lemme \ref{lemma-weakly-ass-finite-ring-map}.
\end{proof}

\begin{lemma}
\label{lemma-localize-weakly-ass}
Soit $R$ un anneau. Soit $M$ un $R$-module.
Soit $S \subset R$ une partie multiplicative.
Par l'injection canonique $\Spec(S^{-1}R) \to \Spec(R)$,
on a $\text{WeakAss}_R(S^{-1}M) = \text{WeakAss}_{S^{-1}R}(S^{-1}M)$
et
$$
\text{WeakAss}(M) \cap \Spec(S^{-1}R) = \text{WeakAss}(S^{-1}M).
$$
\end{lemma}

\begin{proof}
Supposons $m \in S^{-1}M$. Soient $I = \{x \in R \mid xm = 0\}$
et $I' = \{x' \in S^{-1}R \mid x'm = 0\}$. Alors $I' = S^{-1}I$
et $I \cap S = \emptyset$, sauf si $I = R$ (vérifications omises).
Ainsi, les idéaux premiers de $S^{-1}R$ minimaux au-dessus de $I'$ correspondent bijectivement
aux idéaux premiers de $R$ minimaux au-dessus de $I$ et disjoints de $S$. Ceci démontre
l'égalité $\text{WeakAss}_R(S^{-1}M) = \text{WeakAss}_{S^{-1}R}(S^{-1}M)$.
La seconde égalité résulte du
Lemme \ref{lemma-weakly-ass-local},
car, pour $\mathfrak p \in R$ tel que $S \cap \mathfrak p = \emptyset$, on a
$M_{\mathfrak p} = (S^{-1}M)_{S^{-1}\mathfrak p}$.
\end{proof}

\begin{lemma}
\label{lemma-localize-weakly-ass-nonzero-divisors}
Soit $R$ un anneau. Soit $M$ un $R$-module.
Soit $S \subset R$ une partie multiplicative.
Supposons que tout $s \in S$ soit un non-diviseur de zéro sur $M$.
Alors
$$
\text{WeakAss}(M) = \text{WeakAss}(S^{-1}M).
$$
\end{lemma}

\begin{proof}
Puisque, par hypothèse, $M \subset S^{-1}M$, le
Lemme \ref{lemma-weakly-ass} donne
$\text{WeakAss}(M) \subset \text{WeakAss}(S^{-1}M)$.
Réciproquement, soit $n/s \in S^{-1}M$ un élément d'annulateur
$I$, et soit $\mathfrak p$ un idéal premier minimal au-dessus de $I$.
Alors l'annulateur de $n \in M$ est $I$, et $\mathfrak p$ est un idéal premier
minimal au-dessus de $I$.
\end{proof}

\begin{lemma}
\label{lemma-zero-at-weakly-ass-zero}
Soit $R$ un anneau. Soit $M$ un $R$-module. L'application
$$
M
\longrightarrow
\prod\nolimits_{\mathfrak p \in \text{WeakAss}(M)} M_{\mathfrak p}
$$
est injective.
\end{lemma}

\begin{proof}
Soit $x \in M$ un élément du noyau de l'application. Posons
$N = Rx \subset M$. Si $\mathfrak p$ est un idéal premier faiblement associé à $N$,
on a d'une part $\mathfrak p \in \text{WeakAss}(M)$
(Lemme \ref{lemma-weakly-ass})
et, d'autre part, $N_{\mathfrak p} \subset M_{\mathfrak p}$
n'est pas nul. Cette contradiction montre que $\text{WeakAss}(N) = \emptyset$.
Ainsi $N = 0$, c'est-à-dire $x = 0$, d'après le
Lemme \ref{lemma-weakly-ass-zero}.
\end{proof}

\begin{lemma}
\label{lemma-weak-post-bourbaki}
Soit $R \to S$ un homomorphisme d'anneaux.
Soit $N$ un $S$-module.
Supposons que $N$ soit plat comme $R$-module et que
$R$ soit intègre, de corps des fractions $K$.
Alors
$$
\text{WeakAss}_S(N) = \text{WeakAss}_{S \otimes_R K}(N \otimes_R K)
$$
par l'inclusion canonique
$\Spec(S \otimes_R K) \subset \Spec(S)$.
\end{lemma}

\begin{proof}
Remarquons que $S \otimes_R K = (R \setminus \{0\})^{-1}S$ et
$N \otimes_R K = (R \setminus \{0\})^{-1}N$.
Pour tout $x \in R$ non nul, la multiplication par $x$ dans $N$ est injective, car
$N$ est plat sur $R$. Le lemme résulte donc du
Lemme \ref{lemma-localize-weakly-ass-nonzero-divisors}.
\end{proof}

\begin{lemma}
\label{lemma-weakly-ass-change-fields}
Soit $K/k$ une extension de corps. Soit $R$ une $k$-algèbre.
Soit $M$ un $R$-module. Soit $\mathfrak q \subset R \otimes_k K$
un idéal premier au-dessus de $\mathfrak p \subset R$. Si
$\mathfrak q$ est faiblement associé à $M \otimes_k K$,
alors $\mathfrak p$ est faiblement associé à $M$.
\end{lemma}

\begin{proof}
Soit $z \in M \otimes_k K$ un élément tel que $\mathfrak q$
soit minimal au-dessus de l'annulateur $J \subset R \otimes_k K$ de $z$.
Choisissons une sous-extension de type fini $K/L/k$ telle que
$z \in M \otimes_k L$. Comme $R \otimes_k L \to R \otimes_k K$
est plat, on a $J = I(R \otimes_k K)$, où $I \subset R \otimes_k L$
est l'annulateur de $z$ dans le plus petit anneau
(Lemme \ref{lemma-annihilator-flat-base-change}).
Ainsi, $\mathfrak q \cap (R \otimes_k L)$ est minimal au-dessus de $I$ par le
théorème de descente (Lemme \ref{lemma-flat-going-down}).
On se ramène ainsi au cas décrit dans le paragraphe suivant.

\medskip\noindent
Supposons que $K/k$ soit une extension de corps de type fini.
Soient $x_1, \ldots, x_r \in K$ une base de transcendance
de $K$ sur $k$ ; voir Corps, section \ref{fields-section-transcendence}.
Posons $L = k(x_1, \ldots, x_r)$. Écrivons $[K : L] = n$. Alors
$R \otimes_k L \to R \otimes_k K$ est un homomorphisme d'anneaux fini.
Le Lemme \ref{lemma-weakly-ass-finite-ring-map} montre donc que
$\mathfrak q \cap (R \otimes_k L)$
est un idéal premier faiblement associé à $M \otimes_k K$
considéré comme $R \otimes_k L$-module.
Comme $M \otimes_k K \cong (M \otimes_k L)^{\oplus n}$
en tant que $R \otimes_k L$-module, on voit que
$\mathfrak q \cap (R \otimes_k L)$
est un idéal premier faiblement associé à $M \otimes_k L$
(par exemple en utilisant le Lemme \ref{lemma-weakly-ass} et une récurrence).
On se ramène ainsi au cas examiné dans le paragraphe suivant.

\medskip\noindent
Supposons que $K = k(x_1, \ldots, x_r)$ soit une extension transcendante pure.
On peut remplacer $R$ par $R_\mathfrak p$, $M$ par $M_\mathfrak p$
et $\mathfrak q$ par $\mathfrak q(R_\mathfrak p \otimes_k K)$.
Voir le Lemme \ref{lemma-localize-weakly-ass}.
On se ramène ainsi au cas examiné dans le paragraphe suivant.

\medskip\noindent
Supposons que $K = k(x_1, \ldots, x_r)$ soit une extension transcendante pure
et que $R$ soit local d'idéal maximal $\mathfrak p$. Nous affirmons que tout
$f \in R \otimes_k K$, avec $f \not \in \mathfrak p(R \otimes_k K)$,
est un non-diviseur de zéro sur $M \otimes_k K$. En effet, soit
$z \in M \otimes_k K$ un élément.
Il existe un $R$-sous-module de type fini $M' \subset M$ tel que
$z \in M' \otimes_k K$ et tel que $M'$ soit minimal pour
cette propriété : choisissons une base $\{t_\alpha\}$ de $K$ comme
$k$-espace vectoriel, écrivons $z = \sum m_\alpha \otimes t_\alpha$ et prenons pour
$M'$ le $R$-sous-module engendré par les $m_\alpha$.
Si $z \in \mathfrak p(M' \otimes_k K) = \mathfrak p M' \otimes_k K$,
alors $\mathfrak pM' = M'$ et $M' = 0$ d'après le Lemme \ref{lemma-NAK},
ce qui est une contradiction.
Ainsi, $z$ a une image non nulle $\overline{z}$ dans $M'/\mathfrak p M' \otimes_k K$.
Or $R/\mathfrak p \otimes_k K$ est intègre, car c'est une localisation
de $\kappa(\mathfrak p)[x_1, \ldots, x_n]$, et
$M'/\mathfrak p M' \otimes_k K$ est un module libre ; par conséquent,
$f\overline{z} \not = 0$. Cela démontre l'affirmation.

\medskip\noindent
Enfin, choisissons $z \in M \otimes_k K$ tel que $\mathfrak q$
soit minimal au-dessus de l'annulateur $J \subset R \otimes_k K$ de $z$.
Pour $f \in \mathfrak p$, il existe $n \geq 1$ et
$g \in R \otimes_k K$, avec $g \not \in \mathfrak q$, tels que
$g f^n z \in J$, c'est-à-dire $g f^n z = 0$.
(Cela vient de ce que $\mathfrak q$ est au-dessus de $\mathfrak p$
et que $\mathfrak q$ est minimal au-dessus de $J$.)
On a vu plus haut que $g$ est un non-diviseur de zéro, donc $f^n z = 0$.
Ainsi, $\mathfrak p$ est un idéal premier faiblement associé
à $M \otimes_k K$ considéré comme $R$-module.
Comme $M \otimes_k K$ est une somme directe de copies de $M$,
on en déduit, comme précédemment, que $\mathfrak p$ est un idéal premier
faiblement associé à $M$.
\end{proof}





\section{Idéaux premiers immergés}
\label{section-embedded-primes}

\noindent
Voici la définition.

\begin{definition}
\label{definition-embedded-primes}
Soit $R$ un anneau.
Soit $M$ un $R$-module.
\begin{enumerate}
\item  Les idéaux premiers associés à $M$ qui ne sont
pas minimaux parmi les idéaux premiers associés à $M$ sont appelés les
{\it idéaux premiers associés immergés} de $M$.
\item Les {\it idéaux premiers immergés de $R$}
sont les idéaux premiers associés immergés de $R$ considéré comme $R$-module.
\end{enumerate}
\end{definition}

\noindent
Voici une manière de les éliminer.

\begin{lemma}
\label{lemma-remove-embedded-primes}
Soit $R$ un anneau noethérien.
Soit $M$ un $R$-module de type fini.
Considérons l'ensemble des $R$-sous-modules
$$
\{
K \subset M
\mid
\text{Supp}(K)
\text{ rare dans }
\text{Supp}(M)
\}.
$$
Cet ensemble possède un élément maximal $K$, et le quotient
$M' = M/K$ possède les propriétés suivantes :
\begin{enumerate}
\item $\text{Supp}(M) = \text{Supp}(M')$,
\item $M'$ n'a aucun idéal premier associé immergé,
\item pour tout $f \in R$ appartenant à tous les
idéaux premiers associés immergés de $M$, on a $M_f \cong M'_f$.
\end{enumerate}
\end{lemma}

\begin{proof}
Nous utiliserons le Lemme \ref{lemma-finite-ass} et la
Proposition \ref{proposition-minimal-primes-associated-primes}
sans les mentionner davantage.
Notons $\mathfrak q_1, \ldots, \mathfrak q_t$ les
idéaux premiers minimaux du support de $M$. Notons
$\mathfrak p_1, \ldots, \mathfrak p_s$ les
idéaux premiers associés immergés de $M$. Alors
$\text{Ass}(M) = \{\mathfrak q_j, \mathfrak p_i\}$.
Posons
$$
K = \{m \in M \mid \text{Supp}(Rm) \subset \bigcup V(\mathfrak p_i)\}
$$
On voit immédiatement que c'est un sous-module. Puisque $M$ est de type fini sur un
anneau noethérien, $K$ est lui aussi de type fini. Ainsi $\text{Supp}(K)$
est rare dans $\text{Supp}(M)$. Soit $K' \subset M$ un autre sous-module
dont le support est rare dans $\text{Supp}(M)$. Cela signifie
que $K_{\mathfrak q_j} = 0$. Ainsi, si $m \in K'$,
l'image de $m$ dans $M_{\mathfrak q_j}$ est nulle, ce qui
implique à son tour $(Rm)_{\mathfrak q_j} = 0$.
D'autre part, on a $\text{Ass}(Rm) \subset \text{Ass}(M)$.
Le support de $Rm$ est donc contenu dans $\bigcup V(\mathfrak p_i)$.
Ainsi $m \in K$, puis $K' \subset K$, puisque $m$ était arbitraire dans $K'$.

\medskip\noindent
Soit $M' = M/K$. Puisque $K_{\mathfrak q_j}=0$, on a
$M'_{\mathfrak q_j} = M_{\mathfrak q_j}$ pour tout $j$.
Ainsi $M$ et $M'$ ont le même support.

\medskip\noindent
Supposons $\mathfrak q = \text{Ann}(\overline{m}) \in \text{Ass}(M')$,
où $\overline{m} \in M'$ est l'image de $m \in M$.
Alors $m \not \in K$ ; le support de $Rm$ doit donc contenir l'un des
$\mathfrak q_j$. Puisque $M_{\mathfrak q_j} = M'_{\mathfrak q_j}$,
on sait que l'image de $\overline{m}$ dans $M'_{\mathfrak q_j}$ n'est pas nulle.
Ainsi $\mathfrak q \subset \mathfrak q_j$ (en fait, il y a égalité),
ce qui signifie qu'aucun des idéaux premiers associés à $M'$ n'est immergé.

\medskip\noindent
Soit $f$ un élément appartenant à tous les $\mathfrak p_i$.
Alors $D(f) \cap \text{supp}(K) = 0$. Ainsi $M_f = M'_f$
parce que $K_f = 0$.
\end{proof}

\begin{lemma}
\label{lemma-remove-embedded-primes-localize}
Soit $R$ un anneau noethérien.
Soit $M$ un $R$-module de type fini.
Pour tout $f \in R$, on a $(M')_f = (M_f)'$, où
$M \to M'$ et $M_f \to (M_f)'$ sont les quotients
construits dans le Lemme \ref{lemma-remove-embedded-primes}.
\end{lemma}

\begin{proof}
Omis.
\end{proof}

\begin{lemma}
\label{lemma-no-embedded-primes-endos}
Soit $R$ un anneau noethérien.
Soit $M$ un $R$-module de type fini sans idéal premier associé immergé.
Soit $I = \{x \in R \mid xM = 0\}$. Alors l'anneau $R/I$ n'a aucun
idéal premier immergé.
\end{lemma}

\begin{proof}
On peut remplacer $R$ par $R/I$.
On peut donc supposer que tout élément non nul
de $R$ agit non trivialement sur $M$.
D'après le Lemme \ref{lemma-support-closed}, cela implique que
$\Spec(R)$ est égal au support de $M$.
Supposons que $\mathfrak p$ soit un idéal premier immergé de $R$.
Soit $x \in R$ un élément dont l'annulateur est $\mathfrak p$.
Considérons le module non nul $N = xM \subset M$. Il est annulé
par $\mathfrak p$. Tout idéal premier associé $\mathfrak q$ à $N$
contient donc $\mathfrak p$ et est aussi un idéal premier associé à $M$.
Alors $\mathfrak q$ serait un idéal premier associé immergé de
$M$, ce qui contredit l'hypothèse du lemme.
\end{proof}

















\section{Suites régulières}
\label{section-regular-sequences}

\noindent
Dans cette section, nous établissons quelques propriétés élémentaires des suites régulières.

\begin{definition}
\label{definition-regular-sequence}
Soit $R$ un anneau. Soit $M$ un $R$-module. Une suite d'éléments
$f_1, \ldots, f_r$ de $R$ est appelée une {\it suite $M$-régulière}
si les conditions suivantes sont satisfaites :
\begin{enumerate}
\item $f_i$ est un non-diviseur de zéro sur
$M/(f_1, \ldots, f_{i - 1})M$
pour tout $i = 1, \ldots, r$, et
\item le module $M/(f_1, \ldots, f_r)M$ n'est pas nul.
\end{enumerate}
Si $I$ est un idéal de $R$ et si $f_1, \ldots, f_r \in I$,
nous appelons $f_1, \ldots, f_r$ une {\it suite $M$-régulière
dans $I$}. Si $M = R$, nous appelons simplement $f_1, \ldots, f_r$ une
{\it suite régulière} (dans $I$).
\end{definition}

\noindent
Notons bien que la définition dépend de l'ordre
des éléments $f_1, \ldots, f_r$ (voir les exemples ci-dessous). Certains articles ou livres
ne demandent pas que le module $M/(f_1, \ldots, f_r)M$ soit non nul.
Cela présente l'avantage que la propriété d'être une suite régulière est préservée par
localisation. Cependant, nous utiliserons surtout cette définition pour définir la
profondeur d'un module lorsque $R$ est local ; dans ce cas, les $f_i$ doivent
appartenir à l'idéal maximal, condition qui n'est pas préservée lorsqu'on passe
de $R$ à une localisation $R_\mathfrak p$.

\begin{example}
\label{example-global-regular}
Soit $k$ un corps. Dans l'anneau $k[x, y, z]$,
la suite $x, y(1-x), z(1-x)$ est régulière,
mais la suite $y(1-x), z(1-x), x$ ne l'est pas.
\end{example}

\begin{example}
\label{example-local-regular}
Soit $k$ un corps. Considérons l'anneau
$k[x, y, w_0, w_1, w_2, \ldots]/I$
où $I$ est engendré par les $yw_i$, $i = 0, 1, 2, \ldots$, et les
$w_i - xw_{i + 1}$, $i = 0, 1, 2, \ldots$.
La suite $x, y$ est régulière, mais $y$ est un diviseur de zéro.
On peut de plus localiser en l'idéal maximal
$(x, y, w_i)$ et obtenir encore un exemple.
\end{example}

\begin{lemma}
\label{lemma-permute-xi}
Soit $R$ un anneau local noethérien.
Soit $M$ un $R$-module de type fini.
Soit $x_1, \ldots, x_c$ une suite $M$-régulière.
Alors toute permutation des $x_i$ est elle aussi une suite
régulière.
\end{lemma}

\begin{proof}
Traitons d'abord le cas $c = 2$.
Considérons le noyau $K \subset M$ de $x_2 : M \to M$.
Pour tout $z \in K$, on sait que $z = x_1 z'$
pour un certain $z' \in M$, car
$x_2$ est un non-diviseur de zéro sur $M/x_1M$.
Comme $x_1$ est un non-diviseur de zéro sur $M$, on a aussi $x_2 z' = 0$.
Ainsi $x_1 : K \to K$ est surjectif.
Le Lemme de Nakayama \ref{lemma-NAK} donne donc $K = 0$.
Considérons ensuite la multiplication par $x_1$ sur $M/x_2M$.
Si $z \in M$ a pour image un élément $\overline{z} \in M/x_2M$
du noyau de cette application, alors $x_1 z = x_2 y$ pour un certain $y \in M$.
Mais, puisque $x_1, x_2$ est une suite régulière, on a alors
$y = x_1 y'$ pour un certain $y' \in M$. Ainsi $x_1 ( z - x_2 y' ) =0$,
puis $z = x_2 y'$, et donc $\overline{z} = 0$, comme voulu.

\medskip\noindent
Dans le cas général, remarquons que toute permutation est
une composée de transpositions d'indices adjacents.
Il suffit donc de démontrer que
$$
x_1, \ldots, x_{i-2}, x_i, x_{i-1}, x_{i + 1}, \ldots, x_c
$$
est une suite $M$-régulière. Cela résulte du cas que nous venons
de traiter, appliqué au module $M/(x_1, \ldots, x_{i-2})$
et à la suite régulière de longueur $2$ formée par $x_{i-1}, x_i$.
\end{proof}

\begin{lemma}
\label{lemma-flat-increases-depth}
\begin{slogan}
Les homomorphismes locaux plats d'anneaux préservent et reflètent les suites régulières.
\end{slogan}
Soient $R, S$ des anneaux locaux. Soit $R \to S$ un homomorphisme local plat d'anneaux.
Soit $x_1, \ldots, x_r$ une suite dans $R$. Soit $M$ un $R$-module.
Les assertions suivantes sont équivalentes :
\begin{enumerate}
\item $x_1, \ldots, x_r$ est une suite $M$-régulière dans $R$, et
\item les images de $x_1, \ldots, x_r$ dans $S$ forment une suite
$M \otimes_R S$-régulière.
\end{enumerate}
\end{lemma}

\begin{proof}
Cela vient de ce que $R \to S$ est fidèlement plat
d'après le Lemme \ref{lemma-local-flat-ff}.
\end{proof}

\begin{lemma}
\label{lemma-regular-sequence-in-neighbourhood}
Soit $R$ un anneau noethérien. Soit $M$ un $R$-module de type fini.
Soit $\mathfrak p$ un idéal premier. Soit $x_1, \ldots, x_r$ une suite
dans $R$ dont l'image dans $R_{\mathfrak p}$ forme une suite
$M_{\mathfrak p}$-régulière. Alors il existe $g \in R$, avec $g \not \in \mathfrak p$,
tel que l'image de $x_1, \ldots, x_r$ dans $R_g$ forme
une suite $M_g$-régulière.
\end{lemma}

\begin{proof}
Posons
$$
K_i = \Ker\left(x_i : M/(x_1, \ldots, x_{i - 1})M \to
M/(x_1, \ldots, x_{i - 1})M\right).
$$
C'est un $R$-module de type fini dont la localisation en $\mathfrak p$ est
nulle par hypothèse. Il existe donc $g \in R$, avec $g \not \in \mathfrak p$,
tel que $(K_i)_g = 0$ pour tout $i = 1, \ldots, r$. Ce $g$ convient.
\end{proof}

\begin{lemma}
\label{lemma-join-regular-sequences}
Soit $A$ un anneau. Soit $I$ un idéal engendré par une suite régulière
$f_1, \ldots, f_n$ dans $A$. Soient $g_1, \ldots, g_m \in A$ des
éléments dont les images $\overline{g}_1, \ldots, \overline{g}_m$ forment une
suite régulière dans $A/I$. Alors la suite $f_1, \ldots, f_n, g_1, \ldots, g_m$
est régulière dans $A$.
\end{lemma}

\begin{proof}
Cela résulte immédiatement des définitions.
\end{proof}

\begin{lemma}
\label{lemma-regular-sequence-short-exact-sequence}
Soit $R$ un anneau. Soit $0 \to M_1 \to M_2 \to M_3 \to 0$
une suite exacte courte de $R$-modules. Soit $f_1, \ldots, f_r \in R$.
Si $f_1, \ldots, f_r$ est $M_1$-régulière et $M_3$-régulière, alors
$f_1, \ldots, f_r$ est $M_2$-régulière.
\end{lemma}

\begin{proof}
D'après le Lemme \ref{lemma-snake}, si $f_1 : M_1 \to M_1$ et
$f_1 : M_3 \to M_3$ sont injectifs, alors $f_1 : M_2 \to M_2$ l'est aussi,
et l'on obtient une suite exacte courte
$$
0 \to M_1/f_1M_1 \to M_2/f_1M_2 \to M_3/f_1M_3 \to 0
$$
Le lemme en résulte par récurrence sur $r$. Certains détails sont omis.
\end{proof}

\begin{lemma}
\label{lemma-regular-sequence-powers}
Soit $R$ un anneau. Soit $M$ un $R$-module.
Soient $f_1, \ldots, f_r \in R$ et des entiers $e_1, \ldots, e_r > 0$.
Alors la suite $f_1, \ldots, f_r$ est $M$-régulière
si et seulement si $f_1^{e_1}, \ldots, f_r^{e_r}$
est une suite $M$-régulière.
\end{lemma}

\begin{proof}
Nous le démontrons par récurrence sur $r$. Si $r = 1$, cela résulte des
deux faits élémentaires suivants : (a) une puissance d'un non-diviseur de zéro sur $M$
est un non-diviseur de zéro sur $M$, et (b) un diviseur d'un non-diviseur de zéro sur $M$
est un non-diviseur de zéro sur $M$.
Si $r > 1$, la récurrence appliquée à $M/f_1M$ montre que
$f_1, f_2, \ldots, f_r$ est une suite $M$-régulière si et seulement si
$f_1, f_2^{e_2}, \ldots, f_r^{e_r}$ est une suite $M$-régulière.
Il suffit donc de montrer, étant donné $e > 0$, que $f_1^e, f_2, \ldots, f_r$
est une suite $M$-régulière si et seulement si $f_1, \ldots, f_r$
est une suite $M$-régulière. Nous le démontrons
par récurrence sur $e$. Le cas $e = 1$ est trivial. Puisque $f_1$ est un
non-diviseur de zéro sous les deux hypothèses (d'après le cas $r = 1$),
on a une suite exacte courte
$$
0 \to M/f_1M \xrightarrow{f_1^{e - 1}} M/f_1^eM \to M/f_1^{e - 1}M \to 0
$$
Supposons que $f_1, f_2, \ldots, f_r$ soit une suite $M$-régulière.
Alors, par récurrence, les éléments $f_2, \ldots, f_r$ sont des suites régulières sur $M/f_1M$ et
$M/f_1^{e - 1}M$. D'après le
Lemme \ref{lemma-regular-sequence-short-exact-sequence},
$f_2, \ldots, f_r$ est $M/f_1^eM$-régulière. Ainsi $f_1^e, f_2, \ldots, f_r$
est $M$-régulière. Réciproquement, supposons
que $f_1^e, f_2, \ldots, f_r$ soit une suite $M$-régulière. Alors
$f_2 : M/f_1^eM \to M/f_1^eM$ est injectif, donc
$f_2 : M/f_1M \to M/f_1M$ est injectif, puis, par récurrence (!),
$f_2 : M/f_1^{e - 1}M \to M/f_1^{e - 1}M$ est injectif, et donc
$$
0 \to
M/(f_1, f_2)M \xrightarrow{f_1^{e - 1}}
M/(f_1^e, f_2)M \to
M/(f_1^{e - 1}, f_2)M \to 0
$$
est une suite exacte courte d'après le Lemme \ref{lemma-snake}. Cela démontre la
réciproque pour $r = 2$. Si $r > 2$, alors
$f_3 : M/(f_1^e, f_2)M \to M/(f_1^e, f_2)M$ est injectif, donc
$f_3 : M/(f_1, f_2)M \to M/(f_1, f_2)M$ est injectif, et ainsi de suite.
Certains détails sont omis.
\end{proof}

\begin{lemma}
\label{lemma-regular-sequence-in-polynomial-ring}
Soit $R$ un anneau. Soient $f_1, \ldots, f_r \in R$ qui n'engendrent pas
l'idéal unité. Les assertions suivantes sont équivalentes :
\begin{enumerate}
\item toute permutation de $f_1, \ldots, f_r$ est une suite régulière,
\item toute sous-suite de $f_1, \ldots, f_r$ (dans l'ordre donné) est
une suite régulière, et
\item $f_1x_1, \ldots, f_rx_r$ est une suite régulière dans l'anneau de
polynômes $R[x_1, \ldots, x_r]$.
\end{enumerate}
\end{lemma}

\begin{proof}
Il est clair que (1) implique (2). Montrons par récurrence sur $r$
que (2) implique (1). Le cas $r = 1$ est trivial. Pour $r = 2$, il s'agit de dire que si
$a, b \in R$ est une suite régulière et si $b$ est un non-diviseur de zéro, alors
$b, a$ est une suite régulière. Cela est clair, car le noyau de
$a : R/(b) \to R/(b)$ est isomorphe au noyau de $b : R/(a) \to R/(a)$
si $a$ et $b$ sont tous deux des non-diviseurs de zéro. Supposons $r > 2$.
Supposons (2) vérifiée et cherchons à montrer que $f_{\sigma(1)}, \ldots, f_{\sigma(r)}$
est une suite régulière pour une permutation $\sigma$. Nous savons déjà
par récurrence que $f_{\sigma(1)}, \ldots, f_{\sigma(r - 1)}$ est une suite régulière.
Il suffit donc de montrer que $f_s$, où $s = \sigma(r)$,
est un non-diviseur de zéro modulo $f_1, \ldots, \hat f_s, \ldots, f_r$.
Si $s = r$, c'est terminé. Si $s < r$, remarquons que $f_s$ et $f_r$
sont tous deux des non-diviseurs de zéro dans l'anneau
$R/(f_1, \ldots, \hat f_s, \ldots, f_{r - 1})$
(encore par hypothèse de récurrence). Puisque $f_s, f_r$ est une
suite régulière dans cet anneau, le cas des suites de longueur
$2$ montre que $f_r, f_s$ l'est également.

\medskip\noindent
Remarquons que $R[x_1, \ldots, x_r]/(f_1x_1, \ldots, f_ix_i)$, comme $R$-module,
est une somme directe des modules
$$
R/I_E \cdot x_1^{e_1} \ldots x_r^{e_r}
$$
indexés par les multi-indices $E = (e_1, \ldots, e_r)$, où
$I_E$ est l'idéal engendré par les $f_j$ tels que $1 \leq j \leq i$
et $e_j > 0$. Ainsi $f_{i + 1}x_i$ est un non-diviseur de zéro sur ce module si
et seulement si $f_{i + 1}$ est un non-diviseur de zéro sur $R/I_E$ pour tout $E$.
En prenant $E$ dont toutes les composantes sont strictement positives, on voit que $f_{i + 1}$
est un non-diviseur de zéro sur $R/(f_1, \ldots, f_i)$. Ainsi (3) implique (2).
Réciproquement, si (2) est vérifiée, toute sous-suite de
$f_1, \ldots, f_i, f_{i + 1}$ est une suite régulière ;
en particulier, $f_{i + 1}$ est un non-diviseur de zéro sur tous les $R/I_E$.
On voit ainsi que (2) implique (3).
\end{proof}










\section{Suites quasi-régulières}
\label{section-quasi-regular}

\noindent
Nous introduisons la notion de suite quasi-régulière, légèrement plus faible
que celle de suite régulière et plus facile à utiliser.
Soit $R$ un anneau et soient $f_1, \ldots, f_c \in R$.
Posons $J = (f_1, \ldots, f_c)$. Soit $M$ un $R$-module.
Il existe alors une application canonique
\begin{equation}
\label{equation-quasi-regular}
M/JM \otimes_{R/J} R/J[X_1, \ldots, X_c]
\longrightarrow
\bigoplus\nolimits_{n \geq 0} J^nM/J^{n + 1}M
\end{equation}
de $R/J[X_1, \ldots, X_c]$-modules gradués définie par la règle
$$
\overline{m} \otimes X_1^{e_1} \ldots X_c^{e_c} \longmapsto
f_1^{e_1} \ldots f_c^{e_c} m \bmod J^{e_1 + \ldots + e_c + 1}M.
$$
Remarquons que (\ref{equation-quasi-regular}) est toujours surjective.

\begin{definition}
\label{definition-quasi-regular-sequence}
Soit $R$ un anneau.
Soit $M$ un $R$-module.
Une suite d'éléments $f_1, \ldots, f_c$ de $R$ est dite
{\it $M$-quasi-régulière} si (\ref{equation-quasi-regular})
est un isomorphisme. Si $M = R$, nous appelons simplement $f_1, \ldots, f_c$ une
{\it suite quasi-régulière}.
\end{definition}

\noindent
Ainsi, si $f_1, \ldots, f_c$ est une suite quasi-régulière, alors
$$
R/J[X_1, \ldots, X_c] = \bigoplus\nolimits_{n \geq 0} J^n/J^{n + 1}
$$
où $J = (f_1, \ldots, f_c)$. Il est clair que la propriété d'être une suite
quasi-régulière ne dépend pas de l'ordre de $f_1, \ldots, f_c$.

\begin{lemma}
\label{lemma-regular-quasi-regular}
Soit $R$ un anneau.
\begin{enumerate}
\item Une suite régulière $f_1, \ldots, f_c$ de $R$ est une suite
quasi-régulière.
\item Supposons que $M$ soit un $R$-module et que $f_1, \ldots, f_c$
soit une suite $M$-régulière. Alors $f_1, \ldots, f_c$ est une suite
$M$-quasi-régulière.
\end{enumerate}
\end{lemma}

\begin{proof}
Posons $J = (f_1, \ldots, f_c)$.
Démontrons la première assertion par récurrence sur $c$.
Il faut montrer que, pour toute relation
$\sum_{|I| = n} a_I f^I \in J^{n + 1}$ avec $a_I \in R$, on a en fait
$a_I \in J$ pour tout multi-indice $I$. Puisque
tout élément de $J^{n + 1}$ est de la forme $\sum_{|I| = n} b_I f^I$
avec $b_I \in J$, on peut supposer, après avoir remplacé $a_I$ par $a_I - b_I$,
que la relation s'écrit $\sum_{|I| = n} a_I f^I = 0$. On peut la réécrire
sous la forme
$$
\sum\nolimits_{e = 0}^n
\left(
\sum\nolimits_{|I'| = n - e}
a_{I', e} f^{I'}
\right)
f_c^e
=
0
$$
Ici et ci-dessous, les multi-indices « primés » $I'$ sont supposés être de la forme
$I' = (i_1, \ldots, i_{c - 1}, 0)$. Nous allons montrer par
récurrence sur $l \in \{0, \ldots, n\}$
que, si l'on a une relation
$$
\sum\nolimits_{e = 0}^l
\left(
\sum\nolimits_{|I'| = n - e}
a_{I', e} f^{I'}
\right)
f_c^e
=
0
$$
alors $a_{I', e} \in J$ pour tous $I', e$.
En effet, posons $J' = (f_1, \ldots, f_{c-1})$.
Remarquons que $\sum\nolimits_{|I'| = n - l} a_{I', l} f^{I'}$
est envoyé dans $(J')^{n - l + 1}$ par $f_c^{l}$.
L'hypothèse de récurrence (pour la récurrence sur $c$)
montre que $f_c^l a_{I', l} \in J'$.
Comme $f_c$ n'est pas un diviseur de zéro sur $R/J'$ (car $f_1, \ldots, f_c$
est une suite régulière), on en déduit que $a_{I', l} \in J'$.
Cela permet de réécrire le terme
$(\sum\nolimits_{|I'| = n - l} a_{I', l} f^{I'})f_c^l$
 sous la forme $(\sum\nolimits_{|I'| = n - l + 1} f_c b_{I', l - 1}
f^{I'})f_c^{l-1}$. On obtient ainsi une nouvelle relation de la forme
$$
\left(\sum\nolimits_{|I'| = n - l + 1}
(a_{I', l-1} + f_c b_{I', l - 1}) f^{I'}\right)f_c^{l-1}
+
\sum\nolimits_{e = 0}^{l - 2}
\left(
\sum\nolimits_{|I'| = n - e}
a_{I', e} f^{I'}
\right)
f_c^e
=
0
$$
L'hypothèse de récurrence (sur $l$ cette fois) montre maintenant que
tous les $a_{I', l-1} + f_c b_{I', l - 1} \in J$ et que
tous les $a_{I', e} \in J$ pour $e \leq l - 2$. Avec la relation
$a_{I', l} \in J' \subset J$ obtenue ci-dessus, cela achève la démonstration du
pas de récurrence.

\medskip\noindent
La seconde assertion signifie que, pour toute expression formelle
$F = \sum_{|I| = n} m_I X^I$, avec $m_I \in M$ et $\sum m_I f^I
\in J^{n + 1}M$, tous les coefficients $m_I$ appartiennent à $J$.
Cela se démontre exactement comme le résultat correspondant
de la première assertion ci-dessus.
\end{proof}

\begin{lemma}
\label{lemma-flat-base-change-quasi-regular}
Soit $R \to R'$ un homomorphisme plat d'anneaux. Soit $M$ un $R$-module.
Supposons que la suite $f_1, \ldots, f_r \in R$ soit $M$-quasi-régulière.
Alors les images de $f_1, \ldots, f_r$ dans
$R'$ forment une suite $M \otimes_R R'$-quasi-régulière.
\end{lemma}

\begin{proof}
Posons $J = (f_1, \ldots, f_r)$, $J' = JR'$ et $M' = M \otimes_R R'$.
Il faut montrer que l'application canonique
$\mu : R'/J'[X_1, \ldots X_r] \otimes_{R'/J'} M'/J'M' \to
\bigoplus (J')^nM'/(J')^{n + 1}M'$ est un isomorphisme.
Puisque $R \to R'$ est plat, les suites
$0 \to J^nM \to M$ et
$0 \to J^{n + 1}M \to J^nM \to J^nM/J^{n + 1}M \to 0$
restent exactes après tensorisation par $R'$. Cela implique d'abord
$J^nM \otimes_R R' = (J')^nM'$, puis
$(J')^nM'/(J')^{n + 1}M' = J^nM/J^{n + 1}M \otimes_R R'$.
Ainsi $\mu$ est le produit tensoriel de (\ref{equation-quasi-regular}),
qui est un isomorphisme par hypothèse,
avec $\text{id}_{R'}$, ce qui conclut.
\end{proof}

\begin{lemma}
\label{lemma-quasi-regular-sequence-in-neighbourhood}
Soit $R$ un anneau noethérien. Soit $M$ un $R$-module de type fini.
Soit $\mathfrak p$ un idéal premier. Soit $x_1, \ldots, x_c$ une suite
dans $R$ dont l'image dans $R_{\mathfrak p}$ forme une suite
$M_{\mathfrak p}$-quasi-régulière. Alors il existe
$g \in R$, avec $g \not \in \mathfrak p$,
tel que l'image de $x_1, \ldots, x_c$ dans $R_g$ forme
une suite $M_g$-quasi-régulière.
\end{lemma}

\begin{proof}
Considérons le noyau $K$ de l'application (\ref{equation-quasi-regular}).
Comme $M/JM \otimes_{R/J} R/J[X_1, \ldots, X_c]$ est un
$R/J[X_1, \ldots, X_c]$-module de type fini et que $R/J[X_1, \ldots, X_c]$ est
noethérien, $K$ est aussi un $R/J[X_1, \ldots, X_c]$-module de type fini.
Choisissons des générateurs homogènes $k_1, \ldots, k_t \in K$.
Par hypothèse, pour tout $i = 1, \ldots, t$, il existe $g_i \in R$,
avec $g_i \not \in \mathfrak p$, tel que $g_i k_i = 0$.
Alors $g = g_1 \ldots g_t$ convient.
\end{proof}

\begin{lemma}
\label{lemma-truncate-quasi-regular}
Soit $R$ un anneau. Soit $M$ un $R$-module.
Soit $f_1, \ldots, f_c \in R$ une suite $M$-quasi-régulière.
Pour tout $i$, la suite
$\overline{f}_{i + 1}, \ldots, \overline{f}_c$
de $\overline{R} = R/(f_1, \ldots, f_i)$ est une suite
$\overline{M} = M/(f_1, \ldots, f_i)M$-quasi-régulière.
\end{lemma}

\begin{proof}
Il suffit de le démontrer pour $i = 1$. Posons
$\overline{J} = (\overline{f}_2, \ldots, \overline{f}_c) \subset \overline{R}$.
Alors
\begin{align*}
\overline{J}^n\overline{M}/\overline{J}^{n + 1}\overline{M}
& =
(J^nM + f_1M)/(J^{n + 1}M + f_1M) \\
& = J^nM / (J^{n + 1}M + J^nM \cap f_1M).
\end{align*}
Ainsi, pour démontrer le lemme, il suffit de montrer que
$J^{n + 1}M + J^nM \cap f_1M = J^{n + 1}M + f_1J^{n - 1}M$
car cela montrera que
$\bigoplus_{n \geq 0}
\overline{J}^n\overline{M}/\overline{J}^{n + 1}\overline{M}$
est le quotient de
$\bigoplus_{n \geq 0} J^nM/J^{n + 1}M \cong M/JM[X_1, \ldots, X_c]$
par $X_1$. En fait, on a $J^nM \cap f_1M = f_1J^{n - 1}M$.
En effet, si $m \not \in J^{n - 1}M$, alors $f_1m \not \in J^nM$
car $\bigoplus J^nM/J^{n + 1}M$ est l'algèbre de polynômes
$M/J[X_1, \ldots, X_c]$ par hypothèse.
\end{proof}

\begin{lemma}
\label{lemma-quasi-regular-regular}
Soit $(R, \mathfrak m)$ un anneau local noethérien.
Soit $M$ un $R$-module de type fini non nul.
Soit $f_1, \ldots, f_c \in \mathfrak m$ une suite $M$-quasi-régulière.
Alors la suite $f_1, \ldots, f_c$ est $M$-régulière.
\end{lemma}

\begin{proof}
Posons $J = (f_1, \ldots, f_c)$.
Montrons que $f_1$ est un non-diviseur de zéro sur $M$.
Supposons $x \in M$ non nul.
D'après le théorème d'intersection de Krull, il existe un entier $r$
tel que $x \in J^rM$ mais $x \not \in J^{r + 1}M$ ; voir le
Lemme \ref{lemma-intersect-powers-ideal-module-zero}.
Alors $f_1 x \in J^{r + 1}M$ est un élément dont la classe dans
$J^{r + 1}M/J^{r + 2}M$ est non nulle, d'après la structure supposée de
$\bigoplus J^nM/J^{n + 1}M$. Par conséquent $f_1x \not = 0$.

\medskip\noindent
On peut maintenant achever la démonstration par récurrence sur $c$ en utilisant le
Lemme \ref{lemma-truncate-quasi-regular}.
\end{proof}

\begin{remark}[Autres types de suites régulières]
\label{remark-koszul-regular}
Dans l'article \cite{Kabele}, l'auteur étudie deux autres
conditions de régularité pour les suites $x_1, \ldots, x_r$ d'éléments
d'un anneau $R$. On dit précisément que la suite est {\it Koszul-régulière}
si $H_i(K_{\bullet}(R, x_{\bullet})) = 0$ pour $i \geq 1$, où
$K_{\bullet}(R, x_{\bullet})$ est le complexe de Koszul. La suite est
dite {\it $H_1$-régulière} si $H_1(K_{\bullet}(R, x_{\bullet})) = 0$.
On a les implications régulière $\Rightarrow$ Koszul-régulière $\Rightarrow$
$H_1$-régulière $\Rightarrow$ quasi-régulière. Par des exemples, l'auteur montre que
ces implications ne sont en général pas réversibles, même si $R$ est un anneau local
(non noethérien) et si la suite engendre l'idéal maximal de
$R$. Nous présentons ces notions plus en détail dans
Compléments d'algèbre, section \ref{more-algebra-section-koszul-regular}.
\end{remark}

\begin{remark}
\label{remark-join-quasi-regular-sequences}
Soit $k$ un corps. Considérons l'anneau
$$
A = k[x, y, w, z_0, z_1, z_2, \ldots]/
(y^2z_0 - wx, z_0 - yz_1, z_1 - yz_2, \ldots)
$$
Dans cet anneau, $x$ est un non-diviseur de zéro et l'image de $y$ dans
$A/xA$ donne une suite quasi-régulière. Mais il n'est pas vrai que la suite
$x, y$ soit quasi-régulière dans $A$, car $(x, y)/(x, y)^2$
n'est pas libre de rang deux sur $A/(x, y)$ : en effet,
$wx = 0$ dans $(x, y)/(x, y)^2$, tandis que $w$ n'est pas nul dans $A/(x, y)$.
Ainsi, l'analogue du
Lemme \ref{lemma-join-regular-sequences}
n'est pas vrai pour les suites quasi-régulières.
\end{remark}

\begin{lemma}
\label{lemma-quasi-regular-on-quotient}
Soit $R$ un anneau. Soit $J = (f_1, \ldots, f_r)$ un idéal de $R$.
Soit $M$ un $R$-module. Posons $\overline{R} = R/\bigcap_{n \geq 0} J^n$,
$\overline{M} = M/\bigcap_{n \geq 0} J^nM$, et notons
$\overline{f}_i$ l'image de $f_i$ dans $\overline{R}$.
Alors la suite $f_1, \ldots, f_r$ est $M$-quasi-régulière si et seulement si
la suite $\overline{f}_1, \ldots, \overline{f}_r$ est $\overline{M}$-quasi-régulière.
\end{lemma}

\begin{proof}
Cela est vrai parce que
$J^nM/J^{n + 1}M \cong
\overline{J}^n\overline{M}/\overline{J}^{n + 1}\overline{M}$.
\end{proof}





\section{Algèbres d'éclatement}
\label{section-blow-up}

\noindent
Dans cette section, nous présentons quelques observations élémentaires sur l'éclatement.

\begin{definition}
\label{definition-blow-up}
Soit $R$ un anneau.
Soit $I \subset R$ un idéal.
\begin{enumerate}
\item L'{\it algèbre d'éclatement}, ou {\it algèbre de Rees}, associée à
la paire $(R, I)$ est la $R$-algèbre graduée
$$
\text{Bl}_I(R) =
\bigoplus\nolimits_{n \geq 0} I^n =
R \oplus I \oplus I^2 \oplus \ldots
$$
où le facteur direct $I^n$ est placé en degré $n$.
\item Soit $a \in I$ un élément. Notons $a^{(1)}$ l'élément $a$
considéré comme un élément de degré $1$ de l'algèbre de Rees. L'
{\it algèbre d'éclatement affine} $R[\frac{I}{a}]$ est alors l'algèbre
$(\text{Bl}_I(R))_{(a^{(1)})}$ construite dans la section \ref{section-proj}.
\end{enumerate}
\end{definition}

\noindent
Autrement dit, un élément de $R[\frac{I}{a}]$ est représenté par
une expression de la forme $x/a^n$ avec $x \in I^n$. Deux représentants
$x/a^n$ et $y/a^m$ définissent le même élément si et seulement si
$a^k(a^mx - a^ny) = 0$ pour un certain $k \geq 0$.

\begin{lemma}
\label{lemma-affine-blowup}
Soit $R$ un anneau, $I \subset R$ un idéal et $a \in I$.
Soit $R' = R[\frac{I}{a}]$ l'algèbre d'éclatement affine. Alors
\begin{enumerate}
\item l'image de $a$ dans $R'$ est un non-diviseur de zéro,
\item $IR' = aR'$, et
\item $(R')_a = R_a$.
\end{enumerate}
\end{lemma}

\begin{proof}
Cela résulte immédiatement de la description de $R[\frac{I}{a}]$ ci-dessus.
\end{proof}

\begin{lemma}
\label{lemma-blowup-base-change}
Soit $R \to S$ un homomorphisme d'anneaux. Soit $I \subset R$ un idéal
et soit $a \in I$. Posons $J = IS$ et soit $b \in J$ l'image de $a$.
Alors $S[\frac{J}{b}]$ est le quotient de $S \otimes_R R[\frac{I}{a}]$
par l'idéal des éléments annulés par une puissance de $b$.
\end{lemma}

\begin{proof}
Soit $S'$ le quotient de $S \otimes_R R[\frac{I}{a}]$ par ses
éléments de torsion pour les puissances de $b$. L'homomorphisme d'anneaux
$$
S \otimes_R R[\textstyle{\frac{I}{a}}]
\longrightarrow
S[\textstyle{\frac{J}{b}}]
$$
est surjectif et annule la torsion pour les puissances de $a$, puisque $b$ est un non-diviseur de zéro
dans $S[\frac{J}{b}]$. On obtient donc une application surjective $S' \to S[\frac{J}{b}]$.
Pour voir que son noyau est trivial, construisons une application inverse. Soit
$z = y/b^n$ un élément de $S[\frac{J}{b}]$, c'est-à-dire $y \in J^n$.
Écrivons $y = \sum x_is_i$ avec $x_i \in I^n$ et $s_i \in S$.
On envoie $z$ sur la classe de $\sum s_i \otimes x_i/a^n$ dans
$S'$. Cette définition est bonne, car un élément du noyau de l'application
$S \otimes_R I^n \to J^n$ est annulé par $a^n$, donc a une image nulle dans $S'$.
\end{proof}

\begin{example}
\label{example-rees-algebra-polynomial}
Soit $R$ un anneau. Soit $P = R[t_1, \ldots, t_n]$ l'algèbre de polynômes.
Soit $I = (t_1, \ldots, t_n) \subset P$. Avec les notations de la
Définition \ref{definition-blow-up}, il existe un isomorphisme
$$
P[T_1, \ldots, T_n]/(t_iT_j - t_jT_i) \longrightarrow \text{Bl}_I(P)
$$
qui envoie $T_i$ sur $t_i^{(1)}$. Nous laissons au lecteur le soin de montrer que
cette application est bien définie. Puisque $I$ est engendré par $t_1, \ldots, t_n$,
on voit que notre application est surjective. Pour montrer qu'elle est injective,
il faut établir ceci : pour tout $e \geq 1$, le $P$-module $I^e$ est engendré
par les monômes $t^E = t_1^{e_1} \ldots x_n^{e_n}$, pour les multi-indices
$E = (e_1, \ldots, e_n)$ de degré $|E| = e$, soumis uniquement aux
relations $t_i t^E = t_j t^{E'}$ lorsque $|E| = |E'| = e$ et
$e_a + \delta_{a i} = e'_a + \delta_{a j},\ a = 1, \ldots, n$
(delta de Kronecker). Nous omettons les détails.
\end{example}

\begin{example}
\label{example-affine-blowup-algebra-polynomial}
Soit $R$ un anneau. Soit $P = R[t_1, \ldots, t_n]$ l'algèbre de polynômes.
Soit $I = (t_1, \ldots, t_n) \subset P$. Soit $a = t_1$. Avec les notations de la
Définition \ref{definition-blow-up}, il existe un isomorphisme
$$
P[x_2, \ldots, x_n]/(t_1x_2 - t_2, \ldots, t_1x_n - t_n)
\longrightarrow
\textstyle{P[\frac{I}{a}] = P[\frac{I}{t_1}]}
$$
qui envoie $x_i$ sur $t_i/t_1$. Nous laissons au lecteur le soin de montrer que
cette application est bien définie. Puisque $I$ est engendré par $t_1, \ldots, t_n$,
on voit que notre application est surjective. Pour montrer qu'elle est injective,
le lecteur peut observer que la source et le but de notre application sont
sans $t_1$-torsion et que l'application est un isomorphisme après inversion de
$t_1$ ; voir le Lemme \ref{lemma-affine-blowup}. Le lecteur peut aussi
utiliser la description de l'algèbre de Rees donnée dans
l'Exemple \ref{example-rees-algebra-polynomial}.
Nous omettons les détails.
\end{example}

\begin{lemma}
\label{lemma-affine-blowup-quotient-description}
Soit $R$ un anneau. Soit $I = (a_1, \ldots, a_n)$ un idéal de $R$.
Soit $a = a_1$. Il existe alors une surjection
$$
R[x_2, \ldots, x_n]/(a x_2 - a_2, \ldots, a x_n - a_n)
\longrightarrow
\textstyle{R[\frac{I}{a}]}
$$
dont le noyau est la torsion pour les puissances de $a$ dans la source.
\end{lemma}

\begin{proof}
Considérons l'homomorphisme d'anneaux $P = \mathbf{Z}[t_1, \ldots, t_n] \to R$
qui envoie $t_i$ sur $a_i$. Posons $J = (t_1, \ldots, t_n)$.
D'après l'Exemple \ref{example-affine-blowup-algebra-polynomial}, on a
$P[\frac{J}{t_1}] =
P[x_2, \ldots, x_n]/(t_1 x_2 - t_2, \ldots, t_1 x_n - t_n)$.
On conclut en appliquant le Lemme \ref{lemma-blowup-base-change} à l'application $P \to A$.
\end{proof}

\begin{lemma}
\label{lemma-blowup-in-principal}
Soit $R$ un anneau, $I \subset R$ un idéal et $a \in I$.
Posons $R' = R[\frac{I}{a}]$. Si $f \in R$ vérifie $V(f) = V(I)$,
alors l'image de $f$ dans $R'$ est un non-diviseur de zéro et $R'_f = R'_a = R_a$.
\end{lemma}

\begin{proof}
Nous utiliserons les résultats du Lemme \ref{lemma-affine-blowup}
sans les mentionner davantage.
L'hypothèse $V(f) = V(I)$ implique $V(fR') = V(IR') = V(aR')$.
Ainsi $a^n = fb$ et $f^m = ac$ pour certains $b, c \in R'$.
Le lemme en résulte.
\end{proof}

\begin{lemma}
\label{lemma-blowup-add-principal}
Soit $R$ un anneau, $I \subset R$ un idéal, $a \in I$ et $f \in R$.
Posons $R' = R[\frac{I}{a}]$ et $R'' = R[\frac{fI}{fa}]$. Alors
il existe un homomorphisme surjectif de $R$-algèbres $R' \to R''$ dont le noyau
est l'ensemble des éléments de $f$-torsion de $R'$.
\end{lemma}

\begin{proof}
L'application envoie $x/a^n$, pour $x \in I^n$, sur $f^nx/(fa)^n$.
On vérifie immédiatement que cette application est bien définie et surjective.
Puisque $af$ est un non-diviseur de zéro dans $R''$
(Lemme \ref{lemma-affine-blowup}), l'ensemble des éléments de torsion
pour les puissances de $f$ est envoyé sur zéro. Réciproquement, si $x \in R'$
et si $f^n x \not = 0$ pour tout $n > 0$, alors $(af)^n x \not = 0$
pour tout $n$, puisque $a$ est un non-diviseur de zéro dans $R'$. Il s'ensuit
que l'image de $x$ dans $R''$ n'est pas nulle, d'après la description de
$R''$ donnée après la Définition \ref{definition-blow-up}.
\end{proof}

\begin{lemma}
\label{lemma-blowup-reduced}
\begin{slogan}
La propriété d'être réduit est invariante par éclatement
\end{slogan}
Si $R$ est réduit, toute algèbre d'éclatement (affine) de $R$ est réduite.
\end{lemma}

\begin{proof}
Soit $I \subset R$ un idéal et soit $a \in I$. Supposons que $x/a^n$, avec
$x \in I^n$, soit un élément nilpotent de $R[\frac{I}{a}]$. Alors
$(x/a^n)^m = 0$. Ainsi $a^N x^m = 0$ dans $R$ pour un certain $N \geq 0$.
Après avoir augmenté $N$ si nécessaire, on peut supposer $N = me$ pour un certain
$e \geq 0$. Alors $(a^e x)^m = 0$ et, puisque $R$ est réduit, on obtient
$a^e x = 0$. Cela signifie que $x/a^n = 0$ dans $R[\frac{I}{a}]$.
\end{proof}

\begin{lemma}
\label{lemma-blowup-domain}
Soit $R$ un anneau intègre, $I \subset R$ un idéal et $a \in I$ un élément
non nul. Alors l'algèbre d'éclatement affine $R[\frac{I}{a}]$ est intègre.
\end{lemma}

\begin{proof}
Supposons que $x/a^n$, $y/a^m$, avec $x \in I^n$, $y \in I^m$,
soient des éléments de $R[\frac{I}{a}]$ dont le produit est nul.
Alors $a^N x y = 0$ dans $R$. Puisque $R$ est intègre, on en déduit
que $x = 0$ ou $y = 0$.
\end{proof}

\begin{lemma}
\label{lemma-blowup-dominant}
Soit $R$ un anneau. Soit $I \subset R$ un idéal. Soit $a \in I$.
Si $a$ n'appartient à aucun idéal premier minimal de $R$, alors
$\Spec(R[\frac{I}{a}]) \to \Spec(R)$ a une image dense.
\end{lemma}

\begin{proof}
Si $a^k x = 0$ pour $x \in R$, alors $x$ appartient à tous les
idéaux premiers minimaux de $R$ et est donc nilpotent ; voir le
Lemme \ref{lemma-Zariski-topology}.
Ainsi, le noyau de $R \to R[\frac{I}{a}]$ est constitué d'éléments
nilpotents. Le résultat découle donc du
Lemme \ref{lemma-image-dense-generic-points}.
\end{proof}

\begin{lemma}
\label{lemma-valuation-ring-colimit-affine-blowups}
Soit $(R, \mathfrak m)$ un anneau local intègre de corps des fractions $K$.
Soit $R \subset A \subset K$ un anneau de valuation qui domine $R$.
Alors
$$
A = \colim R[\textstyle{\frac{I}{a}}]
$$
est une colimite filtrante d'éclatements affines $R \to R[\frac{I}{a}]$ possédant
les propriétés suivantes :
\begin{enumerate}
\item $a \in I \subset \mathfrak m$,
\item $I$ est de type fini, et
\item l'anneau de la fibre de $R \to R[\frac{I}{a}]$ en $\mathfrak m$
n'est pas nul.
\end{enumerate}
\end{lemma}

\begin{proof}
Toute algèbre d'éclatement $R[\frac{I}{a}]$ est un anneau intègre contenu dans $K$ ; voir le
Lemme \ref{lemma-blowup-domain}. Le lemme dit simplement que $A$ est la
réunion filtrante de celles pour lesquelles $a \in I$ et les propriétés (1), (2), (3) sont satisfaites.
Si $R[\frac{I}{a}] \subset A$ et $R[\frac{J}{b}] \subset A$, alors
on a
$$
R[\textstyle{\frac{I}{a}}] \cup R[\textstyle{\frac{J}{b}}]  \subset
R[\textstyle{\frac{IJ}{ab}}] \subset A
$$
La première inclusion vient de ce que $x/a^n = b^nx/(ab)^n$, et la seconde
de ce que, si $z \in (IJ)^n$, alors $z = \sum x_iy_i$ avec $x_i \in I^n$
et $y_i \in J^n$, et donc $z/(ab)^n = \sum (x_i/a^n)(y_i/b^n)$
appartient à $A$.

\medskip\noindent
Considérons une partie finie $E \subset A$. Écrivons $E = \{e_1, \ldots, e_n\}$.
Choisissons $a \in R$ non nul tel que l'on puisse écrire $e_i = f_i/a$ pour
tout $i = 1, \ldots, n$. Posons $I = (f_1, \ldots, f_n, a)$.
Nous affirmons que $R[\frac{I}{a}] \subset A$. C'est clair, car un élément
de $R[\frac{I}{a}]$ peut être représenté par un polynôme en les éléments
$e_i$. Le lemme résulte immédiatement de cette observation.
\end{proof}










\section{Groupes Ext}
\label{section-ext}

\noindent
Dans cette section, nous faisons un peu d'algèbre homologique
afin d'établir quelques propriétés fondamentales de la
profondeur sur les anneaux locaux noethériens.

\begin{lemma}
\label{lemma-resolution-by-finite-free}
Soit $R$ un anneau. Soit $M$ un $R$-module.
\begin{enumerate}
\item Il existe un complexe exact
$$
\ldots \to F_2 \to F_1 \to F_0 \to M \to 0.
$$
où les $F_i$ sont des $R$-modules libres.
\item Si $R$ est noethérien et si $M$ est de type fini sur $R$, on
peut choisir le complexe de sorte que les $F_i$ soient libres de type fini.
Autrement dit, on peut trouver un complexe exact
$$
\ldots \to R^{\oplus n_2} \to R^{\oplus n_1} \to R^{\oplus n_0} \to M \to 0.
$$
\end{enumerate}
\end{lemma}

\begin{proof}
Expliquons seulement le cas noethérien.
On choisit d'abord une surjection $R^{n_0} \to M$.
Puis, après avoir construit un complexe exact de longueur
$e$, il suffit de choisir une surjection $R^{n_{e + 1}} \to
\Ker(R^{n_e} \to R^{n_{e-1}})$, ce qui est possible
puisque $R$ est noethérien.
\end{proof}

\begin{definition}
\label{definition-finite-free-resolution}
Soit $R$ un anneau. Soit $M$ un $R$-module.
\begin{enumerate}
\item Une {\it résolution} (à gauche) $F_\bullet \to M$ de $M$ est un complexe exact
$$
\ldots \to F_2 \to F_1 \to F_0 \to M \to 0
$$
de $R$-modules.
\item Une {\it résolution de $M$ par des $R$-modules libres} est une résolution
$F_\bullet \to M$ dans laquelle chaque $F_i$ est un $R$-module libre.
\item Une {\it résolution de $M$ par des $R$-modules libres de type fini} est une résolution
$F_\bullet \to M$ dans laquelle chaque $F_i$ est un $R$-module libre de type fini.
\end{enumerate}
\end{definition}

\noindent
Nous employons souvent la notation $F_{\bullet}$ pour désigner un complexe
de $R$-modules
$$
\ldots \to F_i \to F_{i-1} \to \ldots
$$
Dans ce cas, nous notons souvent $d_i$ ou $d_{F, i}$ l'application
$F_i \to F_{i-1}$. Dans cette section, nous supposerons toujours
que $F_0$ est le dernier terme non nul du complexe.
Le {\it $i$-ième groupe d'homologie du complexe} $F_{\bullet}$
est le groupe $H_i = \Ker(d_{F, i})/\Im(d_{F, i + 1})$.
Un {\it morphisme de complexes $\alpha : F_{\bullet} \to G_{\bullet}$}
est donné par des applications $\alpha_i : F_i \to G_i$ telles que
$\alpha_{i-1} \circ d_{F, i} = d_{G, i-1} \circ \alpha_i$.
Un tel morphisme induit en homologie une application $H_i(\alpha) :
H_i(F_{\bullet}) \to H_i(G_{\bullet})$. Si $\alpha, \beta
:  F_{\bullet} \to G_{\bullet}$ sont des morphismes de complexes, une
{\it homotopie} entre $\alpha$ et $\beta$ est donnée par
une famille d'applications $h_i : F_i \to G_{i + 1}$ telles que
$\alpha_i - \beta_i = d_{G, i + 1} \circ h_i +
h_{i-1} \circ d_{F, i}$.
Deux morphismes $\alpha, \beta : F_{\bullet} \to G_{\bullet}$ sont dits
{\it homotopes} s'il existe une homotopie entre $\alpha$
et $\beta$.

\medskip\noindent
Nous emploierons des notations tout à fait analogues pour les complexes
de la forme $F^{\bullet}$ qui s'écrivent
$$
\ldots \to F^i \xrightarrow{d^i} F^{i + 1} \to \ldots
$$
On dispose de morphismes de complexes, d'homotopies, etc.
Dans ce cas, nous posons $H^i(F^{\bullet}) =
\Ker(d^i)/\Im(d^{i - 1})$ et nous l'appelons le
{\it $i$-ième groupe de cohomologie}.

\begin{lemma}
\label{lemma-homotopic-equal-homology}
Deux morphismes de complexes homotopes quelconques induisent les mêmes applications sur les
groupes de (co)homologie.
\end{lemma}

\begin{proof}
Omis.
\end{proof}

\begin{lemma}
\label{lemma-compare-resolutions}
Soit $R$ un anneau. Soit $M \to N$ un homomorphisme de $R$-modules.
Soit $N_\bullet \to N$ une résolution quelconque.
Soit
$$
\ldots \to F_2 \to F_1 \to F_0 \to M
$$
un complexe de $R$-modules dans lequel chaque $F_i$ est un $R$-module libre. Alors
\begin{enumerate}
\item il existe un morphisme de complexes $F_\bullet \to N_\bullet$ tel que
$$
\xymatrix{
F_0 \ar[r] \ar[d] & M \ar[d] \\
N_0 \ar[r] & N
}
$$
soit commutatif, et
\item deux morphismes quelconques $\alpha, \beta : F_\bullet \to N_\bullet$ comme en (1)
sont homotopes.
\end{enumerate}
\end{lemma}

\begin{proof}
Démonstration de (1). Puisque $F_0$ est libre, on peut trouver une application $F_0 \to N_0$
qui relève l'application $F_0 \to M \to N$. On obtient une application induite
$F_1 \to F_0 \to N_0$ dont l'image est contenue dans l'image
de $N_1 \to N_0$. Puisque $F_1$ est libre, on peut la relever
en une application $F_1 \to N_1$. Celle-ci induit à son tour une application
$F_2 \to F_1 \to N_1$ dont l'image dans
$N_0$ est nulle. Comme $N_\bullet$ est exact, on voit que
l'image de cette application est contenue dans l'image
de $N_2 \to N_1$. On peut donc la relever en une application
$F_2 \to N_2$. On poursuit ainsi.

\medskip\noindent
Démonstration de (2). Pour montrer que $\alpha, \beta$ sont homotopes, il suffit
de montrer que la différence $\gamma = \alpha - \beta$ est homotope
à zéro. Remarquons que l'image de $\gamma_0 : F_0 \to N_0$
est contenue dans l'image de $N_1 \to N_0$. On peut donc relever
$\gamma_0$ en une application $h_0 : F_0 \to N_1$. Considérons l'application
$\gamma_1' = \gamma_1 - h_0 \circ d_{F, 1}$. Par le choix de $h_0$,
on voit que l'image de $\gamma_1'$ est contenue dans
le noyau de $N_1 \to N_0$. Comme $N_\bullet$ est exact,
on peut relever $\gamma_1'$ en une application $h_1 : F_1 \to N_2$.
À ce stade, on a $\gamma_1 = h_0 \circ d_{F, 1}
+ d_{N, 2} \circ h_1$. On poursuit ainsi.
\end{proof}

\noindent
Nous sommes maintenant en mesure de définir les groupes
$\Ext^i_R(M, N)$. Choisissons une résolution
$F_{\bullet}$ de $M$ par des $R$-modules libres, voir le Lemme
\ref{lemma-resolution-by-finite-free}. Considérons
le complexe (cohomologique)
$$
\Hom_R(F_\bullet, N) :
\Hom_R(F_0, N) \to
\Hom_R(F_1, N) \to
\Hom_R(F_2, N) \to \ldots
$$
On définit $\Ext^i_R(M, N)$, pour $i \geq 0$, comme le $i$-ième
groupe de cohomologie de ce complexe\footnote{À ce stade,
il conviendrait peut-être davantage de dire ``un'' plutôt
que ``le'' groupe Ext.}. Pour $i < 0$, on pose $\Ext^i_R(M, N) = 0$.
Avant de poursuivre, signalons que
$$
\Ext^0_R(M, N) = \Ker(\Hom_R(F_0, N) \to \Hom_R(F_1, N)) =
\Hom_R(M, N)
$$
car on peut appliquer la partie (1) du Lemme \ref{lemma-hom-exact} à
la suite exacte $F_1 \to F_0 \to M \to 0$.
Le lemme suivant précise
en quel sens cette définition est bien définie.

\begin{lemma}
\label{lemma-ext-welldefined}
Soit $R$ un anneau. Soient $M_1, M_2, N$ des $R$-modules.
Supposons que $F_{\bullet}$ soit une résolution libre du module $M_1$,
et que $G_{\bullet}$ soit une résolution libre du module $M_2$.
Soit $\varphi : M_1 \to M_2$ un homomorphisme de modules.
Soit $\alpha : F_{\bullet} \to G_{\bullet}$
un morphisme de complexes qui induit $\varphi$ sur
$M_1 = \Coker(d_{F, 1}) \to M_2 = \Coker(d_{G, 1})$,
voir le Lemme \ref{lemma-compare-resolutions}.
Alors les applications induites
$$
H^i(\alpha) :
H^i(\Hom_R(F_{\bullet}, N))
\longrightarrow
H^i(\Hom_R(G_{\bullet}, N))
$$
ne dépendent pas du choix de $\alpha$.
Si $\varphi$ est un isomorphisme, toutes les applications
$H^i(\alpha)$ sont des isomorphismes. Si $M_1 = M_2$, $F_\bullet = G_\bullet$, et
si $\varphi$ est l'identité, toutes les applications $H_i(\alpha)$ le sont aussi.
\end{lemma}

\begin{proof}
Un autre morphisme $\beta : F_{\bullet} \to G_{\bullet}$
qui induit $\varphi$ est homotope à $\alpha$ d'après le
Lemme \ref{lemma-compare-resolutions}. Les
applications $\Hom_R(F_\bullet, N) \to
\Hom_R(G_\bullet, N)$ sont donc homotopes.
L'indépendance résulte alors du
Lemme \ref{lemma-homotopic-equal-homology}.

\medskip\noindent
Supposons que $\varphi$ soit un isomorphisme.
Soit $\psi : M_2 \to M_1$ un inverse.
Choisissons $\beta : G_{\bullet} \to F_{\bullet}$
qui induit $\psi :
M_2 = \Coker(d_{G, 1}) \to M_1 = \Coker(d_{F, 1})$,
voir le Lemme \ref{lemma-compare-resolutions}.
Considérons maintenant l'application
$H^i(\alpha) \circ H^i(\beta) =
H^i(\alpha \circ \beta)$. D'après ce qui précède, l'application
$H^i(\alpha \circ \beta)$ est la {\it même}
que l'application $H^i(\text{id}_{G_{\bullet}}) = \text{id}$.
Il en va de même de la composée $H^i(\beta) \circ H^i(\alpha)$.
Ainsi $H^i(\alpha)$ et $H^i(\beta)$ sont inverses l'une de l'autre.
\end{proof}

\begin{lemma}
\label{lemma-long-exact-seq-ext}
Soit $R$ un anneau. Soit $M$ un $R$-module.
Soit $0 \to N' \to N \to N'' \to 0$ une
suite exacte courte. On obtient alors une suite exacte
longue
$$
\begin{matrix}
0
\to \Hom_R(M, N')
\to \Hom_R(M, N)
\to \Hom_R(M, N'')
\\
\phantom{0\ }
\to \Ext^1_R(M, N')
\to \Ext^1_R(M, N)
\to \Ext^1_R(M, N'')
\to \ldots
\end{matrix}
$$
\end{lemma}

\begin{proof}
Choisissons une résolution libre $F_{\bullet} \to M$.
Comme chacun des $F_i$ est libre, on obtient
une suite exacte courte de complexes
$$
0 \to
\Hom_R(F_{\bullet}, N') \to
\Hom_R(F_{\bullet}, N) \to
\Hom_R(F_{\bullet}, N'') \to
0
$$
La suite exacte longue résulte alors du
lemme du serpent appliqué à cette suite.
\end{proof}

\begin{lemma}
\label{lemma-reverse-long-exact-seq-ext}
Soit $R$ un anneau. Soit $N$ un $R$-module.
Soit $0 \to M' \to M \to M'' \to 0$ une
suite exacte courte. On obtient alors une suite exacte
longue
$$
\begin{matrix}
0
\to \Hom_R(M'', N)
\to \Hom_R(M, N)
\to \Hom_R(M', N)
\\
\phantom{0\ }
\to \Ext^1_R(M'', N)
\to \Ext^1_R(M, N)
\to \Ext^1_R(M', N)
\to \ldots
\end{matrix}
$$
\end{lemma}

\begin{proof}
Choisissons des systèmes de générateurs $\{m'_{i'}\}_{i' \in I'}$ et
$\{m''_{i''}\}_{i'' \in I''}$ de $M'$ et de $M''$.
Pour chaque $i'' \in I''$, choisissons un relèvement $\tilde m''_{i''} \in M$
de l'élément $m''_{i''} \in M''$. Posons $F' = \bigoplus_{i' \in I'} R$,
$F'' = \bigoplus_{i'' \in I''} R$ et $F = F' \oplus F''$.
En envoyant les générateurs de ces modules libres sur les générateurs choisis
correspondants, on obtient des homomorphismes surjectifs de $R$-modules $F' \to M'$,
$F'' \to M''$ et $F \to M$. On obtient un morphisme de suites exactes courtes
$$
\begin{matrix}
0 & \to & M' & \to & M & \to & M'' & \to & 0 \\
& & \uparrow & & \uparrow & & \uparrow \\
0 & \to & F' & \to & F & \to & F'' & \to & 0 \\
\end{matrix}
$$
Le lemme du serpent montre que la suite des noyaux
$0 \to K' \to K \to K'' \to 0$ est une suite exacte courte de $R$-modules.
On peut donc poursuivre indéfiniment ce procédé. Autrement dit,
on obtient une suite exacte courte de résolutions qui s'insère dans le diagramme
$$
\begin{matrix}
0 & \to & M' & \to & M & \to & M'' & \to & 0 \\
& & \uparrow & & \uparrow & & \uparrow \\
0 & \to & F_\bullet' & \to & F_\bullet & \to & F_\bullet'' & \to & 0 \\
\end{matrix}
$$
Comme chacune des suites $0 \to F'_n \to F_n \to F''_n \to 0$
est exacte scindée (par construction), on obtient une suite exacte courte de
complexes
$$
0 \to
\Hom_R(F''_{\bullet}, N) \to
\Hom_R(F_{\bullet}, N) \to
\Hom_R(F'_{\bullet}, N) \to
0
$$
en appliquant le foncteur $\Hom_R(-, N)$.
La suite exacte longue résulte alors du
lemme du serpent appliqué à cette suite.
\end{proof}

\begin{lemma}
\label{lemma-annihilate-ext}
Soit $R$ un anneau. Soient $M$, $N$ des $R$-modules.
Tout $x\in R$ tel que $xN = 0$ ou $xM = 0$
annule chacun des modules $\Ext^i_R(M, N)$.
\end{lemma}

\begin{proof}
Choisissons une résolution libre $F_{\bullet}$ de $M$.
Puisque $\Ext^i_R(M, N)$
est défini comme la cohomologie du complexe
$\Hom_R(F_{\bullet}, N)$, le lemme est
clair lorsque $xN = 0$. Si $xM = 0$, alors
la multiplication par $x$ sur $F_{\bullet}$
relève l'application nulle sur $M$. Ainsi, d'après le Lemme
\ref{lemma-ext-welldefined}, elle induit
la même application sur les groupes Ext que l'application
nulle.
\end{proof}

\begin{lemma}
\label{lemma-ext-noetherian}
Soit $R$ un anneau noethérien. Soient $M$, $N$ des $R$-modules de type fini.
Alors $\Ext^i_R(M, N)$ est un $R$-module de type fini pour tout $i$.
\end{lemma}

\begin{proof}
Cela résulte de ce que $\Ext^i_R(M, N)$ se calcule comme les
groupes de cohomologie d'un complexe $\Hom_R(F_\bullet, N)$
où chaque $F_n$ est un $R$-module libre de type fini ; voir le
Lemme \ref{lemma-resolution-by-finite-free}.
\end{proof}




\section{Profondeur}
\label{section-depth}

\noindent
Voici notre définition.

\begin{definition}
\label{definition-depth}
Soit $R$ un anneau, et soit $I \subset R$ un idéal. Soit $M$ un $R$-module de type fini.
La {\it profondeur relativement à $I$} de $M$, notée $\text{depth}_I(M)$, est définie comme suit :
\begin{enumerate}
\item si $IM \not = M$, alors $\text{depth}_I(M)$ est la borne supérieure dans
$\{0, 1, 2, \ldots, \infty\}$ des longueurs des suites $M$-régulières dans $I$,
\item si $IM = M$, on pose $\text{depth}_I(M) = \infty$.
\end{enumerate}
Si $(R, \mathfrak m)$ est local, on appelle simplement $\text{depth}_{\mathfrak m}(M)$ la
{\it profondeur} de $M$.
\end{definition}

\noindent
Explication. D'après la Définition \ref{definition-regular-sequence}, la suite vide
n'est pas une suite régulière sur le module nul, mais, en pratique,
il s'avère commode de fixer la profondeur du module $0$
égale à $+\infty$. Remarquons que, si $I = R$, alors $\text{depth}_I(M) = \infty$
pour tout $R$-module de type fini $M$. Si $I$ est contenu dans le radical de Jacobson
de $R$ (par exemple, si $R$ est local et si $I \subset \mathfrak m_R$), alors
$M \not = 0 \Rightarrow IM \not = M$ d'après le lemme de Nakayama.
Un module $M$ a une profondeur relativement à $I$ égale à $0$ si et seulement si $M$ est non nul et si $I$
ne contient aucun non-diviseur de zéro sur $M$.

\medskip\noindent
L'Exemple \ref{example-global-regular} montre que la profondeur ne se comporte pas
bien, même si l'anneau est noethérien, et l'Exemple
\ref{example-local-regular} montre qu'elle ne se comporte pas
bien si l'anneau est local mais non noethérien.
Nous verrons que la profondeur se comporte bien si l'anneau est local noethérien.

\begin{lemma}
\label{lemma-depth-weak-sequence}
Soit $R$ un anneau, $I \subset R$ un idéal, et $M$ un $R$-module de type fini.
Alors $\text{depth}_I(M)$ est égale à la borne supérieure des longueurs des
suites $f_1, \ldots, f_r \in I$ telles que $f_i$ soit un non-diviseur de zéro
sur $M/(f_1, \ldots, f_{i - 1})M$.
\end{lemma}

\begin{proof}
Supposons que $IM = M$. Alors le Lemme \ref{lemma-NAK} montre qu'il existe
$f \in I$ tel que $f : M \to M$ soit $\text{id}_M$. Ainsi
$f, 0, 0, 0, \ldots$ est une suite infinie de
non-diviseurs de zéro successifs, et les deux quantités coïncident dans ce cas.
Si $IM \not =  M$, une suite comme dans le lemme
est une suite $M$-régulière, et les deux quantités coïncident là encore.
\end{proof}

\begin{lemma}
\label{lemma-bound-depth}
Soit $(R, \mathfrak m)$ un anneau local noethérien.
Soit $M$ un $R$-module de type fini non nul.
Alors $\dim(\text{Supp}(M)) \geq \text{depth}(M)$.
\end{lemma}

\begin{proof}
La démonstration se fait par récurrence sur $\dim(\text{Supp}(M))$.
Si $\dim(\text{Supp}(M)) = 0$, alors
$\text{Supp}(M) = \{\mathfrak m\}$, d'où $\text{Ass}(M) = \{\mathfrak m\}$
(d'après les Lemmes \ref{lemma-ass-support} et \ref{lemma-ass-zero}), et donc
la profondeur de $M$ est nulle, par exemple d'après le
Lemme \ref{lemma-ideal-nonzerodivisor}.
Pour l'étape de récurrence, supposons $\dim(\text{Supp}(M)) > 0$.
Soit $f_1, \ldots, f_d$ une suite d'éléments de $\mathfrak m$
telle que $f_i$ soit un non-diviseur de zéro sur $M/(f_1, \ldots, f_{i - 1})M$.
D'après le Lemme \ref{lemma-depth-weak-sequence}, il suffit de montrer que
$\dim(\text{Supp}(M)) \geq d$. On peut supposer
$d > 0$, sans quoi le lemme est établi. D'après le
Lemme \ref{lemma-one-equation-module},
on a $\dim(\text{Supp}(M/f_1M)) = \dim(\text{Supp}(M)) - 1$.
Par récurrence, on conclut que $\dim(\text{Supp}(M/f_1M)) \geq d - 1$,
comme souhaité.
\end{proof}

\begin{lemma}
\label{lemma-depth-finite-noetherian}
Soit $R$ un anneau noethérien, $I \subset R$ un idéal, et $M$ un
$R$-module non nul de type fini tel que $IM \not = M$. Alors
$\text{depth}_I(M) < \infty$.
\end{lemma}

\begin{proof}
Puisque $M/IM$ est non nul, on peut choisir $\mathfrak p \in \text{Supp}(M/IM)$
d'après le Lemme \ref{lemma-support-zero}. Alors $(M/IM)_\mathfrak p \not = 0$,
ce qui implique $I \subset \mathfrak p$ et, de plus,
$M_\mathfrak p \not = IM_\mathfrak p$, puisque la localisation est exacte.
Soit $f_1, \ldots, f_r \in I$ une suite $M$-régulière.
Alors $M_\mathfrak p/(f_1, \ldots, f_r)M_\mathfrak p$ est
non nul, car $(f_1, \ldots, f_r) \subset I$. La localisation étant
plate, les images de $f_1, \ldots, f_r$ forment une
suite $M_\mathfrak p$-régulière dans $I_\mathfrak p$. Comme cela vaut
pour toute suite $M$-régulière dans $I$, on conclut que
$\text{depth}_I(M) \leq \text{depth}_{I_\mathfrak p}(M_\mathfrak p)$.
Ce dernier membre est $\leq \text{depth}(M_\mathfrak p)$, qui est
$< \infty$ d'après le Lemme \ref{lemma-bound-depth}.
\end{proof}

\begin{lemma}
\label{lemma-depth-ext}
Soit $R$ un anneau local noethérien d'idéal maximal $\mathfrak m$.
Soit $M$ un $R$-module de type fini non nul. Alors $\text{depth}(M)$
est égale au plus petit entier $i$ tel que
$\Ext^i_R(R/\mathfrak m, M)$ soit non nul.
\end{lemma}

\begin{proof}
Notons $\delta(M)$ la profondeur de $M$ et $i(M)$ le
plus petit entier $i$ tel que $\Ext^i_R(R/\mathfrak m, M)$
soit non nul. Nous verrons dans un instant que $i(M) < \infty$.
D'après le Lemme \ref{lemma-ideal-nonzerodivisor}, on a
$\delta(M) = 0$ si et seulement si $i(M) = 0$, car
$\mathfrak m \in \text{Ass}(M)$ signifie exactement
que $i(M) = 0$. Ainsi, si $\delta(M)$ ou $i(M)$ est $> 0$, on peut
choisir $x \in \mathfrak m$ tel que (a) $x$ soit un non-diviseur de zéro
sur $M$, et (b) $\text{depth}(M/xM) = \delta(M) - 1$.
Considérons la suite exacte longue
de groupes Ext associée à la suite exacte courte
$0 \to M \to M \to M/xM \to 0$ par le Lemme \ref{lemma-long-exact-seq-ext} :
$$
\begin{matrix}
0
\to \Hom_R(\kappa, M)
\to \Hom_R(\kappa, M)
\to \Hom_R(\kappa, M/xM)
\\
\phantom{0\ }
\to \Ext^1_R(\kappa, M)
\to \Ext^1_R(\kappa, M)
\to \Ext^1_R(\kappa, M/xM)
\to \ldots
\end{matrix}
$$
Puisque $x \in \mathfrak m$, toutes les applications $\Ext^i_R(\kappa, M)
\to \Ext^i_R(\kappa, M)$ sont nulles ; voir le
Lemme \ref{lemma-annihilate-ext}.
Il est donc clair que $i(M/xM) = i(M) - 1$. Une récurrence sur
$\delta(M)$ achève la démonstration.
\end{proof}

\begin{lemma}
\label{lemma-depth-in-ses}
Soit $R$ un anneau local noethérien. Soit $0 \to N' \to N \to N'' \to 0$
une suite exacte courte de $R$-modules non nuls de type fini.
\begin{enumerate}
\item
$\text{depth}(N) \geq \min\{\text{depth}(N'), \text{depth}(N'')\}$
\item
$\text{depth}(N'') \geq \min\{\text{depth}(N), \text{depth}(N') - 1\}$
\item
$\text{depth}(N') \geq \min\{\text{depth}(N), \text{depth}(N'') + 1\}$
\end{enumerate}
\end{lemma}

\begin{proof}
Utilisons la caractérisation de la profondeur par les groupes Ext
$\Ext^i(\kappa, N)$, voir le Lemme \ref{lemma-depth-ext},
ainsi que la suite exacte longue de cohomologie
$$
\begin{matrix}
0
\to \Hom_R(\kappa, N')
\to \Hom_R(\kappa, N)
\to \Hom_R(\kappa, N'')
\\
\phantom{0\ }
\to \Ext^1_R(\kappa, N')
\to \Ext^1_R(\kappa, N)
\to \Ext^1_R(\kappa, N'')
\to \ldots
\end{matrix}
$$
du Lemme \ref{lemma-long-exact-seq-ext}.
\end{proof}

\begin{lemma}
\label{lemma-depth-drops-by-one}
Soit $R$ un anneau local noethérien et $M$ un $R$-module de type fini non nul.
\begin{enumerate}
\item Si $x \in \mathfrak m$ est un non-diviseur de zéro sur $M$, alors
$\text{depth}(M/xM) = \text{depth}(M) - 1$.
\item Toute suite $M$-régulière $x_1, \ldots, x_r$ peut être prolongée en une
suite $M$-régulière de longueur $\text{depth}(M)$.
\end{enumerate}
\end{lemma}

\begin{proof}
La partie (2) est une conséquence formelle de la partie (1). Soit $x \in R$ comme en (1).
La suite exacte courte $0 \to M \to M \to M/xM \to 0$
et le Lemme \ref{lemma-depth-in-ses} montrent que la profondeur diminue d'au plus 1.
D'autre part, si $x_1, \ldots, x_r \in \mathfrak m$
est une suite régulière pour $M/xM$, alors $x, x_1, \ldots, x_r$
est une suite régulière pour $M$. On voit donc que la profondeur diminue
d'au moins 1.
\end{proof}

\begin{lemma}
\label{lemma-inherit-minimal-primes}
Soit $(R, \mathfrak m)$ un anneau local noethérien et $M$ un $R$-module de type fini.
Soient $x \in \mathfrak m$, $\mathfrak p \in \text{Ass}(M)$, et $\mathfrak q$
minimal au-dessus de $\mathfrak p + (x)$. Alors $\mathfrak q \in \text{Ass}(M/x^nM)$
pour un certain $n \geq 1$.
\end{lemma}

\begin{proof}
Choisissons un sous-module $N \subset M$ tel que $N \cong R/\mathfrak p$.
D'après le lemme d'Artin-Rees (Lemme \ref{lemma-Artin-Rees}),
on peut choisir $n > 0$ tel que $N \cap x^nM \subset xN$.
Soit $\overline{N} \subset M/x^nM$ l'image de $N \to M \to M/x^nM$.
D'après le Lemme \ref{lemma-ass}, il suffit de montrer que
$\mathfrak q \in \text{Ass}(\overline{N})$.
Par notre choix de $n$, il existe une surjection
$\overline{N} \to N/xN = R/\mathfrak p + (x)$,
et donc $\mathfrak q$ appartient au support de $\overline{N}$.
Puisque $\overline{N}$ est annulé par $x^n$ et $\mathfrak p$, on voit que
$\mathfrak q$ est minimal parmi les idéaux premiers du support de $\overline{N}$.
Ainsi $\mathfrak q$ est un idéal premier associé à $\overline{N}$ d'après le
Lemme \ref{lemma-ass-minimal-prime-support}.
\end{proof}

\begin{lemma}
\label{lemma-depth-dim-associated-primes}
Soit $(R, \mathfrak m)$ un anneau local noethérien et $M$ un $R$-module de type fini.
Pour $\mathfrak p \in \text{Ass}(M)$, on a
$\dim(R/\mathfrak p) \geq \text{depth}(M)$.
\end{lemma}

\begin{proof}
Si $\mathfrak m \in \text{Ass}(M)$, il existe un élément non nul
$x \in M$ annulé par tous les éléments de $\mathfrak m$.
Ainsi $\text{depth}(M) = 0$. En particulier, le lemme est vrai dans ce cas.

\medskip\noindent
Si $\text{depth}(M) = 1$, alors le premier paragraphe
montre que $\mathfrak m \not \in \text{Ass}(M)$.
Ainsi $\dim(R/\mathfrak p) \geq 1$ pour tout $\mathfrak p \in \text{Ass}(M)$,
et le lemme est encore vrai dans ce cas.

\medskip\noindent
Démontrons le lemme en général par récurrence sur $\text{depth}(M)$,
que l'on peut et va supposer $> 1$. Choisissons $x \in \mathfrak m$ qui soit
un non-diviseur de zéro sur $M$. Remarquons que $x \not \in \mathfrak p$
(Lemme \ref{lemma-ass-zero-divisors}).
D'après le Lemme \ref{lemma-one-equation}, on a
$\dim(R/\mathfrak p + (x)) = \dim(R/\mathfrak p) - 1$.
Il existe donc un idéal premier $\mathfrak q$ minimal au-dessus de $\mathfrak p + (x)$ tel que
$\dim(R/\mathfrak q) = \dim(R/\mathfrak p) - 1$ (petit argument omis ;
indication : la dimension d'un anneau local noethérien $A$ est le maximum
des dimensions de $A/\mathfrak r$ lorsque $\mathfrak r$ parcourt les idéaux premiers
minimaux de $A$). Choisissons $n$ comme dans le
Lemme \ref{lemma-inherit-minimal-primes}, de sorte que
$\mathfrak q$ soit un idéal premier associé à $M/x^nM$.
On peut appliquer l'hypothèse de récurrence à $M/x^nM$ et $\mathfrak q$,
car $\text{depth}(M/x^nM) = \text{depth}(M) - 1$ d'après le
Lemme \ref{lemma-depth-drops-by-one}. On obtient
$\dim(R/\mathfrak q) \geq \text{depth}(M/x^nM)$, ce qui conclut.
\end{proof}

\begin{lemma}
\label{lemma-depth-localization}
Soit $R$ un anneau local noethérien et $M$ un $R$-module de type fini.
Pour un idéal premier $\mathfrak p \subset R$, on a
$\text{depth}(M_\mathfrak p) + \dim(R/\mathfrak p) \geq \text{depth}(M)$.
\end{lemma}

\begin{proof}
Si $M_\mathfrak p = 0$, alors $\text{depth}(M_\mathfrak p) = \infty$, et
le lemme est établi.
Si $\text{depth}(M) \leq \dim(R/\mathfrak p)$, alors le lemme est vrai.
Si $\text{depth}(M) > \dim(R/\mathfrak p)$, alors $\mathfrak p$ n'est contenu
dans aucun idéal premier associé $\mathfrak q$ à $M$, d'après le
Lemme \ref{lemma-depth-dim-associated-primes}.
On peut donc trouver $x \in \mathfrak p$ qui n'appartient à aucun
idéal premier associé à $M$, d'après le Lemme \ref{lemma-silly} et le
Lemme \ref{lemma-finite-ass}. Alors $x$ est un non-diviseur de zéro
sur $M$ ; voir le Lemme \ref{lemma-ass-zero-divisors}.
Ainsi $\text{depth}(M/xM) = \text{depth}(M) - 1$ et
$\text{depth}(M_\mathfrak p / x M_\mathfrak p) =
\text{depth}(M_\mathfrak p) - 1$, pourvu que $M_\mathfrak p$ soit non nul ;
voir le Lemme \ref{lemma-depth-drops-by-one}.
On conclut donc par récurrence sur $\text{depth}(M)$.
\end{proof}

\begin{lemma}
\label{lemma-depth-goes-down-finite}
Soit $(R, \mathfrak m)$ un anneau local noethérien. Soit $R \to S$
un homomorphisme fini d'anneaux. Soient $\mathfrak m_1, \ldots, \mathfrak m_n$
les idéaux maximaux de $S$. Soit $N$ un $S$-module de type fini.
Alors
$$
\min\nolimits_{i = 1, \ldots, n} \text{depth}(N_{\mathfrak m_i}) =
\text{depth}_\mathfrak m(N)
$$
\end{lemma}

\begin{proof}
D'après les Lemmes \ref{lemma-integral-no-inclusion}, \ref{lemma-integral-going-up},
et le Lemme \ref{lemma-finite-finite-fibres}, les idéaux maximaux de
$S$ sont exactement les idéaux premiers de $S$ au-dessus de $\mathfrak m$, et
ils sont en nombre fini. L'énoncé du lemme a donc
un sens. Nous allons démontrer le lemme par récurrence sur
$k = \min\nolimits_{i = 1, \ldots, n} \text{depth}(N_{\mathfrak m_i})$.
Si $k = 0$, alors $\text{depth}(N_{\mathfrak m_i}) = 0$ pour un certain $i$.
D'après le Lemme \ref{lemma-depth-ext}, cela signifie que
$\mathfrak m_i S_{\mathfrak m_i}$ est un idéal premier associé
à $N_{\mathfrak m_i}$, et donc que $\mathfrak m_i$ est un
idéal premier associé à $N$ (Lemme \ref{lemma-localize-ass}).
D'après le Lemme \ref{lemma-ass-functorial-Noetherian}, on voit que
$\mathfrak m$ est un idéal premier associé à $N$ comme $R$-module.
Ainsi $\text{depth}_\mathfrak m(N) = 0$. Cela établit le cas initial.
Si $k > 0$, alors $\mathfrak m_i \not \in \text{Ass}_S(N)$.
Ainsi $\mathfrak m \not \in \text{Ass}_R(N)$, encore d'après le
Lemme \ref{lemma-ass-functorial-Noetherian}.
On peut donc trouver $f \in \mathfrak m$ qui n'est pas un diviseur de zéro sur
$N$ ; voir le Lemme \ref{lemma-ideal-nonzerodivisor}. D'après le
Lemme \ref{lemma-depth-drops-by-one},
toutes les profondeurs diminuent exactement de $1$ lorsque l'on passe de $N$ à
$N/fN$, et l'hypothèse de récurrence conclut.
\end{proof}




\section{Fonctorialités de Ext}
\label{section-functoriality-ext}

\noindent
Dans cette section, nous examinons brièvement la fonctorialité
de $\Ext$ par rapport au changement d'anneau, etc.
Voici une liste de points à développer.
\begin{enumerate}
\item Étant donnés $R \to R'$, un $R$-module
$M$ et un $R'$-module $N'$,
le $R$-module $\Ext^i_R(M, N')$
possède une structure naturelle de $R'$-module. De plus, il existe un
homomorphisme canonique $R'$-linéaire $\Ext^i_{R'}(M \otimes_R R', N') \to
\Ext^i_R(M, N')$.
\item Étant donnés $R \to R'$ et des $R$-modules $M$, $N$, il existe un
homomorphisme naturel de $R$-modules
$\Ext^i_R(M, N) \to \text{Ext}^i_R(M, N \otimes_R R')$.
\end{enumerate}

\begin{lemma}
\label{lemma-flat-base-change-ext}
Étant donnés un homomorphisme plat d'anneaux $R \to R'$, un $R$-module $M$ et un
$R'$-module $N'$, l'homomorphisme naturel
$$
\Ext^i_{R'}(M \otimes_R R', N') \to \text{Ext}^i_R(M, N')
$$
est un isomorphisme pour $i \geq 0$.
\end{lemma}

\begin{proof}
Choisissons une résolution libre $F_\bullet$ de $M$.
Comme $R \to R'$ est plat, $F_\bullet \otimes_R R'$ est
une résolution libre de $M \otimes_R R'$ sur $R'$.
L'énoncé affirme que l'homomorphisme
$$
\Hom_{R'}(F_\bullet \otimes_R R', N') \to
\Hom_R(F_\bullet, N')
$$
induit un isomorphisme sur les groupes d'homologie, ce qui est vrai puisqu'il
est un isomorphisme de complexes d'après le
Lemme \ref{lemma-adjoint-tensor-restrict}.
\end{proof}






\section{Une application des groupes Ext}
\label{section-ext-application}

\noindent
La voici.

\begin{lemma}
\label{lemma-split-injection-after-completion}
Soit $R$ un anneau noethérien. Soit $I \subset R$ un idéal
contenu dans le radical de Jacobson de $R$.
Soit $N \to M$ un homomorphisme de $R$-modules de type fini.
Supposons qu'il existe des entiers $n$ arbitrairement grands tels que
$N/I^nN \to M/I^nM$ soit une injection scindée.
Alors $N \to M$ est une injection scindée.
\end{lemma}

\begin{proof}
Supposons que $\varphi : N \to M$ satisfasse les hypothèses du lemme.
Remarquons que cela implique $\Ker(\varphi) \subset I^nN$
pour des entiers $n$ arbitrairement grands. Ainsi, d'après le
Lemme \ref{lemma-intersection-powers-ideal-module}, on voit que $\varphi$
est injectif. Posons $Q = M/N$, de sorte que l'on ait une suite exacte courte
$$
0 \to N \to M \to Q \to 0.
$$
Soit
$$
F_2 \xrightarrow{d_2} F_1 \xrightarrow{d_1} F_0 \to Q \to 0
$$
une résolution libre de type fini de $Q$. On peut choisir une application
$\alpha : F_0 \to M$ qui relève l'application $F_0 \to Q$. Elle induit une application
$\beta : F_1 \to N$ telle que $\beta \circ d_2 = 0$. L'extension
ci-dessus est scindée si et seulement s'il existe une application $\gamma : F_0 \to N$
telle que $\beta = \gamma \circ d_1$. Autrement dit, la classe de
$\beta$ dans $\Ext^1_R(Q, N)$ est l'obstruction au scindage
de la suite exacte courte ci-dessus.

\medskip\noindent
Supposons que $n$ soit un entier assez grand tel que $N/I^nN \to M/I^nM$ soit une
injection scindée. Cela implique que
$$
0 \to N/I^nN \to M/I^nM \to Q/I^nQ \to 0.
$$
est encore une suite exacte courte. De plus, la suite
$$
F_1/I^nF_1 \xrightarrow{d_1} F_0/I^nF_0 \to Q/I^nQ \to 0
$$
est encore exacte. En raisonnant comme ci-dessus, on voit que l'application
$\overline{\beta} : F_1/I^nF_1 \to N/I^nN$
induite par $\beta$ est égale à $\overline{\gamma_n} \circ d_1$ pour une certaine
application $\overline{\gamma_n} : F_0/I^nF_0 \to N/I^nN$.
Comme $F_0$ est libre, on peut relever $\overline{\gamma_n}$ en une application
$\gamma_n : F_0 \to N$, et l'on voit alors que
$\beta - \gamma_n \circ d_1$ est une application de $F_1$ dans $I^nN$.
Autrement dit, on conclut que
$$
\beta \in
\Im\Big(\Hom_R(F_0, N) \to \Hom_R(F_1, N)\Big) + I^n\Hom_R(F_1, N).
$$
pour cet entier $n$.

\medskip\noindent
Puisque, par hypothèse, cette propriété vaut pour des entiers $n$ arbitrairement grands,
on conclut que l'image de $\beta$ dans le conoyau de
$\Hom_R(F_0, N) \to \Hom_R(F_1, N)$ est nulle d'après le
Lemme \ref{lemma-intersection-powers-ideal-module}. Ainsi
$\beta$ appartient à l'image de l'application $\Hom_R(F_0, N) \to \Hom_R(F_1, N)$,
comme souhaité.
\end{proof}











\section{Groupes Tor et platitude}
\label{section-tor}

\noindent
Dans cette section, nous utilisons une partie de l'algèbre homologique
développée dans la section précédente pour expliquer ce que sont
les groupes Tor. Supposons donc que $R$ soit un anneau
et que $M$, $N$ soient deux $R$-modules. Choisissons
une résolution $F_\bullet$ de $M$ par des $R$-modules libres.
Voir le Lemme \ref{lemma-resolution-by-finite-free}.
Considérons le complexe homologique
$$
F_\bullet \otimes_R N
:
\ldots
\to F_2 \otimes_R N
\to F_1 \otimes_R N
\to F_0 \otimes_R N
$$
On définit $\text{Tor}^R_i(M, N)$ comme le $i$-ième groupe d'homologie
de ce complexe. Le lemme suivant précise en
quel sens cette définition est bien définie.

\begin{lemma}
\label{lemma-tor-welldefined}
Soit $R$ un anneau. Soient $M_1, M_2, N$ des $R$-modules.
Supposons que $F_\bullet$ soit une résolution libre du
module $M_1$ et que $G_\bullet$ soit une résolution libre du
module $M_2$. Soit $\varphi : M_1 \to M_2$
un homomorphisme de modules. Soit $\alpha : F_\bullet \to G_\bullet$
un morphisme de complexes qui induit $\varphi$ sur
$M_1 = \Coker(d_{F, 1}) \to M_2 = \Coker(d_{G, 1})$,
voir le Lemme \ref{lemma-compare-resolutions}.
Alors les applications induites
$$
H_i(\alpha) :
H_i(F_\bullet \otimes_R N)
\longrightarrow
H_i(G_\bullet \otimes_R N)
$$
ne dépendent pas du choix de $\alpha$. Si $\varphi$
est un isomorphisme, toutes les applications $H_i(\alpha)$ le sont aussi.
Si $M_1 = M_2$, $F_\bullet = G_\bullet$, et
si $\varphi$ est l'identité, toutes les applications $H_i(\alpha)$ le sont aussi.
\end{lemma}

\begin{proof}
La démonstration de ce lemme est identique à celle du Lemme
\ref{lemma-ext-welldefined}.
\end{proof}

\noindent
Ce lemme implique non seulement que les modules Tor sont bien définis,
mais il fournit aussi la fonctorialité des constructions
$(M, N) \mapsto \text{Tor}_i^R(M, N)$ en la première variable. Bien entendu,
la fonctorialité en la seconde variable est évidente. Nous laissons au
lecteur le soin de vérifier que chacun des $\text{Tor}_i^R$ est en fait
un foncteur
$$
\text{Mod}_R \times \text{Mod}_R \to \text{Mod}_R.
$$
Ici $\text{Mod}_R$ désigne la catégorie des $R$-modules et,
pour la définition de la catégorie produit,
voir Catégories, Définition \ref{categories-definition-product-category}.
Plus précisément, étant donnés des morphismes de $R$-modules $M_1 \to M_2$
et $N_1 \to N_2$, on obtient un diagramme commutatif
$$
\xymatrix{
\text{Tor}_i^R(M_1, N_1) \ar[r] \ar[d] &
\text{Tor}_i^R(M_1, N_2) \ar[d] \\
\text{Tor}_i^R(M_2, N_1) \ar[r] &
\text{Tor}_i^R(M_2, N_2) \\
}
$$

\begin{lemma}
\label{lemma-long-exact-sequence-tor}
Soit $R$ un anneau et soit $M$ un $R$-module.
Supposons que $0 \to N' \to N \to N'' \to 0$ soit une suite
exacte courte de $R$-modules. Il existe une suite exacte
longue
$$
\text{Tor}_1^R(M, N')
\to \text{Tor}_1^R(M, N)
\to \text{Tor}_1^R(M, N'')
\to
M \otimes_R N'
\to M \otimes_R N
\to M \otimes_R N''
\to 0
$$
\end{lemma}

\begin{proof}
La démonstration est la même que celle du
Lemme \ref{lemma-long-exact-seq-ext}.
\end{proof}

\noindent
Considérons un bicomplexe homologique de $R$-modules
$$
\xymatrix{
\ldots \ar[r]^d &
A_{2, 0} \ar[r]^d &
A_{1, 0} \ar[r]^d &
A_{0, 0} \\
\ldots \ar[r]^d &
A_{2, 1} \ar[r]^d \ar[u]^\delta &
A_{1, 1} \ar[r]^d \ar[u]^\delta &
A_{0, 1} \ar[u]^\delta \\
\ldots \ar[r]^d &
A_{2, 2} \ar[r]^d \ar[u]^\delta &
A_{1, 2} \ar[r]^d \ar[u]^\delta &
A_{0, 2} \ar[u]^\delta \\
&
\ldots \ar[u]^\delta &
\ldots \ar[u]^\delta &
\ldots \ar[u]^\delta \\
}
$$
Cela signifie que $d_{i, j} : A_{i, j} \to A_{i-1, j}$
et $\delta_{i, j} : A_{i, j} \to A_{i, j-1}$ possèdent les propriétés
suivantes :
\begin{enumerate}
\item Toute composée de deux $d_{i, j}$ est nulle. Autrement dit,
les lignes du bicomplexe sont des complexes.
\item Toute composée de deux $\delta_{i, j}$ est nulle. Autrement dit,
les colonnes du bicomplexe sont des complexes.
\item Pour toute paire $(i, j)$, on a $\delta_{i-1, j} \circ d_{i, j}
= d_{i, j-1} \circ \delta_{i, j}$. Autrement dit, tous les carrés
sont commutatifs.
\end{enumerate}
La bonne construction consiste à associer une suite spectrale à
chacun de ces bicomplexes. Cependant, pour le moment, nous pouvons nous contenter
d'une construction légèrement plus simple.

\medskip\noindent
Plus précisément, pour les seuls besoins de cette section, étant donné un bicomplexe
$(A_{\bullet, \bullet}, d, \delta)$, posons
$R(A)_j = \Coker(A_{1, j} \to A_{0, j})$ et
$U(A)_i = \Coker(A_{i, 1} \to A_{i, 0})$. (Les lettres
$R$ et $U$ évoquent les mots anglais pour « droite » et « haut ».)
On munit $R(A)_\bullet$ d'une structure de complexe
au moyen des applications $\delta$. De même, on munit $U(A)_\bullet$
d'une structure de complexe au moyen des applications $d$.
Autrement dit, on obtient l'immense diagramme commutatif suivant
$$
\xymatrix{
\ldots \ar[r]^d &
U(A)_2 \ar[r]^d &
U(A)_1 \ar[r]^d &
U(A)_0 &
\\
\ldots \ar[r]^d &
A_{2, 0} \ar[r]^d \ar[u] &
A_{1, 0} \ar[r]^d \ar[u] &
A_{0, 0} \ar[r] \ar[u] &
R(A)_0 \\
\ldots \ar[r]^d &
A_{2, 1} \ar[r]^d \ar[u]^\delta &
A_{1, 1} \ar[r]^d \ar[u]^\delta &
A_{0, 1} \ar[r] \ar[u]^\delta &
R(A)_1 \ar[u]^\delta \\
\ldots \ar[r]^d &
A_{2, 2} \ar[r]^d \ar[u]^\delta &
A_{1, 2} \ar[r]^d \ar[u]^\delta &
A_{0, 2} \ar[r] \ar[u]^\delta &
R(A)_2 \ar[u]^\delta \\
&
\ldots \ar[u]^\delta &
\ldots \ar[u]^\delta &
\ldots \ar[u]^\delta &
\ldots \ar[u]^\delta \\
}
$$
(Bien entendu, ce n'est plus un bicomplexe.)
La notion de morphisme $\Phi : (A_{\bullet, \bullet}, d, \delta)
\to (B_{\bullet, \bullet}, d, \delta)$ de bicomplexes
est claire, et il est clair qu'un tel morphisme induit des morphismes de complexes
$R(\Phi) : R(A)_\bullet \to R(B)_\bullet$ et
$U(\Phi) : U(A)_\bullet \to U(B)_\bullet$.

\begin{lemma}
\label{lemma-no-spectral-sequence}
Soit $(A_{\bullet, \bullet}, d, \delta)$ un bicomplexe tel
que
\begin{enumerate}
\item Chaque ligne $A_{\bullet, j}$ soit une résolution de $R(A)_j$.
\item Chaque colonne $A_{i, \bullet}$ soit une résolution de $U(A)_i$.
\end{enumerate}
Il existe alors des isomorphismes canoniques
$$
H_i(R(A)_\bullet)
\cong
H_i(U(A)_\bullet).
$$
Ces isomorphismes sont fonctoriels par rapport aux morphismes
de bicomplexes possédant les propriétés ci-dessus.
\end{lemma}

\begin{proof}
Nous allons montrer que $H_i(R(A)_\bullet))$
et $H_i(U(A)_\bullet)$ sont canoniquement
isomorphes à un troisième groupe, à savoir
$$
\mathbf{H}_i(A) :=
\frac{
\{
(a_{i, 0}, a_{i-1, 1}, \ldots, a_{0, i})
\mid
d(a_{i, 0}) = \delta(a_{i-1, 1}), \ldots,
d(a_{1, i-1}) = \delta(a_{0, i})
\}}
{
\{
d(a_{i + 1, 0}) + \delta(a_{i, 1}),
d(a_{i, 1}) + \delta(a_{i-1, 2}),
\ldots,
d(a_{1, i}) + \delta(a_{0, i + 1})
\}
}
$$
Nous adoptons ici la convention de notation selon laquelle $a_{i, j}$ désigne
un élément de $A_{i, j}$. Autrement dit, un élément de $\mathbf{H}_i$
est représenté par un zigzag, figuré comme suit lorsque $i = 2$
$$
\xymatrix{
a_{2, 0} \ar@{|->}[r] & d(a_{2, 0}) = \delta(a_{1, 1}) & \\
& a_{1, 1} \ar@{|->}[u] \ar@{|->}[r] & d(a_{1, 1}) = \delta(a_{0, 2}) \\
& & a_{0, 2} \ar@{|->}[u] \\
}
$$
Naturellement, on quotient par les zigzags « triviaux », c'est-à-dire par le sous-module
engendré par les éléments de la forme $(0, \ldots, 0, -\delta(a_{t + 1, t-i}),
d(a_{t + 1, t-i}), 0, \ldots, 0)$. Remarquons qu'il existe des
homomorphismes canoniques
$$
\mathbf{H}_i(A) \to H_i(R(A)_\bullet), \quad
(a_{i, 0}, a_{i-1, 1}, \ldots, a_{0, i}) \mapsto
\text{classe de l'image de }a_{0, i}
$$
et
$$
\mathbf{H}_i(A) \to H_i(U(A)_\bullet), \quad
(a_{i, 0}, a_{i-1, 1}, \ldots, a_{0, i}) \mapsto
\text{classe de l'image de }a_{i, 0}
$$

\medskip\noindent
Montrons d'abord que ces homomorphismes sont surjectifs.
Supposons que $\overline{r} \in H_i(R(A)_\bullet)$.
Soit $r \in R(A)_i$ un cycle qui représente la
classe de $\overline{r}$.
Soit $a_{0, i} \in A_{0, i}$ un élément dont
l'image est $r$. Puisque $\delta(r) = 0$,
on voit que $\delta(a_{0, i})$ appartient à
l'image de $d$. Il existe donc un élément
$a_{1, i-1} \in A_{1, i-1}$ tel que
$d(a_{1, i-1}) = \delta(a_{0, i})$. Cela implique à son tour
que $\delta(a_{1, i-1})$ appartient au noyau
de $d$ (car $d(\delta(a_{1, i-1})) = \delta(d(a_{1, i-1}))
= \delta(\delta(a_{0, i})) = 0$). Par exactitude des
lignes, on trouve un élément $a_{2, i-2}$ tel que
$d(a_{2, i-2}) = \delta(a_{1, i-1})$. Et l'on poursuit ainsi
jusqu'à obtenir un zigzag complet. Bien entendu, la surjectivité
de $\mathbf{H}_i \to H_i(U(A))$ se démontre de la même manière.

\medskip\noindent
Pour démontrer l'injectivité, on raisonne exactement de la même manière.
Supposons donc que l'on ait un zigzag
$(a_{i, 0}, a_{i-1, 1}, \ldots, a_{0, i})$
dont l'image dans $H_i(R(A)_\bullet)$ est nulle.
Cela signifie que l'image de $a_{0, i}$ est un élément
de $\Coker(A_{i, 1} \to A_{i, 0})$
qui appartient à l'image de
$\delta : \Coker(A_{i + 1, 1} \to A_{i + 1, 0}) \to
\Coker(A_{i, 1} \to A_{i, 0})$.
Autrement dit, $a_{0, i}$ appartient à l'image de
$\delta \oplus d : A_{0, i + 1} \oplus A_{1, i} \to A_{0, i}$.
La définition des zigzags triviaux montre que
l'on peut modifier notre zigzag par un zigzag trivial et
supposer que $a_{0, i} = 0$. Cela implique immédiatement
que $d(a_{1, i-1}) = 0$. Comme les lignes
sont exactes, il s'ensuit que $a_{1, i-1}$ appartient
à l'image de $d : A_{2, i-1} \to A_{1, i-1}$.
On peut donc modifier de nouveau notre zigzag par un
zigzag trivial et supposer qu'il est de la forme
$(a_{i, 0}, a_{i-1, 1}, \ldots, a_{2, i-2}, 0, 0)$.
En poursuivant ainsi, on obtient l'injectivité voulue.

\medskip\noindent
Si $\Phi : (A_{\bullet, \bullet}, d, \delta)
\to (B_{\bullet, \bullet}, d, \delta)$ est un morphisme
de bicomplexes satisfaisant tous deux aux conditions
du lemme, on obtient clairement un diagramme
commutatif
$$
\xymatrix{
H_i(U(A)_\bullet) \ar[d] &
\mathbf{H}_i(A) \ar[r] \ar[l] \ar[d] &
H_i(R(A)_\bullet) \ar[d] \\
H_i(U(B)_\bullet) &
\mathbf{H}_i(B) \ar[r] \ar[l] &
H_i(R(B)_\bullet) \\
}
$$
Cela démontre la fonctorialité.
\end{proof}

\begin{remark}
\label{remark-signs-double-complex}
L'isomorphisme construit ci-dessus n'est l'isomorphisme « correct » qu'aux signes près.
Une bonne partie de l'algèbre homologique consiste à choisir les signes de
diverses applications et à démontrer la commutativité de diagrammes en faisant intervenir
des signes appropriés. Pour le moment, nous utiliserons simplement l'isomorphisme
donné dans la démonstration ci-dessus et nous nous préoccuperons des signes plus tard.
\end{remark}

\begin{lemma}
\label{lemma-tor-left-right}
Soit $R$ un anneau. Pour tout $i \geq 0$, les foncteurs
$\text{Mod}_R \times \text{Mod}_R \to \text{Mod}_R$,
$(M, N) \mapsto \text{Tor}_i^R(M, N)$ et
$(M, N) \mapsto \text{Tor}_i^R(N, M)$ sont
canoniquement isomorphes.
\end{lemma}

\begin{proof}
Soit $F_\bullet$ une résolution libre du module $M$ et
soit $G_\bullet$ une résolution libre du module $N$.
Considérons le bicomplexe $(A_{i, j}, d, \delta)$ défini
comme suit :
\begin{enumerate}
\item posons $A_{i, j} = F_i \otimes_R G_j$,
\item définissons $d_{i, j} : F_i \otimes_R G_j \to F_{i-1} \otimes G_j$
comme $d_{F, i} \otimes \text{id}$, et
\item définissons $\delta_{i, j} : F_i \otimes_R G_j \to F_i \otimes G_{j-1}$
comme $\text{id} \otimes d_{G, j}$.
\end{enumerate}
Ce bicomplexe est généralement noté simplement $F_\bullet \otimes_R G_\bullet$.

\medskip\noindent
Puisque chaque $G_j$ est libre, donc plat, chaque
ligne du bicomplexe est exacte sauf en degré homologique
$0$. Puisque chaque $F_i$ est libre, donc plat, chaque
colonne du bicomplexe est exacte sauf en degré homologique
$0$. Le bicomplexe satisfait donc aux conditions
du Lemme \ref{lemma-no-spectral-sequence}.

\medskip\noindent
Pour comprendre ce que dit le lemme, calculons $R(A)_\bullet$ et $U(A)_\bullet$.
On a
\begin{eqnarray*}
R(A)_i & = & \Coker(A_{1, i} \to A_{0, i}) \\
& = & \Coker(F_1 \otimes_R G_i \to F_0 \otimes_R G_i) \\
& = & \Coker(F_1 \to F_0) \otimes_R G_i \\
& = & M \otimes_R G_i
\end{eqnarray*}
En fait, ces isomorphismes sont compatibles avec les différentielles
$\delta$, et l'on voit que $R(A)_\bullet = M \otimes_R G_\bullet$
comme complexes homologiques. De manière tout à fait analogue, on voit que
$U(A)_\bullet = F_\bullet \otimes_R N$. On obtient
\begin{eqnarray*}
\text{Tor}_i^R(M, N)
& = & H_i(F_\bullet \otimes_R N) \\
& = & H_i(U(A)_\bullet) \\
& = & H_i(R(A)_\bullet) \\
& = & H_i(M \otimes_R G_\bullet) \\
& = & H_i(G_\bullet \otimes_R M) \\
& = & \text{Tor}_i^R(N, M)
\end{eqnarray*}
Ici, la troisième égalité est le Lemme \ref{lemma-no-spectral-sequence}, et
la cinquième utilise l'isomorphisme $V \otimes W = W \otimes V$
du produit tensoriel.

\medskip\noindent
Fonctorialité. Supposons que l'on ait des $R$-modules $M_\nu$, $N_\nu$,
$\nu = 1, 2$. Soient $\varphi : M_1 \to M_2$ et $\psi : N_1 \to N_2$
des morphismes de $R$-modules.
Supposons que l'on ait des résolutions libres $F_{\nu, \bullet}$
de $M_\nu$ et des résolutions libres $G_{\nu, \bullet}$ de $N_\nu$.
D'après le Lemme \ref{lemma-compare-resolutions}, on peut choisir
des morphismes de complexes $\alpha : F_{1, \bullet} \to F_{2, \bullet}$
et $\beta : G_{1, \bullet} \to G_{2, \bullet}$ compatibles
avec $\varphi$ et $\psi$. Nous affirmons que
la paire $(\alpha, \beta)$ induit un morphisme de
bicomplexes
$$
\alpha \otimes \beta :
F_{1, \bullet} \otimes_R G_{1, \bullet}
\longrightarrow
F_{2, \bullet} \otimes_R G_{2, \bullet}
$$
Il s'agit d'une vérification tout à fait directe : l'application
$F_{1, i} \otimes_R G_{1, j} \to F_{2, i} \otimes_R G_{2, j}$
est donnée par $\alpha_i \otimes \beta_j$, où $\alpha_i$,
resp.\ $\beta_j$, est la composante de degré $i$, resp.\ $j$, de $\alpha$,
resp.\ $\beta$. Le lecteur vérifie aussi aisément que l'application induite
$R(F_{1, \bullet} \otimes_R G_{1, \bullet})_\bullet
\to R(F_{2, \bullet} \otimes_R G_{2, \bullet})_\bullet$
coïncide avec l'application $M_1 \otimes_R G_{1, \bullet}
\to M_2 \otimes_R G_{2, \bullet}$ induite par $\varphi \otimes \beta$.
Il en va de même de l'application induite sur les complexes $U(-)_\bullet$.
L'assertion de fonctorialité résulte donc de l'assertion
de fonctorialité du Lemme \ref{lemma-no-spectral-sequence}.
\end{proof}

\begin{remark}
\label{remark-curiosity-signs-swap}
Un cas intéressant se présente lorsque $M = N$ ci-dessus.
On obtient alors une application canonique $\text{Tor}_i^R(M, M)
\to \text{Tor}_i^R(M, M)$. Remarquons que cette application n'est pas
l'identité, car même lorsque $i = 0$, elle n'est pas
l'identité ! Par exemple, si $V$ est un espace vectoriel de dimension
$n$ sur un corps, alors l'application d'échange $V \otimes_k V \to V \otimes_k V$
possède $(n^2 + n)/2$ valeurs propres égales à $+1$ et $(n^2-n)/2$ valeurs propres
égales à $-1$. En caractéristique $2$, elle n'est même pas diagonalisable.
Remarquons que même changer le signe de l'application ne suffira pas à supprimer
ce phénomène.
\end{remark}

\begin{lemma}
\label{lemma-tor-noetherian}
Soit $R$ un anneau noethérien. Soient $M$, $N$ des $R$-modules de type fini.
Alors $\text{Tor}_p^R(M, N)$ est un $R$-module de type fini pour tout $p$.
\end{lemma}

\begin{proof}
Cela résulte de ce que $\text{Tor}_p^R(M, N)$ se calcule comme les
groupes d'homologie d'un complexe $F_\bullet \otimes_R N$
où chaque $F_n$ est un $R$-module libre de type fini ; voir le
Lemme \ref{lemma-resolution-by-finite-free}.
\end{proof}

\begin{lemma}
\label{lemma-characterize-flat}
Soit $R$ un anneau. Soit $M$ un $R$-module.
Les assertions suivantes sont équivalentes :
\begin{enumerate}
\item Le module $M$ est plat sur $R$.
\item Pour tout $i > 0$, le foncteur $\text{Tor}_i^R(M, -)$ est nul.
\item Le foncteur $\text{Tor}_1^R(M, -)$ est nul.
\item Pour tout idéal $I \subset R$, on a $\text{Tor}_1^R(M, R/I) = 0$.
\item Pour tout idéal de type fini $I \subset R$, on a
$\text{Tor}_1^R(M, R/I) = 0$.
\end{enumerate}
\end{lemma}

\begin{proof}
Supposons $M$ plat. Soit $N$ un $R$-module.
Soit $F_\bullet$ une résolution libre de $N$.
Alors $F_\bullet \otimes_R M$ est une résolution de $N \otimes_R M$,
par platitude de $M$. Tous les groupes Tor supérieurs s'annulent donc.

\medskip\noindent
Il suffit maintenant de montrer que la dernière condition implique que
$M$ est plat. Soit $I \subset R$ un idéal.
Considérons la suite exacte courte
$0 \to I \to R \to R/I \to 0$. Appliquons le
Lemme \ref{lemma-long-exact-sequence-tor}. On obtient une
suite exacte
$$
\text{Tor}_1^R(M, R/I) \to
M \otimes_R I \to
M \otimes_R R \to
M \otimes_R R/I \to
0
$$
Comme, de toute évidence, $M \otimes_R R = M$, on conclut que la
dernière hypothèse implique que $M \otimes_R I \to M$ est
injectif pour tout idéal de type fini $I$.
Ainsi $M$ est plat d'après le Lemme \ref{lemma-flat}.
\end{proof}

\begin{remark}
\label{remark-Tor-ring-mod-ideal}
La démonstration du Lemme \ref{lemma-characterize-flat} montre en fait
que
$$
\text{Tor}_1^R(M, R/I)
=
\Ker(I \otimes_R M \to M).
$$
\end{remark}














\section{Fonctorialités de Tor}
\label{section-functoriality-tor}

\noindent
Dans cette section, nous examinons brièvement la fonctorialité
de $\text{Tor}$ par rapport au changement d'anneau, etc.
Voici une liste de points à développer.
\begin{enumerate}
\item Étant donnés un homomorphisme d'anneaux $R \to R'$, un $R$-module
$M$ et un $R'$-module $N'$,
les $R$-modules $\text{Tor}_i^R(M, N')$ possèdent
une structure naturelle de $R'$-module.
\item Étant donnés un homomorphisme d'anneaux $R \to R'$ et des $R$-modules
$M$, $N$, il existe un homomorphisme naturel de $R$-modules
$\text{Tor}_i^R(M, N) \to \text{Tor}_i^{R'}(M \otimes_R R', N \otimes_R R')$.
\item Étant donnés un homomorphisme d'anneaux $R \to R'$, un $R$-module $M$ et
un $R'$-module $N'$, il existe un homomorphisme naturel
de $R'$-modules
$\text{Tor}_i^R(M, N') \to \text{Tor}_i^{R'}(M \otimes_R R', N')$.
\end{enumerate}

\begin{lemma}
\label{lemma-flat-base-change-tor}
Étant donnés un homomorphisme plat d'anneaux $R \to R'$ et des $R$-modules
$M$, $N$, l'homomorphisme naturel de $R$-modules
$\text{Tor}_i^R(M, N)\otimes_R R'
\to \text{Tor}_i^{R'}(M \otimes_R R', N \otimes_R R')$
est un isomorphisme pour tout $i$.
\end{lemma}

\begin{proof}
Omis. Cela résulte de ce qu'une résolution libre $F_\bullet$ de $M$ sur
$R$ reste exacte après tensorisation par $R'$ sur $R$, et donc que
$(F_\bullet \otimes_R N)\otimes_R R'$ calcule les groupes Tor
sur $R'$.
\end{proof}

\noindent
Le lemme suivant ne semble trouver sa place nulle part ailleurs.

\begin{lemma}
\label{lemma-tor-commutes-filtered-colimits}
Soit $R$ un anneau. Soit $M = \colim M_i$ une colimite filtrante de
$R$-modules. Soit $N$ un $R$-module. Alors
$\text{Tor}_n^R(M, N) = \colim \text{Tor}_n^R(M_i, N)$ pour tout $n$.
\end{lemma}

\begin{proof}
Choisissons une résolution libre $F_\bullet$ de $N$. Alors
$F_\bullet \otimes_R M = \colim F_\bullet \otimes_R M_i$
comme complexes, d'après le Lemme \ref{lemma-tensor-products-commute-with-limits}.
Le résultat découle donc du Lemme \ref{lemma-directed-colimit-exact}.
\end{proof}








\section{Modules projectifs}
\label{section-projective}

\noindent
Quelques lemmes sur les modules projectifs.

\begin{definition}
\label{definition-projective}
Soit $R$ un anneau. Un $R$-module $P$ est {\it projectif} si et seulement si
le foncteur $\Hom_R(P, -) : \text{Mod}_R \to \text{Mod}_R$ est
un foncteur exact.
\end{definition}

\noindent
Le foncteur $\Hom_R(M, - )$ est exact à gauche pour tout $R$-module $M$ ; voir le
Lemme \ref{lemma-hom-exact}.
La condition que $P$ soit projectif signifie donc précisément que, pour toute
surjection de $R$-modules $N \to N'$, l'application
$\Hom_R(P, N) \to \Hom_R(P, N')$ est surjective.

\begin{lemma}
\label{lemma-characterize-projective}
Soit $R$ un anneau. Soit $P$ un $R$-module.
Les assertions suivantes sont équivalentes :
\begin{enumerate}
\item $P$ est projectif,
\item $P$ est un facteur direct d'un $R$-module libre, et
\item $\Ext^1_R(P, M) = 0$ pour tout $R$-module $M$.
\end{enumerate}
\end{lemma}

\begin{proof}
Supposons $P$ projectif. Choisissons une surjection $\pi : F \to P$, où $F$
est un $R$-module libre. Comme $P$ est projectif, il existe
$i \in \Hom_R(P, F)$ tel que $\pi \circ i = \text{id}_P$.
Autrement dit, $F \cong \Ker(\pi) \oplus i(P)$, et l'on voit
que $P$ est un facteur direct de $F$.

\medskip\noindent
Réciproquement, supposons que $P \oplus Q = F$ soit un $R$-module libre.
Remarquons que le module libre $F = \bigoplus_{i \in I} R$ est projectif,
car $\Hom_R(F, M) = \prod_{i \in I} M$ et le foncteur
$M \mapsto \prod_{i \in I} M$ est exact.
Alors $\Hom_R(F, -) = \Hom_R(P, -) \times \Hom_R(Q, -)$
comme foncteurs, et donc $P$ et $Q$ sont tous deux projectifs.

\medskip\noindent
Supposons que $P \oplus Q = F$ soit un $R$-module libre. On dispose alors d'une
résolution libre $F_\bullet$ de la forme
$$
\ldots F \xrightarrow{a} F \xrightarrow{b} F \to P \to 0
$$
où les applications $a, b$ alternent et sont égales aux projecteurs sur
$P$ et $Q$. Le complexe $\Hom_R(F_\bullet, M)$ est donc exact scindé
en degrés $\geq 1$, d'où l'annulation de (3).

\medskip\noindent
Supposons que $\Ext^1_R(P, M) = 0$ pour tout $R$-module $M$.
Choisissons une résolution libre $F_\bullet \to P$. Posons
$M = \Im(F_1 \to F_0) = \Ker(F_0 \to P)$.
Considérons l'élément $\xi \in \Ext^1_R(P, M)$ donné par
la classe de l'application quotient $\pi : F_1 \to M$. Comme $\xi$ est nul,
il existe une application $s : F_0 \to M$ telle que $\pi = s \circ (F_1 \to F_0)$.
Cela signifie clairement que
$$
F_0 = \Ker(s) \oplus \Ker(F_0 \to P) =
P \oplus \Ker(F_0 \to P)
$$
et l'on conclut.
\end{proof}

\begin{lemma}
\label{lemma-characterize-finite-projective-noetherian}
Soit $R$ un anneau noethérien. Soit $P$ un $R$-module de type fini.
Si $\Ext^1_R(P, M) = 0$ pour tout $R$-module de type fini $M$, alors
$P$ est projectif.
\end{lemma}

\noindent
On peut renforcer ce lemme : il existe une version pour les
$R$-modules de présentation finie sans supposer $R$ noethérien. Dans le cas noethérien,
il existe une version où $M$ parcourt tous les modules de longueur finie.

\begin{proof}
Choisissons une surjection $R^{\oplus n} \to P$ de noyau $M$.
Puisque $\Ext^1_R(P, M) = 0$, cette surjection est scindée, et
on conclut par le Lemme \ref{lemma-characterize-projective}.
\end{proof}

\begin{lemma}
\label{lemma-direct-sum-projective}
Une somme directe de modules projectifs est projective.
\end{lemma}

\begin{proof}
Cela résulte de la caractérisation des modules projectifs comme facteurs
directs de modules libres dans le
Lemme \ref{lemma-characterize-projective}.
\end{proof}

\begin{lemma}
\label{lemma-lift-projective-module}
Soit $R$ un anneau. Soit $I \subset R$ un idéal nilpotent. Soit
$\overline{P}$ un $R/I$-module projectif. Il existe alors un
$R$-module projectif $P$ tel que $P/IP \cong \overline{P}$.
\end{lemma}

\begin{proof}
D'après le Lemme \ref{lemma-characterize-projective},
on peut choisir un ensemble $A$ et une décomposition en somme directe
$\bigoplus_{\alpha \in A} R/I = \overline{P} \oplus \overline{K}$
pour un certain $R/I$-module $\overline{K}$. Notons $F = \bigoplus_{\alpha \in A} R$
le $R$-module libre de base $A$. Choisissons un relèvement
$p : F \to F$ du projecteur $\overline{p}$ associé au
facteur direct $\overline{P}$ de $\bigoplus_{\alpha \in A} R/I$.
Remarquons que $p^2 - p \in \text{End}_R(F)$ est un endomorphisme
nilpotent de $F$ (puisque $I$ est nilpotent et que les coefficients matriciels de
$p^2 - p$ appartiennent à $I$ ; plus précisément, si $I^n = 0$, alors $(p^2 - p)^n = 0$).
Ainsi, d'après le Lemme \ref{lemma-lift-idempotents-noncommutative},
on peut modifier le choix de $p$ et supposer que $p$ est un projecteur.
Posons $P = \Im(p)$.
\end{proof}

\begin{lemma}
\label{lemma-lift-finite-projective-module}
Soit $R$ un anneau. Soit $I \subset R$ un idéal localement nilpotent. Soit
$\overline{P}$ un $R/I$-module projectif de type fini. Il existe alors un
$R$-module projectif de type fini $P$ tel que $P/IP \cong \overline{P}$.
\end{lemma}

\begin{proof}
Rappelons que $\overline{P}$ est un facteur direct d'un $R/I$-module libre
$\bigoplus_{\alpha \in A} R/I$ d'après le Lemme \ref{lemma-characterize-projective}.
Comme $\overline{P}$ est de type fini, il s'ensuit que $\overline{P}$ est contenu
dans $\bigoplus_{\alpha \in A'} R/I$ pour une partie finie $A' \subset A$.
On peut donc supposer que l'on a une décomposition en somme directe
$(R/I)^{\oplus n} = \overline{P} \oplus \overline{K}$
pour un certain $n$ et un certain $R/I$-module $\overline{K}$. Choisissons un relèvement
$p \in \text{Mat}(n \times n, R)$ du projecteur $\overline{p}$
associé au facteur direct $\overline{P}$ de $(R/I)^{\oplus n}$.
Remarquons que $p^2 - p \in \text{Mat}(n \times n, R)$ est nilpotent :
comme $I$ est localement nilpotent et que les coefficients matriciels $c_{ij}$ de
$p^2 - p$ appartiennent à $I$, on a $c_{ij}^t = 0$ pour un certain $t > 0$, et
alors $(p^2 - p)^{tn^2} = 0$ (comme on le voit sur les coefficients matriciels).
Ainsi, d'après le Lemme \ref{lemma-lift-idempotents-noncommutative},
on peut modifier le choix de $p$ et supposer que $p$ est un projecteur.
Posons $P = \Im(p)$.
\end{proof}

\begin{lemma}
\label{lemma-lift-projective}
Soit $R$ un anneau.
Soit $I \subset R$ un idéal.
Soit $M$ un $R$-module.
Supposons que
\begin{enumerate}
\item $I$ soit nilpotent,
\item $M/IM$ soit un $R/I$-module projectif,
\item $M$ soit un $R$-module plat.
\end{enumerate}
Alors $M$ est un $R$-module projectif.
\end{lemma}

\begin{proof}
D'après le Lemme \ref{lemma-lift-projective-module}, on peut trouver un
$R$-module projectif $P$ et un isomorphisme $P/IP \to M/IM$. Nous allons montrer
que $M$ est isomorphe à $P$, ce qui achèvera la démonstration. Comme $P$
est projectif, on peut relever l'application $P \to P/IP \to M/IM$ en un homomorphisme de $R$-modules
$P \to M$ qui est un isomorphisme modulo $I$. Puisque $I^n = 0$
pour un certain $n$, on peut utiliser les filtrations
\begin{align*}
0 = I^nM \subset I^{n - 1}M \subset \ldots \subset IM \subset M \\
0 = I^nP \subset I^{n - 1}P \subset \ldots \subset IP \subset P
\end{align*}
pour voir qu'il suffit de montrer que les applications induites
$I^aP/I^{a + 1}P \to I^aM/I^{a + 1}M$ sont bijectives. Comme $P$
et $M$ sont tous deux des $R$-modules plats, on peut identifier cette application à
$$
I^a/I^{a + 1} \otimes_{R/I} P/IP
\longrightarrow
I^a/I^{a + 1} \otimes_{R/I} M/IM
$$
induite par $P \to M$. Puisque l'on a choisi $P \to M$ de sorte que l'application induite
$P/IP \to M/IM$ soit un isomorphisme, on conclut.
\end{proof}

\begin{lemma}
\label{lemma-projective-modulo-two-ideals}
Soit $R$ un anneau. Soient $I, J \subset R$ des idéaux tels que $I \cap J = 0$.
Soit $P$ un $R$-module tel que $P/IP$ soit un $R/I$-module projectif et que
$P/JP$ soit un $R/J$-module projectif. Alors $P$ est un $R$-module projectif.
\end{lemma}

\begin{proof}
Choisissons une surjection $p : F \to P$, où $F$ est un $R$-module libre.
Puisque $P/IP$ est un $R/I$-module projectif, on peut choisir une application
$f : P/IP \to F/IF$ qui soit un inverse à droite de $p \bmod I$.
Considérons l'application $q : F/JF \to F/(I + J)F \times_{P/(I + J)P} P/JP$
induite par l'application quotient $F/JF \to F/(I + J)F$ et par $p$.
Remarquons que $q$ est un homomorphisme surjectif de $R/J$-modules (petit détail omis).
Considérons l'application $f' : P/JP \to F/(I + J)F \times_{P/(I + J)P} P/JP$
induite par $f$ et par l'application identité.
Puisque $P/JP$ est un $R/J$-module projectif et que $q$ est surjectif,
on peut choisir une application $g : P/JP \to F/JF$ telle que $q \circ g = f'$.
Alors $f \bmod I + J = g \bmod I + J$. Puisque $F = F/JF \times_{F/(I + J)F} F/IF$
car $I \cap J = 0$, on conclut que l'on obtient un homomorphisme bien défini
$h = (f, g) : P \to F$ de $R$-modules. L'application $a = p \circ h : P \to P$
se réduit à l'identité modulo $I$ et modulo $J$. Alors $a$ induit aussi
l'identité sur le sous-module $IP$ : si $x = \sum i_\alpha x_\alpha$
avec $i_\alpha \in I$ et $x_\alpha \in P$, alors
$a(x) = \sum i_\alpha a(x_\alpha) = \sum i_\alpha x_\alpha = x$,
car $a(x_\alpha) = x_\alpha + j_\alpha$ avec $j_\alpha \in J$
et $i_\alpha j_\alpha = 0$. Ainsi $a : P \to P$ agit comme l'identité
sur $IP$ et sur $P/IP$, et l'on conclut que $a$ est un automorphisme de $P$
(petit détail omis). Ainsi $P$ est un facteur direct de $F$, donc projectif.
\end{proof}










\section{Modules projectifs de type fini}
\label{section-finite-projective-modules}

\begin{definition}
\label{definition-locally-free}
Soit $R$ un anneau et $M$ un $R$-module.
\begin{enumerate}
\item On dit que $M$ est {\it localement libre} si l'on peut recouvrir $\Spec(R)$ par
des ouverts standards $D(f_i)$, $i \in I$, tels que $M_{f_i}$ soit un
$R_{f_i}$-module libre pour tout $i \in I$.
\item On dit que $M$ est {\it localement libre de type fini} si l'on peut choisir
le recouvrement de sorte que chaque $M_{f_i}$ soit libre de type fini.
\item On dit que $M$ est {\it localement libre de type fini de rang $r$}
si l'on peut choisir le recouvrement de sorte que chaque $M_{f_i}$ soit isomorphe
à $R_{f_i}^{\oplus r}$.
\end{enumerate}
\end{definition}

\noindent
Remarquons qu'un $R$-module localement libre de type fini est automatiquement
de présentation finie d'après le Lemme \ref{lemma-cover}. De plus, si $M$ est un
module localement libre de type fini de rang $r$ sur un anneau $R$ et si
$R$ est non nul, alors $r$ est déterminé de manière unique par le
Lemme \ref{lemma-rank} (car au moins l'un des anneaux
localisés $	R_{f_i}$ est non nul).

\begin{lemma}
\label{lemma-finite-projective}
Soit $R$ un anneau et soit $M$ un $R$-module.
Les assertions suivantes sont équivalentes :
\begin{enumerate}
\item $M$ est de présentation finie et plat sur $R$,
\item $M$ est projectif de type fini,
\item $M$ est un facteur direct d'un $R$-module libre de type fini,
\item $M$ est de présentation finie et,
pour tout $\mathfrak p \in \Spec(R)$, le
module localisé $M_{\mathfrak p}$ est libre,
\item $M$ est de présentation finie et,
pour tout idéal maximal $\mathfrak m \subset R$, le
module localisé $M_{\mathfrak m}$ est libre,
\item $M$ est de type fini et localement libre,
\item $M$ est localement libre de type fini, et
\item $M$ est de type fini, pour tout idéal premier $\mathfrak p$ le module
$M_{\mathfrak p}$ est libre, et la fonction
$$
\rho_M : \Spec(R) \to \mathbf{Z}, \quad
\mathfrak p
\longmapsto
\dim_{\kappa(\mathfrak p)} M \otimes_R \kappa(\mathfrak p)
$$
est localement constante pour la topologie de Zariski.
\end{enumerate}
\end{lemma}

\begin{proof}
Supposons d'abord $M$ projectif de type fini, c'est-à-dire que (2) est vérifiée.
Prenons une surjection $R^n \to M$ et notons $K$ son noyau.
Puisque $M$ est projectif,
$0 \to K \to R^n \to M \to 0$ est scindée.
Ainsi (2) $\Rightarrow$ (3).
L'implication (3) $\Rightarrow$ (2) résulte de ce qu'un
facteur direct d'un module projectif est projectif ; voir le
Lemme \ref{lemma-characterize-projective}.

\medskip\noindent
Supposons (3), de sorte que l'on puisse écrire $K \oplus M \cong R^{\oplus n}$.
Alors $K$ est un facteur direct de $R^n$, et donc de type fini.
Cela montre que $M = R^{\oplus n}/K$ est de présentation finie.
Autrement dit, (3) $\Rightarrow$ (1).

\medskip\noindent
Supposons que $M$ soit de présentation finie et plat, c'est-à-dire que (1) soit vérifiée.
Nous allons démontrer (7). Choisissons un idéal premier $\mathfrak p$ quelconque et
$x_1, \ldots, x_r \in M$ dont les images forment une base de
$M \otimes_R \kappa(\mathfrak p)$. D'après le
lemme de Nakayama (sous la forme du Lemme \ref{lemma-NAK-localization}),
ces éléments engendrent $M_g$ pour un certain $g \in R$, $g \not \in \mathfrak p$.
La surjection correspondante $\varphi : R_g^{\oplus r} \to M_g$
possède les deux propriétés suivantes : (a) $\Ker(\varphi)$ est un
$R_g$-module de type fini (voir le Lemme \ref{lemma-extension})
et (b) $\Ker(\varphi) \otimes \kappa(\mathfrak p) = 0$
par platitude de $M_g$ sur $R_g$ (voir le
Lemme \ref{lemma-flat-tor-zero}).
Ainsi, d'après le lemme de Nakayama encore une fois, il existe $g' \in R_g$ tel que
$\Ker(\varphi)_{g'} = 0$. Autrement dit, $M_{gg'}$ est libre.

\medskip\noindent
Un module localement libre de type fini est de type fini ; voir le
Lemme \ref{lemma-cover},
donc (7) $\Rightarrow$ (6).
Il est clair que (6) $\Rightarrow$ (7) et que (7) $\Rightarrow$ (8).

\medskip\noindent
Un module localement libre de type fini est de présentation finie ; voir le
Lemme \ref{lemma-cover},
donc (7) $\Rightarrow$ (4).
Bien entendu, (4) implique (5).
Puisque l'on peut vérifier la platitude localement (voir le
Lemme \ref{lemma-flat-localization}),
on conclut que (5) implique (1).
À ce stade, on a
$$
\xymatrix{
(2) \ar@{<=>}[r] & (3) \ar@{=>}[r] & (1) \ar@{=>}[r] &
(7)  \ar@{<=>}[r] \ar@{=>}[rd] \ar@{=>}[d] & (6) \\
& & (5) \ar@{=>}[u] & (4) \ar@{=>}[l] & (8)
}
$$

\medskip\noindent
Supposons que $M$ satisfasse (1), (4), (5), (6) et (7).
Nous allons démontrer (3). Il suffit
de montrer que $M$ est projectif. Il faut montrer que $\Hom_R(M, -)$
est exact. Soit $0 \to N'' \to N \to N'\to 0$ une suite exacte courte de
$R$-modules. Il faut montrer que
$0 \to \Hom_R(M, N'') \to \Hom_R(M, N) \to
\Hom_R(M, N') \to 0$ est exacte.
Comme $M$ est localement libre de type fini, il existe un recouvrement
$\Spec(R) = \bigcup D(f_i)$ tel que $M_{f_i}$ soit libre de type fini.
D'après le
Lemme \ref{lemma-hom-from-finitely-presented},
on voit que
$$
0 \to \Hom_R(M, N'')_{f_i} \to \Hom_R(M, N)_{f_i} \to
\Hom_R(M, N')_{f_i} \to 0
$$
est égale à
$0 \to \Hom_{R_{f_i}}(M_{f_i}, N''_{f_i}) \to
\Hom_{R_{f_i}}(M_{f_i}, N_{f_i}) \to
\Hom_{R_{f_i}}(M_{f_i}, N'_{f_i}) \to 0$,
qui est exacte puisque $M_{f_i}$ est libre et que la localisation
$0 \to N''_{f_i} \to N_{f_i} \to N'_{f_i} \to 0$
est exacte (car la localisation est exacte). On voit donc que
$0 \to \Hom_R(M, N'') \to \Hom_R(M, N) \to
\Hom_R(M, N') \to 0$ est exacte d'après le
Lemme \ref{lemma-cover}.

\medskip\noindent
Enfin, supposons que (8) soit vérifiée. Choisissons un idéal maximal $\mathfrak m \subset R$.
Choisissons $x_1, \ldots, x_r \in M$ dont les images forment une base sur $\kappa(\mathfrak m)$ de
$M \otimes_R \kappa(\mathfrak m) = M/\mathfrak mM$. En particulier,
$\rho_M(\mathfrak m) = r$. D'après le
Lemme de Nakayama \ref{lemma-NAK},
il existe $f \in R$, $f \not \in \mathfrak m$, tel que
$x_1, \ldots, x_r$ engendrent $M_f$ sur $R_f$. Comme, par hypothèse,
$\rho_M$ est localement constante, il existe $g \in R$, $g \not \in \mathfrak m$,
tel que $\rho_M$ soit constante de valeur $r$ sur $D(g)$. Nous affirmons que
$$
\Psi : R_{fg}^{\oplus r} \longrightarrow M_{fg}, \quad
(a_1, \ldots, a_r) \longmapsto \sum a_i x_i
$$
est un isomorphisme. Cela montrera que $M$ est localement libre de type
fini, c'est-à-dire que (7) est vérifiée. Pour établir l'affirmation,
il suffit de montrer que l'application induite sur les localisations
$\Psi_{\mathfrak p} : R_{\mathfrak p}^{\oplus r} \to M_{\mathfrak p}$
est un isomorphisme pour tout $\mathfrak p \in D(fg)$ ; voir le
Lemme \ref{lemma-characterize-zero-local}.
Par le choix de $f$, l'application $\Psi_{\mathfrak p}$
est surjective. D'après l'hypothèse (8), on a
$M_{\mathfrak p} \cong R_{\mathfrak p}^{\oplus \rho_M(\mathfrak p)}$,
et, par le choix de $g$, on a $\rho_M(\mathfrak p) = r$.
Ainsi $\Psi_{\mathfrak p}$ détermine une surjection
$R_{\mathfrak p}^{\oplus r} \to
M_{\mathfrak p} \cong R_{\mathfrak p}^{\oplus r}$,
qui est donc un isomorphisme d'après le
Lemme \ref{lemma-fun}.
(Bien entendu, ce dernier fait résulte aussi d'un simple argument matriciel.)
\end{proof}

\begin{lemma}
\label{lemma-finite-projective-reduced}
Soit $R$ un anneau réduit et soit $M$ un $R$-module. Alors les
conditions équivalentes du Lemme \ref{lemma-finite-projective}
sont aussi équivalentes à
\begin{enumerate}
\item[(9)] $M$ est de type fini et la fonction
$\rho_M : \Spec(R) \to \mathbf{Z}$, $\mathfrak p \mapsto
\dim_{\kappa(\mathfrak p)} M \otimes_R \kappa(\mathfrak p)$
est localement constante pour la topologie de Zariski.
\end{enumerate}
\end{lemma}

\begin{proof}
Choisissons un idéal maximal $\mathfrak m \subset R$.
Choisissons $x_1, \ldots, x_r \in M$ dont les images forment une base sur $\kappa(\mathfrak m)$ de
$M \otimes_R \kappa(\mathfrak m) = M/\mathfrak mM$. En particulier,
$\rho_M(\mathfrak m) = r$. D'après le
Lemme de Nakayama \ref{lemma-NAK},
il existe $f \in R$, $f \not \in \mathfrak m$, tel que
$x_1, \ldots, x_r$ engendrent $M_f$ sur $R_f$. Comme, par hypothèse,
$\rho_M$ est localement constante, il existe $g \in R$, $g \not \in \mathfrak m$,
tel que $\rho_M$ soit constante de valeur $r$ sur $D(g)$. Nous affirmons que
$$
\Psi : R_{fg}^{\oplus r} \longrightarrow M_{fg}, \quad
(a_1, \ldots, a_r) \longmapsto \sum a_i x_i
$$
est un isomorphisme. Cela montrera que $M$ est localement libre de type
fini, c'est-à-dire que (7) est vérifiée. Puisque $\Psi$ est surjective, il suffit de
montrer que $\Psi$ est injective. Puisque $R_{fg}$ est réduit,
il suffit de montrer que $\Psi$ est injective après localisation
en tous les idéaux premiers minimaux $\mathfrak p$ de $R_{fg}$ ; voir le
Lemme \ref{lemma-reduced-ring-sub-product-fields}.
Cependant, on sait que $R_\mathfrak p = \kappa(\mathfrak p)$
d'après le Lemme \ref{lemma-minimal-prime-reduced-ring}, et
$\rho_M(\mathfrak p) = r$ ; ainsi $\Psi_\mathfrak p :
R_\mathfrak p^{\oplus r} \to M \otimes_R \kappa(\mathfrak p)$
est un isomorphisme, puisqu'il s'agit d'une application surjective entre espaces vectoriels
de dimension finie et de même dimension.
\end{proof}

\begin{remark}
\label{remark-warning}
Il n'est pas vrai qu'un $R$-module de type fini qui soit
plat sur $R$ soit automatiquement projectif. Un contre-
exemple s'obtient en prenant $R = \mathcal{C}^\infty(\mathbf{R})$
égal à l'anneau des fonctions indéfiniment différentiables sur
$\mathbf{R}$, et $M = R_{\mathfrak m} = R/I$, où
$\mathfrak m = \{f \in R \mid f(0) = 0\}$ et
$I = \{f \in R \mid \exists \epsilon, \epsilon > 0 :
f(x) = 0\ \forall x, |x| < \epsilon\}$.
\end{remark}

\begin{lemma}
\label{lemma-finite-flat-local}
(Attention : voir la Remarque \ref{remark-warning}.)
Supposons que $R$ soit un anneau local et que $M$ soit un $R$-module
plat de type fini. Alors $M$ est libre de type fini.
\end{lemma}

\begin{proof}
Cela résulte du critère équationnel de platitude ; voir le
Lemme \ref{lemma-flat-eq}. Plus précisément, supposons que
$x_1, \ldots, x_r \in M$ aient pour images une base de
$M/\mathfrak mM$. D'après le Lemme de Nakayama \ref{lemma-NAK},
ces éléments engendrent $M$. Nous voulons montrer qu'il
n'existe aucune relation entre les $x_i$. Nous allons plutôt montrer
par récurrence sur $n$ que, si $x_1, \ldots, x_n \in M$
sont linéairement indépendants dans l'espace vectoriel
$M/\mathfrak mM$, alors ils sont indépendants sur $R$.

\medskip\noindent
Le cas initial de la récurrence est celui où l'on a
$x \in M$, $x \not\in \mathfrak mM$ et une relation
$fx = 0$. D'après le critère équationnel, il
existe $y_j \in M$ et $a_j \in R$ tels que
$x = \sum a_j y_j$ et $fa_j = 0$ pour tout $j$.
Puisque $x \not\in \mathfrak mM$, on voit qu'au moins
un $a_j$ est une unité, et donc que $f = 0 $.

\medskip\noindent
Supposons que $\sum f_i x_i$ soit une relation entre $x_1, \ldots, x_n$.
Par le choix des $x_i$, on a $f_i \in \mathfrak m$.
D'après le critère équationnel de platitude, il existe
$a_{ij} \in R$ et $y_j \in M$ tels que
$x_i = \sum a_{ij} y_j$ et $\sum f_i a_{ij} = 0$.
Puisque $x_n \not \in \mathfrak mM$, on voit que
$a_{nj}\not\in \mathfrak m$ pour au moins un $j$.
Puisque $\sum f_i a_{ij} = 0$, on obtient
$f_n = \sum_{i = 1}^{n-1} (-a_{ij}/a_{nj}) f_i$.
La relation $\sum f_i x_i = 0$ peut alors se réécrire
$\sum_{i = 1}^{n-1} f_i( x_i + (-a_{ij}/a_{nj}) x_n) = 0$.
Remarquons que les éléments $x_i + (-a_{ij}/a_{nj}) x_n$ ont pour images
$n-1$ éléments linéairement indépendants de $M/\mathfrak mM$.
Par hypothèse de récurrence, tous les $f_i$, $i \leq n-1$,
doivent être nuls, ainsi que $f_n = \sum_{i = 1}^{n-1} (-a_{ij}/a_{nj}) f_i$.
Cela établit l'étape de récurrence.
\end{proof}

\begin{lemma}
\label{lemma-finite-projective-descends}
Soit $R \to S$ un homomorphisme local plat d'anneaux locaux.
Soit $M$ un $R$-module de type fini. Alors $M$ est projectif de type fini
sur $R$ si et seulement si $M \otimes_R S$ est projectif de type fini
sur $S$.
\end{lemma}

\begin{proof}
D'après le Lemme \ref{lemma-finite-projective}, être projectif de type fini
sur un anneau local revient à être libre de type fini.
Supposons que $M \otimes_R S$ soit un $S$-module libre de type fini.
Choisissons $x_1, \ldots, x_r \in M$ dont les images dans $M/\mathfrak m_RM$
forment une base sur $\kappa(\mathfrak m)$. Alors
$x_1 \otimes 1, \ldots, x_r \otimes 1$
forment une base de $M \otimes_R S$. Cela implique que
l'application $R^{\oplus r} \to M, (a_i) \mapsto \sum a_i x_i$
devient un isomorphisme après tensorisation par $S$.
Par fidèle platitude de $R \to S$, voir le Lemme \ref{lemma-local-flat-ff},
on voit que c'est un isomorphisme.
\end{proof}

\begin{lemma}
\label{lemma-locally-free-semi-local-free}
Soit $R$ un anneau semi-local.
Soit $M$ un module localement libre de type fini.
Si $M$ est de rang constant, alors $M$ est libre. En particulier, si $R$ a
un spectre connexe, alors $M$ est libre.
\end{lemma}

\begin{proof}
Omis. Indications : montrons d'abord que $M/\mathfrak m_iM$ a la
même dimension $d$ pour tous les idéaux maximaux $\mathfrak m_1, \ldots, \mathfrak m_n$
de $R$, en utilisant la constance du rang.
Montrons ensuite qu'il existe des éléments $x_1, \ldots, x_d \in M$
qui forment une base de chaque $M/\mathfrak m_iM$, grâce au théorème
chinois des restes. Montrons enfin que $x_1, \ldots, x_d$ est une base de $M$.
\end{proof}

\noindent
Voici un lemme technique utilisé dans le chapitre sur les groupoïdes.

\begin{lemma}
\label{lemma-semi-local-module-basis-in-submodule}
Soit $R$ un anneau local d'idéal maximal $\mathfrak m$ et de
corps résiduel infini.
Soit $R \to S$ un homomorphisme d'anneaux.
Soit $M$ un $S$-module et soit $N \subset M$ un $R$-sous-module.
Supposons que
\begin{enumerate}
\item $S$ soit semi-local et que $\mathfrak mS$ soit contenu dans le
radical de Jacobson de $S$,
\item $M$ soit un $S$-module libre de type fini, et
\item $N$ engendre $M$ comme $S$-module.
\end{enumerate}
Alors $N$ contient une $S$-base de $M$.
\end{lemma}

\begin{proof}
Supposons $M$ libre de rang $n$. Soit $I \subset S$ le radical de Jacobson.
D'après le Lemme de Nakayama \ref{lemma-NAK}, une suite d'éléments
$m_1, \ldots, m_n$ est une base de $M$ si et seulement si
$\overline{m}_i \in M/IM$ engendrent $M/IM$. On peut donc remplacer
$M$ par $M/IM$, $N$ par $N/(N \cap IM)$, $R$ par $R/\mathfrak m$,
et $S$ par $S/IS$. Dans ce cas, on voit que $S$ est un produit fini
de corps $S = k_1 \times \ldots \times k_r$ et que
$M = k_1^{\oplus n} \times \ldots \times k_r^{\oplus n}$.
Le fait que $N \subset M$ engendre $M$ comme $S$-module
signifie qu'il existe $x_j \in N$ tels qu'une combinaison linéaire
$\sum a_j x_j$, avec $a_j \in S$, ait une composante non nulle dans chaque
facteur $k_i^{\oplus n}$.
Puisque $R = k$ est un corps infini, cela signifie aussi qu'une certaine
combinaison linéaire $y = \sum c_j x_j$, avec $c_j \in k$, a une
composante non nulle dans chaque facteur. Ainsi $y \in N$ engendre un
facteur direct libre $Sy \subset M$. Par récurrence sur $n$, le résultat
vaut pour $M/Sy$ et le sous-module $\overline{N} = N/(N \cap Sy)$.
Autrement dit, il existe $\overline{y}_2, \ldots, \overline{y}_n$
dans $\overline{N}$ qui engendrent librement $M/Sy$. Alors
$y, y_2, \ldots, y_n$ engendrent librement $M$, et l'on conclut.
\end{proof}

\begin{lemma}
\label{lemma-evaluation-map-iso-finite-projective}
Soit $R$ un anneau. Soient $L$, $M$, $N$ des $R$-modules.
L'application canonique
$$
\Hom_R(M, N) \otimes_R L \to \Hom_R(M, N \otimes_R L)
$$
est un isomorphisme si $M$ est projectif de type fini.
\end{lemma}

\begin{proof}
D'après le Lemme \ref{lemma-finite-projective}, on voit que $M$
est de présentation finie et localement libre de type fini.
D'après les Lemmes \ref{lemma-hom-from-finitely-presented} et
\ref{lemma-tensor-product-localization}, la formation du
membre de gauche et du membre de droite de la flèche commute à la
localisation. On peut vérifier que notre application est un isomorphisme
après localisation ; voir le Lemme \ref{lemma-cover}.
On peut donc supposer $M$ libre de type fini. Dans ce cas,
le lemme est immédiat.
\end{proof}


 


\section{Lieux ouverts définis par des homomorphismes de modules}
\label{section-loci-maps}

\noindent
L'ensemble des idéaux premiers où un homomorphisme de modules donné est surjectif ou est un isomorphisme
est parfois ouvert. Dans le cas de modules projectifs de type fini, on peut considérer
le rang de l'homomorphisme.

\begin{lemma}
\label{lemma-map-between-finite}
Soit $R$ un anneau. Soit $\varphi : M \to N$ un homomorphisme de $R$-modules,
où $N$ est un $R$-module de type fini. Alors on a l'égalité
\begin{align*}
U & = \{\mathfrak p \subset R \mid
\varphi_{\mathfrak p} : M_{\mathfrak p} \to N_{\mathfrak p}
\text{ est surjective}\} \\
& = \{\mathfrak p \subset R \mid
\varphi \otimes \kappa(\mathfrak p) :
M \otimes \kappa(\mathfrak p) \to N \otimes \kappa(\mathfrak p)
\text{ est surjective}\}
\end{align*}
et $U$ est un ouvert de $\Spec(R)$. De plus, pour tout $f \in R$
tel que $D(f) \subset U$, l'homomorphisme $M_f \to N_f$ est surjectif.
\end{lemma}

\begin{proof}
L'égalité de la formule affichée résulte du lemme de Nakayama.
Le lemme de Nakayama implique aussi que $U$ est ouvert. Voir le
Lemme \ref{lemma-NAK}, en particulier la partie (3). Si $D(f) \subset U$, alors
$M_f \to N_f$ est surjectif sur toutes les localisations aux idéaux premiers de
$R_f$, et il est donc surjectif d'après le Lemme \ref{lemma-characterize-zero-local}.
\end{proof}

\begin{lemma}
\label{lemma-map-between-finitely-presented}
Soit $R$ un anneau. Soit $\varphi : M \to N$ un homomorphisme de $R$-modules,
où $M$ est de type fini et $N$ de présentation finie. Alors
$$
U = \{\mathfrak p \subset R \mid
\varphi_{\mathfrak p} : M_{\mathfrak p} \to N_{\mathfrak p}
\text{ est un isomorphisme}\}
$$
est un ouvert de $\Spec(R)$.
\end{lemma}

\begin{proof}
Soit $\mathfrak p \in U$. Choisissons une présentation
$N = R^{\oplus n}/\sum_{j = 1, \ldots, m} R k_j$.
Notons $e_i$ l'image dans $N$ du $i$-ième vecteur de base de $R^{\oplus n}$.
Pour chaque $i \in \{1, \ldots, n\}$, choisissons un élément
$m_i \in M_{\mathfrak p}$ tel que $\varphi(m_i) = f_i e_i$
pour un certain $f_i \in R$, $f_i \not \in \mathfrak p$. C'est possible
car $\varphi_{\mathfrak p}$ est un isomorphisme. Posons $f = f_1 \ldots f_n$
et soit $\psi : R_f^{\oplus n} \to M_f$ l'homomorphisme qui envoie
le $i$-ième vecteur de base sur $m_i/f_i$. Remarquons que
$\varphi_f \circ \psi$ est la localisation
en $f$ de l'homomorphisme donné $R^{\oplus n} \to N$.
Comme $\varphi_{\mathfrak p}$ est un isomorphisme, on
voit que $\psi(k_j)$ est un élément de $M$ dont l'image est nulle
dans $M_{\mathfrak p}$. Il existe donc
$g_j \in R$, $g_j \not \in \mathfrak p$, tel que $g_j \psi(k_j) = 0$.
En posant $g = g_1 \ldots g_m$, on voit que $\psi_g$ se factorise par
$N_{fg}$ et donne un homomorphisme $\chi : N_{fg} \to M_{fg}$. Par construction,
$\chi$ est un inverse à droite de $\varphi_{fg}$. Il s'ensuit que
$\chi_\mathfrak p$ est un isomorphisme. D'après le Lemme \ref{lemma-map-between-finite},
il existe $h \in R$, $h \not \in \mathfrak p$, tel que
$\chi_h : N_{fgh} \to M_{fgh}$ soit surjectif.
Ainsi $\varphi_{fgh}$ et $\chi_h$ sont des homomorphismes inverses l'un de l'autre,
ce qui implique $D(fgh) \subset U$, comme souhaité.
\end{proof}

\begin{lemma}
\label{lemma-finitely-presented-localization-free}
Soit $R$ un anneau. Soit $\mathfrak p \subset R$ un idéal premier.
Soit $M$ un $R$-module de présentation finie. Si $M_\mathfrak p$
est libre, alors il existe $f \in R$, $f \not \in \mathfrak p$,
tel que $M_f$ soit un $R_f$-module libre.
\end{lemma}

\begin{proof}
Choisissons une base $x_1, \ldots, x_n \in M_\mathfrak p$. On peut choisir
$f \in R$, $f \not \in \mathfrak p$, tel que $x_i$ soit l'image
d'un certain $y_i \in M_f$. Après avoir remplacé $y_i$ par $f^m y_i$ avec $m \gg 0$,
on peut supposer $y_i \in M$. Cela revient à remplacer $x_1, \ldots, x_n$
par $f^mx_1, \ldots, f^mx_n$, qui est encore une base puisque l'image de $f$ est
une unité dans $R_\mathfrak p$. On obtient donc
un homomorphisme $\varphi = (y_1, \ldots, y_n) : R^{\oplus n} \to M$
de $R$-modules dont la localisation en $\mathfrak p$ est un isomorphisme.
D'après le Lemme \ref{lemma-map-between-finitely-presented},
on peut trouver $f \in R$, $f \not \in \mathfrak p$,
tel que $\varphi_\mathfrak q$ soit un isomorphisme pour tout idéal premier
$\mathfrak q \subset R$ vérifiant $f \not \in \mathfrak q$.
Il résulte alors du Lemme \ref{lemma-characterize-zero-local}
que $\varphi_f$ est un isomorphisme, ce qui achève la démonstration.
\end{proof}

\begin{lemma}
\label{lemma-cokernel-flat}
Soit $R$ un anneau. Soit $\varphi : P_1 \to P_2$ un homomorphisme de
modules projectifs de type fini. Alors
\begin{enumerate}
\item L'ensemble $U$ des idéaux premiers $\mathfrak p \in \Spec(R)$ tels que
$\varphi \otimes \kappa(\mathfrak p)$ soit injectif est ouvert et,
pour tout $f\in R$ tel que $D(f) \subset U$, on a
\begin{enumerate}
\item $P_{1, f} \to P_{2, f}$ est injectif, et
\item le module $\Coker(\varphi)_f$ est projectif de type fini sur $R_f$.
\end{enumerate}
\item L'ensemble $W$ des idéaux premiers $\mathfrak p \in \Spec(R)$ tels que
$\varphi \otimes \kappa(\mathfrak p)$ soit surjectif est ouvert et,
pour tout $f\in R$ tel que $D(f) \subset W$, on a
\begin{enumerate}
\item $P_{1, f} \to P_{2, f}$ est surjectif, et
\item le module $\Ker(\varphi)_f$ est projectif de type fini sur $R_f$.
\end{enumerate}
\item L'ensemble $V$ des idéaux premiers $\mathfrak p \in \Spec(R)$ tels que
$\varphi \otimes \kappa(\mathfrak p)$ soit un isomorphisme est ouvert et,
pour tout $f\in R$ tel que $D(f) \subset V$, l'homomorphisme
$\varphi : P_{1, f} \to P_{2, f}$ est un isomorphisme de modules sur $R_f$.
\end{enumerate}
\end{lemma}

\begin{proof}
Pour démontrer que l'ensemble $U$ est ouvert, on peut travailler localement sur
$\Spec(R)$. On peut donc remplacer $R$ par une localisation convenable
et supposer que $P_1 = R^{n_1}$ et $P_2 = R^{n_2}$ ; voir le Lemme
\ref{lemma-finite-projective}. Dans ce cas, l'injectivité de
$\varphi \otimes \kappa(\mathfrak p)$ équivaut à $n_1 \leq n_2$
et à ce qu'un certain mineur $n_1 \times n_1$, noté $f$, de la matrice de $\varphi$, soit
inversible dans $\kappa(\mathfrak p)$. Ainsi $D(f) \subset U$.
Cet argument montre aussi que $P_{1, \mathfrak p} \to P_{2, \mathfrak p}$
est injectif pour $\mathfrak p \in U$.

\medskip\noindent
Supposons maintenant $D(f) \subset U$. La remarque du paragraphe précédent
et le Lemme \ref{lemma-characterize-zero-local} montrent que
$P_{1, f} \to P_{2, f}$ est injectif, c'est-à-dire que (1)(a) est vérifiée.
D'après le Lemme \ref{lemma-finite-projective}, pour démontrer (1)(b),
il suffit de montrer que $\Coker(\varphi)$ est projectif de type fini
localement sur $D(f)$. Ainsi, comme on l'a vu ci-dessus, on peut
supposer que $P_1 = R^{n_1}$ et $P_2 = R^{n_2}$
et qu'un certain mineur de la matrice de $\varphi$ est inversible dans $R$.
Si le mineur considéré correspond aux $n_1$ premiers
vecteurs de base de $R^{n_2}$, alors les $n_2 - n_1$ derniers vecteurs de base
donnent un homomorphisme $R^{n_2 - n_1} \to R^{n_2} \to \Coker(\varphi)$
dont on vérifie facilement que c'est un isomorphisme.

\medskip\noindent
L'ouverture de $W$ et (2)(a) lorsque $D(f) \subset W$ résultent du
Lemme \ref{lemma-map-between-finite}. Puisque $P_{2, f}$ est projectif
sur $R_f$, on voit que $\varphi_f : P_{1, f} \to P_{2, f}$ possède une section,
et il s'ensuit que $\Ker(\varphi)_f$ est un facteur direct de $P_{2, f}$.
Par conséquent, $\Ker(\varphi)_f$ est projectif de type fini. Ainsi (2)(b) est également vérifiée.

\medskip\noindent
Il est clair que $V = U \cap W$ est ouvert, et l'autre assertion de (3)
résulte de (1)(a) et (2)(a).
\end{proof}






\section{Descente fidèlement plate de la projectivité des modules}
\label{section-ffdescent-projectivity-introduction}

\medskip\noindent
Dans les quelques sections suivantes, nous démontrons, à la suite de
Raynaud et Gruson \cite{GruRay}, que la
projectivité des modules descend le long des homomorphismes d'anneaux fidèlement plats. L'idée de la
démonstration est d'utiliser un dévissage à la Kaplansky \cite{Kaplansky} pour se ramener
au
cas des modules dénombrablement engendrés. Étant donnée une filtration convenable d'un
module $M$, le dévissage permet d'exprimer $M$ comme somme directe des quotients
successifs des sous-modules de la filtration (voir la
section \ref{section-transfinite-devissage}).
Grâce à cette technique, nous démontrons qu'un module projectif est une somme directe de
modules dénombrablement engendrés (Théorème \ref{theorem-projective-direct-sum}). Pour
démontrer la descente de la projectivité des modules dénombrablement engendrés, nous introduisons une
condition « de Mittag-Leffler » sur les modules, montrons qu'un module dénombrablement engendré
est projectif si et seulement s'il est plat et de Mittag-Leffler (Théorème
\ref{theorem-projectivity-characterization}), puis montrons que la propriété
d'être un module de Mittag-Leffler descend (Lemme
\ref{lemma-ffdescent-ML}). Enfin, étant donné un module quelconque $M$ dont le
changement de base par un homomorphisme fidèlement plat d'anneaux est projectif, nous filtrons $M$ par
des sous-modules dont les quotients successifs sont des modules projectifs
dénombrablement engendrés, puis le dévissage montre que $M$ est une somme directe de modules projectifs,
donc est lui-même projectif (Théorème \ref{theorem-ffdescent-projectivity}).

\medskip\noindent
Signalons qu'il y a une erreur dans la démonstration de la descente fidèlement plate
de la projectivité dans \cite{GruRay}. La descente
de la projectivité le long des homomorphismes d'anneaux fidèlement plats y est déduite de la descente de la
projectivité le long d'un type plus général d'homomorphisme d'anneaux
(\cite[exemple 3.1.4(1) de la partie II]{GruRay}).
Cependant, la démonstration de la descente le long de ce type plus général
d'homomorphisme est incorrecte. Dans \cite{G}, Gruson explique ce qui ne va pas,
mais ne fournit pas de correction pour le cas qui nous intéresse.
Combler cette lacune dans la démonstration de la descente fidèlement plate de la projectivité
revient à démontrer que la propriété d'être un module de Mittag-Leffler
descend le long des homomorphismes d'anneaux fidèlement plats. C'est ce que nous faisons dans le
Lemme \ref{lemma-ffdescent-ML}.





\section{Caractérisation de la platitude}
\label{section-characterize-flatness}

\noindent
Dans cette section, nous étudions des critères de platitude. Le résultat principal de
cette section est le théorème de Lazard (Théorème \ref{theorem-lazard}
ci-dessous), qui affirme qu'un module plat est la colimite d'un système filtrant de
modules libres de type fini.
Rappelons au lecteur le « critère équationnel de platitude » ; voir le
Lemme \ref{lemma-flat-eq}.
Il se trouve que l'on peut le transformer en une propriété apparemment bien plus forte.

\begin{lemma}
\label{lemma-flat-factors-free}
Soit $M$ un $R$-module. Les assertions suivantes sont équivalentes :
\begin{enumerate}
\item $M$ est plat.
\item Si $f: R^n \to M$ est un homomorphisme de modules et $x \in \Ker(f)$, alors il existe
des homomorphismes de modules $h: R^n \to R^m$ et $g: R^m \to M$ tels que
$f = g \circ h$ et $x \in \Ker(h)$.
\item Supposons que $f: R^n \to M$ soit un homomorphisme de modules, que $N \subset \Ker(f)$ soit un
sous-module quelconque et que $h: R^n \to R^{m}$ soit un homomorphisme tel que $N \subset \Ker(h)$
et que $f$ se factorise par $h$. Alors, pour tout $x \in \Ker(f)$, on peut trouver un homomorphisme
$h': R^n \to R^{m'}$ tel que $N + Rx \subset \Ker(h')$ et que $f$
se factorise par $h'$.
\item Si $f: R^n \to M$ est un homomorphisme de modules et si $N \subset \Ker(f)$ est un
sous-module de type fini, alors il existe des homomorphismes de modules $h: R^n \to
R^m$ et $g: R^m \to M$ tels que $f = g \circ h$ et $N \subset
\Ker(h)$.
\end{enumerate}
\end{lemma}

\begin{proof}
L'équivalence de (1) et (2) n'est qu'une reformulation du critère équationnel
de platitude\footnote{En fait, un homomorphisme de modules $f : R^n \to M$
correspond au choix d'éléments $x_1, x_2, \ldots, x_n$ de $M$
(à savoir les images des vecteurs de la base canonique $e_1, e_2, \ldots,
e_n$) ; de plus, un élément $x \in \Ker(f)$ correspond à une
relation entre ces $x_1, x_2, \ldots, x_n$ (à savoir la relation
$\sum_i f_i x_i = 0$, où les $f_i$ sont les coordonnées de $x$).
L'homomorphisme de modules $h$ (représenté par une matrice $m \times n$)
correspond à la matrice $(a_{ij})$ du Lemme \ref{lemma-flat-eq},
et les $y_j$ du Lemme \ref{lemma-flat-eq} sont les images des
vecteurs de la base canonique de $R^m$ par $g$.}. Pour montrer que
(2) implique (3), soit $g: R^m \to M$
l'homomorphisme tel que $f$ se factorise sous la forme $f = g \circ h$. D'après (2), choisissons $h'': R^m
\to R^{m'}$ tel que $h''$ annule $h(x)$ et que $g: R^m \to M$
se factorise par $h''$. Alors $h' = h'' \circ h$ convient. L'assertion (3) implique (4)
par récurrence sur le nombre de générateurs de $N \subset \Ker(f)$ dans (4).
Il est clair que (4) implique (2).
\end{proof}

\begin{lemma}
\label{lemma-flat-factors-fp}
Soit $M$ un $R$-module. Alors $M$ est plat si et seulement si la condition
suivante est satisfaite : si $P$ est un $R$-module de présentation finie et $f: P
\to M$ un homomorphisme de modules, alors il existe un $R$-module libre de type fini $F$ et des
homomorphismes de modules $h: P \to F$ et $g: F \to M$ tels que $f = g
\circ h$.
\end{lemma}

\begin{proof}
Ce n'est qu'une reformulation de la condition (4) du
Lemme \ref{lemma-flat-factors-free}.
\end{proof}

\begin{lemma}
\label{lemma-flat-surjective-hom}
Soit $M$ un $R$-module. Alors $M$ est plat si et seulement si la condition
suivante est satisfaite : pour tout $R$-module de présentation finie $P$, si $N \to
M$ est un homomorphisme surjectif de $R$-modules, alors l'homomorphisme induit $\Hom_R(P, N)
\to \Hom_R(P, M)$ est surjectif.
\end{lemma}

\begin{proof}
Supposons d'abord $M$ plat. Il faut montrer que, si $P$ est de présentation finie,
alors toute application $f: P \to M$ se factorise par l'application $N
\to M$. D'après le
Lemme \ref{lemma-flat-factors-fp},
l'application $f$ se factorise par une application $F \to M$, où $F$ est libre de type fini.
Puisque $F$ est libre, cette application se factorise par
$N \to M$. Ainsi $f$ se factorise par $N \to M$.

\medskip\noindent
Réciproquement, supposons que la condition du lemme soit satisfaite. Soit $f: P \to M$
une application issue d'un module de présentation finie $P$. Choisissons un module libre $N$ et une
surjection $N \to M$ sur $M$. Alors $f$ se factorise par $N \to
M$ et, puisque $P$ est de type fini, $f$ se factorise par un sous-module libre
de type fini de $N$. Ainsi $M$ satisfait à la condition du Lemme
\ref{lemma-flat-factors-fp}, et est donc plat.
\end{proof}

\begin{theorem}[Théorème de Lazard]
\label{theorem-lazard}
Soit $M$ un $R$-module. Alors $M$ est plat si et seulement s'il est la colimite
d'un système filtrant de $R$-modules libres de type fini.
\end{theorem}

\begin{proof}
Une colimite d'un système filtrant de modules plats est plate, puisque les colimites
filtrantes sont exactes et commutent au produit tensoriel. Ainsi, si $M$ est la colimite
d'un système filtrant de modules libres de type fini, alors $M$ est plat.

\medskip\noindent
Pour la réciproque, rappelons d'abord que tout module $M$ peut s'écrire comme la
colimite d'un système filtrant de modules de présentation finie, de la manière
suivante. Choisissons une surjection $f: R^I \to M$ pour un certain ensemble $I$, et soit $K$
son noyau. Soit $E$ l'ensemble des couples ordonnés $(J, N)$, où $J$ est une
partie finie de $I$ et $N$ un sous-module de type fini de $R^J \cap K$.
On fait de $E$ un ensemble ordonné filtrant en posant $(J, N) \leq
(J', N')$ si et seulement si $J \subset J'$ et $N \subset N'$. Définissons $M_e =
R^J/N$ pour $e = (J, N)$, et définissons $f_{ee'}: M_e \to M_{e'}$ comme
l'application naturelle pour $e \leq e'$. Alors $(M_e, f_{ee'})$ est un système filtrant, et
les applications naturelles $f_e: M_e \to M$ induisent un isomorphisme
$\colim_{e
\in E} M_e \xrightarrow{\cong} M$.

\medskip\noindent
Supposons maintenant $M$ plat. Soit $I = M \times \mathbf{Z}$, notons $(x_i)$ la
base canonique de $R^{I}$ et prenons, dans la discussion ci-dessus, $f: R^I
\to M$ égal à l'application qui envoie $x_i$ sur la projection de $i$ sur $M$.
Pour démontrer le théorème, il suffit de montrer que les $e \in E$ tels que $M_e$
soit libre forment une partie cofinale de $E$. Soit donc $e = (J, N) \in E$ quelconque.
D'après le Lemme \ref{lemma-flat-factors-fp}, il existe un module libre de type fini $F$ et des applications
$h: R^J/N \to F$ et $g: F \to M$ telles que l'application naturelle
$f_e: R^J/N \to M$ se factorise sous la forme $R^J/N \xrightarrow{h} F
\xrightarrow{g} M$. Nous allons réaliser $F$ comme $M_{e'}$ pour un certain $e' \geq
e$.

\medskip\noindent
Soit $\{ b_1, \ldots, b_n \}$ une base finie de $F$. Choisissons $n$ éléments distincts
$i_1, \ldots, i_n \in I$ tels que $i_{\ell} \notin J$ pour tout $\ell$,
et tels que l'image de $x_{i_{\ell}}$ par $f: R^I \to M$ soit égale
à l'image de $b_{\ell}$ par $g: F \to M$. C'est possible puisque
tout élément de $M$ peut s'écrire sous la forme $f(x_i)$ pour une infinité d'éléments
distincts $i \in I$ (par le choix de $I$). Posons maintenant
$J' = J \cup \{i_1, \ldots , i_n \}$, et définissons $R^{J'}
\to F$ par $x_i \mapsto h(x_i)$ pour $i \in J$ et $x_{i_{\ell}} \mapsto
b_{\ell}$ pour $\ell = 1, \ldots, n$. Soit $N' = \Ker(R^{J'} \to F)$.
Remarquons que :
\begin{enumerate}
\item Le carré
$$
\xymatrix{
R^{J'} \ar[r] \ar@{^{(}->}[d] & F \ar[d]^{g} \\
R^{I} \ar[r]_{f} & M
}
$$
est commutatif,
%$R^{J'} \to F$ factors $f: R^I \to M$,
donc $N' \subset K = \Ker(f)$ ;
\item $R^{J'} \to F$ est une surjection sur un module libre de type fini, elle
est donc scindée, et $N'$ est de type fini ;
\item $J \subset J'$ et $N \subset N'$.
\end{enumerate}
D'après (1) et (2), $e' = (J', N')$ appartient à $E$ ; d'après (3), $e' \geq e$ ; et, par
construction, $M_{e'} = R^{J'}/N' \cong F$ est libre.
\end{proof}

\section{Morphismes de modules universellement injectifs}
\label{section-universally-injective}

\noindent
Nous étudions maintenant les morphismes de modules universellement injectifs, qui
sont en un certain sens complémentaires des modules plats (voir le lemme
\ref{lemma-flat-universally-injective}). Nous suivons la thèse de Lazard
\cite{Autour} ; voir aussi \cite{Lam}.

\begin{definition}
\label{definition-universally-injective}
Soit $f: M \to N$ un morphisme de $R$-modules. On dit que $f$ est
{\it universellement injectif} si, pour tout $R$-module $Q$, le morphisme $f
\otimes_R \text{id}_Q: M \otimes_R Q \to N \otimes_R Q$
est injectif. Une suite $0 \to M_1 \to M_2 \to M_3
\to 0$ de $R$-modules est dite {\it universellement exacte} si elle est exacte
et si $M_1 \to M_2$ est universellement injectif.
\end{definition}

\begin{example}
\label{example-universally-exact}
Exemples de suites universellement exactes.
\begin{enumerate}
\item Une suite exacte courte scindée est universellement exacte, puisque le produit tensoriel
commute aux sommes directes.
\item La colimite d'un système inductif de suites universellement exactes est
universellement exacte. Cela résulte de l'exactitude des colimites filtrantes et
du fait que le produit tensoriel commute aux colimites. En particulier, la
colimite d'un système inductif de suites exactes scindées est universellement exacte. Nous
verrons plus bas que, réciproquement, toute suite universellement exacte s'obtient de
cette manière.
\end{enumerate}
\end{example}

\noindent
Nous donnons maintenant une liste de critères pour qu'une suite exacte courte soit universellement
exacte. Ils sont analogues aux critères de platitude donnés plus haut. Les assertions (3)--(6)
ci-dessous correspondent respectivement aux critères de platitude donnés dans les
lemmes \ref{lemma-flat-eq}, \ref{lemma-flat-factors-free},
\ref{lemma-flat-surjective-hom} et le
théorème \ref{theorem-lazard}.

\begin{theorem}
\label{theorem-universally-exact-criteria}
Soit
$$
0 \to M_1 \xrightarrow{f_1} M_2 \xrightarrow{f_2} M_3 \to 0
$$
une suite exacte de $R$-modules. Les assertions suivantes sont équivalentes :
\begin{enumerate}
\item La suite $0 \to M_1 \to M_2 \to M_3
\to 0$ est universellement exacte.
\item Pour tout $R$-module de présentation finie $Q$, la suite
$$
0 \to M_1 \otimes_R Q \to M_2 \otimes_R Q \to
M_3 \otimes_R Q \to 0
$$
est exacte.
\item Étant donnés des éléments $x_i \in M_1$ $(i = 1, \ldots, n)$, $y_j \in M_2$ $(j = 1,
\ldots, m)$ et $a_{ij} \in R$ $(i = 1, \ldots, n, j = 1, \ldots, m)$ tels que,
pour tout $i$,
$$
f_1(x_i) = \sum\nolimits_j a_{ij} y_j,
$$
il existe $z_j \in M_1$ $(j =1, \ldots, m)$ tels que, pour tout $i$,
$$
x_i = \sum\nolimits_j a_{ij} z_j .
$$
\item Étant donné un diagramme commutatif de morphismes de $R$-modules
$$
\xymatrix{
R^n \ar[r] \ar[d] &  R^m \ar[d] \\
M_1 \ar[r]^{f_1}        &  M_2
}
$$
où $m$ et $n$ sont des entiers, il existe un morphisme $R^m \to M_1$ rendant
commutatif le triangle supérieur.
\item Pour tout $R$-module de présentation finie $P$, le morphisme de $R$-modules
$\Hom_R(P, M_2) \to \Hom_R(P, M_3)$ est surjectif.
\item La suite $0 \to M_1 \to M_2 \to M_3
\to 0$ est la colimite d'un système inductif de suites exactes scindées de
la forme
$$
0 \to M_{1} \to M_{2, i} \to M_{3, i} \to 0
$$
où les $M_{3, i}$ sont de présentation finie.
\end{enumerate}
\end{theorem}

\begin{proof}
Il est clair que (1) implique (2).

\medskip\noindent
Montrons maintenant que (2) implique (3). Soient $f_1(x_i) = \sum_j a_{ij} y_j$ des relations
comme en (3). Soient $(d_j)$ une base de $R^m$, $(e_i)$ une base de $R^n$, et
$R^m \to R^n$ le morphisme donné par $d_j \mapsto \sum_i a_{ij} e_i$.
Soit $Q$ le conoyau de $R^m \to R^n$. En tensorisant
$R^m \to R^n \to Q \to 0$ par le morphisme $f_1: M_1 \to M_2$, on obtient un
diagramme commutatif
$$
\xymatrix{
M_1^{\oplus m} \ar[r] \ar[d] & M_1^{\oplus n} \ar[r] \ar[d] & M_1 \otimes_R Q
\ar[r] \ar[d] & 0 \\
M_2^{\oplus m} \ar[r] & M_2^{\oplus n} \ar[r] & M_2 \otimes_R Q \ar[r] & 0
}
$$
où $M_1^{\oplus m} \to M_1^{\oplus n}$ est donné par
$$
(z_1, \ldots, z_m) \mapsto
(\sum\nolimits_j a_{1j} z_j, \ldots, \sum\nolimits_j a_{nj} z_j),
$$
et où $M_2^{\oplus m} \to M_2^{\oplus n}$ est défini de même. Nous voulons
montrer que $x = (x_1, \ldots, x_n) \in M_1^{\oplus n}$ appartient à l'image de $M_1^{\oplus
m} \to M_1^{\oplus n}$. Par (2), le morphisme $M_1 \otimes Q \to M_2
\otimes Q$ est injectif ; par exactitude de la ligne supérieure, il suffit donc de montrer que
$x$ s'envoie sur $0$ dans $M_2 \otimes Q$ ; par exactitude de la ligne inférieure, il suffit ainsi
de montrer que l'image de $x$ dans $M_2^{\oplus n}$ appartient à l'image de
$M_2^{\oplus m} \to M_2^{\oplus n}$. Cela résulte de l'hypothèse.

\medskip\noindent
La condition (4) n'est que la traduction de (3) sous forme de diagramme.

\medskip\noindent
Montrons maintenant que (4) implique (5). Soit $\varphi : P \to M_3$ un morphisme issu d'un
$R$-module de présentation finie $P$. Nous devons montrer que $\varphi$ se relève en un morphisme
$P \to M_2$. Choisissons une présentation de $P$,
$$
R^n \xrightarrow{g_1} R^m \xrightarrow{g_2} P \to 0.
$$
La liberté de $R^n$ et de $R^m$ permet de construire $h_2: R^m \to M_2$
puis $h_1: R^n \to M_1$ de sorte que le diagramme suivant soit commutatif :
$$
\xymatrix{
         & R^n \ar[r]^{g_1} \ar[d]^{h_1} & R^m \ar[r]^{g_2} \ar[d]^{h_2} & P
\ar[r] \ar[d]^{\varphi} & 0 \\
0 \ar[r] & M_1 \ar[r]^{f_1} & M_2 \ar[r]^{f_2} & M_3 \ar[r] & 0 .
}
$$
Par (4), il existe un morphisme $k_1: R^m \to M_1$ tel que $k_1 \circ g_1 =
h_1$. Définissons maintenant $h'_2: R^m \to M_2$ par $h_2' = h_2 - f_1 \circ k_1$.
Alors
$$
h'_2 \circ g_1 = h_2 \circ g_1 - f_1 \circ k_1 \circ g_1 =
h_2 \circ g_1 - f_1 \circ h_1 = 0 .
$$
Ainsi, par passage au quotient, $h'_2$ définit un morphisme $\varphi': P \to
M_2$ tel que $\varphi' \circ g_2 = h_2'$. Sous forme de diagramme, on a
$$
\xymatrix{
R^m \ar[r]^{g_2} \ar[d]_{h'_2} & P \ar[d]^{\varphi} \ar[dl]_{\varphi'} \\
M_2 \ar[r]^{f_2} & M_3.
}
$$
où le triangle supérieur est commutatif. Nous affirmons que $\varphi'$ est le relèvement voulu,
c'est-à-dire que $f_2 \circ \varphi' = \varphi$. D'après les définitions, on a
$$
f_2 \circ \varphi' \circ g_2 = f_2 \circ h'_2 =
f_2 \circ h_2 - f_2 \circ f_1 \circ k_1 = f_2 \circ h_2 =
\varphi \circ g_2.
$$
Comme $g_2$ est surjectif, ceci achève la démonstration.

\medskip\noindent
Montrons maintenant que (5) implique (6). Écrivons $M_{3}$ comme colimite d'un système inductif
de modules de présentation finie $M_{3, i}$, voir le
lemme \ref{lemma-module-colimit-fp}. Soit $M_{2, i}$ le produit fibré
de $M_{3, i}$ et $M_{2}$ au-dessus de $M_{3}$ ; par définition, c'est le sous-module de
$M_2 \times M_{3, i}$ formé des éléments dont les deux projections sur $M_3$
sont égales. Soit $M_{1, i}$ le noyau de la projection
$M_{2, i} \to M_{3, i}$. On obtient alors un système inductif de suites exactes
$$
0 \to M_{1, i} \to M_{2, i} \to M_{3, i} \to 0,
$$
et, pour chaque $i$, un morphisme de suites exactes
$$
\xymatrix{
0  \ar[r] & M_{1, i} \ar[d] \ar[r]  &  M_{2, i}  \ar[r] \ar[d] & M_{3, i} \ar[d]
\ar[r] & 0 \\
0  \ar[r] & M_{1} \ar[r]  &  M_{2}  \ar[r] & M_{3} \ar[r] & 0
}
$$
compatible au système inductif. Il résulte de la définition du produit fibré
$M_{2, i}$ que le morphisme $M_{1, i} \to M_1$ est un isomorphisme.
 Par (5), il existe un morphisme $M_{3, i} \to M_{2}$ relevant $M_{3, i} \to
M_3$ et, par la propriété universelle du produit fibré, celui-ci donne lieu à une
section de $M_{2, i} \to M_{3, i}$. Ainsi, les suites
$$
0 \to M_{1, i} \to M_{2, i} \to M_{3, i} \to 0
$$
sont scindées. En passant à la colimite, on obtient un diagramme commutatif
$$
\xymatrix{
0  \ar[r] & \colim M_{1, i} \ar[d]^{\cong} \ar[r]  &  \colim M_{2, i}
 \ar[r]
\ar[d] & \colim M_{3, i} \ar[d]^{\cong} \ar[r] & 0 \\
0  \ar[r] & M_{1} \ar[r]  &  M_{2}  \ar[r] & M_{3} \ar[r] & 0
}
$$
à lignes exactes et dont les flèches verticales extérieures sont des isomorphismes. Ainsi, $\colim
M_{2, i}
\to M_2$ est également un isomorphisme, et (6) est satisfaite.

\medskip\noindent
La condition (6) implique (1) d'après
l'exemple \ref{example-universally-exact} (2).
\end{proof}

\noindent
Le théorème précédent montre qu'une suite universellement exacte est toujours une
colimite de suites exactes courtes scindées. Si le conoyau d'un morphisme universellement
injectif est de présentation finie, alors le morphisme lui-même est en fait scindé :

\begin{lemma}
\label{lemma-universally-exact-split}
Soit
$$
0 \to M_1 \to M_2 \to M_3 \to 0
$$
une suite exacte de $R$-modules. Supposons que $M_3$ soit de présentation finie.
Alors
$$
0 \to M_1 \to M_2 \to M_3 \to 0
$$
est universellement exacte si et seulement si elle est scindée.
\end{lemma}

\begin{proof}
Une suite exacte courte scindée est toujours universellement exacte, voir
l'exemple \ref{example-universally-exact}.
Réciproquement, si la suite est universellement exacte, alors le
théorème \ref{theorem-universally-exact-criteria} (5),
appliqué à $P = M_3$, montre que le morphisme $M_2 \to M_3$ admet une section.
\end{proof}

\noindent
Le lemme suivant montre en quel sens les morphismes universellement injectifs sont complémentaires des
modules plats.

\begin{lemma}
\label{lemma-flat-universally-injective}
Soit $M$ un $R$-module. Alors $M$ est plat si et seulement si toute suite exacte
de $R$-modules
$$
0 \to M_1 \to M_2 \to M \to 0
$$
est universellement exacte.
\end{lemma}

\begin{proof}
Cela résulte du lemme \ref{lemma-flat-surjective-hom} et du
théorème \ref{theorem-universally-exact-criteria} (5).
\end{proof}

\begin{example}
\label{example-universally-exact-non-split-non-flat}
Suites universellement exactes non scindées et non plates.
\begin{enumerate}
\item Malgré le
lemme \ref{lemma-universally-exact-split},
il peut exister une suite exacte courte de $R$-modules
$$
0 \to M_1 \to M_2 \to M_3 \to 0
$$
qui soit universellement exacte sans être scindée. Par exemple, prenons $R = \mathbf{Z}$,
$M_1 = \bigoplus_{n=1}^{\infty} \mathbf{Z}$, $M_{2} = \prod_{n =
1}^{\infty} \mathbf{Z}$, et pour $M_{3}$ le conoyau de l'inclusion $M_1
\to M_2$. Alors $M_1, M_2, M_3$ sont tous plats, car ils sont sans torsion
(Compléments d'algèbre, lemme \ref{more-algebra-lemma-dedekind-torsion-free-flat}) ;
par le lemme \ref{lemma-flat-universally-injective},
$$
0 \to M_1 \to M_2 \to M_3 \to 0
$$
est donc universellement exacte. Il ne peut toutefois exister de section $s: M_3 \to
M_2$. En effet, si $x$ est l'image de $(2, 2^2, 2^3, \ldots) \in M_2$ dans $M_3$,
alors tout morphisme de modules $s: M_3 \to M_2$ doit annuler $x$. En effet, $x
\in 2^n M_3$ pour tout $n \geq 1$ ; ainsi, $s(x)$ est divisible par $2^n$ pour tout $n
\geq 1$ et doit donc être $0$.
\item Malgré le lemme \ref{lemma-flat-universally-injective}, il peut exister
une suite exacte courte de $R$-modules
$$
0 \to M_1 \to M_2 \to M_3 \to 0
$$
qui soit universellement exacte alors que $M_1, M_2, M_3$ sont tous non plats. En effet, si $M$
est un module non plat quelconque, il suffit de prendre la suite exacte scindée
$$
0 \to M \to M \oplus M \to M \to 0.
$$
Par exemple, pour $R = \mathbf{Z}$, on peut prendre pour $M$ un module de torsion quelconque.
\item En prenant la somme directe d'une suite exacte comme en (1) et d'une autre comme en (2),
on obtient une suite exacte courte de $R$-modules
$$
0 \to M_1 \to M_2 \to M_3 \to 0
$$
qui est universellement exacte, non scindée, et telle que
$M_1, M_2, M_3$ sont tous non plats.
\end{enumerate}
\end{example}

\begin{lemma}
\label{lemma-ui-flat-domain}
Soit $0 \to M_1 \to M_2 \to M_3 \to 0$ une suite universellement exacte
de $R$-modules, et supposons $M_2$ plat.
Alors $M_1$ et $M_3$ sont plats.
\end{lemma}

\begin{proof}
Soit $0 \to N \to N' \to N'' \to 0$ une suite exacte courte de
$R$-modules. Considérons le diagramme commutatif
$$
\xymatrix{
M_1 \otimes_R N \ar[r] \ar[d] &
M_2 \otimes_R N \ar[r] \ar[d] &
M_3 \otimes_R N \ar[d] \\
M_1 \otimes_R N' \ar[r] \ar[d] &
M_2 \otimes_R N' \ar[r] \ar[d] &
M_3 \otimes_R N' \ar[d] \\
M_1 \otimes_R N'' \ar[r] &
M_2 \otimes_R N'' \ar[r] &
M_3 \otimes_R N''
}
$$
(nous avons omis les $0$ du bord).
Par hypothèse, les lignes sont des suites exactes courtes et la flèche
$M_2 \otimes N \to M_2 \otimes N'$ est injective. Cela implique clairement
que $M_1 \otimes N \to M_1 \otimes N'$ est injective, et l'on voit que $M_1$
est plat. En particulier, les colonnes de gauche et du milieu sont des suites
exactes courtes. Une chasse au diagramme montre alors que la flèche
$M_3 \otimes N \to M_3 \otimes N'$ est injective. Ainsi, $M_3$ est plat.
\end{proof}

\begin{lemma}
\label{lemma-universally-injective-tensor}
Soit $R$ un anneau.
Soit $M \to M'$ un morphisme universellement injectif de $R$-modules.
Alors, pour tout $R$-module $N$, le morphisme $M \otimes_R N \to M' \otimes_R N$
est universellement injectif.
\end{lemma}

\begin{proof}
Démonstration omise.
\end{proof}

\begin{lemma}
\label{lemma-composition-universally-injective}
Soit $R$ un anneau. Toute composée de morphismes universellement injectifs de
$R$-modules est universellement injective.
\end{lemma}

\begin{proof}
Démonstration omise.
\end{proof}

\begin{lemma}
\label{lemma-universally-injective-permanence}
Soit $R$ un anneau. Soient $M \to M'$ et $M' \to M''$ des morphismes de $R$-modules.
Si leur composée $M \to M''$ est universellement injective, alors
$M \to M'$ est universellement injectif.
\end{lemma}

\begin{proof}
Démonstration omise.
\end{proof}

\begin{lemma}
\label{lemma-universally-injective-check-stalks}
Soit $R \to S$ un homomorphisme d'anneaux.
Soit $M \to M'$ un morphisme de $S$-modules.
Les assertions suivantes sont équivalentes :
\begin{enumerate}
\item $M \to M'$ est universellement injectif comme morphisme de $R$-modules,
\item pour tout idéal premier $\mathfrak q$ de $S$, le morphisme
$M_{\mathfrak q} \to M'_{\mathfrak q}$ est universellement injectif
comme morphisme de $R$-modules,
\item pour tout idéal maximal $\mathfrak m$ de $S$, le morphisme
$M_{\mathfrak m} \to M'_{\mathfrak m}$ est universellement injectif
comme morphisme de $R$-modules,
\item pour tout idéal premier $\mathfrak q$ de $S$, le morphisme
$M_{\mathfrak q} \to M'_{\mathfrak q}$ est universellement injectif
comme morphisme de $R_{\mathfrak p}$-modules, où $\mathfrak p$ est l'image
réciproque de $\mathfrak q$ dans $R$, et
\item pour tout idéal maximal $\mathfrak m$ de $S$, le morphisme
$M_{\mathfrak m} \to M'_{\mathfrak m}$ est universellement injectif
comme morphisme de $R_{\mathfrak p}$-modules, où $\mathfrak p$ est l'image
réciproque de $\mathfrak m$ dans $R$.
\end{enumerate}
\end{lemma}

\begin{proof}
Soit $N$ un $R$-module. Soit $\mathfrak q$ un idéal premier de $S$ au-dessus de
l'idéal premier $\mathfrak p$ de $R$. On a alors
$$
(M \otimes_R N)_{\mathfrak q} =
M_{\mathfrak q} \otimes_R N =
M_{\mathfrak q} \otimes_{R_{\mathfrak p}} N_{\mathfrak p}.
$$
La même chose vaut en outre pour $M'$, et la localisation est exacte.
De plus, si $N$ est un $R_{\mathfrak p}$-module, alors $N_{\mathfrak p} = N$.
Ces remarques permettent de démontrer directement les équivalences.

\medskip\noindent
Supposons par exemple que (5) soit satisfaite. Soit
$K = \Ker(M \otimes_R N \to M' \otimes_R N)$. D'après les remarques
précédentes, $K_{\mathfrak m} = 0$ pour tout idéal maximal $\mathfrak m$
de $S$. Ainsi $K = 0$ d'après le
lemme \ref{lemma-characterize-zero-local}.
L'assertion (1) est donc satisfaite. Réciproquement, supposons (1) satisfaite. Prenons un
$\mathfrak q \subset S$ quelconque au-dessus de $\mathfrak p \subset R$.
Prenons un module $N$ quelconque sur $R_{\mathfrak p}$. Par hypothèse,
$\Ker(M \otimes_R N \to M' \otimes_R N) = 0$.
Les formules ci-dessus et le fait que $N = N_{\mathfrak p}$
montrent alors que
$\Ker(M_{\mathfrak q} \otimes_{R_{\mathfrak p}} N \to
M'_{\mathfrak q} \otimes_{R_{\mathfrak p}} N) = 0$. Autrement dit,
(4) est satisfaite. Bien sûr, (4) $\Rightarrow$ (5) est immédiat. Ainsi,
(1), (4) et (5) sont toutes équivalentes.
Nous omettons la démonstration des autres équivalences.
\end{proof}

\begin{lemma}
\label{lemma-universally-injective-localize}
Soit $\varphi : A \to B$ un homomorphisme d'anneaux. Soient $S \subset A$ et
$S' \subset B$ des parties multiplicatives telles que $\varphi(S) \subset S'$.
Soit $M \to M'$ un morphisme de $B$-modules.
\begin{enumerate}
\item Si $M \to M'$ est universellement injectif comme morphisme de $A$-modules,
alors $(S')^{-1}M \to (S')^{-1}M'$ est universellement injectif comme morphisme de
$A$-modules et comme morphisme de $S^{-1}A$-modules.
\item Si $M$ et $M'$ sont des $(S')^{-1}B$-modules, alors $M \to M'$
est universellement injectif comme morphisme de $A$-modules si et seulement s'il
est universellement injectif comme morphisme de $S^{-1}A$-modules.
\end{enumerate}
\end{lemma}

\begin{proof}
On peut démontrer ceci à l'aide du
lemme \ref{lemma-universally-injective-check-stalks},
mais on peut aussi le démontrer directement comme suit.
Supposons $M \to M'$ universellement injectif comme morphisme de $A$-modules.
Soit $Q$ un $A$-module. Alors $Q \otimes_A M \to Q \otimes_A M'$
est injectif. L'exactitude de la localisation montre que
$(S')^{-1}(Q \otimes_A M) \to (S')^{-1}(Q \otimes_A M')$ est injectif.
Comme $(S')^{-1}(Q \otimes_A M) = Q \otimes_A (S')^{-1}M$, et de même pour $M'$,
on voit que
$Q \otimes_A (S')^{-1}M \to Q \otimes_A (S')^{-1}M'$ est injectif ; ainsi,
$(S')^{-1}M \to (S')^{-1}M'$ est universellement injectif comme morphisme de
$A$-modules. Ceci démontre la première partie de (1).
Pour établir (2), on peut utiliser les deux faits suivants : (a) si $Q$ est un
$S^{-1}A$-module, alors $Q \otimes_A S^{-1}A = Q$, c'est-à-dire que tensoriser
par $Q$ sur $A$ revient à tensoriser par $Q$ sur $S^{-1}A$ ;
(b) si $M$ est un $A$-module quelconque sur lequel les éléments de $S$ sont inversibles,
alors $M \otimes_A Q = M \otimes_{S^{-1}A} S^{-1}Q$.
L'assertion (2) en résulte immédiatement.
\end{proof}

\begin{lemma}
\label{lemma-check-universally-injective-into-flat}
Soit $R$ un anneau et soit $M \to M'$ un morphisme de $R$-modules.
Si $M'$ est plat, alors $M \to M'$ est universellement injectif si
et seulement si $M/IM \to M'/IM'$ est injectif pour tout idéal de type
fini $I$ de $R$.
\end{lemma}

\begin{proof}
Il suffit de montrer que $M \otimes_R Q \to M' \otimes_R Q$ est
injectif pour tout $R$-module de type fini $Q$, voir le
théorème \ref{theorem-universally-exact-criteria}.
Alors $Q$ possède une filtration finie
$0 = Q_0 \subset Q_1 \subset \ldots \subset Q_n = Q$
par des sous-modules dont les sous-quotients
sont isomorphes à des modules cycliques $R/I_i$, voir le
lemme \ref{lemma-trivial-filter-finite-module}.
Comme $M'$ est plat, on obtient une filtration
$$
\xymatrix{
M \otimes Q_1 \ar[r] \ar[d] &
M \otimes Q_2 \ar[r] \ar[d] &
\ldots \ar[r] &
M \otimes Q \ar[d] \\
M' \otimes Q_1 \ar@{^{(}->}[r] &
M' \otimes Q_2 \ar@{^{(}->}[r] &
\ldots \ar@{^{(}->}[r] &
M' \otimes Q
}
$$
de $M' \otimes_R Q$ par les sous-modules $M' \otimes_R Q_i$, dont les quotients
successifs sont $M' \otimes_R R/I_i = M'/I_iM'$. Une récurrence simple
montre qu'il suffit de vérifier l'injectivité de $M/I_i M \to M'/I_i M'$.
Notons que l'ensemble des idéaux de type fini $I'_i \subset I_i$
est filtrant. Ainsi, $M/I_iM = \colim M/I'_iM$ est une colimite
filtrante, et de même pour $M'$ ; les morphismes $M/I'_iM \to M'/I'_i M'$ sont
injectifs par hypothèse et, puisque les colimites filtrantes sont exactes
(lemme \ref{lemma-directed-colimit-exact}), on conclut.
\end{proof}

\begin{lemma}
\label{lemma-base-change-along-universally-injective-ring-map}
Soit $R \to S$ un homomorphisme d'anneaux universellement injectif
comme morphisme de $R$-modules. Alors le foncteur $M \mapsto M \otimes_R S$
sur les $R$-modules reflète les morphismes injectifs, surjectifs et les isomorphismes.
\end{lemma}

\begin{proof}
Soit $M \to N$ un morphisme de $R$-modules, de noyau $K$ et de conoyau $Q$.
Si $M \otimes_R S \to N \otimes_R S$ est injectif, alors
l'image de $K \otimes_R S \to M \otimes_R S$ est nulle.
Comme $K \subset K \otimes_R S$ et $M \subset M \otimes_R S$,
on en déduit que $K = 0$.
D'autre part, si $M \otimes_R S \to N \otimes_R S$ est surjectif,
alors $Q \otimes_R S$ est nul (par exactitude à droite du produit tensoriel).
Ainsi, $Q \subset Q \otimes_R S$ est également nul.
\end{proof}

\begin{lemma}
\label{lemma-faithfully-flat-universally-injective}
Soit $R \to S$ un homomorphisme d'anneaux fidèlement plat.
Alors $R \to S$ est universellement injectif comme morphisme de $R$-modules.
En particulier, $R \cap IS = I$ pour tout idéal $I \subset R$.
\end{lemma}

\begin{proof}
Soit $N$ un $R$-module. Nous devons montrer que $N \to N \otimes_R S$ est
injectif. Comme $S$ est fidèlement plat comme $R$-module, il suffit de le démontrer
après tensorisation par $S$. Il suffit donc de montrer que
$N \otimes_R S \to N \otimes_R S \otimes_R S$,
$n \otimes s \mapsto n \otimes 1 \otimes s$ est injectif. C'est bien le cas
car il existe une rétraction, à savoir
$n \otimes s \otimes s' \mapsto n \otimes ss'$.
\end{proof}




\section{Descente pour les modules projectifs de type fini}
\label{section-finite-projective}

\noindent
Dans cette section, nous donnons une démonstration élémentaire du fait que la propriété
d'être un module projectif {\it de type fini} descend le long des homomorphismes
d'anneaux fidèlement plats. La démonstration ne s'applique plus sans l'hypothèse de finitude.
Toutefois, la méthode préfigure celle que nous utiliserons pour démontrer la descente
de la propriété d'être un module projectif {\it dénombrablement engendré} ; voir
les remarques à la fin de cette section.

\begin{lemma}
\label{lemma-finite-projective-again}
Soit $M$ un $R$-module. Alors $M$ est projectif de type fini si et seulement si $M$ est
de présentation finie et plat.
\end{lemma}

\begin{proof}
Ceci fait partie du
lemme \ref{lemma-finite-projective}.
À ce stade, nous pouvons toutefois donner une démonstration plus élégante de l'implication
(1) $\Rightarrow$ (2) de ce lemme, comme suit.
Si $M$ est de présentation finie et plat, choisissons une surjection
$R^n \to M$. D'après le
lemme \ref{lemma-flat-surjective-hom},
appliqué à $P = M$, le morphisme $R^n \to M$ admet une section.
Ainsi, $M$ est un facteur direct d'un module libre, et est donc projectif.
\end{proof}

\noindent
Voici quelques propriétés des modules qui descendent.

\begin{lemma}
\label{lemma-descend-properties-modules}
Soit $R \to S$ un homomorphisme d'anneaux fidèlement plat.
Soit $M$ un $R$-module. Alors
\begin{enumerate}
\item si le $S$-module $M \otimes_R S$ est de type fini, alors
$M$ est de type fini,
\item si le $S$-module $M \otimes_R S$ est de présentation finie, alors
$M$ est de présentation finie,
\item si le $S$-module $M \otimes_R S$ est plat, alors
$M$ est plat, et
\item ajouter ici d'autres propriétés au besoin.
\end{enumerate}
\end{lemma}

\begin{proof}
Supposons $M \otimes_R S$ de type fini. Soient $y_1, \ldots, y_m$ des générateurs
de $M \otimes_R S$ sur $S$. Écrivons $y_j = \sum x_i \otimes f_i$ pour certains
$x_1, \ldots, x_n \in M$. On voit alors que le morphisme
$\varphi : R^{\oplus n} \to M$
a la propriété que
$\varphi \otimes \text{id}_S : S^{\oplus n} \to M \otimes_R S$
est surjectif. Comme $R \to S$ est fidèlement plat, on en déduit que
$\varphi$ est surjectif et que $M$ est de type fini.

\medskip\noindent
Supposons $M \otimes_R S$ de présentation finie. D'après (1),
$M$ est de type fini. Choisissons une surjection $R^{\oplus n} \to M$ et
notons $K$ son noyau. Comme $R \to S$ est plat, $K \otimes_R S$
est le noyau du changement de base $S^{\oplus n} \to M \otimes_R S$.
Puisque $M \otimes_R S$ est de présentation finie, on en déduit que $K \otimes_R S$
est de type fini. D'après (1), $K$ est donc de type fini,
et $M$ est par conséquent de présentation finie.

\medskip\noindent
L'assertion (3) est le
lemme \ref{lemma-flatness-descends}.
\end{proof}

\begin{proposition}
\label{proposition-ffdescent-finite-projectivity}
Soit $R \to S$ un homomorphisme d'anneaux fidèlement plat. Soit $M$ un $R$-module.
Si le $S$-module $M \otimes_R S$ est projectif de type fini, alors $M$ est projectif
de type fini.
\end{proposition}

\begin{proof}
Cela résulte des
lemmes \ref{lemma-finite-projective-again} et
\ref{lemma-descend-properties-modules}.
\end{proof}

\noindent
Les quelques sections suivantes visent à supprimer l'hypothèse de finitude en recourant au
d\'evissage pour se ramener au cas dénombrablement engendré. Dans ce dernier
cas, la stratégie consiste à trouver une caractérisation des modules projectifs
dénombrablement engendrés analogue au
lemme \ref{lemma-finite-projective-again},
puis à démontrer directement que cette caractérisation descend. Pour cela, nous
introduisons la notion de module de Mittag-Leffler et démontrons qu'un module
$M$ dénombrablement engendré est projectif si et seulement s'il est plat et
de Mittag-Leffler (théorème \ref{theorem-projectivity-characterization}). Lorsque $M$
est de type fini, cet énoncé se réduit au lemme
\ref{lemma-finite-projective-again} (puisque, d'après
l'exemple \ref{example-ML} (1),
un module de type fini est de Mittag-Leffler si et seulement s'il est
de présentation finie).

\section{D\'evissage transfini des modules}
\label{section-transfinite-devissage}

\noindent
Dans cette section, nous introduisons une technique de d\'evissage permettant de décomposer un module
en somme directe. Le résultat principal affirme qu'un module projectif est une somme directe
de modules dénombrablement engendrés (théorème \ref{theorem-projective-direct-sum}
ci-dessous). Nous suivons \cite{Kaplansky}.


\begin{definition}
\label{definition-devissage}
Soit $M$ un $R$-module. Un {\it d\'evissage par sommes directes} de $M$ est une famille
de sous-modules $(M_{\alpha})_{\alpha \in S}$, indexée par un ordinal $S$ et
croissante (pour l'inclusion), telle que :
\begin{enumerate}
\item[(0)] $M_0 = 0$ ;
\item[(1)] $M = \bigcup_{\alpha} M_{\alpha}$ ;
\item[(2)] si $\alpha \in S$ est un ordinal limite, alors $M_{\alpha} =
\bigcup_{\beta < \alpha} M_{\beta}$ ;
\item[(3)] si $\alpha + 1 \in S$, alors $M_{\alpha}$ est un facteur direct de
$M_{\alpha + 1}$.
\end{enumerate}
Si de plus
\begin{enumerate}
\item[(4)] $M_{\alpha + 1}/M_{\alpha}$ est dénombrablement engendré pour
$\alpha + 1 \in S$,
\end{enumerate}
alors $(M_{\alpha})_{\alpha \in S}$ est appelé un {\it d\'evissage de Kaplansky}
de $M$.
\end{definition}

\noindent
La terminologie est justifiée par le lemme suivant.

\begin{lemma}
\label{lemma-direct-sum-devissage}
Soit $M$ un $R$-module. Si $(M_{\alpha})_{\alpha \in S}$ est un d\'evissage
par sommes directes de $M$, alors
$M \cong \bigoplus_{\alpha + 1 \in S} M_{\alpha + 1}/M_{\alpha}$.
\end{lemma}

\begin{proof}
Par la propriété (3) d'un d\'evissage par sommes directes, il existe une inclusion
$M_{\alpha + 1}/M_{\alpha} \to M$ pour tout $\alpha \in S$. Considérons le
morphisme
$$
f : \bigoplus\nolimits_{\alpha + 1\in S} M_{\alpha + 1}/M_{\alpha} \to M
$$
donné par la somme de ces inclusions.
Considérons en outre les restrictions
$$
f_{\beta} :
\bigoplus\nolimits_{\alpha + 1 \leq \beta} M_{\alpha + 1}/M_{\alpha}
\longrightarrow
M
$$
pour $\beta\in S$. Une récurrence transfinie sur $S$ montre que l'image de
$f_{\beta}$ est $M_{\beta}$. Pour $\beta=0$, cela résulte de $(0)$. Si $\beta+1$
est un ordinal successeur et si l'assertion vaut pour $\beta$, elle vaut alors pour
$\beta + 1$ par (3). Si enfin $\beta$ est un ordinal limite et que l'assertion vaut pour
$\alpha < \beta$, elle vaut pour $\beta$ par (2). Ainsi, $f$ est surjectif
par (1).

\medskip\noindent
Une récurrence transfinie sur $S$ montre aussi que les restrictions $f_{\beta}$
sont injectives. Pour $\beta = 0$, c'est vrai. Si $\beta+1$ est un
ordinal successeur et que $f_{\beta}$ est injectif, soit $x$ un élément du noyau ;
écrivons $x = (x_{\alpha + 1})_{\alpha + 1 \leq \beta + 1}$ selon ses
composantes $x_{\alpha + 1} \in M_{\alpha + 1}/M_{\alpha}$. Par la propriété (3) et
le fait que l'image de $f_{\beta}$ est $M_{\beta}$, tant
$(x_{\alpha + 1})_{\alpha + 1 \leq \beta}$ que $x_{\beta + 1}$ s'envoient sur $0$.
Ainsi, $x_{\beta+1} = 0$ et, puisque la restriction
$f_{\beta}$ est injective, on a aussi $x_{\alpha + 1} = 0$
pour tout $\alpha + 1 \leq \beta$. Ainsi $x = 0$ et $f_{\beta+1}$ est injectif.
Si $\beta$ est un ordinal limite, considérons un élément $x$ du noyau. Alors $x$
appartient déjà au domaine de $f_{\alpha}$ pour un certain $\alpha < \beta$.
Ainsi $x = 0$, ce qui achève la récurrence.
On conclut que $f$ est injectif, puisque $f_{\beta}$ l'est pour tout $\beta \in S$.
\end{proof}

\begin{lemma}
\label{lemma-Kaplansky-devissage}
Soit $M$ un $R$-module. Alors $M$ est une somme directe de $R$-modules
dénombrablement engendrés si et seulement s'il admet un d\'evissage de Kaplansky.
\end{lemma}

\begin{proof}
Le lemme précédent traite le sens « si ». Réciproquement, supposons $M =
\bigoplus_{i \in I} N_i$, où chaque $N_i$ est un $R$-module dénombrablement engendré.
Munissons $I$ d'un bon ordre afin de pouvoir l'identifier à un ordinal. Alors, en posant $M_i =
\bigoplus_{j < i} N_j$, on obtient un d\'evissage de Kaplansky $(M_i)_{i \in I}$ de
$M$.
\end{proof}

\begin{theorem}
\label{theorem-kaplansky-direct-sum}
Supposons que $M$ soit une somme directe de $R$-modules dénombrablement engendrés. Si $P$ est un
facteur direct de $M$, alors $P$ est lui aussi une somme directe de $R$-modules
dénombrablement engendrés.
\end{theorem}

\begin{proof}
Écrivons $M = P \oplus Q$. Nous allons construire un d\'evissage de Kaplansky
$(M_{\alpha})_{\alpha \in S}$ de $M$ qui, outre les propriétés de définition
(0)--(4), vérifie :
\begin{enumerate}
\item[(5)] Chaque $M_{\alpha}$ est un facteur direct de $M$ ;
\item[(6)] $M_{\alpha} = P_{\alpha} \oplus Q_{\alpha}$, où $P_{\alpha} =P
\cap M_{\alpha}$ et $Q_\alpha = Q \cap M_{\alpha}$.
\end{enumerate}
(Remarque : si les propriétés (0)--(2) sont satisfaites, alors la propriété (3) équivaut en fait à
la propriété (5).)

\medskip\noindent
Pour voir comment ceci implique le théorème, il suffit de montrer que
$(P_{\alpha})_{\alpha \in S}$ forme un d\'evissage de Kaplansky de $P$. Les propriétés
(0), (1) et (2) sont claires. D'après (5) et (6) pour $(M_{\alpha})$, chaque
$P_{\alpha}$ est un facteur direct de $M$. Comme $P_{\alpha} \subset P_{\alpha +
1}$, il s'ensuit que $P_{\alpha}$ est un facteur direct de $P_{\alpha + 1}$ ;
ainsi, (3) vaut pour $(P_{\alpha})$. Pour (4), remarquons que
$$
M_{\alpha + 1}/M_{\alpha} \cong P_{\alpha + 1}/P_{\alpha} \oplus
Q_{\alpha + 1}/Q_{\alpha},
$$
de sorte que $P_{\alpha + 1}/P_{\alpha}$ est dénombrablement engendré, puisque
$M_{\alpha + 1}/M_{\alpha}$ l'est.

\medskip\noindent
Il reste à construire les $M_{\alpha}$. Écrivons $M = \bigoplus_{i \in I} N_i$,
où chaque $N_i$ est un $R$-module dénombrablement engendré. Choisissons un bon ordre
sur $I$. Par récurrence transfinie, nous allons définir une famille croissante
de sous-modules $M_{\alpha}$ de $M$, un pour chaque ordinal $\alpha$, telle que
$M_{\alpha}$ soit la somme directe d'une certaine sous-famille des $N_i$.

\medskip\noindent
Pour $\alpha = 0$, posons $M_{0} = 0$. Si $\alpha$ est un ordinal limite et que
$M_{\beta}$ a été défini pour tout $\beta < \alpha$, posons $M_{\alpha}
= \bigcup_{\beta < \alpha} M_{\beta}$. Puisque chaque $M_{\beta}$, pour $\beta <
\alpha$, est la somme directe d'une sous-famille des $N_i$, il en va de même de
$M_{\alpha}$. Si $\alpha + 1$ est un ordinal successeur et que $M_{\alpha}$ a été
défini, définissons $M_{\alpha + 1}$ comme suit. Si $M_{\alpha} = M$, posons
$M_{\alpha + 1} = M$. Sinon, choisissons le plus petit $j \in I$ tel que $N_j$ ne soit
pas contenu dans $M_{\alpha}$. Nous allons construire une matrice infinie $(x_{mn}),
m, n = 1, 2, 3, \ldots$ telle que :
\begin{enumerate}
\item $N_j$ est contenu dans le sous-module de $M$ engendré par les entrées
$x_{mn}$ ;
\item si l'on écrit une entrée quelconque $x_{k\ell}$ selon ses composantes dans $P$ et
$Q$, $x_{k\ell} = y_{k\ell} + z_{k\ell}$, alors la matrice $(x_{mn})$
contient une famille de générateurs de chaque $N_i$ pour lequel $y_{k\ell}$ ou
$z_{k\ell}$ possède une composante non nulle.
\end{enumerate}
Nous définissons alors $M_{\alpha + 1}$ comme le sous-module de $M$ engendré par
$M_{\alpha}$ et tous les $x_{mn}$ ; d'après la propriété (2) de la matrice $(x_{mn})$,
$M_{\alpha + 1}$ sera la somme directe d'une certaine sous-famille des $N_i$.
Pour construire la matrice $(x_{mn})$, soit $x_{11}, x_{12}, x_{13}, \ldots$
une famille dénombrable de générateurs de $N_j$. Si ensuite
$x_{11} = y_{11} + z_{11}$ est la décomposition selon les composantes dans $P$ et
$Q$, soit $x_{21}, x_{22}, x_{23}, \ldots$ une famille dénombrable
de générateurs de la somme des $N_i$ pour lesquels $y_{11}$ ou $z_{11}$ possède une
composante non nulle. Répétons ce procédé avec $x_{12}$ pour obtenir $x_{31},
x_{32}, \ldots$, la troisième ligne de notre matrice. Répétons-le avec $x_{21}$ pour obtenir la
quatrième ligne, avec $x_{13}$ pour obtenir la cinquième, et ainsi de suite, en parcourant
successivement les antidiagonales comme indiqué ci-dessous :
$$
\left(
\vcenter{
\xymatrix@R=2mm@C=2mm{
x_{11} & x_{12} \ar[dl] & x_{13} \ar[dl] & x_{14} \ar[dl] & \ldots  \\
x_{21} & x_{22} \ar[dl] & x_{23} \ar[dl] & \ldots  \\
x_{31} & x_{32} \ar[dl] & \ldots \\
x_{41} & \ldots \\
\ldots
}
}
\right).
$$

\medskip\noindent
Une récurrence transfinie sur $I$ (utilisant le fait que nous avons construit
$M_{\alpha + 1}$
de façon qu'il contienne $N_j$ pour le plus petit $j$ tel que $N_j$ ne soit pas contenu dans
$M_{\alpha}$) montre que, pour tout $i \in I$, $N_i$ est contenu dans un certain
$M_{\alpha}$. Il existe donc un ordinal $S$ assez grand vérifiant : pour
tout $i \in I$, il existe $\alpha \in S$ tel que $N_i$ soit contenu dans
$M_{\alpha}$. Cela signifie que $(M_{\alpha})_{\alpha \in S}$ vérifie la propriété (1)
d'un d\'evissage de Kaplansky de $M$. La famille $(M_{\alpha})_{\alpha \in S}$
vérifie en outre les autres propriétés de définition, ainsi que (5) et (6) ci-dessus :
les propriétés (0), (2), (4) et (6) sont claires par construction ; la propriété (5) est
vraie parce que chaque $M_{\alpha}$ est par construction une somme directe de certains $N_i$ ;
et (3) résulte de (5) et du fait que $M_{\alpha} \subset M_{\alpha + 1}$.
\end{proof}

\noindent
On obtient comme corollaire le résultat sur les modules projectifs annoncé au début
de la section.

\begin{theorem}
\label{theorem-projective-direct-sum}
\begin{slogan}
Tout module projectif est une somme directe de modules projectifs
dénombrablement engendrés.
\end{slogan}
Si $P$ est un $R$-module projectif, alors $P$ est une somme directe de $R$-modules projectifs
dénombrablement engendrés.
\end{theorem}

\begin{proof}
Un module est projectif si et seulement s'il est un facteur direct d'un module libre ;
le résultat découle donc du théorème \ref{theorem-kaplansky-direct-sum}.
\end{proof}



\section{Modules projectifs sur un anneau local}
\label{section-projective-local-ring}

\noindent
Dans cette section, nous démontrons un très joli résultat :
un module projectif $M$ sur un anneau local est libre
(théorème \ref{theorem-projective-free-over-local-ring} ci-dessous).
Notons que, sous l'hypothèse supplémentaire que $M$ est de type fini, ce résultat est le
lemme \ref{lemma-finite-flat-local}.
En général, on a :

\begin{lemma}
\label{lemma-projective-free}
Soit $R$ un anneau. Alors tout $R$-module projectif est libre si et seulement si
tout $R$-module projectif dénombrablement engendré est libre.
\end{lemma}

\begin{proof}
Cela résulte immédiatement du
théorème \ref{theorem-projective-direct-sum}.
\end{proof}

\noindent
Voici un critère pour qu'un module dénombrablement engendré soit libre.

\begin{lemma}
\label{lemma-freeness-criteria}
Soit $M$ un $R$-module dénombrablement engendré possédant la propriété suivante :
si $M = N \oplus N'$, où $N'$ est un $R$-module libre de type fini, alors
tout élément de $N$ appartient à un facteur direct libre de $N$.
Alors $M$ est libre.
\end{lemma}

\begin{proof}
Soient $x_1, x_2, \ldots$ une famille dénombrable de générateurs de $M$.
Construisons par récurrence des facteurs directs libres de type fini
$F_1, F_2, \ldots$ de $M$ tels que, pour tout $n$,
$F_1 \oplus \ldots \oplus F_n$ soit un facteur direct de $M$
contenant $x_1, \ldots, x_n$. Plus précisément, étant donnés $F_1, \ldots, F_n$
ayant les propriétés voulues, écrivons
$$
M = F_1 \oplus \ldots \oplus F_n \oplus N
$$
et soit $x \in N$ l'image de $x_{n + 1}$.
L'hypothèse de l'énoncé du lemme fournit alors un facteur direct libre $F_{n + 1} \subset N$
contenant $x$.
On peut bien sûr remplacer $F_{n + 1}$ par un facteur direct libre de type fini
de $F_{n + 1}$, ce qui achève l'étape de récurrence.
Alors $M = \bigoplus_{i = 1}^{\infty} F_i$ est libre.
\end{proof}

\begin{lemma}
\label{lemma-projective-freeness-criteria}
Soit $P$ un module projectif sur un anneau local $R$. Alors tout élément de $P$
appartient à un facteur direct libre de $P$.
\end{lemma}

\begin{proof}
Comme $P$ est projectif, il est facteur direct d'un certain $R$-module libre $F$ ; écrivons
$F = P \oplus Q$. Soit $x \in P$ l'élément dont nous voulons montrer qu'il
appartient à un facteur direct libre de $P$. Soit $B$ une base de $F$ telle que
le nombre d'éléments de base nécessaires pour exprimer $x$ soit minimal ; écrivons $x
= \sum_{i=1}^n a_i e_i$ pour certains $e_i \in B$ et $a_i \in R$. Alors aucun $a_j$
ne peut s'écrire comme combinaison linéaire des autres $a_i$ ; en effet, si $a_j =
\sum_{i \neq  j} a_i b_i$ pour certains $b_i \in R$, le remplacement de $e_i$ par $e_i +
b_ie_j$ pour $i \neq j$, les autres éléments de $B$ restant inchangés, donne
une nouvelle base de $F$ dans laquelle $x$ possède une expression plus courte.

\medskip\noindent
Soit $e_i = y_i + z_i, y_i \in P, z_i \in Q$ la décomposition de $e_i$ selon
ses composantes dans $P$ et $Q$. Écrivons $y_i = \sum_{j=1}^{n} b_{ij} e_j + t_i$,
où $t_i$ est une combinaison linéaire d'éléments de $B$ autres que $e_1, \ldots,
e_n$. Pour achever la démonstration, il suffit de montrer que la matrice $(b_{ij})$ est
inversible. Alors, en effet, le morphisme $F \to F$ envoyant $e_i \mapsto y_i$ pour
$i=1, \ldots, n$ et fixant $B \setminus \{e_1, \ldots, e_n\}$ est un
isomorphisme ;
ainsi, $y_1, \ldots, y_n$ et $B \setminus \{e_1, \ldots, e_n\}$
forment une base de $F$. Le sous-module $N$ engendré par $y_1, \ldots, y_n$
est alors un sous-module libre de $P$ ; $N$ est un facteur direct de $P$, car $N \subset P$
et $N$ comme $P$ sont des facteurs directs de $F$ ; enfin $x \in N$, puisque $x \in P$
implique $x = \sum_{i=1}^n a_i e_i = \sum_{i=1}^n a_i y_i$.

\medskip\noindent
Démontrons maintenant que $(b_{ij})$ est inversible. En substituant $y_i = \sum_{j=1}^{n}
b_{ij} e_j + t_i$ dans $\sum_{i=1}^n a_i e_i = \sum_{i=1}^n a_i y_i$ et
en identifiant les coefficients de $e_j$, on obtient $a_j = \sum_{i=1}^n a_i b_{ij}$. Or,
comme observé ci-dessus, le choix de $B$ garantit qu'aucun $a_j$ ne peut s'écrire comme
combinaison linéaire des autres $a_i$. Ainsi, $b_{ij}$ est non inversible pour $i \neq
j$, et $1-b_{ii}$ est non inversible --- donc, en particulier, $b_{ii}$ est inversible --- pour
tout $i$. Une matrice sur un anneau local dont les coefficients diagonaux sont inversibles et
les autres non inversibles est inversible, car son déterminant est inversible.
\end{proof}

\begin{theorem}
\label{theorem-projective-free-over-local-ring}
\begin{slogan}
Les modules projectifs sur les anneaux locaux sont libres.
\end{slogan}
Si $P$ est un module projectif sur un anneau local $R$, alors $P$ est libre.
\end{theorem}

\begin{proof}
Cela résulte des lemmes \ref{lemma-projective-free}, \ref{lemma-freeness-criteria}
et \ref{lemma-projective-freeness-criteria}.
\end{proof}



\section{Systèmes de Mittag-Leffler}
\label{section-mittag-leffler}

\noindent
Cette section a pour but de définir les systèmes de Mittag-Leffler
et d'expliquer l'utilité de cette notion.

\medskip\noindent
Dans la suite, $I$ sera un ensemble préordonné filtrant, voir
Catégories, définition \ref{categories-definition-directed-set}.
Soit $(A_i, \varphi_{ji}: A_j \to A_i)$ un système projectif
d'ensembles ou de modules indexé par $I$, voir
Catégories, définition \ref{categories-definition-directed-system}.
Il s'agit d'un système projectif filtrant, puisque nous avons supposé $I$ filtrant
(Catégories, définition \ref{categories-definition-directed-system}).
Pour tout $i \in I$, les images $\varphi_{ji}(A_j) \subset A_i$, pour $j \geq i$,
forment une famille filtrante décroissante de parties (ou de sous-modules) de $A_i$. Posons
$A'_i = \bigcap_{j \geq i} \varphi_{ji}(A_j)$.
Alors $\varphi_{ji}(A'_j) \subset A'_i$ pour $j \geq i$ ; par restriction,
on obtient donc un système projectif filtrant $(A'_i, \varphi_{ji}|_{A'_j})$.
La construction de la limite d'un système projectif dans la catégorie
des ensembles ou des modules donne $\lim A_i = \lim A'_i$. La condition de Mittag-Leffler
sur $(A_i, \varphi_{ji})$ exige que $A'_i$ soit égal à
$\varphi_{ji}(A_j)$ pour un certain $j \geq i$ (et donc égal à
$\varphi_{ki}(A_k)$ pour tout $k \geq j$) :

\begin{definition}
\label{definition-ML-system}
Soit $(A_i, \varphi_{ji})$ un système projectif filtrant d'ensembles sur $I$. On
dit que $(A_i, \varphi_{ji})$ est {\it de Mittag-Leffler} si, pour
tout $i \in I$, la famille $\varphi_{ji}(A_j) \subset A_i$, pour
$j \geq i$, se stabilise. Explicitement, cela signifie que, pour tout $i \in I$, il
existe $j \geq i$ tel que, pour $k \geq j$, on ait $\varphi_{ki}(A_k) =
\varphi_{ji}( A_j)$. Si $(A_i, \varphi_{ji})$ est un système projectif filtrant
de modules sur un anneau $R$, on dit qu'il est de Mittag-Leffler si le système
projectif d'ensembles sous-jacent est de Mittag-Leffler.
\end{definition}

\begin{example}
\label{example-ML-surjective-maps}
Si $(A_i, \varphi_{ji})$ est un système projectif filtrant d'ensembles ou de modules et si
les morphismes $\varphi_{ji}$ sont surjectifs, alors le système est clairement
de Mittag-Leffler. Réciproquement, supposons $(A_i, \varphi_{ji})$ de Mittag-Leffler.
Soit $A'_i \subset A_i$ l'image stable des $\varphi_{ji}(A_j)$ pour $j \geq
i$. Alors $\varphi_{ji}|_{A'_j}: A'_j \to A'_i$ est surjectif pour
$j \geq i$ et $\lim A_i = \lim A'_i$. Ainsi, la limite du système de Mittag-Leffler
$(A_i, \varphi_{ji})$ peut aussi s'écrire comme la limite d'un système
projectif filtrant sur $I$ à morphismes de transition surjectifs.
\end{example}

\begin{lemma}
\label{lemma-ML-limit-nonempty}
Soit $(A_i, \varphi_{ji})$ un système projectif filtrant sur $I$. Supposons $I$
dénombrable. Si $(A_i, \varphi_{ji})$ est de Mittag-Leffler et que les $A_i$ sont
non vides, alors $\lim A_i$ est non vide.
\end{lemma}

\begin{proof}
Soit $i_1, i_2, i_3, \ldots$ une énumération des éléments de $I$. Définissons
par récurrence une suite d'éléments $j_n \in I$, pour $n = 1, 2, 3, \ldots$, par les
conditions suivantes : $j_1 = i_1$, et $j_n \geq i_n$ et $j_n \geq j_m$ pour $m < n$.
 La suite $j_n$ est alors croissante et forme une partie cofinale de $I$.
Nous pouvons donc supposer $I =\{1, 2, 3, \ldots \}$. L'exemple
\ref{example-ML-surjective-maps} nous ramène alors à montrer que la limite d'un
système projectif d'ensembles non vides, à morphismes surjectifs et indexé par les entiers
strictement positifs, est non vide. Cela résulte de l'axiome du choix.
\end{proof}

\noindent
La condition de Mittag-Leffler sera importante pour nous en raison de la propriété
d'exactitude suivante.

\begin{lemma}
\label{lemma-ML-exact-sequence}
Soit
$$
0 \to A_i \xrightarrow{f_i} B_i \xrightarrow{g_i} C_i \to 0
$$
une suite exacte de systèmes projectifs filtrants de groupes abéliens sur $I$.
Supposons $I$ dénombrable. Si $(A_i)$ est de Mittag-Leffler, alors
$$
0 \to \lim A_i \to \lim B_i \to \lim C_i\to 0
$$
est exacte.
\end{lemma}

\begin{proof}
Le passage à la limite des systèmes projectifs filtrants est exact à gauche ; il suffit donc de
démontrer la surjectivité de $\lim B_i \to \lim C_i$. Soit donc $(c_i) \in \lim
C_i$. Pour tout $i \in I$, posons $E_i = g_i^{-1}(c_i)$, qui est non vide puisque
$g_i: B_i \to C_i$ est surjectif. Le système de morphismes $\varphi_{ji}: B_j
\to B_i$ du système $(B_i)$ induit par restriction des morphismes $E_j \to E_i$ qui
font de $(E_i)$ un système projectif d'ensembles non vides. Il suffit de montrer
que $(E_i)$ est de Mittag-Leffler. Alors le lemme \ref{lemma-ML-limit-nonempty}
montrera que $\lim E_i$ est non vide, et le choix d'un élément de $\lim E_i$
fournira un élément de $\lim B_i$ s'envoyant sur $(c_i)$.

\medskip\noindent
Grâce à l'injection $f_i: A_i \to B_i$, nous considérerons $A_i$ comme une partie de
$B_i$. Puisque $(A_i)$ est de Mittag-Leffler, pour $i \in I$, il existe $j \geq
i$ tel que $\varphi_{ki}(A_k) = \varphi_{ji}(A_j)$ pour $k \geq j$. Nous affirmons
qu'on a aussi $\varphi_{ki}(E_k) = \varphi_{ji}(E_j)$ pour $k \geq j$. On a toujours
$\varphi_{ki}(E_k) \subset \varphi_{ji}(E_j)$ pour $k \geq j$. Pour l'inclusion
réciproque, soit $e_j \in E_j$ ; nous devons trouver $x_k \in E_k$ tel que
$\varphi_{ki}(x_k) = \varphi_{ji}(e_j)$. Soit $e'_k \in E_k$ un élément quelconque,
et posons $e'_j = \varphi_{kj}(e'_k)$. Alors $g_j(e_j - e'_j) = c_j - c_j = 0$ ;
ainsi $e_j - e'_j = a_j \in A_j$. Comme $\varphi_{ki}(A_k) =
\varphi_{ji}(A_j)$, il existe $a_k \in A_k$ tel que $\varphi_{ki}(a_k) =
\varphi_{ji}(a_j)$. Ainsi
$$
\varphi_{ki}(e'_k + a_k) = \varphi_{ji}(e'_j) + \varphi_{ji}(a_j) =
\varphi_{ji}(e_j),
$$
on peut donc prendre $x_k = e'_k + a_k$.
\end{proof}




\section{Systèmes projectifs}
\label{section-inverse-systems}

\noindent
Dans de nombreux articles (et dans cette section), le terme {\it système projectif} est
employé pour désigner un système projectif sur l'ensemble ordonné
$(\mathbf{N}, \geq)$. Nous étudions brièvement de tels systèmes dans cette section.
Ce matériel sera traité plus généralement dans
Homologie, section \ref{homology-section-inverse-systems}.
Supposons donnés un anneau $R$ et une suite de $R$-modules
$$
M_1 \xleftarrow{\varphi_2} M_2 \xleftarrow{\varphi_3} M_3 \leftarrow \ldots
$$
avec les morphismes indiqués. En composant les morphismes successifs, on obtient des morphismes
$\varphi_{ii'} : M_i \to M_{i'}$ pour $i \geq i'$ tels qu'en outre
$\varphi_{ii''} = \varphi_{i'i''} \circ \varphi_{i i'}$ dès que
$i \geq i' \geq i''$. Réciproquement, étant donné le système des morphismes $\varphi_{ii'}$,
on peut poser $\varphi_i = \varphi_{i(i-1)}$ et retrouver les morphismes représentés
ci-dessus. Dans ce cas,
$$
\lim M_i
=
\{(x_i) \in \prod M_i \mid \varphi_i(x_i) = x_{i - 1}, \ i = 2, 3, \ldots\}
$$
voir
Catégories, section \ref{categories-section-limit-sets}.
Comme expliqué dans
Homologie, section \ref{homology-section-inverse-systems},
il s'agit bien d'une limite dans la catégorie des $R$-modules, au sens défini dans
Catégories, section \ref{categories-section-limits}.

\begin{lemma}
\label{lemma-Mittag-Leffler}
Soit $R$ un anneau.
Soient $0 \to K_i \to L_i \to M_i \to 0$ des suites exactes courtes de
$R$-modules, $i \geq 1$, s'insérant dans des morphismes de suites exactes courtes
$$
\xymatrix{
0 \ar[r] &
K_i \ar[r] &
L_i \ar[r] &
M_i \ar[r] &
0 \\
0 \ar[r] &
K_{i + 1} \ar[r] \ar[u] &
L_{i + 1} \ar[r] \ar[u] &
M_{i + 1} \ar[r] \ar[u] &
0}
$$
Si, pour tout $i$, il existe $c = c(i) \geq i$ tel que
$\Im(K_c \to K_i) = \Im(K_j \to K_i)$
pour tout $j \geq c$, alors la suite
$$
0 \to \lim K_i \to \lim L_i \to \lim M_i \to 0
$$
est exacte.
\end{lemma}

\begin{proof}
C'est un cas particulier du lemme plus général
\ref{lemma-ML-exact-sequence}.
\end{proof}



\section{Modules de Mittag-Leffler}
\label{section-mittag-leffler-modules}

\noindent
Un module de Mittag-Leffler est (très grossièrement) un module qui peut s'écrire
comme une colimite filtrante dont le dual est un système de Mittag-Leffler. Afin de
donner une définition précise, il nous faut effectuer quelques préparatifs.

\begin{definition}
\label{definition-ML-inductive-system}
Soit $(M_i, f_{ij})$ un système inductif de $R$-modules. On dit que
$(M_i, f_{ij})$ est un {\it système inductif de Mittag-Leffler} si chaque
$M_i$ est un $R$-module de présentation finie et si, pour tout $R$-module $N$,
le système projectif
$$
(\Hom_R(M_i, N), \Hom_R(f_{ij}, N))
$$
est de Mittag-Leffler.
\end{definition}

\noindent
Nous allons caractériser les $R$-modules qui sont des colimites de
systèmes inductifs de Mittag-Leffler.

\begin{definition}
\label{definition-domination}
Soient $f: M \to N$ et $g: M \to M'$ des morphismes de $R$-modules.
On dit que $g$ {\it domine} $f$ si, pour tout $R$-module $Q$, on a $\Ker(f
\otimes_R \text{id}_Q) \subset \Ker(g \otimes_R \text{id}_Q)$.
\end{definition}

\noindent
Il suffit de vérifier cette condition sur les modules de présentation finie.

\begin{lemma}
\label{lemma-domination-fp}
Soient $f: M \to N$ et $g: M \to M'$ des morphismes de $R$-modules.
Alors $g$ domine $f$ si et seulement si, pour tout $R$-module de présentation finie
$Q$, on a $\Ker(f \otimes_R \text{id}_Q) \subset \Ker(g \otimes_R
\text{id}_Q)$.
\end{lemma}

\begin{proof}
Supposons $\Ker(f \otimes_R \text{id}_Q) \subset \Ker(g \otimes_R
\text{id}_Q)$ pour tout module de présentation finie $Q$. Si $Q$ est un
module quelconque, écrivons $Q = \colim_{i \in I} Q_i$ comme colimite d'un
système
inductif de modules de présentation finie $Q_i$. Alors $\Ker(f \otimes_R
\text{id}_{Q_i}) \subset \Ker(g \otimes_R \text{id}_{Q_i})$ pour
tout $i$. Puisque le passage aux colimites filtrantes est exact et commute au produit
tensoriel, il s'ensuit que $\Ker(f \otimes_R \text{id}_Q) \subset \Ker(g
\otimes_R \text{id}_Q)$.
\end{proof}

\begin{lemma}
\label{lemma-domination-universally-injective}
Soient $f : M \to N$ et $g : M \to M'$ des morphismes de $R$-modules.
Considérons la somme amalgamée de $f$ et $g$,
$$
\xymatrix{
M  \ar[r]_f \ar[d]_g & N \ar[d]^{g'} \\
M' \ar[r]^{f'} & N'
}
$$
Alors $g$ domine $f$ si et seulement si $f'$ est universellement injectif.
\end{lemma}

\begin{proof}
Rappelons que $N'$ est le quotient de $M' \oplus N$ par le sous-module formé des éléments
$(g(x), -f(x))$ pour $x \in M$.
La construction de $N'$ fournit une suite exacte courte
$$
0 \to \Ker(f) \cap \Ker(g) \to \Ker(f) \to \Ker(f')
\to 0.
$$
Comme le produit tensoriel commute aux sommes amalgamées, on a une telle suite exacte
courte
$$
0 \to \Ker(f \otimes \text{id}_Q ) \cap \Ker(g \otimes
\text{id}_Q) \to \Ker(f \otimes \text{id}_Q)
\to \Ker(f' \otimes \text{id}_Q) \to 0
$$
pour tout $R$-module $Q$. Ainsi, $f'$ est universellement injectif si et seulement si
$\Ker(f \otimes \text{id}_Q ) \subset \Ker(g \otimes
\text{id}_Q)$ pour tout $Q$, si et seulement si $g$ domine $f$.
\end{proof}

\noindent
La définition de la domination donnée ci-dessus est parfois reliée à la notion usuelle
de domination de morphismes, comme le montre le lemme suivant.

\begin{lemma}
\label{lemma-domination}
Soient $f: M \to N$ et $g: M \to M'$ des morphismes de $R$-modules.
Supposons $\Coker(f)$ de présentation finie. Alors $g$ domine $f$ si
et
seulement si $g$ se factorise par $f$, c'est-à-dire\ s'il existe un morphisme de modules $h: N
\to M'$ tel que $g = h \circ f$.
\end{lemma}

\begin{proof}
Considérons la somme amalgamée de $f$ et $g$ comme dans l'énoncé du
lemme \ref{lemma-domination-universally-injective}.
La construction de la somme amalgamée montre que
$\Coker(f') = \Coker(f)$ ; ainsi, $\Coker(f')$ est de
présentation finie. D'après le
lemme \ref{lemma-universally-exact-split}, $f'$ est universellement injectif si et
seulement si
$$
0 \to M' \xrightarrow{f'} N' \to \Coker(f') \to 0
$$
est scindée. C'est le cas si et seulement s'il existe un morphisme $h' : N' \to M'$
tel que $h' \circ f' = \text{id}_{M'}$. Par la propriété
universelle de la somme amalgamée, l'existence d'un tel $h'$ équivaut à ce que $g$
se factorise par $f$.
\end{proof}


\begin{proposition}
\label{proposition-ML-characterization}
Soit $M$ un $R$-module. Soit $(M_i, f_{ij})$ un système inductif de $R$-modules de
présentation finie, indexé par $I$, tel que $M = \colim M_i$. Soit
$f_i:
M_i \to M$ le morphisme canonique. Les assertions suivantes sont équivalentes :
\begin{enumerate}
\item Pour tout $R$-module de présentation finie $P$ et tout morphisme de modules $f: P
\to M$, il existe un $R$-module de présentation finie $Q$ et un morphisme de modules
$g: P \to Q$ tels que $g$ et $f$ se dominent mutuellement, c'est-à-dire
$\Ker(f \otimes_R \text{id}_N) = \Ker(g \otimes_R \text{id}_N)$
pour tout $R$-module $N$.
\item Pour tout $i \in I$, il existe $j \geq i$ tel que $f_{ij}: M_i
\to M_j$ domine $f_i: M_i \to M$.
\item Pour tout $i \in I$, il existe $j \geq i$ tel que $f_{ij}: M_i
\to M_j$ se factorise par $f_{ik}: M_i \to M_k$ pour tout $k \geq
i$.
\item Pour tout $R$-module $N$, le système projectif
$(\Hom_R(M_i, N), \Hom_R(f_{ij}, N))$ est de Mittag-Leffler.
\item Pour $N = \prod_{s \in I} M_s$, le système projectif
$(\Hom_R(M_i, N), \Hom_R(f_{ij}, N))$ est de Mittag-Leffler.
\end{enumerate}
\end{proposition}

\begin{proof}
Démontrons d'abord l'équivalence de (1) et (2). Supposons (1) satisfaite et soit $i
\in I$. Au morphisme $f_i: M_i \to M$, on peut associer un $g:
M_i \to Q$ comme en (1). Comme $M_i$ et $Q$ sont de présentation finie,
$\Coker(g)$ l'est aussi. D'après le lemme \ref{lemma-domination}, $f_i : M_i
\to M$ se factorise par $g: M_i \to Q$ ; écrivons $f_i = h \circ g$
pour un certain $h: Q \to M$. Comme $Q$ est de présentation finie, $h$
se factorise par $M_j \to M$ pour un certain $j \geq i$ ; écrivons $h = f_j \circ h'$
pour un certain $h': Q \to M_j$. On obtient au total un diagramme commutatif
$$
\xymatrix{
  &  M  & \\
M_i \ar[dr]_g \ar[ur]^{f_i} \ar[rr]^{f_{ij}} &
& M_j \ar[ul]_{f_j} \\
  & Q  \ar[ur]_{h'} &
}
$$
Ainsi, $f_{ij}$ domine $g$. Or $g$ domine $f_i$, donc $f_{ij}$ domine
$f_i$.

\medskip\noindent
Réciproquement, supposons (2) satisfaite. Soit $P$ de présentation finie et soit $f: P
\to M$ un morphisme de modules. Alors $f$ se factorise par $f_i: M_i \to M$
pour un certain $i \in I$ ; écrivons $f = f_i \circ g'$ pour un certain $g': P \to M_i$.
Choisissons, grâce à (2), un $j \geq i$ tel que $f_{ij}$ domine $f_i$. On a un
diagramme commutatif
$$
\xymatrix{
P \ar[d]_{g'} \ar[r]^{f}           & M  \\
M_i \ar[ur]^{f_i} \ar[r]_{f_{ij}} & M_j \ar[u]_{f_j}
}
$$
Le diagramme et le fait que $f_{ij}$ domine $f_i$ montrent que $f$
et $f_{ij} \circ g'$ se dominent mutuellement. Ainsi, on peut prendre $g = f_{ij} \circ g' :
P \to M_j$.

\medskip\noindent
Montrons ensuite que (2) équivaut à (3). Soit $i \in I$. Il est toujours vrai que
$f_i$ domine $f_{ik}$ pour $k \geq i$, puisque $f_i$ se factorise par
$f_{ik}$. Si (2) est satisfaite, choisissons $j \geq i$ tel que $f_{ij}$ domine
$f_i$. La domination étant transitive, $f_{ij}$ domine alors
$f_{ik}$ pour $k \geq i$. Tous les $M_i$ sont de présentation finie, donc
$\Coker(f_{ik})$ est de présentation finie pour $k \geq i$. D'après le lemme
\ref{lemma-domination}, $f_{ij}$ se factorise par $f_{ik}$ pour tout $k \geq i$.
Ainsi, (2) implique (3). D'autre part, si (3) est satisfaite, alors, pour tout $R$-module
$N$, $f_{ij} \otimes_R \text{id}_N$ se factorise par $f_{ik}
\otimes_R \text{id}_N$ pour $k \geq i$. Ainsi, $\Ker(f_{ik} \otimes_R
\text{id}_N) \subset \Ker(f_{ij} \otimes_R \text{id}_N)$ pour $k
\geq i$. Or $\Ker(f_i \otimes_R \text{id}_N: M_i \otimes_R N
\to M \otimes_R N)$ est la réunion des $\Ker(f_{ik} \otimes_R
\text{id}_N)$ pour $k \geq i$. Ainsi, $\Ker(f_i \otimes_R
\text{id}_N) \subset \Ker(f_{ij} \otimes_R \text{id}_N)$ pour
tout $R$-module $N$, ce qui signifie par définition que $f_{ij}$ domine $f_i$.

\medskip\noindent
Il est immédiat que (3) implique (4), qui implique (5). Montrons que (5) implique (3). Soit
$N = \prod_{s \in I} M_s$. Si (5) est satisfaite, étant donné $i \in I$, choisissons $j \geq i$
tel que
$$
\Im( \Hom(M_j, N) \to  \Hom(M_i, N)) =
\Im( \Hom(M_k, N) \to  \Hom(M_i, N))
$$
pour tout $k \geq j$. En faisant sortir le produit sur $s \in I$ des
$\Hom$
et en examinant les morphismes sur chaque composante du produit, ceci signifie
$$
\Im( \Hom(M_j, M_s) \to  \Hom(M_i, M_s)) =
\Im( \Hom(M_k, M_s) \to   \Hom(M_i, M_s))
$$
pour tous $k \geq j$ et $s \in I$. En prenant $s = j$, on obtient
$$
\Im( \Hom(M_j, M_j) \to  \Hom(M_i, M_j)) =
\Im( \Hom(M_k, M_j) \to   \Hom(M_i, M_j))
$$
pour tout $k \geq j$. Puisque $f_{ij}$ est l'image de
$\text{id} \in  \Hom(M_j, M_j)$ par
$\Hom(M_j, M_j) \to  \Hom(M_i, M_j)$,
cela montre que, pour tout $k \geq j$, il existe $h \in \Hom(M_k, M_j)$
tel que $f_{ij} = h \circ f_{ik}$. Si $j \geq k$, on peut prendre
$h = f_{kj}$. Ainsi, (3) est satisfaite.
\end{proof}

\begin{definition}
\label{definition-mittag-leffler-module}
Soit $M$ un $R$-module. On dit que $M$ est {\it de Mittag-Leffler} si les
conditions équivalentes de la
proposition \ref{proposition-ML-characterization}
sont satisfaites.
\end{definition}

\noindent
En particulier, tout module de présentation finie est de Mittag-Leffler.

\begin{remark}
\label{remark-flat-ML}
Soit $M$ un $R$-module plat. D'après le théorème de Lazard
(théorème \ref{theorem-lazard}),
on peut écrire $M = \colim M_i$ comme la colimite d'un
système inductif $(M_i, f_{ij})$ où les $M_i$ sont des $R$-modules libres
de type fini. Pour que $M$ soit de Mittag-Leffler, il suffit que le système projectif
des duaux $(\Hom_R(M_i, R), \Hom_R(f_{ij}, R))$ soit de
Mittag-Leffler. Cela résulte du critère (4) de la
proposition \ref{proposition-ML-characterization}
et du fait que, pour un $R$-module libre de type fini $F$,
il existe un isomorphisme fonctoriel
$\Hom_R(F, R) \otimes_R N \cong \Hom_R(F, N)$
pour tout $R$-module $N$.
\end{remark}

\begin{lemma}
\label{lemma-tensor-ML-modules}
Si $R$ est un anneau et si $M$, $N$ sont des $R$-modules de Mittag-Leffler,
alors $M \otimes_R N$ est un module de Mittag-Leffler.
\end{lemma}

\begin{proof}
Écrivons $M = \colim_{i \in I} M_i$ et $N = \colim_{j \in J} N_j$
comme des colimites filtrantes de $R$-modules de présentation finie.
Notons $f_{ii'} : M_i \to M_{i'}$ et $g_{jj'} : N_j \to N_{j'}$ les
morphismes de transition. Alors $M_i \otimes_R N_j$ est un
$R$-module de présentation finie (voir le
lemme \ref{lemma-tensor-finiteness}),
et $M \otimes_R N = \colim_{(i, j) \in I \times J} M_i \otimes_R M_j$.
Choisissons $(i, j) \in I \times J$. Par définition d'un module de Mittag-Leffler,
la
proposition \ref{proposition-ML-characterization} (3)
s'applique aux deux systèmes. Autrement dit, il existe $i' \geq i$ et $j' \geq j$
tels que, pour tous $i'' \geq i$ et $j'' \geq j$, il existe des
morphismes $a : M_{i''} \to M_{i'}$ et $b : M_{j''} \to M_{j'}$ tels que
$f_{ii'} = a \circ f_{ii''}$ et $g_{jj'} = b \circ g_{jj''}$.
Il est alors clair que
$a \otimes b : M_{i''} \otimes_R N_{j''} \to M_{i'} \otimes_R N_{j'}$
joue le même rôle pour le système
$(M_i \otimes_R N_j, f_{ii'} \otimes g_{jj'})$.
Ainsi, d'après la caractérisation
de la proposition \ref{proposition-ML-characterization} (3),
on conclut que $M \otimes_R N$ est de Mittag-Leffler.
\end{proof}

\begin{lemma}
\label{lemma-ML-also}
Soient $R$ un anneau et $M$ un $R$-module. Alors $M$ est de Mittag-Leffler si et
seulement si, pour tout $R$-module libre de type fini $F$ et tout morphisme de modules
$f: F \to M$, il existe un $R$-module de présentation finie $Q$
et un morphisme de modules $g : F \to Q$ tels que $g$ et $f$ se dominent mutuellement, c'est-à-dire
$\Ker(f \otimes_R \text{id}_N) = \Ker(g \otimes_R \text{id}_N)$
pour tout $R$-module $N$.
\end{lemma}

\begin{proof}
Comme cette condition est manifestement plus faible que la condition (1) de la
proposition \ref{proposition-ML-characterization},
on voit qu'un module de Mittag-Leffler la satisfait.
Réciproquement, supposons que $M$ satisfasse cette condition et que
$f : P \to M$ soit un morphisme de $R$-modules d'un $R$-module de présentation finie
$P$ vers $M$. Choisissons une surjection $F \to P$, où
$F$ est un $R$-module libre de type fini. Par hypothèse, on peut trouver un morphisme
$F \to Q$, où $Q$ est un $R$-module de présentation finie, tel que
$F \to Q$ et $F \to M$ se dominent mutuellement. En particulier, le noyau
de $F \to Q$ contient le noyau de $F \to P$ ; on obtient donc un
morphisme de $R$-modules $g : P \to Q$ tel que $F \to Q$ soit égal à
la composée $F \to P \to Q$. Soit $N$ un $R$-module quelconque et
considérons le diagramme commutatif
$$
\xymatrix{
F \otimes_R N \ar[d] \ar[r] & Q \otimes_R N \\
P \otimes_R N \ar[ru] \ar[r] & M \otimes_R N
}
$$
Par hypothèse, les noyaux de $F \otimes_R N \to Q \otimes_R N$
et de $F \otimes_R N \to M \otimes_R N$ sont égaux. Comme
$F \otimes_R N \to P \otimes_R N$ est surjectif, les noyaux
de $P \otimes_R N \to Q \otimes_R N$
et de $P \otimes_R N \to M \otimes_R N$ sont donc eux aussi égaux.
\end{proof}

\begin{lemma}
\label{lemma-restrict-ML-modules}
Soit $R \to S$ un homomorphisme d'anneaux fini et de présentation finie.
Soit $M$ un $S$-module.
Si $M$ est un module de Mittag-Leffler sur $S$, alors
$M$ est un module de Mittag-Leffler sur $R$.
\end{lemma}

\begin{proof}
Supposons que $M$ soit un module de Mittag-Leffler sur $S$.
Écrivons $M = \colim M_i$ comme colimite filtrante de
$S$-modules de présentation finie $M_i$. Comme $M$ est de Mittag-Leffler sur $S$, il existe, pour tout
$i$, un indice $j \geq i$ tel que, pour tout $k \geq j$, il existe une factorisation
$f_{ij} = h \circ f_{ik}$ (où $h$ dépend de $i$, du choix de $j$ et de
$k$). Notons que, d'après le
lemme \ref{lemma-finite-finitely-presented-extension},
les modules $M_i$ sont aussi de présentation finie comme $R$-modules. De plus, tous
les morphismes $f_{ij}, f_{ik}, h$ sont des morphismes de $R$-modules. Le
système $(M_i, f_{ij})$ satisfait donc la même condition lorsqu'on le considère comme un système
de $R$-modules. Ainsi, $M$ est de Mittag-Leffler comme $R$-module.
\end{proof}

\begin{lemma}
\label{lemma-mod-ideal-ML-modules}
Soit $R$ un anneau.
Soit $S = R/I$ pour un idéal de type fini $I$.
Soit $M$ un $S$-module.
Alors $M$ est un module de Mittag-Leffler sur $R$ si et seulement si
$M$ est un module de Mittag-Leffler sur $S$.
\end{lemma}

\begin{proof}
Une implication résulte du
lemme \ref{lemma-restrict-ML-modules}.
Pour démontrer l'autre, supposons $M$ de Mittag-Leffler comme $R$-module.
Écrivons $M = \colim M_i$ comme colimite filtrante de
$S$-modules de présentation finie. Comme $I$ est de type fini, l'anneau $S$ est fini et de présentation
finie comme $R$-algèbre ; les modules $M_i$ sont donc de présentation
finie comme $R$-modules, voir le
lemme \ref{lemma-finite-finitely-presented-extension}.
Soit ensuite $N$ un $S$-module quelconque. Notons que, pour tout $i$, on a
$\Hom_R(M_i, N) = \Hom_S(M_i, N)$, puisque $R \to S$ est surjectif.
Ainsi, la condition que le système projectif
$(\Hom_R(M_i, N))_i$ soit de Mittag-Leffler implique que le système
$(\Hom_S(M_i, N))_i$ est de Mittag-Leffler. Par définition, $M$ est donc
de Mittag-Leffler sur $S$.
\end{proof}

\begin{remark}
\label{remark-go-up-ML-modules}
Soit $R \to S$ un homomorphisme d'anneaux fini et de présentation finie.
Soit $M$ un $S$-module qui est de Mittag-Leffler comme $R$-module.
En général, il n'est alors pas vrai que $M$ soit de Mittag-Leffler comme
$S$-module. Supposons par exemple que $S$ soit l'anneau des nombres duaux
sur $R$, c'est-à-dire $S = R \oplus R\epsilon$ avec $\epsilon^2 = 0$. Un
$S$-module consiste alors en un $R$-module $M$ muni d'un endomorphisme
$R$-linéaire de carré nul $\epsilon : M \to M$. Supposons maintenant que $M_0$
soit un $R$-module qui ne soit pas de Mittag-Leffler. Choisissons une présentation
$F_1 \xrightarrow{u} F_0 \to M_0 \to 0$, où $F_1$ et $F_0$ sont des $R$-modules libres.
Posons $M = F_1 \oplus F_0$ et
$$
\epsilon =
\left(
\begin{matrix}
0 & 0 \\
u & 0
\end{matrix}
\right) : M \longrightarrow M.
$$
Alors $M/\epsilon M \cong F_1 \oplus M_0$ n'est pas de Mittag-Leffler sur
$R = S/\epsilon S$, et n'est donc pas de Mittag-Leffler sur $S$ (voir le
lemme \ref{lemma-mod-ideal-ML-modules}).
D'autre part, $M/\epsilon M = M \otimes_S S/\epsilon S$, qui serait
de Mittag-Leffler sur $S$ si $M$ l'était, voir le
lemme \ref{lemma-tensor-ML-modules}.
\end{remark}


\section{Commutation des produits directs avec le produit tensoriel}
\label{section-products-tensor}

\noindent
Soit $M$ un $R$-module et soit $(Q_{\alpha})_{\alpha \in A}$ une famille de
$R$-modules. Il existe alors un morphisme canonique $M \otimes_R \left( \prod_{\alpha
\in A} Q_{\alpha} \right) \to \prod_{\alpha \in A} ( M \otimes_R
Q_{\alpha})$ donné sur les tenseurs purs par $x \otimes (q_{\alpha}) \mapsto (x
\otimes q_{\alpha})$. Ce morphisme n'est pas nécessairement injectif ni surjectif, comme
le montre l'exemple suivant.

\begin{example}
\label{example-Q-not-ML}
Prenons $R = \mathbf{Z}$, $M = \mathbf{Q}$, et considérons la famille $Q_n =
\mathbf{Z}/n$ pour $n \geq 1$. Alors $\prod_n (M \otimes Q_n) = 0$. Toutefois,
il existe une injection $\mathbf{Q} \to M \otimes (\prod_n Q_n)$
obtenue en tensorisant l'injection $\mathbf{Z} \to \prod_n Q_n$ par
$M$ ; ainsi, $M \otimes (\prod_n Q_n)$ est non nul. Le morphisme $M \otimes (\prod_n
Q_n) \to \prod_n (M \otimes Q_n)$ n'est donc pas injectif.

\medskip\noindent
D'autre part, prenons encore $R = \mathbf{Z}$, $M = \mathbf{Q}$, et posons $Q_n
= \mathbf{Z}$ pour $n \geq 1$. L'image de $M \otimes (\prod_n Q_n)
\to \prod_n (M \otimes Q_n) = \prod_n M$ est formée exactement des
suites de la forme $(a_n/m)_{n \geq 1}$, où $a_n \in \mathbf{Z}$ et où $m$ est
un entier non nul. Le morphisme n'est donc pas surjectif.
\end{example}

\noindent
Nous déterminons ci-dessous les conditions précises à imposer à $M$ pour que le morphisme $M
\otimes_R \left( \prod_{\alpha} Q_{\alpha} \right) \to \prod_{\alpha}
(M \otimes_R Q_{\alpha})$ soit surjectif, bijectif ou injectif pour tout
choix de $(Q_{\alpha})_{\alpha \in A}$. Cela est pertinent, car les modules
pour lesquels il est injectif sont exactement les modules de Mittag-Leffler
(proposition \ref{proposition-ML-tensor}). Dans la suite, si $M$ est un
$R$-module et $A$ un ensemble, nous notons $M^A$ le produit $\prod_{\alpha \in A}
M$.

\begin{proposition}
\label{proposition-fg-tensor}
Soit $M$ un $R$-module. Les assertions suivantes sont équivalentes :
\begin{enumerate}
\item $M$ est de type fini.
\item Pour toute famille $(Q_{\alpha})_{\alpha \in A}$ de $R$-modules, le
morphisme canonique $M \otimes_R \left( \prod_{\alpha} Q_{\alpha} \right)
\to \prod_{\alpha} (M \otimes_R Q_{\alpha})$ est surjectif.
\item Pour tout $R$-module $Q$ et tout ensemble $A$, le morphisme canonique $M
\otimes_R Q^{A} \to (M \otimes_R Q)^{A}$ est surjectif.
\item Pour tout ensemble $A$, le morphisme canonique $M \otimes_R R^{A} \to
M^{A}$ est surjectif.
\end{enumerate}
\end{proposition}

\begin{proof}
Montrons d'abord que (1) implique (2). Choisissons une surjection $R^n \to M$ et
considérons le diagramme commutatif
$$
\xymatrix{
R^n \otimes_R (\prod_{\alpha} Q_{\alpha})  \ar[r]^{\cong} \ar[d] &
\prod_{\alpha} (R^n \otimes_R Q_{\alpha}) \ar[d] \\
M \otimes_R (\prod_{\alpha} Q_{\alpha})  \ar[r] & \prod_{\alpha} ( M
\otimes_R Q_{\alpha}).
}
$$
La flèche supérieure est un isomorphisme et les flèches verticales sont surjectives. On
en conclut que la flèche inférieure est surjective.

\medskip\noindent
Il est clair que (2) implique (3), qui implique (4) ; il reste donc à montrer que (4) implique (1).
En fait, pour que (1) soit satisfaite, il suffit que l'élément $d = (x)_{x \in M}$ de
$M^M$ appartienne à l'image du morphisme $f: M \otimes_R R^{M} \to M^M$. Dans
ce cas, $d = \sum_{i = 1}^{n} f(x_i \otimes a_i)$ pour certains $x_i \in M$ et
$a_i \in R^M$. Si, pour $x \in M$, on note $p_x: M^M \to M$ la
projection sur le facteur indexé par $x$, alors
$$
x = p_x(d) = \sum\nolimits_{i = 1}^{n} p_x(f(x_i \otimes a_i)) =
\sum\nolimits_{i=1}^{n} p_x(a_i) x_i.
$$
Ainsi, $x_1, \ldots, x_n$ engendrent $M$.
\end{proof}

\begin{proposition}
\label{proposition-fp-tensor}
Soit $M$ un $R$-module. Les assertions suivantes sont équivalentes :
\begin{enumerate}
\item $M$ est de présentation finie.
\item Pour toute famille $(Q_{\alpha})_{\alpha \in A}$ de $R$-modules, le
morphisme canonique $M \otimes_R \left( \prod_{\alpha} Q_{\alpha} \right)
\to \prod_{\alpha} (M \otimes_R Q_{\alpha})$ est bijectif.
\item Pour tout $R$-module $Q$ et tout ensemble $A$, le morphisme canonique $M
\otimes_R Q^{A} \to (M \otimes_R Q)^{A}$ est bijectif.
\item Pour tout ensemble $A$, le morphisme canonique $M \otimes_R R^{A} \to
M^{A}$ est bijectif.
\end{enumerate}
\end{proposition}

\begin{proof}
Montrons d'abord que (1) implique (2). Choisissons une présentation $R^m \to R^n
\to M$ et considérons le diagramme commutatif
$$
\xymatrix{
R^m \otimes_R (\prod_{\alpha} Q_{\alpha}) \ar[r] \ar[d]^{\cong} & R^n
\otimes_R (\prod_{\alpha} Q_{\alpha}) \ar[r] \ar[d]^{\cong} & M \otimes_R
(\prod_{\alpha} Q_{\alpha}) \ar[r] \ar[d] & 0 \\
\prod_{\alpha} (R^m \otimes_R Q_{\alpha}) \ar[r] & \prod_{\alpha} (R^n
\otimes_R Q_{\alpha}) \ar[r] & \prod_{\alpha} (M \otimes_R Q_{\alpha})
\ar[r] & 0.
}
$$
Les deux premières flèches verticales sont des isomorphismes et les lignes sont exactes. Cela
implique que le morphisme
$M \otimes_R (\prod_{\alpha} Q_{\alpha}) \to
\prod_{\alpha} ( M \otimes_R Q_{\alpha})$
est surjectif et, par une chasse au diagramme, également injectif. Ainsi, (2) est satisfaite.

\medskip\noindent
Il est clair que (2) implique (3), qui implique (4) ; il reste donc à montrer que (4) implique (1).
D'après la proposition \ref{proposition-fg-tensor}, si (4) est satisfaite, on sait déjà que
$M$ est de type fini. On peut donc choisir une surjection $F \to M$,
où $F$ est libre de type fini. Soit $K$ le noyau. Nous devons montrer que $K$ est
de type fini. Pour tout ensemble $A$, on a un diagramme commutatif
$$
\xymatrix{
& K \otimes_R R^A \ar[r] \ar[d]_{f_3} & F \otimes_R R^A \ar[r]
\ar[d]_{f_2}^{\cong} & M \otimes_R R^A \ar[r] \ar[d]_{f_1}^{\cong} & 0 \\
0 \ar[r] & K^A \ar[r] & F^A \ar[r] & M^A \ar[r] & 0 .
}
$$
Le morphisme $f_1$ est un isomorphisme par hypothèse, le morphisme $f_2$ est un isomorphisme
puisque $F$ est libre de type fini, et les lignes sont exactes. Une chasse au diagramme montre
que $f_3$ est surjectif ; la proposition \ref{proposition-fg-tensor} montre donc
que $K$ est de type fini.
\end{proof}

\noindent
Nous aurons besoin du lemme suivant pour la prochaine proposition.

\begin{lemma}
\label{lemma-kernel-tensored-fp}
Soient $M$ un $R$-module, $P$ un $R$-module de présentation finie et $f: P
\to M$ un morphisme. Soit $Q$ un $R$-module et supposons que $x \in \Ker(P
\otimes Q \to M \otimes Q)$. Alors il existe un $R$-module de présentation finie
$P'$ et un morphisme $f': P \to P'$ tels que $f$ se factorise par
$f'$ et que $x \in \Ker(P \otimes Q \to P' \otimes Q)$.
\end{lemma}

\begin{proof}
Écrivons $M$ comme une colimite $M = \colim_{i \in I} M_i$ d'un système filtrant de
modules de présentation finie $M_i$. Puisque $P$ est de présentation finie, le morphisme $f:
P \to M$ se factorise par $M_j \to M$ pour un certain $j \in I$. Après
tensorisation par $Q$, nous avons un diagramme commutatif
$$
\xymatrix{
& M_j \otimes Q \ar[dr] & \\
P \otimes Q \ar[ur] \ar[rr] & & M \otimes Q .
}
$$
L'image $y$ de $x$ dans $M_j \otimes Q$ appartient au noyau de $M_j \otimes Q
\to M \otimes Q$. Puisque $M \otimes Q = \colim_{i \in I} (M_i
\otimes
Q)$, cela signifie que $y$ s'envoie sur $0$ dans $M_{j'} \otimes Q$ pour un certain $j' \geq j$.
Nous pouvons donc prendre $P' = M_{j'}$ et pour $f'$ le composé $P \to M_j
\to M_{j'}$.
\end{proof}

\begin{proposition}
\label{proposition-ML-tensor}
Soit $M$ un $R$-module. Les assertions suivantes sont équivalentes :
\begin{enumerate}
\item $M$ est de Mittag-Leffler.
\item Pour toute famille $(Q_{\alpha})_{\alpha \in A}$ de $R$-modules, le
morphisme canonique $M \otimes_R \left( \prod_{\alpha} Q_{\alpha} \right)
\to \prod_{\alpha} (M \otimes_R Q_{\alpha})$ est injectif.
\end{enumerate}
\end{proposition}

\begin{proof}
Montrons d'abord que (1) implique (2). Supposons que $M$ soit de Mittag-Leffler et soit $x$
dans le noyau de $M \otimes_R (\prod_{\alpha} Q_{\alpha}) \to
\prod_{\alpha} (M \otimes_R Q_{\alpha})$. Écrivons $M$ comme une colimite $M =
\colim_{i \in I} M_i$ d'un système filtrant de modules de présentation finie
$M_i$.
Alors $M \otimes_R (\prod_{\alpha} Q_{\alpha})$ est la colimite des $M_i
\otimes_R (\prod_{\alpha} Q_{\alpha})$. Ainsi, $x$ est l'image d'un élément
$x_i \in M_i \otimes_R (\prod_{\alpha} Q_{\alpha})$. Nous devons montrer que $x_i$
s'envoie sur $0$ dans $M_j \otimes_R (\prod_{\alpha} Q_{\alpha})$ pour un certain $j \geq
i$. Puisque $M$ est de Mittag-Leffler, nous pouvons choisir $j \geq i$ tel que $M_i
\to M_j$ et $M_i \to M$ se dominent mutuellement. Considérons alors
le diagramme commutatif
$$
\xymatrix{
M \otimes_R (\prod_{\alpha} Q_{\alpha}) \ar[r] & \prod_{\alpha} (M
\otimes_R Q_{\alpha}) \\
M_i \otimes_R (\prod_{\alpha} Q_{\alpha}) \ar[r]^{\cong} \ar[d] \ar[u] &
\prod_{\alpha} (M_i \otimes_R Q_{\alpha}) \ar[d] \ar[u] \\
M_j \otimes_R (\prod_{\alpha} Q_{\alpha}) \ar[r]^{\cong} & \prod_{\alpha}
(M_j \otimes_R Q_{\alpha})
}
$$
dont les deux flèches horizontales inférieures sont des isomorphismes, d'après la proposition
\ref{proposition-fp-tensor}. Puisque $x_i$ s'envoie sur $0$ dans $\prod_{\alpha} (M
\otimes_R Q_{\alpha})$, son image dans $\prod_{\alpha} (M_i \otimes_R
Q_{\alpha})$ appartient au noyau du morphisme $\prod_{\alpha} (M_i \otimes_R
Q_{\alpha}) \to \prod_{\alpha} (M \otimes_R Q_{\alpha})$. Mais ce
noyau est égal au noyau de $\prod_{\alpha} (M_i \otimes_R Q_{\alpha})
\to \prod_{\alpha} (M_j \otimes_R Q_{\alpha})$ par le
choix de $j$. Ainsi, $x_i$ s'envoie sur $0$ dans $\prod_{\alpha} (M_j \otimes_R
Q_{\alpha})$, et donc sur $0$ dans $M_j \otimes_R (\prod_{\alpha} Q_{\alpha})$.

\medskip\noindent
Supposons maintenant que (2) soit satisfaite. Nous montrons que $M$ satisfait la formulation (1) de la propriété
de Mittag-Leffler donnée dans la proposition \ref{proposition-ML-characterization}. Soit $f:
P \to M$ un morphisme d'un module de présentation finie $P$ vers $M$. Choisissons
un ensemble $B$ de représentants des classes d'isomorphisme de $R$-modules de présentation finie.
Soit $A$ l'ensemble des couples $(Q, x)$ où $Q \in B$ et $x \in
\Ker(P \otimes Q \to M \otimes Q)$. Pour $\alpha = (Q, x) \in A$, nous
notons $Q_{\alpha}$ le module $Q$ et $x_{\alpha}$ l'élément $x$. Considérons le diagramme
commutatif
$$
\xymatrix{
M \otimes_R (\prod_{\alpha} Q_{\alpha}) \ar[r] &
\prod_{\alpha} (M \otimes_R Q_{\alpha}) \\
P \otimes_R (\prod_{\alpha} Q_{\alpha}) \ar[r]^{\cong} \ar[u] &
\prod_{\alpha} (P \otimes_R Q_{\alpha}) \ar[u] .
}
$$
La flèche supérieure est injective par hypothèse, et la flèche inférieure est un
isomorphisme par la proposition \ref{proposition-fp-tensor}. Soit $x \in P
\otimes_R (\prod_{\alpha} Q_{\alpha})$ l'élément correspondant à
$(x_{\alpha}) \in \prod_{\alpha} (P \otimes_R Q_{\alpha})$ par cet
isomorphisme. Alors $x \in \Ker( P \otimes_R (\prod_{\alpha} Q_{\alpha})
\to M \otimes_R (\prod_{\alpha} Q_{\alpha}))$, puisque la flèche supérieure du
diagramme est injective. Par le lemme \ref{lemma-kernel-tensored-fp}, nous obtenons un
module de présentation finie $P'$ et un morphisme $f': P \to P'$ tels que $f: P
\to M$ se factorise par $f'$ et que $x \in \Ker(P \otimes_R
(\prod_{\alpha} Q_{\alpha}) \to P' \otimes_R (\prod_{\alpha}
Q_{\alpha}))$. Nous avons un diagramme commutatif
$$
\xymatrix{
P' \otimes_R (\prod_{\alpha} Q_{\alpha}) \ar[r]^{\cong} &
\prod_{\alpha} (P' \otimes_R Q_{\alpha}) \\
P \otimes_R (\prod_{\alpha} Q_{\alpha}) \ar[r]^{\cong} \ar[u] &
\prod_{\alpha} (P \otimes_R Q_{\alpha}) \ar[u] .
}
$$
où les flèches supérieure et inférieure sont toutes deux des isomorphismes par la proposition
\ref{proposition-fp-tensor}. Ainsi, puisque $x$ appartient au noyau du morphisme vertical
de gauche, $(x_{\alpha})$ appartient au noyau du morphisme vertical de droite. Cela
signifie que $x_{\alpha} \in \Ker(P \otimes_R Q_{\alpha} \to P' \otimes_R
Q_{\alpha})$ pour tout $\alpha \in A$. Par définition de $A$, cela signifie que
$\Ker(P \otimes_R Q \to P' \otimes_R Q) \supset \Ker(P \otimes_R
Q \to M \otimes_R Q)$ pour tout $Q$ de présentation finie et, puisque $f: P
\to M$ se factorise par $f': P \to P'$, nous avons en fait l'égalité.
Par le lemme \ref{lemma-domination-fp}, $f$ et $f'$ se dominent mutuellement.
\end{proof}

\begin{lemma}
\label{lemma-minimal-contains}
Soit $M$ un module plat de Mittag-Leffler sur $R$. Soit $F$ un $R$-module
et soit $x \in F \otimes_R M$. Alors il existe un plus petit sous-module
$F' \subset F$ tel que $x \in F' \otimes_R M$.
De plus, $F'$ est un $R$-module de type fini.
\end{lemma}

\begin{proof}
Puisque $M$ est plat, nous avons $F' \otimes_R M \subset F \otimes_R M$
si $F' \subset F$ est un sous-module, de sorte que l'énoncé a un sens.
Soit $I = \{F' \subset F \mid x \in F' \otimes_R M\}$ et, pour
$i \in I$, notons $F_i \subset F$ le sous-module correspondant.
Alors $x$ s'envoie sur zéro par le morphisme
$$
F \otimes_R M \longrightarrow \prod (F/F_i \otimes_R M)
$$
d'où, par la proposition \ref{proposition-ML-tensor},
$x$ s'envoie sur zéro par le morphisme
$$
F \otimes_R M \longrightarrow \left(\prod F/F_i\right) \otimes_R M
$$
Puisque $M$ est plat, le noyau de cette flèche est
$(\bigcap F_i) \otimes_R M$, ce qui prouve que $F' = \bigcap F_i$ convient.
Pour voir que $F'$ est un module de type fini, supposons que
$x = \sum_{j = 1, \ldots, m} f_j \otimes m_j$ avec $f_j \in F'$
et $m_j \in M$. Alors $x \in F'' \otimes_R M$, où $F'' \subset F'$ est
le sous-module engendré par $f_1, \ldots, f_m$. Bien entendu,
$F'' = F'$, ce qui donne la dernière assertion.
\end{proof}

\begin{lemma}
\label{lemma-pure-submodule-ML}
Soit $0 \to M_1 \to M_2 \to M_3 \to 0$ une suite
universellement exacte de $R$-modules. Alors :
\begin{enumerate}
\item Si $M_2$ est de Mittag-Leffler, alors $M_1$ est de Mittag-Leffler.
\item Si $M_1$ et $M_3$ sont de Mittag-Leffler, alors $M_2$ est de Mittag-Leffler.
\end{enumerate}
\end{lemma}

\begin{proof}
Pour toute famille $(Q_{\alpha})_{\alpha \in A}$ de $R$-modules, nous avons un
diagramme commutatif
$$
\xymatrix{
0 \ar[r] & M_1 \otimes_R (\prod_{\alpha} Q_{\alpha}) \ar[r] \ar[d] & M_2
\otimes_R (\prod_{\alpha} Q_{\alpha}) \ar[r] \ar[d] & M_3 \otimes_R
(\prod_{\alpha} Q_{\alpha}) \ar[r] \ar[d] & 0 \\
0 \ar[r] & \prod_{\alpha}(M_1 \otimes Q_{\alpha}) \ar[r] & \prod_{\alpha}(M_2
\otimes Q_{\alpha}) \ar[r] & \prod_{\alpha}(M_3 \otimes Q_{\alpha})\ar[r] & 0
}
$$
dont les lignes sont exactes. Ainsi, (1) et (2) résultent de la
proposition \ref{proposition-ML-tensor}.
\end{proof}

\begin{lemma}
\label{lemma-quotient-module-ML}
Soit $M_1 \to M_2 \to M_3 \to 0$ une suite exacte de $R$-modules.
Si $M_1$ est de type fini et $M_2$ est de Mittag-Leffler, alors $M_3$
est de Mittag-Leffler.
\end{lemma}

\begin{proof}
Pour toute famille $(Q_{\alpha})_{\alpha \in A}$ de $R$-modules,
puisque le produit tensoriel est exact à droite, nous avons un diagramme commutatif
$$
\xymatrix{
M_1 \otimes_R (\prod_{\alpha} Q_{\alpha}) \ar[r] \ar[d] & M_2
\otimes_R (\prod_{\alpha} Q_{\alpha}) \ar[r] \ar[d] & M_3 \otimes_R
(\prod_{\alpha} Q_{\alpha}) \ar[r] \ar[d] & 0 \\
\prod_{\alpha}(M_1 \otimes Q_{\alpha}) \ar[r] & \prod_{\alpha}(M_2
\otimes Q_{\alpha}) \ar[r] & \prod_{\alpha}(M_3 \otimes Q_{\alpha})\ar[r] & 0
}
$$
dont les lignes sont exactes. Par la proposition \ref{proposition-fg-tensor},
la flèche verticale de gauche est surjective. Par la
proposition \ref{proposition-ML-tensor}, la flèche verticale du milieu
est injective. Une chasse au diagramme montre que la flèche verticale de droite
est injective. Ainsi, $M_3$ est de Mittag-Leffler par la
proposition \ref{proposition-ML-tensor}.
\end{proof}


\begin{lemma}
\label{lemma-colimit-universally-injective-ML}
Si $M = \colim M_i$ est la colimite d'un système filtrant de $R$-modules de Mittag-Leffler
$M_i$ dont les morphismes de transition sont universellement injectifs, alors $M$ est de
Mittag-Leffler.
\end{lemma}

\begin{proof}
Soit $(Q_{\alpha})_{\alpha \in A}$ une famille de $R$-modules. Nous devons
montrer que $M \otimes_R (\prod Q_\alpha) \to \prod M \otimes_R Q_\alpha$
est injectif, et nous savons que
$M_i \otimes_R (\prod Q_\alpha) \to \prod M_i \otimes_R Q_\alpha$
est injectif pour tout $i$, voir la proposition \ref{proposition-ML-tensor}.
Comme $\otimes$ commute aux colimites filtrantes, il suffit de montrer que
$\prod M_i \otimes_R Q_\alpha \to \prod M \otimes_R Q_\alpha$
est injectif. Cela est clair, car chacun des morphismes
$M_i \otimes_R Q_\alpha \to M \otimes_R Q_\alpha$ est injectif
par l'hypothèse que les morphismes de transition sont universellement injectifs.
\end{proof}

\begin{lemma}
\label{lemma-direct-sum-ML}
Si $M = \bigoplus_{i \in I} M_i$ est une somme directe de $R$-modules, alors $M$ est
de Mittag-Leffler si et seulement si chaque $M_i$ est de Mittag-Leffler.
\end{lemma}

\begin{proof}
Le sens « seulement si » résulte du lemme \ref{lemma-pure-submodule-ML} (1)
et du fait qu'une suite exacte courte scindée est universellement exacte. La
réciproque résulte du lemme \ref{lemma-colimit-universally-injective-ML},
mais nous pouvons aussi raisonner directement comme suit. Remarquons d'abord que si $I$ est
fini, cela résulte du lemme
\ref{lemma-pure-submodule-ML} (2). Pour $I$ quelconque, si tous les $M_i$ sont
de Mittag-Leffler, nous montrons qu'il en est de même de $M$ en vérifiant la condition (1) de la
proposition \ref{proposition-ML-characterization}.
Soit $f: P \to M$ un morphisme d'un module de présentation finie $P$.
Alors $f$ se factorise sous la forme
$P \xrightarrow{f'} \bigoplus_{i' \in I'} M_{i'} \hookrightarrow
\bigoplus_{i \in I} M_i$
pour une certaine partie finie $I'$ de $I$. Par le cas fini,
$\bigoplus_{i' \in I'} M_{i'}$ est de Mittag-Leffler, et il existe donc
un module de présentation finie $Q$ et un morphisme $g: P \to Q$ tels que $g$
et $f'$ se dominent mutuellement. Alors $g$ et $f$ se dominent également mutuellement.
\end{proof}

\begin{lemma}
\label{lemma-flat-ML-over-ML-ring}
Soit $R \to S$ un homomorphisme d'anneaux. Soit $M$ un $S$-module.
Si $S$ est de Mittag-Leffler comme $R$-module, et si $M$ est plat et de Mittag-Leffler
comme $S$-module, alors $M$ est de Mittag-Leffler comme $R$-module.
\end{lemma}

\begin{proof}
Nous déduisons ceci de la caractérisation de la
proposition \ref{proposition-ML-tensor}.
En effet, supposons que les $Q_\alpha$ forment une famille de $R$-modules.
Considérons le composé
$$
\xymatrix{
M \otimes_R \prod_\alpha Q_\alpha =
M \otimes_S S \otimes_R \prod_\alpha Q_\alpha \ar[d] \\
M \otimes_S \prod_\alpha (S \otimes_R Q_\alpha) \ar[d] \\
\prod_\alpha (M \otimes_S S \otimes_R Q_\alpha) =
\prod_\alpha (M \otimes_R Q_\alpha)
}
$$
La première flèche est injective puisque $M$ est plat sur $S$ et que $S$ est
de Mittag-Leffler sur $R$, et la seconde flèche est injective puisque
$M$ est de Mittag-Leffler sur $S$. Ainsi, $M$ est de Mittag-Leffler sur $R$.
\end{proof}


\section{Anneaux cohérents}
\label{section-coherent}

\noindent
Nous utilisons la discussion sur la commutation de $\prod$ et $\otimes$ pour déterminer
les anneaux pour lesquels les produits de modules plats sont plats. Il s'avère qu'il
s'agit des anneaux dits cohérents. La notion de module $\mathcal{O}_X$ cohérent
sur un espace annelé vous est peut-être plus familière, voir
Modules, section \ref{modules-section-coherent}.

\begin{definition}
\label{definition-coherent}
Soit $R$ un anneau. Soit $M$ un $R$-module.
\begin{enumerate}
\item Nous disons que $M$ est un {\it module cohérent} s'il est de type fini
et si tout sous-module de type fini de $M$ est de présentation finie sur
$R$.
\item Nous disons que $R$ est un {\it anneau cohérent} s'il est cohérent comme module
sur lui-même.
\end{enumerate}
\end{definition}

\noindent
Ainsi, un anneau est cohérent si et seulement si tout idéal de type fini est de
présentation finie comme module.

\begin{example}
\label{example-valuation-ring-coherent}
Un anneau de valuation est un anneau cohérent. En effet, tout idéal non nul de type fini
est principal (lemme \ref{lemma-characterize-valuation-ring}),
donc libre puisqu'un anneau de valuation est intègre, et donc de présentation finie.
\end{example}

\noindent
La catégorie des modules cohérents est abélienne.

\begin{lemma}
\label{lemma-coherent}
Soit $R$ un anneau.
\begin{enumerate}
\item Un sous-module de type fini d'un module cohérent est cohérent.
\item Soit $\varphi : N \to M$ un homomorphisme d'un module de type fini
vers un module cohérent. Alors $\Ker(\varphi)$ est de type fini,
$\Im(\varphi)$ est
cohérent et $\Coker(\varphi)$ est cohérent.
\item Soit $\varphi : N \to M$ un homomorphisme de modules cohérents.
Alors $\Ker(\varphi)$ et $\Coker(\varphi)$ sont des modules
cohérents.
\item Étant donnée une suite exacte courte de $R$-modules
$0 \to M_1 \to M_2 \to M_3 \to 0$, si deux d'entre eux sont cohérents,
le troisième l'est aussi.
\end{enumerate}
\end{lemma}

\begin{proof}
La première assertion résulte immédiatement de la définition.

\medskip\noindent
Soit $\varphi : N \to M$ satisfaisant aux hypothèses de (2). Tout d'abord,
$\Im(\varphi)$ est de type fini,
donc cohérent par (1). En particulier, $\Im(\varphi)$ est de présentation finie ;
en appliquant le lemme \ref{lemma-extension} à la suite exacte
$0 \to \Ker(\varphi) \to N \to \Im(\varphi) \to 0$, nous voyons que $\Ker(\varphi)$
est de type fini.
Pour montrer que $\Coker(\varphi)$ est cohérent, soit $E \subset \Coker(\varphi)$
un sous-module de type fini, et soit $E'$ son image réciproque dans $M$.
De la suite
exacte $0 \to \Im(\varphi) \to E' \to E \to 0$ et puisque $\Ker(\varphi)$
est de type fini, nous concluons par le lemme \ref{lemma-extension} que $E' \subset M$ est
de type fini,
donc de présentation finie puisque $M$ est cohérent. La même suite exacte
montre alors que $E$ est de présentation finie, d'où l'assertion.

\medskip\noindent
La partie (3) résulte immédiatement de (1) et (2).

\medskip\noindent
Soit $0 \to M_1 \xrightarrow{i} M_2 \xrightarrow{p} M_3 \to 0$
une suite exacte courte de $R$-modules comme dans (4). Il reste
à montrer que si $M_1$ et $M_3$ sont cohérents,
alors $M_2$ l'est. Par le lemme \ref{lemma-extension},
nous voyons que $M_2$ est de type fini. Soit $N_2 \subset M_2$ un sous-module
de type fini.
Posons $N_3 = p(N_2) \subset M_3$ et $N_1 = i^{-1}(N_2) \subset M_1$.
Nous avons une suite exacte $0 \to N_1 \to N_2 \to N_3 \to 0$. Il est clair que $N_3$
est de type fini (comme quotient de $N_2$), donc de présentation finie (comme sous-module
de type fini de $M_3$). Il résulte du lemme \ref{lemma-extension} (5) que
$N_1$ est de type fini,
donc de présentation finie (comme sous-module de type fini de $M_1$). Nous concluons par le
lemme \ref{lemma-extension} (2) que $N_2$ est de présentation finie.
\end{proof}

\begin{lemma}
\label{lemma-coherent-ring}
Soit $R$ un anneau. Si $R$ est cohérent, alors un module est cohérent
si et seulement s'il est de présentation finie.
\end{lemma}

\begin{proof}
Il est clair qu'un module cohérent est de présentation finie (sur tout anneau).
Réciproquement, si $R$ est cohérent, alors $R^{\oplus n}$ est cohérent, de même que
le conoyau de tout morphisme $R^{\oplus m} \to R^{\oplus n}$, voir le
lemme \ref{lemma-coherent}.
\end{proof}

\begin{lemma}
\label{lemma-Noetherian-coherent}
Un anneau noethérien est un anneau cohérent.
\end{lemma}

\begin{proof}
Par le
lemme \ref{lemma-Noetherian-finite-type-is-finite-presentation},
tout $R$-module de type fini est de présentation finie. En particulier, tout idéal de $R$
est de présentation finie.
\end{proof}

\begin{proposition}
\label{proposition-characterize-coherent}
\begin{reference}
C'est \cite[Theorem 2.1]{Chase}.
\end{reference}
Soit $R$ un anneau. Les assertions suivantes sont équivalentes :
\begin{enumerate}
\item $R$ est cohérent,
\item tout produit de $R$-modules plats est plat, et
\item pour tout ensemble $A$, le module $R^A$ est plat.
\end{enumerate}
\end{proposition}

\begin{proof}
Supposons $R$ cohérent, et soit $Q_\alpha$, $\alpha \in A$, un ensemble de
$R$-modules plats. Nous devons montrer que
$I \otimes_R \prod_\alpha Q_\alpha \to \prod Q_\alpha$ est injectif
pour tout idéal de type fini $I$ de $R$, voir le
lemme \ref{lemma-flat}.
Puisque $R$ est cohérent, $I$ est un $R$-module de présentation finie.
Ainsi, $I \otimes_R \prod_\alpha Q_\alpha = \prod I \otimes_R Q_\alpha$ par la
proposition \ref{proposition-fp-tensor}.
L'injectivité voulue résulte de ce que $I \otimes_R Q_\alpha \to Q_\alpha$
est injectif par platitude de $Q_\alpha$.

\medskip\noindent
L'implication (2) $\Rightarrow$ (3) est triviale.

\medskip\noindent
Supposons que le $R$-module $R^A$ soit plat pour tout ensemble $A$. Soit $I$
un idéal de type fini de $R$. Alors $I \otimes_R R^A \to R^A$
est injectif par hypothèse. Par la
proposition \ref{proposition-fg-tensor}
et la finitude de $I$, l'image est égale à $I^A$. Ainsi,
$I \otimes_R R^A = I^A$ pour tout ensemble $A$, et nous concluons que $I$
est de présentation finie par la
proposition \ref{proposition-fp-tensor}.
\end{proof}







\section{Exemples et contre-exemples de modules de Mittag-Leffler}
\label{section-examples}

\noindent
Nous terminons cette section par quelques exemples et contre-exemples de modules de
Mittag-Leffler.

\begin{example}
\label{example-ML}
Modules de Mittag-Leffler.
\begin{enumerate}
\item Tout module de présentation finie est de Mittag-Leffler. Cela résulte, par
exemple, de la proposition \ref{proposition-ML-characterization} (1). Plus
généralement, un module de type fini est de Mittag-Leffler si et
seulement s'il est de présentation finie. Cela résulte des propositions
\ref{proposition-fg-tensor}, \ref{proposition-fp-tensor} et
\ref{proposition-ML-tensor}.
\item Un module libre est de Mittag-Leffler puisqu'il satisfait la condition (1) de la
proposition \ref{proposition-ML-characterization}.
\item Par l'exemple précédent et le lemme \ref{lemma-direct-sum-ML},
les modules projectifs sont de Mittag-Leffler.
\end{enumerate}
\end{example}

\noindent
Nous voulons aussi ajouter à notre liste d'exemples les anneaux de séries formelles sur un
anneau noethérien $R$. Ce sera une conséquence du lemme suivant.

\begin{lemma}
\label{lemma-flat-ML-criterion}
Soit $M$ un $R$-module plat. Les assertions suivantes sont équivalentes :
\begin{enumerate}
\item $M$ est de Mittag-Leffler, et
\item si $F$ est un $R$-module libre de type fini et si
$x \in F \otimes_R M$, alors il existe un plus petit sous-module $F'$ de $F$
tel que $x \in F' \otimes_R M$.
\end{enumerate}
\end{lemma}

\begin{proof}
L'implication (1) $\Rightarrow$ (2) est un cas particulier du
lemme \ref{lemma-minimal-contains}. Supposons (2).
Par le théorème \ref{theorem-lazard}, nous pouvons écrire $M$ comme la colimite
$M = \colim_{i \in I} M_i$ d'un système filtrant $(M_i, f_{ij})$ de
$R$-modules libres de type fini.
Par la remarque \ref{remark-flat-ML}, il suffit de montrer que le système projectif
$(\Hom_R(M_i, R), \Hom_R(f_{ij}, R))$ est de Mittag-Leffler. Autrement
dit,
fixons $i \in I$ et, pour $j \geq i$, soit $Q_j$ l'image de
$\Hom_R(M_j, R)
\to \Hom_R(M_i, R)$ ; nous devons montrer que les $Q_j$ se stabilisent.

\medskip\noindent
Puisque $M_i$ est libre de type fini, nous pouvons faire l'identification
$\Hom_R(M_i, M_j) = \Hom_R(M_i, R) \otimes_R  M_j$ pour tout $j$.
En utilisant le
fait que les $M_j$ sont libres, il s'ensuit que, pour $j \geq i$, $Q_j$ est le
plus petit sous-module de $\Hom_R(M_i, R)$ tel que $f_{ij} \in Q_j
\otimes_R
M_j$. Sous l'identification $\Hom_R(M_i, M) = \Hom_R(M_i, R)
\otimes_R
M$, le morphisme canonique $f_i: M_i \to M$ appartient à $\Hom_R(M_i, R)
\otimes_R M$. Par l'hypothèse sur $M$, il existe un plus petit sous-module
$Q$ de $\Hom_R(M_i, R)$ tel que $f_i \in Q \otimes_R M$. Nous allons
montrer
que les $Q_j$ se stabilisent à $Q$.

\medskip\noindent
Pour $j \geq i$, nous avons un diagramme commutatif
$$
\xymatrix{
Q_j \otimes_R M_j \ar[r] \ar[d] & \Hom_R(M_i, R) \otimes_R M_j
\ar[d] \\
Q_j \otimes_R M \ar[r] & \Hom_R(M_i, R) \otimes_R M.
}
$$
Puisque $f_{ij} \in Q_j \otimes_R M_j$ s'envoie sur $f_i \in \Hom_R(M_i, R)
\otimes_R M$, il s'ensuit que $f_i \in Q_j \otimes_R M$. Ainsi, par le
choix de $Q$, nous avons $Q \subset Q_j$ pour tout $j \geq i$.

\medskip\noindent
Puisque les $Q_j$ sont décroissants et que $Q \subset Q_j$ pour tout $j \geq i$, pour montrer
que les $Q_j$ se stabilisent à $Q$, il suffit de trouver un $j \geq i$ tel que $Q_j
\subset Q$. Comme élément de
$$
\Hom_R(M_i, R) \otimes_R M = \colim_{j \in J}
(\Hom_R(M_i, R) \otimes_R
M_j),
$$
$f_i$ est la colimite des $f_{ij}$ pour $j \geq i$, et $f_i$ appartient également au
sous-module
$$
\colim_{j \in J} (Q \otimes_R M_j) \subset \colim_{j \in J}
(\Hom_R(M_i, R)
\otimes_R M_j).
$$
Il s'ensuit que, pour un certain $j \geq i$, $f_{ij}$ appartient à $Q \otimes_R M_j$.
Puisque $Q_j$ est le plus petit sous-module de $\Hom_R(M_i, R)$ tel que $f_{ij}
\in
Q_j \otimes_R M_j$, nous concluons que $Q_j\subset Q$.
\end{proof}

\begin{lemma}
\label{lemma-product-over-Noetherian-ring}
Soient $R$ un anneau noethérien et $A$ un ensemble.
Alors $M = R^A$ est un $R$-module plat et de Mittag-Leffler.
\end{lemma}

\begin{proof}
En combinant le
lemme \ref{lemma-Noetherian-coherent}
et la
proposition \ref{proposition-characterize-coherent},
nous voyons que $M$ est plat sur $R$. Nous montrons que $M$ satisfait la condition du
lemme \ref{lemma-flat-ML-criterion}.
Soit $F$ un $R$-module libre de type fini. Si $F'$ est un sous-module quelconque de $F$, alors il
est de présentation finie puisque $R$ est noethérien. Par la
proposition \ref{proposition-fp-tensor},
nous avons donc un diagramme commutatif
$$
\xymatrix{
F' \otimes_R M \ar[r] \ar[d]^{\cong} & F \otimes_R M \ar[d]^{\cong} \\
(F')^A \ar[r] & F^A
}
$$
par lequel nous pouvons identifier le morphisme $F' \otimes_R M \to F \otimes_R M$
à $(F')^A \to F^A$. Ainsi, si $x \in F \otimes_R M$ correspond à
$(x_\alpha) \in F^A$, alors le sous-module $F'$ de $F$ engendré par les
$x_\alpha$ est le plus petit sous-module de $F$ tel que $x \in F' \otimes_R M$.
\end{proof}

\begin{lemma}
\label{lemma-power-series-ML}
Soient $R$ un anneau noethérien et $n$ un entier positif. Alors le $R$-module
$M = R[[t_1, \ldots, t_n]]$ est plat et de Mittag-Leffler.
\end{lemma}

\begin{proof}
Comme $R$-module, nous avons $M = R^A$ pour un ensemble (dénombrable) $A$.
Ce lemme est donc un cas particulier du
lemme \ref{lemma-product-over-Noetherian-ring}.
\end{proof}

\begin{example}
\label{example-not-ML}
Modules qui ne sont pas de Mittag-Leffler.
\begin{enumerate}
\item Par l'exemple \ref{example-Q-not-ML} et la
proposition \ref{proposition-ML-tensor}, $\mathbf{Q}$ n'est pas un
$\mathbf{Z}$-module de Mittag-Leffler.
\item Nous démontrons ci-dessous (théorème \ref{theorem-projectivity-characterization}) que,
pour un module plat et dénombrablement engendré, la projectivité équivaut à être
de Mittag-Leffler. Ainsi, tout module $M$ plat, dénombrablement engendré et non projectif
est un exemple de module qui n'est pas de Mittag-Leffler. Pour un tel exemple, voir la
remarque \ref{remark-warning}.
\item Soit $k$ un corps. Soit $R = k[[x]]$. Le $R$-module
$M = \prod_{n \in \mathbf{N}} R/(x^n)$ n'est pas de Mittag-Leffler.
En effet, considérons l'élément $\xi = (\xi_1, \xi_2, \xi_3, \ldots)$
défini par $\xi_{2^m} = x^{2^{m - 1}}$ et $\xi_n = 0$ sinon, de sorte que
$$
\xi = (0, x, 0, x^2, 0, 0, 0, x^4, 0, 0, 0, 0, 0, 0, 0, x^8, \ldots)
$$
Alors l'annulateur de $\xi$ dans $M/x^{2^m}M$ est engendré par $x^{2^{m - 1}}$
pour $m \gg 0$. Mais si $M$ était de Mittag-Leffler, il existerait un
$R$-module de type fini $Q$ et un élément $\xi' \in Q$ tels que l'annulateur
de $\xi'$ dans $Q/x^l Q$ coïncide avec l'annulateur de $\xi$ dans $M/x^l M$
pour tout $l \geq 1$, voir la
proposition \ref{proposition-ML-characterization} (1).
Vous pouvez alors montrer qu'il existe un entier $a \geq 0$ tel que
l'annulateur de $\xi'$ dans $Q/x^l Q$ soit engendré
soit par $x^a$, soit par $x^{l - a}$ pour tout $l \gg 0$ (selon que
$\xi' \in Q$ est de torsion ou non). En combinant ce qui précède,
on obtiendrait, pour tout $l = 2^m >> 0$, l'égalité $a = l/2$ ou $l - a = l/2$,
ce qui est absurde.
\item Le même argument montre que le complété $(x)$-adique de
$\bigoplus_{n \in \mathbf{N}} R/(x^n)$ n'est pas de Mittag-Leffler sur
$R = k[[x]]$ (indication : $\xi$ est en fait un élément de ce complété).
\item Soit $R = k[a, b]/(a^2, ab, b^2)$. Soit $S$ la $R$-algèbre de présentation finie
présentée par $S = R[t]/(at - b)$. Alors, comme $R$-module,
$S$ est dénombrablement engendré et indécomposable (détails omis).
D'autre part, $R$ est local artinien, donc local complet,
et donc un anneau local hensélien, voir le
lemme \ref{lemma-complete-henselian}.
Si $S$ était de Mittag-Leffler comme $R$-module, alors ce serait une
somme directe de $R$-modules de type fini par le
lemme \ref{lemma-split-ML-henselian}.
Nous concluons donc que $S$ n'est pas de Mittag-Leffler comme $R$-module.
\end{enumerate}
\end{example}




\section{Modules de Mittag-Leffler dénombrablement engendrés}
\label{section-ML-countable}

\noindent
Il s'avère que les modules de Mittag-Leffler dénombrablement engendrés ont une
structure particulièrement simple.

\begin{lemma}
\label{lemma-ML-countable-colimit}
Soit $M$ un $R$-module. Écrivons $M = \colim_{i \in I} M_i$, où $(M_i,
f_{ij})$ est un système filtrant de $R$-modules de présentation finie. Si $M$ est
de Mittag-Leffler et dénombrablement engendré, alors il existe une partie filtrante dénombrable
$I' \subset I$ telle que $M \cong \colim_{i \in I'} M_i$.
\end{lemma}

\begin{proof}
Soit $x_1, x_2, \ldots$ une famille dénombrable de générateurs de $M$. Pour chaque $x_n$,
choisissons $i \in I$ tel que $x_n$ appartienne à l'image du morphisme canonique $f_i:
M_i \to M$ ; soit $I'_{0} \subset I$ l'ensemble de tous ces $i$. Puisque
$M$ est de Mittag-Leffler, pour chaque $i \in I'_{0}$, nous pouvons choisir $j \in I$
tel que $j \geq i$ et que $f_{ij}: M_i \to M_j$ se factorise par
$f_{ik}: M_i \to M_k$ pour tout $k \geq i$ (condition (3) de la proposition
\ref{proposition-ML-characterization}) ; soit $I'_1$ la réunion de $I'_0$ et de
tous ces $j$. Puisque $I'_1$ est dénombrable, nous pouvons l'agrandir en un
ensemble préordonné filtrant dénombrable $I'_{2} \subset I$. Nous pouvons maintenant appliquer à $I'_{2}$
le même procédé qu'à $I'_{0}$ pour obtenir un nouvel ensemble dénombrable $I'_{3} \subset
I$. Puis nous agrandissons $I'_{3}$ en un ensemble préordonné filtrant dénombrable $I'_{4}$. En continuant
ainsi --- en ajoutant un $j$ comme dans la proposition
\ref{proposition-ML-characterization} (3) pour chaque $ i \in I'_{\ell}$ si $\ell$
est pair, et en agrandissant $I'_{\ell}$ en un ensemble préordonné filtrant si $\ell$ est impair --- nous obtenons une
suite de parties $I'_{\ell} \subset I$ pour $\ell \geq 0$. La réunion $I' =
\bigcup I'_{\ell}$ satisfait :
\begin{enumerate}
\item $I'$ est dénombrable et filtrant ;
\item chaque $x_n$ appartient à l'image de $f_i: M_i \to M$ pour un certain $i
\in I'$ ;
\item si $i \in I'$, alors il existe $j \in I'$ tel que $j \geq i$ et que $f_{ij}:
M_i \to M_j$ se factorise par $f_{ik}: M_i \to M_k$ pour tout
$k \in I$ tel que $k \geq i$. En particulier, $\Ker(f_{ik}) \subset \Ker(f_{ij})$
pour $k \geq i$.
\end{enumerate}
Nous affirmons que le morphisme canonique $\colim_{i \in I'} M_i \to
\colim_{i
\in I} M_i = M$ est un isomorphisme. Par (2), il est surjectif. Pour l'injectivité,
supposons que $x \in \colim_{i \in I'} M_i$ s'envoie sur $0$ dans $\colim_{i \in
I} M_i$.
En représentant $x$ par un élément $\tilde{x} \in M_i$ pour un certain $i \in I'$, cela
signifie que $f_{ik}(\tilde{x}) = 0$ pour un certain $k \in I, k \geq i$. Mais, par
(3), il existe $j \in I', j \geq i,$ tel que $f_{ij}(\tilde{x}) = 0$. Ainsi, $x
= 0$ dans $\colim_{i \in I'} M_i$.
\end{proof}

\noindent
Le lemme \ref{lemma-ML-countable-colimit}
implique qu'un module de Mittag-Leffler $M$ dénombrablement engendré sur
$R$ est la colimite d'un système
$$
M_1 \to M_2 \to M_3 \to M_4 \to \ldots
$$
où chaque $M_n$ est un $R$-module de présentation finie. Pour le voir, raisonnons comme dans la
démonstration du
lemme \ref{lemma-ML-limit-nonempty}
pour constater qu'un ensemble préordonné filtrant dénombrable possède une partie cofinale
isomorphe à $(\mathbf{N}, \geq)$. Supposons que
$R = k[x_1, x_2, x_3, \ldots]$ et $M = R/(x_i)$. Alors $M$ est de type
fini, mais n'est pas de présentation finie, et n'est donc pas de Mittag-Leffler (voir
l'exemple \ref{example-ML}, partie (1)).
Bien entendu, on peut écrire $M = \colim_n M_n$ en prenant
$M_n = R/(x_1, \ldots, x_n)$ ; ainsi, la possibilité d'écrire
$M$ comme une telle colimite n'implique pas que $M$ soit de Mittag-Leffler.

\begin{lemma}
\label{lemma-ML-countable}
Soit $R$ un anneau.
Soit $M$ un $R$-module.
Supposons $M$ de Mittag-Leffler et dénombrablement engendré.
Pour tout morphisme de $R$-modules $f : P \to M$, où $P$ est de type fini, il
existe un endomorphisme $\alpha : M \to M$ tel que
\begin{enumerate}
\item $\alpha : M \to M$ se factorise par un $R$-module de présentation finie, et
\item $\alpha \circ f = f$.
\end{enumerate}
\end{lemma}

\begin{proof}
Écrivons $M = \colim_{i \in I} M_i$ comme colimite filtrante de $R$-modules de
présentation finie, avec $I$ dénombrable, voir le
lemme \ref{lemma-ML-countable-colimit}.
Les morphismes de transition sont notés $f_{ij}$ et nous utilisons $f_i : M_i \to M$
pour désigner les morphismes canoniques vers $M$. Posons $N = \prod_{s \in I} M_s$. Notons
$$
M_i^* = \Hom_R(M_i, N) = \prod\nolimits_{s \in I} \Hom_R(M_i, M_s)
$$
de sorte que $(M_i^*)$ est un système projectif de $R$-modules indexé par $I$.
Remarquons que $\Hom_R(M, N) = \lim M_i^*$.
Comme $M$ est de Mittag-Leffler, nous trouvons, pour tout
$i \in I$, un indice $k(i) \geq i$ tel que
$$
E_i := \bigcap\nolimits_{i' \geq i} \Im(M_{i'}^* \to M_i^*)
=
\Im(M_{k(i)}^* \to M_i^*)
$$
Choisissons et fixons $j \in I$ tel que $\Im(P \to M) \subset \Im(M_j \to M)$.
C'est possible puisque $P$ est de type fini. Posons $k = k(j)$.
Soit
$x = (0, \ldots, 0, \text{id}_{M_k}, 0, \ldots, 0) \in M_k^*$ et
remarquons qu'il s'envoie sur $y = (0, \ldots, 0, f_{jk}, 0, \ldots, 0) \in M_j^*$.
Par notre choix de $k$, nous voyons que $y \in E_j$. Par
l'exemple \ref{example-ML-surjective-maps},
les morphismes de transition $E_i \to E_j$ sont surjectifs pour tout $i \geq j$
et $\lim E_i = \lim M_i^* = \Hom_R(M, N)$. Ainsi, le
lemme \ref{lemma-ML-limit-nonempty}
garantit l'existence d'un élément $z \in \Hom_R(M, N)$
qui s'envoie sur $y$ dans $E_j \subset M_j^*$. Soit $z_k$ la $k$-ième composante
de $z$. Alors $z_k : M \to M_k$ est un homomorphisme tel que
$$
\xymatrix{
M \ar[r]_{z_k} & M_k \\
M_j \ar[ru]_{f_{jk}} \ar[u]^{f_j}
}
$$
commute. Soit $\alpha : M \to M$ le composé
$f_k \circ z_k : M \to M_k \to M$.
Alors $\alpha$ se factorise par un module de présentation finie par construction et
$\alpha \circ f_j = f_j$. Puisque l'image de $f$ est contenue dans l'image de
$f_j$, cela implique aussi que $\alpha \circ f = f$.
\end{proof}

\noindent
Nous verrons plus loin (voir le
lemme \ref{lemma-split-ML-henselian})
que le
lemme \ref{lemma-ML-countable}
signifie qu'un module de Mittag-Leffler dénombrablement engendré sur un anneau local
hensélien est une somme directe de modules de présentation finie.




\section{Caractérisation des modules projectifs}
\label{section-characterize-projective}

\noindent
L'objectif de cette section est de démontrer qu'un module est projectif si et seulement s'il
est plat, de Mittag-Leffler et somme directe de modules dénombrablement engendrés
(théorème \ref{theorem-projectivity-characterization} ci-dessous).

\begin{lemma}
\label{lemma-countgen-projective}
Soit $M$ un $R$-module. Si $M$ est plat, de Mittag-Leffler et
dénombrablement engendré, alors $M$ est projectif.
\end{lemma}

\begin{proof}
Par le théorème de Lazard (théorème \ref{theorem-lazard}), nous pouvons écrire $M =
\colim_{i
\in I} M_i$ pour un système filtrant de $R$-modules libres de type fini $(M_i, f_{ij})$
indexé par un ensemble $I$. Par le lemme \ref{lemma-ML-countable-colimit}, nous pouvons supposer
$I$ dénombrable. Soit maintenant
$$
0 \to N_1 \to N_2 \to N_3 \to 0
$$
une suite exacte de $R$-modules. Nous devons montrer que l'application de
$\Hom_R(M, -)$
préserve l'exactitude. Puisque $M_i$ est libre de type fini,
$$
0 \to \Hom_R(M_i, N_1) \to \Hom_R(M_i, N_2) \to
\Hom_R(M_i, N_3) \to 0
$$
est exacte pour tout $i$. Puisque $M$ est de Mittag-Leffler, $(\Hom_R(M_i,
N_{1}))$ est
un système projectif de Mittag-Leffler. Par le lemme \ref{lemma-ML-exact-sequence},
$$
0 \to \lim_{i \in I} \Hom_R(M_i, N_1) \to
\lim_{i \in I} \Hom_R(M_i, N_2) \to
\lim_{i \in I} \Hom_R(M_i, N_3) \to 0
$$
est donc exacte. Mais, pour tout $R$-module $N$, il existe un isomorphisme fonctoriel
$\Hom_R(M, N) \cong \lim_{i \in I} \Hom_R(M_i, N)$, de sorte que
$$
0 \to \Hom_R(M, N_1) \to \Hom_R(M, N_2) \to
\Hom_R(M, N_3) \to 0
$$
est exacte.
\end{proof}

\begin{remark}
\label{remark-characterize-projective}
Le lemme \ref{lemma-countgen-projective} n'est plus vrai sans l'hypothèse
d'engendrement dénombrable. Par exemple, le $\mathbf Z$-module $M =
\mathbf{Z}[[x]]$ est plat et de Mittag-Leffler, mais n'est pas projectif. Il est
de Mittag-Leffler par le lemme \ref{lemma-power-series-ML}. Les sous-groupes des groupes abéliens
libres sont libres ; un $\mathbf Z$-module projectif est donc en fait libre, et il en va de même
de ses sous-modules. Ainsi, pour montrer que $M$ n'est pas projectif, il suffit de construire
un sous-module non libre. Fixons un nombre premier $p$ et considérons le sous-module $N$
formé des séries formelles $f(x) = \sum a_i x^i$ telles que, pour tout entier $m
\geq 1$, $p^m$ divise $a_i$ pour tous les $i$ sauf un nombre fini. Alors $\sum a_i p^i
x^i$ appartient à $N$ pour tous $a_i \in \mathbf{Z}$, et $N$ n'est donc pas dénombrable. Ainsi, si
$N$ était libre, il serait de rang non dénombrable, et la dimension de $N/pN$ sur
$\mathbf{Z}/p$ ne serait pas dénombrable. Ce n'est pas le cas, car les éléments $x^i \in
N/pN$, pour $i \geq 0$, engendrent $N/pN$.
\end{remark}

\begin{theorem}
\label{theorem-projectivity-characterization}
Soit $M$ un $R$-module. Alors $M$ est projectif si et seulement si
\begin{enumerate}
\item $M$ est plat,
\item $M$ est de Mittag-Leffler,
\item $M$ est une somme directe de $R$-modules dénombrablement engendrés.
\end{enumerate}
\end{theorem}

\begin{proof}
Supposons d'abord que $M$ soit projectif. Alors $M$ est un facteur direct d'un module
libre, de sorte que $M$ est plat et de Mittag-Leffler, puisque ces propriétés passent aux facteurs
directs. Par le théorème de Kaplansky (théorème \ref{theorem-projective-direct-sum}),
$M$ satisfait (3).

\medskip\noindent
Réciproquement, supposons que $M$ satisfasse (1)-(3). Puisque les propriétés d'être plat et de Mittag-Leffler
passent aux facteurs directs, $M$ est une somme directe de $R$-modules plats, de Mittag-Leffler et
dénombrablement engendrés.
Le lemme \ref{lemma-countgen-projective}
implique que $M$ est une somme directe de modules projectifs. Ainsi, $M$ est projectif.
\end{proof}

\begin{lemma}
\label{lemma-ML-ui-descent}
Soit $f: M \to N$ un morphisme universellement injectif de $R$-modules. Supposons que
$M$ soit une somme directe de $R$-modules dénombrablement engendrés, et supposons que $N$ soit plat
et de Mittag-Leffler. Alors $M$ est projectif.
\end{lemma}

\begin{proof}
Par les
lemmes \ref{lemma-ui-flat-domain} et
\ref{lemma-pure-submodule-ML},
$M$ est plat et de Mittag-Leffler ; la conclusion résulte donc du théorème
\ref{theorem-projectivity-characterization}.
\end{proof}

\begin{lemma}
\label{lemma-universally-injective-submodule-powerseries}
Soient $R$ un anneau noethérien et $M$ un $R$-module. Supposons que $M$ soit une
somme directe de $R$-modules dénombrablement engendrés, et supposons qu'il existe un
morphisme universellement injectif $M \to R[[t_1, \ldots, t_n]]$ pour un certain $n$.
Alors $M$ est projectif.
\end{lemma}

\begin{proof}
Cela résulte des
lemmes \ref{lemma-ML-ui-descent} et
\ref{lemma-power-series-ML}.
\end{proof}



\section{Montée de propriétés des modules}
\label{section-ascent}

\noindent
Toutes les propriétés d'un module figurant dans le théorème
\ref{theorem-projectivity-characterization} montent le long d'homomorphismes d'anneaux quelconques :

\begin{lemma}
\label{lemma-ascend-properties-modules}
Soit $R \to S$ un homomorphisme d'anneaux. Soit $M$ un $R$-module. Alors :
\begin{enumerate}
\item Si $M$ est plat, alors le $S$-module $M \otimes_R S$ est plat.
\item Si $M$ est de Mittag-Leffler, alors le $S$-module $M \otimes_R S$ est
de Mittag-Leffler.
\item Si $M$ est une somme directe de $R$-modules dénombrablement engendrés, alors le
$S$-module $M \otimes_R S$ est une somme directe de $S$-modules dénombrablement engendrés.
\item Si $M$ est projectif, alors le $S$-module $M \otimes_R S$ est projectif.
\end{enumerate}
\end{lemma}

\begin{proof}
Toutes ces assertions sont évidentes, sauf (2). Pour celle-ci, utilisons la formulation (3) de la propriété
de Mittag-Leffler donnée dans la proposition \ref{proposition-ML-characterization} ainsi que le
fait que la tensorisation commute aux colimites. On peut aussi
utiliser la caractérisation de la proposition \ref{proposition-ML-tensor}.
\end{proof}



\section{Descente de propriétés des modules}
\label{section-descent}

\noindent
Nous abordons la descente fidèlement plate des propriétés du théorème
\ref{theorem-projectivity-characterization} qui caractérisent la projectivité.
En présence de la platitude, la propriété d'être un module de Mittag-Leffler
descend :

\begin{lemma}
\label{lemma-ffdescent-ML}
\begin{reference}
Courriel de Juan Pablo Acosta Lopez daté du 20 décembre 2014.
\end{reference}
Soit $R \to S$ un homomorphisme d'anneaux fidèlement plat. Soit $M$ un $R$-module. Si le
$S$-module $M \otimes_R S$ est de Mittag-Leffler, alors $M$ est de Mittag-Leffler.
\end{lemma}

\begin{proof}
Écrivons $M = \colim_{i\in I} M_i$ comme colimite filtrante
de $R$-modules de présentation finie $M_i$.
En utilisant la proposition \ref{proposition-ML-characterization}, nous voyons que nous
devons montrer que, pour tout $i \in I$, il existe $i \leq j$, $j\in I$,
tel que $M_i\rightarrow M_j$ domine $M_i\rightarrow M$.

\medskip\noindent
Soit $N$ la somme amalgamée
$$
\xymatrix{
M_i  \ar[r] \ar[d] & M_j \ar[d] \\
M \ar[r] & N
}
$$
Le lemme équivaut alors à l'existence de $j$ tel que
$M_j\rightarrow N$ soit universellement injectif, voir le
lemme \ref{lemma-domination-universally-injective}.
Observons que sa tensorisation par $S$,
$$
\xymatrix{
M_i\otimes_R S  \ar[r] \ar[d] & M_j\otimes_R S \ar[d] \\
M\otimes_R S \ar[r] & N\otimes_R S
}
$$
est un diagramme de somme amalgamée. Ainsi, puisque
$M \otimes_R S = \colim_{i\in I} M_i \otimes_R S$
exprime $M\otimes_R S$ comme une colimite de $S$-modules de
présentation finie, et puisque $M\otimes_R S$ est de Mittag-Leffler,
il existe $j \geq i$ tel que $M_j\otimes_R S\rightarrow N\otimes_R S$
soit universellement injectif. Comme $R\rightarrow S$ est fidèlement plat,
nous concluons que $M_j\rightarrow N$ est lui aussi universellement injectif.
\end{proof}

\begin{lemma}
\label{lemma-ffdescent-countable}
Soit $R \to S$ un homomorphisme d'anneaux fidèlement plat. Soit $M$ un $R$-module. Si
le $S$-module $M \otimes_R S$ est dénombrablement engendré, alors $M$
est dénombrablement engendré.
\end{lemma}

\begin{proof}
Supposons que $M \otimes_R S$ soit engendré par les éléments
$y_i$, $i = 1, 2, 3, \ldots$. Écrivons $y_i = \sum_{j = 1, \ldots, n_i}
x_{ij} \otimes s_{ij}$ pour certains $n_i \geq 0$, $x_{ij} \in M$
et $s_{ij} \in S$. Notons $M' \subset M$ le sous-module engendré
par la famille dénombrable des éléments $x_{ij}$. Alors
$M' \otimes_R S \to M \otimes_R S$ est surjectif puisque son
image contient les générateurs $y_i$. Comme $S$ est fidèlement plat
sur $R$, nous concluons que $M' = M$, comme voulu.
\end{proof}

\noindent
À ce stade, la descente fidèlement plate des modules projectifs
dénombrablement engendrés s'obtient aisément.

\begin{lemma}
\label{lemma-ffdescent-countable-projectivity}
Soit $R \to S$ un homomorphisme d'anneaux fidèlement plat. Soit $M$ un $R$-module.
Si le $S$-module $M \otimes_R S$ est dénombrablement engendré et projectif,
alors $M$ est dénombrablement engendré et projectif.
\end{lemma}

\begin{proof}
Cela résulte des lemmes \ref{lemma-descend-properties-modules},
\ref{lemma-ffdescent-ML} et \ref{lemma-ffdescent-countable}, et du
théorème \ref{theorem-projectivity-characterization}.
\end{proof}

\noindent
Il ne reste qu'à utiliser un dévissage pour ramener la descente de la projectivité dans le
cas général au cas dénombrablement engendré. Commençons par deux lemmes simples.

\begin{lemma}
\label{lemma-lift-countably-generated-submodule}
Soient $R \to S$ un homomorphisme d'anneaux, $M$ un $R$-module et $Q$ un
$S$-sous-module dénombrablement engendré de $M \otimes_R S$. Alors il existe un
$R$-sous-module dénombrablement engendré $P$ de $M$ tel que
$\Im(P \otimes_R S \to M \otimes_R S)$ contienne $Q$.
\end{lemma}

\begin{proof}
Soient $y_1, y_2, \ldots$ des générateurs de $Q$ et écrivons $y_j = \sum_k x_{jk}
\otimes s_{jk}$ pour certains $x_{jk} \in M$ et $s_{jk} \in S$. Prenons alors pour $P$
le sous-module de $M$ engendré par les $x_{jk}$.
\end{proof}

\begin{lemma}
\label{lemma-adapted-submodule}
Soient $R \to S$ un homomorphisme d'anneaux et $M$ un $R$-module. Supposons que $M
\otimes_R S = \bigoplus_{i \in I} Q_i$ soit une somme directe de
$S$-modules dénombrablement engendrés $Q_i$. Si $N$ est un sous-module dénombrablement engendré de $M$, alors
il existe un sous-module dénombrablement engendré $N'$ de $M$ tel que $N' \supset N$
et $\Im(N' \otimes_R S \to M \otimes_R S) =
\bigoplus_{i \in I'} Q_i$ pour une certaine partie $I' \subset I$.
\end{lemma}

\begin{proof}
Soit $N'_0 = N$. Nous construisons par récurrence une suite croissante de sous-modules
dénombrablement engendrés $N'_{\ell} \subset M$ pour $\ell = 0, 1, 2, \ldots$
telle que : si $I'_{\ell}$ est l'ensemble des $i \in I$ tels que la projection de
$\Im(N'_{\ell} \otimes_R S \to M \otimes_R S)$ sur $Q_i$ est
non nulle, alors $\Im(N'_{\ell + 1} \otimes_R S \to M \otimes_R
S)$ contient $Q_i$ pour tout $i \in I'_{\ell}$. Pour construire $N'_{\ell + 1}$
à partir de $N'_\ell$, soit $Q$ la somme des $Q_i$ (en nombre dénombrable) pour
$i \in I'_{\ell}$, choisissons $P$ comme dans le lemme
\ref{lemma-lift-countably-generated-submodule}, puis posons $N'_{\ell + 1} =
N'_{\ell} + P$. Une fois les $N'_{\ell}$ construits, prenons simplement $N' =
\bigcup_{\ell} N'_{\ell}$ et $I' = \bigcup_{\ell} I'_{\ell}$.
\end{proof}

\begin{theorem}
\label{theorem-ffdescent-projectivity}
Soit $R \to S$ un homomorphisme d'anneaux fidèlement plat. Soit $M$ un $R$-module.
Si le $S$-module $M \otimes_R S$ est projectif, alors $M$ est projectif.
\end{theorem}

\begin{proof}
Nous allons construire un dévissage de Kaplansky de $M$ pour montrer que c'est une
somme directe de modules projectifs, et donc qu'il est projectif. Par le théorème
\ref{theorem-projective-direct-sum}, nous pouvons écrire $M \otimes_R S =
\bigoplus_{i \in I} Q_i$ comme une somme directe de $S$-modules dénombrablement engendrés
$Q_i$. Choisissons un bon ordre sur $M$. Par récurrence transfinie, nous allons
définir une famille croissante de sous-modules $M_{\alpha}$ de $M$, un pour chaque
ordinal $\alpha$, telle que $M_{\alpha} \otimes_R S$ soit une somme directe d'une certaine
partie des $Q_i$.

\medskip\noindent
Pour $\alpha = 0$, posons $M_0 = 0$. Si $\alpha$ est un ordinal limite et si $M_{\beta}$
a été défini pour tout $\beta < \alpha$, définissons alors $M_\alpha =
\bigcup_{\beta < \alpha} M_{\beta}$. Puisque chaque $M_{\beta} \otimes_R S$, pour
$\beta < \alpha$, est une somme directe d'une partie des $Q_i$, il en va de même de
$M_{\alpha} \otimes_R S$. Si $\alpha + 1$ est un ordinal successeur et si
$M_{\alpha}$ a été défini, définissons $M_{\alpha + 1}$ comme suit. Si
$M_{\alpha} = M$, posons $M_{\alpha +1} = M$. Sinon, choisissons le plus petit
$x \in M$ (pour le bon ordre fixé) tel que $x \notin
M_{\alpha}$. Puisque $S$ est plat sur $R$, $(M/M_{\alpha}) \otimes_R S = M
\otimes_R S/M_{\alpha} \otimes_R S$ ; comme $M_{\alpha} \otimes_R S$ est
une somme directe de certains $Q_i$, il en va de même de $(M/M_{\alpha}) \otimes_R
S$. Par le lemme \ref{lemma-adapted-submodule}, nous pouvons trouver un
$R$-sous-module dénombrablement engendré $P$ de $M/M_{\alpha}$ contenant l'image de $x$ dans
$M/M_{\alpha}$ et tel que $P \otimes_R S$ (qui est égal à $\Im(P
\otimes_R S \to M \otimes_R S)$ puisque $S$ est plat sur $R$) soit une
somme directe de certains $Q_i$. Puisque $M \otimes_R S = \bigoplus_{i \in I} Q_i$
est projectif et que la projectivité passe aux facteurs directs, $P \otimes_R S$ est
également projectif. Ainsi, par le lemme \ref{lemma-ffdescent-countable-projectivity},
$P$ est projectif. Enfin, nous définissons $M_{\alpha + 1}$ comme l'image réciproque de $P$
dans $M$, de sorte que $M_{\alpha + 1}/M_{\alpha} = P$ est dénombrablement engendré et
projectif. En particulier, $M_{\alpha}$ est un facteur direct de $M_{\alpha + 1}$,
car la projectivité de $M_{\alpha + 1}/M_{\alpha}$ implique que la suite $0
\to M_{\alpha} \to M_{\alpha + 1} \to
M_{\alpha + 1}/M_{\alpha} \to 0$ est scindée.

\medskip\noindent
Une récurrence transfinie sur $M$ (en utilisant le fait que nous avons construit
$M_{\alpha + 1}$ de façon qu'il contienne le plus petit $x \in M$ non contenu dans
$M_{\alpha}$) montre que
chaque $x \in M$ est contenu dans un certain $M_{\alpha}$. Ainsi, il existe un ordinal $S$
assez grand satisfaisant : pour chaque $x \in M$, il existe $\alpha \in S$ tel
que $x \in M_{\alpha}$. Cela signifie que $(M_{\alpha})_{\alpha \in S}$ satisfait
la propriété (1) d'un dévissage de Kaplansky de $M$. Les autres propriétés sont claires
par construction. Nous concluons que
$M = \bigoplus_{\alpha + 1 \in S} M_{\alpha + 1}/M_{\alpha}$.
Puisque chaque $M_{\alpha + 1}/M_{\alpha}$ est projectif
par construction, $M$ est projectif.
\end{proof}












\section{Complétion}
\label{section-completion}

\noindent
Supposons que $R$ soit un anneau et que $I$ soit un idéal.
Nous définissons le {\it complété de $R$ relativement à $I$}
comme la limite
$$
R^\wedge = \lim_n R/I^n.
$$
Un élément de $R^\wedge$ est donné par une suite
d'éléments $f_n \in R/I^n$ tels que $f_n \equiv f_{n + 1} \bmod I^n$
pour tout $n$. Nous considérerons $R^\wedge$ comme une $R$-algèbre.
De même, si $M$ est un $R$-module, nous définissons le
{\it complété de $M$ relativement à $I$}
comme la limite
$$
M^\wedge = \lim_n M/I^nM.
$$
Un élément de $M^\wedge$ est donné par une suite
d'éléments $m_n \in M/I^nM$ tels que $m_n \equiv m_{n + 1} \bmod I^nM$
pour tout $n$. Nous considérerons $M^\wedge$ comme un $R^\wedge$-module.
Cette description montre clairement qu'il existe
toujours des morphismes canoniques
$$
M \longrightarrow M^\wedge
\quad\text{et}\quad
M \otimes_R R^\wedge \longrightarrow M^\wedge.
$$
De plus, étant donné un morphisme de modules $\varphi : M \to N$, nous obtenons sur les complétés un
morphisme induit $\varphi^\wedge : M^\wedge \to N^\wedge$ qui rend le
diagramme
$$
\xymatrix{
M \ar[r] \ar[d] & N \ar[d] \\
M^\wedge \ar[r] & N^\wedge
}
$$
commutatif. En général, la complétion n'est pas un foncteur exact, voir
Exemples, section \ref{examples-section-completion-not-exact}.
Voici quelques premiers résultats positifs.

\begin{lemma}
\label{lemma-completion-generalities}
Soit $R$ un anneau. Soit $I \subset R$ un idéal.
Soit $\varphi : M \to N$ un morphisme de $R$-modules.
\begin{enumerate}
\item Si $M/IM \to N/IN$ est surjectif, alors $M^\wedge \to N^\wedge$
est surjectif.
\item Si $M \to N$ est surjectif, alors $M^\wedge \to N^\wedge$ est surjectif.
\item Si $0 \to K \to M \to N \to 0$ est une suite exacte courte de
$R$-modules et si $N$ est plat, alors
$0 \to K^\wedge \to M^\wedge \to N^\wedge \to 0$ est une suite exacte courte.
\item Le morphisme $M \otimes_R R^\wedge \to M^\wedge$ est
surjectif pour tout $R$-module de type fini $M$.
\end{enumerate}
\end{lemma}

\begin{proof}
Supposons que $M/IM \to N/IN$ soit surjectif. Alors le morphisme $M/I^nM \to N/I^nN$
est surjectif pour tout $n \geq 1$ par le lemme de Nakayama. Plus précisément,
appliquons la partie (11) du lemme \ref{lemma-NAK} au
morphisme $M/I^nM \to N/I^nN$ sur l'anneau $R/I^n$ et à l'idéal nilpotent
$I/I^n$. Posons $K_n = \{x \in M \mid \varphi(x) \in I^nN\}$.
Nous obtenons ainsi des suites exactes courtes
$$
0 \to K_n/I^nM \to M/I^nM \to N/I^nN \to 0
$$
Nous affirmons que le morphisme canonique $K_{n + 1}/I^{n + 1}M \to K_n/I^nM$
est surjectif. En effet, si $x \in K_n$, écrivons
$\varphi(x) = \sum z_j n_j$ avec $z_j \in I^n$, $n_j \in N$.
Par hypothèse, nous pouvons écrire $n_j = \varphi(m_j) + \sum z_{jk}n_{jk}$
avec $m_j \in M$, $z_{jk} \in I$ et $n_{jk} \in N$. Par conséquent,
$$
\varphi(x - \sum z_j m_j) = \sum z_jz_{jk} n_{jk}.
$$
Cela signifie que $x' = x - \sum z_j m_j \in K_{n + 1}$ s'envoie
sur $x \bmod I^nM$, ce qui prouve l'assertion. Nous pouvons maintenant appliquer le
lemme \ref{lemma-Mittag-Leffler}
au système projectif de suites exactes courtes ci-dessus pour obtenir (1).
La partie (2) est un cas particulier de (1).
Si les hypothèses de (3) sont satisfaites, alors, pour tout $n$, la suite
$$
0 \to K/I^nK \to M/I^nM \to N/I^nN \to 0
$$
est exacte courte par le
lemme \ref{lemma-flat-tor-zero}.
Nous pouvons donc appliquer directement le
lemme \ref{lemma-Mittag-Leffler}
pour conclure que (3) est vraie.
Pour voir (4), choisissons des générateurs $x_i \in M$, $i = 1, \ldots, n$.
Alors le morphisme $R^{\oplus n} \to M$, $(a_1, \ldots, a_n) \mapsto \sum a_ix_i$
est surjectif. Par (2), nous voyons donc que
$(R^\wedge)^{\oplus n} \to M^\wedge$, $(a_1, \ldots, a_n) \mapsto \sum a_ix_i$
est surjectif. L'assertion (4) en résulte.
\end{proof}

\begin{definition}
\label{definition-complete}
Soit $R$ un anneau. Soit $I \subset R$ un idéal.
Soit $M$ un $R$-module. Nous disons que $M$ est {\it $I$-adiquement complet}
si le morphisme
$$
M \longrightarrow M^\wedge = \lim_n M/I^nM
$$
est un isomorphisme\footnote{Cela comprend la condition
$\bigcap I^nM = 0$.}. Nous disons que $R$ est {\it $I$-adiquement complet}
si $R$ est $I$-adiquement complet comme $R$-module.
\end{definition}

\noindent
Il n'est pas vrai que le complété d'un $R$-module $M$ relativement
à $I$ soit $I$-adiquement complet. Pour un exemple, voir
Exemples, section \ref{examples-section-noncomplete-completion}.
Si l'idéal est de type fini, alors le complété est complet.

\begin{lemma}
\label{lemma-hathat-finitely-generated}
\begin{reference}
\cite[Theorem 15]{Matlis}. La démonstration élégante donnée ici provient
d'un courriel de Bjorn Poonen daté du 5 novembre 2016.
\end{reference}
Soit $R$ un anneau. Soit $I$ un idéal de type fini de $R$.
Soit $M$ un $R$-module. Alors
\begin{enumerate}
\item le complété $M^\wedge$ est $I$-adiquement complet, et
\item $I^nM^\wedge = \Ker(M^\wedge \to M/I^nM) = (I^nM)^\wedge$ pour tout
$n \geq 1$.
\end{enumerate}
En particulier, $R^\wedge$ est $I$-adiquement complet,
$I^nR^\wedge = (I^n)^\wedge$, et
$R^\wedge/I^nR^\wedge = R/I^n$.
\end{lemma}

\begin{proof}
Puisque $I$ est de type fini,
$I^n$ est de type fini, disons engendré par $f_1, \ldots, f_r$.
L'application de la partie (2) du lemme \ref{lemma-completion-generalities}
à la surjection $(f_1, \ldots, f_r) : M^{\oplus r} \to I^n M$
donne une surjection
$$
(M^\wedge)^{\oplus r} \xrightarrow{(f_1, \ldots, f_r)} (I^n M)^\wedge =
\lim_{m \geq n} I^n M/I^m M = \Ker(M^\wedge \to M/I^n M).
$$
D'autre part, l'image de
$(f_1, \ldots, f_r) : (M^\wedge)^{\oplus r} \to M^\wedge$
est $I^n M^\wedge$.
Ainsi, $M^\wedge / I^n M^\wedge \simeq M/I^n M$.
En prenant les limites projectives, nous obtenons $(M^\wedge)^\wedge \simeq M^\wedge$ ;
autrement dit, $M^\wedge$ est $I$-adiquement complet.
\end{proof}

\begin{lemma}
\label{lemma-completion-differ-by-torsion}
Soit $R$ un anneau. Soit $I \subset R$ un idéal. Soit
$0 \to M \to N \to Q \to 0$ une suite exacte de
$R$-modules telle que $Q$ soit annulé par une puissance de $I$.
Alors la complétion fournit une suite exacte
$0 \to M^\wedge \to N^\wedge \to Q \to 0$.
\end{lemma}

\begin{proof}
Disons que $I^c Q = 0$. Alors $Q/I^nQ = Q$ pour $n \geq c$.
D'autre part, il est clair que
$I^nM \subset M \cap I^nN \subset I^{n - c}M$ pour $n \geq c$.
Ainsi, $M^\wedge = \lim M/(M \cap I^n N)$. Appliquons le lemme \ref{lemma-Mittag-Leffler}
au système de suites exactes
$$
0 \to M/(M \cap I^n N) \to N/I^n N \to Q \to 0
$$
pour $n \geq c$ afin de conclure.
\end{proof}

\begin{lemma}
\label{lemma-hathat}
\begin{reference}
Tiré d'une note non publiée de Lenstra et de Smit.
\end{reference}
Soit $R$ un anneau. Soit $I \subset R$ un idéal. Soit $M$ un $R$-module.
Notons $K_n = \Ker(M^\wedge \to M/I^nM)$. Alors $M^\wedge$ est $I$-adiquement
complet si et seulement si $K_n$ est égal à $I^nM^\wedge$ pour tout $n \geq 1$.
\end{lemma}

\begin{proof}
Le module $I^n M^\wedge$ est contenu dans $K_n$.
Ainsi, pour tout $n \geq 1$, il existe une suite exacte canonique
$$
0 \to K_n/I^nM^\wedge \to M^\wedge/I^nM^\wedge \to M/I^nM \to 0.
$$
Comme $I^nM^\wedge$ se surjecte sur $I^nM/I^{n + 1}M$, nous voyons que
$K_{n + 1} + I^n M^\wedge = K_n$. Ainsi, le système projectif
$\{K_n/I^n M^\wedge\}_{n \geq 1}$ a des morphismes de transition surjectifs.
Par le
lemme \ref{lemma-Mittag-Leffler},
nous voyons qu'il existe une suite exacte courte
$$
0 \to
\lim_n K_n/I^n M^\wedge \to
(M^\wedge)^\wedge \to
M^\wedge \to 0
$$
Ainsi, $M^\wedge$ est complet si et seulement si $K_n/I^n M^\wedge = 0$
pour tout $n \geq 1$.
\end{proof}

\begin{lemma}
\label{lemma-radical-completion}
Soit $R$ un anneau, soit $I \subset R$ un idéal et soit
$R^\wedge = \lim R/I^n$.
\begin{enumerate}
\item tout élément de $R^\wedge$ qui s'envoie sur une unité de $R/I$ est une unité,
\item tout élément de $1 + I$ s'envoie sur un élément inversible de $R^\wedge$,
\item tout élément de $1 + IR^\wedge$ est inversible dans $R^\wedge$, et
\item les idéaux $IR^\wedge$ et $\Ker(R^\wedge \to R/I)$ sont contenus
dans le radical de Jacobson de $R^\wedge$.
\end{enumerate}
\end{lemma}

\begin{proof}
Soit $x \in R^\wedge$ qui s'envoie sur une unité $x_1$ dans $R/I$.
Alors $x$ s'envoie sur une unité $x_n$ dans $R/I^n$ pour tout $n$, par le
lemme \ref{lemma-locally-nilpotent-unit}.
Ainsi, $y = (x_n^{-1}) \in \lim R/I^n = R^\wedge$ est un inverse de $x$.
Les parties (2) et (3) résultent immédiatement de (1).
La partie (4) résulte de (1) et du lemme \ref{lemma-contained-in-radical}.
\end{proof}

\begin{lemma}
\label{lemma-when-surjective-to-completion}
Soit $A$ un anneau. Soit $I = (f_1, \ldots, f_r)$ un idéal de type
fini. Si $M \to \lim M/f_i^nM$ est surjectif pour
chaque $i$, alors $M \to \lim M/I^nM$ est surjectif.
\end{lemma}

\begin{proof}
Remarquons que $\lim M/I^nM = \lim M/(f_1^n, \ldots, f_r^n)M$, puisque
$I^n \supset (f_1^n, \ldots, f_r^n) \supset I^{rn}$.
Un élément $\xi$ de $\lim M/(f_1^n, \ldots, f_r^n)M$ peut s'écrire formellement
sous la forme
$$
\xi = \sum\nolimits_{n \geq 0} \sum\nolimits_i f_i^n x_{n, i}
$$
avec $x_{n, i} \in M$. Si $M \to \lim M/f_i^nM$ est surjectif, alors il existe
un $x_i \in M$ qui s'envoie sur $\sum x_{n, i} f_i^n$ dans $\lim M/f_i^nM$.
Alors $x = \sum x_i$ s'envoie sur $\xi$ dans $\lim M/I^nM$.
\end{proof}

\begin{lemma}
\label{lemma-complete-by-sub}
Soit $A$ un anneau. Soient $I \subset J \subset A$ des idéaux.
Si $M$ est $J$-adiquement complet et si $I$ est de type fini, alors
$M$ est $I$-adiquement complet.
\end{lemma}

\begin{proof}
Supposons que $M$ soit $J$-adiquement complet et que $I$ soit de type fini.
Nous avons $\bigcap I^nM = 0$ parce que $\bigcap J^nM = 0$. Par le
lemme \ref{lemma-when-surjective-to-completion},
il suffit de démontrer la surjectivité de $M \to \lim M/I^nM$ dans le cas où
$I$ est engendré par un seul élément. Disons $I = (f)$.
Soient $x_n \in M$ tels que $x_{n + 1} - x_n \in f^nM$. Nous devons montrer qu'il existe
un $x \in M$ tel que $x_n - x \in f^nM$ pour tout $n$.
Comme $x_{n + 1} - x_n \in J^nM$ et comme $M$ est $J$-adiquement complet,
il existe un élément $x \in M$ tel que $x_n - x \in J^nM$.
En remplaçant $x_n$ par $x_n - x$, nous pouvons supposer que $x_n \in J^nM$.
Pour terminer la démonstration, nous allons montrer que cela implique $x_n \in I^nM$.
En effet, écrivons $x_n - x_{n + 1} = f^nz_n$.
Alors
$$
x_n = f^n(z_n + fz_{n + 1} + f^2z_{n + 2} + \ldots)
$$
La somme $z_n + fz_{n + 1} + f^2z_{n + 2} + \ldots$ converge dans $M$
puisque $f^c \in J^c$. La somme $f^n(z_n + fz_{n + 1} + f^2z_{n + 2} + \ldots)$
converge dans $M$ vers $x_n$, car
les sommes partielles sont égales à $x_n - x_{n + c}$ et $x_{n + c} \in J^{n + c}M$.
\end{proof}

\begin{lemma}
\label{lemma-change-ideal-completion}
Soit $R$ un anneau.
Soient $I$, $J$ des idéaux de $R$.
Supposons qu'il existe des entiers $c, d > 0$ tels que
$I^c \subset J$ et $J^d \subset I$.
Alors la complétion relativement à $I$ s'identifie à la complétion
relativement à $J$ pour tout $R$-module.
En particulier, un $R$-module $M$ est $I$-adiquement complet
si et seulement s'il est $J$-adiquement complet.
\end{lemma}

\begin{proof}
Considérons le système de morphismes
$M/I^nM \to M/J^{\lfloor n/c \rfloor}M$ et
le système de morphismes $M/J^mM \to M/I^{\lfloor m/d \rfloor}M$
pour obtenir des morphismes mutuellement inverses entre les complétés.
\end{proof}

\begin{lemma}
\label{lemma-quotient-complete}
Soit $R$ un anneau. Soit $I$ un idéal de $R$.
Soit $M$ un module $I$-adiquement complet sur $R$,
et soit $K \subset M$ un $R$-sous-module.
Les assertions suivantes sont équivalentes :
\begin{enumerate}
\item $K = \bigcap (K + I^nM)$ et
\item $M/K$ est $I$-adiquement complet.
\end{enumerate}
\end{lemma}

\begin{proof}
Posons $N = M/K$. Par le
lemme \ref{lemma-completion-generalities},
le morphisme $M = M^\wedge \to N^\wedge$ est surjectif.
Ainsi, $N \to N^\wedge$ est surjectif. Il est facile de voir que le
noyau de $N \to N^\wedge$ est le module $\bigcap (K + I^nM) / K$.
\end{proof}

\begin{lemma}
\label{lemma-when-finite-module-complete-over-complete-ring}
Soit $R$ un anneau. Soit $I$ un idéal de $R$.
Soit $M$ un $R$-module.
Si (a) $R$ est $I$-adiquement complet, (b) $M$ est un $R$-module de type fini,
et (c) $\bigcap I^nM = (0)$, alors $M$ est $I$-adiquement complet.
\end{lemma}

\begin{proof}
Par le lemme \ref{lemma-completion-generalities},
le morphisme $M = M \otimes_R R = M \otimes_R R^\wedge \to M^\wedge$
est surjectif. Le noyau de ce morphisme est $\bigcap I^nM$, donc nul
par hypothèse. Ainsi, $M \cong M^\wedge$ et $M$ est complet.
\end{proof}

\begin{lemma}
\label{lemma-finite-over-complete-ring}
Soit $R$ un anneau. Soit $I \subset R$ un idéal. Soit $M$ un $R$-module.
Supposons que
\begin{enumerate}
\item $R$ soit $I$-adiquement complet,
\item $\bigcap_{n \geq 1} I^nM = (0)$, et
\item $M/IM$ soit un $R/I$-module de type fini.
\end{enumerate}
Alors $M$ est un $R$-module de type fini.
\end{lemma}

\begin{proof}
Soient $x_1, \ldots, x_n \in M$ des éléments dont les images dans $M/IM$ engendrent
$M/IM$ comme $R/I$-module. Notons $M' \subset M$ le $R$-sous-module
engendré par $x_1, \ldots, x_n$. Par le lemme \ref{lemma-completion-generalities},
le morphisme $(M')^\wedge \to M^\wedge$ est surjectif.
Puisque $\bigcap I^nM = 0$, nous voyons en particulier que $\bigcap I^nM' = (0)$.
Par le lemme \ref{lemma-when-finite-module-complete-over-complete-ring},
nous voyons donc que $M'$ est complet, et nous concluons que $M' \to M^\wedge$
est surjectif. Enfin, le noyau de $M \to M^\wedge$ est
nul, puisqu'il est égal à $\bigcap I^nM = (0)$.
Nous concluons donc que $M \cong M' \cong M^\wedge$
est de type fini.
\end{proof}









\section{Complétion pour les anneaux noethériens}
\label{section-completion-noetherian}

\noindent
Dans cette section, nous étudions la complétion relativement à des idéaux dans les
anneaux noethériens.

\begin{lemma}
\label{lemma-completion-tensor}
Soit $I$ un idéal d'un anneau noethérien $R$.
Notons ${}^\wedge$ la complétion relativement à $I$.
\begin{enumerate}
\item Si $K \to N$ est un morphisme injectif de $R$-modules de type fini,
alors le morphisme induit sur les complétés $K^\wedge \to N^\wedge$ est injectif.
\item Si $0 \to K \to N \to M \to 0$ est une suite exacte courte
de $R$-modules de type fini, alors $0 \to K^\wedge \to N^\wedge \to M^\wedge \to 0$
est une suite exacte courte.
\item Si $M$ est un $R$-module de type fini, alors $M^\wedge = M \otimes_R R^\wedge$.
\end{enumerate}
\end{lemma}

\begin{proof}
En posant $M = N/K$, nous constatons que la partie (1) résulte de la partie (2).
Soit $0 \to K \to N \to M \to 0$ comme dans (2).
Pour tout $n$, nous obtenons la suite exacte courte
$$
0 \to K/(I^nN \cap K) \to N/I^nN \to M/I^nM \to 0.
$$
Par le lemme \ref{lemma-Mittag-Leffler},
nous obtenons la suite exacte
$$
0 \to \lim K/(I^nN \cap K) \to N^\wedge \to M^\wedge \to 0.
$$
Par le lemme d'Artin-Rees \ref{lemma-Artin-Rees}, nous pouvons choisir $c$ tel que
$I^nK \subset I^n N \cap K \subset I^{n-c} K$ pour $n \geq c$.
Ainsi, $K^\wedge = \lim K/I^nK = \lim K/(I^nN \cap K)$,
et nous concluons que (2) est vraie.

\medskip\noindent
Soit $M$ comme dans (3), et soit $0 \to K \to R^{\oplus t} \to M \to 0$
une présentation de $M$. Nous obtenons un diagramme commutatif
$$
\xymatrix{
&
K \otimes_R R^\wedge \ar[r] \ar[d] &
R^{\oplus t} \otimes_R R^\wedge \ar[r] \ar[d] &
M \otimes_R R^\wedge \ar[r] \ar[d] &
0 \\
0 \ar[r] &
K^\wedge \ar[r] &
(R^{\oplus t})^\wedge \ar[r] &
M^\wedge \ar[r] & 0
}
$$
La ligne supérieure est exacte, voir la section \ref{section-flat}.
La ligne inférieure est exacte par la partie (2).
Par le lemme \ref{lemma-completion-generalities},
les flèches verticales sont surjectives.
La flèche verticale du milieu est un isomorphisme.
Nous concluons que (3) est satisfaite par le lemme du serpent \ref{lemma-snake}.
\end{proof}

\begin{lemma}
\label{lemma-completion-flat}
Soit $I$ un idéal d'un anneau noethérien $R$.
Notons ${}^\wedge$ la complétion relativement à $I$.
\begin{enumerate}
\item L'homomorphisme d'anneaux $R \to R^\wedge$ est plat.
\item Le foncteur $M \mapsto M^\wedge$ est exact sur la catégorie des
$R$-modules de type fini.
\end{enumerate}
\end{lemma}

\begin{proof}
Considérons $J \otimes_R R^\wedge \to R \otimes_R R^\wedge = R^\wedge$,
où $J$ est un idéal quelconque de $R$.
D'après le lemme \ref{lemma-completion-tensor}, ce morphisme
s'identifie à $J^\wedge \to R^\wedge$, et $J^\wedge \to R^\wedge$
est injectif. La partie (1) résulte du lemme \ref{lemma-flat}.
La partie (2) est une reformulation de la
partie (2) du lemme \ref{lemma-completion-tensor}.
\end{proof}

\begin{lemma}
\label{lemma-completion-faithfully-flat}
Soit $I$ un idéal d'un anneau noethérien $R$. Notons $R^\wedge$
le complété de $R$ relativement à $I$. Si $I$ est contenu
dans le radical de Jacobson de $R$, alors l'homomorphisme d'anneaux $R \to R^\wedge$
est fidèlement plat. En particulier, si $(R, \mathfrak m)$ est un anneau local
noethérien, alors le complété $\lim_n R/\mathfrak m^n$ est fidèlement plat.
\end{lemma}

\begin{proof}
Il est plat par le lemme \ref{lemma-completion-flat}.
Le composé $R \to R^\wedge \to R/I$, où
le dernier morphisme est la projection $R^\wedge \to R/I$, montre que
tout idéal maximal de $R$ appartient à l'image de $\Spec(R^\wedge)
\to \Spec(R)$. Ainsi, le morphisme est fidèlement
plat par le lemme \ref{lemma-ff}.
\end{proof}

\begin{lemma}
\label{lemma-completion-complete}
Soit $R$ un anneau noethérien.
Soit $I$ un idéal de $R$.
Soit $M$ un $R$-module.
Alors le complété $M^\wedge$
de $M$ relativement à $I$ est $I$-adiquement complet,
$I^n M^\wedge = (I^nM)^\wedge$, et $M^\wedge/I^nM^\wedge = M/I^nM$.
\end{lemma}

\begin{proof}
C'est un cas particulier du
lemme \ref{lemma-hathat-finitely-generated},
car $I$ est un idéal de type fini.
\end{proof}

\begin{lemma}
\label{lemma-completion-Noetherian}
Soit $I$ un idéal d'un anneau $R$. Supposons que
\begin{enumerate}
\item $R/I$ soit un anneau noethérien,
\item $I$ soit de type fini.
\end{enumerate}
Alors le complété $R^\wedge$ de $R$ relativement à $I$
est un anneau noethérien complet relativement à $IR^\wedge$.
\end{lemma}

\begin{proof}
Par le
lemme \ref{lemma-hathat-finitely-generated},
nous voyons que $R^\wedge$ est $I$-adiquement complet. Il est donc aussi
$IR^\wedge$-adiquement complet. Puisque $R^\wedge/IR^\wedge = R/I$ est
noethérien, nous voyons qu'après avoir remplacé $R$ par $R^\wedge$, nous pouvons,
en plus des hypothèses (1) et (2), supposer que $R$ est lui aussi $I$-adiquement
complet.

\medskip\noindent
Soient $f_1, \ldots, f_t$ des générateurs de $I$.
Il existe alors une surjection d'anneaux
$R/I[T_1, \ldots, T_t] \to \bigoplus I^n/I^{n + 1}$
qui envoie $T_i$ sur l'élément $\overline{f}_i \in I/I^2$.
Ainsi, $\bigoplus I^n/I^{n + 1}$ est un anneau noethérien.
Soit $J \subset R$ un idéal. Considérons l'idéal
$$
\bigoplus J \cap I^n/J \cap I^{n + 1} \subset \bigoplus I^n/I^{n + 1}.
$$
Soient $\overline{g}_1, \ldots, \overline{g}_m$ des générateurs de cet
idéal. Nous pouvons choisir $\overline{g}_j$ homogène
de degré $d_j$ et choisir $g_j \in J \cap I^{d_j}$ qui s'envoie sur
$\overline{g}_j \in J \cap I^{d_j}/J \cap I^{d_j + 1}$. Nous affirmons
que $g_1, \ldots, g_m$ engendrent $J$.

\medskip\noindent
Soit $x \in J \cap I^n$. Il existe $a_j \in I^{\max(0, n - d_j)}$ tels que
$x - \sum a_j g_j \in J \cap I^{n + 1}$.
La raison en est que $J \cap I^n/J \cap I^{n + 1}$ est égal à
$\sum \overline{g}_j I^{n - d_j}/I^{n - d_j + 1}$ par notre choix
de $g_1, \ldots, g_m$. Ainsi, en partant de $x \in J$, nous pouvons trouver
une suite de vecteurs $(a_{1, n}, \ldots, a_{m, n})_{n \geq 0}$
avec $a_{j, n} \in I^{\max(0, n - d_j)}$ telle que
$$
x =
\sum\nolimits_{n = 0, \ldots, N}
\sum\nolimits_{j = 1, \ldots, m} a_{j, n} g_j \bmod I^{N + 1}
$$
En posant $A_j = \sum_{n \geq 0} a_{j, n}$, nous voyons que
$x = \sum A_j g_j$, puisque $R$ est complet. Ainsi, $J$ est de type fini, et
le résultat est démontré.
\end{proof}

\begin{lemma}
\label{lemma-completion-Noetherian-Noetherian}
Soit $R$ un anneau noethérien.
Soit $I$ un idéal de $R$.
Le complété $R^\wedge$ de $R$ relativement à $I$ est
noethérien.
\end{lemma}

\begin{proof}
C'est une conséquence du
lemme \ref{lemma-completion-Noetherian}.
On peut aussi le voir directement comme suit.
Choisissons des générateurs $f_1, \ldots, f_n$ de $I$.
Considérons le morphisme
$$
R[[x_1, \ldots, x_n]] \longrightarrow R^\wedge,
\quad
x_i \longmapsto f_i.
$$
C'est un homomorphisme d'anneaux bien défini et surjectif
(détails omis).
Puisque $R[[x_1, \ldots, x_n]]$ est noethérien (voir le
lemme \ref{lemma-Noetherian-power-series}), le résultat est démontré.
\end{proof}

\noindent
Supposons que $R \to S$ soit un homomorphisme local d'anneaux locaux $(R, \mathfrak m)$
et $(S, \mathfrak n)$. Soit $S^\wedge$ le complété
de $S$ relativement à $\mathfrak n$. En général, $S^\wedge$ n'est pas
le complété $\mathfrak m$-adique de $S$. Si
$\mathfrak n^t \subset \mathfrak mS$ pour un certain $t \geq 1$,
alors nous avons bien $S^\wedge = \lim S/\mathfrak m^nS$ par le
lemme \ref{lemma-change-ideal-completion}. Dans certains cas, cela
implique même que $S^\wedge$ est de type fini sur $R^\wedge$.

\begin{lemma}
\label{lemma-finite-after-completion}
Soit $R \to S$ un homomorphisme local d'anneaux locaux $(R, \mathfrak m)$
et $(S, \mathfrak n)$. Soient $R^\wedge$, resp.\ $S^\wedge$, le complété
de $R$, resp.\ de $S$, relativement à $\mathfrak m$, resp.\ à $\mathfrak n$.
Si $\mathfrak m$ et $\mathfrak n$ sont de type fini et si
$\dim_{\kappa(\mathfrak m)} S/\mathfrak mS < \infty$, alors
\begin{enumerate}
\item $S^\wedge$ est égal au complété $\mathfrak m$-adique de $S$, et
\item $S^\wedge$ est un $R^\wedge$-module de type fini.
\end{enumerate}
\end{lemma}

\begin{proof}
Nous avons $\mathfrak mS \subset \mathfrak n$, puisque $R \to S$ est un homomorphisme
local.
L'hypothèse $\dim_{\kappa(\mathfrak m)} S/\mathfrak mS < \infty$
implique que $S/\mathfrak mS$ est un anneau artinien, voir le
lemme \ref{lemma-finite-dimensional-algebra}.
Il est donc de dimension $0$, voir le
lemme \ref{lemma-Noetherian-dimension-0},
et par conséquent $\mathfrak n = \sqrt{\mathfrak mS}$.
Cela et le fait que $\mathfrak n$ est de type fini
impliquent que $\mathfrak n^t \subset \mathfrak mS$ pour
un certain $t \geq 1$. Par le
lemme \ref{lemma-change-ideal-completion},
nous voyons que $S^\wedge$ peut être identifié au complété
$\mathfrak m$-adique de $S$. Comme $\mathfrak m$ est de type fini, il résulte du
lemme \ref{lemma-hathat-finitely-generated}
que $S^\wedge$ et $R^\wedge$ sont $\mathfrak m$-adiquement complets.
Nous pouvons alors appliquer le
lemme \ref{lemma-finite-over-complete-ring}
à $S^\wedge$ comme $R^\wedge$-module pour conclure.
\end{proof}

\begin{lemma}
\label{lemma-completion-finite-extension}
Soit $R$ un anneau noethérien. Soit $R \to S$ un homomorphisme d'anneaux fini.
Soit $\mathfrak p \subset R$ un idéal premier et soient
$\mathfrak q_1, \ldots, \mathfrak q_m$ les idéaux premiers de $S$
au-dessus de $\mathfrak p$
(lemme \ref{lemma-finite-finite-fibres}).
Alors
$$
R_\mathfrak p^\wedge \otimes_R S =
(S_\mathfrak p)^\wedge =
S_{\mathfrak q_1}^\wedge \times \ldots \times S_{\mathfrak q_m}^\wedge
$$
où $(S_\mathfrak p)^\wedge$ est le complété relativement à
$\mathfrak p$, et où les anneaux locaux $R_\mathfrak p$ et
$S_{\mathfrak q_i}$ sont complétés relativement à leurs idéaux maximaux.
\end{lemma}

\begin{proof}
Nous pouvons remplacer $R$ par le localisé $R_\mathfrak p$ et
$S$ par $S_\mathfrak p = S \otimes_R R_\mathfrak p$.
Nous pouvons donc supposer que $R$ est un anneau local noethérien et
que $\mathfrak p = \mathfrak m$ est son idéal maximal.
Le complété $\mathfrak q_iS_{\mathfrak q_i}$-adique
$S_{\mathfrak q_i}^\wedge$ est égal au complété $\mathfrak m$-adique
par le lemme \ref{lemma-finite-after-completion}.
Pour tout $n \geq 1$, les idéaux premiers de $S/\mathfrak m^nS$ sont en
correspondance bijective avec les idéaux maximaux
$\mathfrak q_1, \ldots, \mathfrak q_m$ de $S$
(par montée pour $S$ sur $R$, voir le
lemme \ref{lemma-integral-going-up}).
Ainsi,
$S/\mathfrak m^nS = \prod S_{\mathfrak q_i}/\mathfrak m^nS_{\mathfrak q_i}$
par le lemme \ref{lemma-artinian-finite-length}
(en utilisant par exemple la
proposition \ref{proposition-dimension-zero-ring}
pour voir que $S/\mathfrak m^nS$ est artinien).
Ainsi, le complété $\mathfrak m$-adique $S^\wedge$ de $S$ est égal à
$\prod S_{\mathfrak q_i}^\wedge$. Enfin, nous avons
$R^\wedge \otimes_R S = S^\wedge$ par le lemme \ref{lemma-completion-tensor}.
\end{proof}

\begin{lemma}
\label{lemma-split-completed-sequence}
Soit $R$ un anneau. Soit $I \subset R$ un idéal.
Soit $0 \to K \to P \to M \to 0$ une suite exacte courte de
$R$-modules. Si $M$ est plat sur $R$ et si $M/IM$ est un
$R/I$-module projectif, alors la suite des complétés $I$-adiques
$$
0 \to K^\wedge \to P^\wedge \to M^\wedge \to 0
$$
est une suite exacte scindée.
\end{lemma}

\begin{proof}
Puisque $M$ est plat, chacune des suites
$$
0 \to K/I^nK \to P/I^nP \to M/I^nM \to 0
$$
est exacte courte, voir le
lemme \ref{lemma-flat-tor-zero},
et la suite $0 \to K^\wedge \to P^\wedge \to M^\wedge \to 0$
est une suite exacte courte, voir le
lemme \ref{lemma-completion-generalities}.
Il suffit de montrer que nous pouvons trouver des scindages
$s_n : M/I^nM \to P/I^nP$ tels que $s_{n + 1} \bmod I^n = s_n$.
Nous allons construire ces $s_n$ par récurrence sur $n$.
Choisissons un scindage quelconque $s_1$, qui existe puisque $M/IM$ est un $R/I$-module projectif.
Supposons $s_n$ donné pour un certain $n > 0$. Posons
$P_{n + 1} = \{x \in P \mid x \bmod I^nP \in \Im(s_n)\}$.
Le morphisme $\pi : P_{n + 1}/I^{n + 1}P_{n + 1} \to M/I^{n + 1}M$ est surjectif
(détails omis). Puisque $M/I^{n + 1}M$ est projectif comme $R/I^{n + 1}$-module
par le
lemme \ref{lemma-lift-projective},
nous pouvons choisir une section
$t : M/I^{n + 1}M \to P_{n + 1}/I^{n + 1}P_{n + 1}$
de $\pi$. En prenant pour $s_{n + 1}$ le composé de
$t$ et du morphisme canonique $P_{n + 1}/I^{n + 1}P_{n + 1} \to P/I^{n + 1}P$,
nous obtenons le résultat voulu.
\end{proof}

\begin{lemma}
\label{lemma-complete-modulo-nilpotent}
Soit $A$ un anneau noethérien. Soient $I, J \subset A$ des idéaux.
Si $A$ est $I$-adiquement complet et si $A/I$ est $J$-adiquement complet,
alors $A$ est $J$-adiquement complet.
\end{lemma}

\begin{proof}
Soit $B$ le complété $(I + J)$-adique de $A$. Par le
lemme \ref{lemma-completion-flat}, $B/IB$ est le complété $J$-adique de $A/I$,
donc est isomorphe à $A/I$ par hypothèse. De plus, $B$ est $I$-adiquement
complet par le lemme \ref{lemma-complete-by-sub}. Ainsi, $B$ est un
$A$-module de type fini par le lemme \ref{lemma-finite-over-complete-ring}.
Par le lemme de Nakayama (lemme \ref{lemma-NAK}, en utilisant que $I$ est contenu dans le
radical de Jacobson de $A$
par le lemme \ref{lemma-radical-completion}), nous trouvons que $A \to B$ est surjectif.
Le morphisme $A \to B$ est plat par le lemme \ref{lemma-completion-flat}.
L'image de $\Spec(B) \to \Spec(A)$ contient $V(I)$ et, comme $I$
est contenu dans le radical de Jacobson de $A$, nous trouvons que $A \to B$ est fidèlement plat
(lemme \ref{lemma-ff-rings}). Ainsi, $A \to B$ est injectif. Par conséquent, $A$ est
complet relativement à $I + J$, et donc a fortiori complet
relativement à $J$.
\end{proof}









\section{Passage à la limite de modules}
\label{section-limits}

\noindent
Dans cette section, nous étudions ce qui se produit lorsqu'on prend une limite de modules.

\begin{lemma}
\label{lemma-limit-complete-pre}
Soit $I \subset A$ un idéal de type fini d'un anneau.
Soit $(M_n)$ un système projectif de $A$-modules tel que $I^n M_n = 0$.
Alors $M = \lim M_n$ est $I$-adiquement complet.
\end{lemma}

\begin{proof}
Nous avons $M \to M/I^nM \to M_n$. En prenant la limite, nous obtenons
$M \to M^\wedge \to M$. Ainsi, $M$ est un facteur direct de $M^\wedge$.
Puisque $M^\wedge$ est $I$-adiquement complet par le
lemme \ref{lemma-hathat-finitely-generated}, il en va de même de $M$.
\end{proof}

\begin{lemma}
\label{lemma-limit-complete}
Soit $I \subset A$ un idéal de type fini d'un anneau.
Soit $(M_n)$ un système projectif de $A$-modules tel que
$M_n = M_{n + 1}/I^nM_{n + 1}$. Posons $M = \lim M_n$.
Alors $M/I^nM = M_n$ et $M$ est $I$-adiquement complet.
\end{lemma}

\begin{proof}
Par le lemme \ref{lemma-limit-complete-pre}, nous voyons que $M$ est $I$-adiquement
complet. Puisque les morphismes de transition sont surjectifs, les morphismes $M \to M_n$
sont surjectifs. Considérons le système projectif de suites exactes courtes
$$
0 \to N_n \to M \to M_n \to 0
$$
qui définit $N_n$. Puisque $M_n = M_{n + 1}/I^nM_{n + 1}$, le morphisme
$N_{n + 1} + I^nM \to N_n$ est surjectif. Ainsi,
$N_{n + 1}/(N_{n + 1} \cap I^{n + 1}M) \to N_n/(N_n \cap I^nM)$
est surjectif. En prenant la limite projective des suites exactes courtes
$$
0 \to N_n/(N_n \cap I^nM) \to M/I^nM \to M_n \to 0
$$
nous obtenons une suite exacte
$$
0 \to \lim N_n/(N_n \cap I^nM) \to M^\wedge \to M
$$
Puisque $M$ est $I$-adiquement complet, nous concluons que
$\lim N_n/(N_n \cap I^nM) = 0$ et donc, par la surjectivité
des morphismes de transition, que $N_n/(N_n \cap I^nM) = 0$ pour tout $n$.
Ainsi, $M_n = M/I^nM$, comme voulu.
\end{proof}

\begin{lemma}
\label{lemma-finiteness-graded}
Soit $A$ un anneau gradué noethérien. Soit $I \subset A_+$ un idéal
homogène. Soit $(N_n)$ un système projectif de $A$-modules gradués de type fini tel que
$N_n = N_{n + 1}/I^n N_{n + 1}$. Alors il existe un $A$-module gradué de type fini
$N$ tel que $N_n = N/I^nN$ comme modules gradués pour tout $n$.
\end{lemma}

\begin{proof}
Choisissons $r$ et des éléments homogènes $x_{1, 1}, \ldots, x_{1, r} \in N_1$ de
degrés $d_1, \ldots, d_r$ qui engendrent $N_1$. Puisque les morphismes de transition
sont surjectifs, nous pouvons choisir un système compatible d'éléments homogènes
$x_{n, i} \in N_n$ relevant $x_{1, i}$. Par le lemme de Nakayama gradué
(lemme \ref{lemma-graded-NAK}), nous voyons que
$N_n$ est engendré par les éléments $x_{n, 1}, \ldots, x_{n, r}$
placés dans les degrés $d_1, \ldots, d_r$.
Ainsi, pour $m \leq n$, nous voyons que $N_n \to N_n/I^m N_n$
est un isomorphisme en degrés $< \min(d_i) + m$ (puisque $I^mN_n$ est nul
en ces degrés). Le système projectif des composantes de degré $d$
$$
\ldots
= N_{2 + d - \min(d_i), d}
= N_{1 + d - \min(d_i), d}
= N_{d - \min(d_i), d} \to N_{-1 + d - \min(d_i), d} \to \ldots
$$
se stabilise donc comme indiqué. Soit $N$ le $A$-module gradué dont la
composante graduée de degré $d$ est cette valeur stationnaire. En particulier, nous avons les
éléments $x_i = \lim x_{n, i}$ dans $N$. Nous affirmons que les $x_i$ engendrent $N$ :
tout $x \in N_d$ est une combinaison linéaire de $x_1, \ldots, x_r$,
car nous pouvons le vérifier dans $N_{d - \min(d_i), d}$, où cela est vrai
puisque les $x_{d - \min(d_i), i}$ engendrent $N_{d - \min(d_i)}$.
Enfin, le lecteur vérifie que le morphisme surjectif
$N/I^nN \to N_n$ est un isomorphisme en examinant
ce qui se produit dans chaque degré comme ci-dessus. Les détails sont omis.
\end{proof}

\begin{lemma}
\label{lemma-daniel-litt}
Soit $A$ un anneau gradué. Soit $I \subset A_+$ un idéal homogène.
Notons $A' = \lim A/I^n$. Soit $(G_n)$ un système projectif
de $A$-modules gradués tel que $G_n$ soit annulé par $I^n$.
Soit $M$ un $A$-module gradué et soit
$\varphi_n : M \to G_n$ un système compatible de morphismes gradués de
$A$-modules. Si le morphisme induit
$$
\varphi : M \otimes_A A' \longrightarrow \lim G_n
$$
est un isomorphisme, alors $M_d \to \lim G_{n, d}$
est un isomorphisme pour tout $d \in \mathbf{Z}$.
\end{lemma}

\begin{proof}
Par convention, les anneaux gradués sont en degrés $\geq 0$, tandis que les modules gradués
peuvent avoir des composantes non nulles de tout degré, voir la section \ref{section-graded}.
Le morphisme $\varphi$ existe parce que $\lim G_n$ est
un module sur $A'$, puisque $G_n$ est annulé par $I^n$.
Une autre observation utile à garder à l'esprit est que nous avons
$$
\bigoplus\nolimits_{d \in \mathbf{Z}} \lim G_{n, d} \subset
\lim G_n \subset
\prod\nolimits_{d \in \mathbf{Z}} \lim G_{n, d}
$$
où un indice ${\ }_d$ désigne la composante graduée de degré $d$.

\medskip\noindent
Injectivité. Soit $x \in M_d$. Si $x \mapsto 0$ dans $\lim G_{n, d}$,
alors $x \otimes 1 = 0$ dans $M \otimes_A A'$. Nous pouvons alors trouver
un sous-module de type fini $M' \subset M$ tel que $x \in M'$
et que $x \otimes 1$ soit nul dans $M' \otimes_A A'$.
Disons que $M'$ est engendré par des éléments homogènes placés
en degrés $d_1, \ldots, d_r$. Soit $n = d - \min(d_i) + 1$.
Puisque $A'$ possède un morphisme vers $A/I^n$ et puisque
$A \to A/I^n$ est un isomorphisme en degrés $\leq n - 1$,
nous voyons que $M' \to M' \otimes_A A'$ est injectif en
degrés $\leq n - 1$. Ainsi, $x = 0$, comme voulu.

\medskip\noindent
Surjectivité. Soit $y \in \lim G_{n, d}$. Choisissons une somme
finie $\sum x_i \otimes f'_i$ dans $M \otimes_A A'$ qui s'envoie sur $y$.
Nous pouvons supposer $x_i$ homogène, disons de degré $d_i$.
Observons que, bien que $A'$ ne soit pas un anneau gradué, c'est une
limite des anneaux gradués $A/I^nA$ et que, de plus, dans tout
degré donné, les morphismes de transition deviennent finalement des isomorphismes
(voir ci-dessus). Cela donne
$$
A = \bigoplus\nolimits_{d \geq 0} A_d \subset A' \subset
\prod\nolimits_{d \geq 0} A_d
$$
Nous pouvons donc écrire
$$
f'_i = \sum\nolimits_{j = 0, \ldots, d - d_i - 1} f_{i, j} + f_i + g'_i
$$
avec $f_{i, j} \in A_j$, $f_i \in A_{d - d_i}$ et
$g'_i \in A'$ qui s'envoie sur zéro dans $\prod_{j \leq d - d_i} A_j$.
Si maintenant nous calculons $\varphi_n(\sum_{i, j} f_{i, j}x_i) \in G_n$,
nous obtenons une somme d'éléments homogènes de degré $< d$.
Ainsi, $\varphi(\sum x_i \otimes f_{i, j})$ s'envoie sur zéro dans
$\lim G_{n, d}$.
De même, un calcul montre que l'élément $\varphi(\sum x_i \otimes g'_i)$
s'envoie sur zéro dans $\prod_{d' \leq d} \lim G_{n, d'}$.
Puisque nous savons que $\varphi(\sum x_i \otimes f'_i)$
est $y$, nous concluons que $\sum f_ix_i \in M_d$ s'envoie sur $y$, comme voulu.
\end{proof}













\section{Critères de platitude}
\label{section-criteria-flatness}

\noindent
Dans cette section, nous démontrons quelques lemmes techniques importants dans le cas
noethérien. Nous les généraliserons (partiellement) au cas non noethérien
dans la section \ref{section-more-flatness-criteria}.

\begin{lemma}
\label{lemma-mod-injective}
Supposons que $R \to S$ soit un homomorphisme local d'anneaux locaux, avec $S$
noethérien.
Notons $\mathfrak m$ l'idéal maximal de $R$. Soient $M$ un $R$-module plat
et $N$ un $S$-module de type fini. Soit $u : N \to M$ un morphisme de $R$-modules.
Si $\overline{u} : N/\mathfrak m N \to M/\mathfrak m M$
est injectif, alors $u$ est injectif.
Dans ce cas, $M/u(N)$ est plat sur $R$.
\end{lemma}

\begin{proof}
Affirmons d'abord que $u_n : N/{\mathfrak m}^nN \to M/{\mathfrak m}^nM$
est injectif pour tout $n \geq 1$. Procédons par récurrence, le cas initial
étant que $\overline{u} = u_1$ est injectif. Comme $M$
est plat sur $R$, nous avons une suite exacte courte
$0 \to M \otimes_R {\mathfrak m}^n/{\mathfrak m}^{n + 1}
\to M/{\mathfrak m}^{n + 1}M \to M/{\mathfrak m}^n M \to 0$.
De plus, $M \otimes_R {\mathfrak m}^n/{\mathfrak m}^{n + 1}
= M/{\mathfrak m}M \otimes_{R/{\mathfrak m}}
{\mathfrak m}^n/{\mathfrak m}^{n + 1}$. Nous avons
une suite exacte analogue $N \otimes_R {\mathfrak m}^n/{\mathfrak m}^{n + 1}
\to N/{\mathfrak m}^{n + 1}N \to N/{\mathfrak m}^n N \to 0$
pour $N$, sauf que nous n'avons pas le zéro à gauche. Nous avons aussi
$N \otimes_R {\mathfrak m}^n/{\mathfrak m}^{n + 1}
= N/{\mathfrak m}N \otimes_{R/{\mathfrak m}}
{\mathfrak m}^n/{\mathfrak m}^{n + 1}$. Ainsi, le morphisme $u_{n + 1}$ est
injectif, puisque $u_n$ et le morphisme
$\overline{u} \otimes \text{id}_{{\mathfrak m}^n/{\mathfrak m}^{n + 1}}$ le sont.

\medskip\noindent
Par le théorème d'intersection de Krull
(lemme \ref{lemma-intersect-powers-ideal-module-zero})
appliqué à $N$ sur l'anneau $S$ et à l'idéal $\mathfrak mS$,
nous avons $\bigcap \mathfrak m^nN = 0$. L'injectivité
de $u_n$ pour tout $n$ implique donc que $u$ est injectif.

\medskip\noindent
Pour montrer que $M/u(N)$ est plat sur $R$, il suffit de montrer que
$\text{Tor}_1^R(M/u(N), R/I) = 0$ pour tout idéal $I \subset R$,
voir le lemme \ref{lemma-characterize-flat}. De la suite exacte courte
$$
0 \to N \xrightarrow{u} M \to M/u(N) \to 0
$$
et de la platitude de $M$, nous obtenons une suite exacte de groupes Tor
$$
0 \to \text{Tor}_1^R(M/u(N), R/I) \to N/IN \to M/IM
$$
Voir le lemme \ref{lemma-long-exact-sequence-tor}. Il suffit donc de montrer que
$N/IN$ s'injecte dans $M/IM$. Notons que $R/I \to S/IS$ est un homomorphisme
local d'anneaux locaux avec $S/IS$ noethérien, que $N/IN \to M/IM$ est un morphisme
de $R/I$-modules, que $N/IN$ est de type fini sur $S/IS$, que $M/IM$ est plat sur
$R/I$ et que $u \bmod I : N/IN \to M/IM$ est injectif modulo $\mathfrak m$.
Nous pouvons donc appliquer la première partie de la démonstration à $u \bmod I$, ce qui conclut.
\end{proof}

\begin{lemma}
\label{lemma-grothendieck}
Supposons que $R \to S$ soit un homomorphisme d'anneaux plat et local entre anneaux
locaux noethériens. Notons $\mathfrak m$ l'idéal maximal de $R$.
Supposons que $f \in S$ soit un non-diviseur de zéro dans $S/{\mathfrak m}S$.
Alors $S/fS$ est plat sur $R$, et $f$ est un non-diviseur de zéro dans $S$.
\end{lemma}

\begin{proof}
Cela découle directement du lemme \ref{lemma-mod-injective}.
\end{proof}

\begin{lemma}
\label{lemma-grothendieck-regular-sequence}
Supposons que $R \to S$ soit un homomorphisme d'anneaux plat et local entre anneaux
locaux noethériens. Notons $\mathfrak m$ l'idéal maximal de $R$.
Supposons que $f_1, \ldots, f_c$ soit une suite d'éléments de
$S$ telle que les images $\overline{f}_1, \ldots, \overline{f}_c$
forment une suite régulière dans $S/{\mathfrak m}S$.
Alors $f_1, \ldots, f_c$ est une suite régulière dans $S$ et chacun
des quotients $S/(f_1, \ldots, f_i)$ est plat sur $R$.
\end{lemma}

\begin{proof}
Par récurrence et par le lemme \ref{lemma-grothendieck}.
\end{proof}

\begin{lemma}
\label{lemma-free-fibre-flat-free}
Soit $R \to S$ un homomorphisme local d'anneaux
locaux noethériens. Soit $\mathfrak m$ l'idéal
maximal de $R$. Soit $M$ un $S$-module non nul de type fini.
Supposons que (a) $M/\mathfrak mM$
soit un $S/\mathfrak mS$-module libre, et que (b) $M$ soit plat sur $R$.
Alors $M$ est libre et $S$ est plat sur $R$.
\end{lemma}

\begin{proof}
Soit $\overline{x}_1, \ldots, \overline{x}_n$ une base
du module libre $M/\mathfrak mM$. Choisissons
$x_1, \ldots, x_n \in M$, avec $x_i$ s'envoyant sur $\overline{x}_i$. Soit
$u : S^{\oplus n} \to M$ le morphisme qui envoie le $i$-ème vecteur
de la base canonique sur $x_i$. Le lemme \ref{lemma-mod-injective}
montre que $u$ est injectif. D'autre part, par le
lemme de Nakayama \ref{lemma-NAK}, ce morphisme est surjectif. Le
résultat en découle.
\end{proof}

\begin{lemma}
\label{lemma-complex-exact-mod}
Soit $R \to S$ un homomorphisme local d'anneaux locaux
noethériens. Soit $\mathfrak m$ l'idéal maximal de $R$.
Soit $0 \to F_e \to F_{e-1} \to \ldots \to F_0$
un complexe fini de $S$-modules de type fini. Supposons que
chaque $F_i$ soit plat sur $R$, et que le complexe
$0 \to F_e/\mathfrak m F_e \to F_{e-1}/\mathfrak m F_{e-1}
\to \ldots \to F_0 / \mathfrak m F_0$ soit exact.
Alors $0 \to F_e \to F_{e-1} \to \ldots \to F_0$
est exact et, de plus, le module
$\Coker(F_1 \to F_0)$ est plat sur $R$.
\end{lemma}

\begin{proof}
Raisonnons par récurrence sur $e$. Si $e = 1$, c'est exactement le
lemme \ref{lemma-mod-injective}. Si $e > 1$, le
lemme \ref{lemma-mod-injective} montre que $F_e \to F_{e-1}$
est injectif et que $C = \Coker(F_e \to F_{e-1})$
est un $S$-module de type fini plat sur $R$. Nous pouvons donc
appliquer l'hypothèse de récurrence au complexe
$0 \to C \to F_{e-2} \to \ldots \to F_0$.
Nous en déduisons que $C \to F_{e-2}$ est injectif,
d'où l'exactitude du complexe ainsi que
la platitude du conoyau de $F_1 \to F_0$.
\end{proof}

\noindent
Dans la suite de cette section, nous démontrons deux versions de ce que l'on appelle le
« {\it critère local de platitude} ». Notons aussi l'intéressant
lemme \ref{lemma-CM-over-regular-flat} ci-dessous.

\begin{lemma}
\label{lemma-prepare-local-criterion-flatness}
Soit $R$ un anneau local d'idéal maximal $\mathfrak m$
et de corps résiduel $\kappa = R/\mathfrak m$.
Soit $M$ un $R$-module. Si $\text{Tor}_1^R(\kappa, M) = 0$,
alors, pour tout $R$-module $N$ de longueur finie, nous avons
$\text{Tor}_1^R(N, M) = 0$.
\end{lemma}

\begin{proof}
Raisonnons par récurrence sur la longueur de $N$.
Si la longueur de $N$ vaut $1$, alors $N \cong \kappa$
et le résultat est acquis. Si la longueur de $N$ est strictement supérieure à
$1$, nous pouvons insérer $N$ dans une suite exacte courte
$0 \to N' \to N \to N'' \to 0$, où $N'$ et $N''$ sont des
$R$-modules de longueur finie et de longueurs plus petites.
L'annulation de $\text{Tor}_1^R(N, M)$ découle
de celles de $\text{Tor}_1^R(N', M)$
et de $\text{Tor}_1^R(N'', M)$ (hypothèse de récurrence),
ainsi que de la suite exacte longue des groupes Tor, voir le lemme
\ref{lemma-long-exact-sequence-tor}.
\end{proof}

\begin{lemma}[Critère local de platitude]
\label{lemma-local-criterion-flatness}
Soit $R \to S$ un homomorphisme local d'anneaux locaux
noethériens. Soit $\mathfrak m$ l'idéal maximal de $R$,
et soit $\kappa = R/\mathfrak m$.
Soit $M$ un $S$-module de type fini. Si $\text{Tor}_1^R(\kappa, M) = 0$,
alors $M$ est plat sur $R$.
\end{lemma}

\begin{proof}
Soit $I \subset R$ un idéal. D'après le lemme \ref{lemma-flat}, il suffit
de montrer que $I \otimes_R M \to M$ est injectif. Par la remarque
\ref{remark-Tor-ring-mod-ideal}, nous voyons que son noyau est
égal à $\text{Tor}_1^R(M, R/I)$. Le
lemme \ref{lemma-prepare-local-criterion-flatness}
montre que $J \otimes_R M \to M$ est injectif pour tout idéal
de colongueur finie.

\medskip\noindent
Choisissons $n >> 0$ et considérons la suite exacte
courte suivante
$$
0
\to I \cap \mathfrak m^n
\to I \oplus \mathfrak m^n
\to I + \mathfrak m^n
\to 0
$$
C'est une sous-suite de la suite exacte courte
$0 \to R \to R^{\oplus 2} \to R \to 0$. Nous obtenons donc le diagramme
$$
\xymatrix{
(I\cap \mathfrak m^n) \otimes_R M \ar[r] \ar[d] &
I \otimes_R M \oplus \mathfrak m^n \otimes_R M \ar[r] \ar[d] &
(I + \mathfrak m^n) \otimes_R M \ar[d] \\
M \ar[r] &
M \oplus M \ar[r] &
M
}
$$
Notons que $I + \mathfrak m^n$ et $\mathfrak m^n$
sont des idéaux de colongueur finie.
Une chasse au diagramme montre donc que
$\Ker((I \cap \mathfrak m^n)\otimes_R M \to M)
\to \Ker(I \otimes_R M \to M)$
est surjectif. Nous en concluons en particulier que
$K = \Ker(I \otimes_R M \to M)$ est contenu
dans l'image de $(I \cap \mathfrak m^n) \otimes_R M$
dans $I \otimes_R M$. Par le lemme d'Artin-Rees \ref{lemma-Artin-Rees},
nous voyons que $K$ est contenu
dans $\mathfrak m^{n-c}(I \otimes_R M)$ pour un certain $c > 0$
et tout $n >> 0$. Comme $I \otimes_R M$ est un
$S$-module de type fini (!) et que $S$ est noethérien, nous voyons
que cela implique $K = 0$. En effet, ce qui précède implique que
$K$ s'envoie sur zéro dans la complétion $\mathfrak mS$-adique
de $I \otimes_R M$. Mais le morphisme de $S$
vers sa complétion $\mathfrak mS$-adique est fidèlement plat
par le lemme \ref{lemma-completion-faithfully-flat}.
Ainsi $K = 0$, comme souhaité.
\end{proof}

\noindent
Dans ce qui suit, nous rencontrons souvent les conditions
« $M/IM$ est plat sur $R/I$ et $\text{Tor}_1^R(R/I, M) = 0$ ».
Le lemme suivant donne quelques conséquences de ces conditions
(il s'agit d'une généralisation du
lemme \ref{lemma-prepare-local-criterion-flatness}).

\begin{lemma}
\label{lemma-what-does-it-mean}
Soit $R$ un anneau.
Soit $I \subset R$ un idéal.
Soit $M$ un $R$-module.
Si $M/IM$ est plat sur $R/I$ et $\text{Tor}_1^R(R/I, M) = 0$, alors
\begin{enumerate}
\item $M/I^nM$ est plat sur $R/I^n$ pour tout $n \geq 1$, et
\item pour tout module $N$ annulé par $I^m$ pour un certain $m \geq 0$,
nous avons $\text{Tor}_1^R(N, M) = 0$.
\end{enumerate}
En particulier, si $I$ est nilpotent, alors $M$ est plat sur $R$.
\end{lemma}

\begin{proof}
Supposons que $M/IM$ soit plat sur $R/I$ et que $\text{Tor}_1^R(R/I, M) = 0$.
Soit $N$ un $R/I$-module. Choisissons une suite exacte courte
$$
0 \to K \to \bigoplus\nolimits_{i \in I} R/I \to N \to 0
$$
Par la suite exacte longue de $\text{Tor}$ et l'annulation de
$\text{Tor}_1^R(R/I, M)$, nous obtenons
$$
0 \to \text{Tor}_1^R(N, M) \to K \otimes_R M \to
(\bigoplus\nolimits_{i \in I} R/I) \otimes_R M \to N \otimes_R M \to 0
$$
Mais comme $K$, $\bigoplus_{i \in I} R/I$ et $N$ sont tous annulés
par $I$, nous voyons que
\begin{align*}
K \otimes_R M & = K \otimes_{R/I} M/IM, \\
(\bigoplus\nolimits_{i \in I} R/I) \otimes_R M & =
(\bigoplus\nolimits_{i \in I} R/I) \otimes_{R/I} M/IM, \\
N \otimes_R M & = N \otimes_{R/I} M/IM.
\end{align*}
Puisque $M/IM$ est plat sur $R/I$, nous concluons que
$$
0 \to K \otimes_{R/I} M/IM \to
(\bigoplus\nolimits_{i \in I} R/I) \otimes_{R/I} M/IM \to
N \otimes_{R/} M/IM \to 0
$$
est exacte. En combinant ceci avec ce qui précède, nous concluons que
$\text{Tor}_1^R(N, M) = 0$ pour tout $R$-module $N$ annulé par $I$.

\medskip\noindent
Démontrons (2) par récurrence sur $m$. Le cas $m = 1$ a été
traité au paragraphe précédent. Pour $N$ annulé par $I^m$
avec $m > 1$, nous pouvons choisir une suite exacte
$0 \to N' \to N \to N'' \to 0$, où $N'$ et $N''$ sont annulés
par $I^{m - 1}$. Par exemple, on peut prendre $N' = IN$ et $N'' = N/IN$.
Alors la suite exacte
$$
\text{Tor}_1^R(N', M) \to
\text{Tor}_1^R(N, M) \to
\text{Tor}_1^R(N'', M)
$$
et la récurrence prouvent l'annulation voulue.

\medskip\noindent
Enfin, démontrons (1). Étant donné $n \geq 1$, nous devons montrer que $M/I^nM$
est plat sur $R/I^n$. Autrement dit, nous devons montrer que le foncteur
$N \mapsto N \otimes_{R/I^n} M/I^nM$ est exact sur la catégorie des $R$-modules
$N$ annulés par $I^n$. Or, pour un tel $N$, nous avons
$N \otimes_{R/I^n} M/I^nM = N \otimes_R M$. L'annulation de
$\text{Tor}_1$ dans (2) montre que le foncteur $N \mapsto N \otimes_R M$
est exact sur la catégorie des $N$ annulés par une puissance de $I$,
ce qui conclut.
\end{proof}

\begin{lemma}
\label{lemma-what-does-it-mean-again}
Soient $R$ un anneau, $I \subset R$ un idéal et
$M$ un $R$-module.
\begin{enumerate}
\item Si $M/IM$ est plat sur $R/I$ et si $M \otimes_R I/I^2 \to IM/I^2M$
est injectif, alors $M/I^2M$ est plat sur $R/I^2$.
\item Si $M/IM$ est plat sur $R/I$ et si $M \otimes_R I^n/I^{n + 1}
\to I^nM/I^{n + 1}M$ est injectif pour $n = 1, \ldots, k$,
alors $M/I^{k + 1}M$ est plat sur $R/I^{k + 1}$.
\end{enumerate}
\end{lemma}

\begin{proof}
La première assertion découle du
lemme \ref{lemma-what-does-it-mean} appliqué en remplaçant $R$ par $R/I^2$
et $M$ par $M/I^2M$, en utilisant l'égalité
$$
\text{Tor}_1^{R/I^2}(M/I^2M, R/I) =
\Ker(M \otimes_R I/I^2 \to IM/I^2M),
$$
voir la remarque \ref{remark-Tor-ring-mod-ideal}.
La seconde assertion s'obtient de la même manière en raisonnant par récurrence
sur $n$ pour montrer que $M/I^{n + 1}M$ est plat sur $R/I^{n + 1}$ pour
$n = 1, \ldots, k$. Nous utilisons ici l'égalité
$$
\text{Tor}_1^{R/I^{n + 1}}(M/I^{n + 1}M, R/I^n) =
\Ker(M \otimes_R I^n/I^{n + 1} \to I^nM/I^{n + 1}M)
$$
pour tout $n$.
\end{proof}

\begin{lemma}[Variante du critère local]
\label{lemma-variant-local-criterion-flatness}
Soit $R \to S$ un homomorphisme local d'anneaux locaux
noethériens. Soit $I \not = R$ un idéal de $R$.
Soit $M$ un $S$-module de type fini. Si $\text{Tor}_1^R(M, R/I) = 0$
et si $M/IM$ est plat sur $R/I$, alors $M$ est plat sur $R$.
\end{lemma}

\begin{proof}
Première démonstration : par le
lemme \ref{lemma-what-does-it-mean},
nous voyons que $\text{Tor}_1^R(\kappa, M)$ est nul, où $\kappa$
est le corps résiduel de $R$. Le module $M$ est donc
plat sur $R$ par le
lemme \ref{lemma-local-criterion-flatness}.

\medskip\noindent
Seconde démonstration : soit $\mathfrak m$ l'idéal maximal de $R$.
Nous allons montrer que $\mathfrak m \otimes_R M \to M$ est injectif,
puis appliquer le
lemme \ref{lemma-local-criterion-flatness}.
Supposons que $\sum f_i \otimes x_i \in \mathfrak m \otimes_R M$
et que $\sum f_i x_i = 0$ dans $M$. Par le critère équationnel
de platitude, lemme \ref{lemma-flat-eq}, appliqué à $M/IM$
sur $R/I$, nous voyons qu'il existe des $\overline{a}_{ij} \in R/I$
et des $\overline{y}_j \in M/IM$ tels que
$x_i \bmod IM = \sum_j \overline{a}_{ij} \overline{y}_j $
et $0 = \sum_i (f_i \bmod I) \overline{a}_{ij}$.
Soit $a_{ij} \in R$ un relèvement de $\overline{a}_{ij}$ et,
de même, soit $y_j \in M$ un relèvement de $\overline{y}_j$.
Nous voyons alors que
\begin{eqnarray*}
\sum f_i \otimes x_i
& = &
\sum f_i \otimes x_i +
\sum f_ia_{ij} \otimes y_j -
\sum f_i \otimes a_{ij} y_j
\\
& = &
\sum f_i \otimes (x_i - \sum a_{ij} y_j) +
\sum (\sum f_i a_{ij}) \otimes y_j
\end{eqnarray*}
Comme $x_i - \sum a_{ij} y_j \in IM$ et
$\sum f_i a_{ij} \in I$, nous voyons qu'il existe
un élément de $I \otimes_R M$ qui s'envoie sur notre
élément donné $\sum f_i \otimes x_i$ dans $\mathfrak m \otimes_R M$.
Mais $I \otimes_R M \to M$ est injectif par hypothèse (voir la
remarque \ref{remark-Tor-ring-mod-ideal}), ce qui conclut.
\end{proof}

\noindent
En particulier, dans la situation du
lemme \ref{lemma-variant-local-criterion-flatness}, supposons que
$I = (x)$ soit engendré par un seul élément $x$ qui soit
un non-diviseur de zéro dans $R$. Alors $\text{Tor}_1^R(M, R/(x)) = (0)$
si et seulement si $x$ est un non-diviseur de zéro sur $M$.

\begin{lemma}
\label{lemma-flat-module-powers}
Soit $R \to S$ un homomorphisme d'anneaux. Soit $I \subset R$ un idéal.
Soit $M$ un $S$-module. Supposons que
\begin{enumerate}
\item $R$ soit un anneau noethérien,
\item $S$ soit un anneau noethérien,
\item $M$ soit un $S$-module de type fini, et
\item pour tout $n \geq 1$, le module $M/I^n M$ soit plat sur
$R/I^n$.
\end{enumerate}
Alors, pour tout $\mathfrak q \in V(IS)$,
le localisé $M_{\mathfrak q}$ est plat sur $R$.
En particulier, si $S$ est local et si $IS$ est contenu
dans son idéal maximal, alors $M$ est plat sur $R$.
\end{lemma}

\begin{proof}
Nous allons utiliser le
lemme \ref{lemma-variant-local-criterion-flatness}.
Par hypothèse, $M/IM$ est plat sur $R/I$. Il suffit donc de vérifier
que $\text{Tor}_1^R(M, R/I)$ est nul après localisation en $\mathfrak q$. Par la
remarque \ref{remark-Tor-ring-mod-ideal},
ce groupe Tor est égal à $K = \Ker(I \otimes_R M \to M)$.
Nous savons que le noyau de
$I/I^n \otimes_{R/I^n} M/I^nM \to M/I^nM$ est nul pour tout $n \geq 1$.
Tout élément de $K$ s'envoie donc sur zéro dans $I/I^n \otimes_{R/I^n} M/I^nM$.
Comme
$$
I/I^n \otimes_{R/I^n} M/I^nM = I/I^n \otimes_R M =
(I \otimes_R M)/I^{n - 1}(I \otimes_R M)
$$
nous concluons que $K \subset I^{n - 1}(I \otimes_R M)$ pour tout $n \geq 1$.
Par le lemme d'Artin-Rees, et plus précisément par le
lemme \ref{lemma-intersection-powers-ideal-module},
nous concluons que $K_{\mathfrak q} = 0$, comme souhaité.
\end{proof}

\begin{lemma}
\label{lemma-surjective-on-tor-one}
Soient $R \to R' \to R''$ des homomorphismes d'anneaux.
Soit $M$ un $R$-module. Supposons que $M \otimes_R R'$
soit plat sur $R'$. Alors le morphisme naturel
$\text{Tor}_1^R(M, R') \otimes_{R'} R'' \to
\text{Tor}_1^R(M, R'')$ est surjectif.
\end{lemma}

\begin{proof}
Soit $F_\bullet$ une résolution libre de $M$ sur $R$.
Le complexe $F_2 \otimes_R R' \to F_1\otimes_R R' \to F_0 \otimes_R R'$
calcule $\text{Tor}_1^R(M, R')$.
Le complexe $F_2 \otimes_R R'' \to F_1\otimes_R R'' \to F_0 \otimes_R R''$
calcule $\text{Tor}_1^R(M, R'')$. Notons que
$F_i \otimes_R R' \otimes_{R'} R'' = F_i \otimes_R R''$. Posons
$K' = \Ker(F_1\otimes_R R' \to F_0 \otimes_R R')$ et,
de même, $K'' = \Ker(F_1\otimes_R R'' \to F_0 \otimes_R R'')$.
Nous avons donc une suite exacte
$$
0 \to K' \to F_1\otimes_R R' \to F_0 \otimes_R R' \to M \otimes_R R' \to 0.
$$
Puisque, par hypothèse, $M \otimes_R R'$ est plat sur $R'$,
la suite
$$
K' \otimes_{R'} R'' \to
F_1 \otimes_R R'' \to
F_0 \otimes_R R'' \to
M \otimes_R R'' \to 0
$$
reste exacte. Cela signifie que $K' \otimes_{R'} R'' \to K''$
est surjectif. Comme $\text{Tor}_1^R(M, R')$ est un quotient de $K'$ et que
$\text{Tor}_1^R(M, R'')$ est un quotient de $K''$, le résultat en découle.
\end{proof}

\begin{lemma}
\label{lemma-surjective-on-tor-one-trivial}
Soit $R \to R'$ un homomorphisme d'anneaux. Soient $I \subset R$
un idéal et $I' = IR'$. Soit $M$ un $R$-module
et posons $M' = M \otimes_R R'$. Le morphisme naturel
$\text{Tor}_1^R(R'/I', M) \to \text{Tor}_1^{R'}(R'/I', M')$
est surjectif.
\end{lemma}

\begin{proof}
Soit $F_2 \to F_1 \to F_0 \to M \to 0$ une résolution libre de
$M$ sur $R$. Posons $F_i' = F_i \otimes_R R'$. La suite
$F_2' \to F_1' \to F_0' \to M' \to 0$ peut ne plus être exacte
en $F_1'$. Une résolution libre de $M'$ sur $R'$ est donc de la forme
suivante :
$$
F_2' \oplus F_2'' \to F_1' \to F_0' \to M' \to 0
$$
pour un module libre convenable $F_2''$ sur $R'$. Ensuite, notons que
$F_i \otimes_R R'/I' = F_i' / IF_i' = F_i'/I'F_i'$.
Le complexe $F_2'/I'F_2' \to F_1'/I'F_1' \to F_0'/I'F_0'$
calcule donc $\text{Tor}_1^R(M, R'/I')$. D'autre part,
$F_i' \otimes_{R'} R'/I' = F_i'/I'F_i'$ et il en va de même
pour $F_2''$. Ainsi, le complexe
$F_2'/I'F_2' \oplus F_2''/I'F_2'' \to F_1'/I'F_1' \to F_0'/I'F_0'$
calcule $\text{Tor}_1^{R'}(M', R'/I')$. Comme le morphisme vertical
de complexes
$$
\xymatrix{
F_2'/I'F_2' \ar[r] \ar[d] &
F_1'/I'F_1' \ar[r] \ar[d] &
F_0'/I'F_0' \ar[d] \\
F_2'/I'F_2' \oplus F_2''/I'F_2'' \ar[r] &
F_1'/I'F_1' \ar[r] &
F_0'/I'F_0'
}
$$
induit manifestement une surjection en homologie, le résultat en découle.
\end{proof}

\begin{lemma}
\label{lemma-another-variant-local-criterion-flatness}
Soit
$$
\xymatrix{
S \ar[r] & S' \\
R \ar[r] \ar[u] & R' \ar[u]
}
$$
un diagramme commutatif d'homomorphismes locaux d'anneaux locaux noethériens.
Soit $I \subset R$ un idéal propre.
Soit $M$ un $S$-module de type fini.
Notons $I' = IR'$ et $M' = M \otimes_S S'$.
Supposons que
\begin{enumerate}
\item $S'$ soit un localisé du produit tensoriel
$S \otimes_R R'$,
\item $M/IM$ soit plat sur $R/I$,
\item $\text{Tor}_1^R(M, R/I) \to \text{Tor}_1^{R'}(M', R'/I')$
soit nul.
\end{enumerate}
Alors $M'$ est plat sur $R'$.
\end{lemma}

\begin{proof}
Puisque $S'$ est un localisé de $S \otimes_R R'$, nous voyons que
$M'$ est un localisé de $M \otimes_R R'$. Notons que,
par le lemme \ref{lemma-flat-base-change}, le module $M/IM \otimes_{R/I} R'/I'
= M \otimes_R R' /I'(M \otimes_R R')$ est plat sur $R'/I'$. Par conséquent,
$M'/I'M'$ est lui aussi plat sur $R'/I'$, puisque tout localisé d'un module plat
est plat. D'après le lemme \ref{lemma-variant-local-criterion-flatness},
il suffit de montrer que $\text{Tor}_1^{R'}(M', R'/I')$ est nul.
Comme $M'$ est un localisé de $M \otimes_R R'$, la dernière hypothèse
implique qu'il suffit de montrer que
$\text{Tor}_1^R(M, R/I) \otimes_R R'
\to
\text{Tor}_1^{R'}(M \otimes_R R', R'/I')$
est surjectif.

\medskip\noindent
Le lemme \ref{lemma-surjective-on-tor-one-trivial} montre que
$\text{Tor}_1^R(M, R'/I') \to \text{Tor}_1^{R'}(M \otimes_R R', R'/I')$
est surjectif. Il suffit donc maintenant de montrer que
$\text{Tor}_1^R(M, R/I) \otimes_R R'
\to
\text{Tor}_1^R(M, R'/I')$
est surjectif. Cela découle du lemme \ref{lemma-surjective-on-tor-one}
appliqué aux homomorphismes d'anneaux $R \to R/I \to R'/I'$ et au module $M$.
\end{proof}

\noindent
Veuillez comparer le lemme ci-dessous au
lemme \ref{lemma-criterion-flatness-fibre-nilpotent}
(le cas d'un idéal nilpotent) et au
lemme \ref{lemma-criterion-flatness-fibre}
(le cas des algèbres de présentation finie).

\begin{lemma}[Crit\`ere de platitude par fibres ; cas noethérien]
\label{lemma-criterion-flatness-fibre-Noetherian}
Soient $R$, $S$ et $S'$ des anneaux locaux noethériens, et soient $R \to S \to S'$
des homomorphismes locaux d'anneaux. Soit $\mathfrak m \subset R$
l'idéal maximal. Soit $M$ un $S'$-module. Supposons que
\begin{enumerate}
\item le module $M$ soit de type fini sur $S'$ ;
\item le module $M$ soit non nul ;
\item le module $M/\mathfrak m M$
soit un $S/\mathfrak m S$-module plat ;
\item le module $M$ soit un $R$-module plat.
\end{enumerate}
Alors $S$ est plat sur $R$ et $M$ est un $S$-module plat.
\end{lemma}

\begin{proof}
Posons $I = \mathfrak mS \subset S$. Nous voyons que $M/IM$ est un
$S/I$-module plat en vertu de (3). Comme
$\mathfrak m \otimes_R S' \to I \otimes_S S'$ est surjectif, nous voyons
que $\mathfrak m \otimes_R M \to I \otimes_S M$ est également surjectif.
Considérons
$$
\mathfrak m \otimes_R M \to I \otimes_S M \to M.
$$
Comme $M$ est plat sur $R$, la composée est injective,
et les deux flèches le sont donc.
En particulier, $\text{Tor}_1^S(S/I, M) = 0$, voir la
remarque \ref{remark-Tor-ring-mod-ideal}. Par le
lemme \ref{lemma-variant-local-criterion-flatness}, nous concluons
que $M$ est plat sur $S$. Notons que, puisque $M/\mathfrak m_{S'}M$
n'est pas nul par le lemme de Nakayama \ref{lemma-NAK},
nous voyons qu'en fait $M$ est fidèlement plat sur $S$ par le
lemme \ref{lemma-ff} (car cela impose $M/\mathfrak m_SM \not = 0$).

\medskip\noindent
Considérons la suite exacte
$0 \to \mathfrak m \to R \to \kappa \to 0$.
Elle donne une suite exacte
$0 \to \text{Tor}_1^R(\kappa, S) \to \mathfrak m \otimes_R S \to I \to 0$.
Puisque $M$ est plat sur $S$, celle-ci donne une suite exacte
$0 \to \text{Tor}_1^R(\kappa, S)\otimes_S M \to
\mathfrak m \otimes_R M \to I \otimes_S M \to 0$.
D'après ce qui précède, cela implique que $\text{Tor}_1^R(\kappa, S)\otimes_S M = 0$.
Puisque $M$ est fidèlement plat sur $S$, il s'ensuit que
$\text{Tor}_1^R(\kappa, S) = 0$, et nous concluons que
$S$ est plat sur $R$ par le lemme \ref{lemma-local-criterion-flatness}.
\end{proof}

\begin{lemma}
\label{lemma-flatness-scallop}
Soient $A$ un anneau, $M$ un $A$-module et $f \in A$. Si
\begin{enumerate}
\item $f$ est un non-diviseur de zéro sur $A$ et sur $M$,
\item $M_f$ est un $A_f$-module plat, et
\item $M/fM$ est un $A/fA$-module plat,
\end{enumerate}
alors $M$ est un $A$-module plat. Il en va de même en remplaçant « plat » par
« fidèlement plat ».
\end{lemma}

\begin{proof}
Comme $0 \to A \to A \to A/fA \to 0$ et $0 \to M \to M \to M/fM \to 0$ sont
exactes, nous obtenons $\text{Tor}_i^A(M, A/fA) = 0$ pour $i = 1$ et $i = 2$.
Par le lemme \ref{lemma-what-does-it-mean}, nous concluons que
$\text{Tor}_1^A(M, N) = 0$ pour tout $A$-module $N$ annulé par $f$
(ceci utilise la platitude de $M/fM$ sur $A/fA$).
Étant donné un $A$-module $N$ annulé par $f$, nous pouvons
choisir une suite exacte courte $0 \to N' \to F \to N \to 0$ de $A$-modules,
où $F$ est une somme directe de copies de $A/fA$. De la suite exacte
$$
\text{Tor}_2^A(M, F) \to
\text{Tor}_2^A(M, N) \to
\text{Tor}_1^A(M, N')
$$
nous concluons que $\text{Tor}_2^A(M, N) = 0$ pour tout $A$-module
$N$ annulé par $f$. Soit ensuite $K$ un $A$-module arbitraire.
Nous pouvons décomposer le morphisme $f : K \to K$ en deux suites exactes courtes
$$
0 \to K[f] \to K \to K' \to 0
\quad\text{et}\quad
0 \to K' \to K \to K/fK \to 0
$$
où $K' = K/K[f] \cong fK$. En appliquant les suites exactes de Tor, nous
obtenons les suites exactes
$$
\text{Tor}_1^A(K[f], M) \to
\text{Tor}_1^A(K, M) \to
\text{Tor}_1^A(K', M)
$$
et
$$
\text{Tor}_2^A(K/fK, M) \to
\text{Tor}_1^A(K', M) \to
\text{Tor}_1^A(K, M)
$$
En utilisant l'annulation de
$\text{Tor}_1^A(K[f], M)$ et de
$\text{Tor}_2^A(K/fK, M)$,
nous concluons que $f : K \to K$ induit un morphisme injectif sur
$\text{Tor}^A_1(M, K)$. Autrement dit, nous voyons que $\text{Tor}^A_1(M, K)$
est nul si et seulement si $\text{Tor}^A_1(M, K) \otimes_A A_f$ est nul.
Par le lemme \ref{lemma-flat-base-change-tor}, nous avons
$$
\text{Tor}^A_1(M, K) \otimes_A A_f = \text{Tor}^{A_f}_1(M_f, K_f)  = 0
$$
Cette annulation résulte de la platitude de $M_f$ sur $A_f$
(lemme \ref{lemma-characterize-flat}).
Nous concluons que $\text{Tor}_1^A(M, K) = 0$
pour tout $A$-module $K$. Par conséquent,
$M$ est plat sur $A$ par le lemme \ref{lemma-characterize-flat}.

\medskip\noindent
Nous omettons l'argument dans le cas des modules fidèlement plats.
\end{proof}

\begin{lemma}
\label{lemma-flatness-scallop-pre}
Soient $A$ un anneau et $M$ un $A$-module. Soit $I = (f_1, \ldots, f_r)$
un idéal de $A$ engendré par $r \geq 1$ éléments. Si
\begin{enumerate}
\item $M_{f_i}$ est un $A_{f_i}$-module plat pour $i = 1, \ldots, r$,
\item $M/IM$ est un $A/I$-module plat,
\item $\text{Tor}_i^A(M, A/I) = 0$ pour $i = 1, \ldots, r + 1$,
\end{enumerate}
alors $M$ est un $A$-module plat. Il en va de même en remplaçant « plat » par
« fidèlement plat ».
\end{lemma}

\begin{proof}
Le lemme \ref{lemma-what-does-it-mean} montre que
$\text{Tor}_1^A(M, K) = 0$ pour tout $A$-module $K$ annulé par $I$.
Étant donné un $A$-module $K$ annulé par $I$, nous pouvons choisir une suite
exacte courte $0 \to N \to F \to K \to 0$ de $A$-modules, où $F$
est une somme directe de copies de $A/I$. Nous obtenons une suite exacte
$$
\text{Tor}_i^A(M, F) \to
\text{Tor}_i^A(M, K) \to
\text{Tor}_{i - 1}^A(M, N)
$$
Ainsi, en utilisant l'hypothèse (3) et en raisonnant par récurrence sur $i$,
nous concluons que $\text{Tor}_i^A(M, K) = 0$
pour tout $A$-module $K$ annulé par $I$ et $i = 1, \ldots, r + 1$.

\medskip\noindent
Supposons que, pour un certain $1 \leq j \leq r$,
nous ayons montré que $\text{Tor}_i^A(M, K) = 0$
pour tout $A$-module $K$ annulé par $f_1, \ldots, f_j$
et $i = 1, \ldots, j + 1$.
Soit $K$ un $A$-module annulé par $f_1, \ldots, f_{j - 1}$.
Nous pouvons décomposer le morphisme $f_j : K \to K$ en deux suites exactes courtes
$$
0 \to K[f_j] \to K \to K' \to 0
\quad\text{et}\quad
0 \to K' \to K \to K/f_jK \to 0
$$
où $K' = K/K[f_j] \cong f_jK$. Soit $1 \leq i \leq j$.  En appliquant les
suites exactes de Tor, nous obtenons les suites exactes
$$
\text{Tor}_i^A(K[f_j], M) \to
\text{Tor}_i^A(K, M) \to
\text{Tor}_i^A(K', M)
$$
et
$$
\text{Tor}_{i + 1}^A(K/f_jK, M) \to
\text{Tor}_i^A(K', M) \to
\text{Tor}_i^A(K, M)
$$
En utilisant l'annulation de $\text{Tor}_i^A(K[f_j], M)$ et de
$\text{Tor}_{i + 1}^A(K/f_jK, M)$, nous concluons que $f_j : K \to K$
induit un morphisme injectif sur $\text{Tor}^A_i(M, K)$.
Autrement dit, nous voyons que $\text{Tor}^A_i(M, K)$
est nul si et seulement si $\text{Tor}^A_i(M, K) \otimes_A A_{f_j}$ est nul.
Par le lemme \ref{lemma-flat-base-change-tor}, nous avons
$$
\text{Tor}^A_i(M, K) \otimes_A A_{f_j} =
\text{Tor}^{A_{f_j}}_i(M_{f_j}, K_{f_j})  = 0
$$
Cette annulation résulte de la platitude de $M_{f_j}$ sur $A_{f_j}$
(lemme \ref{lemma-characterize-flat}).
Nous concluons que $\text{Tor}_i^A(M, K) = 0$
pour tout $A$-module $K$ annulé par $f_1, \ldots, f_{j - 1}$
et $i = 1, \ldots, j$.
Par récurrence descendante sur $j$, nous concluons que
ceci vaut pour $j = 0$, c'est-à-dire que
$\text{Tor}_1^A(M, K) = 0$
pour tout $A$-module $K$. Par conséquent,
$M$ est plat sur $A$ par le lemme \ref{lemma-characterize-flat}.

\medskip\noindent
Nous omettons l'argument dans le cas des modules fidèlement plats.
\end{proof}








\section{Changement de base et platitude}
\label{section-base-change-flat}

\noindent
Quelques lemmes décrivant ce qu'il advient de la platitude lors
d'un changement de base.

\begin{lemma}
\label{lemma-base-change-flat-up-down}
Soit
$$
\xymatrix{
S \ar[r] & S' \\
R \ar[r] \ar[u] & R' \ar[u]
}
$$
un diagramme commutatif d'homomorphismes locaux d'anneaux locaux.
Supposons que $S'$ soit un localisé du produit tensoriel $S \otimes_R R'$.
Soit $M$ un $S$-module et posons $M' = S' \otimes_S M$.
\begin{enumerate}
\item Si $M$ est plat sur $R$, alors $M'$ est plat sur $R'$.
\item Si $M'$ est plat sur $R'$ et si $R \to R'$ est plat, alors
$M$ est plat sur $R$.
\end{enumerate}
En particulier, nous avons
\begin{enumerate}
\item[(3)] Si $S$ est plat sur $R$, alors $S'$ est plat sur $R'$.
\item[(4)] Si $R' \to S'$ et $R \to R'$ sont plats, alors $S$ est plat sur $R$.
\end{enumerate}
\end{lemma}

\begin{proof}
Démonstration de (1). Si $M$ est plat sur $R$, alors $M \otimes_R R'$
est plat sur $R'$ par le
lemme \ref{lemma-flat-base-change}.
Si $W \subset S \otimes_R R'$ est la partie multiplicative telle que
$W^{-1}(S \otimes_R R') = S'$, alors $M' = W^{-1}(M \otimes_R R')$.
Ainsi, $M'$ est plat sur $R'$ en tant que localisé d'un module plat, voir la
partie (5) du lemme \ref{lemma-flat-localization}. Ceci démontre (1) et,
en particulier, montre que (3) est vraie.

\medskip\noindent
Démonstration de (2). Supposons que $M'$ soit plat sur $R'$ et que $R \to R'$ soit plat.
En appliquant (3) au diagramme réfléchi par rapport à la diagonale nord-ouest,
nous voyons que $S \to S'$ est plat. Ainsi, $S \to S'$ est fidèlement plat par le
lemme \ref{lemma-local-flat-ff}.
Nous allons utiliser le critère du
lemme \ref{lemma-flat} (\ref{item-f-ideal})
pour montrer que $M$ est plat.
Soit $I \subset R$ un idéal. Si $I \otimes_R M \to M$ a un noyau non nul,
il en va de même de $(I \otimes_R M) \otimes_S S' \to M \otimes_S S' = M'$.
Notons que $I \otimes_R R' = IR'$, puisque $R \to R'$ est plat, et que
$$
(I \otimes_R M) \otimes_S S' =
(I \otimes_R R') \otimes_{R'} (M \otimes_S S') =
IR' \otimes_{R'} M'.
$$
La platitude de $M'$ sur $R'$
implique que ce module s'injecte dans $M'$.
Ceci achève la démonstration de (2), et (4) est donc également vraie.
\end{proof}

\noindent
Voici encore une application du critère local de platitude.

\begin{lemma}
\label{lemma-yet-another-variant-local-criterion-flatness}
Considérons un diagramme commutatif d'anneaux locaux et d'homomorphismes locaux
$$
\xymatrix{
S \ar[r] & S' \\
R \ar[r] \ar[u] & R' \ar[u]
}
$$
Soit $M$ un $S$-module de type fini. Supposons que
\begin{enumerate}
\item les flèches horizontales soient des homomorphismes d'anneaux plats,
\item $M$ soit plat sur $R$,
\item $\mathfrak m_R R' = \mathfrak m_{R'}$,
\item $R'$ et $S'$ soient noethériens.
\end{enumerate}
Alors $M' = M \otimes_S S'$ est plat sur $R'$.
\end{lemma}

\begin{proof}
Comme $\mathfrak m_R \subset R$ et que $R \to R'$ est plat, nous obtenons
$\mathfrak m_R \otimes_R R' = \mathfrak m_R R' = \mathfrak m_{R'}$
par l'hypothèse (3). Observons que $M'$ est un $S'$-module de type fini
qui est plat sur $R$ par le lemme \ref{lemma-flatness-descends-more-general}.
Ainsi, $\mathfrak m_R \otimes_R M' \to M'$ est injectif.
Nous obtenons alors
$$
\mathfrak m_R \otimes_R M' =
\mathfrak m_R \otimes_R R' \otimes_{R'} M' =
\mathfrak m_{R'} \otimes_{R'} M'
$$
Ainsi, $\mathfrak m_{R'} \otimes_{R'} M' \to M'$ est injectif.
Ceci montre que $\text{Tor}_1^{R'}(\kappa_{R'}, M') = 0$
(remarque \ref{remark-Tor-ring-mod-ideal}).
Le module $M'$ est donc plat sur $R'$ par le
lemme \ref{lemma-local-criterion-flatness}.
\end{proof}







\section{Critères de platitude sur les anneaux artiniens}
\label{section-flatness-artinian}

\noindent
Nous examinons quelques critères de platitude pour les modules sur les anneaux artiniens.
Notons que l'idéal maximal d'un anneau local artinien est nilpotent,
de sorte que les deux lemmes suivants s'appliquent aux anneaux locaux artiniens.

\begin{lemma}
\label{lemma-local-artinian-basis-when-flat}
Soit $(R, \mathfrak m)$ un anneau local dont l'idéal maximal
$\mathfrak m$ est nilpotent. Soit $M$ un $R$-module plat.
Si $A$ est un ensemble et si $x_\alpha \in M$, $\alpha \in A$, est une famille
d'éléments de $M$, alors les assertions suivantes sont équivalentes :
\begin{enumerate}
\item $\{\overline{x}_\alpha\}_{\alpha \in A}$ est une base
de l'espace vectoriel $M/\mathfrak mM$ sur $R/\mathfrak m$, et
\item $\{x_\alpha\}_{\alpha \in A}$ est une base de $M$ sur $R$.
\end{enumerate}
\end{lemma}

\begin{proof}
L'implication (2) $\Rightarrow$ (1) est immédiate.
Supposons (1). Par le lemme de Nakayama \ref{lemma-NAK},
les éléments $x_\alpha$ engendrent $M$. On obtient alors une suite exacte
courte
$$
0 \to K \to \bigoplus\nolimits_{\alpha \in A} R \to M \to 0
$$
En tensorisant par $R/\mathfrak m$ et en utilisant le lemme \ref{lemma-flat-tor-zero},
nous obtenons $K/\mathfrak mK = 0$. Par le lemme de Nakayama \ref{lemma-NAK},
nous concluons que $K = 0$.
\end{proof}

\begin{lemma}
\label{lemma-local-artinian-characterize-flat}
Soit $R$ un anneau local d'idéal maximal nilpotent. Soit $M$ un $R$-module.
Les assertions suivantes sont équivalentes :
\begin{enumerate}
\item $M$ est plat sur $R$,
\item $M$ est un $R$-module libre, et
\item $M$ est un $R$-module projectif.
\end{enumerate}
\end{lemma}

\begin{proof}
Comme tout module projectif est plat (en tant que facteur direct d'un module libre)
et que tout module libre est projectif, il suffit de montrer qu'un module plat
est libre. Soit $M$ un module plat. Soient $A$ un ensemble et des éléments $x_\alpha \in M$,
$\alpha \in A$, tels que les
$\overline{x_\alpha} \in M/\mathfrak m M$ forment une base sur le corps
résiduel de $R$. Par le
lemme \ref{lemma-local-artinian-basis-when-flat},
les $x_\alpha$ sont une base de $M$ sur $R$, ce qui conclut.
\end{proof}

\begin{lemma}
\label{lemma-lift-basis}
Soit $R$ un anneau.
Soit $I \subset R$ un idéal.
Soit $M$ un $R$-module.
Soient $A$ un ensemble et $x_\alpha \in M$, $\alpha \in A$, une famille
d'éléments de $M$.
Supposons que
\begin{enumerate}
\item $I$ soit nilpotent,
\item $\{\overline{x}_\alpha\}_{\alpha \in A}$ soit une base de $M/IM$ sur
$R/I$, et
\item $\text{Tor}_1^R(R/I, M) = 0$.
\end{enumerate}
Alors $M$ est libre de base $\{x_\alpha\}_{\alpha \in A}$ sur $R$.
\end{lemma}

\begin{proof}
Soient $R$, $I$, $M$, $\{x_\alpha\}_{\alpha \in A}$ comme dans le lemme,
et supposons les hypothèses (1), (2) et (3) satisfaites. Par le
lemme de Nakayama \ref{lemma-NAK},
les éléments $x_\alpha$ engendrent $M$ sur $R$.
L'hypothèse $\text{Tor}_1^R(R/I, M) = 0$ implique que nous avons une suite
exacte courte
$$
0 \to I \otimes_R M \to M \to M/IM \to 0.
$$
Soit $\sum f_\alpha x_\alpha = 0$ une relation dans $M$.
Par le choix des $x_\alpha$, nous voyons que $f_\alpha \in I$.
Nous concluons donc que $\sum f_\alpha \otimes x_\alpha = 0$ dans
$I \otimes_R M$. Le morphisme $I \otimes_R M \to I/I^2 \otimes_{R/I} M/IM$
et le fait que $\{x_\alpha\}_{\alpha \in A}$ soit une base
de $M/IM$ impliquent que $f_\alpha \in I^2$ ! Nous concluons donc qu'il
n'existe aucune relation entre les images des $x_\alpha$ dans
$M/I^2M$. Autrement dit, nous voyons que $M/I^2M$ est libre, de base
les images des $x_\alpha$. En utilisant le morphisme
$I \otimes_R M \to I/I^3 \otimes_{R/I^2} M/I^2M$,
nous concluons ensuite que $f_\alpha \in I^3$ !
Et ainsi de suite. Comme $I^n = 0$ pour un certain $n$ par l'hypothèse (1), le résultat suit.
\end{proof}

\begin{lemma}
\label{lemma-prepare-lift-flatness}
Soit $\varphi : R \to R'$ un homomorphisme d'anneaux.
Soit $I \subset R$ un idéal.
Soit $M$ un $R$-module.
Supposons que
\begin{enumerate}
\item $M/IM$ soit plat sur $R/I$, et
\item $R' \otimes_R M$ soit plat sur $R'$.
\end{enumerate}
Posons $I_2 = \varphi^{-1}(\varphi(I^2)R')$.
Alors $M/I_2M$ est plat sur $R/I_2$.
\end{lemma}

\begin{proof}
Nous pouvons remplacer $R$, $M$ et $R'$ par $R/I_2$, $M/I_2M$ et
$R'/\varphi(I)^2R'$. Alors $I^2 = 0$ et $\varphi$ est injectif. Par le
lemme \ref{lemma-what-does-it-mean}
et puisque $I^2 = 0$, il suffit de montrer que
$\text{Tor}^R_1(R/I, M) = K = \Ker(I \otimes_R M \to M)$ est nul.
Posons $M' = M \otimes_R R'$ et $I' = IR'$.
Par hypothèse, le morphisme $I' \otimes_{R'} M' \to M'$ est injectif.
Ainsi, $K$ s'envoie sur zéro dans
$$
I' \otimes_{R'} M' = I' \otimes_R M = I' \otimes_{R/I} M/IM.
$$
Le morphisme $I \to I'$ est alors un morphisme injectif de $R/I$-modules.
Comme $M/IM$ est plat sur $R/I$, le morphisme
$$
I \otimes_{R/I} M/IM \longrightarrow I' \otimes_{R/I} M/IM
$$
est injectif. Ceci implique que $K$ est nul dans
$I \otimes_R M = I \otimes_{R/I} M/IM$, comme souhaité.
\end{proof}

\begin{lemma}
\label{lemma-lift-flatness}
Soit $\varphi : R \to R'$ un homomorphisme d'anneaux.
Soit $I \subset R$ un idéal.
Soit $M$ un $R$-module.
Supposons que
\begin{enumerate}
\item $I$ soit nilpotent,
\item $R \to R'$ soit injectif,
\item $M/IM$ soit plat sur $R/I$, et
\item $R' \otimes_R M$ soit plat sur $R'$.
\end{enumerate}
Alors $M$ est plat sur $R$.
\end{lemma}

\begin{proof}
Définissons par récurrence $I_1 = I$ et $I_{n + 1} = \varphi^{-1}(\varphi(I_n)^2R')$
pour $n \geq 1$. Notons que, par le
lemme \ref{lemma-prepare-lift-flatness},
nous obtenons que $M/I_nM$ est plat sur $R/I_n$ pour tout $n \geq 1$.
Il est clair que $\varphi(I_{n + 1}) \subset \varphi(I)^{2^n}R'$. Comme
$I$ est nilpotent, nous voyons que $\varphi(I_n) = 0$ pour un certain $n$. Puisque
$\varphi$ est injectif, nous concluons que $I_n = 0$ pour un certain $n$, ce qui
achève la démonstration.
\end{proof}

\noindent
Voici la version artinienne locale du critère local de platitude.

\begin{lemma}
\label{lemma-artinian-variant-local-criterion-flatness}
Soit $R$ un anneau local artinien. Soit $M$ un $R$-module.
Soit $I \subset R$ un idéal propre. Les assertions suivantes sont
équivalentes :
\begin{enumerate}
\item $M$ est plat sur $R$, et
\item $M/IM$ est plat sur $R/I$ et $\text{Tor}_1^R(R/I, M) = 0$.
\end{enumerate}
\end{lemma}

\begin{proof}
L'implication (1) $\Rightarrow$ (2) découle immédiatement des
définitions. Supposons que $M/IM$ soit plat sur $R/I$ et que
$\text{Tor}_1^R(R/I, M) = 0$. Par le
lemme \ref{lemma-local-artinian-characterize-flat},
ceci implique que $M/IM$ est libre sur $R/I$. Choisissons un ensemble $A$
et des éléments $x_\alpha \in M$ dont les images dans $M/IM$ forment
une base. Par le
lemme \ref{lemma-lift-basis},
nous concluons que $M$ est libre et, en particulier, plat.
\end{proof}

\noindent
Il se trouve que la platitude descend par tout homomorphisme injectif
dont la source est un anneau artinien.

\begin{lemma}
\label{lemma-descent-flatness-injective-map-artinian-rings}
Soit $R \to S$ un homomorphisme d'anneaux. Soit $M$ un $R$-module.
Supposons que
\begin{enumerate}
\item $R$ soit artinien,
\item $R \to S$ soit injectif, et
\item $M \otimes_R S$ soit un $S$-module plat.
\end{enumerate}
Alors $M$ est un $R$-module plat.
\end{lemma}

\begin{proof}
Première démonstration : soit $I \subset R$ le radical de Jacobson de $R$.
Alors $I$ est nilpotent et $M/IM$ est plat sur $R/I$, puisque $R/I$
est un produit de corps, voir la
section \ref{section-artinian}.
Par conséquent, $M$ est plat par une application du
lemme \ref{lemma-lift-flatness}.

\medskip\noindent
Seconde démonstration : par le
lemme \ref{lemma-artinian-finite-length},
nous pouvons écrire $R = \prod R_i$ comme produit fini d'anneaux locaux
artiniens. Ceci induit des décompositions en produits analogues pour $R$ et $S$.
Nous nous ramenons donc au cas où $R$ est local artinien (détails omis).

\medskip\noindent
Supposons que $R \to S$ et $M$ soient comme dans le lemme, satisfassent (1), (2) et (3),
et qu'en outre $R$ soit local d'idéal maximal $\mathfrak m$.
Soient $A$ un ensemble et $x_\alpha \in A$ des éléments tels que
$\overline{x}_\alpha$ forme une base de $M/\mathfrak mM$
sur $R/\mathfrak m$. Par le
lemme de Nakayama \ref{lemma-NAK},
nous voyons que les éléments $x_\alpha$ engendrent $M$ comme $R$-module.
Posons $N = S \otimes_R M$ et $I = \mathfrak mS$.
Alors $\{1 \otimes x_\alpha\}_{\alpha \in A}$ est une famille d'éléments
de $N$ qui forme une base de $N/IN$. De plus, puisque $N$ est plat sur
$S$, nous avons $\text{Tor}_1^S(S/I, N) = 0$. Le
lemme \ref{lemma-lift-basis}
permet donc de conclure que $N$ est libre de base $\{1 \otimes x_\alpha\}_{\alpha \in A}$.
L'injectivité de $R \to S$ garantit alors qu'il ne peut exister de
relation non triviale entre les $x_\alpha$ à coefficients dans $R$.
\end{proof}

\noindent
Veuillez comparer le lemme ci-dessous au
lemme \ref{lemma-criterion-flatness-fibre-Noetherian}
(le cas des anneaux locaux noethériens), au
lemme \ref{lemma-criterion-flatness-fibre}
(le cas des algèbres de présentation finie), et au
lemme \ref{lemma-criterion-flatness-fibre-locally-nilpotent}
(le cas des idéaux localement nilpotents).

\begin{lemma}[Crit\`ere de platitude par fibres : cas nilpotent]
\label{lemma-criterion-flatness-fibre-nilpotent}
Soit
$$
\xymatrix{
S \ar[rr] & & S' \\
& R \ar[lu] \ar[ru]
}
$$
un diagramme commutatif dans la catégorie des anneaux.
Soient $I \subset R$ un idéal nilpotent et $M$ un $S'$-module. Supposons que
\begin{enumerate}
\item le module $M/IM$ soit un $S/IS$-module plat ;
\item le module $M$ soit un $R$-module plat.
\end{enumerate}
Alors $M$ est un $S$-module plat et $S_{\mathfrak q}$ est plat sur $R$
pour tout $\mathfrak q \subset S$ tel que $M \otimes_S \kappa(\mathfrak q)$
soit non nul.
\end{lemma}

\begin{proof}
Puisque $M$ est plat sur $R$, la tensorisation par la suite exacte
courte $0 \to I \to R \to R/I \to 0$ donne une suite exacte courte
$$
0 \to I \otimes_R M \to M \to M/IM \to 0.
$$
Notons que $I \otimes_R M \to IS \otimes_S M$ est surjectif. Combiné avec
ce qui précède, cela signifie que les deux morphismes de
$$
I \otimes_R M \to IS \otimes_S M \to M
$$
sont injectifs. Ainsi, $\text{Tor}_1^S(IS, M) = 0$ (voir la
remarque \ref{remark-Tor-ring-mod-ideal}),
et nous concluons que $M$ est un $S$-module plat par le
lemme \ref{lemma-what-does-it-mean}.
Pour terminer, nous devons montrer que $S_{\mathfrak q}$ est plat sur
$R$ pour tout idéal premier $\mathfrak q \subset S$ tel que
$M \otimes_S \kappa(\mathfrak q)$ soit non nul. Ceci découle du
lemme \ref{lemma-ff} et du lemme \ref{lemma-flat-permanence}.
\end{proof}








\section{Qu'est-ce qui rend un complexe exact ?}
\label{section-complex-exact}

\noindent
Une partie de ce matériau se trouve dans l'article \cite{WhatExact}
de Buchsbaum et Eisenbud.

\begin{situation}
\label{situation-complex}
Ici, $R$ est un anneau, et nous avons un complexe
$$
0
\to
R^{n_e}
\xrightarrow{\varphi_e}
R^{n_{e-1}}
\xrightarrow{\varphi_{e-1}}
\ldots
\xrightarrow{\varphi_{i + 1}}
R^{n_i}
\xrightarrow{\varphi_i}
R^{n_{i-1}}
\xrightarrow{\varphi_{i-1}}
\ldots
\xrightarrow{\varphi_1}
R^{n_0}
$$
Autrement dit, nous imposons $\varphi_i \circ \varphi_{i + 1} = 0$
pour $i = 1, \ldots, e - 1$.
\end{situation}

\begin{lemma}
\label{lemma-add-trivial-complex}
Supposons que $R$ soit un anneau. Soit
$$
\ldots
\xrightarrow{\varphi_{i + 1}}
R^{n_i}
\xrightarrow{\varphi_i}
R^{n_{i-1}}
\xrightarrow{\varphi_{i-1}}
\ldots
$$
un complexe de $R$-modules libres de rang fini. Supposons que, pour un certain $i$,
un coefficient matriciel du morphisme $\varphi_i$ soit inversible.
Alors le complexe affiché est isomorphe à la somme directe d'un complexe
$$
\ldots \to
R^{n_{i + 2}} \xrightarrow{\varphi_{i + 2}}
R^{n_{i + 1}} \to
R^{n_i - 1} \to
R^{n_{i - 1} - 1} \to
R^{n_{i - 2}} \xrightarrow{\varphi_{i - 2}}
R^{n_{i - 3}} \to
\ldots
$$
et du complexe $\ldots \to 0 \to R \to R \to 0 \to \ldots$,
où le morphisme $R \to R$ est l'identité.
\end{lemma}

\begin{proof}
L'hypothèse signifie, après un changement de base de
$R^{n_i}$ et de $R^{n_{i-1}}$, que le premier vecteur de la base
de $R^{n_i}$ est envoyé par $\varphi_i$ sur le premier vecteur de la base
de $R^{n_{i-1}}$. Notons $e_j$ le
$j$-ème vecteur de la base de $R^{n_i}$ et $f_k$ le $k$-ème
vecteur de la base de $R^{n_{i-1}}$. Écrivons $\varphi_i(e_j)
= \sum a_{jk} f_k$. Ainsi, $a_{1k} = 0$ sauf si $k = 1$,
et $a_{11} = 1$. Changeons de nouveau la base de $R^{n_i}$
en posant $e'_j = e_j - a_{j1} e_1$ pour $j > 1$.
Après ce changement de coordonnées, nous avons $a_{j1} = 0$
pour $j > 1$. Notons que l'image
de $R^{n_{i + 1}} \to R^{n_i}$ est contenue dans le
sous-module engendré par les $e_j$, $j > 1$. Notons également
que $R^{n_{i-1}} \to R^{n_{i-2}}$ doit annuler
$f_1$, puisqu'il appartient à l'image. Ces conditions
et la forme de la matrice $(a_{jk})$ de $\varphi_i$
impliquent le lemme.
\end{proof}

\noindent
Dans la situation \ref{situation-complex}, nous disons qu'un complexe de la forme
$$
0 \to \ldots \to 0 \to R \xrightarrow{1} R \to 0 \to \ldots \to 0
$$
ou de la forme
$$
0 \to \ldots \to 0 \to R
$$
est {\it trivial}. Plus précisément, nous disons que
$0 \to R^{n_e} \to R^{n_{e-1}} \to \ldots \to R^{n_0}$
est trivial s'il existe soit un $e \geq i \geq 1$
tel que $n_i = n_{i - 1} = 1$, $\varphi_i = \text{id}_R$, et que
$n_j = 0$ pour $j \not \in \{i, i - 1\}$, soit
$n_0 = 1$ et $n_i = 0$ pour $i > 0$.
Le lemme précédent dit clairement que
tout complexe fini de modules libres de rang fini sur un anneau local est, à des
sommes directes de complexes triviaux près, identique à un complexe
dont tous les morphismes ont tous leurs coefficients matriciels dans
l'idéal maximal.

\begin{lemma}
\label{lemma-exact-depth-zero-local}
Dans la situation \ref{situation-complex}, supposons que $R$ soit
un anneau local noethérien d'idéal maximal $\mathfrak m$.
Supposons que $\mathfrak m \in \text{Ass}(R)$, autrement dit que
$R$ soit de profondeur $0$. Supposons que
$0 \to R^{n_e} \to R^{n_{e-1}} \to \ldots \to R^{n_0}$
soit exact en $R^{n_e}, \ldots, R^{n_1}$.
Alors le complexe est isomorphe à une somme directe de complexes
triviaux.
\end{lemma}

\begin{proof}
Choisissons $x \in R$, $x \not = 0$, tel que $\mathfrak m x = 0$.
Soit $i$ le plus grand indice tel que $n_i > 0$.
Si $i = 0$, alors l'assertion est vraie. Si
$i > 0$, notons $f_1$ le premier vecteur de la base de $R^{n_i}$.
Comme $xf_1$ ne s'envoie pas sur zéro en raison de
l'exactitude du complexe, nous en déduisons qu'un coefficient
matriciel du morphisme $R^{n_i} \to R^{n_{i - 1}}$
n'appartient pas à $\mathfrak m$.
Le lemme \ref{lemma-add-trivial-complex} nous permet alors
de diminuer $n_e + \ldots + n_1$. Une récurrence achève la démonstration.
\end{proof}

\begin{lemma}
\label{lemma-exact-artinian-local}
Dans la situation \ref{situation-complex}, soit $R$ un anneau local artinien.
Supposons que $0 \to R^{n_e} \to R^{n_{e-1}} \to \ldots \to R^{n_0}$
soit exact en $R^{n_e}, \ldots, R^{n_1}$. Alors le complexe est isomorphe
à une somme directe de complexes triviaux.
\end{lemma}

\begin{proof}
C'est un cas particulier du lemme \ref{lemma-exact-depth-zero-local},
car tout anneau local artinien est de profondeur $0$.
\end{proof}

\noindent
Nous définissons ci-dessous le rang d'un morphisme de modules libres de rang fini.
Ce n'est qu'une définition possible du rang. C'est
simplement celle qui convient dans cette section ; il en
existe d'autres qui peuvent être plus commodes dans d'autres situations.

\begin{definition}
\label{definition-rank}
Soit $R$ un anneau. Supposons que $\varphi : R^m \to R^n$ soit un morphisme
de modules libres de rang fini.
\begin{enumerate}
\item Le {\it rang} de $\varphi$ est le plus grand $r$ tel que
$\wedge^r \varphi : \wedge^r R^m \to \wedge^r R^n$ soit non nul.
\item Nous notons $I(\varphi) \subset R$ l'idéal engendré par
les mineurs $r \times r$ de la matrice de $\varphi$, où $r$
est le rang défini ci-dessus.
\end{enumerate}
\end{definition}

\noindent
Le rang de $\varphi : R^m \to R^n$ vaut $0$
si et seulement si $\varphi = 0$, et dans ce cas $I(\varphi) = R$.

\begin{lemma}
\label{lemma-trivial-case-exact}
Dans la situation \ref{situation-complex}, supposons que le complexe soit
isomorphe à une somme directe de complexes triviaux. Alors
nous avons
\begin{enumerate}
\item les morphismes $\varphi_i$ ont pour rang
$r_i = n_i - n_{i + 1} + \ldots + (-1)^{e-i-1} n_{e-1} + (-1)^{e-i} n_e$,
\item pour tout $i$, $1 \leq i \leq e - 1$, nous avons
$\text{rank}(\varphi_{i + 1}) + \text{rank}(\varphi_i) = n_i$,
\item chaque $I(\varphi_i) = R$.
\end{enumerate}
\end{lemma}

\begin{proof}
Nous pouvons supposer que le complexe est la somme directe de complexes
triviaux. Pour chaque $i$, nous pouvons alors répartir les vecteurs de la base
canonique de $R^{n_i}$ entre ceux qui s'envoient sur un vecteur de la base
de $R^{n_{i-1}}$ et ceux qui s'envoient sur zéro (et ces derniers
sont atteints par des vecteurs de la base de $R^{n_{i + 1}}$ si $i > 0$).
En raisonnant par
récurrence descendante à partir de $i = e$, il est facile de prouver qu'il y
a $r_{i + 1}$ vecteurs de la base de $R^{n_i}$ qui s'envoient
sur zéro, et $r_i$ qui s'envoient sur des vecteurs de la base de
$R^{n_{i-1}}$. Le résultat en découle.
\end{proof}

\begin{lemma}
\label{lemma-div-x-exact-one-less}
Dans la situation \ref{situation-complex}, supposons que $R$ soit
un anneau local d'idéal maximal $\mathfrak m$.
Supposons que $0 \to R^{n_e} \to R^{n_{e-1}} \to \ldots \to R^{n_0}$
soit exact en $R^{n_e}, \ldots, R^{n_1}$.
Soit $x \in \mathfrak m$ un non-diviseur de zéro. Le complexe
$0 \to (R/xR)^{n_e} \to \ldots \to (R/xR)^{n_1}$
est exact en $(R/xR)^{n_e}, \ldots, (R/xR)^{n_2}$.
\end{lemma}

\begin{proof}
Notons $F_\bullet$ le complexe de termes $F_i = R^{n_i}$
et dont la différentielle est donnée par $\varphi_i$. Nous avons alors une suite
exacte courte de complexes
$$
0 \to F_\bullet \xrightarrow{x} F_\bullet \to F_\bullet/xF_\bullet \to 0
$$
En appliquant le lemme du serpent, nous obtenons une suite exacte longue
$$
H_i(F_\bullet) \xrightarrow{x} H_i(F_\bullet) \to
H_i(F_\bullet/xF_\bullet) \to H_{i - 1}(F_\bullet)
\xrightarrow{x} H_{i - 1}(F_\bullet)
$$
Le lemme en découle.
\end{proof}

\begin{lemma}[Lemme d'acyclicité]
\label{lemma-acyclic}
\begin{reference}
\cite[Lemma 1.8]{Peskine-Szpiro}
\end{reference}
Soit $R$ un anneau local noethérien.
Soit $0 \to M_e \to M_{e-1} \to \ldots \to M_0$
un complexe de $R$-modules de type fini.
Supposons que $\text{depth}(M_i) \geq i$.
Soit $i$ le plus grand indice où le complexe n'est
pas exact en $M_i$. Si $i > 0$, alors
$\Ker(M_i \to M_{i-1})/\Im(M_{i + 1} \to M_i)$
est de profondeur $\geq 1$.
\end{lemma}

\begin{proof}
Soit $H = \Ker(M_i \to M_{i-1})/\Im(M_{i + 1} \to M_i)$ le
groupe d'homologie considéré.
Nous pouvons décomposer le complexe en suites exactes courtes
$0 \to M_e \to M_{e-1} \to K_{e-2} \to 0$,
$0 \to K_j \to M_j \to K_{j-1} \to 0$, pour $i + 2 \leq j \leq e-2 $,
$0 \to K_{i + 1} \to M_{i + 1} \to B_i \to 0$,
$0 \to K_i \to M_i \to M_{i-1}$, et
$0 \to B_i \to K_i \to H \to 0$.
Nous remontons ces complexes pour
prouver les assertions sur les profondeurs, en utilisant de façon répétée le
lemme \ref{lemma-depth-in-ses}.
Tout d'abord, puisque $\text{depth}(M_e) \geq e$
et que $\text{depth}(M_{e-1}) \geq e-1$, nous en déduisons
que $\text{depth}(K_{e-2}) \geq e - 1$. À ce stade, les
suites $0 \to K_j \to M_j \to K_{j-1} \to 0$ pour $i + 2 \leq j \leq e-2 $
impliquent de même que $\text{depth}(K_{j-1}) \geq j$ pour
$i + 2 \leq j \leq e-2$. La suite
$0 \to K_{i + 1} \to M_{i + 1} \to B_i \to 0$
montre alors que $\text{depth}(B_i) \geq i + 1$. La suite
$0 \to K_i \to M_i \to M_{i-1}$ montre que $\text{depth}(K_i) \geq 1$,
puisque $M_i$ est de profondeur $\geq i \geq 1$ par hypothèse.
La suite $0 \to B_i \to K_i \to H \to 0$
implique alors le résultat.
\end{proof}

\begin{proposition}
\label{proposition-what-exact}
\begin{reference}
\cite[Corollary 1]{WhatExact}
\end{reference}
Dans la situation \ref{situation-complex}, supposons que $R$ soit
un anneau local noethérien. Les assertions suivantes sont équivalentes :
\begin{enumerate}
\item $0 \to R^{n_e} \to R^{n_{e-1}} \to \ldots \to R^{n_0}$
est exact en $R^{n_e}, \ldots, R^{n_1}$, et
\item pour tout $i$, $1 \leq i \leq e$,
les deux conditions suivantes sont satisfaites :
\begin{enumerate}
\item $\text{rank}(\varphi_i) = r_i$, où
$r_i = n_i - n_{i + 1} + \ldots + (-1)^{e-i-1} n_{e-1} + (-1)^{e-i} n_e$,
\item $I(\varphi_i) = R$, ou bien $I(\varphi_i)$ contient une
suite régulière de longueur $i$.
\end{enumerate}
\end{enumerate}
\end{proposition}

\begin{proof}
Si, pour un certain $i$, un coefficient matriciel de $\varphi_i$
n'appartient pas à $\mathfrak m$, nous appliquons le lemme \ref{lemma-add-trivial-complex}.
Il est facile de voir que la proposition pour un complexe et
pour le même complexe auquel on a ajouté un complexe trivial
sont équivalentes. Nous pouvons donc supposer que tous les coefficients matriciels
de chaque $\varphi_i$ sont des éléments de l'idéal maximal.
Nous pouvons également supposer que $e \geq 1$.

\medskip\noindent
Supposons que le complexe soit exact en $R^{n_e}, \ldots, R^{n_1}$.
Soit $\mathfrak q \in \text{Ass}(R)$.
Notons que l'anneau $R_{\mathfrak q}$ est de profondeur $0$
et que le complexe reste exact après localisation en $\mathfrak q$.
Nous appliquons les lemmes \ref{lemma-exact-depth-zero-local} et
\ref{lemma-trivial-case-exact} au complexe localisé
sur $R_{\mathfrak q}$. Nous concluons que
$\varphi_{i, \mathfrak q}$ est de rang $r_i$ pour tout $i$.
Comme $R \to \bigoplus_{\mathfrak q \in \text{Ass}(R)} R_\mathfrak q$
est injectif (lemme \ref{lemma-zero-at-ass-zero}), nous concluons que
$\varphi_i$ est de rang $r_i$ sur $R$, par la définition du rang donnée
dans la définition \ref{definition-rank}. Nous voyons donc que
$I(\varphi_i)_\mathfrak q = I(\varphi_{i, \mathfrak q})$,
puisque les rangs ne changent pas. Comme tous les idéaux
$I(\varphi_i)_{\mathfrak q}$, $e \geq i \geq 1$,
sont égaux à $R_{\mathfrak q}$ (par les lemmes cités ci-dessus),
nous concluons qu'aucun des idéaux $I(\varphi_i)$ n'est contenu dans $\mathfrak q$.
Ceci implique que $I(\varphi_e)I(\varphi_{e-1})\ldots I(\varphi_1)$
n'est contenu dans aucun des idéaux premiers associés
à $R$. Par le lemme \ref{lemma-silly}, nous pouvons choisir
$x \in I(\varphi_e)I(\varphi_{e - 1})\ldots I(\varphi_1)$,
$x \not \in \mathfrak q$ pour tout $\mathfrak q \in \text{Ass}(R)$.
Observons que $x$ est un non-diviseur de zéro (lemme \ref{lemma-ass-zero-divisors}).
D'après le lemme \ref{lemma-div-x-exact-one-less},
le complexe $0 \to (R/xR)^{n_e} \to \ldots \to (R/xR)^{n_1}$ est exact
en $(R/xR)^{n_e}, \ldots, (R/xR)^{n_2}$. Par récurrence
sur $e$, tous les idéaux $I(\varphi_i)/xR$ possèdent une suite
régulière de longueur $i - 1$. Ceci prouve que $I(\varphi_i)$
contient une suite régulière de longueur $i$.

\medskip\noindent
Supposons que (2)(a) et (2)(b) soient satisfaites. Nous allons prouver (1) par
récurrence sur $\dim(R)$. Si $\dim(R) = 0$, alors nous devons avoir
$I(\varphi_i) = R$ pour $1 \leq i \leq e$ par (2)(b). Comme les
coefficients de $\varphi_i$ sont contenus dans l'idéal maximal,
cela ne peut se produire que si $r_i = 0$ pour tout $i$. Par (2)(a), nous
concluons que $e = 0$, et (1) est satisfaite. Supposons que $\dim(R) > 0$.
Affirmons que, pour tout idéal premier $\mathfrak p \subset R$, les conditions
(2)(a) et (2)(b) sont satisfaites pour le complexe
$0 \to R_\mathfrak p^{n_e} \to R_\mathfrak p^{n_{e - 1}} \to \ldots \to
R_\mathfrak p^{n_0}$, de morphismes $\varphi_{i, \mathfrak p}$,
sur $R_\mathfrak p$. En effet, comme $I(\varphi_i)$ contient un
non-diviseur de zéro, l'image de $I(\varphi_i)$ dans $R_\mathfrak p$
est non nulle. Ceci implique que le rang de $\varphi_{i, \mathfrak p}$
est le même que celui de $\varphi_i$ : le rang, tel que défini ci-dessus,
d'une matrice $\varphi$ sur un anneau $R$ ne peut que diminuer lors du passage
à une $R$-algèbre $R'$, et cela se produit si et seulement si $I(\varphi)$
s'envoie sur zéro dans $R'$. Ainsi, (2)(a) est satisfaite. Cela étant dit,
nous savons que $I(\varphi_{i, \mathfrak p}) = I(\varphi_i)_\mathfrak p$,
et nous voyons que (2)(b) est également préservée par localisation.
Par récurrence sur la dimension de $R$, nous pouvons supposer que le complexe
est exact après localisation en tout idéal premier non maximal $\mathfrak p$ de $R$.
Ainsi, $\Ker(\varphi_i)/\Im(\varphi_{i + 1})$ a son support contenu dans
$\{\mathfrak m\}$ et, s'il est non nul, est donc de profondeur $0$.
Comme $I(\varphi_i) \subset \mathfrak m$ pour tout $i$, en vertu
de ce qui a été dit dans le premier paragraphe de la démonstration, nous
voyons que (2)(b) implique $\text{depth}(R) \geq e$.
Par le lemme \ref{lemma-acyclic}, nous voyons
que le complexe est exact en $R^{n_e}, \ldots, R^{n_1}$,
ce qui achève la démonstration.
\end{proof}

\begin{remark}
\label{remark-what-exact}
Si, dans la proposition \ref{proposition-what-exact}, les conditions équivalentes (1)
et (2) sont satisfaites, alors il existe un $j$ tel que $I(\varphi_i) = R$
si et seulement si $i \geq j$. Comme dans la démonstration de la proposition, il suffit
de le voir lorsque toutes les matrices ont leurs coefficients dans l'idéal maximal
$\mathfrak m$ de $R$. Dans ce cas, nous voyons que $I(\varphi_j) = R$
si et seulement si $\varphi_j = 0$. Mais si $\varphi_j = 0$, alors nous obtenons
des complexes exacts arbitrairement longs
$0 \to R^{n_e} \to R^{n_{e-1}} \to \ldots
\to R^{n_j} \to 0 \to 0 \to \ldots \to 0$,
et la proposition montre donc que $I(\varphi_i)$, pour $i > j$,
doit être égal à $R$ (sinon, c'est un idéal propre d'un anneau local
noethérien contenant des suites régulières arbitrairement longues, ce qui est impossible).
\end{remark}











\section{Modules de Cohen-Macaulay}
\label{section-CM}

\noindent
Nous montrons ici que les modules de Cohen-Macaulay possèdent de bonnes propriétés. Nous remettons
à plus tard l'emploi des groupes Ext pour établir le lien avec la dualité, entre autres.

\begin{definition}
\label{definition-CM}
Soit $R$ un anneau local noethérien.
Soit $M$ un $R$-module de type fini.
Nous disons que $M$ est {\it de Cohen-Macaulay}
si $\dim(\text{Supp}(M)) = \text{depth}(M)$.
\end{definition}

\noindent
Un premier objectif sera d'établir la proposition \ref{proposition-CM-module}.
Nous y parvenons par une suite (peut-être non standard) de lemmes élémentaires
ne faisant presque appel à aucun des résultats antérieurs sur la profondeur. Introduisons
quelques notations.

\medskip\noindent
Soit $R$ un anneau local noethérien. Soient $M$ un module
de Cohen-Macaulay et $f_1, \ldots, f_d$
une suite $M$-régulière, avec $d = \dim(\text{Supp}(M))$.
Nous disons que $g \in \mathfrak m$ est {\it bon relativement à
$(M, f_1, \ldots, f_d)$} si, pour tout $i = 0, 1, \ldots, d-1$,
nous avons $\dim (\text{Supp}(M) \cap V(g, f_1, \ldots, f_i))
= d - i - 1$. Ceci équivaut à la condition
$\dim(\text{Supp}(M/(f_1, \ldots, f_i)M) \cap V(g)) =
d - i - 1$ pour $i = 0, 1, \ldots, d - 1$.

\begin{lemma}
\label{lemma-good-element}
Notations et hypothèses comme ci-dessus. Si $g$ est bon relativement à
$(M, f_1, \ldots, f_d)$, alors (a) $g$ est un non-diviseur de zéro sur $M$,
et (b) $M/gM$ est de Cohen-Macaulay, de suite régulière
maximale $f_1, \ldots, f_{d - 1}$.
\end{lemma}

\begin{proof}
Nous démontrons le lemme par récurrence sur $d$.
Si $d = 0$, alors $M$ est de type fini et il n'existe aucun cas
auquel le lemme s'applique.
Si $d = 1$, nous devons montrer que $g : M \to M$ est
injectif. Le noyau $K$ a pour support $\{\mathfrak m\}$,
car, par hypothèse, $\dim \text{Supp}(M) \cap V(g) = 0$.
Ainsi, $K$ est de longueur finie. L'injectivité de $f_1 : K \to K$
implique donc que son image a la même longueur que $K$, et par conséquent que
$f_1 K = K$, ce qui, par le lemme de Nakayama \ref{lemma-NAK}, implique $K = 0$.
De plus, $\dim \text{Supp}(M/gM) = 0$, si bien que $M/gM$ est de Cohen-Macaulay
et de profondeur $0$.

\medskip\noindent
Supposons que $d > 1$. Observons que $g$ est bon pour $(M/f_1M, f_2, \ldots, f_d)$,
comme on le voit aisément d'après la définition. Par récurrence, nous avons que
(a) $g$ est un non-diviseur de zéro sur $M/f_1M$ et que
(b) $M/(g, f_1)M$ est de Cohen-Macaulay, de suite régulière maximale
$f_2, \ldots, f_{d - 1}$. Par le
lemme \ref{lemma-permute-xi},
nous voyons que $g, f_1$ est une suite $M$-régulière.
Ainsi, $g$ est un non-diviseur de zéro sur $M$ et
$f_1, \ldots, f_{d - 1}$ est une suite $M/gM$-régulière.
\end{proof}

\begin{lemma}
\label{lemma-CM-one-g}
Soit $R$ un anneau local noethérien.
Soit $M$ un module de Cohen-Macaulay sur $R$.
Supposons que $g \in \mathfrak m$ vérifie $\dim(\text{Supp}(M) \cap V(g))
= \dim(\text{Supp}(M)) - 1$. Alors (a) $g$ est un non-diviseur de zéro sur $M$,
et (b) $M/gM$ est de Cohen-Macaulay, de profondeur diminuée d'une unité.
\end{lemma}

\begin{proof}
Choisissons une suite $M$-régulière $f_1, \ldots, f_d$, avec
$d = \dim(\text{Supp}(M))$. Si $g$ est bon relativement à
$(M, f_1, \ldots, f_d)$, le lemme \ref{lemma-good-element} conclut.
En particulier, le lemme est vrai si $d = 1$. (Le cas $d = 0$ ne
se produit pas.) Supposons que $d > 1$. Choisissons un élément $h \in R$ tel que
(\romannumeral1) $h$ soit bon relativement à $(M, f_1, \ldots, f_d)$,
et que (\romannumeral2) $\dim(\text{Supp}(M) \cap V(h, g)) = d - 2$.
Pour voir que $h$ existe, soit $\{\mathfrak q_j\}$ l'ensemble (fini) des
idéaux premiers minimaux des fermés $\text{Supp}(M)$,
$\text{Supp}(M)\cap V(f_1, \ldots, f_i)$, $i = 1, \ldots, d - 1$,
et $\text{Supp}(M) \cap V(g)$. Aucun de ces $\mathfrak q_j$
n'est égal à $\mathfrak m$ ; nous pouvons donc trouver $h \in \mathfrak m$,
$h \not \in \mathfrak q_j$, par le lemme \ref{lemma-silly}. Il est clair
que $h$ satisfait (\romannumeral1) et (\romannumeral2). Du
lemme \ref{lemma-good-element}, nous déduisons que
$M/hM$ est de Cohen-Macaulay. Par (\romannumeral2), nous voyons que la paire
$(M/hM, g)$ satisfait l'hypothèse de récurrence. Ainsi,
$M/(h, g)M$ est de Cohen-Macaulay et $g : M/hM \to M/hM$
est injectif. Par le lemme \ref{lemma-permute-xi}, nous voyons
que $g : M \to M$ et $h : M/gM \to M/gM$
sont injectifs. Combiné au fait que $M/(g, h)M$
est de Cohen-Macaulay, ceci achève la démonstration.
\end{proof}

\begin{proposition}
\label{proposition-CM-module}
Soit $R$ un anneau local noethérien, d'idéal maximal $\mathfrak m$.
Soit $M$ un module de Cohen-Macaulay sur $R$, dont le support est de dimension $d$.
Supposons que $g_1, \ldots, g_c$ soient des éléments de
$\mathfrak m$ tels que $\dim(\text{Supp}(M/(g_1, \ldots, g_c)M))
= d - c$. Alors $g_1, \ldots, g_c$ est une suite $M$-régulière,
et peut être prolongée en une suite $M$-régulière maximale.
\end{proposition}

\begin{proof}
Soit $Z = \text{Supp}(M) \subset \Spec(R)$.
Par le lemme \ref{lemma-one-equation}, dans la chaîne
$Z \supset Z \cap V(g_1) \supset \ldots \supset Z \cap V(g_1, \ldots, g_c)$,
chaque étape diminue la dimension d'au plus $1$. Par hypothèse,
chaque étape diminue donc la dimension d'exactement $1$. Nous pouvons ainsi
appliquer successivement le lemme \ref{lemma-CM-one-g} aux modules
$M/(g_1, \ldots, g_i)$ et à l'élément $g_{i + 1}$.

\medskip\noindent
Pour prolonger $g_1, \ldots, g_c$ d'un élément si $c < d$, il suffit de
choisir un élément $g_{c + 1} \in \mathfrak m$ qui n'appartient
à aucun des idéaux premiers minimaux, en nombre fini, de $Z \cap V(g_1, \ldots, g_c)$,
en utilisant le lemme \ref{lemma-silly}.
\end{proof}

\noindent
La proposition \ref{proposition-CM-module} étant démontrée, poursuivons le
développement de la théorie standard.

\begin{lemma}
\label{lemma-nonzerodivisor-on-CM}
Soit $R$ un anneau local noethérien d'idéal maximal $\mathfrak m$.
Soit $M$ un $R$-module de type fini. Soit $x \in \mathfrak m$ un
non-diviseur de zéro sur $M$. Alors $M$ est de Cohen-Macaulay si et seulement
si $M/xM$ est de Cohen-Macaulay.
\end{lemma}

\begin{proof}
Par le lemme \ref{lemma-depth-drops-by-one}, nous avons
$\text{depth}(M/xM) = \text{depth}(M)-1$.
Par le lemme \ref{lemma-one-equation-module},
nous avons $\dim(\text{Supp}(M/xM)) = \dim(\text{Supp}(M)) - 1$.
\end{proof}

\begin{lemma}
\label{lemma-CM-over-quotient}
Soit $R \to S$ un homomorphisme surjectif d'anneaux locaux noethériens.
Soit $N$ un $S$-module de type fini. Alors $N$ est de Cohen-Macaulay comme $S$-module
si et seulement si $N$ est de Cohen-Macaulay comme $R$-module.
\end{lemma}

\begin{proof}
Démonstration omise.
\end{proof}

\begin{lemma}
\label{lemma-CM-ass-minimal-support}
\begin{reference}
\cite[Chapter 0, Proposition 16.5.4]{EGA}
\end{reference}
Soit $R$ un anneau local noethérien. Soit $M$ un $R$-module de Cohen-Macaulay
de type fini. Si $\mathfrak p \in \text{Ass}(M)$, alors
$\dim(R/\mathfrak p) = \dim(\text{Supp}(M))$ et $\mathfrak p$
est un idéal premier minimal du support de $M$.
En particulier, $M$ n'a aucun idéal premier associé immergé.
\end{lemma}

\begin{proof}
Par le lemme \ref{lemma-depth-dim-associated-primes}, nous avons
$\text{depth}(M) \leq \dim(R/\mathfrak p)$.
Bien sûr, $\dim(R/\mathfrak p) \leq \dim(\text{Supp}(M))$,
car $\mathfrak p \in \text{Supp}(M)$ (lemme \ref{lemma-ass-support}).
Nous avons donc égalité dans les deux inégalités, puisque $M$ est de Cohen-Macaulay.
Alors $\mathfrak p$ doit être minimal dans $\text{Supp}(M)$, sans quoi
nous aurions $\dim(R/\mathfrak p) < \dim(\text{Supp}(M))$.
Enfin, les idéaux premiers minimaux du support de $M$ sont égaux
aux éléments minimaux de $\text{Ass}(M)$
(proposition \ref{proposition-minimal-primes-associated-primes}) ;
ainsi, $M$ n'a aucun idéal premier associé immergé
(définition \ref{definition-embedded-primes}).
\end{proof}

\begin{definition}
\label{definition-maximal-CM}
Soit $R$ un anneau local noethérien.
Un module $M$ de type fini sur $R$ est appelé {\it module de Cohen-Macaulay
maximal} si $\text{depth}(M) = \dim(R)$.
\end{definition}

\noindent
Autrement dit, un module de Cohen-Macaulay maximal sur un anneau local
noethérien est un module de type fini ayant la plus grande profondeur possible sur cet anneau.
De manière équivalente, un module de Cohen-Macaulay maximal sur un anneau local
noethérien $R$ est un module de Cohen-Macaulay dont la dimension est égale à celle
de l'anneau. En particulier, si $M$ est un $R$-module de Cohen-Macaulay tel que
$\Spec(R) = \text{Supp}(M)$, alors $M$ est de Cohen-Macaulay maximal.
Les deux lemmes suivants portent donc sur les modules de Cohen-Macaulay maximaux.

\begin{lemma}
\label{lemma-maximal-chain-maximal-CM}
\begin{slogan}
Dans un anneau local de Cohen-Macaulay, toute chaîne maximale d'idéaux premiers est de
longueur égale à la dimension.
\end{slogan}
Soit $R$ un anneau local noethérien. Supposons qu'il existe un
module de Cohen-Macaulay $M$ tel que $\Spec(R) = \text{Supp}(M)$.
Alors toute chaîne maximale d'idéaux premiers $\mathfrak p_0 \subset
\mathfrak p_1 \subset \ldots \subset \mathfrak p_n$
est de longueur $n = \dim(R)$.
\end{lemma}

\begin{proof}
Nous allons le démontrer par récurrence sur $\dim(R)$. Si $\dim(R) = 0$,
l'assertion est claire. Supposons que $\dim(R) > 0$. Alors $n > 0$.
Choisissons un élément $x \in \mathfrak p_1$, avec $x$ n'appartenant
à aucun des idéaux premiers minimaux de $R$, et en particulier
$x \not \in \mathfrak p_0$. (Voir le lemme \ref{lemma-silly}.)
Alors $\dim(R/xR) = \dim(R) - 1$ par le lemme \ref{lemma-one-equation}.
Le module $M/xM$ est de Cohen-Macaulay sur $R/xR$ par la
proposition \ref{proposition-CM-module} et le
lemme \ref{lemma-CM-over-quotient}.
Le support de $M/xM$ est $\Spec(R/xR)$ par le
lemme \ref{lemma-support-quotient}.
Après avoir remplacé $x$ par $x^n$ pour un certain $n$,
nous pouvons supposer que $\mathfrak p_1$ est un idéal premier associé à $M/xM$, voir le
lemme \ref{lemma-inherit-minimal-primes}.
Par le lemme \ref{lemma-CM-ass-minimal-support},
nous concluons que $\mathfrak p_1/(x)$ est un idéal premier minimal de $R/xR$.
Il s'ensuit que la chaîne
$\mathfrak p_1/(x) \subset \ldots \subset \mathfrak p_n/(x)$
est une chaîne maximale d'idéaux premiers de $R/xR$.
Par récurrence, nous trouvons que cette chaîne
est de longueur $\dim(R/xR) = \dim(R) - 1$, comme souhaité.
\end{proof}

\begin{lemma}
\label{lemma-dim-formula-maximal-CM}
Supposons que $R$ soit un anneau local noethérien. Supposons qu'il existe un
module de Cohen-Macaulay $M$ tel que $\Spec(R) = \text{Supp}(M)$. Alors, pour
un idéal premier $\mathfrak p \subset R$, nous avons
$$
\dim(R) = \dim(R_{\mathfrak p}) + \dim(R/\mathfrak p).
$$
\end{lemma}

\begin{proof}
Cela découle immédiatement du lemme \ref{lemma-maximal-chain-maximal-CM}.
\end{proof}

\begin{lemma}
\label{lemma-localize-CM-module}
Supposons que $R$ soit un anneau local noethérien. Soit $M$ un module de Cohen-Macaulay
sur $R$. Pour tout idéal premier $\mathfrak p \subset R$, le
module $M_{\mathfrak p}$ est de Cohen-Macaulay sur $R_\mathfrak p$.
\end{lemma}

\begin{proof}
Nous pouvons supposer, et supposons, que $\mathfrak p \not = \mathfrak m$ et que $M$ est non nul.
Choisissons une chaîne maximale d'idéaux premiers $\mathfrak p = \mathfrak p_c \subset
\mathfrak p_{c - 1} \subset \ldots \subset \mathfrak p_1 \subset \mathfrak m$.
Si nous prouvons le résultat pour $M_{\mathfrak p_1}$ sur $R_{\mathfrak p_1}$,
alors le lemme découlera d'une récurrence sur $c$. Nous pouvons donc supposer qu'il
n'existe aucun idéal premier strictement compris entre $\mathfrak p$ et $\mathfrak m$.
Notons que
$\dim(\text{Supp}(M_\mathfrak p)) \leq \dim(\text{Supp}(M)) - 1$,
car toute chaîne d'idéaux premiers du support de $M_\mathfrak p$
peut être prolongée d'un idéal premier supplémentaire (à savoir $\mathfrak m$) dans le
support de $M$. D'autre part, nous avons
$\text{depth}(M_\mathfrak p) \geq \text{depth}(M) - \dim(R/\mathfrak p) =
\text{depth}(M) - 1$ par le lemme \ref{lemma-depth-localization} et notre
choix de $\mathfrak p$.
Ainsi, $\text{depth}(M_\mathfrak p) \geq \dim(\text{Supp}(M_\mathfrak p))$,
comme souhaité (l'autre inégalité est le lemme \ref{lemma-bound-depth}).
\end{proof}

\begin{definition}
\label{definition-module-CM}
Soit $R$ un anneau noethérien. Soit $M$ un $R$-module de type fini.
Nous disons que $M$ est {\it de Cohen-Macaulay} si $M_\mathfrak p$ est un module
de Cohen-Macaulay sur $R_\mathfrak p$ pour tout idéal premier $\mathfrak p$ de $R$.
\end{definition}

\noindent
Par le lemme \ref{lemma-localize-CM-module}, il suffit de le vérifier
aux idéaux maximaux de $R$.

\begin{lemma}
\label{lemma-maximal-CM-polynomial-algebra}
Soit $R$ un anneau noethérien. Soit $M$ un module de Cohen-Macaulay
sur $R$. Alors $M \otimes_R R[x_1, \ldots, x_n]$ est un module
de Cohen-Macaulay sur $R[x_1, \ldots, x_n]$.
\end{lemma}

\begin{proof}
Par récurrence sur le nombre de variables, il suffit de le démontrer pour
$M[x] = M \otimes_R R[x]$ sur $R[x]$. Soit $\mathfrak m \subset R[x]$
un idéal maximal, et soit $\mathfrak p = R \cap \mathfrak m$.
Soit $f_1, \ldots, f_d$ une suite $M_\mathfrak p$-régulière dans l'idéal
maximal de $R_{\mathfrak p}$, de longueur $d = \dim(\text{Supp}(M_{\mathfrak p}))$.
Notons que, puisque $R[x]$ est plat sur $R$, le localisé
$R[x]_{\mathfrak m}$ est plat sur $R_{\mathfrak p}$.
Ainsi, par le lemme \ref{lemma-flat-increases-depth}, la suite
$f_1, \ldots, f_d$ est une suite $M[x]_{\mathfrak m}$-régulière de longueur $d$
dans $R[x]_{\mathfrak m}$. Le quotient
$$
Q = M[x]_{\mathfrak m}/(f_1, \ldots, f_d)M[x]_{\mathfrak m} =
M_{\mathfrak p}/(f_1, \ldots, f_d)M_{\mathfrak p}
\otimes_{R_\mathfrak p} R[x]_{\mathfrak m}
$$
a pour support les idéaux premiers situés au-dessus de $\mathfrak p$,
car $R_\mathfrak p \to R[x]_\mathfrak m$ est plat et
le support de $M_{\mathfrak p}/(f_1, \ldots, f_d)M_{\mathfrak p}$
est égal à $\{\mathfrak p\}$ (détails omis ; indication : cela découle des
lemmes \ref{lemma-annihilator-flat-base-change} et
\ref{lemma-support-closed}). La dimension vaut donc $1$.
Pour achever la démonstration, il suffit de trouver un $f \in \mathfrak m$
qui soit un non-diviseur de zéro sur $Q$. Comme $\mathfrak m$ est
un idéal maximal, l'extension de corps
$\kappa(\mathfrak m)/\kappa(\mathfrak p)$
est finie (théorème \ref{theorem-nullstellensatz}).
Nous pouvons donc trouver $f \in \mathfrak m$ qui,
vu comme un polynôme en $x$, a son coefficient
dominant hors de $\mathfrak p$. Un tel $f$ agit comme
un non-diviseur de zéro sur
$$
M_{\mathfrak p}/(f_1, \ldots, f_d)M_{\mathfrak p} \otimes_R R[x] =
\bigoplus\nolimits_{n \geq 0}
M_{\mathfrak p}/(f_1, \ldots, f_d)M_{\mathfrak p} \cdot x^n
$$
et agit donc comme un non-diviseur de zéro sur $Q$.
\end{proof}















\section{Anneaux de Cohen-Macaulay}
\label{section-CM-ring}

\noindent
La plupart des résultats de cette section sont des cas particuliers de ceux
de la section \ref{section-CM}.

\begin{definition}
\label{definition-local-ring-CM}
Un anneau local noethérien $R$ est dit {\it de Cohen-Macaulay}
s'il est de Cohen-Macaulay comme module sur lui-même.
\end{definition}

\noindent
Notons que cela équivaut à demander l'existence
d'une suite $R$-régulière $x_1, \ldots, x_d$ de l'idéal
maximal telle que $R/(x_1, \ldots, x_d)$ soit de dimension $0$.
Nous dirons généralement simplement « suite régulière », et non
« suite $R$-régulière ».

\begin{lemma}
\label{lemma-reformulate-CM}
\begin{slogan}
Dans les anneaux locaux de Cohen-Macaulay, les suites régulières sont caractérisées
par le fait de découper un objet de la bonne dimension.
\end{slogan}
Soit $R$ un anneau local noethérien de Cohen-Macaulay, d'idéal
maximal $\mathfrak m $. Soient $x_1, \ldots, x_c \in \mathfrak m$ des
éléments. Alors
$$
x_1, \ldots, x_c
\text{ est une suite régulière }
\Leftrightarrow
\dim(R/(x_1, \ldots, x_c)) = \dim(R) - c
$$
Dans ce cas,
$x_1, \ldots, x_c$ peut être prolongée en
une suite régulière de longueur $\dim(R)$, et chaque quotient
$R/(x_1, \ldots, x_i)$ est un anneau de Cohen-Macaulay de dimension
$\dim(R) - i$.
\end{lemma}

\begin{proof}
C'est un cas particulier de la proposition \ref{proposition-CM-module}.
\end{proof}

\begin{lemma}
\label{lemma-maximal-chain-CM}
Soit $R$ un anneau local noethérien.
Supposons que $R$ soit de Cohen-Macaulay de dimension $d$.
Toute chaîne maximale d'idéaux premiers $\mathfrak p_0 \subset
\mathfrak p_1 \subset \ldots \subset \mathfrak p_n$
est de longueur $n = d$.
\end{lemma}

\begin{proof}
C'est un cas particulier du lemme \ref{lemma-maximal-chain-maximal-CM}.
\end{proof}

\begin{lemma}
\label{lemma-CM-dim-formula}
Supposons que $R$ soit un anneau local noethérien de Cohen-Macaulay de dimension $d$.
Pour tout idéal premier $\mathfrak p \subset R$, nous avons
$$
\dim(R) = \dim(R_{\mathfrak p}) + \dim(R/\mathfrak p).
$$
\end{lemma}

\begin{proof}
Cela découle immédiatement du lemme \ref{lemma-maximal-chain-CM}.
(C'est aussi un cas particulier du lemme \ref{lemma-dim-formula-maximal-CM}.)
\end{proof}

\begin{lemma}
\label{lemma-localize-CM}
Supposons que $R$ soit un anneau local de Cohen-Macaulay.
Pour tout idéal premier $\mathfrak p \subset R$, l'anneau
$R_{\mathfrak p}$ est lui aussi de Cohen-Macaulay.
\end{lemma}

\begin{proof}
C'est un cas particulier du lemme \ref{lemma-localize-CM-module}.
\end{proof}

\begin{definition}
\label{definition-ring-CM}
Un anneau noethérien $R$ est dit {\it de Cohen-Macaulay} si tous
ses anneaux locaux sont de Cohen-Macaulay.
\end{definition}

\begin{lemma}
\label{lemma-CM-polynomial-algebra}
Supposons que $R$ soit un anneau noethérien de Cohen-Macaulay.
Toute algèbre polynomiale sur $R$ est de Cohen-Macaulay.
\end{lemma}

\begin{proof}
C'est un cas particulier du lemme \ref{lemma-maximal-CM-polynomial-algebra}.
\end{proof}

\begin{lemma}
\label{lemma-dimension-shift}
Soit $R$ un anneau local noethérien de Cohen-Macaulay de dimension $d$.
Soit $0 \to K \to R^{\oplus n} \to M \to 0$ une suite exacte de
$R$-modules. Alors soit $M = 0$, soit $\text{depth}(K) > \text{depth}(M)$, soit
$\text{depth}(K) = \text{depth}(M) = d$.
\end{lemma}

\begin{proof}
C'est un cas particulier du lemme \ref{lemma-depth-in-ses}.
\end{proof}

\begin{lemma}
\label{lemma-mcm-resolution}
Soit $R$ un anneau local noethérien de Cohen-Macaulay de dimension $d$.
Soit $M$ un $R$-module de type fini et de profondeur $e$.
Il existe un complexe exact
$$
0 \to K \to F_{d-e-1} \to \ldots \to F_0 \to M \to 0
$$
où chaque $F_i$ est libre de rang fini et $K$ est de Cohen-Macaulay maximal.
\end{lemma}

\begin{proof}
Cela découle immédiatement de la définition et du lemme \ref{lemma-dimension-shift}.
\end{proof}

\begin{lemma}
\label{lemma-find-sequence-image-regular}
Soit $\varphi : A \to B$ un homomorphisme d'anneaux locaux.
Supposons que $B$ soit noethérien et de Cohen-Macaulay, et que
$\mathfrak m_B = \sqrt{\varphi(\mathfrak m_A) B}$. Alors il existe
une suite d'éléments $f_1, \ldots, f_{\dim(B)}$ de $A$
telle que $\varphi(f_1), \ldots, \varphi(f_{\dim(B)})$ soit une
suite régulière dans $B$.
\end{lemma}

\begin{proof}
Par récurrence sur $\dim(B)$, il suffit de prouver ceci : si $\dim(B) \geq 1$, alors nous
pouvons trouver un élément $f$ de $A$ qui s'envoie sur un non-diviseur de zéro dans $B$.
Par le
lemme \ref{lemma-reformulate-CM},
il suffit de trouver $f \in A$ dont l'image dans $B$ n'est contenue dans aucun
des idéaux premiers minimaux, en nombre fini, $\mathfrak q_1, \ldots, \mathfrak q_r$
de $B$. L'hypothèse
$\mathfrak m_B = \sqrt{\varphi(\mathfrak m_A) B}$
montre que $\mathfrak m_A \not \subset \varphi^{-1}(\mathfrak q_i)$.
Nous pouvons donc trouver $f$ grâce au
lemme \ref{lemma-silly}.
\end{proof}








\section{Anneaux caténaires}
\label{section-catenary}

\noindent
Comparer avec Topologie, section \ref{topology-section-catenary-spaces}.

\begin{definition}
\label{definition-catenary}
On dit qu'un anneau $R$ est {\it caténaire} si, pour toute paire d'idéaux premiers
$\mathfrak p \subset \mathfrak q$, il existe un entier qui majore les
longueurs de toutes les chaînes finies d'idéaux premiers
$\mathfrak p = \mathfrak p_0 \subset \mathfrak p_1 \subset \ldots \subset
\mathfrak p_e = \mathfrak q$ et si toutes les chaînes maximales de ce type ont la même longueur.
\end{definition}

\begin{lemma}
\label{lemma-catenary}
Un anneau $R$ est caténaire si et seulement si l'espace topologique
$\Spec(R)$ est caténaire (voir
Topologie, définition \ref{topology-definition-catenary}).
\end{lemma}

\begin{proof}
Cela résulte immédiatement de la définition et de la caractérisation des
fermés irréductibles dans le lemme \ref{lemma-irreducible}.
\end{proof}

\noindent
En général, il n'est pas vrai qu'une $R$-algèbre de type fini
soit caténaire lorsque $R$ l'est. Nous posons donc la
définition suivante.

\begin{definition}
\label{definition-universally-catenary}
Un anneau noethérien $R$ est dit {\it universellement caténaire}
si toute $R$-algèbre de type fini est caténaire.
\end{definition}

\noindent
Nous nous restreignons aux anneaux noethériens, car il n'est pas clair que
cette définition soit la bonne pour les anneaux non noethériens.
D'après le lemme \ref{lemma-quotient-catenary}, pour vérifier qu'un anneau
noethérien $R$ est universellement caténaire, il suffit de vérifier que
toute algèbre polynomiale $R[x_1, \ldots, x_n]$ est caténaire.

\begin{lemma}
\label{lemma-localization-catenary}
Toute localisation d'un anneau caténaire est caténaire.
Toute localisation d'un anneau noethérien universellement caténaire
est universellement caténaire.
\end{lemma}

\begin{proof}
Soient $A$ un anneau et $S \subset A$ une partie multiplicative.
La description de $\Spec(S^{-1}A)$ dans le lemme \ref{lemma-spec-localization}
montre que, si $A$ est caténaire, alors $S^{-1}A$ l'est aussi. Si $S^{-1}A \to C$
est de type fini, alors $C = S^{-1}B$ pour un homomorphisme d'anneaux de type fini
$A \to B$. Ainsi, si $A$ est noethérien et universellement caténaire, alors
$B$ est caténaire et nous voyons que $C$ est également caténaire. Cela prouve le lemme.
\end{proof}

\begin{lemma}
\label{lemma-universally-catenary}
Soit $A$ un anneau noethérien universellement caténaire.
Toute $A$-algèbre essentiellement de type fini sur $A$
est universellement caténaire.
\end{lemma}

\begin{proof}
Si $B$ est une $A$-algèbre de type fini, alors $B$ est noethérien
d'après le lemme \ref{lemma-Noetherian-permanence}. Toute
$B$-algèbre de type fini est une $A$-algèbre de type fini et est donc caténaire
puisque, par hypothèse, $A$ est universellement caténaire. Ainsi, $B$
est universellement caténaire. Toute localisation de $B$ est universellement
caténaire d'après le lemme \ref{lemma-localization-catenary}, ce qui
achève la preuve.
\end{proof}

\begin{lemma}
\label{lemma-catenary-check-local}
Soit $R$ un anneau. Les assertions suivantes sont équivalentes :
\begin{enumerate}
\item $R$ est caténaire,
\item $R_\mathfrak p$ est caténaire pour tout idéal premier $\mathfrak p$,
\item $R_\mathfrak m$ est caténaire pour tout idéal maximal $\mathfrak m$.
\end{enumerate}
Supposons $R$ noethérien. Les assertions suivantes sont équivalentes :
\begin{enumerate}
\item $R$ est universellement caténaire,
\item $R_\mathfrak p$ est universellement caténaire pour tout idéal premier
$\mathfrak p$,
\item $R_\mathfrak m$ est universellement caténaire pour tout idéal maximal
$\mathfrak m$.
\end{enumerate}
\end{lemma}

\begin{proof}
L'implication (1) $\Rightarrow$ (2) découle du
lemme \ref{lemma-localization-catenary} dans les deux cas.
L'implication (2) $\Rightarrow$ (3) est immédiate dans les deux cas.
Supposons que $R_\mathfrak m$ soit caténaire pour tout idéal maximal
$\mathfrak m$ de $R$. Si $\mathfrak p \subset \mathfrak q$ sont des idéaux premiers
de $R$, choisissons un idéal maximal vérifiant $\mathfrak q \subset \mathfrak m$.
Les chaînes d'idéaux premiers entre $\mathfrak p$ et $\mathfrak q$ sont
en correspondance bijective avec les chaînes d'idéaux premiers entre
$\mathfrak pR_\mathfrak m$ et $\mathfrak qR_\mathfrak m$ ; nous voyons donc que
$R$ est caténaire. Supposons $R$ noethérien et $R_\mathfrak m$
universellement caténaire pour tout idéal maximal $\mathfrak m$ de $R$.
Soit $R \to S$ un homomorphisme d'anneaux de type fini. Soit $\mathfrak q$ un idéal premier
de $S$ au-dessus de l'idéal premier $\mathfrak p \subset R$.
Choisissons un idéal maximal vérifiant $\mathfrak p \subset \mathfrak m$ dans $R$.
Alors $R_\mathfrak p$ est une localisation de $R_\mathfrak m$, donc est
universellement caténaire d'après le lemme \ref{lemma-localization-catenary}.
Ainsi, $S_\mathfrak p$ est caténaire comme anneau de type fini sur $R_\mathfrak p$.
Par conséquent, $S_\mathfrak q$ est caténaire comme localisation. Ainsi, $S$ est caténaire
d'après le premier cas traité ci-dessus.
\end{proof}

\begin{lemma}
\label{lemma-quotient-catenary}
Tout quotient d'un anneau caténaire est caténaire.
Tout quotient d'un anneau noethérien universellement caténaire est
universellement caténaire.
\end{lemma}

\begin{proof}
Soient $A$ un anneau et $I \subset A$ un idéal.
La description de $\Spec(A/I)$ dans le lemme \ref{lemma-spec-closed}
montre que, si $A$ est caténaire, alors $A/I$ l'est aussi.
La seconde assertion est un cas particulier du
lemme \ref{lemma-universally-catenary}.
\end{proof}

\begin{lemma}
\label{lemma-catenary-check-irreducible}
Soit $R$ un anneau noethérien.
\begin{enumerate}
\item $R$ est caténaire si et seulement si $R/\mathfrak p$ est caténaire
pour tout idéal premier minimal $\mathfrak p$.
\item $R$ est universellement caténaire si et seulement si $R/\mathfrak p$ est
universellement caténaire pour tout idéal premier minimal $\mathfrak p$.
\end{enumerate}
\end{lemma}

\begin{proof}
Si $\mathfrak a \subset \mathfrak b$ est une inclusion d'idéaux premiers de $R$,
alors nous pouvons trouver un idéal premier minimal $\mathfrak p \subset \mathfrak a$,
et la première assertion est claire. Nous omettons la preuve de la seconde.
\end{proof}

\begin{lemma}
\label{lemma-CM-ring-catenary}
Un anneau de Cohen-Macaulay noethérien est universellement caténaire.
Plus généralement, si $R$ est un anneau noethérien et $M$ un
$R$-module de Cohen-Macaulay tel que $\text{Supp}(M) = \Spec(R)$,
alors $R$ est universellement caténaire.
\end{lemma}

\begin{proof}
Puisqu'une algèbre polynomiale sur $R$ est de Cohen-Macaulay
d'après le lemme \ref{lemma-CM-polynomial-algebra},
il suffit de montrer qu'un anneau de Cohen-Macaulay est
caténaire.
Soit $R$ un anneau de Cohen-Macaulay et soient $\mathfrak p \subset \mathfrak q$
des idéaux premiers de $R$. Par définition, $R_{\mathfrak q}$ et
$R_{\mathfrak p}$ sont de Cohen-Macaulay.
Prenons une chaîne maximale d'idéaux premiers
$\mathfrak p = \mathfrak p_0 \subset \mathfrak p_1 \subset
\ldots \subset \mathfrak p_n = \mathfrak q$.
Choisissons ensuite une chaîne maximale d'idéaux premiers
$\mathfrak q_0 \subset \mathfrak q_1 \subset \ldots \subset
\mathfrak q_m = \mathfrak p$. D'après le
lemme \ref{lemma-maximal-chain-CM},
nous avons $n + m = \dim(R_{\mathfrak q})$. Et nous avons
$m = \dim(R_{\mathfrak p})$ d'après le même lemme.
Ainsi, $n = \dim(R_{\mathfrak q}) - \dim(R_{\mathfrak p})$
est indépendant des choix.

\medskip\noindent
Pour prouver l'assertion plus générale, raisonnons exactement comme ci-dessus, mais
en utilisant les lemmes \ref{lemma-maximal-CM-polynomial-algebra}
et \ref{lemma-maximal-chain-maximal-CM}.
\end{proof}

\begin{lemma}
\label{lemma-catenary-Noetherian-local}
Soit $(A, \mathfrak m)$ un anneau local noethérien. Les assertions suivantes sont équivalentes :
\begin{enumerate}
\item $A$ est caténaire, et
\item $\mathfrak p \mapsto \dim(A/\mathfrak p)$ est une fonction de dimension
sur $\Spec(A)$.
\end{enumerate}
\end{lemma}

\begin{proof}
Si $A$ est caténaire, alors $\Spec(A)$ possède une fonction de dimension $\delta$ d'après
Topologie, lemme \ref{topology-lemma-locally-dimension-function}
(et le lemme \ref{lemma-catenary}). Nous pouvons supposer que $\delta(\mathfrak m) = 0$.
Nous voyons alors que
$$
\delta(\mathfrak p) = \text{codim}(V(\mathfrak m), V(\mathfrak p)) =
\dim(A/\mathfrak p)
$$
d'après Topologie, lemme \ref{topology-lemma-dimension-function-catenary}.
Nous voyons ainsi que (1) implique (2). L'implication réciproque
découle de
Topologie, lemme \ref{topology-lemma-dimension-function-catenary}
également.
\end{proof}














\section{Anneaux locaux réguliers}
\label{section-regular}

\noindent
Les anneaux locaux réguliers sont définis dans la définition \ref{definition-regular-local}.
Il n'est pas si facile de montrer que toutes les localisations premières d'un anneau local
régulier sont régulières. En fait, une bonne partie des résultats développés jusqu'ici
vise à prouver ce fait. Voir la
proposition \ref{proposition-finite-gl-dim-regular}, et
remonter la chaîne des références.

\begin{lemma}
\label{lemma-regular-graded}
Soit $(R, \mathfrak m, \kappa)$ un anneau local régulier de dimension $d$.
L'anneau gradué $\bigoplus \mathfrak m^n / \mathfrak m^{n + 1}$
est isomorphe à l'algèbre polynomiale graduée
$\kappa[X_1, \ldots, X_d]$.
\end{lemma}

\begin{proof}
Soit $x_1, \ldots, x_d$ un système minimal de générateurs
de l'idéal maximal $\mathfrak m$, voir la
définition \ref{definition-regular-local}.
Il existe une surjection $\kappa[X_1, \ldots, X_d]
\to \bigoplus \mathfrak m^n/\mathfrak m^{n + 1}$,
qui envoie $X_i$ sur la classe de $x_i$ dans $\mathfrak m/\mathfrak m^2$.
Puisque $d(R) = d$ d'après la
proposition \ref{proposition-dimension},
nous savons que le polynôme
numérique $n \mapsto \dim_\kappa \mathfrak m^n/\mathfrak m^{n + 1}$
est de degré $d - 1$. D'après le lemme \ref{lemma-quotient-smaller-d}, nous
concluons que la surjection $\kappa[X_1, \ldots, X_d]
\to \bigoplus \mathfrak m^n/\mathfrak m^{n + 1}$ est un isomorphisme.
\end{proof}

\begin{lemma}
\label{lemma-regular-domain}
Tout anneau local régulier est intègre.
\end{lemma}

\begin{proof}
Nous allons utiliser l'égalité $\bigcap \mathfrak m^n = 0$
donnée par le lemme \ref{lemma-intersect-powers-ideal-module-zero}.
Soient $f, g \in R$ tels que $fg = 0$.
Supposons que $f \in \mathfrak m^a$ et
$g \in \mathfrak m^b$, avec $a, b$ maximaux.
Puisque $fg = 0 \in \mathfrak m^{a + b + 1}$,
le résultat du lemme \ref{lemma-regular-graded} montre
que soit $f \in \mathfrak m^{a + 1}$, soit
$g \in \mathfrak m^{b + 1}$. Contradiction.
\end{proof}

\begin{lemma}
\label{lemma-regular-ring-CM}
Soit $R$ un anneau local régulier et soit
$x_1, \ldots, x_d$ un système minimal de générateurs
de l'idéal maximal $\mathfrak m$. Alors
$x_1, \ldots, x_d$ est une suite régulière, et
chaque $R/(x_1, \ldots, x_c)$ est un anneau local régulier
de dimension $d - c$. En particulier, $R$ est de Cohen-Macaulay.
\end{lemma}

\begin{proof}
Remarquons que $R/x_1R$ est un anneau local noethérien de dimension $\geq d - 1$
d'après le lemme \ref{lemma-one-equation}, tandis que $x_2, \ldots, x_d$
engendrent l'idéal maximal. Il s'agit donc d'un anneau local régulier par définition.
Puisque $R$ est intègre d'après le lemme \ref{lemma-regular-domain},
$x_1$ est un non-diviseur de zéro.
\end{proof}

\begin{lemma}
\label{lemma-regular-quotient-regular}
Soit $R$ un anneau local régulier. Soit $I \subset R$ un idéal
tel que $R/I$ soit lui aussi un anneau local régulier. Alors
il existe un système minimal de générateurs $x_1, \ldots, x_d$
de l'idéal maximal $\mathfrak m$ de $R$ tel que
$I = (x_1, \ldots, x_c)$ pour un certain $0 \leq c \leq d$.
\end{lemma}

\begin{proof}
Écrivons $\dim(R) = d$ et $\dim(R/I) = d - c$.
Notons $\overline{\mathfrak m} = \mathfrak m/I$ l'idéal
maximal de $R/I$. Soit $\kappa = R/\mathfrak m$. Nous avons
$$
\dim_\kappa((I + \mathfrak m^2)/\mathfrak m^2) =
\dim_\kappa(\mathfrak m/\mathfrak m^2)
- \dim(\overline{\mathfrak m}/\overline{\mathfrak m}^2) = d - (d - c) = c
$$
par définition d'un anneau local régulier. Nous pouvons donc choisir
$x_1, \ldots, x_c \in I$ dont les images dans $\mathfrak m/\mathfrak m^2$
sont linéairement indépendantes, et les compléter par
$x_{c + 1}, \ldots, x_d$ pour obtenir un système minimal de générateurs de
$\mathfrak m$. Le morphisme induit $R/(x_1, \ldots, x_c) \to R/I$ est une
surjection entre anneaux locaux réguliers de même dimension
(lemme \ref{lemma-regular-ring-CM}). Il s'ensuit que
le noyau est nul, c'est-à-dire que $I = (x_1, \ldots, x_c)$. En effet, sinon,
nous aurions $\dim(R/I) < \dim(R/(x_1, \ldots, x_c))$ d'après les
lemmes \ref{lemma-regular-domain} et \ref{lemma-one-equation}.
\end{proof}

\begin{lemma}
\label{lemma-free-mod-x}
Soit $R$ un anneau local noethérien.
Soit $x \in \mathfrak m$.
Soit $M$ un $R$-module fini tel que
$x$ soit un non-diviseur de zéro sur $M$ et que
$M/xM$ soit libre sur $R/xR$.
Alors $M$ est libre sur $R$.
\end{lemma}

\begin{proof}
Soient $m_1, \ldots, m_r$ des éléments de $M$ qui s'envoient sur
une $R/xR$-base de $M/xM$. D'après le lemme de Nakayama \ref{lemma-NAK},
$m_1, \ldots, m_r$ engendrent $M$. Si $\sum a_i m_i = 0$
est une relation, alors $a_i \in xR$ pour tout $i$. Ainsi,
$a_i = b_i x$ pour certains $b_i \in R$. Par conséquent,
le noyau $K$ de $R^r \to M$ vérifie $xK = K$
et est donc nul d'après le lemme de Nakayama.
\end{proof}

\begin{lemma}
\label{lemma-regular-mcm-free}
Soit $R$ un anneau local régulier.
Tout module de Cohen-Macaulay maximal sur $R$ est libre.
\end{lemma}

\begin{proof}
Soit $M$ un module de Cohen-Macaulay maximal sur $R$.
Soit $x \in \mathfrak m$ un élément d'une suite régulière
engendrant $\mathfrak m$. Alors $x$ est un non-diviseur de zéro
sur $M$ d'après la proposition \ref{proposition-CM-module}, et
$M/xM$ est un module de Cohen-Macaulay maximal sur $R/xR$.
Par récurrence sur $\dim(R)$, nous voyons que $M/xM$ est libre.
Le lemme \ref{lemma-free-mod-x} permet de conclure.
\end{proof}

\begin{lemma}
\label{lemma-regular-mod-x}
Supposons que $R$ soit un anneau local noethérien.
Soit $x \in \mathfrak m$ un non-diviseur de zéro
tel que $R/xR$ soit un anneau local régulier. Alors $R$ est un anneau local régulier.
Plus généralement, si $x_1, \ldots, x_r$ est une suite régulière dans $R$
telle que $R/(x_1, \ldots, x_r)$ soit un anneau local régulier, alors
$R$ est un anneau local régulier.
\end{lemma}

\begin{proof}
C'est vrai parce que $x$, avec les relèvements d'un système
minimal de générateurs de l'idéal maximal de $R/xR$, fournit
$\dim(R)$ générateurs de $\mathfrak m$.
Utiliser le lemme \ref{lemma-one-equation}.
La dernière assertion découle de la première et d'une récurrence.
\end{proof}

\begin{lemma}
\label{lemma-colimit-regular}
Soit $(R_i, \varphi_{ii'})$ un système inductif d'anneaux locaux dont les
morphismes de transition sont des homomorphismes locaux d'anneaux. Si chaque $R_i$ est un anneau local régulier et si
$R = \colim R_i$ est noethérien, alors $R$ est un anneau local régulier.
\end{lemma}

\begin{proof}
Soit $\mathfrak m \subset R$ l'idéal maximal ; c'est la colimite
des idéaux maximaux $\mathfrak m_i \subset R_i$.
Nous prouvons le lemme par récurrence sur $d = \dim \mathfrak m/\mathfrak m^2$.
Si $d = 0$, alors $R = R/\mathfrak m$ est un corps et $R$ est un anneau local régulier.
Si $d > 0$, choisissons $x \in \mathfrak m$, $x \not \in \mathfrak m^2$.
Pour un certain $i$, nous pouvons trouver $x_i \in \mathfrak m_i$ qui s'envoie sur $x$.
Remarquons que $R/xR = \colim_{i' \geq i} R_{i'}/x_iR_{i'}$ est un anneau
local noethérien. D'après le
lemme \ref{lemma-regular-ring-CM},
nous voyons que $R_{i'}/x_iR_{i'}$ est un anneau local régulier.
Ainsi, par récurrence, nous voyons
que $R/xR$ est un anneau local régulier. Puisque chaque $R_i$ est intègre
(lemme \ref{lemma-regular-graded}), nous voyons que $R$ est intègre.
Ainsi, $x$ est un non-diviseur de zéro et nous concluons que $R$ est
un anneau local régulier d'après le lemme \ref{lemma-regular-mod-x}.
\end{proof}





\section{Épimorphismes d'anneaux}
\label{section-epimorphism}

\noindent
Dans toute catégorie, il existe une notion d'{\it épimorphisme}.
Une partie de ces résultats provient de \cite{Autour} et de \cite{Mazet}.

\begin{lemma}
\label{lemma-epimorphism}
Soit $R \to S$ un homomorphisme d'anneaux. Les assertions suivantes sont équivalentes :
\begin{enumerate}
\item $R \to S$ est un épimorphisme,
\item les deux homomorphismes d'anneaux $S \to S \otimes_R S$ sont égaux,
\item l'un ou l'autre des homomorphismes d'anneaux $S \to S \otimes_R S$ est un isomorphisme, et
\item l'homomorphisme d'anneaux $S \otimes_R S \to S$ est un isomorphisme.
\end{enumerate}
\end{lemma}

\begin{proof}
Omis.
\end{proof}

\begin{lemma}
\label{lemma-composition-epimorphism}
La composée de deux épimorphismes d'anneaux est un épimorphisme.
\end{lemma}

\begin{proof}
Omis. Indication : c'est vrai dans toute catégorie.
\end{proof}

\begin{lemma}
\label{lemma-base-change-epimorphism}
Si $R \to S$ est un épimorphisme d'anneaux et $R \to R'$ un homomorphisme d'anneaux quelconque,
alors $R' \to R' \otimes_R S$ est un épimorphisme.
\end{lemma}

\begin{proof}
Omis. Indication : c'est vrai dans toute catégorie possédant des sommes amalgamées.
\end{proof}

\begin{lemma}
\label{lemma-permanence-epimorphism}
Si $A \to B \to C$ sont des homomorphismes d'anneaux et $A \to C$ un épimorphisme, alors
$B \to C$ l'est aussi.
\end{lemma}

\begin{proof}
Omis. Indication : c'est vrai dans toute catégorie.
\end{proof}

\noindent
Cela signifie en particulier que, si $R \to S$ est un épimorphisme d'image
$\overline{R} \subset S$, alors $\overline{R} \to S$ est un épimorphisme.
Ainsi, lorsque nous prouvons des résultats sur les épimorphismes, nous pouvons souvent supposer
que l'homomorphisme est injectif. Le lemme suivant signifie en particulier que
toute localisation est un épimorphisme.

\begin{lemma}
\label{lemma-epimorphism-local}
Soit $R \to S$ un homomorphisme d'anneaux. Les assertions suivantes sont équivalentes :
\begin{enumerate}
\item $R \to S$ est un épimorphisme, et
\item $R_{\mathfrak p} \to S_{\mathfrak p}$ est un épimorphisme pour
tout idéal premier $\mathfrak p$ de $R$.
\end{enumerate}
\end{lemma}

\begin{proof}
Puisque $S_{\mathfrak p} = R_{\mathfrak p} \otimes_R S$ (voir le
lemme \ref{lemma-tensor-localization}),
nous voyons que (1) implique (2) d'après le
lemme \ref{lemma-base-change-epimorphism}.
Réciproquement, supposons que (2) soit vérifiée. Soient $a, b : S \to A$ deux homomorphismes d'anneaux
de $S$ vers un anneau $A$ qui coïncident sur l'homomorphisme $R \to S$. Par hypothèse, nous voyons
que, pour tout idéal premier $\mathfrak p$ de $R$, les homomorphismes induits
$a_{\mathfrak p}, b_{\mathfrak p} : S_{\mathfrak p} \to A_{\mathfrak p}$ sont
égaux. Ainsi, $a = b$ puisque $A \subset \prod_{\mathfrak p} A_{\mathfrak p}$, voir le
lemme \ref{lemma-characterize-zero-local}.
\end{proof}

\begin{lemma}
\label{lemma-finite-epimorphism-surjective}
\begin{slogan}
Un homomorphisme d'anneaux est surjectif si et seulement si c'est un épimorphisme fini.
\end{slogan}
Soit $R \to S$ un homomorphisme d'anneaux. Les assertions suivantes sont équivalentes :
\begin{enumerate}
\item $R \to S$ est un épimorphisme fini, et
\item $R \to S$ est surjectif.
\end{enumerate}
\end{lemma}

\begin{proof}
(Ce lemme semble avoir été redémontré de nombreuses fois dans la littérature, et
possède de nombreuses preuves différentes.)
Il est clair qu'un homomorphisme d'anneaux surjectif est un épimorphisme.
Supposons que $R \to S$ soit un homomorphisme d'anneaux fini tel que
$S \otimes_R S \to S$ soit un isomorphisme. Notre but est de montrer que
$R \to S$ est surjectif. Supposons que $S/R$ ne soit pas nul.
La suite exacte $R \to S \to S/R \to 0$
donne une suite exacte
$$
R \otimes_R S \to S \otimes_R S \to S/R \otimes_R S \to 0.
$$
Notre hypothèse implique que la première flèche est un isomorphisme ; ainsi,
nous concluons que $S/R \otimes_R S = 0$. Donc aussi $S/R \otimes_R S/R = 0$. D'après le
lemme \ref{lemma-trivial-filter-finite-module},
il existe une surjection de $R$-modules $S/R \to R/I$ pour un certain idéal propre
$I \subset R$. Il existe donc une
surjection $S/R \otimes_R S/R \to R/I \otimes_R R/I = R/I \not = 0$,
ce qui est une contradiction.
\end{proof}

\begin{lemma}
\label{lemma-faithfully-flat-epimorphism}
Un épimorphisme fidèlement plat est un isomorphisme.
\end{lemma}

\begin{proof}
Cela résulte clairement de la partie (3) du
lemme \ref{lemma-epimorphism},
car l'homomorphisme $S \to S \otimes_R S$ est l'homomorphisme $R \to S$ tensorisé avec $S$.
\end{proof}

\begin{lemma}
\label{lemma-epimorphism-over-field}
Si $k \to S$ est un épimorphisme et $k$ un corps, alors $S = k$ ou $S = 0$.
\end{lemma}

\begin{proof}
Cela résulte clairement du
lemme \ref{lemma-faithfully-flat-epimorphism}
(puisque toute algèbre non nulle sur $k$ est fidèlement plate), ou
en montrant directement que $R \to R \otimes_k R$ ne peut pas être
surjectif sauf si $\dim_k(R) \leq 1$.
\end{proof}

\begin{lemma}
\label{lemma-epimorphism-injective-spec}
Soit $R \to S$ un épimorphisme d'anneaux. Alors
\begin{enumerate}
\item $\Spec(S) \to \Spec(R)$ est injectif, et
\item pour $\mathfrak q \subset S$ au-dessus de $\mathfrak p \subset R$,
nous avons $\kappa(\mathfrak p) = \kappa(\mathfrak q)$.
\end{enumerate}
\end{lemma}

\begin{proof}
Soit $\mathfrak p$ un idéal premier de $R$. La fibre de l'homomorphisme est le spectre
de l'anneau fibre $S \otimes_R \kappa(\mathfrak p)$. D'après le
lemme \ref{lemma-base-change-epimorphism},
l'homomorphisme $\kappa(\mathfrak p) \to S \otimes_R \kappa(\mathfrak p)$
est un épimorphisme ; ainsi, d'après le
lemme \ref{lemma-epimorphism-over-field},
nous avons soit $S \otimes_R \kappa(\mathfrak p) = 0$,
soit $S \otimes_R \kappa(\mathfrak p) = \kappa(\mathfrak p)$,
ce qui prouve (1) et (2).
\end{proof}

\begin{lemma}
\label{lemma-relations}
Soit $R$ un anneau.
Soient $M$, $N$ des $R$-modules.
Soit $\{x_i\}_{i \in I}$ un système de générateurs de $M$.
Soit $\{y_j\}_{j \in J}$ un système de générateurs de $N$.
Soit $\{m_j\}_{j \in J}$ une famille d'éléments de $M$ telle que $m_j = 0$
pour presque tout $j$.
Alors
$$
\sum\nolimits_{j \in J} m_j \otimes y_j = 0 \text{ dans } M \otimes_R N
$$
équivaut à ce qui suit :
il existe des $a_{i, j} \in R$, avec $a_{i, j} = 0$ pour presque toute
paire $(i, j)$, tels que
\begin{align*}
m_j & = \sum\nolimits_{i \in I} a_{i, j} x_i \quad\text{pour tout } j \in J, \\
0 & = \sum\nolimits_{j \in J} a_{i, j} y_j \quad\text{pour tout } i \in I.
\end{align*}
\end{lemma}

\begin{proof}
La suffisance est immédiate. Supposons que
$\sum_{j \in J} m_j \otimes y_j = 0$.
Considérons la suite exacte courte
$$
0 \to K \to \bigoplus\nolimits_{j \in J} R \to N \to 0
$$
dans laquelle le $j$-ième vecteur de base de $\bigoplus\nolimits_{j \in J} R$ s'envoie
sur $y_j$. Tensorisons-la avec $M$ pour obtenir la suite exacte
$$
K \otimes_R M \to \bigoplus\nolimits_{j \in J} M \to N \otimes_R M \to 0.
$$
L'hypothèse implique qu'il existe des éléments $k_i \in K$ tels que
$\sum k_i \otimes x_i$ s'envoie sur l'élément $(m_j)_{j \in J}$ du terme médian.
En écrivant $k_i = (a_{i, j})_{j \in J}$, nous obtenons le résultat voulu.
\end{proof}

\begin{lemma}
\label{lemma-kernel-difference-projections}
Soit $\varphi : R \to S$ un homomorphisme d'anneaux.
Soit $g \in S$. Les assertions suivantes sont équivalentes :
\begin{enumerate}
\item $g \otimes 1 = 1 \otimes g$ dans $S \otimes_R S$, et
\item il existe $n \geq 0$, des éléments $y_i, z_j \in S$
et des $x_{i, j} \in R$ pour $1 \leq i, j \leq n$ tels que
\begin{enumerate}
\item $g = \sum_{i, j \leq n} x_{i, j} y_i z_j$,
\item pour tout $j$, nous avons $\sum x_{i, j}y_i \in \varphi(R)$, et
\item pour tout $i$, nous avons $\sum x_{i, j}z_j \in \varphi(R)$.
\end{enumerate}
\end{enumerate}
\end{lemma}

\begin{proof}
Il est clair que (2) implique (1). Réciproquement, supposons que
$g \otimes 1 = 1 \otimes g$. Choisissons un système de générateurs $\{s_i\}_{i \in I}$
de $S$ comme $R$-module avec $0, 1 \in I$, $s_0 = 1$ et $s_1 = g$.
Appliquons le
lemme \ref{lemma-relations}
à la relation $g \otimes s_0 + (-1) \otimes s_1 = 0$.
Nous voyons qu'il existe des $a_{i, j} \in R$ tels que
$g = \sum_i a_{i, 0} s_i$, $-1 = \sum_i a_{i, 1} s_i$ et, pour
$j \not = 0, 1$, $0 = \sum_i a_{i, j} s_i$, et tels qu'en outre,
pour tout $i$, nous ayons $\sum_j a_{i, j}s_j = 0$.
Nous avons alors
$$
\sum\nolimits_{i, j \not = 0} a_{i, j} s_i s_j = -g + a_{0, 0}
$$
et, pour tout $j \not = 0$, nous avons
$\sum_{i \not = 0} a_{i, j}s_i \in R$. Cela prouve que $-g + a_{0, 0}$
peut s'écrire comme dans (2). Il s'ensuit que $g$ peut s'écrire
comme dans (2). Les détails sont omis.
Indication : montrer que l'ensemble des éléments de $S$ qui possèdent une
expression comme dans (2) forme une $R$-sous-algèbre de $S$.
\end{proof}

\begin{remark}
\label{remark-matrices-associated-to-elements-epicenter}
Soit $R \to S$ un homomorphisme d'anneaux. L'ensemble des éléments
$g \in S$ tels que $g \otimes 1 = 1 \otimes g$ est parfois appelé l'
{\it épicentre} de $S$. C'est une $R$-algèbre. La construction du
lemme \ref{lemma-kernel-difference-projections}
donne, pour chaque $g$ dans l'épicentre, une factorisation matricielle
$$
(g) = Y X Z
$$
avec $X \in \text{Mat}(n \times n, R)$,
$Y \in \text{Mat}(1 \times n, S)$ et
$Z \in \text{Mat}(n \times 1, S)$. En effet, soient $x_{i, j}, y_i, z_j$
comme dans la partie (2) du lemme. Posons $X = (x_{i, j})$, soit $y$ le
vecteur ligne dont les coefficients sont les $y_i$, et soit $z$ le vecteur colonne
dont les coefficients sont les $z_j$. Avec ces notations, les conditions (b) et (c) du
lemme \ref{lemma-kernel-difference-projections}
signifient exactement que $Y X \in \text{Mat}(1 \times n, R)$,
$X Z \in \text{Mat}(n \times 1, R)$.
Il s'avère très commode de considérer le triplet de
matrices $(X, YX, XZ)$. Étant donnés $n \in \mathbf{N}$ et un triplet
$(P, U, V)$, nous disons que $(P, U, V)$ est un {\it $n$-triplet associé à $g$}
s'il existe une factorisation matricielle comme ci-dessus telle que
$P = X$, $U = YX$ et $V = XZ$.
\end{remark}

\begin{lemma}
\label{lemma-epimorphism-cardinality}
Soit $R \to S$ un épimorphisme d'anneaux.
Alors le cardinal de $S$ est au plus celui de $R$.
En formule : $|S| \leq |R|$.
\end{lemma}

\begin{proof}
La condition que $R \to S$ soit un épimorphisme signifie que tout $g \in S$
vérifie $g \otimes 1 = 1 \otimes g$, voir le
lemme \ref{lemma-epimorphism}.
Nous allons utiliser les notations introduites dans la
remarque \ref{remark-matrices-associated-to-elements-epicenter}.
Supposons que $g, g' \in S$ et que $(P, U, V)$ soit un $n$-triplet
associé à la fois à $g$ et à $g'$. Nous affirmons alors que
$g = g'$. En effet, écrivons $(P, U, V) = (X, YX, XZ)$ pour une factorisation
matricielle $(g) = YXZ$ de $g$, et écrivons $(P, U, V) = (X', Y'X', X'Z')$
pour une factorisation matricielle $(g') = Y'X'Z'$ de $g'$.
Nous voyons alors que
$$
(g) = YXZ = UZ = Y'X'Z = Y'PZ = Y'XZ = Y'V = Y'X'Z' = (g')
$$
et donc que $g = g'$. Cela implique que le cardinal de $S$ est majoré
par le nombre de triplets possibles, dont le cardinal est au plus
$\sup_{n \in \mathbf{N}} |R|^n$. Si $R$ est infini, celui-ci est
au plus $|R|$, voir \cite[Ch. I, 10.13]{Kunen}.

\medskip\noindent
Si $R$ est un anneau fini, l'argument ci-dessus prouve seulement que $S$
est au plus dénombrable. En fait, dans ce cas, $R$ est artinien et
l'homomorphisme $R \to S$ est surjectif. Nous omettons la preuve de ce cas.
\end{proof}

\begin{lemma}
\label{lemma-epimorphism-modules}
Pour un homomorphisme d'anneaux $R \to S$, les assertions suivantes sont équivalentes :
\begin{enumerate}
\item $R \to S$ est un épimorphisme d'anneaux,
\item pour tous $S$-modules $N_1, N_2$, nous avons
$\Hom_S(N_1, N_2) = \Hom_R(N_1, N_2)$, et
\item le foncteur de restriction des scalaires $\text{Mod}_S \to \text{Mod}_R$
est pleinement fidèle.
\end{enumerate}
\end{lemma}

\begin{proof}
Remarquons que (2) et (3) sont équivalentes par définition.

\medskip\noindent
Supposons (1), soient $N_1, N_2$ des $S$-modules et soit $\varphi : N_1 \to N_2$
un homomorphisme $R$-linéaire. Pour tout $x \in N_1$, considérons l'homomorphisme
$S \otimes_R S \to N_2$ défini par la règle
$g \otimes g' \mapsto g\varphi(g'x)$. Puisque les deux homomorphismes $S \to S \otimes_R S$
sont des isomorphismes (lemme \ref{lemma-epimorphism}), nous concluons que
$g \varphi(g'x) = gg'\varphi(x) = \varphi(gg' x)$. Ainsi, $\varphi$
est $S$-linéaire.

\medskip\noindent
Supposons (2). Soit $N_1 = S \otimes_R S$ considéré comme un $S$-module
par l'action à gauche de $S$, et soit $N_2 = S \otimes_R S$
muni de l'action à droite de $S$. Puisque $N_1 = N_2$ comme $R$-modules,
(2) montre que les deux structures de $S$-module sur $S \otimes_R S$
coïncident. Cela implique (1) d'après le lemme \ref{lemma-epimorphism}.
\end{proof}



\section{Idéaux purs}
\label{section-pure-ideals}

\noindent
Les résultats de cette section sont abordés dans de nombreux articles, voir par exemple
\cite{Lazard}, \cite{Bkouche} et \cite{DeMarco}.

\begin{definition}
\label{definition-pure-ideal}
Soit $R$ un anneau. Nous disons qu'un idéal $I \subset R$ est {\it pur}
si l'anneau quotient $R/I$ est plat sur $R$.
\end{definition}

\begin{lemma}
\label{lemma-pure}
Soit $R$ un anneau.
Soit $I \subset R$ un idéal.
Les assertions suivantes sont équivalentes :
\begin{enumerate}
\item $I$ est pur,
\item pour tout idéal $J \subset R$, nous avons $J \cap I = IJ$,
\item pour tout idéal de type fini $J \subset R$, nous avons
$J \cap I = JI$,
\item pour tout $x \in R$, nous avons $(x) \cap I = xI$,
\item pour tout $x \in I$, nous avons $x = yx$ pour un certain $y \in I$,
\item pour tous $x_1, \ldots, x_n \in I$, il existe un
$y \in I$ tel que $x_i = yx_i$ pour tout $i = 1, \ldots, n$,
\item pour tout idéal premier $\mathfrak p$ de $R$, nous avons
$IR_{\mathfrak p} = 0$ ou $IR_{\mathfrak p} = R_{\mathfrak p}$,
\item $\text{Supp}(I) = \Spec(R) \setminus V(I)$,
\item $I$ est le noyau de l'homomorphisme $R \to (1 + I)^{-1}R$,
\item $R/I \cong S^{-1}R$ comme $R$-algèbres pour une certaine partie
multiplicative $S$ de $R$, et
\item $R/I \cong (1 + I)^{-1}R$ comme $R$-algèbres.
\end{enumerate}
\end{lemma}

\begin{proof}
Pour tout idéal $J$ de $R$, nous avons la suite exacte courte
$0 \to J \to R \to R/J \to 0$. En tensorisant avec $R/I$, nous obtenons
une suite exacte $J \otimes_R R/I \to R/I \to R/I + J \to 0$
et $J \otimes_R R/I = J/JI$. Ainsi, l'
équivalence de (1), (2) et (3) découle du
lemme \ref{lemma-flat}. En outre, ces assertions impliquent (4).

\medskip\noindent
L'implication (4) $\Rightarrow$ (5) est triviale.
Supposons (5) et soient $x_1, \ldots, x_n \in I$.
Choisissons $y_i \in I$ tel que $x_i = y_ix_i$.
Soit $y \in I$ l'élément tel que
$1 - y = \prod_{i = 1, \ldots, n} (1 - y_i)$.
Alors $x_i = yx_i$ pour tout $i = 1, \ldots, n$.
Ainsi, (6) est vérifiée, et il s'ensuit que (5) $\Leftrightarrow$ (6).

\medskip\noindent
Supposons (5). Soit $x \in I$. Alors $x = yx$ pour un certain $y \in I$.
Ainsi, $x(1 - y) = 0$, ce qui montre que $x$ s'envoie sur zéro dans
$(1 + I)^{-1}R$. Bien entendu, le noyau de l'homomorphisme
$R \to (1 + I)^{-1}R$ est toujours contenu dans $I$. Nous
voyons donc que (5) implique (9). Supposons (9). Alors, pour tout $x \in I$,
nous voyons que $x(1 - y) = 0$ pour un certain $y \in I$.
Autrement dit, $x = yx$. Nous concluons que (5) équivaut à (9).

\medskip\noindent
Supposons (5). Soit $\mathfrak p$ un idéal premier de $R$.
Si $\mathfrak p \not \in V(I)$, alors $IR_{\mathfrak p} = R_{\mathfrak p}$.
Si $\mathfrak p \in V(I)$, autrement dit, si $I \subset \mathfrak p$,
alors, pour $x \in I$, nous avons $x(1 - y) = 0$ pour un certain $y \in I$, ce qui implique que
$x$ s'envoie sur zéro dans $R_{\mathfrak p}$, c'est-à-dire que $IR_{\mathfrak p} = 0$.
Nous voyons donc que (7) est vérifiée.

\medskip\noindent
Supposons (7). Alors $(R/I)_{\mathfrak p}$ vaut soit $0$, soit $R_{\mathfrak p}$
pour tout idéal premier $\mathfrak p$ de $R$. Ainsi, d'après le
lemme \ref{lemma-flat-localization},
nous voyons que (1) est vérifiée. À ce stade, nous voyons que toutes les assertions
(1) -- (7) et (9) sont équivalentes.

\medskip\noindent
Comme $IR_{\mathfrak p} = I_{\mathfrak p}$, nous voyons que (7) implique (8).
Enfin, si (8) est vérifiée, cela signifie exactement que $I_{\mathfrak p}$
est le module nul si et seulement si $\mathfrak p \in V(I)$, ce qui
exprime clairement que (7) est vérifiée. Ainsi, (1) -- (9) sont équivalentes.

\medskip\noindent
Supposons que (1) -- (9) soient vérifiées. Alors $R/I \subset (1 + I)^{-1}R$ d'après (9), et
l'homomorphisme $R/I \to (1 + I)^{-1}R$ est également surjectif d'après la description
des localisations aux idéaux premiers fournie par (7). Ainsi, (11) est vérifiée.

\medskip\noindent
L'implication (11) $\Rightarrow$ (10) est triviale.
Et (10) implique que (1) est vérifiée, car une localisation de
$R$ est plate sur $R$, voir le
lemme \ref{lemma-flat-localization}.
\end{proof}

\begin{lemma}
\label{lemma-pure-ideal-determined-by-zero-set}
\begin{slogan}
Les idéaux purs sont déterminés par leur lieu d'annulation.
\end{slogan}
Soit $R$ un anneau.
Si $I, J \subset R$ sont des idéaux purs, alors $V(I) = V(J)$
implique $I = J$.
\end{lemma}

\begin{proof}
Par exemple, d'après la propriété (7) du
lemme \ref{lemma-pure},
nous voyons que
$I = \Ker(R \to \prod_{\mathfrak p \in V(I)} R_{\mathfrak p})$
peut être retrouvé à partir du fermé qui lui est associé.
\end{proof}

\begin{lemma}
\label{lemma-pure-open-closed-specializations}
Soit $R$ un anneau. La règle
$I \mapsto V(I)$
détermine une bijection
$$
\{I \subset R \text{ pur}\}
\leftrightarrow
\{Z \subset \Spec(R)\text{ fermé et stable par générisation}\}
$$
\end{lemma}

\begin{proof}
Soit $I$ un idéal pur. Puisque $R \to R/I$ est plat, le théorème de descente
montre que les générisations se relèvent le long de l'homomorphisme $\Spec(R/I) \to \Spec(R)$.
Ainsi, $V(I)$ est stable par générisation. Cela montre que l'application
est bien définie. D'après le
lemme \ref{lemma-pure-ideal-determined-by-zero-set},
l'application est injective. Supposons que
$Z \subset \Spec(R)$ soit fermé et stable par générisation.
Soit $J \subset R$ l'idéal radical tel que $Z = V(J)$.
Soit $I = \{x \in R : x \in xJ\}$. Remarquons que $I$ est un idéal :
si $x, y \in I$, alors il existe $f, g \in J$ tels que
$x = xf$ et $y = yg$. Alors
$$
x + y = (x + y)(f + g - fg)
$$
La vérification est laissée au lecteur.
Nous affirmons que $I$ est pur et que $V(I) = V(J)$.
Si cette affirmation est vraie, alors l'application du lemme est surjective et
le lemme est démontré.

\medskip\noindent
Remarquons que $I \subset J$, de sorte que $V(J) \subset V(I)$.
Soit $I \subset \mathfrak p$, ce dernier idéal étant premier. Considérons la partie multiplicative
$S = (R \setminus \mathfrak p)(1 + J)$. La définition de
$I$ et l'inclusion $I \subset \mathfrak p$ montrent que $0 \not \in S$.
Nous pouvons donc trouver un idéal premier $\mathfrak q$ de $R$ disjoint
de $S$, voir les
lemmes \ref{lemma-localization-zero} et
\ref{lemma-spec-localization}.
Ainsi, $\mathfrak q \subset \mathfrak p$ et
$\mathfrak q \cap (1 + J) = \emptyset$.
Cela implique que $\mathfrak q + J$ est un idéal propre de $R$.
Soit $\mathfrak m$ un idéal maximal contenant $\mathfrak q + J$.
Nous obtenons alors
$\mathfrak m \in V(J)$, et donc $\mathfrak q \in V(J) = Z$
puisque $Z$ est supposé stable par générisation.
Cela implique à son tour que $\mathfrak p \in V(J)$, car
$\mathfrak q \subset \mathfrak p$. Ainsi, nous voyons que $V(I) = V(J)$.

\medskip\noindent
Enfin, puisque $V(I) = V(J)$ (et que $J$ est radical), nous voyons que $J = \sqrt{I}$.
Choisissons $x \in I$ ; par définition, $x = xy$ pour un certain $y \in J$.
Alors $x = xy = xy^2 = \ldots = xy^n$. Puisque $y^n \in I$ pour un certain $n > 0$,
nous concluons que la propriété (5) du
lemme \ref{lemma-pure}
est vérifiée et nous voyons que $I$ est bien pur.
\end{proof}

\begin{lemma}
\label{lemma-finitely-generated-pure-ideal}
Soit $R$ un anneau. Soit $I \subset R$ un idéal.
Les assertions suivantes sont équivalentes :
\begin{enumerate}
\item $I$ est pur et de type fini,
\item $I$ est engendré par un idempotent,
\item $I$ est pur et $V(I)$ est ouvert, et
\item $R/I$ est un $R$-module projectif.
\end{enumerate}
\end{lemma}

\begin{proof}
Si (1) est vérifiée, alors $I = I \cap I = I^2$ d'après le
lemme \ref{lemma-pure}.
Ainsi, $I$ est engendré par un idempotent d'après le
lemme \ref{lemma-ideal-is-squared-union-connected}.
Ainsi, (1) $\Rightarrow$ (2).
Si (2) est vérifiée, alors $I = (e)$ et $R = (1 - e) \oplus (e)$ comme
$R$-module ; ainsi, $R/I$ est plat, $I$ est pur et $V(I) = D(1 - e)$
est ouvert. Donc (2) $\Rightarrow$ (1) $+$ (3).
Enfin, supposons (3). Alors $V(I)$ est ouvert et fermé, donc
$V(I) = D(1 - e)$ pour un certain idempotent $e$ de $R$, voir le
lemme \ref{lemma-disjoint-decomposition}.
L'idéal $J = (e)$ est un idéal pur tel que $V(J) = V(I)$ ; ainsi,
$I = J$ d'après le
lemme \ref{lemma-pure-ideal-determined-by-zero-set}.
Nous voyons ainsi que (3) $\Rightarrow$ (2). D'après le
lemme \ref{lemma-finite-projective},
nous voyons que (4) équivaut à l'assertion que $I$ est pur
et que $R/I$ est de présentation finie. En outre, $R/I$ est de présentation finie
si et seulement si $I$ est de type fini, voir le lemme \ref{lemma-extension}.
Ainsi, (4) équivaut à (1).
\end{proof}

\noindent
Nous pouvons utiliser ce qui précède pour caractériser les anneaux sur lesquels tout module
plat fini est de présentation finie.

\begin{lemma}
\label{lemma-finite-flat-module-finitely-presented}
Soit $R$ un anneau. Les assertions suivantes sont équivalentes :
\begin{enumerate}
\item tout $Z \subset \Spec(R)$ qui est fermé et stable par
générisation est également ouvert, et
\item tout $R$-module plat fini est localement libre de type fini.
\end{enumerate}
\end{lemma}

\begin{proof}
Si tout $R$-module plat fini est localement libre de type fini, alors le support
de $R/I$, où $I$ est un idéal pur, est ouvert. Ainsi, l'implication
(2) $\Rightarrow$ (1) découle du
lemme \ref{lemma-pure-ideal-determined-by-zero-set}.

\medskip\noindent
Pour la réciproque, supposons que $R$ vérifie (1).
Soit $M$ un $R$-module plat fini.
Le support $Z = \text{Supp}(M)$ de $M$ est fermé, voir le
lemme \ref{lemma-support-closed}.
D'autre part, si $\mathfrak p \subset \mathfrak p'$, alors, d'après le
lemme \ref{lemma-finite-flat-local},
le module $M_{\mathfrak p'}$ est libre, et
$M_{\mathfrak p} = M_{\mathfrak p'} \otimes_{R_{\mathfrak p'}} R_{\mathfrak p}$
Ainsi,
$\mathfrak p' \in \text{Supp}(M) \Rightarrow \mathfrak p \in \text{Supp}(M)$,
autrement dit, le support est stable par générisation.
Comme $R$ vérifie (1), nous voyons que le support de $M$ est ouvert et fermé.
Supposons que $M$ soit engendré par $r$ éléments $m_1, \ldots, m_r$.
Les modules $\wedge^i(M)$, $i = 1, \ldots, r$, sont eux aussi des $R$-modules
plats finis, car $\wedge^i(M)_{\mathfrak p} = \wedge^i(M_{\mathfrak p})$
est libre sur $R_{\mathfrak p}$. Remarquons que
$\text{Supp}(\wedge^{i + 1}(M)) \subset \text{Supp}(\wedge^i(M))$.
Nous voyons donc qu'il existe une décomposition
$$
\Spec(R) = U_0 \amalg U_1 \amalg \ldots \amalg U_r
$$
en parties ouvertes et fermées telle que le support de
$\wedge^i(M)$ soit $U_r \cup \ldots \cup U_i$ pour tout $i = 0, \ldots, r$.
Soit $\mathfrak p$ un idéal premier de $R$, et supposons que $\mathfrak p \in U_i$.
Remarquons que
$\wedge^i(M) \otimes_R \kappa(\mathfrak p) =
\wedge^i(M \otimes_R \kappa(\mathfrak p))$.
Ainsi, après avoir éventuellement renuméroté $m_1, \ldots, m_r$, nous pouvons supposer que
$m_1, \ldots, m_i$ engendrent $M \otimes_R \kappa(\mathfrak p)$. D'après le
lemme de Nakayama \ref{lemma-NAK},
nous obtenons une surjection
$$
R_f^{\oplus i} \longrightarrow M_f, \quad
(a_1, \ldots, a_i) \longmapsto \sum a_im_i
$$
pour un certain $f \in R$, $f \not \in \mathfrak p$. Nous pouvons aussi supposer que
$D(f) \subset U_i$. Cela signifie que $\wedge^i(M_f) = \wedge^i(M)_f$
est un $R_f$-module plat dont le support est tout $\Spec(R_f)$.
D'après ce qui précède, il est engendré par un seul élément, à savoir
$m_1 \wedge \ldots \wedge m_i$. Ainsi, $\wedge^i(M)_f \cong R_f/J$
pour un certain idéal pur $J \subset R_f$ tel que $V(J) = \Spec(R_f)$.
Cela signifie clairement que $J = (0)$, voir le
lemme \ref{lemma-pure-ideal-determined-by-zero-set}.
Ainsi, $m_1 \wedge \ldots \wedge m_i$ est une base de
$\wedge^i(M_f)$, et il s'ensuit que l'homomorphisme affiché est injectif
aussi bien que surjectif. Cela prouve que $M$ est localement libre de type fini,
comme voulu.
\end{proof}







\section{Anneaux de dimension globale finie}
\label{section-ring-finite-gl-dim}

\noindent
Le lemme suivant est souvent utilisé pour comparer différentes
résolutions projectives d'un module donné.

\begin{lemma}[Lemme de Schanuel]
\label{lemma-Schanuel}
Soit $R$ un anneau. Soit $M$ un $R$-module.
Supposons que
$$
0 \to K \xrightarrow{c_1} P_1 \xrightarrow{p_1} M \to 0
\quad\text{et}\quad
0 \to L \xrightarrow{c_2} P_2 \xrightarrow{p_2} M \to 0
$$
soient deux suites exactes courtes, avec $P_i$ projectif.
Alors $K \oplus P_2 \cong L \oplus P_1$. Plus précisément,
il existe un diagramme commutatif
$$
\xymatrix{
0 \ar[r] &
K \oplus P_2 \ar[r]_{(c_1, \text{id})} \ar[d] &
P_1 \oplus P_2 \ar[r]_{(p_1, 0)} \ar[d] &
M \ar[r] \ar@{=}[d] &
0 \\
0 \ar[r] &
P_1 \oplus L \ar[r]^{(\text{id}, c_2)} &
P_1 \oplus P_2 \ar[r]^{(0, p_2)} &
M \ar[r] &
0
}
$$
dont les flèches verticales sont des isomorphismes.
\end{lemma}

\begin{proof}
Considérons le module $N$ défini par la suite exacte courte
$0 \to N \to P_1 \oplus P_2 \to M \to 0$,
où le dernier homomorphisme est la somme des deux homomorphismes
$P_i \to M$. Il est facile de voir que la projection
$N \to P_1$ est surjective de noyau $L$, et que
$N \to P_2$ est surjective de noyau $K$.
Puisque les $P_i$ sont projectifs, nous avons $N \cong K \oplus P_2
\cong L \oplus P_1$. Cela prouve la première assertion.

\medskip\noindent
Pour prouver la seconde assertion (et prouver de nouveau la première), choisissons
$a : P_1 \to P_2$ et $b : P_2 \to P_1$ tels que
$p_1 = p_2 \circ a$ et $p_2 = p_1 \circ b$. C'est possible
parce que $P_1$ et $P_2$ sont projectifs. Nous obtenons alors un diagramme commutatif
$$
\xymatrix{
0 \ar[r] &
K \oplus P_2 \ar[r]_{(c_1, \text{id})} &
P_1 \oplus P_2 \ar[r]_{(p_1, 0)} &
M \ar[r] &
0 \\
0 \ar[r] &
N \ar[r] \ar[d] \ar[u] &
P_1 \oplus P_2 \ar[r]_{(p_1, p_2)}
\ar[d]_S \ar[u]^T &
M \ar[r] \ar@{=}[d] \ar@{=}[u] &
0 \\
0 \ar[r] &
P_1 \oplus L \ar[r]^{(\text{id}, c_2)} &
P_1 \oplus P_2 \ar[r]^{(0, p_2)} &
M \ar[r] &
0
}
$$
où $T$ et $S$ sont donnés par les matrices
$$
S = \left(
\begin{matrix}
\text{id} & 0 \\
a & \text{id}
\end{matrix}
\right)
\quad\text{et}\quad
T = \left(
\begin{matrix}
\text{id} & b \\
0 & \text{id}
\end{matrix}
\right)
$$
Alors $S$, $T$ et les homomorphismes $N \to P_1 \oplus L$
et $N \to K \oplus P_2$ sont des isomorphismes, comme voulu.
\end{proof}

\begin{definition}
\label{definition-finite-proj-dim}
Soit $R$ un anneau. Soit $M$ un $R$-module. Nous disons que $M$ est de
{\it dimension projective finie} s'il possède une résolution de longueur finie
par des $R$-modules projectifs. La longueur minimale d'une telle
résolution est appelée la {\it dimension projective}
de $M$.
\end{definition}

\noindent
Il est clair que la dimension projective de $M$ est $0$ si et
seulement si $M$ est un module projectif.
Le lemme suivant explique dans quelle mesure la dimension
projective est indépendante du choix d'une
résolution projective.

\begin{lemma}
\label{lemma-independent-resolution}
Soit $R$ un anneau. Supposons que $M$ soit un $R$-module de dimension
projective $d$. Supposons que $F_e \to F_{e-1} \to \ldots \to F_0 \to M \to 0$
soit exacte, avec les $F_i$ projectifs et $e \geq d - 1$.
Alors le noyau de $F_e \to F_{e-1}$ est projectif
(ou bien le noyau de $F_0 \to M$ est projectif dans le cas où
$e = 0$).
\end{lemma}

\begin{proof}
Nous le prouvons par récurrence sur $d$. Si $d = 0$, alors
$M$ est projectif. Dans ce cas, il existe une décomposition
$F_0 = \Ker(F_0 \to M) \oplus M$, et donc
$\Ker(F_0 \to M)$ est projectif. Cela achève
la preuve si $e = 0$ ; si $e > 0$, en remplaçant
$M$ par $\Ker(F_0 \to M)$, nous diminuons $e$.

\medskip\noindent
Supposons maintenant $d > 0$.
Soit $0 \to P_d \to P_{d-1} \to \ldots \to P_0 \to M \to 0$
une résolution de longueur minimale, avec les $P_i$ projectifs.
D'après le
lemme de Schanuel \ref{lemma-Schanuel},
nous avons
$P_0 \oplus \Ker(F_0 \to M) \cong F_0 \oplus \Ker(P_0 \to M)$.
Cela prouve le cas $d = 1$, $e = 0$, car le membre de
droite est alors $F_0 \oplus P_1$, qui est projectif. Nous pouvons donc désormais
supposer $e > 0$. Le module $F_0 \oplus \Ker(P_0 \to M)$ possède la
résolution projective finie
$$
0 \to P_d \to P_{d-1} \to \ldots \to
P_2 \to P_1 \oplus F_0 \to \Ker(P_0 \to M) \oplus F_0 \to 0
$$
de longueur $d - 1$. En appliquant l'hypothèse de récurrence à la suite exacte
$$
F_e \to F_{e-1} \to \ldots \to F_2 \to P_0 \oplus F_1 \to
P_0 \oplus \Ker(F_0 \to M) \to 0
$$
de longueur $e - 1$, nous concluons que $\Ker(F_e \to F_{e - 1})$
est projectif (si $e \geq 2$),
ou que $\Ker(F_1 \oplus P_0 \to F_0 \oplus P_0)$ est projectif.
Cela implique le lemme.
\end{proof}

\begin{lemma}
\label{lemma-what-kind-of-resolutions}
Soit $R$ un anneau. Soit $M$ un $R$-module. Soit $d \geq 0$.
Les assertions suivantes sont équivalentes :
\begin{enumerate}
\item $M$ est de dimension projective $\leq d$,
\item il existe une résolution
$0 \to P_d \to P_{d - 1} \to \ldots \to P_0 \to M \to 0$
avec les $P_i$ projectifs,
\item pour une certaine résolution
$\ldots \to P_2 \to P_1 \to P_0 \to M \to 0$ avec les
$P_i$ projectifs, $\Ker(P_{d - 1} \to P_{d - 2})$
est projectif si $d \geq 2$, ou $\Ker(P_0 \to M)$ est projectif si
$d = 1$, ou $M$ est projectif si $d = 0$,
\item pour toute résolution
$\ldots \to P_2 \to P_1 \to P_0 \to M \to 0$ avec les
$P_i$ projectifs, $\Ker(P_{d - 1} \to P_{d - 2})$
est projectif si $d \geq 2$, ou $\Ker(P_0 \to M)$ est projectif si
$d = 1$, ou $M$ est projectif si $d = 0$.
\end{enumerate}
\end{lemma}

\begin{proof}
L'équivalence de (1) et (2) est la définition de la dimension
projective, voir la définition \ref{definition-finite-proj-dim}.
Nous avons (2) $\Rightarrow$ (4) d'après le lemme \ref{lemma-independent-resolution}.
Les implications (4) $\Rightarrow$ (3) et (3) $\Rightarrow$ (2) sont
immédiates.
\end{proof}

\begin{lemma}
\label{lemma-what-kind-of-resolutions-local}
Soit $R$ un anneau local. Soit $M$ un $R$-module. Soit $d \geq 0$.
Les conditions équivalentes (1) -- (4) du
lemme \ref{lemma-what-kind-of-resolutions}
sont également équivalentes à
\begin{enumerate}
\item[(5)] il existe une résolution
$0 \to P_d \to P_{d - 1} \to \ldots \to P_0 \to M \to 0$
avec les $P_i$ libres.
\end{enumerate}
\end{lemma}

\begin{proof}
Cela découle du lemme \ref{lemma-what-kind-of-resolutions} et du
théorème \ref{theorem-projective-free-over-local-ring}.
\end{proof}

\begin{lemma}
\label{lemma-what-kind-of-resolutions-Noetherian}
Soit $R$ un anneau noethérien. Soit $M$ un $R$-module fini.
Soit $d \geq 0$. Les conditions équivalentes (1) -- (4) du
lemme \ref{lemma-what-kind-of-resolutions}
sont également équivalentes à
\begin{enumerate}
\item[(6)] il existe une résolution
$0 \to P_d \to P_{d - 1} \to \ldots \to P_0 \to M \to 0$
avec les $P_i$ projectifs finis.
\end{enumerate}
\end{lemma}

\begin{proof}
Choisissons une résolution $\ldots \to F_2 \to F_1 \to F_0 \to M \to 0$
avec les $F_i$ libres finis (lemme \ref{lemma-resolution-by-finite-free}).
D'après le lemme \ref{lemma-what-kind-of-resolutions}, nous voyons que
$P_d = \Ker(F_{d - 1} \to F_{d - 2})$ est projectif, du moins si $d \geq 2$.
Alors $P_d$ est un $R$-module fini puisque $R$ est noethérien et que
$P_d \subset F_{d - 1}$, qui est libre fini.
Ainsi, $0 \to P_d \to F_{d - 1} \to \ldots \to F_1 \to F_0 \to M \to 0$
est la résolution voulue.
\end{proof}

\begin{lemma}
\label{lemma-what-kind-of-resolutions-Noetherian-local}
Soit $R$ un anneau local noethérien. Soit $M$ un $R$-module fini.
Soit $d \geq 0$. Les conditions équivalentes (1) -- (4) du
lemme \ref{lemma-what-kind-of-resolutions},
la condition (5) du lemme \ref{lemma-what-kind-of-resolutions-local},
et la condition (6) du lemme \ref{lemma-what-kind-of-resolutions-Noetherian}
sont également équivalentes à
\begin{enumerate}
\item[(7)] il existe une résolution
$0 \to F_d \to F_{d - 1} \to \ldots \to F_0 \to M \to 0$
avec les $F_i$ libres finis.
\end{enumerate}
\end{lemma}

\begin{proof}
Cela découle des lemmes \ref{lemma-what-kind-of-resolutions},
\ref{lemma-what-kind-of-resolutions-local} et
\ref{lemma-what-kind-of-resolutions-Noetherian},
ainsi que du fait qu'un module projectif fini sur un anneau local
est libre fini, voir le lemme \ref{lemma-finite-projective}.
\end{proof}

\begin{lemma}
\label{lemma-projective-dimension-ext}
Soit $R$ un anneau. Soit $M$ un $R$-module. Soit $n \geq 0$.
Les assertions suivantes sont équivalentes :
\begin{enumerate}
\item $M$ est de dimension projective $\leq n$,
\item $\Ext^i_R(M, N) = 0$ pour tout $R$-module $N$ et tout
$i \geq n + 1$, et
\item $\Ext^{n + 1}_R(M, N) = 0$ pour tout $R$-module $N$.
\end{enumerate}
\end{lemma}

\begin{proof}
Supposons (1). Choisissons une résolution libre $F_\bullet \to M$ de $M$. Notons
$d_e : F_e \to F_{e - 1}$. D'après le
lemme \ref{lemma-independent-resolution},
nous voyons que $P_e = \Ker(d_e)$ est projectif pour $e \geq n - 1$.
Cela implique que $F_e \cong P_e \oplus P_{e - 1}$ pour $e \geq n$,
où $d_e$ envoie isomorphiquement le facteur $P_{e - 1}$ sur $P_{e - 1}$
dans $F_{e - 1}$. Ainsi, pour tout $R$-module $N$, le complexe
$\Hom_R(F_\bullet, N)$ est exact scindé en degrés $\geq n + 1$.
Donc (2) est vérifiée. L'implication (2) $\Rightarrow$ (3) est triviale.

\medskip\noindent
Supposons que (3) soit vérifiée. Si $n = 0$, alors $M$ est projectif d'après le
lemme \ref{lemma-characterize-projective},
et nous voyons que (1) est vérifiée. Si $n > 0$, choisissons un $R$-module libre $F$
et une surjection $F \to M$ de noyau $K$. D'après le
lemme \ref{lemma-reverse-long-exact-seq-ext}
et l'annulation de $\Ext_R^i(F, N)$ pour tout $i > 0$ donnée par la partie (1),
nous voyons que $\Ext_R^n(K, N) = 0$ pour tout $R$-module $N$.
Ainsi, par récurrence, nous voyons que $K$ est de dimension projective $\leq n - 1$.
Alors $M$ est de dimension projective $\leq n$, car toute résolution projective finie
de $K$ donne une résolution projective de longueur augmentée de un
de $M$ en ajoutant $F$ au début.
\end{proof}

\begin{lemma}
\label{lemma-exact-sequence-projective-dimension}
Soit $R$ un anneau. Soit $0 \to M' \to M \to M'' \to 0$ une suite
exacte courte de $R$-modules.
\begin{enumerate}
\item Si $M$ est de dimension projective $\leq n$ et $M''$
de dimension projective $\leq n + 1$, alors $M'$ est de dimension projective
$\leq n$.
\item Si $M'$ et $M''$ sont de dimension projective
$\leq n$, alors $M$ est de dimension projective $\leq n$.
\item Si $M'$ est de dimension projective $\leq n$ et
$M$ de dimension projective $\leq n + 1$, alors
$M''$ est de dimension projective $\leq n + 1$.
\end{enumerate}
\end{lemma}

\begin{proof}
Combiner la caractérisation de la dimension projective du
lemme \ref{lemma-projective-dimension-ext}
avec la suite exacte longue de groupes Ext du
lemme \ref{lemma-reverse-long-exact-seq-ext}.
\end{proof}

\begin{definition}
\label{definition-finite-gl-dim}
Soit $R$ un anneau. On dit que l'anneau
$R$ est de {\it dimension globale finie}
s'il existe un entier $n$ tel que
tout $R$-module possède une résolution par des
$R$-modules projectifs de longueur au plus $n$.
Le plus petit tel $n$ est alors appelé la {\it dimension globale}
de $R$.
\end{definition}

\noindent
L'argument de la preuve du lemme suivant se trouve
dans l'article \cite{Auslander} d'Auslander.

\begin{lemma}
\label{lemma-colimit-projective-dimension}
Soit $R$ un anneau. Supposons donné un module $M = \bigcup_{e \in E} M_e$
où les $M_e$ sont des sous-modules bien ordonnés par inclusion. Supposons que les quotients
$M_e/\bigcup\nolimits_{e' < e} M_{e'}$ soient de dimension projective $\leq n$.
Alors $M$ est de dimension projective $\leq n$.
\end{lemma}

\begin{proof}
Nous allons le prouver par récurrence sur $n$.

\medskip\noindent
Cas initial : $n = 0$. Alors $P_e = M_e/\bigcup_{e' < e} M_{e'}$ est projectif.
Nous pouvons donc choisir une section $P_e \to M_e$ de la projection $M_e \to P_e$.
Nous affirmons que l'homomorphisme induit $\psi : \bigoplus_{e \in E} P_e \to M$ est un
isomorphisme. En effet, si $x = \sum x_e \in \bigoplus P_e$ est non nul,
soit $e_{max}$ maximal tel que $x_{e_{max}}$ soit non nul ;
nous concluons que $y = \psi(x) = \psi(\sum x_e)$ est non nul parce que
$y \in M_{e_{max}}$ a pour image non nulle $x_{e_{max}}$ dans $P_{e_{max}}$.
D'autre part, soit $y \in M$. Alors $y \in M_e$ pour un certain $e$.
Nous montrons que $y \in \Im(\psi)$ par récurrence transfinie sur $e$.
Soit $x_e \in P_e$ l'image de $y$. Alors
$y - \psi(x_e) \in \bigcup_{e' < e} M_{e'}$.
Par hypothèse de récurrence, nous concluons que $y - \psi(x_e) \in \Im(\psi)$,
et donc que $y \in \Im(\psi)$. L'affirmation est ainsi vraie et
$\psi$ est un isomorphisme. Nous concluons que $M$ est projectif comme
somme directe de modules projectifs, voir le
lemme \ref{lemma-direct-sum-projective}.

\medskip\noindent
Si $n > 0$, alors, pour $e \in E$, nous notons $F_e$ le $R$-module libre
sur l'ensemble des éléments de $M_e$. Nous avons alors un système de
suites exactes courtes
$$
0 \to K_e \to F_e \to M_e \to 0
$$
indexé par l'ensemble bien ordonné $E$. Remarquons que les homomorphismes de transition
$F_{e'} \to F_e$ et $K_{e'} \to K_e$ sont eux aussi injectifs.
Posons $F = \bigcup F_e$ et $K = \bigcup K_e$. Alors
$$
0 \to
K_e/\bigcup\nolimits_{e' < e} K_{e'} \to
F_e/\bigcup\nolimits_{e' < e} F_{e'} \to
M_e/\bigcup\nolimits_{e' < e} M_{e'} \to 0
$$
est aussi une suite exacte courte de $R$-modules, et
$F_e/\bigcup_{e' < e} F_{e'}$ est le $R$-module libre sur l'
ensemble des éléments de $M_e$ qui ne sont pas contenus dans $\bigcup_{e' < e} M_{e'}$.
Ainsi, d'après le
lemme \ref{lemma-exact-sequence-projective-dimension},
nous voyons que la dimension projective de $K_e/\bigcup_{e' < e} K_{e'}$
est au plus $n - 1$. Par récurrence, nous concluons que $K$ est de dimension
projective au plus $n - 1$. Donc $M$ est de dimension projective au plus
$n$, ce qui conclut.
\end{proof}

\begin{lemma}
\label{lemma-finite-gl-dim}
Soit $R$ un anneau. Les assertions suivantes sont équivalentes :
\begin{enumerate}
\item $R$ est de dimension globale finie $\leq n$,
\item tout $R$-module fini est de dimension projective $\leq n$, et
\item tout $R$-module cyclique $R/I$ est de dimension projective $\leq n$.
\end{enumerate}
\end{lemma}

\begin{proof}
Il est clair que (1) $\Rightarrow$ (2) et (2) $\Rightarrow$ (3).
Supposons (3). Choisissons un ensemble $E \subset M$ de générateurs de $M$.
Choisissons un bon ordre sur $E$. Pour $e \in E$, notons
$M_e$ le sous-module de $M$ engendré par les éléments $e' \in E$
tels que $e' \leq e$. Alors $M = \bigcup_{e \in E} M_e$.
Remarquons que, pour chaque $e \in E$, le quotient
$$
M_e/\bigcup\nolimits_{e' < e} M_{e'}
$$
est soit nul, soit engendré par un élément ; il est donc de dimension
projective $\leq n$ d'après (3). D'après le lemme \ref{lemma-colimit-projective-dimension},
cela signifie que $M$ est de dimension projective $\leq n$.
\end{proof}

\begin{lemma}
\label{lemma-localize-finite-gl-dim}
Soit $R$ un anneau. Soit $M$ un $R$-module.
Soit $S \subset R$ une partie multiplicative.
\begin{enumerate}
\item Si $M$ est de dimension projective $\leq n$, alors $S^{-1}M$ est de
dimension projective $\leq n$ sur $S^{-1}R$.
\item Si $R$ est de dimension globale finie $\leq n$, alors
$S^{-1}R$ est de dimension globale finie $\leq n$.
\end{enumerate}
\end{lemma}

\begin{proof}
Soit $0 \to P_n \to P_{n - 1} \to \ldots \to P_0 \to M \to 0$
une résolution projective. Comme la localisation est exacte, voir la
proposition \ref{proposition-localization-exact},
et comme chaque $S^{-1}P_i$ est un $S^{-1}R$-module projectif, voir le
lemme \ref{lemma-ascend-properties-modules},
nous voyons que $0 \to S^{-1}P_n \to \ldots \to S^{-1}P_0 \to S^{-1}M \to 0$
est une résolution projective de $S^{-1}M$. Cela prouve (1).
Soit $M'$ un $S^{-1}R$-module.
Remarquons que $M' = S^{-1}M'$.
Nous voyons donc que (2) découle de (1).
\end{proof}










\section{Anneaux réguliers et dimension globale}
\label{section-regular-finite-gl-dim}

\noindent
Nous pouvons utiliser les résultats sur les anneaux de dimension globale finie
pour donner une autre caractérisation des anneaux locaux réguliers.

\begin{proposition}
\label{proposition-regular-finite-gl-dim}
Soit $R$ un anneau local régulier de dimension $d$.
Tout $R$-module fini $M$ de profondeur $e$ possède une résolution
libre finie
$$
0 \to F_{d-e} \to \ldots \to F_0 \to M \to 0.
$$
En particulier, un anneau local régulier est de dimension globale $\leq d$.
\end{proposition}

\begin{proof}
La première partie découle du lemme \ref{lemma-regular-mcm-free}
et du lemme \ref{lemma-mcm-resolution}. La dernière partie en découle
avec le lemme \ref{lemma-finite-gl-dim}.
\end{proof}

\begin{lemma}
\label{lemma-finite-gl-dim-primes}
Soit $R$ un anneau noethérien. Soit $n \geq 0$ un entier.
Alors $R$ est de dimension globale finie $\leq n$ si et
seulement si, pour tout idéal maximal $\mathfrak m$ de $R$,
l'anneau $R_{\mathfrak m}$ est de dimension globale $\leq n$.
\end{lemma}

\begin{proof}
Nous avons vu, dans le lemme \ref{lemma-localize-finite-gl-dim},
que, si $R$ est de dimension globale finie $n$,
alors toutes les localisations $R_{\mathfrak m}$
sont de dimension globale finie au plus $n$.
Réciproquement, supposons que tous les $R_{\mathfrak m}$
soient de dimension globale $\leq n$. Soit $M$ un
$R$-module fini. Soit
$0 \to K_n \to F_{n-1} \to \ldots \to F_0 \to M \to 0$
une résolution avec les $F_i$ libres finis.
Alors $K_n$ est un $R$-module fini.
D'après le
lemme \ref{lemma-independent-resolution}
et l'hypothèse, tous les modules $K_n \otimes_R R_{\mathfrak m}$
sont projectifs. Ainsi, d'après le
lemme \ref{lemma-finite-projective},
le module $K_n$ est projectif fini.
\end{proof}

\begin{lemma}
\label{lemma-length-resolution-residue-field}
Supposons que $R$ soit un anneau local noethérien
d'idéal maximal $\mathfrak m$ et de
corps résiduel $\kappa$. Alors
la dimension projective de $\kappa$ est
$\geq \dim_\kappa \mathfrak m / \mathfrak m^2$.
\end{lemma}

\begin{proof}
Soient $x_1 , \ldots, x_n$ des éléments de $\mathfrak m$
dont les images dans $\mathfrak m / \mathfrak m^2$ forment une base.
Considérons le {\it complexe de Koszul} sur $x_1, \ldots, x_n$.
C'est le complexe
$$
0 \to \wedge^n R^n \to \wedge^{n-1} R^n \to \wedge^{n-2} R^n \to
\ldots \to \wedge^i R^n \to \ldots \to R^n \to R
$$
dont les homomorphismes sont donnés par
$$
e_{j_1} \wedge \ldots \wedge e_{j_i}
\longmapsto
\sum_{a = 1}^i (-1)^{a + 1} x_{j_a} e_{j_1} \wedge \ldots
\wedge \hat e_{j_a} \wedge \ldots \wedge e_{j_i}
$$
Il est facile de voir qu'il s'agit d'un complexe $K_{\bullet}(R, x_{\bullet})$.
Remarquons que le conoyau du dernier homomorphisme de $K_{\bullet}(R, x_{\bullet})$
est $\kappa$ d'après la partie (8) du lemme \ref{lemma-NAK}.

\medskip\noindent
Si $\kappa$ est de dimension projective finie $d$, nous pouvons trouver
une résolution $F_{\bullet} \to \kappa$ par des $R$-modules libres finis
de longueur $d$
(lemme \ref{lemma-what-kind-of-resolutions-Noetherian-local}).
D'après le lemme \ref{lemma-add-trivial-complex},
nous pouvons supposer que tous les homomorphismes du complexe $F_{\bullet}$
ont la propriété que $\Im(F_i \to F_{i-1})
\subset \mathfrak m F_{i-1}$, car supprimer un facteur trivial
de la résolution peut tout au plus raccourcir celle-ci.
D'après le lemme \ref{lemma-compare-resolutions}, nous pouvons trouver un morphisme
de complexes $\alpha : K_{\bullet}(R, x_{\bullet}) \to F_{\bullet}$
induisant l'identité sur $\kappa$. Nous allons prouver par récurrence
que les homomorphismes $\alpha_i : \wedge^i R^n = K_i(R, x_{\bullet}) \to F_i$
ont la propriété que
$\alpha_i \otimes \kappa : \wedge^i \kappa^n \to F_i \otimes \kappa$
sont injectifs. Cela montre que $F_n \not = 0$ et donc que $d \geq n$,
comme voulu.

\medskip\noindent
Le résultat est clair pour $i = 0$, car la composée
$R \xrightarrow{\alpha_0} F_0 \to \kappa$ est non nulle.
Remarquons que $F_0$ doit être de rang $1$, puisque
sinon l'homomorphisme $F_1 \to F_0$, dont le conoyau est une seule
copie de $\kappa$, ne peut pas avoir son image contenue dans $\mathfrak m F_0$.

\medskip\noindent
Vérifions ensuite le cas $i = 1$, que nous jugeons instructif ;
le lecteur peut l'omettre, puisque l'étape de récurrence déduira le cas $i = 1$
du cas $i = 0$. Nous avons vu ci-dessus que
$F_0 = R$ et que l'homomorphisme $F_1 \to F_0 = R$ a pour image $\mathfrak m$.
Nous avons un diagramme commutatif
$$
\begin{matrix}
R^n & = & K_1(R, x_{\bullet}) & \to & K_0(R, x_{\bullet}) & = & R \\
& & \downarrow & & \downarrow & & \downarrow \\
& & F_1 & \to & F_0 & = & R
\end{matrix}
$$
où la flèche verticale la plus à droite est donnée par la multiplication
par une unité. Nous voyons donc que l'image de la composée
$R^n \to F_1 \to F_0 = R$ est elle aussi égale à $\mathfrak m$.
Ainsi, l'homomorphisme $R^n \otimes \kappa \to F_1 \otimes \kappa$
doit être injectif puisque $\dim_\kappa (\mathfrak m / \mathfrak m^2) = n$.

\medskip\noindent
Soit $i \geq 1$ et supposons l'injectivité de $\alpha_j \otimes \kappa$ prouvée
pour tout $j \leq i - 1$. Considérons le diagramme commutatif
$$
\begin{matrix}
\wedge^i R^n & = & K_i(R, x_{\bullet}) & \to & K_{i-1}(R, x_{\bullet})
& = & \wedge^{i-1} R^n \\
& & \downarrow & & \downarrow & & \\
& & F_i & \to & F_{i-1} & &
\end{matrix}
$$
Nous savons que $\wedge^{i-1} \kappa^n \to F_{i-1} \otimes \kappa$
est injectif. Cela prouve que
$\wedge^{i-1} \kappa^n \otimes_{\kappa} \mathfrak m/\mathfrak m^2
\to F_{i-1} \otimes \mathfrak m/\mathfrak m^2$ est injectif.
En outre, par notre choix du complexe, $F_i$ s'envoie dans
$\mathfrak mF_{i-1}$, et il en va de même pour le complexe de Koszul.
Nous obtenons donc un diagramme commutatif
$$
\begin{matrix}
\wedge^i \kappa^n & \to &
\wedge^{i-1} \kappa^n \otimes \mathfrak m/\mathfrak m^2 \\
\downarrow & & \downarrow \\
F_i \otimes \kappa & \to & F_{i-1} \otimes \mathfrak m/\mathfrak m^2
\end{matrix}
$$
À ce stade, il suffit de vérifier que l'homomorphisme
$\wedge^i \kappa^n \to
\wedge^{i-1} \kappa^n \otimes \mathfrak m/\mathfrak m^2$
est injectif, ce qui se vérifie directement.
\end{proof}

\begin{lemma}
\label{lemma-dim-gl-dim}
Soit $R$ un anneau local noethérien.
Supposons que le corps résiduel $\kappa$ ait une dimension
projective finie $n$ sur $R$.
Alors $\dim(R) \geq n$.
\end{lemma}

\begin{proof}
Soit $F_{\bullet}$ une résolution finie de $\kappa$ par des
$R$-modules libres finis (lemme \ref{lemma-what-kind-of-resolutions-Noetherian-local}).
D'après le lemme \ref{lemma-add-trivial-complex},
nous pouvons supposer que tous les homomorphismes du complexe $F_{\bullet}$
ont la propriété que $\Im(F_i \to F_{i-1})
\subset \mathfrak m F_{i-1}$, car supprimer un facteur trivial
de la résolution peut tout au plus raccourcir celle-ci.
Écrivons $F_n \not = 0$ et $F_i = 0$ pour $i > n$, de sorte que
la dimension projective de $\kappa$ soit $n$.
D'après la proposition \ref{proposition-what-exact}, nous voyons que
$\text{depth}_{I(\varphi_n)}(R) \geq n$, puisque $I(\varphi_n)$
ne peut pas être égal à $R$ par notre choix du complexe.
Ainsi, d'après le lemme \ref{lemma-bound-depth}, nous avons aussi $\dim(R) \geq n$.
\end{proof}

\begin{proposition}
\label{proposition-finite-gl-dim-regular}
Soit $(R, \mathfrak m, \kappa)$ un anneau local noethérien.
Les assertions suivantes sont équivalentes :
\begin{enumerate}
\item $\kappa$ est de dimension projective finie comme $R$-module,
\item $R$ est de dimension globale finie,
\item $R$ est un anneau local régulier.
\end{enumerate}
De plus, dans ce cas, la dimension globale de $R$ est égale à
$\dim(R) = \dim_\kappa(\mathfrak m/\mathfrak m^2)$.
\end{proposition}

\begin{proof}
Nous avons (3) $\Rightarrow$ (2) d'après la
proposition \ref{proposition-regular-finite-gl-dim}.
L'implication (2) $\Rightarrow$ (1) est triviale.
Supposons (1). D'après les lemmes \ref{lemma-length-resolution-residue-field}
et \ref{lemma-dim-gl-dim}, nous voyons que
$\dim(R) \geq \dim_\kappa(\mathfrak m /\mathfrak m^2)$.
Ainsi, $R$ est régulier, voir la définition
\ref{definition-regular-local} et la discussion qui la précède.
Supposons que les conditions équivalentes (1) -- (3) soient vérifiées.
D'après la proposition \ref{proposition-regular-finite-gl-dim},
la dimension globale de $R$ est au plus $\dim(R)$ et, d'après le
lemme \ref{lemma-length-resolution-residue-field},
elle est au moins $\dim_\kappa(\mathfrak m/\mathfrak m^2)$.
L'égalité annoncée est donc vérifiée.
\end{proof}

\begin{lemma}
\label{lemma-localization-of-regular-local-is-regular}
Un anneau local noethérien $R$ est régulier si et seulement si
il est de dimension globale finie. Dans ce cas,
$R_{\mathfrak p}$ est un anneau local régulier pour tout idéal premier $\mathfrak p$.
\end{lemma}

\begin{proof}
D'après les propositions \ref{proposition-finite-gl-dim-regular} et
\ref{proposition-regular-finite-gl-dim},
nous voyons qu'un anneau local noethérien est régulier si et seulement si
il est de dimension globale finie. En outre, toute localisation
$R_{\mathfrak p}$ est de dimension globale finie,
voir le lemme \ref{lemma-localize-finite-gl-dim},
et est donc un anneau local régulier.
\end{proof}

\noindent
D'après le lemme \ref{lemma-localization-of-regular-local-is-regular},
il est raisonnable de poser la définition suivante,
car elle n'entre pas en conflit avec la définition antérieure
d'un anneau local régulier.

\begin{definition}
\label{definition-regular}
Un anneau noethérien $R$ est dit {\it régulier}
si toutes ses localisations $R_{\mathfrak p}$ aux idéaux premiers sont des
anneaux locaux réguliers.
\end{definition}

\noindent
Il suffit d'exiger que les anneaux locaux aux idéaux maximaux soient réguliers.
Remarquons que cela ne revient pas à demander que $R$ soit de dimension
globale finie, même en supposant $R$ noethérien. Cela vient
de ce qu'il existe un exemple d'anneau noethérien régulier
qui n'est pas de dimension globale finie, précisément parce qu'
il n'est pas de dimension finie.

\begin{lemma}
\label{lemma-finite-gl-dim-finite-dim-regular}
Soit $R$ un anneau noethérien.
Les assertions suivantes sont équivalentes :
\begin{enumerate}
\item $R$ est de dimension globale finie $n$,
\item $R$ est un anneau régulier de dimension $n$,
\item il existe un entier $n$ tel que
toutes les localisations $R_{\mathfrak m}$ aux idéaux maximaux
soient régulières de dimension $\leq n$, avec égalité pour au moins
un $\mathfrak m$, et
\item il existe un entier $n$ tel que
toutes les localisations $R_{\mathfrak p}$ aux idéaux premiers
soient régulières de dimension $\leq n$, avec égalité pour au moins
un $\mathfrak p$.
\end{enumerate}
\end{lemma}

\begin{proof}
Cela découle de la discussion ci-dessus. Plus précisément, cela s'obtient
en combinant la définition \ref{definition-regular} avec le
lemme \ref{lemma-finite-gl-dim-primes}
et la proposition \ref{proposition-finite-gl-dim-regular}.
\end{proof}

\begin{lemma}
\label{lemma-flat-under-regular}
Soit $R \to S$ un homomorphisme local d'anneaux locaux noethériens.
Supposons que $R \to S$ soit plat et que $S$ soit régulier.
Alors $R$ est régulier.
\end{lemma}

\begin{proof}
Soit $\mathfrak m \subset R$ l'idéal maximal,
et soit $\kappa = R/\mathfrak m$ le corps résiduel.
Soit $d = \dim S$.
Choisissons une résolution quelconque $F_\bullet \to \kappa$
où chaque $F_i$ est un $R$-module libre fini. Posons
$K_d = \Ker(F_{d - 1} \to F_{d - 2})$.
Par platitude de $R \to S$, le complexe
$0 \to K_d \otimes_R S \to F_{d - 1} \otimes_R S \to \ldots
\to F_0 \otimes_R S \to \kappa \otimes_R S \to 0$
reste exact. Puisque la dimension globale de $S$
est $d$, voir la proposition \ref{proposition-regular-finite-gl-dim},
nous voyons que $K_d \otimes_R S$ est un $S$-module libre fini
(voir aussi le lemme \ref{lemma-independent-resolution}).
D'après le lemme \ref{lemma-finite-projective-descends}, nous voyons
que $K_d$ est un $R$-module libre fini.
Ainsi, $\kappa$ est de dimension projective finie et $R$ est régulier d'après la
proposition \ref{proposition-finite-gl-dim-regular}.
\end{proof}








\section{Auslander-Buchsbaum}
\label{section-Auslander-Buchsbaum}

\noindent
Le résultat suivant se trouve dans \cite{Auslander-Buchsbaum}.

\begin{proposition}
\label{proposition-Auslander-Buchsbaum}
Soit $R$ un anneau local noethérien. Soit $M$ un $R$-module fini non nul
de dimension projective finie $\text{pd}_R(M)$. Alors nous avons
$$
\text{depth}(R) = \text{pd}_R(M) + \text{depth}(M)
$$
\end{proposition}

\begin{proof}
Nous le prouvons par récurrence sur $\text{depth}(M)$. Le cas le plus intéressant
est celui où $\text{depth}(M) = 0$. Dans ce cas, soit
$$
0 \to R^{n_e} \to R^{n_{e-1}} \to \ldots \to R^{n_0} \to M \to 0
$$
une résolution libre finie minimale, de sorte que $e = \text{pd}_R(M)$.
D'après le lemme \ref{lemma-add-trivial-complex}, nous pouvons supposer que tous les coefficients
matriciels des homomorphismes du complexe sont contenus dans l'idéal maximal
de $R$. Alors, d'une part, la
proposition \ref{proposition-what-exact} montre que
$\text{depth}(R) \geq e$. D'autre part, en décomposant la longue
suite exacte en suites exactes courtes
\begin{align*}
0 \to R^{n_e} \to R^{n_{e - 1}} \to K_{e - 2} \to 0,\\
0 \to K_{e - 2} \to R^{n_{e - 2}} \to K_{e - 3} \to 0,\\
\ldots,\\
0 \to K_0 \to R^{n_0} \to M \to 0
\end{align*}
nous voyons, à l'aide du lemme \ref{lemma-depth-in-ses}, que
\begin{align*}
\text{depth}(K_{e - 2}) \geq \text{depth}(R) - 1,\\
\text{depth}(K_{e - 3}) \geq \text{depth}(R) - 2,\\
\ldots,\\
\text{depth}(K_0) \geq \text{depth}(R) - (e - 1),\\
\text{depth}(M) \geq \text{depth}(R) - e
\end{align*}
et, puisque $\text{depth}(M) = 0$, nous concluons que $\text{depth}(R) \leq e$.
Cela achève la preuve du cas $\text{depth}(M) = 0$.

\medskip\noindent
Étape de récurrence. Si $\text{depth}(M) > 0$, choisissons $x \in \mathfrak m$
qui soit un non-diviseur de zéro à la fois sur $M$ et sur $R$. C'est possible, car
soit $\text{pd}_R(M) > 0$ et $\text{depth}(R) > 0$ d'après la
proposition \ref{proposition-what-exact} précitée, soit $\text{pd}_R(M) = 0$, auquel
cas $M$ est libre fini et donc aussi $\text{depth}(R) = \text{depth}(M) > 0$.
Ainsi, $\text{depth}(R \oplus M) > 0$ d'après le lemme \ref{lemma-depth-in-ses}
(par exemple), et nous pouvons trouver $x \in \mathfrak m$ qui soit un non-diviseur de zéro
à la fois sur $R$ et sur $M$. Soit
$$
0 \to R^{n_e} \to R^{n_{e-1}} \to \ldots \to R^{n_0} \to M \to 0
$$
une résolution minimale comme ci-dessus. Une application du lemme du serpent
montre que
$$
0 \to (R/xR)^{n_e} \to (R/xR)^{n_{e-1}} \to \ldots \to (R/xR)^{n_0} \to
M/xM \to 0
$$
est elle aussi une résolution minimale. Ainsi, $\text{pd}_R(M) = \text{pd}_{R/xR}(M/xM)$.
D'après le lemme \ref{lemma-depth-drops-by-one}, nous avons
$\text{depth}(R/xR) = \text{depth}(R) - 1$ et
$\text{depth}(M/xM) = \text{depth}(M) - 1$.
Jusqu'ici, toutes les profondeurs étaient prises comme $R$-modules, mais nous observons que
$\text{depth}_R(M/xM) = \text{depth}_{R/xR}(M/xM)$, et de même pour $R/xR$.
Par hypothèse de récurrence, nous voyons que la formule
d'Auslander-Buchsbaum est vérifiée pour $M/xM$ sur $R/xR$. Puisque les
profondeurs de $R/xR$ et de $M/xM$ ont toutes deux diminué de un, tandis que la dimension
projective n'a pas changé, nous concluons.
\end{proof}









\section{Homomorphismes et dimension}
\label{section-homomorphism-dimension}

\noindent
Cette section contient une collection de résultats élémentaires reliant les
dimensions d'anneaux lorsqu'il existe des homomorphismes entre eux.

\begin{lemma}
\label{lemma-dimension-going-up}
Supposons que $R \to S$ soit un homomorphisme d'anneaux vérifiant soit la montée, voir la
définition \ref{definition-going-up-down}, soit la descente,
voir la définition \ref{definition-going-up-down}.
Supposons en outre que $\Spec(S) \to \Spec(R)$
soit surjectif. Alors $\dim(R) \leq \dim(S)$.
\end{lemma}

\begin{proof}
Supposons la montée.
Prenons une chaîne quelconque $\mathfrak p_0 \subset \mathfrak p_1 \subset \ldots
\subset \mathfrak p_e$ d'idéaux premiers de $R$.
Par surjectivité, nous pouvons choisir un idéal premier $\mathfrak q_0$ s'envoyant
sur $\mathfrak p_0$. Par montée, nous pouvons le prolonger en une chaîne
de longueur $e$ d'idéaux premiers $\mathfrak q_i$ au-dessus des
$\mathfrak p_i$. Ainsi, $\dim(S) \geq \dim(R)$.
Le cas de la descente est exactement semblable.
Voir aussi Topologie, lemme \ref{topology-lemma-dimension-specializations-lift},
pour une version purement topologique.
\end{proof}

\begin{lemma}
\label{lemma-going-up-maximal-on-top}
Supposons que $R \to S$ soit un homomorphisme d'anneaux possédant la propriété de montée,
voir la définition \ref{definition-going-up-down}. Si
$\mathfrak q \subset S$ est un idéal maximal,
alors l'image réciproque de $\mathfrak q$ dans $R$
est elle aussi un idéal maximal.
\end{lemma}

\begin{proof}
Trivial.
\end{proof}

\begin{lemma}
\label{lemma-integral-dim-up}
Supposons que $R \to S$ soit un homomorphisme d'anneaux tel que $S$ soit entier sur $R$.
Alors $\dim (R) \geq \dim(S)$, et tout point fermé de $\Spec(S)$
s'envoie sur un point fermé de $\Spec(R)$.
\end{lemma}

\begin{proof}
Cela découle immédiatement des lemmes \ref{lemma-integral-no-inclusion} et
\ref{lemma-going-up-maximal-on-top}
et des définitions.
\end{proof}

\begin{lemma}
\label{lemma-integral-sub-dim-equal}
Supposons $R \subset S$ et $S$ entier sur $R$.
Alors $\dim(R) = \dim(S)$.
\end{lemma}

\begin{proof}
Cela résulte de la combinaison des lemmes
\ref{lemma-integral-going-up},
\ref{lemma-integral-overring-surjective},
\ref{lemma-dimension-going-up} et
\ref{lemma-integral-dim-up}.
\end{proof}

\begin{definition}
\label{definition-fibre}
Supposons que $R \to S$ soit un homomorphisme d'anneaux.
Soit $\mathfrak q \subset S$ un idéal premier au-dessus
de l'idéal premier $\mathfrak p$ de $R$.
L'{\it anneau local de la fibre en $\mathfrak q$}
est l'anneau local
$$
S_{\mathfrak q}/\mathfrak pS_{\mathfrak q}
=
(S/\mathfrak pS)_{\mathfrak q}
=
(S \otimes_R \kappa(\mathfrak p))_{\mathfrak q}
$$
\end{definition}

\begin{lemma}
\label{lemma-dimension-base-fibre-total}
Soit $R \to S$ un homomorphisme d'anneaux noethériens.
Soit $\mathfrak q \subset S$ un idéal premier au-dessus
de l'idéal premier $\mathfrak p$. Alors
$$
\dim(S_{\mathfrak q})
\leq
\dim(R_{\mathfrak p})
+
\dim(S_{\mathfrak q}/\mathfrak pS_{\mathfrak q}).
$$
\end{lemma}

\begin{proof}
Nous utilisons la caractérisation de la dimension de la
proposition \ref{proposition-dimension}.
Soient $x_1, \ldots, x_d$ des éléments de $\mathfrak p$
engendrant un idéal de définition de $R_{\mathfrak p}$, avec
$d = \dim(R_{\mathfrak p})$.
Soient $y_1, \ldots, y_e$ des éléments de $\mathfrak q$
engendrant un idéal de définition de
$S_{\mathfrak q}/\mathfrak pS_{\mathfrak q}$,
avec $e = \dim(S_{\mathfrak q}/\mathfrak pS_{\mathfrak q})$.
Il est clair que $S_{\mathfrak q}/(x_1, \ldots, x_d, y_1, \ldots, y_e)$
a un idéal maximal nilpotent. Ainsi,
$x_1, \ldots, x_d, y_1, \ldots, y_e$ engendrent un idéal de définition
de $S_{\mathfrak q}$.
\end{proof}

\begin{lemma}
\label{lemma-dimension-base-fibre-equals-total}
Soit $R \to S$ un homomorphisme d'anneaux noethériens.
Soit $\mathfrak q \subset S$ un idéal premier au-dessus
de l'idéal premier $\mathfrak p$. Supposons que la propriété de descente soit vérifiée
pour $R \to S$ (par exemple, si $R \to S$ est plat, voir le
lemme \ref{lemma-flat-going-down}). Alors
$$
\dim(S_{\mathfrak q})
=
\dim(R_{\mathfrak p})
+
\dim(S_{\mathfrak q}/\mathfrak pS_{\mathfrak q}).
$$
\end{lemma}

\begin{proof}
D'après le lemme \ref{lemma-dimension-base-fibre-total},
nous avons une inégalité
$\dim(S_{\mathfrak q}) \leq
\dim(R_{\mathfrak p}) + \dim(S_{\mathfrak q}/\mathfrak pS_{\mathfrak q})$.
Pour obtenir l'égalité, choisissons une chaîne d'idéaux premiers
$\mathfrak pS \subset \mathfrak q_0 \subset \mathfrak q_1 \subset \ldots
\subset \mathfrak q_d = \mathfrak q$, avec
$d = \dim(S_{\mathfrak q}/\mathfrak pS_{\mathfrak q})$.
D'autre part, choisissons une chaîne d'idéaux premiers
$\mathfrak p_0 \subset \mathfrak p_1 \subset \ldots \subset \mathfrak p_e
= \mathfrak p$, avec $e = \dim(R_{\mathfrak p})$.
Par le théorème de descente, nous pouvons choisir
$\mathfrak q_{-1} \subset \mathfrak q_0$ au-dessus de
$\mathfrak p_{e-1}$. Puis nous pouvons choisir
$\mathfrak q_{-2} \subset \mathfrak q_{-1}$ au-dessus de
$\mathfrak p_{e-2}$. En poursuivant par récurrence, nous obtenons finalement
une chaîne
$\mathfrak q_{-e} \subset \ldots \subset \mathfrak q_d$
de longueur $e + d$.
\end{proof}

\begin{lemma}
\label{lemma-flat-over-regular-with-regular-fibre}
Soit $R \to S$ un homomorphisme local d'anneaux locaux noethériens.
Supposons que
\begin{enumerate}
\item $R$ soit régulier,
\item $S/\mathfrak m_RS$ soit régulier, et
\item $R \to S$ soit plat.
\end{enumerate}
Alors $S$ est régulier.
\end{lemma}

\begin{proof}
D'après le lemme \ref{lemma-dimension-base-fibre-equals-total},
nous avons
$\dim(S) = \dim(R) + \dim(S/\mathfrak m_RS)$.
Choisissons des générateurs $x_1, \ldots, x_d \in \mathfrak m_R$, avec $d = \dim(R)$,
et choisissons $y_1, \ldots, y_e \in \mathfrak m_S$
qui engendrent l'idéal maximal de $S/\mathfrak m_RS$, avec
$e = \dim(S/\mathfrak m_RS)$. Nous voyons alors que
$x_1, \ldots, x_d, y_1, \ldots, y_e$ sont des éléments qui engendrent
l'idéal maximal de $S$ et que $e + d = \dim(S)$.
\end{proof}

\noindent
Le lemme ci-dessous servira plus loin à montrer que les anneaux de type fini sur
un corps sont de Cohen-Macaulay si et seulement s'ils sont quasi-finis et plats sur
un anneau polynomial. C'est une réciproque partielle du
lemme \ref{lemma-CM-over-regular-flat}.

\begin{lemma}
\label{lemma-finite-flat-over-regular-CM}
Soit $R \to S$ un homomorphisme local d'anneaux locaux noethériens.
Supposons $R$ de Cohen-Macaulay.
Si $S$ est fini et plat sur $R$, ou si $S$ est plat sur $R$ et si
$\dim(S) \leq \dim(R)$, alors $S$ est de Cohen-Macaulay et $\dim(R) = \dim(S)$.
\end{lemma}

\begin{proof}
Soit $x_1, \ldots, x_d \in \mathfrak m_R$ une suite régulière
de longueur $d = \dim(R)$. D'après le lemme \ref{lemma-flat-increases-depth},
elle s'envoie sur une suite régulière dans $S$.
Ainsi, $S$ est de Cohen-Macaulay si $\dim(S) \leq d$. C'est le cas
si $S$ est fini et plat sur $R$, d'après le lemme \ref{lemma-integral-sub-dim-equal}.
Et, dans le second cas, nous l'avons supposé.
\end{proof}








\section{La formule de dimension}
\label{section-dimension-formula}

\noindent
Rappelons les définitions d'un anneau caténaire (définition \ref{definition-catenary})
et d'un anneau universellement caténaire (définition \ref{definition-universally-catenary}).

\begin{lemma}
\label{lemma-dimension-formula}
Soit $R \to S$ un homomorphisme d'anneaux.
Soit $\mathfrak q$ un idéal premier de $S$ au-dessus de l'idéal premier $\mathfrak p$ de $R$.
Supposons que
\begin{enumerate}
\item $R$ est noethérien,
\item $R \to S$ est de type fini,
\item $R$, $S$ sont des anneaux intègres, et
\item $R \subset S$.
\end{enumerate}
Alors nous avons
$$
\text{hauteur}(\mathfrak q)
\leq
\text{hauteur}(\mathfrak p) + \text{trdeg}_R(S)
- \text{trdeg}_{\kappa(\mathfrak p)} \kappa(\mathfrak q)
$$
avec égalité si $R$ est universellement caténaire.
\end{lemma}

\begin{proof}
Supposons que $R \subset S' \subset S$, où $S'$ est une sous-algèbre de type fini
sur $R$ de $S$. Dans ce cas, posons $\mathfrak q' = S' \cap \mathfrak q$.
Le lemme pour les homomorphismes d'anneaux $R \to S'$ et $S' \to S$ implique le
lemme pour $R \to S$ par additivité du degré de transcendance dans les tours
d'extensions de corps (Corps, lemme \ref{fields-lemma-transcendence-degree-tower}).
Nous pouvons donc raisonner par récurrence sur le nombre de générateurs
de $S$ sur $R$ et nous ramener au cas où $S$ est engendré par
un élément sur $R$.

\medskip\noindent
Cas I : $S = R[x]$ est une algèbre polynomiale sur $R$.
Dans ce cas, nous avons $\text{trdeg}_R(S) = 1$.
De plus, $R \to S$ est plat et donc
$$
\dim(S_{\mathfrak q}) =
\dim(R_{\mathfrak p}) + \dim(S_{\mathfrak q}/\mathfrak pS_{\mathfrak q})
$$
voir le lemme \ref{lemma-dimension-base-fibre-equals-total}.
Soit $\mathfrak r = \mathfrak pS$. Alors
$\text{trdeg}_{\kappa(\mathfrak p)} \kappa(\mathfrak q) = 1$
équivaut à $\mathfrak q = \mathfrak r$, et implique que
$\dim(S_{\mathfrak q}/\mathfrak pS_{\mathfrak q}) = 0$.
De même, $\text{trdeg}_{\kappa(\mathfrak p)} \kappa(\mathfrak q) = 0$
équivaut à avoir une inclusion stricte
$\mathfrak r \subset \mathfrak q$, ce qui implique que
$\dim(S_{\mathfrak q}/\mathfrak pS_{\mathfrak q}) = 1$.
Nous avons ainsi terminé le cas I, avec égalité dans tous les cas.

\medskip\noindent
Cas II : $S = R[x]/\mathfrak n$ avec $\mathfrak n \not = 0$.
Dans ce cas, nous avons $\text{trdeg}_R(S) = 0$.
Notons $\mathfrak q' \subset R[x]$ l'idéal premier correspondant à $\mathfrak q$.
Nous avons ainsi
$$
S_{\mathfrak q} = (R[x])_{\mathfrak q'}/\mathfrak n(R[x])_{\mathfrak q'}
$$
D'après le cas précédent, nous avons
$\dim((R[x])_{\mathfrak q'}) =
\dim(R_{\mathfrak p}) + 1
- \text{trdeg}_{\kappa(\mathfrak p)} \kappa(\mathfrak q)$.
Puisque $\mathfrak n \not = 0$, nous voyons que la dimension de
$S_{\mathfrak q}$ diminue d'au moins un, voir le
lemme \ref{lemma-one-equation},
ce qui démontre l'inégalité du lemme.
Pour voir l'égalité lorsque $R$ est universellement caténaire, remarquons que
$\mathfrak n \subset R[x]$ est un idéal premier de hauteur un puisqu'il correspond
à un idéal premier non nul de $F[x]$, où $F$ est le corps des fractions de $R$.
Par conséquent, toute chaîne maximale d'idéaux premiers de
$S_\mathfrak q = R[x]_{\mathfrak q'}/\mathfrak nR[x]_{\mathfrak q'}$
correspond à une chaîne maximale d'idéaux premiers
de longueur supérieure de 1 entre $\mathfrak q'$ et $(0)$ dans $R[x]$.
Si $R$ est universellement caténaire, elles ont toutes la même longueur, égale à
la hauteur de $\mathfrak q'$. Cela montre que
$\dim(S_\mathfrak q) = \dim(R[x]_{\mathfrak q'}) - 1$
et implique l'égalité voulue.
\end{proof}

\noindent
Le lemme suivant affirme que les morphismes génériquement finis
ont tendance à être quasi-finis en codimension $1$.

\begin{lemma}
\label{lemma-finite-in-codim-1}
Soit $A \to B$ un homomorphisme d'anneaux.
Supposons que
\begin{enumerate}
\item $A \subset B$ est une extension d'anneaux intègres,
\item l'extension induite des corps de fractions est finie,
\item $A$ est noethérien, et
\item $A \to B$ est de type fini.
\end{enumerate}
Soit $\mathfrak p \subset A$ un idéal premier de hauteur $1$.
Il existe alors au plus un nombre fini d'idéaux premiers de $B$
au-dessus de $\mathfrak p$, et ils sont tous de hauteur $1$.
\end{lemma}

\begin{proof}
D'après la formule de dimension (lemme \ref{lemma-dimension-formula}),
pour tout idéal premier $\mathfrak q$ au-dessus de $\mathfrak p$, nous avons
$$
\dim(B_{\mathfrak q}) \leq
\dim(A_{\mathfrak p}) - \text{trdeg}_{\kappa(\mathfrak p)} \kappa(\mathfrak q).
$$
Comme l'anneau intègre $B_\mathfrak q$ possède au moins $2$ idéaux premiers, nous voyons que
$\dim(B_{\mathfrak q}) \geq 1$. Nous en concluons que
$\dim(B_{\mathfrak q}) = 1$ et que l'extension
$\kappa(\mathfrak p) \subset \kappa(\mathfrak q)$ est algébrique.
Ainsi, $\mathfrak q$ définit un point fermé de sa fibre
$\Spec(B \otimes_A \kappa(\mathfrak p))$, voir le
lemme \ref{lemma-finite-residue-extension-closed}.
Puisque $B \otimes_A \kappa(\mathfrak p)$ est un anneau noethérien,
la fibre $\Spec(B \otimes_A \kappa(\mathfrak p))$
est un espace topologique noethérien, voir le
lemme \ref{lemma-Noetherian-topology}.
Un espace topologique noethérien constitué de points fermés
est fini, voir par exemple
Topologie, lemme \ref{topology-lemma-Noetherian}.
\end{proof}









\section{Dimension des algèbres de type fini sur les corps}
\label{section-dimension-finite-type-algebras}

\noindent
Dans cette section, nous calculons la dimension d'un anneau de polynômes sur
un corps. Nous démontrons également que la dimension d'une algèbre intègre
de type fini sur un corps est la dimension de ses anneaux locaux aux idéaux
maximaux. Nous établirons le lien avec le degré de transcendance
sur le corps de base dans la
section \ref{section-dimension-finite-type-algebras-reprise}.

\begin{lemma}
\label{lemma-dim-affine-space}
Soit $\mathfrak m$ un idéal maximal de $k[x_1, \ldots, x_n]$.
L'idéal $\mathfrak m$ est engendré par $n$ éléments.
La dimension de $k[x_1, \ldots, x_n]_{\mathfrak m}$ est $n$.
Ainsi, $k[x_1, \ldots, x_n]_{\mathfrak m}$ est un anneau local régulier
de dimension $n$.
\end{lemma}

\begin{proof}
D'après le théorème des zéros de Hilbert
(théorème \ref{theorem-nullstellensatz}),
nous savons que le corps résiduel $\kappa = \kappa(\mathfrak m)$ est
une extension finie de $k$. Notons $\alpha_i \in \kappa$ l'
image de $x_i$. Notons $\kappa_i = k(\alpha_1, \ldots, \alpha_i)
\subset \kappa$, $i = 1, \ldots, n$, et $\kappa_0 = k$.
Remarquons que $\kappa_i = k[\alpha_1, \ldots, \alpha_i]$
d'après la théorie des corps. Définissons par récurrence des éléments
$f_i \in \mathfrak m \cap k[x_1, \ldots, x_i]$
comme suit : soit $P_i(T) \in \kappa_{i-1}[T]$
le polynôme minimal unitaire de $\alpha_i $ sur $\kappa_{i-1}$.
Soit $Q_i(T) \in k[x_1, \ldots, x_{i-1}][T]$ un relèvement unitaire de $P_i(T)$
(de même degré). Posons $f_i = Q_i(x_i)$.
Remarquons que si $d_i = \deg_T(P_i) = \deg_T(Q_i) = \deg_{x_i}(f_i)$,
alors $d_1d_2\ldots d_i = [\kappa_i : k]$ d'après
Corps, lemmes \ref{fields-lemma-multiplicativity-degrees} et
\ref{fields-lemma-degree-minimal-polynomial}.

\medskip\noindent
Nous affirmons que, pour tout $i = 0, 1, \ldots, n$, il existe un
isomorphisme
$$
\psi_i : k[x_1, \ldots, x_i] /(f_1, \ldots, f_i) \cong \kappa_i.
$$
Par construction, le composé
$k[x_1, \ldots, x_i] \to k[x_1, \ldots, x_n] \to \kappa$
a pour image $\kappa_i$, et $f_1, \ldots, f_i$ appartiennent
au noyau. Cela donne un homomorphisme surjectif.
Nous démontrons par récurrence que $\psi_i$ est injectif. C'est clair pour $i = 0$.
Supposons l'énoncé vrai pour $i$ et démontrons-le pour $i + 1$.
L'extension d'anneaux $k[x_1, \ldots, x_i]/(f_1, \ldots, f_i) \to
k[x_1, \ldots, x_{i + 1}]/(f_1, \ldots, f_{i + 1})$
est engendrée par $1$ élément sur un corps et soumise à une
équation irréductible. D'après la théorie élémentaire des corps,
$k[x_1, \ldots, x_{i + 1}]/(f_1, \ldots, f_{i + 1})$
est un corps, et $\psi_i$ est donc injectif.

\medskip\noindent
Cela implique que $\mathfrak m = (f_1, \ldots, f_n)$.
De plus, nous en concluons également que
$$
k[x_1, \ldots, x_n]/(f_1, \ldots, f_i)
\cong
\kappa_i[x_{i + 1}, \ldots, x_n].
$$
Ainsi, $(f_1, \ldots, f_i)$ est un idéal premier. Par conséquent,
$$
(0) \subset (f_1) \subset (f_1, f_2) \subset \ldots \subset
(f_1, \ldots, f_n) = \mathfrak m
$$
est une chaîne d'idéaux premiers de longueur $n$. Le lemme en résulte.
\end{proof}

\begin{proposition}
\label{proposition-finite-gl-dim-polynomial-ring}
Une algèbre polynomiale à $n$ variables sur un corps est un anneau régulier.
Elle est de dimension globale $n$. Toutes ses localisations en des idéaux maximaux
sont des anneaux locaux réguliers de dimension $n$.
\end{proposition}

\begin{proof}
D'après le lemme \ref{lemma-dim-affine-space},
toutes les localisations $k[x_1, \ldots, x_n]_{\mathfrak m}$
en des idéaux maximaux sont des anneaux locaux réguliers de dimension $n$. Ainsi,
nous concluons par le lemme \ref{lemma-finite-gl-dim-finite-dim-regular}.
\end{proof}

\begin{lemma}
\label{lemma-dimension-height-polynomial-ring}
Soit $k$ un corps.
Soit $\mathfrak p \subset \mathfrak q \subset k[x_1, \ldots, x_n]$
un couple d'idéaux premiers.
Toute chaîne maximale d'idéaux premiers entre $\mathfrak p$ et $\mathfrak q$
est de longueur $\text{hauteur}(\mathfrak q) - \text{hauteur}(\mathfrak p)$.
\end{lemma}

\begin{proof}
D'après la
proposition \ref{proposition-finite-gl-dim-polynomial-ring},
tout anneau local de $k[x_1, \ldots, x_n]$ est régulier.
Tous les anneaux locaux sont donc de Cohen-Macaulay, voir le
lemme \ref{lemma-regular-ring-CM}.
Les anneaux locaux aux idéaux maximaux sont de dimension $n$, donc
toute chaîne maximale d'idéaux premiers de $k[x_1, \ldots, x_n]$
est de longueur $n$, voir le
lemme \ref{lemma-maximal-chain-CM}.
Ainsi, toute chaîne maximale d'idéaux premiers entre $(0)$ et $\mathfrak p$
est de longueur $\text{hauteur}(\mathfrak p)$, voir le
lemme \ref{lemma-CM-dim-formula}
par exemple.
En rassemblant ces faits, nous obtenons l'assertion du lemme.
\end{proof}

\begin{lemma}
\label{lemma-dimension-spell-it-out}
Soit $k$ un corps.
Soit $S$ une $k$-algèbre de type fini qui est un anneau intègre.
Alors $\dim(S) = \dim(S_{\mathfrak m})$ pour tout idéal
maximal $\mathfrak m$ de $S$. Autrement dit, toute chaîne maximale
d'idéaux premiers est de longueur égale à la dimension de $S$.
\end{lemma}

\begin{proof}
Écrivons $S = k[x_1, \ldots, x_n]/\mathfrak p$.
D'après la proposition \ref{proposition-finite-gl-dim-polynomial-ring} et le
lemme \ref{lemma-dimension-height-polynomial-ring},
toutes les chaînes maximales d'idéaux premiers de $S$ (qui se terminent nécessairement
par un idéal maximal) sont de longueur $n - \text{hauteur}(\mathfrak p)$.
Ce nombre est donc la dimension de $S$ et de $S_{\mathfrak m}$
pour tout idéal maximal $\mathfrak m$ de $S$.
\end{proof}

\noindent
Rappelons que nous avons défini la
dimension $\dim_x(X)$ d'un espace topologique $X$ en un point $x$
dans Topologie, définition \ref{topology-definition-Krull}.

\begin{lemma}
\label{lemma-dimension-at-a-point-finite-type-over-field}
Soit $k$ un corps.
Soit $S$ une $k$-algèbre de type fini.
Soit $X = \Spec(S)$.
Soit $\mathfrak p \subset S$ un idéal premier et soit
$x \in X$ le point correspondant.
Les nombres suivants sont égaux :
\begin{enumerate}
\item $\dim_x(X)$,
\item $\max \dim(Z)$, où le maximum porte sur les
composantes irréductibles $Z$ de $X$ qui passent par $x$, et
\item $\min \dim(S_{\mathfrak m})$, où le minimum
porte sur les idéaux maximaux $\mathfrak m$ tels que
$\mathfrak p \subset \mathfrak m$.
\end{enumerate}
\end{lemma}

\begin{proof}
Soit $X = \bigcup_{i \in I} Z_i$ la décomposition de $X$ en
ses composantes irréductibles. Elles sont en nombre fini
(voir les
lemmes \ref{lemma-obvious-Noetherian} et \ref{lemma-Noetherian-topology}).
Soit $I' = \{i \mid x \in Z_i\}$, et soit
$T = \bigcup_{i \not \in I'} Z_i$. Alors $U = X \setminus T$
est un ouvert de $X$ qui contient le point $x$.
Le nombre figurant en (2) vaut $\max_{i \in I'} \dim(Z_i)$.
Pour tout ouvert $W \subset U$ tel que $x \in W$,
les composantes irréductibles de $W$ sont les ensembles irréductibles
$W_i = Z_i \cap W$ pour $i \in I'$, et $x$ appartient
à chacun d'eux.
Remarquons que chaque $W_i$, $i \in I'$, contient un point fermé, car
$X$ est jacobsonien, voir la section \ref{section-ring-jacobson}.
Puisque $W_i \subset Z_i$, nous avons $\dim(W_i) \leq \dim(Z_i)$.
L'existence d'un point fermé implique, par le lemme
\ref{lemma-dimension-spell-it-out}, qu'il existe une chaîne de
fermés irréductibles de longueur égale à $\dim(Z_i)$ dans l'ouvert $W_i$.
Ainsi, $\dim(W_i) = \dim(Z_i)$ pour tout $i \in I'$. Par conséquent, $\dim(W)$
est égal au nombre figurant en (2). Cela démontre que (1) $ = $ (2).

\medskip\noindent
Soit $\mathfrak m \supset \mathfrak p$ un idéal maximal quelconque
contenant $\mathfrak p$. Soit $x_0 \in X$ le point
correspondant. Tout d'abord, $x_0$ appartient à toutes les
composantes irréductibles $Z_i$, $i \in I'$. Notons $\mathfrak q_i$
les idéaux premiers minimaux de $S$ correspondant aux
composantes irréductibles $Z_i$. Pour chaque $i$ tel que
$x_0 \in Z_i$ (ce qui équivaut à $\mathfrak m \supset \mathfrak q_i$),
nous avons une surjection
$$
S_{\mathfrak m} \longrightarrow
S_\mathfrak m/\mathfrak q_i S_\mathfrak m =(S/\mathfrak q_i)_{\mathfrak m}
$$
De plus, les idéaux premiers $\mathfrak q_i S_\mathfrak m$ ainsi obtenus
épuisent les idéaux premiers
minimaux de l'anneau local noethérien $S_{\mathfrak m}$, voir le
lemme \ref{lemma-irreducible-components-containing-x}.
Nous concluons, en utilisant le lemme \ref{lemma-dimension-spell-it-out},
que la dimension de $S_{\mathfrak m}$ est le
maximum des dimensions des $Z_i$ qui passent par $x_0$.
Pour terminer la démonstration du lemme, il suffit de montrer que
nous pouvons choisir $x_0$ de sorte que $x_0 \in Z_i \Rightarrow i \in I'$.
Comme $S$ est jacobsonien (ainsi que nous l'avons vu ci-dessus),
il suffit de montrer que $V(\mathfrak p) \setminus T$
(où $T$ est comme ci-dessus) est non vide. Cela est clair, car il
contient le point $x$ (c'est-à-dire $\mathfrak p$).
\end{proof}

\begin{lemma}
\label{lemma-dimension-closed-point-finite-type-field}
Soit $k$ un corps.
Soit $S$ une $k$-algèbre de type fini.
Soit $X = \Spec(S)$.
Soit $\mathfrak m \subset S$ un idéal maximal et soit
$x \in X$ le point fermé associé.
Alors $\dim_x(X) = \dim(S_{\mathfrak m})$.
\end{lemma}

\begin{proof}
C'est un cas particulier du
lemme \ref{lemma-dimension-at-a-point-finite-type-over-field}.
\end{proof}

\begin{lemma}
\label{lemma-disjoint-decomposition-CM-algebra}
Soit $k$ un corps.
Soit $S$ une $k$-algèbre de type fini.
Supposons que $S$ est de Cohen-Macaulay.
Alors $\Spec(S) = \coprod T_d$ est une union disjointe finie de
parties ouvertes et fermées $T_d$, où $T_d$ est équidimensionnelle
(voir Topologie, définition \ref{topology-definition-equidimensional})
de dimension $d$. De manière équivalente, $S$ est un produit d'anneaux
$S_d$, $d = 0, \ldots, \dim(S)$, tels que tout idéal maximal
$\mathfrak m$ de $S_d$ est de hauteur $d$.
\end{lemma}

\begin{proof}
L'équivalence des deux énoncés résulte du
lemme \ref{lemma-disjoint-implies-product}.
Soit $\mathfrak m \subset S$ un idéal maximal.
Toute chaîne maximale d'idéaux premiers de $S_{\mathfrak m}$ est
de même longueur, égale à $\dim(S_{\mathfrak m})$, voir le
lemme \ref{lemma-maximal-chain-CM}. Ainsi, les composantes irréductibles
qui passent par le point correspondant à $\mathfrak m$
sont toutes de dimension égale à $\dim(S_{\mathfrak m})$, voir le
lemme \ref{lemma-dimension-spell-it-out}.
Puisque $\Spec(S)$ est un espace topologique jacobsonien,
l'intersection
de deux quelconques de ses composantes irréductibles contient un point fermé si elle est non vide,
voir les
lemmes \ref{lemma-finite-type-field-Jacobson} et
\ref{lemma-jacobson}.
Nous avons donc montré que deux composantes irréductibles quelconques
qui se rencontrent sont de même dimension. Le lemme en résulte
facilement, ainsi que du fait que $\Spec(S)$
n'a qu'un nombre fini de composantes irréductibles (voir les
lemmes \ref{lemma-obvious-Noetherian} et \ref{lemma-Noetherian-topology}).
\end{proof}
















\section{Normalisation de Noether}
\label{section-Noether-normalization}

\noindent
Dans cette section, nous démontrons des variantes du lemme de normalisation de Noether.
L'ingrédient essentiel que nous utiliserons est contenu dans les deux lemmes suivants.

\begin{lemma}
\label{lemma-helper}
Soit $n \in \mathbf{N}$.
Soit $N$ un ensemble fini non vide
de multi-indices $\nu = (\nu_1, \ldots, \nu_n)$.
Pour $e = (e_1, \ldots, e_n)$, nous posons $e \cdot \nu = \sum e_i\nu_i$.
Alors, pour $e_1 \gg e_2 \gg \ldots \gg e_{n-1} \gg e_n$, nous avons :
si $\nu, \nu' \in N$, alors
$$
(e \cdot \nu = e \cdot \nu')
\Leftrightarrow
(\nu = \nu')
$$
\end{lemma}

\begin{proof}
Écrivons $N = \{\nu_j\}$ avec $\nu_j = (\nu_{j1}, \ldots, \nu_{jn})$.
Soit $A_i = \max_j \nu_{ji} - \min_j \nu_{ji}$. Si, pour tout $i$, nous avons
$e_{i - 1} > A_ie_i + A_{i + 1}e_{i + 1} + \ldots + A_ne_n$, alors
le lemme est vrai. Supposons en effet que $e \cdot (\nu - \nu') = 0$. Alors, pour
$n \ge 2$,
$$
e_1(\nu_1 - \nu'_1) = \sum\nolimits_{i = 2}^n e_i(\nu'_i - \nu_i).
$$
Nous pouvons supposer que $(\nu_1 - \nu'_1) \ge 0$. Si $(\nu_1 - \nu'_1) > 0$, alors
$$
e_1(\nu_1 - \nu'_1) \ge e_1 >
A_2e_2 + \ldots + A_ne_n \ge
\sum\nolimits_{i = 2}^n e_i|\nu'_i - \nu_i| \ge
\sum\nolimits_{i = 2}^n e_i(\nu'_i - \nu_i).
$$
Cette contradiction implique que $\nu'_1 = \nu_1$. Par
récurrence, $\nu'_i = \nu_i$ pour $2 \le i \le n$.
\end{proof}

\begin{lemma}
\label{lemma-helper-polynomial}
Soit $R$ un anneau. Soit $g \in R[x_1, \ldots, x_n]$ un élément
qui n'est pas constant, c'est-à-dire $g \not \in R$.
Pour $e_1 \gg e_2 \gg \ldots \gg e_{n-1} \gg e_n = 1$, le polynôme
$$
g(x_1 + x_n^{e_1}, x_2 + x_n^{e_2}, \ldots, x_{n - 1} + x_n^{e_{n - 1}}, x_n)
=
ax_n^d + \text{termes de degré inférieur en }x_n
$$
où $d > 0$ et $a \in R$ est l'un des coefficients non nuls de $g$.
\end{lemma}

\begin{proof}
Écrivons $g = \sum_{\nu \in N} a_\nu x^\nu$, où les $a_\nu \in R$ sont non nuls.
Ici, $N$ est un ensemble fini de multi-indices comme dans le
lemme \ref{lemma-helper},
et $x^\nu = x_1^{\nu_1} \ldots x_n^{\nu_n}$.
Remarquons que le terme dominant de
$$
(x_1 + x_n^{e_1})^{\nu_1} \ldots (x_{n-1} + x_n^{e_{n-1}})^{\nu_{n-1}}
x_n^{\nu_n}
\quad\text{est}\quad
x_n^{e_1\nu_1 + \ldots + e_{n-1}\nu_{n-1} + \nu_n}.
$$
Le lemme résulte donc du
lemme \ref{lemma-helper},
qui garantit qu'il existe exactement un terme non nul $a_\nu x^\nu$ de $g$
qui donne naissance au terme dominant de
$g(x_1 + x_n^{e_1}, x_2 + x_n^{e_2}, \ldots, x_{n - 1} + x_n^{e_{n - 1}},
x_n)$, c'est-à-dire $a = a_\nu$ pour l'unique $\nu \in N$ tel que
$e \cdot \nu$ soit maximal.
\end{proof}

\begin{lemma}
\label{lemma-one-relation}
Soit $k$ un corps.
Soit $S = k[x_1, \ldots, x_n]/I$ pour un idéal propre $I$.
Si $I \not = 0$, alors il existe $y_1, \ldots, y_{n-1} \in k[x_1, \ldots, x_n]$
tels que $S$ est fini sur $k[y_1, \ldots, y_{n-1}]$. De plus, nous pouvons
choisir les $y_i$ dans la sous-$\mathbf{Z}$-algèbre de $k[x_1, \ldots, x_n]$
engendrée par $x_1, \ldots, x_n$.
\end{lemma}

\begin{proof}
Choisissons $f \in I$, $f\not = 0$. Il suffit de démontrer le lemme
pour $k[x_1, \ldots, x_n]/(f)$, puisque $S$ est un quotient de cet anneau.
Nous prendrons $y_i = x_i - x_n^{e_i}$, $i = 1, \ldots, n-1$,
pour des entiers $e_i$ convenables. Quand cela fonctionne-t-il ? Il suffit
de montrer que $\overline{x_n} \in k[x_1, \ldots, x_n]/(f)$
est entier sur l'anneau $k[y_1, \ldots, y_{n-1}]$. L'
équation de $\overline{x_n}$ sur cet anneau est
$$
f(y_1 + x_n^{e_1}, \ldots, y_{n-1} + x_n^{e_{n-1}}, x_n) = 0.
$$
Nous avons donc terminé si nous pouvons montrer qu'il existe des entiers $e_i$ tels
que le coefficient dominant par rapport à $x_n$ de l'équation
ci-dessus est un élément non nul de $k$. On peut l'obtenir, par exemple,
en choisissant $e_1 \gg e_2 \gg \ldots \gg e_{n-1}$, voir le
lemme \ref{lemma-helper-polynomial}.
\end{proof}

\begin{lemma}
\label{lemma-Noether-normalization}
\begin{slogan}
Normalisation de Noether
\end{slogan}
Soit $k$ un corps. Soit $S = k[x_1, \ldots, x_n]/I$ pour un idéal $I$.
Si $I \neq (1)$, il existe $r\geq 0$, et
$y_1, \ldots, y_r \in k[x_1, \ldots, x_n]$
tels que (a) l'homomorphisme $k[y_1, \ldots, y_r] \to S$ est injectif,
où la source est l'anneau de polynômes en $y_1, \ldots, y_r$,
et (b) l'homomorphisme $k[y_1, \ldots, y_r] \to S$ est fini.
Dans ce cas, l'entier $r$ est la dimension de $S$.
De plus, nous pouvons choisir les $y_i$ dans la
sous-$\mathbf{Z}$-algèbre de $k[x_1, \ldots, x_n]$
engendrée par $x_1, \ldots, x_n$.
\end{lemma}

\begin{proof}
Raisonnons par récurrence sur $n$, le cas $n = 0$ étant trivial.
Si $I = 0$, prenons $r = n$ et $y_i = x_i$.
Si $I \not = 0$, choisissons $y_1, \ldots, y_{n-1}$
comme dans le lemme \ref{lemma-one-relation}. Soit
$S' \subset S$ le sous-anneau engendré par
les images des $y_i$. Par récurrence, nous pouvons
choisir $r$ et $z_1, \ldots, z_r \in k[y_1, \ldots, y_{n-1}]$
de sorte que (a), (b) soient vérifiées pour $k[z_1, \ldots, z_r]
\to S'$. Puisque $S' \to S$ est injectif et fini,
nous voyons que (a), (b) sont vérifiées pour $k[z_1, \ldots, z_r]
\to S$. La dernière assertion résulte du lemme
\ref{lemma-integral-sub-dim-equal}.
\end{proof}

\begin{lemma}
\label{lemma-Noether-normalization-at-point}
Soit $k$ un corps.
Soit $S$ une $k$-algèbre de type fini et notons $X = \Spec(S)$.
Soit $\mathfrak q$ un idéal premier de $S$, et soit $x \in X$ le point
correspondant. Il existe $g \in S$, $g \not \in \mathfrak q$,
tel que $\dim(S_g) = \dim_x(X) =: d$ et tel qu'
il existe un homomorphisme fini injectif $k[y_1, \ldots, y_d] \to S_g$.
\end{lemma}

\begin{proof}
Remarquons que, par définition, $\dim_x(X)$ est le minimum
des dimensions des $S_g$ pour $g \in S$, $g \not \in \mathfrak q$, c'est-à-dire
que le minimum est atteint. Le lemme résulte donc du
lemme \ref{lemma-Noether-normalization}.
\end{proof}

\begin{lemma}
\label{lemma-refined-Noether-normalization}
Soit $k$ un corps. Soit $\mathfrak q \subset k[x_1, \ldots, x_n]$
un idéal premier. Posons $r = \text{trdeg}_k\ \kappa(\mathfrak q)$.
Alors il existe un homomorphisme fini d'anneaux
$\varphi : k[y_1, \ldots, y_n] \to k[x_1, \ldots, x_n]$ tel
que $\varphi^{-1}(\mathfrak q) = (y_{r + 1}, \ldots, y_n)$.
\end{lemma}

\begin{proof}
Raisonnons par récurrence sur $n$. Le cas $n = 0$ est clair. Supposons $n > 0$.
Si $r = n$, alors $\mathfrak q = (0)$ et le résultat est clair.
Choisissons un élément non nul $f \in \mathfrak q$. Bien sûr, $f$ n'est pas constant.
Après avoir appliqué un automorphisme de la forme
$$
k[x_1, \ldots, x_n] \longrightarrow k[x_1, \ldots, x_n],
\quad
x_n \mapsto x_n,
\quad
x_i \mapsto x_i + x_n^{e_i}\ (i < n)
$$
nous pouvons supposer que $f$ est unitaire en $x_n$ sur $k[x_1, \ldots, x_n]$, voir le
lemme \ref{lemma-helper-polynomial}. Ainsi, l'homomorphisme d'anneaux
$$
k[y_1, \ldots, y_n] \longrightarrow k[x_1, \ldots, x_n],
\quad
y_n \mapsto f,
\quad
y_i \mapsto x_i\ (i < n)
$$
est fini. De plus, $y_n \in \mathfrak q \cap k[y_1, \ldots, y_n]$ par
construction. Ainsi,
$\mathfrak q \cap k[y_1, \ldots, y_n] = \mathfrak pk[y_1, \ldots, y_n] + (y_n)$,
où $\mathfrak p \subset k[y_1, \ldots, y_{n - 1}]$ est un idéal premier.
Remarquons que $\kappa(\mathfrak p) \subset \kappa(\mathfrak q)$ est finie, et
donc $r = \text{trdeg}_k\ \kappa(\mathfrak p)$.
Appliquons l'hypothèse de récurrence au couple
$(k[y_1, \ldots, y_{n - 1}], \mathfrak p)$ ; nous obtenons un homomorphisme fini d'anneaux
$k[z_1, \ldots, z_{n - 1}] \to k[y_1, \ldots, y_{n - 1}]$ tel que
$\mathfrak p \cap k[z_1, \ldots, z_{n - 1}] = (z_{r + 1}, \ldots, z_{n - 1})$.
Nous prolongeons l'homomorphisme d'anneaux
$k[z_1, \ldots, z_{n - 1}] \to k[y_1, \ldots, y_{n - 1}]$
en un homomorphisme d'anneaux
$k[z_1, \ldots, z_n] \to k[y_1, \ldots, y_n]$
en envoyant $z_n$ sur $y_n$.
Le composé des homomorphismes d'anneaux
$$
k[z_1, \ldots, z_n] \to k[y_1, \ldots, y_n] \to k[x_1, \ldots, x_n]
$$
résout le problème.
\end{proof}

\begin{lemma}
\label{lemma-Noether-normalization-over-a-domain}
Soit $R \to S$ un homomorphisme injectif d'anneaux de type fini. Supposons que $R$ est intègre.
Alors il existe un entier $d$ et une factorisation
$$
R \to R[y_1, \ldots, y_d] \to S' \to S
$$
par des homomorphismes injectifs tels que $S'$ soit fini sur $R[y_1, \ldots, y_d]$
et que $S'_f \cong S_f$ pour un élément non nul $f \in R$.
\end{lemma}

\begin{proof}
Choisissons $x_1, \ldots, x_n \in S$ qui engendrent $S$ sur $R$.
Soit $K$ le corps des fractions de $R$ et $S_K = S \otimes_R K$.
D'après le lemme \ref{lemma-Noether-normalization},
nous pouvons trouver $y_1, \ldots, y_d \in S$ tels que $K[y_1, \ldots, y_d] \to S_K$
soit un homomorphisme fini injectif. Remarquons que $y_i \in S$, car nous pouvons choisir les
$y_j$ dans la $\mathbf{Z}$-algèbre engendrée par $x_1, \ldots, x_n$.
Comme un homomorphisme fini d'anneaux est entier (voir le
lemme \ref{lemma-finite-is-integral}),
nous pouvons trouver des polynômes unitaires $P_i \in K[y_1, \ldots, y_d][T]$ tels que
$P_i(x_i) = 0$ dans $S_K$. Soit $f \in R$ un élément non nul tel que
$fP_i \in R[y_1, \ldots, y_d][T]$ pour tout $i$. Alors $fP_i(x_i)$
s'envoie sur zéro dans $S_K$. Par conséquent, après avoir remplacé $f$ par un autre
élément non nul de $R$, nous pouvons aussi supposer que $fP_i(x_i)$ est nul dans $S$.
Posons $x_i' = fx_i$ et soit $S' \subset S$ la sous-$R$-algèbre engendrée par
$y_1, \ldots, y_d$ et $x'_1, \ldots, x'_n$. Remarquons que $x'_i$ est entier sur
$R[y_1, \ldots, y_d]$, car nous avons $Q_i(x_i') = 0$, où
$Q_i = f^{\deg_T(P_i)}P_i(T/f)$ est un polynôme unitaire en
$T$ à coefficients dans $R[y_1, \ldots, y_d]$ par notre choix de $f$.
Ainsi, $R[y_1, \ldots, y_d] \subset S'$ est fini d'après le
lemme \ref{lemma-characterize-finite-in-terms-of-integral}.
Puisque $S' \subset S$, nous avons $S'_f \subset S_f$ (la localisation est exacte).
D'autre part, les éléments $x_i = x'_i/f$ de $S'_f$ engendrent $S_f$
sur $R_f$, et donc $S'_f \to S_f$ est surjectif. Par conséquent,
$S'_f \cong S_f$, ce qui termine la démonstration.
\end{proof}





\section{Dimension des algèbres de type fini sur les corps, reprise}
\label{section-dimension-finite-type-algebras-reprise}

\noindent
Cette section est une suite de la
section \ref{section-dimension-finite-type-algebras}.
Dans cette section, nous établissons le lien entre la dimension et le
degré de transcendance sur le corps de base pour les anneaux intègres de type fini
sur un corps.

\begin{lemma}
\label{lemma-dimension-prime-polynomial-ring}
Soit $k$ un corps.
Soit $S$ une $k$-algèbre de type fini qui est un anneau intègre.
Soit $K$ le corps des fractions de $S$.
Soit $r = \text{trdeg}(K/k)$ le degré de transcendance de $K$ sur $k$.
Alors $\dim(S) = r$. De plus, l'anneau local de $S$ en tout idéal
maximal est de dimension $r$.
\end{lemma}

\begin{proof}
Nous pouvons écrire $S = k[x_1, \ldots, x_n]/\mathfrak p$.
D'après le lemme \ref{lemma-dimension-height-polynomial-ring},
tous les anneaux locaux de $S$ en des idéaux maximaux ont la même dimension.
Appliquons le lemme \ref{lemma-Noether-normalization}.
Nous obtenons un homomorphisme fini injectif d'anneaux
$$
k[y_1, \ldots, y_d] \to S
$$
avec $d = \dim(S)$. Il est clair que $k(y_1, \ldots, y_d) \subset K$
est une extension finie, ce qui termine la démonstration.
\end{proof}

\begin{lemma}
\label{lemma-tr-deg-specialization}
Soit $k$ un corps. Soit $S$ une $k$-algèbre de type fini.
Soient $\mathfrak q \subset \mathfrak q' \subset S$ des idéaux
premiers distincts. Alors
$\text{trdeg}_k\ \kappa(\mathfrak q') < \text{trdeg}_k\ \kappa(\mathfrak q)$.
\end{lemma}

\begin{proof}
D'après le
lemme \ref{lemma-dimension-prime-polynomial-ring},
nous avons $\dim V(\mathfrak q) = \text{trdeg}_k\ \kappa(\mathfrak q)$,
et de même pour $\mathfrak q'$. Le résultat en découle, car l'inclusion stricte
$V(\mathfrak q') \subset V(\mathfrak q)$
implique une inégalité stricte entre les dimensions.
\end{proof}

\noindent
Le lemme suivant généralise le
lemme \ref{lemma-dimension-closed-point-finite-type-field}.

\begin{lemma}
\label{lemma-dimension-at-a-point-finite-type-field}
Soit $k$ un corps.
Soit $S$ une $k$-algèbre de type fini.
Soit $X = \Spec(S)$.
Soit $\mathfrak p \subset S$ un idéal premier,
et soit $x \in X$ le point correspondant.
Alors nous avons
$$
\dim_x(X) = \dim(S_{\mathfrak p}) + \text{trdeg}_k\ \kappa(\mathfrak p).
$$
\end{lemma}

\begin{proof}
D'après le lemme \ref{lemma-dimension-prime-polynomial-ring}, nous savons que
$r = \text{trdeg}_k\ \kappa(\mathfrak p)$ est égal à la
dimension de $V(\mathfrak p)$.
Choisissons une chaîne maximale quelconque d'idéaux premiers
$\mathfrak p \subset \mathfrak p_1 \subset \ldots \subset \mathfrak p_r$
qui commence par $\mathfrak p$ dans $S$.
Elle est de longueur $r$ d'après le lemme \ref{lemma-dimension-spell-it-out}.
Soient $\mathfrak q_j$, $j \in J$, les idéaux premiers
minimaux de $S$ qui sont contenus dans $\mathfrak p$.
Ils correspondent $1-1$ aux idéaux premiers minimaux de $S_{\mathfrak p}$
par la règle $\mathfrak q_j \mapsto \mathfrak q_jS_{\mathfrak p}$.
D'après le lemme \ref{lemma-dimension-at-a-point-finite-type-over-field},
nous savons que $\dim_x(X)$ est égal
au maximum des dimensions des anneaux $S/\mathfrak q_j$.
Pour chaque $j$, choisissons une chaîne maximale d'idéaux premiers
$\mathfrak q_j \subset \mathfrak p'_1 \subset \ldots \subset \mathfrak p'_{s(j)}
= \mathfrak p$.
Alors $\dim(S_{\mathfrak p}) = \max_{j \in J} s(j)$.
Or chaque chaîne
$$
\mathfrak q_j \subset \mathfrak p'_1 \subset \ldots \subset
\mathfrak p'_{s(j)} = \mathfrak p \subset
\mathfrak p_1 \subset \ldots \subset \mathfrak p_r
$$
est une chaîne maximale dans $S/\mathfrak q_j$ et, d'après ce qui précède,
nous avons
$\dim_x(X) = \max_{j \in J} r + s(j)$.
Le lemme en résulte.
\end{proof}

\noindent
Le lemme suivant affirme que la codimension d'un spectre de type fini
dans un autre est la différence des hauteurs.

\begin{lemma}
\label{lemma-codimension}
Soit $k$ un corps.
Soit $S' \to S$ une surjection de $k$-algèbres de type fini.
Soit $\mathfrak p \subset S$ un idéal premier,
et soit $\mathfrak p'$ l'idéal premier correspondant de $S'$.
Soient $X = \Spec(S)$, resp.\ $X' = \Spec(S')$,
et soient $x \in X$, resp. $x'\in X'$, les points correspondant
à $\mathfrak p$, resp.\ $\mathfrak p'$.
Alors
$$
\dim_{x'} X' - \dim_x X =
\text{hauteur}(\mathfrak p') - \text{hauteur}(\mathfrak p).
$$
\end{lemma}

\begin{proof}
C'est une conséquence immédiate du lemme \ref{lemma-dimension-at-a-point-finite-type-field}.
\end{proof}

\begin{lemma}
\label{lemma-dimension-preserved-field-extension}
Soit $k$ un corps.
Soit $S$ une $k$-algèbre de type fini.
Soit $K/k$ une extension de corps.
Alors $\dim(S) = \dim(K \otimes_k S)$.
\end{lemma}

\begin{proof}
D'après le lemme \ref{lemma-Noether-normalization},
il existe un homomorphisme fini injectif
$k[y_1, \ldots, y_d] \to S$ avec $d = \dim(S)$.
Puisque $K$ est plat sur $k$, nous obtenons aussi un homomorphisme fini injectif
$K[y_1, \ldots, y_d] \to K \otimes_k S$.
Le résultat découle du lemme \ref{lemma-integral-sub-dim-equal}.
\end{proof}

\begin{lemma}
\label{lemma-dimension-at-a-point-preserved-field-extension}
Soit $k$ un corps.
Soit $S$ une $k$-algèbre de type fini.
Posons $X = \Spec(S)$.
Soit $K/k$ une extension de corps.
Posons $S_K = K \otimes_k S$, et $X_K = \Spec(S_K)$.
Soit $\mathfrak q \subset S$ un idéal premier correspondant à $x \in X$,
et soit $\mathfrak q_K \subset S_K$ un idéal premier correspondant
à $x_K \in X_K$ au-dessus de $\mathfrak q$.
Alors $\dim_x X = \dim_{x_K} X_K$.
\end{lemma}

\begin{proof}
Choisissons une présentation $S = k[x_1, \ldots, x_n]/I$.
Cela donne une présentation
$K \otimes_k S = K[x_1, \ldots, x_n]/(K \otimes_k I)$.
Soit $\mathfrak q_K' \subset K[x_1, \ldots, x_n]$,
resp.\ $\mathfrak q' \subset k[x_1, \ldots, x_n]$, les
idéaux premiers correspondants. Considérons le diagramme commutatif
suivant d'anneaux locaux noethériens :
$$
\xymatrix{
K[x_1, \ldots, x_n]_{\mathfrak q_K'} \ar[r] &
(K \otimes_k S)_{\mathfrak q_K} \\
k[x_1, \ldots, x_n]_{\mathfrak q'} \ar[r] \ar[u] &
S_{\mathfrak q} \ar[u]
}
$$
Les deux flèches verticales sont plates, car ce sont des localisations des
homomorphismes plats d'anneaux $S \to S_K$ et
$k[x_1, \ldots, x_n] \to K[x_1, \ldots, x_n]$.
De plus, les flèches verticales ont les mêmes anneaux fibres.
Par conséquent, il résulte du
lemme \ref{lemma-dimension-base-fibre-equals-total} que
$\text{hauteur}(\mathfrak q') - \text{hauteur}(\mathfrak q)
= \text{hauteur}(\mathfrak q_K') - \text{hauteur}(\mathfrak q_K)$.
Notons $x' \in X' = \Spec(k[x_1, \ldots, x_n])$
et $x'_K \in X'_K = \Spec(K[x_1, \ldots, x_n])$
les points correspondant à $\mathfrak q'$ et
$\mathfrak q_K'$. D'après le lemme \ref{lemma-codimension} et ce que nous avons montré
ci-dessus, nous avons
\begin{eqnarray*}
n - \dim_x X & = & \dim_{x'} X' - \dim_x X \\
& = & \text{hauteur}(\mathfrak q') - \text{hauteur}(\mathfrak q) \\
& = & \text{hauteur}(\mathfrak q_K') - \text{hauteur}(\mathfrak q_K) \\
& = & \dim_{x'_K} X'_K - \dim_{x_K} X_K \\
& = & n - \dim_{x_K} X_K
\end{eqnarray*}
et le lemme en résulte.
\end{proof}

\begin{lemma}
\label{lemma-inequalities-under-field-extension}
Soit $k$ un corps. Soit $S$ une $k$-algèbre de type fini.
Soit $K/k$ une extension de corps. Posons $S_K = K \otimes_k S$.
Soit $\mathfrak q \subset S$ un idéal premier et soit $\mathfrak q_K \subset S_K$
un idéal premier au-dessus de $\mathfrak q$. Alors
$$
\dim (S_K \otimes_S \kappa(\mathfrak q))_{\mathfrak q_K} =
\dim (S_K)_{\mathfrak q_K} - \dim S_\mathfrak q =
\text{trdeg}_k \kappa(\mathfrak q) - \text{trdeg}_K \kappa(\mathfrak q_K)
$$
De plus, étant donné $\mathfrak q$, nous pouvons toujours choisir $\mathfrak q_K$ de sorte
que le nombre ci-dessus soit nul.
\end{lemma}

\begin{proof}
Observons que $S_\mathfrak q \to (S_K)_{\mathfrak q_K}$ est un homomorphisme
local plat d'anneaux locaux noethériens, de fibre spéciale
$(S_K \otimes_S \kappa(\mathfrak q))_{\mathfrak q_K}$. La première
égalité résulte donc du lemme \ref{lemma-dimension-base-fibre-equals-total}.
La seconde égalité résulte du fait que nous avons
$\dim_x X = \dim_{x_K} X_K$, avec les notations du
lemme \ref{lemma-dimension-at-a-point-preserved-field-extension},
et que nous avons
$\dim_x X = \dim S_\mathfrak q + \text{trdeg}_k \kappa(\mathfrak q)$
d'après le lemme \ref{lemma-dimension-at-a-point-finite-type-field},
et de même pour $\dim_{x_K} X_K$.
Si nous choisissons $\mathfrak q_K$ minimal au-dessus de $\mathfrak q S_K$, alors
la dimension de l'anneau fibre sera nulle.
\end{proof}







\section{Dimension des algèbres graduées sur un corps}
\label{section-dimension-graded}

\noindent
Voici un résultat élémentaire.

\begin{lemma}
\label{lemma-dimension-graded}
Soit $k$ un corps. Soit $S$ une $k$-algèbre graduée engendrée sur $k$
par un nombre fini d'éléments de degré $1$.
Supposons $S_0 = k$. Soit $P(T) \in \mathbf{Q}[T]$ le polynôme
tel que $\dim(S_d) = P(d)$ pour tout $d \gg 0$. Voir la
proposition \ref{proposition-graded-hilbert-polynomial}.
Alors
\begin{enumerate}
\item L'idéal non pertinent $S_{+}$ est un idéal maximal $\mathfrak m$.
\item Tout idéal premier minimal de $S$ est un idéal homogène et est contenu
dans $S_{+} = \mathfrak m$.
\item Nous avons $\dim(S) = \deg(P) + 1 = \dim_x\Spec(S)$
(avec la convention que $\deg(0) = -1$),
où $x$ est le point correspondant à l'idéal maximal
$S_{+} = \mathfrak m$.
\item La fonction de Hilbert de l'anneau local $R = S_{\mathfrak m}$
est égale à la fonction de Hilbert de $S$.
\end{enumerate}
\end{lemma}

\begin{proof}
La première assertion est évidente.
La seconde résulte du lemme \ref{lemma-graded-ring-minimal-prime}.
D'après (2), toute composante irréductible passe par $x$.
Ainsi, nous avons $\dim(S) = \dim_x\Spec(S) = \dim(S_\mathfrak m)$
d'après le lemme \ref{lemma-dimension-at-a-point-finite-type-over-field}. Puisque
$\mathfrak m^d/\mathfrak m^{d + 1} \cong
\mathfrak m^dS_\mathfrak m/\mathfrak m^{d + 1}S_\mathfrak m$,
nous voyons que la fonction de Hilbert de l'anneau local $S_\mathfrak m$
est égale à la
fonction de Hilbert de $S$, ce qui donne (4). Nous déduisons la dernière égalité
de (3) de la proposition \ref{proposition-dimension}.
\end{proof}













\section{Platitude générique}
\label{section-generic-flatness}

\noindent
En substance, cela affirme qu'une algèbre de type fini sur un anneau intègre devient
plate après inversion d'un seul élément de l'anneau intègre.
Il existe plusieurs versions de ce résultat (par ordre croissant de
force).

\begin{lemma}
\label{lemma-generic-flatness-Noetherian}
Soit $R \to S$ un homomorphisme d'anneaux.
Soit $M$ un $S$-module.
Supposons que
\begin{enumerate}
\item $R$ est noethérien,
\item $R$ est intègre,
\item $R \to S$ est de type fini, et
\item $M$ est un $S$-module de type fini.
\end{enumerate}
Alors il existe un élément non nul $f \in R$ tel que
$M_f$ soit un $R_f$-module libre.
\end{lemma}

\begin{proof}
Soit $K$ le corps des fractions de $R$. Posons $S_K = K \otimes_R S$. C'
est une algèbre de type fini sur $K$. Nous allons raisonner par récurrence sur
$d = \dim(S_K)$ (qui est fini, par exemple par la normalisation de Noether, voir la
section \ref{section-Noether-normalization}).
Fixons $d \geq 0$.
Supposons que le lemme soit vrai dans tous les cas où $\dim(S_K) < d$.

\medskip\noindent
Supposons donnés $R \to S$ et $M$ comme dans le lemme, avec $\dim(S_K) = d$. D'après le
lemme \ref{lemma-filter-Noetherian-module},
il existe une filtration
$0 \subset M_1 \subset M_2 \subset \ldots \subset M_n = M$
telle que $M_i/M_{i - 1}$ soit isomorphe à $S/\mathfrak q$
pour un idéal premier $\mathfrak q$ de $S$. Remarquons que
$\dim((S/\mathfrak q)_K) \leq \dim(S_K)$. Remarquons aussi qu'une extension de
modules libres est libre (voir la notion élémentaire \ref{item-extension-free}).
Nous pouvons donc supposer que $M = S$ et que $S$ est un anneau intègre de type fini sur $R$.

\medskip\noindent
Si $R \to S$ a un noyau non trivial, choisissons un élément non nul $f \in R$ dans
ce noyau. Dans ce cas, $S_f = 0$ et le lemme est vrai. (Il s'agit en réalité
du cas $d = -1$, qui amorce la récurrence.) Nous pouvons donc
supposer que $R \to S$ est une extension de type fini d'anneaux intègres noethériens.

\medskip\noindent
Appliquons le lemme \ref{lemma-Noether-normalization-over-a-domain}
et remplaçons $R$ par $R_f$ (où $f$ est comme dans le lemme) pour obtenir une
factorisation
$$
R \subset R[y_1, \ldots, y_d] \subset S
$$
où la seconde extension est finie. Choisissons $z_1, \ldots, z_r \in S$ qui
forment une base du corps des fractions de $S$ sur le corps des fractions de
$R[y_1, \ldots, y_d]$. Cela donne une suite exacte courte
$$
0 \to
R[y_1, \ldots, y_d]^{\oplus r} \xrightarrow{(z_1, \ldots, z_r)}
S \to N \to 0
$$
Par construction, $N$ est un $R[y_1, \ldots, y_d]$-module fini dont le
support ne contient pas le point générique $(0)$ de
$\Spec(R[y_1, \ldots, y_d])$. D'après le
lemme \ref{lemma-support-closed},
il existe un élément non nul $g \in R[y_1, \ldots, y_d]$ tel que
$g$ annule $N$ ; nous pouvons donc voir $N$ comme un module fini sur
$S' = R[y_1, \ldots, y_d]/(g)$. Puisque $\dim(S'_K) < d$, il existe par récurrence
un élément non nul $f \in R$ tel que $N_f$ soit un
$R_f$-module libre. Puisque
$(R[y_1, \ldots, y_d])_f \cong R_f[y_1, \ldots, y_d]$ est également libre,
nous concluons par le fait déjà mentionné qu'une extension
de modules libres est libre.
\end{proof}

\begin{lemma}
\label{lemma-generic-flatness-finitely-presented}
\begin{slogan}
Liberté générique.
\end{slogan}
Soit $R \to S$ un homomorphisme d'anneaux.
Soit $M$ un $S$-module.
Supposons que
\begin{enumerate}
\item $R$ est intègre,
\item $R \to S$ est de présentation finie, et
\item $M$ est un $S$-module de présentation finie.
\end{enumerate}
Alors il existe un élément non nul $f \in R$ tel que
$M_f$ soit un $R_f$-module libre.
\end{lemma}

\begin{proof}
Écrivons $S = R[x_1, \ldots, x_n]/(g_1, \ldots, g_m)$.
Pour $g \in R[x_1, \ldots, x_n]$, notons $\overline{g}$ son image dans $S$.
Nous pouvons écrire $M = S^{\oplus t}/\sum Sn_i$ pour certains $n_i \in S^{\oplus t}$.
Écrivons $n_i = (\overline{g}_{i1}, \ldots, \overline{g}_{it})$ pour certains
$g_{ij} \in R[x_1, \ldots, x_n]$. Soit $R_0 \subset R$ le sous-anneau
engendré par tous les coefficients de tous les éléments
$g_i, g_{ij} \in R[x_1, \ldots, x_n]$. Définissons
$S_0 = R_0[x_1, \ldots, x_n]/(g_1, \ldots, g_m)$.
Définissons $M_0 = S_0^{\oplus t}/\sum S_0n_i$.
Alors $R_0$ est un anneau intègre de type fini sur $\mathbf{Z}$, et donc
noethérien (voir le
lemme \ref{lemma-Noetherian-permanence}).
De plus, via l'injection $R_0 \to R$, nous avons $S \cong R \otimes_{R_0} S_0$
et $M \cong R \otimes_{R_0} M_0$. En appliquant le
lemme \ref{lemma-generic-flatness-Noetherian},
nous obtenons un élément non nul $f \in R_0$ tel que $(M_0)_f$ soit un
$(R_0)_f$-module libre. Ainsi, $M_f = R_f \otimes_{(R_0)_f} (M_0)_f$
est un $R_f$-module libre.
\end{proof}

\begin{lemma}
\label{lemma-generic-flatness}
Soit $R \to S$ un homomorphisme d'anneaux.
Soit $M$ un $S$-module.
Supposons que
\begin{enumerate}
\item $R$ est intègre,
\item $R \to S$ est de type fini, et
\item $M$ est un $S$-module de type fini.
\end{enumerate}
Alors il existe un élément non nul $f \in R$ tel que
\begin{enumerate}
\item[(a)] $M_f$ et $S_f$ soient libres comme $R_f$-modules, et
\item[(b)] $S_f$ soit une $R_f$-algèbre de présentation finie et $M_f$ un
$S_f$-module de présentation finie.
\end{enumerate}
\end{lemma}

\begin{proof}
Nous démontrons d'abord le lemme pour $S = R[x_1, \ldots, x_n]$, puis
nous en déduisons le résultat général.

\medskip\noindent
Supposons $S = R[x_1, \ldots, x_n]$.
Choisissons des éléments $m_1, \ldots, m_t$ qui engendrent $M$. Cela donne
une suite exacte courte
$$
0 \to N \to S^{\oplus t} \xrightarrow{(m_1, \ldots, m_t)} M \to 0.
$$
Notons $K$ le corps des fractions de $R$. Notons
$S_K = K \otimes_R S = K[x_1, \ldots, x_n]$, et de même
$N_K = K \otimes_R N$, $M_K = K \otimes_R M$.
Comme $R \to K$ est plat, la suite reste exacte après tensorisation par $K$.
Comme $S_K = K[x_1, \ldots, x_n]$ est un anneau noethérien (voir le
lemme \ref{lemma-Noetherian-permanence}),
nous pouvons trouver un nombre fini d'éléments $n'_1, \ldots, n'_s \in N_K$
qui l'engendrent. Choisissons $n_1, \ldots, n_r \in N$ tels que
$n'_i = \sum a_{ij}n_j$ pour certains $a_{ij} \in K$. Posons
$$
M' = S^{\oplus t}/\sum\nolimits_{i = 1, \ldots, r} Sn_i
$$
Par construction, $M'$ est un $S$-module de présentation finie, et
il existe une surjection $M' \to M$ qui induit un isomorphisme
$M'_K \cong M_K$. Nous pouvons appliquer le
lemme \ref{lemma-generic-flatness-finitely-presented}
à $R \to S$ et $M'$, et nous trouvons un élément $f \in R$ tel que $M'_f$
soit un $R_f$-module libre. Ainsi, $M'_f \to M_f$ est une surjection de
modules sur l'anneau intègre $R_f$, dont la source est un module libre,
et qui devient un isomorphisme après tensorisation par $K$.
Elle est donc injective, puisque $M'_f \subset M'_K$, car $M'_f$
est libre et $R_f$ est un anneau intègre de corps des fractions $K$.
Ainsi, $M'_f \to M_f$ est un isomorphisme, et le résultat est démontré.

\medskip\noindent
Dans le cas général, choisissons une surjection $R[x_1, \ldots, x_n] \to S$.
Considérons $S$ et $M$ comme des modules finis sur $R[x_1, \ldots, x_n]$.
D'après le cas particulier démontré ci-dessus, il existe un élément non nul $f \in R$
tel que $S_f$ et $M_f$ soient tous deux libres comme $R_f$-modules et de présentation
finie comme $R_f[x_1, \ldots, x_n]$-modules. Cela implique clairement que
$S_f$ est une $R_f$-algèbre de présentation finie et que $M_f$ est un
$S_f$-module de présentation finie.
\end{proof}

\noindent
Soit $R \to S$ un homomorphisme d'anneaux. Soit $M$ un $S$-module. Considérons
la condition suivante sur un élément $f \in R$ :
\begin{equation}
\label{equation-flat-and-finitely-presented}
\left\{
\begin{matrix}
S_f & \text{est de présentation finie sur }R_f\\
M_f & \text{est de présentation finie comme }S_f\text{-module}\\
S_f, M_f & \text{sont libres comme }R_f\text{-modules}
\end{matrix}
\right.
\end{equation}
Nous définissons
\begin{equation}
\label{equation-good-locus}
U(R \to S, M)
=
\bigcup\nolimits_{f \in
R\text{ tel que }(\ref{equation-flat-and-finitely-presented})}
D(f)
\end{equation}
qui est un ouvert de $\Spec(R)$.

\begin{lemma}
\label{lemma-generic-flatness-locus-extension}
Soit $R \to S$ un homomorphisme d'anneaux.
Soit $0 \to M_1 \to M_2 \to M_3 \to 0$ une suite exacte courte
de $S$-modules.
Alors
$$
U(R \to S, M_1) \cap U(R \to S, M_3) \subset U(R \to S, M_2).
$$
\end{lemma}

\begin{proof}
Soit $u \in U(R \to S, M_1) \cap U(R \to S, M_3)$. Choisissons
$f_1, f_3 \in R$ tels que $u \in D(f_1)$, $u \in D(f_3)$ et
que (\ref{equation-flat-and-finitely-presented}) soit vérifiée pour
$f_1$ et $M_1$, ainsi que pour $f_3$ et $M_3$. Posons alors $f = f_1f_3$.
Nous avons $u \in D(f)$, et (\ref{equation-flat-and-finitely-presented})
est vérifiée pour $f$ et à la fois $M_1$ et $M_3$. Une extension de modules libres
est libre, et une extension de modules de présentation finie est de présentation finie
(lemme \ref{lemma-extension}). Nous voyons donc que
(\ref{equation-flat-and-finitely-presented}) est vérifiée pour $f$ et $M_2$.
Ainsi, $u \in U(R \to S, M_2)$, ce qui termine la démonstration.
\end{proof}

\begin{lemma}
\label{lemma-generic-flatness-locus-localize}
Soit $R \to S$ un homomorphisme d'anneaux.
Soit $M$ un $S$-module.
Soit $f \in R$.
Avec l'identification $\Spec(R_f) = D(f)$, nous avons
$U(R_f \to S_f, M_f) = D(f) \cap U(R \to S, M)$.
\end{lemma}

\begin{proof}
Supposons que $u \in U(R_f \to S_f, M_f)$. Il existe alors un
élément $g \in R_f$ tel que $u \in D(g)$ et que
(\ref{equation-flat-and-finitely-presented})
soit vérifiée pour le couple $((R_f)_g \to (S_f)_g, (M_f)_g)$.
Écrivons $g = a/f^n$ pour un certain $a \in R$. Posons $h = af$.
Alors $R_h = (R_f)_g$, $S_h = (S_f)_g$ et $M_h = (M_f)_g$.
De plus, $u \in D(h)$. Ainsi, $u \in U(R \to S, M)$.
Réciproquement, supposons que $u \in D(f) \cap U(R \to S, M)$.
Il existe alors un élément $g \in R$ tel que $u \in D(g)$ et que
(\ref{equation-flat-and-finitely-presented})
soit vérifiée pour le couple $(R_g \to S_g, M_g)$.
Il est alors clair que (\ref{equation-flat-and-finitely-presented})
est également vérifiée pour le couple
$(R_{fg} \to S_{fg}, M_{fg}) = ((R_f)_g \to (S_f)_g, (M_f)_g)$.
Ainsi, $u \in U(R_f \to S_f, M_f)$, ce qui termine la démonstration.
\end{proof}

\begin{lemma}
\label{lemma-generic-flatness-locus-reduce}
Soit $R \to S$ un homomorphisme d'anneaux.
Soit $M$ un $S$-module.
Soit $U \subset \Spec(R)$ un ouvert dense.
Supposons qu'il existe un recouvrement $U = \bigcup_{i \in I} D(f_i)$ par des
ouverts tel que $U(R_{f_i} \to S_{f_i}, M_{f_i})$ soit dense dans
$D(f_i)$ pour tout $i \in I$. Alors $U(R \to S, M)$ est dense dans
$\Spec(R)$.
\end{lemma}

\begin{proof}
Au vu du
lemme \ref{lemma-generic-flatness-locus-localize},
il s'agit d'une assertion purement topologique. En effet, d'après ce lemme,
nous voyons que $U(R \to S, M) \cap D(f_i)$ est dense dans $D(f_i)$
pour tout $i \in I$. D'après
Topologie, lemme \ref{topology-lemma-nowhere-dense-local},
nous voyons que $U(R \to S, M) \cap U$ est dense dans $U$.
Puisque $U$ est dense dans $\Spec(R)$, nous concluons que $U(R \to S, M)$
est dense dans $\Spec(R)$.
\end{proof}

\begin{lemma}
\label{lemma-generic-flatness-reduced}
Soit $R \to S$ un homomorphisme d'anneaux. Soit $M$ un $S$-module.
Supposons que
\begin{enumerate}
\item $R \to S$ est de type fini,
\item $M$ est un $S$-module fini, et
\item $R$ est réduit.
\end{enumerate}
Alors il existe une partie $U \subset \Spec(R)$ telle que
\begin{enumerate}
\item $U$ soit ouverte et dense dans $\Spec(R)$,
\item pour tout $u \in U$, il existe $f \in R$
tel que $u \in D(f) \subset U$ et que nous ayons
\begin{enumerate}
\item $M_f$ et $S_f$ sont libres sur $R_f$,
\item $S_f$ est une $R_f$-algèbre de présentation finie, et
\item $M_f$ est un $S_f$-module de présentation finie.
\end{enumerate}
\end{enumerate}
\end{lemma}

\begin{proof}
Remarquons que le lemme équivaut à dire que l'ouvert
$U(R \to S, M)$, voir l'équation (\ref{equation-good-locus}), est dense
dans $\Spec(R)$. Nous démontrons d'abord le lemme pour
$S = R[x_1, \ldots, x_n]$, puis
nous en déduisons le résultat général.

\medskip\noindent
Démonstration du cas $S = R[x_1, \ldots, x_n]$ et où $M$ est un module fini quelconque
sur $S$. Remarquons que, dans ce cas, $S_f = R_f[x_1, \ldots, x_n]$
est libre et de présentation finie sur $R_f$ ; nous n'avons donc pas
à nous préoccuper des conditions qui concernent $S$,
mais seulement de celles qui concernent $M$. Nous raisonnons par récurrence sur $n$.

\medskip\noindent
Il existe une filtration finie
$$
0 \subset M_1 \subset M_2 \subset \ldots \subset M_t = M
$$
telle que $M_i/M_{i - 1} \cong S/J_i$ pour un certain idéal $J_i \subset S$, voir le
lemme \ref{lemma-trivial-filter-finite-module}. Puisqu'
une intersection finie d'ouverts denses est un ouvert dense,
il résulte du
lemme \ref{lemma-generic-flatness-locus-extension}
qu'il suffit de démontrer le lemme pour chacun des modules $R/J_i$.
Nous pouvons donc supposer que $M = S/J$ pour un certain idéal $J$ de
$S = R[x_1, \ldots, x_n]$.

\medskip\noindent
Soit $I \subset R$ l'idéal engendré par les coefficients des
éléments de $J$. Soit $U_1 = \Spec(R) \setminus V(I)$ et
soit
$$
U_2 = \Spec(R) \setminus \overline{U_1}.
$$
Il est alors clair que $U = U_1 \cup U_2$ est dense dans $\Spec(R)$.
Soit $f \in R$ un élément tel que soit (a) $D(f) \subset U_1$, soit
(b) $D(f) \subset U_2$. Si, pour tout $f$ de ce type,
le lemme est vrai pour le couple $(R_f \to R_f[x_1, \ldots, x_n], M_f)$,
alors le
lemme \ref{lemma-generic-flatness-locus-reduce}
montre que $U(R \to S, M)$ est dense dans $\Spec(R)$.
Nous pouvons donc supposer soit (a) $I = R$, soit (b) $V(I) = \Spec(R)$.

\medskip\noindent
Dans le cas (b), nous avons en fait $I = 0$, puisque $R$ est réduit ! Ainsi, $J = 0$,
$M = S$, et le lemme est vrai dans ce cas.

\medskip\noindent
Dans le cas (a), il faut travailler un peu plus. Remarquons que tout élément
de $I$ est en fait le coefficient d'un monôme d'un élément de $J$, car
l'ensemble des coefficients des éléments de $J$ forme un idéal (détails omis).
Nous trouvons donc un élément
$$
g = \sum\nolimits_{K \in E} a_K x^K \in J
$$
où $E$ est un ensemble fini de multi-indices $K = (k_1, \ldots, k_n)$,
avec au moins un coefficient $a_{K_0}$ qui est une unité dans $R$. En fait, nous pouvons
en trouver un dont un coefficient est égal à $1$, puisque $1 \in I$ dans le cas (a).
Soit $m = \#\{K \in E \mid a_K \text{ n'est pas une unité}\}$.
Remarquons que $0 \leq m \leq \# E - 1$.
Nous raisonnons par récurrence sur $m$.

\medskip\noindent
Le cas $m = 0$. Dans ce cas, tous les coefficients $a_K$, $K \in E$,
de $g$ sont des unités, et $E \not = \emptyset$.
Si $E = \{K_0\}$ est un singleton et $K_0 = (0, \ldots, 0)$, alors $g$
est une unité et $J = S$, de sorte que le résultat est certainement vrai. (Cela se produit
en particulier lorsque $n = 0$, et fournit le cas initial de la récurrence
sur $n$.) Si l'on n'a pas $E = \{(0, \ldots, 0)\}$, alors au moins un $K$ n'est pas
égal à $(0, \ldots, 0)$, c'est-à-dire $g \not \in R$. Nous employons alors
l'astuce usuelle de la normalisation de Noether. Plus précisément, nous considérons
$$
G(y_1, \ldots, y_n) =
g(y_1 + y_n^{e_1}, y_2 + y_n^{e_2}, \ldots, y_{n - 1} + y_n^{e_{n - 1}}, y_n)
$$
avec $0 \ll e_{n -1} \ll e_{n - 2} \ll \ldots \ll e_1$. D'après le
lemme \ref{lemma-helper-polynomial},
le polynôme $G(y_1, \ldots, y_n)$ en $y_n$
est de la forme
$$
a_K y_n^{k_n + \sum_{i = 1, \ldots, n - 1} e_i k_i} +
\text{termes de degré inférieur en }y_n
$$
Puisque $a_K$ est une unité, nous concluons que $M = R[x_1, \ldots, x_n]/J$ est
fini sur $R[y_1, \ldots, y_{n - 1}]$. Ainsi,
$U(R \to R[x_1, \ldots, x_n], M) = U(R \to R[y_1, \ldots, y_{n - 1}], M)$,
et nous concluons par récurrence sur $n$.

\medskip\noindent
Le cas $m > 0$. Choisissons un multi-indice $K \in E$ tel que $a_K$ ne soit
pas une unité. Comme précédemment, posons
$U_1 = \Spec(R_{a_K}) = \Spec(R) \setminus V(a_K)$
et posons
$$
U_2 = \Spec(R) \setminus \overline{U_1}.
$$
Il est alors clair que $U = U_1 \cup U_2$ est dense dans $\Spec(R)$.
Soit $f \in R$ un élément tel que soit (a) $D(f) \subset U_1$, soit
(b) $D(f) \subset U_2$. Si, pour tout $f$ de ce type,
le lemme est vrai pour le couple $(R_f \to R_f[x_1, \ldots, x_n], M_f)$,
alors le
lemme \ref{lemma-generic-flatness-locus-reduce}
montre que $U(R \to S, M)$ est dense dans $\Spec(R)$.
Nous pouvons donc supposer soit (a) $a_KR = R$, soit (b) $V(a_K) = \Spec(R)$.
Dans le cas (a), le nombre $m$ diminue, puisque $a_K$ est devenu une unité.
Dans le cas (b), puisque $R$ est réduit, nous concluons que $a_K = 0$.
Ainsi, l'ensemble $E$ diminue, et le nombre $m$ diminue lui aussi.
Dans les deux cas, nous concluons par récurrence sur $m$.

\medskip\noindent
Nous avons maintenant démontré le lemme dans le cas $S = R[x_1, \ldots, x_n]$.
Supposons que $(R \to S, M)$ soit un couple quelconque satisfaisant les conditions du
lemme. Choisissons une surjection $R[x_1, \ldots, x_n] \to S$. Remarquons que,
avec les notations introduites dans (\ref{equation-good-locus}), nous avons
$$
U(R \to S, M) =
U(R \to R[x_1, \ldots, x_n], S)
\cap
U(R \to R[x_1, \ldots, x_n], M)
$$
Puisque nous venons de démontrer que les deux ouverts du membre de droite sont denses,
l'ouvert du membre de gauche est lui aussi dense.
\end{proof}






\section{Autour de Krull-Akizuki}
\label{section-krull-akizuki}

\noindent
Une application du théorème de Krull-Akizuki consiste à montrer qu'il existe de nombreux
anneaux de valuation discrète. Plus généralement, nous montrons dans cette section
comment construire des anneaux de valuation discrète qui dominent des anneaux
locaux noethériens.

\medskip\noindent
Nous montrons d'abord comment dominer un anneau local noethérien intègre
par un anneau local noethérien intègre de dimension $1$, en éclatant
l'idéal maximal.

\begin{lemma}
\label{lemma-dominate-by-dimension-1}
Soit $R$ un anneau local noethérien intègre de corps des fractions $K$.
Supposons que $R$ ne soit pas un corps.
Alors il existe un anneau intermédiaire $R \subset R' \subset K$ qui vérifie les conditions suivantes :
\begin{enumerate}
\item $R'$ est un anneau local noethérien de dimension $1$,
\item $R \to R'$ est un homomorphisme local d'anneaux, c'est-à-dire que $R'$ domine $R$, et
\item $R \to R'$ est essentiellement de type fini.
\end{enumerate}
\end{lemma}

\begin{proof}
Choisissons un anneau de valuation $A \subset K$ qui domine $R$ (il en
existe un d'après le lemme \ref{lemma-dominate}).
Notons $v$ la valuation correspondante.
Soit $x_1, \ldots, x_r$ un système minimal de générateurs de
l'idéal maximal $\mathfrak m$ de $R$. Nous pouvons supposer, et nous le faisons, que
$v(x_r) = \min\{v(x_1), \ldots, v(x_r)\}$. Considérons l'anneau
$$
S = R[x_1/x_r, x_2/x_r, \ldots, x_{r - 1}/x_r] \subset K.
$$
Remarquons que $\mathfrak mS = x_rS$ est un idéal principal.
Remarquons que $S \subset A$ et que $v(x_r) > 0$ ; nous voyons donc
que $x_rS \not = S$. Choisissons un idéal premier minimal $\mathfrak q$
au-dessus de $x_rS$. Alors $\text{hauteur}(\mathfrak q) = 1$ d'après le
lemme \ref{lemma-minimal-over-1},
et $\mathfrak q$ est au-dessus de $\mathfrak m$. Nous
voyons ainsi que $R' = S_{\mathfrak q}$ est une solution.
\end{proof}

\begin{lemma}[Koll\'ar]
\label{lemma-hart-serre-loc-thm}
\begin{reference}
Ceci est tiré d'un article à paraître de
J\'anos Koll\'ar intitulé « Variants of normality for Noetherian schemes ».
\end{reference}
Soit $(R, \mathfrak m)$ un anneau local noethérien.
Alors exactement l'une des possibilités suivantes est réalisée :
\begin{enumerate}
\item $(R, \mathfrak m)$ est artinien,
\item $(R, \mathfrak m)$ est régulier de dimension $1$,
\item $\text{profondeur}(R) \geq 2$, ou
\item il existe un homomorphisme fini d'anneaux $R \to R'$ qui n'est pas
un isomorphisme, dont le noyau et le conoyau sont annulés par une puissance
de $\mathfrak m$, tel que $\mathfrak m$ ne soit pas un idéal premier associé
de $R'$ et que $R' \not = 0$.
\end{enumerate}
\end{lemma}

\begin{proof}
Observons que $(R, \mathfrak m)$ n'est pas artinien si et seulement si
$V(\mathfrak m) \subset \Spec(R)$ est rare. Voir la
proposition \ref{proposition-dimension-zero-ring}. Nous le supposons désormais.

\medskip\noindent
Soit $J \subset R$ le plus grand idéal annulé par une puissance de $\mathfrak m$.
Si $J \not = 0$, alors $R \to R/J$ montre que $(R, \mathfrak m)$
relève du cas (4).

\medskip\noindent
Sinon, $J = 0$. En particulier, $\mathfrak m$ n'est pas un idéal premier associé
de $R$, et le
lemme \ref{lemma-ideal-nonzerodivisor} montre qu'il existe un non-diviseur de zéro $x \in \mathfrak m$. Si $\mathfrak m$
n'est pas un idéal premier associé de $R/xR$, alors $\text{profondeur}(R) \geq 2$
d'après le même lemme. Il ne reste donc que le cas où il existe
$y \in R$, $y \not \in xR$, tel que $y \mathfrak m \subset xR$.

\medskip\noindent
Si $y \mathfrak m \subset x \mathfrak m$, nous pouvons considérer
l'application $\varphi : \mathfrak m \to \mathfrak m$, $f \mapsto yf/x$
(bien définie puisque $x$ est un non-diviseur de zéro). Par l'astuce du déterminant
du lemme \ref{lemma-charpoly-module}, il existe un polynôme unitaire
$P$ à coefficients dans $R$ tel que $P(\varphi) = 0$.
Nous concluons que $P(y/x) = 0$ dans $R_x$.
Soit $R' \subset R_x$ l'anneau engendré par
$R$ et $y/x$. Alors $R \subset R'$, et $R'/R$ est un $R$-module de type fini
annulé par une puissance de $\mathfrak m$. Ainsi, $R$ relève du cas (4).

\medskip\noindent
Sinon, il existe $t \in \mathfrak m$ tel que $y t = u x$
pour une certaine unité $u$ de $R$. Après avoir remplacé $t$ par $u^{-1}t$,
nous obtenons $yt = x$. En particulier, $y$ est un non-diviseur de zéro.
Pour tout $t' \in \mathfrak m$, nous avons $y t' = x s$ pour un certain $s \in R$.
Ainsi, $y (t' - s t ) = x s - x s = 0$.
Comme $y$ n'est pas un diviseur de zéro, cela implique $t' = ts$, et donc
$\mathfrak m = (t)$. Ainsi, $(R, \mathfrak m)$ est régulier de dimension 1.

\medskip\noindent
L'argument mené jusqu'ici montre que tout $R$ relève de l'un des 4 cas.
Pour terminer, nous devons montrer que ces propriétés
s'excluent deux à deux. Nous montrerons seulement que (2) et (3)
excluent chacune (4) ; les autres cas sont laissés
au lecteur.

\medskip\noindent
Supposons que $R$ soit régulier de dimension $1$ et que $R \to R'$ soit un homomorphisme fini
d'anneaux dont le noyau et le conoyau sont annulés par une puissance de $\mathfrak m$,
et que $\mathfrak m$ ne soit pas un idéal premier associé de $R'$. Alors $R'$
est un $R$-module de type fini de profondeur au moins $1$, et est donc libre d'après le
lemme \ref{lemma-regular-mcm-free}. Puisque $R \to R'$ est un isomorphisme
au point générique, nous concluons que $R'$ doit être libre de rang $1$
comme $R$-module. Alors $R' \to \text{End}_R(R') \cong R$ est l'inverse
de l'application $R \to R'$. Ainsi, le cas (4) est impossible.

\medskip\noindent
Supposons que $R$ soit de profondeur $\geq 2$ et que $R \to R'$ soit un homomorphisme fini
d'anneaux dont le noyau et le conoyau sont annulés par une puissance
de $\mathfrak m$, et que $\mathfrak m$ ne soit pas un idéal premier associé de $R'$.
Alors $R \to R'$ est nécessairement injectif (car le noyau serait
à la fois de profondeur $0$ et de profondeur $\geq 1$), et le conoyau est de profondeur
au moins $1$ (d'après le lemme \ref{lemma-depth-in-ses}
et le fait que $R'$ est de profondeur $\geq 1$ par hypothèse) ;
il ne peut donc être supporté sur
$\{\mathfrak m\}$, à moins d'être lui aussi $0$. Ainsi, le cas (4) est impossible.
\end{proof}

\begin{lemma}
\label{lemma-nonregular-dimension-one}
Soit $R$ un anneau local d'idéal maximal $\mathfrak m$.
Supposons que $R$ soit noethérien, de dimension $1$, et que
$\dim(\mathfrak m/\mathfrak m^2) > 1$. Alors il existe
un homomorphisme d'anneaux $R \to R'$ tel que
\begin{enumerate}
\item $R \to R'$ soit fini,
\item $R \to R'$ ne soit pas un isomorphisme,
\item le noyau et le conoyau de $R \to R'$ soient annulés par
une puissance de $\mathfrak m$, et
\item $\mathfrak m$ ne soit pas un idéal premier associé de $R'$.
\end{enumerate}
\end{lemma}

\begin{proof}
Cela résulte du lemme \ref{lemma-hart-serre-loc-thm}
et du fait que $R$ n'est ni artinien, ni régulier, et
n'est pas de profondeur $\geq 2$ (ce dernier point parce que la
profondeur ne dépasse pas la dimension, d'après le
lemme \ref{lemma-bound-depth}).
\end{proof}

\begin{example}
\label{example-nonreduced}
Considérons l'anneau local noethérien
$$
R = k[[x, y]]/(y^2)
$$
Il est de dimension 1 et de Cohen-Macaulay.
Un exemple d'extension comme dans le
lemme \ref{lemma-nonregular-dimension-one}
est l'extension
$$
k[[x, y]]/(y^2) \subset k[[x, z]]/(z^2), \ \ y \mapsto xz
$$
autrement dit, elle s'obtient en adjoignant $y/x$ à $R$.
Répéter la construction $n > 1$ fois revient à
adjoindre l'élément $y/x^n$.
\end{example}

\begin{example}
\label{example-bad-dvr-char-p}
Soit $k$ un corps de caractéristique $p > 0$ tel que $k$
soit de degré infini sur son sous-corps $k^p$ des puissances $p$-ièmes.
Par exemple, $k = \mathbf{F}_p(t_1, t_2, t_3, \ldots)$.
Considérons l'anneau
$$
A =
\left\{
\sum a_i x^i \in k[[x]] \text{ tel que }
[k^p(a_0, a_1, a_2, \ldots) : k^p] < \infty
\right\}
$$
Alors $A$ est un anneau de valuation discrète, et son complété est
$A^\wedge = k[[x]]$. Remarquons que l'inclusion $A \subset k[[x]]$
induit une extension infinie et purement inséparable des corps des fractions.
Choisissons un $f \in k[[x]]$, $f \not \in A$. Soit $R = A[f] \subset k[[x]]$.
Alors $R$ est un anneau local noethérien intègre de dimension $1$ dont le
complété $R^\wedge$ est non réduit (à méditer !).
\end{example}

\begin{remark}
\label{remark-resolution-dim-1}
Supposons que $R$ soit un anneau semi-local noethérien intègre de dimension $1$.
S'il existe un idéal maximal $\mathfrak m \subset R$ tel que
$R_{\mathfrak m}$ ne soit pas régulier, nous pouvons appliquer le
lemme \ref{lemma-nonregular-dimension-one} à $(R, \mathfrak m)$
pour obtenir une extension finie d'anneaux $R \subset R_1$.
(Par exemple, on peut faire en sorte que $\Spec(R_1) \to \Spec(R)$
soit l'éclatement de $\Spec(R)$ le long de l'idéal $\mathfrak m$.)
Bien sûr, $R_1$ est un anneau semi-local noethérien intègre
de dimension $1$ ayant le même corps des fractions que $R$. Si $R_1$ n'est pas un
anneau semi-local régulier, nous pouvons répéter la construction pour
obtenir $R_1 \subset R_2$. Nous obtenons ainsi une suite
$$
R \subset R_1 \subset R_2 \subset R_3 \subset \ldots
$$
d'extensions finies d'anneaux, qui peut s'arrêter si $R_n$ est régulier pour
un certain $n$. La résolution des singularités serait l'affirmation
que $R_n$ finit effectivement par être régulier. En réalité, ce n'est pas
le cas. En effet, il existe en caractéristique $0$ un
anneau local noethérien intègre $A$ de dimension $1$ dont le complété est non réduit,
voir \cite[Proposition 3.1]{Ferrand-Raynaud} ou la
section \ref{examples-section-local-completion-nonreduced} du chapitre Exemples.
Pour un exemple en caractéristique $p > 0$, voir
l'exemple \ref{example-bad-dvr-char-p}.
Comme la formation de l'éclatement est compatible avec la complétion, il
est facile de voir que la suite ne se stabilise jamais.
Voir \cite{Bennett} pour une discussion (principalement en caractéristique positive).
En revanche, si le complété de $R$ en chacun de ses idéaux maximaux
est réduit, alors le procédé s'arrête (insérer ici une référence
future).
\end{remark}

\begin{lemma}
\label{lemma-characterize-dvr}
Soit $A$ un anneau. Les propriétés suivantes sont équivalentes.
\begin{enumerate}
\item L'anneau $A$ est un anneau de valuation discrète.
\item L'anneau $A$ est un anneau de valuation noethérien qui n'est pas un corps.
\item L'anneau $A$ est un anneau local régulier de dimension $1$.
\item L'anneau $A$ est un anneau local noethérien intègre dont l'idéal maximal
$\mathfrak m$ est engendré par un seul élément non nul.
\item L'anneau $A$ est un anneau local noethérien intègre et normal de dimension $1$.
\end{enumerate}
Dans ce cas, si $\pi$ est un générateur de l'idéal maximal de
$A$, alors tout élément non nul de $A$ s'écrit de manière unique
$u\pi^n$, où $u \in A$ est une unité.
\end{lemma}

\begin{proof}
L'équivalence de (1) et (2) est le
lemme \ref{lemma-valuation-ring-Noetherian-discrete}.
De plus, dans la preuve du lemme \ref{lemma-valuation-ring-Noetherian-discrete},
nous avons vu que si $A$ est un anneau de valuation discrète, alors $A$ est un anneau principal, d'où (3).
Remarquons qu'un anneau local régulier est intègre (voir le
lemme \ref{lemma-regular-domain}). Cela étant, l'équivalence de (3) et (4)
résulte de la théorie de la dimension, voir la section \ref{section-dimension}.

\medskip\noindent
Supposons (3) et soit $\pi$ un générateur de l'idéal maximal $\mathfrak m$.
Pour tout $n \geq 0$, nous avons
$\dim_{A/\mathfrak m} \mathfrak m^n/\mathfrak m^{n + 1} = 1$,
car il est engendré par $\pi^n$ (et ne peut être nul).
En particulier, $\mathfrak m^n = (\pi^n)$, et
l'anneau gradué $\bigoplus \mathfrak m^n/\mathfrak m^{n + 1}$
est isomorphe à l'anneau de polynômes $A/\mathfrak m[T]$.
Pour $x \in A \setminus \{0\}$, posons
$v(x) = \max\{n \mid x \in \mathfrak m^n\}$.
Autrement dit, $x = u \pi^{v(x)}$ avec $u \in A^*$.
D'après les remarques ci-dessus, nous avons $v(xy) = v(x) + v(y)$
pour tous $x, y \in A \setminus \{0\}$. Nous l'étendons au corps des fractions
$K$ de $A$ en posant $v(a/b) = v(a) - v(b)$ (ce qui est bien défini par la multiplicativité
démontrée ci-dessus). Il est alors clair que $A$ est l'ensemble des éléments de
$K$ de valuation $\geq 0$. Nous voyons donc que $A$ est un
anneau de valuation d'après le lemme \ref{lemma-valuation-valuation-ring}.

\medskip\noindent
Un anneau de valuation est un anneau intègre normal d'après le lemme \ref{lemma-valuation-ring-normal}.
Nous voyons donc que les conditions équivalentes (1) -- (3) impliquent
(5). Supposons que (5) soit vérifiée et, en vue d'une contradiction, que
$\mathfrak m$ ne puisse pas être engendré par exactement $1$ élément.
Alors le lemme \ref{lemma-nonregular-dimension-one} implique qu'il existe un homomorphisme fini
d'anneaux $A \to A'$ qui est un isomorphisme après inversion
de tout élément non nul de $\mathfrak m$, mais n'est pas un isomorphisme.
En particulier, nous pouvons identifier $A'$ à un sous-anneau du corps des fractions de $A$.
Puisque $A \to A'$ est fini, il est entier (voir le
lemme \ref{lemma-finite-is-integral}).
Puisque $A$ est normal, nous obtenons $A = A'$, contradiction.
\end{proof}

\begin{definition}
\label{definition-uniformizer}
Soit $A$ un anneau de valuation discrète. Une {\it uniformisante} est un élément
$\pi \in A$ qui engendre l'idéal maximal de $A$.
\end{definition}

\noindent
D'après le lemme \ref{lemma-characterize-dvr}, deux uniformisantes quelconques d'un anneau
de valuation discrète sont associées.

\begin{lemma}
\label{lemma-finite-length}
Soit $R$ un anneau intègre de corps des fractions $K$.
Soit $M$ un $R$-sous-module de $K^{\oplus r}$.
Supposons que $R$ soit local noethérien de dimension $1$.
Pour tout élément non nul $x \in R$, nous avons $\text{longueur}_R(R/xR) < \infty$
et
$$
\text{longueur}_R(M/xM) \leq r \cdot \text{longueur}_R(R/xR).
$$
\end{lemma}

\begin{proof}
Si $x$ est une unité, le résultat est vrai. Nous pouvons donc supposer que
$x \in \mathfrak m$, l'idéal maximal de $R$. Puisque $x$ n'est pas
nul et que $R$ est intègre, nous avons $\dim(R/xR) = 0$, et donc
$R/xR$ est de longueur finie. Considérons $M \subset K^{\oplus r}$ comme
dans le lemme. Nous pouvons supposer que les éléments de $M$ engendrent
$K^{\oplus r}$ comme $K$-espace vectoriel, après avoir remplacé $K^{\oplus r}$
par un sous-espace plus petit si nécessaire.

\medskip\noindent
Supposons d'abord que $M$ soit un $R$-module de type fini. Dans ce cas, nous pouvons chasser les
dénominateurs et supposer $M \subset R^{\oplus r}$. Puisque
$M$ engendre $K^{\oplus r}$ comme espace vectoriel, nous voyons que
$R^{\oplus r}/M$ est de longueur finie. En particulier, il existe
un entier $c \geq 0$ tel que $x^cR^{\oplus r} \subset M$.
Remarquons que $M \supset xM \supset x^2M \supset \ldots$ est une suite de
modules dont les quotients successifs sont tous isomorphes à $M/xM$. Ainsi,
nous voyons que
$$
n \text{longueur}_R(M/xM) = \text{longueur}_R(M/x^nM).
$$
Le même argument pour $M = R^{\oplus r}$ montre que
$$
n \text{longueur}_R(R^{\oplus r}/xR^{\oplus r}) =
\text{longueur}_R(R^{\oplus r}/x^nR^{\oplus r}).
$$
Par notre choix de $c$ ci-dessus, nous voyons que $x^nM$ est compris entre
$x^n R^{\oplus r}$ et $x^{n + c}R^{\oplus r}$. Cela donne facilement
$$
r(n + c) \text{longueur}_R(R/xR)
\geq
n \text{longueur}_R(M/xM)
\geq
r (n - c) \text{longueur}_R(R/xR)
$$
Ainsi, dans le cas de type fini, nous obtenons en fait le résultat du lemme avec
égalité.

\medskip\noindent
Supposons maintenant que $M$ ne soit pas de type fini. Supposons que la longueur de $M/xM$ soit
$\geq k$ pour un certain entier naturel $k$. Nous pouvons alors trouver
$$
0 \subset N_0 \subset N_1 \subset N_2 \subset \ldots N_k \subset M/xM
$$
avec $N_i \not = N_{i + 1}$ pour $i = 0, \ldots k - 1$.
Choisissons un élément $m_i \in M$ dont la classe modulo $xM$ appartient
à $N_i$, mais pas à $N_{i - 1}$, pour $i = 1, \ldots, k$.
Considérons le $R$-module de type fini $M' = Rm_1 + \ldots + Rm_k \subset M$.
Soit $N'_i \subset M'/xM'$ l'image réciproque de $N_i$.
Il est clair que $N'_i \not =N'_{i + 1}$ par notre choix de $m_i$.
Nous voyons donc que $\text{longueur}_R(M'/xM') \geq k$. D'après le
cas de type fini, nous concluons que $k \leq r\text{longueur}_R(R/xR)$,
comme voulu.
\end{proof}

\noindent
Voici une première application.

\begin{lemma}
\label{lemma-finite-extension-residue-fields-dimension-1}
Soit $R \to S$ un homomorphisme entre anneaux intègres qui induit une
injection de corps des fractions $K \subset L$. Si $R$ est un anneau local
noethérien de dimension $1$ et si $[L : K] < \infty$, alors
\begin{enumerate}
\item tout idéal premier $\mathfrak n_i$ de $S$ au-dessus de
l'idéal maximal $\mathfrak m$ de $R$ est maximal,
\item ces idéaux sont en nombre fini, et
\item $[\kappa(\mathfrak n_i) : \kappa(\mathfrak m)] < \infty$ pour tout $i$.
\end{enumerate}
\end{lemma}

\begin{proof}
Choisissons un élément non nul $x \in \mathfrak m$. Appliquons le lemme \ref{lemma-finite-length}
au sous-module $S \subset L \cong K^{\oplus n}$, où $n = [L : K]$.
Ainsi, l'anneau $S/xS$ est de longueur finie sur $R$. Il s'ensuit que
$S/\mathfrak m S$ est de longueur finie sur $\kappa(\mathfrak m)$.
Autrement dit, $\dim_{\kappa(\mathfrak m)} S/\mathfrak m S$
est finie (lemme \ref{lemma-dimension-is-length}). Ainsi, $S/\mathfrak mS$
est artinien (lemme \ref{lemma-finite-dimensional-algebra}). Les
résultats de structure sur les anneaux artiniens impliquent (1) et (2), voir
par exemple le lemme \ref{lemma-artinian-finite-length}.
Le point (3) résulte de la finitude établie ci-dessus.
\end{proof}

\begin{lemma}
\label{lemma-finite-length-global}
Soit $R$ un anneau intègre de corps des fractions $K$.
Soit $M$ un $R$-sous-module de $K^{\oplus r}$.
Supposons que $R$ soit noethérien de dimension $1$.
Pour tout élément non nul $x \in R$, nous avons
$\text{longueur}_R(M/xM) < \infty$.
\end{lemma}

\begin{proof}
Puisque $R$ est de dimension $1$, l'élément $x$ n'appartient qu'à un nombre fini
d'idéaux premiers $\mathfrak m_i$, $i = 1, \ldots, n$, qui sont tous maximaux. Puisque $R$ est
noethérien, nous voyons que $R/xR$ est artinien et que
$R/xR = \prod_{i = 1, \ldots, n} (R/xR)_{\mathfrak m_i}$ d'après la
proposition \ref{proposition-dimension-zero-ring} et le
lemme \ref{lemma-artinian-finite-length}. De même, $M/xM$
se décompose en produit $M/xM = \prod (M/xM)_{\mathfrak m_i}$
de ses localisés aux $\mathfrak m_i$. D'après le lemme \ref{lemma-finite-length}
appliqué à $M_{\mathfrak m_i}$, considéré comme $R_{\mathfrak m_i}$-module, chacun des
$M_{\mathfrak m_i}/xM_{\mathfrak m_i} = (M/xM)_{\mathfrak m_i}$
est de longueur finie sur $R_{\mathfrak m_i}$. Ainsi, $M/xM$
est de longueur finie sur $R$, car ce qui précède implique que $M/xM$
possède une filtration finie par des $R$-sous-modules dont les quotients
successifs sont isomorphes aux corps résiduels $\kappa(\mathfrak m_i)$.
\end{proof}

\begin{lemma}[Krull-Akizuki]
\label{lemma-krull-akizuki}
Soit $R$ un anneau intègre de corps des fractions $K$.
Soit $L/K$ une extension finie de corps.
Supposons que $R$ soit noethérien et que $\dim(R) = 1$.
Dans ce cas, tout anneau $A$ tel que $R \subset A \subset L$ est
noethérien.
\end{lemma}

\begin{proof}
Soit $I \subset A$ un idéal non nul. D'après le
lemme \ref{lemma-intersection-not-zero},
nous pouvons trouver un élément non nul $x \in I \cap R$.
Nous obtenons alors $I/xA \subset A/xA$. D'après le lemme \ref{lemma-finite-length-global},
le $R$-module $A/xA$ est de longueur finie comme $R$-module. Ainsi,
$I/xA$ est de longueur finie comme $R$-module. Par conséquent, $I$ est de type fini
comme idéal de $A$.
\end{proof}

\begin{lemma}
\label{lemma-exists-dvr}
Soit $R$ un anneau local noethérien intègre de corps des fractions $K$.
Supposons que $R$ ne soit pas un corps.
Soit $L/K$ une extension de corps de type fini.
Alors il existe un anneau de valuation discrète $A$ de corps des fractions
$L$ qui domine $R$.
\end{lemma}

\begin{proof}
Si $L$ n'est pas fini sur $K$, choisissons une base de transcendance
$x_1, \ldots, x_r$ de $L$ sur $K$ et remplaçons $R$ par
$R[x_1, \ldots, x_r]$ localisé en l'idéal maximal
engendré par $\mathfrak m_R$ et $x_1, \ldots, x_r$.
Nous pouvons donc supposer que $K \subset L$ est une extension finie.

\medskip\noindent
D'après le lemme \ref{lemma-dominate-by-dimension-1}, nous pouvons supposer $\dim(R) = 1$.

\medskip\noindent
Soit $A \subset L$ la clôture intégrale de $R$ dans $L$.
D'après le lemme \ref{lemma-krull-akizuki}, elle est noethérienne.
D'après le lemme \ref{lemma-integral-overring-surjective}, il
existe un idéal premier $\mathfrak q \subset A$ au-dessus
de l'idéal maximal de $R$.
D'après le lemme \ref{lemma-characterize-dvr}, l'anneau $A_{\mathfrak q}$ est un anneau de
valuation discrète qui domine $R$, comme souhaité.
\end{proof}








\section{Factorisation}
\label{section-factoring}

\noindent
Voici quelques notions et les relations entre elles qui sont généralement enseignées
dans un cours d'algèbre de première année de licence.

\begin{definition}
\label{definition-irreducible-prime-element}
Soit $R$ un anneau intègre.
\begin{enumerate}
\item Des éléments $x, y \in R$ sont dits {\it associés} s'il
existe une unité $u \in R^*$ telle que $x = uy$.
\item Un élément $x \in R$ est dit {\it irréductible}
s'il est non nul, n'est pas une unité et si, pour toute écriture $x = yz$ avec
$y, z \in R$, $y$ est soit une unité, soit associé à $x$.
\item Un élément $x \in R$ est dit {\it premier} si l'idéal
engendré par $x$ est un idéal premier.
\end{enumerate}
\end{definition}

\begin{lemma}
\label{lemma-easy-divisibility}
Soit $R$ un anneau intègre. Soient $x, y \in R$.
Alors $x$ et $y$ sont associés si et seulement si $(x) = (y)$.
\end{lemma}

\begin{proof}
Si $x = uy$ pour une unité $u \in R$, alors $(x) \subset (y)$ et
$y = u^{-1}x$, donc aussi $(y) \subset (x)$. Réciproquement, supposons que
$(x) = (y)$. Alors $x = fy$ et $y = gx$ pour certains $f, g \in A$.
Alors $x = fg x$ et, puisque $R$ est un anneau intègre, $fg = 1$. Ainsi,
$x$ et $y$ sont associés.
\end{proof}

\begin{lemma}
\label{lemma-factorization-exists}
Soit $R$ un anneau intègre. Considérons les conditions suivantes :
\begin{enumerate}
\item L'anneau $R$ vérifie la condition de chaîne ascendante sur les
idéaux principaux.
\item Tout élément $a \in R$ non nul qui n'est pas une unité
admet une factorisation $a = b_1 \ldots b_k$
où chaque $b_i$ est un élément irréductible de $R$.
\end{enumerate}
Alors (1) implique (2).
\end{lemma}

\begin{proof}
Soit $x$ un élément non nul, qui n'est pas une unité et qui n'admet pas de
factorisation en
irréductibles. Posons $x_1 = x$. Nous pouvons écrire $x = yz$, où aucun des deux facteurs
$y$ et $z$ n'est irréductible ni inversible.
Alors soit $y$ n'admet pas de factorisation
en irréductibles, auquel cas nous posons $x_2 = y$, soit $z$ n'admet pas
de factorisation en irréductibles, auquel cas nous posons $x_2 = z$.
En poursuivant ainsi, nous obtenons une suite
$$
\ldots | x_3 | x_2 | x_1
$$
d'éléments de $R$ telle que $x_n/x_{n + 1}$ n'est pas une unité.
Cela donne une suite strictement croissante d'idéaux principaux
$(x_1) \subset (x_2) \subset (x_3) \subset \ldots$, ce qui achève la démonstration.
\end{proof}

\begin{definition}
\label{definition-UFD}
Un {\it anneau factoriel}, aussi appelé domaine à factorisation unique ({\it UFD}),
est un anneau intègre $R$ tel que,
si $x \in R$ est non nul et n'est pas une unité, alors $x$ admet une factorisation
en irréductibles, et si
$$
x = a_1 \ldots a_m = b_1 \ldots b_n
$$
sont des factorisations en irréductibles, alors $n = m$ et
il existe une permutation $\sigma : \{1, \ldots, n\} \to \{1, \ldots, n\}$
telle que $a_i$ et $b_{\sigma(i)}$ sont associés.
\end{definition}

\begin{lemma}
\label{lemma-characterize-UFD}
Soit $R$ un anneau intègre. Supposons que tout élément non nul qui n'est pas une unité se factorise en
irréductibles. Alors $R$ est un anneau factoriel si et seulement si tout
élément irréductible est premier.
\end{lemma}

\begin{proof}
Supposons que $R$ soit un anneau factoriel et soit $x \in R$ un élément irréductible.
Supposons que $ab \in (x)$, c'est-à-dire $ab = cx$. Choisissons des factorisations
$a = a_1 \ldots a_n$, $b = b_1 \ldots b_m$ et $c = c_1 \ldots c_r$.
Par unicité de la factorisation
$$
a_1 \ldots a_n b_1 \ldots b_m = c_1 \ldots c_r x
$$
nous constatons que $x$ est associé à l'un des éléments
$a_1, \ldots, b_m$. Autrement dit, soit $a \in (x)$, soit $b \in (x)$,
et nous concluons que $x$ est premier.

\medskip\noindent
Supposons que tout élément irréductible soit premier. Nous devons démontrer que
deux factorisations en irréductibles ne diffèrent que par l'ordre des facteurs
et leur remplacement par des éléments associés. Supposons que $a_1 \ldots a_m = b_1 \ldots b_n$, où les $a_i$
et les $b_j$ sont irréductibles. Puisque $a_1$ est premier, nous voyons que
$b_j \in (a_1)$ pour un certain $j$. Après renumérotation, nous pouvons supposer que
$b_1 \in (a_1)$. Alors $b_1 = a_1 u$ et, puisque $b_1$ est irréductible,
nous voyons que $u$ est une unité. Par conséquent, $a_1$ et $b_1$ sont associés
et $a_2 \ldots a_n = ub_2\ldots b_m$. Par récurrence sur $n + m$,
nous voyons que $n = m$ et que $a_i$ est associé à $b_{\sigma(i)}$ pour
$i = 2, \ldots, n$, ce qu'il fallait démontrer.
\end{proof}

\begin{lemma}
\label{lemma-characterize-UFD-height-1}
Soit $R$ un anneau intègre noethérien. Alors $R$ est un anneau factoriel si et seulement si tout
idéal premier de hauteur $1$ est principal.
\end{lemma}

\begin{proof}
Supposons que $R$ soit un anneau factoriel et soit $\mathfrak p$ un idéal premier de hauteur 1.
Prenons $x \in \mathfrak p$ non nul et soit $x = a_1 \ldots a_n$
une factorisation en irréductibles.
Puisque $\mathfrak p$ est premier, nous voyons que $a_i \in \mathfrak p$
pour un certain $i$. Par le lemme \ref{lemma-characterize-UFD}, l'idéal
$(a_i)$ est premier. Puisque $\mathfrak p$ est de hauteur $1$, nous concluons
que $(a_i) = \mathfrak p$.

\medskip\noindent
Supposons que tout idéal premier de hauteur $1$ soit principal. Puisque $R$ est noethérien,
tout élément $x$ non nul qui n'est pas une unité admet une factorisation en irréductibles,
voir le lemme \ref{lemma-factorization-exists}. Il suffit de démontrer
qu'un élément irréductible $x$ est premier, voir le lemme \ref{lemma-characterize-UFD}.
Soit $(x) \subset \mathfrak p$ un idéal premier minimal au-dessus de $(x)$. Alors
$\mathfrak p$ est de hauteur $1$ par le lemme \ref{lemma-minimal-over-1}.
Par hypothèse, $\mathfrak p = (y)$. Ainsi, $x = yz$ et $z$ est une unité
puisque $x$ est irréductible. Donc $(x) = (y)$ et nous voyons que $x$ est premier.
\end{proof}

\begin{lemma}[Critère de factorialité de Nagata]
\label{lemma-invert-prime-elements}
\begin{reference}
\cite[lemme 2]{Nagata-UFD}
\end{reference}
Soit $A$ un anneau intègre. Soit $S \subset A$ une partie multiplicative
engendrée par des éléments premiers. Soit $x \in A$ irréductible. Alors
\begin{enumerate}
\item l'image de $x$ dans $S^{-1}A$ est irréductible ou une unité, et
\item $x$ est premier si et seulement si l'image de $x$ dans $S^{-1}A$ est
un élément premier ou une unité dans $S^{-1}A$.
\end{enumerate}
De plus, $A$ est alors un anneau factoriel si et seulement si tout élément non nul qui n'est pas une unité
de $A$ admet une factorisation en irréductibles et $S^{-1}A$ est un anneau factoriel.
\end{lemma}

\begin{proof}
Écrivons $x = \alpha \beta$ pour $\alpha, \beta \in S^{-1}A$. Alors
$\alpha = a/s$ et $\beta = b/s'$ pour $a, b \in A$, $s, s' \in S$.
Nous obtenons donc $ss'x = ab$. Par hypothèse, nous pouvons écrire
$ss' = p_1 \ldots p_r$ pour certains éléments premiers $p_i$.
Pour chaque $i$, l'élément $p_i$ divise soit $a$, soit $b$.
En divisant, nous obtenons une factorisation $x = a' b'$ et
$a = s'' a'$, $b = s''' b'$ pour certains $s'', s''' \in S$.
Comme $x$ est irréductible, soit $a'$, soit $b'$ est une unité.
Il s'ensuit que l'un des deux éléments $\alpha$, $\beta$ est une unité.
Cela démontre (1).

\medskip\noindent
Supposons que $x$ soit premier. Alors $A/(x)$ est un anneau intègre. Par conséquent,
$S^{-1}A/xS^{-1}A = S^{-1}(A/(x))$ est un anneau intègre ou nul.
Ainsi, l'image de $x$ est un élément premier ou une unité.

\medskip\noindent
Supposons que l'image de $x$ dans $S^{-1}A$ soit une unité.
Alors $y x = s$ pour un $s \in S$ et un $y \in A$. Par hypothèse,
$s = p_1 \ldots p_r$, où $p_i$ est un élément premier. Pour chaque $i$, soit
$p_i$ divise $y$, soit $p_i$ divise $x$. Dans le second cas,
$p_i$ et $x$ sont associés (puisque $x$ est irréductible) ; la conclusion s'ensuit.
Mais si le premier cas se produit pour tout $i = 1, \ldots, r$, alors
$x$ est une unité, ce qui est une contradiction.

\medskip\noindent
Supposons que l'image de $x$ dans $S^{-1}A$ soit un élément premier.
Supposons que $a, b \in A$ et $ab \in (x)$. Alors $sa = xy$ ou
$sb = xy$ pour certains $s \in S$ et $y \in A$. Supposons que le premier
cas se produise. Par hypothèse,
$s = p_1 \ldots p_r$, où $p_i$ est un élément premier. Pour chaque $i$, soit
$p_i$ divise $y$, soit $p_i$ divise $x$. Dans le second cas,
$p_i$ et $x$ sont associés (puisque $x$ est irréductible) ; la conclusion s'ensuit.
Si le premier cas se produit pour tout $i = 1, \ldots, r$, alors
$a \in (x)$, ce qu'il fallait démontrer. Ceci achève la démonstration de (2).

\medskip\noindent
La dernière assertion du lemme découle de (1) et (2)
et du lemme \ref{lemma-characterize-UFD}.
\end{proof}

\begin{lemma}
\label{lemma-UFD-ascending-chain-condition-principal-ideals}
Un anneau factoriel vérifie la condition de chaîne ascendante sur les idéaux
principaux.
\end{lemma}

\begin{proof}
Considérons une chaîne ascendante $(a_1) \subset (a_2) \subset (a_3) \subset \ldots$
d'idéaux principaux de $R$. Écrivons $a_1 = p_1^{e_1} \ldots p_r^{e_r}$
où les $p_i$ sont premiers. Nous voyons alors que $a_n$ est associé à
$p_1^{c_1} \ldots p_r^{c_r}$ pour certains $0 \leq c_i \leq e_i$.
Puisqu'il n'y a qu'un nombre fini de possibilités, la chaîne est stationnaire.
\end{proof}

\begin{lemma}
\label{lemma-factoring-in-polynomial}
Soit $R$ un anneau intègre. Supposons que $R$ vérifie la condition de chaîne ascendante
sur les idéaux principaux. Alors la même propriété vaut pour un anneau de polynômes
sur $R$.
\end{lemma}

\begin{proof}
Considérons une chaîne ascendante $(f_1) \subset (f_2) \subset (f_3) \subset \ldots$
d'idéaux principaux de $R[x]$. Puisque $f_{n + 1}$ divise $f_n$, nous voyons
que les degrés décroissent dans la suite. Ainsi, $f_n$
a un degré fixe $d \geq 0$ pour tout $n \gg 0$. Soit $a_n$ le
coefficient dominant de $f_n$. La condition $f_n \in (f_{n + 1})$
implique que $a_{n + 1}$ divise $a_n$ pour tout $n$.
Par notre hypothèse sur $R$, nous voyons que $a_{n + 1}$ et $a_n$
sont associés pour tout $n$ assez grand (lemme \ref{lemma-easy-divisibility}).
Ainsi, pour $n$ grand, nous voyons que $f_n = u f_{n + 1}$, où
$u \in R$ (pour des raisons de degré) est une unité (puisque $a_n$ et $a_{n + 1}$
sont associés).
\end{proof}

\begin{lemma}
\label{lemma-polynomial-ring-UFD}
Un anneau de polynômes sur un anneau factoriel est un anneau factoriel. En particulier, si $k$ est un corps,
alors $k[x_1, \ldots, x_n]$ est un anneau factoriel.
\end{lemma}

\begin{proof}
Soit $R$ un anneau factoriel. Alors $R$ vérifie la condition de chaîne ascendante sur les
idéaux principaux
(lemme \ref{lemma-UFD-ascending-chain-condition-principal-ideals}),
donc $R[x]$ vérifie la condition de chaîne ascendante sur les idéaux principaux
(lemme \ref{lemma-factoring-in-polynomial}), et tout
élément de $R[x]$ admet donc une factorisation en irréductibles
(lemme \ref{lemma-factorization-exists}).
Soit $S \subset R$ la partie multiplicative
engendrée par les éléments premiers. Puisque tout élément non inversible de $R$ est un produit
d'éléments premiers, nous voyons que $K = S^{-1}R$ est le corps des fractions
de $R$. Observons que tout élément premier de $R$ reste premier
dans $R[x]$ et que $S^{-1}(R[x]) = S^{-1}R[x] = K[x]$ est un
anneau factoriel (et même un anneau principal). Nous pouvons donc appliquer le
lemme \ref{lemma-invert-prime-elements} pour conclure.
\end{proof}

\begin{lemma}
\label{lemma-UFD-normal}
Un anneau factoriel est normal.
\end{lemma}

\begin{proof}
Soit $R$ un anneau factoriel. Soit $x$ un élément du corps des fractions de
$R$ qui est entier sur $R$. Écrivons $x^d - a_1 x^{d - 1} - \ldots - a_d = 0$
avec $a_i \in R$. Nous pouvons écrire
$x = u p_1^{e_1} \ldots p_r^{e_r}$ où $u$ est une unité, $e_i \in \mathbf{Z}$,
et $p_1, \ldots, p_r$ sont des éléments irréductibles qui ne sont pas associés.
Pour démontrer le lemme, nous devons montrer que $e_i \geq 0$. Sinon, disons $e_1 < 0$,
alors, pour $N \gg 0$, nous obtenons
$$
u^d p_2^{de_2 + N} \ldots p_r^{de_r + N} =
p_1^{-de_1}p_2^N \ldots p_r^N(
\sum\nolimits_{i = 1, \ldots, d} a_i x^{d - i} ) \in (p_1)
$$
ce qui contredit l'unicité de la factorisation dans $R$.
\end{proof}

\begin{definition}
\label{definition-PID}
Un {\it anneau principal} (en abrégé {\it PID})
est un anneau intègre $R$ tel que tout idéal est un idéal principal.
\end{definition}

\begin{lemma}
\label{lemma-PID-UFD}
Un anneau principal est un anneau factoriel.
\end{lemma}

\begin{proof}
Comme un anneau principal est noethérien, cela résulte du
lemme \ref{lemma-characterize-UFD-height-1}.
\end{proof}

\begin{definition}
\label{definition-dedekind-domain}
Un {\it anneau de Dedekind} est un anneau intègre $R$ tel que tout
idéal non nul $I \subset R$ puisse s'écrire comme un produit
$$
I = \mathfrak p_1 \ldots \mathfrak p_r
$$
d'idéaux premiers non nuls, de manière unique à l'ordre des facteurs $\mathfrak p_i$ près.
\end{definition}

\begin{lemma}
\label{lemma-PID-dedekind}
Un anneau principal est un anneau de Dedekind.
\end{lemma}

\begin{proof}
Soit $R$ un anneau principal. Puisque tout idéal non nul de $R$ est principal,
et que $R$ est un anneau factoriel (lemme \ref{lemma-PID-UFD}), cela résulte
du fait que tout élément irréductible de $R$ est premier
(lemme \ref{lemma-characterize-UFD}), chaque factorisation d'un élément
induisant alors une factorisation de l'idéal principal correspondant en idéaux premiers.
\end{proof}

\begin{lemma}
\label{lemma-product-ideals-principal}
\begin{slogan}
Un produit d'idéaux est un module inversible si et seulement si ses deux facteurs le sont.
\end{slogan}
Soit $A$ un anneau. Soient $I$ et $J$ des idéaux non nuls de $A$
tels que $IJ = (f)$ pour un non-diviseur de zéro $f \in A$. Alors $I$ et $J$ sont
des idéaux de type fini et sont, de plus, localement libres de type fini de rang $1$ comme $A$-modules.
\end{lemma}

\begin{proof}
Il suffit de montrer que $I$ et $J$ sont des $A$-modules localement libres de type fini
de rang $1$, voir le lemme \ref{lemma-finite-projective}.
Pour ce faire, écrivons $f = \sum_{i = 1, \ldots, n} x_i y_i$ avec $x_i \in I$
et $y_i \in J$. Nous pouvons
aussi écrire $x_i y_i = a_i f$ pour un certain $a_i \in A$.
Puisque $f$ est un non-diviseur de zéro, nous voyons que $\sum a_i = 1$.
Il suffit donc de montrer que les modules $I_{a_i}$ et
$J_{a_i}$ sont libres de rang $1$ sur $A_{a_i}$. Après avoir remplacé $A$ par
$A_{a_i}$, nous concluons que $f = xy$ pour certains $x \in I$ et $y \in J$.
Notons que $x$ et $y$ sont tous deux des non-diviseurs de zéro. Nous affirmons que
$I = (x)$ et $J = (y)$, ce qui achève la démonstration. En effet, si $x' \in I$,
alors $x'y = af = axy$ pour un certain $a \in A$. Ainsi, $x' = ax$ ; le même argument, avec les rôles intervertis, achève la démonstration.
\end{proof}

\begin{lemma}
\label{lemma-characterize-Dedekind}
Soit $R$ un anneau. Les conditions suivantes sont équivalentes :
\begin{enumerate}
\item $R$ est un anneau de Dedekind,
\item $R$ est un anneau intègre noethérien et, pour tout idéal maximal non nul
$\mathfrak m$, l'anneau local $R_{\mathfrak m}$ est un anneau de valuation discrète, et
\item $R$ est un anneau intègre noethérien et normal, et $\dim(R) \leq 1$.
\end{enumerate}
\end{lemma}

\begin{proof}
Supposons (1). L'argument n'est pas trivial, car nous n'avons pas supposé que $R$
était noethérien dans notre définition d'un anneau de Dedekind. Soit
$\mathfrak p \subset R$ un idéal premier. Observons que
$\mathfrak p \not = \mathfrak p^2$ par unicité
des factorisations de la définition. Choisissons $x \in \mathfrak p$
tel que $x \not \in \mathfrak p^2$. Soit $y \in \mathfrak p$ un
second élément (par exemple $y = 0$).
Écrivons $(x, y) = \mathfrak p_1 \ldots \mathfrak p_r$.
Puisque $(x, y) \subset \mathfrak p$, au moins l'un des idéaux premiers
$\mathfrak p_i$ est contenu dans $\mathfrak p$.
Mais puisque $x \not \in \mathfrak p^2$, il y en a au plus un.
Ainsi, exactement l'un des idéaux $\mathfrak p_1, \ldots, \mathfrak p_r$ est contenu
dans $\mathfrak p$, disons $\mathfrak p_1 \subset \mathfrak p$.
Nous concluons que $(x, y)R_\mathfrak p = \mathfrak p_1R_\mathfrak p$
est premier pour tout choix de $y$. Nous affirmons que
$(x)R_\mathfrak p = \mathfrak pR_\mathfrak p$. En effet,
choisissons $y \in \mathfrak p$. D'après ce qui précède, appliqué à $y^2$, nous voyons
que $(x, y^2)R_\mathfrak p$ est premier.
Par conséquent, $y \in (x, y^2)R_\mathfrak p$, c'est-à-dire $y = ax + by^2$ dans
$R_\mathfrak p$. Ainsi, $(1 - by)y = ax \in (x)R_\mathfrak p$, c'est-à-dire
$y \in (x)R_\mathfrak p$, comme il fallait le montrer.

\medskip\noindent
En écrivant à nouveau $(x) = \mathfrak p_1 \ldots \mathfrak p_r$ avec
$\mathfrak p_1 \subset \mathfrak p$, nous concluons que
$\mathfrak p_1 R_\mathfrak p = \mathfrak p R_\mathfrak p$, c'est-à-dire
$\mathfrak p_1 = \mathfrak p$. De plus, $\mathfrak p_1 = \mathfrak p$ est
un idéal de type fini de $R$ par le
lemme \ref{lemma-product-ideals-principal}.
Nous concluons que $R$ est noethérien par le lemme \ref{lemma-cohen}.
De plus, il s'ensuit que $R_\mathfrak m$ est un anneau de
valuation discrète pour tout idéal premier $\mathfrak p$, voir le
lemme \ref{lemma-characterize-dvr}.

\medskip\noindent
L'équivalence de (2) et (3) résulte des
lemmes \ref{lemma-normality-is-local} et
\ref{lemma-characterize-dvr}. Supposons que (2) et (3) soient satisfaites.
Soit $I \subset R$ un idéal. Nous allons construire une factorisation
de $I$. Si $I$ est premier, il n'y a rien à démontrer.
Sinon, choisissons $I \subset \mathfrak p$, où $\mathfrak p \subset R$ est un idéal maximal.
Posons $J = \{x \in R \mid x \mathfrak p \subset I\}$.
Nous affirmons que $J \mathfrak p = I$. Il suffit de le vérifier après localisation
aux idéaux maximaux $\mathfrak m$ de $R$ (la formation de $J$ commute à la
localisation et nous utilisons le lemme \ref{lemma-characterize-zero-local}).
Alors, soit $\mathfrak p R_\mathfrak m = R_\mathfrak m$ et le résultat
est clair, soit $\mathfrak p R_\mathfrak m = \mathfrak m R_\mathfrak m$.
Dans ce dernier cas, $\mathfrak p R_\mathfrak m = (\pi)$ et le cas où
$\mathfrak p$ est principal est immédiat. Par récurrence noethérienne,
l'idéal $J$ admet une factorisation et nous obtenons la factorisation souhaitée
de $I$. Nous omettons la démonstration de l'unicité de la factorisation.
\end{proof}

\noindent
Le lemme suivant est une variante du théorème de Krull-Akizuki.

\begin{lemma}
\label{lemma-integral-closure-Dedekind}
Soit $A$ un anneau intègre noethérien de dimension $1$ et de corps des fractions $K$.
Soit $L/K$ une extension finie. Soit $B$ la
clôture intégrale de $A$ dans $L$. Alors $B$ est un anneau de Dedekind et
$\Spec(B) \to \Spec(A)$ est surjective, a des fibres finies et
induit des extensions finies de corps résiduels.
\end{lemma}

\begin{proof}
D'après le théorème de Krull-Akizuki (lemme \ref{lemma-krull-akizuki}),
l'anneau $B$ est noethérien. Par le lemme \ref{lemma-integral-sub-dim-equal},
$\dim(B) = 1$. Ainsi, $B$ est un anneau de Dedekind par le
lemme \ref{lemma-characterize-Dedekind}.
La surjectivité de l'application sur les spectres résulte du
lemme \ref{lemma-integral-overring-surjective}.
Les deux dernières assertions résultent du
lemme \ref{lemma-finite-extension-residue-fields-dimension-1}.
\end{proof}







\section{Ordres d'annulation}
\label{section-orders-of-vanishing}

\begin{lemma}
\label{lemma-ord-additive}
Soit $R$ un anneau semi-local noethérien de dimension $1$.
Si $a, b \in R$ sont des non-diviseurs de zéro, alors
$$
\text{longueur}_R(R/(ab)) =
\text{longueur}_R(R/(a)) +
\text{longueur}_R(R/(b))
$$
et ces longueurs sont finies.
\end{lemma}

\begin{proof}
Cette finitude a été établie au lemme \ref{lemma-finite-length-global}.
L'additivité résulte de l'existence d'une suite exacte courte
$0 \to R/(a) \to R/(ab) \to R/(b) \to 0$, où la première flèche
est la multiplication par $b$. (On utilise l'additivité de la longueur ;
voir le lemme \ref{lemma-length-additive}.)
\end{proof}

\begin{definition}
\label{definition-ord}
Supposons que $K$ soit un corps et que $R \subset K$ soit un sous-anneau
local\footnote{On pourrait aussi donner cette définition lorsque $R$ est seulement
semi-local, mais ce n'est sans doute jamais vraiment la situation recherchée !}
noethérien de dimension $1$, de corps des fractions $K$.
Dans ce cas, on définit l'{\it ordre d'annulation relatif à $R$}
$$
\text{ord}_R : K^* \longrightarrow \mathbf{Z}
$$
par la règle
$$
\text{ord}_R(x) = \text{longueur}_R(R/(x))
$$
si $x \in R$, et l'on pose
$\text{ord}_R(x/y) = \text{ord}_R(x) - \text{ord}_R(y)$
pour $x, y \in R$ tous deux non nuls.
\end{definition}

\noindent
On peut utiliser l'ordre d'annulation pour comparer des réseaux dans un
espace vectoriel. En voici la définition.

\begin{definition}
\label{definition-lattice}
Soit $R$ un anneau intègre local noethérien de dimension $1$, de
corps des fractions $K$. Soit $V$ un espace vectoriel de dimension finie sur $K$.
Un {\it réseau dans $V$} est un sous-$R$-module de type fini $M \subset V$ tel
que $V = K \otimes_R M$.
\end{definition}

\noindent
La condition $V = K \otimes_R M$ signifie que $M$ contient une
base de l'espace vectoriel $V$. Signalons que, dans de nombreux ouvrages,
la notion de réseau n'est définie que lorsque
l'anneau $R$ est un anneau de valuation discrète. Si $R$ est un anneau de valuation
discrète, tout réseau est un $R$-module libre, ce qui n'est pas nécessairement le cas
en général.

\begin{lemma}
\label{lemma-compare-lattices}
Soit $R$ un anneau intègre local noethérien de dimension $1$, de
corps des fractions $K$. Soit $V$ un espace vectoriel de dimension finie sur $K$.
\begin{enumerate}
\item Si $M$ est un réseau dans $V$ et si $M \subset M' \subset V$
est un sous-$R$-module de $V$ contenant $M$,
alors les assertions suivantes sont équivalentes :
\begin{enumerate}
\item $M'$ est un réseau,
\item $\text{longueur}_R(M'/M)$ est finie, et
\item $M'$ est de type fini.
\end{enumerate}
\item Si $M$ est un réseau dans $V$ et si $M' \subset M$ est un sous-$R$-module
de $M$, alors $M'$ est un réseau si et seulement si
$\text{longueur}_R(M/M')$ est finie.
\item Si $M$, $M'$ sont des réseaux dans $V$, alors il en est de même de
$M \cap M'$ et $M + M'$.
\item Si $M \subset M' \subset M'' \subset V$ sont des réseaux dans $V$,
alors
$$
\text{longueur}_R(M''/M) =
\text{longueur}_R(M'/M) +
\text{longueur}_R(M''/M').
$$
\item Si $M$, $M'$, $N$, $N'$ sont des réseaux dans $V$ et si
$N \subset M \cap M'$, $M + M' \subset N'$, alors on a
\begin{eqnarray*}
& & \text{longueur}_R(M/M \cap M') - \text{longueur}_R(M'/M \cap M')\\
& = &
\text{longueur}_R(M/N) - \text{longueur}_R(M'/N) \\
& = &
\text{longueur}_R(M + M' / M') - \text{longueur}_R(M + M'/M) \\
& = &
\text{longueur}_R(N' / M') - \text{longueur}_R(N'/M)
\end{eqnarray*}
\end{enumerate}
\end{lemma}

\begin{proof}
Démonstration de (1). Supposons (1)(a) vérifiée. Soient $y_1, \ldots, y_m$ des générateurs de $M'$.
Alors chaque $y_i = x_i/f_i$ avec $x_i \in M$ et
$f_i \in R$ non nul.
On en déduit que $f_1 \ldots f_m M' \subset M$.
Puisque $R$ est un anneau local noethérien de dimension $1$,
on a $\mathfrak m^n \subset (f_1 \ldots f_m)$
pour un certain $n$ (par exemple en combinant les
lemmes \ref{lemma-one-equation} et la
proposition \ref{proposition-dimension-zero-ring}, ou en combinant les
lemmes \ref{lemma-finite-length} et \ref{lemma-length-infinite}).
En d'autres termes, $\mathfrak m^nM' \subset M$ pour un certain $n$.
Par conséquent,
$\text{longueur}(M'/M) < \infty$ d'après le lemme \ref{lemma-length-finite},
autrement dit, (1)(b) est vérifiée.
Supposons (1)(b). Alors $M'/M$ est un $R$-module de type fini
(voir le lemme \ref{lemma-finite-length-finite}).
Par conséquent, $M'$ est un $R$-module de type fini puisqu'il est une extension de $R$-modules de type fini.
Donc (1)(c) est vérifiée. L'implication
(1)(c) $\Rightarrow$ (1)(a) résulte de la remarque qui suit la
définition \ref{definition-lattice}.

\medskip\noindent
Démonstration de (2). Supposons que
$M$ soit un réseau dans $V$ et que $M' \subset M$ soit un sous-$R$-module.
Nous avons vu en (1) que, si $M'$ est un réseau, alors
$\text{longueur}_R(M/M') < \infty$. Réciproquement, supposons que
$\text{longueur}_R(M/M') < \infty$. Alors $M'$ est de type fini
puisque $R$ est noethérien. De plus, pour un certain $n$, on a
$\mathfrak m^n M \subset M'$ (lemme \ref{lemma-length-infinite}).
Il s'ensuit
que $M'$ contient une base de $V$ et que $M'$ est un réseau.

\medskip\noindent
Démonstration de (3). Supposons que $M$, $M'$ soient des réseaux dans $V$.
Puisque $R$ est noethérien, le sous-module $M \cap M'$ de $M$ est de type fini.
Comme $M$ est un réseau, nous pouvons trouver
$x_1, \ldots, x_n \in M$ qui forment une $K$-base de
$V$. Puisque $M'$ est un réseau, nous pouvons écrire $x_i = y_i/f_i$ avec
$y_i \in M'$ et $f_i \in R$. Ainsi $f_ix_i \in M \cap M'$. Par conséquent,
$M \cap M'$ est lui aussi un réseau.
Le fait que $M + M'$ soit un réseau résulte du point (1).

\medskip\noindent
Le point (4) résulte de l'additivité des longueurs
(lemme \ref{lemma-length-additive})
et de la suite exacte
$$
0 \to M'/M \to M''/M \to M''/M' \to 0
$$
Le point (5) s'obtient en appliquant le point (4) à plusieurs reprises.
\end{proof}

\begin{definition}
\label{definition-distance}
Soit $R$ un anneau intègre local noethérien de dimension $1$, de
corps des fractions $K$. Soit $V$ un espace vectoriel de dimension finie sur $K$.
Soient $M$, $M'$ deux réseaux dans $V$. La {\it distance de
$M$ à $M'$} est l'entier
$$
d(M, M') = \text{longueur}_R(M/M \cap M') - \text{longueur}_R(M'/M \cap M')
$$
défini au point (5) du lemme \ref{lemma-compare-lattices}.
\end{definition}

\noindent
En particulier, si $M' \subset M$, alors
$d(M, M') = \text{longueur}_R(M/M')$.

\begin{lemma}
\label{lemma-properties-distance-function}
Soit $R$ un anneau intègre local noethérien de dimension $1$, de
corps des fractions $K$. Soit $V$ un espace vectoriel de dimension finie sur $K$.
Cette fonction de distance vérifie
$$
d(M, M'') = d(M, M') + d(M', M'')
$$
quels que soient les trois réseaux $M$, $M'$, $M''$ de $V$.
En particulier, on a $d(M, M') = - d(M', M)$.
\end{lemma}

\begin{proof}
La démonstration est omise.
\end{proof}

\begin{lemma}
\label{lemma-order-vanishing-determinant}
Soit $R$ un anneau intègre local noethérien de dimension $1$, de
corps des fractions $K$. Soit $V$ un espace vectoriel de dimension finie sur $K$.
Soit $\varphi : V \to V$ un isomorphisme $K$-linéaire.
Pour tout réseau $M \subset V$, on a
$$
d(M, \varphi(M)) = \text{ord}_R(\det(\varphi))
$$
\end{lemma}

\begin{proof}
L'entier $d(M, \varphi(M))$ est indépendant du choix
du réseau $M$. En effet, soit $M'$ un autre
réseau. Alors
\begin{eqnarray*}
d(M, \varphi(M)) & = & d(M, M') + d(M', \varphi(M)) \\
& = & d(M, M') + d(\varphi(M'), \varphi(M)) + d(M', \varphi(M'))
\end{eqnarray*}
Puisque $\varphi$ est un isomorphisme, on a
$d(\varphi(M'), \varphi(M)) = d(M', M) = -d(M, M')$, et donc
$d(M, \varphi(M)) = d(M', \varphi(M'))$. De plus, les deux membres de
l'égalité du lemme, considérés comme fonctions de $\varphi$, sont additifs pour la composition, c'est-à-dire que
$$
\text{ord}_R(\det(\varphi \circ \psi))
=
\text{ord}_R(\det(\varphi))
+
\text{ord}_R(\det(\psi))
$$
et l'on a aussi
\begin{eqnarray*}
d(M, \varphi(\psi((M))) & = &
d(M, \psi(M)) + d(\psi(M), \varphi(\psi(M))) \\
& = & d(M, \psi(M)) + d(M, \varphi(M))
\end{eqnarray*}
par l'indépendance établie ci-dessus. Il suffit donc de démontrer le lemme
pour un système de générateurs de $\text{GL}(V)$. Choisissons un isomorphisme
$K^{\oplus n} \cong V$. Alors $\text{GL}(V) = \text{GL}_n(K)$ est
engendré par les matrices élémentaires $E$.
Le résultat est immédiat lorsque $E$ est la matrice identité.
Si $E = E_{ij}(\lambda)$ avec $i \not = j$, $\lambda \in K$,
$\lambda \not = 0$, par exemple
$$
E_{12}(\lambda) =
\left(
\begin{matrix}
1 & \lambda & \ldots \\
0 & 1 & \ldots \\
\ldots & \ldots & \ldots
\end{matrix}
\right)
$$
alors, après un changement de base, on obtient $E_{12}(1)$.
Le résultat est immédiat pour $E = E_{12}(1)$ en prenant pour réseau
$R^{\oplus n} \subset K^{\oplus n}$. Enfin, si $E = E_i(a)$,
où $a \in K^*$, par exemple
$$
E_1(a) =
\left(
\begin{matrix}
a & 0 & \ldots \\
0 & 1 & \ldots \\
\ldots & \ldots & \ldots
\end{matrix}
\right)
$$
alors $E_1(a)(R^{\oplus b}) = aR \oplus R^{\oplus n - 1}$ et
on voit immédiatement que $d(R^{\oplus n}, aR \oplus R^{\oplus n - 1})
= \text{ord}_R(a)$, d'où le résultat.
\end{proof}

\begin{lemma}
\label{lemma-finite-extension-dim-1}
Soit $A \to B$ un homomorphisme d'anneaux. Supposons que
\begin{enumerate}
\item $A$ est un anneau intègre local noethérien de dimension $1$,
\item $A \subset B$ est une inclusion finie entre anneaux intègres.
\end{enumerate}
Soit $L/K$ l'extension finie correspondante de leurs corps des fractions respectifs.
Soit $y \in L^*$ et posons $x = \text{Nm}_{L/K}(y)$.
Dans cette situation, $B$ est semi-local.
Soient $\mathfrak m_i$, $i = 1, \ldots, n$, les idéaux maximaux de $B$.
Alors
$$
\text{ord}_A(x) =
\sum\nolimits_i
[\kappa(\mathfrak m_i) : \kappa(\mathfrak m_A)]
\text{ord}_{B_{\mathfrak m_i}}(y)
$$
où $\text{ord}$ désigne l'ordre défini dans la définition \ref{definition-ord}.
\end{lemma}

\begin{proof}
L'anneau $B$ est semi-local d'après le lemme \ref{lemma-finite-in-codim-1}.
Écrivons $y = b/b'$ avec $b, b' \in B$.
Grâce à l'additivité de $\text{ord}$ et à la multiplicativité de
$\text{Nm}$, il suffit de démontrer le lemme pour
$y = b$ ou $y = b'$. Autrement dit, on peut supposer que $y \in B$.
Dans ce cas, le membre de droite de la formule est
$$
\sum [\kappa(\mathfrak m_i) : \kappa(\mathfrak m_A)]
\text{longueur}_{B_{\mathfrak m_i}}((B/yB)_{\mathfrak m_i})
$$
D'après le lemme \ref{lemma-pushdown-module}, cette expression est égale à
$\text{longueur}_A(B/yB)$. D'après le lemme \ref{lemma-order-vanishing-determinant},
on a
$$
\text{longueur}_A(B/yB) = d(B, yB) =
\text{ord}_A(\det\nolimits_K(L \xrightarrow{y} L)).
$$
Puisque $x = \text{Nm}_{L/K}(y) = \det\nolimits_K(L \xrightarrow{y} L)$
par définition, le lemme est démontré.
\end{proof}














\section{Morphismes quasi-finis}
\label{section-quasi-finite}

\noindent
Considérons un homomorphisme d'anneaux $R \to S$ de type fini.
Le morphisme $\Spec(S) \to \Spec(R)$ est quasi-fini
en un point si ce point est isolé dans sa fibre.
Cela signifie que la fibre est de dimension zéro en ce point.
Dans cette section, nous étudions les propriétés fondamentales de cette
notion importante mais technique. On trouvera des résultats plus avancés
dans la section suivante.

\begin{lemma}
\label{lemma-isolated-point}
Soit $k$ un corps.
Soit $S$ une $k$-algèbre de type fini.
Soit $\mathfrak q$ un idéal premier de $S$.
Les assertions suivantes sont équivalentes :
\begin{enumerate}
\item $\mathfrak q$ est un point isolé de $\Spec(S)$,
\item $S_{\mathfrak q}$ est fini sur $k$,
\item il existe $g \in S$, $g \not\in \mathfrak q$, tel que
$D(g) = \{ \mathfrak q \}$,
\item $\dim_{\mathfrak q} \Spec(S) = 0$,
\item $\mathfrak q$ est un point fermé de $\Spec(S)$ et
$\dim(S_{\mathfrak q}) = 0$, et
\item l'extension de corps $\kappa(\mathfrak q)/k$ est finie
et $\dim(S_{\mathfrak q}) = 0$.
\end{enumerate}
Dans ce cas, $S = S_{\mathfrak q} \times S'$ pour une
$k$-algèbre $S'$ de type fini. De plus, l'élément $g$
du point (3) vérifie $S_{\mathfrak q} = S_g$.
\end{lemma}

\begin{proof}
Supposons que $\mathfrak q$ soit un point isolé de $\Spec(S)$, c'est-à-dire que
$\{\mathfrak q\}$ soit ouvert dans $\Spec(S)$.
Puisque $\Spec(S)$ est un espace de Jacobson (voir les
lemmes \ref{lemma-finite-type-field-Jacobson} et
\ref{lemma-jacobson}),
on voit que $\mathfrak q$ est un point fermé. Par conséquent,
$\{\mathfrak q\}$ est ouvert et fermé dans $\Spec(S)$.
D'après les
lemmes \ref{lemma-disjoint-decomposition} et
\ref{lemma-disjoint-implies-product}, on peut
écrire $S = S_1 \times S_2$, où $\mathfrak q$
correspond à l'unique point de $\Spec(S_1)$.
Ainsi, $S_1 = S_{\mathfrak q}$ est un anneau de dimension
zéro et de type fini sur $k$. Il est donc fini sur $k$,
par exemple d'après le lemme \ref{lemma-Noether-normalization}.
Nous avons démontré que (1) implique (2).

\medskip\noindent
Supposons que $S_{\mathfrak q}$ soit fini sur $k$.
Alors $S_{\mathfrak q}$ est un anneau local artinien ; voir le
lemme \ref{lemma-finite-dimensional-algebra}. Ainsi,
$\Spec(S_{\mathfrak q}) = \{\mathfrak qS_{\mathfrak q}\}$ d'après le
lemme \ref{lemma-artinian-finite-length}.
Considérons la suite exacte $0 \to K \to S \to S_{\mathfrak q}
\to Q \to 0$. Il est clair que $K_{\mathfrak q} = Q_{\mathfrak q} = 0$.
De plus, $K$ est un $S$-module de type fini puisque $S$ est noethérien, et
$Q$ est un $S$-module de type fini puisque $S_{\mathfrak q}$ est fini sur $k$.
Il existe donc $g \in S$, $g \not \in \mathfrak q$, tel que
$K_g = Q_g = 0$. Ainsi, $S_{\mathfrak q} = S_g$ et
$D(g) = \{ \mathfrak q \}$. Nous avons démontré que (2) implique (3).

\medskip\noindent
Supposons que $D(g) = \{ \mathfrak q \}$. Comme $D(g)$ est ouvert par
définition de la topologie sur $\Spec(S)$, on voit que
$\mathfrak q$ est un point isolé de $\Spec(S)$.
Nous avons démontré que (3) implique (1).
Autrement dit, (1), (2) et (3) sont équivalentes.

\medskip\noindent
Supposons que $\dim_{\mathfrak q} \Spec(S) = 0$. Cela signifie qu'il
existe un voisinage ouvert de $\mathfrak q$ dans $\Spec(S)$
qui est de dimension zéro. Il existe alors un voisinage ouvert de la
forme $D(g)$ qui est de dimension zéro. Puisque $S_g$ est noethérien,
on en déduit que $S_g$ est artinien et que
$D(g) = \Spec(S_g)$ est un ensemble discret fini ; voir la
proposition \ref{proposition-dimension-zero-ring}.
Ainsi, $\mathfrak q$ est un point isolé de $D(g)$ et,
en appliquant l'équivalence de (1) et (2) ci-dessus à
$\mathfrak qS_g \subset S_g$, on voit que
$S_{\mathfrak q} = (S_g)_{\mathfrak qS_g}$ est fini sur $k$.
Par conséquent, (4) implique (2). Il est clair que (1) implique (4).
Ainsi, (1) -- (4) sont toutes équivalentes.

\medskip\noindent
Le lemme \ref{lemma-dimension-closed-point-finite-type-field}
donne l'implication (5) $\Rightarrow$ (4).
L'implication (4) $\Rightarrow$ (6) résulte du
lemme \ref{lemma-dimension-at-a-point-finite-type-field}.
L'implication (6) $\Rightarrow$ (5) résulte du
lemme \ref{lemma-finite-residue-extension-closed}.
À ce stade, nous savons que (1) -- (6) sont équivalentes.

\medskip\noindent
Les deux assertions finales du lemme ont été obtenues au
cours de la démonstration de l'équivalence de (1), (2) et (3) ci-dessus.
\end{proof}

\begin{lemma}
\label{lemma-isolated-point-fibre}
\begin{slogan}
Conditions équivalentes pour les points isolés dans les fibres
\end{slogan}
Soit $R \to S$ un homomorphisme d'anneaux de type fini.
Soit $\mathfrak q \subset S$ un idéal premier au-dessus de
$\mathfrak p \subset R$. Soit $F = \Spec(S \otimes_R \kappa(\mathfrak p))$
la fibre de $\Spec(S) \to \Spec(R)$ ; voir la
remarque \ref{remark-fundamental-diagram}.
Notons $\overline{\mathfrak q} \in F$ le point correspondant à
$\mathfrak q$. Les assertions suivantes sont équivalentes :
\begin{enumerate}
\item $\overline{\mathfrak q}$ est un point isolé de $F$,
\item $S_{\mathfrak q}/\mathfrak pS_{\mathfrak q}$ est fini sur
$\kappa(\mathfrak p)$,
\item il existe $g \in S$, $g \not \in \mathfrak q$, tel que
le seul point de $D(g)$ dont l'image est $\mathfrak p$ soit $\mathfrak q$,
\item $\dim_{\overline{\mathfrak q}}(F) = 0$,
\item $\overline{\mathfrak q}$ est un point fermé de $F$ et
$\dim(S_{\mathfrak q}/\mathfrak pS_{\mathfrak q}) = 0$, et
\item l'extension de corps $\kappa(\mathfrak q)/\kappa(\mathfrak p)$
est finie et $\dim(S_{\mathfrak q}/\mathfrak pS_{\mathfrak q}) = 0$.
\end{enumerate}
\end{lemma}

\begin{proof}
Remarquons que $S_{\mathfrak q}/\mathfrak pS_{\mathfrak q} =
(S \otimes_R \kappa(\mathfrak p))_{\overline{\mathfrak q}}$.
De plus, $S \otimes_R \kappa(\mathfrak p)$ est de type fini sur
$\kappa(\mathfrak p)$.
Ces conditions correspondent exactement à celles du
lemme \ref{lemma-isolated-point}
pour la $\kappa(\mathfrak p)$-algèbre $S \otimes_R \kappa(\mathfrak p)$
et l'idéal premier $\overline{\mathfrak q}$ ; elles sont donc équivalentes.
\end{proof}

\begin{definition}
\label{definition-quasi-finite}
Soit $R \to S$ un homomorphisme d'anneaux de type fini.
Soit $\mathfrak q \subset S$ un idéal premier.
\begin{enumerate}
\item Si les conditions équivalentes du lemme \ref{lemma-isolated-point-fibre}
sont satisfaites, on dit que $R \to S$ est {\it quasi-fini en $\mathfrak q$}.
\item On dit qu'un homomorphisme d'anneaux $A \to B$ est {\it quasi-fini}
s'il est de type fini et quasi-fini en tout idéal premier de $B$.
\end{enumerate}
\end{definition}

\begin{lemma}
\label{lemma-quasi-finite-above-p}
Soit $R \to S$ un homomorphisme d'anneaux de type fini
et soit $\mathfrak p$ un idéal premier de $R$.
Alors les assertions suivantes sont équivalentes :
\begin{enumerate}
\item $R \to S$ est quasi-fini en tout idéal premier de $S$ au-dessus de $\mathfrak p$,
\item $S \otimes_R \kappa(\mathfrak p)$ est une
$\kappa(\mathfrak p)$-algèbre finie, et
\item $\Spec(S \otimes_R \kappa(\mathfrak p))$ est un ensemble fini.
\end{enumerate}
\end{lemma}

\begin{proof}
La condition (1) signifie que la topologie de
$\Spec(S \otimes_R \kappa(\mathfrak p))$ est discrète (puisque tout point
est ouvert). L'équivalence de (1), (2) et (3) résulte donc du
lemme \ref{lemma-finite-type-algebra-finite-nr-primes}.
\end{proof}

\begin{lemma}
\label{lemma-quasi-finite}
Soit $R \to S$ un homomorphisme d'anneaux de type fini. Alors $R \to S$ est quasi-fini
si et seulement si, pour tout idéal premier $\mathfrak p \subset R$, l'anneau
$S \otimes_R \kappa(\mathfrak p)$ est fini sur $\kappa(\mathfrak p)$.
\end{lemma}

\begin{proof}
Cela résulte immédiatement du lemme plus général
\ref{lemma-quasi-finite-above-p}.
\end{proof}

\begin{lemma}
\label{lemma-quasi-finite-local}
Soit $R \to S$ un homomorphisme d'anneaux de type fini. Soit $\mathfrak q \subset S$
un idéal premier au-dessus de $\mathfrak p \subset R$. Soient
$f \in R$, $f \not \in \mathfrak p$ et $g \in S$, $g \not \in \mathfrak q$.
Alors $R \to S$ est quasi-fini en $\mathfrak q$ si et seulement si
$R_f \to S_{fg}$ est quasi-fini en $\mathfrak qS_{fg}$.
\end{lemma}

\begin{proof}
La fibre de $\Spec(S_{fg}) \to \Spec(R_f)$ est homéomorphe
à un ouvert de la fibre de $\Spec(S) \to \Spec(R)$.
Le lemme résulte donc de la condition (1) du
lemme \ref{lemma-isolated-point-fibre}.
\end{proof}

\begin{lemma}
\label{lemma-four-rings}
Soit
$$
\xymatrix{
S \ar[r] & S' & &
\mathfrak q \ar@{-}[r] & \mathfrak q' \\
R \ar[u] \ar[r] &  R' \ar[u] & &
\mathfrak p \ar@{-}[r] \ar@{-}[u] & \mathfrak p' \ar@{-}[u]
}
$$
un diagramme commutatif d'anneaux, avec les idéaux premiers indiqués.
Supposons que $R \to S$ soit de type fini et que $S \otimes_R R' \to S'$ soit surjectif.
Si $R \to S$ est quasi-fini en $\mathfrak q$, alors
$R' \to S'$ est quasi-fini en $\mathfrak q'$.
\end{lemma}

\begin{proof}
Écrivons $S \otimes_R \kappa(\mathfrak p) = S_1 \times S_2$,
avec $S_1$ fini sur $\kappa(\mathfrak p)$ et tel que
$\mathfrak q$ corresponde à un idéal premier de $S_1$, comme dans le
lemme \ref{lemma-isolated-point}. Cette décomposition en produit
induit une décomposition en produit correspondante pour toute
$S \otimes_R \kappa(\mathfrak p)$-algèbre. En particulier,
on obtient $S' \otimes_{R'} \kappa(\mathfrak p') = S'_1 \times S'_2$.
Puisque $S \otimes_R R' \to S'$ est surjectif, l'homomorphisme canonique
$(S \otimes_R \kappa(\mathfrak p)) \otimes_{\kappa(\mathfrak p)}
\kappa(\mathfrak p') \to S' \otimes_{R'} \kappa(\mathfrak p')$
est surjectif, et par conséquent
$S_i \otimes_{\kappa(\mathfrak p)} \kappa(\mathfrak p') \to S'_i$
est surjectif. Il s'ensuit que $S'_1$ est fini sur $\kappa(\mathfrak p')$.
L'homomorphisme $S' \otimes_{R'} \kappa(\mathfrak p') \to
\kappa(\mathfrak q')$ se factorise par $S_1'$
(c.-à-d.\ qu'il annule le facteur $S_2'$)
car l'homomorphisme $S \otimes_R \kappa(\mathfrak p) \to
\kappa(\mathfrak q)$ se factorise par $S_1$
(c.-à-d.\ qu'il annule le facteur $S_2$). Ainsi,
$\mathfrak q'$ correspond à un point de
$\Spec(S_1')$ dans la décomposition en union disjointe
de la fibre : $\Spec(S' \otimes_{R'} \kappa(\mathfrak p'))
= \Spec(S_1') \amalg \Spec(S_2')$ ; voir le
lemme \ref{lemma-spec-product}.
Puisque $S_1'$ est fini sur un corps, c'est un anneau artinien,
et $\Spec(S_1')$ est donc un ensemble discret fini.
(Voir la proposition \ref{proposition-dimension-zero-ring}.)
On conclut que $\mathfrak q'$ est isolé dans sa fibre,
ce qui donne le résultat.
\end{proof}

\begin{lemma}
\label{lemma-quasi-finite-composition}
La composée d'homomorphismes d'anneaux quasi-finis est quasi-finie.
\end{lemma}

\begin{proof}
Supposons que $A \to B$ et $B \to C$ soient des homomorphismes d'anneaux quasi-finis. D'après le
lemme \ref{lemma-compose-finite-type},
on voit que $A \to C$ est de type fini.
Soit $\mathfrak r \subset C$ un idéal premier de $C$ au-dessus de
$\mathfrak q \subset B$ et de $\mathfrak p \subset A$. Puisque $A \to B$ et
$B \to C$ sont quasi-finis en $\mathfrak q$ et $\mathfrak r$ respectivement,
il existe alors $b \in B$ et $c \in C$ tels que $\mathfrak q$ soit
le seul idéal premier de $D(b)$ au-dessus de $\mathfrak p$ et, de même,
$\mathfrak r$ soit le seul idéal premier de $D(c)$ au-dessus de $\mathfrak q$.
Si $c' \in C$ est l'image de $b \in B$, alors $\mathfrak r$ est le seul
idéal premier de $D(cc')$ au-dessus de $\mathfrak p$.
Par conséquent, $A \to C$ est quasi-fini en $\mathfrak r$.
\end{proof}

\begin{lemma}
\label{lemma-quasi-finite-base-change}
Soit $R \to S$ un homomorphisme d'anneaux de type fini.
Soit $R \to R'$ un homomorphisme d'anneaux quelconque. Posons $S' = R' \otimes_R S$.
\begin{enumerate}
\item L'ensemble
$\{\mathfrak q' \mid R' \to S' \text{ quasi-fini en }\mathfrak q'\}$
est l'image réciproque de l'ensemble correspondant de $\Spec(S)$
par le morphisme canonique $\Spec(S') \to \Spec(S)$.
\item Si $\Spec(R') \to \Spec(R)$ est surjectif,
alors $R \to S$ est quasi-fini si et seulement si $R' \to S'$ est quasi-fini.
\item Tout changement de base d'un homomorphisme d'anneaux quasi-fini est quasi-fini.
\end{enumerate}
\end{lemma}

\begin{proof}
Soit $\mathfrak p' \subset R'$ un idéal premier au-dessus de $\mathfrak p \subset R$.
Alors l'anneau de la fibre $S' \otimes_{R'} \kappa(\mathfrak p')$ est le
changement de base de l'anneau de la fibre $S \otimes_R \kappa(\mathfrak p)$
par l'extension de corps $\kappa(\mathfrak p) \to \kappa(\mathfrak p')$.
La première assertion résulte donc de l'invariance de la dimension
par extension du corps de base
(lemme \ref{lemma-dimension-at-a-point-preserved-field-extension})
et du lemme \ref{lemma-isolated-point}.
La stabilité des morphismes quasi-finis par changement de base résulte de
ce qui précède et de la stabilité de la propriété d'être de type fini par changement de base.
La deuxième assertion en résulte
car l'hypothèse implique que, pour tout idéal premier $\mathfrak q \subset S$, on peut
trouver un idéal premier $\mathfrak q' \subset S'$ au-dessus de lui.
\end{proof}

\begin{lemma}
\label{lemma-quasi-finite-permanence}
Soient $A \to B$ et $B \to C$ des homomorphismes d'anneaux tels que $A \to C$
soit de type fini. Soit $\mathfrak r$ un idéal premier de $C$ au-dessus de
$\mathfrak q \subset B$ et de $\mathfrak p \subset A$.
Si $A \to C$ est quasi-fini en $\mathfrak r$, alors
$B \to C$ est quasi-fini en $\mathfrak r$.
\end{lemma}

\begin{proof}
Remarquons que $B \to C$ est de type fini
(lemme \ref{lemma-compose-finite-type}),
de sorte que l'énoncé a un sens.
Utilisons la caractérisation (3) du lemme \ref{lemma-isolated-point-fibre}.
Si $A \to C$ est quasi-fini en $\mathfrak r$, alors
il existe $c \in C$ tel que
$$
\{\mathfrak r' \subset C \text{ au-dessus de }\mathfrak p\} \cap D(c) =
\{\mathfrak{r}\}.
$$
Puisque les idéaux premiers $\mathfrak r' \subset C$ au-dessus de $\mathfrak q$
forment un sous-ensemble des idéaux premiers $\mathfrak r' \subset C$ au-dessus de
$\mathfrak p$, on conclut que $B \to C$ est quasi-fini en $\mathfrak r$.
\end{proof}

\begin{lemma}
\label{lemma-generically-finite}
Soit $R \to S$ un homomorphisme d'anneaux de type fini.
Soit $\mathfrak p \subset R$ un idéal premier minimal.
Supposons qu'il n'y ait qu'un nombre fini d'idéaux premiers de $S$
au-dessus de $\mathfrak p$. Il existe alors
$g \in R$, $g \not \in \mathfrak p$, tel que le
homomorphisme d'anneaux $R_g \to S_g$ soit fini.
\end{lemma}

\begin{proof}
Soient $x_1, \ldots, x_n$ des générateurs de $S$ sur $R$.
D'après le lemme \ref{lemma-quasi-finite-above-p}, l'hypothèse signifie que
$S \otimes_R \kappa(\mathfrak p) = S_{\mathfrak p}/\mathfrak pS_{\mathfrak p}$
est une $\kappa(\mathfrak p)$-algèbre finie.
On peut donc trouver des polynômes unitaires
$P_i \in R_{\mathfrak p}[X]$ tels que l'image de $P_i(x_i)$ soit nulle
dans $S_{\mathfrak p}/\mathfrak pS_{\mathfrak p}$.
Puisque $\mathfrak p$ est un idéal premier minimal,
$\mathfrak pR_{\mathfrak p}$ est un idéal localement nilpotent ; voir le
lemme \ref{lemma-minimal-prime-reduced-ring}.
Par conséquent, $\mathfrak pS_{\mathfrak p}$ est un idéal localement nilpotent ; voir le
lemme \ref{lemma-locally-nilpotent}. Ainsi,
il existe des $e_i \geq 1$ tels que $P(x_i)^{e_i} = 0$
dans $S_{\mathfrak p}$. Choisissons $g_1 \in R$, $g_1 \not \in \mathfrak p$,
tel que les coefficients de $P_i$ appartiennent à $R[1/g_1]$ pour tout $i$.
Choisissons ensuite $g_2 \in R$, $g_2 \not \in \mathfrak p$,
tel que $P(x_i)^{e_i} = 0$ dans $S_{g_1g_2}$. En posant $g = g_1g_2$,
on obtient le résultat.
\end{proof}







\section{Théorème principal de Zariski}
\label{section-Zariski}

\noindent
Dans cette section, notre objectif est de démontrer la version algébrique du
théorème principal de Zariski. Ce théorème servira de fondement à de nombreux
développements ultérieurs de la théorie des schémas et de leurs
morphismes dans la suite du projet Champs.

\medskip\noindent
Soit $R \to S$ un homomorphisme d'anneaux de type fini.
Notre objectif dans cette section est de montrer que l'ensemble
des points de $\Spec(S)$ en lesquels cet homomorphisme est quasi-fini
est {\it ouvert} (théorème \ref{theorem-main-theorem}).
En fait, nous verrons qu'il existe
un homomorphisme d'anneaux fini $R \to S'$ tel qu'en un certain sens le lieu de quasi-finitude
de $S/R$ soit ouvert dans $\Spec(S')$
(mais nous ne le démontrerons pas dans ce chapitre d'algèbre, car nous n'y
développons pas le langage des schémas -- lorsque $R \to S$
est quasi-fini, voir le lemme \ref{lemma-quasi-finite-open-integral-closure}).
Ces assertions sont assez délicates à démontrer ; nous procédons
par une longue suite de lemmes sur les extensions entières et
finies d'anneaux. On trouvera ces résultats
dans \cite{Henselian} et \cite{Peskine}. Les notes
de Thierry Coquand nous ont également été utiles.

\begin{lemma}
\label{lemma-make-integral-trivial}
Soit $\varphi : R \to S$ un homomorphisme d'anneaux.
Supposons que $t \in S$ satisfasse la
relation $\varphi(a_0) + \varphi(a_1)t + \ldots + \varphi(a_n) t^n = 0$.
Alors $\varphi(a_n)t$ est entier sur $R$.
\end{lemma}

\begin{proof}
En effet, multiplions l'égalité
$\varphi(a_0) + \varphi(a_1)t + \ldots + \varphi(a_n) t^n = 0$
par $\varphi(a_n)^{n-1}$ et écrivons-la sous la forme
$\varphi(a_0 a_n^{n-1}) +
\varphi(a_1 a_n^{n-2}) (\varphi(a_n)t) +
\ldots +
(\varphi(a_n) t)^n = 0$.
\end{proof}

\noindent
Le lemme suivant est, en un certain sens, le lemme clé de cette
section.

\begin{lemma}
\label{lemma-make-integral-trick}
Soit $R$ un anneau. Soit $\varphi : R[x] \to S$ un
homomorphisme d'anneaux. Soit $t \in S$.
Supposons (a) que $t$ soit entier sur $R[x]$,
et (b) qu'il existe un polynôme unitaire $p \in R[x]$ tel que
$t \varphi(p) \in \Im(\varphi)$. Alors il
existe $q \in R[x]$ tel que $t - \varphi(q)$
soit entier sur $R$.
\end{lemma}

\begin{proof}
Écrivons $t \varphi(p) = \varphi(r)$ pour un certain $r \in R[x]$.
Par division euclidienne, écrivons $r = qp + r'$ avec
$q, r' \in R[x]$ et $\deg(r') < \deg(p)$. On peut remplacer
$t$ par $t - \varphi(q)$, qui est encore entier sur
$R[x]$, et l'on obtient alors $t \varphi(p) = \varphi(r')$.
Dans l'anneau localisé $S_t$, cette égalité s'écrit
$\varphi(p) - (1/t) \varphi(r') = 0$.
Cela implique que $\varphi(x)$ définit un élément de la
localisation $S_t$ qui est entier sur
$\varphi(R)[1/t] \subset S_t$. D'autre part,
$t$ est entier sur le sous-anneau $\varphi(R)[\varphi(x)] \subset S$.
En combinant ces faits, on conclut que $t$ est entier sur
le sous-anneau $\varphi(R)[1/t] \subset S_t$ ; voir le lemme
\ref{lemma-integral-transitive}. Autrement dit,
il existe une égalité de la forme
$$
t^d + \sum\nolimits_{i < d}
\left(\sum\nolimits_{j = 0, \ldots, n_i} \varphi(r_{i, j})/t^j\right) t^i = 0
$$
dans $S_t$, où $r_{i, j} \in R$. Cela signifie que
$t^{d + N} +
\sum_{i < d} \sum_{j = 0, \ldots, n_i} \varphi(r_{i, j}) t^{i + N - j} = 0$
dans $S$ pour un certain $N$ assez grand. Autrement dit,
$t$ est entier sur $R$.
\end{proof}

\begin{lemma}
\label{lemma-combine-lemmas}
Soit $R$ un anneau. Soit $\varphi : R[x] \to S$ un
homomorphisme d'anneaux. Soit $t \in S$. Supposons que $t$ soit entier
sur $R[x]$. Soit $p \in R[x]$, $p = a_0 + a_1x + \ldots +
a_k x^k$ tel que $t \varphi(p) \in \Im(\varphi)$.
Alors il existe $q \in R[x]$ et $n \geq 0$
tels que $\varphi(a_k)^n t - \varphi(q)$ soit entier
sur $R$.
\end{lemma}

\begin{proof}
Soient $R'$ et $S'$ les localisés de $R$ et $S$, respectivement en $a_k$ et en son image.
Soit $\varphi' : R'[x] \to S'$ l'homomorphisme localisé induit par $\varphi$.
Soit $t' \in S'$ l'image de $t$. Posons $p' = p/a_k \in R'[x]$.
Alors $t' \varphi'(p') \in \Im(\varphi')$, puisque $t \varphi(p) \in \Im(\varphi)$.
Comme $p'$ est unitaire, le
lemme \ref{lemma-make-integral-trick} donne un $q' \in R'[x]$
tel que $t' - \varphi'(q')$ soit entier sur $R'$.
On peut choisir $n \geq 0$ et un élément $q \in R[x]$
tels que $a_k^n q'$ soit l'image de $q$.
Alors $\varphi(a_k)^n t - \varphi(q)$ est un élément de $S$ dont
l'image dans $S'$ est entière sur $R'$.
D'après le lemme \ref{lemma-integral-closure-localize},
il existe $m \geq 0$ tel que
$\varphi(a_k)^m(\varphi(a_k)^n t - \varphi(q))$ soit entier sur $R$.
Ainsi, $\varphi(a_k)^{m + n}t - \varphi(a_k^m q)$
est entier sur $R$, ce qui donne le résultat.
\end{proof}

\begin{situation}
\label{situation-one-transcendental-element}
Soit $R$ un anneau.
Soit $\varphi : R[x] \to S$ un homomorphisme d'anneaux fini.
Soit
$$
J = \{ g \in S \mid gS \subset \Im(\varphi)\}
$$
l'``idéal conducteur'' de $\varphi$.
Supposons que le sous-anneau $\varphi(R) \subset S$ soit intégralement clos dans $S$.
\end{situation}

\begin{lemma}
\label{lemma-leading-coefficient-in-J}
Plaçons-nous dans la situation \ref{situation-one-transcendental-element}.
Supposons que $u \in S$, que $a_0, \ldots, a_k \in R$ et que
$u \varphi(a_0 + a_1x + \ldots + a_k x^k) \in J$.
Alors il existe $m \geq 0$ tel que
$u \varphi(a_k)^m \in J$.
\end{lemma}

\begin{proof}
Supposons que $S$ soit engendré par $t_1, \ldots, t_n$
comme $R[x]$-module. Dans ce cas,
$J = \{ g \in S \mid gt_i \in \Im(\varphi)\text{ pour tout }i\}$.
Notons que chaque élément $u t_i$ est entier sur
$R[x]$ ; voir le lemme \ref{lemma-finite-is-integral}.
On a $\varphi(a_0 + a_1x + \ldots + a_k x^k) u t_i \in
\Im(\varphi)$. D'après le lemme \ref{lemma-combine-lemmas}, pour
chaque $i$, il existe un entier $n_i$ et un élément
$q_i \in R[x]$ tels que $\varphi(a_k^{n_i}) u t_i - \varphi(q_i)$
soit entier sur $R$. Par hypothèse, cet élément appartient à $\varphi(R)$,
et donc $\varphi(a_k^{n_i}) u t_i \in \Im(\varphi)$.
Il s'ensuit que $m = \max\{n_1, \ldots, n_n\}$ convient.
\end{proof}

\begin{lemma}
\label{lemma-all-coefficients-in-J}
Plaçons-nous dans la situation \ref{situation-one-transcendental-element}.
Supposons que $u \in S$, que $a_0, \ldots, a_k \in R$ et que
$u \varphi(a_0 + a_1x + \ldots + a_k x^k) \in \sqrt{J}$.
Alors $u \varphi(a_i) \in \sqrt{J}$ pour tout $i$.
\end{lemma}

\begin{proof}
Sous les hypothèses du lemme, on a
$u^n \varphi(a_0 + a_1x + \ldots + a_k x^k)^n \in J$ pour
un certain $n \geq 1$. D'après le lemme \ref{lemma-leading-coefficient-in-J},
on en déduit que $u^n \varphi(a_k^{nm}) \in J$ pour un certain $m \geq 1$.
Ainsi, $u \varphi(a_k) \in \sqrt{J}$, et donc
$u \varphi(a_0 + a_1x + \ldots + a_k x^k) - u \varphi(a_k x^k) =
u \varphi(a_0 + a_1x + \ldots + a_{k-1} x^{k-1}) \in \sqrt{J}$.
On conclut par récurrence sur $k$.
\end{proof}

\noindent
Ce lemme suggère la définition suivante.

\begin{definition}
\label{definition-strongly-transcendental}
Étant donnés une inclusion d'anneaux $R \subset S$ et
un élément $x \in S$, on dit que $x$ est
{\it fortement transcendant sur $R$} si,
dès que $u(a_0 + a_1 x + \ldots + a_k x^k) = 0$
avec $u \in S$ et $a_i \in R$, alors
on a $ua_i = 0$ pour tout $i$.
\end{definition}

\noindent
Notons que, si $S$ est un anneau intègre, cela revient à
dire que $x$, considéré comme élément du corps des fractions de
$S$, est transcendant sur le corps des fractions de $R$.

\begin{lemma}
\label{lemma-reduced-strongly-transcendental-minimal-prime}
Supposons que $R \subset S$ soit une inclusion d'anneaux réduits
et que $x \in S$ soit fortement transcendant sur $R$.
Soit $\mathfrak q \subset S$ un idéal premier minimal
et soit $\mathfrak p = R \cap \mathfrak q$.
Alors l'image de $x$ dans $S/\mathfrak q$ est fortement
transcendante sur le sous-anneau $R/\mathfrak p$.
\end{lemma}

\begin{proof}
Supposons que $u(a_0 + a_1x + \ldots + a_k x^k) \in \mathfrak q$.
D'après le lemme \ref{lemma-minimal-prime-reduced-ring},
l'anneau local $S_{\mathfrak q}$ est un corps,
et donc $u(a_0 + a_1x + \ldots + a_k x^k) $ est nul
dans $S_{\mathfrak q}$. Ainsi, $uu'(a_0 + a_1x + \ldots + a_k x^k) = 0$
pour un certain $u' \in S$, $u' \not\in \mathfrak q$.
Puisque $x$ est fortement transcendant sur $R$, on obtient
$uu'a_i = 0$ pour tout $i$. Cela implique à son tour
que $ua_i \in \mathfrak q$.
\end{proof}

\begin{lemma}
\label{lemma-domains-transcendental-not-quasi-finite}
Supposons que $R\subset S$ soit une inclusion d'anneaux intègres et
soit $x \in S$. Supposons que $x$ soit (fortement) transcendant sur $R$
et que $S$ soit fini sur $R[x]$. Alors l'homomorphisme $R \to S$ n'est
quasi-fini en aucun idéal premier de $S$.
\end{lemma}

\begin{proof}
Commençons par le cas où $R$ est normal ; voir la
définition \ref{definition-ring-normal}.
D'après le lemme \ref{lemma-polynomial-ring-normal},
on voit que $R[x]$ est normal.
Prenons un idéal premier $\mathfrak q \subset S$
et posons $\mathfrak p = R \cap \mathfrak q$.
Supposons que l'extension $\kappa(\mathfrak p)
\subset \kappa(\mathfrak q)$ soit finie.
Ce serait le cas si l'homomorphisme $R \to S$ était
quasi-fini en $\mathfrak q$.
Soit $\mathfrak r = R[x] \cap \mathfrak q$.
Alors, puisque $\kappa(\mathfrak p)
\subset \kappa(\mathfrak r) \subset \kappa(\mathfrak q)$,
on voit que l'extension $\kappa(\mathfrak p)
\subset \kappa(\mathfrak r)$ est également finie.
Ainsi, l'inclusion $\mathfrak r \supset \mathfrak p R[x]$
est stricte. Par le théorème de descente appliqué à $R[x] \subset S$,
voir la proposition \ref{proposition-going-down-normal-integral},
on trouve un idéal premier $\mathfrak q' \subset \mathfrak q$
au-dessus de l'idéal premier $\mathfrak pR[x]$. Par conséquent,
la fibre $\Spec(S \otimes_R \kappa(\mathfrak p))$
contient un point distinct de $\mathfrak q$,
à savoir $\mathfrak q'$, dont l'adhérence contient $\mathfrak q$ ; ainsi,
$\mathfrak q$ n'est pas isolé dans sa fibre.

\medskip\noindent
Si $R$ n'est pas normal, introduisons une chaîne $R \subset R' \subset K$
dans laquelle $R'$ désigne la clôture intégrale de $R$ dans son corps des fractions
$K$. De même, introduisons $S \subset S' \subset L$, où $S'$ est le sous-anneau du
corps des fractions $L$ de $S$ engendré par $R'$ et
$S$. Notons que, par construction, l'homomorphisme $S \otimes_R R'
\to S'$ est surjectif. Cela implique que $R'[x] \subset S'$
est fini. De plus, l'inclusion $S \subset S'$
induit sur $\Spec$ un morphisme surjectif ; voir le
lemme \ref{lemma-integral-overring-surjective}.
On conclut à l'aide du lemme \ref{lemma-four-rings} et du cas normal
que nous venons de traiter.
\end{proof}

\begin{lemma}
\label{lemma-reduced-strongly-transcendental-not-quasi-finite}
Supposons que $R \subset S$ soit une inclusion d'anneaux réduits.
Supposons que $x \in S$ soit fortement transcendant sur $R$
et que $S$ soit fini sur $R[x]$. Alors l'homomorphisme $R \to S$ n'est
quasi-fini en aucun idéal premier de $S$.
\end{lemma}

\begin{proof}
Soit $\mathfrak q \subset S$ un idéal premier quelconque.
Choisissons un idéal premier minimal $\mathfrak q' \subset \mathfrak q$.
D'après les lemmes
\ref{lemma-reduced-strongly-transcendental-minimal-prime} et
\ref{lemma-domains-transcendental-not-quasi-finite},
l'extension $R/(R \cap \mathfrak q') \subset
S/\mathfrak q'$ n'est pas quasi-finie en l'idéal premier correspondant
à $\mathfrak q$. D'après le lemme \ref{lemma-four-rings},
l'homomorphisme $R \to S$ n'est pas quasi-fini
en $\mathfrak q$.
\end{proof}

\begin{lemma}
\label{lemma-quasi-finite-monogenic}
Soit $R$ un anneau. Soit $S = R[x]/I$.
Soit $\mathfrak q \subset S$ un idéal premier.
Supposons que l'homomorphisme $R \to S$ soit quasi-fini en $\mathfrak q$.
Soit $S' \subset S$ la clôture intégrale de $R$ dans $S$.
Alors il existe un élément
$g \in S'$, $g \not\in \mathfrak q$ tel que
$S'_g \cong S_g$.
\end{lemma}

\begin{proof}
Soit $\mathfrak p$ l'image de $\mathfrak q$ dans $\Spec(R)$.
Il existe $f \in I$, $f = a_nx^n + \ldots + a_0$ tel que
$a_i \not \in \mathfrak p$ pour un certain $i$. En effet, sinon, l'anneau de la fibre
$S \otimes_R \kappa(\mathfrak p)$ serait $\kappa(\mathfrak p)[x]$
et l'homomorphisme ne serait quasi-fini en aucun idéal premier situé
au-dessus de $\mathfrak p$. On en déduit qu'il existe une relation
$b_m x^m + \ldots + b_0 = 0$ avec $b_j \in S'$, $j = 0, \ldots, m$
et $b_j \not \in \mathfrak q \cap S'$ pour un certain $j$.
Nous démontrons le lemme par récurrence sur $m$. Le cas initial
$m = 0$ ne se présente pas (car les assertions $b_0 = 0$ et
$b_0 \not \in \mathfrak q$ sont contradictoires).

\medskip\noindent
Supposons que $b_m \not \in \mathfrak q$. Dans ce cas, $x$ est entier
sur $S'_{b_m}$ ; en fait, $b_mx \in S'$ d'après le
lemme \ref{lemma-make-integral-trivial}.
Ainsi, l'homomorphisme injectif $S'_{b_m} \to S_{b_m}$ est également surjectif, donc
c'est l'isomorphisme voulu.

\medskip\noindent
Supposons que $b_m \in \mathfrak q$. Dans ce cas, on a $b_mx \in S'$ d'après le
lemme \ref{lemma-make-integral-trivial}.
Posons $b'_{m - 1} = b_mx + b_{m - 1}$. Alors
$$
b'_{m - 1}x^{m - 1} + b_{m - 2}x^{m - 2} + \ldots + b_0 = 0
$$
Puisque $b'_{m - 1}$ est congru à $b_{m - 1}$ modulo $S' \cap \mathfrak q$,
on voit que l'un des éléments
$b'_{m - 1}, b_{m - 2}, \ldots, b_0$ n'appartient toujours pas à $S' \cap \mathfrak q$.
On conclut donc par récurrence sur $m$.
\end{proof}

\begin{theorem}[Théorème principal de Zariski]
\label{theorem-main-theorem}
Soit $R$ un anneau. Soit $R \to S$ l'homomorphisme structural d'une $R$-algèbre de type fini.
Soit $S' \subset S$ la clôture intégrale de $R$ dans $S$.
Soit $\mathfrak q \subset S$ un idéal premier de $S$.
Si l'homomorphisme $R \to S$ est quasi-fini en $\mathfrak q$, alors
il existe $g \in S'$, $g \not \in \mathfrak q$
tel que $S'_g \cong S_g$.
\end{theorem}

\begin{proof}
Il existe un nombre fini d'éléments
$x_1, \ldots, x_n \in S$ tels que $S$ soit fini
sur la sous-$R$-algèbre engendrée par $x_1, \ldots, x_n$. (Par
exemple, on peut prendre des générateurs de $S$ sur $R$.) Nous démontrons le théorème
par récurrence sur le plus petit entier $n$ ayant cette propriété.

\medskip\noindent
Le cas $n = 0$ est trivial, car alors $S' = S$ ;
voir le lemme \ref{lemma-finite-is-integral}.

\medskip\noindent
Supposons que $n = 1$. On peut remplacer $R$ par sa clôture intégrale dans $S$
(le lemme \ref{lemma-quasi-finite-permanence} garantit que l'homomorphisme $R \to S$
reste quasi-fini en $\mathfrak q$). On peut donc supposer que
l'inclusion $R \subset S$ identifie l'anneau de base à un sous-anneau intégralement clos dans $S$, autrement dit que $R = S'$.
Considérons l'homomorphisme $\varphi : R[x] \to S$, $x \mapsto x_1$.
(Nous verrons ci-dessous que $\varphi$ n'est pas injectif.)
Par hypothèse, $\varphi$ est fini. Nous sommes donc dans la situation
\ref{situation-one-transcendental-element}.
Soit $J \subset S$ l'``idéal conducteur'' défini
dans la situation \ref{situation-one-transcendental-element}.
Considérons le diagramme
$$
\xymatrix{
R[x] \ar[r] & S \ar[r] & S/\sqrt{J} & R/(R \cap \sqrt{J})[x] \ar[l]
\\
& R \ar[lu] \ar[r] \ar[u] & R/(R \cap \sqrt{J}) \ar[u] \ar[ru] &
}
$$
D'après le lemme \ref{lemma-all-coefficients-in-J},
l'image de $x$ dans le quotient $S/\sqrt{J}$
est fortement transcendante sur $R/ (R \cap \sqrt{J})$.
Ainsi, d'après le lemme \ref{lemma-reduced-strongly-transcendental-not-quasi-finite},
l'homomorphisme d'anneaux $R/ (R \cap \sqrt{J}) \to S/\sqrt{J}$
n'est quasi-fini en aucun idéal premier de $S/\sqrt{J}$.
D'après le lemme \ref{lemma-four-rings}, on en déduit que $\mathfrak q$
n'appartient pas à $V(J) \subset \Spec(S)$.
Il existe donc un élément $s \in J$
tel que $s \not\in \mathfrak q$. Par définition de $J$, on peut écrire
$s = \varphi(f)$ pour un polynôme $f \in R[x]$.
Soit $I = \Ker(\varphi : R[x] \to S)$. Puisque $\varphi(f) \in J$,
on obtient $(R[x]/I)_f \cong S_{\varphi(f)}$. De plus, $s \not \in \mathfrak q$
signifie que $f \not \in \varphi^{-1}(\mathfrak q)$. Ainsi,
$\varphi^{-1}(\mathfrak q)/I$ est un idéal premier de $R[x]/I$
en lequel l'homomorphisme $R \to R[x]/I$ est quasi-fini ; voir le
lemme \ref{lemma-quasi-finite-local}. Notons que $R$ est intégralement clos
dans $R[x]/I$ puisque $R$ est intégralement clos dans $S$. D'après le
lemme \ref{lemma-quasi-finite-monogenic},
il existe un élément $h \in R$, $h \not \in R \cap \mathfrak q$
tel que $R_h \cong (R[x]/I)_h$. Ainsi,
$(R[x]/I)_{fh} = S_{\varphi(fh)}$ est isomorphe à une
localisation principale $R_{h'}$ de $R$ pour un certain
$h' \in R$, $h' \not \in \mathfrak q$.

\medskip\noindent
Supposons que $n > 1$. Considérons le sous-anneau $R' \subset S$
qui est la clôture intégrale de $R[x_1, \ldots, x_{n-1}]$
dans $S$. D'après le lemme \ref{lemma-quasi-finite-permanence}, l'extension
$S/R'$ est quasi-finie en $\mathfrak q$. Notons également
que $S$ est fini sur $R'[x_n]$.
D'après le cas $n = 1$ ci-dessus, il existe $g' \in R'$
tel que $g' \not \in \mathfrak q$ et que
$(R')_{g'} \cong S_{g'}$. À ce stade, on ne peut pas
appliquer l'hypothèse de récurrence à $R \to R'$, car $R'$ peut ne pas être de type fini sur $R$.
Puisque $S$ est de type fini sur $R$, on en déduit en particulier
que $(R')_{g'}$ est de type fini sur $R$. Disons que
les éléments $g'$ et $y_1/(g')^{n_1}, \ldots, y_N/(g')^{n_N}$,
où $y_i \in R'$,
engendrent $(R')_{g'}$ sur $R$. Soit $R''$ la sous-$R$-algèbre
de $R'$ engendrée par $x_1, \ldots, x_{n-1}, y_1, \ldots, y_N, g'$.
Elle vérifie $(R'')_{g'} \cong S_{g'}$. La surjectivité
résulte du choix des $y_i$ ; l'injectivité résulte de
$R'' \subset R'$ et de l'exactitude de la localisation. Notons que
$R''$ est fini sur $R[x_1, \ldots, x_{n-1}]$ en raison
du choix de $R'$ ; voir le lemme \ref{lemma-characterize-integral}.
Soit $\mathfrak q'' = R'' \cap \mathfrak q$. Puisque
$(R'')_{\mathfrak q''} = S_{\mathfrak q}$, on voit que
l'homomorphisme $R \to R''$ est quasi-fini en $\mathfrak q''$ ; voir le
lemme \ref{lemma-isolated-point-fibre}.
Appliquons l'hypothèse de récurrence à $R \to R''$, à $\mathfrak q''$
et à $x_1, \ldots, x_{n-1} \in R''$ ; on obtient un sous-anneau
$R''' \subset R''$ entier sur $R$ et un
élément $g'' \in R'''$, $g'' \not \in \mathfrak q''$
tel que $(R''')_{g''} \cong (R'')_{g''}$. Écrivons l'image de $g'$ dans
$(R'')_{g''}$ sous la forme $g'''/(g'')^n$ pour un certain $g''' \in  R'''$.
Posons $g = g''g''' \in R'''$. Il est alors clair que $g \not\in
\mathfrak q$ et que $(R''')_g \cong S_g$. Comme, par construction,
on a $R''' \subset S'$, on a également $S'_g \cong S_g$, ce qui donne le résultat.
\end{proof}

\begin{lemma}
\label{lemma-quasi-finite-open}
Soit $R \to S$ un homomorphisme d'anneaux de type fini.
L'ensemble des points $\mathfrak q$ de $\Spec(S)$ en lesquels
$S/R$ est quasi-fini est ouvert dans $\Spec(S)$.
\end{lemma}

\begin{proof}
Soit $\mathfrak q \subset S$ un idéal premier en lequel l'homomorphisme d'anneaux
est quasi-fini. D'après le théorème \ref{theorem-main-theorem},
il existe un homomorphisme d'anneaux entier $R \to S'$, $S' \subset S$,
et un élément $g \in S'$, $g\not \in \mathfrak q$ tel que
$S'_g \cong S_g$. Puisque $S$, et donc $S_g$, sont de type fini
sur $R$, on peut trouver un nombre fini d'éléments
$y_1, \ldots, y_N$ de $S'$ tels que $S''_g \cong S_g$,
où $S'' \subset S'$ est la sous-$R$-algèbre engendrée
par $g, y_1, \ldots, y_N$. Puisque $S''$ est fini sur $R$
(voir le lemme \ref{lemma-characterize-integral}), on voit que
$S''$ est quasi-fini sur $R$ (voir le lemme \ref{lemma-quasi-finite}).
On voit facilement que cela implique que $S''_g$ est quasi-fini sur $R$,
par exemple parce que la propriété d'être quasi-fini en un idéal premier ne dépend
que de l'anneau local en cet idéal. Ainsi, $S_g$ est quasi-fini
sur $R$. Le même raisonnement montre que l'homomorphisme $R \to S$ est quasi-fini
en tout idéal premier de $S$ appartenant à $D(g)$.
\end{proof}

\begin{lemma}
\label{lemma-quasi-finite-open-integral-closure}
Soit $R \to S$ un homomorphisme d'anneaux de type fini.
Supposons que $S$ soit quasi-fini sur $R$.
Soit $S' \subset S$ la clôture intégrale de $R$ dans $S$. Alors
\begin{enumerate}
\item Le morphisme $\Spec(S) \to \Spec(S')$ est un homéomorphisme
sur un ouvert,
\item si $g \in S'$ et si $D(g)$ est contenu dans l'image
du morphisme, alors $S'_g \cong S_g$, et
\item il existe une $R$-algèbre finie $S'' \subset S'$
telle que (1) et (2) soient vérifiées pour l'homomorphisme d'anneaux
$S'' \to S$.
\end{enumerate}
\end{lemma}

\begin{proof}
Puisque $S/R$ est quasi-fini, on peut appliquer le
théorème \ref{theorem-main-theorem} à
chaque point $\mathfrak q$ de $\Spec(S)$.
Comme $\Spec(S)$ est quasi-compact (voir le
lemme \ref{lemma-quasi-compact}), on peut choisir
un nombre fini de $g_i \in S'$, $i = 1, \ldots, n$
tels que $S'_{g_i} = S_{g_i}$ et que
$g_1, \ldots, g_n$ engendrent l'idéal unité de $S$
(autrement dit, les ouverts principaux de $\Spec(S)$ associés
à $g_1, \ldots, g_n$ recouvrent tout $\Spec(S)$).

\medskip\noindent
Supposons que $D(g) \subset \Spec(S')$
soit contenu dans l'image. Alors $D(g) \subset \bigcup D(g_i)$.
Autrement dit, $g_1, \ldots, g_n$ engendrent l'idéal unité de
$S'_g$. Notons que $S'_{gg_i} \cong S_{gg_i}$ par notre choix
des $g_i$. Par conséquent, $S'_g \cong S_g$ d'après le lemme \ref{lemma-cover}.

\medskip\noindent
Construisons une algèbre finie $S'' \subset S'$ comme
au point (3). Pour cela, notons que chaque $S'_{g_i} \cong S_{g_i}$
est une $R$-algèbre de type fini. Pour chaque $i$, choisissons
des éléments $y_{ij} \in S'$ tels que chaque
$S'_{g_i}$ soit engendrée comme $R$-algèbre par $1/g_i$
et les éléments $y_{ij}$. Posons alors $S''$
égal à la sous-$R$-algèbre de $S'$ engendrée par tous les $g_i$
et tous les $y_{ij}$. Nous omettons les détails.
\end{proof}




\section{Applications du théorème principal de Zariski}
\label{section-apply-main-theorem}

\noindent
Voici une application immédiate qui caractérise les homomorphismes finis
parmi les homomorphismes quasi-finis d'anneaux semi-locaux de dimension $1$ :
ce sont ceux pour lesquels l'égalité est toujours vérifiée dans la
formule du lemme \ref{lemma-finite-extension-dim-1}.

\begin{lemma}
\label{lemma-quasi-finite-extension-dim-1}
Soit $A \subset B$ une extension d'anneaux intègres. Supposons que
\begin{enumerate}
\item $A$ soit un anneau local noethérien de dimension $1$,
\item l'homomorphisme $A \to B$ soit de type fini, et
\item l'extension induite $L/K$ des corps des fractions soit finie.
\end{enumerate}
Alors $B$ est semi-local.
Soit $x \in \mathfrak m_A$, $x \not = 0$.
Soient $\mathfrak m_i$, $i = 1, \ldots, n$,
les idéaux maximaux de $B$.
Alors
$$
[L : K]\text{ord}_A(x)
\geq
\sum\nolimits_i
[\kappa(\mathfrak m_i) : \kappa(\mathfrak m_A)]
\text{ord}_{B_{\mathfrak m_i}}(x)
$$
où $\text{ord}$ est défini comme dans la définition \ref{definition-ord}.
On a égalité si et seulement si l'homomorphisme $A \to B$ est fini.
\end{lemma}

\begin{proof}
L'anneau $B$ est semi-local d'après le lemme \ref{lemma-finite-in-codim-1}.
Soit $B'$ la clôture intégrale de $A$ dans $B$. D'après le
lemme \ref{lemma-quasi-finite-open-integral-closure},
on peut trouver une sous-$A$-algèbre finie $C \subset B'$ telle que,
en posant $\mathfrak n_i = C \cap \mathfrak m_i$, on ait
$C_{\mathfrak n_i} \cong B_{\mathfrak m_i}$ et que les idéaux premiers
$\mathfrak n_1, \ldots, \mathfrak n_n$ soient deux à deux distincts.
L'anneau $C$ est semi-local d'après le lemme \ref{lemma-finite-in-codim-1}.
Soient $\mathfrak p_j$, $j = 1, \ldots, m$, les autres idéaux
maximaux de $C$ (les ``points manquants''). D'après le
lemme \ref{lemma-finite-extension-dim-1}, on a
$$
\text{ord}_A(x^{[L : K]}) =
\sum\nolimits_i
[\kappa(\mathfrak n_i) : \kappa(\mathfrak m_A)]
\text{ord}_{C_{\mathfrak n_i}}(x)
+
\sum\nolimits_j
[\kappa(\mathfrak p_j) : \kappa(\mathfrak m_A)]
\text{ord}_{C_{\mathfrak p_j}}(x)
$$
d'où l'inégalité. En cas d'égalité, on en déduit que
$m = 0$ (il n'y a pas de ``points manquants''). Ainsi, $C \subset B$ est une inclusion
d'anneaux semi-locaux qui induit une bijection sur les idéaux maximaux et
un isomorphisme sur toutes les localisations aux idéaux maximaux. Ainsi, si $b \in B$,
alors $I = \{x \in C \mid xb \in C\}$ est un idéal de $C$ qui n'est
contenu dans aucun idéal maximal de $C$ ; par conséquent, $I = C$,
donc $b \in C$. Ainsi, $B = C$ et $B$ est fini sur $A$.
\end{proof}

\noindent
Voici une application plus classique du théorème principal de Zariski à la
structure des homomorphismes locaux d'anneaux locaux.

\begin{lemma}
\label{lemma-essentially-finite-type-fibre-dim-zero}
Soit $(R, \mathfrak m_R) \to (S, \mathfrak m_S)$ un homomorphisme local
d'anneaux locaux. Supposons que
\begin{enumerate}
\item $R \to S$ soit essentiellement de type fini,
\item l'extension $\kappa(\mathfrak m_R) \subset \kappa(\mathfrak m_S)$ soit finie, et
\item $\dim(S/\mathfrak m_RS) = 0$.
\end{enumerate}
Alors $S$ est la localisation d'une $R$-algèbre finie.
\end{lemma}

\begin{proof}
Soit $S'$ une $R$-algèbre de type fini telle que $S = S'_{\mathfrak q'}$
pour un certain idéal premier $\mathfrak q'$ de $S'$. D'après la
définition \ref{definition-quasi-finite},
on voit que l'homomorphisme $R \to S'$ est quasi-fini en $\mathfrak q'$.
Après avoir remplacé $S'$ par $S'_{g'}$ pour un certain
$g' \in S'$, $g' \not \in \mathfrak q'$, on peut supposer que $R \to S'$ est
quasi-fini ; voir le
lemme \ref{lemma-quasi-finite-open}.
Alors, d'après le
lemme \ref{lemma-quasi-finite-open-integral-closure},
il existe une $R$-algèbre finie $S''$ et des éléments
$g' \in S'$, $g' \not \in \mathfrak q'$ et $g'' \in S''$
tels que $S'_{g'} \cong S''_{g''}$ comme $R$-algèbres.
Cela démontre le lemme.
\end{proof}

\begin{lemma}
\label{lemma-completion-at-quasi-finite-prime}
Soit $R \to S$ un homomorphisme d'anneaux et soit $\mathfrak q$ un idéal premier de $S$
situé au-dessus d'un idéal premier $\mathfrak p$ de $R$. Si
\begin{enumerate}
\item $R$ est noethérien,
\item l'homomorphisme $R \to S$ est de type fini, et
\item l'homomorphisme $R \to S$ est quasi-fini en $\mathfrak q$,
\end{enumerate}
alors $R_\mathfrak p^\wedge \otimes_R S = S_\mathfrak q^\wedge \times B$
pour une certaine $R_\mathfrak p^\wedge$-algèbre $B$.
\end{lemma}

\begin{proof}
Il existe une $R$-algèbre finie $S' \subset S$ et un élément
$g \in S'$, $g \not \in \mathfrak q' = S' \cap \mathfrak q$
tels que $S'_g = S_g$ et, en particulier,
$S'_{\mathfrak q'} = S_\mathfrak q$ ; voir le
lemme \ref{lemma-quasi-finite-open-integral-closure}.
On a
$$
R_\mathfrak p^\wedge \otimes_R S' = (S'_{\mathfrak q'})^\wedge \times B'
$$
d'après le lemme \ref{lemma-completion-finite-extension}. Remarquons que, dans cette
décomposition en produit, $g$ s'envoie sur un couple $(u, b')$ où
$u \in (S'_{\mathfrak q'})^\wedge$ est une unité puisque $g \not \in \mathfrak q'$.
La décomposition en produit de $R_\mathfrak p^\wedge \otimes_R S'$
induit une décomposition en produit
$$
R_\mathfrak p^\wedge \otimes_R S = A \times B
$$
Puisque $S'_g = S_g$, on a également
$(R_\mathfrak p^\wedge \otimes_R S')_g = (R_\mathfrak p^\wedge \otimes_R S)_g$,
et puisque $g \mapsto (u, b')$, où $u$ est une unité, on voit que
$(S'_{\mathfrak q'})^\wedge = A$. Comme l'isomorphisme
$S'_{\mathfrak q'} = S_\mathfrak q$ induit un isomorphisme des
complétés, on obtient également $A = S_\mathfrak q^\wedge$.
La démonstration est terminée, à ceci près qu'il faut vérifier
que l'homomorphisme induit
$\phi : R_\mathfrak p^\wedge \otimes_R S \to A = S_\mathfrak q^\wedge$
est l'homomorphisme naturel. Pour les éléments de la forme $x \otimes 1$,
où $x \in R_\mathfrak p^\wedge$, c'est clair puisque l'homomorphisme naturel
$R_\mathfrak p^\wedge \to S_\mathfrak q^\wedge$ se factorise par
$(S'_{\mathfrak q'})^\wedge$. Pour les éléments de la forme $1 \otimes y$,
où $y \in S$, on peut remarquer que, pour un certain $n \geq 1$, l'élément
$g^ny$ est l'image d'un certain $y' \in S'$. Ainsi, $\phi(1 \otimes g^ny)$
est l'image de $y'$ par la composée
$S' \to (S'_{\mathfrak q'})^\wedge \to S_\mathfrak q^\wedge$. Cette image est
égale à celle de $g^ny$ par l'homomorphisme $S \to S_\mathfrak q^\wedge$.
Puisque $g$ s'envoie sur une unité,
il s'ensuit que $\phi(1 \otimes y)$ est bien l'image de $y$
par l'homomorphisme $S \to S_\mathfrak q^\wedge$.
\end{proof}




































\section{Dimension des fibres}
\label{section-dimension-fibres}

\noindent
Nous étudions le comportement des dimensions des fibres à l'aide du
théorème principal de Zariski. Rappelons que nous avons défini la
dimension $\dim_x(X)$ d'un espace topologique $X$ en un point $x$
dans Topologie, définition \ref{topology-definition-Krull}.

\begin{definition}
\label{definition-relative-dimension}
Supposons que $R \to S$ soit de type fini, et soit
$\mathfrak q \subset S$ un idéal premier au-dessus d'un idéal premier
$\mathfrak p$ de $R$.
Nous définissons la {\it dimension relative
de $S/R$ en $\mathfrak q$}, notée
$\dim_{\mathfrak q}(S/R)$, comme la dimension
de $\Spec(S \otimes_R \kappa(\mathfrak p))$
au point correspondant à $\mathfrak q$. Nous désignons par
$\dim(S/R)$ la borne supérieure des $\dim_{\mathfrak q}(S/R)$
lorsque $\mathfrak q$ parcourt tous les idéaux premiers. On l'appelle la
{\it dimension relative de} $S/R$.
\end{definition}

\noindent
En particulier, $R \to S$ est quasi-fini en $\mathfrak q$ si et
seulement si $\dim_{\mathfrak q}(S/R) = 0$. Le lemme suivant
est essentiellement une reformulation du théorème principal de Zariski.

\begin{lemma}
\label{lemma-quasi-finite-over-polynomial-algebra}
Soit $R \to S$ un homomorphisme d'anneaux de type fini.
Soit $\mathfrak q \subset S$ un idéal premier.
Supposons que $\dim_{\mathfrak q}(S/R) = n$.
Il existe un élément $g \in S$, $g \not\in \mathfrak q$
tel que $S_g$ soit quasi-fini sur une
algèbre de polynômes $R[t_1, \ldots, t_n]$.
\end{lemma}

\begin{proof}
L'anneau $\overline{S} = S \otimes_R \kappa(\mathfrak p)$ est
de type fini sur $\kappa(\mathfrak p)$.
Soit $\overline{\mathfrak q}$ l'idéal premier de $\overline{S}$
correspondant à $\mathfrak q$.
Par définition de
la dimension d'un espace topologique en un point, il existe
un ouvert $U \subset \Spec(\overline{S})$ tel que
$\overline{q} \in U$ et $\dim(U) = n$.
Puisque la topologie sur $\Spec(\overline{S})$ est
induite par la topologie sur $\Spec(S)$ (voir la
remarque \ref{remark-fundamental-diagram}), nous pouvons trouver
un élément $g \in S$, $g \not \in \mathfrak q$, dont l'image
$\overline{g} \in \overline{S}$ est telle que
$D(\overline{g}) \subset U$.
Ainsi, après remplacement de $S$ par $S_g$, on voit que
$\dim(\overline{S}) = n$.

\medskip\noindent
Choisissons ensuite des générateurs $x_1, \ldots, x_N$ de $S$ comme $R$-algèbre. D'après le
lemme \ref{lemma-Noether-normalization},
il existe des éléments $y_1, \ldots, y_n$ dans la sous-algèbre sur $\mathbf{Z}$ de $S$
engendrée par $x_1, \ldots, x_N$ tels que l'homomorphisme
$R[t_1, \ldots, t_n] \to S$, $t_i \mapsto y_i$ a pour propriété
que $\kappa(\mathfrak p)[t_1\ldots, t_n] \to \overline{S}$
est fini. En particulier, $S$ est quasi-fini sur $R[t_1, \ldots, t_n]$
en $\mathfrak q$. Par conséquent, d'après le lemme \ref{lemma-quasi-finite-open},
nous pouvons remplacer $S$ par $S_g$ pour un certain $g\in S$, $g \not \in \mathfrak q$
tel que $R[t_1, \ldots, t_n] \to S$ soit quasi-fini.
\end{proof}

\begin{lemma}
\label{lemma-refined-quasi-finite-over-polynomial-algebra}
Soit $R \to S$ un homomorphisme d'anneaux. Soit $\mathfrak q \subset S$
un idéal premier au-dessus de l'idéal premier $\mathfrak p$ de $R$.
Supposons que
\begin{enumerate}
\item $R \to S$ soit de type fini,
\item $\dim_{\mathfrak q}(S/R) = n$, et
\item $\text{trdeg}_{\kappa(\mathfrak p)}\kappa(\mathfrak q) = r$.
\end{enumerate}
Alors il existe $f \in R$, $f \not \in \mathfrak p$,
$g \in S$, $g \not\in \mathfrak q$, et un homomorphisme d'anneaux quasi-fini
$$
\varphi : R_f[x_1, \ldots, x_n] \longrightarrow S_g
$$
tel que $\varphi^{-1}(\mathfrak qS_g) =
(\mathfrak p, x_{r + 1}, \ldots, x_n)R_f[x_1, \ldots, x_n]$
\end{lemma}

\begin{proof}
Après remplacement de $S$ par une localisation principale, nous pouvons supposer qu'il
existe un homomorphisme d'anneaux quasi-fini $\varphi : R[t_1, \ldots, t_n] \to S$, voir le
lemme \ref{lemma-quasi-finite-over-polynomial-algebra}.
Posons $\mathfrak q' = \varphi^{-1}(\mathfrak q)$.
Soit $\overline{\mathfrak q}' \subset \kappa(\mathfrak p)[t_1, \ldots, t_n]$
l'idéal premier correspondant à $\mathfrak q'$. D'après le
lemme \ref{lemma-refined-Noether-normalization},
il existe un homomorphisme d'anneaux fini
$\kappa(\mathfrak p)[x_1, \ldots, x_n] \to
\kappa(\mathfrak p)[t_1, \ldots, t_n]$
tel que l'image réciproque de $\overline{\mathfrak q}'$ soit
$(x_{r + 1}, \ldots, x_n)$. Soit
$\overline{h}_i \in \kappa(\mathfrak p)[t_1, \ldots, t_n]$
l'image de $x_i$. Il existe un élément
$f \in R$, $f \not \in \mathfrak p$
et des $h_i \in R_f[t_1, \ldots, t_n]$ dont les images sont $\overline{h}_i$
dans $\kappa(\mathfrak p)[t_1, \ldots, t_n]$. Alors l'homomorphisme d'anneaux
$$
R_f[x_1, \ldots, x_n] \longrightarrow R_f[t_1, \ldots, t_n]
$$
devient fini après extension des scalaires à $\kappa(\mathfrak p)$.
En particulier, $R_f[t_1, \ldots, t_n]$ est quasi-fini sur
$R_f[x_1, \ldots, x_n]$ en $\mathfrak q'R_f[t_1, \ldots, t_n]$.
Ainsi, d'après le
lemme \ref{lemma-quasi-finite-open},
il existe $g \in R_f[t_1, \ldots, t_n]$,
$g \not \in \mathfrak q'R_f[t_1, \ldots, t_n]$
tel que $R_f[x_1, \ldots, x_n] \to R_f[t_1, \ldots, t_n, 1/g]$
soit quasi-fini. On voit donc que le composé
$$
R_f[x_1, \ldots, x_n] \longrightarrow
R_f[t_1, \ldots, t_n, 1/g] \longrightarrow S_{\varphi(g)}
$$
est quasi-fini, d'où le résultat.
\end{proof}

\begin{lemma}
\label{lemma-dimension-inequality-quasi-finite}
Soit $R \to S$ un homomorphisme d'anneaux de type fini.
Soit $\mathfrak q \subset S$ un idéal premier au-dessus de $\mathfrak p \subset R$.
Si $R \to S$ est quasi-fini en $\mathfrak q$, alors
$\dim(S_{\mathfrak q}) \leq \dim(R_{\mathfrak p})$.
\end{lemma}

\begin{proof}
Si $R_{\mathfrak p}$ est noethérien
(et donc $S_{\mathfrak q}$ noethérien puisqu'il est essentiellement
de type fini sur $R_{\mathfrak p}$),
cela résulte immédiatement du
lemme \ref{lemma-dimension-base-fibre-total} et des
définitions. Dans le cas général, soit $S'$ la clôture intégrale
de $R_\mathfrak p$ dans $S_\mathfrak p$.
Par le théorème principal de Zariski \ref{theorem-main-theorem},
on a $S_{\mathfrak q} = S'_{\mathfrak q'}$ pour un certain
$\mathfrak q' \subset S'$ au-dessus de $\mathfrak q$.
D'après le lemme \ref{lemma-integral-dim-up}, on a
$\dim(S') \leq \dim(R_\mathfrak p)$ et donc, a fortiori,
$\dim(S_\mathfrak q) = \dim(S'_{\mathfrak q'}) \leq \dim(R_\mathfrak p)$.
\end{proof}

\begin{lemma}
\label{lemma-dimension-quasi-finite-over-polynomial-algebra}
\begin{slogan}
Un revêtement quasi-fini de l'espace affine de dimension n est de dimension au plus n.
\end{slogan}
Soit $k$ un corps. Soit $S$ une $k$-algèbre de type fini.
Supposons qu'il existe un homomorphisme de $k$-algèbres quasi-fini
$k[t_1, \ldots, t_n] \subset S$. Alors $\dim(S) \leq n$.
\end{lemma}

\begin{proof}
D'après le lemme \ref{lemma-dim-affine-space}, la dimension de
tout anneau local de $k[t_1, \ldots, t_n]$ est au plus $n$.
Le résultat découle donc du
lemme \ref{lemma-dimension-inequality-quasi-finite}.
\end{proof}

\begin{lemma}
\label{lemma-dimension-fibres-bounded-open-upstairs}
Soit $R \to S$ un homomorphisme d'anneaux de type fini.
Soit $\mathfrak q \subset S$ un idéal premier.
Supposons que $\dim_{\mathfrak q}(S/R) = n$.
Il existe un voisinage ouvert $V$ de $\mathfrak q$
dans $\Spec(S)$ tel que
$\dim_{\mathfrak q'}(S/R) \leq n$ pour tout $\mathfrak q' \in V$.
\end{lemma}

\begin{proof}
D'après le lemme \ref{lemma-quasi-finite-over-polynomial-algebra},
on voit que l'on peut supposer $S$ quasi-fini sur
une algèbre de polynômes $R[t_1, \ldots, t_n]$. En considérant
les fibres, on se ramène au
lemme \ref{lemma-dimension-quasi-finite-over-polynomial-algebra}.
\end{proof}

\noindent
Autrement dit, le lemme affirme que l'ensemble des points où la
dimension de la fibre est $\leq n$ est ouvert dans $\Spec(S)$.
Le lemme suivant affirme que la formation de cet ouvert commute au
changement de base.
Si l'homomorphisme d'anneaux est de présentation finie, cet ensemble est
un ouvert quasi-compact (voir ci-dessous).

\begin{lemma}
\label{lemma-dimension-fibres-bounded-open-upstairs-base-change}
Soit $R \to S$ un homomorphisme d'anneaux de type fini.
Soit $R \to R'$ un homomorphisme d'anneaux quelconque.
Posons $S' = R' \otimes_R S$ et notons $f : \Spec(S') \to \Spec(S)$
l'application correspondante entre les spectres.
Soit $n \geq 0$.
L'image réciproque
$f^{-1}(\{\mathfrak q \in \Spec(S) \mid
\dim_{\mathfrak q}(S/R) \leq n\})$
est égale à
$\{\mathfrak q' \in \Spec(S') \mid
\dim_{\mathfrak q'}(S'/R') \leq n\}$.
\end{lemma}

\begin{proof}
La condition s'exprime en termes des dimensions
des anneaux des fibres, qui sont de type fini sur un corps.
En la combinant avec le
lemme \ref{lemma-dimension-at-a-point-preserved-field-extension},
on obtient le lemme.
\end{proof}

\begin{lemma}
\label{lemma-dimension-fibres-bounded-quasi-compact-open-upstairs}
Soit $R \to S$ un homomorphisme d'anneaux de présentation finie.
Soit $n \geq 0$. L'ensemble
$$
V_n = \{\mathfrak q \in \Spec(S) \mid \dim_{\mathfrak q}(S/R) \leq n\}
$$
est un ouvert quasi-compact de $\Spec(S)$.
\end{lemma}

\begin{proof}
Il est ouvert d'après le lemme \ref{lemma-dimension-fibres-bounded-open-upstairs}.
Soit $S = R[x_1, \ldots, x_n]/(f_1, \ldots, f_m)$ une présentation de
$S$. Soit $R_0$ la sous-algèbre sur $\mathbf{Z}$ de $R$ engendrée par les
coefficients des polynômes $f_i$.
Soit $S_0 = R_0[x_1, \ldots, x_n]/(f_1, \ldots, f_m)$.
Alors $S = R \otimes_{R_0} S_0$. D'après le
lemme \ref{lemma-dimension-fibres-bounded-open-upstairs-base-change},
$V_n$ est l'image réciproque d'un ouvert $V_{0, n}$ par l'application continue
quasi-compacte $\Spec(S) \to \Spec(S_0)$. Comme
$S_0$ est noethérien, on voit que $V_{0, n}$ est quasi-compact.
\end{proof}

\begin{lemma}
\label{lemma-finite-type-domain-over-valuation-ring-dim-fibres}
Soit $R$ un anneau de valuation de corps résiduel $k$ et de corps
des fractions $K$. Soit $S$ un anneau intègre contenant $R$ tel que
$S$ soit de type fini sur $R$. Si $S \otimes_R k$ n'est pas l'anneau
nul, alors
$$
\dim(S \otimes_R k) = \dim(S \otimes_R K)
$$
En fait, $\Spec(S \otimes_R k)$ est équidimensionnel.
\end{lemma}

\begin{proof}
Il suffit de montrer que $\dim_{\mathfrak q}(S/k)$ est égale
à $\dim(S \otimes_R K)$ pour tout idéal premier $\mathfrak q$ de
$S$ contenant $\mathfrak m_RS$. Choisissons un tel idéal premier. D'après le
lemme \ref{lemma-dimension-fibres-bounded-open-upstairs},
l'inégalité $\dim_{\mathfrak q}(S/k) \geq \dim(S \otimes_R K)$
est vérifiée. Posons $n = \dim_{\mathfrak q}(S/k)$. D'après le
lemme \ref{lemma-quasi-finite-over-polynomial-algebra},
après remplacement de $S$ par $S_g$ pour un certain $g \in S$, $g \not \in \mathfrak q$,
il existe un homomorphisme d'anneaux quasi-fini
$R[t_1, \ldots, t_n] \to S$. Si $\dim(S \otimes_R K) < n$,
alors $K[t_1, \ldots, t_n] \to S \otimes_R K$ a un noyau non nul.
Écrivons $f = \sum a_I t_1^{i_1}\ldots t_n^{i_n}$. Après avoir divisé
$f$ par un coefficient non nul de $f$ de valuation minimale, nous pouvons
supposer que $f\in R[t_1, \ldots, t_n]$ et qu'un certain $a_I$ a une image non nulle
dans $k$. Ainsi, l'homomorphisme d'anneaux $k[t_1, \ldots, t_n] \to S \otimes_R k$
a un noyau non nul, ce qui entraîne que $\dim(S \otimes_R k) < n$.
Contradiction.
\end{proof}
























\section{Algèbres et modules de présentation finie}
\label{section-finite-presentation}

\noindent
Dans cette section, nous examinons quelques résultats classiques dont
l'élément essentiel est une hypothèse de type fini ou de
présentation finie.

\begin{lemma}
\label{lemma-finite-type-descends}
Soit $R \to S$ un homomorphisme d'anneaux.
Soit $R \to R'$ un homomorphisme d'anneaux fidèlement plat.
Posons $S' = R'\otimes_R S$.
Alors $R \to S$ est de type fini si et seulement si $R' \to S'$
est de type fini.
\end{lemma}

\begin{proof}
Il est clair que, si $R \to S$ est de type fini, alors $R' \to S'$
est de type fini. Supposons que $R' \to S'$ soit de type fini.
Supposons que $y_1, \ldots, y_m$ engendrent $S'$ sur $R'$.
Écrivons $y_j = \sum_i a_{ij} \otimes x_{ji}$ pour certains
$a_{ij} \in R'$ et $x_{ji} \in S$. Soit $A \subset S$
la sous-$R$-algèbre engendrée par les $x_{ij}$.
Par platitude, on a $A' := R' \otimes_R A \subset S'$ et,
par construction, $y_j \in A'$. Ainsi, $A' = S'$.
Par fidèle platitude, $A = S$.
\end{proof}

\begin{lemma}
\label{lemma-finite-presentation-descends}
Soit $R \to S$ un homomorphisme d'anneaux.
Soit $R \to R'$ un homomorphisme d'anneaux fidèlement plat.
Posons $S' = R'\otimes_R S$.
Alors $R \to S$ est de présentation finie si et seulement si $R' \to S'$
est de présentation finie.
\end{lemma}

\begin{proof}
Il est clair que, si $R \to S$ est de présentation finie, alors $R' \to S'$
est de présentation finie. Supposons que $R' \to S'$ soit de présentation finie.
D'après le lemme \ref{lemma-finite-type-descends}, on voit
que $R \to S$ est de type fini. Écrivons $S = R[x_1, \ldots, x_n]/I$.
Par platitude, $S' = R'[x_1, \ldots, x_n]/R'\otimes I$.
Supposons que $g_1, \ldots, g_m$ engendrent $R'\otimes I$ sur $R'[x_1, \ldots, x_n]$.
Écrivons $g_j = \sum_i a_{ij} \otimes f_{ji}$ pour certains
$a_{ij} \in R'$ et $f_{ji} \in I$. Soit $J \subset I$
l'idéal engendré par les $f_{ij}$.
Par platitude, on a $R' \otimes_R J \subset R'\otimes_R I$, et
ce sont tous deux des idéaux de $R'[x_1, \ldots, x_n]$.
Par construction, $g_j \in R' \otimes_R J$. Ainsi,
$R' \otimes_R J = R'\otimes_R I$.
Par fidèle platitude, $J = I$.
\end{proof}

\begin{lemma}
\label{lemma-construct-fp-module}
Soit $R$ un anneau.
Soit $I \subset R$ un idéal.
Soit $S \subset R$ une partie multiplicative.
Posons $R' = S^{-1}(R/I) = S^{-1}R/S^{-1}I$.
\begin{enumerate}
\item Pour tout $R'$-module de type fini $M'$, il existe un
$R$-module de type fini $M$ tel que $S^{-1}(M/IM) \cong M'$.
\item Pour tout $R'$-module de présentation finie $M'$, il existe un
$R$-module de présentation finie $M$ tel que $S^{-1}(M/IM) \cong M'$.
\end{enumerate}
\end{lemma}

\begin{proof}
Démonstration de (1). Choisissons une suite exacte courte
$0 \to K' \to (R')^{\oplus n} \to M' \to 0$.
Soit $K \subset R^{\oplus n}$ l'image réciproque de
$K'$ par l'homomorphisme $R^{\oplus n} \to (R')^{\oplus n}$.
Alors $M = R^{\oplus n}/K$ convient.

\medskip\noindent
Démonstration de (2).
Choisissons une présentation $(R')^{\oplus m} \to (R')^{\oplus n} \to M' \to 0$.
Supposons que le premier homomorphisme soit donné par la matrice
$A' = (a'_{ij})$ et que le second soit déterminé par des générateurs
$x'_i \in M'$, $i = 1, \ldots, n$. Comme $R' = S^{-1}(R/I)$, on peut choisir
$s \in S$ et une matrice $A = (a_{ij})$ à coefficients dans $R$
tels que $a'_{ij} = a_{ij} / s \bmod S^{-1}I$. Soit $M$ le
$R$-module de présentation finie ayant pour présentation
$R^{\oplus m} \to R^{\oplus n} \to M \to 0$,
où le premier homomorphisme est donné par la matrice $A$ et le second
est déterminé par des générateurs $x_i \in M$, $i = 1, \ldots, n$.
Alors l'homomorphisme $M \to M'$, $x_i \mapsto x'_i$, induit un isomorphisme
$S^{-1}(M/IM) \cong M'$.
\end{proof}

\begin{lemma}
\label{lemma-construct-fp-module-from-localization}
Soit $R$ un anneau.
Soit $S \subset R$ une partie multiplicative.
Soit $M$ un $R$-module.
\begin{enumerate}
\item Si $S^{-1}M$ est un $S^{-1}R$-module de type fini, alors il
existe un $R$-module de type fini $M'$ et un homomorphisme $M' \to M$ qui induit un
isomorphisme $S^{-1}M' \to S^{-1}M$.
\item Si $S^{-1}M$ est un $S^{-1}R$-module de présentation finie,
alors il existe un $R$-module $M'$ de présentation finie
et un homomorphisme $M' \to M$ qui induit un isomorphisme
$S^{-1}M' \to S^{-1}M$.
\end{enumerate}
\end{lemma}

\begin{proof}
Démonstration de (1). Soient $x_1, \ldots, x_n \in M$ des éléments qui engendrent
$S^{-1}M$ comme $S^{-1}R$-module. Soit $M'$ le
sous-$R$-module de $M$ engendré par $x_1, \ldots, x_n$.

\medskip\noindent
Démonstration de (2). Soient $x_1, \ldots, x_n \in M$ des éléments qui engendrent
$S^{-1}M$ comme $S^{-1}R$-module. Posons
$K = \Ker(R^{\oplus n} \to M)$, où l'homomorphisme est donné par
la règle $(a_1, \ldots, a_n) \mapsto \sum a_i x_i$. D'après le
lemme \ref{lemma-extension},
on voit que $S^{-1}K$ est un $S^{-1}R$-module de type fini.
D'après (1), on peut trouver un sous-module de type fini $K' \subset K$
tel que $S^{-1}K' = S^{-1}K$. Prenons
$M' = \Coker(K' \to R^{\oplus n})$.
\end{proof}

\begin{lemma}
\label{lemma-construct-fp-module-from-stalk}
Soit $R$ un anneau.
Soit $\mathfrak p \subset R$ un idéal premier.
Soit $M$ un $R$-module.
\begin{enumerate}
\item Si $M_{\mathfrak p}$ est un $R_{\mathfrak p}$-module de type fini, alors il
existe un $R$-module de type fini $M'$ et un homomorphisme $M' \to M$ qui induit un
isomorphisme $M'_{\mathfrak p} \to M_{\mathfrak p}$.
\item Si $M_{\mathfrak p}$ est un $R_{\mathfrak p}$-module de présentation finie,
alors il existe un $R$-module $M'$ de présentation finie
et un homomorphisme $M' \to M$ qui induit un isomorphisme
$M'_{\mathfrak p} \to M_{\mathfrak p}$.
\end{enumerate}
\end{lemma}

\begin{proof}
C'est un cas particulier du
lemme \ref{lemma-construct-fp-module-from-localization}.
\end{proof}

\begin{lemma}
\label{lemma-local-isomorphism}
Soit $\varphi : R \to S$ un homomorphisme d'anneaux. Soit $\mathfrak q \subset S$
un idéal premier au-dessus de $\mathfrak p \subset R$. Supposons que
\begin{enumerate}
\item $S$ soit de présentation finie sur $R$,
\item $\varphi$ induise un isomorphisme $R_\mathfrak p \cong S_\mathfrak q$.
\end{enumerate}
Alors il existe $f \in R$, $f \not \in \mathfrak p$ et une
$R_f$-algèbre $C$ tels que $S_f \cong R_f \times C$ comme $R_f$-algèbres.
\end{lemma}

\begin{proof}
Écrivons $S = R[x_1, \ldots, x_n]/(g_1, \ldots, g_m)$. Soit $a_i \in R_\mathfrak p$
un élément qui s'envoie sur l'image de $x_i$ dans $S_\mathfrak q$.
Écrivons $a_i = b_i/f$ pour un certain $f \in R$, $f \not \in \mathfrak p$.
Après avoir remplacé $R$ par $R_f$ et $x_i$ par $x_i - a_i$, on peut
supposer que $S = R[x_1, \ldots, x_n]/(g_1, \ldots, g_m)$ et
que $x_i$ s'envoie sur zéro dans $S_\mathfrak q$. Alors, si $c_j$ désigne
le terme constant de $g_j$, on conclut que $c_j$ s'envoie sur zéro
dans $R_\mathfrak p$. Après un nouveau remplacement de $R$, on peut
supposer que les termes constants $c_j$ des $g_j$ sont nuls.
On obtient ainsi un homomorphisme de $R$-algèbres $S \to R$, $x_i \mapsto 0$, dont le
noyau est l'idéal $(x_1, \ldots, x_n)$.

\medskip\noindent
On a les isomorphismes $R_\mathfrak p \to S_\mathfrak q \to R_\mathfrak p$,
et $S \to R$ envoie $x_i$ sur zéro. On doit donc avoir
$S_\mathfrak q = R_\mathfrak p[x_1, \ldots, x_n]/(x_1, \ldots, x_n)$
et, a fortiori, $S_\mathfrak q = S_\mathfrak p/(x_1, \ldots, x_n)S_\mathfrak p$.
Cela signifie que l'idéal de type fini $(x_1, \ldots, x_n)S_\mathfrak p$
est pur dans $S_\mathfrak p$ ; voir la définition \ref{definition-pure-ideal}.
Ainsi, $(x_1, \ldots, x_n)S_\mathfrak p$ est engendré par un idempotent
$e$ de $S_\mathfrak p$ d'après le lemme \ref{lemma-finitely-generated-pure-ideal}.
Après avoir remplacé $R \to S$ par $R_f \to S_f$ pour un certain
$f \in R$, $f \not \in \mathfrak p$, on peut trouver un idempotent
$e' \in S$ qui s'envoie sur $e$.
Alors $e'S$ et $(x_1, \ldots, x_n)S$ sont des idéaux de type fini
qui deviennent égaux dans $S_\mathfrak p$. Ainsi, après avoir remplacé $R \to S$ par
$R_f \to S_f$ pour un certain $f \in R$, $f \not \in \mathfrak p$, on peut supposer que
$e'S = (x_1, \ldots, x_n)S$. Poser $C = e'S$ achève la démonstration.
\end{proof}

\begin{lemma}
\label{lemma-isomorphic-local-rings}
Soit $R$ un anneau.
Soient $S$, $S'$ de présentation finie sur $R$.
Soient $\mathfrak q \subset S$ et $\mathfrak q' \subset S'$
des idéaux premiers. Si $S_{\mathfrak q} \cong S'_{\mathfrak q'}$ comme
$R$-algèbres, alors il existe $g \in S$, $g \not \in \mathfrak q$
et $g' \in S'$, $g' \not \in \mathfrak q'$, tels que
$S_g \cong S'_{g'}$ comme $R$-algèbres.
\end{lemma}

\begin{proof}
Soit $\psi : S_{\mathfrak q} \to S'_{\mathfrak q'}$ l'isomorphisme
de l'hypothèse du lemme.
Écrivons $S = R[x_1, \ldots, x_n]/(f_1, \ldots, f_r)$ et
$S' = R[y_1, \ldots, y_m]/J$.
Pour chaque $i = 1, \ldots, n$, choisissons une fraction
$h_i/g_i$ avec $h_i, g_i \in R[y_1, \ldots, y_m]$
avec $g_i \bmod J$ n'appartenant pas à $\mathfrak q'$ ; cette fraction représente
l'image de $x_i$ par $\psi$. Après avoir remplacé
$S'$ par $S'_{g_1 \ldots g_n}$ et introduit
$R[y_1, \ldots, y_m, y_{m + 1}]$ (en envoyant $y_{m + 1}$ sur $1/(g_1\ldots g_n)$),
on peut supposer que $\psi(x_i)$ est l'image d'un certain
$h_i \in R[y_1, \ldots, y_m]$. Considérons les éléments
$f_j(h_1, \ldots, h_n) \in R[y_1, \ldots, y_m]$.
Puisque $\psi$ annule chaque $f_j$,
il existe $g \in R[y_1, \ldots, y_m]$, $g \bmod J \not \in \mathfrak q'$
tel que $g f_j(h_1, \ldots, h_n) \in J$ pour chaque $j = 1, \ldots, r$.
Après avoir remplacé $S'$ par $S'_g$ et
introduit $R[y_1, \ldots, y_m, y_{m + 1}]$ comme précédemment, on peut supposer que
$f_j(h_1, \ldots, h_n) \in J$. On obtient ainsi un homomorphisme d'anneaux
$S \to S'$, $x_i \mapsto h_i$, qui induit $\psi$ sur les anneaux locaux.
D'après le lemme \ref{lemma-compose-finite-type},
l'homomorphisme $S \to S'$ est de présentation finie.
D'après le lemme \ref{lemma-local-isomorphism},
on peut supposer que $S' = S \times C$. Ainsi, en inversant dans $S'$
l'idempotent correspondant au facteur $C$, on obtient le résultat.
\end{proof}

\begin{lemma}
\label{lemma-finite-type-mod-nilpotent}
Soit $R$ un anneau. Soit $I \subset R$ un idéal nilpotent.
Soit $S$ une $R$-algèbre telle que $R/I \to S/IS$
soit de type fini. Alors $R \to S$ est de type fini.
\end{lemma}

\begin{proof}
Choisissons $s_1, \ldots, s_n \in S$ dont les images dans $S/IS$ engendrent
$S/IS$ comme algèbre sur $R/I$. D'après la partie (11) du lemme \ref{lemma-NAK},
on voit que l'homomorphisme de $R$-algèbres $R[x_1, \ldots, x_n] \to S$, $x_i \mapsto s_i$,
est surjectif, ce qui donne le résultat.
\end{proof}

\begin{lemma}
\label{lemma-surjective-mod-locally-nilpotent}
Soit $R$ un anneau. Soit $I \subset R$ un idéal localement nilpotent.
Soit $S \to S'$ un homomorphisme de $R$-algèbres tel que $S \to S'/IS'$ soit surjectif
et que $S'$ soit de type fini sur $R$. Alors $S \to S'$ est surjectif.
\end{lemma}

\begin{proof}
Écrivons $S' = R[x_1, \ldots, x_m]/K$ pour un certain idéal $K$. Par hypothèse, il
existe des $g_j = x_j + \sum \delta_{j, J} x^J \in R[x_1, \ldots, x_n]$ avec
$\delta_{j, J} \in I$ et $g_j \bmod K \in \Im(S \to S')$.
Il suffit donc de montrer que $g_1, \ldots, g_m$ engendrent
$R[x_1, \ldots, x_n]$. Soit $R_0 \subset R$ une
$\mathbf{Z}$-sous-algèbre de type fini de $R$ contenant notamment les $\delta_{j, J}$.
Alors $R_0 \cap I$ est un idéal nilpotent (d'après le
lemme \ref{lemma-Noetherian-power}). Il s'ensuit que
$R_0[x_1, \ldots, x_n]$ est engendré par $g_1, \ldots, g_m$ (car
$x_j \mapsto g_j$ définit un automorphisme de $R_0[x_1, \ldots, x_m]$ ;
nous omettons les détails). Comme $R$ est la réunion des sous-anneaux $R_0$, le résultat suit.
\end{proof}

\begin{lemma}
\label{lemma-isomorphism-modulo-ideal}
Soit $R$ un anneau. Soit $I \subset R$ un idéal. Soit $S \to S'$
un homomorphisme de $R$-algèbres. Soit $IS \subset \mathfrak q \subset S$
un idéal
premier. Supposons que
\begin{enumerate}
\item $S \to S'$ soit surjectif,
\item $S_\mathfrak q/IS_\mathfrak q \to S'_\mathfrak q/IS'_\mathfrak q$
soit un isomorphisme,
\item $S$ soit de type fini sur $R$,
\item $S'$ soit de présentation finie sur $R$, et
\item $S'_\mathfrak q$ soit plat sur $R$.
\end{enumerate}
Alors $S_g \to S'_g$ est un isomorphisme pour un certain
$g \in S$, $g \not \in \mathfrak q$.
\end{lemma}

\begin{proof}
Soit $J = \Ker(S \to S')$. D'après le
lemme \ref{lemma-compose-finite-type},
$J$ est un idéal de type fini. Comme $S'_\mathfrak q$ est plat
sur $R$, on voit que
$J_\mathfrak q/IJ_\mathfrak q \subset S_\mathfrak q/IS_{\mathfrak q}$
(appliquer le lemme \ref{lemma-flat-tor-zero} à $0 \to J \to S \to S' \to 0$).
D'après l'hypothèse (2), on voit que $J_\mathfrak q/IJ_\mathfrak q$ est nul.
D'après le lemme de Nakayama (lemme \ref{lemma-NAK}),
il existe $g \in S$, $g \not \in \mathfrak q$, tel
que $J_g = 0$. Ainsi, $S_g \cong S'_g$, comme souhaité.
\end{proof}

\begin{lemma}
\label{lemma-isomorphism-modulo-locally-nilpotent}
Soit $R$ un anneau. Soit $I \subset R$ un idéal. Soit $S \to S'$
un homomorphisme de $R$-algèbres. Supposons que
\begin{enumerate}
\item $I$ soit localement nilpotent,
\item $S/IS \to S'/IS'$ soit un isomorphisme,
\item $S$ soit de type fini sur $R$,
\item $S'$ soit de présentation finie sur $R$, et
\item $S'$ soit plat sur $R$.
\end{enumerate}
Alors $S \to S'$ est un isomorphisme.
\end{lemma}

\begin{proof}
D'après le lemme \ref{lemma-surjective-mod-locally-nilpotent}, l'homomorphisme
$S \to S'$ est surjectif. Comme $I$ est localement nilpotent, les
idéaux $IS$ et $IS'$ le sont aussi (lemme \ref{lemma-locally-nilpotent}). Ainsi,
tout idéal premier $\mathfrak q$ de $S$ contient $IS$ et, trivialement,
$S_\mathfrak q/IS_\mathfrak q \cong S'_\mathfrak q/IS'_\mathfrak q$.
Le lemme \ref{lemma-isomorphism-modulo-ideal} s'applique donc,
et l'on voit que $S_\mathfrak q \to S'_\mathfrak q$ est un
isomorphisme pour tout idéal premier $\mathfrak q \subset S$.
Il s'ensuit que $S \to S'$ est injectif, par exemple d'après le
lemme \ref{lemma-characterize-zero-local}.
\end{proof}





\section{Colimites et morphismes de présentation finie}
\label{section-colimits-flat}

\noindent
Dans cette section, nous démontrons quelques lemmes préliminaires
qui nous aideront par la suite à établir un résultat au moyen d'une
réduction noethérienne absolue.
Dans Catégories, section \ref{categories-section-directed-colimits},
nous étudions les colimites filtrantes en général.
Voici un exemple de cette notion très générale.

\begin{lemma}
\label{lemma-ring-colimit-fp-category}
Soit $R \to A$ un morphisme d'anneaux. Considérons la catégorie
$\mathcal{I}$ de tous les diagrammes de morphismes de $R$-algèbres
$A' \to A$, où $A'$ est de présentation finie sur $R$. Alors
$\mathcal{I}$ est filtrante et la colimite des $A'$ suivant
$\mathcal{I}$ est isomorphe à $A$.
\end{lemma}
\begin{proof}
La catégorie\footnote{Afin d'éviter les difficultés ensemblistes, nous ne
considérons que les $A' \to A$ tels que $A'$ soit un quotient de
$R[x_1, x_2, x_3, \ldots]$.}
$\mathcal{I}$ est non vide, puisque $R \to R$ en est un objet.
Considérons deux objets $A' \to A$ et $A'' \to A$ de $\mathcal{I}$.
Alors $A' \otimes_R A'' \to A$ appartient à
$\mathcal{I}$ (utiliser les Lemmes \ref{lemma-compose-finite-type} et
\ref{lemma-base-change-finiteness}). Les morphismes d'anneaux
$A' \to A' \otimes_R A''$ et $A'' \to A' \otimes_R A''$
définissent des flèches de $\mathcal{I}$, ce qui démontre la deuxième
propriété définissant une catégorie filtrante ; voir
Catégories, Définition \ref{categories-definition-directed}.
Enfin, supposons donnés deux morphismes $\sigma, \tau : A' \to A''$
dans $\mathcal{I}$. Si $x_1, \ldots, x_r \in A'$ engendrent
$A'$ comme $R$-algèbre, nous pouvons considérer
$A''' = A''/(\sigma(x_i) - \tau(x_i))$.
C'est une $R$-algèbre de présentation finie, et le morphisme de
$R$-algèbres donné $A'' \to A$ se factorise par la surjection
$\nu : A'' \to A'''$. Ainsi, $\nu$ est un morphisme de $\mathcal{I}$
qui égalise $\sigma$ et $\tau$, comme souhaité.

\medskip\noindent
Le fait que notre catégorie d'indices soit cofiltrante signifie que nous
pouvons calculer la valeur de $B = \colim_{A' \to A} A'$ dans la catégorie
des ensembles (certains détails sont omis ; comparer avec la discussion de
Catégories, section \ref{categories-section-directed-colimits}).
Pour voir que $B \to A$ est surjectif, pour tout $a \in A$ nous pouvons
utiliser $R[x] \to A$, $x \mapsto a$, et constater que $a$ appartient à
l'image de $B \to A$. Réciproquement, si $b \in B$ s'envoie sur zéro dans
$A$, nous pouvons trouver $A' \to A$ dans $\mathcal{I}$ et $a' \in A'$
qui s'envoie sur $b$. Alors $A'/(a') \to A$ appartient lui aussi à
$\mathcal{I}$, et le morphisme $A' \to B$ se factorise sous la forme
$A' \to A'/(a') \to B$, ce qui montre que $b = 0$, comme souhaité.
\end{proof}
\noindent
Il est souvent plus commode de considérer les colimites indexées par des
ensembles préordonnés. Soit $(\Lambda, \geq)$ un ensemble préordonné.
Un système d'anneaux sur $\Lambda$ est constitué d'un anneau
$R_\lambda$ pour chaque $\lambda \in \Lambda$ et d'un morphisme
$R_\lambda \to R_\mu$ dès que $\lambda \leq \mu$.
Ces morphismes doivent satisfaire la règle selon laquelle
$R_\lambda \to R_\mu \to R_\nu$ est égal au morphisme
$R_\lambda \to R_\nu$ pour tous $\lambda \leq \mu \leq \nu$.
Voir Catégories, section \ref{categories-section-posets-limits}.
Nous supposerons souvent que $(I, \leq)$ est {\it filtrant},
ce qui signifie que $\Lambda$ est non vide et que, étant donnés
$\lambda, \mu \in \Lambda$, il existe $\nu \in \Lambda$ tel que
$\lambda \leq \nu$ et $\mu \leq \nu$.
Rappelons que, dans ce cas, la colimite $\colim_\lambda R_\lambda$
est parfois appelée « limite directe »
(mais nous n'emploierons pas cette terminologie).

\medskip\noindent
Notons que Catégories, Lemme \ref{categories-lemma-directed-category-system},
nous dit que les colimites indexées par des catégories filtrantes reviennent
au même que les colimites indexées par des ensembles filtrants.

\begin{lemma}
\label{lemma-ring-colimit-fp}
Soit $R \to A$ un morphisme d'anneaux. Il existe un système filtrant
$A_\lambda$ de $R$-algèbres de présentation finie tel que
$A = \colim_\lambda A_\lambda$. Si $A$ est de type fini sur $R$, on peut
faire en sorte que tous les morphismes de transition du système des
$A_\lambda$ soient surjectifs.
\end{lemma}

\begin{proof}
Une première démonstration consiste à déduire l'énoncé du
Lemme \ref{lemma-ring-colimit-fp-category} et de
Catégories, Lemme \ref{categories-lemma-directed-category-system}.

\medskip\noindent
Seconde démonstration.
Comparer avec la démonstration du Lemme \ref{lemma-module-colimit-fp}.
Considérons une partie finie quelconque $S \subset A$ et une famille finie
$E$ de relations polynomiales entre les éléments de $S$.
Ainsi, à chaque $s \in S$ correspond $x_s \in A$, et chaque $e \in E$
consiste en un polynôme $f_e \in R[X_s; s\in S]$ tel que
$f_e(x_s) = 0$. Posons
$A_{S, E} = R[X_s; s\in S]/(f_e; e\in E)$ ; c'est une $R$-algèbre de
présentation finie. Il existe des morphismes canoniques
$A_{S, E} \to A$.
Si $S \subset S'$ et si les éléments de $E$ correspondent, par le morphisme
$R[X_s; s \in S] \to R[X_s; s\in S']$, à une partie de $E'$, il existe
alors un morphisme évident $A_{S, E} \to A_{S', E'}$ qui commute aux
morphismes vers $A$. Ainsi, en prenant pour $\Lambda$ l'ensemble des couples
$(S, E)$ ordonné par inclusion comme ci-dessus, on obtient un ensemble
partiellement ordonné filtrant. Il est clair que la colimite de ce système
filtrant est $A$.

\medskip\noindent
Pour la dernière assertion, supposons $A = R[x_1, \ldots, x_n]/I$.
Considérons alors la partie $\Lambda' \subset \Lambda$ formée des systèmes
$(S, E)$ ci-dessus pour lesquels $S = \{x_1, \ldots, x_n\}$.
On voit aisément que l'on a encore
$A = \colim_{\lambda' \in \Lambda'} A_{\lambda'}$.
De plus, les morphismes de transition sont manifestement surjectifs.
\end{proof}
\noindent
Il se trouve que l'on peut caractériser comme suit les morphismes d'anneaux
de présentation finie. Cela signifie en un certain sens que les algèbres
de présentation finie sont les objets « compacts » de la catégorie des
$R$-algèbres.

\begin{lemma}
\label{lemma-characterize-finite-presentation}
Soit $\varphi : R \to S$ un morphisme d'anneaux. Les assertions suivantes
sont équivalentes :
\begin{enumerate}
\item $\varphi$ est de présentation finie,
\item pour tout système filtrant $A_\lambda$ de $R$-algèbres, l'application
$$
\colim_\lambda \Hom_R(S, A_\lambda) \longrightarrow
\Hom_R(S, \colim_\lambda A_\lambda)
$$
est bijective, et
\item pour tout système filtrant $A_\lambda$ de $R$-algèbres, l'application
$$
\colim_\lambda \Hom_R(S, A_\lambda) \longrightarrow
\Hom_R(S, \colim_\lambda A_\lambda)
$$
est surjective.
\end{enumerate}
\end{lemma}
\begin{proof}
Supposons (1) et écrivons
$S = R[x_1, \ldots, x_n] / (f_1, \ldots, f_m)$.
Posons $A = \colim A_\lambda$. Remarquons qu'un homomorphisme de
$R$-algèbres $S \to A$ ou $S \to A_\lambda$ est déterminé par les images
de $x_1, \ldots, x_n$. Il est donc clair que
$\colim_\lambda \Hom_R(S, A_\lambda) \to \Hom_R(S, A)$
est injectif. Pour voir qu'il est surjectif, soit $\chi : S \to A$
un homomorphisme de $R$-algèbres. Chaque $x_i$ s'envoie alors sur un
élément de l'image d'un certain $A_{\lambda_i}$.
Nous pouvons choisir $\mu \geq \lambda_i$, $i = 1, \ldots, n$, et supposer
que $\chi(x_i)$ est l'image de $y_i \in A_\mu$ pour
$i = 1, \ldots, n$. Considérons
$z_j = f_j(y_1, \ldots, y_n) \in A_\mu$.
Comme $\chi$ est un homomorphisme, l'image de $z_j$ dans
$A = \colim_\lambda A_\lambda$ est nulle. Il existe donc
$\mu_j \geq \mu$ tel que $z_j$ s'envoie sur zéro dans $A_{\mu_j}$.
Choisissons $\nu \geq \mu_j$, $j = 1, \ldots, m$. Les images de
$z_1, \ldots, z_m$ sont alors nulles dans $A_\nu$. Cela signifie
exactement que les $y_i$ s'envoient sur des éléments
$y'_i \in A_\nu$ qui satisfont les relations
$f_j(y'_1, \ldots, y'_n) = 0$. Nous obtenons donc un morphisme d'anneaux
$S \to A_\nu$. Cela montre que (1) implique (2).

\medskip\noindent
Il est clair que (2) implique (3). Supposons (3).
D'après le Lemme \ref{lemma-ring-colimit-fp}, nous pouvons écrire
$S = \colim_\lambda S_\lambda$, où $S_\lambda$ est de présentation finie
sur $R$. Le morphisme identité se factorise alors sous la forme
$$
S \to S_\lambda \to S
$$
pour un certain $\lambda$. Il s'ensuit que $S$ est de présentation finie
sur $S_\lambda$ d'après le Lemme \ref{lemma-compose-finite-type}, partie (4),
appliqué à $S \to S_\lambda \to S$. En appliquant la partie (2) du même
lemme à $R \to S_\lambda \to S$, nous concluons que $S$ est de
présentation finie sur $R$.
\end{proof}
\noindent
À l'aide des résultats élémentaires qui précèdent, nous pouvons donner le
critère suivant pour qu'une algèbre $A$ soit une colimite filtrante
d'algèbres d'un type donné.

\begin{lemma}
\label{lemma-when-colimit}
Soit $R \to \Lambda$ un morphisme d'anneaux. Soit $\mathcal{E}$ un ensemble
de $R$-algèbres tel que chaque $A \in \mathcal{E}$ soit de présentation
finie sur $R$. Les deux assertions suivantes sont équivalentes :
\begin{enumerate}
\item $\Lambda$ est une colimite filtrante d'éléments de $\mathcal{E}$, et
\item pour tout morphisme de $R$-algèbres $A \to \Lambda$, avec $A$ de
présentation finie sur $R$, on peut trouver une factorisation
$A \to B \to \Lambda$ avec $B \in \mathcal{E}$.
\end{enumerate}
\end{lemma}

\begin{proof}
Supposons que $\mathcal{I} \to \mathcal{E}$, $i \mapsto A_i$, soit un
diagramme filtrant tel que $\Lambda = \colim_i A_i$.
Soit $A \to \Lambda$ un morphisme de $R$-algèbres, avec $A$ de
présentation finie sur $R$. L'application du
Lemme \ref{lemma-characterize-finite-presentation} fournit une
factorisation $A \to A_i \to \Lambda$. Ainsi, (1) implique (2).

\medskip\noindent
Considérons la catégorie $\mathcal{I}$ du
Lemme \ref{lemma-ring-colimit-fp-category}.
D'après Catégories, Lemme \ref{categories-lemma-cofinal-in-filtered},
la sous-catégorie pleine $\mathcal{J}$ formée des $A \to \Lambda$ tels que
$A \in \mathcal{E}$ est cofinale dans $\mathcal{I}$ et est une catégorie
filtrante. Alors $\Lambda$ est aussi la colimite suivant $\mathcal{J}$,
d'après Catégories, Lemme \ref{categories-lemma-cofinal}.
\end{proof}
\noindent
Mais il y a davantage. Plus précisément, étant donné
$R = \colim_\lambda R_\lambda$, on voit que la catégorie des $R$-modules
de présentation finie est équivalente à la limite des catégories des
$R_\lambda$-modules de présentation finie. Il en va de même des catégories
des $R$-algèbres de présentation finie.

\begin{lemma}
\label{lemma-module-map-property-in-colimit}
Soient $A$ un anneau et $M, N$ des $A$-modules.
Supposons que $R = \colim_{i \in I} R_i$ soit une colimite filtrante de
$A$-algèbres.
\begin{enumerate}
\item Si $M$ est un $A$-module de type fini et si $u, u' : M \to N$ sont des
morphismes de $A$-modules tels que
$u \otimes 1 = u' \otimes 1 : M \otimes_A R \to N \otimes_A R$,
alors, pour un certain $i$, on a
$u \otimes 1 = u' \otimes 1 : M \otimes_A R_i \to N \otimes_A R_i$.
\item Si $N$ est un $A$-module de type fini et si $u : M \to N$ est un morphisme
de $A$-modules tel que
$u \otimes 1 : M \otimes_A R \to N \otimes_A R$ soit surjectif, alors,
pour un certain $i$, le morphisme
$u \otimes 1 : M \otimes_A R_i \to N \otimes_A R_i$ est surjectif.
\item Si $N$ est un $A$-module de présentation finie et si
$v : N \otimes_A R \to M \otimes_A R$ est un morphisme de $R$-modules,
alors il existe $i$ et un morphisme de $R_i$-modules
$v_i : N \otimes_A R_i \to M \otimes_A R_i$ tels que
$v = v_i \otimes 1$.
\item Si $M$ est un $A$-module de type fini, si $N$ est un $A$-module de
présentation finie et si $u : M \to N$ est un morphisme de $A$-modules
tel que $u \otimes 1 : M \otimes_A R \to N \otimes_A R$ soit un
isomorphisme, alors, pour un certain $i$, le morphisme
$u \otimes 1 : M \otimes_A R_i \to N \otimes_A R_i$ est un isomorphisme.
\end{enumerate}
\end{lemma}

\begin{proof}
Pour démontrer (1), supposons que $u$ soit comme dans (1) et soient
$x_1, \ldots, x_m \in M$ des générateurs. Comme
$N \otimes_A R = \colim_i N \otimes_A R_i$, nous pouvons choisir
$i \in I$ tel que $u(x_j) \otimes 1 = u'(x_j) \otimes 1$ dans
$M \otimes_A R_i$, $j = 1, \ldots, m$. Pour un tel $i$, on a
$u \otimes 1 = u' \otimes 1 : M \otimes_A R_i \to N \otimes_A R_i$.

\medskip\noindent
Pour démontrer (2), supposons $u \otimes 1$ surjectif et soient
$y_1, \ldots, y_m \in N$ des générateurs. Comme
$N \otimes_A R = \colim_i N \otimes_A R_i$, nous pouvons choisir
$i \in I$ et $z_j \in M \otimes_A R_i$, $j = 1, \ldots, m$, dont les
images dans $N \otimes_A R$ sont égales à $y_j \otimes 1$.
Pour un tel $i$, le morphisme
$u \otimes 1 : M \otimes_A R_i \to N \otimes_A R_i$ est surjectif.

\medskip\noindent
Pour démontrer (3), soient $y_1, \ldots, y_m \in N$ des générateurs.
Soit $K = \Ker(A^{\oplus m} \to N)$, où le morphisme est donné par la règle
$(a_1, \ldots, a_m) \mapsto \sum a_j x_j$. Soient
$k_1, \ldots, k_t$ des générateurs de $K$. Écrivons
$k_s = (k_{s1}, \ldots, k_{sm})$.
Comme $M \otimes_A R = \colim_i M \otimes_A R_i$, nous pouvons choisir
$i \in I$ et $z_j \in M \otimes_A R_i$, $j = 1, \ldots, m$, dont les
images dans $M \otimes_A R$ sont égales à $v(y_j \otimes 1)$.
Nous voulons utiliser les $z_j$ pour définir le morphisme
$v_i : N \otimes_A R_i \to M \otimes_A R_i$.
Puisque
$K \otimes_A R_i \to R_i^{\oplus m} \to N \otimes_A R_i \to 0$
est une présentation, il suffit de vérifier que
$\xi_s = \sum_j k_{sj}z_j$ est nul dans $M \otimes_A R_i$ pour chaque
$s = 1, \ldots, t$. Ce n'est pas nécessairement le cas, mais, puisque
l'image de $\xi_s$ dans $M \otimes_A R$ est nulle, cela le devient après
avoir augmenté quelque peu $i$.

\medskip\noindent
Pour démontrer (4), supposons que $u \otimes 1$ soit un isomorphisme, que
$M$ soit de type fini et que $N$ soit de présentation finie.
Soit $v : N \otimes_A R \to M \otimes_A R$ un inverse de $u \otimes 1$.
La partie (3) fournit un morphisme
$v_i : N \otimes_A R_i \to M \otimes_A R_i$ pour un certain $i$.
La partie (1) montre qu'après avoir augmenté $i$, on a
$v_i \circ (u \otimes 1) = \text{id}_{M \otimes_R R_i}$ et
$(u \otimes 1) \circ v_i = \text{id}_{N \otimes_R R_i}$.
\end{proof}
\begin{lemma}
\label{lemma-colimit-category-fp-modules}
Supposons que $R = \colim_{\lambda \in \Lambda} R_\lambda$ soit une
colimite filtrante d'anneaux. Alors la catégorie des $R$-modules de
présentation finie est la colimite des catégories des $R_\lambda$-modules
de présentation finie. Plus précisément :
\begin{enumerate}
\item Étant donné un $R$-module $M$ de présentation finie, il existe
$\lambda \in \Lambda$ et un $R_\lambda$-module $M_\lambda$ de présentation
finie tels que $M \cong M_\lambda \otimes_{R_\lambda} R$.
\item Étant donnés $\lambda \in \Lambda$, des $R_\lambda$-modules de
présentation finie $M_\lambda, N_\lambda$ et un morphisme de $R$-modules
$\varphi : M_\lambda \otimes_{R_\lambda} R \to
N_\lambda \otimes_{R_\lambda} R$, il existe $\mu \geq \lambda$ et un
morphisme de $R_\mu$-modules
$\varphi_\mu : M_\lambda \otimes_{R_\lambda} R_\mu \to
N_\lambda \otimes_{R_\lambda} R_\mu$
tels que $\varphi = \varphi_\mu \otimes 1_R$.
\item Étant donnés $\lambda \in \Lambda$, des $R_\lambda$-modules de
présentation finie $M_\lambda, N_\lambda$ et des morphismes de
$R_\lambda$-modules $\varphi, \psi : M_\lambda \to N_\lambda$ tels que
$\varphi \otimes 1_R = \psi \otimes 1_R$, on a
$\varphi \otimes 1_{R_\mu} = \psi \otimes 1_{R_\mu}$ pour un certain
$\mu \geq \lambda$.
\end{enumerate}
\end{lemma}

\begin{proof}
Pour démontrer (1), choisissons une présentation
$R^{\oplus m} \to R^{\oplus n} \to M \to 0$.
Supposons que le premier morphisme soit donné par la matrice $A = (a_{ij})$.
Nous pouvons choisir $\lambda \in \Lambda$ et une matrice
$A_\lambda = (a_{\lambda, ij})$ à coefficients dans $R_\lambda$ qui
s'envoie sur $A$ dans $R$.
Il suffit alors de prendre pour $M_\lambda$ le $R_\lambda$-module ayant pour
présentation
$R_\lambda^{\oplus m} \to R_\lambda^{\oplus n} \to M_\lambda \to 0$,
où le premier morphisme est donné par $A_\lambda$.

\medskip\noindent
Les parties (2) et (3) résultent du
Lemme \ref{lemma-module-map-property-in-colimit}.
\end{proof}
\begin{lemma}
\label{lemma-algebra-map-property-in-colimit}
Soient $A$ un anneau et $B, C$ des $A$-algèbres.
Supposons que $R = \colim_{i \in I} R_i$ soit une colimite filtrante de
$A$-algèbres.
\begin{enumerate}
\item Si $B$ est une $A$-algèbre de type fini et si
$u, u' : B \to C$ sont des morphismes de $A$-algèbres tels que
$u \otimes 1 = u' \otimes 1 : B \otimes_A R \to C \otimes_A R$,
alors, pour un certain $i$, on a
$u \otimes 1 = u' \otimes 1 : B \otimes_A R_i \to C \otimes_A R_i$.
\item Si $C$ est une $A$-algèbre de type fini et si $u : B \to C$ est un
morphisme de $A$-algèbres tel que
$u \otimes 1 : B \otimes_A R \to C \otimes_A R$ soit surjectif, alors,
pour un certain $i$, le morphisme
$u \otimes 1 : B \otimes_A R_i \to C \otimes_A R_i$ est surjectif.
\item Si $C$ est de présentation finie sur $A$ et si
$v : C \otimes_A R \to B \otimes_A R$ est un morphisme de $R$-algèbres,
alors il existe $i$ et un morphisme de $R_i$-algèbres
$v_i : C \otimes_A R_i \to B \otimes_A R_i$ tels que
$v = v_i \otimes 1$.
\item Si $B$ est une $A$-algèbre de type fini, si $C$ est une $A$-algèbre
de présentation finie et si
$u \otimes 1 : B \otimes_A R \to C \otimes_A R$ est un isomorphisme,
alors, pour un certain $i$, le morphisme
$u \otimes 1 : B \otimes_A R_i \to C \otimes_A R_i$ est un isomorphisme.
\end{enumerate}
\end{lemma}
\begin{proof}
Pour démontrer (1), supposons que $u$ soit comme dans (1) et soient
$x_1, \ldots, x_m \in B$ des générateurs. Comme
$B \otimes_A R = \colim_i B \otimes_A R_i$, nous pouvons choisir
$i \in I$ tel que $u(x_j) \otimes 1 = u'(x_j) \otimes 1$ dans
$B \otimes_A R_i$, $j = 1, \ldots, m$. Pour un tel $i$, on a
$u \otimes 1 = u' \otimes 1 : B \otimes_A R_i \to C \otimes_A R_i$.

\medskip\noindent
Pour démontrer (2), supposons $u \otimes 1$ surjectif et soient
$y_1, \ldots, y_m \in C$ des générateurs. Comme
$B \otimes_A R = \colim_i B \otimes_A R_i$, nous pouvons choisir
$i \in I$ et $z_j \in B \otimes_A R_i$, $j = 1, \ldots, m$, dont les
images dans $C \otimes_A R$ sont égales à $y_j \otimes 1$.
Pour un tel $i$, le morphisme
$u \otimes 1 : B \otimes_A R_i \to C \otimes_A R_i$ est surjectif.

\medskip\noindent
Pour démontrer (3), soient $c_1, \ldots, c_m \in C$ des générateurs.
Soit $K = \Ker(A[x_1, \ldots, x_m] \to N)$, où le morphisme est donné par
la règle $x_j \mapsto \sum c_j$. Soient $f_1, \ldots, f_t$ des
générateurs de l'idéal $K$ de $A[x_1, \ldots, x_m]$.
Nous considérons $f_j = f_j(x_1, \ldots, x_m)$ comme un polynôme.
Comme $B \otimes_A R = \colim_i B \otimes_A R_i$, nous pouvons choisir
$i \in I$ et $z_j \in B \otimes_A R_i$, $j = 1, \ldots, m$, dont les
images dans $B \otimes_A R$ sont égales à $v(c_j \otimes 1)$.
Nous voulons utiliser les $z_j$ pour définir un morphisme
$v_i : C \otimes_A R_i \to B \otimes_A R_i$.
Puisque
$K \otimes_A R_i \to R_i[x_1, \ldots, x_m] \to
C \otimes_A R_i \to 0$
est une présentation, il suffit de vérifier que
$\xi_s = f_j(z_1, \ldots, z_m)$ est nul dans $B \otimes_A R_i$ pour chaque
$s = 1, \ldots, t$. Ce n'est pas nécessairement le cas, mais, puisque
l'image de $\xi_s$ dans $B \otimes_A R$ est nulle, cela le devient après
avoir augmenté quelque peu $i$.

\medskip\noindent
Pour démontrer (4), supposons que $u \otimes 1$ soit un isomorphisme, que
$B$ soit une $A$-algèbre de type fini et que $C$ soit une $A$-algèbre de
présentation finie. Soit
$v : B \otimes_A R \to C \otimes_A R$ un inverse de $u \otimes 1$.
Soit $v_i : C \otimes_A R_i \to B \otimes_A R_i$ comme dans la partie (3).
La partie (1) montre qu'après avoir augmenté $i$, on a
$v_i \circ (u \otimes 1) = \text{id}_{B \otimes_R R_i}$ et
$(u \otimes 1) \circ v_i = \text{id}_{C \otimes_R R_i}$.
\end{proof}
\begin{lemma}
\label{lemma-colimit-category-fp-algebras}
Supposons que $R = \colim_{\lambda \in \Lambda} R_\lambda$ soit une
colimite filtrante d'anneaux. Alors la catégorie des $R$-algèbres de
présentation finie est la colimite des catégories des $R_\lambda$-algèbres
de présentation finie. Plus précisément :
\begin{enumerate}
\item Étant donnée une $R$-algèbre $A$ de présentation finie, il existe
$\lambda \in \Lambda$ et une $R_\lambda$-algèbre $A_\lambda$ de
présentation finie telles que
$A \cong A_\lambda \otimes_{R_\lambda} R$.
\item Étant donnés $\lambda \in \Lambda$, des $R_\lambda$-algèbres de
présentation finie $A_\lambda, B_\lambda$ et un morphisme de $R$-algèbres
$\varphi : A_\lambda \otimes_{R_\lambda} R \to
B_\lambda \otimes_{R_\lambda} R$, il existe $\mu \geq \lambda$ et un
morphisme de $R_\mu$-algèbres
$\varphi_\mu : A_\lambda \otimes_{R_\lambda} R_\mu \to
B_\lambda \otimes_{R_\lambda} R_\mu$
tels que $\varphi = \varphi_\mu \otimes 1_R$.
\item Étant donnés $\lambda \in \Lambda$, des $R_\lambda$-algèbres de
présentation finie $A_\lambda, B_\lambda$ et des morphismes de
$R_\lambda$-algèbres
$\varphi_\lambda, \psi_\lambda : A_\lambda \to B_\lambda$ tels que
$\varphi \otimes 1_R = \psi \otimes 1_R$, on a
$\varphi \otimes 1_{R_\mu} = \psi \otimes 1_{R_\mu}$ pour un certain
$\mu \geq \lambda$.
\end{enumerate}
\end{lemma}

\begin{proof}
Pour démontrer (1), choisissons une présentation
$A = R[x_1, \ldots, x_n]/(f_1, \ldots, f_m)$.
Nous pouvons choisir $\lambda \in \Lambda$ et des éléments
$f_{\lambda, j} \in R_\lambda[x_1, \ldots, x_n]$ qui s'envoient sur
$f_j \in R[x_1, \ldots, x_n]$.
Il suffit alors de poser
$A_\lambda =
R_\lambda[x_1, \ldots, x_n]/(f_{\lambda, 1}, \ldots, f_{\lambda, m})$.

\medskip\noindent
Les parties (2) et (3) résultent du
Lemme \ref{lemma-algebra-map-property-in-colimit}.
\end{proof}
\begin{lemma}
\label{lemma-limit-no-condition-local}
Supposons que $R \to S$ soit un homomorphisme local d'anneaux locaux.
Il existe un ensemble filtrant $(\Lambda, \leq)$ et un système
d'homomorphismes locaux $R_\lambda \to S_\lambda$ d'anneaux locaux tels que :
\begin{enumerate}
\item La colimite du système $R_\lambda \to S_\lambda$ est égale à
$R \to S$.
\item Chaque $R_\lambda$ est essentiellement de type fini sur
$\mathbf{Z}$.
\item Chaque $S_\lambda$ est essentiellement de type fini sur
$R_\lambda$.
\end{enumerate}
\end{lemma}

\begin{proof}
Notons $\varphi : R  \to S$ le morphisme d'anneaux.
Soient $\mathfrak m \subset R$ l'idéal maximal de $R$ et
$\mathfrak n \subset S$ l'idéal maximal de $S$. Posons
$$
\Lambda = \{
(A, B)
\mid
A \subset R, B \subset S, \# A < \infty, \# B < \infty, \varphi(A) \subset B
\}.
$$
Nous prenons pour ordre partiel la relation d'inclusion. Pour chaque
$\lambda = (A, B) \in \Lambda$, soit $R'_\lambda$ la sous-$\mathbf{Z}$-algèbre
engendrée par les $a \in A$, et soit $S'_\lambda$ la sous-$\mathbf{Z}$-algèbre
engendrée par les $b$, $b \in B$.
Soit $R_\lambda$ la localisation de $R'_\lambda$ en l'idéal premier
$R'_\lambda \cap \mathfrak m$, et soit $S_\lambda$ la localisation de
$S'_\lambda$ en l'idéal premier $S'_\lambda \cap \mathfrak n$.
Voici le diagramme :
$$
\xymatrix{
B \ar[r] &
S'_\lambda \ar[r] &
S_\lambda \ar[r] &
S \\
A \ar[r] \ar[u] &
R'_\lambda \ar[r] \ar[u] &
R_\lambda \ar[r] \ar[u] &
R \ar[u]
}.
$$
Les morphismes de transition sont clairs. Nous laissons au lecteur la
démonstration des autres assertions.
\end{proof}
\begin{lemma}
\label{lemma-limit-essentially-finite-type}
Supposons que $R \to S$ soit un homomorphisme local d'anneaux locaux.
Supposons $S$ essentiellement de type fini sur $R$.
Il existe alors un ensemble filtrant $(\Lambda, \leq)$ et un système
d'homomorphismes locaux $R_\lambda \to S_\lambda$ d'anneaux locaux tels que :
\begin{enumerate}
\item La colimite du système $R_\lambda \to S_\lambda$ est égale à
$R \to S$.
\item Chaque $R_\lambda$ est essentiellement de type fini sur
$\mathbf{Z}$.
\item Chaque $S_\lambda$ est essentiellement de type fini sur
$R_\lambda$.
\item Pour chaque $\lambda \leq \mu$, le morphisme
$S_\lambda \otimes_{R_\lambda} R_\mu \to S_\mu$
présente $S_\mu$ comme la localisation d'un quotient de
$S_\lambda \otimes_{R_\lambda} R_\mu$.
\end{enumerate}
\end{lemma}
\begin{proof}
Notons $\varphi : R  \to S$ le morphisme d'anneaux.
Soient $\mathfrak m \subset R$ l'idéal maximal de $R$ et
$\mathfrak n \subset S$ l'idéal maximal de $S$.
Soient $x_1, \ldots, x_n \in S$ des éléments tels que $S$ soit une
localisation de la sous-$R$-algèbre de $S$ engendrée par
$x_1, \ldots, x_n$. Autrement dit, $S$ est un quotient d'une localisation
de l'anneau de polynômes $R[x_1, \ldots, x_n]$.

\medskip\noindent
Soit $\Lambda = \{ A \subset R \mid \# A < \infty\}$ l'ensemble des
parties finies de $R$. Nous prenons pour ordre partiel la relation
d'inclusion. Pour chaque $\lambda = A \in \Lambda$, soit $R'_\lambda$
la sous-$\mathbf{Z}$-algèbre engendrée par les $a \in A$, et soit
$S'_\lambda$ la sous-$\mathbf{Z}$-algèbre engendrée par les $\varphi(a)$,
$a \in A$, et par les éléments $x_1, \ldots, x_n$.
Soit $R_\lambda$ la localisation de $R'_\lambda$ en l'idéal premier
$R'_\lambda \cap \mathfrak m$, et soit $S_\lambda$ la localisation de
$S'_\lambda$ en l'idéal premier $S'_\lambda \cap \mathfrak n$.
Voici le diagramme :
$$
\xymatrix{
\varphi(A) \amalg \{x_i\} \ar[r] &
S'_\lambda \ar[r] &
S_\lambda \ar[r] &
S \\
A \ar[r] \ar[u] &
R'_\lambda \ar[r] \ar[u] &
R_\lambda \ar[r] \ar[u] &
R \ar[u]
}
$$
Il est clair que si $A \subset B$ correspond à $\lambda \leq \mu$ dans
$\Lambda$, il existe des morphismes canoniques
$R_\lambda \to R_\mu$ et $S_\lambda \to S_\mu$, et l'on obtient ainsi
un système sur l'ensemble filtrant $\Lambda$.

\medskip\noindent
L'assertion $R = \colim R_\lambda$ est claire : tous les morphismes
$R_\lambda \to R$ sont injectifs et tout élément de $R$ finit par
appartenir à l'image. Le même argument s'applique à
$S = \colim S_\lambda$. Les assertions (2) et (3) sont vraies par
construction. La dernière assertion résulte du fait que les morphismes
$S'_\lambda \otimes_{R'_\lambda} R'_\mu \to S'_\mu$
sont manifestement surjectifs.
\end{proof}
\begin{lemma}
\label{lemma-limit-essentially-finite-presentation}
Supposons que $R \to S$ soit un homomorphisme local d'anneaux locaux.
Supposons $S$ essentiellement de présentation finie sur $R$.
Il existe alors un ensemble filtrant $(\Lambda, \leq)$ et un système
d'homomorphismes locaux $R_\lambda \to S_\lambda$ d'anneaux locaux tels que :
\begin{enumerate}
\item La colimite du système $R_\lambda \to S_\lambda$ est égale à
$R \to S$.
\item Chaque $R_\lambda$ est essentiellement de type fini sur
$\mathbf{Z}$.
\item Chaque $S_\lambda$ est essentiellement de type fini sur
$R_\lambda$.
\item Pour chaque $\lambda \leq \mu$, le morphisme
$S_\lambda \otimes_{R_\lambda} R_\mu \to S_\mu$
présente $S_\mu$ comme la localisation de
$S_\lambda \otimes_{R_\lambda} R_\mu$ en un idéal premier.
\end{enumerate}
\end{lemma}

\begin{proof}
Par hypothèse, nous pouvons choisir un isomorphisme
$\Phi : (R[x_1, \ldots, x_n]/I)_{\mathfrak q} \to S$,
où $I \subset R[x_1, \ldots, x_n]$ est un idéal de type fini et
$\mathfrak q \subset R[x_1, \ldots, x_n]/I$ est premier.
(Notons que $R \cap \mathfrak q$ est égal à l'idéal maximal
$\mathfrak m$ de $R$.) Choisissons aussi des générateurs
$f_1, \ldots, f_m \in I$ de l'idéal $I$.
Écrivons de façon quelconque $R$ comme une colimite
$R = \colim R_\lambda$ indexée par un ensemble filtrant
$(\Lambda, \leq )$, chaque $R_\lambda$ étant local et essentiellement de
type fini sur $\mathbf{Z}$. Il existe $\lambda_0 \in \Lambda$ tel que
$f_j$ soit l'image d'un certain
$f_{j, \lambda_0} \in R_{\lambda_0}[x_1, \ldots, x_n]$.
Pour tout $\lambda \geq \lambda_0$, notons
$f_{j, \lambda} \in R_{\lambda}[x_1, \ldots, x_n]$ l'image de
$f_{j, \lambda_0}$. Nous obtenons ainsi un système de morphismes d'anneaux
$$
R_\lambda[x_1, \ldots, x_n]/(f_{1, \lambda}, \ldots, f_{m, \lambda})
\to
R[x_1, \ldots, x_n]/(f_1, \ldots, f_m) \to S
$$
Posons $\mathfrak q_\lambda$ égal à l'image réciproque de $\mathfrak q$.
Posons
$S_\lambda = (R_\lambda[x_1, \ldots, x_n]/
(f_{1, \lambda}, \ldots, f_{m, \lambda}))_{\mathfrak q_\lambda}$.
Nous laissons au lecteur le soin de vérifier que ce choix convient.
\end{proof}
\begin{remark}
\label{remark-suitable-systems-limits}
Supposons que $R \to S$ soit un homomorphisme local d'anneaux locaux,
essentiellement de présentation finie. Prenons un système quelconque
$(\Lambda, \leq)$, $R_\lambda \to S_\lambda$, possédant les propriétés
énumérées dans le Lemme \ref{lemma-limit-essentially-finite-type}.
Il peut arriver que ce soit le « mauvais » système, c'est-à-dire que la
propriété (4) du Lemme \ref{lemma-limit-essentially-finite-presentation}
ne soit pas satisfaite. En voici un exemple.
Soit $k = \mathbf{F}_2$. Considérons l'anneau
$$
R = \text{localisation de }
k[z, y_1, y_2, \ldots]/(y_i^2 - zy_{i + 1})
\text{ en }(z, y_1, y_2, \ldots)
$$
Posons $S = R/zR$. Comme système, prenons $\Lambda = \mathbf{N}$ et
$$
R_n = \text{localisation de }
k[z, y_1, \ldots, y_n]/(\{y_i^2 - zy_{i + 1}\}_{i \leq n-1})
\text{ en }(z, y_1, \ldots, y_n)
$$
et $S_n = R_n/(z, y_n^2)$. Aucun des morphismes
$S_n \otimes_{R_n} R_{n + 1} \to S_{n + 1}$
n'est une localisation (c'est-à-dire, en l'espèce, un isomorphisme),
puisque $1 \otimes y_{n + 1}^2$ s'envoie sur zéro.
Si l'on prend plutôt $S_n' = R_n/zR_n$, alors les morphismes
$S'_n \otimes_{R_n} R_{n + 1} \to S'_{n + 1}$ sont des isomorphismes.
La morale de cette remarque est qu'il faut choisir les systèmes avec
un peu de soin.
\end{remark}
\begin{lemma}
\label{lemma-limit-module-essentially-finite-presentation}
Supposons que $R \to S$ soit un homomorphisme local d'anneaux locaux et que
$S$ soit essentiellement de présentation finie sur $R$.
Soit $M$ un $S$-module de présentation finie.
Il existe alors un ensemble filtrant $(\Lambda, \leq)$, un système
d'homomorphismes locaux $R_\lambda \to S_\lambda$ d'anneaux locaux et des
$S_\lambda$-modules $M_\lambda$ tels que :
\begin{enumerate}
\item La colimite du système $R_\lambda \to S_\lambda$ est égale à
$R \to S$. La colimite du système des $M_\lambda$ est $M$.
\item Chaque $R_\lambda$ est essentiellement de type fini sur
$\mathbf{Z}$.
\item Chaque $S_\lambda$ est essentiellement de type fini sur
$R_\lambda$.
\item Chaque $M_\lambda$ est de type fini sur $S_\lambda$.
\item Pour chaque $\lambda \leq \mu$, le morphisme
$S_\lambda \otimes_{R_\lambda} R_\mu \to S_\mu$
présente $S_\mu$ comme la localisation de
$S_\lambda \otimes_{R_\lambda} R_\mu$ en un idéal premier.
\item Pour chaque $\lambda \leq \mu$, le morphisme
$M_\lambda \otimes_{S_\lambda} S_\mu \to M_\mu$
est un isomorphisme.
\end{enumerate}
\end{lemma}

\begin{proof}
Comme dans la démonstration du
Lemme \ref{lemma-limit-essentially-finite-presentation}, nous pouvons
d'abord écrire $R = \colim R_\lambda$ comme une colimite filtrante de
$\mathbf{Z}$-algèbres locales essentiellement de type fini.
Ensuite, nous pouvons supposer qu'il existe $\lambda_1 \in \Lambda$ et des
$f_{j, \lambda_1} \in R_{\lambda_1}[x_1, \ldots, x_n]$ tels que
$$
S =
\colim_{\lambda \geq \lambda_1} S_\lambda, \text{ avec }
S_\lambda =
(R_\lambda[x_1, \ldots, x_n]/
(f_{1, \lambda}, \ldots, f_{m, \lambda}))_{\mathfrak q_\lambda}
$$
Choisissons une présentation
$$
S^{\oplus s} \to S^{\oplus t} \to M \to 0
$$
de $M$ sur $S$. Soit $A \in \text{Mat}(t \times s, S)$ la matrice de la
présentation. Pour un certain $\lambda_2 \in \Lambda$,
$\lambda_2 \geq \lambda_1$, nous pouvons trouver une matrice
$A_{\lambda_2} \in \text{Mat}(t \times s, S_{\lambda_2})$ qui s'envoie sur
$A$. Pour tout $\lambda \geq \lambda_2$, posons
$M_\lambda = \Coker(S_\lambda^{\oplus s} \xrightarrow{A_\lambda}
S_\lambda^{\oplus t})$. Nous laissons au lecteur le soin de vérifier que
ce choix convient.
\end{proof}
\begin{lemma}
\label{lemma-limit-no-condition}
Supposons que $R \to S$ soit un morphisme d'anneaux.
Il existe alors un ensemble filtrant $(\Lambda, \leq)$ et un système de
morphismes d'anneaux $R_\lambda \to S_\lambda$ tels que :
\begin{enumerate}
\item La colimite du système $R_\lambda \to S_\lambda$ est égale à
$R \to S$.
\item Chaque $R_\lambda$ est de type fini sur $\mathbf{Z}$.
\item Chaque $S_\lambda$ est de type fini sur $R_\lambda$.
\end{enumerate}
\end{lemma}

\begin{proof}
C'est la version non locale du Lemme \ref{lemma-limit-no-condition-local}.
La démonstration est analogue et laissée au lecteur.
\end{proof}
\begin{lemma}
\label{lemma-limit-integral}
Supposons que $R \to S$ soit un morphisme d'anneaux et que $S$ soit entier
sur $R$. Il existe alors un ensemble filtrant $(\Lambda, \leq)$ et un
système de morphismes d'anneaux $R_\lambda \to S_\lambda$ tels que :
\begin{enumerate}
\item La colimite du système $R_\lambda \to S_\lambda$ est égale à
$R \to S$.
\item Chaque $R_\lambda$ est de type fini sur $\mathbf{Z}$.
\item Chaque $S_\lambda$ est fini sur $R_\lambda$.
\end{enumerate}
\end{lemma}

\begin{proof}
Considérons l'ensemble $\Lambda$ des couples $(E, F)$ où $E \subset R$ et
$F \subset S$ sont des parties finies, et où tout élément $f \in F$ est
racine d'un polynôme unitaire $P(X) \in R[X]$ dont les coefficients
appartiennent à $E$. Posons $(E, F) \leq (E', F')$ si $E \subset E'$ et
$F \subset F'$.
Étant donné $\lambda = (E, F) \in \Lambda$, soit
$R_\lambda \subset R$ la sous-$\mathbf{Z}$-algèbre de $R$ engendrée par
$E$, et soit $S_\lambda \subset S$ la sous-$\mathbf{Z}$-algèbre engendrée
par $F$ et par l'image de $E$ dans $S$.
Il est clair que $R = \colim R_\lambda$. On a
$S = \colim S_\lambda$, puisque tout élément de $S$ est entier sur $S$.
Les morphismes d'anneaux $R_\lambda \to S_\lambda$ sont finis d'après le
Lemme \ref{lemma-characterize-finite-in-terms-of-integral} et parce que
$S_\lambda$ est engendrée sur $R_\lambda$ par les éléments de $F$, qui sont
entiers sur $R_\lambda$ en vertu de notre condition sur les couples
$(E, F)$. Le lemme en résulte.
\end{proof}
\begin{lemma}
\label{lemma-limit-finite-type}
Supposons que $R \to S$ soit un morphisme d'anneaux et que $S$ soit de type
fini sur $R$. Il existe alors un ensemble filtrant $(\Lambda, \leq)$ et un
système de morphismes d'anneaux $R_\lambda \to S_\lambda$ tels que :
\begin{enumerate}
\item La colimite du système $R_\lambda \to S_\lambda$ est égale à
$R \to S$.
\item Chaque $R_\lambda$ est de type fini sur $\mathbf{Z}$.
\item Chaque $S_\lambda$ est de type fini sur $R_\lambda$.
\item Pour chaque $\lambda \leq \mu$, le morphisme
$S_\lambda \otimes_{R_\lambda} R_\mu \to S_\mu$
présente $S_\mu$ comme un quotient de
$S_\lambda \otimes_{R_\lambda} R_\mu$.
\end{enumerate}
\end{lemma}

\begin{proof}
C'est la version non locale du
Lemme \ref{lemma-limit-essentially-finite-type}.
La démonstration est analogue et laissée au lecteur.
\end{proof}
\begin{lemma}
\label{lemma-limit-finite-presentation}
Supposons que $R \to S$ soit un morphisme d'anneaux et que $S$ soit de
présentation finie sur $R$. Il existe alors un ensemble filtrant
$(\Lambda, \leq)$ et un système de morphismes d'anneaux
$R_\lambda \to S_\lambda$ tels que :
\begin{enumerate}
\item La colimite du système $R_\lambda \to S_\lambda$ est égale à
$R \to S$.
\item Chaque $R_\lambda$ est de type fini sur $\mathbf{Z}$.
\item Chaque $S_\lambda$ est de type fini sur $R_\lambda$.
\item Pour chaque $\lambda \leq \mu$, le morphisme
$S_\lambda \otimes_{R_\lambda} R_\mu \to S_\mu$
est un isomorphisme.
\end{enumerate}
\end{lemma}

\begin{proof}
C'est la version non locale du
Lemme \ref{lemma-limit-essentially-finite-presentation}.
La démonstration est analogue et laissée au lecteur.
\end{proof}
\begin{lemma}
\label{lemma-limit-module-finite-presentation}
Supposons que $R \to S$ soit un morphisme d'anneaux et que $S$ soit de
présentation finie sur $R$. Soit $M$ un $S$-module de présentation finie.
Il existe alors un ensemble filtrant $(\Lambda, \leq)$, un système de
morphismes d'anneaux $R_\lambda \to S_\lambda$ et des $S_\lambda$-modules
$M_\lambda$ tels que :
\begin{enumerate}
\item La colimite du système $R_\lambda \to S_\lambda$ est égale à
$R \to S$. La colimite du système des $M_\lambda$ est $M$.
\item Chaque $R_\lambda$ est de type fini sur $\mathbf{Z}$.
\item Chaque $S_\lambda$ est de type fini sur $R_\lambda$.
\item Chaque $M_\lambda$ est de type fini sur $S_\lambda$.
\item Pour chaque $\lambda \leq \mu$, le morphisme
$S_\lambda \otimes_{R_\lambda} R_\mu \to S_\mu$
est un isomorphisme.
\item Pour chaque $\lambda \leq \mu$, le morphisme
$M_\lambda \otimes_{S_\lambda} S_\mu \to M_\mu$
est un isomorphisme.
\end{enumerate}
En particulier, pour tout $\lambda \in \Lambda$, on a
$$
M = M_\lambda \otimes_{S_\lambda} S
= M_\lambda \otimes_{R_\lambda} R.
$$
\end{lemma}

\begin{proof}
C'est la version non locale du
Lemme \ref{lemma-limit-module-essentially-finite-presentation}.
La démonstration est analogue et laissée au lecteur.
\end{proof}








\section{Autres critères de platitude}
\label{section-more-flatness-criteria}

\noindent
Le lemme suivant est souvent utilisé en géométrie algébrique pour montrer
qu'un morphisme fini d'une surface normale vers une surface lisse est plat.
Il constitue une réciproque partielle du
Lemme \ref{lemma-finite-flat-over-regular-CM}, puisqu'un homomorphisme local
injectif et fini d'anneaux satisfait assurément la condition (3).

\begin{lemma}
\label{lemma-CM-over-regular-flat}
\begin{slogan}
Platitude miraculeuse
\end{slogan}
Soit $R \to S$ un homomorphisme local d'anneaux locaux noethériens.
Supposons que
\begin{enumerate}
\item $R$ soit régulier,
\item $S$ soit de Cohen-Macaulay,
\item $\dim(S) = \dim(R) + \dim(S/\mathfrak m_R S)$.
\end{enumerate}
Alors $R \to S$ est plat.
\end{lemma}

\begin{proof}
Raisonnons par récurrence sur $\dim(R)$. Le cas $\dim(R) = 0$ est trivial,
car $R$ est alors un corps. Supposons $\dim(R) > 0$. D'après (3), cela
implique que $\dim(S) > 0$. Soient
$\mathfrak q_1, \ldots, \mathfrak q_r$ les idéaux premiers minimaux de $S$.
Notons que $\mathfrak q_i \not \supset \mathfrak m_R S$, puisque
$$
\dim(S/\mathfrak q_i) = \dim(S) > \dim(S/\mathfrak m_R S)
$$
la première égalité résultant du Lemme \ref{lemma-maximal-chain-CM}, et
l'inégalité de (3). Ainsi,
$\mathfrak p_i = R \cap \mathfrak q_i$ n'est pas égal à $\mathfrak m_R$.
Choisissons $x \in \mathfrak m_R$, $x \not \in \mathfrak m_R^2$ et
$x \not \in \mathfrak p_i$ ; voir le Lemme \ref{lemma-silly}.
On voit donc que $x$ n'appartient à aucun idéal premier minimal de $S$.
Par conséquent, $x$ est un non-diviseur de zéro sur $S$ d'après (2) ; voir
le Lemme \ref{lemma-reformulate-CM}. En outre, $S/xS$ est de Cohen-Macaulay
et $\dim(S/xS) = \dim(S) - 1$.
D'après (1) et le Lemme \ref{lemma-regular-ring-CM}, l'anneau $R/xR$ est
régulier et $\dim(R/xR) = \dim(R) - 1$.
L'hypothèse de récurrence montre que $R/xR \to S/xS$ est plat.
On conclut alors par le Lemme \ref{lemma-variant-local-criterion-flatness}
et la remarque qui le suit.
\end{proof}

\begin{lemma}
\label{lemma-flat-over-regular}
Soit $R \to S$ un homomorphisme d'anneaux locaux noethériens.
Supposons que $R$ soit un anneau local régulier et qu'un système régulier
de paramètres s'envoie sur une suite régulière dans $S$. Alors
$R \to S$ est plat.
\end{lemma}

\begin{proof}
Supposons que $x_1, \ldots, x_d$ soit un système de paramètres de $R$
qui s'envoie sur une suite régulière dans $S$. Notons que
$S/(x_1, \ldots, x_d)S$ est plat sur $R/(x_1, \ldots, x_d)$, puisque ce
dernier est un corps. Puis $x_d$ est un non-diviseur de zéro dans
$S/(x_1, \ldots, x_{d - 1})S$ ; par conséquent,
$S/(x_1, \ldots, x_{d - 1})S$ est plat sur
$R/(x_1, \ldots, x_{d - 1})$ d'après le critère local de platitude (voir le
Lemme \ref{lemma-variant-local-criterion-flatness} et les remarques qui le
suivent). Puis $x_{d - 1}$ est un non-diviseur de zéro dans
$S/(x_1, \ldots, x_{d - 2})S$ ; ainsi,
$S/(x_1, \ldots, x_{d - 2})S$ est plat sur
$R/(x_1, \ldots, x_{d - 2})$ d'après le critère local de platitude (voir le
Lemme \ref{lemma-variant-local-criterion-flatness} et les remarques qui le
suivent). On continue ainsi jusqu'à conclure que $S$ est plat sur $R$.
\end{proof}

\noindent
Le lemme suivant est la clef qui permet de montrer que les résultats sur les
modules de présentation finie sur des anneaux de présentation finie sur un
anneau de base découlent des résultats correspondants sur les modules de type fini
dans le cas noethérien.

\begin{lemma}
\label{lemma-colimit-eventually-flat}
Soient $R \to S$, $M$, $\Lambda$, $R_\lambda \to S_\lambda$ et
$M_\lambda$ comme dans le
Lemme \ref{lemma-limit-module-essentially-finite-presentation}.
Supposons que $M$ soit plat sur $R$. Alors, pour un certain
$\lambda \in \Lambda$, le module $M_\lambda$ est plat sur $R_\lambda$.
\end{lemma}

\begin{proof}
Choisissons $\lambda \in \Lambda$ et considérons
$$
\text{Tor}_1^{R_\lambda}(M_\lambda, R_\lambda/\mathfrak m_\lambda)
=
\Ker(\mathfrak m_\lambda \otimes_{R_\lambda} M_\lambda
\to M_\lambda).
$$
Voir la Remarque \ref{remark-Tor-ring-mod-ideal}. Le second membre montre
qu'il s'agit d'un $S_\lambda$-module de type fini (car $S_\lambda$ est
noethérien et les modules en question sont de type fini).
Soient $\xi_1, \ldots, \xi_n$ des générateurs.
Comme $M$ est plat sur $R$, on a
$0 = \Ker(\mathfrak m_\lambda R \otimes_R M \to M)$.
Puisque $\otimes$ commute aux colimites, il existe
$\lambda' \geq \lambda$ tel que chaque $\xi_i$ s'envoie sur zéro dans
$\mathfrak m_{\lambda}R_{\lambda'} \otimes_{R_{\lambda'}} M_{\lambda'}$.
On voit donc que
$$
\text{Tor}_1^{R_\lambda}(M_\lambda, R_\lambda/\mathfrak m_\lambda)
\longrightarrow
\text{Tor}_1^{R_{\lambda'}}(M_{\lambda'},
R_{\lambda'}/\mathfrak m_{\lambda}R_{\lambda'})
$$
est nul. Notons que
$M_\lambda \otimes_{R_\lambda} R_\lambda/\mathfrak m_\lambda$
est plat sur $R_\lambda/\mathfrak m_\lambda$, puisque ce dernier anneau
est un corps. Nous pouvons donc appliquer le
Lemme \ref{lemma-another-variant-local-criterion-flatness} et conclure que
$M_{\lambda'}$ est plat sur $R_{\lambda'}$.
\end{proof}

\noindent
À l'aide du lemme précédent, nous pouvons commencer à redémontrer les
résultats de la section \ref{section-criteria-flatness} dans le cas non
noethérien.

\begin{lemma}
\label{lemma-mod-injective-general}
Supposons que $R \to S$ soit un homomorphisme local d'anneaux locaux.
Notons $\mathfrak m$ l'idéal maximal de $R$.
Soit $u : M \to N$ un morphisme de $S$-modules. Supposons que
\begin{enumerate}
\item $S$ soit essentiellement de présentation finie sur $R$,
\item $M$ et $N$ soient de présentation finie sur $S$,
\item $N$ soit plat sur $R$, et
\item $\overline{u} : M/\mathfrak mM \to N/\mathfrak mN$ soit injectif.
\end{enumerate}
Alors $u$ est injectif et $N/u(M)$ est plat sur $R$.
\end{lemma}

\begin{proof}
D'après le Lemme \ref{lemma-limit-module-essentially-finite-presentation}
et sa démonstration, on peut trouver un système de morphismes locaux
$R_\lambda \to S_\lambda$, ainsi que des morphismes de $S_\lambda$-modules
$u_\lambda : M_\lambda \to N_\lambda$, qui satisfont, pour $N$ comme pour
$M$, aux conclusions (1) -- (6) de ce lemme, et tels que la colimite des
morphismes $u_\lambda$ soit $u$.
D'après le Lemme \ref{lemma-colimit-eventually-flat}, nous pouvons supposer
que $N_\lambda$ est plat sur $R_\lambda$ pour tout $\lambda$ suffisamment
grand. Notons $\mathfrak m_\lambda \subset R_\lambda$ l'idéal maximal et
$\kappa_\lambda = R_\lambda / \mathfrak m_\lambda$, et notons
resp.\ $\kappa = R/\mathfrak m$ les corps résiduels.

\medskip\noindent
Considérons le morphisme
$$
\Psi_\lambda :
M_\lambda/\mathfrak m_\lambda M_\lambda \otimes_{\kappa_\lambda} \kappa
\longrightarrow
M/\mathfrak m M.
$$
Puisque $S_\lambda/\mathfrak m_\lambda S_\lambda$ est essentiellement de
type fini sur le corps $\kappa_\lambda$, le produit tensoriel
$S_\lambda/\mathfrak m_\lambda S_\lambda
\otimes_{\kappa_\lambda} \kappa$ est essentiellement de type fini sur
$\kappa$. C'est donc un anneau noethérien, et nous en concluons que le noyau
de $\Psi_\lambda$ est de type fini.
Comme $M/\mathfrak m M$ est la colimite du système des
$M_\lambda/\mathfrak m_\lambda M_\lambda$ et que $\kappa$ est la colimite
des corps $\kappa_\lambda$, il existe $\lambda' > \lambda$ tel que le noyau
de $\Psi_\lambda$ soit engendré par le noyau de
$$
\Psi_{\lambda, \lambda'} :
M_\lambda/\mathfrak m_\lambda M_\lambda
\otimes_{\kappa_\lambda}
\kappa_{\lambda'}
\longrightarrow
M_{\lambda'}/\mathfrak m_{\lambda'} M_{\lambda'}.
$$
Par construction, il existe une partie multiplicative
$W \subset S_\lambda \otimes_{R_\lambda} R_{\lambda'}$ telle que
$S_{\lambda'} = W^{-1}(S_\lambda \otimes_{R_\lambda} R_{\lambda'})$ et
$$
W^{-1}(M_\lambda/\mathfrak m_\lambda M_\lambda
\otimes_{\kappa_\lambda}
\kappa_{\lambda'})
=
M_{\lambda'}/\mathfrak m_{\lambda'} M_{\lambda'}.
$$
Supposons maintenant que $x$ soit un élément du noyau de
$$
\Psi_{\lambda'} :
M_{\lambda'}/\mathfrak m_{\lambda'} M_{\lambda'}
\otimes_{\kappa_{\lambda'}} \kappa
\longrightarrow
M/\mathfrak m M.
$$
Alors, pour un certain $w \in W$, on a
$wx \in M_\lambda/\mathfrak m_\lambda M_\lambda \otimes \kappa$.
Ainsi $wx \in \Ker(\Psi_\lambda)$. Par conséquent, $wx$ est une
combinaison linéaire d'éléments du noyau de
$\Psi_{\lambda, \lambda'}$. Il s'ensuit que
$wx = 0$ dans $M_{\lambda'}/\mathfrak m_{\lambda'} M_{\lambda'}
\otimes_{\kappa_{\lambda'}} \kappa$, donc $x = 0$ puisque $w$ est
inversible dans $S_{\lambda'}$.
Nous concluons que le noyau de $\Psi_{\lambda'}$ est nul pour tout
$\lambda'$ suffisamment grand !

\medskip\noindent
Le résultat du paragraphe précédent nous permet de supposer que le noyau
de $\Psi_\lambda$ est nul pour tout $\lambda$ suffisamment grand, ce qui
implique que le morphisme
$M_\lambda/\mathfrak m_\lambda M_\lambda \to M/\mathfrak m M$
est injectif. Combiné à l'injectivité de $\overline{u}$, cela implique
formellement que
$\overline{u_\lambda} : M_\lambda/\mathfrak m_\lambda M_\lambda
\to N_\lambda/\mathfrak m_\lambda N_\lambda$ est lui aussi injectif.
D'après le Lemme \ref{lemma-mod-injective}, nous concluons que, pour tout
$\lambda$ suffisamment grand, le morphisme $u_\lambda$ est injectif et que
$N_\lambda/u_\lambda(M_\lambda)$ est plat sur $R_\lambda$.
Le lemme en résulte.
\end{proof}

\begin{lemma}
\label{lemma-grothendieck-general}
Supposons que $R \to S$ soit un homomorphisme local d'anneaux locaux.
Notons $\mathfrak m$ l'idéal maximal de $R$. Supposons que
\begin{enumerate}
\item $S$ soit essentiellement de présentation finie sur $R$,
\item $S$ soit plat sur $R$, et
\item $f \in S$ soit un non-diviseur de zéro dans $S/{\mathfrak m}S$.
\end{enumerate}
Alors $S/fS$ est plat sur $R$, et $f$ est un non-diviseur de zéro dans $S$.
\end{lemma}

\begin{proof}
Cela résulte directement du Lemme \ref{lemma-mod-injective-general}.
\end{proof}

\begin{lemma}
\label{lemma-grothendieck-regular-sequence-general}
Supposons que $R \to S$ soit un homomorphisme local d'anneaux locaux.
Notons $\mathfrak m$ l'idéal maximal de $R$. Supposons que
\begin{enumerate}
\item $R \to S$ soit essentiellement de présentation finie,
\item $R \to S$ soit plat, et
\item $f_1, \ldots, f_c$ soit une suite d'éléments de $S$ telle que les
images $\overline{f}_1, \ldots, \overline{f}_c$ forment une suite régulière
dans $S/{\mathfrak m}S$.
\end{enumerate}
Alors $f_1, \ldots, f_c$ est une suite régulière dans $S$, et chacun des
quotients $S/(f_1, \ldots, f_i)$ est plat sur $R$.
\end{lemma}

\begin{proof}
Cela résulte d'une récurrence et du Lemme \ref{lemma-grothendieck-general}.
\end{proof}

\noindent
Voici la version du critère local de platitude pour les morphismes locaux
d'anneaux qui sont localement de présentation finie.

\begin{lemma}
\label{lemma-variant-local-criterion-flatness-general}
Soit $R \to S$ un homomorphisme local d'anneaux locaux.
Soit $I \not = R$ un idéal de $R$. Soit $M$ un $S$-module. Supposons que
\begin{enumerate}
\item $S$ soit essentiellement de présentation finie sur $R$,
\item $M$ soit de présentation finie sur $S$,
\item $\text{Tor}_1^R(M, R/I) = 0$, et
\item $M/IM$ soit plat sur $R/I$.
\end{enumerate}
Alors $M$ est plat sur $R$.
\end{lemma}

\begin{proof}
Soient $\Lambda$, $R_\lambda \to S_\lambda$ et $M_\lambda$ comme dans le
Lemme \ref{lemma-limit-module-essentially-finite-presentation}.
Notons $I_\lambda \subset R_\lambda$ l'image réciproque de $I$.
Dans ce cas, le système constitué par
$R/I \to S/IS$, $M/IM$, $R_\lambda \to S_\lambda/I_\lambda S_\lambda$
et $M_\lambda/I_\lambda M_\lambda$ satisfait lui aussi les conclusions du
Lemme \ref{lemma-limit-module-essentially-finite-presentation}.
D'après le Lemme \ref{lemma-colimit-eventually-flat}, nous pouvons donc
supposer, après avoir restreint l'ensemble d'indices $\Lambda$, que
$M_\lambda/I_\lambda M_\lambda$ est plat pour tout $\lambda$.
Choisissons un $\lambda$ et considérons
$$
\text{Tor}_1^{R_\lambda}(M_\lambda, R_\lambda/I_\lambda)
=
\Ker(I_\lambda \otimes_{R_\lambda} M_\lambda
\to M_\lambda).
$$
Voir la Remarque \ref{remark-Tor-ring-mod-ideal}. Le second membre montre
qu'il s'agit d'un $S_\lambda$-module de type fini (car $S_\lambda$ est
noethérien et les modules en question sont de type fini).
Soient $\xi_1, \ldots, \xi_n$ des générateurs.
Comme $\text{Tor}_1^R(M, R/I) = 0$ et puisque $\otimes$ commute aux
colimites, il existe $\lambda' \geq \lambda$ tel que chaque $\xi_i$
s'envoie sur zéro dans
$\text{Tor}_1^{R_{\lambda'}}(M_{\lambda'},
R_{\lambda'}/I_{\lambda'})$.
Le composé des morphismes
$$
\xymatrix{
R_{\lambda'} \otimes_{R_\lambda}
\text{Tor}_1^{R_\lambda}(M_\lambda, R_\lambda/I_\lambda)
\ar[d]^{\text{surjectif d'après le lemme \ref{lemma-surjective-on-tor-one}}} \\
\text{Tor}_1^{R_\lambda}(M_\lambda, R_{\lambda'}/I_\lambda R_{\lambda'})
\ar[d]^{\text{surjectif après localisation d'après le
lemme \ref{lemma-surjective-on-tor-one-trivial}}} \\
\text{Tor}_1^{R_{\lambda'}}(M_{\lambda'}, R_{\lambda'}/I_\lambda R_{\lambda'})
\ar[d]^{\text{surjectif d'après le lemme \ref{lemma-surjective-on-tor-one}}} \\
\text{Tor}_1^{R_{\lambda'}}(M_{\lambda'}, R_{\lambda'}/I_{\lambda'}).
}
$$
est surjectif après une localisation, pour les raisons indiquées.
La localisation est nécessaire parce que $M_{\lambda'}$ n'est pas égal à
$M_\lambda \otimes_{R_\lambda} R_{\lambda'}$. En effet, il est égal à
$M_\lambda \otimes_{S_\lambda} S_{\lambda'}$, et $S_{\lambda'}$ est la
localisation de $S_{\lambda} \otimes_{R_\lambda} R_{\lambda'}$, d'où
l'assertion après localisation (ou tensorisation par $S_{\lambda'}$).
Notons que le Lemme \ref{lemma-surjective-on-tor-one} s'applique à la
première et à la troisième flèche, parce que
$M_\lambda/I_\lambda M_\lambda$ est plat sur $R_\lambda/I_\lambda$ et que
$M_{\lambda'}/I_\lambda M_{\lambda'}$ est plat sur
$R_{\lambda'}/I_\lambda R_{\lambda'}$, car c'est un changement de base du
module plat $M_\lambda/I_\lambda M_\lambda$.
Le composé envoie les générateurs $\xi_i$ sur zéro, comme nous l'avons
expliqué ci-dessus. Nous concluons enfin que
$\text{Tor}_1^{R_{\lambda'}}(M_{\lambda'},
R_{\lambda'}/I_{\lambda'})$ est nul. Cela implique que $M_{\lambda'}$
est plat sur $R_{\lambda'}$ d'après le
Lemme \ref{lemma-variant-local-criterion-flatness}.
\end{proof}

\noindent
Comparer le lemme ci-dessous au
Lemme \ref{lemma-criterion-flatness-fibre-Noetherian}
(cas des anneaux locaux noethériens) et au
Lemme \ref{lemma-criterion-flatness-fibre-nilpotent}
(cas d'un idéal nilpotent dans la base).

\begin{lemma}[Crit\`ere de platitude par fibres]
\label{lemma-criterion-flatness-fibre}
Soient $R$, $S$ et $S'$ des anneaux locaux, et soient
$R \to S \to S'$ des homomorphismes locaux d'anneaux.
Soit $M$ un $S'$-module. Soit $\mathfrak m \subset R$ l'idéal maximal.
Supposons que
\begin{enumerate}
\item les morphismes d'anneaux $R \to S$ et $R \to S'$ soient
essentiellement de présentation finie ;
\item le module $M$ soit de présentation finie sur $S'$ ;
\item le module $M$ ne soit pas nul ;
\item le module $M/\mathfrak mM$ soit un $S/\mathfrak mS$-module plat ;
\item le module $M$ soit un $R$-module plat.
\end{enumerate}
Alors $S$ est plat sur $R$ et $M$ est un $S$-module plat.
\end{lemma}

\begin{proof}
Comme dans la démonstration du
Lemme \ref{lemma-limit-essentially-finite-presentation}, nous pouvons
d'abord écrire $R = \colim R_\lambda$ comme une colimite filtrante de
$\mathbf{Z}$-algèbres locales essentiellement de type fini.
Notons $\mathfrak p_\lambda$ l'idéal maximal de $R_\lambda$.
Ensuite, nous pouvons supposer que, pour un certain
$\lambda_1 \in \Lambda$, il existe des
$f_{j, \lambda_1} \in R_{\lambda_1}[x_1, \ldots, x_n]$ tels que
$$
S =
\colim_{\lambda \geq \lambda_1} S_\lambda, \text{ avec }
S_\lambda =
(R_\lambda[x_1, \ldots, x_n]/
(f_{1, \lambda}, \ldots, f_{u, \lambda}))_{\mathfrak q_\lambda}
$$
Pour un certain $\lambda_2 \in \Lambda$, $\lambda_2 \geq \lambda_1$, il
existe des
$g_{j, \lambda_2} \in R_{\lambda_2}[x_1, \ldots, x_n, y_1, \ldots, y_m]$
d'images
$\overline{g}_{j, \lambda_2} \in S_{\lambda_2}[y_1, \ldots, y_m]$
telles que
$$
S' =
\colim_{\lambda \geq \lambda_2} S'_\lambda, \text{ avec }
S'_\lambda =
(S_\lambda[y_1, \ldots, y_m]/
(\overline{g}_{1, \lambda}, \ldots,
\overline{g}_{v, \lambda}))_{\overline{\mathfrak q}'_\lambda}
$$
Notons que cela implique aussi
$$
S'_\lambda =
(R_\lambda[x_1, \ldots, x_n, y_1, \ldots, y_m]/
(g_{1, \lambda}, \ldots, g_{v, \lambda}))_{\mathfrak q'_\lambda}
$$
Choisissons une présentation
$$
(S')^{\oplus s} \to (S')^{\oplus t} \to M \to 0
$$
de $M$ sur $S'$. Soit $A \in \text{Mat}(t \times s, S')$ la matrice de
la présentation. Pour un certain $\lambda_3 \in \Lambda$,
$\lambda_3 \geq \lambda_2$, nous pouvons trouver une matrice
$A_{\lambda_3} \in \text{Mat}(t \times s, S_{\lambda_3})$ qui s'envoie sur
$A$. Pour tout $\lambda \geq \lambda_3$, posons
$M_\lambda = \Coker((S'_\lambda)^{\oplus s} \xrightarrow{A_\lambda}
(S'_\lambda)^{\oplus t})$.

\medskip\noindent
Avec ces choix, pour tous $\lambda_3 \leq \lambda \leq \mu$, le morphisme
$S_\lambda \otimes_{R_{\lambda}} R_\mu \to S_\mu$ est une localisation,
$S'_\lambda \otimes_{S_{\lambda}} S_\mu \to S'_\mu$ est une localisation,
et le morphisme $M_\lambda \otimes_{S'_\lambda} S'_\mu \to M_\mu$ est un
isomorphisme. Cela implique aussi que
$S'_\lambda \otimes_{R_{\lambda}} R_\mu \to S'_\mu$ est une localisation.
Ainsi, puisque $M$ est plat sur $R$, le
Lemme \ref{lemma-colimit-eventually-flat} montre que, pour tout $\lambda$
assez grand, le module $M_\lambda$ est plat sur $R_\lambda$.
Notons en outre que
$
\mathfrak m = \colim \mathfrak p_\lambda
$,
$
S/\mathfrak mS = \colim S_\lambda/\mathfrak p_\lambda S_\lambda
$,
$
S'/\mathfrak mS' = \colim S'_\lambda/\mathfrak p_\lambda S'_\lambda
$,
et
$
M/\mathfrak mM = \colim M_\lambda/\mathfrak p_\lambda M_\lambda
$. De plus, pour tous $\lambda_3 \leq \lambda \leq \mu$, les propriétés
énumérées ci-dessus montrent que
$$
S'_\lambda/\mathfrak p_\lambda S'_\lambda
\otimes_{S_{\lambda}/\mathfrak p_\lambda S_\lambda}
S_\mu/\mathfrak p_\mu S_\mu
\longrightarrow
S'_\mu/\mathfrak p_\mu S'_\mu
$$
est une localisation, et que le morphisme
$$
M_\lambda / \mathfrak p_\lambda M_\lambda
\otimes_{S'_\lambda/\mathfrak p_\lambda S'_\lambda}
S'_\mu /\mathfrak p_\mu S'_\mu
\longrightarrow
M_\mu/\mathfrak p_\mu M_\mu
$$
est un isomorphisme. Par conséquent, le système
$(S_\lambda/\mathfrak p_\lambda S_\lambda \to
S'_\lambda/\mathfrak p_\lambda S'_\lambda,
M_\lambda/\mathfrak p_\lambda M_\lambda)$ est lui aussi un système comme
dans le Lemme \ref{lemma-limit-module-essentially-finite-presentation}.
Nous pouvons appliquer une nouvelle fois le
Lemme \ref{lemma-colimit-eventually-flat}, car $M/\mathfrak m M$ est supposé
plat sur $S/\mathfrak mS$, et nous voyons que
$M_\lambda/\mathfrak p_\lambda M_\lambda$ est plat sur
$S_\lambda/\mathfrak p_\lambda S_\lambda$ pour tout $\lambda$ assez grand.
Ainsi, pour $\lambda$ assez grand, les données
$R_\lambda \to S_\lambda \to S'_\lambda, M_\lambda$ satisfont les
hypothèses du Lemme \ref{lemma-criterion-flatness-fibre-Noetherian}.
Choisissons un tel $\lambda$. Alors
$S = S_\lambda \otimes_{R_\lambda} R$ est plat sur $R$, et
$M = M_\lambda \otimes_{S_\lambda} S$ est plat sur $S$ (puisqu'un
changement de base d'un module plat est plat).
\end{proof}

\noindent
Le résultat suivant est une conséquence facile du « crit\`ere de platitude
par fibres », Lemme \ref{lemma-criterion-flatness-fibre}. Pour d'autres
résultats de ce genre, voir Compléments sur la platitude,
section \ref{flat-section-introduction}.

\begin{lemma}
\label{lemma-criterion-flatness-fibre-fp-over-ft}
Soient $R$, $S$ et $S'$ des anneaux locaux, et soient
$R \to S \to S'$ des homomorphismes locaux d'anneaux.
Soit $M$ un $S'$-module. Soit $\mathfrak m \subset R$ l'idéal maximal.
Supposons que
\begin{enumerate}
\item $R \to S'$ soit essentiellement de présentation finie,
\item $R \to S$ soit essentiellement de type fini,
\item $M$ soit de présentation finie sur $S'$,
\item $M$ ne soit pas nul,
\item $M/\mathfrak mM$ soit un $S/\mathfrak mS$-module plat, et
\item $M$ soit un $R$-module plat.
\end{enumerate}
Alors $S$ est essentiellement de présentation finie et plat sur $R$, et
$M$ est un $S$-module plat.
\end{lemma}

\begin{proof}
Comme $S$ est essentiellement de présentation finie sur $R$, nous pouvons
écrire $S = C_{\overline{\mathfrak q}}$ pour une certaine $R$-algèbre $C$
de type fini. Écrivons $C = R[x_1, \ldots, x_n]/I$.
Notons $\mathfrak q \subset R[x_1, \ldots, x_n]$ l'idéal premier
correspondant à $\overline{\mathfrak q}$. Alors $S = B/J$, où
$B = R[x_1, \ldots, x_n]_{\mathfrak q}$ est essentiellement de présentation
finie sur $R$ et $J = IB$. Nous pouvons trouver
$f_1, \ldots, f_k \in J$ tels que les images
$\overline{f}_i \in B/\mathfrak mB$ engendrent l'image $\overline{J}$ de
$J$ dans l'anneau noethérien $B/\mathfrak mB$.
Il existe donc des idéaux de type fini $J' \subset J$ tels que
$B/J' \to B/J$ induise un isomorphisme
$$
(B/J') \otimes_R R/\mathfrak m \longrightarrow
B/J \otimes_R R/\mathfrak m = S/\mathfrak mS.
$$
Pour tout $J'$ comme ci-dessus, le
Lemme \ref{lemma-criterion-flatness-fibre} s'applique aux morphismes
d'anneaux
$$
R \longrightarrow B/J' \longrightarrow S'
$$
et au module $M$. Nous concluons donc que $B/J'$ est plat sur $R$ pour tout
choix de $J'$ comme ci-dessus. Maintenant, si $J' \subset J' \subset J$
sont deux idéaux de type fini comme ci-dessus, nous concluons que
$B/J' \to B/J''$ est un morphisme surjectif entre des $R$-algèbres plates
et essentiellement de présentation finie, qui est un isomorphisme modulo
$\mathfrak m$. Par conséquent, le Lemme \ref{lemma-mod-injective-general}
implique que $B/J' = B/J''$, c'est-à-dire $J' = J''$.
Cela signifie clairement que $J$ est de type fini, c'est-à-dire que $S$ est
essentiellement de présentation finie sur $R$. Nous pouvons donc appliquer
le Lemme \ref{lemma-criterion-flatness-fibre} à $R \to S \to S'$, ce qui
achève la démonstration.
\end{proof}

\begin{lemma}[Crit\`ere de platitude par fibres : cas localement nilpotent]
\label{lemma-criterion-flatness-fibre-locally-nilpotent}
Soit
$$
\xymatrix{
S \ar[rr] & & S' \\
& R \ar[lu] \ar[ru]
}
$$
un diagramme commutatif dans la catégorie des anneaux.
Soient $I \subset R$ un idéal localement nilpotent et $M$ un $S'$-module.
Supposons que
\begin{enumerate}
\item $R \to S$ soit de type fini,
\item $R \to S'$ soit de présentation finie,
\item $M$ soit un $S'$-module de présentation finie,
\item $M/IM$ soit plat comme $S/IS$-module, et
\item $M$ soit plat comme $R$-module.
\end{enumerate}
Alors $M$ est un $S$-module plat, et $S_\mathfrak q$ est plat et
essentiellement de présentation finie sur $R$ pour tout
$\mathfrak q \subset S$ tel que
$M \otimes_S \kappa(\mathfrak q)$ soit non nul.
\end{lemma}

\begin{proof}
Si $M \otimes_S \kappa(\mathfrak q)$ est non nul, alors
$S' \otimes_S \kappa(\mathfrak q)$ est non nul ; il existe donc un idéal
premier $\mathfrak q' \subset S'$ au-dessus de $\mathfrak q$
(Lemme \ref{lemma-in-image}). Soit $\mathfrak p \subset R$ l'image de
$\mathfrak q$ dans $\Spec(R)$. Alors $I \subset \mathfrak p$, puisque $I$
est localement nilpotent ; ainsi $M/\mathfrak p M$ est plat sur
$S/\mathfrak pS$. Nous pouvons donc appliquer le
Lemme \ref{lemma-criterion-flatness-fibre-fp-over-ft} à
$R_\mathfrak p \to S_\mathfrak q \to S'_{\mathfrak q'}$ et à
$M_{\mathfrak q'}$. Nous concluons que $M_{\mathfrak q'}$ est plat sur $S$
et que $S_\mathfrak q$ est plat et essentiellement de présentation finie
sur $R$. Comme $\mathfrak q'$ était un idéal premier arbitraire de $S'$,
nous voyons aussi que $M$ est plat sur $S$
(Lemme \ref{lemma-flat-localization}).
\end{proof}














\section{Ouverture du lieu de platitude}
\label{section-open-flat}

\noindent
Nous utilisons le Lemme \ref{lemma-colimit-eventually-flat} pour nous ramener
au cas noethérien. Celui-ci est traité à l'aide de la caractérisation des
complexes exacts donnée dans la section \ref{section-complex-exact}.

\begin{lemma}
\label{lemma-CM-dim-finite-type}
Soit $k$ un corps. Soit $S$ une $k$-algèbre de type fini.
Soient $f_1, \ldots, f_i$ des éléments de $S$.
Supposons que $S$ soit de Cohen-Macaulay, équidimensionnelle de dimension
$d$, et que $\dim V(f_1, \ldots, f_i) \leq d - i$.
Alors il y a égalité, et $f_1, \ldots, f_i$ forment une suite régulière
dans $S_{\mathfrak q}$ pour tout idéal premier $\mathfrak q$ de
$V(f_1, \ldots, f_i)$.
\end{lemma}

\begin{proof}
Si $S$ est de Cohen-Macaulay et équidimensionnelle de dimension $d$, alors
$\dim(S_{\mathfrak m}) = d$ pour tout idéal maximal $\mathfrak m$ de $S$ ;
voir le Lemme \ref{lemma-disjoint-decomposition-CM-algebra}.
D'après la Proposition \ref{proposition-CM-module}, pour tout idéal maximal
$\mathfrak m \in V(f_1, \ldots, f_i)$, la suite est régulière dans
$S_{\mathfrak m}$ et l'anneau local
$S_{\mathfrak m}/(f_1, \ldots, f_i)$ est de Cohen-Macaulay, de dimension
$d - i$. Cela signifie en fait que $S/(f_1, \ldots, f_i)$ est de
Cohen-Macaulay et équidimensionnelle de dimension $d - i$.
\end{proof}

\begin{lemma}
\label{lemma-open-regular-sequence}
Soit $R \to S$ un morphisme d'anneaux de type fini. Soit $d$ un entier tel
que toutes les fibres $S \otimes_R \kappa(\mathfrak p)$ soient de
Cohen-Macaulay et équidimensionnelles de dimension $d$.
Soient $f_1, \ldots, f_i$ des éléments de $S$. L'ensemble
$$
\{ \mathfrak q \in V(f_1, \ldots, f_i)
\mid f_1, \ldots, f_i
\text{ forment une suite régulière dans }
S_{\mathfrak q}/\mathfrak p S_{\mathfrak q}
\text{ où }\mathfrak p = R \cap \mathfrak q
\}
$$
est ouvert dans $V(f_1, \ldots, f_i)$.
\end{lemma}

\begin{proof}
Écrivons $\overline{S} = S/(f_1, \ldots, f_i)$.
Supposons que $\mathfrak q$ appartienne à l'ensemble défini dans le lemme,
et soit $\mathfrak p$ l'idéal premier correspondant de $R$.
Nous utiliserons la dimension relative définie dans la
Définition \ref{definition-relative-dimension}.
Notons d'abord que
$d = \dim_{\mathfrak q}(S/R) =
\dim(S_{\mathfrak q}/\mathfrak pS_{\mathfrak q}) +
\text{trdeg}_{\kappa(\mathfrak p)}\ \kappa(\mathfrak q)$
d'après le Lemme \ref{lemma-dimension-at-a-point-finite-type-field}.
Puisque $f_1, \ldots, f_i$ forment une suite régulière dans l'anneau local
noethérien $S_{\mathfrak q}/\mathfrak pS_{\mathfrak q}$, le
Lemme \ref{lemma-one-equation} donne
$\dim(\overline{S}_{\mathfrak q}/\mathfrak p\overline{S}_{\mathfrak q})
= \dim(S_{\mathfrak q}/\mathfrak pS_{\mathfrak q}) - i$.
Nous en concluons que
$\dim_{\mathfrak q}(\overline{S}/R)
= \dim(\overline{S}_{\mathfrak q}/\mathfrak p\overline{S}_{\mathfrak q}) +
\text{trdeg}_{\kappa(\mathfrak p)}\ \kappa(\mathfrak q)
= d - i$ d'après le
Lemme \ref{lemma-dimension-at-a-point-finite-type-field}.
D'après le Lemme \ref{lemma-dimension-fibres-bounded-open-upstairs}, on a
$\dim_{\mathfrak q'}(\overline{S}/R) \leq d - i$
pour tout $\mathfrak q' \in V(f_1, \ldots, f_i) = \Spec(\overline{S})$
dans un voisinage de $\mathfrak q$. Ainsi, après avoir remplacé $S$ par
$S_g$ pour un certain $g \in S$, $g \not \in \mathfrak q$, nous pouvons
supposer que l'inégalité vaut pour tout $\mathfrak q'$.
Le résultat découle alors du Lemme \ref{lemma-CM-dim-finite-type}.
\end{proof}

\begin{lemma}
\label{lemma-exact-on-fibres-open}
Soit $R \to S$ un morphisme d'anneaux. Considérons un complexe homologique
fini de $S$-modules libres de type fini :
$$
F_{\bullet} :
0
\to
S^{n_e}
\xrightarrow{\varphi_e}
S^{n_{e-1}}
\xrightarrow{\varphi_{e-1}}
\ldots
\xrightarrow{\varphi_{i + 1}}
S^{n_i}
\xrightarrow{\varphi_i}
S^{n_{i-1}}
\xrightarrow{\varphi_{i-1}}
\ldots
\xrightarrow{\varphi_1}
S^{n_0}
$$
Pour tout idéal premier $\mathfrak q$ de $S$, considérons le complexe
$\overline{F}_{\bullet, \mathfrak q} =
F_{\bullet, \mathfrak q} \otimes_R \kappa(\mathfrak p)$,
où $\mathfrak p$ est l'image réciproque de $\mathfrak q$ dans $R$.
Supposons $R$ noethérien et supposons qu'il existe un entier $d$ tel que
$R \to S$ soit de type fini et plat, à fibres
$S \otimes_R \kappa(\mathfrak p)$ de Cohen-Macaulay et de dimension $d$.
L'ensemble
$$
\{\mathfrak q \in \Spec(S) \mid
\overline{F}_{\bullet, \mathfrak q}\text{ est exact}\}
$$
est ouvert dans $\Spec(S)$.
\end{lemma}

\begin{proof}
Soit $\mathfrak q$ un élément de l'ensemble défini dans le lemme.
Nous allons utiliser la Proposition \ref{proposition-what-exact} pour montrer
qu'il existe $g \in S$, $g \not \in \mathfrak q$, tel que $D(g)$ soit
contenu dans l'ensemble défini dans le lemme. Autrement dit, nous allons
montrer qu'après avoir remplacé $S$ par $S_g$, l'ensemble du lemme est
$\Spec(S)$ tout entier. Au cours de la démonstration, nous remplacerons donc
un nombre fini de fois $S$ par une telle localisation.
Rappelons que la Proposition \ref{proposition-what-exact} caractérise
l'exactitude des complexes, lorsque l'anneau est local, en fonction des rangs
des morphismes $\varphi_i$ et des idéaux $I(\varphi_i)$.
Traitons d'abord la condition sur les rangs. Posons
$r_i = n_i - n_{i + 1} + \ldots + (-1)^{e - i} n_e$.
Notons que $r_i + r_{i + 1} = n_i$ et que $r_i$ est le rang attendu de
$\varphi_i$ dans le cas exact.

\medskip\noindent
D'après le Lemme \ref{lemma-complex-exact-mod}, si
$\overline{F}_{\bullet, \mathfrak q}$ est exact, alors le complexe localisé
$F_{\bullet, \mathfrak q}$ est exact. En particulier, le complexe
$F_\bullet$ devient exact après localisation par un élément
$g \in S$, $g \not \in \mathfrak q$. Dans ce cas, la
Proposition \ref{proposition-what-exact}, appliquée à toutes les
localisations de $S$ en des idéaux premiers, implique que tous les mineurs
d'ordre $(r_i + 1) \times (r_i + 1)$ de $\varphi_i$ sont nuls.
On voit donc que le rang de $\varphi_i$ est au plus $r_i$.

\medskip\noindent
Notons $I_i \subset S$ l'idéal engendré par les mineurs d'ordre
$r_i \times r_i$ de la matrice de $\varphi_i$.
D'après la Proposition \ref{proposition-what-exact}, le complexe
$\overline{F}_{\bullet, \mathfrak q}$ est exact si et seulement si, pour
tout $1 \leq i \leq e$, on a soit
$(I_i)_{\mathfrak q} = S_{\mathfrak q}$, soit que
$(I_i)_{\mathfrak q}$ contient une suite
$S_{\mathfrak q}/\mathfrak p S_{\mathfrak q}$-régulière de longueur $i$.
En effet, par le choix de $r_i$ ci-dessus et la borne sur les rangs des
$\varphi_i$, c'est la seule façon dont les conditions de la
Proposition \ref{proposition-what-exact} peuvent être satisfaites.

\medskip\noindent
Si $(I_i)_{\mathfrak q} = S_{\mathfrak q}$, après avoir localisé $S$ en un
élément $g \not\in \mathfrak q$, nous pouvons supposer que $I_i = S$.
C'est manifestement une condition ouverte.

\medskip\noindent
Si $(I_i)_{\mathfrak q} \not = S_{\mathfrak q}$, il existe une suite
$f_1, \ldots, f_i \in (I_i)_{\mathfrak q}$ qui forme une suite régulière
dans $S_{\mathfrak q}/\mathfrak pS_{\mathfrak q}$.
Notons que, pour tout idéal premier $\mathfrak q' \subset S$ tel que
$(f_1, \ldots, f_i) \not \subset \mathfrak q'$, on a
$(I_i)_{\mathfrak q'} = S_{\mathfrak q'}$.
Le résultat découle donc du Lemme \ref{lemma-open-regular-sequence}.
\end{proof}


\begin{theorem}
\label{theorem-openness-flatness}
Soit $R$ un anneau. Soit $R \to S$ un morphisme d'anneaux de présentation
finie. Soit $M$ un $S$-module de présentation finie. L'ensemble
$$
\{ \mathfrak q \in \Spec(S) \mid
M_{\mathfrak q}\text{ est plat sur }R\}
$$
est ouvert dans $\Spec(S)$.
\end{theorem}

\begin{proof}
Soit $\mathfrak q \in \Spec(S)$ un idéal premier.
Soit $\mathfrak p \subset R$ l'image réciproque de $\mathfrak q$ dans $R$.
Notons que $M_{\mathfrak q}$ est plat sur $R$ si et seulement s'il est plat
sur $R_{\mathfrak p}$. Supposons que $M_{\mathfrak q}$ soit plat sur $R$.
Nous affirmons qu'il existe $g \in S$, $g \not \in \mathfrak q$, tel que
$M_g$ soit plat sur $R$.

\medskip\noindent
Nous nous ramenons d'abord au cas où $R$ et $S$ sont de type fini sur
$\mathbf{Z}$. Choisissons un ensemble filtrant $\Lambda$ et un système
$(R_\lambda \to S_\lambda, M_\lambda)$ comme dans le
Lemme \ref{lemma-limit-module-finite-presentation}.
Posons $\mathfrak p_\lambda$ égal à l'image réciproque de $\mathfrak p$
dans $R_\lambda$, et $\mathfrak q_\lambda$ égal à l'image réciproque de
$\mathfrak q$ dans $S_\lambda$. Le système
$$
((R_\lambda)_{\mathfrak p_\lambda},
(S_\lambda)_{\mathfrak q_\lambda},
(M_\lambda)_{\mathfrak q_{\lambda}})
$$
est alors un système comme dans le
Lemme \ref{lemma-limit-module-essentially-finite-presentation}.
Par conséquent, le Lemme \ref{lemma-colimit-eventually-flat} montre que,
pour un certain $\lambda$, le module $M_\lambda$ est plat sur $R_\lambda$
en l'idéal premier $\mathfrak q_{\lambda}$.
Supposons que nous sachions démontrer notre assertion pour le système
$(R_\lambda \to S_\lambda, M_\lambda, \mathfrak q_{\lambda})$.
Autrement dit, supposons que nous puissions trouver $g \in S_\lambda$,
$g \not\in \mathfrak q_\lambda$, tel que $(M_\lambda)_g$ soit plat sur
$R_\lambda$. D'après le Lemme \ref{lemma-limit-module-finite-presentation},
on a $M = M_\lambda \otimes_{R_\lambda} R$, et donc aussi
$M_g = (M_\lambda)_g \otimes_{R_\lambda} R$.
Le Lemme \ref{lemma-flat-base-change} donne alors l'assertion pour le
système $(R \to S, M, \mathfrak q)$.

\medskip\noindent
Nous pouvons désormais supposer que $R$ et $S$ sont de type fini sur
$\mathbf{Z}$. Nous pouvons écrire $S$ comme quotient d'un anneau de
polynômes $R[x_1, \ldots, x_n]$. Bien entendu, nous pouvons remplacer $S$
par $R[x_1, \ldots, x_n]$ et supposer que $S$ est un anneau de polynômes
sur $R$. En particulier, $R \to S$ est plat et tous les anneaux des fibres
$S \otimes_R \kappa(\mathfrak p)$ sont de dimension globale $n$.

\medskip\noindent
Choisissons une résolution $F_\bullet$ de $M$ sur $S$, chaque $F_i$ étant
libre de type fini ; voir le Lemme \ref{lemma-resolution-by-finite-free}.
Soit $K_n = \Ker(F_{n-1} \to F_{n-2})$. Notons que
$(K_n)_{\mathfrak q}$ est plat sur $R$, puisque chaque $F_i$ est plat sur
$R$ et en vertu de l'hypothèse sur $M$ ; voir le
Lemme \ref{lemma-flat-ses}. En outre, la suite
$$
0 \to
K_n/\mathfrak p K_n \to
F_{n-1}/ \mathfrak p F_{n-1} \to
\ldots \to
F_0 / \mathfrak p F_0 \to
M/\mathfrak p M \to
0
$$
est exacte après localisation en $\mathfrak q$, en raison de l'annulation
de $\text{Tor}_i^{R_\mathfrak p}(\kappa(\mathfrak p), M_{\mathfrak q})$.
Comme la dimension globale de
$S_\mathfrak q/\mathfrak p S_{\mathfrak q}$ est $n$, nous concluons que
$K_n / \mathfrak p K_n$, localisé en $\mathfrak q$, est un module libre
de type fini sur $S_\mathfrak q/\mathfrak p S_{\mathfrak q}$.
D'après le Lemme \ref{lemma-free-fibre-flat-free},
$(K_n)_{\mathfrak q}$ est libre sur $S_{\mathfrak q}$.
En particulier, il existe $g \in S$, $g \not \in \mathfrak q$, tel que
$(K_n)_g$ soit libre de type fini sur $S_g$.

\medskip\noindent
D'après le Lemme \ref{lemma-exact-on-fibres-open}, il existe une localisation
supplémentaire $S_g$ telle que le complexe
$$
0 \to K_n \to F_{n-1} \to \ldots \to F_0
$$
soit exact sur {\it toutes les fibres} de $R \to S$.
D'après le Lemme \ref{lemma-complex-exact-mod}, cela implique que le conoyau
de $F_1 \to F_0$ est plat. Cela démontre le théorème dans le cas noethérien.
\end{proof}


\section{Ouverture des lieux de Cohen-Macaulay}
\label{section-CM-open}

\noindent
Dans cette section, nous caractérisons la propriété de Cohen-Macaulay des
algèbres de type fini en termes de platitude. Nous utilisons ensuite cette
caractérisation pour montrer que l'ensemble des points où une telle algèbre
est de Cohen-Macaulay est ouvert.

\begin{lemma}
\label{lemma-where-CM}
Soit $S$ une algèbre de type fini sur un corps $k$.
Soit $\varphi : k[y_1, \ldots, y_d] \to S$ un morphisme d'anneaux
quasi-fini. Comme sous-ensembles de $\Spec(S)$, on a
$$
\{ \mathfrak q \mid
S_{\mathfrak q} \text{ plat sur }k[y_1, \ldots, y_d]\}
=
\{ \mathfrak q \mid
S_{\mathfrak q} \text{ CM et }\dim_{\mathfrak q}(S/k) = d\}
$$
Pour la notation, voir la Définition \ref{definition-relative-dimension}.
\end{lemma}

\begin{proof}
Soit $\mathfrak q \subset S$ un idéal premier. Notons
$\mathfrak p = k[y_1, \ldots, y_d] \cap \mathfrak q$.
Notons que l'on a toujours
$\dim(S_{\mathfrak q}) \leq \dim(k[y_1, \ldots, y_d]_{\mathfrak p})$,
par exemple d'après le Lemme \ref{lemma-dimension-inequality-quasi-finite}.
En outre, l'extension de corps
$\kappa(\mathfrak q)/\kappa(\mathfrak p)$ est finie et donc
$\text{trdeg}_k(\kappa(\mathfrak p)) =
\text{trdeg}_k(\kappa(\mathfrak q))$.

\medskip\noindent
Soit $\mathfrak q$ un élément du membre de gauche.
Le Lemme \ref{lemma-finite-flat-over-regular-CM} s'applique alors, et nous
concluons que $S_{\mathfrak q}$ est de Cohen-Macaulay et que
$\dim(S_{\mathfrak q}) = \dim(k[y_1, \ldots, y_d]_{\mathfrak p})$.
Combiné à l'égalité des degrés de transcendance ci-dessus et au
Lemme \ref{lemma-dimension-at-a-point-finite-type-field}, cela implique que
$\dim_{\mathfrak q}(S/k) = d$. Ainsi, $\mathfrak q$ appartient au membre
de droite.

\medskip\noindent
Soit $\mathfrak q$ un élément du membre de droite.
L'égalité des degrés de transcendance ci-dessus, l'hypothèse
$\dim_{\mathfrak q}(S/k) = d$ et le
Lemme \ref{lemma-dimension-at-a-point-finite-type-field} donnent
$\dim(S_{\mathfrak q}) = \dim(k[y_1, \ldots, y_d]_{\mathfrak p})$.
Le Lemme \ref{lemma-CM-over-regular-flat} s'applique donc, et nous voyons que
$\mathfrak q$ appartient au membre de gauche.
\end{proof}

\begin{lemma}
\label{lemma-finite-type-over-field-CM-open}
Soit $S$ une algèbre de type fini sur un corps $k$.
L'ensemble des idéaux premiers $\mathfrak q$ tels que $S_{\mathfrak q}$
soit de Cohen-Macaulay est ouvert dans $S$.
\end{lemma}

\noindent
Ce lemme est un cas particulier du
Lemme \ref{lemma-finite-presentation-flat-CM-locus-open} ci-dessous ; on
peut donc, si on le souhaite, passer directement à la démonstration de ce
dernier.

\begin{proof}
Soit $\mathfrak q \subset S$ un idéal premier tel que $S_{\mathfrak q}$
soit de Cohen-Macaulay. Il faut montrer qu'il existe
$g \in S$, $g \not \in \mathfrak q$, tel que l'anneau $S_g$ soit de
Cohen-Macaulay. Pour tout $g \in S$, $g \not \in \mathfrak q$, nous pouvons
remplacer $S$ par $S_g$ et $\mathfrak q$ par $\mathfrak qS_g$.
En combinant cette observation aux
Lemmes \ref{lemma-Noether-normalization-at-point} et
\ref{lemma-dimension-at-a-point-finite-type-field}, nous pouvons supposer
qu'il existe un morphisme d'anneaux fini injectif
$k[y_1, \ldots, y_d] \to S$ tel que
$d = \dim(S_{\mathfrak q}) + \text{trdeg}_k(\kappa(\mathfrak q))$.
Posons $\mathfrak p = k[y_1, \ldots, y_d] \cap \mathfrak q$.
Par construction, $\mathfrak q$ appartient au membre de droite de l'égalité
affichée du Lemme \ref{lemma-where-CM}. Il appartient donc aussi au membre
de gauche.

\medskip\noindent
D'après le Théorème \ref{theorem-openness-flatness}, il existe
$g \in S$, $g \not \in \mathfrak q$, tel que l'anneau $S_g$ soit plat sur
$k[y_1, \ldots, y_d]$. En appliquant de nouveau l'égalité du
Lemme \ref{lemma-where-CM}, nous concluons que tous les anneaux locaux de
$S_g$ sont de Cohen-Macaulay, comme souhaité.
\end{proof}

\begin{lemma}
\label{lemma-generic-CM}
Soit $k$ un corps. Soit $S$ une $k$-algèbre de type fini.
L'ensemble des idéaux premiers de Cohen-Macaulay forme un ouvert dense
$U \subset \Spec(S)$.
\end{lemma}

\begin{proof}
L'ensemble est ouvert d'après le
Lemme \ref{lemma-finite-type-over-field-CM-open}.
Il contient tous les idéaux premiers minimaux $\mathfrak q \subset S$,
puisque l'anneau local $S_{\mathfrak q}$ en un idéal premier minimal est de
dimension zéro, et donc de Cohen-Macaulay.
\end{proof}

\begin{lemma}
\label{lemma-finite-presentation-flat-CM-locus-open}
Soit $R$ un anneau. Soit $R \to S$ de présentation finie et plat.
Pour tout $d \geq 0$, l'ensemble
$$
\left\{
\begin{matrix}
\mathfrak q \in \Spec(S)
\text{ tel que, en posant }\mathfrak p = R \cap \mathfrak q
\text{ l'anneau de la fibre}\\
S_{\mathfrak q}/\mathfrak pS_{\mathfrak q}
\text{ est de Cohen-Macaulay}
\text{ et } \dim_{\mathfrak q}(S/R) = d
\end{matrix}
\right\}
$$
est ouvert dans $\Spec(S)$.
\end{lemma}

\begin{proof}
Soit $\mathfrak q$ un élément de l'ensemble indiqué, et soit $\mathfrak p$
l'idéal premier correspondant de $R$. Il faut trouver
$g \in S$, $g \not \in \mathfrak q$, tel que tous les anneaux des fibres
de $R \to S_g$ soient de Cohen-Macaulay.
Au cours de la démonstration, nous pourrons remplacer, un nombre fini de
fois, $S$ par $S_g$ pour un $g \in S$, $g \not \in \mathfrak q$.
Ainsi, d'après le
Lemme \ref{lemma-quasi-finite-over-polynomial-algebra}, nous pouvons supposer
qu'il existe un morphisme d'anneaux quasi-fini
$R[t_1, \ldots, t_d] \to S$ avec
$d = \dim_{\mathfrak q}(S/R)$.
Soit $\mathfrak q' = R[t_1, \ldots, t_d] \cap \mathfrak q$.
D'après le Lemme \ref{lemma-where-CM}, le morphisme d'anneaux
$$
R[t_1, \ldots, t_d]_{\mathfrak q'} /
\mathfrak p R[t_1, \ldots, t_d]_{\mathfrak q'}
\longrightarrow
S_{\mathfrak q}/\mathfrak p S_{\mathfrak q}
$$
est plat. Par conséquent, le crit\`ere de platitude par fibres,
Lemme \ref{lemma-criterion-flatness-fibre}, montre que
$R[t_1, \ldots, t_d]_{\mathfrak q'} \to S_{\mathfrak q}$ est plat.
D'après le Théorème \ref{theorem-openness-flatness}, il existe donc
$g \in S$, $g \not \in \mathfrak q$, tel que le morphisme d'anneaux
$R[t_1, \ldots, t_d] \to S_g$ soit plat. En remplaçant $S$ par $S_g$, nous
voyons que, pour tout idéal premier $\mathfrak r \subset S$, en posant
$\mathfrak r' = R[t_1, \ldots, t_d] \cap \mathfrak r$ et
$\mathfrak p' = R \cap \mathfrak r$, l'homomorphisme local
$R[t_1, \ldots, t_d]_{\mathfrak r'} \to S_{\mathfrak r}$ est plat.
Il en est donc de même du changement de base
$$
R[t_1, \ldots, t_d]_{\mathfrak r'} /
\mathfrak p' R[t_1, \ldots, t_d]_{\mathfrak r'}
\longrightarrow
S_{\mathfrak r}/\mathfrak p' S_{\mathfrak r}
$$
qui est plat. Le Lemme \ref{lemma-where-CM}, appliqué avec
$k = \kappa(\mathfrak p')$, montre donc que $\mathfrak r$ appartient à
l'ensemble du lemme, comme souhaité.
\end{proof}

\begin{lemma}
\label{lemma-generic-CM-flat-finite-presentation}
Soit $R$ un anneau. Soit $R \to S$ plat et de présentation finie.
L'ensemble des idéaux premiers $\mathfrak q$ tels que l'anneau de la fibre
$S_{\mathfrak q} \otimes_R \kappa(\mathfrak p)$, où
$\mathfrak p = R \cap \mathfrak q$, soit de Cohen-Macaulay est ouvert et
dense dans chaque fibre de $\Spec(S) \to \Spec(R)$.
\end{lemma}

\begin{proof}
Cet ensemble, noté $W$, est ouvert d'après le
Lemme \ref{lemma-finite-presentation-flat-CM-locus-open}.
Il est dense dans les fibres, car l'intersection de $W$ avec une fibre est
l'ensemble correspondant de cette fibre, auquel s'applique le
Lemme \ref{lemma-generic-CM}.
\end{proof}

\begin{lemma}
\label{lemma-extend-field-CM-locus}
Soit $k$ un corps. Soit $S$ une $k$-algèbre de type fini.
Soit $K/k$ une extension de corps, et posons $S_K = K \otimes_k S$.
Soit $\mathfrak q \subset S$ un idéal premier de $S$.
Soit $\mathfrak q_K \subset S_K$ un idéal premier de $S_K$ au-dessus de
$\mathfrak q$. Alors $S_{\mathfrak q}$ est de Cohen-Macaulay si et
seulement si $(S_K)_{\mathfrak q_K}$ est de Cohen-Macaulay.
\end{lemma}

\begin{proof}
Au cours de la démonstration, nous pouvons remplacer, un nombre fini de fois,
$S$ par $S_g$ pour tout $g \in S$, $g \not \in \mathfrak q$.
Par conséquent, en utilisant le
Lemme \ref{lemma-Noether-normalization-at-point}, nous pouvons supposer que
$\dim(S) = \dim_{\mathfrak q}(S/k) =: d$ et trouver un morphisme fini
injectif $k[x_1, \ldots, x_d] \to S$.
Notons que celui-ci induit aussi, par changement de base, un morphisme fini
injectif $K[x_1, \ldots, x_d] \to S_K$.
D'après le Lemme \ref{lemma-dimension-at-a-point-preserved-field-extension},
on a $\dim_{\mathfrak q_K}(S_K/K) = d$.
Posons $\mathfrak p = k[x_1, \ldots, x_d] \cap \mathfrak q$ et
$\mathfrak p_K = K[x_1, \ldots, x_d] \cap \mathfrak q_K$.
Considérons le diagramme commutatif suivant d'anneaux locaux noethériens :
$$
\xymatrix{
S_{\mathfrak q} \ar[r] &
(S_K)_{\mathfrak q_K} \\
k[x_1, \ldots, x_d]_{\mathfrak p} \ar[r] \ar[u] &
K[x_1, \ldots, x_d]_{\mathfrak p_K} \ar[u]
}
$$
D'après le Lemme \ref{lemma-where-CM}, il faut montrer que la flèche
verticale de gauche est plate si et seulement si celle de droite l'est.
Comme la flèche du bas est plate, cette équivalence résulte du
Lemme \ref{lemma-base-change-flat-up-down}.
\end{proof}

\begin{lemma}
\label{lemma-CM-locus-commutes-base-change}
Soit $R$ un anneau. Soit $R \to S$ de type fini.
Soit $R \to R'$ un morphisme d'anneaux. Posons $S' = R' \otimes_R S$.
Notons $f : \Spec(S') \to \Spec(S)$ l'application associée au morphisme
d'anneaux $S \to S'$. Soit $W$ l'ensemble des idéaux premiers
$\mathfrak q$ tels que l'anneau de la fibre
$S_{\mathfrak q} \otimes_R \kappa(\mathfrak p)$,
$\mathfrak p = R \cap \mathfrak q$, soit de Cohen-Macaulay, et notons
$W'$ l'ensemble analogue pour $S'/R'$. Alors $W' = f^{-1}(W)$.
\end{lemma}

\begin{proof}
Cela résulte immédiatement du Lemme \ref{lemma-extend-field-CM-locus} et
des définitions.
\end{proof}

\begin{lemma}
\label{lemma-relative-dimension-CM}
Soit $R$ un anneau. Soit $R \to S$ un morphisme d'anneaux qui est (a) plat,
(b) de présentation finie et (c) à fibres de Cohen-Macaulay.
Alors on peut écrire $S = S_0 \times \ldots \times S_n$ comme un produit de
$R$-algèbres $S_d$ de sorte que chaque $S_d$ satisfasse (a), (b) et (c) et
ait toutes ses fibres équidimensionnelles de dimension $d$.
\end{lemma}

\begin{proof}
Pour chaque entier $d$, notons $W_d \subset \Spec(S)$ l'ensemble défini
dans le Lemme \ref{lemma-finite-presentation-flat-CM-locus-open}.
On a manifestement $\Spec(S) = \coprod W_d$, et chaque $W_d$ est ouvert
d'après le lemme que nous venons de citer. Le résultat découle donc du
Lemme \ref{lemma-disjoint-implies-product}.
\end{proof}


\section{Différentielles}
\label{section-differentials}

\noindent
Dans cette section, nous définissons le module des différentielles d'un
morphisme d'anneaux.

\begin{definition}
\label{definition-derivation}
Soit $\varphi : R \to S$ un morphisme d'anneaux et soit $M$ un $S$-module.
Une {\it dérivation}, ou plus précisément une
{\it $R$-dérivation} à valeurs dans $M$, est une application $D : S \to M$
qui est additive, s'annule sur les éléments de $\varphi(R)$ et satisfait
la {\it règle de Leibniz} : $D(ab) = aD(b) + bD(a)$.
\end{definition}

\noindent
Notons que $D(ra) = rD(a)$ si $r \in R$ et $a \in S$. Une définition
équivalente est qu'une $R$-dérivation est une application $R$-linéaire
$D : S \to M$ qui satisfait la règle de Leibniz.
L'ensemble de toutes les $R$-dérivations forme un $S$-module : étant données
deux $R$-dérivations $D, D'$, leur somme $D + D' : S \to M$,
$a \mapsto D(a)+D'(a)$, est une $R$-dérivation, et, étant donnés une
$R$-dérivation $D$ et un élément $c\in S$, le multiple scalaire
$cD : S \to M$, $a \mapsto cD(a)$, est une $R$-dérivation. Nous notons ce
$S$-module
$$
\text{Der}_R(S, M).
$$
En outre, si $\alpha : M \to N$ est un morphisme de $S$-modules, alors la
composée $\alpha \circ D$ est une $R$-dérivation à valeurs dans $N$.
Ainsi, l'association $M \mapsto \text{Der}_R(S, M)$ est un foncteur
covariant.

\medskip\noindent
Considérons le morphisme suivant de $S$-modules libres :
$$
\bigoplus\nolimits_{(a, b)\in S^2} S[(a, b)] \oplus
\bigoplus\nolimits_{(f, g)\in S^2} S[(f, g)] \oplus
\bigoplus\nolimits_{r\in R} S[r]
\longrightarrow
\bigoplus\nolimits_{a\in S} S[a]
$$
défini par les règles
$$
[(a, b)] \longmapsto [a + b] - [a] - [b],\quad
[(f, g)] \longmapsto [fg] -f[g] - g[f],\quad
[r]      \longmapsto [\varphi(r)]
$$
avec les notations évidentes. Soit $\Omega_{S/R}$ le conoyau de ce morphisme.
Il existe une application $\text{d} : S \to \Omega_{S/R}$ qui envoie $a$
sur la classe $\text{d}a$ de $[a]$ dans le conoyau. C'est une $R$-dérivation
en vertu des relations imposées à $\Omega_{S/R}$ ; autrement dit,
$$
\text{d}(a + b) = \text{d}a + \text{d}b, \quad
\text{d}(fg) = f\text{d}g + g\text{d}f, \quad
\text{d}\varphi(r) = 0
$$
où $a,b,f,g \in S$ et $r \in R$.

\begin{definition}
\label{definition-differentials}
Le couple $(\Omega_{S/R}, \text{d})$ est appelé le {\it module des
différentielles de K\"ahler}, ou le {\it module des différentielles}
de $S$ sur $R$.
\end{definition}

\begin{lemma}
\label{lemma-universal-omega}
\begin{slogan}
Les morphismes partant du module des différentielles sont les mêmes que
les dérivations.
\end{slogan}
Le module des différentielles de $S$ sur $R$ possède la propriété
universelle suivante. L'application
$$
\Hom_S(\Omega_{S/R}, M)
\longrightarrow
\text{Der}_R(S, M), \quad
\alpha
\longmapsto
\alpha \circ \text{d}
$$
est un isomorphisme de foncteurs.
\end{lemma}

\begin{proof}
Par définition, une $R$-dérivation est une règle qui associe à chaque
$a \in S$ un élément $D(a) \in M$. Ainsi, $D$ donne naissance à un
morphisme $[D] : \bigoplus S[a] \to M$. Or les conditions requises pour
être une $R$-dérivation signifient exactement que $[D]$ annule l'image du
morphisme dans la présentation de $\Omega_{S/R}$ affichée ci-dessus.
\end{proof}

\begin{lemma}
\label{lemma-trivial-differential-surjective}
Supposons que $R \to S$ soit surjectif.
Alors $\Omega_{S/R} = 0$.
\end{lemma}

\begin{proof}
On peut le voir soit parce que toutes les $R$-dérivations doivent
manifestement être nulles, soit parce que le morphisme de la présentation
de $\Omega_{S/R}$ est surjectif.
\end{proof}

\noindent
Supposons que
\begin{equation}
\label{equation-functorial-omega}
\vcenter{
\xymatrix{
S \ar[r]_\varphi
&
S'
\\
R \ar[r]^\psi \ar[u]^\alpha
&
R' \ar[u]_\beta
}
}
\end{equation}
soit un diagramme commutatif d'anneaux. Dans ce cas, il existe un morphisme
naturel entre les modules des différentielles qui s'insère dans le diagramme
commutatif
$$
\xymatrix{
\Omega_{S/R} \ar[r] &
\Omega_{S'/R'}
\\
S \ar[u]^{\text{d}} \ar[r]^{\varphi}
&
S' \ar[u]_{\text{d}}
}
$$
Pour construire ce morphisme, il suffit d'utiliser le morphisme évident
entre les présentations de $\Omega_{S/R}$ et $\Omega_{S'/R'}$, à savoir
\begin{equation}
\label{equation-map-presentations}
\vcenter{
\xymatrix{
\bigoplus S'[(a', b')]
\oplus
\bigoplus S'[(f', g')]
\oplus
\bigoplus S'[r'] \ar[r]
&
\bigoplus S' [a'] \\
 \\
\bigoplus S[(a, b)]
\oplus
\bigoplus S[(f, g)]
\oplus
\bigoplus S[r] \ar[r]
\ar[uu]^{
\begin{matrix}
[(a, b)] \mapsto [(\varphi(a), \varphi(b))] \\
[(f, g)] \mapsto [(\varphi(f), \varphi(g))] \\
[r]\mapsto [\psi(r)]
\end{matrix}
} &
\bigoplus S[a] \ar[uu]_{[a] \mapsto [\varphi(a)]}
}
}
\end{equation}
Le résultat est simplement que $f\text{d}g \in \Omega_{S/R}$ est envoyé sur
$\varphi(f)\text{d}\varphi(g)$.


\begin{lemma}
\label{lemma-colimit-differentials}
Soit $I$ un ensemble filtrant.
Soit $(R_i \to S_i, \varphi_{ii'})$ un système de morphismes d'anneaux
indexé par $I$ ; voir
Catégories, section \ref{categories-section-posets-limits}.
Alors on a
$$
\Omega_{S/R} =
\colim_i \Omega_{S_i/R_i}.
$$
où $R \to S = \colim (R_i \to S_i)$.
\end{lemma}

\begin{proof}
Cela résulte clairement de la présentation définissant $\Omega_{S/R}$ et
de la fonctorialité décrite ci-dessus.
\end{proof}

\begin{lemma}
\label{lemma-differential-surjective}
Dans le diagramme (\ref{equation-functorial-omega}), supposons que
$S \to S'$ soit surjectif de noyau $I \subset S$.
Alors $\Omega_{S/R} \to \Omega_{S'/R'}$ est surjectif et son noyau est
engendré comme $S$-module par les éléments $\text{d}a$, où $a \in S$ est
tel que $\varphi(a) \in \beta(R')$.
(Cela comprend en particulier les éléments $\text{d}(i)$, $i \in I$.)
\end{lemma}

\begin{proof}[Première démonstration]
Considérons le morphisme de présentations
(\ref{equation-map-presentations}). Le morphisme vertical de droite entre
modules libres est manifestement surjectif. Le morphisme est donc surjectif.
Supposons qu'un élément $\eta$ de $\Omega_{S/R}$ soit envoyé sur zéro dans
$\Omega_{S'/R'}$. Écrivons $\eta$ comme l'image de $\sum s_i[a_i]$ pour des
$s_i, a_i \in S$. Nous voyons alors que
$\sum \varphi(s_i)[\varphi(a_i)]$ est l'image d'un élément
$$
\theta = \sum s_j'[a_j', b_j'] + \sum s_k'[f_k', g_k'] + \sum s_l'[r_l']
$$
du coin supérieur gauche du diagramme.
Puisque $\varphi$ est surjectif, les termes $s_j'[a_j', b_j']$ et
$s_k'[f_k', g_k']$ sont dans l'image d'éléments du coin inférieur droit.
Ainsi, en modifiant $\eta$ et $\theta$ par soustraction des images de ces
éléments, nous pouvons supposer que $\theta = \sum s_l'[r_l']$.
Autrement dit, nous voyons que $\sum \varphi(s_i)[\varphi(a_i)]$ est de la
forme $\sum s'_l [\beta(r'_l)]$.
Ensuite, nous pouvons supposer qu'il existe $a' \in S'$ tel que
$a' = \varphi(a_i)$ pour tout $i$ et $a' = \beta(r_l')$ pour tout $l$.
Cela résulte clairement de la décomposition en somme directe du coin
supérieur droit du diagramme. Choisissons $a \in S$ tel que
$\varphi(a) = a'$. Nous pouvons alors écrire $a_i = a + x_i$ pour des
$x_i \in I$. En utilisant les relations permises, nous pouvons donc
supposer que tous les $a_i$ sont égaux à $a$. Mais nous pouvons alors
supposer que notre élément est de la forme $s[a]$. Nous savons encore que
$\varphi(s)[a'] = \sum \varphi(s_l')[\beta(r_l')]$.
Ainsi, soit $\varphi(s) = 0$ et nous avons terminé, soit
$a' = \varphi(a)$ appartient à l'image de $\beta$ et nous avons également
terminé.
\end{proof}

\begin{proof}[Seconde démonstration]
Nous utiliserons sans autre mention la propriété universelle des modules
des différentielles donnée dans le Lemme \ref{lemma-universal-omega}.

\medskip\noindent
Dans (\ref{equation-functorial-omega}), posons $R'' = S \times_{S'} R'$.
Nous obtenons alors le diagramme suivant :
$$
\xymatrix{
S \ar[r] & S \ar[r] & S' \\
R \ar[r] \ar[u] & R'' \ar[r] \ar[u] & R' \ar[u]
}
$$
Soit $M$ un $S$-module. Il résulte immédiatement des définitions qu'une
$R$-dérivation $D : S \to M$ est une $R''$-dérivation si et seulement si
elle annule les éléments de l'image de $R'' \to S$. La propriété universelle
traduit cela en l'énoncé suivant : le morphisme naturel
$\Omega_{S/R} \to \Omega_{S/R''}$ est surjectif et son noyau est engendré
comme $S$-module par l'image de $R''$.

\medskip\noindent
D'après le paragraphe précédent, il suffit de montrer que
$\Omega_{S/R} \to \Omega_{S'/R'}$ est un isomorphisme lorsque $S \to S'$
est surjectif et $R = S \times_{S'} R'$. Soit $M'$ un $S'$-module.
Observons que toute $R'$-dérivation $D' : S' \to M'$ donne une
$R$-dérivation par précomposition avec $S \to S'$. Réciproquement,
supposons que $M$ soit un $S$-module et que $D : S \to M$ soit une
$R$-dérivation. Si $i \in I$, alors il existe $a \in R$ tel que
$\alpha(a) = i$ (puisque $R = S \times_{S'} R'$). Il s'ensuit que
$D(i) = 0$, et donc que $0 = D(is) = iD(s)$ pour tout $s \in S$. Ainsi,
l'image de $D$ est contenue dans le sous-module $M' \subset M$ des éléments
annulés par $I$ et, de plus, l'application induite $S \to M'$ se factorise
par une $R'$-dérivation $S' \to M'$. C'est un exercice d'utiliser la
propriété universelle pour voir que cela signifie que
$\Omega_{S/R} \to \Omega_{S'/R'}$ est un isomorphisme ; nous omettons les
détails.
\end{proof}

\begin{lemma}
\label{lemma-exact-sequence-differentials}
Soient $A \to B \to C$ des morphismes d'anneaux.
Il existe alors une suite exacte canonique
$$
C \otimes_B \Omega_{B/A} \to
\Omega_{C/A} \to
\Omega_{C/B} \to 0
$$
de $C$-modules.
\end{lemma}

\begin{proof}
Nous obtenons un diagramme (\ref{equation-functorial-omega}) en posant
$R = A$, $S = C$, $R' = B$ et $S' = C$.
D'après le Lemme \ref{lemma-differential-surjective}, le morphisme
$\Omega_{C/A} \to \Omega_{C/B}$ est surjectif, et son noyau est engendré
par les éléments $\text{d}(c)$, où $c \in C$ appartient à l'image de
$B \to C$. Le lemme en résulte.
\end{proof}

\begin{lemma}
\label{lemma-differentials-localize}
Soit $\varphi : A \to B$ un morphisme d'anneaux.
\begin{enumerate}
\item Si $S \subset A$ est une partie multiplicative dont l'image est
formée d'éléments inversibles de $B$, alors
$\Omega_{B/A} = \Omega_{B/S^{-1}A}$.
\item Si $S \subset B$ est une partie multiplicative, alors
$S^{-1}\Omega_{B/A} = \Omega_{S^{-1}B/A}$.
\end{enumerate}
\end{lemma}

\begin{proof}
Pour démontrer l'égalité de (1), il suffit de montrer que toute
$A$-dérivation $D : B \to M$ annule les éléments $\varphi(s)^{-1}$.
Cela résulte clairement de la règle de Leibniz appliquée à
$1 = \varphi(s) \varphi(s)^{-1}$.
Pour montrer (2), notons qu'il existe un morphisme évident
$S^{-1}\Omega_{B/A} \to \Omega_{S^{-1}B/A}$.
Pour montrer qu'il s'agit d'un isomorphisme, il suffit de montrer qu'il
existe une $A$-dérivation $\text{d}'$ de $S^{-1}B$ à valeurs dans
$S^{-1}\Omega_{B/A}$. Pour la définir, posons simplement
$\text{d}'(b/s) = (1/s)\text{d}b - (1/s^2)b\text{d}s$.
Nous omettons les détails.
\end{proof}

\begin{lemma}
\label{lemma-differential-seq}
Dans le diagramme (\ref{equation-functorial-omega}), supposons que
$S \to S'$ soit surjectif de noyau $I \subset S$, et supposons que
$R' = R$.
Il existe alors une suite exacte canonique de $S'$-modules
$$
I/I^2
\longrightarrow
\Omega_{S/R} \otimes_S S'
\longrightarrow
\Omega_{S'/R}
\longrightarrow
0
$$
Le morphisme de gauche est caractérisé par la règle qui envoie
$f \in I$ sur $\text{d}f \otimes 1$.
\end{lemma}

\begin{proof}
Le terme du milieu est $\Omega_{S/R} \otimes_S S/I$.
Pour $f \in I$, notons $\overline{f}$ l'image de $f$ dans $I/I^2$.
Pour montrer que l'application
$\overline{f} \mapsto \text{d}f \otimes 1$ est bien définie, il suffit de
vérifier que $\text{d} f_1f_2 \otimes 1 = 0$ si $f_1, f_2 \in I$.
Cela résulte clairement de la règle de Leibniz
$\text{d} f_1f_2 \otimes 1 = (f_1 \text{d}f_2 + f_2 \text{d} f_1 )\otimes 1 =
\text{d}f_2 \otimes f_1 + \text{d}f_1 \otimes f_2 = 0$.
Un calcul analogue montre que cette application est
$S' = S/I$-linéaire.

\medskip\noindent
Le morphisme $\Omega_{S/R} \otimes_S S' \to \Omega_{S'/R}$ est le morphisme
$S'$-linéaire canonique associé au morphisme $S$-linéaire
$\Omega_{S/R} \to \Omega_{S'/R}$. Il est surjectif parce que
$\Omega_{S/R} \to \Omega_{S'/R}$ est surjectif d'après le
Lemme \ref{lemma-differential-surjective}.

\medskip\noindent
La composée des deux morphismes est nulle parce que $\text{d}f$ est envoyé
sur zéro dans $\Omega_{S'/R}$ pour $f \in I$. Notons que l'exactitude dit
simplement que le noyau de $\Omega_{S/R} \to \Omega_{S'/R}$ est engendré
comme sous-$S$-module par le sous-module $I\Omega_{S/R}$ et par les éléments
$\text{d}f$, avec $f \in I$. Nous savons par le
Lemme \ref{lemma-differential-surjective} que ce noyau est engendré par les
éléments $\text{d}(a)$ où $\varphi(a) = \beta(r)$ pour un certain $r \in R$.
Mais alors $a = \alpha(r) + a - \alpha(r)$, de sorte que
$\text{d}(a) = \text{d}(a - \alpha(r))$. De plus,
$a - \alpha(r) \in I$ puisque $\varphi(a - \alpha(r)) =
\varphi(a) - \varphi(\alpha(r)) = \beta(r) - \beta(r) = 0$.
Nous concluons que les éléments $\text{d}f$ avec $f \in I$ engendrent déjà
le noyau comme $S$-module, comme souhaité.
\end{proof}

\begin{lemma}
\label{lemma-differential-seq-split}
Dans le diagramme (\ref{equation-functorial-omega}), supposons que
$S \to S'$ soit surjectif de noyau $I \subset S$ et supposons que
$R' = R$. Supposons en outre qu'il existe un morphisme de $R$-algèbres
$S' \to S$ qui soit un inverse à droite de $S \to S'$.
Alors la suite exacte de $S'$-modules du
Lemme \ref{lemma-differential-seq} devient une suite exacte courte
$$
0 \longrightarrow
I/I^2
\longrightarrow
\Omega_{S/R} \otimes_S S'
\longrightarrow
\Omega_{S'/R}
\longrightarrow
0
$$
qui est même une suite exacte courte scindée.
\end{lemma}

\begin{proof}
Soit $\beta : S' \to S$ l'inverse à droite de la surjection
$\alpha : S \to S'$. Considérons l'application
$$
D : S \longrightarrow I/I^2,
\quad
x \longmapsto x - \beta(\alpha(x))
$$
Il est facile de montrer que $D$ est une $R$-dérivation (nous l'omettons).
En outre, $x D(s) = 0$ si $x \in I, s \in S$. Par conséquent, la propriété
universelle montre que $D$ induit un morphisme
$\tau : \Omega_{S/R} \otimes_S S' \to I/I^2$.
Nous omettons la vérification que c'est un inverse à gauche de
$\text{d} : I/I^2  \to \Omega_{S/R} \otimes_S S'$. D'où le résultat.
\end{proof}

\begin{lemma}
\label{lemma-differential-mod-power-ideal}
Soit $R \to S$ un morphisme d'anneaux. Soit $I \subset S$ un idéal.
Soit $n \geq 1$ un entier. Posons $S' = S/I^{n + 1}$.
Le morphisme $\Omega_{S/R} \to \Omega_{S'/R}$ induit un isomorphisme
$$
\Omega_{S/R} \otimes_S S/I^n
\longrightarrow
\Omega_{S'/R} \otimes_{S'} S/I^n.
$$
\end{lemma}

\begin{proof}
Cela résulte du Lemme \ref{lemma-differential-seq} et du fait que
$\text{d}(I^{n + 1}) \subset I^n\Omega_{S/R}$ par la règle de Leibniz pour
$\text{d}$.
\end{proof}

\begin{lemma}
\label{lemma-differentials-base-change}
Supposons donnés des morphismes d'anneaux $R \to R'$ et $R \to S$.
Posons $S' = S \otimes_R R'$, de sorte que nous obtenions un diagramme
(\ref{equation-functorial-omega}). Alors le morphisme canonique défini
ci-dessus induit un isomorphisme
$\Omega_{S/R} \otimes_R R' = \Omega_{S'/R'}$.
\end{lemma}

\begin{proof}
Notons $\text{d}' : S' = S \otimes_R R' \to \Omega_{S/R} \otimes_R R'$
l'application
$\text{d}'( \sum a_i \otimes x_i ) = \sum \text{d}(a_i) \otimes x_i$.
Elle existe parce que l'application
$S \times R' \to \Omega_{S/R} \otimes_R R'$,
$(a, x)\mapsto \text{d}a \otimes_R x$, est $R$-bilinéaire.
C'est une $R'$-dérivation, comme on le vérifie par un calcul simple.
Nous allons montrer que $(\Omega_{S/R} \otimes_R R', \text{d}')$ satisfait
la propriété universelle. Soit $D : S' \to M'$ une $R'$-dérivation à valeurs
dans un $S'$-module. La composée $S \to S' \to M'$ est une $R$-dérivation ;
nous obtenons donc un morphisme $S$-linéaire
$\varphi_D : \Omega_{S/R} \to M'$. Nous pouvons le tensoriser par $R'$ et
obtenir le morphisme $\varphi'_D :
\Omega_{S/R} \otimes_R R' \to M'$, $\varphi'_D(\eta \otimes x) =
x\varphi_D(\eta)$. Il est clair que $D = \varphi'_D \circ \text{d}'$.
\end{proof}

\noindent
Le morphisme de multiplication $S \otimes_R S \to S$ est le morphisme de
$R$-algèbres qui envoie $a \otimes b$ sur $ab$ dans $S$. C'est aussi un
morphisme de $S$-algèbres si l'on considère $S \otimes_R S$ comme une
$S$-algèbre par l'un ou l'autre des morphismes $S \to S \otimes_R S$.

\begin{lemma}
\label{lemma-differentials-diagonal}
Soit $R \to S$ un morphisme d'anneaux. Soit
$J = \Ker(S \otimes_R S \to S)$ le noyau du morphisme de multiplication.
Il existe un isomorphisme canonique de $S$-modules
$\Omega_{S/R} \to J/J^2$,
$a \text{d} b \mapsto a \otimes b - ab \otimes 1$.
\end{lemma}

\begin{proof}[Première démonstration]
Appliquons le Lemme \ref{lemma-differential-seq-split} au diagramme commutatif
$$
\xymatrix{
S \otimes_R S \ar[r] & S \\
S \ar[r] \ar[u] & S \ar[u]
}
$$
où la flèche verticale de gauche est $a \mapsto a \otimes 1$. Nous obtenons
la suite exacte
$0 \to J/J^2 \to
\Omega_{S \otimes_R S/S} \otimes_{S \otimes_R S} S \to \Omega_{S/S} \to 0$.
D'après le Lemme \ref{lemma-trivial-differential-surjective}, le terme
$\Omega_{S/S}$ est $0$, et nous obtenons un isomorphisme entre les deux
autres termes. D'après le Lemme \ref{lemma-differentials-base-change}, on a
$\Omega_{S \otimes_R S/S} = \Omega_{S/R} \otimes_S (S \otimes_R S)$,
car $S \to S \otimes_R S$ est le changement de base de $R \to S$, et donc
$$
\Omega_{S \otimes_R S/S} \otimes_{S \otimes_R S} S =
\Omega_{S/R} \otimes_S (S \otimes_R S) \otimes_{S \otimes_R S} S =
\Omega_{S/R}
$$
Nous omettons la vérification que le morphisme est donné par la règle du
lemme.
\end{proof}

\begin{proof}[Seconde démonstration]
Montrons d'abord que la règle
$a \text{d} b \mapsto a \otimes b - ab \otimes 1$ est bien définie.
Pour cela, nous devons montrer que $\text{d}r$ et
$a\text{d}b + b \text{d}a - d(ab)$ sont envoyés sur zéro.
Le premier l'est parce que $r \otimes 1 - 1 \otimes r = 0$ par définition
du produit tensoriel. Le second l'est parce que
$$
(a \otimes b - ab \otimes 1) +
(b \otimes a - ba \otimes 1) -
(1 \otimes ab - ab \otimes 1)
=
(a \otimes 1 - 1\otimes a)(1\otimes b - b \otimes 1)
$$
appartient à $J^2$.

\medskip\noindent
Construisons un morphisme dans l'autre sens.
Nous pouvons considérer $S \to S \otimes_R S$, $a \mapsto a \otimes 1$,
comme le changement de base de $R \to S$. Ainsi, d'après le
Lemme \ref{lemma-differentials-base-change}, on a
$\Omega_{S \otimes_R S/S} = \Omega_{S/R} \otimes_S (S \otimes_R S)$.
À ce stade, la suite du Lemme \ref{lemma-differential-seq} donne un morphisme
$$
J/J^2  \to \Omega_{S \otimes_R S/ S} \otimes_{S \otimes_R S} S
= (\Omega_{S/R} \otimes_S (S \otimes_R S))\otimes_{S \otimes_R S} S
= \Omega_{S/R}.
$$
Nous laissons au lecteur le soin de vérifier qu'il est l'inverse du
morphisme ci-dessus.
\end{proof}





\begin{lemma}
\label{lemma-differentials-polynomial-ring}
Si $S = R[x_1, \ldots, x_n]$, alors $\Omega_{S/R}$ est un $S$-module libre
fini de base $\text{d}x_1, \ldots, \text{d}x_n$.
\end{lemma}

\begin{proof}
Montrons d'abord que $\text{d}x_1, \ldots, \text{d}x_n$ engendrent
$\Omega_{S/R}$ comme $S$-module. Pour cela, montrons que $\text{d}g$ peut
s'écrire comme une somme $\sum g_i \text{d}x_i$ pour tout
$g \in R[x_1, \ldots, x_n]$. Nous raisonnons par récurrence sur le degré
(total) de $g$. C'est clair si le degré de $g$ est $0$, car alors
$\text{d}g = 0$. Si le degré de $g$ est $> 0$, nous pouvons écrire $g$
sous la forme $c + \sum g_i x_i$ avec $c\in R$ et
$\deg(g_i) < \deg(g)$.
La règle de Leibniz donne
$\text{d}g = \sum g_i \text{d} x_i + \sum x_i \text{d}g_i$,
et la récurrence permet de conclure.

\medskip\noindent
Considérons la $R$-dérivation $\partial / \partial x_i :
R[x_1, \ldots, x_n] \to R[x_1, \ldots, x_n]$. (Nous laissons au lecteur le
soin de la définir ; sa propriété caractéristique est que
$\partial / \partial x_i (x_j) = \delta_{ij}$.)
Par la propriété universelle, elle correspond à un morphisme de $S$-modules
$l_i : \Omega_{S/R} \to R[x_1, \ldots, x_n]$ qui envoie $\text{d}x_i$
sur $1$ et $\text{d}x_j$ sur $0$ pour $j \not = i$.
Il est donc clair qu'il n'existe aucune relation $S$-linéaire entre les
éléments $\text{d}x_1, \ldots, \text{d}x_n$.
\end{proof}

\begin{lemma}
\label{lemma-differentials-finitely-presented}
Supposons que $R \to S$ soit de présentation finie.
Alors $\Omega_{S/R}$ est un $S$-module de présentation finie.
\end{lemma}

\begin{proof}
Écrivons $S = R[x_1, \ldots, x_n]/(f_1, \ldots, f_m)$.
Écrivons $I = (f_1, \ldots, f_m)$. D'après le
Lemme \ref{lemma-differential-seq}, il existe une suite exacte de $S$-modules
$$
I/I^2
\to
\Omega_{R[x_1, \ldots, x_n]/R} \otimes_{R[x_1, \ldots, x_n]} S
\to
\Omega_{S/R}
\to
0
$$
Le résultat découle du fait que $I/I^2$ est un $S$-module de type fini
(engendré par les images des $f_i$) et que le terme du milieu est libre de
type fini d'après le
Lemme \ref{lemma-differentials-polynomial-ring}.
\end{proof}

\begin{lemma}
\label{lemma-differentials-finitely-generated}
Supposons que $R \to S$ soit de type fini.
Alors $\Omega_{S/R}$ est un $S$-module de type fini.
\end{lemma}

\begin{proof}
La démonstration est très semblable à celle du
Lemme \ref{lemma-differentials-finitely-presented}, mais plus facile.
\end{proof}


\section{Le complexe de de Rham}
\label{section-de-rham-complex}

\noindent
Soit $A \to B$ un morphisme d'anneaux. Notons
$\text{d} : B \to \Omega_{B/A}$ le module des différentielles muni de sa
$A$-dérivation universelle construite dans la
section \ref{section-differentials}.
Soit $\Omega_{B/A}^i = \wedge^i_B(\Omega_{B/A})$, pour $i \geq 0$, la
$i$-ième puissance extérieure, comme dans la
section \ref{section-tensor-algebra}.
Le {\it complexe de de Rham de $B$ sur $A$} est le complexe
$$
\Omega_{B/A}^0 \to \Omega_{B/A}^1 \to \Omega_{B/A}^2 \to \ldots
$$
muni des différentielles $A$-linéaires construites et décrites ci-dessous.

\medskip\noindent
L'application $\text{d} : \Omega^0_{B/A} \to \Omega^1_{B/A}$ est la
dérivation universelle $\text{d} : B \to \Omega_{B/A}$.
Observons qu'elle est bien $A$-linéaire.

\medskip\noindent
Pour $p \geq 1$, nous affirmons qu'il existe une unique application
$A$-linéaire $\text{d} : \Omega_{B/A}^p \to \Omega_{B/A}^{p + 1}$ telle que
\begin{equation}
\label{equation-rule}
\text{d}\left(b_0\text{d}b_1 \wedge \ldots \wedge \text{d}b_p\right) =
\text{d}b_0 \wedge \text{d}b_1 \wedge \ldots \wedge \text{d}b_p
\end{equation}
Rappelons que $\Omega_{B/A}$ est engendré comme $B$-module par les éléments
$\text{d}b$. Ainsi, $\Omega^p_{B/A}$ est engendré comme $A$-module par les
éléments $b_0\text{d}b_1 \wedge \ldots \wedge \text{d}b_p$, et il s'ensuit
que l'application $\text{d} : \Omega^p_{B/A} \to \Omega^{p + 1}_{B/A}$,
si elle existe, est unique.

\medskip\noindent
Construction de $\text{d} : \Omega_{B/A}^1 \to \Omega_{B/A}^2$.
D'après la Définition \ref{definition-differentials}, les éléments
$\text{d}b$ engendrent librement $\Omega_{B/A}$ comme $B$-module, modulo
les relations $\text{d}a = 0$ pour $a \in A$ et
$\text{d}(b' + b'') = \text{d}b' + \text{d}b''$ et
$\text{d}(b'b'') = b'\text{d}b'' + b''\text{d}b'$ pour $b', b'' \in B$.
Par conséquent, pour montrer que la règle
$$
\sum b'_i \text{d}b_i \longmapsto \sum \text{d}b'_i \wedge \text{d}b_i
$$
est bien définie, nous devons montrer que les éléments
$$
b\text{d}a,
\quad\text{et}\quad
b\text{d}(b' + b'') - b\text{d}b' - b\text{d}b''
\quad\text{et}\quad
b\text{d}(b'b'') - bb'\text{d}b'' - bb''\text{d}b'
$$
pour $a \in A$ et $b, b', b'' \in B$ sont envoyés sur zéro.
Cela résulte clairement d'un calcul direct utilisant la règle de Leibniz
pour $\text{d}$.

\medskip\noindent
Observons que la composée
$\Omega^0_{B/A} \to \Omega^1_{B/A} \to \Omega^2_{B/A}$ est nulle puisque
$\text{d}(\text{d}(b)) = \text{d}(1 \text{d}b) =
\text{d}(1) \wedge \text{d}(b) = 0 \wedge \text{d}b = 0$.
Ici, $\text{d}(1) = 0$ car $1 \in B$ appartient à l'image de $A \to B$.
Nous utiliserons ce fait ci-dessous.

\medskip\noindent
Construction de $\text{d} : \Omega_{B/A}^p \to \Omega_{B/A}^{p + 1}$
pour $p \geq 2$. Nous allons montrer que l'application $A$-linéaire
$$
\gamma : \Omega^1_{B/A} \otimes_A \ldots \otimes_A \Omega^1_{B/A}
\longrightarrow
\Omega_{B/A}^{p + 1}
$$
définie par la formule
$$
\omega_1 \otimes \ldots \otimes \omega_p
\longmapsto
\sum (-1)^{i + 1}
\omega_1 \wedge \ldots \wedge \text{d}(\omega_i) \wedge \ldots \wedge \omega_p
$$
se factorise par la surjection naturelle
$\Omega^1_{B/A} \otimes_A \ldots \otimes_A \Omega^1_{B/A} \to \Omega^p_{B/A}$
pour donner l'application souhaitée
$\text{d} : \Omega^p_{B/A} \to \Omega^{p + 1}_{B/A}$.
D'après le Lemme \ref{lemma-present-wedge}, le noyau de
$\Omega^1_{B/A} \otimes_A \ldots \otimes_A \Omega^1_{B/A} \to \Omega^p_{B/A}$
est engendré comme $A$-module par les éléments
$\omega_1 \otimes \ldots \otimes \omega_p$ tels que
$\omega_i = \omega_j$ pour certains $i \not = j$, et par les éléments
$\omega_1 \otimes \ldots \otimes f\omega_i \otimes \ldots \otimes \omega_p -
\omega_1 \otimes \ldots \otimes f\omega_j \otimes \ldots \otimes \omega_p$
pour un certain $f \in B$.
Un calcul direct montre que le premier type d'éléments est envoyé sur $0$
par $\gamma$ ; autrement dit, $\gamma$ est alternée.
Pour terminer, nous devons montrer que
$$
\gamma(
\omega_1 \otimes \ldots \otimes f\omega_i \otimes \ldots \otimes \omega_p)
=
\gamma(
\omega_1 \otimes \ldots \otimes f\omega_j \otimes \ldots \otimes \omega_p)
$$
pour $f \in B$. Par $A$-linéarité et par la propriété d'alternance, il suffit
de le montrer pour $p = 2$, $i = 1$, $j = 2$,
$\omega_1 = b \text{d}b'$ et $\omega_2 = c \text{d} c'$ avec
$b, b', c, c' \in B$.
Nous devons donc montrer que
\begin{align*}
& \text{d}(fb) \wedge \text{d}b' \wedge c \text{d}c'
- fb \text{d}b' \wedge \text{d}c \wedge \text{d}c' \\
& =
\text{d}b \wedge \text{d}b' \wedge fc\text{d}c'
- b \text{d}b' \wedge \text{d}(fc) \wedge \text{d}c'
\end{align*}
c'est-à-dire que
$$
(c \text{d}(fb) + fb \text{d}c - fc \text{d}b - b \text{d}(fc))
\wedge \text{d}b' \wedge \text{d}c' = 0.
$$
Cela découle de la règle de Leibniz. Observons que la valeur de $\gamma$
sur l'élément
$b_0\text{d}b_1 \otimes \text{d}b_2 \otimes \ldots \otimes \text{d}b_p$
est $\text{d}b_0 \wedge \text{d}b_1 \wedge \ldots \wedge \text{d}b_p$,
et donc que (\ref{equation-rule}) est satisfaite pour l'application
$\text{d} : \Omega^p_{B/A} \to \Omega^{p + 1}_{B/A}$ ainsi obtenue.

\medskip\noindent
Enfin, puisque $\Omega^p_{B/A}$ est engendré additivement par les éléments
$b_0\text{d}b_1 \wedge \ldots \wedge \text{d}b_p$ et puisque
$\text{d}(b_0\text{d}b_1 \wedge \ldots \wedge \text{d}b_p) =
\text{d}b_0 \wedge \ldots \wedge \text{d}b_p$, nous voyons exactement de
la même manière que la composée
$
\Omega^p_{B/A} \to
\Omega^{p + 1}_{B/A} \to
\Omega^{p + 2}_{B/A}
$
est nulle pour $p \geq 1$. Ainsi, le complexe de de Rham est bien un
complexe.

\medskip\noindent
Étant donné un anneau $R$, nous posons $\Omega_R = \Omega_{R/\mathbf{Z}}$.
On l'appelle parfois le module absolu des différentielles de $R$ ; cette
terminologie a un sens : si $\Omega_R$ est le module des différentielles
pour lequel nous imposons seulement la règle de Leibniz, sans imposer
l'annulation de $\text{d}1$, alors la règle de Leibniz donne
$\text{d}1 = \text{d}(1 \cdot 1) =
1 \text{d}1 + 1 \text{d}1 = 2 \text{d}1$, et donc
$\text{d}1 = 0$ dans $\Omega_R$. Dans ce cas, le
{\it complexe de de Rham absolu de $R$} est le complexe correspondant
$$
\Omega_R^0 \to \Omega_R^1 \to \Omega_R^2 \to \ldots
$$
où nous posons $\Omega^i_R = \Omega^i_{R/\mathbf{Z}}$, et ainsi de suite.

\medskip\noindent
Supposons donné un diagramme commutatif d'anneaux
$$
\xymatrix{
B \ar[r] & B' \\
A \ar[r] \ar[u] & A' \ar[u]
}
$$
Il existe un morphisme naturel de complexes de de Rham
$$
\Omega^\bullet_{B/A} \longrightarrow \Omega^\bullet_{B'/A'}
$$
En degré $0$, il s'agit du morphisme $B \to B'$ ; en degré $1$, il s'agit
du morphisme $\Omega_{B/A} \to \Omega_{B'/A'}$ construit dans la
section \ref{section-differentials} ; et, pour $p \geq 2$, il s'agit du
morphisme induit
$\Omega^p_{B/A} = \wedge^p_B(\Omega_{B/A}) \to
\wedge^p_{B'}(\Omega_{B'/A'}) = \Omega^p_{B'/A'}$.
La compatibilité aux différentielles découle de la caractérisation des
différentielles par la formule (\ref{equation-rule}).

\begin{lemma}
\label{lemma-de-rham-complex-base-change}
Supposons donnés des morphismes d'anneaux $A \to A'$ et $A \to B$.
Posons $B' = B \otimes_A A'$, de sorte que nous obtenions un diagramme
comme ci-dessus. Alors le morphisme canonique défini ci-dessus induit un
isomorphisme
$\Omega^\bullet_{B/A} \otimes_A A' = \Omega^\bullet_{B'/A'}$
de complexes.
\end{lemma}

\begin{proof}
Cela découle du Lemme \ref{lemma-differentials-base-change} et du fait que
la formation des puissances extérieures commute au changement de base.
\end{proof}

\begin{lemma}
\label{lemma-de-rham-complex}
Soit $A \to B$ un morphisme d'anneaux. Soit
$\pi : \Omega_{B/A} \to \Omega$ un morphisme surjectif de $B$-modules.
Notons $\text{d} : B \to \Omega$ la composée de $\pi$ avec la dérivation
universelle $\text{d}_{B/A} : B \to \Omega_{B/A}$.
Posons $\Omega^i = \wedge_B^i(\Omega)$.
Supposons que le noyau de $\pi$ soit engendré, comme $B$-module, par des
éléments $\omega \in \Omega_{B/A}$ tels que
$\text{d}_{B/A}(\omega) \in \Omega_{B/A}^2$ soit envoyé sur zéro dans
$\Omega^2$. Il existe alors un complexe de de Rham
$$
\Omega^0 \to \Omega^1 \to \Omega^2 \to \ldots
$$
dont la différentielle est définie par la règle
$$
\text{d} : \Omega^p \to \Omega^{p + 1},\quad
\text{d}\left(f_0\text{d}f_1 \wedge \ldots \wedge \text{d}f_p\right) =
\text{d}f_0 \wedge \text{d}f_1 \wedge \ldots \wedge \text{d}f_p
$$
\end{lemma}

\begin{proof}
Nous allons montrer qu'il existe un diagramme commutatif
$$
\xymatrix{
\Omega_{B/A}^0 \ar[d] \ar[r]_{\text{d}_{B/A}} &
\Omega_{B/A}^1 \ar[d]_\pi \ar[r]_{\text{d}_{B/A}} &
\Omega_{B/A}^2 \ar[d]_{\wedge^2\pi} \ar[r]_{\text{d}_{B/A}} &
\ldots \\
\Omega^0 \ar[r]^{\text{d}} &
\Omega^1 \ar[r]^{\text{d}} &
\Omega^2 \ar[r]^{\text{d}} &
\ldots
}
$$
La description de l'application $\text{d}$ découlera de la construction,
donnée ci-dessus, des différentielles
$\text{d}_{B/A} : \Omega^p_{B/A} \to \Omega^{p + 1}_{B/A}$ du complexe
de de Rham de $B$ sur $A$.
Puisque la flèche verticale la plus à gauche est un isomorphisme, nous avons
le premier carré. Comme $\pi$ est surjectif, pour obtenir le deuxième carré,
il suffit de montrer que $\text{d}_{B/A}$ envoie le noyau de $\pi$ dans le
noyau de $\wedge^2\pi$. Par hypothèse, tout élément du noyau de $\pi$ est de
la forme $\sum b_i\omega_i$ avec $\pi(\omega_i) = 0$ et
$\wedge^2\pi(\text{d}_{B/A}(\omega_i)) = 0$.
La règle de Leibniz pour $\text{d}_{B/A}$ donne
$\text{d}_{B/A}(\sum b_i\omega_i) = \sum b_i \text{d}_{B/A}(\omega_i) +
\sum \text{d}_{B/A}(b_i) \wedge \omega_i$. Cet élément est donc envoyé sur
zéro par $\wedge^2\pi$.

\medskip\noindent
Pour $i > 1$, notons que $\wedge^i \pi$ est surjectif et que son noyau est
l'image de $\Ker(\pi) \wedge \Omega^{i - 1}_{B/A}
\to \Omega_{B/A}^i$. Pour $\omega_1 \in \Ker(\pi)$ et
$\omega_2 \in \Omega^{i - 1}_{B/A}$, on a
$$
\text{d}_{B/A}(\omega_1 \wedge \omega_2) =
\text{d}_{B/A}(\omega_1) \wedge \omega_2 -
\omega_1 \wedge \text{d}_{B/A}(\omega_2)
$$
qui appartient au noyau de $\wedge^{i + 1}\pi$ d'après ce que nous venons
de démontrer. Nous obtenons donc le $(i + 1)$-ième carré du diagramme
ci-dessus. Cela achève la démonstration.
\end{proof}


\section{Opérateurs différentiels d'ordre fini}
\label{section-differential-operators}

\noindent
Dans cette section, nous introduisons les opérateurs différentiels d'ordre
fini.

\begin{definition}
\label{definition-differential-operators}
Soit $R \to S$ un morphisme d'anneaux. Soient $M$, $N$ des $S$-modules.
Soit $k \geq 0$ un entier. Nous définissons par récurrence un
{\it opérateur différentiel $D : M \to N$ d'ordre $k$} comme une application
$R$-linéaire telle que, pour tout $g \in S$, l'application
$m \mapsto D(gm) - gD(m)$ soit un opérateur différentiel d'ordre $k - 1$.
Pour le cas initial $k = 0$, nous définissons un opérateur différentiel
d'ordre $0$ comme une application $S$-linéaire.
\end{definition}

\noindent
Si $D : M \to N$ est un opérateur différentiel d'ordre $k$, alors, pour tout
$g \in S$, l'application $gD$ est un opérateur différentiel d'ordre $k$.
La somme de deux opérateurs différentiels d'ordre $k$ en est encore un.
L'ensemble de tous ces opérateurs
$$
\text{Diff}^k(M, N) = \text{Diff}^k_{S/R}(M, N)
$$
est donc un $S$-module. On a
$$
\text{Diff}^0(M, N) \subset
\text{Diff}^1(M, N) \subset
\text{Diff}^2(M, N) \subset \ldots
$$

\begin{lemma}
\label{lemma-composition-differential-operators}
Soit $R \to S$ un morphisme d'anneaux. Soient $L, M, N$ des $S$-modules.
Si $D : L \to M$ et $D' : M \to N$ sont des opérateurs différentiels
d'ordres $k$ et $k'$, alors $D' \circ D$ est un opérateur différentiel
d'ordre $k + k'$.
\end{lemma}

\begin{proof}
Soit $g \in S$. L'application qui envoie $x \in L$ sur
$$
D'(D(gx)) - gD'(D(x)) = D'(D(gx)) - D'(gD(x)) + D'(gD(x)) - gD'(D(x))
$$
est une somme de deux composées d'opérateurs différentiels d'ordres
inférieurs. Le lemme découle donc d'une récurrence sur $k + k'$.
\end{proof}

\begin{lemma}
\label{lemma-module-principal-parts}
Soit $R \to S$ un morphisme d'anneaux. Soit $M$ un $S$-module.
Soit $k \geq 0$. Il existe un $S$-module $P^k_{S/R}(M)$ et un isomorphisme
canonique
$$
\text{Diff}^k_{S/R}(M, N) = \Hom_S(P^k_{S/R}(M), N)
$$
fonctoriel en le $S$-module $N$.
\end{lemma}

\noindent
Par le lemme de Yoneda, cela nous dit qu'il existe un opérateur différentiel
« universel » $D_{univ} : M \to P^k_{S/R}(M)$ d'ordre $k$ tel que, pour tout
opérateur différentiel $D : M \to N$ d'ordre $k$, il existe une unique
application $S$-linéaire $\alpha : P^k_{S/R}(M) \to N$ telle que
$D = D_{univ} \circ \alpha$. Le couple
$(P^k_{S/R}(M), D_{univ})$ est unique à isomorphisme unique près.

\begin{proof}
L'existence de $P^k_{S/R}(M)$ découle d'arguments généraux de théorie des
catégories (insérer ici une référence future), mais nous allons également
en donner une construction. Posons $F = \bigoplus_{m \in M} S[m]$, où
$[m]$ est un symbole désignant l'élément de base du facteur correspondant
à $m$. Étant donné un opérateur différentiel $D : M \to N$, nous obtenons
une application $S$-linéaire $L_D : F \to N$ qui envoie $[m]$ sur $D(m)$.
Si $D$ est d'ordre $0$, alors $L_D$ annule les éléments
$$
[m + m'] - [m] - [m'],\quad
g_0[m] - [g_0m]
$$
où $g_0 \in S$ et $m, m' \in M$.
Si $D$ est d'ordre $1$, alors $L_D$ annule les éléments
$$
[m + m'] - [m] - [m'],\quad
f[m] - [fm], \quad
g_0g_1[m] - g_0[g_1m] - g_1[g_0m] + [g_1g_0m]
$$
où $f \in R$, $g_0, g_1 \in S$ et $m \in M$.
Si $D$ est d'ordre $k$, alors $L_D$ annule les éléments
$[m + m'] - [m] - [m']$, $f[m] - [fm]$, ainsi que les éléments
$$
g_0g_1\ldots g_k[m] - \sum g_0 \ldots \hat g_i \ldots g_k[g_im] + \ldots
+(-1)^{k + 1}[g_0\ldots g_km]
$$
Réciproquement, si $L : F \to N$ est une application $S$-linéaire qui
annule tous les éléments énumérés dans la phrase précédente, alors
$m \mapsto L([m])$ est un opérateur différentiel d'ordre $k$.
Nous voyons ainsi que $P^k_{S/R}(M)$ est le quotient de $F$ par le
sous-module engendré par ces éléments.
\end{proof}

\begin{definition}
\label{definition-module-principal-parts}
Soit $R \to S$ un morphisme d'anneaux. Soit $M$ un $S$-module. Le module
$P^k_{S/R}(M)$ construit dans le Lemme \ref{lemma-module-principal-parts}
est appelé le {\it module des parties principales d'ordre $k$} de $M$.
\end{definition}

\noindent
Notons que les inclusions
$$
\text{Diff}^0(M, N) \subset
\text{Diff}^1(M, N) \subset
\text{Diff}^2(M, N) \subset \ldots
$$
correspondent, par le lemme de Yoneda (Catégories,
Lemme \ref{categories-lemma-yoneda}), aux surjections
$$
\ldots \to P^2_{S/R}(M) \to P^1_{S/R}(M) \to P^0_{S/R}(M) = M
$$

\begin{example}
\label{example-derivations-and-differential-operators}
Soit $R \to S$ un morphisme d'anneaux et soit $N$ un $S$-module.
Observons que $\text{Diff}^1(S, N) = \text{Der}_R(S, N) \oplus N$.
En effet, si $D : S \to N$ est un opérateur différentiel d'ordre $1$, alors
$\sigma_D : S \to N$, défini par $\sigma_D(g) := D(g) - gD(1)$, est une
$R$-dérivation et $D = \sigma_D + \lambda_{D(1)}$, où
$\lambda_x : S \to N$ est l'application linéaire qui envoie $g$ sur $gx$.
Il s'ensuit que $P^1_{S/R} = \Omega_{S/R} \oplus S$ par la propriété
universelle de $\Omega_{S/R}$.
\end{example}

\begin{lemma}
\label{lemma-sequence-of-principal-parts}
Soit $R \to S$ un morphisme d'anneaux. Soit $M$ un $S$-module. Il existe
une suite exacte courte canonique
$$
0 \to \Omega_{S/R} \otimes_S M \to P^1_{S/R}(M) \to M \to 0
$$
fonctorielle en $M$, appelée la {\it suite des parties principales}.
\end{lemma}

\begin{proof}
Le morphisme $P^1_{S/R}(M) \to M$ est donné ci-dessus.
Soit $N$ un $S$-module et soit $D : M \to N$ un opérateur différentiel
d'ordre $1$. Pour $m \in M$, l'application
$$
g \longmapsto D(gm) - gD(m)
$$
est une $R$-dérivation $S \to N$ par les axiomes des opérateurs
différentiels d'ordre $1$. Elle correspond donc à une application linéaire
$D_m : \Omega_{S/R} \to N$ déterminée par la règle
$a\text{d}b \mapsto aD(bm) - abD(m)$
(voir le Lemme \ref{lemma-universal-omega}). L'application
$$
\Omega_{S/R} \times M \longrightarrow N,\quad
(\eta, m) \longmapsto D_m(\eta)
$$
est $S$-bilinéaire (nous omettons les détails) et détermine donc une
application $S$-linéaire
$$
\sigma_D : \Omega_{S/R} \otimes_S M \to N
$$
Nous obtenons ainsi une application
$\text{Diff}^1(M, N) \to \Hom_S(\Omega_{S/R} \otimes_S M, N)$,
$D \mapsto \sigma_D$, fonctorielle en $N$. Par le lemme de Yoneda, cela
correspond à un morphisme
$\Omega_{S/R} \otimes_S M \to P^1_{S/R}(M)$. Il résulte immédiatement de
la construction que ce morphisme est fonctoriel en $M$. La suite
$$
\Omega_{S/R} \otimes_S M \to P^1_{S/R}(M) \to M \to 0
$$
est exacte parce que, pour tout module $N$, la suite
$$
0 \to \Hom_S(M, N) \to
\text{Diff}^1(M, N) \to
\Hom_S(\Omega_{S/R} \otimes_S M, N)
$$
est exacte par inspection.

\medskip\noindent
Pour voir que $\Omega_{S/R} \otimes_S M \to P^1_{S/R}(M)$ est injectif,
raisonnons comme suit. Choisissons une suite exacte
$$
0 \to M' \to F \to M \to 0
$$
où $F$ est un $S$-module libre. Elle induit une suite exacte
$$
0 \to \text{Diff}^1(M, N) \to \text{Diff}^1(F, N) \to \text{Diff}^1(M', N)
$$
pour tout $N$. Cela démontre que, dans le diagramme commutatif
$$
\xymatrix{
0 \ar[r] &
\Omega_{S/R} \otimes_S M' \ar[r] \ar[d] &
P^1_{S/R}(M') \ar[r] \ar[d] &
M' \ar[r] \ar[d] &
0 \\
0 \ar[r] &
\Omega_{S/R} \otimes_S F \ar[r] \ar[d] &
P^1_{S/R}(F) \ar[r] \ar[d] &
F \ar[r] \ar[d] &
0 \\
0 \ar[r] &
\Omega_{S/R} \otimes_S M \ar[r] \ar[d] &
P^1_{S/R}(M) \ar[r] \ar[d] &
M \ar[r] \ar[d] &
0 \\
& 0 & 0 & 0
}
$$
la colonne du milieu est exacte. La colonne de gauche est exacte par
l'exactitude à droite de $\Omega_{S/R} \otimes_S -$. D'après le lemme du
serpent (voir la section \ref{section-snake}), il suffit de démontrer
l'exactitude à gauche pour le module libre $F$.
En utilisant le fait que $P^1_{S/R}(-)$ commute aux sommes directes, nous
nous ramenons au cas $M = S$. Ce cas découle de la discussion de
l'Exemple \ref{example-derivations-and-differential-operators}.
\end{proof}

\begin{remark}
\label{remark-functoriality-principal-parts}
Supposons donnés un diagramme commutatif d'anneaux
$$
\xymatrix{
B \ar[r] & B' \\
A \ar[u] \ar[r] & A' \ar[u]
}
$$
un $B$-module $M$, un $B'$-module $M'$ et une application $B$-linéaire
$M \to M'$. Nous obtenons alors un système compatible de morphismes de
modules
$$
\xymatrix{
\ldots \ar[r] &
P^2_{B'/A'}(M') \ar[r] &
P^1_{B'/A'}(M') \ar[r] &
P^0_{B'/A'}(M') \\
\ldots \ar[r] &
P^2_{B/A}(M) \ar[r] \ar[u] &
P^1_{B/A}(M) \ar[r] \ar[u] &
P^0_{B/A}(M) \ar[u]
}
$$
Ces morphismes sont compatibles avec toute composition ultérieure de
morphismes de ce type. La façon la plus simple de le voir est d'utiliser
la description des modules $P^k_{B/A}(M)$ par générateurs et relations dans
la démonstration du Lemme \ref{lemma-module-principal-parts}, mais on peut
aussi le voir directement à partir de la propriété universelle de ces
modules. En outre, ces morphismes sont compatibles avec les suites exactes
courtes du Lemme \ref{lemma-sequence-of-principal-parts}.
\end{remark}

\begin{lemma}
\label{lemma-principal-parts-base-change}
Supposons donnés des morphismes d'anneaux $A \to A'$ et $A \to B$, ainsi
qu'un $B$-module $M$. Posons $B' = B \otimes_A A'$ et considérons
$M' = M \otimes_A A'$ comme un $B'$-module.
Le morphisme de la Remarque \ref{remark-functoriality-principal-parts}
induit un isomorphisme
$P^k_{B/A}(M) \otimes_A A' = P^k_{B'/A'}(M')$.
\end{lemma}

\begin{proof}
Soit $D_{univ} : M \to P^k_{B/A}(M)$ l'opérateur différentiel universel.
L'application $A'$-linéaire induite
$D_{univ} \otimes 1 : M' = M \otimes_A A' \to P^k_{B/A}(M) \otimes_A A'$
est aisément vérifiée être un opérateur différentiel d'ordre $k$ pour
$M'/B'/A'$. Nous allons montrer que $D_{univ} \otimes 1$ satisfait la
propriété universelle. Soit $D' : M' \to N'$ un opérateur différentiel
d'ordre $k$ à valeurs dans un $B'$-module $N'$. Alors la composée
$D : M \to M' \to N'$ est un opérateur différentiel à valeurs dans $N'$,
considéré comme un $B$-module. Il existe donc une application $B$-linéaire
$\gamma : P^k_{B/A}(M) \to N'$ telle que $D = \gamma \circ D_{univ}$.
Le lecteur vérifiera aisément que l'application $B'$-linéaire induite
$\gamma' : P^k_{B/A}(M) \otimes_A A' \to N'$ satisfait
$\gamma' \circ (D_{univ} \otimes 1) = D'$.
Nous omettons la démonstration de l'unicité de $\gamma'$.
\end{proof}

\begin{lemma}
\label{lemma-principal-parts-diagonal}
Soit $R \to S$ un morphisme d'anneaux. Soit $M$ un $S$-module.
Soit $J = \Ker(S \otimes_R S \to S)$ le noyau du morphisme de
multiplication. Il existe un isomorphisme canonique de $S$-modules
$$
P^k_{S/R}(M) \longrightarrow (S \otimes_R M)/J^{k + 1}(S \otimes_R M)
$$
où $s \in S$ agit sur la cible par multiplication par $s \otimes 1$, et tel
que l'opérateur différentiel universel d'ordre $k$ corresponde à
l'application donnée par $m \mapsto \text{classe de }1 \otimes m$.
\end{lemma}

\begin{proof}
Considérons l'application $T : M \to S \otimes M$,
$m \mapsto 1 \otimes m$. Puisque $T$ est $R$-linéaire et puisque
\begin{align*}
g_0g_1\ldots g_k T(m)
- \sum g_0 \ldots \hat g_i \ldots g_k T(g_im) + \ldots
+(-1)^{k + 1}T(g_0\ldots g_km) \\
=
\prod_{i = 0,\ldots,k} (g_i \otimes 1 - 1 \otimes g_i) 1 \otimes m
\end{align*}
pour $m \in M$ et $g_0, \ldots, g_k \in S$, nous concluons que la règle
$m \mapsto \text{classe de }1 \otimes m$ de l'énoncé du lemme est un
opérateur différentiel d'ordre $k$ (voir la démonstration du
Lemme \ref{lemma-module-principal-parts}).
Par la propriété universelle de $P^k_{S/R}(M)$, nous obtenons la flèche de
l'énoncé du lemme. D'autre part, si $D : M \to N$ est un opérateur
différentiel d'ordre $k$, nous pouvons considérer l'application $S$-linéaire
$L : S \otimes_R M \to N$, $g \otimes m \mapsto gD(m)$.
Le lecteur vérifie, par un calcul semblable à celui ci-dessus et en utilisant
le fait que $J$ est engendré par des éléments de la forme
$g \otimes 1 - 1 \otimes g$, que $L$ annule
$J^{k + 1}(S \otimes_R M)$. En particulier, si nous appliquons cela à
l'opérateur différentiel universel
$D_{univ} : M \to P^k_{S/R}(M)$, nous obtenons une application $S$-linéaire
$(S \otimes_R M)/J^{k + 1}(S \otimes_R M) \to P^k_{S/R}(M)$.
Nous laissons au lecteur le soin de vérifier que cette application est
l'inverse de celle de l'énoncé du lemme.
\end{proof}

\begin{lemma}
\label{lemma-differentials-de-rham-complex-order-1}
Soit $A \to B$ un morphisme d'anneaux. Les différentielles
$\text{d} : \Omega^i_{B/A}  \to \Omega^{i + 1}_{B/A}$ sont des opérateurs
différentiels d'ordre $1$.
\end{lemma}

\begin{proof}
Étant donné $b \in B$, nous devons montrer que
$\text{d} \circ b - b \circ \text{d}$ est un opérateur linéaire.
Nous devons donc montrer que
$$
\text{d} \circ b \circ b' - b \circ \text{d} \circ b' -
b' \circ \text{d} \circ b + b' \circ b \circ \text{d} = 0
$$
Pour le voir, il suffit de vérifier cette égalité sur des générateurs
additifs de $\Omega^i_{B/A}$. Il suffit donc de montrer que
$$
\text{d}(bb'b_0\text{d}b_1 \wedge \ldots \wedge \text{d}b_i) -
b\text{d}(b'b_0\text{d}b_1 \wedge \ldots \wedge \text{d}b_i) -
b'\text{d}(bb_0\text{d}b_1 \wedge \ldots \wedge \text{d}b_i) +
bb'\text{d}(b_0\text{d}b_1 \wedge \ldots \wedge \text{d}b_i)
$$
est nul. C'est un agréable calcul utilisant la règle de Leibniz, que nous
laissons au lecteur.
\end{proof}

\begin{lemma}
\label{lemma-check-differential-operators}
Soit $A \to B$ un morphisme d'anneaux. Soient $g_i \in B$, $i \in I$, un
ensemble de générateurs de $B$ comme $A$-algèbre. Soient $M, N$ des
$B$-modules. Soit $D : M \to N$ une application $A$-linéaire. Pour montrer
que $D$ est un opérateur différentiel d'ordre $k$, il suffit de montrer que
$D \circ g_i - g_i \circ D$ est un opérateur différentiel d'ordre $k - 1$
pour $i \in I$.
\end{lemma}

\begin{proof}
En effet, nous affirmons que l'ensemble des éléments $g \in B$ tels que
$D \circ g - g \circ D$ soit un opérateur différentiel d'ordre $k - 1$ est
une sous-$A$-algèbre de $B$. Cela découle des relations
$$
D \circ (g + g') - (g + g') \circ D =
(D \circ g - g \circ D) + (D \circ g' - g' \circ D)
$$
et
$$
D \circ gg' - gg' \circ D =
(D \circ g - g \circ D) \circ g' + g \circ (D \circ g' - g' \circ D)
$$
À proprement parler, pour conclure dans le cas des produits, nous utilisons
aussi le Lemme \ref{lemma-composition-differential-operators}.
\end{proof}

\begin{lemma}
\label{lemma-invert-system-differential-operators}
Soit $A \to B$ un morphisme d'anneaux. Soient $M, N$ des $B$-modules.
Soit $S \subset B$ une partie multiplicative. Tout opérateur différentiel
$D : M \to N$ d'ordre $k$ se prolonge de manière unique en un opérateur
différentiel $E : S^{-1}M \to S^{-1}N$ d'ordre $k$.
\end{lemma}

\begin{proof}
Raisonnons par récurrence sur $k$. Si $k = 0$, alors $D$ est $B$-linéaire,
et la fonctorialité de la localisation fournit le prolongement.
Étant donné $b \in B$, l'opérateur
$L_b : m \mapsto D(bm) - bD(m)$ est d'ordre $k - 1$. Par récurrence, il
possède donc un unique prolongement en un opérateur différentiel
$E_b : S^{-1}M \to S^{-1}N$ d'ordre $k - 1$.
En outre, un calcul montre que
$L_{b'b} = L_{b'} \circ b + b' \circ L_b$ ; l'unicité donne donc
$E_{b'b} = E_{b'} \circ b + b' \circ E_b$.
De même, nous obtenons $E_{b'} \circ b - b \circ E_{b'} =
E_b \circ b' - b' \circ E_b$.
Pour $m \in M$ et $g \in S$, posons maintenant
$$
E(m/g) = (1/g)(D(m) - E_g(m/g))
$$
Pour montrer que cette définition est indépendante du choix du représentant,
il suffit de montrer que, pour $g' \in S$, l'emploi du représentant
$g'm/g'g$ donne le même résultat. Nous calculons
\begin{align*}
(1/g'g)(D(g'm) - E_{g'g}(g'm/gg'))
& =
(1/gg')(g'D(m) + E_{g'}(m) - E_{g'g}(g'm/gg')) \\
& =
(1/g'g)(g'D(m) - g' E_g(m/g))
\end{align*}
qui est le même qu'auparavant. Il est clair que $E$ est $R$-linéaire puisque
$D$ et $E_g$ sont $R$-linéaires. En prenant $g = 1$ et en utilisant
$E_1 = 0$, nous voyons que $E$ prolonge $D$.
D'après le Lemme \ref{lemma-check-differential-operators}, il suffit
maintenant de montrer que $E \circ b - b \circ E$ pour $b \in B$ et
$E \circ 1/g' - 1/g' \circ E$ pour $g' \in S$ sont des opérateurs
différentiels d'ordre $k - 1$ afin de montrer que $E$ est un opérateur
différentiel d'ordre $k$. Pour le premier, choisissons un élément $m/g$ de
$S^{-1}M$ et observons que
\begin{align*}
E(b m/g) - bE(m/g)
& =
(1/g)(D(bm) - bD(m) - E_g(bm/g) + bE_g(m/g)) \\
& =
(1/g)(L_b(m) - E_b(m) + gE_b(m/g)) \\
& =
E_b(m/g)
\end{align*}
qui est un opérateur différentiel d'ordre $k - 1$. Enfin, on a
\begin{align*}
E(m/g'g) - (1/g')E(m/g)
& =
(1/g'g)(D(m) - E_{g'g}(m/g'g)) -
(1/g'g)(D(m) - E_g(m/g)) \\
& =
-(1/g')E_{g'}(m/g'g)
\end{align*}
qui est également un opérateur différentiel d'ordre $k - 1$, comme composée
d'applications linéaires (multiplication par $1/g'$ et changement de signe)
et de $E_{g'}$. Nous omettons la démonstration de l'unicité.
\end{proof}

\begin{lemma}
\label{lemma-base-change-differential-operators}
Soient $R \to A$ et $R \to B$ des morphismes d'anneaux. Soient $M$ et $M'$
des $A$-modules. Soit $D : M \to M'$ un opérateur différentiel d'ordre $k$
relativement à $R \to A$. Soit $N$ un $B$-module quelconque. Alors
l'application
$$
D \otimes \text{id}_N : M \otimes_R N \to M' \otimes_R N
$$
est un opérateur différentiel d'ordre $k$ relativement à
$B \to A \otimes_R B$.
\end{lemma}

\begin{proof}
Il est clair que $D' = D \otimes \text{id}_N$ est $B$-linéaire.
D'après le Lemme \ref{lemma-check-differential-operators}, il suffit de
montrer que
$$
D' \circ a \otimes 1 -  a \otimes 1 \circ D' =
(D \circ a - a \circ D) \otimes \text{id}_N
$$
est un opérateur différentiel d'ordre $k - 1$, ce qui découle d'une
récurrence sur $k$.
\end{proof}


\section{Le complexe cotangent naïf}
\label{section-netherlander}

\noindent
Soit $R \to S$ un morphisme d'anneaux. Notons $R[S]$ l'anneau de polynômes
dont les variables sont les éléments $s \in S$. Notons $[s] \in R[S]$ la
variable correspondant à $s \in S$. Ainsi, $R[S]$ est un $R$-module libre
sur les éléments de base $[s_1] \ldots [s_n]$, où
$s_1, \ldots, s_n$ parcourt toutes les suites non ordonnées d'éléments de
$S$. Il existe une surjection canonique
\begin{equation}
\label{equation-canonical-presentation}
R[S] \longrightarrow S,\quad [s] \longmapsto s
\end{equation}
dont nous notons le noyau $I \subset R[S]$. Une observation simple montre
que $I$ est engendré par les éléments
$[s + s'] - [s] - [s']$, $[s][s'] - [ss']$ et $[r] - r$.
D'après le Lemme \ref{lemma-differential-seq}, il existe un morphisme
canonique
\begin{equation}
\label{equation-naive-cotangent-complex}
I/I^2 \longrightarrow \Omega_{R[S]/R} \otimes_{R[S]} S
\end{equation}
dont le conoyau est canoniquement isomorphe à $\Omega_{S/R}$. Observons que
le $S$-module $\Omega_{R[S]/R} \otimes_{R[S]} S$ est libre sur les
générateurs $\text{d}[s]$.

\begin{definition}
\label{definition-naive-cotangent-complex}
Soit $R \to S$ un morphisme d'anneaux. Le {\it complexe cotangent naïf}
$\NL_{S/R}$ est le complexe de chaînes
(\ref{equation-naive-cotangent-complex})
$$
\NL_{S/R} = \left(I/I^2 \longrightarrow \Omega_{R[S]/R} \otimes_{R[S]} S\right)
$$
où $I/I^2$ est placé en degré (homologique) $1$ et
$\Omega_{R[S]/R} \otimes_{R[S]} S$ en degré $0$. Nous noterons
$H_1(L_{S/R}) = H_1(\NL_{S/R})$\footnote{Ce module est parfois noté
$\Gamma_{S/R}$ dans la littérature.} l'homologie en degré $1$.
\end{definition}

\noindent
Avant de poursuivre, disons quelques mots du véritable complexe cotangent
(Complexe cotangent, section \ref{cotangent-section-cotangent-ring-map}).
Étant donné un morphisme d'anneaux $R \to S$, il existe une $R$-algèbre
simpliciale canonique $P_\bullet$ dont les termes sont des algèbres de
polynômes et qui est munie d'une équivalence d'homotopie canonique
$$
P_\bullet \longrightarrow S
$$
Le complexe cotangent $L_{S/R}$ de $S$ sur $R$ est défini comme le complexe
de chaînes associé au module simplicial
$$
\Omega_{P_\bullet/R} \otimes_{P_\bullet} S
$$
Le complexe cotangent naïf défini ci-dessus est canoniquement isomorphe à la
troncation $\tau_{\leq 1}L_{S/R}$ (voir
Homologie, section \ref{homology-section-truncations}, et
Complexe cotangent, section \ref{cotangent-section-surjections}).
En particulier, on a bien
$H_1(\NL_{S/R}) = H_1(L_{S/R})$ ; notre définition est donc compatible avec
celle qui utilise le complexe cotangent. En outre,
$H_0(L_{S/R}) = H_0(\NL_{S/R}) = \Omega_{S/R}$, comme nous l'avons vu
ci-dessus.

\medskip\noindent
Soit $R \to S$ un morphisme d'anneaux. Une {\it présentation de $S$ sur
$R$} est une surjection $\alpha : P \to S$ de $R$-algèbres, où $P$ est une
algèbre de polynômes (en un ensemble de variables). Souvent, lorsque $S$ est
de type fini sur $R$, nous l'indiquerons en disant : « Soit
$R[x_1, \ldots, x_n] \to S$ une présentation de $S/R$ », ou « Soit
$0 \to I \to R[x_1, \ldots, x_n] \to S \to 0$ une présentation de $S/R$ »
si nous voulons indiquer que $I$ est le noyau de la présentation.
Notons que le morphisme $R[S] \to S$ utilisé pour définir le complexe
cotangent naïf est un exemple de présentation.

\medskip\noindent
Notons que toute présentation $\alpha$ donne un complexe de chaînes à deux
termes de $S$-modules
$$
\NL(\alpha) : I/I^2 \longrightarrow \Omega_{P/R} \otimes_P S.
$$
Ici, le terme $I/I^2$ est placé en degré $1$ et le terme
$\Omega_{P/R} \otimes S$ en degré $0$. La classe de $f \in I$ dans $I/I^2$
est envoyée sur $\text{d}f \otimes 1$ dans $\Omega_{P/R} \otimes S$.
Le conoyau de ce complexe est canoniquement $\Omega_{S/R}$ ; voir le
Lemme \ref{lemma-differential-seq}. Nous appelons le complexe $\NL(\alpha)$
le {\it complexe cotangent naïf associé à la présentation
$\alpha : P \to S$ de $S/R$}. Notons que, si $P = R[S]$ muni de sa
surjection canonique sur $S$, nous retrouvons $\NL_{S/R}$.
Si $P = R[x_1, \ldots, x_n]$, nous utiliserons parfois la notation
$I/I^2 \to \bigoplus_{i = 1, \ldots, n} S\text{d}x_i$ pour désigner ce
complexe.

\medskip\noindent
Supposons donné un diagramme commutatif
\begin{equation}
\label{equation-functoriality-NL}
\vcenter{
\xymatrix{
S \ar[r]_{\phi} & S' \\
R \ar[r] \ar[u] & R' \ar[u]
}
}
\end{equation}
d'anneaux. Soit $\alpha : P \to S$ une présentation de $S$ sur $R$, et
soit $\alpha' : P' \to S'$ une présentation de $S'$ sur $R'$.
Un {\it morphisme de présentations de $\alpha : P \to S$ vers
$\alpha' : P' \to S'$} est, par définition, un morphisme de $R$-algèbres
$$
\varphi : P \to P'
$$
tel que $\phi \circ \alpha = \alpha' \circ \varphi$. Notons que, dans ce
cas, $\varphi(I) \subset I'$, où $I = \Ker(\alpha)$ et
$I' = \Ker(\alpha')$. Ainsi, $\varphi$ induit un morphisme de $S$-modules
$I/I^2 \to I'/(I')^2$ et, par fonctorialité des différentielles, également
un morphisme de $S$-modules
$\Omega_{P/R} \otimes S \to \Omega_{P'/R'} \otimes S'$.
Ces morphismes sont compatibles avec les différentielles de $\NL(\alpha)$
et $\NL(\alpha')$, et nous obtenons un morphisme de complexes cotangents
naïfs
$$
\NL(\alpha) \longrightarrow \NL(\alpha').
$$
Il est souvent commode de considérer le morphisme induit
$\NL(\alpha) \otimes_S S' \to \NL(\alpha')$.

\medskip\noindent
Dans le cas particulier où $P = R[S]$ et $P' = R'[S']$, le morphisme
$\phi : S \to S'$ induit un morphisme canonique d'anneaux
$\varphi : P \to P'$ par la règle $[s] \mapsto [\phi(s)]$.
La construction ci-dessus détermine donc des morphismes canoniques (!) de
complexes de chaînes
$$
\NL_{S/R} \longrightarrow \NL_{S'/R'},\quad\text{et}\quad
\NL_{S/R} \otimes_S S' \longrightarrow \NL_{S'/R'}
$$
associés au diagramme (\ref{equation-functoriality-NL}). Notons que cette
construction est compatible à la composition : étant donné un diagramme
commutatif
$$
\xymatrix{
S \ar[r]_{\phi} & S' \ar[r]_{\phi'} & S'' \\
R \ar[r] \ar[u] & R' \ar[u] \ar[r] & R'' \ar[u]
}
$$
nous voyons que la composée de
$$
\NL_{S/R} \longrightarrow \NL_{S'/R'} \longrightarrow \NL_{S''/R''}
$$
est le morphisme $\NL_{S/R} \to \NL_{S''/R''}$ donné par le carré extérieur.

\medskip\noindent
Il se trouve que $\NL(\alpha)$ est homotopiquement équivalent à $\NL_{S/R}$
et que les morphismes construits ci-dessus sont bien définis à homotopie
près (les homotopies de morphismes de complexes sont discutées dans
Homologie, section \ref{homology-section-complexes}, mais nous explicitons
également dans la démonstration du lemme ci-dessous le sens exact de ces
énoncés).

\begin{lemma}
\label{lemma-NL-homotopy}
Supposons donné un diagramme (\ref{equation-functoriality-NL}).
Soient $\alpha : P \to S$ et $\alpha' : P' \to S'$ des présentations.
\begin{enumerate}
\item Il existe un morphisme de présentations de $\alpha$ vers $\alpha'$.
\item Deux morphismes de présentations quelconques induisent des morphismes
de complexes homotopes $\NL(\alpha) \to \NL(\alpha')$.
\item La construction est compatible aux compositions de morphismes de
présentations (voir la démonstration pour l'énoncé exact).
\item Si $R \to R'$ et $S \to S'$ sont des isomorphismes, alors, pour tout
morphisme $\varphi$ de présentations de $\alpha$ vers $\alpha'$, le morphisme
induit $\NL(\alpha) \to \NL(\alpha')$ est une équivalence d'homotopie et un
quasi-isomorphisme.
\end{enumerate}
En particulier, en comparant $\alpha$ à la présentation canonique
(\ref{equation-canonical-presentation}), nous concluons qu'il existe un
quasi-isomorphisme $\NL(\alpha) \to \NL_{S/R}$ bien défini à homotopie près
et compatible avec toutes les fonctorialités (à homotopie près).
\end{lemma}

\begin{proof}
Puisque $P$ est une algèbre de polynômes sur $R$, nous pouvons écrire
$P = R[x_a, a \in A]$ pour un certain ensemble $A$.
Comme $\alpha'$ est surjectif, nous pouvons choisir, pour tout $a \in A$,
un élément $f_a \in P'$ tel que
$\alpha'(f_a) = \phi(\alpha(x_a))$. Soit
$\varphi : P = R[x_a, a \in A] \to P'$ l'unique morphisme de $R$-algèbres
tel que $\varphi(x_a) = f_a$. Cela donne le morphisme de (1).

\medskip\noindent
Soient $\varphi$ et $\varphi'$ des morphismes de présentations de $\alpha$
vers $\alpha'$. Soient $I = \Ker(\alpha)$ et $I' = \Ker(\alpha')$.
Nous devons construire le morphisme diagonal $h$ dans le diagramme
$$
\xymatrix{
I/I^2 \ar[r]^-{\text{d}}
\ar@<1ex>[d]^{\varphi'_1} \ar@<-1ex>[d]_{\varphi_1}
&
\Omega_{P/R} \otimes_P S
\ar@<1ex>[d]^{\varphi'_0} \ar@<-1ex>[d]_{\varphi_0}
\ar[ld]_h
\\
I'/(I')^2 \ar[r]^-{\text{d}}
&
\Omega_{P'/R'} \otimes_{P'} S'
}
$$
où les morphismes verticaux sont induits par $\varphi$, $\varphi'$ et tels
que
$$
\varphi_1 - \varphi'_1 = h \circ \text{d}
\quad\text{et}\quad
\varphi_0 - \varphi'_0 = \text{d} \circ h
$$
Considérons l'application $\varphi - \varphi' : P \to P'$. Puisque
$\varphi$ et $\varphi'$ sont toutes deux compatibles avec $\alpha$ et
$\alpha'$, nous obtenons $\varphi - \varphi' : P \to I'$.
Cela implique que $\varphi, \varphi' : P \to P'$ induisent la même structure
de $P$-module sur $I'/(I')^2$, puisque
$\varphi(p)i' - \varphi'(p)i' = (\varphi - \varphi')(p)i' \in (I')^2$.
De plus, $\varphi - \varphi'$ est $R$-linéaire et
$$
(\varphi - \varphi')(fg) =
\varphi(f)(\varphi - \varphi')(g) + (\varphi - \varphi')(f)\varphi'(g)
$$
L'application induite $D : P \to I'/(I')^2$ est donc une $R$-dérivation.
Nous obtenons ainsi un morphisme canonique
$h : \Omega_{P/R} \otimes_P S \to I'/(I')^2$ tel que
$D = h \circ \text{d}$.
Un calcul (omis) montre que $h$ est l'homotopie souhaitée.

\medskip\noindent
Supposons que nous ayons un diagramme commutatif
$$
\xymatrix{
S \ar[r]_{\phi} & S' \ar[r]_{\phi'} & S'' \\
R \ar[r] \ar[u] & R' \ar[u] \ar[r] & R'' \ar[u]
}
$$
et que
\begin{enumerate}
\item $\alpha : P \to S$,
\item $\alpha' : P' \to S'$, et
\item $\alpha'' : P'' \to S''$
\end{enumerate}
soient des présentations. Supposons que
\begin{enumerate}
\item $\varphi : P \to P'$ soit un morphisme de présentations de
$\alpha$ vers $\alpha'$, et
\item $\varphi' : P' \to P''$ soit un morphisme de présentations de
$\alpha'$ vers $\alpha''$.
\end{enumerate}
Il est alors immédiat que
$\varphi' \circ \varphi : P \to P''$ est un morphisme de présentations de
$\alpha$ vers $\alpha''$ et que le morphisme induit
$\NL(\alpha) \to \NL(\alpha'')$ de complexes cotangents naïfs est la
composée des morphismes $\NL(\alpha) \to \NL(\alpha')$ et
$\NL(\alpha') \to \NL(\alpha'')$ induits par $\varphi$ et $\varphi'$.

\medskip\noindent
Dans le cas simple des complexes à deux termes, un quasi-isomorphisme est
simplement un morphisme qui induit un isomorphisme à la fois sur le conoyau
et sur le noyau des morphismes entre les termes. Notons que des morphismes
homotopes de complexes à deux termes (au sens expliqué ci-dessus)
définissent les mêmes morphismes sur le noyau et le conoyau. Ainsi, si
$\varphi$ est un morphisme d'une présentation $\alpha$ de $S$ sur $R$ vers
elle-même, alors le morphisme induit
$\NL(\alpha) \to \NL(\alpha)$ est un quasi-isomorphisme, car il est homotope
à l'identité d'après (2). Pour démontrer (4) en toute généralité,
considérons un morphisme $\varphi'$ de $\alpha'$ vers $\alpha$, qui existe
d'après (1). Les composées
$\NL(\alpha) \to \NL(\alpha') \to \NL(\alpha)$ et
$\NL(\alpha') \to \NL(\alpha) \to \NL(\alpha')$ sont homotopes aux
morphismes identités d'après (3) ; ces morphismes sont donc, par définition,
des équivalences d'homotopie. Il s'ensuit formellement que les deux
morphismes $\NL(\alpha) \to \NL(\alpha')$ et
$\NL(\alpha') \to \NL(\alpha)$ sont des quasi-isomorphismes.
Nous omettons certains détails.
\end{proof}

\begin{lemma}
\label{lemma-NL-polynomial-algebra}
Soit $A \to B$ une algèbre de polynômes. Alors $\NL_{B/A}$ est
homotopiquement équivalent au complexe de chaînes
$(0 \to \Omega_{B/A})$, où $\Omega_{B/A}$ est en degré $0$.
\end{lemma}

\begin{proof}
Cela découle du Lemme \ref{lemma-NL-homotopy} et du fait que
$\text{id}_B : B \to B$ est une présentation de $B$ sur $A$ de noyau nul.
\end{proof}

\noindent
Le lemme suivant fournit une partie de la motivation pour introduire le
complexe cotangent naïf. Le complexe cotangent étend ceci en une véritable
suite exacte longue de cohomologie. Si $B \to C$ est une intersection
complète locale, on peut prolonger la suite par un zéro à gauche ; voir
Compléments d'algèbre,
Lemme \ref{more-algebra-lemma-transitive-lci-at-end}.

\begin{lemma}[Suite de Jacobi-Zariski]
\label{lemma-exact-sequence-NL}
Soient $A \to B \to C$ des morphismes d'anneaux. Choisissons une
présentation $\alpha : A[x_s, s \in S] \to B$ de noyau $I$.
Choisissons une présentation $\beta : B[y_t, t \in T] \to C$ de noyau $J$.
Soit $\gamma : A[x_s, y_t] \to C$ la présentation induite de $C$, de noyau
$K$. Nous obtenons alors un diagramme commutatif canonique
$$
\xymatrix{
0 \ar[r] &
\Omega_{A[x_s]/A} \otimes C \ar[r] &
\Omega_{A[x_s, y_t]/A} \otimes C \ar[r] &
\Omega_{B[y_t]/B} \otimes C \ar[r] &
0 \\
&
I/I^2 \otimes C \ar[r] \ar[u] &
K/K^2 \ar[r] \ar[u] &
J/J^2 \ar[r] \ar[u] &
0
}
$$
à lignes exactes. Nous obtenons la suite exacte suivante de groupes
d'homologie
$$
H_1(\NL_{B/A} \otimes_B C) \to
H_1(L_{C/A}) \to
H_1(L_{C/B}) \to
C \otimes_B \Omega_{B/A} \to
\Omega_{C/A} \to
\Omega_{C/B} \to 0
$$
de $C$-modules, qui prolonge la suite du
Lemme \ref{lemma-exact-sequence-differentials}.
Si $\text{Tor}_1^B(\Omega_{B/A}, C) = 0$ et
$\text{Tor}_2^B(\Omega_{B/A}, C) = 0$, alors
$H_1(\NL_{B/A} \otimes_B C) = H_1(L_{B/A}) \otimes_B C$.
\end{lemma}

\begin{proof}
Nous omettons la définition précise des morphismes.
L'exactitude de la ligne supérieure découle du fait que les
$\text{d}x_s$, $\text{d}y_t$ forment une base du module du milieu.
Le morphisme $\gamma$ se factorise en
$$
A[x_s, y_t] \to B[y_t] \to C
$$
où la première flèche est surjective et la seconde est égale à $\beta$.
Nous voyons donc que $K \to J$ est surjectif.
En outre, le noyau de la première flèche affichée est
$IA[x_s, y_t]$. Par conséquent, $I/I^2 \otimes C$ se surjecte sur le noyau
de $K/K^2 \to J/J^2$. Enfin, nous pouvons utiliser le
Lemme \ref{lemma-NL-homotopy} pour identifier les termes aux groupes
d'homologie des complexes cotangents naïfs.

\medskip\noindent
La dernière assertion est un énoncé d'algèbre homologique.
Rappelons que $\NL_{B/A} = (N^{-1} \to N^0)$ est un complexe à deux termes
de $B$-modules, où $N^0$ est libre, et que ses modules de cohomologie sont
$H^0 = \Omega_{B/A}$ et $H^{-1} = H_1(L_{B/A})$.
Notons $M \subset N^0$ l'image de la différentielle.
Si $\text{Tor}_1^B(H^0, C) = 0$, alors nous avons une suite exacte
$$
0 \to M \otimes_B C \to N^0 \otimes_B C \to H^0 \otimes_B C \to 0
$$
Puisque $N^0$ est libre, nous voyons aussi que
$\text{Tor}_2^B(H^0, C) =
\text{Tor}_1^B(M, C)$. Ainsi, si $\text{Tor}_2^B(H^0, C) = 0$, nous avons
également une suite exacte
$$
0 \to H^{-1} \otimes_B C \to N^{-1} \otimes_B C \to M \otimes_B C \to 0
$$
En rassemblant ces faits, nous voyons que, si
$\text{Tor}_1^B(H^0, C) = 0$ et $\text{Tor}_2^B(H^0, C) = 0$, alors
$H^{-1} \otimes_B C$ est le noyau de
$N^{-1} \otimes_B C \to N^0 \otimes_B C$, comme souhaité.
\end{proof}

\begin{remark}
\label{remark-composition-homotopy-equivalent-to-zero}
Soient $A \to B$ et $\phi : B \to C$ des morphismes d'anneaux.
Alors la composée $\NL_{B/A} \to \NL_{C/A} \to \NL_{C/B}$ est
homotopiquement équivalente à zéro. En effet, cette composée est la
fonctorialité du complexe cotangent naïf pour le carré
$$
\xymatrix{
B \ar[r]_\phi & C \\
A \ar[r] \ar[u] & B \ar[u]
}
$$
Écrivons $J = \Ker(B[C] \to C)$. Une homotopie explicite est donnée par le
morphisme $\Omega_{A[B]/A} \otimes_{A[B]} B \to J/J^2$ qui envoie l'élément
de base $\text{d}[b]$ sur la classe de $[\phi(b)] - b$ dans $J/J^2$.
\end{remark}

\begin{lemma}
\label{lemma-NL-surjection}
Soit $A \to B$ un morphisme surjectif d'anneaux, de noyau $I$.
Alors $\NL_{B/A}$ est homotopiquement équivalent au complexe de chaînes
$(I/I^2 \to 0)$, où $I/I^2$ est en degré $1$. En particulier,
$H_1(L_{B/A}) = I/I^2$.
\end{lemma}

\begin{proof}
Cela découle du Lemme \ref{lemma-NL-homotopy} et du fait que $A \to B$ est
une présentation de $B$ sur $A$.
\end{proof}

\begin{lemma}
\label{lemma-application-NL}
Soient $A \to B \to C$ des morphismes d'anneaux. Supposons que $A \to C$
soit surjectif (donc $B \to C$ l'est aussi). Notons
$I = \Ker(A \to C)$ et $J = \Ker(B \to C)$. Alors la suite
$$
I/I^2 \to J/J^2 \to \Omega_{B/A} \otimes_B B/J \to 0
$$
est exacte.
\end{lemma}

\begin{proof}
Cela découle du Lemme \ref{lemma-exact-sequence-NL} et de la description des
complexes cotangents naïfs $\NL_{C/B}$ et $\NL_{C/A}$ dans le
Lemme \ref{lemma-NL-surjection}.
\end{proof}

\begin{lemma}[Changement de base plat]
\label{lemma-change-base-NL}
Soit $R \to S$ un morphisme d'anneaux. Soit $\alpha : P \to S$ une
présentation. Soit $R \to R'$ un morphisme plat d'anneaux.
Soit $\alpha' : P \otimes_R R' \to S' = S \otimes_R R'$ la présentation
induite. Alors
$\NL(\alpha) \otimes_R R' = \NL(\alpha) \otimes_S S' = \NL(\alpha')$.
En particulier, le morphisme canonique
$$
\NL_{S/R} \otimes_S S' \longrightarrow \NL_{S \otimes_R R'/R'}
$$
est une équivalence d'homotopie si $R \to R'$ est plat.
\end{lemma}

\begin{proof}
Cela est vrai parce que
$\Ker(\alpha') = R' \otimes_R \Ker(\alpha)$, puisque $R \to R'$ est plat.
\end{proof}

\begin{lemma}
\label{lemma-colimits-NL}
Soit $R_\lambda \to S_\lambda$ un système de morphismes d'anneaux indexé
par l'ensemble filtrant $\Lambda$.
Posons $R = \colim R_\lambda$ et $S = \colim S_\lambda$.
Alors $\NL_{S/R} = \colim \NL_{S_\lambda/R_\lambda}$.
\end{lemma}

\begin{proof}
Rappelons que $\NL_{S/R}$ est le complexe
$I/I^2 \to \bigoplus_{s \in S} S\text{d}[s]$, où $I \subset R[S]$ est le
noyau de la présentation canonique $R[S] \to S$.
Il est maintenant clair que $R[S] = \colim R_\lambda[S_\lambda]$ et, de
même, que $I = \colim I_\lambda$, où
$I_\lambda = \Ker(R_\lambda[S_\lambda] \to S_\lambda)$.
Le lemme est donc clair.
\end{proof}

\begin{lemma}
\label{lemma-NL-of-localization}
Si $S \subset A$ est une partie multiplicative de $A$, alors
$\NL_{S^{-1}A/A}$ est homotopiquement équivalent au complexe nul.
\end{lemma}

\begin{proof}
Puisque $A \to S^{-1}A$ est plat, nous voyons que
$\NL_{S^{-1}A/A} \otimes_A S^{-1}A \to \NL_{S^{-1}A/S^{-1}A}$ est une
équivalence d'homotopie par changement de base plat
(Lemme \ref{lemma-change-base-NL}). Comme la source de la flèche est
isomorphe à $\NL_{S^{-1}A/A}$ et que sa cible est nulle (d'après le
Lemme \ref{lemma-NL-surjection}), le résultat en découle.
\end{proof}

\begin{lemma}
\label{lemma-NL-localize-bottom}
Soit $S \subset A$ une partie multiplicative de $A$.
Soit $S^{-1}A \to B$ un morphisme d'anneaux.
Alors $\NL_{B/A} \to \NL_{B/S^{-1}A}$ est une équivalence d'homotopie.
\end{lemma}

\begin{proof}
Choisissons une présentation $\alpha : P \to B$ de $B$ sur $A$.
Alors $\beta : S^{-1}P \to B$ est une présentation de $B$ sur $S^{-1}A$.
Un calcul direct montre que $\NL(\alpha) = \NL(\beta)$, ce qui démontre le
lemme puisque le complexe cotangent naïf est bien défini à homotopie près
d'après le Lemme \ref{lemma-NL-homotopy}.
\end{proof}

\begin{lemma}
\label{lemma-principal-localization-NL}
\begin{slogan}
La formation du complexe cotangent naïf commute à la localisation en un
élément.
\end{slogan}
Soit $A \to B$ un morphisme d'anneaux. Soit $g \in B$.
Supposons que $\alpha : P \to B$ soit une présentation de noyau $I$.
Une présentation de $B_g$ sur $A$ est alors le morphisme
$$
\beta : P[x] \longrightarrow B_g
$$
qui prolonge $\alpha$ et envoie $x$ sur $1/g$.
Le noyau $J$ de $\beta$ est engendré par $I$ et par l'élément $f x - 1$,
où $f \in P$ est un élément envoyé sur $g \in B$ par $\alpha$.
Dans cette situation, on a
\begin{enumerate}
\item $J/J^2 = (I/I^2)_g \oplus B_g (f x - 1)$,
\item $\Omega_{P[x]/A} \otimes_{P[x]} B_g =
\Omega_{P/A} \otimes_P B_g \oplus B_g \text{d}x$,
\item $\NL(\beta) \cong
\NL(\alpha) \otimes_B B_g \oplus (B_g \xrightarrow{g} B_g)$.
\end{enumerate}
Le morphisme canonique $\NL_{B/A} \otimes_B B_g \to \NL_{B_g/A}$ est donc
une équivalence d'homotopie.
\end{lemma}

\begin{proof}
Puisque $P[x]/(I, fx - 1) = B[x]/(gx - 1) = B_g$, nous obtenons l'assertion
selon laquelle $I$ et $fx - 1$ engendrent $J$. Considérons le diagramme
commutatif
$$
\xymatrix{
0 \ar[r] &
\Omega_{P/A} \otimes B_g \ar[r] &
\Omega_{P[x]/A} \otimes B_g \ar[r] &
\Omega_{B[x]/B} \otimes B_g \ar[r] &
0 \\
&
(I/I^2)_g \ar[r] \ar[u] &
J/J^2 \ar[r] \ar[u] &
(gx - 1)/(gx - 1)^2 \ar[r] \ar[u] &
0
}
$$
dont les lignes sont exactes d'après le
Lemme \ref{lemma-exact-sequence-NL}.
Le $B_g$-module $\Omega_{B[x]/B} \otimes B_g$ est libre de rang $1$ sur
$\text{d}x$. L'élément $\text{d}x$ du $B_g$-module
$\Omega_{P[x]/A} \otimes B_g$ fournit un scindage de la ligne supérieure.
L'élément $gx - 1 \in (gx - 1)/(gx - 1)^2$ est envoyé sur
$g\text{d}x$ dans $\Omega_{B[x]/B} \otimes B_g$ ; par conséquent,
$(gx - 1)/(gx - 1)^2$ est libre de rang $1$ sur $B_g$.
(On peut aussi le voir en remarquant que $gx - 1$ est un non-diviseur de
zéro dans $B[x]$, parce que c'est un polynôme dont le terme constant est
inversible, et que tout non-diviseur de zéro donne une suite quasi-régulière
de longueur $1$ d'après le Lemme \ref{lemma-regular-quasi-regular}.)

\medskip\noindent
Démontrons que $(I/I^2)_g \to J/J^2$ est injectif. Considérons le morphisme
de $P$-algèbres
$$
\pi : P[x] \to (P/I^2)_f = P_f/I_f^2
$$
qui envoie $x$ sur $1/f$. Puisque $J$ est engendré par $I$ et $fx - 1$,
nous voyons que $\pi(J) \subset (I/I^2)_f = (I/I^2)_g$. Comme cet idéal est
de carré nul, on a $\pi(J^2) = 0$.
Si $a \in I$ est envoyé sur un élément de $J^2$ dans $J$, alors
$\pi(a) = 0$, ce qui implique que $a$ est envoyé sur zéro dans $I_f/I_f^2$.
Cela démontre l'injectivité souhaitée.

\medskip\noindent
Nous avons donc une suite exacte courte de complexes à deux termes
$$
0 \to \NL(\alpha) \otimes_B B_g \to \NL(\beta)
\to (B_g \xrightarrow{g} B_g) \to 0
$$
Une telle suite exacte courte peut toujours être scindée dans la catégorie
des complexes. Dans notre cas particulier, nous pouvons choisir les
scindages
$$
J/J^2 = (I/I^2)_g \oplus B_g (fx - 1)\quad\text{et}\quad
\Omega_{P[x]/A} \otimes B_g = \Omega_{P/A} \otimes B_g \oplus
B_g (g^{-2}\text{d}f + \text{d}x)
$$
Cela fonctionne parce que
$\text{d}(fx - 1) = x\text{d}f + f \text{d}x =
g(g^{-2}\text{d}f + \text{d}x)$ dans $\Omega_{P[x]/A} \otimes B_g$.
\end{proof}

\begin{lemma}
\label{lemma-localize-NL}
Soit $A \to B$ un morphisme d'anneaux. Soit $S \subset B$ une partie
multiplicative. Le morphisme canonique
$\NL_{B/A} \otimes_B S^{-1}B \to \NL_{S^{-1}B/A}$ est un quasi-isomorphisme.
\end{lemma}

\begin{proof}
On a $S^{-1}B = \colim_{g \in S} B_g$, où nous considérons $S$ comme un
ensemble filtrant (ordonné par divisibilité) ; voir le
Lemme \ref{lemma-localization-colimit}.
D'après le Lemme \ref{lemma-principal-localization-NL}, chacun des morphismes
$\NL_{B/A} \otimes_B B_g \to \NL_{B_g/A}$ est un quasi-isomorphisme.
Le lemme découle du Lemme \ref{lemma-colimits-NL}.
\end{proof}

\begin{lemma}
\label{lemma-sum-two-terms}
Soit $R$ un anneau.
Soient $A_1 \to A_0$ et $B_1 \to B_0$ des complexes à deux termes.
Supposons qu'il existe des morphismes de complexes
$\varphi : A_\bullet \to B_\bullet$ et
$\psi : B_\bullet \to A_\bullet$ tels que
$\varphi \circ \psi$ et $\psi \circ \varphi$ soient homotopes aux
morphismes identités. Alors
$A_1 \oplus B_0 \cong B_1 \oplus A_0$ comme $R$-modules.
\end{lemma}

\begin{proof}
Choisissons une application $h : A_0 \to A_1$ telle que
$$
\text{id}_{A_1} - \psi_1 \circ \varphi_1 = h \circ d_A
\text{ et }
\text{id}_{A_0} - \psi_0 \circ \varphi_0 = d_A \circ h.
$$
De même, choisissons une application $h' : B_0 \to B_1$ telle que
$$
\text{id}_{B_1} - \varphi_1 \circ \psi_1 = h' \circ d_B
\text{ et }
\text{id}_{B_0} - \varphi_0 \circ \psi_0 = d_B \circ h'.
$$
Un calcul immédiat montre que
$$
\left(
\begin{matrix}
\text{id}_{A_1} & -h' \circ \psi_1 + h \circ \psi_0 \\
0 & \text{id}_{B_0}
\end{matrix}
\right)
=
\left(
\begin{matrix}
\psi_1 & h \\
-d_B & \varphi_0
\end{matrix}
\right)
\left(
\begin{matrix}
\varphi_1 & - h' \\
d_A & \psi_0
\end{matrix}
\right)
$$
Cela montre que les deux matrices du membre de droite sont inversibles et
démontre le lemme.
\end{proof}

\begin{lemma}
\label{lemma-conormal-module}
Soit $R \to S$ un morphisme d'anneaux de type fini.
Pour toutes présentations $\alpha : R[x_1, \ldots, x_n] \to S$ et
$\beta : R[y_1, \ldots, y_m] \to S$, on a
$$
I/I^2 \oplus S^{\oplus m} \cong J/J^2 \oplus S^{\oplus n}
$$
comme $S$-modules, où $I = \Ker(\alpha)$ et $J = \Ker(\beta)$.
\end{lemma}

\begin{proof}
Voir les Lemmes \ref{lemma-NL-homotopy} et \ref{lemma-sum-two-terms}.
\end{proof}

\begin{lemma}
\label{lemma-conormal-module-localize}
Soit $R \to S$ un morphisme d'anneaux de type fini.
Soit $g \in S$. Pour toutes présentations
$\alpha : R[x_1, \ldots, x_n] \to S$ et
$\beta : R[y_1, \ldots, y_m] \to S_g$, on a
$$
(I/I^2)_g \oplus S^{\oplus m}_g \cong J/J^2 \oplus S_g^{\oplus n}
$$
comme $S_g$-modules, où $I = \Ker(\alpha)$ et $J = \Ker(\beta)$.
\end{lemma}

\begin{proof}
Soit $\beta' : R[x_1, \ldots, x_n, x] \to S_g$ la présentation du
Lemme \ref{lemma-principal-localization-NL} construite à partir de $\alpha$.
Nous savons alors que $\NL(\alpha) \otimes_S S_g$ est homotopiquement
équivalent à $\NL(\beta')$.
Nous savons que $\NL(\beta)$ et $\NL(\beta')$ sont homotopiquement
équivalents d'après le Lemme \ref{lemma-NL-homotopy}. Nous concluons que
$\NL(\alpha) \otimes_S S_g$ est homotopiquement équivalent à
$\NL(\beta)$. Enfin, nous appliquons le
Lemme \ref{lemma-conormal-module}.
\end{proof}


\section{Intersections complètes locales}
\label{section-lci}

\noindent
Être une intersection complète locale est une propriété intrinsèque d'un
anneau local noethérien. Cela sera discuté dans
Algèbre à puissances divisées, section \ref{dpa-section-lci}.
Cependant, pour le moment, nous définissons seulement cette propriété pour
les algèbres de type fini sur un corps.

\begin{definition}
\label{definition-lci-field}
Soit $k$ un corps.
Soit $S$ une $k$-algèbre de type fini.
\begin{enumerate}
\item Nous disons que $S$ est une {\it intersection complète globale sur
$k$} s'il existe une présentation
$S = k[x_1, \ldots, x_n]/(f_1, \ldots, f_c)$ telle que
$\dim(S) = n - c$.
\item Nous disons que $S$ est une {\it intersection complète locale sur
$k$} s'il existe un recouvrement $\Spec(S) = \bigcup D(g_i)$ tel que chacun
des anneaux $S_{g_i}$ soit une intersection complète globale sur $k$.
\end{enumerate}
Nous adopterons aussi la convention que l'anneau nul est une intersection
complète globale sur $k$.
\end{definition}

\noindent
Supposons que $S$ soit une intersection complète globale
$S = k[x_1, \ldots, x_n]/(f_1, \ldots, f_c)$ comme dans la
Définition \ref{definition-lci-field}.
Pour un idéal maximal $\mathfrak m \subset k[x_1, \ldots, x_n]$, on a
$\dim(k[x_1, \ldots, x_n]_\mathfrak m) = n$
(Lemme \ref{lemma-dim-affine-space}).
Si $(f_1, \ldots, f_c) \subset \mathfrak m$, alors le
Lemme \ref{lemma-one-equation} donne $\dim(S_\mathfrak m) \geq n - c$.
Comme $\dim(S) = n - c$ d'après la Définition \ref{definition-lci-field},
nous concluons que $\dim(S_\mathfrak m) = n - c$ pour tout idéal maximal de
$S$ et que $\Spec(S)$ est équidimensionnel
(Topologie, Définition \ref{topology-definition-equidimensional}) de
dimension $n - c$ ; voir le
Lemme \ref{lemma-dimension-at-a-point-finite-type-over-field}.
Nous utiliserons souvent ce fait sans autre mention.

\begin{lemma}
\label{lemma-localize-lci}
Soit $k$ un corps.
Soit $S$ une $k$-algèbre de type fini.
Soit $g \in S$.
\begin{enumerate}
\item Si $S$ est une intersection complète globale, alors $S_g$ l'est aussi.
\item Si $S$ est une intersection complète locale, alors $S_g$ l'est aussi.
\end{enumerate}
\end{lemma}

\begin{proof}
La seconde assertion découle immédiatement de la première.
Démontrons la première. Si $S_g$ est l'anneau nul, elle est vraie.
Supposons $S_g$ non nul.
Écrivons $S = k[x_1, \ldots, x_n]/(f_1, \ldots, f_c)$ avec
$n - c = \dim(S)$ comme dans la Définition \ref{definition-lci-field}.
D'après les remarques qui suivent la définition,
$\dim(S_g) = n - c$. Soit $g' \in k[x_1, \ldots, x_n]$ un élément dont la
classe résiduelle correspond à $g$. Alors
$S_g =  k[x_1, \ldots, x_n, x_{n + 1}]/(f_1, \ldots, f_c, x_{n + 1}g' - 1)$,
comme souhaité.
\end{proof}

\begin{lemma}
\label{lemma-lci-CM}
Soit $k$ un corps. Soit $S$ une $k$-algèbre de type fini.
Si $S$ est une intersection complète locale, alors $S$ est un anneau de
Cohen-Macaulay.
\end{lemma}

\begin{proof}
Choisissons un idéal premier maximal $\mathfrak m$ de $S$.
Nous devons montrer que $S_\mathfrak m$ est de Cohen-Macaulay.
Par hypothèse, nous pouvons supposer
$S = k[x_1, \ldots, x_n]/(f_1, \ldots, f_c)$ avec
$\dim(S) = n - c$. Soit
$\mathfrak m' \subset k[x_1, \ldots, x_n]$ l'idéal maximal correspondant à
$\mathfrak m$. D'après la
Proposition \ref{proposition-finite-gl-dim-polynomial-ring}, l'anneau local
$k[x_1, \ldots, x_n]_{\mathfrak m'}$ est régulier local de dimension $n$.
En particulier, il est de Cohen-Macaulay d'après le
Lemme \ref{lemma-regular-ring-CM}.
En appliquant $c$ fois le Lemme \ref{lemma-one-equation}, l'anneau local
$S_{\mathfrak m} = k[x_1, \ldots, x_n]_{\mathfrak m'}/(f_1, \ldots, f_c)$
est de dimension $\geq n - c$. Par hypothèse,
$\dim(S_{\mathfrak m}) \leq n - c$. Nous obtenons donc l'égalité.
Cela implique que $f_1, \ldots, f_c$ est une suite régulière dans
$k[x_1, \ldots, x_n]_{\mathfrak m'}$ et que $S_{\mathfrak m}$ est de
Cohen-Macaulay ; voir la Proposition \ref{proposition-CM-module}.
\end{proof}

\noindent
Le lemme suivant est la clef technique du reste de cette section. Un aspect
important de ce lemme est que nous pouvons choisir n'importe quelle
présentation de l'anneau $S$, mais que la condition (1) ne dépend pas de ce
choix.

\begin{lemma}
\label{lemma-lci}
Soit $k$ un corps.
Soit $S$ une $k$-algèbre de type fini.
Soit $\mathfrak q$ un idéal premier de $S$.
Choisissons une présentation quelconque $S = k[x_1, \ldots, x_n]/I$.
Soit $\mathfrak q'$ l'idéal premier de $k[x_1, \ldots, x_n]$ correspondant
à $\mathfrak q$. Posons
$c = \text{height}(\mathfrak q') - \text{height}(\mathfrak q)$ ; autrement
dit, $\dim_{\mathfrak q}(S) = n - c$
(voir le Lemme \ref{lemma-codimension}). Les assertions suivantes sont
équivalentes :
\begin{enumerate}
\item Il existe $g \in S$, $g \not \in \mathfrak q$, tel que $S_g$ soit une
intersection complète globale sur $k$.
\item L'idéal $I_{\mathfrak q'} \subset
k[x_1, \ldots, x_n]_{\mathfrak q'}$ peut être engendré par $c$ éléments.
\item Le module conormal $(I/I^2)_{\mathfrak q}$ peut être engendré par $c$
éléments sur $S_{\mathfrak q}$.
\item Le module conormal $(I/I^2)_{\mathfrak q}$ est un
$S_{\mathfrak q}$-module libre de rang $c$.
\item L'idéal $I_{\mathfrak q'}$ peut être engendré par une suite régulière
dans l'anneau local régulier $k[x_1, \ldots, x_n]_{\mathfrak q'}$.
\end{enumerate}
Dans ce cas, tout choix de $c$ éléments de $I_{\mathfrak q'}$ qui engendrent
$I_{\mathfrak q'}/\mathfrak q'I_{\mathfrak q'}$ forme une suite régulière
dans l'anneau local $k[x_1, \ldots, x_n]_{\mathfrak q'}$.
\end{lemma}

\begin{proof}
Posons $R = k[x_1, \ldots, x_n]_{\mathfrak q'}$. C'est un anneau local de
Cohen-Macaulay de dimension $\text{height}(\mathfrak q')$ ; voir par exemple
le Lemme \ref{lemma-lci-CM}. En outre,
$\overline{R} = R/IR = R/I_{\mathfrak q'} = S_{\mathfrak q}$ est un quotient
de dimension $\text{height}(\mathfrak q)$.
Soient $f_1, \ldots, f_c \in I_{\mathfrak q'}$ des éléments qui engendrent
$(I/I^2)_{\mathfrak q}$. D'après le Lemme \ref{lemma-NAK},
$f_1, \ldots, f_c$ engendrent $I_{\mathfrak q'}$.
Comme les dimensions concordent, nous concluons par la
Proposition \ref{proposition-CM-module} que $f_1, \ldots, f_c$ est une suite
régulière dans $R$.
D'après le Lemme \ref{lemma-regular-quasi-regular},
$(I/I^2)_{\mathfrak q}$ est libre.
Ces arguments montrent que (2), (3) et (4) sont équivalentes, qu'elles
impliquent la dernière assertion du lemme et, par conséquent, qu'elles
impliquent (5).

\medskip\noindent
Si (5) est satisfaite, disons que $I_{\mathfrak q'}$ est engendré par une
suite régulière de longueur $e$. Alors, par la théorie de la dimension,
$\text{height}(\mathfrak q) = \dim(S_{\mathfrak q}) =
\dim(k[x_1, \ldots, x_n]_{\mathfrak q'}) - e =
\text{height}(\mathfrak q') - e$ ; voir la
section \ref{section-dimension}. Nous concluons que $e = c$.
Ainsi, (5) implique (2).

\medskip\noindent
Conservons les notations introduites dans le premier paragraphe.
Pour chaque $f_i$, nous pouvons trouver
$d_i \in k[x_1, \ldots, x_n]$, $d_i \not \in \mathfrak q'$, tel que
$f_i' = d_i f_i \in k[x_1, \ldots, x_n]$.
On a encore $I_{\mathfrak q'} = (f_1', \ldots, f_c')R$.
Il existe donc $g' \in k[x_1, \ldots, x_n]$,
$g' \not \in \mathfrak q'$, tel que $I_{g'} = (f_1', \ldots, f_c')$.
Choisissons en outre $g'' \in k[x_1, \ldots, x_n]$,
$g'' \not \in \mathfrak q'$, tel que
$\dim(S_{g''}) = \dim_{\mathfrak q} \Spec(S)$.
D'après le Lemme \ref{lemma-codimension}, cette dimension est égale à
$n - c$. Enfin, soit $g$ l'image de $g'g''$ dans $S$.
Nous voyons alors que
$$
S_g \cong k[x_1, \ldots, x_n, x_{n + 1}]
/
(f_1', \ldots, f_c', x_{n + 1}g'g'' - 1)
$$
et, par notre choix de $g''$, cet anneau est de dimension $n - c$.
C'est donc une intersection complète globale.
Ainsi, chacune des assertions (2), (3) et (4) implique (1).

\medskip\noindent
Supposons (1). Soit
$S_g \cong k[y_1, \ldots, y_m]/(f_1, \ldots, f_t)$ une présentation de
$S_g$ comme intersection complète globale.
Écrivons $J = (f_1, \ldots, f_t)$. Soit
$\mathfrak q'' \subset k[y_1, \ldots, y_m]$ l'idéal premier correspondant à
$\mathfrak qS_g$. Notons que
$t = m - \dim(S_g) =
\text{height}(\mathfrak q'') - \text{height}(\mathfrak q)$ ; voir le
Lemme \ref{lemma-codimension} pour la dernière égalité.
Comme on l'a vu dans la démonstration du Lemme \ref{lemma-lci-CM} (et aussi
ci-dessus), les éléments $f_1, \ldots, f_t$ forment une suite régulière dans
l'anneau local $k[y_1, \ldots, y_m]_{\mathfrak q''}$.
D'après le Lemme \ref{lemma-regular-quasi-regular},
$(J/J^2)_{\mathfrak q}$ est libre de rang $t$.
D'après le Lemme \ref{lemma-conormal-module-localize}, on a
$$
J/J^2 \oplus S_g^n \cong (I/I^2)_g \oplus S_g^m
$$
Ainsi, $(I/I^2)_{\mathfrak q}$ est libre de rang
$t + n - m = m - \dim(S_g) + n - m = n - \dim(S_g) =
\text{height}(\mathfrak q') - \text{height}(\mathfrak q) = c$.
Nous obtenons donc (4).
\end{proof}

\noindent
Le résultat du Lemme \ref{lemma-lci} suggère la définition suivante.

\begin{definition}
\label{definition-lci-local-ring}
Soit $k$ un corps. Soit $S$ une $k$-algèbre locale essentiellement de type
fini sur $k$. Nous disons que $S$ est une {\it intersection complète (sur
$k$)} s'il existe une $k$-algèbre locale $R$ et des éléments
$f_1, \ldots, f_c \in \mathfrak m_R$ tels que
\begin{enumerate}
\item $R$ soit essentiellement de type fini sur $k$,
\item $R$ soit un anneau local régulier,
\item $f_1, \ldots, f_c$ forment une suite régulière dans $R$, et
\item $S \cong R/(f_1, \ldots, f_c)$ comme $k$-algèbres.
\end{enumerate}
\end{definition}

\noindent
D'après le théorème de structure de Cohen (voir le
Théorème \ref{theorem-cohen-structure-theorem}), tout anneau local
noethérien complet s'écrit comme quotient d'un anneau local régulier complet.
Nous pouvons donc utiliser la définition ci-dessus pour définir la notion
d'anneau d'intersection complète pour tout anneau local noethérien complet.
Nous en discuterons dans Algèbre à puissances divisées,
section \ref{dpa-section-lci}. En attendant, le lemme suivant montre qu'une
telle définition a un sens.

\begin{lemma}
\label{lemma-ci-well-defined}
Soient $A \to B \to C$ des homomorphismes locaux surjectifs d'anneaux.
Supposons que $A$ et $B$ soient des anneaux locaux réguliers. Les assertions
suivantes sont équivalentes :
\begin{enumerate}
\item $\Ker(A \to C)$ est engendré par une suite régulière,
\item $\Ker(A \to C)$ est engendré par $\dim(A) - \dim(C)$ éléments,
\item $\Ker(B \to C)$ est engendré par une suite régulière, et
\item $\Ker(B \to C)$ est engendré par $\dim(B) - \dim(C)$ éléments.
\end{enumerate}
\end{lemma}

\begin{proof}
Un anneau local régulier est de Cohen-Macaulay ; voir le
Lemme \ref{lemma-regular-ring-CM}. D'où les équivalences
(1) $\Leftrightarrow$ (2) et (3) $\Leftrightarrow$ (4) ; voir la
Proposition \ref{proposition-CM-module}.
D'après le Lemme \ref{lemma-regular-quotient-regular}, l'idéal
$\Ker(A \to B)$ peut être engendré par $\dim(A) - \dim(B)$ éléments.
Nous voyons donc que (4) implique (2).

\medskip\noindent
Il reste à montrer que (1) implique (4). Raisonnons par récurrence sur
$\dim(A) - \dim(B)$. Le cas $\dim(A) - \dim(B) = 0$ est trivial.
Supposons $\dim(A) > \dim (B)$.
Écrivons $I = \Ker(A \to C)$ et $J = \Ker(A \to B)$.
Notons que $J \subset I$. Notre hypothèse dit que le nombre minimal de
générateurs de $I$ est $\dim(A) - \dim(C)$.
Soit $\mathfrak m \subset A$ l'idéal maximal. Considérons les applications
$$
J/ \mathfrak m J \to  I / \mathfrak m I \to \mathfrak m /\mathfrak m^2
$$
D'après le Lemme \ref{lemma-regular-quotient-regular} et sa démonstration,
la composée est injective. Prenons un élément $x \in J$ qui ne soit pas nul
dans $J /\mathfrak mJ$. D'après ce qui précède et le lemme de Nakayama,
$x$ appartient à un système minimal de générateurs de $I$.
Nous pouvons donc remplacer $A$ par $A/xA$ et $I$ par $I/xA$, ce qui diminue
à la fois $\dim(A)$ et le nombre minimal de générateurs de $I$ de $1$.
Le résultat en découle.
\end{proof}

\begin{lemma}
\label{lemma-lci-local}
Soit $k$ un corps. Soit $S$ une $k$-algèbre locale essentiellement de type
fini sur $k$. Les assertions suivantes sont équivalentes :
\begin{enumerate}
\item $S$ est une intersection complète sur $k$,
\item pour toute surjection $R \to S$, où $R$ est un anneau local régulier
essentiellement de présentation finie sur $k$, l'idéal $\Ker(R \to S)$ peut
être engendré par une suite régulière,
\item pour une certaine surjection $R \to S$, où $R$ est un anneau local
régulier essentiellement de présentation finie sur $k$, l'idéal
$\Ker(R \to S)$ peut être engendré par $\dim(R) - \dim(S)$ éléments,
\item il existe une intersection complète globale $A$ sur $k$ et un idéal
premier $\mathfrak a$ de $A$ tels que $S \cong A_{\mathfrak a}$, et
\item il existe une intersection complète locale $A$ sur $k$ et un idéal
premier $\mathfrak a$ de $A$ tels que $S \cong A_{\mathfrak a}$.
\end{enumerate}
\end{lemma}

\begin{proof}
Il est clair que (2) implique (1), que (1) implique (3), et aussi que (4)
implique (5). Montrons que (3) implique (4). Nous supposons donc qu'il existe
une surjection $R \to S$, où $R$ est un anneau local régulier essentiellement
de présentation finie sur $k$, telle que l'idéal $\Ker(R \to S)$ puisse être
engendré par $\dim(R) - \dim(S)$ éléments.
Nous pouvons écrire $R = (k[x_1, \ldots, x_n]/J)_{\mathfrak q}$ pour un
certain $J \subset k[x_1, \ldots, x_n]$ et un certain idéal premier
$\mathfrak q \subset k[x_1, \ldots, x_n]$ tel que $J \subset \mathfrak q$.
Soit $I \subset k[x_1, \ldots, x_n]$ le noyau du morphisme
$k[x_1, \ldots, x_n] \to S$, de sorte que
$S \cong (k[x_1, \ldots, x_n]/I)_{\mathfrak q}$.
Par hypothèse, $(I/J)_{\mathfrak q}$ est engendré par
$\dim(R) - \dim(S)$ éléments. D'après le
Lemme \ref{lemma-ci-well-defined}, nous concluons que
$I_{\mathfrak q}$ peut être engendré par
$\dim(k[x_1, \ldots, x_n]_{\mathfrak q}) - \dim(S)$ éléments.
Le Lemme \ref{lemma-lci} montre que, pour un certain
$g \in k[x_1, \ldots, x_n]$, $g \not \in \mathfrak q$, l'algèbre
$(k[x_1, \ldots, x_n]/I)_g$ est une intersection complète globale et que
$S$ est isomorphe à l'un de ses anneaux locaux.

\medskip\noindent
Pour achever la démonstration, nous devons montrer que (5) implique (2).
Supposons (5) et soit $\pi : R \to S$ une surjection, où $R$ est une
$k$-algèbre locale régulière essentiellement de type fini sur $k$.
Par hypothèse, $S = A_{\mathfrak a}$ pour une certaine intersection complète
locale $A$ sur $k$.
Choisissons une présentation
$R = (k[y_1, \ldots, y_m]/J)_{\mathfrak q}$ avec
$J \subset \mathfrak q \subset k[y_1, \ldots, y_m]$.
Nous pouvons supposer, et supposons, que $J$ est le noyau du morphisme
$k[y_1, \ldots, y_m] \to R$. Soit
$I \subset k[y_1, \ldots, y_m]$ le noyau du morphisme
$k[y_1, \ldots, y_m] \to S = A_{\mathfrak a}$.
Alors $J \subset I$ et $(I/J)_{\mathfrak q}$ est le noyau de la surjection
$\pi : R \to S$. Ainsi,
$S = (k[y_1, \ldots, y_m]/I)_{\mathfrak q}$.

\medskip\noindent
D'après le Lemme \ref{lemma-isomorphic-local-rings}, il existe
$g \in A$, $g \not \in \mathfrak a$, et
$g' \in k[y_1, \ldots, y_m]$, $g' \not \in \mathfrak q$, tels que
$A_g \cong (k[y_1, \ldots, y_m]/I)_{g'}$.
Après avoir remplacé $A$ par $A_g$ et $k[y_1, \ldots, y_m]$ par
$k[y_1, \ldots, y_{m + 1}]$, nous pouvons supposer que
$A \cong k[y_1, \ldots, y_m]/I$. Considérons les morphismes surjectifs
d'anneaux locaux
$$
k[y_1, \ldots, y_m]_{\mathfrak q} \to R \to S.
$$
Nous devons montrer que le noyau de $R \to S$ est engendré par une suite
régulière. D'après le Lemme \ref{lemma-lci},
$k[y_1, \ldots, y_m]_{\mathfrak q} \to A_{\mathfrak a} = S$ possède cette
propriété (car $A$ est une intersection complète locale sur $k$).
Le résultat découle du Lemme \ref{lemma-ci-well-defined}.
\end{proof}

\begin{lemma}
\label{lemma-lci-at-prime}
Soit $k$ un corps. Soit $S$ une $k$-algèbre de type fini.
Soit $\mathfrak q$ un idéal premier de $S$. Les assertions suivantes sont
équivalentes :
\begin{enumerate}
\item L'anneau local $S_{\mathfrak q}$ est un anneau d'intersection complète
(Définition \ref{definition-lci-local-ring}).
\item Il existe $g \in S$, $g \not \in \mathfrak q$, tel que $S_g$ soit une
intersection complète locale sur $k$.
\item Il existe $g \in S$, $g \not \in \mathfrak q$, tel que $S_g$ soit une
intersection complète globale sur $k$.
\item Pour toute présentation $S = k[x_1, \ldots, x_n]/I$, où
$\mathfrak q' \subset k[x_1, \ldots, x_n]$ correspond à $\mathfrak q$,
l'une quelconque des conditions équivalentes (1) -- (5) du
Lemme \ref{lemma-lci} est satisfaite.
\end{enumerate}
\end{lemma}

\begin{proof}
C'est une combinaison des Lemmes \ref{lemma-lci} et
\ref{lemma-lci-local}, ainsi que des définitions.
\end{proof}

\begin{lemma}
\label{lemma-lci-global}
Soit $k$ un corps. Soit $S$ une $k$-algèbre de type fini.
Les assertions suivantes sont équivalentes :
\begin{enumerate}
\item L'anneau $S$ est une intersection complète locale sur $k$.
\item Tous les anneaux locaux de $S$ sont des anneaux d'intersection
complète sur $k$.
\item Toutes les localisations de $S$ en des idéaux maximaux sont des
anneaux d'intersection complète sur $k$.
\end{enumerate}
\end{lemma}

\begin{proof}
Cela découle du Lemme \ref{lemma-lci-at-prime}, du fait que $\Spec(S)$ est
quasi-compact et des définitions.
\end{proof}

\noindent
Le lemme suivant affirme que la propriété d'être une intersection complète
est préservée par changement du corps de base (en un sens fort).

\begin{lemma}
\label{lemma-lci-field-change-local}
Soit $K/k$ une extension de corps.
Soit $S$ une algèbre de type fini sur $k$.
Soit $\mathfrak q_K$ un idéal premier de $S_K = K \otimes_k S$, et soit
$\mathfrak q$ l'idéal premier correspondant de $S$.
Alors $S_{\mathfrak q}$ est une intersection complète sur $k$
(Définition \ref{definition-lci-local-ring}) si et seulement si
$(S_K)_{\mathfrak q_K}$ est une intersection complète sur $K$.
\end{lemma}

\begin{proof}
Choisissons une présentation $S = k[x_1, \ldots, x_n]/I$.
Elle donne une présentation $S_K = K[x_1, \ldots, x_n]/I_K$, où
$I_K = K \otimes_k I$. Soient
$\mathfrak q_K' \subset K[x_1, \ldots, x_n]$, respectivement
$\mathfrak q' \subset k[x_1, \ldots, x_n]$, les idéaux premiers
correspondants. Nous allons montrer que les conditions équivalentes du
Lemme \ref{lemma-lci} sont satisfaites pour le couple
$(S = k[x_1, \ldots, x_n]/I, \mathfrak q)$ si et seulement si elles le sont
pour le couple $(S_K = K[x_1, \ldots, x_n]/I_K, \mathfrak q_K)$.
Le lemme en découlera (voir le Lemme \ref{lemma-lci-at-prime}).

\medskip\noindent
D'après le Lemme \ref{lemma-dimension-at-a-point-preserved-field-extension},
on a $\dim_{\mathfrak q} S = \dim_{\mathfrak q_K} S_K$.
Ainsi, l'entier $c$ qui intervient dans le Lemme \ref{lemma-lci} est le même
pour le couple $(S = k[x_1, \ldots, x_n]/I, \mathfrak q)$ que pour le couple
$(S_K = K[x_1, \ldots, x_n]/I_K, \mathfrak q_K)$.
D'autre part, on a
\begin{eqnarray*}
I \otimes_{k[x_1, \ldots, x_n]} \kappa(\mathfrak q')
\otimes_{\kappa(\mathfrak q')} \kappa(\mathfrak q_K')
& = &
I \otimes_{k[x_1, \ldots, x_n]} \kappa(\mathfrak q_K') \\
& = &
I \otimes_{k[x_1, \ldots, x_n]} K[x_1, \ldots, x_n]
\otimes_{K[x_1, \ldots, x_n]} \kappa(\mathfrak q_K') \\
& = &
(K \otimes_k I) \otimes_{K[x_1, \ldots, x_n]} \kappa(\mathfrak q_K') \\
& = &
I_K \otimes_{K[x_1, \ldots, x_n]} \kappa(\mathfrak q'_K).
\end{eqnarray*}
Par conséquent,
$\dim_{\kappa(\mathfrak q')}
I \otimes_{k[x_1, \ldots, x_n]} \kappa(\mathfrak q') =
\dim_{\kappa(\mathfrak q'_K)}
I_K \otimes_{K[x_1, \ldots, x_n]} \kappa(\mathfrak q_K')$.
Il découle donc du Lemme de Nakayama \ref{lemma-NAK} que le nombre minimal
de générateurs de $I_{\mathfrak q'}$ est égal au nombre minimal de
générateurs de $(I_K)_{\mathfrak q'_K}$.
Le lemme découle alors de la caractérisation (2) du
Lemme \ref{lemma-lci}.
\end{proof}





\begin{lemma}
\label{lemma-lci-field-change}
Soit $k \to K$ une extension de corps.
Soit $S$ une $k$-algèbre de type fini.
Alors $S$ est une intersection complète locale sur $k$ si et seulement si
$S \otimes_k K$ est une intersection complète locale sur $K$.
\end{lemma}

\begin{proof}
Cela résulte de la combinaison des Lemmes
\ref{lemma-lci-global} et \ref{lemma-lci-field-change-local}.
Mais nous donnons également ici une autre démonstration
(fondée sur les mêmes principes).

\medskip\noindent
Posons $S' = S \otimes_k K$. Soit
$\alpha : k[x_1, \ldots, x_n] \to S$ une présentation de noyau $I$.
Soit $\alpha' : K[x_1, \ldots, x_n] \to S'$ la présentation induite,
de noyau $I'$.

\medskip\noindent
Supposons que $S$ soit une intersection complète locale.
Choisissons un idéal premier $\mathfrak q \subset S'$. Notons
$\mathfrak q'$ l'idéal premier correspondant de $K[x_1, \ldots, x_n]$,
$\mathfrak p$ l'idéal premier correspondant de $S$, et
$\mathfrak p'$ l'idéal premier correspondant de $k[x_1, \ldots, x_n]$.
Considérons le diagramme suivant d'anneaux locaux noethériens :
$$
\xymatrix{
S'_{\mathfrak q} &  K[x_1, \ldots, x_n]_{\mathfrak q'} \ar[l] \\
S_{\mathfrak p}\ar[u] &  k[x_1, \ldots, x_n]_{\mathfrak p'} \ar[u] \ar[l]
}
$$
D'après le Lemme \ref{lemma-lci}, nous savons que $S_{\mathfrak p}$
est défini par une certaine suite régulière $f_1, \ldots, f_c$ dans
$k[x_1, \ldots, x_n]_{\mathfrak p'}$. Puisque la flèche verticale de
droite est plate, nous voyons que les images de $f_1, \ldots, f_c$
forment une suite régulière dans $K[x_1, \ldots, x_n]_{\mathfrak q'}$.
Comme le produit tensoriel par $K$ sur $k$ est un foncteur exact, on a
$S'_{\mathfrak q} = K[x_1, \ldots, x_n]_{\mathfrak q'}/(f_1, \ldots, f_c)$.
En appliquant à nouveau le Lemme \ref{lemma-lci}, nous voyons donc que $S'$
est une intersection complète locale dans un voisinage de $\mathfrak q$.
Comme $\mathfrak q$ était arbitraire, nous voyons que $S'$ est une
intersection complète locale sur $K$.

\medskip\noindent
Supposons que $S'$ soit une intersection complète locale.
Choisissons un idéal maximal $\mathfrak m$ de $S$. Notons $\mathfrak m'$
l'idéal maximal correspondant de $k[x_1, \ldots, x_n]$.
Notons $\kappa = \kappa(\mathfrak m)$ le corps résiduel.
D'après la Remarque \ref{remark-fundamental-diagram}, les idéaux premiers de
$S'$ situés au-dessus de $\mathfrak m$ correspondent aux idéaux premiers de
$K \otimes_k \kappa$. D'après le théorème de Hilbert-Nullstellensatz
\ref{theorem-nullstellensatz}, on a $[\kappa : k] < \infty$.
Ainsi, $K \otimes_k \kappa$ est non nul et de dimension finie sur $K$.
$K \otimes_k \kappa$ possède donc un nombre fini $> 0$ d'idéaux premiers,
qui sont tous maximaux et dont chacun a un corps résiduel fini sur $K$
(voir la section \ref{section-artinian}).
Il existe donc un nombre fini $> 0$ d'idéaux premiers
$\mathfrak n \subset S'$ situés au-dessus de $\mathfrak m$ ; chacun d'eux
est maximal et possède un corps résiduel fini sur $K$. Choisissons-en un,
disons $\mathfrak n \subset S'$, et soit
$\mathfrak n' \subset K[x_1, \ldots, x_n]$ l'idéal premier correspondant de
$K[x_1, \ldots, x_n]$.
Notons que, puisque $V(\mathfrak mS')$ est fini, $\mathfrak n$ en est un
point fermé isolé, et nous en déduisons que $\mathfrak mS'_{\mathfrak n}$
est un idéal de définition de $S'_{\mathfrak n}$. Cela implique que
$\dim(S_{\mathfrak m}) = \dim(S'_{\mathfrak n})$, par exemple d'après le
Lemme \ref{lemma-dimension-base-fibre-equals-total}.
(On peut aussi le voir à l'aide du
Lemme \ref{lemma-dimension-at-a-point-preserved-field-extension}.)
Considérons le diagramme correspondant d'anneaux locaux noethériens :
$$
\xymatrix{
S'_{\mathfrak n} &  K[x_1, \ldots, x_n]_{\mathfrak n'} \ar[l] \\
S_{\mathfrak m}\ar[u] &  k[x_1, \ldots, x_n]_{\mathfrak m'} \ar[u] \ar[l]
}
$$
D'après le Lemme \ref{lemma-change-base-NL}, on a
$\NL(\alpha) \otimes_S S' = \NL(\alpha')$, et en particulier
$I'/(I')^2 = I/I^2 \otimes_S S'$. Par conséquent,
$(I/I^2)_{\mathfrak m} \otimes_{S_{\mathfrak m}} \kappa$
et
$(I'/(I')^2)_{\mathfrak n} \otimes_{S'_{\mathfrak n}} \kappa(\mathfrak n)$
ont la même dimension. Puisque $(I'/(I')^2)_{\mathfrak n}$
est libre de rang $n - \dim S'_{\mathfrak n}$, nous en déduisons que
$(I/I^2)_{\mathfrak m}$ peut être engendré par
$n - \dim S'_{\mathfrak n} = n - \dim S_{\mathfrak m}$ éléments.
D'après le Lemme \ref{lemma-lci}, nous voyons que $S$ est une intersection
complète locale dans un voisinage de $\mathfrak m$.
Comme $\mathfrak m$ était un idéal maximal quelconque, nous concluons que
$S$ est une intersection complète locale.
\end{proof}

\noindent
Nous terminons par un lemme que nous utiliserons plus tard pour montrer que,
étant donnés des morphismes d'anneaux $T \to A \to B$ tels que $B$ soit
syntomique sur $T$ et que $B$ soit syntomique sur $A$, alors $A$ est
syntomique sur $T$.

\begin{lemma}
\label{lemma-lci-permanence-initial}
Soit
$$
\xymatrix{
B & S \ar[l] \\
A \ar[u] & R \ar[l] \ar[u]
}
$$
un carré commutatif d'anneaux locaux. Supposons que
\begin{enumerate}
\item $R$ et $\overline{S} = S/\mathfrak m_R S$ soient des anneaux locaux
réguliers,
\item $A = R/I$ et $B = S/J$ pour certains idéaux $I$, $J$,
\item $J \subset S$ et
$\overline{J} = J/\mathfrak m_R \cap J \subset \overline{S}$
soient engendrés par des suites régulières, et
\item $A \to B$ et $R \to S$ soient plats.
\end{enumerate}
Alors $I$ est engendré par une suite régulière.
\end{lemma}

\begin{proof}
Posons $\overline{B} = B/\mathfrak m_RB = B/\mathfrak m_AB$, de sorte que
$\overline{B} = \overline{S}/\overline{J}$.
Soient $f_1, \ldots, f_{\overline{c}} \in J$ des éléments tels que
$\overline{f}_1, \ldots, \overline{f}_{\overline{c}} \in \overline{J}$
forment une suite régulière engendrant $\overline{J}$.
Notons que $\overline{c} = \dim(\overline{S}) - \dim(\overline{B})$ ;
voir le Lemme \ref{lemma-ci-well-defined}.
D'après le Lemme \ref{lemma-grothendieck-regular-sequence}, l'anneau
$S/(f_1, \ldots, f_{\overline{c}})$ est plat sur $R$.
Ainsi, $S/(f_1, \ldots, f_{\overline{c}}) + IS$ est plat sur $A$.
Le morphisme $S/(f_1, \ldots, f_{\overline{c}}) + IS \to B$ est donc une
surjection de $S/IS$-modules de type fini plats sur $A$, qui est un
isomorphisme modulo $\mathfrak m_A$ et, par conséquent, un isomorphisme
d'après le Lemme \ref{lemma-mod-injective}. Autrement dit,
$J = (f_1, \ldots, f_{\overline{c}}) + IS$.

\medskip\noindent
D'après le Lemme \ref{lemma-ci-well-defined}, à nouveau, l'idéal $J$ est
engendré par une suite régulière de $c = \dim(S) - \dim(B)$ éléments.
Ainsi, $J/\mathfrak m_SJ$ est un espace vectoriel de dimension $c$.
D'après la description de $J$ ci-dessus, il existe
$g_1, \ldots, g_{c - \overline{c}} \in I$ tels que $J$ soit engendré par
$f_1, \ldots, f_{\overline{c}}, g_1, \ldots, g_{c - \overline{c}}$
(utiliser le Lemme de Nakayama \ref{lemma-NAK}). Considérons l'anneau
$A' = R/(g_1, \ldots, g_{c - \overline{c}})$ et la surjection
$A' \to A$. Nous voyons d'après ce qui précède que
$B = S/(f_1, \ldots, f_{\overline{c}}, g_1, \ldots, g_{c - \overline{c}})$
est plat sur $A'$ (puisque $S/(f_1, \ldots, f_{\overline{c}})$ est plat
sur $R$). Ainsi, $A' \to B$ est injectif (car il est fidèlement plat ;
voir le Lemme \ref{lemma-local-flat-ff}).
Comme ce morphisme se factorise par $A$, on obtient $A' = A$.
Notons que $\dim(B) = \dim(A) + \dim(\overline{B})$, et que
$\dim(S) = \dim(R) + \dim(\overline{S})$ ; voir le
Lemme \ref{lemma-dimension-base-fibre-equals-total}.
Ainsi, $c - \overline{c} = \dim(R) -\dim(A)$ par un calcul élémentaire.
Par conséquent, $I = (g_1, \ldots, g_{c - \overline{c}})$ est engendré
par une suite régulière d'après le Lemme \ref{lemma-ci-well-defined}.
\end{proof}




\section{Morphismes syntomiques}
\label{section-syntomic}

\noindent
Les morphismes d'anneaux syntomiques sont les morphismes d'anneaux plats de
présentation finie dont toutes les fibres sont des intersections complètes
locales. Nous étudions les morphismes généraux d'intersection complète locale
dans Compléments d'algèbre,
section \ref{more-algebra-section-lci}.

\begin{definition}
\label{definition-lci}
Un morphisme d'anneaux $R \to S$ est dit {\it syntomique}, ou bien nous disons
que $S$ est une {\it intersection complète locale plate sur $R$}, s'il est
plat, de présentation finie, et si tous ses anneaux fibres
$S \otimes_R \kappa(\mathfrak p)$ sont des intersections complètes locales ;
voir la Définition \ref{definition-lci-field}.
\end{definition}

\noindent
Il est clair qu'une algèbre sur un corps est syntomique sur ce corps si et
seulement si elle est une intersection complète locale. Voici une propriété
agréable de cette définition.

\begin{lemma}
\label{lemma-syntomic-descends}
\begin{slogan}
La propriété d'être syntomique est locale pour la topologie fpqc sur la base.
\end{slogan}
Soit $R \to S$ un morphisme d'anneaux.
Soit $R \to R'$ un morphisme d'anneaux fidèlement plat.
Posons $S' = R'\otimes_R S$.
Alors $R \to S$ est syntomique si et seulement si $R' \to S'$ est syntomique.
\end{lemma}

\begin{proof}
D'après le Lemme \ref{lemma-finite-presentation-descends} et le
Lemme \ref{lemma-flatness-descends}, cela vaut pour la propriété d'être plat
et pour celle d'être de présentation finie.
Le morphisme $\Spec(R') \to \Spec(R)$ est surjectif ;
voir le Lemme \ref{lemma-ff-rings}. Il suffit donc de montrer, pour des idéaux
premiers $\mathfrak p' \subset R'$ situés au-dessus de
$\mathfrak p \subset R$, que $S \otimes_R \kappa(\mathfrak p)$ est une
intersection complète locale si et seulement si
$S' \otimes_{R'} \kappa(\mathfrak p')$ est une intersection complète locale.
Notons que
$S' \otimes_{R'} \kappa(\mathfrak p') =
S \otimes_R \kappa(\mathfrak p)
\otimes_{\kappa(\mathfrak p)} \kappa(\mathfrak p')$.
Le Lemme \ref{lemma-lci-field-change} s'applique donc.
\end{proof}

\begin{lemma}
\label{lemma-base-change-syntomic}
Tout changement de base d'un morphisme syntomique est syntomique.
\end{lemma}

\begin{proof}
Cela est vrai pour la platitude, pour la propriété d'être de présentation
finie et pour celle d'avoir des intersections complètes locales comme fibres,
d'après les Lemmes \ref{lemma-flat-base-change},
\ref{lemma-compose-finite-type} et \ref{lemma-lci-field-change}.
\end{proof}

\begin{lemma}
\label{lemma-local-syntomic}
Soit $R \to S$ un morphisme d'anneaux.
Supposons donnés des éléments $g_1, \ldots g_m \in S$ qui engendrent l'idéal
unité et tels que chacun des morphismes $R \to S_{g_i}$ soit syntomique.
Alors $R \to S$ est syntomique.
\end{lemma}

\begin{proof}
Cela est vrai pour la platitude et pour la propriété d'être de présentation
finie d'après les Lemmes \ref{lemma-flat-localization} et
\ref{lemma-cover-upstairs}.
La propriété d'avoir des anneaux fibres qui soient des intersections complètes
locales est locale sur $S$ par sa définition même ; voir la
Définition \ref{definition-lci-field}.
\end{proof}

\begin{definition}
\label{definition-relative-global-complete-intersection}
Soit $R \to S$ un morphisme d'anneaux. Nous disons que $R \to S$ est une
{\it intersection complète globale relative} s'il existe une présentation
$S = R[x_1, \ldots, x_n]/(f_1, \ldots, f_c)$ et si toute fibre non vide de
$\Spec(S) \to \Spec(R)$ est de dimension $n - c$.
Nous dirons ``soit
$S = R[x_1, \ldots, x_n]/(f_1, \ldots, f_c)$ une intersection complète
globale relative'' pour désigner cette situation.
\end{definition}

\noindent
Le lemme suivant est parfois utile pour trouver des présentations globales.

\begin{lemma}
\label{lemma-huber}
Soit $S$ une $R$-algèbre de présentation finie qui possède une présentation
$S = R[x_1, \ldots, x_n]/I$ telle que $I/I^2$ soit libre sur $S$. Alors
$S$ possède une présentation $S = R[y_1, \ldots, y_m]/(f_1, \ldots, f_c)$
telle que $(f_1, \ldots, f_c)/(f_1, \ldots, f_c)^2$ soit libre, avec pour
base les classes de $f_1, \ldots, f_c$.
\end{lemma}

\begin{proof}
Notons que $I$ est un idéal de type fini d'après le
Lemme \ref{lemma-finite-presentation-independent}.
Soient $f_1, \ldots, f_c \in I$ des éléments qui s'envoient sur une base de
$I/I^2$. D'après le lemme de Nakayama (Lemme \ref{lemma-NAK}),
il existe $g \in 1 + I$ tel que
$$
g \cdot I \subset (f_1, \ldots, f_c)
$$
et $I_g \cong (f_1, \ldots, f_c)_g$. Nous voyons donc que
$$
S \cong R[x_1, \ldots, x_n]/(f_1, \ldots, f_c)[1/g]
\cong R[x_1, \ldots, x_n, x_{n + 1}]/(f_1, \ldots, f_c, gx_{n + 1} - 1)
$$
comme souhaité. Il en résulte que
$f_1, \ldots, f_c,gx_{n + 1} - 1$ forment une base de
$(f_1, \ldots, f_c, gx_{n + 1} - 1)/(f_1, \ldots, f_c, gx_{n + 1} - 1)^2$,
par exemple en appliquant le
Lemme \ref{lemma-principal-localization-NL}.
\end{proof}



\begin{example}
\label{example-factor-polynomials}
Soient $n , m \geq 1$ des entiers. Considérons le morphisme d'anneaux
\begin{eqnarray*}
R = \mathbf{Z}[a_1, \ldots, a_{n + m}]
& \longrightarrow &
S = \mathbf{Z}[b_1, \ldots, b_n, c_1, \ldots, c_m] \\
a_1 & \longmapsto & b_1 + c_1 \\
a_2 & \longmapsto & b_2 + b_1 c_1 + c_2 \\
\ldots & \ldots & \ldots \\
a_{n + m} & \longmapsto & b_n c_m
\end{eqnarray*}
Autrement dit, c'est l'unique morphisme d'anneaux de polynômes indiqué tel
que la factorisation de polynômes
$$
x^{n + m} + a_1 x^{n + m - 1} + \ldots + a_{n + m}
=
(x^n + b_1 x^{n - 1} + \ldots + b_n)
(x^m + c_1 x^{m - 1} + \ldots + c_m)
$$
soit satisfaite. Notons que $S$ est engendré par $n + m$ éléments sur $R$
(à savoir les $b_i, c_j$) et qu'il y a $n + m$ équations
(à savoir $a_k = a_k(b_i, c_j)$). Pour montrer que $S$ est une intersection
complète globale relative sur $R$, il suffit de prouver que toutes les fibres
sont de dimension $0$.

\medskip\noindent
Pour cela, soit $R \to k$ un morphisme d'anneaux à valeurs dans un corps $k$.
Disons que $a_i$ s'envoie sur $\alpha_i \in k$.
Considérons l'anneau fibre $S_k = k \otimes_R S$. Soit $k \to K$ une
extension de corps. Se donner un morphisme de $k$-algèbres de $S_k \to K$
revient à trouver
$\beta_1, \ldots, \beta_n, \gamma_1, \ldots, \gamma_m \in K$ tels que
$$
x^{n + m} + \alpha_1 x^{n + m - 1} + \ldots + \alpha_{n + m}
=
(x^n + \beta_1 x^{n - 1} + \ldots + \beta_n)
(x^m + \gamma_1 x^{m - 1} + \ldots + \gamma_m).
$$
Il n'existe donc qu'un nombre fini de tels $n + m$-uplets dans $K$.
Cela prouve que toutes les fibres n'ont qu'un nombre fini de points fermés
(utiliser par exemple le théorème de Hilbert-Nullstellensatz pour voir qu'ils
correspondent tous à des solutions dans $\overline{k}$), et donc que
$R \to S$ est une intersection complète globale relative.

\medskip\noindent
Une autre manière de raisonner consiste à montrer que
$\mathbf{Z}[a_1, \ldots, a_{n + m}] \to
\mathbf{Z}[b_1, \ldots, b_n, c_1, \ldots, c_m]$ est en fait aussi un
morphisme d'anneaux {\it fini}. En effet, d'après le
Lemme \ref{lemma-polynomials-divide}, chacun des $b_i, c_j$ est entier sur
$R$, et donc $R \to S$ est fini d'après le
Lemme \ref{lemma-characterize-integral}.
\end{example}

\begin{example}
\label{example-roots-universal-polynomial}
Considérons le morphisme d'anneaux
\begin{eqnarray*}
R = \mathbf{Z}[a_1, \ldots, a_n]
& \longrightarrow &
S = \mathbf{Z}[\alpha_1, \ldots, \alpha_n] \\
a_1 & \longmapsto &
\alpha_1 + \ldots + \alpha_n \\
\ldots & \ldots & \ldots \\
a_n & \longmapsto & \alpha_1 \ldots \alpha_n
\end{eqnarray*}
Autrement dit, c'est l'unique morphisme d'anneaux de polynômes indiqué tel que
$$
x^n + a_1 x^{n - 1} + \ldots + a_n
=
\prod\nolimits_{i = 1}^n (x + \alpha_i)
$$
soit satisfait dans $\mathbf{Z}[\alpha_i, x]$. On peut aussi dire que $a_i$
s'envoie sur la $i$-ième fonction symétrique élémentaire en
$\alpha_1, \ldots, \alpha_n$. D'après la théorie usuelle des polynômes
symétriques élémentaires (détails omis), l'anneau $S$ est un $R$-module libre
de type fini dont une base est constituée des éléments
$\alpha_1^{e_1} \alpha_2^{e_2} \ldots \alpha_n^{e_n}$ avec
$0 \leq e_i \leq n - i$. A fortiori, les anneaux fibres de
$R \to S$ sont finis et sont donc de dimension $0$.
D'autre part, $S$ est engendré par $n$ éléments sur $R$ soumis à $n$
équations. Ainsi, $S$ est une intersection complète globale relative sur
$R$. Comme le rang de $S$ sur $R$ est positif, nous voyons aussi que $S$ est
fidèlement plat sur $R$.
\end{example}

\begin{lemma}
\label{lemma-base-change-relative-global-complete-intersection}
Soit $S = R[x_1, \ldots, x_n]/(f_1, \ldots, f_c)$ une intersection complète
globale relative
(Définition \ref{definition-relative-global-complete-intersection}).
\begin{enumerate}
\item Pour tout $R \to R'$, le changement de base
$R' \otimes_R S = R'[x_1, \ldots, x_n]/(f_1, \ldots, f_c)$ est une
intersection complète globale relative.
\item Pour tout $g \in S$ qui est l'image d'un
$h \in R[x_1, \ldots, x_n]$, l'anneau
$S_g = R[x_1, \ldots, x_n, x_{n + 1}]/(f_1, \ldots, f_c, hx_{n + 1} - 1)$
est une intersection complète globale relative.
\item Si $R \to S$ se factorise sous la forme $R \to R_f \to S$ pour un
certain $f \in R$, alors l'anneau
$S = R_f[x_1, \ldots, x_n]/(f_1, \ldots, f_c)$ est une intersection complète
globale relative sur $R_f$.
\end{enumerate}
\end{lemma}

\begin{proof}
D'après le Lemme \ref{lemma-dimension-preserved-field-extension}, les fibres
d'un changement de base ont la même dimension que les fibres du morphisme
initial. De plus,
$R' \otimes_R R[x_1, \ldots, x_n]/(f_1, \ldots, f_c)
= R'[x_1, \ldots, x_n]/(f_1, \ldots, f_c)$. Ainsi, (1) en résulte.
La démonstration de (2) repose sur le fait que la localisation en un élément
peut s'écrire
$S_g \cong S[x_{n + 1}]/(gx_{n + 1} - 1)$.
L'assertion (3) découle de (1), car, sous les hypothèses de (3), on a
$R_f \otimes_R S \cong S$.
\end{proof}

\begin{lemma}
\label{lemma-localize-relative-complete-intersection}
Soit $R$ un anneau. Soit
$S = R[x_1, \ldots, x_n]/(f_1, \ldots, f_c)$.
Dans chacun des cas suivants, nous allons trouver
$h \in R[x_1, \ldots, x_n]$ qui s'envoie sur $g \in S$ tel que
$$
S_g = R[x_1, \ldots, x_n, x_{n + 1}]/(f_1, \ldots, f_c, hx_{n + 1} - 1)
$$
soit une intersection complète globale relative avec une présentation comme
dans la Définition \ref{definition-relative-global-complete-intersection} :
\begin{enumerate}
\item Soit $I \subset R$ un idéal. Si les fibres de
$\Spec(S/IS) \to \Spec(R/I)$ sont de dimension $n - c$, alors nous pouvons
trouver $(h, g)$ comme ci-dessus de sorte que $g$ s'envoie sur $1 \in S/IS$.
\item Soit $\mathfrak p \subset R$ un idéal premier. Si
$\dim(S \otimes_R \kappa(\mathfrak p)) = n - c$, alors nous pouvons trouver
$(h, g)$ comme ci-dessus de sorte que $g$ s'envoie sur une unité de
$S \otimes_R \kappa(\mathfrak p)$.
\item Soit $\mathfrak q \subset S$ un idéal premier situé au-dessus de
$\mathfrak p \subset R$. Si $\dim_{\mathfrak q}(S/R) = n - c$, alors nous
pouvons trouver $(h, g)$ comme ci-dessus de sorte que
$g \not \in \mathfrak q$.
\end{enumerate}
\end{lemma}

\begin{proof}
Pour (1). D'après le Lemme \ref{lemma-dimension-fibres-bounded-open-upstairs},
il existe un ouvert $W \subset \Spec(S)$ contenant $V(IS)$ tel que toutes les
fibres de $W \to \Spec(R)$ soient de dimension $\leq n - c$.
Écrivons $W = \Spec(S) \setminus V(J)$. Alors
$V(J) \cap V(IS) = \emptyset$ ; nous pouvons donc trouver $g \in J$ qui
s'envoie sur $1 \in S/IS$.
Soit $h \in R[x_1, \ldots, x_n]$ un antécédent quelconque de $g$.

\medskip\noindent
Pour (2). D'après le Lemme \ref{lemma-dimension-fibres-bounded-open-upstairs},
il existe un ouvert $W \subset \Spec(S)$ contenant
$\Spec(S \otimes_R \kappa(\mathfrak p))$
tel que toutes les fibres de $W \to \Spec(R)$ soient de dimension
$\leq n - c$.
Écrivons $W = \Spec(S) \setminus V(J)$. Alors
$V(J \cdot S \otimes_R \kappa(\mathfrak p)) = \emptyset$.
Nous pouvons donc trouver $g \in J$ qui s'envoie sur une unité de
$S \otimes_R \kappa(\mathfrak p)$ (détails omis).
Soit $h \in R[x_1, \ldots, x_n]$ un antécédent quelconque de $g$.

\medskip\noindent
Pour (3). D'après le Lemme
\ref{lemma-dimension-fibres-bounded-open-upstairs}, il existe
$g \in S$, $g \not \in \mathfrak q$, tel que toutes les fibres non vides de
$R \to S_g$ soient de dimension $\leq n - c$.
Soit $h \in R[x_1, \ldots, x_n]$ un élément quelconque qui s'envoie sur $g$.
\end{proof}

\noindent
Le lemme suivant affirme que nous pouvons effectuer une approximation
noethérienne absolue pour les intersections complètes globales relatives.

\begin{lemma}
\label{lemma-relative-global-complete-intersection-Noetherian}
Soit $R$ un anneau. Soit
$S = R[x_1, \ldots, x_n]/(f_1, \ldots, f_c)$ une intersection complète
globale relative
(Définition \ref{definition-relative-global-complete-intersection}).
Il existe une sous-$\mathbf{Z}$-algèbre $R_0 \subset R$ de type fini telle que
$f_i \in R_0[x_1, \ldots, x_n]$ et telle que
$$
S_0 = R_0[x_1, \ldots, x_n]/(f_1, \ldots, f_c)
$$
soit une intersection complète globale relative.
\end{lemma}

\begin{proof}
Soit $R_0 \subset R$ la $\mathbf{Z}$-algèbre de $R$ engendrée par tous les
coefficients des polynômes $f_1, \ldots, f_c$. Posons
$S_0 = R_0[x_1, \ldots, x_n]/(f_1, \ldots, f_c)$.
Il est clair que $S = R \otimes_{R_0} S_0$.
Choisissons un idéal premier $\mathfrak q \subset S$ et notons
$\mathfrak p \subset R$, $\mathfrak q_0 \subset S_0$ et
$\mathfrak p_0 \subset R_0$ les idéaux premiers au-dessus desquels il se
trouve. Puisque
$\dim (S \otimes_R \kappa(\mathfrak p) ) = n - c$,
on a aussi
$\dim (S_0 \otimes_{R_0} \kappa(\mathfrak p_0)) = n - c$ ;
voir le Lemme \ref{lemma-dimension-preserved-field-extension}.
D'après le Lemme \ref{lemma-dimension-fibres-bounded-open-upstairs}, il existe
$g \in S_0$, $g \not \in \mathfrak q_0$, tel que toutes les fibres non vides
de $R_0 \to (S_0)_g$ soient de dimension $\leq n - c$.
Comme $\mathfrak q$ était arbitraire et que $\Spec(S)$ est quasi-compact, nous
pouvons trouver un nombre fini d'éléments
$g_1, \ldots, g_m \in S_0$ tels que (a), pour $j = 1, \ldots, m$, les fibres
non vides de $R_0 \to (S_0)_{g_j}$ soient de dimension $\leq n - c$ et (b),
l'image de $\Spec(S) \to \Spec(S_0)$ soit contenue dans
$D(g_1) \cup \ldots \cup D(g_m)$.
Autrement dit, les images de $g_1, \ldots, g_m$ dans
$S = R \otimes_{R_0} S_0$ engendrent l'idéal unité. Après avoir agrandi $R_0$,
nous pouvons supposer que $g_1, \ldots, g_m$ engendrent l'idéal unité dans
$S_0$. D'après (a), toutes les fibres non vides de $R_0 \to S_0$ sont de
dimension $\leq n - c$, et nous concluons.
\end{proof}

\begin{lemma}
\label{lemma-relative-global-complete-intersection-conormal}
Soit $R$ un anneau. Soit
$S = R[x_1, \ldots, x_n]/(f_1, \ldots, f_c)$ une intersection complète
globale relative (Définition
\ref{definition-relative-global-complete-intersection}). Pour tout idéal
premier $\mathfrak q$ de $S$, soit $\mathfrak q'$ l'idéal premier
correspondant de $R[x_1, \ldots, x_n]$. Alors
\begin{enumerate}
\item $f_1, \ldots, f_c$ forment une suite régulière dans l'anneau local
$R[x_1, \ldots, x_n]_{\mathfrak q'}$,
\item chacun des anneaux
$R[x_1, \ldots, x_n]_{\mathfrak q'}/(f_1, \ldots, f_i)$ est plat sur $R$, et
\item le $S$-module $(f_1, \ldots, f_c)/(f_1, \ldots, f_c)^2$ est libre, avec
pour base les éléments $f_i \bmod (f_1, \ldots, f_c)^2$.
\end{enumerate}
\end{lemma}

\begin{proof}
Supposons $R$ noethérien. Soit $\mathfrak p = R \cap \mathfrak q'$.
D'après le Lemme \ref{lemma-lci}, par exemple, nous voyons que
$f_1, \ldots, f_c$ forment une suite régulière dans l'anneau local
$R[x_1, \ldots, x_n]_{\mathfrak q'} \otimes_R \kappa(\mathfrak p)$.
De plus, l'anneau local $R[x_1, \ldots, x_n]_{\mathfrak q'}$ est plat sur
$R_{\mathfrak p}$. Comme $R$, et donc
$R[x_1, \ldots, x_n]_{\mathfrak q'}$, est noethérien, nous déduisons du
Lemme \ref{lemma-grothendieck-regular-sequence} que (1) et (2) sont vraies.

\medskip\noindent
Soit maintenant $R$ quelconque. Écrivons
$R = \colim_{\lambda \in \Lambda} R_\lambda$ comme la colimite filtrante de
sous-$\mathbf{Z}$-algèbres de type fini (comparer avec la
section \ref{section-colimits-flat}). Nous pouvons supposer que
$f_1, \ldots, f_c \in R_\lambda[x_1, \ldots, x_n]$ pour tout $\lambda$.
Soit $R_0 \subset R$ comme dans le
Lemme \ref{lemma-relative-global-complete-intersection-Noetherian}.
Nous pouvons alors supposer que $R_0 \subset R_\lambda$ pour tout $\lambda$.
Il en résulte que
$S_\lambda = R_\lambda[x_1, \ldots, x_n]/(f_1, \ldots, f_c)$ est une
intersection complète globale relative (comme changement de base de $S_0$
par $R_0 \to R_\lambda$ ; voir le
Lemme \ref{lemma-base-change-relative-global-complete-intersection}).
Notons $\mathfrak p_\lambda$, $\mathfrak q_\lambda$,
$\mathfrak q'_\lambda$ les idéaux premiers de $R_\lambda$, $S_\lambda$,
$R_\lambda[x_1, \ldots, x_n]$ induits par $\mathfrak p$, $\mathfrak q$,
$\mathfrak q'$. Avec ces notations, (1) et (2) sont vraies pour tout
$\lambda$. Puisque
$$
R[x_1, \ldots, x_n]_{\mathfrak q'}/(f_1, \ldots, f_i)
=
\colim R_\lambda[x_1, \ldots, x_n]_{\mathfrak q_\lambda'}/(f_1, \ldots, f_i)
$$
nous déduisons la platitude dans (2) sur $R$ du
Lemme \ref{lemma-colimit-rings-flat}.
Comme on a
\begin{align*}
R[x_1, \ldots, x_n]_{\mathfrak q'}/(f_1, \ldots, f_i)
\xrightarrow{f_{i + 1}}
R[x_1, \ldots, x_n]_{\mathfrak q'}/(f_1, \ldots, f_i) \\
=
\colim
\left(
R_\lambda[x_1, \ldots, x_n]_{\mathfrak q_\lambda'}/(f_1, \ldots, f_i)
\xrightarrow{f_{i + 1}}
R_\lambda[x_1, \ldots, x_n]_{\mathfrak q_\lambda'}/(f_1, \ldots, f_i)
\right)
\end{align*}
et comme les colimites filtrantes sont exactes
(Lemme \ref{lemma-directed-colimit-exact}), nous concluons que (1) est vraie.

\medskip\noindent
Démonstration de (3). Notons $N$ le $S$-module
$(f_1, \ldots, f_c)/(f_1, \ldots, f_c)^2$ et $e_i \in N$ l'image de $f_i$.
D'après le Lemme \ref{lemma-regular-quasi-regular} et (1), nous savons que
$e_1, \ldots, e_c$ est une base de $N_\mathfrak q$ pour tout idéal premier
$\mathfrak q$ de $S$. Le Lemme \ref{lemma-characterize-zero-local} permet de
conclure que (3) est vraie.
\end{proof}

\begin{lemma}
\label{lemma-relative-global-complete-intersection}
Une intersection complète globale relative est syntomique, c'est-à-dire
plate.
\end{lemma}

\begin{proof}
Soit $R \to S$ une intersection complète globale relative.
Les fibres sont des intersections complètes globales, et $S$ est de
présentation finie sur $R$.
La seule chose qui reste à démontrer est donc que $R \to S$ est plat.
C'est vrai d'après (2) du
Lemme \ref{lemma-relative-global-complete-intersection-conormal}.
\end{proof}

\begin{lemma}
\label{lemma-adjoin-roots}
Supposons que $A$ soit un anneau et que
$P(x) = x^n + b_1 x^{n-1} + \ldots + b_n \in A[x]$ soit un polynôme unitaire
sur $A$. Alors il existe une extension d'anneaux syntomique, libre de type
fini et fidèlement plate $A \subset A'$ telle que
$P(x) = \prod_{i = 1, \ldots, n} (x - \beta_i)$ pour certains
$\beta_i \in A'$.
\end{lemma}

\begin{proof}
Prenons $A' = A \otimes_R S$, où $R$ et $S$ sont comme dans
l'Exemple \ref{example-roots-universal-polynomial}, où $R \to A$ envoie
$a_i$ sur $b_i$, et posons $\beta_i = -1 \otimes \alpha_i$.
Observons que $R \to S$ est fidèlement plat, libre de type fini et syntomique
d'après le Lemme \ref{lemma-relative-global-complete-intersection}.
Ces propriétés sont héritées par le changement de base $A \to A'$ ; certains
détails sont omis.
\end{proof}

\begin{lemma}
\label{lemma-syntomic}
Soit $R \to S$ un morphisme d'anneaux.
Soit $\mathfrak q \subset S$ un idéal premier situé au-dessus de l'idéal
premier $\mathfrak p$ de $R$.
Les assertions suivantes sont équivalentes :
\begin{enumerate}
\item Il existe un élément $g \in S$, $g \not \in \mathfrak q$ tel que
$R \to S_g$ soit syntomique.
\item Il existe un élément $g \in S$, $g \not \in \mathfrak q$ tel que
$S_g$ soit une intersection complète globale relative sur $R$.
\item Il existe un élément $g \in S$, $g \not \in \mathfrak q$, tel que
$R \to S_g$ soit de présentation finie, que le morphisme d'anneaux locaux
$R_{\mathfrak p} \to S_{\mathfrak q}$ soit plat et que l'anneau local
$S_{\mathfrak q}/\mathfrak pS_{\mathfrak q}$ soit un anneau d'intersection
complète sur $\kappa(\mathfrak p)$ (voir la
Définition \ref{definition-lci-local-ring}).
\end{enumerate}
\end{lemma}

\begin{proof}
L'implication (1) $\Rightarrow$ (3) est le
Lemme \ref{lemma-lci-at-prime}.
L'implication (2) $\Rightarrow$ (1) est le
Lemme \ref{lemma-relative-global-complete-intersection}.
Il reste à montrer que (3) implique (2).

\medskip\noindent
Supposons (3). Après avoir remplacé $S$ par $S_g$ pour un certain
$g \in S$, $g\not\in \mathfrak q$, nous pouvons supposer que $S$ est de
présentation finie sur $R$.
Choisissons une présentation $S = R[x_1, \ldots, x_n]/I$. Soit
$\mathfrak q' \subset R[x_1, \ldots, x_n]$ l'idéal premier correspondant à
$\mathfrak q$. Écrivons $\kappa(\mathfrak p) = k$.
Notons que $S \otimes_R k = k[x_1, \ldots, x_n]/\overline{I}$, où
$\overline{I} \subset k[x_1, \ldots, x_n]$ est l'idéal engendré par l'image
de $I$. Soit
$\overline{\mathfrak q}' \subset k[x_1, \ldots, x_n]$ l'idéal premier
engendré par l'image de $\mathfrak q'$.
D'après le Lemme \ref{lemma-lci-at-prime}, les conditions équivalentes du
Lemme \ref{lemma-lci} sont satisfaites pour $\overline{I}$ et
$\overline{\mathfrak q}'$.
Disons que la dimension de
$\overline{I}_{\overline{\mathfrak q}'}/
\overline{\mathfrak q}'\overline{I}_{\overline{\mathfrak q}'}$
sur $\kappa(\overline{\mathfrak q}')$ est $c$.
Choisissons $f_1, \ldots, f_c \in I$ qui s'envoient sur une base de cet
espace vectoriel. Les images $\overline{f}_j \in \overline{I}$ engendrent
$\overline{I}_{\overline{\mathfrak q}'}$
(d'après le Lemme \ref{lemma-lci}).
Posons $S' = R[x_1, \ldots, x_n]/(f_1, \ldots, f_c)$. Soit $J$ le noyau de
la surjection $S' \to S$. Comme $S$ est de présentation finie, $J$ est un
idéal de type fini (Lemme \ref{lemma-compose-finite-type}).
Considérons la suite exacte courte
$$
0 \to J \to S' \to S \to 0
$$
Comme $S_\mathfrak q$ est plat sur $R$, nous voyons que
$J_{\mathfrak q'} \otimes_R k \to S'_{\mathfrak q'} \otimes_R k$
est injectif (Lemme \ref{lemma-flat-tor-zero}).
Or, par construction, $S'_{\mathfrak q'} \otimes_R k$ s'envoie
isomorphiquement sur $S_\mathfrak q \otimes_R k$. Nous concluons donc que
$J_{\mathfrak q'} \otimes_R k =
J_{\mathfrak q'}/\mathfrak pJ_{\mathfrak q'} = 0$. D'après le lemme de
Nakayama (Lemme \ref{lemma-NAK}), il existe donc
$g \in R[x_1, \ldots, x_n]$, $g \not \in \mathfrak q'$ tel que $J_g = 0$.
Autrement dit, $S'_g \cong S_g$. Après une localisation supplémentaire, nous
voyons que $S'$ (et donc $S$) devient une intersection complète globale
relative d'après le
Lemme \ref{lemma-localize-relative-complete-intersection}, comme souhaité.
\end{proof}

\begin{lemma}
\label{lemma-syntomic-presentation-ideal-mod-squares}
Soit $R$ un anneau. Soit $S = R[x_1, \ldots, x_n]/I$ pour un certain idéal
$I$ de type fini. Si $g \in S$ est tel que $S_g$ soit syntomique sur $R$,
alors $(I/I^2)_g$ est un $S_g$-module projectif de type fini.
\end{lemma}

\begin{proof}
D'après le Lemme \ref{lemma-syntomic}, il existe un nombre fini d'éléments
$g_1, \ldots, g_m \in S$ qui engendrent l'idéal unité dans $S_g$ et tels que
chacun des $S_{gg_j}$ soit une intersection complète globale relative sur
$R$. Comme il suffit de montrer que $(I/I^2)_{gg_j}$ est projectif de type
fini, voir le Lemme \ref{lemma-finite-projective}, nous pouvons supposer que
$S_g$ est une intersection complète globale relative.
Dans ce cas, le résultat découle des Lemmes
\ref{lemma-conormal-module-localize} et
\ref{lemma-relative-global-complete-intersection-conormal}.
\end{proof}

\begin{lemma}
\label{lemma-composition-syntomic}
Soient $R \to S$ et $S \to S'$ des morphismes d'anneaux.
\begin{enumerate}
\item Si $R \to S$ et $S \to S'$ sont syntomiques, alors $R \to S'$ est
syntomique.
\item Si $R \to S$ et $S \to S'$ sont des intersections complètes globales
relatives, alors $R \to S'$ est une intersection complète globale relative.
\end{enumerate}
\end{lemma}

\begin{proof}
Démonstration de (2). Supposons que $R \to S$ et $S \to S'$ soient des
intersections complètes globales relatives et que nous ayons des présentations
$S =  R[x_1, \ldots, x_n]/(f_1, \ldots, f_c)$ et
$S' = S[y_1, \ldots, y_m]/(h_1, \ldots, h_d)$ comme dans la
Définition \ref{definition-relative-global-complete-intersection}.
Alors
$$
S' \cong
R[x_1, \ldots, x_n, y_1, \ldots, y_m]/(f_1, \ldots, f_c, h'_1, \ldots, h'_d)
$$
pour certains relèvements
$h_j' \in R[x_1, \ldots, x_n, y_1, \ldots, y_m]$ des $h_j$.
Il suffit donc de borner les dimensions des anneaux fibres.
Nous pouvons ainsi supposer que $R = k$ est un corps.
Dans ce cas, nous avons un anneau, à savoir $S$, qui est de type fini sur $k$
et équidimensionnel de dimension $n - c$, ainsi qu'un morphisme d'anneaux de
type fini $S \to S'$ dont tous les anneaux fibres non vides sont
équidimensionnels de dimension $m - d$. Alors, par exemple d'après le
Lemme \ref{lemma-dimension-base-fibre-total} appliqué aux localisations en les
idéaux maximaux de $S'$, nous voyons que
$\dim(S') \leq n - c + m - d$, comme souhaité.

\medskip\noindent
Nous allons ramener la partie (1) à la partie (2). Supposons que $R \to S$ et
$S \to S'$ soient syntomiques. Soit $\mathfrak q' \subset S$ un idéal premier
situé au-dessus de $\mathfrak q \subset S$. D'après le
Lemme \ref{lemma-syntomic}, il existe $g' \in S'$,
$g' \not \in \mathfrak q'$, tel que $S \to S'_{g'}$ soit une intersection
complète globale relative. De même, nous trouvons $g \in S$,
$g \not \in \mathfrak q$, tel que $R \to S_g$ soit une intersection complète
globale relative.
D'après le Lemme \ref{lemma-base-change-relative-global-complete-intersection},
le morphisme d'anneaux $S_g \to S_{gg'}$ est une intersection complète
globale relative. D'après la partie (2), nous voyons que
$R \to S_{gg'}$ est une intersection complète globale relative et que
$gg' \not \in \mathfrak q'$.
Comme $\mathfrak q'$ était arbitraire, en combinant les
Lemmes \ref{lemma-syntomic} et \ref{lemma-local-syntomic}, nous voyons que
$R \to S'$ est syntomique (cela utilise aussi que le spectre de $S'$ est
quasi-compact ; voir le Lemme \ref{lemma-quasi-compact}).
\end{proof}

\noindent
Le lemme suivant sera amélioré plus tard ; voir Lissage des morphismes
d'anneaux, Proposition \ref{smoothing-proposition-lift-smooth}.

\begin{lemma}
\label{lemma-lift-syntomic}
Soit $R$ un anneau et soit $I \subset R$ un idéal.
Soit $R/I \to \overline{S}$ un morphisme syntomique.
Alors il existe des éléments $\overline{g}_i \in \overline{S}$ qui engendrent
l'idéal unité de $\overline{S}$ et tels que chacun des
$\overline{S}_{g_i} \cong S_i/IS_i$ pour une certaine intersection complète
globale relative $S_i$ sur $R$.
\end{lemma}

\begin{proof}
D'après le Lemme \ref{lemma-syntomic}, nous trouvons une famille d'éléments
$\overline{g}_i \in \overline{S}$ qui engendrent l'idéal unité de
$\overline{S}$ et tels que chacun des $\overline{S}_{g_i}$ soit une
intersection complète globale relative sur $R/I$.
Nous pouvons donc supposer que $\overline{S}$ est une intersection complète
globale relative.
Écrivons
$\overline{S} =
(R/I)[x_1, \ldots, x_n]/(\overline{f}_1, \ldots, \overline{f}_c)$
comme dans la Définition \ref{definition-relative-global-complete-intersection}.
Choisissons $f_1, \ldots, f_c \in R[x_1, \ldots, x_n]$ relevant
$\overline{f}_1, \ldots, \overline{f}_c$.
Posons $S = R[x_1, \ldots, x_n]/(f_1, \ldots, f_c)$.
Notons que $S/IS \cong \overline{S}$.
D'après le Lemme \ref{lemma-localize-relative-complete-intersection}, nous
pouvons trouver $g \in S$ qui s'envoie sur $1$ dans $\overline{S}$ et tel que
$S_g$ soit une intersection complète globale relative sur $R$.
Comme $\overline{S} \cong S_g/IS_g$, cela achève la démonstration.
\end{proof}


\section{Morphismes d'anneaux lisses}
\label{section-smooth}

\noindent
Motivons la définition d'un morphisme d'anneaux lisse par un exemple.
Supposons que $R$ soit un anneau et que $S = R[x, y]/(f)$ pour un certain
$f \in R[x, y]$ non nul. Il existe alors une suite exacte
$$
S \to
S\text{d}x \oplus S\text{d}y \to
\Omega_{S/R} \to 0
$$
où la première flèche envoie $1$ sur
$\frac{\partial f}{\partial x} \text{d}x +
\frac{\partial f}{\partial y} \text{d}y$ ; voir la
section \ref{section-netherlander}.
Nous concluons que $\Omega_{S/R}$ est localement libre de rang $1$ si les
dérivées partielles de $f$ engendrent l'idéal unité dans $S$.
Dans ce cas, $S$ est lisse de dimension relative $1$ sur $R$.
Mais il peut arriver que $\Omega_{S/R}$ soit localement libre de rang $2$, à
savoir lorsque les deux dérivées partielles de $f$ sont nulles. Cela se
produit par exemple si, pour un nombre premier $p$, on a $p = 0$ dans $R$ et
$f = x^p + y^p$. Ici, $R \to S$ est une intersection complète globale
relative de dimension relative $1$ qui n'est pas lisse.
Ainsi, pour vérifier qu'un morphisme d'anneaux est lisse, il ne suffit pas de
vérifier si le module des différentielles est libre. La condition correcte est
la suivante.

\begin{definition}
\label{definition-smooth}
Un morphisme d'anneaux $R \to S$ est {\it lisse} s'il est de présentation
finie et si le complexe cotangent naïf $\NL_{S/R}$ est quasi-isomorphe à un
$S$-module projectif de type fini placé en degré $0$ : cela signifie que
$H_1(\NL_{S/R}) = 0$ et que $\Omega_{S/R}$ est un $S$-module projectif de
type fini.
\end{definition}

\noindent
Soit $R \to S$ un morphisme d'anneaux. D'après le
Lemme \ref{lemma-NL-homotopy}, pour toute présentation
$\alpha : P \to S$, on a $H_1(\NL(\alpha)) = H_1(\NL_{S/R})$ et
$H_0(\NL(\alpha)) = \Omega_{S/R}$. Ainsi, si $R \to S$ est lisse, alors,
pour toute surjection $\alpha : R[x_1, \ldots, x_n] \to S$ de noyau $I$, le
morphisme
$$
I/I^2
\longrightarrow
\Omega_{R[x_1, \ldots, x_n]/R} \otimes_{R[x_1, \ldots, x_n]} S
$$
est injectif et son conoyau est un $S$-module projectif de type fini.
La flèche affichée est donc une injection scindée. Autrement dit,
$\bigoplus_{i = 1}^n S \text{d}x_i \cong I/I^2 \oplus \Omega_{S/R}$
comme $S$-modules, et $I/I^2$ est lui aussi un $S$-module projectif de type
fini. Réciproquement, si $R \to S$ est de présentation finie et si nous avons
une surjection $\alpha : R[x_1, \ldots, x_n] \to S$ de noyau $I$ telle que la
flèche affichée soit une injection scindée, alors $R \to S$ est lisse.

\begin{lemma}
\label{lemma-localize-smooth}
Soit $R \to S$ un morphisme d'anneaux lisse.
Toute localisation $S_g$ est lisse sur $R$.
Si $f \in R$ s'envoie sur un élément inversible de $S$, alors
$R_f \to S$ est lisse.
\end{lemma}

\begin{proof}
D'après le Lemme \ref{lemma-localize-NL}, le complexe cotangent naïf de $S_g$
sur $R$ est le changement de base du complexe cotangent naïf de $S$ sur $R$.
L'hypothèse est que le complexe cotangent naïf de $S/R$ est $\Omega_{S/R}$ et
que c'est un $S$-module projectif de type fini. Il en va donc de même de son
changement de base. Ainsi, $S_g$ est lisse sur $R$.

\medskip\noindent
La seconde assertion découle de la même manière du
Lemme \ref{lemma-NL-localize-bottom}.
\end{proof}

\begin{lemma}
\label{lemma-base-change-smooth}
\begin{slogan}
La lissité est préservée par changement de base
\end{slogan}
Soit $R \to S$ un morphisme d'anneaux lisse.
Soit $R \to R'$ un morphisme d'anneaux quelconque.
Alors le changement de base $R' \to S' = R' \otimes_R S$ est lisse.
\end{lemma}

\begin{proof}
Soit $\alpha : R[x_1, \ldots, x_n] \to S$ une présentation de noyau $I$.
Soit $\alpha' : R'[x_1, \ldots, x_n] \to R' \otimes_R S$ la présentation
induite. Posons $I' = \Ker(\alpha')$.
Puisque $0 \to I \to R[x_1, \ldots, x_n] \to S \to 0$ est exacte, la suite
$R' \otimes_R I \to R'[x_1, \ldots, x_n] \to R' \otimes_R S \to 0$
est exacte. Ainsi, $R' \otimes_R I \to I'$ est surjectif.
D'après la Définition \ref{definition-smooth}, il existe une suite exacte
courte
$$
0 \to I/I^2 \to
\Omega_{R[x_1, \ldots, x_n]/R} \otimes_{R[x_1, \ldots, x_n]} S \to
\Omega_{S/R} \to
0
$$
et le $S$-module $\Omega_{S/R}$ est projectif de type fini.
En particulier, $I/I^2$ est un facteur direct de
$\Omega_{R[x_1, \ldots, x_n]/R} \otimes_{R[x_1, \ldots, x_n]} S$.
Considérons le diagramme commutatif
$$
\xymatrix{
R' \otimes_R (I/I^2) \ar[r] \ar[d] &
R' \otimes_R (\Omega_{R[x_1, \ldots, x_n]/R} \otimes_{R[x_1, \ldots, x_n]} S)
\ar[d] \\
I'/(I')^2 \ar[r] &
\Omega_{R'[x_1, \ldots, x_n]/R'}
\otimes_{R'[x_1, \ldots, x_n]} (R' \otimes_R S)
}
$$
Comme le morphisme vertical de droite est un isomorphisme, nous voyons, par ce
qui précède, que le morphisme vertical de gauche est injectif et surjectif.
Nous concluons donc que $\NL(\alpha')$ est quasi-isomorphe à
$\Omega_{S'/R'} \cong S' \otimes_S \Omega_{S/R}$ placé en degré $0$.
Ce module est projectif de type fini, car c'est le changement de base d'un
module projectif de type fini.
\end{proof}

\begin{lemma}
\label{lemma-smooth-over-field}
Soit $k$ un corps.
Soit $S$ une $k$-algèbre lisse.
Alors $S$ est une intersection complète locale.
\end{lemma}

\begin{proof}
D'après les Lemmes \ref{lemma-base-change-smooth} et
\ref{lemma-lci-field-change}, il suffit de le démontrer lorsque $k$ est
algébriquement clos. Choisissons une présentation
$\alpha : k[x_1, \ldots, x_n] \to S$ de noyau $I$. Soit $\mathfrak m$ un
idéal maximal de $S$, et soit $\mathfrak m' \supset I$ l'idéal maximal
correspondant de $k[x_1, \ldots, x_n]$.
Nous allons montrer que la condition (5) du
Lemme \ref{lemma-lci} est satisfaite (avec $\mathfrak m$ à la place de
$\mathfrak q$).
Nous pouvons écrire $\mathfrak m' = (x_1 - a_1, \ldots, x_n - a_n)$ pour
certains $a_i \in k$, car $k$ est algébriquement clos ; voir le
Théorème \ref{theorem-nullstellensatz}.
Puisque, par hypothèse, $k \to S$ est lisse, le morphisme de $S$-modules
$\text{d} : I/I^2 \to \bigoplus_{i = 1}^n S \text{d}x_i$
est une injection scindée. Le morphisme correspondant
$I/\mathfrak m' I \to \bigoplus \kappa(\mathfrak m') \text{d}x_i$
est donc injectif. Écrivons
$\dim_{\kappa(\mathfrak m')}(I/\mathfrak m' I) = c$ et choisissons
$f_1, \ldots, f_c \in I$ dont les images forment une base sur
$\kappa(\mathfrak m')$ de $I/\mathfrak m' I$.
D'après le Lemme de Nakayama \ref{lemma-NAK}, nous voyons que
$f_1, \ldots, f_c$ engendrent $I_{\mathfrak m'}$ sur
$k[x_1, \ldots, x_n]_{\mathfrak m'}$. Considérons le diagramme commutatif
$$
\xymatrix{
I \ar[r] \ar[d] & I/I^2 \ar[rr] \ar[d] & &
I/\mathfrak m'I \ar[d] \\
\Omega_{k[x_1, \ldots, x_n]/k} \ar[r] &
\bigoplus S\text{d}x_i \ar[rr]^{\text{d}x_i \mapsto x_i - a_i} & &
\mathfrak m'/(\mathfrak m')^2
}
$$
(démonstration de la commutativité omise). Le morphisme vertical du milieu est
celui qui définit le complexe cotangent naïf de $\alpha$. Notons que la flèche
horizontale inférieure de droite induit un isomorphisme
$\bigoplus \kappa(\mathfrak m') \text{d}x_i \to \mathfrak m'/(\mathfrak m')^2$.
Ainsi, nos générateurs $f_1, \ldots, f_c$ de $I_{\mathfrak m'}$ s'envoient sur
une famille d'éléments de $k[x_1, \ldots, x_n]_{\mathfrak m'}$ dont les
classes dans $\mathfrak m'/(\mathfrak m')^2$ sont linéairement indépendantes
sur $\kappa(\mathfrak m')$. Ils forment donc une suite régulière dans l'anneau
$k[x_1, \ldots, x_n]_{\mathfrak m'}$ d'après le
Lemme \ref{lemma-regular-ring-CM}.
Cela vérifie la condition (5) du Lemme \ref{lemma-lci} ; par conséquent,
$S_g$ est une intersection complète globale sur $k$ pour un certain
$g \in S$, $g \not \in \mathfrak m$. Comme cela vaut pour tout idéal maximal
de $S$, nous concluons que $S$ est une intersection complète locale sur $k$.
\end{proof}

\begin{definition}
\label{definition-standard-smooth}
Soit $R$ un anneau. Étant donnés des entiers $n \geq c \geq 0$ et
$f_1, \ldots, f_c \in R[x_1, \ldots, x_n]$, nous disons que
$$
R[x_1, \ldots, x_n]/(f_1, \ldots, f_c)
$$
est une {\it algèbre lisse standard sur $R$} si le polynôme
$$
g =
\det
\left(
\begin{matrix}
\partial f_1/\partial x_1 &
\partial f_2/\partial x_1 &
\ldots &
\partial f_c/\partial x_1 \\
\partial f_1/\partial x_2 &
\partial f_2/\partial x_2 &
\ldots &
\partial f_c/\partial x_2 \\
\ldots & \ldots & \ldots & \ldots \\
\partial f_1/\partial x_c &
\partial f_2/\partial x_c &
\ldots &
\partial f_c/\partial x_c
\end{matrix}
\right)
$$
s'envoie sur un élément inversible de
$R[x_1, \ldots, x_n]/(f_1, \ldots, f_c)$.
Nous disons qu'une $R$-algèbre $S$ est {\it lisse standard}, ou que le
morphisme d'anneaux $R \to S$ est {\it lisse standard}, s'il existe
$n \geq c \geq 0$ et $f_1, \ldots, f_c \in R[x_1, \ldots, x_n]$ tels que
$R[x_1, \ldots, x_n]/(f_1, \ldots, f_c)$ soit une algèbre lisse standard sur
$R$ et que $S$ soit isomorphe à
$R[x_1, \ldots, x_n]/(f_1, \ldots, f_c)$ comme $R$-algèbre.
\end{definition}

\begin{lemma}
\label{lemma-standard-smooth}
Soit
$S = R[x_1, \ldots, x_n]/(f_1, \ldots, f_c) = R[x_1, \ldots, x_n]/I$
une algèbre lisse standard. Alors
\begin{enumerate}
\item le morphisme d'anneaux $R \to S$ est lisse,
\item le $S$-module $\Omega_{S/R}$ est libre sur
$\text{d}x_{c + 1}, \ldots, \text{d}x_n$,
\item le $S$-module $I/I^2$ est libre sur les classes de
$f_1, \ldots, f_c$,
\item pour tout $g \in S$, le morphisme d'anneaux $R \to S_g$ est lisse
standard,
\item pour tout morphisme d'anneaux $R \to R'$, le changement de base
$R' \to R'\otimes_R S$ est lisse standard,
\item si $f \in R$ s'envoie sur un élément inversible de $S$, alors
$R_f \to S$ est lisse standard, et
\item l'anneau $S$ est une intersection complète globale relative sur $R$.
\end{enumerate}
\end{lemma}

\begin{proof}
Considérons le complexe cotangent naïf de la présentation donnée
$$
(f_1, \ldots, f_c)/(f_1, \ldots, f_c)^2
\longrightarrow
\bigoplus\nolimits_{i = 1}^n S \text{d}x_i
$$
Composons ce morphisme avec la projection sur les $c$ premiers facteurs
directs de la somme directe. D'après la définition d'une algèbre lisse
standard, les classes $f_i \bmod (f_1, \ldots, f_c)^2$ s'envoient sur une
base de $\bigoplus_{i = 1}^c S\text{d}x_i$. Nous concluons que
$(f_1, \ldots, f_c)/(f_1, \ldots, f_c)^2$ est libre de rang $c$, avec pour
base les éléments $f_i \bmod (f_1, \ldots, f_c)^2$, et que l'homologie en
degré $0$, c'est-à-dire $\Omega_{S/R}$, du complexe cotangent naïf est un
$S$-module libre dont une base est formée par les images de
$\text{d}x_{c + j}$, $j = 1, \ldots, n - c$.
Cela prouve en particulier que $R \to S$ est lisse.

\medskip\noindent
Les démonstrations de (4) et (6) sont omises. Voir toutefois l'exemple
ci-dessous et la démonstration du
Lemme \ref{lemma-base-change-relative-global-complete-intersection}.

\medskip\noindent
Soit $\varphi : R \to R'$ un morphisme d'anneaux quelconque.
Notons $S' = R'[x_1, \ldots, x_n]/(f_1^\varphi, \ldots, f_c^\varphi)$, où
$f^\varphi$ est le polynôme obtenu à partir de
$f \in R[x_1, \ldots, x_n]$ en appliquant $\varphi$ à tous les coefficients.
Alors $S' \cong R' \otimes_R S$.
De plus, le déterminant de la Définition \ref{definition-standard-smooth}
pour $S'/R'$ est égal à $g^\varphi$. Son image dans $S'$ est donc l'image de
$g$ par $R[x_1, \ldots, x_n] \to S \to S'$, et est par conséquent
inversible. Cela prouve (5).

\medskip\noindent
Pour prouver (7), il suffit de montrer que
$S \otimes_R \kappa(\mathfrak p)$ est de dimension $n - c$ pour tout idéal
premier $\mathfrak p \subset R$.
D'après (5), il suffit de montrer que toute algèbre lisse standard
$k[x_1, \ldots, x_n]/(f_1, \ldots, f_c)$ sur un corps $k$ est de dimension
$n - c$. Nous savons déjà que
$k[x_1, \ldots, x_n]/(f_1, \ldots, f_c)$ est une intersection complète
locale d'après le Lemme \ref{lemma-smooth-over-field}.
Ainsi, puisque $I/I^2$ est libre de rang $c$, nous voyons que
$k[x_1, \ldots, x_n]/(f_1, \ldots, f_c)$ est de dimension $n - c$, par
exemple d'après le Lemme \ref{lemma-lci}.
\end{proof}

\begin{example}
\label{example-make-standard-smooth}
Soit $R$ un anneau.
Soient $f_1, \ldots, f_c \in R[x_1, \ldots, x_n]$.
Soit
$$
h =
\det
\left(
\begin{matrix}
\partial f_1/\partial x_1 &
\partial f_2/\partial x_1 &
\ldots &
\partial f_c/\partial x_1 \\
\partial f_1/\partial x_2 &
\partial f_2/\partial x_2 &
\ldots &
\partial f_c/\partial x_2 \\
\ldots & \ldots & \ldots & \ldots \\
\partial f_1/\partial x_c &
\partial f_2/\partial x_c &
\ldots &
\partial f_c/\partial x_c
\end{matrix}
\right).
$$
Posons
$S = R[x_1, \ldots, x_{n + 1}]/(f_1, \ldots, f_c, x_{n + 1}h - 1)$.
C'est un exemple d'algèbre lisse standard, si ce n'est que la présentation est
mauvaise et que les variables devraient être rangées dans l'ordre suivant :
$x_1, \ldots, x_c, x_{n + 1}, x_{c + 1}, \ldots, x_n$.
\end{example}

\begin{lemma}
\label{lemma-compose-standard-smooth}
Une composée de morphismes d'anneaux lisses standard est lisse standard.
\end{lemma}

\begin{proof}
Supposons que $R \to S$ et $S \to S'$ soient lisses standard. Choisissons des
présentations
$S =  R[x_1, \ldots, x_n]/(f_1, \ldots, f_c)$
et
$S' = S[y_1, \ldots, y_m]/(g_1, \ldots, g_d)$.
Choisissons des éléments
$g_j' \in R[x_1, \ldots, x_n, y_1, \ldots, y_m]$ qui s'envoient sur les
$g_j$. Nous voyons ainsi que
$S' = R[x_1, \ldots, x_n, y_1, \ldots, y_m]/
(f_1, \ldots, f_c, g'_1, \ldots, g'_d)$.
Pour montrer que $S'$ est lisse standard, il suffit de vérifier que le
déterminant
$$
\det
\left(
\begin{matrix}
\partial f_1/\partial x_1 &
\ldots &
\partial f_c/\partial x_1 &
\partial g_1/\partial x_1 &
\ldots &
\partial g_d/\partial x_1 \\
\ldots &
\ldots &
\ldots &
\ldots &
\ldots &
\ldots  \\
\partial f_1/\partial x_c &
\ldots &
\partial f_c/\partial x_c &
\partial g_1/\partial x_c &
\ldots &
\partial g_d/\partial x_c \\
0 &
\ldots &
0 &
\partial g_1/\partial y_1 &
\ldots &
\partial g_d/\partial y_1 \\
\ldots &
\ldots &
\ldots &
\ldots &
\ldots &
\ldots  \\
0 &
\ldots &
0 &
\partial g_1/\partial y_d &
\ldots &
\partial g_d/\partial y_d
\end{matrix}
\right)
$$
est inversible dans $S'$. C'est clair, puisqu'il est le produit des deux
déterminants qui sont inversibles par hypothèse.
\end{proof}

\begin{lemma}
\label{lemma-smooth-syntomic}
Soit $R \to S$ un morphisme d'anneaux lisse.
Il existe un recouvrement ouvert de $\Spec(S)$ par des ouverts standards
$D(g)$ tel que chacun des $S_g$ soit lisse standard sur $R$.
En particulier, $R \to S$ est syntomique.
\end{lemma}

\begin{proof}
Choisissons une présentation $\alpha : R[x_1, \ldots, x_n] \to S$ de noyau
$I = (f_1, \ldots, f_m)$. Pour toute partie
$E \subset \{1, \ldots, m\}$, considérons l'ouvert $U_E$ sur lequel les
classes $f_e, e\in E$ engendrent librement le $S$-module projectif de type
fini $I/I^2$ ; voir le Lemme \ref{lemma-cokernel-flat}.
Nous pouvons recouvrir $\Spec(S)$ par des ouverts standards $D(g)$ dont chacun
est entièrement contenu dans l'un des ouverts $U_E$. Pour un tel $g$, nous
considérons la présentation
$$
\beta : R[x_1, \ldots, x_n, x_{n + 1}] \longrightarrow S_g
$$
qui envoie $x_{n + 1}$ sur $1/g$. Posant $J = \Ker(\beta)$, nous utilisons le
Lemme \ref{lemma-principal-localization-NL} pour voir que
$J/J^2 \cong (I/I^2)_g \oplus S_g$ est libre.
Nous pouvons remplacer, et remplaçons, $S$ par $S_g$. Alors, en utilisant le
Lemme \ref{lemma-huber}, nous pouvons supposer que nous avons une présentation
$\alpha : R[x_1, \ldots, x_n] \to S$ de noyau
$I = (f_1, \ldots, f_c)$ telle que $I/I^2$ soit libre sur les classes de
$f_1, \ldots, f_c$.

\medskip\noindent
En utilisant la présentation $\alpha$ obtenue à la fin du paragraphe
précédent, nous répétons plus ou moins cet argument avec les éléments de base
$\text{d}x_1, \ldots, \text{d}x_n$ de
$\Omega_{R[x_1, \ldots, x_n]/R}$.
En effet, pour toute partie $E \subset \{1, \ldots, n\}$ de cardinal $c$,
nous pouvons considérer l'ouvert $U_E$ de $\Spec(S)$ sur lequel la
différentielle de $\NL(\alpha)$, composée avec la projection
$$
S^{\oplus c} \cong I/I^2
\longrightarrow
\Omega_{R[x_1, \ldots, x_n]/R} \otimes_{R[x_1, \ldots, x_n]} S
\longrightarrow
\bigoplus\nolimits_{i \in E} S\text{d}x_i
$$
est un isomorphisme. À nouveau, nous pouvons trouver un recouvrement de
$\Spec(S)$ par des ouverts standards $D(g)$ (en nombre fini) tel que chaque
$D(g)$ soit entièrement contenu dans l'un des ouverts $U_E$.
En renumérotant, nous pouvons supposer que $E = \{1, \ldots, c\}$.
Pour un $g$ tel que $D(g) \subset U_E$, nous considérons la présentation
$$
\beta : R[x_1, \ldots, x_n, x_{n + 1}] \to S_g
$$
qui envoie $x_{n + 1}$ sur $1/g$. Posant $J = \Ker(\beta)$, nous déduisons du
Lemme \ref{lemma-principal-localization-NL} que
$J = (f_1, \ldots, f_c, fx_{n + 1} - 1)$, où $\alpha(f) = g$, et que la
composée
$$
J/J^2 \longrightarrow
\Omega_{R[x_1, \ldots, x_{n + 1}]/R} \otimes_{R[x_1, \ldots, x_{n + 1}]} S_g
\longrightarrow
\bigoplus\nolimits_{i = 1}^c S_g\text{d}x_i \oplus S_g \text{d}x_{n + 1}
$$
est un isomorphisme. En réordonnant les coordonnées sous la forme
$x_1, \ldots, x_c, x_{n + 1}, x_{c + 1}, \ldots, x_n$, nous concluons que
$S_g$ est lisse standard sur $R$, comme souhaité.

\medskip\noindent
Cela achève la démonstration, car les algèbres lisses standard sont
syntomiques (Lemmes \ref{lemma-standard-smooth} et
\ref{lemma-relative-global-complete-intersection}) et la propriété d'être
syntomique sur $R$ est locale sur $S$
(Lemme \ref{lemma-local-syntomic}).
\end{proof}

\begin{definition}
\label{definition-smooth-at-prime}
Soit $R \to S$ un morphisme d'anneaux.
Soit $\mathfrak q$ un idéal premier de $S$.
Nous disons que $R \to S$ est {\it lisse en $\mathfrak q$} s'il existe
$g \in S$, $g \not \in \mathfrak q$, tel que
$R \to S_g$ soit lisse.
\end{definition}

\noindent
Pour les morphismes d'anneaux de présentation finie, on peut caractériser
cette propriété comme suit.

\begin{lemma}
\label{lemma-smooth-at-point}
Soit $R \to S$ de présentation finie. Soit $\mathfrak q$ un idéal premier de
$S$. Les conditions suivantes sont équivalentes :
\begin{enumerate}
\item $R \to S$ est lisse en $\mathfrak q$,
\item $H_1(L_{S/R})_\mathfrak q = 0$ et
$\Omega_{S/R, \mathfrak q}$ est un $S_\mathfrak q$-module libre de type fini,
\item $H_1(L_{S/R})_\mathfrak q = 0$ et
$\Omega_{S/R, \mathfrak q}$ est un $S_\mathfrak q$-module projectif, et
\item $H_1(L_{S/R})_\mathfrak q = 0$ et
$\Omega_{S/R, \mathfrak q}$ est un $S_\mathfrak q$-module plat.
\end{enumerate}
\end{lemma}

\begin{proof}
Nous utiliserons sans autre mention le fait que la formation du complexe
cotangent naïf commute à la localisation ; voir la
section \ref{section-netherlander}, en particulier le
Lemme \ref{lemma-localize-NL}.
Notons que $\Omega_{S/R}$ est un $S$-module de présentation finie ; voir le
Lemme \ref{lemma-differentials-finitely-presented}. Ainsi (2), (3) et (4) sont
équivalentes par le Lemme \ref{lemma-finite-projective}.
Il est clair que (1) implique les conditions équivalentes (2), (3) et (4).
Supposons (2). En écrivant $S_\mathfrak q$ comme la colimite de localisations
principales, le Lemme \ref{lemma-colimit-category-fp-modules} montre qu'il
existe $g \in S$, $g \not \in \mathfrak q$, tel que
$(\Omega_{S/R})_g$ soit libre de type fini. Choisissons une présentation
$\alpha : R[x_1, \ldots, x_n] \to S$ de noyau $I$. Nous pouvons travailler
avec $\NL(\alpha)$ au lieu de $\NL_{S/R}$ ; voir le
Lemme \ref{lemma-NL-homotopy}. La surjection
$$
\Omega_{R[x_1, \ldots, x_n]/R} \otimes_{R[x_1, \ldots, x_n]} S
\to \Omega_{S/R} \to 0
$$
admet un inverse à droite après inversion de $g$, car $(\Omega_{S/R})_g$ est
projectif. L'image de
$\text{d} : (I/I^2)_g \to
\Omega_{R[x_1, \ldots, x_n]/R} \otimes_{R[x_1, \ldots, x_n]} S_g$
est donc un facteur direct, et ce morphisme admet lui aussi un inverse à
droite. Nous en déduisons que $H_1(L_{S/R})_g$ est un quotient de $(I/I^2)_g$.
En particulier, $H_1(L_{S/R})_g$ est un $S_g$-module de type fini.
L'annulation de $H_1(L_{S/R})_{\mathfrak q}$ entraîne donc celle de
$H_1(L_{S/R})_{gg'}$ pour un certain $g' \in S$,
$g' \not \in \mathfrak q$. Alors $R \to S_{gg'}$ est lisse par définition.
\end{proof}

\begin{lemma}
\label{lemma-locally-smooth}
\begin{slogan}
Un morphisme d'anneaux est lisse si et seulement s'il est lisse en tout idéal
premier de la cible
\end{slogan}
Soit $R \to S$ un morphisme d'anneaux.
Alors $R \to S$ est lisse si et seulement si $R \to S$ est lisse
en tout idéal premier $\mathfrak q$ de $S$.
\end{lemma}

\begin{proof}
L'implication directe est triviale. Supposons que $R \to S$ soit lisse en tout
idéal premier $\mathfrak q$ de $S$. Comme $\Spec(S)$ est quasi-compact ; voir
le Lemme \ref{lemma-quasi-compact}, il existe un recouvrement fini
$\Spec(S) = \bigcup D(g_i)$ tel que chaque $S_{g_i}$ soit lisse.
Le Lemme \ref{lemma-cover-upstairs} implique alors que $S$ est de présentation
finie sur $R$. D'après le Lemme \ref{lemma-localize-NL},
$\NL_{S/R} \otimes_S S_{g_i}$ est quasi-isomorphe à un $S_{g_i}$-module
projectif de type fini placé en degré $0$. Le
Lemme \ref{lemma-finite-projective} implique alors que $\NL_{S/R}$ est
quasi-isomorphe à un $S$-module projectif de type fini placé en degré $0$.
\end{proof}

\begin{lemma}
\label{lemma-compose-smooth}
La composée de morphismes d'anneaux lisses est lisse.
\end{lemma}

\begin{proof}
On peut le démontrer de nombreuses façons. L'une d'elles consiste à utiliser
le lemme du serpent (Lemme \ref{lemma-snake}), la suite de Jacobi--Zariski
(Lemme \ref{lemma-exact-sequence-NL}), ainsi que la caractérisation des modules
projectifs comme facteurs directs de modules libres
(Lemme \ref{lemma-characterize-projective}). Une autre démonstration s'obtient
en combinant les Lemmes \ref{lemma-smooth-syntomic},
\ref{lemma-compose-standard-smooth} et \ref{lemma-locally-smooth}.
\end{proof}

\begin{lemma}
\label{lemma-product-smooth}
Soit $R$ un anneau. Soit $S = S' \times S''$ un produit de $R$-algèbres.
Alors $S$ est lisse sur $R$ si et seulement si $S'$ et $S''$ sont toutes deux
lisses sur $R$.
\end{lemma}

\begin{proof}
Omis. Indications : par le Lemme \ref{lemma-locally-smooth}, on peut vérifier
la lissité idéal premier par idéal premier. Comme $\Spec(S)$ est la réunion
disjointe de $\Spec(S')$ et de $\Spec(S'')$ par le
Lemme \ref{lemma-spec-product}, la lissité de $R \to S$ en $\mathfrak q$
correspond soit à la lissité de $R \to S'$ en l'idéal premier correspondant,
soit à celle de $R \to S''$ en l'idéal premier correspondant.
\end{proof}

\begin{lemma}
\label{lemma-relative-global-complete-intersection-smooth}
Soit $R$ un anneau. Soit $S = R[x_1, \ldots, x_n]/(f_1, \ldots, f_c)$
une intersection complète globale relative.
Soit $\mathfrak q \subset S$ un idéal premier. Alors $R \to S$ est lisse en
$\mathfrak q$ si et seulement s'il existe une partie
$I \subset \{1, \ldots, n\}$ de cardinal $c$ telle que le polynôme
$$
g_I = \det (\partial f_j/\partial x_i)_{j = 1, \ldots, c, \ i \in I}.
$$
ne s'envoie pas sur un élément de $\mathfrak q$.
\end{lemma}

\begin{proof}
Le Lemme \ref{lemma-relative-global-complete-intersection-conormal} montre que
le complexe cotangent naïf associé à la présentation donnée de $S$ est le
complexe
$$
\bigoplus\nolimits_{j = 1}^c S \cdot f_j
\longrightarrow
\bigoplus\nolimits_{i = 1}^n S \cdot \text{d}x_i, \quad
f_j \longmapsto \sum \frac{\partial f_j}{\partial x_i} \text{d}x_i.
$$
Les mineurs maximaux de la matrice qui donne ce morphisme sont précisément les
polynômes $g_I$.

\medskip\noindent
Supposons que $g_I$ s'envoie sur $g \in S$, avec
$g \not \in \mathfrak q$. Alors l'algèbre $S_g$ est lisse sur $R$.
En effet, son complexe cotangent naïf est quasi-isomorphe au complexe ci-dessus
localisé en $g$ ; voir le Lemme \ref{lemma-localize-NL}. Par construction, il
est quasi-isomorphe à un module libre de rang $n - c$ en degré $0$.

\medskip\noindent
Réciproquement, supposons que tous les $g_I$ appartiennent finalement à
$\mathfrak q$. Dans ce cas, le complexe ci-dessus tensorisé avec
$\kappa(\mathfrak q)$ n'est pas de rang maximal ; il n'existe donc aucune
localisation par un élément $g \in S$, $g \not \in \mathfrak q$, dans laquelle
ce morphisme devienne une injection scindée. De nouveau, le
Lemme \ref{lemma-localize-NL} montre qu'aucune telle localisation n'est lisse
sur $R$.
\end{proof}

\begin{lemma}
\label{lemma-flat-fibre-smooth}
Soit $R \to S$ un morphisme d'anneaux.
Soit $\mathfrak q \subset S$ un idéal premier au-dessus de l'idéal premier
$\mathfrak p$ de $R$. Supposons que
\begin{enumerate}
\item il existe $g \in S$, $g \not\in \mathfrak q$, tel que
$R \to S_g$ soit de présentation finie,
\item l'homomorphisme d'anneaux locaux
$R_{\mathfrak p} \to S_{\mathfrak q}$ soit plat,
\item la fibre $S \otimes_R \kappa(\mathfrak p)$ soit lisse sur
$\kappa(\mathfrak p)$ en l'idéal premier correspondant à $\mathfrak q$.
\end{enumerate}
Alors $R \to S$ est lisse en $\mathfrak q$.
\end{lemma}

\begin{proof}
Les Lemmes \ref{lemma-syntomic} et \ref{lemma-smooth-over-field} montrent qu'il
existe $g \in S$ tel que $S_g$ soit une intersection complète globale relative.
En remplaçant $S$ par $S_g$, nous pouvons supposer que
$S = R[x_1, \ldots, x_n]/(f_1, \ldots, f_c)$ est une intersection complète
globale relative. Pour toute partie $I \subset \{1, \ldots, n\}$ de cardinal
$c$, considérons le polynôme
$g_I = \det (\partial f_j/\partial x_i)_{j = 1, \ldots, c, i \in I}$
du Lemme \ref{lemma-relative-global-complete-intersection-smooth}.
Notons que l'image $\overline{g}_I$ de $g_I$ dans l'anneau de polynômes
$\kappa(\mathfrak p)[x_1, \ldots, x_n]$ est le déterminant des dérivées
partielles des images $\overline{f}_j$ des $f_j$ dans l'anneau
$\kappa(\mathfrak p)[x_1, \ldots, x_n]$. Le lemme résulte donc de l'application
du Lemme \ref{lemma-relative-global-complete-intersection-smooth} à la fois à
$R \to S$ et à $\kappa(\mathfrak p) \to S \otimes_R \kappa(\mathfrak p)$.
\end{proof}

\noindent
Notons que les ensembles $U, V$ du lemme suivant sont ouverts par définition.

\begin{lemma}
\label{lemma-flat-base-change-locus-smooth}
Soit $R \to S$ un morphisme d'anneaux de présentation finie.
Soit $R \to R'$ un morphisme d'anneaux plat.
Notons $S' = R' \otimes_R S$ le changement de base.
Soit $U \subset \Spec(S)$ l'ensemble des idéaux premiers auxquels
$R \to S$ est lisse. Soit $V \subset \Spec(S')$ l'ensemble des idéaux premiers
auxquels $R' \to S'$ est lisse. Alors $V$ est l'image réciproque de $U$ par le
morphisme $f : \Spec(S') \to \Spec(S)$.
\end{lemma}

\begin{proof}
Le Lemme \ref{lemma-change-base-NL} montre que
$\NL_{S/R} \otimes_S S'$ est homotopiquement équivalent à $\NL_{S'/R'}$.
Cela implique déjà que $f^{-1}(U) \subset V$.

\medskip\noindent
Soit $\mathfrak q' \subset S'$ un idéal premier au-dessus de
$\mathfrak q \subset S$. Supposons $\mathfrak q' \in V$.
Il faut montrer que $\mathfrak q \in U$.
Comme $S \to S'$ est plat, $S_{\mathfrak q} \to S'_{\mathfrak q'}$ est
fidèlement plat (Lemme \ref{lemma-local-flat-ff}). L'annulation de
$H_1(L_{S'/R'})_{\mathfrak q'}$ implique donc celle de
$H_1(L_{S/R})_{\mathfrak q}$. En appliquant le
Lemme \ref{lemma-finite-projective-descends} au $S_{\mathfrak q}$-module
$(\Omega_{S/R})_{\mathfrak q}$ et au morphisme
$S_{\mathfrak q} \to S'_{\mathfrak q'}$, nous voyons que
$(\Omega_{S/R})_{\mathfrak q}$ est projectif. Ainsi $R \to S$ est lisse en
$\mathfrak q$ par le Lemme \ref{lemma-smooth-at-point}.
\end{proof}

\begin{lemma}
\label{lemma-smooth-field-change-local}
Soit $K/k$ une extension de corps.
Soit $S$ une algèbre de type fini sur $k$.
Soit $\mathfrak q_K$ un idéal premier de $S_K = K \otimes_k S$, et soit
$\mathfrak q$ l'idéal premier correspondant de $S$.
Alors $S$ est lisse sur $k$ en $\mathfrak q$ si et seulement si
$S_K$ est lisse en $\mathfrak q_K$ sur $K$.
\end{lemma}

\begin{proof}
C'est un cas particulier du Lemme \ref{lemma-flat-base-change-locus-smooth}.
\end{proof}

\begin{lemma}
\label{lemma-lift-smooth}
Soit $R$ un anneau et soit $I \subset R$ un idéal.
Soit $R/I \to \overline{S}$ un morphisme d'anneaux lisse.
Alors il existe des éléments $\overline{g}_i \in \overline{S}$ qui engendrent
l'idéal unité de $\overline{S}$, tels que chacun des
$\overline{S}_{g_i} \cong S_i/IS_i$ pour un certain anneau $S_i$ lisse
(standard) sur $R$.
\end{lemma}

\begin{proof}
Le Lemme \ref{lemma-smooth-syntomic} fournit une collection d'éléments
$\overline{g}_i \in \overline{S}$ qui engendrent l'idéal unité de
$\overline{S}$, telle que chacun des $\overline{S}_{g_i}$ soit lisse standard
sur $R/I$. Nous pouvons donc supposer que $\overline{S}$ est lisse standard
sur $R/I$. Écrivons
$\overline{S} =
(R/I)[x_1, \ldots, x_n]/(\overline{f}_1, \ldots, \overline{f}_c)$
comme dans la Définition \ref{definition-standard-smooth}.
Choisissons $f_1, \ldots, f_c \in R[x_1, \ldots, x_n]$ relevant
$\overline{f}_1, \ldots, \overline{f}_c$. Posons
$S = R[x_1, \ldots, x_n, x_{n + 1}]/(f_1, \ldots, f_c, x_{n + 1}\Delta - 1)$,
où $\Delta = \det(\frac{\partial f_j}{\partial x_i})_{i, j = 1, \ldots, c}$,
comme dans l'Exemple \ref{example-make-standard-smooth}.
Cela démontre le lemme.
\end{proof}






\section{Morphismes formellement lisses}
\label{section-formally-smooth}

\noindent
Dans cette section, nous définissons les morphismes d'anneaux formellement
lisses. Il s'avérera qu'un morphisme d'anneaux de présentation finie est
formellement lisse si et seulement s'il est lisse ; voir la
Proposition \ref{proposition-smooth-formally-smooth}.

\begin{definition}
\label{definition-formally-smooth}
Soit $R \to S$ un morphisme d'anneaux.
Nous disons que $S$ est {\it formellement lisse sur $R$} si, pour tout
diagramme commutatif aux flèches pleines
$$
\xymatrix{
S \ar[r] \ar@{-->}[rd] & A/I \\
R \ar[r] \ar[u] & A \ar[u]
}
$$
où $I \subset A$ est un idéal de carré nul, il existe une flèche en pointillés
qui rend le diagramme commutatif.
\end{definition}

\begin{lemma}
\label{lemma-base-change-fs}
Soit $R \to S$ un morphisme d'anneaux formellement lisse.
Soit $R \to R'$ un morphisme d'anneaux quelconque.
Alors le changement de base $S' = R' \otimes_R S$ est formellement lisse sur
$R'$.
\end{lemma}

\begin{proof}
Considérons un diagramme aux flèches pleines
$$
\xymatrix{
S \ar[r] \ar@{-->}[rrd] & R' \otimes_R S \ar[r] \ar@{-->}[rd] & A/I \\
R  \ar[u] \ar[r] & R' \ar[r] \ar[u] & A \ar[u]
}
$$
comme dans la Définition \ref{definition-formally-smooth}.
Par hypothèse, la flèche en pointillés la plus longue existe. La propriété
universelle du produit tensoriel fournit alors la plus courte.
\end{proof}

\begin{lemma}
\label{lemma-compose-formally-smooth}
La composée de morphismes d'anneaux formellement lisses est formellement lisse.
\end{lemma}

\begin{proof}
Omis. (Indication : c'est entièrement formel et résulte de la considération
d'un diagramme convenable.)
\end{proof}

\begin{lemma}
\label{lemma-polynomial-ring-formally-smooth}
Un anneau de polynômes sur $R$ est formellement lisse sur $R$.
\end{lemma}

\begin{proof}
Supposons donné un diagramme comme dans la
Définition \ref{definition-formally-smooth}, avec $S = R[x_j; j \in J]$.
Il existe alors une flèche en pointillés : il suffit de choisir des relèvements
$a_j \in A$ des éléments de $A/I$ sur lesquels les éléments $x_j$ sont envoyés
par la flèche horizontale supérieure.
\end{proof}

\begin{lemma}
\label{lemma-characterize-formally-smooth}
Soit $R \to S$ un morphisme d'anneaux.
Soit $P \to S$ un morphisme surjectif de $R$-algèbres d'un anneau de polynômes
$P$ sur $S$. Notons $J \subset P$ son noyau. Alors $R \to S$ est formellement
lisse si et seulement s'il existe un morphisme de $R$-algèbres
$\sigma : S \to P/J^2$ qui soit un inverse à droite de la surjection
$P/J^2 \to S$.
\end{lemma}

\begin{proof}
Supposons que $R \to S$ soit formellement lisse.
Considérons le diagramme commutatif
$$
\xymatrix{
S \ar[r] \ar@{-->}[rd] & P/J \\
R \ar[r] \ar[u] &  P/J^2\ar[u]
}
$$
Par hypothèse, la flèche en pointillés existe. Cela démontre l'existence de
$\sigma$.

\medskip\noindent
Réciproquement, supposons donné un $\sigma$ comme dans le lemme.
Considérons un diagramme aux flèches pleines
$$
\xymatrix{
S \ar[r] \ar@{-->}[rd] & A/I \\
R \ar[r] \ar[u] & A \ar[u]
}
$$
comme dans la Définition \ref{definition-formally-smooth}.
Puisque $P$ est formellement lisse par le
Lemme \ref{lemma-polynomial-ring-formally-smooth}, il existe un homomorphisme de
$R$-algèbres $\psi : P \to A$ qui relève le morphisme $P \to S \to A/I$.
Clairement, $\psi(J) \subset I$ et, comme $I^2 = 0$, nous en déduisons que
$\psi(J^2) = 0$. Ainsi $\psi$ se factorise en
$\overline{\psi} : P/J^2 \to A$. La flèche en pointillés recherchée est la
composée $\overline{\psi} \circ \sigma : S \to A$.
\end{proof}

\begin{remark}
\label{remark-lemma-characterize-formally-smooth}
Le Lemme \ref{lemma-characterize-formally-smooth} reste valable, plus
généralement, dès que $P$ est formellement lisse sur $R$.
\end{remark}

\begin{lemma}
\label{lemma-characterize-formally-smooth-again}
Soit $R \to S$ un morphisme d'anneaux.
Soit $P \to S$ un morphisme surjectif de $R$-algèbres d'un anneau de polynômes
$P$ sur $S$. Notons $J \subset P$ son noyau. Alors $R \to S$ est formellement
lisse si et seulement si la suite
$$
0 \to J/J^2 \to \Omega_{P/R} \otimes_P S \to \Omega_{S/R} \to 0
$$
du Lemme \ref{lemma-differential-seq} est une suite exacte scindée.
\end{lemma}

\begin{proof}
Supposons $S$ formellement lisse sur $R$. D'après le
Lemme \ref{lemma-characterize-formally-smooth}, cela signifie qu'il existe un
morphisme de $R$-algèbres $S \to P/J^2$ qui est un inverse à droite du morphisme
canonique $P/J^2 \to S$. Par le
Lemme \ref{lemma-differential-mod-power-ideal}, nous avons
$\Omega_{P/R} \otimes_P S = \Omega_{(P/J^2)/R} \otimes_{P/J^2} S$.
Le Lemme \ref{lemma-differential-seq-split} montre que la suite est scindée.

\medskip\noindent
Supposons que la suite exacte du lemme soit exacte scindée.
Choisissons un scindage $\sigma : \Omega_{S/R} \to
\Omega_{P/R} \otimes_P S$. Pour chaque $\lambda \in S$, choisissons
$x_\lambda \in P$ qui s'envoie sur $\lambda$. Puis, pour chaque
$\lambda \in S$, choisissons $f_\lambda \in J$ tel que
$$
\text{d}f_\lambda = \text{d}x_\lambda - \sigma(\text{d}\lambda)
$$
dans le terme médian de la suite exacte.
Nous affirmons que $s : \lambda \mapsto x_\lambda - f_\lambda \mod J^2$
est un homomorphisme de $R$-algèbres $s : S \to P/J^2$.
Pour le démontrer, nous utiliserons à plusieurs reprises que si $h \in J$ et
$\text{d}h = 0$ dans $\Omega_{P/R} \otimes_R S$, alors $h \in J^2$.
Soient $\lambda, \mu \in S$.
Alors $\sigma(\text{d}\lambda + \text{d}\mu - \text{d}(\lambda + \mu)) = 0$.
Cela implique
$$
\text{d}(x_\lambda + x_\mu - x_{\lambda + \mu}
- f_\lambda - f_\mu + f_{\lambda + \mu}) = 0
$$
et donc $x_\lambda + x_\mu - x_{\lambda + \mu}
- f_\lambda - f_\mu + f_{\lambda + \mu} \in J^2$, ce qui signifie à son tour
que $s(\lambda) + s(\mu) = s(\lambda + \mu)$.
De même,
$\sigma(\lambda \text{d}\mu + \mu \text{d}\lambda - \text{d}\lambda \mu) = 0$,
ce qui implique
$$
\mu(\text{d}x_\lambda - \text{d}f_\lambda) +
\lambda(\text{d}x_\mu - \text{d}f_\mu) -
\text{d}x_{\lambda\mu} + \text{d}f_{\lambda\mu} = 0
$$
dans le terme médian de la suite exacte.
De plus,
$$
\text{d}(x_\lambda x_\mu) =
x_\lambda \text{d}x_\mu + x_\mu \text{d}x_\lambda =
\lambda \text{d}x_\mu + \mu \text{d} x_\lambda
$$
dans le terme médian, là encore. Ensemble, ces équations signifient que
$x_\lambda x_\mu - x_{\lambda\mu}
- \mu f_\lambda - \lambda f_\mu + f_{\lambda\mu} \in J^2$ ; ainsi
$(x_\lambda - f_\lambda)(x_\mu - f_\mu) -
(x_{\lambda\mu} - f_{\lambda\mu}) \in J^2$, puisque
$f_\lambda f_\mu \in J^2$, ce qui signifie que
$s(\lambda)s(\mu) = s(\lambda\mu)$.
Si $\lambda \in R$, alors $\text{d}\lambda = 0$ et nous voyons que
$\text{d}f_\lambda = \text{d}x_\lambda$, d'où
$\lambda - x_\lambda + f_\lambda \in J^2$, puis
$s(\lambda) = \lambda$, comme souhaité. Nous pouvons maintenant appliquer le
Lemme \ref{lemma-characterize-formally-smooth} pour conclure que $S/R$ est
formellement lisse.
\end{proof}




\begin{proposition}
\label{proposition-characterize-formally-smooth}
Soit $R \to S$ un morphisme d'anneaux. Considérons une $R$-algèbre $P$
formellement lisse et une surjection $P \to S$ de noyau $J$. Les conditions
suivantes sont équivalentes :
\begin{enumerate}
\item $S$ est formellement lisse sur $R$,
\item pour un certain $P \to S$ comme ci-dessus, il existe une section de
$P/J^2 \to S$,
\item pour tout $P \to S$ comme ci-dessus, il existe une section de
$P/J^2 \to S$,
\item pour un certain $P \to S$ comme ci-dessus, la suite
$0 \to J/J^2 \to \Omega_{P/R} \otimes S \to \Omega_{S/R} \to 0$ est exacte
scindée,
\item pour tout $P \to S$ comme ci-dessus, la suite
$0 \to J/J^2 \to \Omega_{P/R} \otimes S \to \Omega_{S/R} \to 0$ est exacte
scindée, et
\item le complexe cotangent naïf $\NL_{S/R}$ est quasi-isomorphe à un
$S$-module projectif placé en degré $0$ : cela signifie que
$H_1(\NL_{S/R}) = 0$ et que $\Omega_{S/R}$ est un $S$-module projectif.
\end{enumerate}
\end{proposition}

\begin{proof}
Il est clair que (1) implique (3), qui implique (2) ; voir la première partie
de la démonstration du Lemme \ref{lemma-characterize-formally-smooth}.
Il est également vrai que (3) implique (5), qui implique (4), et que (2)
implique (4) ; voir la première partie de la démonstration du
Lemme \ref{lemma-characterize-formally-smooth-again}.
Enfin, le Lemme \ref{lemma-characterize-formally-smooth-again}, appliqué à la
surjection canonique $R[S] \to S$ (\ref{equation-canonical-presentation}),
montre que (1) implique (6).

\medskip\noindent
Supposons (4) et démontrons (6). Considérons la suite du
Lemme \ref{lemma-exact-sequence-NL} associée aux morphismes d'anneaux
$R \to P \to S$. Par l'implication (1) $\Rightarrow$ (6) démontrée ci-dessus,
nous voyons que $\NL_{P/R} \otimes_R S$ est quasi-isomorphe à
$\Omega_{P/R} \otimes_P S$ placé en degré $0$.
Ainsi $H_1(\NL_{P/R} \otimes_P S) = 0$. Puisque $P \to S$ est surjectif,
$\NL_{S/P}$ est homotopiquement équivalent à $J/J^2$ placé en degré $1$
(Lemme \ref{lemma-NL-surjection}). Nous obtenons donc la suite exacte
$0 \to H_1(L_{S/R}) \to J/J^2 \to \Omega_{P/R} \otimes_P S \to
\Omega_{S/R} \to 0$.
Par hypothèse, $H_1(L_{S/R}) = 0$ et $\Omega_{S/R}$ est un $S$-module
projectif. Ainsi (6) en résulte.

\medskip\noindent
Enfin, démontrons que (6) implique (1). L'hypothèse signifie que le complexe
$J/J^2 \to \Omega_{P/R} \otimes S$, où $P = R[S]$ et $P \to S$ est la
surjection canonique (\ref{equation-canonical-presentation}), est
quasi-isomorphe à un $S$-module projectif placé en degré $0$.
Le Lemme \ref{lemma-characterize-formally-smooth-again} montre donc que $S$
est formellement lisse sur $R$.
\end{proof}

\begin{lemma}
\label{lemma-ses-formally-smooth}
Soient $A \to B \to C$ des morphismes d'anneaux. Supposons que $B \to C$
soit formellement lisse. Alors la suite
$$
0 \to \Omega_{B/A} \otimes_B C \to \Omega_{C/A} \to \Omega_{C/B} \to 0
$$
du Lemme \ref{lemma-exact-sequence-differentials} est une suite exacte courte
scindée.
\end{lemma}

\begin{proof}
Cela résulte de la
Proposition \ref{proposition-characterize-formally-smooth} et du
Lemme \ref{lemma-exact-sequence-NL}.
\end{proof}

\begin{lemma}
\label{lemma-differential-seq-formally-smooth}
Soient $A \to B \to C$ des morphismes d'anneaux, avec $A \to C$ formellement
lisse et $B \to C$ surjectif de noyau $J \subset B$.
Alors la suite exacte
$$
0 \to J/J^2 \to \Omega_{B/A} \otimes_B C \to \Omega_{C/A} \to 0
$$
du Lemme \ref{lemma-differential-seq} est exacte scindée.
\end{lemma}

\begin{proof}
Cela résulte de la
Proposition \ref{proposition-characterize-formally-smooth}, du
Lemme \ref{lemma-exact-sequence-NL} et du
Lemme \ref{lemma-differential-seq}.
\end{proof}

\begin{lemma}
\label{lemma-application-NL-formally-smooth}
Soient $A \to B \to C$ des morphismes d'anneaux. Supposons que $A \to C$ soit
surjectif (de sorte que $B \to C$ le soit aussi) et que $A \to B$ soit
formellement lisse. Notons $I = \Ker(A \to C)$ et $J = \Ker(B \to C)$.
Alors la suite
$$
0 \to I/I^2 \to J/J^2 \to \Omega_{B/A} \otimes_B B/J \to 0
$$
du Lemme \ref{lemma-application-NL} est exacte scindée.
\end{lemma}

\begin{proof}
Comme $A \to B$ est formellement lisse, il existe un morphisme d'anneaux
$\sigma : B \to A/I^2$ dont la composée avec $A \to B$ est le morphisme
quotient $A \to A/I^2$. Alors $\sigma$ induit un morphisme
$J/J^2 \to I/I^2$ inverse du morphisme $I/I^2 \to J/J^2$.
\end{proof}

\begin{lemma}
\label{lemma-lift-formal-smoothness}
Soit $R \to S$ un morphisme d'anneaux.
Soit $I \subset R$ un idéal. Supposons que
\begin{enumerate}
\item $I^2 = 0$,
\item $R \to S$ est plat, et
\item $R/I \to S/IS$ est formellement lisse.
\end{enumerate}
Alors $R \to S$ est formellement lisse.
\end{lemma}

\begin{proof}
Supposons (1), (2) et (3).
Soit $P = R[\{x_t\}_{t \in T}] \to S$ une surjection de $R$-algèbres de
noyau $J$. Ainsi $0 \to J \to P \to S \to 0$ est une suite exacte courte de
$R$-modules plats. Cela implique $I \otimes_R S = IS$,
$I \otimes_R P = IP$ et $I \otimes_R J = IJ$, ainsi que $J \cap IP = IJ$.
Nous utiliserons dans toute la démonstration que
$$
\Omega_{(S/IS)/(R/I)} = \Omega_{S/R} \otimes_S (S/IS)
= \Omega_{S/R} \otimes_R R/I = \Omega_{S/R} / I\Omega_{S/R}
$$
et de même pour $P$ (voir le Lemme \ref{lemma-differentials-base-change}).
Par le Lemme \ref{lemma-characterize-formally-smooth-again}, la suite
\begin{equation}
\label{equation-split}
0 \to J/(IJ + J^2) \to
\Omega_{P/R} \otimes_P S/IS \to
\Omega_{S/R} \otimes_S S/IS \to 0
\end{equation}
est exacte scindée. Bien entendu, le terme médian est
$\bigoplus_{t \in T} S/IS \text{d}x_t$. Choisissons un scindage
$\sigma : \Omega_{P/R} \otimes_P S/IS \to J/(IJ + J^2)$.
Pour chaque $t \in T$, choisissons un élément $f_t \in J$ qui s'envoie sur
$\sigma(\text{d}x_t)$ dans $J/(IJ + J^2)$. Cela détermine un unique morphisme
de $S$-modules
$$
\tilde \sigma : \Omega_{P/R} \otimes_R S
= \bigoplus S\text{d}x_t \longrightarrow J/J^2
$$
tel que $\tilde\sigma(\text{d}x_t) = f_t$.
Comme $\sigma$ est une section de $\text{d}$, la différence
$$
\Delta = \text{id}_{J/J^2} - \tilde \sigma \circ \text{d}
$$
est un endomorphisme de $J/J^2 \to J/J^2$ dont l'image est contenue dans
$(IJ + J^2)/J^2$. En particulier, $\Delta((IJ + J^2)/J^2) = 0$, car
$I^2 = 0$. Cela signifie que $\Delta$ se factorise en
$$
J/J^2 \to J/(IJ + J^2) \xrightarrow{\overline{\Delta}}
(IJ + J^2)/J^2 \to J/J^2
$$
où $\overline{\Delta}$ est un morphisme de $S/IS$-modules.
En utilisant de nouveau le fait que la suite (\ref{equation-split}) est
scindée, nous pouvons trouver un morphisme de $S/IS$-modules
$\overline{\delta} : \Omega_{P/R} \otimes_P S/IS \to (IJ + J^2)/J^2$
tel que $\overline{\delta} \circ d$ soit égal à $\overline{\Delta}$.
De la même manière que ci-dessus, le morphisme $\overline{\delta}$ détermine
un morphisme de $S$-modules
$\delta : \Omega_{P/R} \otimes_P S \to J/J^2$.
Après remplacement de $\tilde \sigma$ par $\tilde \sigma + \delta$, un calcul
simple montre que $\Delta = 0$. Autrement dit, $\tilde \sigma$ est une section
de $J/J^2 \to \Omega_{P/R} \otimes_P S$.
Par le Lemme \ref{lemma-characterize-formally-smooth-again}, nous concluons que
$R \to S$ est formellement lisse.
\end{proof}

\begin{proposition}
\label{proposition-smooth-formally-smooth}
Soit $R \to S$ un morphisme d'anneaux. Les conditions suivantes sont
équivalentes :
\begin{enumerate}
\item $R \to S$ est de présentation finie et formellement lisse,
\item $R \to S$ est lisse.
\end{enumerate}
\end{proposition}

\begin{proof}
Cela résulte de la Proposition \ref{proposition-characterize-formally-smooth}
et de la Définition \ref{definition-smooth}.
(Notons que $\Omega_{S/R}$ est un $S$-module de présentation finie si
$R \to S$ est de présentation finie ; voir le
Lemme \ref{lemma-differentials-finitely-presented}.)
\end{proof}


\begin{lemma}
\label{lemma-finite-presentation-fs-Noetherian}
Soit $R \to S$ un morphisme d'anneaux lisse. Alors il existe un sous-anneau
$R_0 \subset R$ de type fini sur $\mathbf{Z}$ et un morphisme d'anneaux lisse
$R_0 \to S_0$ tels que $S \cong R \otimes_{R_0} S_0$.
\end{lemma}

\begin{proof}
Nous allons utiliser le fait que la lissité équivaut à la présentation finie
et à la lissité formelle ; voir la
Proposition \ref{proposition-smooth-formally-smooth}.
Écrivons $S = R[x_1, \ldots, x_n]/(f_1, \ldots, f_m)$ et notons
$I = (f_1, \ldots, f_m)$.
Choisissons un inverse à droite
$\sigma : S \to R[x_1, \ldots, x_n]/I^2$ de la projection sur $S$, comme dans
le Lemme \ref{lemma-characterize-formally-smooth}.
Choisissons $h_i \in R[x_1, \ldots, x_n]$ tels que
$\sigma(x_i \bmod I) = h_i \bmod I^2$. Puisque $x_i - h_i \in I$, il existe
des $b_{ij} \in R[x_1, \ldots, x_n]$ tels que
$$
x_i - h_i = \sum\nolimits_j b_{ij} f_j
$$
Le fait que $\sigma$ soit un homomorphisme de $R$-algèbres
$R[x_1, \ldots, x_n]/I \to R[x_1, \ldots, x_n]/I^2$ équivaut à la condition
$$
f_j(h_1, \ldots, h_n) = \sum\nolimits_{j_1 j_2} a_{j_1 j_2} f_{j_1} f_{j_2}
$$
pour certains $a_{kl} \in R[x_1, \ldots, x_n]$.
Soit $R_0 \subset R$ le sous-anneau engendré sur $\mathbf{Z}$ par tous les
coefficients des polynômes $f_j, h_i, a_{kl}, b_{ij}$.
Posons $S_0 = R_0[x_1, \ldots, x_n]/(f_1, \ldots, f_m)$, avec
$I_0 = (f_1, \ldots, f_m)$. Comme la seconde équation affichée est satisfaite
dans $R_0[x_1, \ldots, x_n]$, nous pouvons prendre pour
$\sigma_0 : S_0 \to R_0[x_1, \ldots, x_n]/I_0^2$ le morphisme de
$R_0$-algèbres défini par la règle $x_i \mapsto h_i \bmod I_0^2$.
Comme la première équation affichée est satisfaite dans
$R_0[x_1, \ldots, x_n]$, nous voyons que $\sigma_0$ est un inverse à droite de
la projection
$R_0[x_1, \ldots, x_n] / I_0^2 \to R_0[x_1, \ldots, x_n] / I_0 = S_0$.
Ainsi, par le Lemme \ref{lemma-characterize-formally-smooth}, l'anneau $S_0$
est formellement lisse sur $R_0$.
\end{proof}

\begin{lemma}
\label{lemma-smooth-descends-through-colimit}
Soit $A = \colim A_i$ une colimite filtrante d'anneaux. Soit $A \to B$ un
morphisme d'anneaux lisse. Il existe un $i$ et un morphisme d'anneaux lisse
$A_i \to B_i$ tels que $B = B_i \otimes_{A_i} A$.
\end{lemma}

\begin{proof}
Cela résulte du Lemme \ref{lemma-finite-presentation-fs-Noetherian}, car
$R_0 \to A$ se factorisera par $A_i$ pour un certain $i$, d'après le
Lemme \ref{lemma-characterize-finite-presentation}.
\end{proof}

\begin{lemma}
\label{lemma-descent-formally-smooth}
Soit $R \to S$ un morphisme d'anneaux. Soit $R \to R'$ un morphisme d'anneaux
fidèlement plat. Posons $S' = S \otimes_R R'$. Alors $R \to S$ est
formellement lisse si et seulement si $R' \to S'$ est formellement lisse.
\end{lemma}

\begin{proof}
Si $R \to S$ est formellement lisse, alors $R' \to S'$ l'est par le
Lemme \ref{lemma-base-change-fs}.
Pour démontrer la réciproque, supposons $R' \to S'$ formellement lisse.
Notons que $N \otimes_R R' = N \otimes_S S'$ pour tout $S$-module $N$.
En particulier, $S \to S'$ est lui aussi fidèlement plat.
Choisissons un anneau de polynômes $P = R[\{x_i\}_{i \in I}]$ et une surjection
de $R$-algèbres $P \to S$ de noyau $J$. Notons que
$P' = P \otimes_R R'$ est une algèbre de polynômes sur $R'$. Comme
$R \to R'$ est plat, le noyau $J'$ de la surjection $P' \to S'$ est
$J \otimes_R R'$. Par conséquent, la suite exacte scindée (voir le
Lemme \ref{lemma-characterize-formally-smooth-again})
$$
0 \to J'/(J')^2 \to \Omega_{P'/R'} \otimes_{P'} S' \to \Omega_{S'/R'} \to 0
$$
est le changement de base par $S \to S'$ de la suite correspondante
$$
J/J^2 \to \Omega_{P/R} \otimes_P S \to \Omega_{S/R} \to 0
$$
voir le Lemme \ref{lemma-differential-seq}.
Comme $S \to S'$ est fidèlement plat, nous en déduisons deux choses :
(1) cette suite (sans ${}'$) est elle aussi exacte et (2)
$\Omega_{S/R}$ est un $S$-module projectif. En effet, $\Omega_{S'/R'}$ est
projectif comme facteur direct du module libre
$\Omega_{P'/R'} \otimes_{P'} S'$, et
$\Omega_{S/R} \otimes_S {S'} = \Omega_{S'/R'}$ d'après ce que nous venons de
dire. Ainsi (2) résulte de la descente de la projectivité par les morphismes
d'anneaux fidèlement plats ; voir le
Théorème \ref{theorem-ffdescent-projectivity}.
La suite $0 \to J/J^2 \to \Omega_{P/R} \otimes_P S \to \Omega_{S/R} \to 0$
est donc elle aussi exacte, et nous concluons en appliquant de nouveau le
Lemme \ref{lemma-characterize-formally-smooth-again}.
\end{proof}

\noindent
Il s'avère que les morphismes d'anneaux lisses satisfont la propriété forte de
relèvement suivante.

\begin{lemma}
\label{lemma-smooth-strong-lift}
Soit $R \to S$ un morphisme d'anneaux lisse. Étant donné un diagramme
commutatif aux flèches pleines
$$
\xymatrix{
S \ar[r] \ar@{-->}[rd] & A/I \\
R \ar[r] \ar[u] & A \ar[u]
}
$$
où $I \subset A$ est un idéal localement nilpotent, il existe une flèche en
pointillés qui rend le diagramme commutatif.
\end{lemma}

\begin{proof}
Par le Lemme \ref{lemma-finite-presentation-fs-Noetherian}, nous pouvons
prolonger le diagramme en un diagramme commutatif
$$
\xymatrix{
S_0 \ar[r] & S \ar[r] \ar@{-->}[rd] & A/I \\
R_0 \ar[r] \ar[u] & R \ar[r] \ar[u] & A \ar[u]
}
$$
où $R_0 \to S_0$ est lisse, $R_0$ est de type fini sur $\mathbf{Z}$ et
$S = S_0 \otimes_{R_0} R$. Soient $x_1, \ldots, x_n \in S_0$ des générateurs
de $S_0$ sur $R_0$. Soient $a_1, \ldots, a_n$ des éléments de $A$ qui
s'envoient sur les mêmes éléments de $A/I$ que $x_1, \ldots, x_n$.
Notons $A_0 \subset A$ le sous-anneau engendré par l'image de $R_0$ et les
éléments $a_1, \ldots, a_n$. Posons $I_0 = A_0 \cap I$. Alors
$A_0/I_0 \subset A/I$, et $S_0 \to A/I$ est à valeurs dans $A_0/I_0$.
Il suffit donc de trouver la flèche en pointillés dans le diagramme
$$
\xymatrix{
S_0 \ar[r] \ar@{-->}[rd] & A_0/I_0 \\
R_0 \ar[r] \ar[u] & A_0 \ar[u]
}
$$
Par construction, l'anneau $A_0$ est de type fini sur $\mathbf{Z}$.
Ainsi $A_0$ est noethérien, donc $I_0$ est nilpotent ; voir le
Lemme \ref{lemma-Noetherian-power}.
Disons que $I_0^n = 0$. Par la
Proposition \ref{proposition-smooth-formally-smooth}, nous pouvons relever
successivement le morphisme de $R_0$-algèbres $S_0 \to A_0/I_0$ en
$S_0 \to A_0/I_0^2$, $S_0 \to A_0/I_0^3$, $\ldots$, et finalement en
$S_0 \to A_0/I_0^n = A_0$.
\end{proof}





\section{Lissité et différentielles}
\label{section-smooth-differential}

\noindent
Quelques résultats sur les différentielles et les morphismes d'anneaux lisses.

\begin{lemma}
\label{lemma-triangle-differentials-smooth}
Étant donnés des morphismes d'anneaux $A \to B \to C$, avec $B \to C$ lisse,
la suite
$$
0 \to C \otimes_B \Omega_{B/A} \to \Omega_{C/A} \to \Omega_{C/B} \to 0
$$
du Lemme \ref{lemma-exact-sequence-differentials} est exacte.
\end{lemma}

\begin{proof}
Cela résulte du Lemme \ref{lemma-ses-formally-smooth}, plus général, car un
morphisme d'anneaux lisse est formellement lisse ; voir la
Proposition \ref{proposition-smooth-formally-smooth}.
Mais cela résulte aussi directement du
Lemme \ref{lemma-exact-sequence-NL}, puisque $H_1(L_{C/B}) = 0$ fait partie de
la définition de la lissité de $B \to C$.
\end{proof}

\begin{lemma}
\label{lemma-differential-seq-smooth}
Soient $A \to B \to C$ des morphismes d'anneaux, avec $A \to C$ lisse et
$B \to C$ surjectif de noyau $J \subset B$.
Alors la suite exacte
$$
0 \to J/J^2 \to \Omega_{B/A} \otimes_B C \to \Omega_{C/A} \to 0
$$
du Lemme \ref{lemma-differential-seq} est exacte scindée.
\end{lemma}

\begin{proof}
Cela résulte du Lemme \ref{lemma-differential-seq-formally-smooth}, plus
général, car un morphisme d'anneaux lisse est formellement lisse ; voir la
Proposition \ref{proposition-smooth-formally-smooth}.
\end{proof}

\begin{lemma}
\label{lemma-application-NL-smooth}
Soient $A \to B \to C$ des morphismes d'anneaux. Supposons que $A \to C$ soit
surjectif (de sorte que $B \to C$ le soit aussi) et que $A \to B$ soit lisse.
Notons $I = \Ker(A \to C)$ et $J = \Ker(B \to C)$.
Alors la suite
$$
0 \to I/I^2 \to J/J^2 \to \Omega_{B/A} \otimes_B B/J \to 0
$$
du Lemme \ref{lemma-application-NL} est exacte.
\end{lemma}

\begin{proof}
Cela résulte du Lemme \ref{lemma-application-NL-formally-smooth}, plus général,
car un morphisme d'anneaux lisse est formellement lisse ; voir la
Proposition \ref{proposition-smooth-formally-smooth}.
\end{proof}

\begin{lemma}
\label{lemma-section-smooth}
\begin{slogan}
Si $R$ est un facteur direct de $S$ et si $S$ est lisse sur $R$, alors le
complété $I$-adique de $S$ est souvent un anneau de séries formelles sur $R$,
où $I$ est le noyau du morphisme de projection de $S$ sur $R$.
\end{slogan}
Soit $\varphi : R \to S$ un morphisme d'anneaux lisse.
Soit $\sigma : S \to R$ un inverse à gauche de $\varphi$.
Posons $I = \Ker(\sigma)$. Alors
\begin{enumerate}
\item $I/I^2$ est un $R$-module localement libre de type fini, et
\item si $I/I^2$ est libre, alors $S^\wedge \cong R[[t_1, \ldots, t_d]]$
comme $R$-algèbres, où $S^\wedge$ est le complété $I$-adique de $S$.
\end{enumerate}
\end{lemma}

\begin{proof}
En appliquant le Lemme \ref{lemma-differential-seq-split} à
$R \to S \to R$, nous voyons que
$I/I^2 = \Omega_{S/R} \otimes_{S, \sigma} R$.
Puisque, par définition d'un morphisme lisse, le module $\Omega_{S/R}$ est
localement libre de type fini sur $S$, nous en déduisons (1).
Si $I/I^2$ est libre, choisissons $f_1, \ldots, f_d \in I$ dont les images dans
$I/I^2$ forment une $R$-base. Considérons le morphisme de $R$-algèbres défini
par
$$
\Psi : R[[x_1, \ldots, x_d]] \longrightarrow S^\wedge, \quad
x_i \longmapsto f_i.
$$
Notons $P = R[[x_1, \ldots, x_d]]$ et $J = (x_1, \ldots, x_d) \subset P$.
Notons $\Psi_n : P/J^n \to S/I^n$ le morphisme induit entre les anneaux
quotients. Remarquons que $S/I^2 = \varphi(R) \oplus I/I^2$. Ainsi $\Psi_2$
est un isomorphisme. Notons $\sigma_2 : S/I^2 \to P/J^2$ l'inverse de
$\Psi_2$. Nous allons démontrer par récurrence sur $n$ que, pour tout $n > 2$,
il existe un inverse $\sigma_n : S/I^n \to P/J^n$ de $\Psi_n$. En effet,
comme $S$ est formellement lisse sur $R$ (par la
Proposition \ref{proposition-smooth-formally-smooth}), nous voyons que, dans le
diagramme aux flèches pleines
$$
\xymatrix{
S \ar@{..>}[r] \ar[rd]_{\sigma_{n - 1}} & P/J^n \ar[d] \\
 & P/J^{n - 1}
}
$$
de $R$-algèbres, nous pouvons compléter la flèche en pointillés par un certain
morphisme de $R$-algèbres $\tau : S \to P/J^n$ qui rende le diagramme
commutatif. Celui-ci induit un morphisme de $R$-algèbres
$\overline{\tau} : S/I^n \to P/J^n$ égal à $\sigma_{n - 1}$ modulo
$J^{n - 1}$. Par construction, le morphisme $\Psi_n$ est surjectif, et
$\overline{\tau} \circ \Psi_n$ est maintenant un endomorphisme de
$R$-algèbres de $P/J^n$ qui envoie $x_i$ sur $x_i + \delta_{i, n}$, avec
$\delta_{i, n} \in J^{n - 1}/J^n$. Il s'ensuit que $\Psi_n$ est un
isomorphisme et possède donc un inverse $\sigma_n$.
Cela démontre le lemme.
\end{proof}







\section{Algèbres lisses sur les corps}
\label{section-smooth-over-field}

\noindent
Avertissement : en général, les deux lemmes suivants ne sont pas valables sur
les corps non parfaits.

\begin{lemma}
\label{lemma-rank-omega}
Soit $k$ un corps algébriquement clos.
Soit $S$ une $k$-algèbre de type fini.
Soit $\mathfrak m \subset S$ un idéal maximal.
Alors
$$
\dim_{\kappa(\mathfrak m)} \Omega_{S/k} \otimes_S \kappa(\mathfrak m)
=
\dim_{\kappa(\mathfrak m)} \mathfrak m/\mathfrak m^2.
$$
\end{lemma}

\begin{proof}
Considérons la suite exacte
$$
\mathfrak m/\mathfrak m^2 \to
\Omega_{S/k} \otimes_S \kappa(\mathfrak m) \to
\Omega_{\kappa(\mathfrak m)/k} \to 0
$$
du Lemme \ref{lemma-differential-seq}. Nous voulons montrer que le premier
morphisme est un isomorphisme. Comme $k$ est algébriquement clos, la composée
$k \to \kappa(\mathfrak m)$ est un isomorphisme par le
Théorème \ref{theorem-nullstellensatz}.
Ainsi, la surjection $S \to \kappa(\mathfrak m)$ admet une section comme
morphisme de $k$-algèbres, et le Lemme \ref{lemma-differential-seq-split}
montre que la suite ci-dessus est exacte à gauche. Puisque
$\Omega_{\kappa(\mathfrak m)/k} = 0$, nous avons terminé.
\end{proof}

\begin{lemma}
\label{lemma-characterize-smooth-kbar}
Soit $k$ un corps algébriquement clos.
Soit $S$ une $k$-algèbre de type fini.
Soit $\mathfrak m \subset S$ un idéal maximal.
Les conditions suivantes sont équivalentes :
\begin{enumerate}
\item L'anneau $S_{\mathfrak m}$ est un anneau local régulier.
\item Nous avons
$\dim_{\kappa(\mathfrak m)} \Omega_{S/k} \otimes_S \kappa(\mathfrak m)
\leq \dim(S_{\mathfrak m})$.
\item Nous avons
$\dim_{\kappa(\mathfrak m)} \Omega_{S/k} \otimes_S \kappa(\mathfrak m)
= \dim(S_{\mathfrak m})$.
\item Il existe $g \in S$, $g \not \in \mathfrak m$, tel que $S_g$ soit lisse
sur $k$. Autrement dit, $S/k$ est lisse en $\mathfrak m$.
\end{enumerate}
\end{lemma}

\begin{proof}
Notons que (1), (2) et (3) sont équivalentes par le
Lemme \ref{lemma-rank-omega} et la Définition \ref{definition-regular}.

\medskip\noindent
Supposons que $S$ soit lisse en $\mathfrak m$.
Le Lemme \ref{lemma-smooth-syntomic} montre que $S_g$ est lisse standard sur
$k$ pour un $g \in S$, $g \not \in \mathfrak m$, convenable.
Par le Lemme \ref{lemma-standard-smooth}, $\Omega_{S_g/k}$ est donc libre de
rang $\dim(S_g)$. Le Lemme \ref{lemma-rank-omega} montre alors que
$\dim(S_{\mathfrak m}) = \dim (\mathfrak m/\mathfrak m^2)$ ; autrement dit,
$S_\mathfrak m$ est régulier.

\medskip\noindent
Réciproquement, supposons que $S_{\mathfrak m}$ soit régulier.
Soit $d = \dim(S_{\mathfrak m}) = \dim \mathfrak m/\mathfrak m^2$.
Choisissons une présentation $S = k[x_1, \ldots, x_n]/I$ telle que $x_i$
s'envoie sur un élément de $\mathfrak m$ pour tout $i$. Autrement dit,
$\mathfrak m'' = (x_1, \ldots, x_n)$ est l'idéal maximal correspondant de
$k[x_1, \ldots, x_n]$. Notons que nous avons une suite exacte courte
$$
I/\mathfrak m''I \to \mathfrak m''/(\mathfrak m'')^2
\to \mathfrak m/(\mathfrak m)^2 \to 0
$$
Choisissons $c = n - d$ éléments $f_1, \ldots, f_c \in I$ dont les images
dans $\mathfrak m''/(\mathfrak m'')^2$ engendrent le noyau du morphisme vers
$\mathfrak m/\mathfrak m^2$. C'est clairement possible. Notons
$J = (f_1, \ldots, f_c)$. Ainsi $J \subset I$.
Notons $S' = k[x_1, \ldots, x_n]/J$, de sorte qu'il existe une surjection
$S' \to S$. Notons $\mathfrak m' = \mathfrak m''S'$ l'idéal maximal
correspondant de $S'$. Nous avons donc
$$
\xymatrix{
k[x_1, \ldots, x_n] \ar[r] & S' \ar[r] & S \\
\mathfrak m'' \ar[u] \ar[r] & \mathfrak m' \ar[r] \ar[u] &
\mathfrak m \ar[u]
}
$$
Par notre choix de $J$, la suite exacte
$$
J/\mathfrak m''J \to \mathfrak m''/(\mathfrak m'')^2
\to \mathfrak m'/(\mathfrak m')^2 \to 0
$$
montre que $\dim( \mathfrak m'/(\mathfrak m')^2 ) = d$.
Comme $S'_{\mathfrak m'}$ se surjecte sur $S_{\mathfrak m}$, nous voyons que
$\dim(S_{\mathfrak m'}) \geq d$. La discussion précédant la
Définition \ref{definition-regular-local} permet donc de conclure que
$S'_{\mathfrak m'}$ est lui aussi régulier de dimension $d$. Comme $S'$ était
découpé par $c = n - d$ équations, nous concluons qu'il existe
$g' \in S'$, $g' \not \in \mathfrak m'$, tel que $S'_{g'}$ soit une
intersection complète globale sur $k$ ; voir le Lemme \ref{lemma-lci}.
De plus, le morphisme $S'_{\mathfrak m'} \to S_{\mathfrak m}$ est une
surjection de domaines locaux noethériens de même dimension, et donc un
isomorphisme. Ainsi $S' \to S$ est surjectif, son noyau est de type fini, et il
devient un isomorphisme après localisation en $\mathfrak m'$. Nous pouvons donc
trouver $g' \in S'$, $g \not \in \mathfrak m'$, tel que
$S'_{g'} \to S_{g'}$ soit un isomorphisme. En définitive, nous concluons
qu'après remplacement de $S$ par une localisation principale, nous pouvons
supposer que $S$ est une intersection complète globale.

\medskip\noindent
À ce stade, nous pouvons écrire $S = k[x_1, \ldots, x_n]/(f_1, \ldots, f_c)$,
avec $\dim S = n - c$. Rappelons que le complexe cotangent naïf de cette
algèbre est donné par
$$
\bigoplus S \cdot f_j
\to
\bigoplus S \cdot \text{d}x_i
$$
voir le Lemme \ref{lemma-relative-global-complete-intersection-conormal}.
Par le Lemme \ref{lemma-relative-global-complete-intersection-smooth}, pour
montrer que $S$ est lisse en $\mathfrak m$, nous devons montrer que l'un des
mineurs $c \times c$, notés $g_I$, de la matrice « $A$ » qui donne le morphisme
ci-dessus ne s'annule pas en $\mathfrak m$. Par le
Lemme \ref{lemma-rank-omega}, la matrice $A \bmod \mathfrak m$ est de rang $c$.
Nous avons donc terminé.
\end{proof}
\begin{lemma}
\label{lemma-characterize-smooth-over-field}
Soit $k$ un corps quelconque.
Soit $S$ une $k$-algèbre de type fini.
Soit $X = \Spec(S)$.
Soit $\mathfrak q \subset S$ un idéal premier correspondant à $x \in X$.
Les conditions suivantes sont équivalentes :
\begin{enumerate}
\item La $k$-algèbre $S$ est lisse en $\mathfrak q$ sur $k$.
\item Nous avons
$\dim_{\kappa(\mathfrak q)} \Omega_{S/k} \otimes_S \kappa(\mathfrak q)
\leq \dim_x X$.
\item Nous avons
$\dim_{\kappa(\mathfrak q)} \Omega_{S/k} \otimes_S \kappa(\mathfrak q)
= \dim_x X$.
\end{enumerate}
De plus, dans ce cas, l'anneau local $S_{\mathfrak q}$ est régulier.
\end{lemma}

\begin{proof}
Si $S$ est lisse en $\mathfrak q$ sur $k$, alors il existe
$g \in S$, $g \not \in \mathfrak q$, tel que $S_g$ soit lisse standard sur
$k$ ; voir le Lemme \ref{lemma-smooth-syntomic}. Le module des différentielles
d'une algèbre lisse standard sur $k$ est libre, de rang égal à la dimension ;
voir le Lemme \ref{lemma-standard-smooth} (utiliser qu'une intersection
complète globale relative sur un corps a pour dimension le nombre de variables
moins le nombre d'équations). Nous voyons ainsi que (1) implique (3). Pour
achever la démonstration du lemme, il suffit de montrer que (2) implique (1),
et qu'il implique que $S_{\mathfrak q}$ est régulier.

\medskip\noindent
Supposons (2). Par le lemme de Nakayama \ref{lemma-NAK}, nous voyons que
$\Omega_{S/k, \mathfrak q}$ peut être engendré par
$\leq \dim_x X$ éléments. Nous pouvons remplacer $S$ par $S_g$ pour un certain
$g \in S$, $g \not \in \mathfrak q$, de sorte que $\Omega_{S/k}$ soit engendré
par au plus $\dim_x X$ éléments.
Soit $K/k$ une extension de corps algébriquement close telle qu'il existe un
morphisme de $k$-algèbres $\psi : \kappa(\mathfrak q) \to K$.
Considérons $S_K = K \otimes_k S$. Soit $\mathfrak m \subset S_K$ l'idéal
maximal correspondant à la surjection
$$
\xymatrix{
S_K = K \otimes_k S \ar[r] &
K \otimes_k \kappa(\mathfrak q)
\ar[r]^-{\text{id}_K \otimes \psi} &
K.
}
$$
Notons que $\mathfrak m \cap S = \mathfrak q$ ; autrement dit,
$\mathfrak m$ est au-dessus de $\mathfrak q$.
Par le Lemme \ref{lemma-dimension-at-a-point-preserved-field-extension}, la
dimension de $X_K = \Spec(S_K)$ au point correspondant à $\mathfrak m$ est
$\dim_x X$. Par le
Lemme \ref{lemma-dimension-closed-point-finite-type-field}, elle est égale à
$\dim((S_K)_{\mathfrak m})$.
Par le Lemme \ref{lemma-differentials-base-change}, le module des
différentielles de $S_K$ sur $K$ est le changement de base de
$\Omega_{S/k}$ ; il est donc lui aussi engendré par au plus
$\dim_x X = \dim((S_K)_{\mathfrak m})$ éléments. Par le
Lemme \ref{lemma-characterize-smooth-kbar}, nous voyons que $S_K$ est lisse en
$\mathfrak m$ sur $K$. Le Lemme \ref{lemma-flat-base-change-locus-smooth}
implique alors que $S$ est lisse en $\mathfrak q$ sur $k$.
Cela démontre (1). De plus, nous savons par le
Lemme \ref{lemma-characterize-smooth-kbar} que l'anneau local
$(S_K)_{\mathfrak m}$ est régulier. Comme
$S_{\mathfrak q} \to (S_K)_{\mathfrak m}$ est plat, nous concluons du
Lemme \ref{lemma-flat-under-regular} que $S_{\mathfrak q}$ est régulier.
\end{proof}

\noindent
Le lemme suivant se généralise considérablement (de plusieurs façons
différentes).

\begin{lemma}
\label{lemma-computation-differential}
Soit $k$ un corps.
Soit $R$ un anneau local noethérien contenant $k$.
Supposons que le corps résiduel $\kappa = R/\mathfrak m$ soit une extension
séparable de type fini de $k$. Alors le morphisme
$$
\text{d} :
\mathfrak m/\mathfrak m^2
\longrightarrow
\Omega_{R/k} \otimes_R \kappa(\mathfrak m)
$$
est injectif.
\end{lemma}

\begin{proof}
Nous pouvons remplacer $R$ par $R/\mathfrak m^2$. Nous pouvons donc supposer
que $\mathfrak m^2 = 0$. Par hypothèse, nous pouvons écrire
$\kappa = k(\overline{x}_1, \ldots, \overline{x}_r, \overline{y})$, où
$\overline{x}_1, \ldots, \overline{x}_r$ est une base de transcendance de
$\kappa$ sur $k$, et $\overline{y}$ est séparable algébrique sur
$k(\overline{x}_1, \ldots, \overline{x}_r)$. Soit $P(T)$ dans
$k(\overline{x}_1, \ldots, \overline{x}_r)[T]$ le polynôme minimal de
$\overline{y}$ sur $k(\overline{x}_1, \ldots, \overline{x}_r)$.
Alors $P(\overline{y}) = 0$, mais $P'(\overline{y}) \not = 0$. Choisissons des
relèvements quelconques $x_i \in R$ des éléments
$\overline{x}_i \in \kappa$. Cela donne un diagramme commutatif
$$
\xymatrix{
R \ar[r] & \kappa \\
& k(\overline{x}_1, \ldots, \overline{x}_r) \ar[lu]^\varphi \ar[u]
}
$$
de $k$-algèbres. Nous voulons prolonger la flèche ascendante gauche
$\varphi$ en un morphisme de $k$-algèbres de $\kappa$ vers $R$. Pour ce faire,
choisissons un $y \in R$ quelconque relevant $\overline{y}$. Pour voir qu'il
définit un morphisme de $k$-algèbres défini sur
$\kappa \cong k(\overline{x}_1, \ldots, \overline{x}_r)[T]/(P)$, il suffit de
montrer que nous pouvons choisir $y$ tel que $P^\varphi(y) = 0$. Sinon, pour
$\delta \in \mathfrak m$, nous calculons
$$
P^\varphi(y + \delta) = P^\varphi(y) + (P')^\varphi(y)\delta
$$
car $\mathfrak m^2 = 0$. Puisque $(P')^\varphi(y)$ est une unité de $R$, nous
voyons que nous pouvons ajuster notre choix comme souhaité.
Cela montre que $R \cong \kappa \oplus \mathfrak m$ comme $k$-algèbres !
Un calcul direct de $\Omega_{\kappa \oplus \mathfrak m/k}$ ou une application
du Lemme \ref{lemma-differential-seq-split} achève alors la démonstration.
\end{proof}

\begin{lemma}
\label{lemma-separable-smooth}
Soit $k$ un corps.
Soit $S$ une $k$-algèbre de type fini.
Soit $\mathfrak q \subset S$ un idéal premier.
Supposons $\kappa(\mathfrak q)$ séparable sur $k$.
Les conditions suivantes sont équivalentes :
\begin{enumerate}
\item L'algèbre $S$ est lisse en $\mathfrak q$ sur $k$.
\item L'anneau $S_{\mathfrak q}$ est régulier.
\end{enumerate}
\end{lemma}

\begin{proof}
Notons $R = S_{\mathfrak q}$, son idéal maximal $\mathfrak m$ et son corps
résiduel $\kappa$.
Le Lemme \ref{lemma-computation-differential} et
\ref{lemma-differential-seq} montrent qu'il existe une suite exacte courte
$$
0 \to \mathfrak m/\mathfrak m^2 \to
\Omega_{R/k} \otimes_R \kappa \to
\Omega_{\kappa/k} \to 0
$$
Notons que $\Omega_{R/k} = \Omega_{S/k, \mathfrak q}$ ; voir le
Lemme \ref{lemma-differentials-localize}.
De plus, puisque $\kappa$ est séparable sur $k$, nous avons
$\dim_{\kappa} \Omega_{\kappa/k} = \text{trdeg}_k(\kappa)$.
Nous obtenons donc
$$
\dim_{\kappa} \Omega_{R/k} \otimes_R \kappa
=
\dim_\kappa \mathfrak m/\mathfrak m^2 + \text{trdeg}_k (\kappa)
\geq
\dim R + \text{trdeg}_k (\kappa)
=
\dim_{\mathfrak q} S
$$
(voir le Lemme \ref{lemma-dimension-at-a-point-finite-type-field} pour la
dernière égalité), avec égalité si et seulement si $R$ est régulier.
Nous concluons donc en appliquant le
Lemme \ref{lemma-characterize-smooth-over-field}.
\end{proof}
\begin{lemma}
\label{lemma-characteristic-zero}
Soit $R \to S$ un morphisme de $\mathbf{Q}$-algèbres.
Soit $f \in S$ tel que $\Omega_{S/R} = S \text{d}f \oplus C$ pour un certain
sous-$S$-module $C$. Alors
\begin{enumerate}
\item $f$ n'est pas nilpotent, et
\item si $S$ est un anneau local noethérien, alors $f$ est non-diviseur de zéro
dans $S$.
\end{enumerate}
\end{lemma}

\begin{proof}
Pour $a \in S$, écrivons $\text{d}(a) = \theta(a)\text{d}f + c(a)$, avec
$\theta(a) \in S$ et $c(a) \in C$.
Considérons la $R$-dérivation $S \to S$, $a \mapsto \theta(a)$.
Notons que $\theta(f) = 1$.

\medskip\noindent
Si $f^n = 0$, avec $n > 1$ minimal, alors
$0 = \theta(f^n) = n f^{n - 1}$, ce qui contredit la minimalité de $n$.
Nous concluons que $f$ n'est pas nilpotent.

\medskip\noindent
Supposons $fa = 0$. Si $f$ est une unité, alors $a = 0$ et nous avons terminé.
Supposons que $f$ ne soit pas une unité. Alors
$0 = \theta(fa) = f\theta(a) + a$ par la règle de Leibniz, et donc
$a \in (f)$. Supposons avoir démontré par récurrence que
$fa = 0 \Rightarrow a \in (f^n)$. En écrivant $a = f^nb$, nous obtenons
$0 = \theta(f^{n + 1}b) = (n + 1)f^nb + f^{n + 1}\theta(b)$.
Ainsi $a = f^n b = -f^{n + 1}\theta(b)/(n + 1) \in (f^{n + 1})$.
Comme, dans l'anneau local noethérien $S$, nous avons
$\bigcap (f^n) = 0$ ; voir le
Lemme \ref{lemma-intersect-powers-ideal-module-zero}, nous avons terminé.
\end{proof}

\noindent
Le résultat suivant est probablement assez inutile dans les applications.

\begin{lemma}
\label{lemma-characteristic-zero-local-smooth}
Soit $k$ un corps de caractéristique $0$.
Soit $S$ une $k$-algèbre de type fini.
Soit $\mathfrak q \subset S$ un idéal premier.
Les conditions suivantes sont équivalentes :
\begin{enumerate}
\item L'algèbre $S$ est lisse en $\mathfrak q$ sur $k$.
\item Le $S_{\mathfrak q}$-module $\Omega_{S/k, \mathfrak q}$ est libre (de
type fini).
\item L'anneau $S_{\mathfrak q}$ est régulier.
\end{enumerate}
\end{lemma}

\begin{proof}
En caractéristique zéro, toute extension de corps est séparable ; l'équivalence
de (1) et (3) résulte donc du Lemme \ref{lemma-separable-smooth}.
De plus, (1) implique (2) par définition des algèbres lisses.
Supposons que $\Omega_{S/k, \mathfrak q}$ soit libre sur $S_{\mathfrak q}$.
Nous allons utiliser les notations et observations de la démonstration du
Lemme \ref{lemma-separable-smooth}. Ainsi, $R = S_{\mathfrak q}$, d'idéal
maximal $\mathfrak m$ et de corps résiduel $\kappa$.
Notre but est de démontrer que $R$ est régulier.

\medskip\noindent
Si $\mathfrak m/\mathfrak m^2 = 0$, alors $\mathfrak m = 0$ et
$R \cong \kappa$. Ainsi $R$ est régulier et nous avons terminé.

\medskip\noindent
Si $\mathfrak m/ \mathfrak m^2 \not = 0$, choisissons un
$f \in \mathfrak m$ quelconque dont l'image dans
$\mathfrak m/ \mathfrak m^2$ n'est pas nulle. Par le
Lemme \ref{lemma-computation-differential}, nous voyons que $\text{d}f$ a une
image non nulle dans $\Omega_{R/k}/\mathfrak m\Omega_{R/k}$.
Par hypothèse, $\Omega_{R/k} = \Omega_{S/k, \mathfrak q}$ est libre de type
fini et, par le lemme de Nakayama \ref{lemma-NAK}, nous voyons donc que
$\text{d}f$ engendre un facteur direct. Nous appliquons le
Lemme \ref{lemma-characteristic-zero} pour en déduire que $f$ est un
non-diviseur de zéro dans $R$.
En outre, le Lemme \ref{lemma-differential-seq} donne une suite exacte
$$
(f)/(f^2) \to \Omega_{R/k} \otimes_R R/fR \to \Omega_{(R/fR)/k} \to 0
$$
Cela implique que $\Omega_{(R/fR)/k}$ est lui aussi libre de type fini.
Par récurrence, nous voyons donc que $R/fR$ est un anneau local régulier.
Puisque $f \in \mathfrak m$ était un non-diviseur de zéro, nous concluons que
$R$ est régulier ; voir le Lemme \ref{lemma-regular-mod-x}.
\end{proof}

\begin{example}
\label{example-characteristic-p}
Le Lemme \ref{lemma-characteristic-zero-local-smooth} n'est pas valable en
caractéristique $p > 0$.
Les exemples standards sont les morphismes d'anneaux
$$
\mathbf{F}_p \longrightarrow \mathbf{F}_p[x]/(x^p)
$$
dont le module des différentielles est libre mais qui ne sont clairement pas
lisses, et le morphisme d'anneaux ($p > 2$)
$$
\mathbf{F}_p(t) \to \mathbf{F}_p(t)[x, y]/(x^p + y^2 + \alpha)
$$
qui n'est pas lisse en l'idéal premier
$\mathfrak q = (y, x^p + \alpha)$, mais qui est régulier.
\end{example}

\noindent
À l'aide de ce qui précède, nous pouvons caractériser la lissité au point
générique en termes d'extensions de corps.

\begin{lemma}
\label{lemma-smooth-at-generic-point}
Soit $R \to S$ un morphisme d'anneaux injectif et de type fini, où $R$ et $S$
sont des domaines. Alors $R \to S$ est lisse en $\mathfrak q = (0)$ si et
seulement si l'extension induite $L/K$ des corps de fractions est séparable.
\end{lemma}

\begin{proof}
Supposons $R \to S$ lisse en $(0)$. Nous pouvons remplacer $S$ par $S_g$ pour
un certain $g \in S$ non nul et supposer que $R \to S$ est lisse.
Alors $K \to S \otimes_R K$ est lisse
(Lemme \ref{lemma-base-change-smooth}). De plus, pour toute extension de corps
$K'/K$, le morphisme d'anneaux $K' \to S \otimes_R K'$ est lui aussi lisse.
Le Lemme \ref{lemma-characterize-smooth-over-field} montre donc que
$S \otimes_R K'$ est un anneau régulier, en particulier réduit.
Il s'ensuit que $S \otimes_R K$ est géométriquement réduit sur $K$.
Ainsi $L$ est géométriquement réduit sur $K$ ; voir le
Lemme \ref{lemma-geometrically-reduced-permanence}.
Ainsi $L/K$ est séparable par le
Lemme \ref{lemma-characterize-separable-field-extensions}.

\medskip\noindent
Réciproquement, supposons que $L/K$ soit séparable.
Nous pouvons supposer que $R \to S$ est de présentation finie ; voir le
Lemme \ref{lemma-generic-finite-presentation}.
Il suffit de démontrer que $K \to S \otimes_R K$ est lisse en $(0)$ ; voir le
Lemme \ref{lemma-flat-base-change-locus-smooth}.
Cela résulte du Lemme \ref{lemma-separable-smooth}, du fait qu'un corps est un
anneau régulier, et de l'hypothèse que $L/K$ est séparable.
\end{proof}







\section{Morphismes d'anneaux lisses dans le cas noethérien}
\label{section-smooth-Noetherian}

\begin{definition}
\label{definition-small-extension}
Soit $\varphi : B' \to B$ un morphisme d'anneaux.
Nous disons que $\varphi$ est une {\it petite extension} si $B'$ et $B$ sont des
anneaux locaux artiniens, si $\varphi$ est surjectif et si
$I = \Ker(\varphi)$ est de longueur $1$ comme $B'$-module.
\end{definition}

\noindent
Clairement, cela signifie que $I^2 = 0$ et que $I = (x)$ pour un certain
$x \in B'$ tel que $\mathfrak m' x = 0$, où $\mathfrak m' \subset B'$ est
l'idéal maximal.

\begin{lemma}
\label{lemma-smooth-test-artinian}
Soit $R \to S$ un morphisme d'anneaux. Soit $\mathfrak q$ un idéal premier de
$S$ au-dessus de $\mathfrak p \subset R$. Supposons que $R$ soit noethérien et
que $R \to S$ soit de type fini.
Les conditions suivantes sont équivalentes :
\begin{enumerate}
\item $R \to S$ est lisse en $\mathfrak q$,
\item pour toute surjection de $R$-algèbres locales
$(B', \mathfrak m') \to (B, \mathfrak m)$ telle que
$\Ker(B' \to B)$ soit de carré nul, et pour tout diagramme commutatif aux
flèches pleines
$$
\xymatrix{
S \ar[r] \ar@{-->}[rd] & B \\
R \ar[r] \ar[u] & B' \ar[u]
}
$$
tel que $\mathfrak q = S \cap \mathfrak m$, il existe une flèche en pointillés
qui rend le diagramme commutatif,
\item comme dans (2), mais où $B' \to B$ parcourt les petites extensions, et
\item comme dans (2), mais où $B' \to B$ parcourt les petites extensions telles
que, de plus, $S \to B$ induise un isomorphisme
$\kappa(\mathfrak q) \cong \kappa(\mathfrak m)$.
\end{enumerate}
\end{lemma}

\begin{proof}
Supposons (1). Cela signifie qu'il existe $g \in S$,
$g \not \in \mathfrak q$, tel que $R \to S_g$ soit lisse. Par la
Proposition \ref{proposition-smooth-formally-smooth}, nous savons que
$R \to S_g$ est formellement lisse. Notons que, dans tout diagramme comme en
(2), le morphisme $S \to B$ se factorise automatiquement par
$S_{\mathfrak q}$, et a fortiori par $S_g$. La lissité formelle de $S_g$ sur
$R$ nous donne un morphisme $S_g \to B'$ qui s'insère dans un diagramme
semblable avec $S_g$ dans le coin supérieur gauche. Sa composée avec
$S \to S_g$ donne la flèche recherchée. Autrement dit, nous avons montré que
(1) implique (2).

\medskip\noindent
Clairement, (2) implique (3), et (3) implique (4).

\medskip\noindent
Supposons (4). Nous allons montrer que (1) est satisfaite, ce qui achèvera la
démonstration du lemme. Choisissons une présentation
$S = R[x_1, \ldots, x_n]/(f_1, \ldots, f_m)$.
C'est possible, car $S$ est de type fini sur $R$ et donc de présentation finie
(voir le Lemme \ref{lemma-Noetherian-finite-type-is-finite-presentation}).
Posons $I = (f_1, \ldots, f_m)$.
Considérons le complexe cotangent naïf
$$
\text{d} : I/I^2
\longrightarrow
\bigoplus\nolimits_{j = 1}^n S\text{d}x_j
$$
de cette présentation (voir la section \ref{section-netherlander}).
Il suffit de montrer que, lorsqu'on localise ce complexe en $\mathfrak q$, le
morphisme devient une injection scindée ; voir le
Lemme \ref{lemma-smooth-at-point}.
Notons $S' = R[x_1, \ldots, x_n]/I^2$.
Par le Lemme \ref{lemma-differential-mod-power-ideal}, nous avons
$$
S \otimes_{S'} \Omega_{S'/R} =
S \otimes_{R[x_1, \ldots, x_n]} \Omega_{R[x_1, \ldots, x_n]/R} =
\bigoplus\nolimits_{j = 1}^m S\text{d}x_j.
$$
Ainsi, le morphisme
$$
\text{d} :
I/I^2
\longrightarrow
S \otimes_{S'} \Omega_{S'/R}
$$
est le même que celui du complexe cotangent naïf ci-dessus. En particulier, la
vérité de l'assertion que nous cherchons à démontrer ne dépend que des trois
anneaux $R \to S' \to S$.
Soit $\mathfrak q' \subset R[x_1, \ldots, x_n]$ l'idéal premier correspondant
à $\mathfrak q$. Puisque la localisation commute à la formation des modules de
différentielles (Lemme \ref{lemma-differentials-localize}), il suffit de montrer
que le morphisme
\begin{equation}
\label{equation-target-map}
\text{d} :
I_{\mathfrak q'}/I_{\mathfrak q'}^2
\longrightarrow
S_{\mathfrak q} \otimes_{S'_{\mathfrak q'}} \Omega_{S'_{\mathfrak q'}/R}
\end{equation}
provenant de $R \to S'_{\mathfrak q'} \to S_{\mathfrak q}$ soit une injection
scindée.

\medskip\noindent
Soit $N \in \mathbf{N}$ un entier.
Considérons l'anneau
$$
B'_N = S'_{\mathfrak q'} / (\mathfrak q')^N S'_{\mathfrak q'}
= (S'/(\mathfrak q')^N S')_{\mathfrak q'}
$$
et son quotient $B_N = B'_N/IB'_N$. Notons que
$B_N \cong S_{\mathfrak q}/\mathfrak q^NS_{\mathfrak q}$.
Observons que $B'_N$ est un anneau local artinien, puisqu'il est le quotient
d'un anneau local noethérien par une puissance de son idéal maximal.
Considérons une filtration du noyau $I_N$ de $B'_N \to B_N$ par des
sous-$B'_N$-modules
$$
0 \subset J_{N, 1} \subset J_{N, 2} \subset \ldots \subset J_{N, n(N)} = I_N
$$
telle que chaque quotient successif $J_{N, i}/J_{N, i - 1}$ soit de longueur
$1$. (Une telle filtration existe puisque $B'_N$ est artinien.)
Cela donne une suite de petites extensions
$$
B'_N  \to B'_N/J_{N, 1} \to B'_N/J_{N, 2} \to \ldots \to
B'_N/J_{N, n(N)} = B'_N/I_N
= B_N = S_{\mathfrak q}/\mathfrak q^NS_{\mathfrak q}
$$
En appliquant successivement la condition (4) à ces petites extensions, à
partir du morphisme $S \to B_N$, nous voyons qu'il existe un diagramme
commutatif
$$
\xymatrix{
S \ar[r] \ar[rd] & B_N  \\
R \ar[r] \ar[u] & B'_N \ar[u]
}
$$
Clairement, le morphisme d'anneaux $S \to B'_N$ se factorise en
$S \to S_{\mathfrak q} \to B'_N$, où
$S_{\mathfrak q} \to B'_N$ est un homomorphisme local d'anneaux locaux.
De plus, comme l'idéal maximal de $B'_N$, élevé à la puissance $N$, est nul,
nous concluons que $S_{\mathfrak q} \to B'_N$ se factorise par
$S_{\mathfrak q}/(\mathfrak q)^NS_{\mathfrak q} = B_N$. Autrement dit, nous
avons montré que, pour tout $N \in \mathbf{N}$, la surjection de $R$-algèbres
$B'_N \to B_N$ admet un scindage.

\medskip\noindent
Considérons la présentation
$$
I_N \to B_N \otimes_{B'_N} \Omega_{B'_N/R} \to \Omega_{B_N/R} \to 0
$$
provenant de la surjection $B'_N \to B_N$ de noyau $I_N$ (voir le
Lemme \ref{lemma-differential-seq}). D'après ce qui précède, le morphisme de
$R$-algèbres $B'_N \to B_N$ admet un inverse à droite. Ainsi, par le
Lemme \ref{lemma-differential-seq-split}, la suite ci-dessus est exacte
scindée ! Pour tout $N$, le morphisme
$$
I_N \longrightarrow B_N \otimes_{B'_N} \Omega_{B'_N/R}
$$
est donc une injection scindée. La suite de la démonstration s'obtient en
explicitant exactement ce que cela signifie. Notons que
$$
I_N = I_{\mathfrak q'}/
(I_{\mathfrak q'}^2 + (\mathfrak q')^N \cap I_{\mathfrak q'})
$$
Par Artin--Rees (Lemme \ref{lemma-Artin-Rees}), il existe $c \geq 0$ tel que
$$
S_{\mathfrak q}/\mathfrak q^{N - c}S_{\mathfrak q}
\otimes_{S_{\mathfrak q}} I_N =
S_{\mathfrak q}/\mathfrak q^{N - c}S_{\mathfrak q}
\otimes_{S_{\mathfrak q}}
I_{\mathfrak q'}/I_{\mathfrak q'}^2
$$
pour tout $N \geq c$ (ces produits tensoriels ne sont qu'une manière élaborée
de diviser par $\mathfrak q^{N - c}$). Nous pouvons bien sûr supposer
$c \geq 1$.
Par le Lemme \ref{lemma-differential-mod-power-ideal}, nous voyons que
$$
S'_{\mathfrak q'}/(\mathfrak q')^{N - c}S'_{\mathfrak q'}
\otimes_{S'_{\mathfrak q'}} \Omega_{B'_N/R} =
S'_{\mathfrak q'}/(\mathfrak q')^{N - c}S'_{\mathfrak q'}
\otimes_{S'_{\mathfrak q'}} \Omega_{S'_{\mathfrak q'}/R}
$$
nous pouvons encore tensoriser cette égalité par
$B_N = S_{\mathfrak q}/\mathfrak q^N$ pour voir que
$$
S_{\mathfrak q}/\mathfrak q^{N - c}S_{\mathfrak q}
\otimes_{S'_{\mathfrak q'}} \Omega_{B'_N/R} =
S_{\mathfrak q}/\mathfrak q^{N - c}S_{\mathfrak q}
\otimes_{S'_{\mathfrak q'}} \Omega_{S'_{\mathfrak q'}/R}.
$$
Puisqu'une injection scindée reste une injection scindée après tensorisation
par n'importe quel module, nous voyons que
$$
S_{\mathfrak q}/\mathfrak q^{N - c}S_{\mathfrak q}
\otimes_{S_{\mathfrak q}}
(\ref{equation-target-map}) =
S_{\mathfrak q}/\mathfrak q^{N - c}S_{\mathfrak q}
\otimes_{S_{\mathfrak q}/\mathfrak q^N S_{\mathfrak q}}
(I_N \longrightarrow B_N \otimes_{B'_N} \Omega_{B'_N/R})
$$
est une injection scindée pour tout $N \geq c$. Par le
Lemme \ref{lemma-split-injection-after-completion}, nous voyons que
(\ref{equation-target-map}) est une injection scindée. Cela achève la
démonstration.
\end{proof}











\section{Aperçu des résultats sur les morphismes d'anneaux lisses}
\label{section-smooth-overview}

\noindent
Voici une liste de résultats sur les morphismes d'anneaux lisses que nous avons
démontrés dans les sections précédentes. Pour des énoncés et des définitions
plus précis, consulter les références indiquées.
\begin{enumerate}
\item Un morphisme d'anneaux $R \to S$ est lisse s'il est de présentation finie
et si le complexe cotangent naïf de $S/R$ est quasi-isomorphe à un $S$-module
projectif de type fini en degré $0$ ; voir la
Définition \ref{definition-smooth}.
\item Si $S$ est lisse sur $R$, alors $\Omega_{S/R}$ est un $S$-module
projectif de type fini ; voir la discussion qui suit la
Définition \ref{definition-smooth}.
\item La propriété d'être lisse est locale sur $S$ ; voir le
Lemme \ref{lemma-locally-smooth}.
\item La propriété d'être lisse est stable par changement de base ; voir le
Lemme \ref{lemma-base-change-smooth}.
\item La propriété d'être lisse est stable par composition ; voir le
Lemme \ref{lemma-compose-smooth}.
\item Un morphisme d'anneaux lisse est syntomique, donc en particulier plat ;
voir le Lemme \ref{lemma-smooth-syntomic}.
\item Un morphisme d'anneaux plat et de présentation finie dont les anneaux
fibres sont lisses est lisse ; voir le Lemme \ref{lemma-flat-fibre-smooth}.
\item Un morphisme d'anneaux de présentation finie $R \to S$ est lisse si et
seulement s'il est formellement lisse ; voir la
Proposition \ref{proposition-smooth-formally-smooth}.
\item Si $R \to S$ est un morphisme d'anneaux de type fini et si $R$ est
noethérien, alors, pour vérifier que $R \to S$ est lisse, il suffit de vérifier
la propriété de relèvement de la lissité formelle le long des petites
extensions d'anneaux locaux artiniens ; voir le
Lemme \ref{lemma-smooth-test-artinian}.
\item Un morphisme d'anneaux lisse $R \to S$ est le changement de base d'un
morphisme d'anneaux lisse $R_0 \to S_0$, où $R_0$ est de type fini sur
$\mathbf{Z}$ ; voir le
Lemme \ref{lemma-finite-presentation-fs-Noetherian}.
\item La formation de l'ensemble des points où un morphisme d'anneaux est lisse
commute au changement de base plat ; voir le
Lemme \ref{lemma-flat-base-change-locus-smooth}.
\item Si $S$ est de type fini sur un corps algébriquement clos $k$, et si
$\mathfrak m \subset S$ est un idéal maximal, alors les conditions suivantes
sont équivalentes :
\begin{enumerate}
\item $S$ est lisse sur $k$ dans un voisinage de $\mathfrak m$,
\item $S_{\mathfrak m}$ est un anneau local régulier,
\item $\dim(S_{\mathfrak m}) =
\dim_{\kappa(m)} \Omega_{S/k} \otimes_S \kappa(\mathfrak m)$.
\end{enumerate}
voir le Lemme \ref{lemma-characterize-smooth-kbar}.
\item Si $S$ est de type fini sur un corps $k$, et si
$\mathfrak q \subset S$ est un idéal premier, alors les conditions suivantes
sont équivalentes :
\begin{enumerate}
\item $S$ est lisse sur $k$ dans un voisinage de $\mathfrak q$,
\item $\dim_{\mathfrak q}(S/k) =
\dim_{\kappa(\mathfrak q)} \Omega_{S/k} \otimes_S \kappa(\mathfrak q)$.
\end{enumerate}
voir le Lemme \ref{lemma-characterize-smooth-over-field}.
\item Si $S$ est lisse sur un corps, alors tous ses anneaux locaux sont
réguliers ; voir le Lemme \ref{lemma-characterize-smooth-over-field}.
\item Si $S$ est de type fini sur un corps $k$, si
$\mathfrak q \subset S$ est un idéal premier, si l'extension de corps
$\kappa(\mathfrak q)/k$ est séparable et si $S_{\mathfrak q}$ est régulier,
alors $S$ est lisse sur $k$ en $\mathfrak q$ ; voir le
Lemme \ref{lemma-separable-smooth}.
\item Si $S$ est de type fini sur un corps $k$, si $k$ est de caractéristique
$0$, si $\mathfrak q \subset S$ est un idéal premier et si
$\Omega_{S/k, \mathfrak q}$ est libre, alors $S$ est lisse sur $k$ en
$\mathfrak q$ ; voir le Lemme \ref{lemma-characteristic-zero-local-smooth}.
\end{enumerate}
Certains de ces résultats ont été démontrés à l'aide de la notion de morphisme
d'anneaux lisse standard ; voir la Définition \ref{definition-standard-smooth}.
C'est l'analogue, pour les morphismes syntomiques, de la notion de morphisme
d'intersection complète globale relative. C'est aussi la façon la plus simple
de construire des exemples.




















\section{Morphismes d'anneaux étales}
\label{section-etale}

\noindent
Un morphisme d'anneaux étale est un morphisme d'anneaux lisse dont la dimension
relative est nulle. Cela équivaut à la définition un peu plus directe suivante.

\begin{definition}
\label{definition-etale}
Soit $R \to S$ un morphisme d'anneaux. Nous disons que $R \to S$ est
{\it étale} s'il est de présentation finie et si le complexe cotangent naïf
$\NL_{S/R}$ est quasi-isomorphe à zéro : cela signifie que
$H_1(\NL_{S/R}) = 0$ et $\Omega_{S/R} = 0$. Étant donné un idéal premier
$\mathfrak q$ de $S$, nous disons que $R \to S$ est
{\it étale en $\mathfrak q$} s'il existe $g \in S$,
$g \not \in \mathfrak q$, tel que $R \to S_g$ soit étale.
\end{definition}

\noindent
En particulier, nous voyons que $\Omega_{S/R} = 0$ si $S$ est étale sur $R$.
Si $R \to S$ est lisse, alors $R \to S$ est étale si et seulement si
$\Omega_{S/R} = 0$.
Nos résultats sur les morphismes d'anneaux lisses fournissent automatiquement
toute une série de résultats pour les morphismes étales. Nous les résumons dans
le Lemme \ref{lemma-etale} ci-dessous. Mais auparavant, nous démontrons que
{\it tout} morphisme d'anneaux étale est lisse standard.

\begin{lemma}
\label{lemma-etale-standard-smooth}
Tout morphisme d'anneaux étale est lisse standard. Plus précisément, si
$R \to S$ est étale, alors il existe une présentation
$S = R[x_1, \ldots, x_n]/(f_1, \ldots, f_n)$ telle que l'image de
$\det(\partial f_j/\partial x_i)$ soit inversible dans $S$.
\end{lemma}

\begin{proof}
Soit $R \to S$ étale. Choisissons une présentation
$S = R[x_1, \ldots, x_n]/I$. Comme $R \to S$ est étale, nous savons que
$$
\text{d} :
I/I^2
\longrightarrow
\bigoplus\nolimits_{i = 1, \ldots, n} S\text{d}x_i
$$
est un isomorphisme ; en particulier, $I/I^2$ est un $S$-module libre.
Ainsi, par le Lemme \ref{lemma-huber}, nous pouvons supposer (après avoir
éventuellement changé la présentation) que $I = (f_1, \ldots, f_c)$ et que les
classes $f_i \bmod I^2$ forment une base de $I/I^2$. Le fait que le morphisme
affiché ci-dessus soit un isomorphisme implique immédiatement que $c = n$ et
que $\det(\partial f_j/\partial x_i)$ est inversible dans $S$.
\end{proof}

\begin{lemma}
\label{lemma-etale}
Résultats sur les morphismes d'anneaux étales.
\begin{enumerate}
\item Le morphisme d'anneaux $R \to R_f$ est étale pour tout anneau $R$ et tout
$f \in R$.
\item Les composées de morphismes d'anneaux étales sont étales.
\item Un changement de base d'un morphisme d'anneaux étale est étale.
\item La propriété d'être étale est locale : étant donnés un morphisme
d'anneaux $R \to S$ et des éléments $g_1, \ldots, g_m \in S$ qui engendrent
l'idéal unité, tels que $R \to S_{g_j}$ soit étale pour
$j = 1, \ldots, m$, alors $R \to S$ est étale.
\item Étant donné $R \to S$ de présentation finie et un morphisme d'anneaux
plat $R \to R'$, posons $S' = R' \otimes_R S$. L'ensemble des idéaux premiers
où $R' \to S'$ est étale est l'image réciproque, par
$\Spec(S') \to \Spec(S)$, de l'ensemble des idéaux premiers où $R \to S$ est
étale.
\item Un morphisme d'anneaux étale est syntomique, donc en particulier plat.
\item Si $S$ est de type fini sur un corps $k$, alors $S$ est étale sur $k$ si
et seulement si $\Omega_{S/k} = 0$.
\item Tout morphisme d'anneaux étale $R \to S$ est le changement de base d'un
morphisme d'anneaux étale $R_0 \to S_0$, où $R_0$ est de type fini sur
$\mathbf{Z}$.
\item Soit $A = \colim A_i$ une colimite filtrante d'anneaux.
Soit $A \to B$ un morphisme d'anneaux étale. Alors il existe un morphisme
d'anneaux étale $A_i \to B_i$ pour un certain $i$, tel que
$B \cong A \otimes_{A_i} B_i$.
\item Soit $A$ un anneau. Soit $S$ une partie multiplicative de $A$.
Soit $S^{-1}A \to B'$ étale. Alors il existe un morphisme d'anneaux étale
$A \to B$ tel que $B' \cong S^{-1}B$.
\item Soit $A$ un anneau. Soit $B = B' \times B''$ un produit de
$A$-algèbres. Alors $B$ est étale sur $A$ si et seulement si $B'$ et $B''$
sont toutes deux étales sur $A$.
\end{enumerate}
\end{lemma}

\begin{proof}
Dans chaque cas, nous utilisons le résultat correspondant pour les morphismes
d'anneaux lisses, auquel nous ajoutons un petit argument pour montrer que
$\Omega_{S/R}$ est nul.

\medskip\noindent
Démonstration de (1). Le morphisme d'anneaux $R \to R_f$ est lisse et
$\Omega_{R_f/R} = 0$.

\medskip\noindent
Démonstration de (2). La composée $A \to C$ des morphismes lisses
$A \to B$ et $B \to C$ est lisse ; voir le
Lemme \ref{lemma-compose-smooth}. Par le
Lemme \ref{lemma-exact-sequence-differentials}, nous voyons que
$\Omega_{C/A}$ est nul, puisque $\Omega_{C/B}$ et $\Omega_{B/A}$ sont tous deux
nuls.

\medskip\noindent
Démonstration de (3). Soit $R \to S$ étale, et soit $R \to R'$ quelconque.
Alors $R' \to S' = R' \otimes_R S$ est lisse ; voir le
Lemme \ref{lemma-base-change-smooth}. Puisque
$\Omega_{S'/R'} = S' \otimes_S \Omega_{S/R}$ par le
Lemme \ref{lemma-differentials-base-change}, nous concluons que
$\Omega_{S'/R'} = 0$. Ainsi $R' \to S'$ est étale.

\medskip\noindent
Démonstration de (4). Supposons les hypothèses de (4). Par le
Lemme \ref{lemma-locally-smooth}, nous voyons que $R \to S$ est lisse.
Nous savons également que
$\Omega_{S_{g_i}/R} = (\Omega_{S/R})_{g_i} = 0$ pour tout $i$.
Alors $\Omega_{S/R} = 0$ ; voir le Lemme \ref{lemma-cover}.

\medskip\noindent
Démonstration de (5). Le résultat pour les morphismes lisses est le
Lemme \ref{lemma-flat-base-change-locus-smooth}.
Dans la démonstration de ce lemme, nous avons utilisé que
$\NL_{S/R} \otimes_S S'$ est homotopiquement équivalent à $\NL_{S'/R'}$.
Nous sommes donc ramenés à montrer que, si $M$ est un $S$-module de
présentation finie, l'ensemble des idéaux premiers $\mathfrak q'$ de $S'$ tels
que $(M \otimes_S S')_{\mathfrak q'} = 0$ est l'image réciproque de l'ensemble
des idéaux premiers $\mathfrak q$ de $S$ tels que $M_{\mathfrak q} = 0$.
Cela résulte du Lemme \ref{lemma-support-base-change}.

\medskip\noindent
Démonstration de (6). Cela résulte directement du résultat correspondant pour
les morphismes d'anneaux lisses (Lemme \ref{lemma-smooth-syntomic}).

\medskip\noindent
Démonstration de (7). Cela résulte du
Lemme \ref{lemma-characterize-smooth-over-field} et des définitions.

\medskip\noindent
Démonstration de (8). Le
Lemme \ref{lemma-finite-presentation-fs-Noetherian} donne le résultat pour les
morphismes d'anneaux lisses. Le morphisme d'anneaux lisse obtenu
$R_0 \to S_0$ satisfait aux hypothèses du
Lemme \ref{lemma-relative-dimension-CM} ; nous pouvons donc remplacer $S_0$ par
le facteur de dimension relative $0$ sur $R_0$.

\medskip\noindent
Démonstration de (9). Cela résulte de (8), car $R_0 \to A$ se factorisera par
$A_i$ pour un certain $i$, d'après le
Lemme \ref{lemma-characterize-finite-presentation}.

\medskip\noindent
Démonstration de (10). Cela résulte de (9), (1) et (2), puisque $S^{-1}A$ est
une colimite filtrante de localisations principales de $A$.

\medskip\noindent
Démonstration de (11). Utilisons le Lemme \ref{lemma-product-smooth} pour
obtenir le résultat pour la lissité, puis le fait que $\Omega_{B/A}$ est nul si
et seulement si $\Omega_{B'/A}$ et $\Omega_{B''/A}$ sont tous deux nuls.
\end{proof}


















\noindent
Examinons maintenant plus en détail ce que signifie être étale sur un corps.

\begin{lemma}
\label{lemma-etale-over-field}
Soit $k$ un corps. Un morphisme d'anneaux $k \to S$ est étale si et seulement
si $S$ est isomorphe, comme $k$-algèbre, à un produit fini d'extensions
séparables finies de $k$.
\end{lemma}

\begin{proof}
Nous utiliserons sans autre mention le fait suivant : si
$S = S_1 \times \ldots \times S_n$ est un produit fini de $k$-algèbres, alors
$S$ est étale sur $k$ si et seulement si chaque $S_i$ est étale sur $k$.
Voir le Lemme \ref{lemma-etale}, partie (11).

\medskip\noindent
Si $k'/k$ est une extension de corps séparable finie, nous pouvons écrire
$k' = k(\alpha) \cong k[x]/(f)$. Ici, $f$ est le polynôme minimal de l'élément
$\alpha$. Comme $k'$ est séparable sur $k$, nous avons
$\gcd(f, f') = 1$. Cela implique que
$\text{d} : k'\cdot f \to k' \cdot \text{d}x$ est un isomorphisme.
Ainsi $k \to k'$ est étale. Par conséquent, si $S$ est un produit fini
d'extensions séparables finies de $k$, alors $S$ est étale sur $k$.

\medskip\noindent
Réciproquement, supposons que $k \to S$ soit étale. Alors $S$ est lisse sur
$k$ et $\Omega_{S/k} = 0$. Par le
Lemme \ref{lemma-characterize-smooth-over-field}, nous voyons que
$\dim_\mathfrak m \Spec(S) = 0$ pour tout idéal maximal $\mathfrak m$ de $S$.
Ainsi $\dim(S) = 0$. Par la Proposition \ref{proposition-dimension-zero-ring},
nous obtenons que $S$ est un produit fini d'anneaux locaux artiniens.
Par le Lemme \ref{lemma-characterize-smooth-over-field}, déjà utilisé, ces
anneaux locaux sont des corps. Nous pouvons donc supposer que $S = k'$ est un
corps. Par le théorème des zéros de Hilbert
(Théorème \ref{theorem-nullstellensatz}), nous voyons que l'extension $k'/k$
est finie. La lissité de $k \to k'$ implique, par le
Lemme \ref{lemma-smooth-at-generic-point}, que $k'/k$ est une extension
séparable, ce qui achève la démonstration.
\end{proof}

\begin{lemma}
\label{lemma-etale-at-prime}
Soit $R \to S$ un morphisme d'anneaux.
Soit $\mathfrak q \subset S$ un idéal premier au-dessus de $\mathfrak p$ dans
$R$. Si $S/R$ est étale en $\mathfrak q$, alors
\begin{enumerate}
\item nous avons $\mathfrak p S_{\mathfrak q} = \mathfrak qS_{\mathfrak q}$,
qui est l'idéal maximal de l'anneau local $S_{\mathfrak q}$, et
\item l'extension de corps $\kappa(\mathfrak q)/\kappa(\mathfrak p)$ est
séparable finie.
\end{enumerate}
\end{lemma}

\begin{proof}
Nous pouvons d'abord remplacer $S$ par $S_g$ pour un certain
$g \in S$, $g \not \in \mathfrak q$, et supposer que $R \to S$ est étale.
Le lemme résulte alors du Lemme \ref{lemma-etale-over-field}, en explicitant le
fait que $S \otimes_R \kappa(\mathfrak p)$ est étale sur
$\kappa(\mathfrak p)$.
\end{proof}

\begin{lemma}
\label{lemma-etale-quasi-finite}
Un morphisme d'anneaux étale est quasi-fini.
\end{lemma}

\begin{proof}
Soit $R \to S$ un morphisme d'anneaux étale. Par définition, $R \to S$ est de
type fini. Pour tout idéal premier $\mathfrak p \subset R$, l'anneau fibre
$S \otimes_R \kappa(\mathfrak p)$ est étale sur $\kappa(\mathfrak p)$ ; c'est
donc un produit fini de corps séparables finis sur $\kappa(\mathfrak p)$, et il
est en particulier fini sur $\kappa(\mathfrak p)$. Ainsi $R \to S$ est
quasi-fini par le Lemme \ref{lemma-quasi-finite}.
\end{proof}

\begin{lemma}
\label{lemma-characterize-etale}
Soit $R \to S$ un morphisme d'anneaux. Soit $\mathfrak q$ un idéal premier de
$S$ au-dessus d'un idéal premier $\mathfrak p$ de $R$. Si
\begin{enumerate}
\item $R \to S$ est de présentation finie,
\item $R_{\mathfrak p} \to S_{\mathfrak q}$ est plat,
\item $\mathfrak p S_{\mathfrak q}$ est l'idéal maximal de l'anneau local
$S_{\mathfrak q}$, et
\item l'extension de corps $\kappa(\mathfrak q)/\kappa(\mathfrak p)$ est
séparable finie,
\end{enumerate}
alors $R \to S$ est étale en $\mathfrak q$.
\end{lemma}

\begin{proof}
Appliquons le Lemme \ref{lemma-isolated-point-fibre} pour trouver
$g \in S$, $g \not \in \mathfrak q$, tel que $\mathfrak q$ soit l'unique idéal
premier de $S_g$ au-dessus de $\mathfrak p$.
Nous remplaçons $S$ par $S_g$. Alors $S \otimes_R \kappa(\mathfrak p)$ possède
un unique idéal premier ; c'est donc un anneau local, et il est donc égal à
$S_{\mathfrak q}/\mathfrak pS_{\mathfrak q}
\cong \kappa(\mathfrak q)$.
Par le Lemme \ref{lemma-flat-fibre-smooth}, il existe
$g \in S$, $g \not \in \mathfrak q$, tel que $R \to S_g$ soit lisse.
Remplaçons encore $S$ par $S_g$ ; nous pouvons supposer que $R \to S$ est
lisse. Par le Lemme \ref{lemma-smooth-syntomic}, nous pouvons même supposer que
$R \to S$ est lisse standard, disons
$S = R[x_1, \ldots, x_n]/(f_1, \ldots, f_c)$. Puisque
$S \otimes_R \kappa(\mathfrak p) = \kappa(\mathfrak q)$ est de dimension $0$,
nous concluons que $n = c$, c'est-à-dire que $R \to S$ est étale.
\end{proof}

\noindent
Voici un phénomène tout à fait nouveau.

\begin{lemma}
\label{lemma-map-between-etale}
Soient $R \to S$ et $R \to S'$ étales.
Alors tout morphisme de $R$-algèbres $S' \to S$ est étale.
\end{lemma}

\begin{proof}
Notons tout d'abord que $S' \to S$ est de présentation finie par le
Lemme \ref{lemma-compose-finite-type}.
Soit $\mathfrak q \subset S$ un idéal premier au-dessus des idéaux premiers
$\mathfrak q' \subset S'$ et $\mathfrak p \subset R$.
Par le Lemme \ref{lemma-etale-at-prime}, le morphisme d'anneaux
$S'_{\mathfrak q'}/\mathfrak p S'_{\mathfrak q'} \to
S_{\mathfrak q}/\mathfrak p S_{\mathfrak q}$ est un morphisme entre extensions
séparables finies de $\kappa(\mathfrak p)$. Il est en particulier plat.
Le Lemme \ref{lemma-criterion-flatness-fibre} montre donc que
$S'_{\mathfrak q'} \to S_{\mathfrak q}$ est plat. Ainsi $S' \to S$ est plat.
De plus, ce qui précède montre aussi que $\mathfrak q'S_{\mathfrak q}$ est
l'idéal maximal de $S_{\mathfrak q}$ et que l'extension de corps résiduels de
$S'_{\mathfrak q'} \to S_{\mathfrak q}$ est séparable finie.
Le Lemme \ref{lemma-characterize-etale} permet donc de conclure que
$S' \to S$ est étale en $\mathfrak q$. Comme la propriété d'être étale est
locale (voir le Lemme \ref{lemma-etale}), nous avons terminé.
\end{proof}

\begin{lemma}
\label{lemma-surjective-flat-finitely-presented}
Soit $\varphi : R \to S$ un morphisme d'anneaux. Si $R \to S$ est surjectif,
plat et de présentation finie, alors il existe un idempotent $e \in R$ tel que
$S = R_e$.
\end{lemma}

\begin{proof}[Première démonstration]
Soit $I$ le noyau de $\varphi$.
Le Lemme \ref{lemma-finite-presentation-independent} montre que $I$ est de type
fini, puisque $\varphi$ est de présentation finie.
De plus, comme $S$ est plat sur $R$, tensoriser la suite exacte
$0 \to I \to R \to S \to 0$ sur $R$ avec $S$ donne $I/I^2 = 0$.
Nous concluons maintenant par le
Lemme \ref{lemma-ideal-is-squared-union-connected}.
\end{proof}

\begin{proof}[Seconde démonstration]
Puisque $\Spec(S) \to \Spec(R)$ est un homéomorphisme sur un sous-ensemble
fermé (voir le Lemme \ref{lemma-spec-closed}) et est ouvert (voir la
Proposition \ref{proposition-fppf-open}), nous voyons que son image est $D(e)$
pour un certain idempotent $e \in R$ (voir le
Lemme \ref{lemma-disjoint-decomposition}). Ainsi $R_e \to S$ induit une
bijection sur les spectres. Ce morphisme induit maintenant un isomorphisme sur
tous les anneaux locaux, par exemple par les
Lemmes \ref{lemma-finite-flat-local} et \ref{lemma-NAK}.
Il s'ensuit que $R_e \to S$ est aussi injectif ; voir, par exemple, le
Lemme \ref{lemma-characterize-zero-local}.
\end{proof}

\begin{lemma}
\label{lemma-lift-etale}
\begin{slogan}
Les morphismes d'anneaux étales se relèvent le long des surjections d'anneaux
\end{slogan}
Soit $R$ un anneau et soit $I \subset R$ un idéal.
Soit $R/I \to \overline{S}$ un morphisme d'anneaux étale.
Alors il existe un morphisme d'anneaux étale $R \to S$ tel que
$\overline{S} \cong S/IS$ comme $R/I$-algèbres.
\end{lemma}

\begin{proof}
Par le Lemme \ref{lemma-etale-standard-smooth}, nous pouvons écrire
$\overline{S} =
(R/I)[x_1, \ldots, x_n]/(\overline{f}_1, \ldots, \overline{f}_n)$
comme dans la Définition \ref{definition-standard-smooth}, avec
$\overline{\Delta} =
\det(\frac{\partial \overline{f}_i}{\partial x_j})_{i, j = 1, \ldots, n}$
inversible dans $\overline{S}$. Il suffit de choisir des relèvements $f_i$ et
de poser
$S = R[x_1, \ldots, x_n, x_{n+1}]/(f_1, \ldots, f_n, x_{n + 1}\Delta - 1)$,
où $\Delta = \det(\frac{\partial f_i}{\partial x_j})_{i, j = 1, \ldots, n}$,
comme dans l'Exemple \ref{example-make-standard-smooth}.
Cela démontre le lemme.
\end{proof}















\begin{lemma}
\label{lemma-lift-etale-infinitesimal}
Considérons un diagramme commutatif
$$
\xymatrix{
0 \ar[r] &
J \ar[r] &
B' \ar[r] &
B \ar[r] & 0 \\
0 \ar[r] &
I \ar[r] \ar[u] &
A' \ar[r] \ar[u] &
A \ar[r] \ar[u] & 0
}
$$
à lignes exactes, où $B' \to B$ et $A' \to A$ sont des morphismes d'anneaux
surjectifs dont les noyaux sont des idéaux de carré nul. Si $A \to B$ est étale
et si $J = I \otimes_A B$, alors $A' \to B'$ est étale.
\end{lemma}

\begin{proof}
Par le Lemme \ref{lemma-lift-etale}, il existe un morphisme d'anneaux étale
$A' \to C$ tel que $C/IC = B$. Alors $A' \to C$ est formellement lisse (par la
Proposition \ref{proposition-smooth-formally-smooth}) ; nous obtenons donc un
morphisme de $A'$-algèbres $\varphi : C \to B'$.
Comme $A' \to C$ est plat, nous avons
$I \otimes_A B = I \otimes_A C/IC = IC$.
Ainsi, l'hypothèse $J = I \otimes_A B$ implique que $\varphi$ induit un
isomorphisme $IC \to J$ et un isomorphisme $C/IC \to B'/IB'$ ; par conséquent,
$\varphi$ est un isomorphisme.
\end{proof}

\begin{example}
\label{example-factor-polynomials-etale}
Soient $n , m \geq 1$ des entiers. Considérons le morphisme d'anneaux
\begin{eqnarray*}
R = \mathbf{Z}[a_1, \ldots, a_{n + m}]
& \longrightarrow &
S = \mathbf{Z}[b_1, \ldots, b_n, c_1, \ldots, c_m] \\
a_1 & \longmapsto & b_1 + c_1 \\
a_2 & \longmapsto & b_2 + b_1 c_1 + c_2 \\
\ldots & \ldots & \ldots \\
a_{n + m} & \longmapsto & b_n c_m
\end{eqnarray*}
de l'Exemple \ref{example-factor-polynomials}.
Écrivons symboliquement
$$
S = R[b_1, \ldots, c_m]/(\{a_k(b_i, c_j) - a_k\}_{k = 1, \ldots, n + m})
$$
où, par exemple, $a_1(b_i, c_j) = b_1 + c_1$.
La matrice des dérivées partielles est
$$
\left(
\begin{matrix}
1 & c_1 & \ldots & c_m & 0 & \ldots & \ldots & 0 \\
0 & 1 & c_1 & \ldots & c_m & 0 & \ldots & 0 \\
\ldots & \ldots & \ldots & \ldots & \ldots & \ldots & \ldots & \ldots \\
0 & \ldots & 0 & 1 & c_1 & c_2 & \ldots & c_m \\
1 & b_1 & \ldots & b_{n - 1} & b_n & 0 & \ldots & 0 \\
0 & 1 & b_1 & \ldots & b_{n - 1} & b_n & \ldots & 0 \\
\ldots & \ldots & \ldots & \ldots & \ldots & \ldots & \ldots & \ldots \\
0 & \ldots & \ldots & 0 & 1 & b_1 & \ldots & b_n
\end{matrix}
\right)
$$
Le déterminant $\Delta$ de cette matrice est mieux connu sous le nom de
{\it résultant} des polynômes
$g = x^n + b_1 x^{n - 1} + \ldots + b_n$ et
$h = x^m + c_1 x^{m - 1} + \ldots + c_m$, et la matrice ci-dessus est appelée
{\it matrice de Sylvester} associée à $g, h$.
Sous forme de formule, $\Delta = \text{Res}_x(g, h)$. La matrice de Sylvester
est la transposée de la matrice du morphisme linéaire
\begin{eqnarray*}
S[x]_{< m} \oplus S[x]_{< n} & \longrightarrow & S[x]_{< n + m} \\
a \oplus b & \longmapsto & ag + bh
\end{eqnarray*}
Soit $\mathfrak q \subset S$ un idéal premier quelconque. D'après ce qui
précède, les conditions suivantes sont équivalentes :
\begin{enumerate}
\item $R \to S$ est étale en $\mathfrak q$,
\item $\Delta = \text{Res}_x(g, h) \not \in \mathfrak q$,
\item les images $\overline{g}, \overline{h} \in \kappa(\mathfrak q)[x]$ des
polynômes $g, h$ sont premiers entre eux dans $\kappa(\mathfrak q)[x]$.
\end{enumerate}
L'équivalence de (2) et (3) vaut parce que l'image de la matrice de Sylvester
dans $\text{Mat}(n + m, \kappa(\mathfrak q))$ possède un noyau si et seulement
si les polynômes $\overline{g}, \overline{h}$ ont un facteur commun.
Nous concluons que le morphisme d'anneaux
$$
R \longrightarrow S[\frac{1}{\Delta}] = S[\frac{1}{\text{Res}_x(g, h)}]
$$
est étale.
\end{example}


\begin{lemma}
\label{lemma-factor-mod-lift-etale}
Soit $R$ un anneau. Soit $f \in R[x]$ un polynôme unitaire. Soit
$\mathfrak p$ un idéal premier de $R$. Soit
$f \bmod \mathfrak p = \overline{g} \overline{h}$ une factorisation de
l'image de $f$ dans $\kappa(\mathfrak p)[x]$.
Si $\gcd(\overline{g}, \overline{h}) = 1$, alors il existe
\begin{enumerate}
\item un morphisme d'anneaux étale $R \to R'$,
\item un idéal premier $\mathfrak p' \subset R'$ au-dessus de $\mathfrak p$, et
\item une factorisation $f = g h$ dans $R'[x]$
\end{enumerate}
tels que
\begin{enumerate}
\item $\kappa(\mathfrak p) = \kappa(\mathfrak p')$,
\item $\overline{g} = g \bmod \mathfrak p'$,
$\overline{h} = h \bmod \mathfrak p'$, et
\item les polynômes $g, h$ engendrent l'idéal unité dans $R'[x]$.
\end{enumerate}
\end{lemma}

\begin{proof}
Supposons
$\overline{g} = \overline{b}_0 x^n + \overline{b}_1 x^{n - 1} + \ldots
+ \overline{b}_n$, et
$\overline{h} = \overline{c}_0 x^m + \overline{c}_1 x^{m - 1} + \ldots
+ \overline{c}_m$, avec
$\overline{b}_0, \overline{c}_0 \in \kappa(\mathfrak p)$ non nuls.
Après avoir localisé $R$ en un certain élément de $R$ qui n'appartient pas à
$\mathfrak p$, nous pouvons supposer que $\overline{b}_0$ est l'image d'un
élément inversible $b_0 \in R$. En remplaçant
$\overline{g}$ par $\overline{g}/b_0$ et $\overline{h}$ par
$b_0\overline{h}$, nous nous ramenons au cas où $\overline{g}$,
$\overline{h}$ sont unitaires (vérification omise).
Disons que
$\overline{g} = x^n + \overline{b}_1 x^{n - 1} + \ldots + \overline{b}_n$,
et $\overline{h} = x^m + \overline{c}_1 x^{m - 1} + \ldots + \overline{c}_m$.
Écrivons $f = x^{n + m} + a_1 x^{n + m - 1} + \ldots + a_{n + m}$.
Considérons le produit fibré
$$
R' = R \otimes_{\mathbf{Z}[a_1, \ldots, a_{n + m}]}
\mathbf{Z}[b_1, \ldots, b_n, c_1, \ldots, c_m]
$$
où le morphisme $\mathbf{Z}[a_k] \to \mathbf{Z}[b_i, c_j]$ est celui des
Exemples \ref{example-factor-polynomials} et
\ref{example-factor-polynomials-etale}. Par construction, il existe un
morphisme de $R$-algèbres
$$
R' = R \otimes_{\mathbf{Z}[a_1, \ldots, a_{n + m}]}
\mathbf{Z}[b_1, \ldots, b_n, c_1, \ldots, c_m]
\longrightarrow
\kappa(\mathfrak p)
$$
qui envoie $b_i$ sur $\overline{b}_i$ et $c_j$ sur $\overline{c}_j$.
Notons $\mathfrak p' \subset R'$ le noyau de ce morphisme.
Comme, par hypothèse, les polynômes $\overline{g}, \overline{h}$ sont premiers
entre eux, nous voyons que l'élément
$\Delta = \text{Res}_x(g, h) \in \mathbf{Z}[b_i, c_j]$
(voir l'Exemple \ref{example-factor-polynomials-etale}) ne s'envoie pas sur zéro
dans $\kappa(\mathfrak p)$ par le morphisme affiché.
Nous concluons que $R \to R'$ est étale en $\mathfrak p'$.
En fait, une solution au problème posé dans le lemme est donnée par le
morphisme d'anneaux $R \to R'[1/\Delta]$ et l'idéal premier
$\mathfrak p' R'[1/\Delta]$. Comme $\text{Res}_x(f, g)$ est inversible dans cet
anneau, la matrice de Sylvester est inversible sur $R'[1/\Delta]$ et donc
$1 = a g +  b h$ pour certains $a, b \in R'[1/\Delta][x]$ ; voir
l'Exemple \ref{example-factor-polynomials-etale}.
\end{proof}




\section{Structure locale des morphismes d'anneaux étales}
\label{section-etale-local-structure}

\noindent
Le Lemme \ref{lemma-etale-standard-smooth} nous indique qu'il n'est pas
vraiment pertinent de définir un morphisme étale standard comme un morphisme
lisse standard de dimension relative $0$.
Comme modèle d'un morphisme étale, nous prenons l'exemple donné par une
extension de corps finie séparable $k'/k$.
En effet, nous pouvons toujours trouver un élément $\alpha \in k'$ tel que
$k' = k(\alpha)$ et tel que le polynôme minimal $f(x) \in k[x]$ de $\alpha$
ait une dérivée $f'$ première avec $f$.

\begin{definition}
\label{definition-standard-etale}
Soit $R$ un anneau. Soient $g , f  \in R[x]$.
Supposons que $f$ soit unitaire et que la dérivée $f'$ soit inversible dans
la localisation $R[x]_g/(f)$.
Dans ce cas, le morphisme d'anneaux $R \to R[x]_g/(f)$ est dit
{\it étale standard}.
\end{definition}

\noindent
Dans la Proposition \ref{proposition-etale-locally-standard}, nous montrons
que tout morphisme d'anneaux étale est localement étale standard.

\begin{lemma}
\label{lemma-standard-etale}
Soit $R \to R[x]_g/(f)$ un morphisme étale standard.
\begin{enumerate}
\item Le morphisme d'anneaux $R \to R[x]_g/(f)$ est étale.
\item Pour tout morphisme d'anneaux $R \to R'$, le changement de base
$R' \to R'[x]_g/(f)$ du morphisme étale standard
$R \to R[x]_g/(f)$ est étale standard.
\item Toute localisation principale de $R[x]_g/(f)$ est étale standard sur $R$.
\item Une composée de morphismes étales standard n'est {\bf pas} étale standard
en général.
\end{enumerate}
\end{lemma}

\begin{proof}
Omis. Voici un exemple pour (4).
Le morphisme d'anneaux $\mathbf{F}_2 \to \mathbf{F}_{2^2}$ est étale standard.
Le morphisme d'anneaux
$\mathbf{F}_{2^2} \to \mathbf{F}_{2^2} \times \mathbf{F}_{2^2}
\times \mathbf{F}_{2^2} \times \mathbf{F}_{2^2}$ est étale standard.
Mais le morphisme d'anneaux
$\mathbf{F}_2 \to \mathbf{F}_{2^2} \times \mathbf{F}_{2^2}
\times \mathbf{F}_{2^2} \times \mathbf{F}_{2^2}$ n'est pas étale standard.
\end{proof}

\noindent
Les morphismes étales standard constituent un moyen commode de produire des
morphismes étales. En voici un exemple.

\begin{lemma}
\label{lemma-make-etale-map-prescribed-residue-field}
Soit $R$ un anneau.
Soit $\mathfrak p$ un idéal premier de $R$.
Soit $L/\kappa(\mathfrak p)$ une extension de corps finie séparable.
Il existe un morphisme d'anneaux étale $R \to R'$, ainsi qu'un idéal premier
$\mathfrak p'$ au-dessus de $\mathfrak p$, tels que l'extension de corps
$\kappa(\mathfrak p')/\kappa(\mathfrak p)$ soit isomorphe à
$\kappa(\mathfrak p) \subset L$.
\end{lemma}

\begin{proof}
Par le théorème de l'élément primitif, nous pouvons écrire
$L = \kappa(\mathfrak p)[\alpha]$. Soit
$\overline{f} \in \kappa(\mathfrak p)[x]$ le polynôme minimal de $\alpha$
(en particulier, il est unitaire).
Après avoir remplacé $\alpha$ par $c\alpha$ pour un certain
$c \in R$, $c\not \in \mathfrak p$, nous pouvons supposer que tous les
coefficients de $\overline{f}$ appartiennent à l'image de
$R \to \kappa(\mathfrak p)$ (vérification omise).
Nous pouvons donc trouver un polynôme unitaire $f \in R[x]$ qui s'envoie sur
$\overline{f}$ dans $\kappa(\mathfrak p)[x]$.
Comme $\kappa(\mathfrak p) \subset L$ est séparable, nous voyons que
$\gcd(\overline{f}, \overline{f}') = 1$.
Il existe donc un élément $\gamma \in L$ tel que
$\overline{f}'(\alpha) \gamma = 1$. Nous obtenons ainsi un morphisme de
$R$-algèbres
\begin{eqnarray*}
R[x, 1/f']/(f) & \longrightarrow & L \\
x & \longmapsto & \alpha \\
1/f' & \longmapsto & \gamma
\end{eqnarray*}
Le membre de gauche est une algèbre étale standard $R'$ sur $R$, et le noyau
du morphisme d'anneaux fournit l'idéal premier recherché.
\end{proof}

\begin{proposition}
\label{proposition-etale-locally-standard}
Soit $R \to S$ un morphisme d'anneaux. Soit $\mathfrak q \subset S$ un idéal
premier. Si $R \to S$ est étale en $\mathfrak q$, alors il existe
$g \in S$, $g \not \in \mathfrak q$, tel que $R \to S_g$ soit étale standard.
\end{proposition}

\begin{proof}
La démonstration suivante est un peu indirecte, et il existe peut-être des
moyens de la raccourcir.

\medskip\noindent
Étape 1. Par la Définition \ref{definition-etale}, il existe
$g \in S$, $g \not \in \mathfrak q$, tel que $R \to S_g$ soit étale.
Nous pouvons donc supposer que $S$ est étale sur $R$.

\medskip\noindent
Étape 2. Par le Lemme \ref{lemma-etale}, il existe un morphisme d'anneaux étale
$R_0 \to S_0$, avec $R_0$ de type fini sur $\mathbf{Z}$, et un morphisme
d'anneaux $R_0 \to R$ tels que $R = R \otimes_{R_0} S_0$.
Notons $\mathfrak q_0$ l'idéal premier de $S_0$ correspondant à
$\mathfrak q$. Si nous démontrons le résultat pour
$(R_0 \to S_0, \mathfrak q_0)$, alors il en découle pour
$(R \to S, \mathfrak q)$ par changement de base.
Nous pouvons donc supposer que $R$ est noethérien.

\medskip\noindent
Étape 3.
Remarquons que $R \to S$ est quasi-fini par le
Lemme \ref{lemma-etale-quasi-finite}.
Par le Lemme \ref{lemma-quasi-finite-open-integral-closure}, il existe un
morphisme d'anneaux fini $R \to S'$, un morphisme de $R$-algèbres
$S' \to S$ et un élément $g' \in S'$ dont l'image vérifie
$g' \not \in \mathfrak q$ et tel que $S' \to S$ induise un isomorphisme
$S'_{g'} \cong S_{g'}$.
(Remarquons que, bien entendu, $S'$ n'est en général pas étale sur $R$.)
Nous pouvons donc supposer que (a) $R$ est noethérien, (b) $R \to S$ est fini
et (c) $R \to S$ est étale en $\mathfrak q$
(mais plus nécessairement étale en tout idéal premier).

\medskip\noindent
Étape 4. Soit $\mathfrak p \subset R$ l'idéal premier correspondant à
$\mathfrak q$. Considérons l'anneau fibre
$S \otimes_R \kappa(\mathfrak p)$. C'est une algèbre finie sur
$\kappa(\mathfrak p)$. Il est donc artinien
(voir le Lemme \ref{lemma-finite-dimensional-algebra}), et par conséquent
un produit fini d'anneaux locaux
$$
S \otimes_R \kappa(\mathfrak p) = \prod\nolimits_{i = 1}^n A_i
$$
voir la Proposition \ref{proposition-dimension-zero-ring}. L'un des facteurs,
disons $A_1$, est l'anneau local
$S_{\mathfrak q}/\mathfrak pS_{\mathfrak q}$, qui est isomorphe à
$\kappa(\mathfrak q)$ ; voir le Lemme \ref{lemma-etale-at-prime}.
Les autres facteurs correspondent aux autres idéaux premiers, disons
$\mathfrak q_2, \ldots, \mathfrak q_n$, de $S$ au-dessus de $\mathfrak p$.

\medskip\noindent
Étape 5. Nous pouvons choisir un élément non nul
$\alpha \in \kappa(\mathfrak q)$ qui engendre l'extension de corps finie
séparable $\kappa(\mathfrak q)/\kappa(\mathfrak p)$
(ainsi, même si l'extension de corps est triviale, nous n'autorisons pas
$\alpha = 0$).
Remarquons que, pour tout $\lambda \in \kappa(\mathfrak p)^*$, l'élément
$\lambda \alpha$ engendre lui aussi $\kappa(\mathfrak q)$ sur
$\kappa(\mathfrak p)$. Considérons l'élément
$$
\overline{t} =
(\alpha, 0, \ldots, 0) \in
\prod\nolimits_{i = 1}^n A_i =
S \otimes_R \kappa(\mathfrak p).
$$
Après avoir éventuellement remplacé $\alpha$ par $\lambda \alpha$ comme
ci-dessus, nous pouvons supposer que $\overline{t}$ est l'image de
$t \in S$. Soit $I \subset R[x]$ le noyau du morphisme de $R$-algèbres
$R[x] \to S$ qui envoie $x$ sur $t$. Posons $S' = R[x]/I$,
de sorte que $S' \subset S$. Voici un diagramme
$$
\xymatrix{
R[x] \ar[r] & S' \ar[r] & S \\
R \ar[u] \ar[ru] \ar[rru] & &
}
$$
Par construction, les idéaux premiers $\mathfrak q_j$, $j \geq 2$, de $S$
se trouvent tous au-dessus de l'idéal premier $(\mathfrak p, x)$ de $R[x]$,
tandis que l'idéal premier $\mathfrak q$ se trouve au-dessus d'un idéal
premier différent de $R[x]$, puisque $\alpha \not = 0$.

\medskip\noindent
Étape 6. Notons $\mathfrak q' \subset S'$ l'idéal premier de $S'$
correspondant à $\mathfrak q$. D'après ce qui précède, $\mathfrak q$ est
le seul idéal premier de $S$ au-dessus de $\mathfrak q'$. Nous voyons donc que
$S_{\mathfrak q} = S_{\mathfrak q'}$ ; voir le
Lemme \ref{lemma-unique-prime-over-localize-below}
(nous disposons de la propriété de montée pour $S' \to S$ par le
Lemme \ref{lemma-integral-going-up}, puisque $S' \to S$ est fini car
$R \to S$ est fini).
Il s'ensuit que $S'_{\mathfrak q'} \to S_{\mathfrak q}$ est fini et injectif,
comme localisation du morphisme d'anneaux fini injectif $S' \to S$.
Considérons les morphismes d'anneaux locaux
$$
R_{\mathfrak p} \to S'_{\mathfrak q'} \to S_{\mathfrak q}
$$
Le second morphisme est fini et injectif. Nous avons
$S_{\mathfrak q}/\mathfrak pS_{\mathfrak q} = \kappa(\mathfrak q)$ ;
voir le Lemme \ref{lemma-etale-at-prime}.
Donc, a fortiori,
$S_{\mathfrak q}/\mathfrak q'S_{\mathfrak q} = \kappa(\mathfrak q)$.
Puisque
$$
\kappa(\mathfrak p) \subset \kappa(\mathfrak q') \subset \kappa(\mathfrak q)
$$
et puisque $\alpha$ appartient à l'image de $\kappa(\mathfrak q')$ dans
$\kappa(\mathfrak q)$, nous concluons que
$\kappa(\mathfrak q') = \kappa(\mathfrak q)$.
Ainsi, par le Lemme de Nakayama \ref{lemma-NAK} appliqué au morphisme de
$S'_{\mathfrak q'}$-modules $S'_{\mathfrak q'} \to S_{\mathfrak q}$,
le morphisme $S'_{\mathfrak q'} \to S_{\mathfrak q}$ est surjectif.
Autrement dit,
$S'_{\mathfrak q'} \cong S_{\mathfrak q}$.

\medskip\noindent
Étape 7. Par le Lemme \ref{lemma-isomorphic-local-rings}, il existe
$g \in S$, $g \not \in \mathfrak q$, et
$g' \in S'$, $g' \not \in \mathfrak q'$, tels que
$S'_{g'} \cong S_g$. Comme $R$ est noethérien, l'anneau $S'$ est fini
sur $R$, car c'est un sous-$R$-module du $R$-module fini $S$.
Après avoir remplacé $S$ par $S'$, nous pouvons donc supposer que
(a) $R$ est noethérien, (b) $S$ est fini sur $R$, (c)
$S$ est étale sur $R$ en $\mathfrak q$ et (d) $S = R[x]/I$.

\medskip\noindent
Étape 8. Considérons l'anneau
$S \otimes_R \kappa(\mathfrak p) = \kappa(\mathfrak p)[x]/\overline{I}$,
où $\overline{I} = I \cdot \kappa(\mathfrak p)[x]$ est l'idéal engendré par
$I$ dans $\kappa(\mathfrak p)[x]$. Comme $\kappa(\mathfrak p)[x]$ est un
anneau principal, nous savons que
$\overline{I} = (\overline{h})$ pour un certain polynôme unitaire
$\overline{h} \in \kappa(\mathfrak p)[x]$.
Après avoir remplacé $\overline{h}$ par
$\lambda \cdot \overline{h}$ pour un certain
$\lambda \in \kappa(\mathfrak p)$, nous pouvons supposer que
$\overline{h}$ est l'image d'un certain $h \in I \subset R[x]$.
(Le problème est que nous ne savons pas si nous pouvons choisir $h$ unitaire.)
De plus, comme à l'Étape 4, nous savons que
$S \otimes_R \kappa(\mathfrak p) = A_1 \times \ldots \times A_n$, avec
$A_1 = \kappa(\mathfrak q)$ une extension finie séparable de
$\kappa(\mathfrak p)$ et $A_2, \ldots, A_n$ locaux. Cela implique que
$$
\overline{h} = \overline{h}_1 \overline{h}_2^{e_2} \ldots \overline{h}_n^{e_n}
$$
pour certains polynômes unitaires irréductibles deux à deux premiers entre eux
$\overline{h}_i \in \kappa(\mathfrak p)[x]$ et certains
$e_2, \ldots, e_n \geq 1$. Ici, la numérotation est choisie de telle sorte que
$A_i = \kappa(\mathfrak p)[x]/(\overline{h}_i^{e_i})$ comme
$\kappa(\mathfrak p)[x]$-algèbres. Remarquons que $\overline{h}_1$ est le
polynôme minimal de $\alpha \in \kappa(\mathfrak q)$ et est donc un polynôme
séparable (sa dérivée est première avec lui-même).

\medskip\noindent
Étape 9. Soit $m \in I$ un élément unitaire ; un tel élément existe parce que
l'extension d'anneaux $R \to R[x]/I$ est finie, donc entière.
Notons $\overline{m}$ l'image dans $\kappa(\mathfrak p)[x]$.
Nous pouvons factoriser
$$
\overline{m} = \overline{k}
\overline{h}_1^{d_1} \overline{h}_2^{d_2} \ldots \overline{h}_n^{d_n}
$$
pour un certain $d_1 \geq 1$, des $d_j \geq e_j$,
$j = 2, \ldots, n$, et un
$\overline{k} \in \kappa(\mathfrak p)[x]$ premier avec tous les
$\overline{h}_i$. Posons $f = m^l + h$, où
$l \deg(m) > \deg(h)$ et $l \geq 2$.
Alors $f$ est unitaire comme polynôme sur $R$. De plus, l'image
$\overline{f}$ de $f$ dans $\kappa(\mathfrak p)[x]$ se factorise sous la forme
$$
\overline{f} =
\overline{h}_1 \overline{h}_2^{e_2} \ldots \overline{h}_n^{e_n}
+
\overline{k}^l \overline{h}_1^{ld_1} \overline{h}_2^{ld_2}
\ldots \overline{h}_n^{ld_n}
=
\overline{h}_1(\overline{h}_2^{e_2} \ldots \overline{h}_n^{e_n}
+
\overline{k}^l
\overline{h}_1^{ld_1 - 1} \overline{h}_2^{ld_2} \ldots \overline{h}_n^{ld_n})
= \overline{h}_1 \overline{w}
$$
avec $\overline{w}$ un polynôme premier avec $\overline{h}_1$.
Posons $g = f'$ (la dérivée par rapport à $x$).

\medskip\noindent
Étape 10. Le morphisme d'anneaux $R[x] \to S = R[x]/I$ possède les propriétés
suivantes : (1) il envoie $f$ sur zéro, et
(2) il envoie $g$ sur un élément de $S \setminus \mathfrak q$.
La première assertion est claire, puisque $f$ est un élément de $I$.
Pour la seconde, il suffit de montrer que $g$ ne s'envoie pas sur zéro dans
$\kappa(\mathfrak q) = \kappa(\mathfrak p)[x]/(\overline{h}_1)$.
L'image de $g$ dans $\kappa(\mathfrak p)[x]$ est la dérivée de
$\overline{f}$. Ainsi, (2) est clair parce que
$$
\overline{g} =
\frac{\text{d}\overline{f}}{\text{d}x} =
\overline{w}\frac{\text{d}\overline{h}_1}{\text{d}x} +
\overline{h}_1\frac{\text{d}\overline{w}}{\text{d}x},
$$
$\overline{w}$ est premier avec $\overline{h}_1$ et
$\overline{h}_1$ est séparable.

\medskip\noindent
Étape 11.
Nous concluons que $\varphi : R[x]/(f) \to S$ est un morphisme d'anneaux
surjectif, que $R[x]_g/(f)$ est étale sur $R$
(parce qu'il est étale standard ; voir le Lemme \ref{lemma-standard-etale})
et que $\varphi(g) \not \in \mathfrak q$.
Choisissons un élément $g' \in R[x]/(f)$ tel que
$\varphi(g') \not \in \mathfrak q$ lui aussi et que
$S_{\varphi(g')}$ soit étale sur $R$ (un tel élément existe puisque $S$ est
étale sur $R$ en $\mathfrak q$). Alors le morphisme d'anneaux
$R[x]_{gg'}/(f) \to S_{\varphi(gg')}$ est un morphisme surjectif d'algèbres
étales sur $R$. Il est donc étale par le
Lemme \ref{lemma-map-between-etale}.
C'est donc une localisation par le
Lemme \ref{lemma-surjective-flat-finitely-presented}.
Ainsi, une localisation de $S$ en un élément n'appartenant pas à
$\mathfrak q$ est isomorphe à une localisation d'une algèbre étale standard
sur $R$, ce que nous voulions démontrer.
\end{proof}

\noindent
Les deux lemmes suivants affirment que la topologie étale est moins fine que
la topologie engendrée par les recouvrements de Zariski et les morphismes
finis plats. Ils devraient être omis lors d'une première lecture.

\begin{lemma}
\label{lemma-standard-etale-finite-flat-Zariski}
Soit $R \to S$ un morphisme étale standard.
Il existe un morphisme d'anneaux $R \to S'$ possédant les propriétés suivantes :
\begin{enumerate}
\item $R \to S'$ est fini, de présentation finie et plat
(autrement dit, $S'$ est un $R$-module projectif de type fini),
\item $\Spec(S') \to \Spec(R)$ est surjectif,
\item pour tout idéal premier $\mathfrak q \subset S$ au-dessus de
$\mathfrak p \subset R$ et tout idéal premier
$\mathfrak q' \subset S'$ au-dessus de $\mathfrak p$, il existe
$g' \in S'$, $g' \not \in \mathfrak q'$, tel que le morphisme d'anneaux
$R \to S'_{g'}$ se factorise par un morphisme
$\varphi : S \to S'_{g'}$ avec
$\varphi^{-1}(\mathfrak q'S'_{g'}) = \mathfrak q$.
\end{enumerate}
\end{lemma}

\begin{proof}
Soit $S = R[x]_g/(f)$ une présentation de $S$ comme dans la
Définition \ref{definition-standard-etale}.
Écrivons $f = x^n + a_1 x^{n - 1} + \ldots + a_n$, avec $a_i \in R$.
Par le Lemme \ref{lemma-adjoin-roots}, il existe un morphisme d'anneaux
fini libre et fidèlement plat $R \to S'$ tel que
$f = \prod (x - \alpha_i)$ pour certains $\alpha_i \in S'$.
Ainsi, $R \to S'$ satisfait les conditions (1) et (2).
Soit $\mathfrak q \subset R[x]/(f)$ un idéal premier avec
$g \not \in \mathfrak q$ (c'est-à-dire qu'il correspond à un idéal premier
de $S$). Soit $\mathfrak p = R \cap \mathfrak q$, et soit
$\mathfrak q' \subset S'$ un idéal premier au-dessus de $\mathfrak p$.
Remarquons qu'il existe $n$ morphismes de $R$-algèbres
\begin{eqnarray*}
\varphi_i : R[x]/(f) & \longrightarrow & S' \\
x & \longmapsto & \alpha_i
\end{eqnarray*}
Pour achever la démonstration, nous devons montrer que, pour un certain $i$,
(a) l'image de $\varphi_i(g)$ dans $\kappa(\mathfrak q')$ n'est pas nulle,
et (b) $\varphi_i^{-1}(\mathfrak q') = \mathfrak q$.
Nous pourrons alors simplement prendre $g' = \varphi_i(g)$ et
$\varphi = \varphi_i$ pour cet $i$.

\medskip\noindent
Notons $\overline{f}$ l'image de $f$ dans $\kappa(\mathfrak p)[x]$.
Remarquons que, comme point de
$\Spec(\kappa(\mathfrak p)[x]/(\overline{f}))$, l'idéal premier
$\mathfrak q$ correspond à un facteur irréductible $f_1$ de
$\overline{f}$. De plus, $g \not \in \mathfrak q$ signifie que $f_1$ ne
divise pas l'image $\overline{g}$ de $g$ dans $\kappa(\mathfrak p)[x]$.
Notons $\overline{\alpha}_1, \ldots, \overline{\alpha}_n$ les images de
$\alpha_1, \ldots, \alpha_n$ dans $\kappa(\mathfrak q')$.
Remarquons que le polynôme $\overline{f}$ se scinde complètement dans
$\kappa(\mathfrak q')[x]$, à savoir
$$
\overline{f} = \prod\nolimits_i (x - \overline{\alpha}_i)
$$
De plus, $\varphi_i(g)$ se réduit à
$\overline{g}(\overline{\alpha}_i)$.
Il s'ensuit que nous pouvons choisir $i$ tel que
$f_1(\overline{\alpha}_i) = 0$ et
$\overline{g}(\overline{\alpha}_i) \not = 0$.
Pour cet $i$, les propriétés (a) et (b) sont satisfaites.
Certains détails sont omis.
\end{proof}

\begin{lemma}
\label{lemma-etale-finite-flat-zariski}
Soit $R \to S$ un morphisme d'anneaux.
Supposons que
\begin{enumerate}
\item $R \to S$ soit étale, et
\item $\Spec(S) \to \Spec(R)$ soit surjectif.
\end{enumerate}
Alors il existe un morphisme d'anneaux $R \to S'$ tel que
\begin{enumerate}
\item $R \to S'$ soit fini, de présentation finie et plat
(autrement dit, il est projectif de type fini comme $R$-module),
\item $\Spec(S') \to \Spec(R)$ soit surjectif,
\item pour tout idéal premier $\mathfrak q' \subset S'$, il existe
$g' \in S'$, $g' \not \in \mathfrak q'$, tel que le morphisme d'anneaux
$R \to S'_{g'}$ se factorise sous la forme
$R \to S \to S'_{g'}$.
\end{enumerate}
\end{lemma}

\begin{proof}
Par la Proposition \ref{proposition-etale-locally-standard} et la
quasi-compacité de $\Spec(S)$ (voir le Lemme \ref{lemma-quasi-compact}),
nous pouvons trouver des éléments $g_1, \ldots, g_n \in S$ qui engendrent
l'idéal unité de $S$ et tels que chacun des $R \to S_{g_i}$ soit étale
standard. Si nous démontrons le lemme pour le morphisme d'anneaux
$R \to \prod_{i = 1, \ldots, n} S_{g_i}$, alors le lemme en découle pour
le morphisme d'anneaux $R \to S$.
Nous pouvons donc supposer que
$S = \prod_{i = 1, \ldots, n} S_i$ est un produit fini de morphismes
étales standard.

\medskip\noindent
Pour chaque $i$, choisissons un morphisme d'anneaux $R \to S_i'$ comme dans le
Lemme \ref{lemma-standard-etale-finite-flat-Zariski}, adapté au morphisme
étale standard $R \to S_i$.
Posons $S' = S_1' \otimes_R \ldots \otimes_R S_n'$ ; ci-dessous, nous
utiliserons sans autre mention les morphismes de $R$-algèbres
$S_i' \to S'$. Nous affirmons que cela convient.
Les propriétés (1) et (2) sont immédiates.
Pour la propriété (3), supposons que $\mathfrak q' \subset S'$ soit un idéal
premier. Notons $\mathfrak p$ son image dans $\Spec(R)$.
Choisissons $i \in \{1, \ldots, n\}$ tel que $\mathfrak p$ appartienne à
l'image de $\Spec(S_i) \to \Spec(R)$ ; cela est possible par hypothèse.
Soit $\mathfrak q_i' \subset S_i'$ l'image de $\mathfrak q'$ dans le spectre
de $S_i'$. Par construction de $S'_i$, il existe
$g'_i \in S_i'$ tel que $R \to (S_i')_{g_i'}$ se factorise sous la forme
$R \to S_i \to (S_i')_{g_i'}$. Par conséquent,
$R \to S'_{g_i'}$ se factorise également sous la forme
$$
R \to S_i \to (S_i')_{g_i'} \to S'_{g_i'}
$$
comme désiré.
\end{proof}







\section{Structure locale étale des morphismes d'anneaux quasi-finis}
\label{section-etale-local-quasi-finite}

\noindent
Les lemmes suivants affirment, en gros, qu'après une extension étale, un
morphisme d'anneaux quasi-fini devient fini.
Pour faciliter l'interprétation des résultats, rappelons que le lieu où un
morphisme d'anneaux de type fini est quasi-fini est ouvert
(voir le Lemme \ref{lemma-quasi-finite-open}) et que la formation de ce lieu
commute aux changements de base arbitraires
(voir le Lemme \ref{lemma-quasi-finite-base-change}).

\begin{lemma}
\label{lemma-produce-finite}
Soient $R \to S' \to S$ des morphismes d'anneaux.
Soit $\mathfrak p \subset R$ un idéal premier.
Soit $g \in S'$ un élément.
Supposons que
\begin{enumerate}
\item $R \to S'$ soit entier,
\item $R \to S$ soit de type fini,
\item $S'_g \cong S_g$, et
\item $g$ soit inversible dans $S' \otimes_R \kappa(\mathfrak p)$.
\end{enumerate}
Alors il existe $f \in R$, $f \not \in \mathfrak p$, tel que
$R_f \to S_f$ soit fini.
\end{lemma}

\begin{proof}
Par hypothèse, l'image $T$ de $V(g) \subset \Spec(S')$ par le morphisme
$\Spec(S') \to \Spec(R)$ ne contient pas $\mathfrak p$.
Par la section \ref{section-going-up}, et en particulier le
Lemme \ref{lemma-going-up-closed}, nous voyons que $T$ est fermé.
Choisissons $f \in R$, $f \not \in \mathfrak p$, tel que
$T \cap D(f) = \emptyset$. Nous voyons alors que $g$ devient inversible dans
$S'_f$. Par conséquent, $S'_f \cong S_f$.
Ainsi, $S_f$ est à la fois de type fini et entier sur $R_f$, donc fini.
\end{proof}

\begin{lemma}
\label{lemma-etale-makes-quasi-finite-finite-one-prime}
Soit $R \to S$ un morphisme d'anneaux.
Soit $\mathfrak q \subset S$ un idéal premier au-dessus de l'idéal premier
$\mathfrak p \subset R$.
Supposons que $R \to S$ soit de type fini et quasi-fini en $\mathfrak q$.
Alors il existe
\begin{enumerate}
\item un morphisme d'anneaux étale $R \to R'$,
\item un idéal premier $\mathfrak p' \subset R'$ au-dessus de $\mathfrak p$,
\item une décomposition en produit
$$
R' \otimes_R S = A \times B
$$
\end{enumerate}
possédant les propriétés suivantes :
\begin{enumerate}
\item $\kappa(\mathfrak p) = \kappa(\mathfrak p')$,
\item $R' \to A$ est fini,
\item $A$ possède exactement un idéal premier $\mathfrak r$ au-dessus de
$\mathfrak p'$,
\item $\mathfrak r$ est au-dessus de $\mathfrak q$, et
\item $B$ ne possède aucun idéal premier au-dessus de $\mathfrak q$ et de
$\mathfrak p'$.
\end{enumerate}
\end{lemma}

\begin{proof}
Soit $S' \subset S$ la clôture intégrale de $R$ dans $S$.
Soit $\mathfrak q' = S' \cap \mathfrak q$.
Par le Théorème principal de Zariski \ref{theorem-main-theorem}, il existe
$g \in S'$, $g \not \in \mathfrak q'$, tel que $S'_g \cong S_g$.
Considérons les anneaux fibres
$F = S \otimes_R \kappa(\mathfrak p)$ et
$F' = S' \otimes_R \kappa(\mathfrak p)$.
Notons $\overline{\mathfrak q}'$ l'idéal premier de $F'$ correspondant à
$\mathfrak q'$. Comme $F'$ est entier sur $\kappa(\mathfrak p)$, nous voyons
que $\overline{\mathfrak q}'$ est un point fermé de $\Spec(F')$ ;
voir le Lemme \ref{lemma-integral-over-field}.
Remarquons que $\mathfrak q$ définit un point fermé isolé
$\overline{\mathfrak q}$ de $\Spec(F)$
(voir la Définition \ref{definition-quasi-finite}).
Comme $S'_g \cong S_g$, nous avons $F'_g \cong F_g$ ; ainsi,
$\overline{\mathfrak q}$ et $\overline{\mathfrak q}'$ possèdent des
voisinages ouverts isomorphes dans $\Spec(F)$ et $\Spec(F')$.
Nous concluons que l'ensemble
$\{\overline{\mathfrak q}'\} \subset \Spec(F')$ est ouvert.
En combinant cela avec le fait que $\mathfrak q'$ est fermé
(démontré ci-dessus), nous concluons que $\overline{\mathfrak q}'$ définit
lui aussi un point fermé isolé de $\Spec(F')$.

\medskip\noindent
Une petite remarque supplémentaire est que, sous le morphisme
$\Spec(F) \to \Spec(F')$, le point $\overline{\mathfrak q}$ est l'unique
point envoyé sur $\overline{\mathfrak q}'$. Cela découle de la discussion
ci-dessus.

\medskip\noindent
Par le Lemme \ref{lemma-disjoint-implies-product}, nous pouvons écrire
$F' = F'_1 \times F'_2$, avec
$\Spec(F'_1) = \{\overline{\mathfrak q}'\}$.
Comme $F' = S' \otimes_R \kappa(\mathfrak p)$, il existe
$s' \in S'$ qui s'envoie sur l'élément
$(r, 0) \in F'_1 \times F'_2 = F'$ pour un certain
$r \in R$, $r \not \in \mathfrak p$.
En fait, ce que nous utiliserons au sujet de $s'$, c'est qu'il s'agit d'un
élément de $S'$ qui n'appartient pas à $\mathfrak q'$ et qui appartient à tout
autre idéal premier au-dessus de $\mathfrak p$.

\medskip\noindent
Soit $f(x) \in R[x]$ un polynôme unitaire tel que $f(s') = 0$.
Notons $\overline{f} \in \kappa(\mathfrak p)[x]$ son image.
Nous pouvons le factoriser sous la forme
$\overline{f} = x^e \overline{h}$, où
$\overline{h}(0) \not = 0$. Après avoir remplacé $f$ par $x f$ si nécessaire,
nous pouvons supposer que $e \geq 1$.
Par le Lemme \ref{lemma-factor-mod-lift-etale}, nous pouvons trouver une
extension d'anneaux étale $R \to R'$, un idéal premier
$\mathfrak p'$ au-dessus de $\mathfrak p$ et une factorisation
$f = h i$ dans $R'[x]$, tels que
$\kappa(\mathfrak p) = \kappa(\mathfrak p')$,
$\overline{h} = h \bmod \mathfrak p'$,
$x^e = i \bmod \mathfrak p'$, et que nous puissions écrire
$a h + b i = 1$ dans $R'[x]$ (pour des $a, b$ convenables).

\medskip\noindent
Considérons les éléments $h(s'), i(s') \in R' \otimes_R S'$.
Par construction, nous avons $h(s')i(s') = f(s') = 0$.
D'autre part, ils engendrent l'idéal unité puisque
$a(s')h(s') + b(s')i(s') = 1$.
Nous voyons donc que $R' \otimes_R S'$ est le produit des localisations en
ces éléments :
$$
R' \otimes_R S'
=
(R' \otimes_R S')_{i(s')}
\times
(R' \otimes_R S')_{h(s')}
=
S'_1 \times S'_2
$$
De plus, cette décomposition en produit est compatible avec celle que nous
avons trouvée pour l'anneau fibre $F'$ ; cela résulte de nos choix de
$s', i, h$, qui garantissent que $\overline{\mathfrak q}'$ est le seul idéal
premier de $F'$ qui ne contient pas l'image de $i(s')$ dans $F'$.
Nous utilisons ici le fait que l'anneau fibre de
$R'\otimes_R S'$ sur $R'$ en $\mathfrak p'$ est le même que $F'$, puisque
$\kappa(\mathfrak p) = \kappa(\mathfrak p')$.
Il s'ensuit que $S'_1$ possède exactement un idéal premier,
disons $\mathfrak r'$, au-dessus de $\mathfrak p'$, et que cet idéal premier
est au-dessus de $\mathfrak q'$.
Par conséquent, l'élément $g \in S'$ s'envoie sur un élément de $S'_1$
qui n'appartient pas à $\mathfrak r'$.

\medskip\noindent
Le changement de base $R'\otimes_R S$ hérite d'une décomposition en produit
semblable :
$$
R' \otimes_R S
=
(R' \otimes_R S)_{i(s')}
\times
(R' \otimes_R S)_{h(s')}
=
S_1 \times S_2
$$
Il résulte de ce qui précède que $S_1$ possède exactement un idéal premier,
disons $\mathfrak r$, au-dessus de $\mathfrak p'$
(considérer l'anneau fibre comme ci-dessus), et que cet idéal premier est
au-dessus de $\mathfrak q$.

\medskip\noindent
Nous pouvons maintenant appliquer le Lemme \ref{lemma-produce-finite} aux
morphismes d'anneaux $R' \to S'_1 \to S_1$, à l'idéal premier
$\mathfrak p'$ et à l'élément $g$ pour voir qu'après avoir remplacé $R'$ par
une localisation principale, nous pouvons supposer que $S_1$ est fini sur
$R'$, comme désiré.
\end{proof}

\begin{lemma}
\label{lemma-etale-makes-quasi-finite-finite}
Soit $R \to S$ un morphisme d'anneaux.
Soit $\mathfrak p \subset R$ un idéal premier.
Supposons que $R \to S$ soit de type fini.
Alors il existe
\begin{enumerate}
\item un morphisme d'anneaux étale $R \to R'$,
\item un idéal premier $\mathfrak p' \subset R'$ au-dessus de $\mathfrak p$,
\item une décomposition en produit
$$
R' \otimes_R S = A_1 \times \ldots \times A_n \times B
$$
\end{enumerate}
possédant les propriétés suivantes :
\begin{enumerate}
\item nous avons $\kappa(\mathfrak p) = \kappa(\mathfrak p')$,
\item chaque $A_i$ est fini sur $R'$,
\item chaque $A_i$ possède exactement un idéal premier $\mathfrak r_i$
au-dessus de $\mathfrak p'$, et
\item $R' \to B$ n'est quasi-fini en aucun idéal premier au-dessus de
$\mathfrak p'$.
\end{enumerate}
\end{lemma}

\begin{proof}
Notons $F = S \otimes_R \kappa(\mathfrak p)$ l'anneau fibre de $S/R$ en
l'idéal premier $\mathfrak p$. Comme $F$ est de type fini sur
$\kappa(\mathfrak p)$, il est noethérien, et $\Spec(F)$ possède donc un nombre
fini de points fermés isolés. S'il n'existe aucun point fermé isolé,
c'est-à-dire aucun idéal premier $\mathfrak q$ de $S$ au-dessus de
$\mathfrak p$ tel que $S/R$ soit quasi-fini en $\mathfrak q$, alors le lemme
est satisfait. S'il existe au moins un tel idéal premier $\mathfrak q$,
nous pouvons appliquer le
Lemme \ref{lemma-etale-makes-quasi-finite-finite-one-prime}.
Cela donne un diagramme
$$
\xymatrix{
S \ar[r] & R'\otimes_R S \ar@{=}[r] & A_1 \times B' \\
R \ar[r] \ar[u] & R' \ar[u] \ar[ru]
}
$$
comme dans ledit lemme. Comme les corps résiduels en $\mathfrak p$ et
$\mathfrak p'$ sont les mêmes, les anneaux fibres de $S/R$ et de
$(A_1 \times B')/R'$ sont les mêmes.
Ainsi, par récurrence sur le nombre de points fermés isolés de la fibre, nous
pouvons supposer que le lemme est satisfait pour $R' \to B'$ et
$\mathfrak p'$. Nous obtenons donc un morphisme d'anneaux étale
$R' \to R''$, un idéal premier $\mathfrak p'' \subset R''$ et une
décomposition
$$
R'' \otimes_{R'} B' = A_2 \times \ldots \times A_n \times B
$$
Nous omettons de vérifier que le morphisme d'anneaux $R \to R''$, l'idéal
premier $\mathfrak p''$ et la décomposition qui en résulte
$$
R'' \otimes_R S = (R'' \otimes_{R'} A_1) \times
A_2 \times \ldots \times A_n \times B
$$
constituent une solution au problème posé dans le lemme.
\end{proof}

\begin{lemma}
\label{lemma-etale-makes-quasi-finite-finite-variant}
Soit $R \to S$ un morphisme d'anneaux.
Soit $\mathfrak p \subset R$ un idéal premier.
Supposons que $R \to S$ soit de type fini.
Alors il existe
\begin{enumerate}
\item un morphisme d'anneaux étale $R \to R'$,
\item un idéal premier $\mathfrak p' \subset R'$ au-dessus de $\mathfrak p$,
\item une décomposition en produit
$$
R' \otimes_R S = A_1 \times \ldots \times A_n \times B
$$
\end{enumerate}
possédant les propriétés suivantes :
\begin{enumerate}
\item chaque $A_i$ est fini sur $R'$,
\item chaque $A_i$ possède exactement un idéal premier $\mathfrak r_i$
au-dessus de $\mathfrak p'$,
\item les extensions de corps finies
$\kappa(\mathfrak r_i)/\kappa(\mathfrak p')$ sont purement inséparables, et
\item $R' \to B$ n'est quasi-fini en aucun idéal premier au-dessus de
$\mathfrak p'$.
\end{enumerate}
\end{lemma}

\begin{proof}
La stratégie de la démonstration consiste à effectuer deux extensions
d'anneaux étales : nous contrôlons d'abord les corps résiduels, puis nous
appliquons le Lemme \ref{lemma-etale-makes-quasi-finite-finite}.

\medskip\noindent
Notons $F = S \otimes_R \kappa(\mathfrak p)$ l'anneau fibre de $S/R$ en
l'idéal premier $\mathfrak p$.
Comme dans la démonstration du
Lemme \ref{lemma-etale-makes-quasi-finite-finite}, il existe un nombre fini
d'idéaux premiers, disons $\mathfrak q_1, \ldots, \mathfrak q_n$, de $S$
au-dessus de $R$, auxquels le morphisme d'anneaux
$R \to S$ est quasi-fini.
Soit $\kappa(\mathfrak p) \subset L_i \subset \kappa(\mathfrak q_i)$ le
sous-corps tel que $\kappa(\mathfrak p) \subset L_i$ soit séparable et que
l'extension de corps $\kappa(\mathfrak q_i)/L_i$ soit purement inséparable.
Soit $L/\kappa(\mathfrak p)$ une extension galoisienne finie dans laquelle
$L_i$ se plonge pour $i = 1, \ldots, n$.
Par le Lemme \ref{lemma-make-etale-map-prescribed-residue-field}, nous pouvons
trouver une extension d'anneaux étale $R \to R'$, ainsi qu'un idéal premier
$\mathfrak p'$ au-dessus de $\mathfrak p$, tels que l'extension de corps
$\kappa(\mathfrak p')/\kappa(\mathfrak p)$ soit isomorphe à
$\kappa(\mathfrak p) \subset L$.
Ainsi, l'anneau fibre de $R' \otimes_R S$ en $\mathfrak p'$ est isomorphe à
$F \otimes_{\kappa(\mathfrak p)} L$.
Les idéaux premiers au-dessus de $\mathfrak q_i$ correspondent aux idéaux
premiers de
$\kappa(\mathfrak q_i) \otimes_{\kappa(\mathfrak p)} L$, qui est un produit
de corps purement inséparables sur $L$, par notre choix de $L$ et la théorie
élémentaire des corps.
Ce sont également les seuls idéaux premiers au-dessus de $\mathfrak p'$
auxquels $R' \to R' \otimes_R S$ est quasi-fini, par le
Lemme \ref{lemma-quasi-finite-base-change}.
Après avoir remplacé $R$ par $R'$, $\mathfrak p$ par $\mathfrak p'$ et
$S$ par $R' \otimes_R S$, nous pouvons donc supposer que, pour tous les idéaux
premiers $\mathfrak q$ au-dessus de $\mathfrak p$ auxquels $S/R$ est
quasi-fini, les extensions de corps
$\kappa(\mathfrak q)/\kappa(\mathfrak p)$ sont purement inséparables.

\medskip\noindent
Appliquons ensuite le Lemme \ref{lemma-etale-makes-quasi-finite-finite}.
Le résultat est celui que nous souhaitons, puisque les extensions de corps ne
changent pas sous cette extension d'anneaux étale.
\end{proof}






\section{Homomorphismes locaux}
\label{section-local-homomorphisms}

\noindent
Quelques lemmes qui ne trouvent naturellement place dans aucune section.
Le premier affirme, en termes vagues, qu'un morphisme étale d'anneaux locaux
est un isomorphisme modulo toutes les puissances d'un idéal principal engendré
par un élément non inversible.

\begin{lemma}
\label{lemma-lindel}
\begin{reference}
\cite[Lemme à la page 321]{Lindel}, \cite[Lemme 4.1.5]{KC}
\end{reference}
Soit $(R, \mathfrak m_R) \to (S, \mathfrak m_S)$ un homomorphisme local
d'anneaux locaux. Supposons que $S$ soit la localisation d'une extension
d'anneaux étale de $R$ et que
$\kappa(\mathfrak m_R) \to \kappa(\mathfrak m_S)$ soit un isomorphisme.
Alors il existe $t \in \mathfrak m_R$ tel que
$R/t^nR \to S/t^nS$ soit un isomorphisme pour tout $n \geq 1$.
\end{lemma}

\begin{proof}
Écrivons $S = T_{\mathfrak q}$ pour une certaine $R$-algèbre étale $T$ et un
idéal premier $\mathfrak q \subset T$ au-dessus de $\mathfrak m_R$.
Par la Proposition \ref{proposition-etale-locally-standard}, nous pouvons
supposer que $R \to T$ est étale standard.
Écrivons $T = R[x]_g/(f)$ comme dans la
Définition \ref{definition-standard-etale}.
Par notre hypothèse sur les corps résiduels, nous pouvons choisir
$a \in R$ tel que $x$ et $a$ aient la même image dans
$\kappa(\mathfrak q) = \kappa(\mathfrak m_S) = \kappa(\mathfrak m_R)$.
Après avoir remplacé $x$ par $x - a$, nous pouvons alors supposer que
$\mathfrak q$ est engendré par $x$ et $\mathfrak m_R$ dans $T$.
En particulier, $t = f(0) \in \mathfrak m_R$.
Nous allons montrer que $t = f(0)$ convient.

\medskip\noindent
Écrivons $f = x^d + \sum_{i = 1, \ldots, d - 1} a_i x^i + t$.
Comme $R \to T$ est étale standard, nous trouvons que $a_1$ est une unité de
$R$ : la dérivée de $f$ est inversible dans $T$ et, en particulier,
n'appartient pas à $\mathfrak q$.
Soit
$h = a_1 + a_2 x + \ldots + a_{d - 1} x^{d - 2} + x^{d - 1} \in R[x]$,
de sorte que $f = t + xh$ dans $R[x]$.
Nous voyons que $h \not \in \mathfrak q$ et pouvons donc remplacer $T$ par
$R[x]_{hg}/(f)$. Après ce remplacement, nous voyons que
$$
T/tT = (R/tR)[x]_{hg}/(f) = (R/tR)[x]_{hg}/(xh) = (R/tR)[x]_{hg}/(x)
$$
est un quotient de $R/tR$.
Par le Lemme \ref{lemma-surjective-mod-locally-nilpotent}, nous concluons que
$R/t^nR \to T/t^nT$ est surjectif pour tout $n \geq 1$.
D'autre part, nous savons que le morphisme local plat d'anneaux
$R/t^nR \to S/t^nS$ se factorise par $R/t^nR \to T/t^nT$ pour tout $n$ ;
ces morphismes sont donc également injectifs (un homomorphisme local plat
d'anneaux locaux est fidèlement plat et donc injectif ; voir les
Lemmes \ref{lemma-local-flat-ff} et
\ref{lemma-faithfully-flat-universally-injective}).
Comme $S$ est la localisation de $T$, nous voyons que $S/t^nS$ est la
localisation de $T/t^nT = R/t^nR$ en un idéal premier au-dessus de l'idéal
maximal, mais cet anneau est déjà local, ce qui achève la démonstration.
\end{proof}

\begin{lemma}
\label{lemma-etale-under-finite-flat}
Soit $(R, \mathfrak m_R) \to (S, \mathfrak m_S)$ un homomorphisme local
d'anneaux locaux. Supposons que $S$ soit la localisation d'une extension
d'anneaux étale de $R$. Alors il existe un morphisme d'anneaux
fini, de présentation finie et fidèlement plat $R \to S'$ tel que, pour tout
idéal maximal $\mathfrak m'$ de $S'$, il existe une factorisation
$$
R \to S \to S'_{\mathfrak m'}.
$$
Il s'agit d'une factorisation du morphisme d'anneaux
$R \to S'_{\mathfrak m'}$.
\end{lemma}

\begin{proof}
Écrivons $S = T_{\mathfrak q}$ pour une certaine $R$-algèbre étale $T$.
Par la Proposition \ref{proposition-etale-locally-standard}, nous pouvons
supposer que $T$ est étale standard.
Appliquons le Lemme \ref{lemma-standard-etale-finite-flat-Zariski} au
morphisme d'anneaux $R \to T$ pour obtenir $R \to S'$.
Alors, en particulier, pour tout idéal maximal $\mathfrak m'$ de $S'$, nous
obtenons une factorisation $\varphi : T \to S'_{g'}$ pour un certain
$g' \not \in \mathfrak m'$ tel que
$\mathfrak q = \varphi^{-1}(\mathfrak m'S'_{g'})$.
Ainsi, $\varphi$ induit le morphisme local d'anneaux recherché
$S \to S'_{\mathfrak m'}$.
\end{proof}





\section{Clôture intégrale et changement de base lisse}
\label{section-integral-closure-smooth-base-change}

\begin{lemma}
\label{lemma-trick}
Soit $R$ un anneau.
Soit $f \in R[x]$ un polynôme unitaire.
Soit $R \to B$ un morphisme d'anneaux.
Si $h \in B[x]/(f)$ est entier sur $R$, alors l'élément $f' h$ peut s'écrire
$f'h = \sum_i b_i x^i$, avec des $b_i \in B$ entiers sur $R$.
\end{lemma}

\begin{proof}
Disons que $h^e + r_1 h^{e - 1} + \ldots + r_e = 0$ dans l'anneau
$B[x]/(f)$, avec $r_i \in R$.
Il existe une extension d'anneaux finie libre $B \subset B'$ telle que
$f = (x - \alpha_1) \ldots (x - \alpha_d)$ pour certains
$\alpha_i \in B'$ ; voir le Lemme \ref{lemma-adjoin-roots}.
Remarquons que chaque $\alpha_i$ est entier sur $R$.
Nous pouvons représenter
$h = h_0 + h_1 x + \ldots + h_{d - 1} x^{d - 1}$, avec $h_i \in B$.
Alors le fait universel suivant est valable :
$$
f' h =
\sum\nolimits_{i = 1, \ldots, d}
h(\alpha_i)
(x - \alpha_1) \ldots \widehat{(x - \alpha_i)} \ldots (x - \alpha_d)
$$
dans $B'[x]/(f)$. On le démontre en évaluant les deux membres aux points
$\alpha_i$ sur l'anneau
$B_{univ} = \mathbf{Z}[\alpha_i, h_j]$ (certains détails sont omis).
Comme, par hypothèse, $h$ satisfait
$h^e + r_1 h^{e - 1} + \ldots + r_e = 0$ dans l'anneau $B[x]/(f)$,
nous voyons que
$$
h(\alpha_i)^e + r_1 h(\alpha_i)^{e - 1} + \ldots + r_e = 0
$$
dans $B'$. Ainsi, $h(\alpha_i)$ est entier sur $R$.
En utilisant la formule ci-dessus, nous voyons que
$f'h \equiv \sum_{j = 0, \ldots, d - 1} b'_j x^j$ dans $B'[x]/(f)$,
avec des $b'_j \in B'$ entiers sur $R$.
Cependant, puisque $f' h \in B[x]/(f)$ et que
$1, x, \ldots, x^{d - 1}$ constituent une $B'$-base de $B'[x]/(f)$,
nous voyons que $b'_j \in B$, comme désiré.
\end{proof}

\begin{lemma}
\label{lemma-integral-closure-commutes-etale}
Soit $R \to S$ un morphisme d'anneaux étale.
Soit $R \to B$ un morphisme d'anneaux quelconque.
Soit $A \subset B$ la clôture intégrale de $R$ dans $B$.
Soit $A' \subset S \otimes_R B$ la clôture intégrale de $S$ dans
$S \otimes_R B$. Alors le morphisme canonique
$S \otimes_R A \to A'$ est un isomorphisme.
\end{lemma}

\begin{proof}
Le morphisme $S \otimes_R A \to A'$ est injectif parce que
$A \subset B$ et que $R \to S$ est plat.
Nous allons utiliser à plusieurs reprises le fait que la formation de la
clôture intégrale commute à la localisation ; voir le
Lemme \ref{lemma-integral-closure-localize}.
Nous pouvons donc localiser sur $S$, par le Lemme \ref{lemma-cover}
(le critère qui permet de vérifier si un morphisme de $S$-modules est un
isomorphisme).
Ainsi, nous pouvons supposer que
$S = R[x]_g/(f) = (R[x]/(f))_g$ est étale standard sur $R$ ;
voir la Proposition \ref{proposition-etale-locally-standard}.
En localisant une fois de plus, nous voyons que $A'$ est $(A'')_g$, où
$A''$ est la clôture intégrale de $R[x]/(f)$ dans $B[x]/(f)$.
Supposons que $a \in A''$. Il suffit de montrer que $a$ appartient à
$S \otimes_R A$. Par le Lemme \ref{lemma-trick}, nous voyons que
$f' a = \sum a_i x^i$, avec $a_i \in A$.
Comme $f'$ est inversible dans $S$ (par définition d'un morphisme d'anneaux
étale standard), nous concluons que $a \in S \otimes_R A$, comme désiré.
\end{proof}

\begin{example}
\label{example-fourier}
Soit $p$ un nombre premier. L'extension d'anneaux
$$
R = \mathbf{Z}[1/p] \subset
R' = \mathbf{Z}[1/p][x]/(x^{p - 1} + \ldots + x + 1)
$$
possède la propriété suivante : pour $d < p$, il existe des éléments
$\alpha_0, \ldots, \alpha_{d - 1} \in R'$ tels que
$$
\prod\nolimits_{0 \leq i < j < d} (\alpha_i - \alpha_j)
$$
soit une unité de $R'$. En effet, prenons $\alpha_i$ égal à la classe de
$x^i$ dans $R'$ pour $i = 0, \ldots, p - 1$. Nous avons alors
$$
T^p - 1 = \prod\nolimits_{i = 0, \ldots, p - 1} (T - \alpha_i)
$$
dans $R'[T]$. En effet, l'anneau
$\mathbf{Q}[x]/(x^{p - 1} + \ldots + x + 1)$ est un corps parce que le
polynôme cyclotomique $x^{p - 1} + \ldots + x + 1$ est irréductible sur
$\mathbf{Q}$ et que les $\alpha_i$ sont des racines deux à deux distinctes de
$T^p - 1$, d'où l'égalité.
En dérivant les deux membres puis en substituant $T = \alpha_i$, nous obtenons
$$
p \alpha_i^{p - 1}
=
(\alpha_i - \alpha_1) \ldots
\widehat{(\alpha_i - \alpha_i)} \ldots
(\alpha_i - \alpha_1)
$$
et nous voyons que cet élément est inversible dans $R'$.
\end{example}

\begin{lemma}
\label{lemma-integral-closure-commutes-smooth}
Soit $R \to S$ un morphisme d'anneaux lisse.
Soit $R \to B$ un morphisme d'anneaux quelconque.
Soit $A \subset B$ la clôture intégrale de $R$ dans $B$.
Soit $A' \subset S \otimes_R B$ la clôture intégrale de $S$ dans
$S \otimes_R B$. Alors le morphisme canonique
$S \otimes_R A \to A'$ est un isomorphisme.
\end{lemma}

\begin{proof}
En raisonnant comme dans la démonstration du
Lemme \ref{lemma-integral-closure-commutes-etale}, nous pouvons localiser sur
$S$. Nous pouvons donc supposer que $R \to S$ est un morphisme d'anneaux lisse
standard ; voir le Lemme \ref{lemma-smooth-syntomic}.
Par définition d'un morphisme d'anneaux lisse standard, nous voyons que $S$
est étale sur un anneau de polynômes $R[x_1, \ldots, x_n]$.
Comme nous avons vu le résultat dans le cas d'une extension d'anneaux étale
(Lemme \ref{lemma-integral-closure-commutes-etale}), cela nous ramène au cas
où $S = R[x]$. Nous devons donc montrer
$$
f = \sum b_i x^i
\text{ entier sur }R[x]
\Leftrightarrow
\text{chaque }b_i\text{ est entier sur }R.
$$
L'implication de droite à gauche vaut parce que l'ensemble des éléments de
$B[x]$ entiers sur $R[x]$ est un anneau
(Lemme \ref{lemma-integral-closure-is-ring}) et contient $x$.

\medskip\noindent
Supposons que $f \in B[x]$ soit entier sur $R[x]$ et supposons que
$f = \sum_{i < d} b_i x^i$ soit de degré $< d$.
Comme la clôture intégrale et la localisation commutent, il suffit de montrer
qu'il existe des nombres premiers distincts $p, q$ tels que chaque $b_i$ soit
entier à la fois sur $R[1/p]$ et sur $R[1/q]$.
Nous pouvons donc trouver une extension d'anneaux finie libre
$R \subset R'$ telle que $R'$ contienne
$\alpha_1, \ldots, \alpha_d$ et que
$\prod_{i < j} (\alpha_i - \alpha_j)$ soit une unité de $R'$ ; voir
l'Exemple \ref{example-fourier}.
Dans ce cas, nous avons l'égalité universelle
$$
f
=
\sum_i
f(\alpha_i)
\frac{(x - \alpha_1) \ldots \widehat{(x - \alpha_i)} \ldots (x - \alpha_d)}
{(\alpha_i - \alpha_1) \ldots \widehat{(\alpha_i - \alpha_i)} \ldots
(\alpha_i - \alpha_d)}.
$$
De plus, les éléments $f(\alpha_i)$ sont entiers sur $R'$, puisque
$(R' \otimes_R B)[x] \to R' \otimes_R B$, $h \mapsto h(\alpha_i)$,
est un morphisme d'anneaux.
Nous voyons donc que les coefficients de $f$ dans
$(R' \otimes_R B)[x]$ sont entiers sur $R'$.
Comme $R'$ est fini sur $R$ (donc entier sur $R$), nous voyons qu'ils sont
également entiers sur $R$, comme désiré.
\end{proof}

\begin{lemma}
\label{lemma-integral-closure-commutes-colim-smooth}
Soient $R \to S$ et $R \to B$ des morphismes d'anneaux.
Soit $A \subset B$ la clôture intégrale de $R$ dans $B$.
Soit $A' \subset S \otimes_R B$ la clôture intégrale de $S$ dans
$S \otimes_R B$. Si $S$ est une colimite filtrante de $R$-algèbres lisses,
alors le morphisme canonique $S \otimes_R A \to A'$ est un isomorphisme.
\end{lemma}

\begin{proof}
Cela résulte du fait immédiat que la formation des produits tensoriels et
celle des clôtures intégrales commutent aux colimites filtrantes, ainsi que du
Lemme \ref{lemma-integral-closure-commutes-smooth}.
\end{proof}






\section{Morphismes formellement non ramifiés}
\label{section-formally-unramified}

\noindent
Il s'avère logiquement plus efficace de définir la notion de morphisme
formellement non ramifié avant d'introduire celle de morphisme formellement
étale.

\begin{definition}
\label{definition-formally-unramified}
Soit $R \to S$ un morphisme d'anneaux.
Nous disons que $S$ est {\it formellement non ramifié sur $R$} si, pour tout
diagramme commutatif dont les flèches pleines sont données
$$
\xymatrix{
S \ar[r] \ar@{-->}[rd] & A/I \\
R \ar[r] \ar[u] & A \ar[u]
}
$$
où $I \subset A$ est un idéal de carré nul, il existe au plus une flèche
pointillée rendant le diagramme commutatif.
\end{definition}

\begin{lemma}
\label{lemma-base-change-formally-unramified}
Soit $R \to S$ un morphisme formellement non ramifié.
Soit $R \to R'$ un morphisme d'anneaux quelconque.
Alors le changement de base $S' = R' \otimes_R S$ est formellement non
ramifié sur $R'$.
\end{lemma}

\begin{proof}
Soit donné un diagramme dont les flèches pleines sont données
$$
\xymatrix{
S \ar[r] \ar@{-->}[rrd] & R' \otimes_R S \ar[r] \ar@{-->}[rd] & A/I \\
R  \ar[u] \ar[r] & R' \ar[r] \ar[u] & A \ar[u]
}
$$
comme dans la Définition \ref{definition-formally-unramified}.
Par hypothèse, il existe au plus une flèche pointillée longue.
Par la propriété universelle du produit tensoriel, nous concluons qu'il existe
au plus une flèche pointillée courte.
\end{proof}

\begin{lemma}
\label{lemma-characterize-formally-unramified}
Soit $R \to S$ un morphisme d'anneaux.
Les conditions suivantes sont équivalentes :
\begin{enumerate}
\item $R \to S$ est formellement non ramifié,
\item le module des différentielles $\Omega_{S/R}$ est nul.
\end{enumerate}
\end{lemma}

\begin{proof}
Soit $J = \Ker(S \otimes_R S \to S)$ le noyau du morphisme de multiplication.
Soit $A_{univ} = S \otimes_R S/J^2$. Rappelons que
$I_{univ} = J/J^2$ est isomorphe à $\Omega_{S/R}$ ; voir le
Lemme \ref{lemma-differentials-diagonal}.
De plus, les deux morphismes de $R$-algèbres
$\sigma_1, \sigma_2 : S \to A_{univ}$,
$\sigma_1(s) = s \otimes 1 \bmod J^2$ et
$\sigma_2(s) = 1 \otimes s \bmod J^2$ diffèrent par la dérivation universelle
$\text{d} : S \to \Omega_{S/R} = I_{univ}$.

\medskip\noindent
Supposons que $R \to S$ soit formellement non ramifié.
Nous voyons alors que $\sigma_1 = \sigma_2$.
Ainsi, $\text{d}(s) = 0$ pour tout $s \in S$.
Par conséquent, $\Omega_{S/R} = 0$.

\medskip\noindent
Supposons que $\Omega_{S/R} = 0$.
Soient $A, I, R \to A, S \to A/I$ un diagramme dont les flèches pleines sont
données, comme dans la Définition \ref{definition-formally-unramified}.
Soient $\tau_1, \tau_2 : S \to A$ deux flèches pointillées rendant le
diagramme commutatif. Considérons le morphisme de $R$-algèbres
$A_{univ} \to A$ défini par la règle
$s_1 \otimes s_2 \mapsto \tau_1(s_1)\tau_2(s_2)$.
Nous omettons de vérifier qu'il est bien défini.
Comme $A_{univ} \cong S$, puisque
$I_{univ} = \Omega_{S/R} = 0$, nous concluons que $\tau_1 = \tau_2$.
\end{proof}

\begin{lemma}
\label{lemma-formally-unramified-local}
Soit $R \to S$ un morphisme d'anneaux.
Les conditions suivantes sont équivalentes :
\begin{enumerate}
\item $R \to S$ est formellement non ramifié,
\item $R \to S_{\mathfrak q}$ est formellement non ramifié pour tout idéal
premier $\mathfrak q$ de $S$, et
\item $R_{\mathfrak p} \to S_{\mathfrak q}$ est formellement non ramifié pour
tout idéal premier $\mathfrak q$ de $S$, avec
$\mathfrak p = R \cap \mathfrak q$.
\end{enumerate}
\end{lemma}

\begin{proof}
Nous avons vu dans le
Lemme \ref{lemma-characterize-formally-unramified} que (1) équivaut à
$\Omega_{S/R} = 0$. De même, par le
Lemme \ref{lemma-differentials-localize}, nous voyons que (2) et (3)
équivalent à $(\Omega_{S/R})_{\mathfrak q} = 0$ pour tout
$\mathfrak q$. L'équivalence résulte donc du
Lemme \ref{lemma-characterize-zero-local}.
\end{proof}

\begin{lemma}
\label{lemma-formally-unramified-localize}
Soit $A \to B$ un morphisme d'anneaux formellement non ramifié.
\begin{enumerate}
\item Pour une partie multiplicative $S \subset A$,
$S^{-1}A \to S^{-1}B$ est formellement non ramifié.
\item Pour une partie multiplicative $S \subset B$,
$A \to S^{-1}B$ est formellement non ramifié.
\end{enumerate}
\end{lemma}

\begin{proof}
Cela résulte du Lemme \ref{lemma-formally-unramified-local}.
(On peut aussi le déduire du
Lemme \ref{lemma-characterize-formally-unramified}, combiné au
Lemme \ref{lemma-differentials-localize}.)
\end{proof}

\begin{lemma}
\label{lemma-colimit-formally-unramified}
Soit $R$ un anneau. Soit $I$ un ensemble ordonné filtrant.
Soit $(S_i, \varphi_{ii'})$ un système de $R$-algèbres sur $I$.
Si chaque $R \to S_i$ est formellement non ramifié, alors
$S = \colim_{i \in I} S_i$ est formellement non ramifié sur $R$
\end{lemma}

\begin{proof}
Considérons un diagramme comme dans la
Définition \ref{definition-formally-unramified}.
Par hypothèse, il existe au plus un morphisme de $R$-algèbres
$S_i \to A$ relevant les composées $S_i \to S \to A/I$.
Comme tout élément de $S$ appartient à l'image de l'un des morphismes
$S_i \to S$, nous voyons qu'il existe au plus un morphisme $S \to A$
s'insérant dans le diagramme.
\end{proof}



\section{Modules conormaux et épaississements universels}
\label{section-conormal}

\noindent
Il s'avère que l'on peut définir le premier voisinage infinitésimal non
seulement pour une immersion fermée de schémas, mais déjà pour tout morphisme
formellement non ramifié. Cela repose sur le fait algébrique suivant.

\begin{lemma}
\label{lemma-universal-thickening}
Soit $R \to S$ un morphisme d'anneaux formellement non ramifié.
Il existe une surjection de $R$-algèbres $S' \to S$ dont le noyau est un idéal
de carré nul et qui possède la propriété universelle suivante : étant donné
un diagramme commutatif quelconque
$$
\xymatrix{
S \ar[r]_a & A/I \\
R \ar[r]^b \ar[u] & A \ar[u]
}
$$
où $I \subset A$ est un idéal de carré nul, il existe un unique morphisme de
$R$-algèbres $a' : S' \to A$ tel que $S' \to A \to A/I$ soit égal à
$S' \to S \to A/I$.
\end{lemma}

\begin{proof}
Choisissons un ensemble de générateurs $z_i \in S$, $i \in I$, de $S$ comme
$R$-algèbre.
Soit $P = R[\{x_i\}_{i \in I}]$ l'anneau de polynômes en les générateurs
$x_i$, $i \in I$. Considérons le morphisme de $R$-algèbres $P \to S$ qui
envoie $x_i$ sur $z_i$. Soit $J = \Ker(P \to S)$.
Considérons le morphisme
$$
\text{d} : J/J^2 \longrightarrow \Omega_{P/R} \otimes_P S
$$
voir le Lemme \ref{lemma-differential-seq}.
Il est surjectif puisque $\Omega_{S/R} = 0$ par hypothèse ; voir le
Lemme \ref{lemma-characterize-formally-unramified}.
Remarquons que $\Omega_{P/R}$ est libre sur les $\text{d}x_i$, et que le
module $\Omega_{P/R} \otimes_P S$ est donc libre sur $S$.
Nous pouvons ainsi choisir un scindage de la surjection ci-dessus et écrire
$$
J/J^2 = K \oplus \Omega_{P/R} \otimes_P S
$$
Soit $J^2 \subset J' \subset J$ l'idéal de $P$ tel que $J'/J^2$ soit le
second facteur direct de la décomposition ci-dessus.
Posons $S' = P/J'$. Nous obtenons une suite exacte courte
$$
0 \to J/J' \to S' \to S \to 0
$$
et nous voyons que $J/J' \cong K$ est un idéal de carré nul de $S'$.
Par conséquent,
$$
\xymatrix{
S \ar[r]_1 & S \\
R \ar[r] \ar[u] & S' \ar[u]
}
$$
est un diagramme comme ci-dessus.
En fait, nous affirmons qu'il s'agit d'un objet initial de la catégorie de ces
diagrammes. En effet, soit $(I \subset A, a, b)$ un diagramme arbitraire.
Nous pouvons choisir un morphisme de $R$-algèbres $\beta : P \to A$ tel que
$$
\xymatrix{
S \ar[r]_1 & S \ar[r]_a & A/I \\
R \ar[r] \ar@/_/[rr]_b \ar[u] & P \ar[u] \ar[r]^\beta & A \ar[u]
}
$$
soit commutatif.
Il n'est pas nécessairement vrai que $\beta(J') = 0$ ; autrement dit,
$\beta$ ne se factorise pas nécessairement par $S' = P/J'$.
En revanche, il est clair que $\beta(J') \subset I$ et, puisque
$\beta(J) \subset I$ et $I^2 = 0$, nous avons $\beta(J^2) = 0$.
Ainsi, l'« obstruction » à l'existence d'un morphisme de
$(J/J' \subset S', 1, R \to S')$ vers $(I \subset A, a, b)$ est le morphisme
$S$-linéaire correspondant $\overline{\beta} : J'/J^2 \to I$.
Le choix de $\beta$ réside dans le choix des $\beta(x_i)$.
Un choix différent de $\beta$, disons $\beta'$, s'obtient en prenant
$\beta'(x_i) = \beta(x_i) + \delta_i$, avec $\delta_i \in I$.
Dans ce cas, pour $g \in J'$, nous obtenons
$$
\beta'(g) =
\beta(g) + \sum\nolimits_i \delta_i \beta(\frac{\partial g}{\partial x_i})
$$
Comme le morphisme
$\text{d}|_{J'/J^2} : J'/J^2 \to \Omega_{P/R} \otimes_P S$ donné par
$g \mapsto \frac{\partial g}{\partial x_i}\text{d}x_i \otimes 1$ est un
isomorphisme par construction, nous voyons qu'il existe un choix unique de
$\delta_i \in I$ tel que $\beta'(g) = 0$ pour tout $g \in J'$.
(Plus précisément, $\delta_i$ est $-\overline{\beta}(g)$, où
$g \in J'/J^2$ est l'unique élément envoyé sur
$\text{d}x_i \otimes 1$ par $\text{d}|_{J'/J^2}$.)
L'unicité de la solution implique celle requise dans le lemme.
\end{proof}

\noindent
Dans la situation du Lemme \ref{lemma-universal-thickening}, le morphisme de
$R$-algèbres $S' \to S$ est unique à isomorphisme unique près.

\begin{definition}
\label{definition-universal-thickening}
Soit $R \to S$ un morphisme d'anneaux formellement non ramifié.
\begin{enumerate}
\item L'{\it épaississement universel du premier ordre} de $S$ sur $R$ est la
surjection de $R$-algèbres $S' \to S$ du
Lemme \ref{lemma-universal-thickening}.
\item Le {\it module conormal} de $R \to S$ est le noyau $I$ de
l'épaississement universel du premier ordre $S' \to S$, considéré comme un
$S$-module.
\end{enumerate}
Dans cette situation, nous notons souvent le module conormal {\it $C_{S/R}$}.
\end{definition}

\begin{lemma}
\label{lemma-universal-thickening-quotient}
Soit $I \subset R$ un idéal d'un anneau.
L'épaississement universel du premier ordre de $R/I$ sur $R$ est la
surjection $R/I^2 \to R/I$. Le module conormal de $R/I$ sur $R$ est
$C_{(R/I)/R} = I/I^2$.
\end{lemma}

\begin{proof}
Omis.
\end{proof}

\begin{lemma}
\label{lemma-universal-thickening-localize}
Soit $A \to B$ un morphisme d'anneaux formellement non ramifié.
Soit $\varphi : B' \to B$ l'épaississement universel du premier ordre de
$B$ sur $A$.
\begin{enumerate}
\item Soit $S \subset A$ une partie multiplicative.
Alors $S^{-1}B' \to S^{-1}B$ est l'épaississement universel du premier ordre
de $S^{-1}B$ sur $S^{-1}A$. En particulier,
$S^{-1}C_{B/A} = C_{S^{-1}B/S^{-1}A}$.
\item Soit $S \subset B$ une partie multiplicative.
Alors $S' = \varphi^{-1}(S)$ est une partie multiplicative de $B'$, et
$(S')^{-1}B' \to S^{-1}B$ est l'épaississement universel du premier ordre de
$S^{-1}B$ sur $A$. En particulier,
$S^{-1}C_{B/A} = C_{S^{-1}B/A}$.
\end{enumerate}
Remarquons que le lemme a un sens en vertu du
Lemme \ref{lemma-formally-unramified-localize}.
\end{lemma}

\begin{proof}
Avec les notations et hypothèses de (1), soit
$(S^{-1}B)' \to S^{-1}B$ l'épaississement universel du premier ordre de
$S^{-1}B$ sur $S^{-1}A$.
Remarquons que $S^{-1}B' \to S^{-1}B$ est une surjection de
$S^{-1}A$-algèbres dont le noyau est de carré nul.
Par définition, nous obtenons donc un morphisme
$(S^{-1}B)' \to S^{-1}B'$ compatible avec les morphismes vers $S^{-1}B$.
Considérons un diagramme commutatif quelconque
$$
\xymatrix{
B \ar[r] & S^{-1}B \ar[r] & D/I \\
A \ar[r] \ar[u] & S^{-1}A \ar[r] \ar[u] & D \ar[u]
}
$$
où $I \subset D$ est un idéal de carré nul.
Comme $B'$ est l'épaississement universel du premier ordre de $B$ sur $A$,
nous obtenons un morphisme de $A$-algèbres $B' \to D$.
Mais il est clair que l'image de $S$ dans $D$ est envoyée sur des éléments
inversibles de $D$, et nous obtenons donc un morphisme compatible
$S^{-1}B' \to D$.
En appliquant cela à $D = (S^{-1}B)'$, nous voyons que nous obtenons un
morphisme $S^{-1}B' \to (S^{-1}B)'$.
Nous omettons de vérifier que ce morphisme est inverse du morphisme décrit
ci-dessus.

\medskip\noindent
Avec les notations et hypothèses de (2), soit
$(S^{-1}B)' \to S^{-1}B$ l'épaississement universel du premier ordre de
$S^{-1}B$ sur $A$.
Remarquons que $(S')^{-1}B' \to S^{-1}B$ est une surjection de $A$-algèbres
dont le noyau est de carré nul.
Par définition, nous obtenons donc un morphisme
$(S^{-1}B)' \to (S')^{-1}B'$ compatible avec les morphismes vers
$S^{-1}B$. Considérons un diagramme commutatif quelconque
$$
\xymatrix{
B \ar[r] & S^{-1}B \ar[r] & D/I \\
A \ar[r] \ar[u] & A \ar[r] \ar[u] & D \ar[u]
}
$$
où $I \subset D$ est un idéal de carré nul.
Comme $B'$ est l'épaississement universel du premier ordre de $B$ sur $A$,
nous obtenons un morphisme de $A$-algèbres $B' \to D$.
Mais il est clair que l'image de $S'$ dans $D$ est envoyée sur des éléments
inversibles de $D$, et nous obtenons donc un morphisme compatible
$(S')^{-1}B' \to D$.
En appliquant cela à $D = (S^{-1}B)'$, nous voyons que nous obtenons un
morphisme $(S')^{-1}B' \to (S^{-1}B)'$.
Nous omettons de vérifier que ce morphisme est inverse du morphisme décrit
ci-dessus.
\end{proof}

\begin{lemma}
\label{lemma-differentials-universal-thickening}
Soient $R \to A  \to B$ des morphismes d'anneaux.
Supposons que $A \to B$ soit formellement non ramifié.
Soit $B' \to B$ l'épaississement universel du premier ordre de $B$ sur $A$.
Alors $B'$ est formellement non ramifié sur $A$, et le morphisme canonique
$\Omega_{A/R} \otimes_A B \to \Omega_{B'/R} \otimes_{B'} B$ est un
isomorphisme.
\end{lemma}

\begin{proof}
Nous allons utiliser la construction de $B'$ donnée dans la démonstration du
Lemme \ref{lemma-universal-thickening}, bien qu'il doive en principe être
possible de déduire formellement ces résultats de la définition.
Plus précisément, nous choisissons une présentation
$B = P/J$, où $P = A[x_i]$ est un anneau de polynômes sur $A$.
Ensuite, nous choisissons des éléments $f_i \in J$ tels que
$\text{d}f_i = \text{d}x_i \otimes 1$ dans
$\Omega_{P/A} \otimes_P B$.
Ces choix étant faits, nous avons $B' = P/J'$, avec
$J' = (f_i) + J^2$ ; voir la démonstration du
Lemme \ref{lemma-universal-thickening}.

\medskip\noindent
Considérons la suite exacte canonique
$$
J'/(J')^2 \to \Omega_{P/A} \otimes_P B' \to \Omega_{B'/A} \to 0
$$
voir le Lemme \ref{lemma-differential-seq}.
Par construction, les classes des $f_i \in J'$ s'envoient sur des éléments
du module $\Omega_{P/A} \otimes_P B'$ qui l'engendrent modulo
$J'/J^2$, par construction. Comme $J'/J^2$ est un idéal nilpotent, nous voyons
que ces éléments engendrent le module tout entier (par le
Lemme de Nakayama \ref{lemma-NAK}).
Cela démontre que $\Omega_{B'/A} = 0$ et donc que $B'$ est formellement non
ramifié sur $A$ ; voir le
Lemme \ref{lemma-characterize-formally-unramified}.

\medskip\noindent
Comme $P$ est un anneau de polynômes sur $A$, nous avons
$\Omega_{P/R} = \Omega_{A/R} \otimes_A P \oplus \bigoplus P\text{d}x_i$.
Nous allons utiliser cette décomposition.
Considérons la suite exacte suivante :
$$
J'/(J')^2 \to
\Omega_{P/R} \otimes_P B' \to
\Omega_{B'/R} \to 0
$$
voir le Lemme \ref{lemma-differential-seq}.
Nous pouvons la tensoriser avec $B$ et obtenir la suite exacte
$$
J'/(J')^2 \otimes_{B'} B \to
\Omega_{P/R} \otimes_P B \to
\Omega_{B'/R} \otimes_{B'} B \to 0
$$
Si nous nous rappelons que $J' = (f_i) + J^2$, nous voyons que le premier
morphisme annule le sous-module $J^2/(J')^2$.
En termes de la décomposition en somme directe
$\Omega_{P/R} \otimes_P B =
\Omega_{A/R} \otimes_A B \oplus \bigoplus B\text{d}x_i $ donnée,
nous voyons que le sous-module
$(f_i)/(J')^2 \otimes_{B'} B$ s'envoie isomorphiquement sur le facteur direct
$\bigoplus B\text{d}x_i$. Ce qui reste de cette suite exacte est donc un
isomorphisme
$\Omega_{A/R} \otimes_A B \to \Omega_{B'/R} \otimes_{B'} B$,
comme désiré.
\end{proof}









\section{Morphismes formellement \'etales}
\label{section-formally-etale}

\begin{definition}
\label{definition-formally-etale}
Soit $R \to S$ un morphisme d'anneaux.
On dit que $S$ est {\it formellement \'etale sur $R$} si, pour tout
diagramme commutatif dont les flèches pleines sont données
$$
\xymatrix{
S \ar[r] \ar@{-->}[rd] & A/I \\
R \ar[r] \ar[u] & A \ar[u]
}
$$
où $I \subset A$ est un idéal de carré nul, il existe
une unique flèche pointillée rendant le diagramme commutatif.
\end{definition}

\noindent
Il est clair qu'un morphisme d'anneaux est formellement \'etale si et seulement s'il
est à la fois formellement lisse et formellement non ramifié.

\begin{lemma}
\label{lemma-base-change-formally-etale}
Soit $R \to S$ un morphisme formellement \'etale.
Soit $R \to R'$ un morphisme d'anneaux quelconque.
Alors le changement de base $S' = R' \otimes_R S$ est formellement \'etale sur $R'$.
\end{lemma}

\begin{proof}
Cela résulte des lemmes \ref{lemma-base-change-fs} et
\ref{lemma-base-change-formally-unramified}.
\end{proof}

\begin{lemma}
\label{lemma-formally-etale-etale}
Soit $R \to S$ un morphisme d'anneaux de présentation finie.
Les assertions suivantes sont équivalentes :
\begin{enumerate}
\item $R \to S$ est formellement \'etale,
\item $R \to S$ est \'etale.
\end{enumerate}
\end{lemma}

\begin{proof}
Supposons que $R \to S$ soit formellement \'etale.
Alors $R \to S$ est lisse par la
proposition \ref{proposition-smooth-formally-smooth}.
D'après le lemme \ref{lemma-characterize-formally-unramified},
nous avons $\Omega_{S/R} = 0$.
Ainsi, $R \to S$ est \'etale par définition.

\medskip\noindent
Supposons que $R \to S$ soit \'etale.
Alors $R \to S$ est formellement lisse par la
proposition \ref{proposition-smooth-formally-smooth}.
D'après le lemme \ref{lemma-characterize-formally-unramified},
il est formellement non ramifié. Ainsi, $R \to S$ est formellement \'etale.
\end{proof}

\begin{lemma}
\label{lemma-colimit-formally-etale}
Soit $R$ un anneau. Soit $I$ un ensemble ordonné filtrant.
Soit $(S_i, \varphi_{ii'})$ un système de $R$-algèbres
indexé par $I$. Si chaque $R \to S_i$ est formellement \'etale, alors
$S = \colim_{i \in I} S_i$ est formellement \'etale sur $R$.
\end{lemma}

\begin{proof}
Considérons un diagramme comme dans la définition \ref{definition-formally-etale}.
Par hypothèse, nous obtenons d'uniques morphismes de $R$-algèbres $S_i \to A$ relevant
les composées $S_i \to S \to A/I$. Ils sont donc compatibles
avec les morphismes de transition $\varphi_{ii'}$ et définissent un relèvement
$S \to A$. Cela prouve l'existence.
L'unicité est claire en se restreignant à chaque $S_i$.
\end{proof}

\begin{lemma}
\label{lemma-localization-formally-etale}
Soit $R$ un anneau. Soit $S \subset R$ une partie multiplicative quelconque.
Alors le morphisme d'anneaux $R \to S^{-1}R$ est formellement \'etale.
\end{lemma}

\begin{proof}
Soit $I \subset A$ un idéal de carré nul. Ce que nous affirmons
ici est qu'étant donné un morphisme d'anneaux $\varphi : R \to A$ tel que
$\varphi(f) \mod I$ soit inversible pour tout $f \in S$, nous avons aussi que
$\varphi(f)$ est inversible dans $A$ pour tout $f \in S$. Cela est vrai parce que
$A^*$ est l'image réciproque de $(A/I)^*$ par le morphisme canonique
$A \to A/I$.
\end{proof}

\begin{lemma}
\label{lemma-formally-etale-lift-infinitesimal}
Soit $R \to S$ un morphisme d'anneaux. Soit $J \subset S$ un idéal tel
que $R \to S/J$ soit surjectif ; soit $I \subset R$ le noyau.
Si $R \to S$ est formellement \'etale, alors $R/I^n \to S/J^n$ est un
isomorphisme pour tout $n$, et
$\bigoplus I^n/I^{n + 1} \to \bigoplus J^n/J^{n + 1}$
est un isomorphisme d'anneaux gradués.
\end{lemma}

\begin{proof}
En utilisant la propriété de relèvement par récurrence, nous obtenons des flèches pointillées
$$
\xymatrix{
S \ar[r] \ar@{-->}[rd] & S/J = R/I \\
R \ar[r] \ar[u] & R/I^2 \ar[u]
}
\quad
\xymatrix{
S \ar[r] \ar@{-->}[rd] & R/I^2 \\
R \ar[r] \ar[u] & R/I^3 \ar[u]
}
\quad
\xymatrix{
S \ar[r] \ar@{-->}[rd] & R/I^3 \\
R \ar[r] \ar[u] & R/I^4 \ar[u]
}
$$
Les morphismes correspondants $S/J^n \to R/I^n$ sont des isomorphismes
puisque les composées $S/J^n \to R/I^n \to S/J^n$ sont
(par récurrence) l'identité, par l'unicité dans la propriété de
relèvement des morphismes d'anneaux formellement \'etales.
\end{proof}

\begin{lemma}
\label{lemma-formally-etale-omega}
Soient $R \to S \to S'$ des morphismes d'anneaux. Soient $J$, resp.\ $J'$,
les noyaux des morphismes de multiplication
$S \otimes_R S \to S$, resp.\ $S' \otimes_{R'} S' \to S'$.
Si $S \to S'$ est formellement \'etale, alors le morphisme
$$
S' \otimes_S \left((S \otimes_R S)/J^{k + 1}\right)
\longrightarrow
(S' \otimes_R S')/(J')^{k + 1}
$$
est un isomorphisme pour tout $k \geq 0$. En particulier, le morphisme
$S' \otimes_S \Omega_{S/R} \to \Omega_{S'/R}$ est un isomorphisme.
\end{lemma}

\begin{proof}
Observons que $S' \otimes_S (S \otimes_R S) = S' \otimes_R S$
et que l'idéal $J$ engendre dans $S' \otimes_R S$ le noyau $I$ de la
surjection $S' \otimes_R S \to S'$. Par conséquent, le
membre de gauche est égal à $S' \otimes_R S/I^{k + 1}$.
Le morphisme $S' \otimes_R S \to S' \otimes_R S'$ est formellement \'etale
par le lemme \ref{lemma-base-change-formally-etale}. Nous concluons donc
que la flèche affichée dans l'énoncé du lemme est un isomorphisme
par le lemme \ref{lemma-formally-etale-lift-infinitesimal}.
La dernière assertion résulte de ceci et du fait que
$\Omega_{S/R} = J/J^2$ et $\Omega_{S'/R} = J'/(J')^2$ d'après le
lemme \ref{lemma-differentials-diagonal}.
\end{proof}

\begin{lemma}
\label{lemma-formally-etale-principal-parts}
Soient $R \to S \to S'$ des morphismes d'anneaux, avec $S \to S'$ formellement \'etale
(par exemple \'etale).
Soit $M$ un $S$-module. Posons $M' = S' \otimes_R M$. Alors nous avons
$$
S' \otimes_S P^k_{S/R}(M) = P^k_{S'/R}(M')
$$
Il s'ensuit que, pour tout $S$-module $N$ et tout
opérateur différentiel d'ordre fini $D : M \to N$,
il existe un unique prolongement $D' : M' \to S' \otimes_S N$
de $D$ en un opérateur différentiel (du même ordre ou d'ordre inférieur).
\end{lemma}

\begin{proof}
Soient $J$ et $J'$ comme dans l'énoncé du
lemme \ref{lemma-formally-etale-omega}. Alors nous avons
\begin{align*}
S' \otimes_S P^k_{S/R}(M)
& =
S' \otimes_S \left((S \otimes_R M)/J^{k + 1}(S \otimes_R M)\right) \\
& =
S' \otimes_S \left((S \otimes_R S)/J^{k + 1}\right) \otimes_S M \\
& =
\left(S' \otimes_R S'/(J')^{k + 1}\right) \otimes_S M \\
& =
\left(S' \otimes_R S'/(J')^{k + 1}\right) \otimes_{S'} M' \\
& =
S' \otimes_R M'/(J')^{k + 1}(S' \otimes_R M') \\
& =
P^k_{S'/R'}(M')
\end{align*}
La première et la dernière égalité proviennent du
lemme \ref{lemma-principal-parts-diagonal}.
La troisième égalité est le lemme \ref{lemma-formally-etale-omega}.
La dernière assertion est vraie parce que, si $D$ correspond à l'application
linéaire $\gamma : P^k_{S/R}(M) \to N$, alors nous pouvons faire
correspondre à $D' : M' \to S' \otimes_R N$ l'application linéaire
$1 \otimes \gamma : S' \otimes_R P^k_{S/R}(M) \to S' \otimes_R N$.
\end{proof}

\begin{remark}
\label{remark-formally-etale-differential-operators}
Soient $R \to S \to S'$ des morphismes d'anneaux, avec $S \to S'$ formellement \'etale
(par exemple \'etale). Soient $M_i$, $i = 1, 2, 3$, des $S$-modules,
et soient $D_i : M_i \to M_{i + 1}$, $i = 1, 2$, des opérateurs
différentiels d'ordre fini. Alors, si $D'_i : M'_i \to M'_{i + 1}$, $i = 1, 2$,
sont les prolongements des $D_i$ à $M'_i = S' \otimes_S M_i$ comme dans le
lemme \ref{lemma-formally-etale-principal-parts}, alors $D'_2 \circ D'_1$
est le prolongement de $D_2 \circ D_1$. En particulier, si $M$ est un $S$-module,
alors $M' = S' \otimes_S M$ est un module sur la $S$-algèbre
$\text{Diff}_{S/R}(M, M)$.
\end{remark}





\section{Morphismes d'anneaux non ramifiés}
\label{section-unramified}

\noindent
Notre définition d'un morphisme d'anneaux non ramifié est celle de \cite{Henselian}.
Ce que nous appelons morphisme d'anneaux G-non ramifié est ce qu'EGA appelle morphisme non ramifié.

\begin{definition}
\label{definition-unramified}
Soit $R \to S$ un morphisme d'anneaux.
\begin{enumerate}
\item On dit que $R \to S$ est {\it non ramifié} si $R \to S$ est de
type fini et $\Omega_{S/R} = 0$.
\item On dit que $R \to S$ est {\it G-non ramifié} si $R \to S$ est de
présentation finie et $\Omega_{S/R} = 0$.
\item Étant donné un idéal premier $\mathfrak q$ de $S$, on dit que $S$ est
{\it non ramifié en $\mathfrak q$} s'il existe un
$g \in S$, $g \not \in \mathfrak q$, tel que $R \to S_g$ soit non ramifié.
\item Étant donné un idéal premier $\mathfrak q$ de $S$, on dit que $S$ est
{\it G-non ramifié en $\mathfrak q$} s'il existe un
$g \in S$, $g \not \in \mathfrak q$, tel que $R \to S_g$ soit G-non ramifié.
\end{enumerate}
\end{definition}

\noindent
Bien sûr, un morphisme G-non ramifié est non ramifié.

\begin{lemma}
\label{lemma-formally-unramified-unramified}
Soit $R \to S$ un morphisme d'anneaux. Les assertions suivantes sont équivalentes :
\begin{enumerate}
\item $R \to S$ est formellement non ramifié et de type fini, et
\item $R \to S$ est non ramifié.
\end{enumerate}
De plus, les assertions suivantes sont également équivalentes :
\begin{enumerate}
\item $R \to S$ est formellement non ramifié et de présentation finie, et
\item $R \to S$ est G-non ramifié.
\end{enumerate}
\end{lemma}

\begin{proof}
Cela résulte du lemme \ref{lemma-characterize-formally-unramified}
et des définitions.
\end{proof}

\begin{lemma}
\label{lemma-unramified}
Propriétés des morphismes d'anneaux non ramifiés et G-non ramifiés.
\begin{enumerate}
\item Le changement de base d'un morphisme d'anneaux non ramifié est non ramifié.
Le changement de base d'un morphisme d'anneaux G-non ramifié est G-non ramifié.
\item La composée de morphismes d'anneaux non ramifiés est non ramifiée.
La composée de morphismes d'anneaux G-non ramifiés est G-non ramifiée.
\item Toute localisation principale $R \to R_f$ est G-non ramifiée et
non ramifiée.
\item Si $I \subset R$ est un idéal, alors $R \to R/I$ est non ramifié.
Si $I \subset R$ est un idéal de type fini, alors $R \to R/I$ est
G-non ramifié.
\item Un morphisme d'anneaux \'etale est G-non ramifié et non ramifié.
\item Si $R \to S$ est de type fini (resp.\ de présentation finie),
si $\mathfrak q \subset S$ est un idéal premier et si $(\Omega_{S/R})_{\mathfrak q} = 0$,
alors $R \to S$ est non ramifié (resp.\ G-non ramifié) en $\mathfrak q$.
\item Si $R \to S$ est de type fini (resp.\ de présentation finie),
si $\mathfrak q \subset S$ est un idéal premier et si
$\Omega_{S/R} \otimes_S \kappa(\mathfrak q) = 0$, alors
$R \to S$ est non ramifié (resp.\ G-non ramifié) en $\mathfrak q$.
\item Si $R \to S$ est de type fini (resp.\ de présentation finie),
si $\mathfrak q \subset S$ est un idéal premier au-dessus de $\mathfrak p \subset R$ et si
$(\Omega_{S \otimes_R \kappa(\mathfrak p)/\kappa(\mathfrak p)})_{\mathfrak q}
= 0$, alors $R \to S$ est non ramifié (resp.\ G-non ramifié) en $\mathfrak q$.
\item Si $R \to S$ est de type fini (resp.\ de présentation finie),
si $\mathfrak q \subset S$ est un idéal premier au-dessus de $\mathfrak p \subset R$ et si
$(\Omega_{S \otimes_R \kappa(\mathfrak p)/\kappa(\mathfrak p)})
\otimes_{S \otimes_R \kappa(\mathfrak p)} \kappa(\mathfrak q) = 0$,
alors $R \to S$ est non ramifié (resp.\ G-non ramifié) en $\mathfrak q$.
\item Si $R \to S$ est un morphisme d'anneaux, si $g_1, \ldots, g_m \in S$ engendrent
l'idéal unité et si $R \to S_{g_j}$ est non ramifié (resp.\ G-non ramifié) pour
$j = 1, \ldots, m$, alors $R \to S$ est non ramifié (resp.\ G-non ramifié).
\item Si $R \to S$ est un morphisme d'anneaux qui est non ramifié (resp.\ G-non ramifié)
en tout idéal premier de $S$, alors $R \to S$ est non ramifié (resp.\ G-non ramifié).
\item Si $R \to S$ est G-non ramifié, alors il existe une
$\mathbf{Z}$-algèbre $R_0$ de type fini, un morphisme d'anneaux G-non ramifié
$R_0 \to S_0$ et un morphisme d'anneaux $R_0 \to R$ tels que
$S = R \otimes_{R_0} S_0$.
\item Si $R \to S$ est non ramifié, alors il existe une
$\mathbf{Z}$-algèbre $R_0$ de type fini, un morphisme d'anneaux non ramifié
$R_0 \to S_0$ et un morphisme d'anneaux $R_0 \to R$ tels que $S$ soit un quotient de
$R \otimes_{R_0} S_0$.
\end{enumerate}
\end{lemma}

\begin{proof}
Nous démontrons chaque point dans l'ordre.

\medskip\noindent
Pour (1). Cela résulte des lemmes \ref{lemma-differentials-base-change}
et \ref{lemma-base-change-finiteness}.

\medskip\noindent
Pour (2). Cela résulte des lemmes \ref{lemma-exact-sequence-differentials}
et \ref{lemma-base-change-finiteness}.

\medskip\noindent
Pour (3). Cela résulte d'un calcul direct de $\Omega_{R_f/R}$, que nous omettons.

\medskip\noindent
Pour (4). Nous avons $\Omega_{(R/I)/R} = 0$, voir le
lemme \ref{lemma-trivial-differential-surjective},
et le morphisme d'anneaux $R \to R/I$
est de type fini. Si $I$ est un idéal de type fini, alors $R \to R/I$
est de présentation finie.

\medskip\noindent
Pour (5). Voir la discussion suivant la définition \ref{definition-etale}.

\medskip\noindent
Pour (6). Dans ce cas, $\Omega_{S/R}$ est un $S$-module de type fini (voir le
lemme \ref{lemma-differentials-finitely-generated}) et il existe donc
un $g \in S$, $g \not \in \mathfrak q$, tel que
$(\Omega_{S/R})_g = 0$. D'après le lemme \ref{lemma-differentials-localize},
cela signifie que $\Omega_{S_g/R} = 0$ et donc que $R \to S_g$ est
non ramifié, comme désiré.

\medskip\noindent
Pour (7). Le lemme de Nakayama (lemme \ref{lemma-NAK}) montre que
la condition est équivalente à celle de (6).

\medskip\noindent
Pour (8) et (9). Ces assertions sont équivalentes de la même manière que (6) et (7)
le sont. De plus,
$\Omega_{S \otimes_R \kappa(\mathfrak p)/\kappa(\mathfrak p)} =
\Omega_{S/R} \otimes_S (S \otimes_R \kappa(\mathfrak p))$ d'après le
lemme \ref{lemma-differentials-base-change}.
Nous voyons donc que (9) est équivalente à (7), puisque les
$\kappa(\mathfrak q)$-espaces vectoriels qui y figurent sont canoniquement
isomorphes.

\medskip\noindent
Pour (10). Cela résulte des lemmes \ref{lemma-cover}
et \ref{lemma-differentials-localize}.

\medskip\noindent
Pour (11). Cela résulte de (6) et (7) et du fait que le spectre de $S$
est quasi-compact.

\medskip\noindent
Pour (12). Écrivons $S = R[x_1, \ldots, x_n]/(g_1, \ldots, g_m)$.
Comme $\Omega_{S/R} = 0$, nous pouvons écrire
$$
\text{d}x_i = \sum h_{ij}\text{d}g_j + \sum a_{ijk}g_j\text{d}x_k
$$
dans $\Omega_{R[x_1, \ldots, x_n]/R}$
pour certains $h_{ij}, a_{ijk} \in R[x_1, \ldots, x_n]$.
Choisissons une
$\mathbf{Z}$-sous-algèbre $R_0 \subset R$ de type fini contenant tous les coefficients des
polynômes $g_i, h_{ij}, a_{ijk}$. Posons
$S_0 = R_0[x_1, \ldots, x_n]/(g_1, \ldots, g_m)$. Cela convient.

\medskip\noindent
Pour (13). Écrivons $S = R[x_1, \ldots, x_n]/I$.
Comme $\Omega_{S/R} = 0$, nous pouvons écrire
$$
\text{d}x_i = \sum h_{ij}\text{d}g_{ij} + \sum g'_{ik}\text{d}x_k
$$
dans $\Omega_{R[x_1, \ldots, x_n]/R}$
pour certains $h_{ij} \in R[x_1, \ldots, x_n]$ et $g_{ij}, g'_{ik} \in I$.
Choisissons une $\mathbf{Z}$-sous-algèbre $R_0 \subset R$ de type fini
contenant tous les coefficients des
polynômes $g_{ij}, h_{ij}, g'_{ik}$. Posons
$S_0 = R_0[x_1, \ldots, x_n]/(g_{ij}, g'_{ik})$. Cela convient.
\end{proof}

\begin{lemma}
\label{lemma-diagonal-unramified}
Soit $R \to S$ un morphisme d'anneaux.
Si $R \to S$ est non ramifié, alors il existe un idempotent
$e \in S \otimes_R S$ tel que $S \otimes_R S \to S$ soit isomorphe à
$S \otimes_R S \to (S \otimes_R S)_e$.
\end{lemma}

\begin{proof}
Soit $J = \Ker(S \otimes_R S \to S)$. Par hypothèse,
$J/J^2 = 0$, voir le
lemme \ref{lemma-differentials-diagonal}.
Comme $S$ est de type fini sur $R$, nous
voyons que $J$ est de type fini, à savoir engendré par les
$x_i \otimes 1 - 1 \otimes x_i$, où les $x_i$ engendrent $S$ sur $R$.
On conclut par le lemme \ref{lemma-ideal-is-squared-union-connected}.
\end{proof}

\begin{lemma}
\label{lemma-unramified-at-prime}
Soit $R \to S$ un morphisme d'anneaux.
Soit $\mathfrak q \subset S$ un
idéal premier au-dessus de $\mathfrak p$ dans $R$.
Si $S/R$ est non ramifié en $\mathfrak q$, alors
\begin{enumerate}
\item nous avons $\mathfrak p S_{\mathfrak q} = \mathfrak qS_{\mathfrak q}$,
et cet idéal est l'idéal maximal de l'anneau local $S_{\mathfrak q}$, et
\item l'extension de corps $\kappa(\mathfrak q)/\kappa(\mathfrak p)$
est finie séparable.
\end{enumerate}
\end{lemma}

\begin{proof}
Nous pouvons d'abord remplacer $S$ par $S_g$ pour un $g \in S$, $g \not \in \mathfrak q$,
et supposer que $R \to S$ est non ramifié.
Le changement de base $S \otimes_R \kappa(\mathfrak p)$
est non ramifié sur $\kappa(\mathfrak p)$ d'après le
lemme \ref{lemma-unramified}.
D'après le
lemme \ref{lemma-characterize-smooth-over-field},
il est lisse, donc \'etale, sur $\kappa(\mathfrak p)$.
Ainsi, $F = S \otimes_R \kappa(\mathfrak p)$ est un produit fini
d'extensions de corps finies séparables de $\kappa(\mathfrak p)$, voir le
lemme \ref{lemma-etale-over-field}. Avec les notations et résultats de la
remarque \ref{remark-local-ring-fibre}, nous trouvons que
$F_{\overline{\mathfrak q}} = S_\mathfrak q/\mathfrak pS_\mathfrak q$
est égal à $\kappa(\mathfrak q)$. Cela implique le lemme.
\end{proof}

\begin{lemma}
\label{lemma-unramified-quasi-finite}
Soit $R \to S$ un morphisme d'anneaux de type fini.
Soit $\mathfrak q$ un idéal premier de $S$.
Si $R \to S$ est non ramifié en $\mathfrak q$, alors
$R \to S$ est quasi-fini en $\mathfrak q$.
En particulier, un morphisme d'anneaux non ramifié est quasi-fini.
\end{lemma}

\begin{proof}
Un morphisme d'anneaux non ramifié est de type fini.
Il est donc clair que la seconde assertion résulte de la première.
Pour voir la première assertion, appliquons la caractérisation du
lemme \ref{lemma-isolated-point-fibre}, partie (2), en utilisant le
lemme \ref{lemma-unramified-at-prime}.
\end{proof}

\begin{lemma}
\label{lemma-characterize-unramified}
Soit $R \to S$ un morphisme d'anneaux. Soit $\mathfrak q$ un idéal premier de $S$
au-dessus d'un idéal premier $\mathfrak p$ de $R$. Si
\begin{enumerate}
\item $R \to S$ est de type fini,
\item $\mathfrak p S_{\mathfrak q}$ est l'idéal maximal
de l'anneau local $S_{\mathfrak q}$, et
\item l'extension de corps $\kappa(\mathfrak q)/\kappa(\mathfrak p)$
est finie séparable,
\end{enumerate}
alors $R \to S$ est non ramifié en $\mathfrak q$.
\end{lemma}

\begin{proof}
D'après le lemme \ref{lemma-unramified} (8), il suffit de montrer que
$\Omega_{S \otimes_R \kappa(\mathfrak p) / \kappa(\mathfrak p)}$
est nul après localisation en $\mathfrak q$. Nous pouvons donc remplacer $S$
par $S \otimes_R \kappa(\mathfrak p)$ et $R$ par $\kappa(\mathfrak p)$.
Autrement dit, nous pouvons supposer que $R = k$ est un corps et que $S$
est une $k$-algèbre de type fini.
Dans ce cas, les hypothèses impliquent que
$S_{\mathfrak q} \cong \kappa(\mathfrak q)$.
Ainsi, $(\Omega_{S/k})_{\mathfrak q} = \Omega_{S_\mathfrak q/k} =
\Omega_{\kappa(\mathfrak q)/k}$ est nul, comme désiré (la
première égalité est le lemme \ref{lemma-differentials-localize}).
\end{proof}

\begin{lemma}
\label{lemma-etale-flat-unramified-finite-presentation}
Soit $R \to S$ un morphisme d'anneaux. Les assertions suivantes sont équivalentes :
\begin{enumerate}
\item $R \to S$ est \'etale,
\item $R \to S$ est plat et G-non ramifié, et
\item $R \to S$ est plat, non ramifié et de présentation finie.
\end{enumerate}
\end{lemma}

\begin{proof}
Les assertions (2) et (3) sont équivalentes par définition.
L'implication (1) $\Rightarrow$ (3) résulte du
fait que les morphismes d'anneaux \'etales sont de présentation finie,
du lemme \ref{lemma-etale} (platitude des morphismes \'etales) et du
lemme \ref{lemma-unramified} (les morphismes \'etales sont non ramifiés).
Réciproquement, la caractérisation des morphismes d'anneaux \'etales du
lemme \ref{lemma-characterize-etale}
et la structure des morphismes d'anneaux non ramifiés du
lemme \ref{lemma-unramified-at-prime}
montrent que (3) implique (1). (On utilise ici que $R \to S$
est \'etale si $R \to S$ est \'etale en tout idéal premier $\mathfrak q \subset S$,
voir le lemme \ref{lemma-etale}.)
\end{proof}

\begin{lemma}
\label{lemma-characterize-etale-over-polynomial-ring}
Soit $k$ un corps. Soit
$$
\varphi : k[x_1, \ldots, x_n] \to A, \quad x_i \longmapsto a_i
$$
un morphisme d'anneaux de type fini. Alors $\varphi$ est \'etale si et seulement si les
deux conditions suivantes sont vérifiées : (a) les anneaux locaux de $A$ aux idéaux maximaux
sont de dimension $n$, et (b) les éléments $\text{d}(a_1), \ldots, \text{d}(a_n)$
engendrent $\Omega_{A/k}$ comme $A$-module.
\end{lemma}

\begin{proof}
Supposons (a) et (b). La condition (b) implique que
$\Omega_{A/k[x_1, \ldots, x_n]} = 0$ et donc que $\varphi$ est non ramifié.
Il suffit ainsi de prouver que $\varphi$ est plat, voir le
lemme \ref{lemma-etale-flat-unramified-finite-presentation}.
Soit $\mathfrak m \subset A$ un idéal maximal.
Posons $X = \Spec(A)$ et notons $x \in X$ le point fermé correspondant
à $\mathfrak m$. Alors $\dim(A_\mathfrak m)$ est $\dim_x X$, voir le
lemme \ref{lemma-dimension-closed-point-finite-type-field}.
Ainsi, d'après le lemme \ref{lemma-characterize-smooth-over-field},
si (a) et (b) sont vérifiées, alors $A_\mathfrak m$ est
un anneau local régulier pour tout idéal maximal $\mathfrak m$. Ensuite,
$k[x_1, \ldots, x_n]_{\varphi^{-1}(\mathfrak m)} \to A_\mathfrak m$
est plat d'après le lemme \ref{lemma-CM-over-regular-flat}
(et le fait qu'un anneau local régulier est de Cohen-Macaulay, voir le
lemme \ref{lemma-regular-ring-CM}).
Ainsi, $\varphi$ est plat d'après le lemme \ref{lemma-flat-localization}.

\medskip\noindent
Supposons que $\varphi$ soit \'etale. Alors $\Omega_{A/k[x_1, \ldots, x_n]} = 0$
et donc (b) est vérifiée. D'autre part, les morphismes d'anneaux \'etales sont plats
(lemme \ref{lemma-etale}) et quasi-finis
(lemme \ref{lemma-etale-quasi-finite}).
Ainsi, pour tout idéal maximal $\mathfrak m$ de $A$, nous pouvons appliquer le
lemme \ref{lemma-dimension-base-fibre-equals-total} à
$k[x_1, \ldots, x_n]_{\varphi^{-1}(\mathfrak m)} \to A_\mathfrak m$
pour voir que $\dim(A_\mathfrak m) = n$ et donc que (a) est vérifiée.
\end{proof}




\section{Structure locale des morphismes d'anneaux non ramifiés}
\label{section-local-structure-unramified}

\noindent
Un morphisme non ramifié est localement (en un sens convenable) la composée
d'une immersion fermée et d'un morphisme \'etale. Les fondements algébriques
de ce fait sont examinés dans cette section.

\begin{proposition}
\label{proposition-unramified-locally-standard}
Soit $R \to S$ un morphisme d'anneaux. Soit $\mathfrak q \subset S$ un idéal premier.
Si $R \to S$ est non ramifié en $\mathfrak q$, alors il existe
\begin{enumerate}
\item un $g \in S$, $g \not \in \mathfrak q$,
\item un morphisme d'anneaux \'etale standard $R \to S'$, et
\item un morphisme surjectif de $R$-algèbres $S' \to S_g$.
\end{enumerate}
\end{proposition}

\begin{proof}
Cette démonstration est la « même » que celle de la
proposition \ref{proposition-etale-locally-standard}.
Elle est quelque peu indirecte et il existe peut-être des moyens de
la raccourcir.

\medskip\noindent
Étape 1. D'après la définition \ref{definition-unramified},
il existe un $g \in S$, $g \not \in \mathfrak q$,
tel que $R \to S_g$ soit non ramifié. Nous pouvons donc supposer que $S$ est
non ramifié sur $R$.

\medskip\noindent
Étape 2. D'après le lemme \ref{lemma-unramified},
il existe un morphisme d'anneaux non ramifié $R_0 \to S_0$,
avec $R_0$ de type fini sur $\mathbf{Z}$, et un morphisme d'anneaux
$R_0 \to R$ tel que $S$ soit un quotient de $R \otimes_{R_0} S_0$. Notons
$\mathfrak q_0$ l'idéal premier de $S_0$ correspondant à $\mathfrak q$.
Si nous montrons le résultat pour $(R_0 \to S_0, \mathfrak q_0)$, alors le
résultat pour $(R \to S, \mathfrak q)$ s'en déduit par changement de base. Ainsi,
nous pouvons supposer que $R$ est noethérien.

\medskip\noindent
Étape 3.
Notons que $R \to S$ est quasi-fini d'après le
lemme \ref{lemma-unramified-quasi-finite}.
D'après le lemme \ref{lemma-quasi-finite-open-integral-closure},
il existe un morphisme d'anneaux fini $R \to S'$, un morphisme de $R$-algèbres
$S' \to S$ et un élément $g' \in S'$ satisfaisant
$g' \not \in \mathfrak q$ et tel que $S' \to S$ induise
un isomorphisme $S'_{g'} \cong S_{g'}$.
(Notons que $S'$ n'est pas nécessairement non ramifié sur $R$.)
Nous pouvons donc supposer que (a) $R$ est noethérien, (b) $R \to S$ est fini
et (c) $R \to S$ est non ramifié en $\mathfrak q$
(mais plus nécessairement non ramifié en tous les idéaux premiers).

\medskip\noindent
Étape 4. Soit $\mathfrak p \subset R$ l'idéal premier correspondant
à $\mathfrak q$. Considérons l'anneau fibre
$S \otimes_R \kappa(\mathfrak p)$. C'est une algèbre finie sur
$\kappa(\mathfrak p)$. Il est donc artinien
(voir le lemme \ref{lemma-finite-dimensional-algebra}) et
est ainsi un produit fini d'anneaux locaux
$$
S \otimes_R \kappa(\mathfrak p) = \prod\nolimits_{i = 1}^n A_i
$$
voir la proposition \ref{proposition-dimension-zero-ring}. L'un des facteurs,
disons $A_1$, est l'anneau local $S_{\mathfrak q}/\mathfrak pS_{\mathfrak q}$,
qui est isomorphe à $\kappa(\mathfrak q)$,
voir le lemme \ref{lemma-unramified-at-prime}. Les autres facteurs correspondent aux
autres idéaux premiers, disons $\mathfrak q_2, \ldots, \mathfrak q_n$, de
$S$ au-dessus de $\mathfrak p$.

\medskip\noindent
Étape 5. Nous pouvons choisir un élément non nul $\alpha \in \kappa(\mathfrak q)$ qui
engendre l'extension de corps finie séparable
$\kappa(\mathfrak q)/\kappa(\mathfrak p)$ (ainsi, même si
l'extension de corps est triviale, nous n'autorisons pas $\alpha = 0$).
Notons que, pour tout $\lambda \in \kappa(\mathfrak p)^*$, l'élément
$\lambda \alpha$ engendre lui aussi $\kappa(\mathfrak q)$
sur $\kappa(\mathfrak p)$. Considérons l'élément
$$
\overline{t} =
(\alpha, 0, \ldots, 0) \in
\prod\nolimits_{i = 1}^n A_i =
S \otimes_R \kappa(\mathfrak p).
$$
Après avoir éventuellement remplacé $\alpha$ par $\lambda \alpha$ comme ci-dessus,
nous pouvons supposer que $\overline{t}$ est l'image d'un $t \in S$.
Soit $I \subset R[x]$ le noyau du morphisme de $R$-algèbres
$R[x] \to S$ qui envoie $x$ sur $t$. Posons $S' = R[x]/I$,
de sorte que $S' \subset S$. Voici un diagramme :
$$
\xymatrix{
R[x] \ar[r] & S' \ar[r] & S \\
R \ar[u] \ar[ru] \ar[rru] & &
}
$$
Par construction, les idéaux premiers $\mathfrak q_j$, $j \geq 2$, de $S$ se
trouvent tous au-dessus de l'idéal premier $(\mathfrak p, x)$ de $R[x]$, tandis que
l'idéal premier $\mathfrak q$ se trouve au-dessus d'un autre idéal premier de $R[x]$,
car $\alpha \not = 0$.

\medskip\noindent
Étape 6. Notons $\mathfrak q' \subset S'$ l'idéal premier de $S'$
correspondant à $\mathfrak q$. D'après ce qui précède, $\mathfrak q$ est
le seul idéal premier de $S$ au-dessus de $\mathfrak q'$. Nous voyons donc que
$S_{\mathfrak q} = S_{\mathfrak q'}$, voir le
lemme \ref{lemma-unique-prime-over-localize-below} (nous disposons de la
montée pour $S' \to S$ d'après le lemme \ref{lemma-integral-going-up},
puisque $S' \to S$ est fini car $R \to S$ est fini).
Il s'ensuit que $S'_{\mathfrak q'} \to S_{\mathfrak q}$ est fini
et injectif en tant que localisation du morphisme d'anneaux fini injectif
$S' \to S$. Considérons les morphismes d'anneaux locaux
$$
R_{\mathfrak p} \to S'_{\mathfrak q'} \to S_{\mathfrak q}
$$
Le second morphisme est fini et injectif. Nous avons
$S_{\mathfrak q}/\mathfrak pS_{\mathfrak q} = \kappa(\mathfrak q)$,
voir le lemme \ref{lemma-unramified-at-prime}.
Donc, a fortiori,
$S_{\mathfrak q}/\mathfrak q'S_{\mathfrak q} = \kappa(\mathfrak q)$.
Puisque
$$
\kappa(\mathfrak p) \subset \kappa(\mathfrak q') \subset \kappa(\mathfrak q)
$$
et puisque $\alpha$ appartient à l'image de $\kappa(\mathfrak q')$ dans
$\kappa(\mathfrak q)$,
nous concluons que $\kappa(\mathfrak q') = \kappa(\mathfrak q)$.
Ainsi, d'après le lemme de Nakayama \ref{lemma-NAK} appliqué au morphisme de
$S'_{\mathfrak q'}$-modules $S'_{\mathfrak q'} \to S_{\mathfrak q}$,
le morphisme $S'_{\mathfrak q'} \to S_{\mathfrak q}$ est surjectif.
Autrement dit,
$S'_{\mathfrak q'} \cong S_{\mathfrak q}$.

\medskip\noindent
Étape 7. D'après le lemme \ref{lemma-isomorphic-local-rings}, il existe
$g \in S$, $g \not \in \mathfrak q$, et
$g' \in S'$, $g' \not \in \mathfrak q'$, tels que $S'_{g'} \cong S_g$.
Comme $R$ est noethérien, l'anneau $S'$ est fini sur $R$,
car c'est un $R$-sous-module
du $R$-module fini $S$. Ainsi, après avoir remplacé $S$ par $S'$, nous pouvons
supposer que (a) $R$ est noethérien, (b) $S$ est fini sur $R$, (c)
$S$ est non ramifié sur $R$ en $\mathfrak q$, et (d) $S = R[x]/I$.

\medskip\noindent
Étape 8. Considérons l'anneau
$S \otimes_R \kappa(\mathfrak p) = \kappa(\mathfrak p)[x]/\overline{I}$,
où $\overline{I} = I \cdot \kappa(\mathfrak p)[x]$ est l'idéal engendré
par $I$ dans $\kappa(\mathfrak p)[x]$. Comme $\kappa(\mathfrak p)[x]$ est un anneau principal,
nous savons que $\overline{I} = (\overline{h})$ pour un certain polynôme unitaire
$\overline{h} \in \kappa(\mathfrak p)$. Après avoir remplacé $\overline{h}$
par $\lambda \cdot \overline{h}$ pour un certain $\lambda \in \kappa(\mathfrak p)$,
nous pouvons supposer que $\overline{h}$ est l'image d'un $h \in R[x]$.
(Le problème est que nous ne savons pas si nous pouvons choisir $h$ unitaire.)
De plus, comme à l'étape 4, nous savons que
$S \otimes_R \kappa(\mathfrak p) = A_1 \times \ldots \times A_n$, où
$A_1 = \kappa(\mathfrak q)$ est une extension finie séparable de
$\kappa(\mathfrak p)$ et $A_2, \ldots, A_n$ sont locaux. Cela implique
que
$$
\overline{h} = \overline{h}_1 \overline{h}_2^{e_2} \ldots \overline{h}_n^{e_n}
$$
pour certains polynômes unitaires irréductibles deux à deux premiers entre eux
$\overline{h}_i \in \kappa(\mathfrak p)[x]$ et certains entiers
$e_2, \ldots, e_n \geq 1$. La numérotation est choisie de sorte que
$A_i = \kappa(\mathfrak p)[x]/(\overline{h}_i^{e_i})$ comme
$\kappa(\mathfrak p)[x]$-algèbres. Notons que $\overline{h}_1$ est
le polynôme minimal de $\alpha \in \kappa(\mathfrak q)$ et est donc
un polynôme séparable (sa dérivée est première avec lui-même).

\medskip\noindent
Étape 9. Soit $m \in I$ un élément unitaire ; un tel élément existe
parce que l'extension d'anneaux $R \to R[x]/I$ est finie, donc entière.
Notons $\overline{m}$ son image dans $\kappa(\mathfrak p)[x]$.
Nous pouvons factoriser
$$
\overline{m} = \overline{k}
\overline{h}_1^{d_1} \overline{h}_2^{d_2} \ldots \overline{h}_n^{d_n}
$$
pour un certain $d_1 \geq 1$, des $d_j \geq e_j$, $j = 2, \ldots, n$, et
un $\overline{k} \in \kappa(\mathfrak p)[x]$ premier avec tous les $\overline{h}_i$.
Posons $f = m^l + h$, où $l \deg(m) > \deg(h)$ et $l \geq 2$.
Alors $f$ est unitaire comme polynôme sur $R$. De plus, l'image $\overline{f}$
de $f$ dans $\kappa(\mathfrak p)[x]$ se factorise comme
$$
\overline{f} =
\overline{h}_1 \overline{h}_2^{e_2} \ldots \overline{h}_n^{e_n}
+
\overline{k}^l \overline{h}_1^{ld_1} \overline{h}_2^{ld_2}
\ldots \overline{h}_n^{ld_n}
=
\overline{h}_1(\overline{h}_2^{e_2} \ldots \overline{h}_n^{e_n}
+
\overline{k}^l
\overline{h}_1^{ld_1 - 1} \overline{h}_2^{ld_2} \ldots \overline{h}_n^{ld_n})
= \overline{h}_1 \overline{w}
$$
avec $\overline{w}$ un polynôme premier avec $\overline{h}_1$.
Posons $g = f'$ (la dérivée par rapport à $x$).

\medskip\noindent
Étape 10. Le morphisme d'anneaux $R[x] \to S = R[x]/I$ a les propriétés suivantes :
(1) il envoie $f$ sur zéro, et
(2) il envoie $g$ sur un élément de $S \setminus \mathfrak q$.
La première assertion est claire puisque $f$ appartient à $I$.
Pour la seconde assertion, il nous suffit de montrer que $g$ ne
s'envoie pas sur zéro dans
$\kappa(\mathfrak q) = \kappa(\mathfrak p)[x]/(\overline{h}_1)$.
L'image de $g$ dans $\kappa(\mathfrak p)[x]$ est la dérivée
de $\overline{f}$. Ainsi, (2) est claire parce que
$$
\overline{g} =
\frac{\text{d}\overline{f}}{\text{d}x} =
\overline{w}\frac{\text{d}\overline{h}_1}{\text{d}x} +
\overline{h}_1\frac{\text{d}\overline{w}}{\text{d}x},
$$
$\overline{w}$ est premier avec $\overline{h}_1$ et
$\overline{h}_1$ est séparable.

\medskip\noindent
Étape 11.
Nous concluons que $\varphi : R[x]/(f) \to S$ est un morphisme d'anneaux surjectif,
que $R[x]_g/(f)$ est \'etale sur $R$ (parce qu'il est \'etale standard,
voir le lemme \ref{lemma-standard-etale}) et que $\varphi(g) \not \in \mathfrak q$.
Ainsi, le morphisme $(R[x]/(f))_g \to S_{\varphi(g)}$ est la
surjection recherchée.
\end{proof}



\begin{lemma}
\label{lemma-etale-makes-unramified-closed-at-prime}
Soit $R \to S$ un morphisme d'anneaux.
Soit $\mathfrak q$ un idéal premier de $S$ au-dessus de $\mathfrak p \subset R$.
Supposons que $R \to S$ soit de type fini et non ramifié en $\mathfrak q$.
Alors il existe
\begin{enumerate}
\item un morphisme d'anneaux \'etale $R \to R'$,
\item un idéal premier $\mathfrak p' \subset R'$ au-dessus de $\mathfrak p$,
\item une décomposition en produit
$$
R' \otimes_R S = A \times B
$$
\end{enumerate}
ayant les propriétés suivantes :
\begin{enumerate}
\item $R' \to A$ est surjectif, et
\item $\mathfrak p'A$ est un idéal premier de $A$ au-dessus de $\mathfrak p'$ et
de $\mathfrak q$.
\end{enumerate}
\end{lemma}

\begin{proof}
Nous pouvons remplacer $(R \to S, \mathfrak p, \mathfrak q)$
par tout changement de base $(R' \to R'\otimes_R S, \mathfrak p', \mathfrak q')$
par un morphisme d'anneaux \'etale $R \to R'$, avec un idéal premier $\mathfrak p'$
au-dessus de $\mathfrak p$ et un choix de $\mathfrak q'$ au-dessus
à la fois de $\mathfrak q$ et de $\mathfrak p'$. Notons aussi qu'étant donnés
$R \to R'$ et $\mathfrak p'$, on peut toujours trouver un $\mathfrak q'$ convenable.

\medskip\noindent
L'hypothèse que $R \to S$ est de type fini signifie que nous pouvons appliquer le
lemme \ref{lemma-etale-makes-quasi-finite-finite-variant}. Nous pouvons donc
supposer que $S = A_1 \times \ldots \times A_n \times B$, que
chaque $R \to A_i$ est fini avec exactement un idéal premier $\mathfrak r_i$
au-dessus de $\mathfrak p$ tel que
$\kappa(\mathfrak p) \subset \kappa(\mathfrak r_i)$ soit purement inséparable,
et que $R \to B$ ne soit quasi-fini en aucun idéal premier au-dessus de $\mathfrak p$.
Alors, clairement, $\mathfrak q = \mathfrak r_i$ pour un certain $i$, puisqu'un
morphisme non ramifié est quasi-fini
(voir le lemme \ref{lemma-unramified-quasi-finite}).
Disons que $\mathfrak q = \mathfrak r_1$.
D'après le lemme \ref{lemma-unramified-at-prime}, nous voyons que
$\kappa(\mathfrak r_1)/\kappa(\mathfrak p)$
est séparable, donc est l'extension de corps triviale, et que
$\mathfrak p(A_1)_{\mathfrak r_1}$ est l'idéal maximal.
De plus, d'après le lemme \ref{lemma-unique-prime-over-localize-below}
(qui s'applique à $R \to A_1$ parce qu'un morphisme d'anneaux fini satisfait la montée d'après le
lemme \ref{lemma-integral-going-up}),
nous avons $(A_1)_{\mathfrak r_1} = (A_1)_{\mathfrak p}$.
Il résulte du lemme de Nakayama \ref{lemma-NAK}
que le morphisme d'anneaux locaux
$R_{\mathfrak p} \to (A_1)_{\mathfrak p} = (A_1)_{\mathfrak r_1}$
est surjectif. Comme $A_1$ est fini sur $R$, nous voyons qu'il
existe un $f \in R$, $f \not \in \mathfrak p$, tel que
$R_f \to (A_1)_f$ soit surjectif. Après avoir remplacé $R$ par $R_f$, on conclut.
\end{proof}

\begin{lemma}
\label{lemma-etale-makes-unramified-closed}
\begin{slogan}
Dans un morphisme d'anneaux non ramifié, on peut séparer les points d'une fibre
en passant à un voisinage \'etale.
\end{slogan}
Soit $R \to S$ un morphisme d'anneaux.
Soit $\mathfrak p$ un idéal premier de $R$.
Si $R \to S$ est non ramifié, alors il existe
\begin{enumerate}
\item un morphisme d'anneaux \'etale $R \to R'$,
\item un idéal premier $\mathfrak p' \subset R'$ au-dessus de $\mathfrak p$,
\item une décomposition en produit
$$
R' \otimes_R S = A_1 \times \ldots \times A_n \times B
$$
\end{enumerate}
ayant les propriétés suivantes :
\begin{enumerate}
\item $R' \to A_i$ est surjectif,
\item $\mathfrak p'A_i$ est un idéal premier de $A_i$ au-dessus de $\mathfrak p'$, et
\item il n'existe aucun idéal premier de $B$ au-dessus de $\mathfrak p'$.
\end{enumerate}
\end{lemma}

\begin{proof}
Nous pouvons appliquer le lemme \ref{lemma-etale-makes-quasi-finite-finite-variant}.
Ainsi, après un changement de base \'etale,
nous pouvons supposer que $S = A_1 \times \ldots \times A_n \times B$,
que chaque $R \to A_i$ est fini avec exactement un idéal premier $\mathfrak r_i$
au-dessus de $\mathfrak p$ tel que
$\kappa(\mathfrak p) \subset \kappa(\mathfrak r_i)$ soit purement inséparable,
et que $R \to B$ ne soit quasi-fini en aucun idéal premier au-dessus de $\mathfrak p$.
Comme $R \to S$ est quasi-fini (voir le
lemme \ref{lemma-unramified-quasi-finite}),
nous voyons qu'il n'existe aucun idéal premier de $B$ au-dessus de $\mathfrak p$.
D'après le lemme \ref{lemma-unramified-at-prime}, nous voyons que
$\kappa(\mathfrak r_i)/\kappa(\mathfrak p)$
est séparable, donc est l'extension de corps triviale, et que
$\mathfrak p(A_i)_{\mathfrak r_i}$ est l'idéal maximal.
De plus, d'après le lemme \ref{lemma-unique-prime-over-localize-below}
(qui s'applique à $R \to A_i$ parce qu'un morphisme d'anneaux fini satisfait la montée d'après le
lemme \ref{lemma-integral-going-up}),
nous avons $(A_i)_{\mathfrak r_i} = (A_i)_{\mathfrak p}$.
Il résulte du lemme de Nakayama \ref{lemma-NAK}
que le morphisme d'anneaux locaux
$R_{\mathfrak p} \to (A_i)_{\mathfrak p} = (A_i)_{\mathfrak r_i}$
est surjectif. Comme $A_i$ est fini sur $R$, nous voyons qu'il
existe un $f \in R$, $f \not \in \mathfrak p$, tel que
$R_f \to (A_i)_f$ soit surjectif. Après avoir remplacé $R$ par $R_f$, on conclut.
\end{proof}








\section{Anneaux locaux henséliens}
\label{section-henselian}

\noindent
Dans cette section, nous examinons brièvement la notion d'anneau local hensélien.
Soit $(R, \mathfrak m, \kappa)$ un anneau local.
Pour $a \in R$, nous notons $\overline{a}$ l'image de $a$ dans $\kappa$.
Pour un polynôme $f \in R[T]$, nous notons souvent $\overline{f}$
l'image de $f$ dans $\kappa[T]$.
Étant donné un polynôme $f \in R[T]$, nous notons $f'$ la dérivée
de $f$ par rapport à $T$. Notons que $\overline{f}' = \overline{f'}$.

\begin{definition}
\label{definition-henselian}
Soit $(R, \mathfrak m, \kappa)$ un anneau local.
\begin{enumerate}
\item On dit que $R$ est {\it hensélien} si, pour tout $f \in R[T]$ unitaire et
toute racine $a_0 \in \kappa$ de $\overline{f}$ telle que
$\overline{f'}(a_0) \not = 0$,
il existe un $a \in R$ tel que $f(a) = 0$ et
$a_0 = \overline{a}$.
\item On dit que $R$ est {\it strictement hensélien} si $R$ est hensélien
et si son corps résiduel est séparablement clos.
\end{enumerate}
\end{definition}

\noindent
Notons que la condition $\overline{f'}(a_0) \not = 0$ équivaut à la
condition que $a_0$ soit une racine simple du polynôme $\overline{f}$.
En fait, elle implique que le relèvement $a \in R$, s'il existe, est unique.

\begin{lemma}
\label{lemma-uniqueness}
Soit $(R, \mathfrak m, \kappa)$ un anneau local.
Soit $f \in R[T]$. Soient $a, b \in R$ tels que $f(a) = f(b) = 0$,
$a = b \bmod \mathfrak m$, et $f'(a) \not \in \mathfrak m$.
Alors $a = b$.
\end{lemma}

\begin{proof}
Écrivons $f(x + y) - f(x) = f'(x)y + g(x, y) y^2$ dans $R[x, y]$ (cela est possible,
comme on le voit en développant $f(x + y)$ ; nous omettons les détails).
Nous voyons alors que $0 = f(b) - f(a) = f(a + (b - a)) - f(a) =
f'(a)(b - a) + c (b - a)^2$ pour un certain $c \in R$. Par hypothèse,
$f'(a)$ est une unité de $R$. Ainsi, $(b - a)(1 + f'(a)^{-1}c(b - a)) = 0$.
Par hypothèse, $b - a \in \mathfrak m$, donc $1 + f'(a)^{-1}c(b - a)$
est une unité de $R$. Ainsi, $b - a = 0$ dans $R$.
\end{proof}

\noindent
Voici la caractérisation des anneaux locaux henséliens.

\begin{lemma}
\label{lemma-characterize-henselian}
\begin{slogan}
Caractérisations des anneaux locaux henséliens
\end{slogan}
Soit $(R, \mathfrak m, \kappa)$ un anneau local.
Les assertions suivantes sont équivalentes :
\begin{enumerate}
\item $R$ est hensélien,
\item pour tout $f \in R[T]$ et toute racine $a_0 \in \kappa$
de $\overline{f}$ telle que $\overline{f'}(a_0) \not = 0$,
il existe un $a \in R$ tel que $f(a) = 0$ et
$a_0 = \overline{a}$,
\item pour tout $f \in R[T]$ unitaire et toute factorisation
$\overline{f} = g_0 h_0$ avec $\gcd(g_0, h_0) = 1$, il
existe une factorisation $f = gh$ dans $R[T]$ telle que
$g_0 = \overline{g}$ et $h_0 = \overline{h}$,
\item pour tout $f \in R[T]$ unitaire et toute factorisation
$\overline{f} = g_0 h_0$ avec $\gcd(g_0, h_0) = 1$, il
existe une factorisation $f = gh$ dans $R[T]$ telle que
$g_0 = \overline{g}$ et $h_0 = \overline{h}$ et, de plus,
$\deg_T(g) = \deg_T(g_0)$,
\item pour tout $f \in R[T]$ et toute factorisation
$\overline{f} = g_0 h_0$ avec $\gcd(g_0, h_0) = 1$, il
existe une factorisation $f = gh$ dans $R[T]$ telle que
$g_0 = \overline{g}$ et $h_0 = \overline{h}$,
\item pour tout $f \in R[T]$ et toute factorisation
$\overline{f} = g_0 h_0$ avec $\gcd(g_0, h_0) = 1$, il
existe une factorisation $f = gh$ dans $R[T]$ telle que
$g_0 = \overline{g}$ et $h_0 = \overline{h}$ et
telle que, de plus, $\deg_T(g) = \deg_T(g_0)$,
\item pour tout morphisme d'anneaux \'etale $R \to S$ et tout idéal premier
$\mathfrak q$ de $S$ au-dessus de $\mathfrak m$ avec
$\kappa = \kappa(\mathfrak q)$, il existe une rétraction
$\tau : S \to R$ de $R \to S$,
\item pour tout morphisme d'anneaux \'etale $R \to S$ et tout idéal premier
$\mathfrak q$ de $S$ au-dessus de $\mathfrak m$ avec
$\kappa = \kappa(\mathfrak q)$, il existe une unique rétraction
$\tau : S \to R$ de $R \to S$ telle que
$\mathfrak q = \tau^{-1}(\mathfrak m)$,
\item toute $R$-algèbre finie est un produit d'anneaux locaux,
\item toute $R$-algèbre finie est un produit fini d'anneaux locaux,
\item toute $R$-algèbre $S$ de type fini peut s'écrire
$A \times B$, avec $R \to A$ fini
et $R \to B$ non quasi-fini en tout idéal premier au-dessus de $\mathfrak m$,
\item toute $R$-algèbre $S$ de type fini peut s'écrire
$A \times B$, avec $R \to A$ fini,
de sorte que toute composante irréductible de $\Spec(B \otimes_R \kappa)$
soit de dimension $\geq 1$, et
\item toute $R$-algèbre quasi-finie $S$ peut s'écrire
$S = A \times B$, avec $R \to A$ fini, de sorte que $B \otimes_R \kappa = 0$.
\end{enumerate}
\end{lemma}

\begin{proof}
Voici une liste des implications les plus faciles :
\begin{enumerate}
\item 2$\Rightarrow$1 parce que, dans (2), nous considérons tous les polynômes et
dans (1) seulement ceux qui sont unitaires,
\item 5$\Rightarrow$3 parce que, dans (5), nous considérons tous les polynômes et
dans (3) seulement ceux qui sont unitaires,
\item 6$\Rightarrow$4 parce que, dans (6), nous considérons tous les polynômes et
dans (4) seulement ceux qui sont unitaires,
\item 4$\Rightarrow$3 est évidente,
\item 6$\Rightarrow$5 est évidente,
\item 8$\Rightarrow$7 est évidente,
\item 10$\Rightarrow$9 est évidente,
\item 11$\Leftrightarrow$12 par définition du fait d'être quasi-fini en un idéal premier,
\item 11$\Rightarrow$13 par définition du fait d'être quasi-fini,
\end{enumerate}

\noindent
Démonstration de 1$\Rightarrow$8. Supposons (1).
Soit $R \to S$ \'etale, et soit $\mathfrak q \subset S$
un idéal premier tel que $\kappa(\mathfrak q) \cong \kappa$. D'après la
proposition \ref{proposition-etale-locally-standard},
nous pouvons trouver un $g \in S$, $g \not \in \mathfrak q$, tel que
$R \to S_g$ soit \'etale standard. Après avoir remplacé $S$ par $S_g$, nous pouvons supposer
que $S = R[t]_g/(f)$ est \'etale standard (détails omis).
Puisque l'idéal premier $\mathfrak q$ a pour corps résiduel $\kappa$, il
correspond à une racine $a_0$ de $\overline{f}$ qui n'est pas une
racine de $\overline{g}$. Par définition d'une algèbre \'etale standard,
cela signifie aussi que $\overline{f'}(a_0) \not = 0$.
Puisque $f$ est également unitaire, à nouveau par définition d'une algèbre \'etale standard,
nous pouvons utiliser le fait que $R$ est hensélien pour conclure qu'il existe un $a \in R$
tel que $a_0 = \overline{a}$ et $f(a) = 0$. Cela implique que
$g(a)$ est une unité de $R$, et nous obtenons le morphisme recherché
$\tau : S = R[t]_g/(f) \to R$ par la règle $t \mapsto a$. Par construction,
$\tau^{-1}(\mathfrak m) = \mathfrak q$. D'après le lemme \ref{lemma-uniqueness},
le morphisme $\tau$ est unique. Cela prouve (8).

\medskip\noindent
Démonstration de 7$\Rightarrow$8. (Ceci est vraiment sans importance et devrait être
omis.) Supposons (7) vérifiée et supposons que $R \to S$ soit \'etale.
Soient $\mathfrak q_1, \ldots, \mathfrak q_r$
les autres idéaux premiers de $S$ au-dessus de $\mathfrak m$.
Nous pouvons alors trouver un $g \in S$, $g \not \in \mathfrak q$, tel que
$g \in \mathfrak q_i$ pour $i = 1, \ldots, r$.
En effet, nous pouvons faire valoir que
$\bigcap_{i=1}^{r} \mathfrak{q}_{i} \not\subset \mathfrak{q}$
car, sinon,
$\mathfrak{q}_{i} \subset \mathfrak{q}$
pour un certain $i$, ce qui est impossible puisque la fibre d'un
morphisme \'etale est discrète (utiliser par exemple le
lemme \ref{lemma-etale-over-field}).
Appliquons (7) au morphisme d'anneaux \'etale
$R \to S_g$ et à l'idéal premier $\mathfrak qS_g$. Nous obtenons une rétraction
$\tau_g : S_g \to R$ telle que la composée $\tau : S \to S_g \to R$
ait la propriété $\tau^{-1}(\mathfrak m) = \mathfrak q$.
Nous omettons les détails.

\medskip\noindent
Démonstration de 8$\Rightarrow$11. Supposons (8) et soit $R \to S$ un morphisme
d'anneaux de type fini. Appliquons le
lemme \ref{lemma-etale-makes-quasi-finite-finite}.
Nous trouvons un morphisme d'anneaux \'etale $R \to R'$ et un idéal premier
$\mathfrak m' \subset R'$ au-dessus de $\mathfrak m$, avec
$\kappa = \kappa(\mathfrak m')$, tels que
$R' \otimes_R S = A' \times B'$, où $A'$ est fini sur $R'$
et $B'$ n'est quasi-fini sur $R'$ en aucun idéal premier au-dessus de $\mathfrak m'$.
Appliquons (8) pour obtenir une rétraction $\tau : R' \to R$ avec
$\mathfrak m = \tau^{-1}(\mathfrak m')$. Utilisons alors
$$
S = (S \otimes_R R') \otimes_{R', \tau} R
= (A' \times B') \otimes_{R', \tau} R
= (A' \otimes_{R', \tau} R)  \times  (B' \otimes_{R', \tau} R)
$$
ce qui donne une décomposition comme dans (11).

\medskip\noindent
Démonstration de 8$\Rightarrow$10. Supposons (8) et soit $R \to S$ un morphisme
d'anneaux fini. Appliquons le
lemme \ref{lemma-etale-makes-quasi-finite-finite}.
Nous trouvons un morphisme d'anneaux \'etale $R \to R'$ et un idéal premier
$\mathfrak m' \subset R'$ au-dessus de $\mathfrak m$, avec
$\kappa = \kappa(\mathfrak m')$, tels que
$R' \otimes_R S = A'_1 \times \ldots \times A'_n \times B'$,
où $A'_i$ est fini sur $R'$ et possède exactement un idéal premier au-dessus de
$\mathfrak m'$, tandis que $B'$ n'est quasi-fini sur $R'$ en aucun idéal premier
au-dessus de $\mathfrak m'$.
Appliquons (8) pour obtenir une rétraction $\tau : R' \to R$ avec
$\mathfrak m' = \tau^{-1}(\mathfrak m)$. Alors nous obtenons
\begin{align*}
S & = (S \otimes_R R') \otimes_{R', \tau} R \\
& = (A'_1 \times \ldots \times A'_n \times B') \otimes_{R', \tau} R \\
& = (A'_1 \otimes_{R', \tau} R)  \times
\ldots \times (A'_1 \otimes_{R', \tau} R) \times
(B' \otimes_{R', \tau} R) \\
& = A_1 \times \ldots \times A_n \times B
\end{align*}
Le facteur $B$ est fini sur $R$, mais $R \to B$
n'est quasi-fini en aucun idéal premier au-dessus de $\mathfrak m$. Par conséquent,
$B = 0$. Les facteurs $A_i$ sont des $R$-algèbres finies possédant exactement
un idéal premier au-dessus de $\mathfrak m$ ; ce sont donc des anneaux locaux.
Cela prouve que $S$ est un produit fini d'anneaux locaux.

\medskip\noindent
Démonstration de 9$\Rightarrow$10. Cela est vrai parce que, si $S$ est fini sur l'anneau
local $R$, alors il possède au plus un nombre fini d'idéaux maximaux. En effet, par la
montée pour $R \to S$, les idéaux maximaux de $S$ se trouvent tous au-dessus de
$\mathfrak m$, et $S/\mathfrak mS$ est artinien, donc possède un nombre fini d'idéaux premiers.

\medskip\noindent
Démonstration de 10$\Rightarrow$1. Supposons (10). Soit $f \in R[T]$ un polynôme
unitaire et soit $a_0 \in \kappa$ une racine simple de $\overline{f}$.
Alors $S = R[T]/(f)$ est une $R$-algèbre finie. En appliquant (10),
nous obtenons que $S = A_1 \times \ldots \times A_r$ est un produit fini de
$R$-algèbres locales. En particulier, nous voyons que
$S/\mathfrak mS = \prod A_i/\mathfrak mA_i$ est la décomposition
de $\kappa[T]/(\overline{f})$ en produit d'anneaux locaux.
Cela signifie que l'un des facteurs, disons $A_1/\mathfrak mA_1$,
est le quotient $\kappa[T]/(\overline{f}) \to \kappa[T]/(T - a_0)$.
Comme $A_1$ est un facteur direct du $R$-module libre de type fini $S$, c'est
lui-même un $R$-module libre de type fini. Comme $A_1/\mathfrak mA_1$ est un
$\kappa$-espace vectoriel de dimension 1, nous voyons que $A_1 \cong R$ comme
$R$-module. Cela signifie clairement que $R \to A_1$ est un isomorphisme.
Soit $a \in R$ l'image de $T$ par le morphisme
$R[T] \to S \to A_1 \to R$. Alors $f(a) = 0$ et $\overline{a} = a_0$,
comme désiré.

\medskip\noindent
Démonstration de 13$\Rightarrow$1. Supposons (13). Soit $f \in R[T]$ un polynôme
unitaire et soit $a_0 \in \kappa$ une racine simple de $\overline{f}$.
Alors $S_1 = R[T]/(f)$ est une $R$-algèbre finie. Soit $g \in R[T]$
un élément quelconque tel que $\overline{g} = \overline{f}/(T - a_0)$.
Alors $S = (S_1)_g$ est une $R$-algèbre quasi-finie telle que
$S \otimes_R \kappa \cong \kappa[T]_{\overline{g}}/(\overline{f})
\cong \kappa[T]/(T - a_0) \cong \kappa$.
En appliquant (13) à $S$, nous obtenons $S = A \times B$, où $A$ est fini sur $R$ et
$B \otimes_R \kappa = 0$. En particulier, nous voyons que
$\kappa \cong S/\mathfrak mS = A/\mathfrak mA$.
Comme $A$ est un facteur direct de la $R$-algèbre plate $S$, nous voyons
qu'il est fini plat, donc libre sur $R$.
Comme $A/\mathfrak mA$ est un
$\kappa$-espace vectoriel de dimension 1, nous voyons que $A \cong R$ comme
$R$-module. Cela signifie clairement que $R \to A$ est un isomorphisme.
Soit $a \in R$ l'image de $T$ par le morphisme
$R[T] \to S \to A \to R$. Alors $f(a) = 0$ et $\overline{a} = a_0$,
comme désiré.

\medskip\noindent
Démonstration de 8$\Rightarrow$2. Supposons (8). Soit $f \in R[T]$ un
polynôme quelconque et soit $a_0 \in \kappa$ une racine simple. Alors
l'algèbre $S = R[T]_{f'}/(f)$ est \'etale sur $R$.
Soit $\mathfrak q \subset S$ l'idéal premier
engendré par $\mathfrak m$ et $T - b$, où $b \in R$ est un
élément quelconque tel que $\overline{b} = a_0$. Appliquons (8) à $S$ et
$\mathfrak q$ pour obtenir $\tau : S \to R$.
Alors l'image $\tau(T) = a \in R$ convient dans (2).

\medskip\noindent
À ce stade, nous voyons que (1), (2), (7), (8), (9), (10), (11), (12), (13) sont
toutes équivalentes. La plus faible des assertions (3), (4), (5) et (6)
est (3), et la plus forte est (6). Il nous reste donc à prouver que
(3) implique (1) et que (1) implique (6).

\medskip\noindent
Démonstration de 3$\Rightarrow$1. Supposons (3). Soit $f \in R[T]$ unitaire et
soit $a_0 \in \kappa$ une racine simple de $\overline{f}$. Cela donne
une factorisation $\overline{f} = (T - a_0)h_0$ avec $h_0(a_0) \not = 0$,
donc $\gcd(T - a_0, h_0) = 1$. Appliquons (3) pour obtenir une factorisation
$f = gh$ avec $\overline{g} = T - a_0$ et $\overline{h} = h_0$.
Posons $S = R[T]/(f)$, qui est une $R$-algèbre libre de type fini. Nous noterons
aussi $g$, $h$ les images de $g$ et $h$ dans $S$. Alors
$gS + hS = S$ d'après le
lemme de Nakayama \ref{lemma-NAK},
puisque l'égalité est vraie modulo $\mathfrak m$. Comme $gh = f = 0$ dans $S$,
cela implique aussi que $gS \cap hS = 0$. Le théorème chinois des restes
donne donc $S = S/(g) \times S/(h)$. Cela implique que
$A = S/(g)$ est un facteur direct d'un $R$-module libre de type fini, donc est
libre de type fini. De plus, le rang de $A$ est $1$, car
$A/\mathfrak mA = \kappa[T]/(T - a_0)$. Le morphisme $R \to A$
est donc un isomorphisme. En prenant pour $a \in R$ l'image de $T$
par les morphismes $R[T] \to S \to A \to R$, nous obtenons un élément de $R$
tel que $f(a) = 0$ et $\overline{a} = a_0$.

\medskip\noindent
Démonstration de 1$\Rightarrow$6. Supposons (1), ou, de façon équivalente, toutes les
assertions (1), (2), (7), (8), (9), (10), (11), (12), (13).
Soit $f \in R[T]$ un polynôme.
Supposons que $\overline{f} = g_0h_0$ soit une factorisation avec
$\gcd(g_0, h_0) = 1$. Nous pouvons supposer, et supposons, que $g_0$ est unitaire.
Considérons $S = R[T]/(f)$. Puisque nous
disposons de cette factorisation, nous voyons que les coefficients de
$f$ engendrent l'idéal unité de $R$.
Cela implique que $S$ a des fibres finies sur $R$, et est donc
quasi-fini sur $R$. Cela implique également que $S$ est plat sur $R$ d'après le
lemme \ref{lemma-grothendieck-general}.
En combinant (13) et (10), nous pouvons écrire
$S = A_1 \times \ldots \times A_n \times B$,
où chaque $A_i$ est local et fini sur $R$, et
$B \otimes_R \kappa = 0$. Après avoir réordonné les facteurs $A_1, \ldots, A_n$,
nous pouvons supposer que
$$
\kappa[T]/(g_0) =
A_1/\mathfrak m A_1 \times \ldots \times A_r/\mathfrak mA_r,
\ \kappa[T]/(h_0) =
A_{r + 1}/\mathfrak mA_{r + 1} \times \ldots \times A_n/\mathfrak mA_n
$$
comme quotients de $\kappa[T]$. La $R$-algèbre finie plate
$A = A_1 \times \ldots \times A_r$ est libre comme $R$-module, voir le
lemme \ref{lemma-finite-flat-local}.
Son rang est $\deg_T(g_0)$. Soit $g \in R[T]$ le polynôme caractéristique
de l'opérateur $R$-linéaire $T : A \to A$. Alors $g$ est un polynôme unitaire
de degré $\deg_T(g) = \deg_T(g_0)$ et, de plus, $\overline{g} = g_0$.
Par Cayley-Hamilton
(lemme \ref{lemma-charpoly}),
nous voyons que $g(T_A) = 0$, où $T_A$ désigne
l'image de $T$ dans $A$. Nous obtenons donc un morphisme surjectif bien défini
$R[T]/(g) \to A$, qui est un isomorphisme d'après le
lemme de Nakayama \ref{lemma-NAK}. Le morphisme $R[T] \to A$ se factorise
par $R[T]/(f)$ par construction ; nous pouvons donc écrire $f = gh$ pour
un certain $h$. Cela achève la démonstration.
\end{proof}

\begin{lemma}
\label{lemma-finite-over-henselian}
Soit $(R, \mathfrak m, \kappa)$ un anneau local hensélien.
\begin{enumerate}
\item Si $R \to S$ est un morphisme d'anneaux fini, alors $S$ est
un produit fini d'anneaux locaux henséliens, chacun fini sur $R$.
\item Si $R \to S$ est un morphisme d'anneaux fini et si $S$ est local, alors
$S$ est un anneau local hensélien et $R \to S$ est un morphisme (fini) d'anneaux locaux.
\item Si $R \to S$ est un morphisme d'anneaux de type fini, et si $\mathfrak q$ est
un idéal premier de $S$ au-dessus de $\mathfrak m$ auquel $R \to S$ est quasi-fini,
alors $S_{\mathfrak q}$ est hensélien et fini sur $R$.
\item Si $R \to S$ est quasi-fini, alors $S_{\mathfrak q}$ est hensélien
et fini sur $R$ pour tout idéal premier $\mathfrak q$ au-dessus de $\mathfrak m$.
\end{enumerate}
\end{lemma}

\begin{proof}
La partie (2) implique la partie (1), puisque $S$, comme dans la partie (1), est un produit fini
de ses localisations aux idéaux premiers au-dessus de $\mathfrak m$ d'après le
lemme \ref{lemma-characterize-henselian}, partie (10).
La partie (2) résulte aussi du lemme \ref{lemma-characterize-henselian}, partie (10),
puisque toute $S$-algèbre finie est également une $R$-algèbre finie (bien sûr,
tout morphisme d'anneaux fini entre anneaux locaux est local).

\medskip\noindent
Soient $R \to S$ et $\mathfrak q$ comme dans (3). Écrivons $S = A \times B$,
où $A$ est fini sur $R$ et $B$ n'est quasi-fini sur $R$ en aucun
idéal premier au-dessus de $\mathfrak m$, voir le
lemme \ref{lemma-characterize-henselian}, partie (11).
Ainsi, $S_\mathfrak q$ est une localisation de $A$ en un idéal maximal,
et nous déduisons (3) de (1). La partie (4) résulte de la partie (3).
\end{proof}

\begin{lemma}
\label{lemma-mop-up}
Soit $(R, \mathfrak m, \kappa)$ un anneau local hensélien.
Toute $R$-algèbre $S$ de type fini peut s'écrire
$S = A_1 \times \ldots \times A_n \times B$, où les $A_i$ sont locaux
et finis sur $R$, et où $R \to B$ n'est quasi-fini en aucun
idéal premier de $B$ au-dessus de $\mathfrak m$.
\end{lemma}

\begin{proof}
C'est une combinaison des parties (11) et (10) du
lemme \ref{lemma-characterize-henselian}.
\end{proof}

\begin{lemma}
\label{lemma-mop-up-strictly-henselian}
Soit $(R, \mathfrak m, \kappa)$ un anneau local strictement hensélien.
Toute $R$-algèbre $S$ de type fini peut s'écrire
$S = A_1 \times \ldots \times A_n \times B$, où les $A_i$ sont locaux
et finis sur $R$, où $\kappa \subset \kappa(\mathfrak m_{A_i})$ est
finie purement inséparable, et où $R \to B$ n'est quasi-fini
en aucun idéal premier de $B$ au-dessus de $\mathfrak m$.
\end{lemma}

\begin{proof}
Écrivons d'abord $S = A_1 \times \ldots \times A_n \times B$ comme dans le
lemme \ref{lemma-mop-up}.
L'extension de corps $\kappa(\mathfrak m_{A_i})/\kappa$
est finie et $\kappa$ est séparablement clos ; elle
est donc finie purement inséparable.
\end{proof}

\begin{lemma}
\label{lemma-henselian-cat-finite-etale}
Soit $(R, \mathfrak m, \kappa)$ un anneau local hensélien.
La catégorie des extensions finies étales d'anneaux $R \to S$ est
équivalente à la catégorie des algèbres finies étales
$\kappa \to \overline{S}$ par le foncteur $S \mapsto S/\mathfrak mS$.
\end{lemma}

\begin{proof}
Notons $\mathcal{C} \to \mathcal{D}$ le foncteur entre les catégories
de l'énoncé.
Supposons que $R \to S$ soit fini étale. Alors nous pouvons écrire
$$
S = A_1 \times \ldots \times A_n
$$
avec les $A_i$ locaux et finis étales sur $S$, en utilisant soit le
lemme \ref{lemma-mop-up},
soit le
lemme \ref{lemma-characterize-henselian}, partie (10).
En particulier, $A_i/\mathfrak mA_i$ est une extension de corps finie
séparable de $\kappa$, voir le
lemme \ref{lemma-etale-at-prime}.
Nous voyons ainsi que tout objet de $\mathcal{C}$ et de
$\mathcal{D}$ se décompose canoniquement en composantes irréductibles
qui se correspondent par le foncteur donné.
Supposons ensuite que $S_1$, $S_2$ soient finis étales sur $R$, et que
$\kappa_1 = S_1/\mathfrak mS_1$ et $\kappa_2 = S_2/\mathfrak mS_2$
soient des corps (finis séparables sur $\kappa$). Alors $S_1 \otimes_R S_2$
est fini étale sur $R$, et nous pouvons écrire
$$
S_1 \otimes_R S_2 = A_1 \times \ldots \times A_n
$$
comme précédemment. Nous voyons alors que $\Hom_R(S_1, S_2)$ s'identifie
à l'ensemble des indices $i \in \{1, \ldots, n\}$ tels que
$S_2 \to A_i$ soit un isomorphisme. Pour le voir, utiliser que, pour tout morphisme
de $R$-algèbres $\varphi : S_1 \to S_2$, le morphisme
$\varphi \times 1 : S_1 \otimes_R S_2 \to S_2$
est surjectif, et est donc égal à la projection sur l'un des facteurs $A_i$.
Mais nous voyons exactement de la même manière que
$\Hom_\kappa(\kappa_1, \kappa_2)$ s'identifie à
l'ensemble des indices $i \in \{1, \ldots, n\}$ tels que
$\kappa_2 \to A_i/\mathfrak mA_i$ soit un isomorphisme.
D'après la discussion ci-dessus, ces ensembles d'indices coïncident, et nous concluons
que notre foncteur est pleinement fidèle.
Enfin, soit $\kappa'/\kappa$ une extension de corps finie
séparable. D'après le
lemme \ref{lemma-make-etale-map-prescribed-residue-field},
il existe un morphisme d'anneaux \'etale $R \to S$ et un idéal premier
$\mathfrak q$ de $S$ au-dessus de $\mathfrak m$ tels que
$\kappa \subset \kappa(\mathfrak q)$ soit isomorphe à l'extension donnée.
D'après la partie (1), nous pouvons écrire
$S = A_1 \times \ldots \times A_n \times B$.
Comme $R \to S$ est quasi-fini, nous voyons qu'il n'existe aucun
idéal premier de $B$ au-dessus de $\mathfrak m$. Ainsi, $S_{\mathfrak q}$ est
égal à $A_i$ pour un certain $i$. Par conséquent, $R \to A_i$ est fini étale
et produit l'extension de corps résiduels donnée. Le foncteur est donc
essentiellement surjectif, ce qui conclut la démonstration.
\end{proof}

\begin{lemma}
\label{lemma-unramified-over-strictly-henselian}
Soit $(R, \mathfrak m, \kappa)$ un anneau local strictement hensélien.
Soit $R \to S$ un morphisme d'anneaux non ramifié. Alors
$$
S = A_1 \times \ldots \times A_n \times B
$$
où chaque $R \to A_i$ est surjectif et où aucun idéal premier de $B$ ne se trouve
au-dessus de $\mathfrak m$.
\end{lemma}

\begin{proof}
Écrivons d'abord $S = A_1 \times \ldots \times A_n \times B$ comme dans le
lemme \ref{lemma-mop-up}.
Nous voyons maintenant que $R \to A_i$ est fini non ramifié et que $A_i$ est local.
L'idéal maximal de $A_i$ est donc $\mathfrak mA_i$, et son
corps résiduel $A_i / \mathfrak m A_i$ est une extension finie
séparable de $\kappa$, voir le
lemme \ref{lemma-unramified-at-prime}.
Cependant, la condition que $R$ soit strictement hensélien signifie que
$\kappa$ est séparablement clos, donc
$\kappa = A_i / \mathfrak m A_i$. D'après le
lemme de Nakayama \ref{lemma-NAK},
nous concluons que $R \to A_i$ est surjectif, comme désiré.
\end{proof}

\begin{lemma}
\label{lemma-complete-henselian}
\begin{slogan}
Les anneaux locaux complets sont henséliens par la méthode de Newton
\end{slogan}
Soit $(R, \mathfrak m, \kappa)$ un anneau local complet, voir la
définition \ref{definition-complete-local-ring}.
Alors $R$ est hensélien.
\end{lemma}

\begin{proof}
Soit $f \in R[T]$ unitaire.
Notons $f_n \in R/\mathfrak m^{n + 1}[T]$ son image.
Notons $f'_n$ la dérivée de $f_n$ par rapport à $T$.
Soit $a_0 \in \kappa$ une racine simple de $f_0$. Nous la relevons
par récurrence en une solution de $f$ sur $R$ comme suit :
supposons donné $a_n \in R/\mathfrak m^{n + 1}$ tel que
$a_n \bmod \mathfrak m = a_0$ et $f_n(a_n) = 0$. Choisissons un
élément quelconque $b \in R/\mathfrak m^{n + 2}$ tel que
$a_n = b \bmod \mathfrak m^{n + 1}$. Alors
$f_{n + 1}(b) \in \mathfrak m^{n + 1}/\mathfrak m^{n + 2}$.
Posons
$$
a_{n + 1} = b - f_{n + 1}(b)/f'_{n + 1}(b)
$$
(méthode de Newton). Cela a un sens, car
$f'_{n + 1}(b) \in R/\mathfrak m^{n + 2}$
est inversible par la condition sur $a_0$. Nous calculons alors
$f_{n + 1}(a_{n + 1}) = f_{n + 1}(b) - f_{n + 1}(b) = 0$
dans $R/\mathfrak m^{n + 2}$. Comme le système d'éléments
$a_n \in R/\mathfrak m^{n + 1}$ ainsi construit est compatible,
nous obtenons un élément
$a \in \lim R/\mathfrak m^n = R$ (nous utilisons ici que $R$ est complet).
De plus, $f(a) = 0$, puisqu'il s'envoie sur zéro dans chaque
$R/\mathfrak m^n$.
Enfin, $\overline{a} = a_0$, ce qui conclut.
\end{proof}

\begin{lemma}
\label{lemma-local-dimension-zero-henselian}
\begin{slogan}
Les anneaux locaux de dimension zéro sont henséliens.
\end{slogan}
Soit $(R, \mathfrak m)$ un anneau local de dimension $0$.
Alors $R$ est hensélien.
\end{lemma}

\begin{proof}
Soit $R \to S$ un morphisme d'anneaux fini. D'après le
lemme \ref{lemma-characterize-henselian},
il suffit de montrer que $S$ est un produit d'anneaux locaux. D'après le
lemme \ref{lemma-finite-finite-fibres},
$S$ possède un nombre fini d'idéaux premiers
$\mathfrak m_1, \ldots, \mathfrak m_r$,
qui se trouvent tous au-dessus de $\mathfrak m$. Il n'existe aucune inclusion entre ces
idéaux premiers, voir le
lemme \ref{lemma-integral-no-inclusion} ;
ils sont donc tous maximaux. Tout élément de
$\mathfrak m_1 \cap \ldots \cap \mathfrak m_r$ est nilpotent d'après le
lemme \ref{lemma-Zariski-topology}.
Il s'ensuit que $S$ est le produit des localisations de $S$ aux idéaux premiers
$\mathfrak m_i$, d'après le
lemme \ref{lemma-product-local}.
\end{proof}

\noindent
Le lemme suivant sera la clé des propriétés d'unicité et de fonctorialité
de l'hensélisation et de l'hensélisation stricte.

\begin{lemma}
\label{lemma-map-into-henselian}
Soit $R \to S$ un morphisme d'anneaux, avec $S$ local hensélien.
Étant donnés
\begin{enumerate}
\item un morphisme d'anneaux \'etale $R \to A$,
\item un idéal premier $\mathfrak q$ de $A$ au-dessus de
$\mathfrak p = R \cap \mathfrak m_S$,
\item un morphisme de $\kappa(\mathfrak p)$-algèbres
$\tau : \kappa(\mathfrak q) \to S/\mathfrak m_S$,
\end{enumerate}
il existe un unique morphisme de $R$-algèbres $f : A \to S$
tel que $\mathfrak q = f^{-1}(\mathfrak m_S)$ et que
$f$ induise le morphisme $\tau$ sur les corps résiduels.
\end{lemma}

\begin{proof}
Considérons $A \otimes_R S$. C'est une algèbre \'etale sur $S$, voir le
lemme \ref{lemma-etale}. De plus, le noyau
$$
\mathfrak q' = \Ker(A \otimes_R S \to
\kappa(\mathfrak q) \otimes_{\kappa(\mathfrak p)} \kappa(\mathfrak m_S)
\xrightarrow{\tau \otimes 1}
\kappa(\mathfrak m_S))
$$
est un idéal premier au-dessus de $\mathfrak m_S$ dont le corps résiduel est égal
à celui de $S$. Ainsi, d'après le lemme \ref{lemma-characterize-henselian},
il existe une unique rétraction $\sigma : A \otimes_R S \to S$
avec $\sigma^{-1}(\mathfrak m_S) = \mathfrak q'$.
Prenons pour $f$ la composée $A \to A \otimes_R S \to S$.
Nous omettons la vérification des propriétés de $f$ ; l'unicité
de $f$ résulte de celle de $\sigma$ (détails omis).
\end{proof}

\begin{lemma}
\label{lemma-strictly-henselian-solutions}
Soit $\varphi : R \to S$ un morphisme local
d'anneaux locaux strictement henséliens.
Soient $P_1, \ldots, P_n \in R[x_1, \ldots, x_n]$ des polynômes tels que
$R[x_1, \ldots, x_n]/(P_1, \ldots, P_n)$ soit \'etale sur $R$.
Alors le morphisme
$$
R^n \longrightarrow S^n, \quad
(h_1, \ldots, h_n) \longmapsto (\varphi(h_1), \ldots, \varphi(h_n))
$$
induit une bijection entre
$$
\{
(r_1, \ldots, r_n) \in R^n
\mid
P_i(r_1, \ldots, r_n) = 0, \ i = 1, \ldots, n
\}
$$
et
$$
\{
(s_1, \ldots, s_n) \in S^n
\mid
P^\varphi_i(s_1, \ldots, s_n) = 0, \ i = 1, \ldots, n
\}
$$
où les $P^\varphi_i \in S[x_1, \ldots, x_n]$ sont les images des $P_i$
par $\varphi$.
\end{lemma}

\begin{proof}
Le premier ensemble de solutions est canoniquement isomorphe à l'ensemble
$$
\Hom_R(R[x_1, \ldots, x_n]/(P_1, \ldots, P_n), R).
$$
Comme $R$ est hensélien, le morphisme $R \to R/\mathfrak m_R$ induit une bijection
entre cet ensemble et l'ensemble des solutions dans le
corps résiduel $R/\mathfrak m_R$, voir le
lemme \ref{lemma-characterize-henselian}.
La même chose est vraie pour $S$.
Puisque $R[x_1, \ldots, x_n]/(P_1, \ldots, P_n)$ est \'etale sur $R$
et que $R/\mathfrak m_R$ est séparablement clos, nous voyons que
$R/\mathfrak m_R[x_1, \ldots, x_n]/(\overline{P}_1, \ldots, \overline{P}_n)$
est un produit fini de copies de $R/\mathfrak m_R$, où $\overline{P}_i$
est l'image de $P_i$ dans $R/\mathfrak m_R[x_1, \ldots, x_n]$. Par conséquent, le
produit tensoriel
$$
R/\mathfrak m_R[x_1, \ldots, x_n]/(\overline{P}_1, \ldots, \overline{P}_n)
\otimes_{R/\mathfrak m_R} S/\mathfrak m_S
=
S/\mathfrak m_S[x_1, \ldots, x_n]/
(\overline{P}^\varphi_1, \ldots, \overline{P}^\varphi_n)
$$
est lui aussi un produit fini de copies de $S/\mathfrak m_S$ avec le même
ensemble d'indices. Cela prouve le lemme.
\end{proof}

\begin{lemma}
\label{lemma-split-ML-henselian}
Soit $R$ un anneau local hensélien.
Tout module de Mittag-Leffler à engendrement dénombrable sur $R$ est une somme
directe de $R$-modules de présentation finie.
\end{lemma}

\begin{proof}
Soit $M$ un $R$-module à engendrement dénombrable et de Mittag-Leffler.
Nous affirmons que, pour tout élément $x \in M$, il existe une décomposition en somme
directe $M = N \oplus K$ avec $x \in N$, le module
$N$ de présentation finie et $K$ de Mittag-Leffler.

\medskip\noindent
Supposons l'affirmation vraie. Choisissons des générateurs
$x_1, x_2, x_3, \ldots$ de $M$. D'après l'affirmation, nous pouvons trouver
par récurrence des décompositions en somme directe
$$
M = N_1 \oplus N_2 \oplus \ldots \oplus N_n \oplus K_n
$$
avec les $N_i$ de présentation finie,
$x_1, \ldots, x_n \in N_1 \oplus \ldots \oplus N_n$, et $K_n$ de Mittag-Leffler.
En répétant cela indéfiniment, nous voyons que $M = \bigoplus N_i$.

\medskip\noindent
Il nous reste à démontrer l'affirmation. Soit $x \in M$. D'après le
lemme \ref{lemma-ML-countable},
il existe un endomorphisme $\alpha : M \to M$
tel que $\alpha$ se factorise par un module de présentation finie et
$\alpha (x) = x$. Disons que $\alpha$ se factorise comme
$$
\xymatrix{
M \ar[r]^\pi & P \ar[r]^i & M
}
$$
Posons $a = \pi \circ \alpha \circ i : P \to P$, de sorte que
$i \circ a \circ \pi = \alpha^3$. D'après le
lemme \ref{lemma-charpoly-module},
il existe un polynôme unitaire $P \in R[T]$ tel que $P(a) = 0$.
Notons que cela implique formellement que $\alpha^2 P(\alpha) = 0$.
Nous pouvons donc considérer $M$ comme un module sur $R[T]/(T^2P)$.
Supposons que $x \not = 0$. Alors $\alpha(x) = x$ implique que
$0 = \alpha^2P(\alpha)x = P(1)x$, donc que $P(1) = 0$ dans $R/I$, où
$I = \{r \in R \mid rx = 0\}$ est l'annulateur de $x$.
Comme $x \not = 0$, nous voyons que $I \subset \mathfrak m_R$, donc que
$1$ est une racine de $\overline{P} = P \bmod \mathfrak m_R \in R/\mathfrak m_R[T]$.
Comme $R$ est hensélien, nous pouvons trouver une factorisation
$$
T^2P = (T^2 Q_1) Q_2
$$
pour certains $Q_1, Q_2 \in R[T]$ avec
$Q_2 = (T - 1)^e \bmod \mathfrak m_R R[T]$ et
$Q_1(1) \not = 0 \bmod \mathfrak m_R$, voir le
lemme \ref{lemma-characterize-henselian}.
Soient $N = \Im(\alpha^2Q_1(\alpha) : M \to M)$ et
$K = \Im(Q_2(\alpha) : M \to M)$. Comme $T^2Q_1$ et
$Q_2$ engendrent l'idéal unité de $R[T]$, nous obtenons une décomposition en somme
directe $M = N \oplus K$. De plus, $Q_2$ agit par zéro sur $N$, et
$T^2Q_1$ agit par zéro sur $K$. Notons que $N$ est un quotient de $P$,
donc est de type fini. De plus, $x \in N$, parce que
$\alpha^2Q_1(\alpha)x = Q_1(1)x$ et que $Q_1(1)$ est une unité de $R$. D'après le
lemme \ref{lemma-direct-sum-ML},
les modules $N$ et $K$ sont de Mittag-Leffler. Enfin, le module de type fini
$N$ est de présentation finie, puisqu'un module de Mittag-Leffler de type fini
est de présentation finie, voir
l'exemple \ref{example-ML}, partie (1).
\end{proof}


\section{Colimites filtrantes de morphismes d'anneaux \'etales}
\label{section-ind-etale}

\noindent
Cette section prépare la section sur les morphismes d'anneaux ind-\'etales
(Cohomologie pro-\'etale, section \ref{proetale-section-ind-etale}).
Les résultats seront également utiles pour démontrer des propriétés d'unicité de
l'hensélisation et de l'hensélisation stricte d'un anneau local.

\begin{lemma}
\label{lemma-base-change-colimit-etale}
Soient $R \to A$ et $R \to R'$ des morphismes d'anneaux. Si $A$ est une colimite
filtrante de $R$-algèbres \'etales, alors $R' \otimes_R A$
est une colimite filtrante de $R'$-algèbres \'etales.
\end{lemma}

\begin{proof}
Cela est vrai parce que les colimites commutent aux produits tensoriels
et que les morphismes d'anneaux \'etales sont préservés par changement de base
(lemme \ref{lemma-etale}).
\end{proof}

\begin{lemma}
\label{lemma-composition-colimit-etale}
Soient $A \to B \to C$ des morphismes d'anneaux. Si $B$ est une colimite
filtrante de $A$-algèbres \'etales et si $C$ est une colimite filtrante
de $B$-algèbres \'etales, alors $C$ est une colimite filtrante
de $A$-algèbres \'etales.
\end{lemma}

\begin{proof}
Nous utiliserons le critère du lemme \ref{lemma-when-colimit}.
Soit $A \to P \to C$ une factorisation de $A \to C$
avec $P$ de présentation finie sur $A$.
Écrivons $B = \colim_{i \in I} B_i$, où $I$ est un ensemble ordonné filtrant et
où $B_i$ est une $A$-algèbre \'etale.
Écrivons $C = \colim_{j \in J} C_j$, où $J$ est un ensemble ordonné filtrant et
où $C_j$ est une $B$-algèbre \'etale.
Nous pouvons factoriser $P \to C$ sous la forme $P \to C_j \to C$ pour un
certain $j$, d'après le lemme \ref{lemma-characterize-finite-presentation}.
D'après le lemme \ref{lemma-etale}, nous pouvons trouver
un $i \in I$ et un morphisme d'anneaux \'etale $B_i \to C'_j$
tels que $C_j = B \otimes_{B_i} C'_j$.
Alors $C_j = \colim_{i' \geq i} B_{i'} \otimes_{B_i} C'_j$,
et nous voyons à nouveau que $P \to C_j$ se factorise sous la forme
$P \to B_{i'} \otimes_{B_i} C'_j \to C$.
Le morphisme $A \to C' = B_{i'} \otimes_{B_i} C'_j$ est \'etale, car les
composées et les produits tensoriels de morphismes d'anneaux \'etales
sont \'etales. Nous avons donc factorisé $P \to C$ sous la forme
$P \to C' \to C$, avec $C'$ \'etale sur $A$, et le critère du
lemme \ref{lemma-when-colimit} s'applique.
\end{proof}

\begin{lemma}
\label{lemma-colimit-colimit-etale}
Soit $R$ un anneau. Soit $A = \colim A_i$ une colimite filtrante
de $R$-algèbres telle que chaque $A_i$ soit une colimite filtrante de
$R$-algèbres \'etales. Alors $A$ est une colimite filtrante de
$R$-algèbres \'etales.
\end{lemma}

\begin{proof}
Écrivons $A_i = \colim_{j \in J_i} A_j$, où $J_i$ est un ensemble ordonné filtrant et
où $A_j$ est une $R$-algèbre \'etale.
Pour tous $i \leq i'$ et $j \in J_i$, il existe un
$j' \in J_{i'}$ et un morphisme de $R$-algèbres
$\varphi_{jj'} : A_j \to A_{j'}$ rendant commutatif le diagramme
$$
\xymatrix{
A_i \ar[r] & A_{i'} \\
A_j \ar[u] \ar[r]^{\varphi_{jj'}} & A_{j'} \ar[u]
}
$$
Cela est vrai parce que $R \to A_j$ est de présentation finie,
de sorte que le lemme \ref{lemma-characterize-finite-presentation} s'applique.
Soit $\mathcal{J}$ la catégorie dont les objets sont
$\coprod_{i \in I} J_i$ et dont les morphismes sont les triplets
$(j, j', \varphi_{jj'})$ ci-dessus (avec la loi de composition évidente).
Alors $\mathcal{J}$ est une catégorie filtrante et
$A = \colim_\mathcal{J} A_j$. Nous omettons les détails.
\end{proof}

\begin{lemma}
\label{lemma-colimit-colimit-etale-better}
Soit $I$ un ensemble ordonné filtrant. Soit
$i \mapsto (R_i \to A_i)$ un système de flèches d'anneaux indexé par $I$.
Posons $R = \colim R_i$ et $A = \colim A_i$. Si chaque $A_i$ est une
colimite filtrante de $R_i$-algèbres \'etales, alors $A$ est une colimite filtrante
de $R$-algèbres \'etales.
\end{lemma}

\begin{proof}
Cela est vrai parce que $A = A \otimes_R R = \colim A_i \otimes_{R_i} R$ ;
nous pouvons donc appliquer le lemme \ref{lemma-colimit-colimit-etale},
puisque $R \to A_i \otimes_{R_i} R$ est une colimite filtrante de
morphismes d'anneaux \'etales d'après le
lemme \ref{lemma-base-change-colimit-etale}.
\end{proof}

\begin{lemma}
\label{lemma-colimits-of-etale}
Soit $R$ un anneau. Soit $A \to B$ un morphisme de $R$-algèbres.
Si $A$ et $B$ sont des colimites filtrantes de $R$-algèbres \'etales, alors
$B$ est une colimite filtrante de $A$-algèbres \'etales.
\end{lemma}

\begin{proof}
Écrivons $A = \colim A_i$ et $B = \colim B_j$ comme colimites filtrantes, avec
les $A_i$ et les $B_j$ \'etales sur $R$. Pour tout $i$, nous pouvons trouver un
$j$ tel que $A_i \to B$ se factorise par $B_j$, voir le
lemme \ref{lemma-characterize-finite-presentation}.
La factorisation $A_i \to B_j$ est \'etale d'après le
lemme \ref{lemma-map-between-etale}.
Comme $A \to A \otimes_{A_i} B_j$ est \'etale (lemme \ref{lemma-etale}),
il suffit de prouver que $B = \colim A \otimes_{A_i} B_j$, où la
colimite est prise sur les paires $(i, j)$ et les factorisations
$A_i \to B_j \to B$ de $A_i \to B$ (c'est un système filtrant ; détails omis).
Cela est clair parce que les colimites commutent aux produits tensoriels,
et donc $\colim A \otimes_{A_i} B_j = A \otimes_A B = B$.
\end{proof}

\begin{lemma}
\label{lemma-map-into-henselian-colimit}
Soit $R \to S$ un morphisme d'anneaux, avec $S$ local hensélien. Étant donnés
\begin{enumerate}
\item une $R$-algèbre $A$ qui est une colimite filtrante de $R$-algèbres \'etales,
\item un idéal premier $\mathfrak q$ de $A$ au-dessus de
$\mathfrak p = R \cap \mathfrak m_S$,
\item un morphisme de $\kappa(\mathfrak p)$-algèbres
$\tau : \kappa(\mathfrak q) \to S/\mathfrak m_S$,
\end{enumerate}
il existe un unique morphisme de $R$-algèbres $f : A \to S$
tel que $\mathfrak q = f^{-1}(\mathfrak m_S)$ et
$f \bmod \mathfrak q = \tau$.
\end{lemma}

\begin{proof}
Écrivons $A = \colim A_i$ comme colimite filtrante de $R$-algèbres \'etales.
Posons $\mathfrak q_i = A_i \cap \mathfrak q$. Nous obtenons
$f_i : A_i \to S$ en appliquant le lemme \ref{lemma-map-into-henselian}.
Posons $f = \colim f_i$.
\end{proof}

\begin{lemma}
\label{lemma-uniqueness-henselian}
Soit $R$ un anneau. Étant donné un diagramme commutatif de morphismes d'anneaux
$$
\xymatrix{
S \ar[r] & K \\
R \ar[u] \ar[r] & S' \ar[u]
}
$$
où $S$, $S'$ sont locaux henséliens, où $S$, $S'$ sont des colimites filtrantes
de $R$-algèbres \'etales, où $K$ est un corps, et où les flèches $S \to K$ et
$S' \to K$ identifient $K$ au corps résiduel de $S$ et de $S'$,
il existe un unique isomorphisme de $R$-algèbres $S \to S'$
compatible avec les morphismes vers $K$.
\end{lemma}

\begin{proof}
Cela résulte immédiatement du lemme \ref{lemma-map-into-henselian-colimit}.
\end{proof}

\noindent
Le lemme suivant ne concerne pas, à proprement parler, les colimites de
morphismes d'anneaux \'etales.

\begin{lemma}
\label{lemma-colimit-henselian}
Une colimite filtrante d'anneaux locaux (strictement) henséliens suivant des
morphismes locaux est (strictement) hensélienne.
\end{lemma}

\begin{proof}
Le lemme \ref{categories-lemma-directed-category-system} de Catégories
dit qu'il s'agit en réalité simplement d'une colimite d'anneaux
locaux (strictement) henséliens indexée par un ensemble ordonné filtrant.
Soit $(R_i, \varphi_{ii'})$ un tel système, avec chaque $\varphi_{ii'}$
local. Alors $R = \colim_i R_i$ est local, et
son corps résiduel $\kappa$ est $\colim \kappa_i$
(argument omis). Il est facile de voir que $\colim \kappa_i$
est séparablement clos si chaque $\kappa_i$ l'est ;
il suffit donc de prouver que $R$ est hensélien si chaque $R_i$ est hensélien.
Supposons que $f \in R[T]$ soit unitaire et que $a_0 \in \kappa$ soit
une racine simple de $\overline{f}$. Alors, pour un $i$ assez grand,
il existe un $f_i \in R_i[T]$ s'envoyant sur $f$ et un
$a_{0, i} \in \kappa_i$ s'envoyant sur $a_0$. Puisque
$\overline{f_i}(a_{0, i}) \in \kappa_i$,
resp.\ $\overline{f_i'}(a_{0, i}) \in \kappa_i$, s'envoie sur
$0 = \overline{f}(a_0) \in \kappa$,
resp.\ $0 \not = \overline{f'}(a_0) \in \kappa$,
nous concluons que $a_{0, i}$ est une racine simple de $\overline{f_i}$.
Comme $R_i$ est hensélien, nous pouvons trouver $a_i \in R_i$ tel que
$f_i(a_i) = 0$ et $a_{0, i} = \overline{a_i}$.
Alors l'image $a \in R$ de $a_i$ est la solution recherchée.
Ainsi, $R$ est hensélien.
\end{proof}



\section{Hensélisation et hensélisation stricte}
\label{section-henselization}

\noindent
Dans cette section, nous construisons l'hensélisation. Nous encourageons le lecteur
à garder à l'esprit l'unicité déjà démontrée dans le
lemme \ref{lemma-uniqueness-henselian}
et le comportement fonctoriel signalé dans le
lemme \ref{lemma-map-into-henselian-colimit}
pendant la lecture de ces résultats.

\begin{lemma}
\label{lemma-henselization}
Soit $(R, \mathfrak m, \kappa)$ un anneau local. Il existe un
morphisme local d'anneaux $R \to R^h$ ayant les propriétés suivantes :
\begin{enumerate}
\item $R^h$ est hensélien,
\item $R^h$ est une colimite filtrante de $R$-algèbres \'etales,
\item $\mathfrak m R^h$ est l'idéal
maximal de $R^h$, et
\item $\kappa = R^h/\mathfrak m R^h$.
\end{enumerate}
\end{lemma}

\begin{proof}
Considérons la catégorie des paires $(S, \mathfrak q)$ où $R \to S$ est un
morphisme d'anneaux \'etale et où $\mathfrak q$ est un idéal premier de $S$
au-dessus de $\mathfrak m$ avec $\kappa = \kappa(\mathfrak q)$. Un morphisme de paires
$(S, \mathfrak q) \to (S', \mathfrak q')$ est donné par un morphisme de
$R$-algèbres $\varphi : S \to S'$ tel que
$\varphi^{-1}(\mathfrak q') = \mathfrak q$.
Posons
$$
R^h = \colim_{(S, \mathfrak q)} S.
$$
Montrons que la catégorie des paires est filtrante, voir
Catégories, définition \ref{categories-definition-directed}.
La catégorie contient la paire $(R, \mathfrak m)$ et n'est donc pas vide,
ce qui prouve la partie (1) de
Catégories, définition \ref{categories-definition-directed}.
Pour toute paire $(S, \mathfrak q)$, l'idéal premier $\mathfrak q$
est maximal de corps résiduel $\kappa$, puisque la composée
$\kappa \to S/\mathfrak q \to \kappa(\mathfrak q)$ est un isomorphisme.
Supposons que $(S, \mathfrak q)$ et $(S', \mathfrak q')$ soient deux objets. Posons
$S'' = S \otimes_R S'$ et
$\mathfrak q'' = \mathfrak qS'' + \mathfrak q'S''$.
Alors $S''/\mathfrak q'' = S/\mathfrak q \otimes_R S'/\mathfrak q' = \kappa$
d'après ce qui précède. De plus, $R \to S''$ est \'etale d'après le
lemme \ref{lemma-etale}.
Cela prouve la partie (2) de
Catégories, définition \ref{categories-definition-directed}.
Supposons ensuite que
$\varphi, \psi : (S, \mathfrak q) \to (S', \mathfrak q')$
soient deux morphismes de paires. Alors $\varphi$, $\psi$, et
$S' \otimes_R S' \to S'$ sont des morphismes d'anneaux \'etales d'après le
lemme \ref{lemma-map-between-etale}.
Considérons
$$
S'' = (S' \otimes_{\varphi, S, \psi} S')
\otimes_{S' \otimes_R S'} S'
$$
avec l'idéal premier
$$
\mathfrak q'' =
(\mathfrak q' \otimes S' + S' \otimes \mathfrak q') \otimes S'
+
(S' \otimes_{\varphi, S, \psi} S') \otimes \mathfrak q'
$$
En raisonnant comme ci-dessus (le changement de base d'un morphisme \'etale est
\'etale, et la composée de morphismes \'etales est \'etale), nous voyons que
$S''$ est \'etale sur $R$. De plus, le morphisme canonique
$S' \to S''$ (utilisant par exemple le facteur le plus à droite)
égalise $\varphi$ et $\psi$. Cela prouve la partie (3) de
Catégories, définition \ref{categories-definition-directed}.
Nous concluons donc que $R^h$ est constitué de triplets
$(S, \mathfrak q, f)$ avec $f \in S$, et que deux tels triplets
$(S, \mathfrak q, f)$, $(S', \mathfrak q', f')$
définissent le même élément de $R^h$ si et seulement s'il existe
une paire $(S'', \mathfrak q'')$ et des morphismes de paires
$\varphi : (S, \mathfrak q) \to (S'', \mathfrak q'')$
et
$\varphi' : (S', \mathfrak q') \to (S'', \mathfrak q'')$
tels que $\varphi(f) = \varphi'(f')$.

\medskip\noindent
Supposons que $x \in R^h$. Représentons $x$ par un triplet
$(S, \mathfrak q, f)$.
Soient $\mathfrak q_1, \ldots, \mathfrak q_r$ les autres idéaux premiers de $S$
au-dessus de $\mathfrak m$. Alors $\mathfrak q \not \subset \mathfrak q_i$,
puisque nous avons vu ci-dessus que $\mathfrak q$ est maximal.
Ainsi, comme $\mathfrak q$ est un idéal premier,
nous pouvons trouver un $g \in S$, $g \not \in \mathfrak q$, tel que
$g \in \mathfrak q_i$ pour $i = 1, \ldots, r$. Considérons le morphisme de
paires $(S, \mathfrak q) \to (S_g, \mathfrak qS_g)$.
De cette manière, nous voyons que nous pouvons toujours supposer que $x$
est donné par un triplet $(S, \mathfrak q, f)$ où
$\mathfrak q$ est le seul idéal premier de $S$ au-dessus de $\mathfrak m$,
c'est-à-dire $\sqrt{\mathfrak mS} = \mathfrak q$. Mais, puisque
$R \to S$ est \'etale, nous avons
$\mathfrak mS_{\mathfrak q} = \mathfrak qS_{\mathfrak q}$, voir le
lemme \ref{lemma-etale-at-prime}.
Nous obtenons donc en fait $\mathfrak mS = \mathfrak q$.

\medskip\noindent
Supposons que $x \not \in \mathfrak mR^h$.
Représentons $x$ par un triplet $(S, \mathfrak q, f)$ avec
$\mathfrak mS = \mathfrak q$.
Alors $f \not \in \mathfrak mS$, c'est-à-dire $f \not \in \mathfrak q$.
Ainsi, $(S, \mathfrak q) \to (S_f, \mathfrak qS_f)$ est un morphisme
de paires tel que l'image de $f$ devienne inversible.
Par conséquent, $x$ est inversible, d'inverse représenté par le triplet
$(S_f, \mathfrak qS_f, 1/f)$. Nous concluons que $R^h$ est un anneau
local d'idéal maximal $\mathfrak mR^h$. Le corps résiduel est
$\kappa$, puisque nous pouvons définir $R^h/\mathfrak mR^h \to \kappa$
en envoyant un triplet $(S, \mathfrak q, f)$ sur la classe
résiduelle de $f$ modulo $\mathfrak q$.

\medskip\noindent
Il nous reste à montrer que $R^h$ est hensélien.
Supposons donc que $P \in R^h[T]$ soit un polynôme
unitaire et que $a_0 \in \kappa$ soit une racine simple de
la réduction $\overline{P} \in \kappa[T]$.
Nous pouvons alors trouver une paire $(S, \mathfrak q)$ telle que
$P$ soit l'image d'un polynôme unitaire $Q \in S[T]$.
Posons $S' = S[T]/(Q)$ et soit $\mathfrak q' \subset S'$
l'idéal maximal $\mathfrak q' = \mathfrak qS' + (T - a')S'$,
où $a' \in S$ est un élément quelconque relevant $a_0$.
Par construction, $S \to S'$ est \'etale en $\mathfrak q'$
et $\kappa = \kappa(\mathfrak q')$.
Choisissons $g \in S'$, $g \not \in \mathfrak q'$, tel que
$S'' = S'_g$ soit \'etale sur $S$. Alors
$(S, \mathfrak q) \to (S'', \mathfrak q'S'')$ est un morphisme
de paires. Le triplet $(S'', \mathfrak q'S'', \text{classe de }T)$
détermine maintenant un élément $a \in R^h$ ayant les propriétés
$P(a) = 0$ et $\overline{a} = a_0$, comme désiré.
\end{proof}

\begin{lemma}
\label{lemma-strict-henselization}
Soit $(R, \mathfrak m, \kappa)$ un anneau local.
Soit $\kappa \subset \kappa^{sep}$ une clôture algébrique séparable.
Il existe un diagramme commutatif
$$
\xymatrix{
\kappa \ar[r] & \kappa \ar[r] & \kappa^{sep} \\
R \ar[r] \ar[u] & R^h \ar[r] \ar[u] & R^{sh} \ar[u]
}
$$
ayant les propriétés suivantes :
\begin{enumerate}
\item le morphisme $R^h \to R^{sh}$ est local,
\item $R^{sh}$ est strictement hensélien,
\item $R^{sh}$ est une colimite filtrante de $R$-algèbres \'etales,
\item $\mathfrak m R^{sh}$ est l'idéal
maximal de $R^{sh}$, et
\item $\kappa^{sep} = R^{sh}/\mathfrak m R^{sh}$.
\end{enumerate}
\end{lemma}

\begin{proof}
Cela se démontre exactement comme le
lemme \ref{lemma-henselization}.
La seule différence est qu'au lieu de paires, on utilise des triplets
$(S, \mathfrak q, \alpha)$, où $R \to S$ est \'etale,
$\mathfrak q$ est un idéal premier de $S$ au-dessus de $\mathfrak m$, et
$\alpha : \kappa(\mathfrak q) \to \kappa^{sep}$ est un plongement
d'extensions de $\kappa$.
\end{proof}

\begin{definition}
\label{definition-henselization}
Soit $(R, \mathfrak m, \kappa)$ un anneau local.
\begin{enumerate}
\item Le morphisme local d'anneaux $R \to R^h$ construit dans le
lemme \ref{lemma-henselization}
est appelé l'{\it hensélisation} de $R$.
\item Étant donnée une clôture algébrique séparable
$\kappa \subset \kappa^{sep}$, le morphisme local d'anneaux
$R \to R^{sh}$ construit dans le
lemme \ref{lemma-strict-henselization}
est appelé l'{\it hensélisation stricte de $R$ relativement à
$\kappa \subset \kappa^{sep}$}.
\item Un morphisme local d'anneaux $R \to R^{sh}$ est appelé une
{\it hensélisation stricte} de $R$ s'il est isomorphe à l'un des
morphismes locaux d'anneaux construits dans le
lemme \ref{lemma-strict-henselization}.
\end{enumerate}
\end{definition}

\noindent
Les morphismes $R \to R^h \to R^{sh}$ sont des morphismes plats d'anneaux locaux.
D'après le lemme \ref{lemma-uniqueness-henselian}, les $R$-algèbres
$R^h$ et $R^{sh}$ sont bien définies à isomorphisme unique près par les conditions
qu'elles soient locales henséliennes, colimites filtrantes de $R$-algèbres \'etales,
de corps résiduels respectifs $\kappa$ et $\kappa^{sep}$.
Dans le reste de cette section, nous examinons principalement la fonctorialité des
hensélisations (strictes).
Nous examinerons des résultats plus subtils concernant
la relation entre $R$ et son hensélisation dans
Plus sur l'algèbre, section \ref{more-algebra-section-permanence-henselization}.

\begin{remark}
\label{remark-construct-sh-from-h}
Nous pouvons aussi construire $R^{sh}$ à partir de $R^h$. En effet, pour toute
sous-extension finie séparable $\kappa^{sep}/\kappa'/\kappa$,
il existe une unique extension locale finie étale d'anneaux
$R^h \subset R^h(\kappa')$ (à isomorphisme unique près)
dont l'extension de corps résiduels reproduit l'extension donnée, voir le
lemme \ref{lemma-henselian-cat-finite-etale}.
Nous pouvons donc poser
$$
R^{sh} =
\bigcup\nolimits_{\kappa \subset \kappa' \subset \kappa^{sep}}
R^h(\kappa')
$$
Les flèches de ce système, compatibles avec les flèches au niveau
des corps résiduels, existent d'après le
lemme \ref{lemma-henselian-cat-finite-etale}.
Cela produit un anneau local hensélien d'après le
lemme \ref{lemma-colimit-henselian},
puisque chacun des anneaux
$R^h(\kappa')$ est hensélien d'après le
lemme \ref{lemma-finite-over-henselian}.
Par construction, l'extension de corps résiduels induite par
$R^h \to R^{sh}$ est l'extension de corps $\kappa^{sep}/\kappa$.
Ainsi, $R^{sh}$ ainsi construit est strictement hensélien.
D'après le lemme \ref{lemma-composition-colimit-etale}, la $R$-algèbre
$R^{sh}$ est une colimite d'algèbres \'etales sur $R$. Ainsi, l'unicité du
lemme \ref{lemma-uniqueness-henselian} montre que $R^{sh}$
est l'hensélisation stricte.
\end{remark}

\begin{lemma}
\label{lemma-henselian-functorial-prepare}
Soit $R \to S$ un morphisme local d'anneaux locaux.
Soit $S \to S^h$ l'hensélisation.
Soit $R \to A$ un morphisme d'anneaux \'etale, et soit $\mathfrak q$
un idéal premier de $A$ au-dessus de $\mathfrak m_R$
tel que $R/\mathfrak m_R \cong \kappa(\mathfrak q)$.
Alors il existe un unique morphisme d'anneaux
$f : A \to S^h$ s'insérant dans le diagramme commutatif
$$
\xymatrix{
A \ar[r]_f & S^h \\
R \ar[u] \ar[r] & S \ar[u]
}
$$
tel que $f^{-1}(\mathfrak m_{S^h}) = \mathfrak q$.
\end{lemma}

\begin{proof}
C'est un cas particulier du lemme \ref{lemma-map-into-henselian}.
\end{proof}

\begin{lemma}
\label{lemma-henselian-functorial}
Soit $R \to S$ un morphisme local d'anneaux locaux.
Soient $R \to R^h$ et $S \to S^h$ les hensélisations.
Il existe un unique morphisme local d'anneaux $R^h \to S^h$ s'insérant
dans le diagramme commutatif
$$
\xymatrix{
R^h \ar[r]_f & S^h \\
R \ar[u] \ar[r] & S \ar[u]
}
$$
\end{lemma}

\begin{proof}
Cela découle immédiatement du lemme \ref{lemma-map-into-henselian-colimit}.
\end{proof}

\noindent
Voici une construction légèrement différente de l'hensélisation.

\begin{lemma}
\label{lemma-henselization-different}
Soit $R$ un anneau.
Soit $\mathfrak p \subset R$ un idéal premier.
Considérons la catégorie des paires $(S, \mathfrak q)$ où
$R \to S$ est \'etale et $\mathfrak q$ est un idéal premier au-dessus de $\mathfrak p$
tel que $\kappa(\mathfrak p) = \kappa(\mathfrak q)$.
Cette catégorie est filtrante et
$$
(R_{\mathfrak p})^h = \colim_{(S, \mathfrak q)} S
= \colim_{(S, \mathfrak q)} S_{\mathfrak q}
$$
canoniquement.
\end{lemma}

\begin{proof}
Un morphisme de paires $(S, \mathfrak q) \to (S', \mathfrak q')$
est donné par un morphisme de $R$-algèbres $\varphi : S \to S'$ tel que
$\varphi^{-1}(\mathfrak q') = \mathfrak q$.
Montrons que la catégorie des paires est filtrante, voir
Catégories, définition \ref{categories-definition-directed}.
La catégorie contient la paire $(R, \mathfrak p)$ et n'est donc pas vide,
ce qui prouve la partie (1) de
Catégories, définition \ref{categories-definition-directed}.
Supposons que $(S, \mathfrak q)$ et $(S', \mathfrak q')$ soient deux paires.
Notons que $\mathfrak q$, resp.\ $\mathfrak q'$ correspondent à des idéaux premiers
des anneaux fibres $S \otimes \kappa(\mathfrak p)$,
resp.\ $S' \otimes \kappa(\mathfrak p)$, de corps résiduels
$\kappa(\mathfrak p)$ ; ils correspondent donc à des idéaux maximaux de
$S \otimes \kappa(\mathfrak p)$, resp.\ $S' \otimes \kappa(\mathfrak p)$.
Posons $S'' = S \otimes_R S'$. D'après ce qui précède, il existe un unique
idéal premier $\mathfrak q'' \subset S''$ au-dessus de $\mathfrak q$ et de
$\mathfrak q'$ dont le corps résiduel est $\kappa(\mathfrak p)$.
Le morphisme d'anneaux $R \to S''$ est \'etale d'après le
lemme \ref{lemma-etale}.
Cela prouve la partie (2) de
Catégories, définition \ref{categories-definition-directed}.
Supposons ensuite que
$\varphi, \psi : (S, \mathfrak q) \to (S', \mathfrak q')$
soient deux morphismes de paires. Alors $\varphi$, $\psi$, et
$S' \otimes_R S' \to S'$ sont des morphismes d'anneaux \'etales d'après le
lemme \ref{lemma-map-between-etale}. Considérons
$$
S'' = (S' \otimes_{\varphi, S, \psi} S')
\otimes_{S' \otimes_R S'} S'
$$
En raisonnant comme ci-dessus (le changement de base d'un morphisme \'etale est
\'etale, et la composée de morphismes \'etales est \'etale), nous voyons que $S''$ est \'etale sur $R$. L'anneau
fibre de $S''$ au-dessus de $\mathfrak p$ est
$$
F'' = (F' \otimes_{\varphi, F, \psi} F')
\otimes_{F' \otimes_{\kappa(\mathfrak p)} F'} F'
$$
où $F', F$ sont les anneaux fibres de $S'$ et $S$. Puisque $\varphi$ et
$\psi$ sont des morphismes de paires, le morphisme $F' \to \kappa(\mathfrak p)$
correspondant à $\mathfrak p'$ se prolonge en un morphisme $F'' \to \kappa(\mathfrak p)$
et correspond à son tour à un idéal premier $\mathfrak q'' \subset S''$
dont le corps résiduel est $\kappa(\mathfrak p)$.
Le morphisme canonique $S' \to S''$ (utilisant par exemple le facteur le plus à droite)
est un morphisme de paires $(S', \mathfrak q') \to (S'', \mathfrak q'')$
qui égalise $\varphi$ et $\psi$. Cela prouve la partie (3) de
Catégories, définition \ref{categories-definition-directed}.
Nous concluons donc que la catégorie est filtrante.

\medskip\noindent
Rappelons que, dans la preuve du
lemme \ref{lemma-henselization},
nous avons construit $(R_{\mathfrak p})^h$ comme la colimite correspondante,
mais en partant de $R_{\mathfrak p}$ et de son idéal maximal
$\mathfrak pR_{\mathfrak p}$. Or, étant donnée une paire $(S, \mathfrak q)$
pour $(R, \mathfrak p)$, nous obtenons une paire
$(S_{\mathfrak p}, \mathfrak qS_{\mathfrak p})$ pour
$(R_{\mathfrak p}, \mathfrak pR_{\mathfrak p})$.
De plus, dans cette situation,
$$
S_{\mathfrak p} = \colim_{f \in R, f \not \in \mathfrak p} S_f.
$$
Ainsi, pour montrer les égalités
du lemme, il suffit de montrer que toute paire $(S_{loc}, \mathfrak q_{loc})$
pour $(R_{\mathfrak p}, \mathfrak pR_{\mathfrak p})$ est de la forme
$(S_{\mathfrak p}, \mathfrak qS_{\mathfrak p})$ pour une certaine paire
$(S, \mathfrak q)$ au-dessus de $(R, \mathfrak p)$ (certains détails sont omis).
Cela découle du
lemme \ref{lemma-etale}.
\end{proof}

\begin{lemma}
\label{lemma-henselian-functorial-improve}
Soit $R \to S$ un morphisme d'anneaux. Soit $\mathfrak q \subset S$ un idéal premier au-dessus
de $\mathfrak p \subset R$. Soient $R \to R^h$ et $S \to S^h$ les
hensélisations de $R_\mathfrak p$ et $S_\mathfrak q$. Le morphisme local d'anneaux
$R^h \to S^h$ du lemme \ref{lemma-henselian-functorial} identifie $S^h$
à l'hensélisation de $R^h \otimes_R S$ en l'unique idéal premier
au-dessus de $\mathfrak m^h$ et de $\mathfrak q$.
\end{lemma}

\begin{proof}
Le lemme \ref{lemma-henselization-different} montre que $R^h$, resp.\ $S^h$
sont des colimites filtrantes de $R$-algèbres, resp.\ de $S$-algèbres, \'etales.
Ainsi, $R^h \otimes_R S$ est une colimite filtrante de
$S$-algèbres \'etales $A_i$ (lemme \ref{lemma-etale}). D'après le
lemme \ref{lemma-colimits-of-etale}, $S^h$ est une
colimite filtrante de $R^h \otimes_R S$-algèbres \'etales.
Comme, de plus, $S^h$ est un anneau local hensélien de corps résiduel
égal à $\kappa(\mathfrak q)$, l'assertion découle du résultat d'unicité du
lemme \ref{lemma-uniqueness-henselian}.
\end{proof}

\begin{lemma}
\label{lemma-strictly-henselian-functorial-prepare}
Soit $\varphi : R \to S$ un morphisme local d'anneaux locaux.
Soit $S/\mathfrak m_S \subset \kappa^{sep}$ une clôture algébrique séparable.
Soit $S \to S^{sh}$ l'hensélisation stricte de $S$
relativement à $S/\mathfrak m_S \subset \kappa^{sep}$.
Soit $R \to A$ un morphisme d'anneaux \'etale, et soit $\mathfrak q$
un idéal premier de $A$ au-dessus de $\mathfrak m_R$.
Étant donné un diagramme commutatif
$$
\xymatrix{
\kappa(\mathfrak q) \ar[r]_{\phi} & \kappa^{sep} \\
R/\mathfrak m_R \ar[r]^{\varphi} \ar[u] & S/\mathfrak m_S \ar[u]
}
$$
il existe un unique morphisme d'anneaux
$f : A \to S^{sh}$ s'insérant dans le diagramme commutatif
$$
\xymatrix{
A \ar[r]_f & S^{sh} \\
R \ar[u] \ar[r]^{\varphi} & S \ar[u]
}
$$
tel que $f^{-1}(\mathfrak m_{S^h}) = \mathfrak q$ et que le
morphisme induit $\kappa(\mathfrak q) \to \kappa^{sep}$ soit celui qui est donné.
\end{lemma}

\begin{proof}
C'est un cas particulier du lemme \ref{lemma-map-into-henselian}.
\end{proof}

\begin{lemma}
\label{lemma-strictly-henselian-functorial}
Soit $R \to S$ un morphisme local d'anneaux locaux.
Choisissons des clôtures algébriques séparables
$R/\mathfrak m_R \subset \kappa_1^{sep}$
et
$S/\mathfrak m_S \subset \kappa_2^{sep}$.
Soient $R \to R^{sh}$ et $S \to S^{sh}$ les hensélisations
strictes correspondantes. Étant donné un diagramme commutatif
$$
\xymatrix{
\kappa_1^{sep} \ar[r]_{\phi} & \kappa_2^{sep} \\
R/\mathfrak m_R \ar[r]^{\varphi} \ar[u] & S/\mathfrak m_S \ar[u]
}
$$
il existe un unique morphisme local d'anneaux $R^{sh} \to S^{sh}$ s'insérant
dans le diagramme commutatif
$$
\xymatrix{
R^{sh} \ar[r]_f & S^{sh} \\
R \ar[u] \ar[r] & S \ar[u]
}
$$
et induisant $\phi$ sur les corps résiduels de
$R^{sh}$ et $S^{sh}$.
\end{lemma}

\begin{proof}
Cela découle immédiatement du lemme \ref{lemma-map-into-henselian-colimit}.
\end{proof}

\begin{lemma}
\label{lemma-strict-henselization-different}
Soit $R$ un anneau.
Soit $\mathfrak p \subset R$ un idéal premier.
Soit $\kappa(\mathfrak p) \subset \kappa^{sep}$ une
clôture algébrique séparable.
Considérons la catégorie des triplets $(S, \mathfrak q, \phi)$
où $R \to S$ est \'etale, $\mathfrak q$ est un idéal premier au-dessus de $\mathfrak p$,
et $\phi : \kappa(\mathfrak q) \to \kappa^{sep}$ est un morphisme de
$\kappa(\mathfrak p)$-algèbres. Cette catégorie est filtrante et
$$
(R_{\mathfrak p})^{sh} =
\colim_{(S, \mathfrak q, \phi)} S =
\colim_{(S, \mathfrak q, \phi)} S_{\mathfrak q}
$$
canoniquement.
\end{lemma}

\begin{proof}
Un morphisme de triplets $(S, \mathfrak q, \phi) \to (S', \mathfrak q', \phi')$
est donné par un morphisme de $R$-algèbres $\varphi : S \to S'$ tel que
$\varphi^{-1}(\mathfrak q') = \mathfrak q$ et que
$\phi' \circ \varphi = \phi$.
Montrons que la catégorie des triplets est filtrante, voir
Catégories, définition \ref{categories-definition-directed}.
La catégorie contient le triplet
$(R, \mathfrak p, \kappa(\mathfrak p) \subset \kappa^{sep})$
et n'est donc pas vide, ce qui prouve la partie (1) de
Catégories, définition \ref{categories-definition-directed}.
Supposons que $(S, \mathfrak q, \phi)$ et $(S', \mathfrak q', \phi')$
soient deux triplets.
Notons que $\mathfrak q$, resp.\ $\mathfrak q'$ correspondent à des idéaux premiers
des anneaux fibres $S \otimes \kappa(\mathfrak p)$,
resp.\ $S' \otimes \kappa(\mathfrak p)$, dont les corps résiduels sont
des extensions séparables finies de $\kappa(\mathfrak p)$, et que $\phi$, resp.\ $\phi'$,
correspondent à des morphismes vers $\kappa^{sep}$. Ces données correspondent donc à des
morphismes de $\kappa(\mathfrak p)$-algèbres
$$
\phi : S \otimes_R \kappa(\mathfrak p) \longrightarrow \kappa^{sep},
\quad
\phi' : S' \otimes_R \kappa(\mathfrak p) \longrightarrow \kappa^{sep}.
$$
Posons $S'' = S \otimes_R S'$. En combinant les morphismes ci-dessus, nous obtenons un unique
morphisme de $\kappa(\mathfrak p)$-algèbres
$$
\phi'' = \phi \otimes \phi' :
S'' \otimes_R \kappa(\mathfrak p)
\longrightarrow
\kappa^{sep}
$$
dont le noyau correspond à un idéal premier $\mathfrak q'' \subset S''$
au-dessus de $\mathfrak q$ et de $\mathfrak q'$, et dont le corps résiduel
s'envoie par $\phi''$ dans le corps composé de
$\phi(\kappa(\mathfrak q))$ et $\phi'(\kappa(\mathfrak q'))$ dans
$\kappa^{sep}$. Le morphisme d'anneaux $R \to S''$ est \'etale d'après le
lemme \ref{lemma-etale}.
Ainsi, $(S'', \mathfrak q'', \phi'')$ est un triplet dominant à la fois
$(S, \mathfrak q, \phi)$ et $(S', \mathfrak q', \phi')$.
Cela prouve la partie (2) de
Catégories, définition \ref{categories-definition-directed}.
Supposons ensuite que
$\varphi, \psi : (S, \mathfrak q, \phi) \to (S', \mathfrak q', \phi')$
soient deux morphismes de triplets. Alors $\varphi$, $\psi$, et
$S' \otimes_R S' \to S'$ sont des morphismes d'anneaux \'etales d'après le
lemme \ref{lemma-map-between-etale}.
Considérons
$$
S'' = (S' \otimes_{\varphi, S, \psi} S')
\otimes_{S' \otimes_R S'} S'
$$
En raisonnant comme ci-dessus (le changement de base d'un morphisme \'etale est
\'etale, et la composée de morphismes \'etales est \'etale), nous voyons que $S''$ est \'etale sur $R$. L'anneau
fibre de $S''$ au-dessus de $\mathfrak p$ est
$$
F'' = (F' \otimes_{\varphi, F, \psi} F')
\otimes_{F' \otimes_{\kappa(\mathfrak p)} F'} F'
$$
où $F', F$ sont les anneaux fibres de $S'$ et $S$. Puisque $\varphi$ et
$\psi$ sont des morphismes de triplets, le morphisme $\phi' : F' \to \kappa^{sep}$
se prolonge en un morphisme $\phi'' : F'' \to \kappa^{sep}$
qui correspond à son tour à un idéal premier $\mathfrak q'' \subset S''$.
Le morphisme canonique $S' \to S''$ (utilisant par exemple le facteur le plus à droite)
est un morphisme de triplets
$(S', \mathfrak q', \phi') \to (S'', \mathfrak q'', \phi'')$
qui égalise $\varphi$ et $\psi$. Cela prouve la partie (3) de
Catégories, définition \ref{categories-definition-directed}.
Nous concluons donc que la catégorie est filtrante.

\medskip\noindent
Il reste à montrer que la colimite $R_{colim}$ du système
est égale à l'hensélisation stricte
de $R_{\mathfrak p}$ relativement à $\kappa^{sep}$. Pour cela, notons que
le système des triplets $(S, \mathfrak q, \phi)$ contient comme sous-système
les paires $(S, \mathfrak q)$ du
lemme \ref{lemma-henselization-different}.
Ainsi, $R_{colim}$ contient $R_{\mathfrak p}^h$ d'après ce lemme.
De plus, il est clair que $R_{\mathfrak p}^h \subset R_{colim}$
est une colimite filtrante d'extensions d'anneaux \'etales.
Il s'ensuit que $R_{colim}$ est hensélien d'après les
lemmes \ref{lemma-finite-over-henselian} et
\ref{lemma-colimit-henselian}.
Enfin, d'après le
lemme \ref{lemma-make-etale-map-prescribed-residue-field},
le corps résiduel de $R_{colim}$ est égal à
$\kappa^{sep}$. Nous concluons donc que $R_{colim}$ est strictement hensélien
et est ainsi égal à l'hensélisation stricte de $R_{\mathfrak p}$, comme souhaité.
Certains détails sont omis.
\end{proof}

\begin{lemma}
\label{lemma-strictly-henselian-functorial-improve}
Soit $R \to S$ un morphisme d'anneaux. Soit $\mathfrak q \subset S$ un idéal premier au-dessus
de $\mathfrak p \subset R$. Choisissons des clôtures algébriques séparables
$\kappa(\mathfrak p) \subset \kappa_1^{sep}$
et
$\kappa(\mathfrak q) \subset \kappa_2^{sep}$.
Soient $R^{sh}$ et $S^{sh}$ les hensélisations strictes
correspondantes de $R_\mathfrak p$ et $S_\mathfrak q$.
Étant donné un diagramme commutatif
$$
\xymatrix{
\kappa_1^{sep} \ar[r]_{\phi} & \kappa_2^{sep} \\
\kappa(\mathfrak p) \ar[r]^{\varphi} \ar[u] & \kappa(\mathfrak q) \ar[u]
}
$$
le morphisme local d'anneaux $R^{sh} \to S^{sh}$ du
lemme \ref{lemma-strictly-henselian-functorial} identifie $S^{sh}$
à l'hensélisation stricte de $R^{sh} \otimes_R S$ en un idéal premier
au-dessus de $\mathfrak q$ et de l'idéal maximal
$\mathfrak m^{sh} \subset R^{sh}$.
\end{lemma}

\begin{proof}
La preuve est identique à celle du
lemme \ref{lemma-henselian-functorial-improve},
à ceci près qu'elle utilise le
lemme \ref{lemma-strict-henselization-different}
au lieu du
lemme \ref{lemma-henselization-different}.
\end{proof}

\begin{lemma}
\label{lemma-sh-from-h-map}
Soit $R \to S$ un morphisme d'anneaux. Soit $\mathfrak q \subset S$ un idéal premier
au-dessus de $\mathfrak p \subset R$ tel que
$\kappa(\mathfrak p) \to \kappa(\mathfrak q)$ soit un isomorphisme.
Choisissons une clôture algébrique séparable $\kappa^{sep}$ de
$\kappa(\mathfrak p) = \kappa(\mathfrak q)$.
Alors
$$
(S_\mathfrak q)^{sh} =
(S_\mathfrak q)^h \otimes_{(R_\mathfrak p)^h} (R_\mathfrak p)^{sh}
$$
\end{lemma}

\begin{proof}
Cela découle de l'autre construction de l'hensélisation stricte
d'un anneau local donnée dans la remarque \ref{remark-construct-sh-from-h}, ainsi que du
fait que les corps résiduels sont égaux. Certains détails sont omis.
\end{proof}














\section{Hensélisation et morphismes d'anneaux quasi-finis}
\label{section-henselization-quasi-finite}

\noindent
Dans cette section, nous démontrons quelques résultats concernant les morphismes fonctoriels
entre hensélisés (stricts) pour les morphismes d'anneaux quasi-finis.

\begin{lemma}
\label{lemma-quasi-finite-henselization}
Soit $R \to S$ un morphisme d'anneaux.
Soit $\mathfrak q$ un idéal premier de $S$ au-dessus de $\mathfrak p$ dans $R$.
Supposons que $R \to S$ soit quasi-fini en $\mathfrak q$.
Le diagramme commutatif
$$
\xymatrix{
R_{\mathfrak p}^h \ar[r] & S_{\mathfrak q}^h \\
R_{\mathfrak p} \ar[u] \ar[r] & S_{\mathfrak q} \ar[u]
}
$$
du
lemme \ref{lemma-henselian-functorial}
identifie $S_{\mathfrak q}^h$ à la localisation de
$R_{\mathfrak p}^h \otimes_{R_{\mathfrak p}} S_{\mathfrak q}$
en l'idéal premier engendré par $\mathfrak q$. De plus, le morphisme
d'anneaux $R_{\mathfrak p}^h \to S_{\mathfrak q}^h$ est fini.
\end{lemma}

\begin{proof}
Notons que $R_{\mathfrak p}^h \otimes_R S$ est quasi-fini sur
$R_{\mathfrak p}^h$ en l'idéal premier correspondant à $\mathfrak q$, voir le
lemme \ref{lemma-four-rings}. Ainsi, la localisation $S'$ de
$R_{\mathfrak p}^h \otimes_{R_{\mathfrak p}} S_{\mathfrak q}$ est hensélienne
et finie sur $R_{\mathfrak p}^h$, voir le
lemme \ref{lemma-finite-over-henselian}. En tant que localisation, $S'$ est une limite
inductive filtrante d'algèbres \'etales sur
$R_{\mathfrak p}^h \otimes_{R_{\mathfrak p}} S_{\mathfrak q}$.
Le lemme \ref{lemma-henselian-functorial-improve} montre que
$S_\mathfrak q^h$ est le hensélisé de
$R_{\mathfrak p}^h \otimes_{R_{\mathfrak p}} S_{\mathfrak q}$.
Ainsi, $S' = S_\mathfrak q^h$ d'après le résultat d'unicité du
lemme \ref{lemma-uniqueness-henselian}.
\end{proof}

\begin{lemma}
\label{lemma-quotient-henselization}
\begin{slogan}
L'hensélisation est compatible au passage au quotient.
\end{slogan}
Soit $R$ un anneau local et soit $R^h$ son hensélisé.
Soit $I \subset \mathfrak m_R$.
Alors $R^h/IR^h$ est le hensélisé de $R/I$.
\end{lemma}

\begin{proof}
C'est un cas particulier du
lemme \ref{lemma-quasi-finite-henselization}.
\end{proof}

\begin{lemma}
\label{lemma-quasi-finite-strict-henselization}
Soit $R \to S$ un morphisme d'anneaux.
Soit $\mathfrak q$ un idéal premier de $S$ au-dessus de $\mathfrak p$ dans $R$.
Supposons que $R \to S$ soit quasi-fini en $\mathfrak q$.
Soit $\kappa_2^{sep}/\kappa(\mathfrak q)$ une clôture séparable,
et notons $\kappa_1^{sep} \subset \kappa_2^{sep}$ le sous-corps
des éléments algébriques séparables sur $\kappa(\mathfrak p)$
(Corps, lemme \ref{fields-lemma-separable-first}).
Le diagramme commutatif
$$
\xymatrix{
R_{\mathfrak p}^{sh} \ar[r] & S_{\mathfrak q}^{sh} \\
R_{\mathfrak p} \ar[u] \ar[r] & S_{\mathfrak q} \ar[u]
}
$$
du lemme \ref{lemma-strictly-henselian-functorial}
identifie $S_{\mathfrak q}^{sh}$ à la localisation de
$R_{\mathfrak p}^{sh} \otimes_{R_{\mathfrak p}} S_{\mathfrak q}$
en l'idéal premier qui est le noyau du morphisme
$$
R_{\mathfrak p}^{sh} \otimes_{R_{\mathfrak p}} S_{\mathfrak q}
\longrightarrow
\kappa_1^{sep} \otimes_{\kappa(\mathfrak p)} \kappa(\mathfrak q)
\longrightarrow
\kappa_2^{sep}
$$
De plus, le morphisme d'anneaux $R_{\mathfrak p}^{sh} \to S_{\mathfrak q}^{sh}$
est un homomorphisme local et fini d'anneaux locaux dont l'extension de corps résiduels
est l'extension $\kappa_2^{sep}/\kappa_1^{sep}$, laquelle est à la fois finie et
purement inséparable.
\end{lemma}

\begin{proof}
Comme $R \to S$ est quasi-fini en $\mathfrak q$, l'extension
$\kappa(\mathfrak q)/\kappa(\mathfrak p)$ est finie, voir la
définition \ref{definition-quasi-finite} et le
lemme \ref{lemma-isolated-point-fibre}.
Ainsi, $\kappa_1^{sep}$ est une clôture séparable de
$\kappa(\mathfrak p)$ (détail mineur omis).
En particulier, le lemme \ref{lemma-strictly-henselian-functorial}
s'applique bien. Ensuite, le corps composé de $\kappa(\mathfrak q)$ et de
$\kappa_1^{sep}$ dans $\kappa_2^{sep}$ est séparablement
clos, et est donc égal à $\kappa_2^{sep}$. Nous concluons que
$\kappa_2^{sep}/\kappa_1^{sep}$ est finie. Par construction,
l'extension $\kappa_2^{sep}/\kappa_1^{sep}$ est purement inséparable.
Le morphisme d'anneaux $R_{\mathfrak p}^{sh} \to S_{\mathfrak q}^{sh}$ est
bien local et induit l'extension de corps résiduels
$\kappa_2^{sep}/\kappa_1^{sep}$, qui est bien finie et purement inséparable.

\medskip\noindent
Notons que $R_{\mathfrak p}^{sh} \otimes_R S$ est quasi-fini sur
$R_{\mathfrak p}^{sh}$ en l'idéal premier $\mathfrak q'$
donné dans l'énoncé du lemme, voir le
lemme \ref{lemma-four-rings}. Ainsi, la localisation $S'$ de
$R_{\mathfrak p}^{sh} \otimes_{R_{\mathfrak p}} S_{\mathfrak q}$
en $\mathfrak q'$ est hensélienne et finie sur $R_{\mathfrak p}^{sh}$, voir le
lemme \ref{lemma-finite-over-henselian}.
Notons que le corps résiduel de $S'$ est $\kappa_2^{sep}$,
car le morphisme $\kappa_1^{sep} \otimes_{\kappa(\mathfrak p)} \kappa(\mathfrak q)
\to \kappa_2^{sep}$ est surjectif d'après la discussion du paragraphe
précédent.
En outre, en tant que localisation, $S'$ est une limite inductive filtrante d'algèbres \'etales sur
$R_{\mathfrak p}^{sh} \otimes_{R_{\mathfrak p}} S_{\mathfrak q}$.
Le lemme \ref{lemma-strictly-henselian-functorial-improve}
montre que $S_{\mathfrak q}^{sh}$ est le hensélisé strict de
$R_{\mathfrak p}^{sh} \otimes_{R_{\mathfrak p}} S_{\mathfrak q}$
en $\mathfrak q'$. Ainsi, $S' = S_\mathfrak q^{sh}$ d'après le résultat d'unicité du
lemme \ref{lemma-uniqueness-henselian}.
\end{proof}

\begin{lemma}
\label{lemma-quotient-strict-henselization}
Soit $R$ un anneau local et soit $R^{sh}$ son hensélisé strict.
Soit $I \subset \mathfrak m_R$.
Alors $R^{sh}/IR^{sh}$ est un hensélisé strict de $R/I$.
\end{lemma}

\begin{proof}
C'est un cas particulier du
lemme \ref{lemma-quasi-finite-strict-henselization}.
\end{proof}

\begin{lemma}
\label{lemma-local-tensor-with-integral}
Soient $A \to B$ et $A \to C$ des homomorphismes locaux d'anneaux locaux.
Si $A \to C$ est entier et si l'une des extensions
$\kappa(\mathfrak m_C)/\kappa(\mathfrak m_A)$ ou
$\kappa(\mathfrak m_B)/\kappa(\mathfrak m_A)$ est purement
inséparable, alors $D = B \otimes_A C$ est un anneau local et
$B \to D$ et $C \to D$ sont locaux.
\end{lemma}

\begin{proof}
Tout idéal maximal de $D$ est au-dessus de l'idéal maximal de $B$ par le
théorème de montée pour le morphisme d'anneaux entier
$B \to D$ (lemme \ref{lemma-integral-going-up}).
Or $D/\mathfrak m_B D = \kappa(\mathfrak m_B) \otimes_A C =
\kappa(\mathfrak m_B) \otimes_{\kappa(\mathfrak m_A)} C/\mathfrak m_A C$.
Le spectre de $C/\mathfrak m_A C$ est réduit à un
seul point, à savoir $\mathfrak m_C$. Ainsi, le spectre de
$D/\mathfrak m_B D$ est le même que celui de
$\kappa(\mathfrak m_B) \otimes_{\kappa(\mathfrak m_A)} \kappa(\mathfrak m_C)$,
qui est réduit à un seul point puisque, par hypothèse, l'une des extensions
$\kappa(\mathfrak m_C)/\kappa(\mathfrak m_A)$ ou
$\kappa(\mathfrak m_B)/\kappa(\mathfrak m_A)$ est purement
inséparable. Cela prouve que $D$ est local et que les morphismes d'anneaux
$B \to D$ et $C \to D$ sont locaux.
\end{proof}

\begin{lemma}
\label{lemma-base-change-strict-henselization-quasi-finite}
Soient $A \to B$ et $A \to C$ des morphismes d'anneaux. Soit $\kappa$ un corps
séparablement clos et soit $B \otimes_A C \to \kappa$
un homomorphisme d'anneaux. Notons
$$
\xymatrix{
B^{sh} \ar[r] & (B \otimes_A C)^{sh} \\
A^{sh} \ar[u] \ar[r] & C^{sh} \ar[u]
}
$$
les morphismes correspondants entre hensélisés stricts (voir la démonstration). Si
\begin{enumerate}
\item $A \to B$ est quasi-fini en l'idéal premier
$\mathfrak p_B = \Ker(B \to \kappa)$, ou
\item $B$ est une limite inductive filtrante de $A$-algèbres quasi-finies, ou
\item $B_{\mathfrak p_B}$ est une limite inductive filtrante d'algèbres quasi-finies
sur $A_{\mathfrak p_A}$, ou
\item $B$ est entier sur $A$,
\end{enumerate}
alors $B^{sh} \otimes_{A^{sh}} C^{sh} \to (B \otimes_A C)^{sh}$
est un isomorphisme.
\end{lemma}

\begin{proof}
Posons $D = B \otimes_A C$. Notons
$\mathfrak p_A = \Ker(A \to \kappa)$ et de même
pour $\mathfrak p_B$, $\mathfrak p_C$ et $\mathfrak p_D$.
Notons $\kappa_A \subset \kappa$ la clôture séparable
de $\kappa(\mathfrak p_A)$ dans $\kappa$
et de même pour $\kappa_B$, $\kappa_C$ et $\kappa_D$.
Notons $A^{sh}$ le hensélisé strict de $A_{\mathfrak p_A}$
construit à l'aide de la clôture séparable
$\kappa_A/\kappa(\mathfrak p_A)$. Faisons de même pour
$B^{sh}$, $C^{sh}$ et $D^{sh}$. Nous obtenons le diagramme commutatif
du lemme par fonctorialité du
lemme \ref{lemma-strictly-henselian-functorial}.

\medskip\noindent
Considérons le morphisme
$$
c : B^{sh} \otimes_{A^{sh}} C^{sh} \to D^{sh} = (B \otimes_A C)^{sh}
$$
que fournit le diagramme commutatif. Si $A \to B$ est quasi-fini en
$\mathfrak p_B = \Ker(B \to \kappa)$, alors le morphisme d'anneaux $C \to D$ est
quasi-fini en $\mathfrak p_D$ d'après le lemme \ref{lemma-four-rings}. Ainsi, d'après le
lemme \ref{lemma-quasi-finite-strict-henselization}
(et le lemme \ref{lemma-base-change-integral}),
le morphisme d'anneaux $c$ est un homomorphisme d'algèbres finies sur $C^{sh}$ et
$$
B^{sh} = (B \otimes_A A^{sh})_{\mathfrak q}
\quad\text{et}\quad
D^{sh} = (D \otimes_C C^{sh})_{\mathfrak r} =
(B \otimes_A C^{sh})_{\mathfrak r}
$$
pour certains idéaux premiers $\mathfrak q$ et $\mathfrak r$. Comme
$$
B^{sh} \otimes_{A^{sh}} C^{sh} =
(B \otimes_A A^{sh})_{\mathfrak q} \otimes_{A^{sh}} C^{sh} =
\text{une localisation de }
B \otimes_A C^{sh}
$$
nous concluons que la source et le but de $c$ sont tous deux
des localisations de $B \otimes_A C^{sh}$ (de manière compatible avec le morphisme).
Il suffit donc de montrer que $B^{sh} \otimes_{A^{sh}} C^{sh}$
est local (détail mineur omis). Cela résulte du
lemme \ref{lemma-local-tensor-with-integral}
et du fait que $A^{sh} \to B^{sh}$ est fini et induit une extension
purement inséparable des corps résiduels, d'après le
lemme \ref{lemma-quasi-finite-strict-henselization} déjà utilisé.
Cela démontre le cas (1) du lemme.

\medskip\noindent
Dans le cas (2), écrivons $B = \colim B_i$ sous la forme d'une limite inductive filtrante de
$A$-algèbres quasi-finies. Nous obtenons alors
$D = \colim D_i$, où $D_i = B_i \otimes_A C$.
Observons que $B^{sh} = \colim B_i^{sh}$. En effet, l'anneau
$\colim B_i^{sh}$ est un anneau local strictement hensélien
d'après le lemme \ref{lemma-colimit-henselian}.
De plus, $\colim B_i^{sh}$ est une limite inductive filtrante de
$B$-algèbres \'etales d'après le lemme \ref{lemma-colimit-colimit-etale-better}.
Enfin, le corps résiduel de $\colim B_i^{sh}$ est une clôture
séparable de $\kappa(\mathfrak p_B)$ (détails omis).
Nous concluons donc que $B^{sh} = \colim B_i^{sh}$, voir la
discussion qui suit la définition \ref{definition-henselization}.
De même, nous avons $D^{sh} = \colim D_i^{sh}$.
On conclut alors en appliquant le cas (1), car
$$
D^{sh} = \colim D_i^{sh} = \colim B_i^{sh} \otimes_{A^{sh}} C^{sh} =
B^{sh} \otimes_{A^{sh}} C^{sh}
$$
puisque les limites inductives filtrantes commutent avec les produits tensoriels.

\medskip\noindent
Cas (3). Nous pouvons remplacer $A$, $B$ et $C$ par leurs localisations
en $\mathfrak p_A$, $\mathfrak p_B$ et $\mathfrak p_C$.
Ainsi, (3) résulte de (2).

\medskip\noindent
Puisqu'un morphisme d'anneaux entier est une limite inductive filtrante de morphismes d'anneaux
finis, nous voyons que (4) résulte également de (2).
\end{proof}






















\section{Critère de normalité de Serre}
\label{section-serre-criterion}

\noindent
Nous introduisons les propriétés suivantes pour les anneaux noethériens.

\begin{definition}
\label{definition-conditions}
Soit $R$ un anneau noethérien.
Soit $k \geq 0$ un entier.
\begin{enumerate}
\item On dit que $R$ possède la propriété {\it $(R_k)$} si, pour tout idéal premier $\mathfrak p$
de hauteur $\leq k$, l'anneau local $R_{\mathfrak p}$ est régulier.
On dit aussi que $R$ est {\it régulier en codimension $\leq k$}.
\item On dit que $R$ possède la propriété {\it $(S_k)$} si, pour tout idéal premier $\mathfrak p$,
l'anneau local $R_{\mathfrak p}$ est de profondeur au moins
$\min\{k, \dim(R_{\mathfrak p})\}$.
\item Soit $M$ un $R$-module de type fini. On dit que $M$ possède la propriété $(S_k)$
si, pour tout idéal premier $\mathfrak p$, le module
$M_{\mathfrak p}$ est de profondeur au moins
$\min\{k, \dim(\text{Supp}(M_{\mathfrak p}))\}$.
\end{enumerate}
\end{definition}

\noindent
Tout anneau noethérien possède la propriété $(S_0)$, de même que tout module de type fini
sur cet anneau. Notre convention selon laquelle la profondeur du module nul vaut $\infty$
(voir la section \ref{section-depth}) et la dimension de l'ensemble vide vaut
$-\infty$ (voir Topologie, section \ref{topology-section-krull-dimension})
garantit que le module nul possède la propriété $(S_k)$ pour tout $k$.

\begin{lemma}
\label{lemma-criterion-no-embedded-primes}
Soit $R$ un anneau noethérien.
Soit $M$ un $R$-module de type fini.
Les assertions suivantes sont équivalentes :
\begin{enumerate}
\item $M$ n'a aucun idéal premier associé immergé, et
\item $M$ possède la propriété $(S_1)$.
\end{enumerate}
\end{lemma}

\begin{proof}
Soit $\mathfrak p$ un idéal premier associé immergé de $M$.
Il existe alors un autre idéal premier associé $\mathfrak q$ de $M$
tel que $\mathfrak p \supset \mathfrak q$. Cela implique en particulier
que $\dim(\text{Supp}(M_{\mathfrak p})) \geq 1$ (puisque $\mathfrak q$
appartient lui aussi au support). D'autre part, $\mathfrak pR_{\mathfrak p}$
est associé à $M_{\mathfrak p}$
(lemme \ref{lemma-associated-primes-localize}), et donc
$\text{profondeur}(M_{\mathfrak p}) = 0$
(voir le lemme \ref{lemma-ideal-nonzerodivisor}).
Autrement dit, la propriété $(S_1)$ n'est pas vérifiée.
Réciproquement, si la propriété $(S_1)$ n'est pas vérifiée, il existe un idéal premier
$\mathfrak p$ tel que $\dim(\text{Supp}(M_{\mathfrak p})) \geq 1$
et $\text{profondeur}(M_{\mathfrak p}) = 0$. Puisque
$\text{profondeur}(M_{\mathfrak p}) = 0$, nous voyons que
$\mathfrak p \in \text{Ass}(M)$ d'après les deux lemmes
\ref{lemma-associated-primes-localize} et \ref{lemma-ideal-nonzerodivisor}.
Comme $\dim(\text{Supp}(M_{\mathfrak p})) \geq 1$, il existe
un idéal premier $\mathfrak q \in \text{Supp}(M)$ tel que
$\mathfrak q \subset \mathfrak p$, $\mathfrak q \not = \mathfrak p$.
Nous pouvons choisir un tel $\mathfrak q$ minimal dans
$\text{Supp}(M)$. Alors, d'après la
proposition \ref{proposition-minimal-primes-associated-primes},
nous avons $\mathfrak q \in \text{Ass}(M)$, et donc
$\mathfrak p$ est un idéal premier associé immergé.
\end{proof}

\begin{lemma}
\label{lemma-criterion-reduced}
\begin{slogan}
Réduit équivaut à R0 et S1.
\end{slogan}
Soit $R$ un anneau noethérien.
Les assertions suivantes sont équivalentes :
\begin{enumerate}
\item $R$ est réduit, et
\item $R$ vérifie les propriétés $(R_0)$ et $(S_1)$.
\end{enumerate}
\end{lemma}

\begin{proof}
Supposons que $R$ soit réduit. Alors $R_{\mathfrak p}$ est un corps pour
tout idéal premier minimal $\mathfrak p$ de $R$, d'après le
lemme \ref{lemma-minimal-prime-reduced-ring}. Ainsi, la propriété $(R_0)$ est vérifiée.
Soit $\mathfrak p$ un idéal premier de hauteur $\geq 1$. Alors $A = R_{\mathfrak p}$
est un anneau local réduit de dimension $\geq 1$. Son idéal maximal
$\mathfrak m$ n'est donc pas un idéal premier associé,
car cela signifierait qu'il existe un $x \in \mathfrak m$
d'annulateur $\mathfrak m$, si bien que $x^2 = 0$. La profondeur de
$A = R_{\mathfrak p}$ est donc au moins un, d'après le lemme \ref{lemma-ass-zero-divisors}.
Cela montre que la propriété $(S_1)$ est vérifiée.

\medskip\noindent
Réciproquement, supposons que $R$ vérifie les propriétés $(R_0)$ et $(S_1)$.
Si $\mathfrak p$ est un idéal premier minimal de $R$, alors
$R_{\mathfrak p}$ est un corps d'après $(R_0)$, et est donc réduit.
Si $\mathfrak p$ n'est pas minimal, alors nous voyons que $R_{\mathfrak p}$
est de profondeur $\geq 1$ d'après $(S_1)$, et nous en déduisons qu'il existe un élément
$t \in \mathfrak pR_{\mathfrak p}$ tel que
$R_{\mathfrak p} \to R_{\mathfrak p}[1/t]$ soit injectif.
Or $R_\mathfrak p[1/t]$ est contenu dans le produit
de ses localisations aux idéaux premiers, voir le
lemme \ref{lemma-characterize-zero-local}.
Il s'ensuit que $R_{\mathfrak p}$ est un sous-anneau d'un produit de
localisations de $R$ en des idéaux premiers satisfaisant $\mathfrak p \supset \mathfrak q$ et
$t \not \in \mathfrak q$. Comme ces idéaux premiers sont de hauteur plus petite,
nous concluons par récurrence sur la hauteur que $R$ est réduit.
\end{proof}

\begin{lemma}[Critère de normalité de Serre]
\label{lemma-criterion-normal}
\begin{reference}
\cite[IV, Theorem 5.8.6]{EGA}
\end{reference}
\begin{slogan}
Normal équivaut à R1 et S2.
\end{slogan}
Soit $R$ un anneau noethérien.
Les assertions suivantes sont équivalentes :
\begin{enumerate}
\item $R$ est un anneau normal, et
\item $R$ vérifie les propriétés $(R_1)$ et $(S_2)$.
\end{enumerate}
\end{lemma}

\begin{proof}
Démonstration de (1) $\Rightarrow$ (2). Supposons que $R$ soit normal, c'est-à-dire que tous
les localisés $R_{\mathfrak p}$ aux idéaux premiers soient des anneaux intègres normaux.
En particulier, nous voyons que $R$ vérifie $(R_0)$ et $(S_1)$ d'après le
lemme \ref{lemma-criterion-reduced}. Il suffit donc de montrer
qu'un anneau local noethérien intègre et normal $R$ de dimension $d$ est
de profondeur $\geq \min(2, d)$ et est régulier si $d = 1$. L'assertion
lorsque $d = 1$ résulte du lemme \ref{lemma-characterize-dvr}.

\medskip\noindent
Soit $R$ un anneau local noethérien intègre et normal, d'idéal maximal
$\mathfrak m$ et de dimension $d \geq 2$. Appliquons le
lemme \ref{lemma-hart-serre-loc-thm} à $R$.
Il est clair que $R$ ne relève ni du cas (1) ni du cas (2)
du lemme.
Considérons $R \to R'$ comme dans le cas (4) du lemme.
Puisque $R$ est intègre, nous avons $R \subset R'$. Comme $\mathfrak m$
n'est pas un idéal premier associé à $R'$, il existe un $x \in \mathfrak m$
qui est un non-diviseur de zéro dans $R'$. Alors $R_x = R'_x$, de sorte que
$R$ et $R'$ sont des anneaux intègres de même corps des fractions. Mais la
finitude de $R \subset R'$ implique que tout élément de $R'$ est entier
sur $R$ (lemme \ref{lemma-finite-is-integral}),
et nous concluons que $R = R'$, puisque $R$ est normal.
Cela signifie que le cas (4) ne se produit pas. Il reste donc la possibilité
(3), c'est-à-dire $\text{profondeur}(R) \geq 2$, comme souhaité.

\medskip\noindent
Démonstration de (2) $\Rightarrow$ (1). Supposons que $R$ vérifie $(R_1)$ et $(S_2)$.
Le lemme \ref{lemma-criterion-reduced} montre que $R$ est
réduit. Il suffit donc de montrer que, si $R$ est un anneau local réduit
noethérien de dimension $d$ vérifiant $(S_2)$ et $(R_1)$,
alors $R$ est un anneau intègre normal. Si $d = 0$, le résultat est clair.
Si $d = 1$, le résultat découle du lemme \ref{lemma-characterize-dvr}.

\medskip\noindent
Soit $R$ un anneau local noethérien réduit, d'idéal maximal
$\mathfrak m$ et de dimension $d \geq 2$, vérifiant $(R_1)$ et
$(S_2)$. D'après le lemme \ref{lemma-characterize-reduced-ring-normal},
il suffit de montrer que $R$ est intégralement clos dans son
anneau total des fractions $Q(R)$. Choisissons $x \in Q(R)$ entier
sur $R$. Alors $R' = R[x]$ est une extension finie de $R$
(lemme \ref{lemma-characterize-finite-in-terms-of-integral}).
Puisque $\dim(R_\mathfrak p) < d$ pour
tout idéal premier non maximal $\mathfrak p \subset R$,
nous avons $R_\mathfrak p = R'_\mathfrak p$ par récurrence.
Le support de $R'/R$ est donc $\{\mathfrak m\}$.
Il s'ensuit que $R'/R$ est annulé par une puissance de $\mathfrak m$
(lemme \ref{lemma-Noetherian-power-ideal-kills-module}).
Le lemme \ref{lemma-hart-serre-loc-thm} contredit alors
l'hypothèse selon laquelle la profondeur de $R$ est $\geq 2 = \min(2, d)$,
et la démonstration est achevée.
\end{proof}

\begin{lemma}
\label{lemma-regular-normal}
Un anneau régulier est normal.
\end{lemma}

\begin{proof}
Soit $R$ un anneau régulier. D'après le
lemme \ref{lemma-criterion-normal},
il suffit de montrer que $R$ vérifie $(R_1)$ et $(S_2)$.
Comme tout anneau local régulier est de Cohen-Macaulay, voir le
lemme \ref{lemma-regular-ring-CM},
il est clair que $R$ vérifie $(S_2)$.
La propriété $(R_1)$ est immédiate.
\end{proof}

\begin{lemma}
\label{lemma-normal-domain-intersection-localizations-height-1}
Soit $R$ un anneau noethérien intègre et normal de corps des fractions $K$. Alors
\begin{enumerate}
\item pour tout $a \in R$ non nul, le quotient $R/aR$ n'a aucun idéal premier associé immergé,
et tous ses idéaux premiers associés sont de hauteur $1$
\item
$$
R = \bigcap\nolimits_{\text{hauteur}(\mathfrak p) = 1} R_{\mathfrak p}
$$
\item Pour tout $x \in K$ non nul, le quotient $R/(R \cap xR)$
n'a aucun idéal premier associé immergé, et tous ses idéaux premiers associés sont de hauteur $1$.
\end{enumerate}
\end{lemma}

\begin{proof}
Le lemme \ref{lemma-criterion-normal} montre que $R$ vérifie $(S_2)$.
Ainsi, pour tout élément $a \in R$ non nul, le quotient $R/aR$ vérifie $(S_1)$
(utiliser par exemple le lemme \ref{lemma-depth-in-ses}).
Par conséquent, $R/aR$ n'a aucun idéal premier associé immergé
(lemme \ref{lemma-criterion-no-embedded-primes}).
Nous en déduisons que les idéaux premiers associés à $R/aR$ sont exactement
les idéaux premiers minimaux $\mathfrak p$ au-dessus de $(a)$, qui sont de hauteur $1$
puisque $a$ n'est pas nul (lemme \ref{lemma-minimal-over-1}). Cela démontre (1).

\medskip\noindent
Ainsi, étant donné $b \in R$, on a $b \in aR$ si et seulement si
$b \in aR_{\mathfrak p}$ pour tout idéal premier minimal $\mathfrak p$
au-dessus de $(a)$ (voir le lemme \ref{lemma-zero-at-ass-zero}).
Ces idéaux premiers sont tous de hauteur $1$, comme on l'a vu ci-dessus, donc
$b/a \in R$ si et seulement si $b/a \in R_{\mathfrak p}$ pour tous les idéaux premiers
de hauteur 1. L'assertion (2) s'ensuit.

\medskip\noindent
Pour (3), écrivons $x = a/b$. Soient $\mathfrak p_1, \ldots, \mathfrak p_r$
les idéaux premiers minimaux au-dessus de $(ab)$. Ils sont tous de hauteur 1 d'après ce qui précède.
Nous avons alors
$R \cap xR = \bigcap_{i = 1, \ldots, r} (R \cap xR_{\mathfrak p_i})$
d'après la partie (2) du lemme. Par conséquent, $R/(R \cap xR)$ est un sous-module de
$\bigoplus R/(R \cap xR_{\mathfrak p_i})$.
Comme $R_{\mathfrak p_i}$ est un anneau de valuation discrète (d'après la propriété $(R_1)$
pour l'anneau noethérien intègre et normal $R$, voir le lemme \ref{lemma-criterion-normal}),
nous avons $xR_{\mathfrak p_i} = \mathfrak p_i^{e_i}R_{\mathfrak p_i}$
pour un certain $e_i \in \mathbf{Z}$. La somme directe est donc égale
à $\bigoplus_{e_i > 0} R/\mathfrak p_i^{(e_i)}$, voir la
définition \ref{definition-symbolic-power}.
D'après le lemme \ref{lemma-symbolic-power-associated},
le seul idéal premier associé au module
$R/\mathfrak p^{(n)}$ est $\mathfrak p$. Ainsi, l'ensemble des idéaux premiers associés
à $R/(R \cap xR)$ est un sous-ensemble de $\{\mathfrak p_i\}$ et il n'existe
aucune relation d'inclusion entre eux. Cela démontre (3).
\end{proof}















\section{Lissité formelle des corps}
\label{section-p-bases}

\noindent
Dans cette section, nous montrons que les extensions de corps sont formellement lisses
si et seulement si elles sont séparables. Cependant, nous démontrons d'abord que les extensions de corps
de type fini sont algébriques séparables si et seulement si
elles sont formellement non ramifiées.

\begin{lemma}
\label{lemma-characterize-separable-algebraic-field-extensions}
Soit $K/k$ une extension de corps de type fini.
Les propriétés suivantes sont équivalentes :
\begin{enumerate}
\item $K$ est une extension finie séparable de $k$,
\item $\Omega_{K/k} = 0$,
\item $K$ est formellement non ramifié sur $k$,
\item $K$ est non ramifié sur $k$,
\item $K$ est formellement \'etale sur $k$,
\item $K$ est \'etale sur $k$.
\end{enumerate}
\end{lemma}

\begin{proof}
L'équivalence de (2) et (3) résulte du
Lemme \ref{lemma-characterize-formally-unramified}.
D'après le Lemme \ref{lemma-etale-over-field},
nous voyons que (1) est équivalent à (6).
La propriété (6) implique (5) et (4), qui impliquent toutes deux à leur tour (3)
(Lemmes \ref{lemma-formally-etale-etale}, \ref{lemma-unramified},
et \ref{lemma-formally-unramified-unramified}).
Il suffit donc de montrer que (2) implique (1).
Choisissons une $k$-sous-algèbre de type fini $A \subset K$
telle que $K$ soit le corps des fractions de l'anneau intègre $A$.
Posons $S = A \setminus \{0\}$.
Puisque $0 = \Omega_{K/k} = S^{-1}\Omega_{A/k}$
(Lemme \ref{lemma-differentials-localize})
et puisque $\Omega_{A/k}$ est de type fini
(Lemme \ref{lemma-differentials-finitely-generated}),
nous pouvons remplacer $A$ par une localisation $A_f$ pour nous ramener au cas
où $\Omega_{A/k} = 0$ (détails omis).
Alors $A$ est non ramifié sur $k$, donc
$K/k$ est une extension finie séparable, par exemple d'après le
Lemme \ref{lemma-unramified-at-prime} appliqué avec $\mathfrak q = (0)$.
\end{proof}

\begin{lemma}
\label{lemma-derivative-zero-pth-power}
Soit $k$ un corps parfait de caractéristique $p > 0$.
Soit $K/k$ une extension de corps.
Soit $a \in K$. Alors $\text{d}a = 0$ dans $\Omega_{K/k}$
si et seulement si $a$ est une puissance $p$-ième.
\end{lemma}

\begin{proof}
D'après le Lemme \ref{lemma-colimit-differentials}, nous voyons qu'il existe un sous-corps
$k \subset L \subset K$ tel que $L/k$
soit une extension de corps de type fini et que
$\text{d}a$ soit nul dans $\Omega_{L/k}$.
Ainsi, nous pouvons supposer que $K$ est une extension de corps de type fini
de $k$.

\medskip\noindent
Choisissons une base de transcendance $x_1, \ldots, x_r \in K$
telle que $K$ soit une extension finie séparable de $k(x_1, \ldots, x_r)$.
Cela est possible par définition ; voir les
Définitions \ref{definition-perfect} et
\ref{definition-separable-field-extension}.
Remarquons que le résultat vaut pour le sous-corps
$k(x_1, \ldots, x_r) \subset K$ ; ce sous-corps est une extension transcendante pure du corps de base.
En effet,
$$
\Omega_{k(x_1, \ldots, x_r)/k} =
\bigoplus\nolimits_{i = 1}^r k(x_1, \ldots, x_r) \text{d}x_i
$$
et toute fonction rationnelle dont toutes les dérivées partielles sont nulles
est une puissance $p$-ième. De plus, nous avons aussi
$$
\Omega_{K/k} =
\bigoplus\nolimits_{i = 1}^r K\text{d}x_i
$$
puisque l'extension $k(x_1, \ldots, x_r) \subset K$ est finie séparable
(calcul omis). Supposons que $a \in K$ soit un élément tel que
$\text{d}a = 0$ dans le module des différentielles. Par notre choix des $x_i$, nous
voyons que le polynôme minimal $P(T) \in k(x_1, \ldots, x_r)[T]$
de $a$ est séparable. Écrivons
$$
P(T) = T^d + \sum\nolimits_{i = 1}^d a_i T^{d - i}
$$
et donc
$$
0 = \text{d}P(a) = \sum\nolimits_{i = 1}^d a^{d - i}\text{d}a_i
$$
dans $\Omega_{K/k}$. D'après la description de
$\Omega_{K/k}$ ci-dessus et le fait que $P$ était le polynôme minimal
de $a$, nous voyons que cela implique $\text{d}a_i = 0$.
Ainsi, $a_i = b_i^p$ pour tout $i$. Par conséquent, d'après
Corps, Lemme \ref{fields-lemma-pth-root},
nous voyons que $a$ est une puissance $p$-ième.
\end{proof}

\begin{lemma}
\label{lemma-size-extension-pth-roots}
Soit $k$ un corps de caractéristique $p > 0$.
Soient $a_1, \ldots, a_n \in k$ des éléments tels que
$\text{d}a_1, \ldots, \text{d}a_n$ soient linéairement indépendants dans
$\Omega_{k/\mathbf{F}_p}$. Alors l'extension de corps
$k(a_1^{1/p}, \ldots, a_n^{1/p})$ est de degré $p^n$ sur $k$.
\end{lemma}

\begin{proof}
Nous procédons par récurrence sur $n$. Si $n = 1$, le résultat est le
Lemme \ref{lemma-derivative-zero-pth-power}.
Pour l'étape de récurrence, supposons que $k(a_1^{1/p}, \ldots, a_{n - 1}^{1/p})$
soit de degré $p^{n - 1}$ sur $k$. Nous devons montrer que $a_n$ ne devient pas
une puissance $p$-ième dans $k(a_1^{1/p}, \ldots, a_{n - 1}^{1/p})$.
Si c'est le cas, nous pouvons écrire
\begin{align*}
a_n & =
\left(\sum\nolimits_{I = (i_1, \ldots, i_{n - 1}),\ 0 \leq i_j \leq p - 1}
\lambda_I a_1^{i_1/p} \ldots a_{n - 1}^{i_{n - 1}/p}\right)^p \\
& = \sum\nolimits_{I = (i_1, \ldots, i_{n - 1}),\ 0 \leq i_j \leq p - 1}
\lambda_I^p a_1^{i_1} \ldots a_{n - 1}^{i_{n - 1}}
\end{align*}
En appliquant $\text{d}$, nous voyons que $\text{d}a_n$ dépend linéairement des
$\text{d}a_i$, $i < n$. C'est une contradiction.
\end{proof}

\begin{lemma}
\label{lemma-separable-differentials}
Soit $k$ un corps de caractéristique $p > 0$.
Les propriétés suivantes sont équivalentes :
\begin{enumerate}
\item l'extension de corps $K/k$ est séparable
(voir la Définition \ref{definition-separable-field-extension}), et
\item l'application
$K \otimes_k \Omega_{k/\mathbf{F}_p} \to \Omega_{K/\mathbf{F}_p}$
est injective.
\end{enumerate}
\end{lemma}

\begin{proof}
Écrivons $K$ comme une limite inductive filtrante $K = \colim_i K_i$ d'extensions de corps
$K_i/k$ de type fini. Par définition, $K$ est séparable si et seulement
si chaque $K_i$ est séparable sur $k$ et, d'après le
Lemme \ref{lemma-colimit-differentials}, nous voyons que
$K \otimes_k \Omega_{k/\mathbf{F}_p} \to \Omega_{K/\mathbf{F}_p}$
est injective si et seulement si chaque application
$K_i \otimes_k \Omega_{k/\mathbf{F}_p} \to \Omega_{K_i/\mathbf{F}_p}$
est injective. Ainsi, nous pouvons supposer que $K/k$ est une extension de corps
de type fini.

\medskip\noindent
Supposons que $K/k$ soit une extension de corps de type fini qui soit
séparable. Choisissons $x_1, \ldots, x_{r + 1} \in K$ comme dans le
Lemme \ref{lemma-generating-finitely-generated-separable-field-extensions}.
Dans ce cas, il existe un polynôme irréductible
$G(X_1, \ldots, X_{r + 1}) \in k[X_1, \ldots, X_{r + 1}]$
tel que $G(x_1, \ldots, x_{r + 1}) = 0$ et tel que
$\partial G/\partial X_{r + 1}$ ne soit pas identiquement nul.
De plus, $K$ est le corps des fractions de l'anneau intègre
$S = K[X_1, \ldots, X_{r + 1}]/(G)$.
Écrivons
$$
G = \sum a_I X^I, \quad X^I = X_1^{i_1}\ldots X_{r + 1}^{i_{r + 1}}.
$$
À l'aide de la présentation de $S$ ci-dessus, nous voyons que
$$
\Omega_{S/\mathbf{F}_p}
=
\frac{
S \otimes_k \Omega_k \oplus
\bigoplus\nolimits_{i = 1, \ldots, r + 1} S\text{d}X_i
}{
\langle
\sum X^I \text{d}a_I + \sum \partial G/\partial X_i \text{d}X_i
\rangle
}
$$
Puisque $\Omega_{K/\mathbf{F}_p}$ est la localisation
du $S$-module $\Omega_{S/\mathbf{F}_p}$ (voir le
Lemme \ref{lemma-differentials-localize}), nous en concluons
que
$$
\Omega_{K/\mathbf{F}_p}
=
\frac{
K \otimes_k \Omega_k \oplus
\bigoplus\nolimits_{i = 1, \ldots, r + 1} K\text{d}X_i
}{
\langle
\sum X^I \text{d}a_I + \sum \partial G/\partial X_i \text{d}X_i
\rangle
}
$$
Maintenant, puisque le polynôme $\partial G/\partial X_{r + 1}$ n'est pas
identiquement nul, nous en concluons que l'application
$K \otimes_k \Omega_{k/\mathbf{F}_p} \to \Omega_{S/\mathbf{F}_p}$
est injective, comme souhaité.

\medskip\noindent
Supposons que $K/k$ soit une extension de corps de type fini
et que
$K \otimes_k \Omega_{k/\mathbf{F}_p} \to \Omega_{K/\mathbf{F}_p}$
soit injective.
(Cette partie de la preuve est identique à l'argument démontrant le
Lemme \ref{lemma-characterize-separable-field-extensions}.)
Soit $x_1, \ldots, x_r$ une base de transcendance de $K$ sur $k$ telle
que le degré d'inséparabilité de l'extension finie
$k(x_1, \ldots, x_r) \subset K$ soit minimal.
Si $K$ est séparable sur $k(x_1, \ldots, x_r)$, alors la conclusion est acquise.
Supposons que ce ne soit pas le cas afin d'obtenir une contradiction.
Il existe alors un élément $\alpha \in K$ qui n'est pas
séparable sur $k(x_1, \ldots, x_r)$. Soit $P(T) \in k(x_1, \ldots, x_r)[T]$
son polynôme minimal. Puisque $\alpha$ n'est pas séparable,
$P$ est en fait un polynôme en $T^p$. Chassons les dénominateurs
pour obtenir un polynôme irréductible
$$
G(X_1, \ldots, X_r, T) = \sum a_{I, i} X^I T^i \in k[X_1, \ldots, X_r, T]
$$
tel que $G(x_1, \ldots, x_r, \alpha) = 0$ dans $K$.
Notons que cela signifie que $k[X_1, \ldots, X_r, T]/(G) \subset K$.
Nous pouvons supposer que, pour un couple $(I_0, i_0)$, le coefficient
$a_{I_0, i_0} = 1$.
Nous affirmons que $\text{d}G/\text{d}X_i$ n'est pas identiquement nulle
pour au moins un $i$. En effet, si ce n'est pas le cas,
$G$ est en fait un polynôme en $X_1^p, \ldots, X_r^p, T^p$.
Cela signifie alors que
$$
\sum\nolimits_{(I, i) \not = (I_0, i_0)} x^I\alpha^i \text{d}a_{I, i}
$$
est nulle dans $\Omega_{K/\mathbf{F}_p}$. Notons qu'il n'existe aucune
relation $k$-linéaire entre les éléments
$$
\{x^I\alpha^i \mid a_{I, i} \not = 0 \text{ et } (I, i) \not = (I_0, i_0)\}
$$
de $K$. L'hypothèse
que $K \otimes_k \Omega_{k/\mathbf{F}_p} \to \Omega_{K/\mathbf{F}_p}$
est injective implique donc que $\text{d}a_{I, i} = 0$
dans $\Omega_{k/\mathbf{F}_p}$ pour tout couple $(I, i)$.
D'après le Lemme \ref{lemma-derivative-zero-pth-power},
nous voyons que chaque $a_{I, i}$ est une puissance $p$-ième, ce qui
implique que $G$ est une puissance $p$-ième, en contradiction avec l'irréductibilité de
$G$. Ainsi,
après renumérotation, nous pouvons supposer que $\text{d}G/\text{d}X_1$ n'est pas nulle.
Nous voyons alors que $x_1$ est algébrique séparable sur
$k(x_2, \ldots, x_r, \alpha)$ et que $x_2, \ldots, x_r, \alpha$
forment une base de transcendance de $K$ sur $k$. Cela signifie que
le degré d'inséparabilité de l'extension finie
$k(x_2, \ldots, x_r, \alpha) \subset K$ est inférieur au
degré d'inséparabilité de l'extension finie
$k(x_1, \ldots, x_r) \subset K$, ce qui est une contradiction.
\end{proof}

\begin{lemma}
\label{lemma-formally-smooth-implies-separable}
Soit $K/k$ une extension de corps.
Si $K$ est formellement lisse sur $k$, alors $K$ est
une extension séparable de $k$.
\end{lemma}

\begin{proof}
Supposons que $K$ soit formellement lisse sur $k$.
D'après le Lemme \ref{lemma-ses-formally-smooth}, nous voyons que
$K \otimes_k \Omega_{k/\mathbf{Z}} \to \Omega_{K/\mathbf{Z}}$
est injective. Ainsi, $K$ est séparable sur $k$ d'après le
Lemme \ref{lemma-separable-differentials}.
\end{proof}

\begin{lemma}
\label{lemma-characterize-formally-smooth-field-extension}
Soit $K/k$ une extension de corps.
Alors $K$ est formellement lisse sur $k$ si et seulement si
$H_1(L_{K/k}) = 0$.
\end{lemma}

\begin{proof}
Cela résulte de la Proposition \ref{proposition-characterize-formally-smooth}
et du fait qu'un espace vectoriel est libre (donc projectif).
\end{proof}

\begin{lemma}
\label{lemma-formally-smooth-extensions-easy}
Soit $K/k$ une extension de corps.
\begin{enumerate}
\item Si $K$ est une extension transcendante pure de $k$, alors
$K$ est formellement lisse sur $k$.
\item Si $K$ est une extension algébrique séparable de $k$, alors $K$ est
formellement lisse sur $k$.
\item Si $K$ est séparable sur $k$, alors $K$ est formellement lisse
sur $k$.
\end{enumerate}
\end{lemma}

\begin{proof}
Pour (1), écrivons $K = k(x_j; j \in J)$. Supposons que
$A$ soit une $k$-algèbre et que $I \subset A$ soit un idéal de
carré nul. Soit $\varphi : K \to A/I$ un morphisme de $k$-algèbres.
Soit $a_j \in A$ un élément tel que $a_j \mod I = \varphi(x_j)$.
Il est alors aisé de voir qu'il existe un unique morphisme de $k$-algèbres
$K \to A$ qui envoie $x_j$ sur $a_j$ et qui se réduit
à $\varphi$ modulo $I$. Ainsi, $k \subset K$ est formellement lisse.

\medskip\noindent
Dans le cas (2), nous voyons que $k \subset K$ est une limite inductive
d'extensions d'anneaux \'etales. Un morphisme d'anneaux \'etale est formellement \'etale
(Lemme \ref{lemma-formally-etale-etale}). Ce cas résulte donc du
Lemme \ref{lemma-colimit-formally-etale} et de l'observation triviale
qu'un morphisme d'anneaux formellement \'etale est formellement lisse.

\medskip\noindent
Dans le cas (3), écrivons $K = \colim K_i$ comme la limite inductive filtrante de ses
sous-extensions de type fini sur $k$. D'après la
Définition \ref{definition-separable-field-extension},
chaque $K_i$ est une extension algébrique séparable d'une extension transcendante pure
de $k$. Ainsi, $K_i/k$ est formellement lisse d'après les cas (1) et (2) et le
Lemme \ref{lemma-compose-formally-smooth}. Donc
$H_1(L_{K_i/k}) = 0$ d'après le
Lemme \ref{lemma-characterize-formally-smooth-field-extension}.
Ainsi, $H_1(L_{K/k}) = 0$ d'après le Lemme \ref{lemma-colimits-NL}.
Donc $K/k$ est formellement lisse, de nouveau d'après le
Lemme \ref{lemma-characterize-formally-smooth-field-extension}.
\end{proof}

\begin{lemma}
\label{lemma-fields-are-formally-smooth}
\begin{slogan}
La lissité formelle équivaut à la séparabilité pour les extensions de corps.
\end{slogan}
Soit $k$ un corps.
\begin{enumerate}
\item Si la caractéristique de $k$ est nulle, alors toute extension de corps
de $k$ est formellement lisse sur $k$.
\item Si la caractéristique de $k$ est $p > 0$, alors $K/k$ est
formellement lisse si et seulement si cette extension de corps est séparable.
\end{enumerate}
\end{lemma}

\begin{proof}
On combine les Lemmes \ref{lemma-formally-smooth-implies-separable} et
\ref{lemma-formally-smooth-extensions-easy}.
\end{proof}

\noindent
Nous rassemblons ici toutes les caractérisations des extensions
de corps séparables.

\begin{proposition}
\label{proposition-characterize-separable-field-extensions}
Soit $K/k$ une extension de corps.
Si la caractéristique de $k$ est nulle, les assertions suivantes sont équivalentes :
\begin{enumerate}
\item $K$ est séparable sur $k$,
\item $K$ est géométriquement réduit sur $k$,
\item $K$ est formellement lisse sur $k$,
\item $H_1(L_{K/k}) = 0$, et
\item l'application $K \otimes_k \Omega_{k/\mathbf{Z}} \to \Omega_{K/\mathbf{Z}}$
est injective.
\end{enumerate}
Si la caractéristique de $k$ est $p > 0$, alors les propriétés suivantes sont
équivalentes :
\begin{enumerate}
\item $K$ est séparable sur $k$,
\item l'anneau $K \otimes_k k^{1/p}$ est réduit,
\item $K$ est géométriquement réduit sur $k$,
\item l'application $K \otimes_k \Omega_{k/\mathbf{F}_p} \to \Omega_{K/\mathbf{F}_p}$
est injective,
\item $H_1(L_{K/k}) = 0$, et
\item $K$ est formellement lisse sur $k$.
\end{enumerate}
\end{proposition}

\begin{proof}
C'est une combinaison des
Lemmes \ref{lemma-characterize-separable-field-extensions},
\ref{lemma-fields-are-formally-smooth},
\ref{lemma-formally-smooth-implies-separable} et
\ref{lemma-separable-differentials}.
\end{proof}

\noindent
Voici une autre caractérisation des extensions de corps séparables
de type fini.

\begin{lemma}
\label{lemma-localization-smooth-separable}
Soit $K/k$ une extension de corps de type fini.
Alors $K$ est séparable sur $k$ si et seulement si $K$ est
la localisation d'une $k$-algèbre lisse.
\end{lemma}

\begin{proof}
Choisissons une $k$-algèbre de type fini $R$ qui soit un anneau intègre dont le
corps des fractions est $K$. Le Lemme \ref{lemma-smooth-at-generic-point}
affirme que $k \to R$ est lisse
en $(0)$ si et seulement si $K/k$ est séparable.
Cela prouve le lemme.
\end{proof}

\begin{lemma}
\label{lemma-colimit-syntomic}
Soit $K/k$ une extension de corps.
Alors $K$ est une limite inductive filtrante d'algèbres qui sont des intersections complètes
globales sur $k$. Si $K/k$ est séparable, alors $K$ est une limite inductive filtrante
d'algèbres lisses sur $k$.
\end{lemma}

\begin{proof}
Supposons que $E \subset K$ soit un sous-ensemble fini. Il suffit de montrer qu'il
existe une $k$-sous-algèbre $A \subset K$ qui contient $E$
et qui est une intersection complète globale (resp.\ lisse) sur $k$.
Le cas séparable/lisse résulte du
Lemme \ref{lemma-localization-smooth-separable}.
Dans le cas général, soit $L \subset K$ le sous-corps engendré par $E$.
Choisissons une base de transcendance $x_1, \ldots, x_d \in L$ sur $k$.
L'extension $L/k(x_1, \ldots, x_d)$ est finie.
Écrivons $L = k(x_1, \ldots, x_d)[y_1, \ldots, y_r]$.
Choisissons par récurrence des polynômes $P_i \in k(x_1, \ldots, x_d)[Y_1, \ldots, Y_r]$
tels que $P_i = P_i(Y_1, \ldots, Y_i)$ soit unitaire en $Y_i$ sur
$k(x_1, \ldots, x_d)[Y_1, \ldots, Y_{i - 1}]$ et dont l'image est le
polynôme minimal de $y_i$ dans
$k(x_1, \ldots, x_d)[y_1, \ldots, y_{i - 1}][Y_i]$.
Il est alors clair que la suite $P_1, \ldots, P_r$ est régulière
dans $k(x_1, \ldots, x_r)[Y_1, \ldots, Y_r]$ et que
$L = k(x_1, \ldots, x_r)[Y_1, \ldots, Y_r]/(P_1, \ldots, P_r)$.
Si $h \in k[x_1, \ldots, x_d]$ est un polynôme tel que
$P_i \in k[x_1, \ldots, x_d, 1/h, Y_1, \ldots, Y_r]$, alors
nous voyons que $P_1, \ldots, P_r$ est une suite régulière dans
$k[x_1, \ldots, x_d, 1/h, Y_1, \ldots, Y_r]$ et
$A = k[x_1, \ldots, x_d, 1/h, Y_1, \ldots, Y_r]/(P_1, \ldots, P_r)$
est une intersection complète globale. Après avoir ajusté notre choix de $h$,
nous pouvons supposer que $E \subset A$, ce qui conclut.
\end{proof}






\section{Construction de morphismes d'anneaux plats}
\label{section-constructing-flat}

\noindent
Le lemme suivant est parfois utile.

\begin{lemma}
\label{lemma-flat-local-given-residue-field}
Soit $(R, \mathfrak m, k)$ un anneau local. Soit $K/k$ une
extension de corps. Il existe un anneau local $(R', \mathfrak m', k')$, un morphisme local plat
d'anneaux $R \to R'$ tel que $\mathfrak m' = \mathfrak mR'$ et tel que
$k'$ soit isomorphe à $K$ comme extension de $k$.
\end{lemma}

\begin{proof}
Supposons que $k' = k(\alpha)$ soit une extension monogène de $k$.
Alors $k'$ est le corps résiduel d'une extension locale plate $R \subset R'$
comme dans le lemme. Plus précisément, si $\alpha$ est transcendant sur $k$, nous prenons
pour $R'$ la localisation de $R[x]$ en l'idéal premier $\mathfrak mR[x]$.
Si $\alpha$ est algébrique de polynôme minimal
$T^d + \sum \overline{\lambda}_iT^{d - i}$, alors nous prenons
$R' = R[T]/(T^d + \sum \lambda_i T^{d - i})$.

\medskip\noindent
Considérons la famille des triplets $(k', R \to R', \phi)$, où
$k \subset k' \subset K$ est un sous-corps,
$R \to R'$ est un morphisme local d'anneaux comme dans le lemme, et
$\phi : R' \to k'$ induit un isomorphisme $R'/\mathfrak mR' \cong k'$
d'extensions de $k$. Ces triplets forment une ``grande'' catégorie $\mathcal{C}$ dont les morphismes
$(k_1, R_1, \phi_1) \to (k_2, R_2, \phi_2)$
sont donnés par des morphismes d'anneaux $\psi : R_1 \to R_2$ tels que
$$
\xymatrix{
R_1 \ar[d]_\psi \ar[r]_{\phi_1} & k_1 \ar[r] & K \ar@{=}[d] \\
R_2 \ar[r]^{\phi_2} & k_2 \ar[r] & K
}
$$
le diagramme soit commutatif. Cela implique que $k_1 \subset k_2$.

\medskip\noindent
Supposons que $I$ soit un ensemble ordonné filtrant et que
$((R_i, k_i, \phi_i), \psi_{ii'})$ soit un système inductif indexé par $I$ ; voir
Catégories, section \ref{categories-section-posets-limits}.
Dans ce cas, nous pouvons considérer
$$
R' = \colim_{i \in I} R_i
$$
C'est un anneau local d'idéal maximal $\mathfrak mR'$ et de
corps résiduel $k' = \bigcup_{i \in I} k_i$. De plus, le morphisme d'anneaux
$R \to R'$ est plat, car c'est une limite inductive de morphismes plats (et les produits
tensoriels commutent aux limites inductives filtrantes).
Ainsi, $(R', k', \phi')$ est un ``majorant'' du système.

\medskip\noindent
Une application presque triviale du lemme de Zorn achèverait la preuve
si $\mathcal{C}$ était un ensemble, mais ce n'en est pas un.
(En fait, on peut faire fonctionner cet argument en trouvant une majoration raisonnable des
cardinaux des anneaux locaux qui interviennent.)
Pour contourner ce problème, choisissons un bon ordre sur $K$.
Pour $x \in K$, notons $K(x)$ le sous-corps de $K$ engendré
par tous les éléments de $K$ qui sont $\leq x$.
Par récurrence transfinie sur $x \in K$, nous allons construire des morphismes d'anneaux
$R \subset R(x)$ comme dans le lemme, avec pour extension de corps résiduels
$K(x)/k$. De plus, la construction assure
que $R(x)$ contient $R(y)$ pour tout $y \leq x$.
En effet, si $x$ a un prédécesseur $x'$, alors $K(x) = K(x')[x]$
et nous pouvons donc prendre pour $R(x') \subset R(x)$ l'extension d'anneaux locaux
construite dans le premier paragraphe de la preuve. Si $x$ n'a pas
de prédécesseur, posons d'abord
$R'(x) = \colim_{x' < x} R(x')$ comme dans le troisième paragraphe
de la preuve. Le corps résiduel de $R'(x)$ est $K'(x) = \bigcup_{x' < x} K(x')$.
Puisque $K(x) = K'(x)[x]$, nous pouvons utiliser la construction du
premier paragraphe de la preuve pour produire $R'(x) \subset R(x)$.
Cela achève la preuve du lemme.
\end{proof}

\begin{lemma}
\label{lemma-colimit-finite-etale-given-residue-field}
Soit $(R, \mathfrak m, k)$ un anneau local. Si $k \subset K$ est une
extension algébrique séparable, alors il existe un ensemble ordonné filtrant $I$ et
un système inductif d'extensions finies \'etales $R \subset R_i$, $i \in I$
d'anneaux locaux tel que $R' = \colim R_i$ ait pour corps résiduel
$K$ (comme extension de $k$).
\end{lemma}

\begin{proof}
Soit $R \subset R'$ l'extension construite dans la preuve du
Lemme \ref{lemma-flat-local-given-residue-field}. Par construction,
$R' = \colim_{\alpha \in A} R_\alpha$, où $A$ est un ensemble bien ordonné,
et les morphismes de transition $R_\alpha \to R_{\alpha + 1}$
sont finis \'etales, tandis que $R_\alpha = \colim_{\beta < \alpha} R_\beta$
si $\alpha$ n'est pas un successeur. Nous allons démontrer le résultat par récurrence
transfinie.

\medskip\noindent
Supposons le résultat vrai pour $R_\alpha$, c'est-à-dire $R_\alpha = \colim R_i$,
où $R_i$ est fini \'etale sur $R$. Puisque
$R_\alpha \to R_{\alpha + 1}$ est fini \'etale,
il existe un $i$ et une extension finie \'etale $R_i \to R_{i, 1}$
tels que $R_{\alpha + 1} = R_\alpha \otimes_{R_i} R_{i, 1}$.
Ainsi, $R_{\alpha + 1} = \colim_{i' \geq i} R_{i'} \otimes_{R_i} R_{i, 1}$,
et le résultat est vrai pour $\alpha + 1$. Supposons que $\alpha$ ne soit pas un successeur
et que le résultat soit vrai pour $R_\beta$ pour tout $\beta < \alpha$.
Tout sous-ensemble fini $E \subset R_\alpha$ est contenu dans $R_\beta$
pour un certain $\beta < \alpha$ et, par hypothèse, $E$ est contenu dans une sous-extension finie \'etale.
Ainsi, le résultat est vrai pour $R_\alpha$.
\end{proof}

\begin{lemma}
\label{lemma-finite-free-given-residue-field-extension}
Soit $R$ un anneau. Soit $\mathfrak p \subset R$ un idéal premier et
soit $L/\kappa(\mathfrak p)$ une extension finie de corps.
Alors il existe un morphisme d'anneaux fini et libre $R \to S$ tel que
$\mathfrak q = \mathfrak pS$ soit premier et que
$\kappa(\mathfrak q)/\kappa(\mathfrak p)$ soit isomorphe à l'extension donnée
$L/\kappa(\mathfrak p)$.
\end{lemma}

\begin{proof}
Nous procédons par récurrence sur le degré de l'extension $\kappa(\mathfrak p) \subset L$.
Si le degré est $1$, nous prenons $R = S$.
En général, s'il existe une sous-extension
$\kappa(\mathfrak p) \subset L' \subset L$, alors la conclusion est acquise par récurrence
sur le degré (en construisant d'abord $R \subset S'$ correspondant
à $L'/\kappa(\mathfrak p)$, puis en construisant $S' \subset S$
correspondant à $L/L'$). Nous pouvons donc supposer que
l'extension $L \supset \kappa(\mathfrak p)$ est engendrée par un seul élément
$\alpha \in L$. Soit $X^d + \sum_{i < d} a_iX^i$ le polynôme minimal
de $\alpha$ sur $\kappa(\mathfrak p)$, de sorte que $a_i \in \kappa(\mathfrak p)$.
Nous pouvons écrire $a_i$ comme l'image
de $f_i/g$ pour certains $f_i, g \in R$ avec $g \not \in \mathfrak p$.
Après avoir remplacé $\alpha$ par $g\alpha$ (et, corrélativement,
remplacé $a_i$ par $g^{d - i}a_i$), nous pouvons supposer que $a_i$ est
l'image d'un élément $f_i \in R$.
Nous prenons alors simplement $S = R[x]/(x^d + \sum f_ix^i)$.
\end{proof}

\begin{lemma}
\label{lemma-cofinal-system-flat}
Soit $A$ un anneau. Posons $\kappa = \max(|A|, \aleph_0)$. Alors toute
$A$-algèbre plate $B$ est la limite inductive filtrante de ses $A$-sous-algèbres plates
$B' \subset B$ telles que $|B'| \leq \kappa$. (Remarquons que $B'$ est
fidèlement plate sur $A$ si $B$ est fidèlement plate sur $A$.)
\end{lemma}

\begin{proof}
Si $B$ est de cardinalité $\leq \kappa$, l'assertion est vraie.
Soit $E \subset B$ une $A$-sous-algèbre telle que $|E| \leq \kappa$.
Nous allons montrer que $E$ est contenue dans une $A$-sous-algèbre plate
$B'$ telle que $|B'| \leq \kappa$. Le lemme s'en déduit, car
(a) tout sous-ensemble fini de $B$ est contenu dans une $A$-sous-algèbre de
cardinalité au plus $\kappa$ et (b) toute paire de $A$-sous-algèbres
de $B$ de cardinalité au plus $\kappa$ est contenue dans une $A$-sous-algèbre
de cardinalité au plus $\kappa$. Détails omis.

\medskip\noindent
Nous allons construire par récurrence une suite de $A$-sous-algèbres
$$
E = E_0 \subset E_1 \subset E_2 \subset \ldots
$$
chacune de cardinalité $\leq \kappa$, et nous montrerons que
$B' = \bigcup E_k$ est plate sur $A$, ce qui achèvera la preuve.

\medskip\noindent
La construction est la suivante. Posons $E_0 = E$.
Étant donné $E_k$ pour $k \geq 0$, considérons l'ensemble $S_k$ des
relations entre des éléments de $E_k$ à coefficients dans $A$.
Ainsi, un élément $s \in S_k$ est la donnée d'un entier $n \geq 1$, de
$a_1, \ldots, a_n \in A$ et de $e_1, \ldots, e_n \in E_k$
tels que $\sum a_i e_i = 0$ dans $E_k$.
La platitude de $A \to B$ implique, d'après le Lemme \ref{lemma-flat-eq},
que, pour tout $s = (n, a_1, \ldots, a_n, e_1, \ldots, e_n) \in S_k$,
nous pouvons choisir
$$
(m_s, b_{s, 1}, \ldots, b_{s, m_s}, a_{s, 11}, \ldots, a_{s, nm_s})
$$
où $m_s \geq 0$ est un entier, $b_{s, j} \in B$, $a_{s, ij} \in A$, et
$$
e_i = \sum\nolimits_j a_{s, ij} b_{s, j}, \forall i,
\quad\text{et}\quad
0 = \sum\nolimits_i a_i a_{s, ij}, \forall j.
$$
Ces choix étant faits, soit $E_{k + 1} \subset B$ la $A$-sous-algèbre
engendrée par
\begin{enumerate}
\item $E_k$ et
\item les éléments $b_{s, 1}, \ldots, b_{s, m_s}$
pour tout $s \in S_k$.
\end{enumerate}
Un argument de théorie des ensembles (omis) montre que $E_{k + 1}$ est de
cardinalité au plus $\kappa$ (on utilise ici que, par récurrence,
$|E_k| \leq \kappa$ et que, par conséquent, la cardinalité de $S_k$ est
également au plus $\kappa$).

\medskip\noindent
Pour montrer que $B' = \bigcup E_k$ est plate sur $A$, considérons
une relation $\sum_{i = 1, \ldots, n} a_i b'_i = 0$ dans $B'$
à coefficients dans $A$. Choisissons $k$ assez grand pour que
$b'_i \in E_k$ pour $i = 1, \ldots, n$. Alors
$(n, a_1, \ldots, a_n, b'_1, \ldots, b'_n) \in S_k$,
et nous voyons donc que la relation est triviale dans $E_{k + 1}$
et, a fortiori, dans $B'$.
Ainsi, $A \to B'$ est plat d'après le Lemme \ref{lemma-flat-eq}.
\end{proof}









\section{Le théorème de structure de Cohen}
\label{section-cohen-structure-theorem}

\noindent
Voici une notion fondamentale d'algèbre commutative.

\begin{definition}
\label{definition-complete-local-ring}
Soit $(R, \mathfrak m)$ un anneau local. Nous disons que $R$ est un
{\it anneau local complet} si le morphisme canonique
$$
R \longrightarrow \lim_n R/\mathfrak m^n
$$
vers le complété de $R$ relativement à $\mathfrak m$ est un
isomorphisme\footnote{Cela inclut la condition
que $\bigcap \mathfrak m^n = (0)$ ; dans certains textes, on l'indique
en disant que $R$ est complet et séparé. Attention : il peut arriver
que le complété $\lim_n R/\mathfrak m^n$ d'un anneau local soit
non complet ; voir
Exemples, Lemme \ref{examples-lemma-noncomplete-completion}.
Cela ne se produit pas lorsque $\mathfrak m$ est de type fini ; voir le
Lemme \ref{lemma-hathat-finitely-generated}, auquel
cas le complété est noethérien ; voir le
Lemme \ref{lemma-completion-Noetherian}.}.
\end{definition}

\noindent
Remarquons qu'un anneau local artinien $R$ est un anneau local complet
car $\mathfrak m_R^n = 0$ pour un certain $n > 0$. Dans cette section,
nous nous concentrons principalement sur les anneaux locaux noethériens complets.

\begin{lemma}
\label{lemma-quotient-complete-local}
Soit $R$ un anneau local noethérien complet.
Tout quotient de $R$ est aussi un anneau local noethérien complet.
Étant donné un morphisme d'anneaux fini $R \to S$, l'anneau $S$ est
un produit d'anneaux locaux noethériens complets.
\end{lemma}

\begin{proof}
L'anneau $S$ est noethérien d'après le Lemme \ref{lemma-Noetherian-permanence}.
En tant que $R$-module, $S$ est complet d'après le Lemme \ref{lemma-completion-tensor}.
Ainsi, $S$ est le produit de ses complétés relativement à ses idéaux maximaux
d'après le Lemme \ref{lemma-completion-finite-extension}.
\end{proof}

\begin{lemma}
\label{lemma-complete-local-ring-Noetherian}
Soit $(R, \mathfrak m)$ un anneau local complet.
Si $\mathfrak m$ est un idéal de type fini, alors
$R$ est noethérien.
\end{lemma}

\begin{proof}
Voir le Lemme \ref{lemma-completion-Noetherian}.
\end{proof}

\begin{definition}
\label{definition-coefficient-ring}
Soit $(R, \mathfrak m)$ un anneau local complet.
Un sous-anneau $\Lambda \subset R$ est
appelé un {\it anneau de coefficients} si les conditions suivantes sont satisfaites :
\begin{enumerate}
\item $\Lambda$ est un anneau local complet d'idéal maximal
$\Lambda \cap \mathfrak m$,
\item le corps résiduel de $\Lambda$ s'envoie isomorphiquement sur le
corps résiduel de $R$, et
\item $\Lambda \cap \mathfrak m = p\Lambda$, où $p$ est la caractéristique
du corps résiduel de $R$.
\end{enumerate}
\end{definition}

\noindent
Faisons quelques remarques sur cette définition. Nous distinguons les
cas suivants :
\begin{enumerate}
\item L'anneau local $R$ contient un corps. Cela se produit
lorsque $\mathbf{Q} \subset R$, ou lorsque $pR = 0$, où $p$ est la
caractéristique de $R/\mathfrak m$. Dans ce cas, un anneau de coefficients
$\Lambda$ est un corps contenu dans $R$ qui s'envoie isomorphiquement sur
$R/\mathfrak m$.
\item La caractéristique de $R/\mathfrak m$ est $p > 0$, mais aucune
puissance de $p$ n'est nulle dans $R$. Dans ce cas, $\Lambda$ est un anneau
de valuation discrète complet d'uniformisante $p$ et de corps résiduel $R/\mathfrak m$.
\item La caractéristique de $R/\mathfrak m$ est $p > 0$ et, pour un certain
$n > 1$, nous avons $p^{n - 1} \not = 0$, $p^n = 0$ dans $R$. Dans ce cas,
$\Lambda$ est un anneau local artinien dont l'idéal maximal est
engendré par $p$ et dont le corps résiduel est $R/\mathfrak m$.
\end{enumerate}
Les anneaux de valuation discrète complets d'uniformisante $p$
ci-dessus jouent un rôle particulier ; donnons-leur le nom suivant.

\begin{definition}
\label{definition-cohen-ring}
Un {\it anneau de Cohen} est un anneau de valuation discrète complet dont
l'uniformisante $p$ est un nombre premier.
\end{definition}

\begin{lemma}
\label{lemma-cohen-rings-exist}
Soit $p$ un nombre premier.
Soit $k$ un corps de caractéristique $p$.
Il existe un anneau de Cohen $\Lambda$ tel que $\Lambda/p\Lambda \cong k$.
\end{lemma}

\begin{proof}
Remarquons d'abord que les entiers $p$-adiques $\mathbf{Z}_p$ forment un anneau de Cohen
de corps résiduel $\mathbf{F}_p$. Soit $k$ un corps arbitraire de caractéristique $p$.
Soit $\mathbf{Z}_p \to R$ un morphisme local plat d'anneaux tel que
$\mathfrak m_R = pR$ et $R/pR = k$ ; voir le
Lemme \ref{lemma-flat-local-given-residue-field}.
D'après le Lemme \ref{lemma-completion-Noetherian}, le complété
$\Lambda = R^\wedge$ est noethérien. Puisque $\Lambda/p\Lambda = R/pR$
est un corps, nous concluons que $\Lambda$ est un
anneau local noethérien complet d'idéal maximal $(p)$
(l'égalité et la complétude résultent du Lemme \ref{lemma-hathat-finitely-generated}).
Pour tout $n$, le morphisme $\mathbf{Z}/p^n\mathbf{Z} \to \Lambda/p^n\Lambda = R/p^nR$
est plat (l'égalité résulte de nouveau du Lemme \ref{lemma-hathat-finitely-generated}).
Ainsi, $\mathbf{Z}_p \to \Lambda$ est plat, par exemple d'après le
Lemme \ref{lemma-flat-module-powers}. Ainsi, $p$ est un
non-diviseur de zéro dans $\Lambda$. Par conséquent, $\Lambda$ est de dimension $\geq 1$
(Lemme \ref{lemma-one-equation}),
et nous concluons que $\Lambda$ est régulier de dimension $1$, c'est-à-dire
un anneau de valuation discrète d'après le Lemme \ref{lemma-characterize-dvr}.
Nous concluons que $\Lambda$ est un anneau de Cohen de corps résiduel $k$.
\end{proof}

\begin{lemma}
\label{lemma-cohen-ring-formally-smooth}
Soit $p > 0$ un nombre premier.
Soit $\Lambda$ un anneau de Cohen dont le corps résiduel est de caractéristique $p$.
Pour tout $n \geq 1$, le morphisme d'anneaux
$$
\mathbf{Z}/p^n\mathbf{Z} \to \Lambda/p^n\Lambda
$$
est formellement lisse.
\end{lemma}

\begin{proof}
Si $n = 1$, cela résulte de la
Proposition \ref{proposition-characterize-separable-field-extensions}.
Pour $n$ quelconque, nous procédons par récurrence sur $n$.
En effet, si $\mathbf{Z}/p^n\mathbf{Z} \to \Lambda/p^n\Lambda$ est
formellement lisse, alors nous pouvons appliquer le Lemme \ref{lemma-lift-formal-smoothness}
au morphisme d'anneaux
$\mathbf{Z}/p^{n + 1}\mathbf{Z} \to \Lambda/p^{n + 1}\Lambda$
et à l'idéal $I = (p^n) \subset \mathbf{Z}/p^{n + 1}\mathbf{Z}$.
\end{proof}

\begin{theorem}[Théorème de structure de Cohen]
\label{theorem-cohen-structure-theorem}
Soit $(R, \mathfrak m)$ un anneau local complet.
\begin{enumerate}
\item $R$ possède un anneau de coefficients (voir la
Définition \ref{definition-coefficient-ring}),
\item si $\mathfrak m$ est un idéal de type fini, alors
$R$ est isomorphe à un quotient
$$
\Lambda[[x_1, \ldots, x_n]]/I
$$
où $\Lambda$ est soit un corps, soit un anneau de Cohen.
\end{enumerate}
\end{theorem}

\begin{proof}
Montrons qu'il existe un anneau de coefficients.
Nous le démontrons d'abord lorsque la caractéristique du corps résiduel $\kappa$
est nulle. Dans ce cas, nous allons démontrer par récurrence
sur $n > 0$ qu'il existe une section
$$
\varphi_n : \kappa \longrightarrow R/\mathfrak m^n
$$
du morphisme canonique $R/\mathfrak m^n \to \kappa = R/\mathfrak m$.
C'est immédiat pour $n = 1$. Si $n > 1$, supposons $\varphi_{n - 1}$ donnée.
L'extension de corps $\kappa/\mathbf{Q}$ est formellement lisse d'après la
Proposition \ref{proposition-characterize-separable-field-extensions}.
Nous pouvons donc trouver la flèche en pointillés
dans le diagramme suivant :
$$
\xymatrix{
R/\mathfrak m^{n - 1} &
R/\mathfrak m^n \ar[l] \\
\kappa \ar[u]^{\varphi_{n - 1}} \ar@{..>}[ru] & \mathbf{Q} \ar[l] \ar[u]
}
$$
Cela démontre l'étape de récurrence. En passant à la limite dans ces morphismes,
$$
\lim_n\ \varphi_n : \kappa \longrightarrow
R = \lim_n\ R/\mathfrak m^n
$$
nous obtenons un morphisme dont l'image est l'anneau de coefficients recherché.

\medskip\noindent
Démontrons ensuite l'existence d'un anneau de coefficients lorsque
la caractéristique du corps résiduel $\kappa$ est $p > 0$.
Choisissons un anneau de Cohen $\Lambda$ tel que $\kappa = \Lambda/p\Lambda$ ;
voir le Lemme \ref{lemma-cohen-rings-exist}. Dans ce cas, nous allons démontrer par
récurrence sur $n > 0$ qu'il existe un morphisme
$$
\varphi_n :
\Lambda/p^n\Lambda
\longrightarrow
R/\mathfrak m^n
$$
dont la composée avec le morphisme de réduction $R/\mathfrak m^n \to \kappa$
donne l'isomorphisme fixé $\Lambda/p\Lambda = \kappa$. C'est immédiat
pour $n = 1$. Si $n > 1$, supposons $\varphi_{n - 1}$ donnée.
Le morphisme d'anneaux $\mathbf{Z}/p^n\mathbf{Z} \to \Lambda/p^n\Lambda$
est formellement lisse d'après le Lemme \ref{lemma-cohen-ring-formally-smooth}.
Nous pouvons donc trouver la flèche en pointillés
dans le diagramme suivant :
$$
\xymatrix{
R/\mathfrak m^{n - 1} &
R/\mathfrak m^n \ar[l] \\
\Lambda/p^n\Lambda \ar[u]^{\varphi_{n - 1}} \ar@{..>}[ru] &
\mathbf{Z}/p^n\mathbf{Z} \ar[l] \ar[u]
}
$$
Cela démontre l'étape de récurrence. En passant à la limite dans ces morphismes,
$$
\lim_n\ \varphi_n :
\Lambda = \lim_n\ \Lambda/p^n\Lambda
\longrightarrow
R = \lim_n\ R/\mathfrak m^n
$$
nous obtenons un morphisme dont l'image est l'anneau de coefficients recherché.

\medskip\noindent
Le dernier énoncé du théorème en résulte aisément. En effet, si
$y_1, \ldots, y_n$ sont des générateurs de l'idéal $\mathfrak m$,
nous pouvons utiliser le morphisme $\Lambda \to R$ que nous venons de construire
pour obtenir un morphisme
$$
\Lambda[[x_1, \ldots, x_n]] \longrightarrow R,
\quad x_i \longmapsto y_i.
$$
Comme les deux membres sont $(x_1, \ldots, x_n)$-adiquement complets,
ce morphisme est surjectif d'après le Lemme \ref{lemma-completion-generalities},
puisqu'il est surjectif modulo $(x_1, \ldots, x_n)$ par
construction.
\end{proof}

\begin{remark}
\label{remark-Noetherian-complete-local-ring-universally-catenary}
Si $k$ est un corps, alors l'anneau de séries formelles $k[[X_1, \ldots, X_d]]$
est un anneau local noethérien complet et régulier de dimension $d$.
Si $\Lambda$ est un anneau de Cohen, alors $\Lambda[[X_1, \ldots, X_d]]$
est un anneau local noethérien complet et régulier de dimension $d + 1$.
Le théorème de structure de Cohen implique donc que tout anneau local
noethérien complet est un quotient d'un anneau local régulier.
En particulier, nous voyons qu'un anneau local noethérien complet est
universellement caténaire ; voir le Lemme \ref{lemma-CM-ring-catenary}
et le Lemme \ref{lemma-regular-ring-CM}.
\end{remark}

\begin{lemma}
\label{lemma-regular-complete-containing-coefficient-field}
Soit $(R, \mathfrak m)$ un anneau local noethérien complet.
Supposons que $R$ soit régulier.
\begin{enumerate}
\item Si $R$ contient soit $\mathbf{F}_p$, soit $\mathbf{Q}$, alors $R$
est isomorphe à un anneau de séries formelles sur son corps résiduel.
\item Si $k$ est un corps et si $k \to R$ est un morphisme d'anneaux induisant
un isomorphisme $k \to R/\mathfrak m$, alors $R$ est isomorphe,
comme $k$-algèbre, à un anneau de séries formelles sur $k$.
\end{enumerate}
\end{lemma}

\begin{proof}
Dans le cas (1), le théorème de structure de Cohen
(Théorème \ref{theorem-cohen-structure-theorem})
fournit un anneau de coefficients qui est nécessairement un corps
et qui s'envoie isomorphiquement sur le corps résiduel. Ainsi,
il suffit de démontrer (2). Dans le cas (2), choisissons
$f_1, \ldots, f_d \in \mathfrak m$ dont les images
forment une base de $\mathfrak m/\mathfrak m^2$, et considérons
le morphisme continu de $k$-algèbres $k[[x_1, \ldots, x_d]] \to R$
qui envoie $x_i$ sur $f_i$. Comme la source et le but sont
$(x_1, \ldots, x_d)$-adiquement complets, ce morphisme est surjectif d'après le
Lemme \ref{lemma-completion-generalities}. D'autre part, il
est nécessairement injectif, car sinon la dimension de
$R$ serait $< d$ d'après le Lemme \ref{lemma-one-equation}.
\end{proof}

\begin{lemma}
\label{lemma-complete-local-Noetherian-domain-finite-over-regular}
Soit $(R, \mathfrak m)$ un anneau local noethérien complet et intègre.
Alors il existe un sous-anneau $R_0 \subset R$ ayant les propriétés suivantes :
\begin{enumerate}
\item $R_0$ est un anneau local régulier et complet,
\item l'inclusion $R_0 \subset R$ est finie et induit un isomorphisme sur les
corps résiduels,
\item $R_0$ est soit isomorphe à $k[[X_1, \ldots, X_d]]$, où $k$
est un corps, soit à $\Lambda[[X_1, \ldots, X_d]]$, où $\Lambda$ est un anneau de Cohen.
\end{enumerate}
\end{lemma}

\begin{proof}
Soit $\Lambda$ un anneau de coefficients de $R$.
Puisque $R$ est intègre, nous voyons que $\Lambda$ est un corps
ou que $\Lambda$ est un anneau de Cohen.

\medskip\noindent
Cas I : $\Lambda = k$ est un corps. Posons $d = \dim(R)$.
Choisissons $x_1, \ldots, x_d \in \mathfrak m$
qui engendrent un idéal de définition $I \subset R$.
(Voir la section \ref{section-dimension}.)
D'après le Lemme \ref{lemma-change-ideal-completion}, nous voyons que $R$
est également $I$-adiquement complet.
Considérons le morphisme $R_0 = k[[X_1, \ldots, X_d]] \to R$
qui envoie $X_i$ sur $x_i$.
Remarquons que $R_0$ est complet relativement à l'idéal
$I_0 = (X_1, \ldots, X_d)$,
et que $R/I_0R \cong R/IR$ est fini sur $k = R_0/I_0$
(car $\dim(R/I) = 0$ ; voir la section \ref{section-dimension}).
Nous concluons donc que $R_0 \to R$ est fini d'après le
Lemme \ref{lemma-finite-over-complete-ring}.
Puisque $\dim(R) = \dim(R_0)$, cela implique que
$R_0 \to R$ est injectif (voir le Lemme \ref{lemma-integral-dim-up}),
et le lemme est démontré.

\medskip\noindent
Cas II : $\Lambda$ est un anneau de Cohen. Posons $d + 1 = \dim(R)$.
Soit $p > 0$ la caractéristique du corps résiduel $k$.
Puisque $R$ est intègre, nous voyons que $p$ est un non-diviseur de zéro dans $R$.
Ainsi, $\dim(R/pR) = d$ ; voir le Lemme \ref{lemma-one-equation}.
Choisissons $x_1, \ldots, x_d \in R$
dont les images engendrent un idéal de définition dans $R/pR$.
Alors $I = (p, x_1, \ldots, x_d)$ est un idéal de définition de $R$.
D'après le Lemme \ref{lemma-change-ideal-completion}, nous voyons que $R$
est également $I$-adiquement complet.
Considérons le morphisme $R_0 = \Lambda[[X_1, \ldots, X_d]] \to R$
qui envoie $X_i$ sur $x_i$.
Remarquons que $R_0$ est complet relativement à l'idéal
$I_0 = (p, X_1, \ldots, X_d)$,
et que $R/I_0R \cong R/IR$ est fini sur $k = R_0/I_0$
(car $\dim(R/I) = 0$ ; voir la section \ref{section-dimension}).
Nous concluons donc que $R_0 \to R$ est fini d'après le
Lemme \ref{lemma-finite-over-complete-ring}.
Puisque $\dim(R) = \dim(R_0)$, cela implique que
$R_0 \to R$ est injectif (voir le Lemme \ref{lemma-integral-dim-up}),
et le lemme est démontré.
\end{proof}





\section{Anneaux japonais}
\label{section-japanese}

\noindent
Dans cette section, nous commençons l'étude de la finitude de la clôture intégrale.

\begin{definition}
\label{definition-N}
\begin{reference}
\cite[Chapter 0, Definition 23.1.1]{EGA}
\end{reference}
Soit $R$ un anneau intègre de corps des fractions $K$.
\begin{enumerate}
\item Nous disons que $R$ est {\it N-1} si la clôture intégrale de $R$ dans $K$
est un $R$-module de type fini.
\item Nous disons que $R$ est {\it N-2} ou {\it japonais} si, pour toute extension
finie $L/K$ de corps, la clôture intégrale de $R$ dans $L$
est finie sur $R$.
\end{enumerate}
\end{definition}

\noindent
Ces notions présentent surtout un intérêt pour les anneaux noethériens,
mais voici un exemple non noethérien.

\begin{example}
\label{example-Japanese-not-Noetherian}
Soit $k$ un corps. L'anneau intègre $R = k[x_1, x_2, x_3, \ldots]$ est N-2,
mais n'est pas noethérien. La raison en est la suivante. Supposons que $R \subset L$
et que le corps $L$ soit une extension finie du corps des fractions de $R$.
Il existe alors un entier $n$ tel que $L$ provienne d'une extension finie
$L_0/k(x_1, \ldots, x_n)$ par adjonction
des éléments (transcendants) $x_{n + 1}, x_{n + 2}$, etc.
Soit $S_0$ la clôture intégrale
de $k[x_1, \ldots, x_n]$ dans $L_0$. D'après la
proposition \ref{proposition-ubiquity-nagata} ci-dessous,
$S_0$ est fini sur $k[x_1, \ldots, x_n]$.
De plus, la clôture intégrale de $R$ dans $L$ est
$S = S_0[x_{n + 1}, x_{n + 2}, \ldots]$ (utiliser le
lemme \ref{lemma-polynomial-domain-normal}) et est
donc finie sur $R$. Le même argument s'applique à
$R = \mathbf{Z}[x_1, x_2, x_3, \ldots]$.
\end{example}

\begin{lemma}
\label{lemma-localize-N}
Soit $R$ un anneau intègre.
Si $R$ est N-1, toute localisation de $R$ l'est aussi.
Il en va de même pour N-2.
\end{lemma}

\begin{proof}
Ces assertions résultent du fait que la formation de la clôture intégrale commute
à la localisation ; voir le lemme \ref{lemma-integral-closure-localize}.
\end{proof}

\begin{lemma}
\label{lemma-Japanese-local}
Soit $R$ un anneau intègre. Soient $f_1, \ldots, f_n \in R$ engendrant
l'idéal unité. Si chacun des anneaux intègres $R_{f_i}$ est N-1, alors $R$ l'est aussi.
Il en va de même pour N-2.
\end{lemma}

\begin{proof}
Supposons que $R_{f_i}$ soit N-2 (ou N-1).
Soit $L$ une extension finie du corps des fractions de $R$ (égale à
ce corps des fractions dans le cas N-1). Soit $S$ la clôture
intégrale de $R$ dans $L$. Le lemme \ref{lemma-integral-closure-localize}
montre que $S_{f_i}$ est la clôture intégrale de $R_{f_i}$ dans $L$.
Par hypothèse, $S_{f_i}$ est donc fini sur $R_{f_i}$.
Le lemme \ref{lemma-cover} montre alors que $S$ est fini sur $R$.
\end{proof}

\begin{lemma}
\label{lemma-quasi-finite-over-Noetherian-japanese}
Soit $R$ un anneau intègre. Soit $R \subset S$ une extension quasi-finie d'anneaux intègres
(par exemple finie). Supposons que $R$ soit N-2 et noethérien. Alors $S$ est N-2.
\end{lemma}

\begin{proof}
Soit $L/K$ l'extension induite des corps des fractions.
Remarquons qu'il s'agit d'une extension finie de corps (par exemple d'après le
lemme \ref{lemma-isolated-point-fibre} (2)
appliqué à la fibre $S \otimes_R K$, et la définition d'un
morphisme d'anneaux quasi-fini).
Soit $S'$ la clôture intégrale de $R$ dans $S$.
Alors $S'$ est contenu dans la clôture intégrale de $R$ dans $L$,
qui est finie sur $R$ par hypothèse. Comme $R$ est noethérien, cela
implique que $S'$ est fini sur $R$.
D'après le lemme \ref{lemma-quasi-finite-open-integral-closure},
il existe des éléments $g_1, \ldots, g_n \in S'$ tels
que $S'_{g_i} \cong S_{g_i}$ et que $g_1, \ldots, g_n$
engendrent l'idéal unité de $S$. Il suffit donc de montrer que
$S'$ est N-2, d'après les lemmes \ref{lemma-localize-N} et \ref{lemma-Japanese-local}.
Nous nous sommes ainsi ramenés au cas où $S$ est fini sur $R$.

\medskip\noindent
Supposons que $R \subset S$ vérifie les hypothèses du lemme et, de plus,
que $S$ soit fini sur $R$. Soit $M$ une extension finie du corps
des fractions de $S$. Alors $M$ est aussi une extension de corps finie
sur $K$, et nous en déduisons que la clôture intégrale $T$ de $R$ dans
$M$ est finie sur $R$. Le lemme \ref{lemma-integral-closure-transitive}
montre que $T$ est aussi la clôture intégrale de $S$ dans $M$, et le résultat suit du
lemme \ref{lemma-integral-permanence}.
\end{proof}

\begin{lemma}
\label{lemma-Laurent-ring-N-1}
Soit $R$ un anneau intègre noethérien.
Si $R[z, z^{-1}]$ est N-1, alors $R$ l'est aussi.
\end{lemma}

\begin{proof}
Soit $R'$ la clôture intégrale de $R$ dans son corps des fractions $K$.
Soit $S'$ la clôture intégrale de $R[z, z^{-1}]$ dans son corps des fractions.
Manifestement, $R' \subset S'$.
Comme $K[z, z^{-1}]$ est un anneau intègre normal, on a $S' \subset K[z, z^{-1}]$.
Supposons que $f_1, \ldots, f_n \in S'$ engendrent $S'$ comme $R[z, z^{-1}]$-module.
Écrivons $f_i = \sum a_{ij}z^j$ (somme finie), avec $a_{ij} \in K$.
Pour tout $x \in R'$, on peut écrire
$$
x = \sum h_i f_i
$$
avec $h_i \in R[z, z^{-1}]$. On voit ainsi que $R'$ est contenu dans le
sous-$R$-module de type fini $\sum Ra_{ij} \subset K$. Comme $R$ est noethérien,
nous en concluons que $R'$ est un $R$-module de type fini.
\end{proof}

\begin{lemma}
\label{lemma-finite-extension-N-2}
Soit $R$ un anneau intègre noethérien, et soit $R \subset S$ une
extension finie d'anneaux intègres. Si $S$ est N-1, alors $R$ l'est aussi.
Si $S$ est N-2, alors $R$ l'est aussi.
\end{lemma}

\begin{proof}
Démonstration omise. (Indication : les clôtures intégrales de $R$ dans des corps d'extension
sont contenues dans les clôtures intégrales de $S$ dans ces corps d'extension.)
\end{proof}

\begin{lemma}
\label{lemma-Noetherian-normal-domain-finite-separable-extension}
Soit $R$ un anneau intègre normal noethérien de corps des fractions $K$.
Soit $L/K$ une extension finie séparable de corps.
Alors la clôture intégrale de $R$ dans $L$ est finie sur $R$.
\end{lemma}

\begin{proof}
Considérons la forme trace
(Corps, définition \ref{fields-definition-trace-pairing})
$$
L \times L \longrightarrow K,
\quad (x, y) \longmapsto \langle x, y\rangle := \text{Trace}_{L/K}(xy).
$$
Comme $L/K$ est séparable, cette forme est non dégénérée
(Corps, lemme \ref{fields-lemma-separable-trace-pairing}).
De plus, si $x \in L$ est entier sur $R$, alors
$\text{Trace}_{L/K}(x)$ appartient à $R$. En effet, le
polynôme minimal de $x$ sur $K$ a ses coefficients dans $R$
(lemme \ref{lemma-minimal-polynomial-normal-domain}),
et $\text{Trace}_{L/K}(x)$ est un
multiple entier de l'un de ces coefficients
(Corps, lemme \ref{fields-lemma-trace-and-norm-from-minimal-polynomial}).
Choisissons $x_1, \ldots, x_n \in L$ entiers sur $R$
et formant une $K$-base de $L$. Alors la clôture intégrale
$S \subset L$ est contenue dans le $R$-module
$$
M = \{y \in L \mid \langle x_i, y\rangle \in R, \ i = 1, \ldots, n\}
$$
L'algèbre linéaire montre que $M \cong R^{\oplus n}$ comme $R$-module.
Ainsi, le sous-module $S \subset R^{\oplus n}$ est un $R$-module de type fini,
puisque $R$ est noethérien.
\end{proof}

\begin{example}
\label{example-bad-invariants}
Le lemme \ref{lemma-Noetherian-normal-domain-finite-separable-extension}
n'est plus valable si l'anneau n'est pas noethérien.
Considérons par exemple l'action de $G = \{+1, -1\}$ sur
$A = \mathbf{C}[x_1, x_2, x_3, \ldots]$, où $-1$ agit en
envoyant $x_i$ sur $-x_i$. L'anneau des invariants $R = A^G$ est
la $\mathbf{C}$-algèbre engendrée par tous les $x_ix_j$. Ainsi,
$R \subset A$ n'est pas une extension finie. Mais $R$ est un anneau intègre normal
de corps des fractions $K = L^G$, le sous-corps des $G$-invariants du corps des fractions
$L$ de $A$. Et, manifestement, $A$ est la clôture intégrale de $R$ dans
$L$.
\end{example}

\noindent
Le lemme suivant peut parfois remplacer le
lemme \ref{lemma-Noetherian-normal-domain-finite-separable-extension}
dans le cas d'extensions purement inséparables.

\begin{lemma}
\label{lemma-Noetherian-normal-domain-insep-extension}
Soit $R$ un anneau intègre normal noethérien de corps des fractions $K$
de caractéristique $p > 0$.
Soit $a \in K$ un élément tel qu'il existe une dérivation
$D : R \to R$ vérifiant $D(a) \not = 0$. Alors la clôture intégrale
de $R$ dans $L = K[x]/(x^p - a)$ est finie sur $R$.
\end{lemma}

\begin{proof}
Après avoir remplacé $x$ par $fx$ et $a$ par $f^pa$, pour un certain $f \in R$,
nous pouvons supposer que $a \in R$. Ainsi, $D(a) \in R$ également. Nous allons montrer
par récurrence sur $i \leq p - 1$ que, si
$$
y = a_0 + a_1x + \ldots + a_i x^i,\quad a_j \in K
$$
est entier sur $R$, alors $D(a)^i a_j \in R$. La clôture intégrale
est donc contenue dans le $R$-module libre de type fini ayant pour base
$D(a)^{-p + 1}x^j$, $j = 0, \ldots, p - 1$. Comme $R$ est noethérien,
cela démontre le lemme.

\medskip\noindent
Si $i = 0$, alors $y = a_0$ est entier sur $R$ si et seulement si $a_0 \in R$,
et l'assertion est vraie. Supposons l'assertion vraie pour un certain $i < p - 1$
et supposons que
$$
y = a_0 + a_1x + \ldots + a_{i + 1} x^{i + 1},\quad a_j \in K
$$
soit entier sur $R$. Alors
$$
y^p = a_0^p + a_1^p a + \ldots + a_{i + 1}^pa^{i + 1}
$$
est un élément de $R$ (puisqu'il appartient à $K$ et est entier sur $R$). En appliquant
$D$, nous obtenons que
$$
(a_1^p + 2a_2^p a + \ldots + (i + 1)a_{i + 1}^p a^i)D(a)
$$
appartient à $R$. Il s'ensuit que
$$
D(a)a_1 + 2D(a) a_2 x + \ldots + (i + 1)D(a) a_{i + 1} x^i
$$
est entier sur $R$. Par hypothèse de récurrence, $D(a)^{i + 1}a_j \in R$
pour $j = 1, \ldots, i + 1$. (Nous utilisons ici que $1, \ldots, i + 1$
sont inversibles.) Par conséquent, $D(a)^{i + 1}a_0$ appartient aussi à $R$, car c'est la
différence de $D(a)^{i + 1}y$ et de $\sum_{j > 0} D(a)^{i + 1}a_jx^j$, qui
sont entiers sur $R$ (puisque $x$ est entier sur $R$, car $a \in R$).
\end{proof}

\begin{lemma}
\label{lemma-domain-char-zero-N-1-2}
Un anneau intègre noethérien dont le corps des fractions est de caractéristique zéro est N-1
si et seulement s'il est N-2 (c'est-à-dire japonais).
\end{lemma}

\begin{proof}
Cela découle immédiatement du
lemme \ref{lemma-Noetherian-normal-domain-finite-separable-extension},
puisque toute extension de corps en caractéristique zéro est séparable.
\end{proof}

\begin{lemma}
\label{lemma-domain-char-p-N-1-2}
Soit $R$ un anneau intègre noethérien de corps des fractions $K$ de
caractéristique $p > 0$. Alors $R$ est N-2 si et seulement si,
pour toute extension finie purement inséparable $L/K$, la clôture
intégrale de $R$ dans $L$ est finie sur $R$.
\end{lemma}

\begin{proof}
Supposons que la clôture intégrale de $R$ dans toute extension de corps
finie purement inséparable de $K$ soit finie.
Soit $L/K$ une extension finie quelconque. Il nous faut montrer que la
clôture intégrale de $R$ dans $L$ est finie sur $R$.
Choisissons une extension finie normale de corps $M/K$
contenant $L$. Comme $R$ est noethérien, il suffit de montrer que
la clôture intégrale de $R$ dans $M$ est finie sur $R$.
D'après le lemme \ref{fields-lemma-normal-case} de la partie Corps,
il existe une sous-extension $M/M_{insep}/K$
telle que $M_{insep}/K$ soit purement inséparable et que $M/M_{insep}$
soit séparable. Par hypothèse, la clôture intégrale $R'$ de $R$ dans
$M_{insep}$ est finie sur $R$. D'après le
lemme \ref{lemma-Noetherian-normal-domain-finite-separable-extension},
la clôture
intégrale $R''$ de $R'$ dans $M$ est finie sur $R'$. Alors $R''$ est fini
sur $R$ d'après le lemme \ref{lemma-finite-transitive}.
Comme $R''$ est aussi la clôture intégrale
de $R$ dans $M$ (voir le lemme \ref{lemma-integral-closure-transitive}), le résultat est démontré.
\end{proof}

\begin{lemma}
\label{lemma-polynomial-ring-N-2}
Soit $R$ un anneau intègre noethérien.
Si $R$ est N-1, alors $R[x]$ est N-1.
Si $R$ est N-2, alors $R[x]$ est N-2.
\end{lemma}

\begin{proof}
Supposons que $R$ soit N-1. Soit $R'$ la clôture intégrale de $R$,
qui est finie sur $R$. Par conséquent, $R'[x]$ est aussi fini sur
$R[x]$. L'anneau $R'[x]$ est normal (voir le
lemme \ref{lemma-polynomial-domain-normal}), donc N-1.
Cela démontre la première assertion.

\medskip\noindent
Pour la seconde assertion, d'après le lemme \ref{lemma-finite-extension-N-2},
il suffit de montrer que $R'[x]$ est N-2. Autrement dit, nous pouvons
supposer que $R$ est un anneau intègre normal N-2. En caractéristique zéro,
le lemme \ref{lemma-domain-char-zero-N-1-2} permet de conclure.
En caractéristique $p > 0$, nous devons montrer que la clôture
intégrale de $R[x]$ dans toute extension finie purement inséparable
$L/K(x)$ est finie, où $K$ est le corps des fractions de $R$. Il
existe une extension finie purement inséparable de corps $L'/K$
et $q = p^e$ tels que $L \subset L'(x^{1/q})$ ; certains détails sont omis.
Comme $R[x]$ est noethérien, il suffit de montrer que la clôture intégrale de $R[x]$
dans $L'(x^{1/q})$ est finie sur $R[x]$. Cette clôture intégrale
est égale à $R'[x^{1/q}]$, où $R \subset R' \subset L'$ est la clôture intégrale
de $R$ dans $L'$.
Comme $R$ est N-2, l'anneau $R'$ est fini sur $R$ et, par suite,
$R'[x^{1/q}]$ est fini sur $R[x]$.
\end{proof}

\begin{lemma}
\label{lemma-openness-normal-locus}
Soit $R$ un anneau intègre noethérien.
S'il existe un élément non nul $f \in R$ tel que $R_f$ soit normal,
alors
$$
U = \{\mathfrak p \in \Spec(R) \mid R_{\mathfrak p} \text{ est normal}\}
$$
est ouvert dans $\Spec(R)$.
\end{lemma}

\begin{proof}
Il est clair que l'ouvert principal $D(f)$ est contenu dans $U$.
Le critère de normalité de Serre, lemme \ref{lemma-criterion-normal}, montre que
$\mathfrak p \not \in U$ implique qu'il existe
$\mathfrak q \subset \mathfrak p$ pour lequel
l'un des cas suivants se présente :
\begin{enumerate}
\item Cas I : $\text{depth}(R_{\mathfrak q}) < 2$
et $\dim(R_{\mathfrak q}) \geq 2$, ou bien
\item Cas II : $R_{\mathfrak q}$ n'est pas régulier
et $\dim(R_{\mathfrak q}) = 1$.
\end{enumerate}
Cela signifie en particulier que $R_{\mathfrak q}$ n'est pas
normal, et donc que $f \in \mathfrak q$. Dans le cas I, on a
$\text{depth}(R_{\mathfrak q}) =
\text{depth}(R_{\mathfrak q}/fR_{\mathfrak q}) + 1$.
Un tel idéal premier $\mathfrak q$ est donc exactement un idéal premier associé
immergé de $R/fR$. Dans le cas II, $\mathfrak q$ est un idéal premier associé
de $R/fR$ de hauteur 1. Il existe donc un ensemble fini $E$
de tels premiers $\mathfrak q$ (voir le lemme \ref{lemma-finite-ass}), et
$$
\Spec(R) \setminus U
=
\bigcup\nolimits_{\mathfrak q \in E} V(\mathfrak q)
$$
comme souhaité.
\end{proof}

\begin{lemma}
\label{lemma-characterize-N-1}
Soit $R$ un anneau intègre noethérien. Alors $R$ est N-1 si et seulement si les deux
conditions suivantes sont satisfaites :
\begin{enumerate}
\item il existe un élément non nul $f \in R$ tel que $R_f$ soit normal, et
\item pour tout idéal maximal $\mathfrak m \subset R$,
l'anneau local $R_{\mathfrak m}$ est N-1.
\end{enumerate}
\end{lemma}

\begin{proof}
Supposons d'abord que $R$ soit N-1. Soit $R'$ la clôture intégrale de $R$ dans son
corps des fractions $K$. Par hypothèse, on peut trouver $x_1, \ldots, x_n$ dans $R'$
qui engendrent $R'$ comme $R$-module. Comme $R' \subset K$, on peut trouver
des éléments non nuls $f_i \in R$ tels que $f_i x_i \in R$. Alors $R_f \cong R'_f$,
où $f = f_1 \ldots f_n$. Ainsi, $R_f$ est normal et la condition (1) est satisfaite.
La condition (2) résulte du lemme \ref{lemma-localize-N}.

\medskip\noindent
Supposons (1) et (2). Soit $K$ le corps des fractions de $R$.
Supposons que $R \subset R' \subset K$ soit une extension finie
de $R$ contenue dans $K$. Remarquons que $R_f = R'_f$, puisque
$R_f$ est déjà normal. Le lemme \ref{lemma-openness-normal-locus} montre donc
que l'ensemble des premiers
$\mathfrak p' \in \Spec(R')$ tels que $R'_{\mathfrak p'}$ ne soit pas normal
est fermé dans $\Spec(R')$. Comme $\Spec(R') \to \Spec(R)$
est fermé, l'image de cet ensemble est fermée dans $\Spec(R)$.
Pour un tel anneau $R'$, notons $Z_{R'} \subset \Spec(R)$ cette image.

\medskip\noindent
Choisissons un idéal maximal $\mathfrak m \subset R$.
Soit $R_{\mathfrak m} \subset R_{\mathfrak m}'$ la clôture
intégrale de l'anneau local dans $K$. Par hypothèse, il s'agit
d'une extension finie d'anneaux. D'après le lemme \ref{lemma-integral-closure-localize},
on peut trouver un nombre fini
d'éléments $x_1, \ldots, x_n \in K$ entiers sur $R$ tels que
$R_{\mathfrak m}'$ soit engendré par $x_1, \ldots, x_n$ sur $R_{\mathfrak m}$.
Soit $R' = R[x_1, \ldots, x_n] \subset K$. Avec ce choix, il est clair
que $\mathfrak m \not \in Z_{R'}$.

\medskip\noindent
Comme $\Spec(R)$ est quasi-compact, ce qui précède montre que nous pouvons
trouver une famille finie $R \subset R'_i \subset K$ telle que
$\bigcap Z_{R'_i} = \emptyset$. Soit $R'$ le sous-anneau de $K$
engendré par tous ces sous-anneaux. Il est fini sur $R$. De plus, $Z_{R'} = \emptyset$.
En effet, tout premier $\mathfrak p'$ est au-dessus d'un premier $\mathfrak p'_i$
tel que $(R'_i)_{\mathfrak p'_i}$ soit normal. Cela implique
que $R'_{\mathfrak p'} = (R'_i)_{\mathfrak p'_i}$ est lui aussi normal.
Ainsi, $R'$ est normal ; autrement dit,
$R'$ est la clôture intégrale de $R$ dans $K$.
\end{proof}

\begin{lemma}[Tate]
\label{lemma-tate-japanese}
\begin{reference}
\cite[Theorem 23.1.3]{EGA}
\end{reference}
Soit $R$ un anneau.
Soit $x \in R$.
Supposons que
\begin{enumerate}
\item $R$ soit un anneau intègre normal noethérien,
\item $R/xR$ soit un anneau intègre et N-2,
\item $R \cong \lim_n R/x^nR$ soit complet pour la topologie $x$-adique.
\end{enumerate}
Alors $R$ est N-2.
\end{lemma}

\begin{proof}
Nous pouvons supposer que $x \not = 0$, car sinon le lemme est trivial.
Soit $K$ le corps des fractions de $R$. Si la caractéristique de $K$
est nulle, le lemme résulte de (1) ; voir le
lemme \ref{lemma-domain-char-zero-N-1-2}. Nous pouvons donc supposer
que la caractéristique de $K$ est $p > 0$ et appliquer le
lemme \ref{lemma-domain-char-p-N-1-2}. Ainsi, pour toute extension de corps $L/K$
finie et purement inséparable, il nous faut donc montrer
que la clôture intégrale $S$ de $R$ dans $L$ est finie sur $R$.

\medskip\noindent
Soit $q$ une puissance de $p$ telle que $L^q \subset K$.
En agrandissant $L$ si nécessaire, nous pouvons supposer qu'il existe
un élément $y \in L$ tel que $y^q = x$. Comme $R \to S$
induit un homéomorphisme entre les spectres (voir le lemme \ref{lemma-p-ring-map}),
il existe un unique idéal premier $\mathfrak q \subset S$ au-dessus
de l'idéal premier $\mathfrak p = xR$. Il est clair que
$$
\mathfrak q = \{f \in S \mid f^q \in \mathfrak p\} = yS
$$
puisque $y^q = x$. Remarquons que $R_{\mathfrak p}$ est un anneau de
valuation discrète d'après le lemme \ref{lemma-characterize-dvr}. D'autre part,
$S_{\mathfrak q}$ est noethérien d'après le théorème de Krull-Akizuki
(lemme \ref{lemma-krull-akizuki}). Nous en concluons alors que
$S_{\mathfrak q}$ est un anneau de valuation discrète, à nouveau d'après le
lemme \ref{lemma-characterize-dvr}.
Le lemme \ref{lemma-finite-extension-residue-fields-dimension-1} montre que
$\kappa(\mathfrak q)/\kappa(\mathfrak p)$ est
une extension finie de corps. La clôture intégrale
$S' \subset \kappa(\mathfrak q)$ de $R/xR$ est donc finie sur
$R/xR$ par l'hypothèse (2). Comme $S/yS \subset S'$, cela implique
que $S/yS$ est fini sur $R$. Remarquons que $S/y^nS$ possède une filtration
finie dont les sous-quotients sont les modules
$y^iS/y^{i + 1}S \cong S/yS$. On voit donc que chaque $S/y^nS$
est fini sur $R$. En particulier, $S/xS$ est fini sur $R$.
De plus, il est clair que $\bigcap x^nS = (0)$, car la puissance $q$-ième d'un élément
de l'intersection appartient à $\bigcap x^nR = (0)$
(lemme \ref{lemma-intersect-powers-ideal-module-zero}).
Nous pouvons donc appliquer le lemme \ref{lemma-finite-over-complete-ring} pour conclure
que $S$ est fini sur $R$, ce qui termine la démonstration.
\end{proof}

\begin{lemma}
\label{lemma-power-series-over-N-2}
Soit $R$ un anneau.
Si $R$ est un anneau intègre noethérien et N-2, alors $R[[x]]$ l'est aussi.
\end{lemma}

\begin{proof}
Remarquons que $R[[x]]$ est noethérien d'après le
lemme \ref{lemma-Noetherian-power-series}.
Soit $R' \supset R$ la clôture intégrale de $R$ dans son corps des
fractions. Comme $R$ est N-2, elle est finie sur $R$. Par conséquent, $R'[[x]]$
est fini sur $R[[x]]$. D'après le
lemme \ref{lemma-power-series-over-Noetherian-normal-domain},
$R'[[x]]$ est un anneau intègre normal.
Appliquons le lemme \ref{lemma-tate-japanese} à
l'élément $x \in R'[[x]]$ : l'anneau $R'[[x]]$ est N-2. Le
lemme \ref{lemma-finite-extension-N-2} montre alors que $R[[x]]$ est N-2.
\end{proof}





\section{Anneaux de Nagata}
\label{section-nagata}

\noindent
Voici la définition.

\begin{definition}
\label{definition-nagata}
Soit $R$ un anneau.
\begin{enumerate}
\item Nous disons que $R$ est {\it universellement japonais} si, pour tout morphisme
d'anneaux de type fini $R \to S$ avec $S$ intègre, l'anneau $S$ est N-2
(c'est-à-dire japonais).
\item Nous disons que $R$ est un {\it anneau de Nagata} si $R$ est noethérien et si,
pour tout idéal premier $\mathfrak p$, l'anneau $R/\mathfrak p$ est N-2.
\end{enumerate}
\end{definition}

\noindent
Il est clair qu'un anneau noethérien universellement japonais est un anneau de Nagata.
Notre objectif est de montrer qu'un anneau de Nagata est universellement japonais. Ce n'est
pas du tout évident et cela demande un certain travail. Mais commençons par un lemme
utile.

\begin{lemma}
\label{lemma-nagata-in-reduced-finite-type-finite-integral-closure}
Soit $R$ un anneau de Nagata.
Soit $R \to S$ un morphisme essentiellement de type fini, avec $S$ réduit.
Alors la clôture intégrale $A$ de $R$ dans $S$ est finie sur $R$.
\end{lemma}

\begin{proof}
Comme $S$ est essentiellement de type fini sur $R$, il est noethérien et
possède un nombre fini d'idéaux premiers minimaux $\mathfrak q_1, \ldots, \mathfrak q_m$ ;
voir le lemme \ref{lemma-Noetherian-irreducible-components}.
Puisque $S$ est réduit, on a $S \subset \prod S_{\mathfrak q_i}$,
et chacun des $S_{\mathfrak q_i} = K_i$ est un corps ; voir les
lemmes \ref{lemma-total-ring-fractions-no-embedded-points}
et \ref{lemma-minimal-prime-reduced-ring}.
Il suffit de montrer que la clôture intégrale
$A_i'$ de $R$ dans chacun des $K_i$ est finie sur $R$.
Cela suffit parce que $R$ est noethérien et que $A \subset \prod A_i'$.
Soit $\mathfrak p_i \subset R$ l'idéal premier de $R$
correspondant à $\mathfrak q_i$.
Comme $S$ est essentiellement de type fini sur $R$, on voit que
$K_i = S_{\mathfrak q_i} = \kappa(\mathfrak q_i)$ est une extension de corps de type
fini de $\kappa(\mathfrak p_i)$. Par conséquent, la clôture algébrique
$L_i$ de $\kappa(\mathfrak p_i)$ dans $K_i$
est finie sur $\kappa(\mathfrak p_i)$ ; voir
Corps, lemme \ref{fields-lemma-algebraic-closure-in-finitely-generated}.
Il est clair que $A_i'$ est la clôture intégrale de $R/\mathfrak p_i$
dans $L_i$, et le résultat découle donc de la définition d'un anneau de Nagata.
\end{proof}

\begin{lemma}
\label{lemma-check-universally-japanese}
Soit $R$ un anneau.
Pour vérifier que $R$ est universellement japonais, il suffit de montrer ceci :
si $R \to S$ est de type fini et si $S$ est intègre, alors $S$ est N-1.
\end{lemma}

\begin{proof}
Supposons en effet satisfaite la condition du lemme.
Soit $R \to S$ un morphisme d'anneaux de type fini, avec $S$ intègre.
Soit $L$ une extension finie du corps des fractions de $S$.
Il existe alors une extension finie d'anneaux $S \subset S' \subset L$
telle que $L$ soit le corps des fractions de $S'$.
Par hypothèse, $S'$ est N-1, et la clôture intégrale
$S''$ de $S'$ dans $L$ est donc finie sur $S'$. Ainsi, $S''$ est fini
sur $S$ (lemme \ref{lemma-finite-transitive})
et $S''$ est la clôture intégrale de $S$ dans $L$
(lemme \ref{lemma-integral-closure-transitive}).
Nous en concluons que $R$ est universellement japonais.
\end{proof}

\begin{lemma}
\label{lemma-universally-japanese}
Si $R$ est universellement japonais, toute algèbre essentiellement de type fini
sur $R$ est universellement japonaise.
\end{lemma}

\begin{proof}
Le cas d'une algèbre de type fini sur $R$ résulte immédiatement de
la définition. Le cas général s'obtient en appliquant le
lemme \ref{lemma-localize-N}.
\end{proof}

\begin{lemma}
\label{lemma-quasi-finite-over-nagata}
Soit $R$ un anneau de Nagata.
Si $R \to S$ est un morphisme d'anneaux quasi-fini (par exemple fini),
alors $S$ est lui aussi un anneau de Nagata.
\end{lemma}

\begin{proof}
Remarquons d'abord que $S$ est noethérien, car $R$ est noethérien et qu'un morphisme
d'anneaux quasi-fini est de type fini.
Soit $\mathfrak q \subset S$ un idéal premier et posons
$\mathfrak p = R \cap \mathfrak q$. Alors
$R/\mathfrak p \subset S/\mathfrak q$ est quasi-fini, et
nous en concluons que $S/\mathfrak q$ est N-2 d'après le
lemme \ref{lemma-quasi-finite-over-Noetherian-japanese},
comme souhaité.
\end{proof}

\begin{lemma}
\label{lemma-nagata-localize}
Toute localisation d'un anneau de Nagata est un anneau de Nagata.
\end{lemma}

\begin{proof}
Cela résulte immédiatement du lemme \ref{lemma-localize-N}.
\end{proof}

\begin{lemma}
\label{lemma-nagata-local}
Soit $R$ un anneau. Soient $f_1, \ldots, f_n \in R$ engendrant
l'idéal unité.
\begin{enumerate}
\item Si chacun des $R_{f_i}$ est universellement japonais, alors $R$ l'est aussi.
\item Si chacun des $R_{f_i}$ est un anneau de Nagata, alors $R$ l'est aussi.
\end{enumerate}
\end{lemma}

\begin{proof}
Soit $\varphi : R \to S$ un morphisme d'anneaux de type fini tel que $S$ soit intègre.
Alors $\varphi(f_1), \ldots, \varphi(f_n)$ engendrent l'idéal unité
de $S$. Par conséquent, si chacun des $S_{f_i} = S_{\varphi(f_i)}$ est N-1, alors
$S$ l'est aussi ; voir le lemme \ref{lemma-Japanese-local}. Cela démontre (1).

\medskip\noindent
Si chacun des $R_{f_i}$ est un anneau de Nagata, chacun des $R_{f_i}$ est noethérien,
et $R$ est donc noethérien ; voir le lemme \ref{lemma-cover}. De plus, si
$\mathfrak p \subset R$ est un idéal premier, chacun des anneaux
$R_{f_i}/\mathfrak pR_{f_i} = (R/\mathfrak p)_{f_i}$ est N-2,
et nous en concluons que $R/\mathfrak p$ est N-2 d'après le
lemme \ref{lemma-Japanese-local}. Cela démontre (2).
\end{proof}

\begin{lemma}
\label{lemma-Noetherian-complete-local-Nagata}
Tout anneau local noethérien complet est un anneau de Nagata.
\end{lemma}

\begin{proof}
Soit $R$ un anneau local noethérien complet.
Soit $\mathfrak p \subset R$ un idéal premier.
Alors $R/\mathfrak p$ est également un anneau local noethérien complet ;
voir le lemme \ref{lemma-quotient-complete-local}.
Il suffit donc de montrer qu'un anneau intègre local noethérien
complet $R$ est N-2. D'après les
lemmes \ref{lemma-quasi-finite-over-Noetherian-japanese}
et \ref{lemma-complete-local-Noetherian-domain-finite-over-regular},
nous nous ramenons au cas $R = k[[X_1, \ldots, X_d]]$, où $k$ est un corps, ou au cas
$R = \Lambda[[X_1, \ldots, X_d]]$, où $\Lambda$ est un anneau de Cohen.

\medskip\noindent
Dans le cas $k[[X_1, \ldots, X_d]]$, le lemme
\ref{lemma-power-series-over-N-2} nous ramène à l'assertion qu'un corps est N-2. C'est clair.
Dans le cas $\Lambda[[X_1, \ldots, X_d]]$, nous nous ramenons à l'assertion
qu'un anneau de Cohen $\Lambda$ est N-2. En appliquant une fois encore le lemme
\ref{lemma-tate-japanese} avec $x = p \in \Lambda$, nous nous ramenons encore
au cas d'un corps. Le résultat est donc démontré.
\end{proof}

\begin{definition}
\label{definition-analytically-unramified}
Soit $(R, \mathfrak m)$ un anneau local noethérien.
Nous disons que $R$ est {\it analytiquement non ramifié} si son complété
$R^\wedge = \lim_n R/\mathfrak m^n$ est réduit.
On dit qu'un idéal premier $\mathfrak p \subset R$ est
{\it analytiquement non ramifié} si $R/\mathfrak p$ est analytiquement
non ramifié.
\end{definition}

\noindent
À ce stade, nous savons que
les assertions suivantes sont vraies pour tout anneau local noethérien $R$ :
le morphisme $R \to R^\wedge$ est un homomorphisme local fidèlement plat d'anneaux
(lemme \ref{lemma-completion-faithfully-flat}).
Le complété $R^\wedge$ est noethérien
(lemme \ref{lemma-completion-Noetherian})
et complet (lemme \ref{lemma-completion-complete}).
Le complété $R^\wedge$ est donc un anneau de Nagata
(lemme \ref{lemma-Noetherian-complete-local-Nagata}).
En outre, nous avons vu dans la section \ref{section-cohen-structure-theorem}
que $R^\wedge$ est
un quotient d'un anneau local régulier
(théorème \ref{theorem-cohen-structure-theorem}), et qu'il est donc
universellement caténaire
(remarque \ref{remark-Noetherian-complete-local-ring-universally-catenary}).

\begin{lemma}
\label{lemma-analytically-unramified-easy}
Soit $(R, \mathfrak m)$ un anneau local noethérien.
\begin{enumerate}
\item Si $R$ est analytiquement non ramifié, alors $R$ est réduit.
\item Si $R$ est analytiquement non ramifié, alors tout idéal premier minimal de
$R$ est analytiquement non ramifié.
\item Si $R$ est réduit, d'idéaux premiers minimaux
$\mathfrak q_1, \ldots, \mathfrak q_t$, et si chaque $\mathfrak q_i$
est analytiquement non ramifié, alors $R$ est analytiquement non ramifié.
\item Si $R$ est analytiquement non ramifié, alors la clôture intégrale
de $R$ dans son anneau total des fractions $Q(R)$ est finie sur $R$.
\item Si $R$ est intègre et analytiquement non ramifié, alors $R$ est N-1.
\end{enumerate}
\end{lemma}

\begin{proof}
Dans cette démonstration, nous utiliserons les remarques qui suivent immédiatement la
définition \ref{definition-analytically-unramified}.
Comme $R \to R^\wedge$ est un homomorphisme local fidèlement plat d'anneaux,
il est injectif, d'où (1).

\medskip\noindent
Soit $\mathfrak q$ un idéal premier minimal de $R$, et supposons $R$
analytiquement non ramifié.
Alors $\mathfrak q$ est un idéal premier associé
à $R$ (voir la
proposition \ref{proposition-minimal-primes-associated-primes}).
Il existe donc $f \in R$
tel que $\{x \in R \mid fx = 0\} = \mathfrak q$.
Remarquons que $(R/\mathfrak q)^\wedge = R^\wedge/\mathfrak q^\wedge$
et que $\{x \in R^\wedge \mid fx = 0\} = \mathfrak q^\wedge$,
car la complétion est exacte (lemme \ref{lemma-completion-flat}).
Si $x \in R^\wedge$ est tel
que $x^2 \in \mathfrak q^\wedge$, alors $fx^2 = 0$, donc
$(fx)^2 = 0$, donc $fx = 0$, donc $x \in \mathfrak q^\wedge$.
Ainsi, $\mathfrak q$ est analytiquement non ramifié et (2) est satisfaite.

\medskip\noindent
Supposons $R$ réduit, d'idéaux premiers minimaux
$\mathfrak q_1, \ldots, \mathfrak q_t$, et chaque $\mathfrak q_i$
analytiquement non ramifié. Alors
$R \to R/\mathfrak q_1 \times \ldots \times R/\mathfrak q_t$ est
injectif. Comme la complétion est exacte (voir le lemme \ref{lemma-completion-flat}),
nous voyons que
$R^\wedge \subset (R/\mathfrak q_1)^\wedge \times \ldots \times
(R/\mathfrak q_t)^\wedge$. L'assertion (3) est donc claire.

\medskip\noindent
Supposons $R$ analytiquement non ramifié.
Soient $\mathfrak p_1, \ldots, \mathfrak p_s$ les idéaux premiers minimaux
de $R^\wedge$. On a alors
$$
Q(R^\wedge) =
R^\wedge_{\mathfrak p_1} \times \ldots \times R^\wedge_{\mathfrak p_s}
$$
où chaque $R^\wedge_{\mathfrak p_i}$ est un corps,
puisque $R^\wedge$ est réduit (voir le
lemme \ref{lemma-total-ring-fractions-no-embedded-points}).
Par conséquent, la clôture intégrale $S$ de $R^\wedge$
dans $Q(R^\wedge)$ est $S = S_1 \times \ldots \times S_s$, où
$S_i$ est la clôture intégrale de $R^\wedge/\mathfrak p_i$ dans son corps des
fractions. D'après le lemme \ref{lemma-Noetherian-complete-local-Nagata}, l'anneau
$S_i$ est fini sur $R^\wedge/\mathfrak p_i$. Ainsi,
$S$ est fini sur $R^\wedge$.
Notons $R'$ la clôture intégrale de $R$ dans $Q(R)$.
Comme $R \to R^\wedge$ est plat, on a
$R' \otimes_R R^\wedge \subset Q(R) \otimes_R R^\wedge \subset Q(R^\wedge)$.
De plus, $R' \otimes_R R^\wedge$ est entier sur $R^\wedge$
(lemme \ref{lemma-base-change-integral}).
Ainsi, $R' \otimes_R R^\wedge \subset S$ est un sous-$R^\wedge$-module.
Comme $R^\wedge$ est noethérien, c'est un $R^\wedge$-module de type fini.
Nous pouvons donc trouver $f_1, \ldots, f_n \in R'$ tels que
$R' \otimes_R R^\wedge$ soit engendré par les éléments $f_i \otimes 1$
comme $R^\wedge$-module.
La fidèle platitude montre que $R'$ est engendré par $f_1, \ldots, f_n$
comme $R$-module. Cela démontre (4).

\medskip\noindent
L'assertion (5) est un cas particulier de l'assertion (4).
\end{proof}

\begin{lemma}
\label{lemma-codimension-1-analytically-unramified}
Soit $R$ un anneau local noethérien.
Soit $\mathfrak p \subset R$ un idéal premier.
Supposons que
\begin{enumerate}
\item $R_{\mathfrak p}$ soit un anneau de valuation discrète, et que
\item $\mathfrak p$ soit analytiquement non ramifié.
\end{enumerate}
Alors, pour tout idéal premier associé $\mathfrak q$ de $R^\wedge/\mathfrak pR^\wedge$,
l'anneau local $(R^\wedge)_{\mathfrak q}$ est un anneau de valuation discrète.
\end{lemma}

\begin{proof}
L'hypothèse (2) signifie que $R^\wedge/\mathfrak pR^\wedge$ est un anneau réduit.
Par conséquent, un idéal premier associé $\mathfrak q \subset R^\wedge$
de $R^\wedge/\mathfrak pR^\wedge$
est la même chose qu'un idéal premier minimal au-dessus de $\mathfrak pR^\wedge$.
On voit en particulier que l'idéal maximal de $(R^\wedge)_{\mathfrak q}$
est $\mathfrak p(R^\wedge)_{\mathfrak q}$.
Choisissons $x \in R$ tel que $xR_{\mathfrak p} = \mathfrak pR_{\mathfrak p}$.
D'après ce qui précède, $x \in (R^\wedge)_{\mathfrak q}$ engendre
l'idéal maximal. Comme $R \to R^\wedge$ est fidèlement plat, on voit que
$x$ est un non-diviseur de zéro dans $(R^\wedge)_{\mathfrak q}$.
Le résultat s'ensuit.
\end{proof}

\begin{lemma}
\label{lemma-criterion-analytically-unramified}
Soit $(R, \mathfrak m)$ un anneau intègre local noethérien.
Soit $x \in \mathfrak m$. Supposons que
\begin{enumerate}
\item $x \not = 0$,
\item $R/xR$ n'ait aucun idéal premier associé immergé, et que
\item pour tout idéal premier associé $\mathfrak p \subset R$
à $R/xR$, on ait
\begin{enumerate}
\item l'anneau local $R_{\mathfrak p}$ est régulier, et
\item $\mathfrak p$ est analytiquement non ramifié.
\end{enumerate}
\end{enumerate}
Alors $R$ est analytiquement non ramifié.
\end{lemma}

\begin{proof}
Soient $\mathfrak p_1, \ldots, \mathfrak p_t$ les idéaux premiers associés
au $R$-module $R/xR$. Puisque $R/xR$ n'a aucun idéal premier associé immergé, on
voit que chaque $\mathfrak p_i$ est de hauteur $1$ et est un idéal premier
minimal au-dessus de $(x)$.
Pour chaque $i$, soient $\mathfrak q_{i1}, \ldots, \mathfrak q_{is_i}$
les idéaux premiers associés au $R^\wedge$-module
$R^\wedge/\mathfrak p_iR^\wedge$.
D'après le lemme \ref{lemma-codimension-1-analytically-unramified},
on voit que $(R^\wedge)_{\mathfrak q_{ij}}$ est régulier.
D'après le lemme \ref{lemma-bourbaki}, on a
$$
\text{Ass}_{R^\wedge}(R^\wedge/xR^\wedge)
=
\bigcup\nolimits_{\mathfrak p \in \text{Ass}_R(R/xR)}
\text{Ass}_{R^\wedge}(R^\wedge/\mathfrak pR^\wedge)
=
\{\mathfrak q_{ij}\}.
$$
Soit $y \in R^\wedge$ tel que $y^2 = 0$.
Comme $(R^\wedge)_{\mathfrak q_{ij}}$ est régulier, et donc intègre
(lemme \ref{lemma-regular-domain}),
on voit que $y$ s'envoie sur zéro dans $(R^\wedge)_{\mathfrak q_{ij}}$.
Il s'ensuit que $y$ s'envoie sur zéro dans $R^\wedge/xR^\wedge$ d'après le
lemme \ref{lemma-zero-at-ass-zero}.
Ainsi, $y = xy'$. Comme $x$ est un non-diviseur de zéro (puisque $R \to R^\wedge$ est plat),
on voit que $(y')^2 = 0$. Nous en concluons que
$y \in \bigcap x^nR^\wedge = (0)$
(lemme \ref{lemma-intersect-powers-ideal-module-zero}).
\end{proof}

\begin{lemma}
\label{lemma-local-nagata-domain-analytically-unramified}
Soit $(R, \mathfrak m)$ un anneau local.
Si $R$ est noethérien, intègre et de Nagata, alors $R$ est
analytiquement non ramifié.
\end{lemma}

\begin{proof}
Raisonnons par récurrence sur $\dim(R)$.
Le cas $\dim(R) = 0$ est trivial. Supposons donc $\dim(R) = d$ et le
lemme vrai pour tous les anneaux intègres noethériens de Nagata de dimension $< d$.

\medskip\noindent
Soit $R \subset S$ la clôture intégrale
de $R$ dans le corps des fractions de $R$. Par hypothèse, $S$ est un
$R$-module de type fini. D'après le lemme \ref{lemma-quasi-finite-over-nagata},
$S$ est un anneau de Nagata. D'après le lemme \ref{lemma-integral-sub-dim-equal}, on a
$\dim(R) = \dim(S)$.
Soient $\mathfrak m_1, \ldots, \mathfrak m_t$ les idéaux
maximaux de $S$. Chacun d'eux est au-dessus de l'idéal maximal $\mathfrak m$
de $R$. De plus,
$$
(\mathfrak m_1 \cap \ldots \cap \mathfrak m_t)^n \subset \mathfrak mS
$$
pour $n$ assez grand, puisque $S/\mathfrak mS$ est artinien.
D'après le lemme \ref{lemma-completion-flat}, $R^\wedge \to S^\wedge$
est injectif et, en combinant le théorème chinois des restes,
lemme \ref{lemma-chinese-remainder}, avec le
lemme \ref{lemma-change-ideal-completion}, on obtient
$S^\wedge = \prod S^\wedge_i$, où $S^\wedge_i$
est le complété de $S$ pour la topologie $\mathfrak m_i$-adique.
Il suffit donc de montrer que $S_{\mathfrak m_i}$ est analytiquement non ramifié.
Autrement dit, nous nous sommes ramenés au cas où $R$ est un anneau intègre
normal noethérien de Nagata.

\medskip\noindent
Supposons que $R$ soit un anneau intègre de Nagata, normal, noethérien et local.
Choisissons un élément non nul $x \in \mathfrak m$ de l'idéal maximal.
Nous allons appliquer le lemme \ref{lemma-criterion-analytically-unramified}.
Nous devons vérifier les conditions (1), (2), (3)(a) et (3)(b).
La condition (1) est claire.
L'anneau $R/xR$ n'a aucun idéal premier associé immergé d'après le
lemme \ref{lemma-normal-domain-intersection-localizations-height-1}.
La condition (2) est donc satisfaite. Le même lemme montre aussi que tout idéal
premier associé $\mathfrak p$ à $R/xR$ est de hauteur $1$.
Ainsi, $R_{\mathfrak p}$ est un anneau intègre normal de dimension $1$,
donc régulier (lemme \ref{lemma-characterize-dvr}). La condition (3)(a) est ainsi satisfaite.
Enfin, la condition (3)(b) résulte de l'hypothèse de récurrence, puisque
$R/\mathfrak p$ est un anneau de Nagata (d'après le lemme \ref{lemma-quasi-finite-over-nagata}
ou directement par définition).
Nous en concluons que $R$ est analytiquement non ramifié.
\end{proof}

\begin{lemma}
\label{lemma-local-nagata-and-analytically-unramified}
Soit $(R, \mathfrak m)$ un anneau local noethérien. Les conditions suivantes
sont équivalentes :
\begin{enumerate}
\item $R$ est un anneau de Nagata,
\item pour tout morphisme fini $R \to S$, avec $S$ intègre, et tout idéal maximal $\mathfrak m' \subset S$,
l'anneau local $S_{\mathfrak m'}$ est analytiquement non ramifié,
\item pour tout homomorphisme local fini $(R, \mathfrak m) \to (S, \mathfrak m')$
avec $S$ intègre, l'anneau $S$ est analytiquement
non ramifié.
\end{enumerate}
\end{lemma}

\begin{proof}
Supposons que $R$ soit un anneau de Nagata et soient $R \to S$ et $\mathfrak m' \subset S$
comme en (2). Alors $S$ est un anneau de Nagata d'après le lemme \ref{lemma-quasi-finite-over-nagata}.
Par conséquent, l'anneau local $S_{\mathfrak m'}$ est un anneau de Nagata
(lemme \ref{lemma-nagata-localize}). Il est donc analytiquement
non ramifié d'après le lemme \ref{lemma-local-nagata-domain-analytically-unramified}.
Il est clair que (2) implique (3).

\medskip\noindent
Supposons (3) satisfaite. Soit $\mathfrak p \subset R$ un idéal premier et
soit $L/\kappa(\mathfrak p)$ une extension finie de corps.
Pour démontrer (1), nous devons montrer que la clôture intégrale de $R/\mathfrak p$
est finie sur $R/\mathfrak p$. Choisissons $x_1, \ldots, x_n \in L$
engendrant $L$ sur $\kappa(\mathfrak p)$. Pour chaque $i$, soit
$P_i(T) = T^{d_i} + a_{i, 1} T^{d_i - 1} + \ldots + a_{i, d_i}$
le polynôme minimal de $x_i$ sur $\kappa(\mathfrak p)$.
Après avoir remplacé $x_i$ par $f_i x_i$ pour un
$f_i \in R$ convenable, $f_i \not \in \mathfrak p$, nous pouvons supposer
$a_{i, j} \in R/\mathfrak p$. En fait, après avoir encore multiplié
par des éléments de $\mathfrak m$, nous pouvons supposer
$a_{i, j} \in \mathfrak m/\mathfrak p \subset R/\mathfrak p$ pour tous $i, j$.
Cela étant fait, posons $S = R/\mathfrak p[x_1, \ldots, x_n] \subset L$.
Alors $S$ est fini sur $R$, intègre, et $S/\mathfrak m S$ est un quotient
de $R/\mathfrak m[T_1, \ldots, T_n]/(T_1^{d_1}, \ldots, T_n^{d_n})$.
Ainsi, $S$ est local. D'après (3), $S$ est analytiquement non ramifié et, d'après le
lemme \ref{lemma-analytically-unramified-easy},
sa clôture intégrale $S'$ dans $L$ est finie sur $S$.
Puisque $S'$ est également la clôture intégrale de $R/\mathfrak p$ dans
$L$, le résultat est démontré.
\end{proof}

\noindent
La proposition suivante affirme notamment qu'une algèbre de type
fini sur un anneau de Nagata est un anneau de Nagata.

\begin{proposition}[Nagata]
\label{proposition-nagata-universally-japanese}
Soit $R$ un anneau. Les conditions suivantes sont équivalentes :
\begin{enumerate}
\item $R$ est un anneau de Nagata,
\item toute $R$-algèbre de type fini est un anneau de Nagata, et
\item $R$ est universellement japonais et noethérien.
\end{enumerate}
\end{proposition}

\begin{proof}
Il est clair que toute algèbre de type fini sur un anneau noethérien universellement
japonais est un anneau de Nagata (c'est-à-dire que la condition (2) est satisfaite). Soit $R$ un anneau de Nagata.
Nous allons montrer que toute $R$-algèbre de type fini $S$ est un anneau de Nagata.
Cela démontrera la proposition.

\medskip\noindent
Étape 1. Il existe une suite de morphismes d'anneaux
$R = R_0 \to R_1 \to R_2 \to \ldots \to R_n = S$ telle que
chaque $R_i \to R_{i + 1}$ soit engendré par un seul élément.
Par récurrence, il suffit donc de démontrer que $S$ est un anneau de Nagata si
$S \cong R[x]/I$.

\medskip\noindent
Étape 2. Soit $\mathfrak q \subset S$ un idéal premier de $S$, et soit
$\mathfrak p \subset R$ l'idéal premier correspondant de $R$.
Nous devons montrer que $S/\mathfrak q$ est N-2. Nous nous sommes donc
ramenés à démontrer l'assertion suivante :
(*) Étant donnés un anneau intègre de Nagata $R$ et une extension monogène $R \subset S$
d'anneaux intègres, alors $S$ est N-2.

\medskip\noindent
Étape 3. Soit $R$ un anneau intègre de Nagata et
soit $R \subset S$ une extension monogène d'anneaux intègres.
Supposons que l'extension induite des corps des fractions de $R$ et de $S$
soit transcendante pure. Dans ce cas, $S = R[x]$. D'après le
lemme \ref{lemma-polynomial-ring-N-2}, on voit que $S$ est N-2.
Nous nous sommes donc ramenés à démontrer l'assertion suivante :
(**) Étant donnés un anneau intègre de Nagata $R$ et une
extension monogène $R \subset S$ d'anneaux intègres
induisant une extension finie des corps des fractions,
alors $S$ est N-2.

\medskip\noindent
Étape 4. Soit $R$ un anneau intègre normal de Nagata et
soit $R \subset S$ une extension monogène d'anneaux intègres
induisant une extension finie $L/K$ des corps des fractions.
Choisissons un élément $x \in S$
qui engendre $S$ comme $R$-algèbre. Soit $M/L$
une extension finie de corps.
Soit $R'$ la clôture intégrale de $R$ dans $M$.
Alors la clôture intégrale $S'$ de $S$ dans $M$ coïncide avec la clôture
intégrale de $R'[x]$ dans $M$.
De plus, le corps des fractions de $R'$ est $M$, et $R \subset R'$
est fini (par la propriété de Nagata de $R$).
Cela implique que $R'$ est un anneau de Nagata
(lemme \ref{lemma-quasi-finite-over-nagata}).
Montrer que $S'$ est fini sur $S$ revient à montrer que
$S'$ est fini sur $R'[x]$. Remplaçons $R$ par $R'$ et $S$ par $R'[x]$
pour nous ramener à l'assertion suivante :
(***) Étant donnés un anneau intègre normal de Nagata $R$, de corps des fractions $K$,
et $x \in K$, l'anneau $S \subset K$ engendré par $R$ et $x$
est N-1.

\medskip\noindent
Étape 5. Soit $R$ un anneau intègre normal de Nagata de corps des fractions $K$.
Soit $x = b/a \in K$. Nous devons montrer que l'anneau $S \subset K$
engendré par $R$ et $x$ est N-1. Remarquons que $S_a \cong R_a$ est normal.
Ainsi, d'après le lemme \ref{lemma-characterize-N-1}, il suffit de montrer que
$S_{\mathfrak m}$ est N-1 pour tout idéal maximal $\mathfrak m$ de $S$.

\medskip\noindent
Sous les hypothèses du paragraphe précédent, choisissons un tel idéal maximal
et posons $\mathfrak n = R \cap \mathfrak m$. L'extension de corps résiduels
$\kappa(\mathfrak m)/\kappa(\mathfrak n)$ est finie
(théorème \ref{theorem-nullstellensatz}) et engendrée par l'image de $x$.
Après avoir remplacé $R$ par $R_a$ pour un certain $a \in R$, $a \not \in \mathfrak n$,
et $S$ par $S_a$, nous pouvons supposer qu'il existe un polynôme unitaire
$f(X) = X^d + \sum_{i = 1, \ldots, d} a_iX^{d - i}$ dans $R[x]$ tel que
$f(x) \in \mathfrak m$ (détails omis).
Soit $K''/K$ une extension finie de corps dans laquelle le polynôme
$f(X)$ se scinde complètement dans $K''[X]$.
Soit $R'$ la clôture intégrale de $R$ dans $K''$.
Soit $S' \subset K''$ le sous-anneau engendré par $R'$ et $x$.
Comme $R$ est un anneau de Nagata, $R'$ est fini sur $R$ et est de Nagata
(lemme \ref{lemma-quasi-finite-over-nagata}).
De plus, $S'$ est fini sur $S$. Si, pour tout idéal maximal
$\mathfrak m'$ de $S'$, l'anneau local $S'_{\mathfrak m'}$ est
N-1, alors $S'_{\mathfrak m}$ est N-1 d'après le
lemme \ref{lemma-characterize-N-1}, ce qui implique à son tour
que $S_{\mathfrak m}$ est N-1 d'après le
lemme \ref{lemma-finite-extension-N-2}.
Après avoir remplacé $R$ par $R'$ et $S$ par $S'$, et $\mathfrak m$ par
l'un quelconque des idéaux maximaux $\mathfrak m'$ au-dessus de $\mathfrak m$,
nous arrivons à la situation où le polynôme $f$ ci-dessus est complètement scindé :
$f(X) = \prod_{i = 1, \ldots, d} (X - a_i)$ avec $a_i \in R$.
Puisque $f(x) \in \mathfrak m$, on a $x - a_i \in \mathfrak m$
pour un certain $i$. Enfin, après avoir remplacé $x$ par $x - a_i$, nous pouvons supposer
que $x \in \mathfrak m$.

\medskip\noindent
Récapitulons : $R$ est un anneau intègre normal de Nagata de corps des fractions $K$,
$x \in K$ et $S$ est le sous-anneau de $K$ engendré par $x$ et $R$ ;
enfin, $\mathfrak m \subset S$ est un idéal maximal tel que $x \in \mathfrak m$.
Nous devons montrer que $S_{\mathfrak m}$ est N-1.

\medskip\noindent
Nous allons montrer que le lemme \ref{lemma-criterion-analytically-unramified}
s'applique à l'anneau local
$S_{\mathfrak m}$ et à l'élément $x$. Il en résultera que
$S_{\mathfrak m}$ est analytiquement non ramifié, et nous
verrons alors qu'il est N-1 d'après le lemme \ref{lemma-analytically-unramified-easy}.

\medskip\noindent
Nous devons vérifier les conditions (1), (2), (3)(a) et (3)(b).
La condition (1) est triviale.
Soit $I = \Ker(R[X] \to S)$, où $X \mapsto x$.
Nous affirmons que $I$ est engendré par toutes les formes linéaires $aX - b$ telles que
$ax = b$ dans $K$. Il est clair que toutes ces formes linéaires appartiennent à $I$.
Si $g = a_d X^d + \ldots a_1 X + a_0 \in I$, alors on voit que
$a_dx$ est entier sur $R$ (lemme \ref{lemma-make-integral-trivial}),
et donc que $b := a_dx \in R$,
puisque $R$ est normal. Alors $g - (a_dX - b)X^{d - 1} \in I$, et le résultat suit par
récurrence sur le degré. Par conséquent,
$$
S/xS = R[X]/(X, I) = R/J
$$
où
$$
J = \{b \in R \mid ax = b \text{ pour un certain }a \in R\} = xR \cap R
$$
D'après le lemme \ref{lemma-normal-domain-intersection-localizations-height-1},
on voit que $S/xS = R/J$ n'a aucun idéal premier associé immergé comme $R$-module, donc comme
$R/J$-module, donc comme $S/xS$-module, donc comme $S$-module.
Cela démontre la condition (2).
Prenons un tel idéal premier associé $\mathfrak q \subset S$ vérifiant
$\mathfrak q \subset \mathfrak m$ (de sorte que ce soit un
idéal premier associé à $S_{\mathfrak m}/xS_{\mathfrak m}$ -- cela n'a aucune
incidence sur les arguments).
Alors $\mathfrak q$ est minimal au-dessus de $xS$ et est donc de hauteur $1$.
La suite d'égalités ci-dessus montre que
$\mathfrak p = R \cap \mathfrak q$ est un idéal premier associé
à $R/J$, et est donc de hauteur $1$
(voir le lemme \ref{lemma-normal-domain-intersection-localizations-height-1}).
Ainsi, $R_{\mathfrak p}$ est un anneau de valuation discrète, et par conséquent
$R_{\mathfrak p} \subset S_{\mathfrak q}$ est une égalité. Cela montre
que $S_{\mathfrak q}$ est régulier. La condition (3)(a) est ainsi satisfaite.
Enfin, $(S/\mathfrak q)_{\mathfrak m}$ est une localisation
de $S/\mathfrak q$, qui est un quotient de $S/xS = R/J$.
Par conséquent, $(S/\mathfrak q)_{\mathfrak m}$ est une localisation d'un
quotient de l'anneau de Nagata $R$, et est donc
de Nagata (lemmes \ref{lemma-quasi-finite-over-nagata}
et \ref{lemma-nagata-localize}),
et donc analytiquement non ramifié
(lemme \ref{lemma-local-nagata-domain-analytically-unramified}).
Cela montre que (3)(b) est satisfaite et termine la démonstration.
\end{proof}

\begin{proposition}
\label{proposition-ubiquity-nagata}
Les types d'anneaux suivants sont de Nagata et, en particulier, universellement japonais :
\begin{enumerate}
\item les corps,
\item les anneaux locaux noethériens complets,
\item $\mathbf{Z}$,
\item les anneaux de Dedekind dont le corps des fractions est de caractéristique zéro,
\item les extensions d'anneaux de type fini de l'un quelconque des anneaux précédents.
\end{enumerate}
\end{proposition}

\begin{proof}
Le cas des anneaux locaux noethériens complets est traité dans le
lemme \ref{lemma-Noetherian-complete-local-Nagata}.
Dans les autres cas, il suffit de vérifier que $R/\mathfrak p$ est N-2 pour tout
idéal premier $\mathfrak p$ de l'anneau. Cela est clair dès que
$R/\mathfrak p$ est un corps, c'est-à-dire lorsque $\mathfrak p$ est maximal.
Dans le cas d'un anneau de Dedekind, il suffit donc d'effectuer la vérification lorsque
$\mathfrak p = (0)$. Mais, puisque nous supposons que le corps des fractions est de
caractéristique zéro, le lemme \ref{lemma-domain-char-zero-N-1-2} s'applique.
\end{proof}

\begin{example}
\label{example-nonjapanese-dvr}
Un anneau de valuation discrète est de Nagata si et seulement s'il est N-2 (car le
quotient par l'idéal maximal est un corps, donc est N-2). L'anneau de valuation discrète
$A$ de l'exemple \ref{example-bad-dvr-char-p} n'est pas de Nagata, c'est-à-dire
qu'il n'est pas N-2. En effet, l'extension finie
$A \subset R = A[f]$ n'est pas N-1. Pour le voir, écrivons $f = \sum a_i x^i$.
Pour tout $n \geq 1$, posons $g_n = \sum_{i < n} a_i x^i \in A$.
Alors $h_n = (f - g_n)/x^n$ est un élément du corps des fractions de $R$
et $h_n^p \in k^p[[x]] \subset A$. La clôture intégrale $R'$
de $R$ contient donc $h_1, h_2, h_3, \ldots$. Or, si $R'$
était fini sur $R$, et donc sur $A$, alors $f  = x^n h_n + g_n$
appartiendrait au sous-module $A + x^nR'$ pour tout $n$. Le lemme
d'Artin--Rees impliquerait alors $f \in A$
(lemme \ref{lemma-intersect-powers-ideal-module-zero}), ce qui est contradictoire.
\end{example}

\begin{lemma}
\label{lemma-nagata-pth-roots}
Soit $(A, \mathfrak m)$ un anneau intègre local noethérien qui est un anneau de Nagata
et dont le corps des fractions est de caractéristique $p$. Si $a \in A$ possède une
racine $p$-ième dans $A^\wedge$, alors $a$ possède une racine $p$-ième dans $A$.
\end{lemma}

\begin{proof}
Soit $\alpha \in A^\wedge$ une racine $p$-ième de $a$. Pour obtenir une contradiction,
supposons que $a$ ne possède pas de racine $p$-ième dans $A$. Alors
$a$ ne possède pas de racine $p$-ième dans le corps des fractions $K$ de $A$.
En effet, si $\beta = b/c$, avec $b, c \in A$ et $c \not = 0$,
vérifie $\beta^p = a$, alors $\alpha = b/c$ dans $A^\wedge$,
ce qui implique que $c$ divise $b$ dans $A$
par fidèle platitude de $A \to A^\wedge$,
et donc que $\beta \in A$, contradiction. Ainsi,
$B = A[x]/(x^p - a)$ est intègre, car $K[x]/(x^p - a)$ est un corps.
Cependant, le complété de $B$ n'est pas réduit en raison de l'existence supposée
de $\alpha$. Cela contredit les résultats précédents, car $B$ est un anneau de Nagata
(proposition \ref{proposition-nagata-universally-japanese})
et est donc analytiquement non ramifié d'après le
lemme \ref{lemma-local-nagata-domain-analytically-unramified}.
\end{proof}






\section{Montée des propriétés}
\label{section-ascending-properties}

\noindent
Dans cette section, nous commençons à démontrer quelques faits algébriques concernant la
« montée » de propriétés des anneaux. Pour traiter le cas de la profondeur des anneaux,
on utilise le résultat suivant sur la montée de la profondeur des modules, voir
\cite[IV, Proposition 6.3.1]{EGA}.

\begin{lemma}
\label{lemma-apply-grothendieck-module}
\begin{reference}
\cite[IV, Proposition 6.3.1]{EGA}
\end{reference}
Nous avons
$$
\text{depth}(M \otimes_R N)
=
\text{depth}(M) + \text{depth}(N/\mathfrak m_RN)
$$
où $R \to S$ est un morphisme local d'anneaux locaux noethériens,
$M$ est un $R$-module de type fini et $N$ est un $S$-module de type fini plat sur $R$.
\end{lemma}

\begin{proof}
Dans l'énoncé et dans la démonstration ci-dessous, la profondeur de $M$ est prise
en tant que $R$-module, celle de $M \otimes_R N$ en tant que $S$-module, et
celle de $N/\mathfrak m_RN$ en tant que $S/\mathfrak m_RS$-module.
Notons $n$ le membre de droite. Supposons d'abord que $n$ soit nul.
Alors nous avons à la fois $\text{depth}(M) = 0$ et
$\text{depth}(N/\mathfrak m_RN) = 0$.
Cela signifie qu'il existe $z \in M$ dont l'annulateur est $\mathfrak m_R$
et $\overline{y} \in N/\mathfrak m_RN$
dont l'annulateur est $\mathfrak m_S/\mathfrak m_RS$.
Soit $y \in N$ un relèvement de $\overline{y}$.
Puisque $N$ est plat sur $R$, l'application $z : R/\mathfrak m_R \to M$
induit une application injective $N/\mathfrak m_RN \to M \otimes_R N$.
L'annulateur de $z \otimes y$ est donc $\mathfrak m_S$.
Ainsi, $\text{depth}(M \otimes_R N) = 0$ également.

\medskip\noindent
Supposons $n > 0$. Si $\text{depth}(N/\mathfrak m_RN) > 0$, nous pouvons choisir
$f \in \mathfrak m_S$ dont l'image $\overline{f} \in S/\mathfrak m_RS$
est un non-diviseur de zéro dans $N/\mathfrak m_RN$.
Alors $\text{depth}(N/\mathfrak m_RN) =
\text{depth}(N/(f, \mathfrak m_R)N) + 1$
d'après le lemme \ref{lemma-depth-drops-by-one}.
D'après le lemme \ref{lemma-mod-injective}, l'élément $f \in S$ est un
non-diviseur de zéro dans $N$ et $N/fN$ est plat sur $R$.
Par récurrence sur $n$, nous avons donc
$$
\text{depth}(M \otimes_R N/fN) =
\text{depth}(M) + \text{depth}(N/(f, \mathfrak m_R)N).
$$
Puisque $N/fN$ est plat sur $R$, la suite
$$
0 \to M \otimes_R N \to M \otimes_R N \to M \otimes_R N/fN \to 0
$$
est exacte, où la première application est la multiplication par $f$
(lemme \ref{lemma-flat-tor-zero}). Par conséquent, d'après le
lemme \ref{lemma-depth-drops-by-one}, nous obtenons
$\text{depth}(M \otimes_R N) = \text{depth}(M \otimes_R N/fN) + 1$
et nous obtenons ainsi l'égalité de la formule du lemme.

\medskip\noindent
Si $n > 0$, mais $\text{depth}(N/\mathfrak m_RN) = 0$,
nous pouvons choisir $f \in \mathfrak m_R$ qui soit un non-diviseur de zéro dans $M$.
Puisque $N$ est plat sur $R$, $f$ est également un non-diviseur de zéro dans
$M \otimes_R N$. Par une nouvelle récurrence sur $n$, nous avons
$$
\text{depth}(M/fM \otimes_R N) =
\text{depth}(M/fM) + \text{depth}(N/\mathfrak m_RN).
$$
Dans ce cas,
$\text{depth}(M \otimes_R N) = \text{depth}(M/fM \otimes_R N) + 1$
et $\text{depth}(M) = \text{depth}(M/fM) + 1$
d'après le lemme \ref{lemma-depth-drops-by-one}, et
nous obtenons ainsi l'égalité de la formule du lemme.
\end{proof}

\begin{lemma}
\label{lemma-apply-grothendieck}
Supposons que $R \to S$ soit un morphisme local plat d'anneaux locaux
noethériens. Alors
$$
\text{depth}(S) = \text{depth}(R) + \text{depth}(S/\mathfrak m_RS).
$$
\end{lemma}

\begin{proof}
C'est un cas particulier du lemme \ref{lemma-apply-grothendieck-module}.
\end{proof}

\begin{lemma}
\label{lemma-CM-goes-up}
Soit $R \to S$ un morphisme local plat d'anneaux locaux noethériens.
Alors les conditions suivantes sont équivalentes :
\begin{enumerate}
\item $S$ est de Cohen-Macaulay, et
\item $R$ et $S/\mathfrak m_RS$ sont de Cohen-Macaulay.
\end{enumerate}
\end{lemma}

\begin{proof}
Cela résulte des définitions et des
lemmes \ref{lemma-apply-grothendieck} et
\ref{lemma-dimension-base-fibre-equals-total}.
\end{proof}

\begin{lemma}
\label{lemma-Sk-goes-up}
Soit $\varphi : R \to S$ un morphisme d'anneaux. Supposons que
\begin{enumerate}
\item $R$ soit noethérien,
\item $S$ soit noethérien,
\item $\varphi$ soit plat,
\item les anneaux des fibres $S \otimes_R \kappa(\mathfrak p)$ vérifient $(S_k)$, et
\item $R$ vérifie la propriété $(S_k)$.
\end{enumerate}
Alors $S$ vérifie la propriété $(S_k)$.
\end{lemma}

\begin{proof}
Soit $\mathfrak q$ un idéal premier de $S$
au-dessus d'un idéal premier $\mathfrak p$ de $R$. D'après le
lemme \ref{lemma-apply-grothendieck}, nous avons
$$
\text{depth}(S_{\mathfrak q}) =
\text{depth}(S_{\mathfrak q}/\mathfrak pS_{\mathfrak q}) +
\text{depth}(R_{\mathfrak p}).
$$
D'autre part, nous avons
$$
\dim(R_{\mathfrak p})
+
\dim(S_{\mathfrak q}/\mathfrak pS_{\mathfrak q})
\geq
\dim(S_{\mathfrak q})
$$
d'après le lemme \ref{lemma-dimension-base-fibre-total}.
(En fait, il y a égalité, d'après le
lemme \ref{lemma-dimension-base-fibre-equals-total},
mais, à strictement parler, nous n'en avons pas besoin.)
Enfin, puisque les anneaux des fibres de ce morphisme
sont supposés vérifier $(S_k)$, nous voyons que
$\text{depth}(S_{\mathfrak q}/\mathfrak pS_{\mathfrak q})
\geq \min(k, \dim(S_{\mathfrak q}/\mathfrak pS_{\mathfrak q}))$.
Ainsi, le lemme résulte de la suite d'inégalités suivante :
\begin{eqnarray*}
\text{depth}(S_{\mathfrak q}) & = &
\text{depth}(S_{\mathfrak q}/\mathfrak pS_{\mathfrak q}) +
\text{depth}(R_{\mathfrak p}) \\
& \geq &
\min(k, \dim(S_{\mathfrak q}/\mathfrak pS_{\mathfrak q})) +
\min(k, \dim(R_{\mathfrak p})) \\
& = &
\min(2k, \dim(S_{\mathfrak q}/\mathfrak pS_{\mathfrak q}) + k,
k + \dim(R_\mathfrak p),
\dim(S_{\mathfrak q}/\mathfrak pS_{\mathfrak q}) +
\dim(R_{\mathfrak p})) \\
& \geq &
\min(k, \dim(S_{\mathfrak q}))
\end{eqnarray*}
ce qu'il fallait démontrer.
\end{proof}

\begin{lemma}
\label{lemma-Rk-goes-up}
Soit $\varphi : R \to S$ un morphisme d'anneaux. Supposons que
\begin{enumerate}
\item $R$ soit noethérien,
\item $S$ soit noethérien,
\item $\varphi$ soit plat,
\item les anneaux des fibres $S \otimes_R \kappa(\mathfrak p)$
vérifient la propriété $(R_k)$, et
\item $R$ vérifie la propriété $(R_k)$.
\end{enumerate}
Alors $S$ vérifie la propriété $(R_k)$.
\end{lemma}

\begin{proof}
Soit $\mathfrak q$ un idéal premier de $S$
au-dessus d'un idéal premier $\mathfrak p$ de $R$.
Supposons que $\dim(S_{\mathfrak q}) \leq k$.
Puisque $\dim(S_{\mathfrak q}) = \dim(R_{\mathfrak p})
+ \dim(S_{\mathfrak q}/\mathfrak pS_{\mathfrak q})$ d'après le
lemme \ref{lemma-dimension-base-fibre-equals-total},
nous voyons que $\dim(R_{\mathfrak p}) \leq k$ et que
$\dim(S_{\mathfrak q}/\mathfrak pS_{\mathfrak q}) \leq k$.
Ainsi, $R_{\mathfrak p}$ et $S_{\mathfrak q}/\mathfrak pS_{\mathfrak q}$
sont réguliers par hypothèse.
Il s'ensuit que $S_{\mathfrak q}$ est régulier d'après le
lemme \ref{lemma-flat-over-regular-with-regular-fibre}.
\end{proof}

\begin{lemma}
\label{lemma-reduced-goes-up-noetherian}
Soit $\varphi : R \to S$ un morphisme d'anneaux. Supposons que
\begin{enumerate}
\item $R$ soit noethérien,
\item $S$ soit noethérien,
\item $\varphi$ soit plat,
\item les anneaux des fibres $S \otimes_R \kappa(\mathfrak p)$ soient réduits,
\item $R$ soit réduit.
\end{enumerate}
Alors $S$ est réduit.
\end{lemma}

\begin{proof}
Pour un anneau noethérien, être réduit équivaut à vérifier les propriétés
$(S_1)$ et $(R_0)$, voir le lemme \ref{lemma-criterion-reduced}.
Nous savons donc que $R$ et les anneaux des fibres vérifient ces propriétés.
Nous pouvons ainsi appliquer les lemmes \ref{lemma-Sk-goes-up} et \ref{lemma-Rk-goes-up}
et nous voyons que $S$ vérifie $(S_1)$ et $(R_0)$, c'est-à-dire qu'il est réduit,
à nouveau d'après le lemme \ref{lemma-criterion-reduced}.
\end{proof}

\begin{lemma}
\label{lemma-reduced-goes-up}
Soit $\varphi : R \to S$ un morphisme d'anneaux. Supposons que
\begin{enumerate}
\item $\varphi$ soit lisse,
\item $R$ soit réduit.
\end{enumerate}
Alors $S$ est réduit.
\end{lemma}

\begin{proof}
Remarquons que $R \to S$ est plat et à fibres régulières (voir la liste des
résultats sur les morphismes lisses d'anneaux à la section \ref{section-smooth-overview}).
En particulier, les fibres sont réduites.
Ainsi, si $R$ est noethérien, alors $S$ est noethérien et le résultat
découle du lemme \ref{lemma-reduced-goes-up-noetherian}.

\medskip\noindent
Dans le cas général, nous pouvons trouver une sous-$\mathbf{Z}$-algèbre
$R_0 \subset R$ de type fini et un morphisme lisse d'anneaux
$R_0 \to S_0$ tels que $S \cong R \otimes_{R_0} S_0$, voir la
remarque (10) de la section \ref{section-smooth-overview}.
Si maintenant $x \in S$ est un élément tel que $x^2 = 0$,
nous pouvons agrandir $R_0$ et supposer que $x$ provient
d'un élément $x_0 \in S_0$. Après avoir agrandi
$R_0$ une fois encore, nous pouvons supposer que $x_0^2 = 0$ dans $S_0$.
Cependant, puisque la sous-algèbre $R_0 \subset R$ est réduite, nous voyons que
$S_0$ est réduit et donc que $x_0 = 0$, comme voulu.
\end{proof}

\begin{lemma}
\label{lemma-normal-goes-up-noetherian}
Soit $\varphi : R \to S$ un morphisme d'anneaux. Supposons que
\begin{enumerate}
\item $R$ soit noethérien,
\item $S$ soit noethérien,
\item $\varphi$ soit plat,
\item les anneaux des fibres $S \otimes_R \kappa(\mathfrak p)$ soient normaux, et
\item $R$ soit normal.
\end{enumerate}
Alors $S$ est normal.
\end{lemma}

\begin{proof}
Pour un anneau noethérien, être normal équivaut à vérifier les propriétés
$(S_2)$ et $(R_1)$, voir le lemme \ref{lemma-criterion-normal}.
Nous savons donc que $R$ et les anneaux des fibres vérifient ces propriétés.
Nous pouvons ainsi appliquer les lemmes \ref{lemma-Sk-goes-up} et \ref{lemma-Rk-goes-up}
et nous voyons que $S$ vérifie $(S_2)$ et $(R_1)$, c'est-à-dire qu'il est normal,
à nouveau d'après le lemme \ref{lemma-criterion-normal}.
\end{proof}

\begin{lemma}
\label{lemma-normal-goes-up}
Soit $\varphi : R \to S$ un morphisme d'anneaux. Supposons que
\begin{enumerate}
\item $\varphi$ soit lisse,
\item $R$ soit normal.
\end{enumerate}
Alors $S$ est normal.
\end{lemma}

\begin{proof}
Remarquons que $R \to S$ est plat et à fibres régulières (voir la liste des
résultats sur les morphismes lisses d'anneaux à la section \ref{section-smooth-overview}).
En particulier, les fibres sont normales. Ainsi, si $R$ est noethérien,
alors $S$ est noethérien et le résultat découle du
lemme \ref{lemma-normal-goes-up-noetherian}.

\medskip\noindent
Passons au cas général. Remarquons d'abord que $R$ est réduit et donc que
$S$ est réduit d'après le lemme \ref{lemma-reduced-goes-up}.
Soit $\mathfrak q$ un idéal premier de $S$ et soit $\mathfrak p$
l'idéal premier correspondant dans $R$. Remarquons que $R_{\mathfrak p}$
est un anneau intègre normal. Nous devons montrer que $S_{\mathfrak q}$ est
un anneau intègre normal. Pour ce faire, nous pouvons remplacer $R$ par $R_{\mathfrak p}$
et $S$ par $S_{\mathfrak p}$. Nous pouvons donc supposer que $R$ est
un anneau intègre normal.

\medskip\noindent
Supposons le morphisme $R \to S$ lisse et $R$ un anneau intègre normal.
Nous pouvons trouver une sous-$\mathbf{Z}$-algèbre de type fini
$R_0 \subset R$ et un morphisme lisse d'anneaux $R_0 \to S_0$ tels
que $S \cong R \otimes_{R_0} S_0$, voir la
remarque (10) de la section \ref{section-smooth-overview}.
Puisque $R_0$ est un anneau intègre de Nagata (voir la proposition \ref{proposition-ubiquity-nagata}),
nous voyons que sa clôture intégrale $R_0'$ est finie sur $R_0$.
De plus, comme $R$ est un anneau intègre normal, il est clair que $R_0' \subset R$.
Nous pouvons donc remplacer $R_0$ par $R_0'$ et $S_0$ par
$R_0' \otimes_{R_0} S_0$ et supposer que $R_0$ est un anneau intègre
normal noethérien. Le premier paragraphe de la démonstration montre alors
que $S_0$ est un anneau normal (il n'est bien sûr pas nécessairement intègre).
De cette manière, nous voyons que $R = \bigcup R_\lambda$
est la réunion d'anneaux intègres normaux noethériens et, parallèlement,
$S = \colim R_\lambda \otimes_{R_0} S_0$ est la limite inductive
d'anneaux normaux. Cela implique que $S$ est un anneau normal.
Nous omettons quelques détails.
\end{proof}

\begin{lemma}
\label{lemma-regular-goes-up}
\begin{slogan}
La régularité se propage le long des morphismes lisses d'anneaux.
\end{slogan}
Soit $\varphi : R \to S$ un morphisme d'anneaux. Supposons que
\begin{enumerate}
\item $\varphi$ soit lisse,
\item $R$ soit un anneau régulier.
\end{enumerate}
Alors $S$ est régulier.
\end{lemma}

\begin{proof}
Cela résulte du lemme \ref{lemma-Rk-goes-up} appliqué à toutes les propriétés $(R_k)$,
le lemme \ref{lemma-characterize-smooth-over-field} permettant de voir que les
hypothèses sont satisfaites.
\end{proof}






\section{Descente des propriétés}
\label{section-descending-properties}

\noindent
Dans cette section, nous commençons par établir quelques résultats algébriques concernant la
« descente » de propriétés des anneaux. Il s'avère qu'il est souvent
« plus facile » de faire descendre des propriétés que de les faire monter. Autrement
dit, les hypothèses portant sur le morphisme d'anneaux $R \to S$ sont souvent plus faibles que
celles du lemme correspondant de la section précédente.
Nous avertissons toutefois le lecteur que les résultats de descente sont souvent
inutiles si l'on ne peut pas également établir les résultats de montée correspondants !
Voici un résultat typique illustrant ce phénomène.

\begin{lemma}
\label{lemma-descent-Noetherian}
Soit $R \to S$ un morphisme d'anneaux.
Supposons que
\begin{enumerate}
\item $R \to S$ soit fidèlement plat, et
\item $S$ soit noethérien.
\end{enumerate}
Alors $R$ est noethérien.
\end{lemma}

\begin{proof}
Soit $I_0 \subset I_1 \subset I_2 \subset \ldots$ une
suite croissante d'idéaux de $R$. Par hypothèse, nous avons
$I_nS = I_{n+1}S = I_{n+2}S = \ldots$ pour un certain $n$.
Par fidèle platitude, l'extension suivie de la contraction laisse tout idéal inchangé,
c'est-à-dire que $I = R \cap IS$ pour tout idéal $I$ de $R$ (lemme
\ref{lemma-faithfully-flat-universally-injective}). Ainsi,
$I_n = I_{n+1} = I_{n+2} = \ldots$, ce qu'il fallait démontrer.
\end{proof}

\begin{lemma}
\label{lemma-descent-reduced}
Soit $R \to S$ un morphisme d'anneaux.
Supposons que
\begin{enumerate}
\item $R \to S$ soit fidèlement plat, et
\item $S$ soit réduit.
\end{enumerate}
Alors $R$ est réduit.
\end{lemma}

\begin{proof}
Cela est clair puisque $R \to S$ est injectif, d'après le
lemme \ref{lemma-faithfully-flat-universally-injective}.
\end{proof}

\begin{lemma}
\label{lemma-descent-normal}
Soit $R \to S$ un morphisme d'anneaux.
Supposons que
\begin{enumerate}
\item $R \to S$ soit fidèlement plat, et
\item $S$ soit un anneau normal.
\end{enumerate}
Alors $R$ est un anneau normal.
\end{lemma}

\begin{proof}
Puisque $S$ est réduit, il s'ensuit que $R$ est réduit.
Soit $\mathfrak p$ un idéal premier de $R$. Nous devons montrer que
$R_{\mathfrak p}$ est un anneau intègre normal. Comme $S_{\mathfrak p}$
est lui aussi fidèlement plat sur $R_{\mathfrak p}$, nous pouvons supposer que
$R$ est local, d'idéal maximal $\mathfrak m$.
Soit $\mathfrak q$ un idéal premier de $S$ au-dessus de $\mathfrak m$.
Nous voyons alors que $R \to S_{\mathfrak q}$ est fidèlement plat
(lemme \ref{lemma-local-flat-ff}).
Nous pouvons donc supposer que $S$ est lui aussi local.
En particulier, $S$ est un anneau intègre normal.
Comme $R \to S$ est fidèlement plat
et que $S$ est un anneau intègre normal, nous voyons que $R$ est un anneau intègre.
Ensuite, supposons que $a/b$ soit entier sur $R$, où $a, b \in R$.
Alors $a/b \in S$ puisque $S$ est normal. Ainsi $a \in bS$.
Cela signifie que l'application $a : R \to R/bR$ devient nulle
après changement de base à $S$. Par fidèle platitude, nous voyons que
$a \in bR$, donc $a/b \in R$. Ainsi $R$ est normal.
\end{proof}

\begin{lemma}
\label{lemma-descent-regular}
Soit $R \to S$ un morphisme d'anneaux.
Supposons que
\begin{enumerate}
\item $R \to S$ soit fidèlement plat, et
\item $S$ soit un anneau régulier.
\end{enumerate}
Alors $R$ est un anneau régulier.
\end{lemma}

\begin{proof}
Nous voyons que $R$ est noethérien d'après le lemme \ref{lemma-descent-Noetherian}.
Soit $\mathfrak p \subset R$ un idéal premier. Choisissons un idéal premier $\mathfrak q \subset S$
au-dessus de $\mathfrak p$. Le lemme \ref{lemma-flat-under-regular}
s'applique à $R_\mathfrak p \to S_\mathfrak q$ et nous en concluons que
$R_\mathfrak p$ est régulier. Comme $\mathfrak p$ était arbitraire, nous voyons que
$R$ est régulier.
\end{proof}

\begin{lemma}
\label{lemma-descent-Sk}
Soit $R \to S$ un morphisme d'anneaux.
Supposons que
\begin{enumerate}
\item $R \to S$ soit fidèlement plat, et
\item $S$ soit noethérien et vérifie la propriété $(S_k)$.
\end{enumerate}
Alors $R$ est noethérien et vérifie la propriété $(S_k)$.
\end{lemma}

\begin{proof}
Nous avons déjà vu que (1) et (2) impliquent que $R$ est noethérien,
voir le lemme \ref{lemma-descent-Noetherian}.
Soit $\mathfrak p \subset R$ un idéal premier.
Choisissons un idéal premier $\mathfrak q \subset S$ au-dessus de $\mathfrak p$
correspondant à un idéal premier minimal de l'anneau de la fibre
$S \otimes_R \kappa(\mathfrak p)$. Alors
$A = R_{\mathfrak p} \to S_{\mathfrak q} = B$ est un homomorphisme local et plat
d'anneaux locaux noethériens, et $\mathfrak m_AB$ est un
idéal de définition de $B$. Ainsi,
$\dim(A) = \dim(B)$ (lemme \ref{lemma-dimension-base-fibre-equals-total}) et
$\text{depth}(A) = \text{depth}(B)$ (lemme \ref{lemma-apply-grothendieck}).
Par conséquent, puisque $B$ vérifie la propriété $(S_k)$, nous
voyons que $A$ vérifie la propriété $(S_k)$.
\end{proof}

\begin{lemma}
\label{lemma-descent-Rk}
Soit $R \to S$ un morphisme d'anneaux. Supposons que
\begin{enumerate}
\item $R \to S$ soit fidèlement plat, et
\item $S$ soit noethérien et vérifie la propriété $(R_k)$.
\end{enumerate}
Alors $R$ est noethérien et vérifie la propriété $(R_k)$.
\end{lemma}

\begin{proof}
Nous avons déjà vu que (1) et (2) impliquent que $R$ est noethérien,
voir le lemme \ref{lemma-descent-Noetherian}.
Soit $\mathfrak p \subset R$ un idéal premier et supposons que
$\dim(R_{\mathfrak p}) \leq k$.
Choisissons un idéal premier $\mathfrak q \subset S$ au-dessus de $\mathfrak p$
correspondant à un idéal premier minimal de l'anneau de la fibre
$S \otimes_R \kappa(\mathfrak p)$. Alors
$A = R_{\mathfrak p} \to S_{\mathfrak q} = B$ est un homomorphisme local et plat
d'anneaux locaux noethériens, et $\mathfrak m_AB$ est un
idéal de définition de $B$. Ainsi,
$\dim(A) = \dim(B)$ (lemme \ref{lemma-dimension-base-fibre-equals-total}).
Comme $S$ vérifie la propriété $(R_k)$, nous en concluons que $B$ est un anneau local régulier.
D'après le lemme \ref{lemma-flat-under-regular}, nous en concluons que $A$ est régulier.
\end{proof}

\begin{lemma}
\label{lemma-descent-nagata}
Soit $R \to S$ un morphisme d'anneaux. Supposons que
\begin{enumerate}
\item $R \to S$ soit lisse et induise une surjection sur les spectres, et
\item $S$ soit un anneau de Nagata.
\end{enumerate}
Alors $R$ est un anneau de Nagata.
\end{lemma}

\begin{proof}
Rappelons qu'un anneau de Nagata est la même chose qu'un anneau noethérien
universellement japonais
(proposition \ref{proposition-nagata-universally-japanese}).
Nous avons déjà vu que $R$ est noethérien dans le
lemme \ref{lemma-descent-Noetherian}.
Soit $R \to A$ un morphisme d'anneaux de type fini, l'anneau cible étant intègre.
D'après le lemme \ref{lemma-check-universally-japanese},
il suffit de vérifier que $A$ est N-1.
Il est clair que $B = A \otimes_R S$ est une $S$-algèbre de type fini
et donc un anneau de Nagata (proposition \ref{proposition-nagata-universally-japanese}).
Puisque $A \to B$ est lisse (lemme \ref{lemma-base-change-smooth}),
nous voyons que $B$ est réduit (lemme \ref{lemma-reduced-goes-up}).
Comme $B$ est noethérien, il ne possède qu'un nombre fini d'idéaux
premiers minimaux $\mathfrak q_1, \ldots, \mathfrak q_t$ (voir le
lemme \ref{lemma-Noetherian-irreducible-components}).
Comme $A \to B$ est plat, chacun de ces idéaux se trouve au-dessus de $(0) \subset A$
(par le théorème de descente, voir le lemme \ref{lemma-flat-going-down}).
L'anneau total des fractions $Q(B)$ est le produit des
$L_i = \kappa(\mathfrak q_i)$ (lemmes
\ref{lemma-total-ring-fractions-no-embedded-points} et
\ref{lemma-minimal-prime-reduced-ring}).
En outre, la clôture intégrale $B'$ de $B$ dans $Q(B)$ est
le produit des clôtures intégrales $B_i'$ des anneaux $B/\mathfrak q_i$
dans les facteurs $L_i$ (comparer avec le
lemme \ref{lemma-characterize-reduced-ring-normal}).
Puisque $B$ est universellement japonais, les
extensions d'anneaux $B/\mathfrak q_i \subset B_i'$ sont finies
et nous en concluons que $B' = \prod B_i'$ est un $B$-module de type fini.
Comme $A \to B$ est plat, nous voyons que tout
non-diviseur de zéro dans $A$ a pour image un non-diviseur de zéro dans $B$.
L'application correspondante
$$
Q(A) \otimes_A B = (A \setminus \{0\})^{-1}A \otimes_A B
= (A \setminus \{0\})^{-1}B \to Q(B)
$$
est injective (nous avons utilisé le lemme \ref{lemma-tensor-localization}).
$A'$ désigne ici la clôture intégrale de l'anneau de départ dans son corps des fractions, et cette application l'envoie dans $B'$. Elle induit une application
$$
A' \otimes_A B \longrightarrow B'
$$
qui est injective (d'après ce qui précède et la platitude de $A \to B$).
Puisque $B'$ est un $B$-module de type fini
et que $B$ est noethérien, nous voyons que $A' \otimes_A B$ est un $B$-module de type fini.
Il existe donc un nombre fini d'éléments $x_i \in A'$ tels que
les éléments $x_i \otimes 1$ engendrent $A' \otimes_A B$ comme $B$-module.
Enfin, par fidèle platitude de $A \to B$, nous concluons que
les $x_i$ engendrent aussi $A'$ comme $A$-module, ce qui achève la démonstration.
\end{proof}

\begin{remark}
\label{remark-universally-catenary-does-not-descend}
La propriété d'être « universellement caténaire » ne descend pas, même
le long de morphismes \'etales d'anneaux. Dans le
chapitre des exemples, à la section \ref{examples-section-non-catenary-Noetherian-local},
on construit un morphisme fini d'anneaux $A \to B$ tel que
$A$ soit un anneau local noethérien non universellement caténaire
et que $B$ soit un anneau semi-local dont les deux idéaux maximaux sont $\mathfrak m$ et $\mathfrak n$.
Les anneaux $B_{\mathfrak m}$ et $B_{\mathfrak n}$ sont réguliers, de dimensions $2$ et $1$
respectivement, et ont chacun le même corps résiduel que celui de $A$.
De plus, $\mathfrak m_A$ engendre l'idéal maximal de chacun des anneaux
$B_{\mathfrak m}$ et $B_{\mathfrak n}$ (ainsi $A \to B$ est à la fois non ramifié
et fini).
D'après le lemme \ref{lemma-etale-makes-unramified-closed},
il existe un morphisme local et \'etale d'anneaux
$A \to A'$ pour lequel $B \otimes_A A' = B_1 \times B_2$,
les morphismes $A' \to B_i$ étant surjectifs.
Cela montre que $A'$ possède deux idéaux premiers minimaux $\mathfrak q_i$
tels que $A'/\mathfrak q_i \cong B_i$. Puisque $B_i$ est un anneau local régulier
(puisqu'il est \'etale sur $B_{\mathfrak m}$ ou sur $B_{\mathfrak n}$),
nous concluons que $A'$ est universellement caténaire.
\end{remark}









\section{Algèbres géométriquement normales}
\label{section-geometrically-normal}

\noindent
Dans cette section, nous donnons quelques applications de la montée et de la
descente des propriétés des anneaux.

\begin{lemma}
\label{lemma-geometrically-normal}
Soit $k$ un corps. Soit $A$ une $k$-algèbre.
Les propriétés suivantes de $A$ sont équivalentes :
\begin{enumerate}
\item $k' \otimes_k A$ est un anneau normal
pour toute extension de corps $k'/k$,
\item $k' \otimes_k A$ est un anneau normal
pour toute extension de corps de type fini $k'/k$,
\item $k' \otimes_k A$ est un anneau normal
pour toute extension de corps finie purement inséparable $k'/k$,
\item $k^{perf} \otimes_k A$ est un anneau normal.
\end{enumerate}
Ici, l'expression « anneau normal » est prise au sens de la définition \ref{definition-ring-normal}.
\end{lemma}

\begin{proof}
Il est clair que (1) $\Rightarrow$ (2) $\Rightarrow$ (3)
et que (1) $\Rightarrow$ (4).

\medskip\noindent
Si $k'/k$ est une extension de corps finie purement inséparable, alors
il existe un plongement $k' \to k^{perf}$ au-dessus de $k$.
Le morphisme d'anneaux $k' \otimes_k A \to k^{perf} \otimes_k A$
est fidèlement plat ; ainsi, $k' \otimes_k A$ est normal si
$k^{perf} \otimes_k A$ est normal, d'après le
lemme \ref{lemma-descent-normal}. On voit ainsi que
(4) $\Rightarrow$ (3).

\medskip\noindent
Supposons (2) et soit $k'/k$ une extension de corps quelconque.
Nous pouvons alors écrire $k' = \colim_i k_i$ sous la forme d'une limite
inductive filtrante d'extensions de corps de type fini. Par conséquent,
$k' \otimes_k A = \colim_i k_i \otimes_k A$
est une limite inductive filtrante d'anneaux normaux. Il s'ensuit
que $k' \otimes_k A$ est un anneau normal d'après le
lemme \ref{lemma-colimit-normal-ring}.
Ainsi, (1) est vérifiée.

\medskip\noindent
Supposons (3) et soit $K/k$ une extension de corps de type fini.
D'après le lemme \ref{lemma-make-separable}, il existe un diagramme
$$
\xymatrix{
K \ar[r] & K' \\
k \ar[u] \ar[r] & k' \ar[u]
}
$$
où $k'/k$ et $K'/K$ sont des extensions de corps finies purement inséparables
telles que $K'/k'$ soit séparable. D'après le
lemme \ref{lemma-localization-smooth-separable},
il existe une $k'$-algèbre lisse $B$ telle que $K'$ soit le
corps des fractions de $B$. Nous pouvons alors raisonner comme suit :
Étape 1 : $k' \otimes_k A$ est un anneau normal, puisque nous avons supposé (3).
Étape 2 : $B \otimes_{k'} k' \otimes_k A$ est un anneau normal, car
$k' \otimes_k A \to B \otimes_{k'} k' \otimes_k A$ est lisse
(lemme \ref{lemma-base-change-smooth})
et la normalité monte le long des morphismes lisses
(lemme \ref{lemma-normal-goes-up}).
Étape 3 : $K' \otimes_{k'} k' \otimes_k A = K' \otimes_k A$ est
un anneau normal, car c'est une localisation d'un anneau normal
(lemme \ref{lemma-localization-normal-ring}).
Étape 4 : Enfin, $K \otimes_k A$ est un anneau normal par descente de la
normalité le long du morphisme d'anneaux fidèlement plat
$K \otimes_k A \to K' \otimes_k A$ (lemme \ref{lemma-descent-normal}).
Ceci démontre le lemme.
\end{proof}

\begin{definition}
\label{definition-geometrically-normal}
Soit $k$ un corps.
Une $k$-algèbre $R$ est dite {\it géométriquement normale} sur $k$ si
les conditions équivalentes du lemme \ref{lemma-geometrically-normal} sont satisfaites.
\end{definition}

\begin{lemma}
\label{lemma-localization-geometrically-normal-algebra}
\begin{slogan}
La localisation préserve la normalité géométrique.
\end{slogan}
Soit $k$ un corps. Une localisation d'une $k$-algèbre géométriquement normale
est géométriquement normale.
\end{lemma}

\begin{proof}
Cela est clair, puisque la normalité d'un anneau se vérifie après localisation en tout
idéal premier.
\end{proof}

\begin{lemma}
\label{lemma-separable-field-extension-geometrically-normal}
Soit $k$ un corps. Soit $K/k$ une extension de corps séparable.
Alors $K$ est géométriquement normal sur $k$.
\end{lemma}

\begin{proof}
Cela résulte du fait que $k^{perf} \otimes_k K$ est un corps.
En effet, cet anneau est réduit d'après le
lemme \ref{lemma-separable-extension-preserves-reducedness}.
D'après le lemme \ref{lemma-perfection} (ou la
définition \ref{definition-perfection}),
l'extension de corps $k^{perf}/k$ est purement inséparable.
Par conséquent, d'après le lemme \ref{lemma-radicial-integral-bijective}, l'anneau
$k^{perf} \otimes_k K$ possède un unique idéal premier. Un anneau réduit ayant
un unique idéal premier est un corps.
\end{proof}

\begin{lemma}
\label{lemma-geometrically-normal-tensor-normal}
Soit $k$ un corps. Soient $A, B$ des $k$-algèbres. Supposons que $A$ soit géométriquement
normale sur $k$ et que $B$ soit un anneau normal. Alors $A \otimes_k B$ est un anneau
normal.
\end{lemma}

\begin{proof}
Soit $\mathfrak r$ un idéal premier de $A \otimes_k B$. Notons
$\mathfrak p$ (resp.\ $\mathfrak q$) l'idéal premier correspondant de $A$
(resp.\ de $B$). Alors $(A \otimes_k B)_{\mathfrak r}$ est une localisation de
$A_{\mathfrak p} \otimes_k B_{\mathfrak q}$. Il suffit donc de démontrer le
résultat pour l'anneau $A_{\mathfrak p} \otimes_k B_{\mathfrak q}$, voir le
lemme \ref{lemma-localization-normal-ring}
et le
lemme \ref{lemma-localization-geometrically-normal-algebra}.
Nous pouvons donc supposer que $A$ et $B$ sont des anneaux intègres.

\medskip\noindent
Supposons que $A$ et $B$ soient des anneaux intègres dont les corps des fractions respectifs sont $K$ et $L$.
Notons que $B$ est la limite inductive filtrante de ses sous-algèbres normales
de type fini sur $k$ (en effet, $k$ est un anneau de Nagata, voir la
proposition \ref{proposition-ubiquity-nagata},
et la clôture intégrale d'une sous-algèbre de type fini sur $k$ est encore
une sous-algèbre de type fini sur $k$, d'après la
proposition \ref{proposition-nagata-universally-japanese}).
D'après le
lemme \ref{lemma-colimit-normal-ring},
on se ramène au cas où $B$ est de type fini sur $k$.

\medskip\noindent
Supposons que $A$ et $B$ soient des anneaux intègres dont les corps des fractions respectifs sont $K$ et $L$,
et que $B$ soit de type fini sur $k$. Dans ce cas, l'anneau $K \otimes_k B$
est de type fini sur $K$, donc noethérien
(lemme \ref{lemma-Noetherian-permanence}).
En particulier, $K \otimes_k B$ possède un nombre fini d'idéaux premiers minimaux
(lemme \ref{lemma-Noetherian-irreducible-components}).
Comme $A \to A \otimes_k B$ est plat, il s'ensuit que $A \otimes_k B$
possède un nombre fini d'idéaux premiers minimaux (par le théorème de descente pour les morphismes d'anneaux plats --
lemme \ref{lemma-flat-going-down}
-- tous ces idéaux sont situés au-dessus de $(0) \subset A$). Il suffit donc de démontrer
que $A \otimes_k B$ est intégralement clos dans son anneau total des fractions
(lemme \ref{lemma-characterize-reduced-ring-normal}).

\medskip\noindent
Nous affirmons que $K \otimes_k B$ et $A \otimes_k L$ sont tous deux des anneaux normaux.
Si tel est le cas, tout élément $x$ de $Q(A \otimes_k B)$ qui est
entier sur $A \otimes_k B$ appartient (d'après le
lemme \ref{lemma-normal-ring-integrally-closed})
à $K \otimes_k B \cap A \otimes_k L = A \otimes_k B$, et la démonstration est terminée.
Puisque $A \otimes_k L$ est un anneau normal par hypothèse, il suffit de
prouver que $K \otimes_k B$ est normal.

\medskip\noindent
Comme $A$ est géométriquement normale sur $k$, on voit que $K$ est géométriquement normal
sur $k$
(lemme \ref{lemma-localization-geometrically-normal-algebra}),
donc $K$ est géométriquement réduit sur $k$.
Ainsi, $K = \bigcup K_i$ est la réunion d'extensions de corps de type fini
de $k$ qui sont géométriquement réduites
(lemme \ref{lemma-subalgebra-separable}).
Chaque $K_i$ est la localisation d'une $k$-algèbre lisse
(lemme \ref{lemma-localization-smooth-separable}).
Ainsi, $K_i \otimes_k B$ est une localisation d'une $B$-algèbre lisse, donc un anneau normal
(lemme \ref{lemma-normal-goes-up}).
Ainsi, $K \otimes_k B$ est un anneau normal
(lemme \ref{lemma-colimit-normal-ring}),
ce qui conclut la démonstration.
\end{proof}

\begin{lemma}
\label{lemma-geometrically-normal-over-separable-algebraic}
Soit $k'/k$ une extension de corps algébrique séparable.
Soit $A$ une algèbre sur $k'$. Alors $A$ est géométriquement normale
sur $k$ si et seulement si elle est géométriquement normale sur $k'$.
\end{lemma}

\begin{proof}
Soit $L/k$ une extension de corps finie purement inséparable.
Alors $L' = k' \otimes_k L$ est un corps (voir les résultats du chapitre
Corps, section \ref{fields-section-algebraic})
et $A \otimes_k L = A \otimes_{k'} L'$. Ainsi, si
$A$ est géométriquement normale sur $k'$, alors $A$ est géométriquement
normale sur $k$.

\medskip\noindent
Supposons que $A$ soit géométriquement normale sur $k$. Soit $K/k'$ une extension de corps.
Alors
$$
K \otimes_{k'} A = (K \otimes_k A) \otimes_{(k' \otimes_k k')} k'
$$
Puisque $k' \otimes_k k' \to k'$ est une localisation d'après le
lemme \ref{lemma-separable-algebraic-diagonal},
on voit que $K \otimes_{k'} A$
est une localisation d'un anneau normal, donc un anneau normal.
\end{proof}



\section{Algèbres géométriquement régulières}
\label{section-geometrically-regular}

\noindent
Soit $k$ un corps.
Soit $A$ une $k$-algèbre noethérienne.
Soit $K/k$ une extension de corps de type fini.
Alors l'anneau $K \otimes_k A$ est lui aussi noethérien, voir le
lemme \ref{lemma-Noetherian-field-extension}.
Le lemme suivant a donc un sens.

\begin{lemma}
\label{lemma-geometrically-regular}
Soit $k$ un corps. Soit $A$ une $k$-algèbre.
Supposons que $A$ soit noethérienne.
Les propriétés suivantes de $A$ sont équivalentes :
\begin{enumerate}
\item $k' \otimes_k A$ est régulier pour toute extension de corps
$k'/k$ de type fini, et
\item $k' \otimes_k A$ est régulier pour toute extension de corps
$k'/k$ finie purement inséparable.
\end{enumerate}
Ici, un anneau régulier est entendu au sens de la définition \ref{definition-regular}.
\end{lemma}

\begin{proof}
Le lemme a un sens d'après les remarques qui le précèdent.
Il est clair que (1) $\Rightarrow$ (2).

\medskip\noindent
Supposons (2) et soit $K/k$ une extension de corps de type fini.
D'après le lemme \ref{lemma-make-separable}, on peut trouver un diagramme
$$
\xymatrix{
K \ar[r] & K' \\
k \ar[u] \ar[r] & k' \ar[u]
}
$$
où $k'/k$ et $K'/K$ sont des extensions de corps finies purement
inséparables telles que $K'/k'$ soit séparable. D'après le
lemme \ref{lemma-localization-smooth-separable},
il existe une $k'$-algèbre lisse $B$ dont $K'$ est le
corps des fractions de $B$. On peut maintenant raisonner comme suit :
Étape 1 : $k' \otimes_k A$ est un anneau régulier puisque nous avons supposé (2).
Étape 2 : $B \otimes_{k'} k' \otimes_k A$ est un anneau régulier, puisque
le morphisme $k' \otimes_k A \to B \otimes_{k'} k' \otimes_k A$ est lisse
(lemme \ref{lemma-base-change-smooth})
et que la régularité monte le long des morphismes lisses
(lemme \ref{lemma-regular-goes-up}).
Étape 3. $K' \otimes_{k'} k' \otimes_k A = K' \otimes_k A$ est
un anneau régulier, puisqu'il s'agit d'un localisé d'un anneau régulier
(cela découle immédiatement de la définition).
Étape 4. Enfin, $K \otimes_k A$ est un anneau régulier par descente de la
régularité le long du morphisme d'anneaux fidèlement plat
$K \otimes_k A \to K' \otimes_k A$ (lemme \ref{lemma-descent-regular}).
Ceci démontre le lemme.
\end{proof}

\begin{definition}
\label{definition-geometrically-regular}
Soit $k$ un corps. Soit $R$ une $k$-algèbre noethérienne.
La $k$-algèbre $R$ est dite {\it géométriquement régulière} sur $k$ si
les conditions équivalentes du lemme \ref{lemma-geometrically-regular} sont satisfaites.
\end{definition}

\noindent
Il résulte immédiatement de la définition que $K \otimes_k R$ est une algèbre
géométriquement régulière sur $K$, pour toute extension de corps $K$ de
type fini sur $k$. Nous verrons plus loin (Compléments d'algèbre, proposition
\ref{more-algebra-proposition-characterization-geometrically-regular})
qu'il suffit de vérifier que $R \otimes_k k'$ est régulier dès que
$k \subset k' \subset k^{1/p}$ (finie).

\begin{lemma}
\label{lemma-geometrically-regular-descent}
\begin{slogan}
La régularité géométrique descend le long des morphismes fidèlement plats d'algèbres
\end{slogan}
Soit $k$ un corps. Soit $A \to B$ un morphisme fidèlement plat de $k$-algèbres.
Si $B$ est géométriquement régulière sur $k$, alors $A$ l'est aussi.
\end{lemma}

\begin{proof}
Supposons que $B$ soit géométriquement régulière sur $k$.
Soit $k'/k$ une extension de corps finie purement inséparable.
Alors $A \otimes_k k' \to B \otimes_k k'$ est fidèlement plat, puisque c'est un
changement de base de $A \to B$ (d'après les
lemmes \ref{lemma-surjective-spec-radical-ideal} et
\ref{lemma-flat-base-change}),
et $B \otimes_k k'$ est régulier d'après notre
hypothèse sur $B$ en tant que $k$-algèbre. Alors $A \otimes_k k'$ est régulier d'après le
lemme \ref{lemma-descent-regular}.
\end{proof}

\begin{lemma}
\label{lemma-geometrically-regular-goes-up}
Soit $k$ un corps. Soit $A \to B$ un morphisme lisse d'anneaux
entre $k$-algèbres. Si $A$ est géométriquement régulière sur $k$,
alors $B$ est géométriquement régulière sur $k$.
\end{lemma}

\begin{proof}
Soit $k'/k$ une extension de corps de type fini.
Alors $A \otimes_k k' \to B \otimes_k k'$ est un morphisme lisse d'anneaux
(lemme \ref{lemma-base-change-smooth}) et $A \otimes_k k'$
est régulier. Par conséquent, $B \otimes_k k'$ est régulier d'après le
lemme \ref{lemma-regular-goes-up}.
\end{proof}

\begin{lemma}
\label{lemma-geometrically-regular-over-subfields}
Soit $k$ un corps. Soit $A$ une algèbre sur $k$.
Soit $k = \colim k_i$ une limite inductive filtrante de sous-corps.
Si $A$ est géométriquement régulière sur chaque $k_i$, alors
$A$ est géométriquement régulière sur $k$.
\end{lemma}

\begin{proof}
Soit $k'/k$ une extension de corps finie purement inséparable.
On peut obtenir $k'$ en adjoignant à $k$ un nombre fini de variables et en
imposant un nombre fini de relations polynomiales. On voit donc qu'il
existe un $i$ et une extension de corps finie purement inséparable
$k_i'/k_i$ tels que $k' = k \otimes_{k_i} k_i'$.
Ainsi, $A \otimes_k k' = A \otimes_{k_i} k_i'$, et le lemme est clair.
\end{proof}

\begin{lemma}
\label{lemma-geometrically-regular-over-separable-algebraic}
Soit $k'/k$ une extension algébrique séparable de corps.
Soit $A$ une algèbre sur $k'$. Alors $A$ est géométriquement
régulière sur $k$ si et seulement si elle est géométriquement régulière sur $k'$.
\end{lemma}

\begin{proof}
Soit $L/k$ une extension de corps finie purement inséparable.
Alors $L' = k' \otimes_k L$ est un corps (voir les résultats du chapitre
Corps, section \ref{fields-section-algebraic})
et $A \otimes_k L = A \otimes_{k'} L'$. Ainsi, si
$A$ est géométriquement régulière sur $k'$, alors $A$ est géométriquement
régulière sur $k$.

\medskip\noindent
Supposons que $A$ soit géométriquement régulière sur $k$. Puisque $k'$
est la limite inductive filtrante des extensions finies de $k$ qu'il contient, on peut
supposer, d'après le lemme \ref{lemma-geometrically-regular-over-subfields},
que $k'/k$ est finie séparable. Considérons les morphismes d'anneaux
$$
k' \to A \otimes_k k' \to A.
$$
Remarquons que $A \otimes_k k'$ est géométriquement régulière sur $k'$
en tant que changement de base de $A$ à $k'$. Remarquons que $A \otimes_k k' \to A$
est le changement de base de $k' \otimes_k k' \to k'$ par le morphisme
$k' \to A$. Puisque $k'/k$ est une extension d'anneaux \'etale, on
voit que $k' \otimes_k k' \to k'$ est \'etale
(lemme \ref{lemma-etale}). Par conséquent, $A$ est
géométriquement régulière sur $k'$ d'après le
lemme \ref{lemma-geometrically-regular-goes-up}.
\end{proof}




\section{Algèbres géométriquement de Cohen-Macaulay}
\label{section-geometrically-CM}

\noindent
Le titre de cette section est quelque peu impropre, puisque les algèbres
de Cohen-Macaulay sont automatiquement géométriquement de Cohen-Macaulay. Plus précisément, voir
le lemme \ref{lemma-extend-field-CM-locus}
et
le lemme \ref{lemma-CM-geometrically-CM}
ci-dessous.

\begin{lemma}
\label{lemma-tensor-fields-CM}
Soit $k$ un corps et soient $K/k$ et $L/k$
deux extensions de corps telles que l'une d'elles soit de type fini.
Alors $K \otimes_k L$ est un anneau noethérien de Cohen-Macaulay.
\end{lemma}

\begin{proof}
L'anneau $K \otimes_k L$ est noethérien d'après
le lemme \ref{lemma-Noetherian-field-extension}.
Supposons que $K$ soit une extension finie de l'extension purement transcendante
$k(t_1, \ldots, t_r)$. Alors
$k(t_1, \ldots, t_r) \otimes_k L \to K \otimes_k L$
est un morphisme d'anneaux fini et libre. D'après
le lemme \ref{lemma-finite-flat-over-regular-CM},
il suffit de montrer que $k(t_1, \ldots, t_r) \otimes_k L$ est de Cohen-Macaulay.
Cela est clair, car il s'agit d'une localisation de l'anneau de polynômes
$L[t_1, \ldots, t_r]$. (Voir par exemple le
lemme \ref{lemma-CM-polynomial-algebra}, qui affirme qu'un anneau de polynômes
est de Cohen-Macaulay.)
\end{proof}

\begin{lemma}
\label{lemma-CM-geometrically-CM}
Soit $k$ un corps. Soit $S$ une $k$-algèbre noethérienne.
Soit $K/k$ une extension de corps de type fini,
et posons $S_K = K \otimes_k S$. Soit $\mathfrak q \subset S$
un idéal premier de $S$. Soit $\mathfrak q_K \subset S_K$ un idéal premier
de $S_K$ au-dessus de $\mathfrak q$. Alors $S_{\mathfrak q}$ est de Cohen-Macaulay
si et seulement si $(S_K)_{\mathfrak q_K}$ est de Cohen-Macaulay.
\end{lemma}

\begin{proof}
D'après
le lemme \ref{lemma-Noetherian-field-extension},
l'anneau $S_K$ est noethérien. Ainsi,
$S_{\mathfrak q} \to (S_K)_{\mathfrak q_K}$ est un morphisme local plat
d'anneaux locaux noethériens. Notons que la fibre
$$
(S_K)_{\mathfrak q_K} / \mathfrak q (S_K)_{\mathfrak q_K}
\cong (\kappa(\mathfrak q) \otimes_k K)_{\mathfrak q'}
$$
est la localisation de l'anneau de Cohen-Macaulay (lemme \ref{lemma-tensor-fields-CM})
$\kappa(\mathfrak q) \otimes_k K$ en un idéal premier convenable
$\mathfrak q'$. Le lemme résulte donc du lemme \ref{lemma-CM-goes-up}.
\end{proof}






\section{Limites inductives et morphismes de présentation finie, II}
\label{section-colimits-finite-presentation}

\noindent
Cette section fait suite à la section \ref{section-colimits-flat}.

\medskip\noindent
Nous commençons par une application de l'ouverture de la platitude.
Elle affirme que nous pouvons approcher les modules plats par des modules plats,
ce qui est utile.

\begin{lemma}
\label{lemma-flat-finite-presentation-limit-flat}
Soit $R \to S$ un morphisme d'anneaux.
Soit $M$ un $S$-module.
Supposons que
\begin{enumerate}
\item $R \to S$ soit de présentation finie,
\item $M$ soit un $S$-module de présentation finie, et
\item $M$ soit plat sur $R$.
\end{enumerate}
Dans ce cas, nous avons les assertions suivantes :
\begin{enumerate}
\item Il existe une $\mathbf{Z}$-algèbre $R_0$ de type fini, un
morphisme d'anneaux $R_0 \to S_0$ de type fini et un $S_0$-module $M_0$ de type fini
tels que $M_0$ soit plat sur $R_0$, ainsi que des morphismes d'anneaux
$R_0 \to R$ et $S_0 \to S$ et un morphisme de $S_0$-modules $M_0 \to M$
tels que $S \cong R \otimes_{R_0} S_0$ et $M = S \otimes_{S_0} M_0$.
\item Si $R = \colim_{\lambda \in \Lambda} R_\lambda$ est présenté
comme une limite inductive filtrante, alors il existe un $\lambda$, un morphisme d'anneaux
$R_\lambda \to S_\lambda$ de présentation finie, et un $S_\lambda$-module
$M_\lambda$ de présentation finie tels que $M_\lambda$ soit plat sur
$R_\lambda$ et que $S = R \otimes_{R_\lambda} S_\lambda$ et
$M = S \otimes_{S_{\lambda}} M_\lambda$.
\item Si
$$
(R \to S, M) =
\colim_{\lambda \in \Lambda}
(R_\lambda \to S_\lambda, M_\lambda)
$$
est présenté comme une limite inductive filtrante telle que
\begin{enumerate}
\item $R_\mu \otimes_{R_\lambda} S_\lambda \to S_\mu$ et
$S_\mu \otimes_{S_\lambda} M_\lambda \to M_\mu$ soient des isomorphismes
pour $\mu \geq \lambda$,
\item $R_\lambda \to S_\lambda$ soit de présentation finie,
\item $M_\lambda$ soit un $S_\lambda$-module de présentation finie,
\end{enumerate}
alors, pour tout $\lambda$ suffisamment grand, le module $M_\lambda$
est plat sur $R_\lambda$.
\end{enumerate}
\end{lemma}

\begin{proof}
Écrivons d'abord $(R \to S, M)$ comme la limite inductive filtrante d'un système
$(R_\lambda \to S_\lambda, M_\lambda)$ comme dans le
lemme \ref{lemma-limit-module-finite-presentation}.
Soit $\mathfrak q \subset S$ un idéal premier.
Soient $\mathfrak p \subset R$, $\mathfrak q_\lambda \subset S_\lambda$
et $\mathfrak p_\lambda \subset R_\lambda$ les idéaux premiers correspondants.
Comme on l'a vu dans la démonstration du théorème \ref{theorem-openness-flatness},
$$
((R_\lambda)_{\mathfrak p_\lambda},
(S_\lambda)_{\mathfrak q_\lambda},
(M_\lambda)_{\mathfrak q_{\lambda}})
$$
est un système comme dans le
lemme \ref{lemma-limit-module-essentially-finite-presentation}, et
le lemme \ref{lemma-colimit-eventually-flat} montre donc
qu'il existe $\lambda_{\mathfrak q} \in \Lambda$
tel que, pour tout $\lambda \geq \lambda_{\mathfrak q}$,
le module $M_\lambda$ soit plat sur
$R_\lambda$ en l'idéal premier $\mathfrak q_{\lambda}$.

\medskip\noindent
D'après le théorème \ref{theorem-openness-flatness},
$U_\lambda \subset \Spec(S_\lambda)$ désigne l'ensemble des idéaux premiers en lesquels $M_\lambda$
est plat sur $R_\lambda$ ; cet ensemble $U_\lambda$ est ouvert.
Notons $V_\lambda \subset \Spec(S)$ l'image réciproque de
$U_\lambda$ par le morphisme $\Spec(S) \to \Spec(S_\lambda)$.
L'argument ci-dessus montre que, pour tout $\mathfrak q \in \Spec(S)$,
il existe $\lambda_{\mathfrak q}$ tel que
$\mathfrak q \in V_\lambda$ pour tout $\lambda \geq \lambda_{\mathfrak q}$.
Comme $\Spec(S)$ est quasi-compact, cela implique qu'il
existe un indice $\lambda_0 \in \Lambda$ tel que
$V_{\lambda_0} = \Spec(S)$.

\medskip\noindent
Le complémentaire $\Spec(S_{\lambda_0}) \setminus U_{\lambda_0}$
est $V(I)$ pour un idéal $I \subset S_{\lambda_0}$. Comme
$V_{\lambda_0} = \Spec(S)$, nous voyons que $IS = S$.
Choisissons $f_1, \ldots, f_r \in I$ et $s_1, \ldots, s_r \in S$ tels
que $\sum f_i s_i = 1$. Puisque $\colim S_\lambda = S$, quitte à
augmenter $\lambda_0$, nous pouvons supposer qu'il existe des
$s_{i, \lambda_0} \in S_{\lambda_0}$ tels que
$\sum f_i s_{i, \lambda_0} = 1$.
Ainsi, pour ce $\lambda_0$, nous avons
$U_{\lambda_0} = \Spec(S_{\lambda_0})$.
Ceci démontre (1).

\medskip\noindent
Démonstration de (2). Soit $(R_0 \to S_0, M_0)$ comme dans (1), et supposons que
$R = \colim R_\lambda$. Comme $R_0$ est une $\mathbf{Z}$-algèbre de type
fini, il existe un $\lambda$ et un morphisme $R_0 \to R_\lambda$ tels
que $R_0 \to R_\lambda \to R$ soit le morphisme donné $R_0 \to R$ (voir le
lemme \ref{lemma-characterize-finite-presentation}).
Alors, l'assertion (2) résulte du choix $S_\lambda = R_\lambda \otimes_{R_0} S_0$
et $M_\lambda = S_\lambda \otimes_{S_0} M_0$.

\medskip\noindent
Enfin, démontrons (3). Soit
$(R_\lambda \to S_\lambda, M_\lambda)$ comme dans (3). Choisissons
$(R_0 \to S_0, M_0)$ et $R_0 \to R$ comme dans (1).
Comme dans la démonstration de (2), il existe $\lambda_0$ et un morphisme d'anneaux
$R_0 \to R_{\lambda_0}$ tels que $R_0 \to R_{\lambda_0} \to R$ soit le
morphisme donné $R_0 \to R$. Comme $S_0$ est de présentation finie sur $R_0$ et que
$S = \colim S_\lambda$, nous voyons que, pour un $\lambda_1 \geq \lambda_0$,
il existe un morphisme de $R_0$-algèbres $S_0 \to S_{\lambda_1}$ tel que la
composée $S_0 \to S_{\lambda_1} \to S$ soit le morphisme donné $S_0 \to S$
(voir le lemme \ref{lemma-characterize-finite-presentation}).
Pour tout $\lambda \geq \lambda_1$, cela donne des morphismes
$$
\Psi_{\lambda} :
R_\lambda \otimes_{R_0} S_0
\longrightarrow
R_\lambda \otimes_{R_{\lambda_1}} S_{\lambda_1}
\cong
S_\lambda
$$
où le dernier isomorphisme résulte de l'hypothèse. Par construction,
$\colim_\lambda \Psi_\lambda$ est un isomorphisme. Ainsi, $\Psi_\lambda$
est un isomorphisme pour tout $\lambda$ suffisamment grand d'après le
lemme \ref{lemma-colimit-category-fp-algebras}.
De même, il existe $\lambda_2 \geq \lambda_1$
et un morphisme de $S_0$-modules $M_0 \to M_{\lambda_2}$ tel que
$M_0 \to M_{\lambda_2} \to M$ soit le morphisme donné
$M_0 \to M$ (voir le lemme \ref{lemma-module-map-property-in-colimit}).
Pour $\lambda \geq \lambda_2$, il existe un morphisme induit
$$
S_\lambda \otimes_{S_0} M_0
\longrightarrow
S_\lambda \otimes_{S_{\lambda_2}} M_{\lambda_2}
\cong
M_\lambda
$$
et, pour $\lambda$ suffisamment grand, ce morphisme est un isomorphisme d'après le
lemme \ref{lemma-colimit-category-fp-modules}.
Ceci implique (3), car $M_0$ est plat sur $R_0$.
\end{proof}

\begin{lemma}
\label{lemma-descend-faithfully-flat-finite-presentation}
Soient $R \to A \to B$ des morphismes d'anneaux.
Supposons $A \to B$ fidèlement plat et de présentation finie.
Alors il existe un diagramme commutatif
$$
\xymatrix{
R \ar[r] \ar@{=}[d] &
A_0 \ar[d] \ar[r] &
B_0 \ar[d] \\
R \ar[r] & A \ar[r] & B
}
$$
où $R \to A_0$ est de présentation finie,
$A_0 \to B_0$ est fidèlement plat et de présentation finie,
et $B = A \otimes_{A_0} B_0$.
\end{lemma}

\begin{proof}
Démontrons d'abord le lemme en remplaçant $R$ par $\mathbf{Z}$.
D'après le lemme \ref{lemma-flat-finite-presentation-limit-flat},
il existe un diagramme
$$
\xymatrix{
A_0 \ar[r] \ar[d] & A \ar[d] \\
B_0 \ar[r] & B
}
$$
où $A_0$ est de type fini sur $\mathbf{Z}$ et $B_0$ est plat et de présentation
finie sur $A_0$, tel que $B = A \otimes_{A_0} B_0$.
Comme $A_0 \to B_0$ est plat et de présentation finie, l'image de
$\Spec(B_0) \to \Spec(A_0)$ est ouverte, voir la
proposition \ref{proposition-fppf-open}. Le complémentaire de cette image
est donc $V(I_0)$ pour un idéal $I_0 \subset A_0$.
Comme $A \to B$ est fidèlement
plat, le morphisme $\Spec(B) \to \Spec(A)$ est surjectif, voir le
lemme \ref{lemma-ff-rings}.
Utilisons maintenant le fait que
le changement de base de l'image est l'image du changement de base.
Ainsi, $I_0A = A$. Choisissons une relation
$\sum f_i r_i = 1$, avec $r_i \in A$, $f_i \in I_0$. Alors, après avoir
agrandi $A_0$ pour qu'il contienne les éléments $r_i$ (et avoir agrandi
$B_0$ en conséquence), nous voyons que $A_0 \to B_0$ est également surjectif
sur les spectres, c'est-à-dire fidèlement plat.

\medskip\noindent
Le lemme est donc vrai dans le cas où $R = \mathbf{Z}$.
Dans le cas général, prenons la solution $A_0' \to B_0'$
que nous venons d'obtenir et posons $A_0 = A_0' \otimes_{\mathbf{Z}} R$,
$B_0 = B_0' \otimes_{\mathbf{Z}} R$.
\end{proof}

\begin{lemma}
\label{lemma-colimit-finite}
Soit $A = \colim_{i \in I} A_i$ une limite inductive filtrante d'anneaux.
Soient $0 \in I$ et $\varphi_0 : B_0 \to C_0$ un morphisme de $A_0$-algèbres.
Supposons que
\begin{enumerate}
\item $A \otimes_{A_0} B_0 \to A \otimes_{A_0} C_0$ soit fini,
\item $C_0$ soit de type fini sur $B_0$.
\end{enumerate}
Alors il existe $i \geq 0$ tel que le morphisme
$A_i \otimes_{A_0} B_0 \to A_i \otimes_{A_0} C_0$
soit fini.
\end{lemma}

\begin{proof}
Soient $x_1, \ldots, x_m$ des générateurs de $C_0$ sur $B_0$.
Choisissons des polynômes unitaires $P_j \in A \otimes_{A_0} B_0[T]$ tels
que $P_j(1 \otimes x_j) = 0$ dans $A \otimes_{A_0} C_0$. Pour un certain
$i \geq 0$, nous pouvons trouver $P_{j, i} \in A_i \otimes_{A_0} B_0[T]$
ayant pour image $P_j$. Comme $\otimes$
commute aux limites inductives, nous voyons que $P_{j, i}(1 \otimes x_j)$ est nul
dans $A_i \otimes_{A_0} C_0$, quitte à augmenter $i$.
Alors cet $i$ convient.
\end{proof}

\begin{lemma}
\label{lemma-colimit-surjective}
Soit $A = \colim_{i \in I} A_i$ une limite inductive filtrante d'anneaux.
Soient $0 \in I$ et $\varphi_0 : B_0 \to C_0$ un morphisme de $A_0$-algèbres.
Supposons que
\begin{enumerate}
\item $A \otimes_{A_0} B_0 \to A \otimes_{A_0} C_0$ soit surjectif,
\item $C_0$ soit de type fini sur $B_0$.
\end{enumerate}
Alors, pour un certain $i \geq 0$, le morphisme
$A_i \otimes_{A_0} B_0 \to A_i \otimes_{A_0} C_0$
est surjectif.
\end{lemma}

\begin{proof}
Soient $x_1, \ldots, x_m$ des générateurs de $C_0$ sur $B_0$.
Choisissons $b_j \in A \otimes_{A_0} B_0$ ayant pour image $1 \otimes x_j$ dans
$A \otimes_{A_0} C_0$. Pour un certain $i \geq 0$, nous pouvons trouver
$b_{j, i} \in A_i \otimes_{A_0} B_0$ ayant pour image $b_j$.
Après avoir augmenté $i$, nous pouvons supposer que $b_{j, i}$ a pour image
$1 \otimes x_j$ dans $A_i \otimes_{A_0} C_0$ pour tout $j = 1, \ldots, m$.
Alors cet $i$ convient.
\end{proof}

\begin{lemma}
\label{lemma-colimit-unramified}
Soit $A = \colim_{i \in I} A_i$ une limite inductive filtrante d'anneaux.
Soient $0 \in I$ et $\varphi_0 : B_0 \to C_0$ un morphisme de $A_0$-algèbres.
Supposons que
\begin{enumerate}
\item $A \otimes_{A_0} B_0 \to A \otimes_{A_0} C_0$ soit non ramifié,
\item $C_0$ soit de type fini sur $B_0$.
\end{enumerate}
Alors, pour un certain $i \geq 0$, le morphisme
$A_i \otimes_{A_0} B_0 \to A_i \otimes_{A_0} C_0$
est non ramifié.
\end{lemma}

\begin{proof}
Posons $B_i = A_i \otimes_{A_0} B_0$, $C_i = A_i \otimes_{A_0} C_0$,
$B = A \otimes_{A_0} B_0$ et $C = A \otimes_{A_0} C_0$.
Soient $x_1, \ldots, x_m$ des générateurs de $C_0$ sur $B_0$.
Alors $\text{d}x_1, \ldots, \text{d}x_m$ engendrent $\Omega_{C_0/B_0}$
sur $C_0$, et leurs images engendrent $\Omega_{C_i/B_i}$ sur $C_i$
(lemmes \ref{lemma-differentials-polynomial-ring} et
\ref{lemma-differential-seq}).
Observons que $0 = \Omega_{C/B} = \colim \Omega_{C_i/B_i}$
(lemme \ref{lemma-colimit-differentials}).
Il existe donc un $i$ tel que $\text{d}x_1, \ldots, \text{d}x_m$
aient une image nulle et, par conséquent, $\Omega_{C_i/B_i} = 0$, comme souhaité.
\end{proof}

\begin{lemma}
\label{lemma-colimit-isomorphism}
Soit $A = \colim_{i \in I} A_i$ une limite inductive filtrante d'anneaux.
Soient $0 \in I$ et $\varphi_0 : B_0 \to C_0$ un morphisme de $A_0$-algèbres.
Supposons que
\begin{enumerate}
\item $A \otimes_{A_0} B_0 \to A \otimes_{A_0} C_0$ soit
un isomorphisme,
\item $B_0 \to C_0$ soit de présentation finie.
\end{enumerate}
Alors, pour un certain $i \geq 0$, le morphisme
$A_i \otimes_{A_0} B_0 \to A_i \otimes_{A_0} C_0$ est
un isomorphisme.
\end{lemma}

\begin{proof}
D'après le lemme \ref{lemma-colimit-surjective}, il existe un $i$ tel que
$A_i \otimes_{A_0} B_0 \to A_i \otimes_{A_0} C_0$ soit
surjectif. Comme le morphisme est de présentation finie,
son noyau est un idéal de type fini. Soient
$g_1, \ldots, g_r \in A_i \otimes_{A_0} B_0$ des générateurs de ce noyau.
Nous pouvons alors choisir $i' \geq i$ tel que les $g_j$ aient une image nulle
dans $A_{i'} \otimes_{A_0} B_0$. Alors
$A_{i'} \otimes_{A_0} B_0 \to A_{i'} \otimes_{A_0} C_0$ est
un isomorphisme.
\end{proof}

\begin{lemma}
\label{lemma-colimit-etale}
Soit $A = \colim_{i \in I} A_i$ une limite inductive filtrante d'anneaux.
Soient $0 \in I$ et $\varphi_0 : B_0 \to C_0$ un morphisme de $A_0$-algèbres.
Supposons que
\begin{enumerate}
\item $A \otimes_{A_0} B_0 \to A \otimes_{A_0} C_0$ soit \'etale,
\item $B_0 \to C_0$ soit de présentation finie.
\end{enumerate}
Alors, pour un certain $i \geq 0$, le morphisme
$A_i \otimes_{A_0} B_0 \to A_i \otimes_{A_0} C_0$
est \'etale.
\end{lemma}

\begin{proof}
Écrivons $C_0 = B_0[x_1, \ldots, x_n]/(f_{1, 0}, \ldots, f_{m, 0})$.
Écrivons $B_i = A_i \otimes_{A_0} B_0$ et $C_i = A_i \otimes_{A_0} C_0$.
Notons que $C_i = B_i[x_1, \ldots, x_n]/(f_{1, i}, \ldots, f_{m, i})$,
où $f_{j, i}$ est l'image de $f_{j, 0}$ dans l'anneau de polynômes
sur $B_i$. Écrivons $B = A \otimes_{A_0} B_0$ et $C = A \otimes_{A_0} C_0$.
Notons que $C = B[x_1, \ldots, x_n]/(f_1, \ldots, f_m)$,
où $f_j$ est l'image de $f_{j, 0}$ dans l'anneau de polynômes
sur $B$. L'hypothèse est que le morphisme
$$
\text{d} :
(f_1, \ldots, f_m)/(f_1, \ldots, f_m)^2
\longrightarrow
\bigoplus C \text{d}x_k
$$
est un isomorphisme. Ainsi, pour $i$ suffisamment grand, nous pouvons trouver des éléments
$$
\xi_{k, i} \in (f_{1, i}, \ldots, f_{m, i})/(f_{1, i}, \ldots, f_{m, i})^2
$$
tels que $\text{d}\xi_{k, i} = \text{d}x_k$ dans $\bigoplus C_i \text{d}x_k$.
De plus, quitte à augmenter $i$, nous voyons que
$\sum (\partial f_{j, i}/\partial x_k) \xi_{k, i} =
f_{j, i} \bmod (f_{1, i}, \ldots, f_{m, i})^2$,
puisque cette égalité est vraie à la limite. Alors cet $i$ convient.
\end{proof}

\begin{lemma}
\label{lemma-colimit-smooth}
Soit $A = \colim_{i \in I} A_i$ une limite inductive filtrante d'anneaux.
Soient $0 \in I$ et $\varphi_0 : B_0 \to C_0$ un morphisme de $A_0$-algèbres.
Supposons que
\begin{enumerate}
\item $A \otimes_{A_0} B_0 \to A \otimes_{A_0} C_0$ soit lisse,
\item $B_0 \to C_0$ soit de présentation finie.
\end{enumerate}
Alors, pour un certain $i \geq 0$, le morphisme
$A_i \otimes_{A_0} B_0 \to A_i \otimes_{A_0} C_0$ est lisse.
\end{lemma}

\begin{proof}
Écrivons $C_0 = B_0[x_1, \ldots, x_n]/(f_{1, 0}, \ldots, f_{m, 0})$.
Écrivons $B_i = A_i \otimes_{A_0} B_0$ et $C_i = A_i \otimes_{A_0} C_0$.
Notons que $C_i = B_i[x_1, \ldots, x_n]/(f_{1, i}, \ldots, f_{m, i})$,
où $f_{j, i}$ est l'image de $f_{j, 0}$ dans l'anneau de polynômes
sur $B_i$. Écrivons $B = A \otimes_{A_0} B_0$ et $C = A \otimes_{A_0} C_0$.
Notons que $C = B[x_1, \ldots, x_n]/(f_1, \ldots, f_m)$,
où $f_j$ est l'image de $f_{j, 0}$ dans l'anneau de polynômes
sur $B$. L'hypothèse est que le morphisme
$$
\text{d} :
(f_1, \ldots, f_m)/(f_1, \ldots, f_m)^2
\longrightarrow
\bigoplus C \text{d}x_k
$$
est une injection scindée. Soient $\xi_k \in (f_1, \ldots, f_m)/(f_1, \ldots, f_m)^2$
des éléments tels que $\sum (\partial f_j/\partial x_k) \xi_k =
f_j \bmod (f_1, \ldots, f_m)^2$. Alors, pour $i$ suffisamment grand, nous pouvons
trouver des éléments
$$
\xi_{k, i} \in (f_{1, i}, \ldots, f_{m, i})/(f_{1, i}, \ldots, f_{m, i})^2
$$
tels que $\sum (\partial f_{j, i}/\partial x_k) \xi_{k, i} =
f_{j, i} \bmod (f_{1, i}, \ldots, f_{m, i})^2$,
puisque cette égalité est vraie à la limite. Alors cet $i$ convient.
\end{proof}

\begin{lemma}
\label{lemma-colimit-lci}
Soit $A = \colim_{i \in I} A_i$ une limite inductive filtrante d'anneaux.
Soient $0 \in I$ et $\varphi_0 : B_0 \to C_0$ un morphisme de $A_0$-algèbres.
Supposons que
\begin{enumerate}
\item $A \otimes_{A_0} B_0 \to A \otimes_{A_0} C_0$ soit
syntomique (resp.\ une intersection complète globale relative),
\item $C_0$ soit de présentation finie sur $B_0$.
\end{enumerate}
Alors il existe $i \geq 0$ tel que le morphisme
$A_i \otimes_{A_0} B_0 \to A_i \otimes_{A_0} C_0$
soit syntomique (resp.\ une intersection complète globale relative).
\end{lemma}

\begin{proof}
Supposons que $A \otimes_{A_0} B_0 \to A \otimes_{A_0} C_0$ soit une intersection
complète globale relative.
D'après le lemme \ref{lemma-relative-global-complete-intersection-Noetherian},
il existe une $\mathbf{Z}$-algèbre $R$ de type fini,
un morphisme d'anneaux $R \to A \otimes_{A_0} B_0$, une intersection
complète globale relative $R \to S$ et un isomorphisme
$$
(A \otimes_{A_0} B_0) \otimes_R S
\longrightarrow
A \otimes_{A_0} C_0
$$
Comme $R$ est de type fini (et donc de présentation finie)
sur $\mathbf{Z}$, il existe un $i$ et un morphisme
$R \to A_i \otimes_{A_0} B_0$ relevant le morphisme $R \to A \otimes_{A_0} B_0$ ;
voir le lemme \ref{lemma-characterize-finite-presentation}.
En utilisant le même lemme, il existe $i' \geq i$ tel que
$(A_i \otimes_{A_0} B_0) \otimes_R S \to A \otimes_{A_0} C_0$
provienne d'un morphisme
$(A_i \otimes_{A_0} B_0) \otimes_R S \to A_{i'} \otimes_{A_0} C_0$.
Nous pouvons donc supposer, après avoir remplacé $i$ par $i'$,
que le morphisme affiché provient d'un morphisme de $A_i \otimes_{A_0} B_0$-algèbres
$$
(A_i \otimes_{A_0} B_0) \otimes_R S
\longrightarrow
A_i \otimes_{A_0} C_0
$$
D'après le lemme \ref{lemma-colimit-isomorphism}, après avoir augmenté $i$, ce
morphisme est un isomorphisme. Ceci achève la démonstration dans ce cas, car le changement
de base d'une intersection complète globale relative est une intersection
complète globale relative d'après le
lemme \ref{lemma-base-change-relative-global-complete-intersection}.

\medskip\noindent
Supposons que $A \otimes_{A_0} B_0 \to A \otimes_{A_0} C_0$ soit syntomique.
Il existe alors des éléments $g_1, \ldots, g_m$ de
$A \otimes_{A_0} C_0$ qui engendrent l'idéal unité tels que
$A \otimes_{A_0} B_0 \to (A \otimes_{A_0} C_0)_{g_j}$ soit une
intersection complète globale relative, voir le lemme \ref{lemma-syntomic}.
Nous pouvons trouver un $i$ et des éléments $g_{i, j} \in A_i \otimes_{A_0} C_0$
ayant pour images les $g_j$. Après avoir augmenté $i$, nous pouvons supposer que
$g_{i, 1}, \ldots, g_{i, m}$ engendrent l'idéal unité
de $A_i \otimes_{A_0} C_0$. Le résultat du paragraphe précédent
implique qu'après avoir augmenté $i$, nous pouvons supposer que les morphismes
$A_i \otimes_{A_0} B_0 \to (A_i \otimes_{A_0} C_0)_{g_{i, j}}$
sont des intersections complètes globales relatives.
Alors $A_i \otimes_{A_0} B_0 \to A_i \otimes_{A_0} C_0$
est syntomique d'après le lemme \ref{lemma-local-syntomic}
(et le lemme \ref{lemma-syntomic}, déjà utilisé).
\end{proof}














\noindent
Le lemme suivant est une application des résultats ci-dessus
qui ne semble trouver naturellement sa place nulle part ailleurs.

\begin{lemma}
\label{lemma-fppf-fpqf}
Soit $R \to S$ un morphisme d'anneaux fidèlement plat et de présentation finie.
Alors il existe un diagramme commutatif
$$
\xymatrix{
S \ar[rr] & & S' \\
& R \ar[lu] \ar[ru]
}
$$
où $R \to S'$ est quasi-fini, fidèlement plat et de présentation finie.
\end{lemma}

\begin{proof}
Dans une première étape, nous ramenons ce lemme au cas où $R$ est de type
fini sur $\mathbf{Z}$.
D'après le lemme \ref{lemma-descend-faithfully-flat-finite-presentation},
il existe un diagramme
$$
\xymatrix{
S_0 \ar[r] & S \\
R_0 \ar[u] \ar[r] & R \ar[u]
}
$$
où $R_0$ est de type fini sur $\mathbf{Z}$
et où $S_0$ est fidèlement plat et de présentation finie sur $R_0$,
tel que $S = R \otimes_{R_0} S_0$.
Si nous démontrons le lemme pour le morphisme d'anneaux $R_0 \to S_0$, alors le lemme
en résulte pour $R \to S$ par changement de base, puisque le changement de base
d'un morphisme d'anneaux quasi-fini est quasi-fini, voir le
lemme \ref{lemma-quasi-finite-base-change}. (Bien entendu, nous
utilisons également le fait que les changements de base de morphismes plats sont plats et que les
changements de base de morphismes de présentation finie sont de présentation finie.)

\medskip\noindent
Supposons que $R \to S$ soit un morphisme d'anneaux fidèlement plat et de présentation finie,
et que $R$ soit noethérien (ce que nous pouvons supposer d'après le paragraphe
précédent). Soit $W \subset \Spec(S)$ l'ouvert du
lemme \ref{lemma-finite-presentation-flat-CM-locus-open}.
Comme $R \to S$ est fidèlement plat, le morphisme $\Spec(S) \to \Spec(R)$
est surjectif, voir le lemme \ref{lemma-ff-rings}.
D'après le lemme \ref{lemma-generic-CM-flat-finite-presentation},
le morphisme $W \to \Spec(R)$ est lui aussi surjectif.
Ainsi, en remplaçant $S$ par un produit $S_{g_1} \times \ldots \times S_{g_m}$,
nous pouvons supposer que $W = \Spec(S)$ ; ici, nous utilisons le fait que $\Spec(R)$
est quasi-compact (lemme \ref{lemma-quasi-compact}) et que le morphisme
$\Spec(S) \to \Spec(R)$ est ouvert
(proposition \ref{proposition-fppf-open}).
Supposons que $\mathfrak p \subset R$ soit un idéal premier. Choisissons un idéal premier
$\mathfrak q \subset S$ situé au-dessus de $\mathfrak p$ et correspondant
à un idéal maximal de l'anneau de la fibre $S \otimes_R \kappa(\mathfrak p)$.
L'anneau local noethérien
$\overline{S}_{\mathfrak q} = S_{\mathfrak q}/\mathfrak pS_{\mathfrak q}$
est de Cohen-Macaulay, disons de dimension $d$. Nous pouvons choisir $f_1, \ldots, f_d$
dans l'idéal maximal de $S_{\mathfrak q}$ dont les images forment une suite régulière
dans $\overline{S}_{\mathfrak q}$. Choisissons un dénominateur commun
$g \in S$, $g \not \in \mathfrak q$, de $f_1, \ldots, f_d$, et considérons
la $R$-algèbre
$$
S' = S_g/(f_1, \ldots, f_d).
$$
Par construction, il existe un idéal premier $\mathfrak q' \subset S'$
au-dessus de $\mathfrak p$ et correspondant à $\mathfrak q$ (par
$S_g \to S'$). Toujours par construction, le morphisme d'anneaux $R \to S'$ est
quasi-fini en $\mathfrak q'$, car l'anneau local
$$
S'_{\mathfrak q'}/\mathfrak pS'_{\mathfrak q'} =
S_{\mathfrak q}/\big((f_1, \ldots, f_d) + \mathfrak pS_{\mathfrak q}\big) =
\overline{S}_{\mathfrak q}/(\overline{f}_1, \ldots, \overline{f}_d)
$$
est de dimension zéro, voir le lemme \ref{lemma-isolated-point-fibre}.
Toujours par construction, $R \to S'$ est de présentation finie.
Enfin, d'après le lemme \ref{lemma-grothendieck-regular-sequence}, le morphisme local d'anneaux
$R_{\mathfrak p} \to S'_{\mathfrak q'}$ est plat (c'est ici que nous
utilisons le fait que $R$ est noethérien). Ainsi, par ouverture de la platitude
(théorème \ref{theorem-openness-flatness}) et ouverture de la quasi-finitude
(lemme \ref{lemma-quasi-finite-open}),
nous pouvons, après avoir remplacé
$g$ par $gg'$ pour un $g' \in S$ convenable, $g' \not \in \mathfrak q$,
supposer que $R \to S'$ est plat et quasi-fini.
L'image de $\Spec(S') \to \Spec(R)$ est ouverte et
contient $\mathfrak p$. Autrement dit, nous avons montré qu'il existe
un anneau $S'$ comme dans l'énoncé du lemme (à l'exception éventuelle
de la fidélité) dont l'image contient tout idéal premier donné.
En utilisant une fois encore la quasi-compacité de $\Spec(R)$,
nous voyons qu'un produit fini de tels anneaux convient.
\end{proof}





\begin{multicols}{2}[\section{Autres chapitres}]
\noindent
Préliminaires
\begin{enumerate}
\item \hyperref[introduction-section-phantom]{Introduction}
\item \hyperref[conventions-section-phantom]{Conventions}
\item \hyperref[sets-section-phantom]{Théorie des ensembles}
\item \hyperref[categories-section-phantom]{Catégories}
\item \hyperref[topology-section-phantom]{Topologie}
\item \hyperref[sheaves-section-phantom]{Faisceaux sur les espaces}
\item \hyperref[sites-section-phantom]{Sites et faisceaux}
\item \hyperref[stacks-section-phantom]{Champs}
\item \hyperref[fields-section-phantom]{Corps}
\item \hyperref[algebra-section-phantom]{Algèbre commutative}
\item \hyperref[brauer-section-phantom]{Groupes de Brauer}
\item \hyperref[homology-section-phantom]{Algèbre homologique}
\item \hyperref[derived-section-phantom]{Catégories dérivées}
\item \hyperref[simplicial-section-phantom]{Méthodes simpliciales}
\item \hyperref[more-algebra-section-phantom]{Compléments d'algèbre}
\item \hyperref[smoothing-section-phantom]{Lissage des morphismes d'anneaux}
\item \hyperref[modules-section-phantom]{Faisceaux de modules}
\item \hyperref[sites-modules-section-phantom]{Modules sur les sites}
\item \hyperref[injectives-section-phantom]{Objets injectifs}
\item \hyperref[cohomology-section-phantom]{Cohomologie des faisceaux}
\item \hyperref[sites-cohomology-section-phantom]{Cohomologie sur les sites}
\item \hyperref[dga-section-phantom]{Algèbre différentielle graduée}
\item \hyperref[dpa-section-phantom]{Algèbre à puissances divisées}
\item \hyperref[sdga-section-phantom]{Faisceaux différentiels gradués}
\item \hyperref[hypercovering-section-phantom]{Hyperrecouvrements}
\end{enumerate}
Schémas
\begin{enumerate}
\setcounter{enumi}{25}
\item \hyperref[schemes-section-phantom]{Schémas}
\item \hyperref[constructions-section-phantom]{Constructions de schémas}
\item \hyperref[properties-section-phantom]{Propriétés des schémas}
\item \hyperref[morphisms-section-phantom]{Morphismes de schémas}
\item \hyperref[coherent-section-phantom]{Cohomologie des schémas}
\item \hyperref[divisors-section-phantom]{Diviseurs}
\item \hyperref[limits-section-phantom]{Limites de schémas}
\item \hyperref[varieties-section-phantom]{Variétés}
\item \hyperref[topologies-section-phantom]{Topologies sur les schémas}
\item \hyperref[descent-section-phantom]{Descente}
\item \hyperref[perfect-section-phantom]{Catégories dérivées de schémas}
\item \hyperref[more-morphisms-section-phantom]{Compléments sur les morphismes}
\item \hyperref[flat-section-phantom]{Compléments sur la platitude}
\item \hyperref[groupoids-section-phantom]{Schémas en groupoïdes}
\item \hyperref[more-groupoids-section-phantom]{Compléments sur les schémas en groupoïdes}
\item \hyperref[etale-section-phantom]{Morphismes étales de schémas}
\end{enumerate}
Thèmes de théorie des schémas
\begin{enumerate}
\setcounter{enumi}{41}
\item \hyperref[chow-section-phantom]{Homologie de Chow}
\item \hyperref[intersection-section-phantom]{Théorie de l'intersection}
\item \hyperref[pic-section-phantom]{Schémas de Picard des courbes}
\item \hyperref[weil-section-phantom]{Théories cohomologiques de Weil}
\item \hyperref[adequate-section-phantom]{Modules adéquats}
\item \hyperref[dualizing-section-phantom]{Complexes dualisants}
\item \hyperref[duality-section-phantom]{Dualité pour les schémas}
\item \hyperref[discriminant-section-phantom]{Discriminants et différentes}
\item \hyperref[derham-section-phantom]{Cohomologie de de Rham}
\item \hyperref[local-cohomology-section-phantom]{Cohomologie locale}
\item \hyperref[algebraization-section-phantom]{Géométrie algébrique et formelle}
\item \hyperref[curves-section-phantom]{Courbes algébriques}
\item \hyperref[resolve-section-phantom]{Résolution des surfaces}
\item \hyperref[models-section-phantom]{Réduction semi-stable}
\item \hyperref[functors-section-phantom]{Foncteurs et morphismes}
\item \hyperref[equiv-section-phantom]{Catégories dérivées de variétés}
\item \hyperref[pione-section-phantom]{Groupes fondamentaux de schémas}
\item \hyperref[etale-cohomology-section-phantom]{Cohomologie étale}
\item \hyperref[crystalline-section-phantom]{Cohomologie cristalline}
\item \hyperref[proetale-section-phantom]{Cohomologie pro-étale}
\item \hyperref[relative-cycles-section-phantom]{Cycles relatifs}
\item \hyperref[more-etale-section-phantom]{Compléments de cohomologie étale}
\item \hyperref[trace-section-phantom]{La formule des traces}
\end{enumerate}
Espaces algébriques
\begin{enumerate}
\setcounter{enumi}{64}
\item \hyperref[spaces-section-phantom]{Espaces algébriques}
\item \hyperref[spaces-properties-section-phantom]{Propriétés des espaces algébriques}
\item \hyperref[spaces-morphisms-section-phantom]{Morphismes d'espaces algébriques}
\item \hyperref[decent-spaces-section-phantom]{Espaces algébriques décents}
\item \hyperref[spaces-cohomology-section-phantom]{Cohomologie des espaces algébriques}
\item \hyperref[spaces-limits-section-phantom]{Limites d'espaces algébriques}
\item \hyperref[spaces-divisors-section-phantom]{Diviseurs sur les espaces algébriques}
\item \hyperref[spaces-over-fields-section-phantom]{Espaces algébriques sur des corps}
\item \hyperref[spaces-topologies-section-phantom]{Topologies sur les espaces algébriques}
\item \hyperref[spaces-descent-section-phantom]{Descente et espaces algébriques}
\item \hyperref[spaces-perfect-section-phantom]{Catégories dérivées d'espaces}
\item \hyperref[spaces-more-morphisms-section-phantom]{Compléments sur les morphismes d'espaces}
\item \hyperref[spaces-flat-section-phantom]{Platitude sur les espaces algébriques}
\item \hyperref[spaces-groupoids-section-phantom]{Groupoïdes dans les espaces algébriques}
\item \hyperref[spaces-more-groupoids-section-phantom]{Compléments sur les groupoïdes dans les espaces}
\item \hyperref[bootstrap-section-phantom]{Amorçage}
\item \hyperref[spaces-pushouts-section-phantom]{Sommes amalgamées d'espaces algébriques}
\end{enumerate}
Thèmes de géométrie
\begin{enumerate}
\setcounter{enumi}{81}
\item \hyperref[spaces-chow-section-phantom]{Groupes de Chow des espaces}
\item \hyperref[groupoids-quotients-section-phantom]{Quotients de groupoïdes}
\item \hyperref[spaces-more-cohomology-section-phantom]{Compléments sur la cohomologie des espaces}
\item \hyperref[spaces-simplicial-section-phantom]{Espaces simpliciaux}
\item \hyperref[spaces-duality-section-phantom]{Dualité pour les espaces}
\item \hyperref[formal-spaces-section-phantom]{Espaces algébriques formels}
\item \hyperref[restricted-section-phantom]{Algébrisation des espaces formels}
\item \hyperref[spaces-resolve-section-phantom]{Résolution des surfaces revisitée}
\end{enumerate}
Théorie des déformations
\begin{enumerate}
\setcounter{enumi}{89}
\item \hyperref[formal-defos-section-phantom]{Théorie formelle des déformations}
\item \hyperref[defos-section-phantom]{Théorie des déformations}
\item \hyperref[cotangent-section-phantom]{Le complexe cotangent}
\item \hyperref[examples-defos-section-phantom]{Problèmes de déformation}
\end{enumerate}
Champs algébriques
\begin{enumerate}
\setcounter{enumi}{93}
\item \hyperref[algebraic-section-phantom]{Champs algébriques}
\item \hyperref[examples-stacks-section-phantom]{Exemples de champs}
\item \hyperref[stacks-sheaves-section-phantom]{Faisceaux sur les champs algébriques}
\item \hyperref[criteria-section-phantom]{Critères de représentabilité}
\item \hyperref[artin-section-phantom]{Axiomes d'Artin}
\item \hyperref[quot-section-phantom]{Espaces de Quot et de Hilbert}
\item \hyperref[stacks-properties-section-phantom]{Propriétés des champs algébriques}
\item \hyperref[stacks-morphisms-section-phantom]{Morphismes de champs algébriques}
\item \hyperref[stacks-limits-section-phantom]{Limites de champs algébriques}
\item \hyperref[stacks-cohomology-section-phantom]{Cohomologie des champs algébriques}
\item \hyperref[stacks-perfect-section-phantom]{Catégories dérivées de champs}
\item \hyperref[stacks-introduction-section-phantom]{Introduction aux champs algébriques}
\item \hyperref[stacks-more-morphisms-section-phantom]{Compléments sur les morphismes de champs}
\item \hyperref[stacks-geometry-section-phantom]{La géométrie des champs}
\end{enumerate}
Thèmes de théorie des espaces de modules
\begin{enumerate}
\setcounter{enumi}{107}
\item \hyperref[moduli-section-phantom]{Champs de modules}
\item \hyperref[moduli-curves-section-phantom]{Espaces de modules de courbes}
\end{enumerate}
Divers
\begin{enumerate}
\setcounter{enumi}{109}
\item \hyperref[examples-section-phantom]{Exemples}
\item \hyperref[exercises-section-phantom]{Exercices}
\item \hyperref[guide-section-phantom]{Guide bibliographique}
\item \hyperref[desirables-section-phantom]{Desiderata}
\item \hyperref[coding-section-phantom]{Style de programmation}
\item \hyperref[obsolete-section-phantom]{Obsolète}
\item \hyperref[fdl-section-phantom]{Licence de documentation libre GNU}
\item \hyperref[index-section-phantom]{Index généré automatiquement}
\end{enumerate}
\end{multicols}


\bibliography{my}
\bibliographystyle{amsalpha}

\end{document}
